% Auto-generated from data/segments.json + layout.json (+ transforms) — do not edit by hand.
% volume: 14Sam03
% mode: reading
% segments: 2919
% transform_rules: 0
% toc_marks: 636

% Volume layout defaults (layout.json ``layout``).
\csromanlayoutapply{1}{1.5}{21.6pt}{5pt}{6.3pt}{65pt}{2.5em}%

% Reading physical-page layout (layout.json ``page_layout_reading_mode``).
\csromanreadingpagelayoutvolume{1}{1.5}{21.6pt}{5pt}{6.3pt}{65pt}{2.5em}%
\csromanreadingpagelayoutenable

\csromanheader{2}{\textbf{สํยุตฺตนิกาย}}
\csromanmatikaanchor{mtk.0}
\csromantocmarknopage{chapter}{มหาวคฺคสํยุตฺตปาฬิ}{mtk.0}
\bookmark[dest={mtk.0},level=0]{มหาวคฺคสํยุตฺตปาฬิ}
\gambhira{\textbf{มหาวคฺคสํยุตฺตปาฬิ}}
\namakkaram{นโม ตสฺส ภควโต อรหโต สมฺมาสมฺพุทฺธสฺส\textbf{.}}
\chapterhead{1. อวิชฺชาวคฺค}
\csromanheader{2}{1. อวิชฺชาสุตฺต}
\proseitem{1}{\csromanfolio{1}เอวํ เม สุตํ\csromandash{}เอกํ สมยํ ภควา สาวตฺถิยํ วิหรติ เชตวเน อนาถปิณฺฑิกสฺส อาราเม\textbf{.}\csromanspacer{} ตตฺร โข ภควา ภิกฺขู อามนฺเตสิ “ภิกฺขโว”ติ\textbf{.} “ภทนฺเต”ติ เต ภิกฺขู ภควโต ปจฺจสฺโสสุํ\textbf{.}\csromanspacer{} ภควา เอตทโวจ\csromandash{}}

\prose{อวิชฺชา ภิกฺขเว ปุพฺพงฺคมา อกุสลานํ ธมฺมานํ สมาปตฺติยา, อนฺวเทว\footnote{อนุเทว (สี, อิ, ก)} อหิริกํ อโนตฺตปฺปํ\textbf{.}\csromanspacer{} อวิชฺชาคตสฺส ภิกฺขเว อวิทฺทสุโน มิจฺฉาทิฏฺฐิ ปโหติ, มิจฺฉาทิฏฺฐิสฺส มิจฺฉาสงฺกปฺโป ปโหติ, มิจฺฉาสงฺกปฺปสฺส มิจฺฉาวาจา ปโหติ, มิจฺฉาวาจสฺส มิจฺฉากมฺมนฺโต ปโหติ, มิจฺฉากมฺมนฺตสฺส มิจฺฉา-อาชีโว ปโหติ, มิจฺฉา อาชีวสฺส มิจฺฉาวายาโม ปโหติ, มิจฺฉาวายามสฺส มิจฺฉาสติ ปโหติ, มิจฺฉาสติสฺส มิจฺฉาสมาธิ ปโหติ\textbf{.}}

\prose{\csromanfolio{2}วิชฺชา จ โข ภิกฺขเว ปุพฺพงฺคมา กุสลานํ ธมฺมานํ สมาปตฺติยา, อนฺวเทว หิโรตฺตปฺปํ.\csromanspacer{} วิชฺชาคตสฺส ภิกฺขเว วิทฺทสุโน สมฺมาทิฏฺฐิ ปโหติ, สมฺมาทิฏฺฐิสฺส สมฺมาสงฺกปฺโป ปโหติ, สมฺมาสงฺกปฺปสฺส สมฺมาวาจา ปโหติ, สมฺมาวาจสฺส สมฺมากมฺมนฺโต ปโหติ, สมฺมากมฺมนฺตสฺส สมฺมา-อาชีโว ปโหติ, สมฺมา-อาชีวสฺส สมฺมาวายาโม ปโหติ, สมฺมาวายามสฺส สมฺมาสติ ปโหติ, สมฺมาสติสฺส สมฺมาสมาธิ ปโหตีติ.\csromanspacer{}\csromanspacer{}ปฐมํ.}

\csromanheader{2}{2. อุปฑฺฒสุตฺต}
\proseitem{2}{เอวํ เม สุตํ\csromandash{}เอกํ สมยํ ภควา สเกฺยสุ วิหรติ นครกํ นาม\footnote{นาครกํ นาม (สี), สกฺกรํ นาม (สฺยา, ก)} สกฺยานํ นิคโม.\csromanspacer{} อถ โข อายสฺมา อานนฺโท เยน ภควา เตนุปสงฺกมิ, อุปสงฺกมิตฺวา ภควนฺตํ อภิวาเทตฺวา เอกมนฺตํ นิสีทิ, เอกมนฺตํ นิสินฺโน โข อายสฺมา อานนฺโท ภควนฺตํ เอตทโวจ “อุปฑฺฒมิทํ ภนฺเต พฺรหฺมจริยํ ยทิทํ กลฺยาณมิตฺตตา กลฺยาณสหายตา กลฺยาณสมฺปวงฺกตา”ติ.}

\prose{มา เหวํ อานนฺท, มา เหวํ อานนฺท, สกลเมวิทํ อานนฺท พฺรหฺมจริยํ ยทิทํ กลฺยาณมิตฺตตา กลฺยาณสหายตา กลฺยาณสมฺปวงฺกตา, กลฺยาณมิตฺตสฺเสตํ อานนฺท ภิกฺขุโน ปาฏิกงฺขํ กลฺยาณสหายสฺส กลฺยาณสมฺปวงฺกสฺส, อริยํ อฏฺฐงฺคิกํ มคฺคํ ภาเวสฺสติ, อริยํ อฏฺฐงฺคิกํ มคฺคํ พหุลีกริสฺสติ.}

\prose{กถญฺจานนฺท ภิกฺขุ กลฺยาณมิตฺโต กลฺยาณสหาโย กลฺยาณสมฺปวงฺโก อริยํ อฏฺฐงฺคิกํ มคฺคํ ภาเวติ, อริยํ อฏฺฐงฺคิกํ มคฺคํ พหุลีกโรติ.\csromanspacer{} อิธานนฺท ภิกฺขุ สมฺมาทิฏฺฐิํ ภาเวติ วิเวกนิสฺสิตํ วิราคนิสฺสิตํ นิโรธนิสฺสิตํ โวสฺสคฺคปริณามิํ, สมฺมาสงฺกปฺปํ ภาเวติ วิเวกนิสฺสิตํ ฯเปฯ สมฺมาวาจํ ภาเวติ ฯเปฯ สมฺมากมฺมนฺตํ ภาเวติ ฯเปฯ สมฺมา-อาชีวํ ภาเวติ ฯเปฯ สมฺมาวายามํ ภาเวติ ฯเปฯ สมฺมาสติํ ภาเวติ ฯเปฯ สมฺมาสมาธิํ ภาเวติ วิเวกนิสฺสิตํ วิราคนิสฺสิตํ นิโรธนิสฺสิตํ โวสฺสคฺคปริณามิํ.\csromanspacer{} เอวํ โข อานนฺท ภิกฺขุ กลฺยาณมิตฺโต กลฺยาณสหาโย กลฺยาณสมฺปวงฺโก อริยํ อฏฺฐงฺคิกํ มคฺคํ ภาเวติ, อริยํ อฏฺฐงฺคิกํ มคฺคํ พหุลีกโรติ.}

\prose{\csromanfolio{3}ตทมินาเปตํ อานนฺท ปริยาเยน เวทิตพฺพํ, ยถา สกลเมวิทํ พฺรหฺมจริยํ ยทิทํ กลฺยาณมิตฺตตา กลฺยาณสหายตา กลฺยาณสมฺป- วงฺกตา.\csromanspacer{} มมํ หิ อานนฺท กลฺยาณมิตฺตํ อาคมฺม ชาติธมฺมา สตฺตา ชาติยา ปริมุจฺจนฺติ, ชราธมฺมา สตฺตา ชราย ปริมุจฺจนฺติ, มรณธมฺมา สตฺตา มรเณน ปริมุจฺจนฺติ, โสกปริเทวทุกฺขโทมนสฺสุปายาสธมฺมา สตฺตา โสกปริเทวทุกฺขโทมนสฺสุปายาเสหิ ปริมุจฺจนฺติ.\csromanspacer{} อิมินา โข เอตํ อานนฺท ปริยาเยน เวทิตพฺพํ, ยถา สกลเมวิทํ พฺรหฺมจริยํ ยทิทํ กลฺยาณมิตฺตตา กลฺยาณสหายตา กลฺยาณสมฺปวงฺกตาติ.\csromanspacer{}\csromanspacer{}ทุติยํ.}

\csromanheader{2}{3. สาริปุตฺตสุตฺต}
\proseitem{3}{สาวตฺถินิทานํ.\csromanspacer{} อถ โข อายสฺมา สาริปุตฺโต เยน ภควา เตนุปสงฺกมิ, อุปสงฺกมิตฺวา ภควนฺตํ อภิวาเทตฺวา เอกมนฺตํ นิสีทิ, เอกมนฺตํ นิสินฺโน โข อายสฺมา สาริปุตฺโต ภควนฺตํ เอตทโวจ “สกลมิทํ ภนฺเต พฺรหฺมจริยํ ยทิทํ กลฺยาณมิตฺตตา กลฺยาณสหายตา กลฺยาณสมฺปวงฺกตา”ติ.}

\prose{สาธุ สาธุ สาริปุตฺต, สกลมิทํ สาริปุตฺต พฺรหฺมจริยํ ยทิทํ กลฺยาณมิตฺตตา กลฺยาณสหายตา กลฺยาณสมฺปวงฺกตา.\csromanspacer{} กลฺยาณมิตฺตสฺเสตํ สาริปุตฺต ภิกฺขุโน ปาฏิกงฺขํ กลฺยาณสหายสฺส กลฺยาณสมฺปวงฺกสฺส, อริยํ อฏฺฐงฺคิกํ มคฺคํ ภาเวสฺสติ, อริยํ อฏฺฐงฺคิกํ มคฺคํ พหุลีกริสฺสติ.\csromanspacer{} กถญฺจ สาริปุตฺต ภิกฺขุ กลฺยาณมิตฺโต กลฺยาณสหาโย กลฺยาณสมฺปวงฺโก อริยํ อฏฺฐงฺคิกํ มคฺคํ ภาเวติ, อริยํ อฏฺฐงฺคิกํ มคฺคํ พหุลีกโรติ.}

\prose{อิธ สาริปุตฺต ภิกฺขุ สมฺมาทิฏฺฐิํ ภาเวติ วิเวกนิสฺสิตํ วิราคนิสฺสิตํ นิโรธนิสฺสิตํ โวสฺสคฺคปริณามิํ ฯเปฯ สมฺมาสมาธิํ ภาเวติ วิเวกนิสฺสิตํ วิราคนิสฺสิตํ นิโรธนิสฺสิตํ โวสฺสคฺคปริณามิํ.\csromanspacer{} เอวํ โข สาริปุตฺต ภิกฺขุ กลฺยาณมิตฺโต กลฺยาณสหาโย กลฺยาณสมฺปวงฺโก อริยํ อฏฺฐงฺคิกํ มคฺคํ ภาเวติ, อริยํ อฏฺฐงฺคิกํ มคฺคํ พหุลีกโรติ.}

\prose{ตทมินาเปตํ สาริปุตฺต ปริยาเยน เวทิตพฺพํ, ยถา สกลมิทํ พฺรหฺมจริยํ ยทิทํ กลฺยาณมิตฺตตา กลฺยาณสหายตา กลฺยาณสมฺปวงฺกตา.\csromanspacer{} มมํ หิ สาริปุตฺต กลฺยาณมิตฺตํ อาคมฺม ชาติธมฺมา สตฺตา \csromanfolio{4}ชาติยา ปริมุจฺจนฺติ, ชราธมฺมา สตฺตา ชราย ปริมุจฺจนฺติ, มรณธมฺมา สตฺตา มรเณน ปริมุจฺจนฺติ, โสกปริเทวทุกฺขโทมนสฺสุปายาสธมฺมา สตฺตา โสกปริเทวทุกฺขโทมนสฺสุปายาเสหิ ปริมุจฺจนฺติ.\csromanspacer{} อิมินา โข เอตํ สาริปุตฺต ปริยาเยน เวทิตพฺพํ, ยถา สกลมิทํ พฺรหฺมจริยํ ยทิทํ กลฺยาณมิตฺตตา กลฺยาณสหายตา กลฺยาณสมฺปวงฺกตาติ.\csromanspacer{}\csromanspacer{}ตติยํ.}

\csromanheader{2}{4. ชาณุสฺโสณิพฺราหฺมณสุตฺต}
\proseitem{4}{สาวตฺถินิทานํ.\csromanspacer{} อถ โข อายสฺมา อานนฺโท ปุพฺพณฺหสมยํ นิวาเสตฺวา ปตฺตจีวรมาทาย สาวตฺถิํ ปิณฺฑาย ปาวิสิ.\csromanspacer{} อทฺทสา โข อายสฺมา อานนฺโท ชาณุสฺโสณิํ พฺราหฺมณํ สพฺพเสเตน วฬวาภิรเถน\footnote{วฬภีรเถน (สี)} สาวตฺถิยา นิยฺยายนฺตํ, เสตา สุทํ อสฺสา ยุตฺตา โหนฺติ เสตาลงฺการา เสโต รโถ เสตปริวาโร เสตา รสฺมิโย เสตา ปโตทลฏฺฐิ เสตํ ฉตฺตํ เสตํ อุณฺหีสํ เสตานิ วตฺถานิ เสตา อุปาหนา เสตาย สุทํ วาลพีชนิยา พีชียติ.\csromanspacer{} ตเมนํ ชโน ทิสฺวา เอวมาห “พฺรหฺมํ วต โภ ยานํ, พฺรหฺมยานรูปํ วต โภ”ติ.}

\prose{อถ โข อายสฺมา อานนฺโท สาวตฺถิยํ ปิณฺฑาย จริตฺวา ปจฺฉาภตฺตํ ปิณฺฑปาตปฏิกฺกนฺโต เยน ภควา เตนุปสงฺกมิ, อุปสงฺกมิตฺวา ภควนฺตํ อภิวาเทตฺวา เอกมนฺตํ นิสีทิ, เอกมนฺตํ นิสินฺโน โข อายสฺมา อานนฺโท ภควนฺตํ เอตทโวจ\csromandash{}}

\prose{อิธาหํ ภนฺเต ปุพฺพณฺหสมยํ นิวาเสตฺวา ปตฺตจีวรมาทาย สาวตฺถิํ ปิณฺฑาย ปาวิสิํ, อทฺทสํ ขฺวาหํ ภนฺเต ชาณุสฺโสณิํ พฺราหฺมณํ สพฺพเสเตน วฬวาภิรเถน สาวตฺถิยา นิยฺยายนฺตํ, เสตา สุทํ อสฺสา ยุตฺตา โหนฺติ เสตาลงฺการา เสโต รโถ เสตปริวาโร เสตา รสฺมิโย เสตา ปโตทลฏฺฐิ เสตํ ฉตฺตํ เสตํ อุณฺหีสํ เสตานิ วตฺถานิ เสตา อุปาหนา เสตาย สุทํ วาลพีชนิยา พีชียติ.\csromanspacer{} ตเมนํ ชโน ทิสฺวา เอวมาห “พฺรหฺมํ วต โภ ยานํ, พฺรหฺมยานรูปํ วต โภ”ติ.\csromanspacer{} สกฺกา นุ โข ภนฺเต อิมสฺมิํ ธมฺมวินเย พฺรหฺมยานํ ปญฺญาเปตุนฺติ.}

\prose{\csromanfolio{5}“สกฺกา อานนฺทา”ติ ภควา อโวจ.\csromanspacer{} อิมสฺเสว โข เอตํ อานนฺท อริยสฺส อฏฺฐงฺคิกสฺส มคฺคสฺส อธิวจนํ พฺรหฺมยานํ อิติปิ ธมฺมยานํ อิติปิ อนุตฺตโร สงฺคามวิชโย อิติปีติ.}

\prose{สมฺมาทิฏฺฐิ อานนฺท ภาวิตา พหุลีกตา ราควินยปริโยสานา โหติ, โทสวินยปริโยสานา โหติ, โมหวินยปริโยสานา โหติ.\csromanspacer{} สมฺมาสงฺกปฺโป อานนฺท ภาวิโต พหุลีกโต ราควินยปริโยสาโน โหติ, โทสวินยปริโยสาโน โหติ, โมหวินยปริโยสาโน โหติ.\csromanspacer{} สมฺมาวาจา อานนฺท ภาวิตา พหุลีกตา ราควินยปริโยสานา โหติ, โทส ฯเปฯ โมหวินยปริโยสานา โหติ.\csromanspacer{} สมฺมากมฺมนฺโต อานนฺท ภาวิโต พหุลีกโต ราควินยปริโยสาโน โหติ, โทส ฯเปฯ โมหวินยปริโยสาโน โหติ.\csromanspacer{} สมฺมา-อาชีโว อานนฺท ภาวิโต พหุลีกโต ราควินยปริโยสาโน โหติ, โทส ฯเปฯ โมหวินยปริโยสาโน โหติ.\csromanspacer{} สมฺมาวายาโม อานนฺท ภาวิโต พหุลีกโต ราควินยปริโยสาโน โหติ, โทส ฯเปฯ โมหวินยปริโยสาโน โหติ.\csromanspacer{} สมฺมาสติ อานนฺท ภาวิตา พหุลีกตา ราควินยปริโยสานา โหติ, โทส ฯเปฯ โมหวินยปริโยสานา โหติ.\csromanspacer{} สมฺมาสมาธิ อานนฺท ภาวิโต พหุลีกโต ราควินยปริโยสาโน โหติ, โทส ฯเปฯ โมหวินยปริโยสาโน โหติ.}

\prose{อิมินา โข เอตํ อานนฺท ปริยาเยน เวทิตพฺพํ, ยถา อิมสฺเสเวตํ อริยสฺส อฏฺฐงฺคิกสฺส มคฺคสฺส อธิวจนํ พฺรหฺมยานํ อิติปิ ธมฺมยานํ อิติปิ อนุตฺตโร สงฺคามวิชโย อิติปีติ.\csromanspacer{} อิทมโวจ ภควา, อิทํ วตฺวาน สุคโต อถาปรํ เอตทโวจ สตฺถา\csromandash{}}

\csromangathameasure{“ยสฺส สทฺธา จ ปญฺญา จ, ธมฺมา ยุตฺตา สทา ธุรํ. \\ หิรี อีสา มโนโยตฺตํ, สติ อารกฺขสารถิ. \\ รโถ สีลปริกฺขาโร, ฌานกฺโข จกฺกวีริโย. \\ อุเปกฺขา ธุรสมาธิ, อนิจฺฉา ปริวารณํ. \\ อพฺยาปาโท อวิหิํสา, วิเวโก ยสฺส อาวุธํ. \\ ติติกฺขาจมฺมสนฺนาโห, โยคกฺเขมาย วตฺตติ. \\ เอตทตฺตนิ สมฺภูตํ, พฺรหฺมยานํ อนุตฺตรํ. \\ นิยฺยนฺติ ธีรา โลกมฺหา, อญฺญทตฺถุ ชยํ ชยนฺ”ติ. จตุตฺถํ.}
\csromangathagroup{\csromangathastanza{“ยสฺส สทฺธา จ ปญฺญา จ, ธมฺมา ยุตฺตา สทา ธุรํ. \\ หิรี อีสา มโนโยตฺตํ, สติ อารกฺขสารถิ.}\csromangathastanza{รโถ สีลปริกฺขาโร, ฌานกฺโข จกฺกวีริโย. \\ อุเปกฺขา ธุรสมาธิ, อนิจฺฉา ปริวารณํ.}\csromangathastanza{อพฺยาปาโท อวิหิํสา, วิเวโก ยสฺส อาวุธํ. \\ ติติกฺขาจมฺมสนฺนาโห\footnote{วมฺมสนฺนาโห (สี)}, โยคกฺเขมาย วตฺตติ.}\csromangathastanza{\csromanfolio{6}เอตทตฺตนิ สมฺภูตํ, พฺรหฺมยานํ อนุตฺตรํ. \\ นิยฺยนฺติ ธีรา โลกมฺหา, อญฺญทตฺถุ ชยํ ชยนฺ”ติ.\csromanspacer{} จตุตฺถํ.}}

\csromanheader{2}{5. กิมตฺถิยสุตฺต}
\proseitem{5}{สาวตฺถินิทานํ.\csromanspacer{} อถ โข สมฺพหุลา ภิกฺขู เยน ภควา เตนุปสงฺกมิํสุ ฯเปฯ เอกมนฺตํ นิสีทิํสุ, เอกมนฺตํ นิสินฺนา โข เต ภิกฺขู ภควนฺตํ เอตทโวจุํ\csromandash{}}

\prose{อิธ โน ภนฺเต อญฺญติตฺถิยา ปริพฺพาชกา อเมฺห เอวํ ปุจฺฉนฺติ “กิมตฺถิยํ อาวุโส สมเณ โคตเม พฺรหฺมจริยํ วุสฺสตี”ติ.\csromanspacer{} เอวํ ปุฏฺฐา มยํ ภนฺเต เตสํ อญฺญติตฺถิยานํ ปริพฺพาชกานํ เอวํ พฺยากโรม “ทุกฺขสฺส โข อาวุโส ปริญฺญตฺถํ ภควติ พฺรหฺมจริยํ วุสฺสตี”ติ.\csromanspacer{} กจฺจิ มยํ ภนฺเต เอวํ ปุฏฺฐา เอวํ พฺยากรมานา วุตฺตวาทิโน เจว ภควโต โหม, น จ ภควนฺตํ อภูเตน อพฺภาจิกฺขาม, ธมฺมสฺส จานุธมฺมํ พฺยากโรม, น จ โกจิ สหธมฺมิโก วาทานุวาโท คารยฺหํ ฐานํ อาคจฺฉตีติ.}

\prose{ตคฺฆ ตุเมฺห ภิกฺขเว เอวํ ปุฏฺฐา เอวํ พฺยากรมานา วุตฺตวาทิโน เจว เม โหถ, น จ มํ อภูเตน อพฺภาจิกฺขถ, ธมฺมสฺส จานุธมฺมํ พฺยากโรถ, น จ โกจิ สหธมฺมิโก วาทานุวาโท คารยฺหํ ฐานํ อาคจฺฉติ.\csromanspacer{} ทุกฺขสฺส หิ ปริญฺญตฺถํ มยิ พฺรหฺมจริยํ วุสฺสติ.\csromanspacer{} สเจ โว ภิกฺขเว อญฺญติตฺถิยา ปริพฺพาชกา เอวํ ปุจฺเฉยฺยุํ “อตฺถิ ปนาวุโส มคฺโค, อตฺถิ ปฏิปทา เอตสฺส ทุกฺขสฺส ปริญฺญายา”ติ.\csromanspacer{} เอวํ ปุฏฺฐา ตุเมฺห ภิกฺขเว เตสํ อญฺญติตฺถิยานํ ปริพฺพาชกานํ เอวํ พฺยากเรยฺยาถ “อตฺถิ โข อาวุโส มคฺโค, อตฺถิ ปฏิปทา เอตสฺส ทุกฺขสฺส ปริญฺญายา”ติ.}

\prose{กตโม จ ภิกฺขเว มคฺโค, กตมา ปฏิปทา เอตสฺส ทุกฺขสฺส ปริญฺญายาติ.\csromanspacer{} อยเมว อริโย อฏฺฐงฺคิโก มคฺโค.\csromanspacer{} เสยฺยถิทํ, สมฺมาทิฏฺฐิ ฯเปฯ สมฺมาสมาธิ.\csromanspacer{} อยํ ภิกฺขเว มคฺโค, อยํ ปฏิปทา เอตสฺส ทุกฺขสฺส ปริญฺญายาติ.\csromanspacer{} เอวํ ปุฏฺฐา ตุเมฺห ภิกฺขเว เตสํ อญฺญติตฺถิยานํ ปริพฺพาชกานํ เอวํ พฺยากเรยฺยาถาติ.\csromanspacer{}\csromanspacer{}ปญฺจมํ.}

\csromanheader{2}{6. ปฐม-อญฺญตรภิกฺขุสุตฺต}
\proseitem{6}{\csromanfolio{7}สาวตฺถินิทานํ.\csromanspacer{} อถ โข อญฺญตโร ภิกฺขุ เยน ภควา เตนุปสงฺกมิ ฯเปฯ เอกมนฺตํ นิสินฺโน โข โส ภิกฺขุ ภควนฺตํ เอตทโวจ\csromandash{}“พฺรหฺมจริยํ พฺรหฺมจริยนฺติ ภนฺเต วุจฺจติ, กตมํ นุ โข ภนฺเต พฺรหฺมจริยํ, กตมํ พฺรหฺมจริยปริโยสานนฺ”ติ.}

\prose{อยเมว โข ภิกฺขุ อริโย อฏฺฐงฺคิโก มคฺโค พฺรหฺมจริยํ.\csromanspacer{} เสยฺยถิทํ, สมฺมาทิฏฺฐิ ฯเปฯ สมฺมาสมาธิ.\csromanspacer{} โย โข ภิกฺขุ ราคกฺขโย โทสกฺขโย โมหกฺขโย, อิทํ พฺรหฺมจริยปริโยสานนฺติ.\csromanspacer{}\csromanspacer{}ฉฏฺฐํ.}

\csromanheader{2}{7. ทุติย-อญฺญตรภิกฺขุสุตฺต}
\proseitem{7}{สาวตฺถินิทานํ.\csromanspacer{} อถ โข อญฺญตโร ภิกฺขุ เยน ภควา เตนุปสงฺกมิ ฯเปฯ เอกมนฺตํ นิสินฺโน โข โส ภิกฺขุ ภควนฺตํ เอตทโวจ\csromandash{}}

\prose{“ราควินโย โทสวินโย โมหวินโย”ติ ภนฺเต วุจฺจติ, กิสฺส นุ โข เอตํ ภนฺเต อธิวจนํ “ราควินโย โทสวินโย โมหวินโย”ติ.\csromanspacer{} นิพฺพานธาตุยา โข เอตํ ภิกฺขุ อธิวจนํ ราควินโย โทสวินโย โมหวินโยติ, อาสวานํ ขโย เตน วุจฺจตีติ.}

\prose{เอวํ วุตฺเต โส ภิกฺขุ ภควนฺตํ เอตทโวจ “อมตํ อมตนฺ”ติ ภนฺเต วุจฺจติ, กตมํ นุ โข ภนฺเต อมตํ, กตโม อมตคามิมคฺโคติ.\csromanspacer{} โย โข ภิกฺขุ ราคกฺขโย โทสกฺขโย โมหกฺขโย, อิทํ วุจฺจติ อมตํ.\csromanspacer{} อยเมว อริโย อฏฺฐงฺคิโก มคฺโค อมตคามิมคฺโค.\csromanspacer{} เสยฺยถิทํ, สมฺมาทิฏฺฐิ ฯเปฯ สมฺมาสมาธีติ.\csromanspacer{}\csromanspacer{}สตฺตมํ.}

\csromanheader{2}{8. วิภงฺคสุตฺต}
\proseitem{8}{สาวตฺถินิทานํ.\csromanspacer{} อริยํ โว ภิกฺขเว อฏฺฐงฺคิกํ มคฺคํ เทเสสฺสามิ วิภชิสฺสามิ, ตํ สุณาถ สาธุกํ มนสิ กโรถ, ภาสิสฺสามีติ. “เอวํ ภนฺเต”ติ โข เต ภิกฺขู ภควโต ปจฺจสฺโสสุํ.\csromanspacer{} ภควา เอตทโวจ\csromandash{}}

\prose{กตโม จ ภิกฺขเว อริโย อฏฺฐงฺคิโก มคฺโค, เสยฺยถิทํ, สมฺมาทิฏฺฐิ ฯเปฯ สมฺมาสมาธิ.\csromanspacer{} กตมา จ ภิกฺขเว สมฺมาทิฏฺฐิ, ยํ โข ภิกฺขเว \csromanfolio{8}ทุกฺเข ญาณํ ทุกฺขสมุทเย ญาณํ ทุกฺขนิโรเธ ญาณํ ทุกฺขนิโรธคามินิยา ปฏิปทาย ญาณํ.\csromanspacer{} อยํ วุจฺจติ ภิกฺขเว สมฺมาทิฏฺฐิ.}

\prose{กตโม จ ภิกฺขเว สมฺมาสงฺกปฺโป, โย โข ภิกฺขเว เนกฺขมฺมสงฺกปฺโป อพฺยาปาทสงฺกปฺโป อวิหิํสาสงฺกปฺโป.\csromanspacer{} อยํ วุจฺจติ ภิกฺขเว สมฺมาสงฺกปฺโป.}

\prose{กตมา จ ภิกฺขเว สมฺมาวาจา, ยา โข ภิกฺขเว มุสาวาทา เวรมณี ปิสุณาย วาจาย เวรมณี ผรุสาย วาจาย เวรมณี สมฺผปฺปลาปา เวรมณี.\csromanspacer{} อยํ วุจฺจติ ภิกฺขเว สมฺมาวาจา.}

\prose{กตโม จ ภิกฺขเว สมฺมากมฺมนฺโต, ยา โข ภิกฺขเว ปาณาติปาตา เวรมณี อทินฺนาทานา เวรมณี อพฺรหฺมจริยา เวรมณี.\csromanspacer{} อยํ วุจฺจติ ภิกฺขเว สมฺมากมฺมนฺโต.}

\prose{กตโม จ ภิกฺขเว สมฺมา-อาชีโว, อิธ ภิกฺขเว อริยสาวโก มิจฺฉา-อาชีวํ ปหาย สมฺมา-อาชีเวน ชีวิตํ กปฺเปติ.\csromanspacer{} อยํ วุจฺจติ ภิกฺขเว สมฺมา-อาชีโว.}

\prose{กตโม จ ภิกฺขเว สมฺมาวายาโม, อิธ ภิกฺขเว ภิกฺขุ อนุปฺปนฺนานํ ปาปกานํ อกุสลานํ ธมฺมานํ อนุปฺปาทาย ฉนฺทํ ชเนติ วายมติ วีริยํ อารภติ, จิตฺตํ ปคฺคณฺหาติ ปทหติ, อุปฺปนฺนานํ ปาปกานํ อกุสลานํ ธมฺมานํ ปหานาย ฉนฺทํ ชเนติ ฯเปฯ อนุปฺปนฺนานํ กุสลานํ ธมฺมานํ อุปฺปาทาย ฉนฺทํ ชเนติ ฯเปฯ อุปฺปนฺนานํ กุสลานํ ธมฺมานํ ฐิติยา อสมฺโมสาย ภิยฺโยภาวาย เวปุลฺลาย ภาวนาย ปาริปูริยา ฉนฺทํ ชเนติ วายมติ วีริยํ อารภติ, จิตฺตํ ปคฺคณฺหาติ ปทหติ.\csromanspacer{} อยํ วุจฺจติ ภิกฺขเว สมฺมาวายาโม.}

\prose{กตมา จ ภิกฺขเว สมฺมาสติ, อิธ ภิกฺขเว ภิกฺขุ กาเย กายานุปสฺสี วิหรติ อาตาปี สมฺปชาโน สติมา วิเนยฺย โลเก อภิชฺฌาโทมนสฺสํ, เวทนาสุ เวทนานุปสฺสี วิหรติ อาตาปี สมฺปชาโน สติมา วิเนยฺย โลเก อภิชฺฌาโทมนสฺสํ, จิตฺเต จิตฺตานุปสฺสี วิหรติ อาตาปี สมฺปชาโน สติมา วิเนยฺย โลเก อภิชฺฌาโทมนสฺสํ, ธมฺเมสุ ธมฺมานุปสฺสี วิหรติ อาตาปี สมฺปชาโน สติมา วิเนยฺย โลเก อภิชฺฌาโทมนสฺสํ.\csromanspacer{} อยํ วุจฺจติ ภิกฺขเว สมฺมาสติ.}

\prose{\csromanfolio{9}กตโม จ ภิกฺขเว สมฺมาสมาธิ, อิธ ภิกฺขเว ภิกฺขุ วิวิจฺเจว กาเมหิ วิวิจฺจ อกุสเลหิ ธมฺเมหิ สวิตกฺกํ สวิจารํ วิเวกชํ ปีติสุขํ ปฐมํ ฌานํ อุปสมฺปชฺช วิหรติ, วิตกฺกวิจารานํ วูปสมา อชฺฌตฺตํ สมฺปสาทนํ เจตโส เอโกทิภาวํ อวิตกฺกํ อวิจารํ สมาธิชํ ปีติสุขํ ทุติยํ ฌานํ อุปสมฺปชฺช วิหรติ, ปีติยา จ วิราคา อุเปกฺขโก จ วิหรติ, สโต จ สมฺปชาโน สุขญฺจ กาเยน ปฏิสํเวเทติ.\csromanspacer{} ยํ ตํ อริยา อาจิกฺขนฺติ “อุเปกฺขโก สติมา สุขวิหารี”ติ ตติยํ ฌานํ อุปสมฺปชฺช วิหรติ.\csromanspacer{} สุขสฺส จ ปหานา ทุกฺขสฺส จ ปหานา ปุพฺเพว โสมนสฺสโทมนสฺสานํ อตฺถงฺคมา อทุกฺขมสุขํ อุเปกฺขาสติปาริสุทฺธิํ จตุตฺถํ ฌานํ อุปสมฺปชฺช วิหรติ.\csromanspacer{} อยํ วุจฺจติ ภิกฺขเว สมฺมาสมาธีติ.\csromanspacer{}\csromanspacer{}อฏฺฐมํ.}

\csromanheader{2}{9. สูกสุตฺต}
\proseitem{9}{สาวตฺถินิทานํ.\csromanspacer{} เสยฺยถาปิ ภิกฺขเว สาลิสูกํ วา ยวสูกํ วา มิจฺฉาปณิหิตํ หตฺเถน วา ปาเทน วา อกฺกนฺตํ หตฺถํ วา ปาทํ วา ภินฺทิสฺสติ\footnote{เภจฺฉติ (ก)}, โลหิตํ วา อุปฺปาเทสฺสตีติ เนตํ ฐานํ วิชฺชติ.\csromanspacer{} ตํ กิสฺส เหตุ, มิจฺฉาปณิหิตตฺตา ภิกฺขเว สูกสฺส.\csromanspacer{} เอวเมว โข ภิกฺขเว โส วต ภิกฺขุ มิจฺฉาปณิหิตาย ทิฏฺฐิยา มิจฺฉาปณิหิตาย มคฺคภาวนาย อวิชฺชํ ภินฺทิสฺสติ, วิชฺชํ อุปฺปาเทสฺสติ, นิพฺพานํ สจฺฉิกริสฺสตีติ เนตํ ฐานํ วิชฺชติ.\csromanspacer{} ตํ กิสฺส เหตุ, มิจฺฉาปณิหิตตฺตา ภิกฺขเว ทิฏฺฐิยา.}

\prose{เสยฺยถาปิ ภิกฺขเว สาลิสูกํ วา ยวสูกํ วา สมฺมาปณิหิตํ หตฺเถน วา ปาเทน วา อกฺกนฺตํ หตฺถํ วา ปาทํ วา ภินฺทิสฺสติ, โลหิตํ วา อุปฺปาเทสฺสตีติ ฐานเมตํ วิชฺชติ.\csromanspacer{} ตํ กิสฺส เหตุ, สมฺมาปณิหิตตฺตา ภิกฺขเว สูกสฺส.\csromanspacer{} เอวเมว โข ภิกฺขเว โส วต ภิกฺขุ สมฺมาปณิหิตาย ทิฏฺฐิยา สมฺมาปณิหิตาย มคฺคภาวนาย อวิชฺชํ ภินฺทิสฺสติ, วิชฺชํ อุปฺปาเทสฺสติ, นิพฺพานํ สจฺฉิกริสฺสตีติ ฐานเมตํ วิชฺชติ.\csromanspacer{} ตํ กิสฺส เหตุ, สมฺมาปณิหิตตฺตา ภิกฺขเว ทิฏฺฐิยา.}

\prose{กถญฺจ ภิกฺขเว ภิกฺขุ สมฺมาปณิหิตาย ทิฏฺฐิยา สมฺมาปณิหิตาย มคฺคภาวนาย อวิชฺชํ ภินฺทติ, วิชฺชํ อุปฺปาเทติ, นิพฺพานํ สจฺฉิกโรตีติ.\csromanspacer{} อิธ ภิกฺขเว ภิกฺขุ สมฺมาทิฏฺฐิํ ภาเวติ วิเวกนิสฺสิตํ วิราคนิสฺสิตํ \csromanfolio{10}นิโรธนิสฺสิตํ โวสฺสคฺคปริณามิํ ฯเปฯ สมฺมาสมาธิํ ภาเวติ วิเวกนิสฺสิตํ วิราคนิสฺสิตํ นิโรธนิสฺสิตํ โวสฺสคฺคปริณามิํ.\csromanspacer{} เอวํ โข ภิกฺขเว ภิกฺขุ สมฺมาปณิหิตาย ทิฏฺฐิยา สมฺมาปณิหิตาย มคฺคภาวนาย อวิชฺชํ ภินฺทติ, วิชฺชํ อุปฺปาเทติ, นิพฺพานํ สจฺฉิกโรตีติ.\csromanspacer{}\csromanspacer{}นวมํ.}

\csromanheader{2}{10. นนฺทิยสุตฺต}
\proseitem{10}{สาวตฺถินิทานํ.\csromanspacer{} อถ โข นนฺทิโย ปริพฺพาชโก เยน ภควา เตนุปสงฺกมิ, อุปสงฺกมิตฺวา ภควตา สทฺธิํ สมฺโมทิ, สมฺโมทนียํ กถํ สารณียํ วีติสาเรตฺวา เอกมนฺตํ นิสีทิ, เอกมนฺตํ นิสินฺโน โข นนฺทิโย ปริพฺพาชโก ภควนฺตํ เอตทโวจ “กติ นุ โข โภ โคตม ธมฺมา ภาวิตา พหุลีกตา นิพฺพานงฺคมา โหนฺติ นิพฺพานปรายนา นิพฺพานปริโยสานา”ติ.}

\prose{อฏฺฐิเม โข นนฺทิย ธมฺมา ภาวิตา พหุลีกตา นิพฺพานงฺคมา โหนฺติ นิพฺพานปรายนา นิพฺพานปริโยสานา.\csromanspacer{} กตเม อฏฺฐ, เสยฺยถิทํ, สมฺมาทิฏฺฐิ ฯเปฯ สมฺมาสมาธิ.\csromanspacer{} อิเม โข นนฺทิย อฏฺฐ ธมฺมา ภาวิตา พหุลีกตา นิพฺพานงฺคมา โหนฺติ นิพฺพานปรายนา นิพฺพานปริโยสานาติ.\csromanspacer{} เอวํ วุตฺเต นนฺทิโย ปริพฺพาชโก ภควนฺตํ เอตทโวจ “อภิกฺกนฺตํ โภ โคตม, อภิกฺกนฺตํ โภ โคตม ฯเปฯ อุปาสกํ มํ ภวํ โคตโม ธาเรตุ อชฺชตคฺเค ปาณุเปตํ สรณํ คตนฺ”ติ.\csromanspacer{}\csromanspacer{}ทสมํ.\csromanspacer{} อวิชฺชาวคฺโค ปฐโม.}

\titlehead{\textbf{ตสฺสุทฺทานํ}}
\csromangathameasure{อวิชฺชญฺจ อุปฑฺฒญฺจ, สาริปุตฺโต จ พฺราหฺมโณ. \\ กิมตฺถิโย จ เทฺว ภิกฺขู, วิภงฺโค สูกนนฺทิยาติ.}
\csromangathagroup{\csromangathastanza{อวิชฺชญฺจ อุปฑฺฒญฺจ, สาริปุตฺโต จ พฺราหฺมโณ. \\ กิมตฺถิโย จ เทฺว ภิกฺขู, วิภงฺโค สูกนนฺทิยาติ.}}

\chapterhead{2. วิหารวคฺค}
\csromanheader{2}{1. ปฐมวิหารสุตฺต}
\proseitem{11}{สาวตฺถินิทานํ.\csromanspacer{} อิจฺฉามหํ ภิกฺขเว อฑฺฒมาสํ ปฏิสลฺลิยิตุํ, นมฺหิ เกนจิ อุปสงฺกมิตพฺโพ อญฺญตฺร เอเกน ปิณฺฑปาตนีหารเกนาติ. “เอวํ \csromanfolio{11}ภนฺเต”ติ โข เต นฺหิกฺขู ภควโต ปฏิสฺสุตฺวา นาสฺสุธ โกจิ ภควนฺตํ อุปสงฺกมติ อญฺญตฺร เอเกน ปิณฺฑปาตนีหารเกน.}

\prose{อถ โข ภควา ตสฺส อฑฺฒมาสสฺส อจฺจเยน ปฏิสลฺลานา วุฏฺฐิโต ภิกฺขู อามนฺเตสิ “เยน สฺวาหํ ภิกฺขเว วิหาเรน ปฐมาภิสมฺพุทฺโธ วิหรามิ, ตสฺส ปเทเสน วิหาสิํ, โส เอวํ ปชานามิ “มิจฺฉาทิฏฺฐิปจฺจยาปิ เวทยิตํ, สมฺมาทิฏฺฐิปจฺจยาปิ เวทยิตํ ฯเปฯ มิจฺฉาสมาธิปจฺจยาปิ เวทยิตํ, สมฺมาสมาธิปจฺจยาปิ เวทยิตํ, ฉนฺทปจฺจยาปิ เวทยิตํ, วิตกฺกปจฺจยาปิ เวทยิตํ, สญฺญาปจฺจยาปิ เวทยิตํ.\csromanspacer{} ฉนฺโท จ อวูปสนฺโต โหติ, วิตกฺโก จ อวูปสนฺโต โหติ, สญฺญา จ อวูปสนฺตา โหติ, ตปฺปจฺจยาปิ เวทยิตํ.\csromanspacer{} ฉนฺโท จ วูปสนฺโต โหติ, วิตกฺโก จ วูปสนฺโต โหติ, สญฺญา จ วูปสนฺตา โหติ, ตปฺปจฺจายาปิ เวทยิตํ.\csromanspacer{} อปฺปตฺตสฺส ปตฺติยา อตฺถิ อายามํ\footnote{วายามํ (สี, สฺยา)}, ตสฺมิมฺปิ ฐาเน อนุปฺปตฺเต ตปฺปจฺจยาปิ เวทยิตนฺ’ติ”.\csromanspacer{}\csromanspacer{}ปฐมํ.}

\csromanheader{2}{2. ทุติยวิหารสุตฺต}
\proseitem{12}{สาวตฺถินิทานํ.\csromanspacer{} อิจฺฉามหํ ภิกฺขเว เตมาสํ ปฏิสลฺลิยิตุํ, นมฺหิ เกนจิ อุปสงฺกมิตพฺโพ อญฺญตฺร เอเกน ปิณฺฑปาตนีหารเกนาติ. “เอวํ ภนฺเต”ติ โข เต ภิกฺขู ภควโต ปฏิสฺสุตฺวา นาสฺสุธ โกจิ ภควนฺตํ อุปสงฺกมติ อญฺญตฺร เอเกน ปิณฺฑปาตนีหารเกน.}

\prose{อถ โข ภควา ตสฺส เตมาสสฺส อจฺจเยน ปฏิสลฺลานา วุฏฺฐิโต ภิกฺขู อามนฺเตสิ “เยน สฺวาหํ ภิกฺขเว วิหาเรน ปฐมาภิสมฺพุทฺโธ วิหรามิ, ตสฺส ปเทเสน วิหาสิํ.\csromanspacer{} โส เอวํ ปชานามิ ‘มิจฺฉาทิฏฺฐิปจฺจยาปิ เวทยิตํ, มิจฺฉาทิฏฺฐิวูปสมปจฺจยาปิ เวทยิตํ, สมฺมาทิฏฺฐิปจฺจยาปิ เวทยิตํ, สมฺมาทิฏฺฐิวูปสมปจฺจยาปิ เวทยิตํ ฯเปฯ มิจฺฉาสมาธิปจฺจยาปิ เวทยิตํ, มิจฺฉาสมาธิวูปสมปจฺจยาปิ เวทยิตํ, สมฺมาสมาธิปจฺจยาปิ เวทยิตํ, สมฺมาสมาธิวูปสมปจฺจยาปิ เวทยิตํ, ฉนฺทปจฺจยาปิ เวทยิตํ, ฉนฺทวูปสมปจฺจยาปิ เวทยิตํ, วิตกฺกปจฺจยาปิ เวทยิตํ, วิตกฺกวูปสมปจฺจยาปิ เวทยิตํ, สญฺญาปจฺจยาปิ เวทยิตํ, สญฺญาวูปสมปจฺจยาปิ เวทยิตํ.\csromanspacer{} ฉนฺโท จ อวูปสนฺโต โหติ, วิตกฺโก จ อวูปสนฺโต โหติ, สญฺญา จ อวูปสนฺตา โหติ, ตปฺปจฺจยาปิ เวทยิตํ.\csromanspacer{} ฉนฺโท จ วูปสนฺโต \csromanfolio{12}โหติ, วิตกฺโก จ วูปสนฺโต โหติ, สญฺญา จ วูปสนฺตา โหติ, ตปฺปจฺจยาปิ เวทยิตํ.\csromanspacer{} อปฺปตฺตสฺส ปตฺติยา อตฺถิ อายามํ\footnote{วายามํ (สี, สฺยา)}, ตสฺมิมฺปิ ฐาเน อนุปฺปตฺเต ตปฺปจฺจยาปิ เวทยิตนฺ’ติ”.\csromanspacer{}\csromanspacer{}ทุติยํ.}

\csromanheader{2}{3. เสกฺขสุตฺต}
\proseitem{13}{สาวตฺถินิทานํ.\csromanspacer{} อถ โข อญฺญตโร ภิกฺขุ เยน ภควา เตนุปสงฺกมิ ฯเปฯ เอกมนฺตํ นิสินฺโน โข โส ภิกฺขุ ภควนฺตํ เอตทโวจ\csromandash{}“เสกฺโข เสกฺโข”ติ ภนฺเต วุจฺจติ.\csromanspacer{} กิตฺตาวตา นุ โข ภนฺเต เสกฺโข โหตีติ.\csromanspacer{} อิธ ภิกฺขุ เสกฺขาย สมฺมาทิฏฺฐิยา สมนฺนาคโต โหติ ฯเปฯ เสกฺเขน สมฺมาสมาธินา สมนฺนาคโต โหติ.\csromanspacer{} เอตฺตาวตา โข ภิกฺขุ เสกฺโข โหตีติ.\csromanspacer{}\csromanspacer{}ตติยํ.}

\csromanheader{2}{4. ปฐม-อุปฺปาทสุตฺต}
\proseitem{14}{สาวตฺถินิทานํ.\csromanspacer{} อฏฺฐิเม ภิกฺขเว ธมฺมา ภาวิตา พหุลีกตา อนุปฺปนฺนา อุปฺปชฺชนฺติ นาญฺญตฺร ตถาคตสฺส ปาตุภาวา อรหโต สมฺมาสมฺพุทฺธสฺส.\csromanspacer{} กตเม อฏฺฐ, เสยฺยถิทํ, สมฺมาทิฏฺฐิ ฯเปฯ สมฺมาสมาธิ.\csromanspacer{} อิเม โข ภิกฺขเว อฏฺฐ ธมฺมา ภาวิตา พหุลีกตา อนุปฺปนฺนา อุปฺปชฺชนฺติ นาญฺญตฺร ตถาคตสฺส ปาตุภาวา อรหโต สมฺมาสมฺพุทฺธสฺสาติ.\csromanspacer{}\csromanspacer{}จตุตฺถํ.}

\csromanheader{2}{5. ทุติย-อุปฺปาทสุตฺต}
\proseitem{15}{สาวตฺถินิทานํ.\csromanspacer{} อฏฺฐิเม ภิกฺขเว ธมฺมา ภาวิตา พหุลีกตา อนุปฺปนฺนา อุปฺปชฺชนฺติ นาญฺญตฺร สุคตวินยา.\csromanspacer{} กตเม อฏฺฐ, เสยฺยถิทํ, สมฺมาทิฏฺฐิ ฯเปฯ สมฺมาสมาธิ.\csromanspacer{} อิเม โข ภิกฺขเว อฏฺฐ ธมฺมา ภาวิตา พหุลีกตา อนุปฺปนฺนา อุปฺปชฺชนฺติ นาญฺญตฺร สุคตวินยาติ.\csromanspacer{}\csromanspacer{}ปญฺจมํ.}

\csromanheader{2}{6. ปฐมปริสุทฺธสุตฺต}
\proseitem{16}{สาวตฺถินิทานํ.\csromanspacer{} อฏฺฐิเม ภิกฺขเว ธมฺมา ปริสุทฺธา ปริโยทาตา อนงฺคณา วิคตูปกฺกิเลสา อนุปฺปนฺนา อุปฺปชฺชนฺติ นาญฺญตฺร ตถาคตสฺส ปาตุภาวา อรหโต สมฺมาสมฺพุทฺธสฺส.\csromanspacer{} กตเม อฏฺฐ, เสยฺยถิทํ, สมฺมาทิฏฺฐิ ฯเปฯ สมฺมาสมาธิ.\csromanspacer{} อิเม โข ภิกฺขเว อฏฺฐ ธมฺมา ปริสุทฺธา \csromanfolio{13}ปริโยทาตา อนงฺคณา วิคตูปกฺกิเลสา อนุปฺปนฺนา อุปฺปชฺชนฺติ นาญฺญตฺร ตถาคตสฺส ปาตุภาวา อรหโต สมฺมาสมฺพุทฺธสฺสาติ.\csromanspacer{}\csromanspacer{}ฉฏฺฐํ.}

\csromanheader{2}{7. ทุติยปริสุทฺธสุตฺต}
\proseitem{17}{สาวตฺถินิทานํ.\csromanspacer{} อฏฺฐิเม ภิกฺขเว ธมฺมา ปริสุทฺธา ปริโยทาตา อนงฺคณา วิคตูปกฺกิเลสา อนุปฺปนฺนา อุปฺปชฺชนฺติ นาญฺญตฺร สุคตวินยา.\csromanspacer{} กตเม อฏฺฐ, เสยฺยถิทํ, สมฺมาทิฏฺฐิ ฯเปฯ สมฺมาสมาธิ.\csromanspacer{} อิเม โข ภิกฺขเว อฏฺฐ ธมฺมา ปริสุทฺธา ปริโยทาตา อนงฺคณา วิคตูปกฺกิเลสา อนุปฺปนฺนา อุปฺปชฺชนฺติ นาญฺญตฺร สุคตวินยาติ.\csromanspacer{}\csromanspacer{}สตฺตมํ.}

\csromanheader{2}{8. ปฐมกุกฺกุฏารามสุตฺต}
\proseitem{18}{เอวํ เม สุตํ\csromandash{}เอกํ สมยํ อายสฺมา จ อานนฺโท อายสฺมา จ ภทฺโท ปาฏลิปุตฺเต วิหรนฺติ กุกฺกุฏาราเม.\csromanspacer{} อถ โข อายสฺมา ภทฺโท สายนฺหสมยํ ปฏิสลฺลานา วุฏฺฐิโต เยนายสฺมา อานนฺโท เตนุปสงฺกมิ, อุปสงฺกมิตฺวา อายสฺมตา อานนฺเทน สทฺธิํ สมฺโมทิ, สมฺโมทนียํ กถํ สารณียํ วีติสาเรตฺวา เอกมนฺตํ นิสีทิ, เอกมนฺตํ นิสินฺโน โข อายสฺมา ภทฺโท อายสฺมนฺตํ อานนฺทํ เอตทโวจ\csromandash{}}

\prose{“อพฺรหฺมจริยํ อพฺรหฺมจริยนฺ”ติ อาวุโส อานนฺท วุจฺจติ, กตมํ นุ โข อาวุโส อพฺรหฺมจริยนฺติ.\csromanspacer{} สาธุ สาธุ อาวุโส ภทฺท, ภทฺทโก โข เต อาวุโส ภทฺท อุมฺมงฺโค, ภทฺทกํ ปฏิภานํ, กลฺยาณี ปริปุจฺฉา, เอวํ หิ ตฺวํ อาวุโส ภทฺท ปุจฺฉสิ “อพฺรหฺมจริยํ อพฺรหฺมจริยนฺติ อาวุโส อานนฺท วุจฺจติ, กตมํ นุ โข อาวุโส อพฺรหฺมจริยนฺ”ติ.\csromanspacer{} เอวมาวุโสติ.\csromanspacer{} อยเมว โข อาวุโส อฏฺฐงฺคิโก มิจฺฉามคฺโค อพฺรหฺมจริยํ.\csromanspacer{} เสยฺยถิทํ, มิจฺฉาทิฏฺฐิ ฯเปฯ มิจฺฉาสมาธีติ.\csromanspacer{}\csromanspacer{}อฏฺฐมํ.}

\csromanheader{2}{9. ทุติยกุกฺกุฏารามสุตฺต}
\proseitem{19}{ปาฏลิปุตฺตนิทานํ. “พฺรหฺมจริยํ พฺรหฺมจริยนฺ”ติ อาวุโส อานนฺท วุจฺจติ, กตมํ นุ โข อาวุโส พฺรหฺมจริยํ, กตมํ พฺรหฺมจริยปริโยสานนฺติ.\csromanspacer{} สาธุ สาธุ อาวุโส ภทฺท, ภทฺทโก โข เต อาวุโส ภทฺท อุมฺมงฺโค, ภทฺทกํ ปฏิภานํ, กลฺยาณี ปริปุจฺฉา, เอวํ หิ ตฺวํ อาวุโส ภทฺท ปุจฺฉสิ “พฺรหฺมจริยํ พฺรหฺมจริยนฺติ อาวุโส อานนฺท วุจฺจติ, กตมํ นุ โข \csromanfolio{14}อาวุโส พฺรหฺมจริยํ, กตมํ พฺรหฺมจริยปริโยสานนฺ”ติ.\csromanspacer{} เอวมาวุโสติ.\csromanspacer{} อยเมว โข อาวุโส อริโย อฏฺฐงฺคิโก มคฺโค พฺรหฺมจริยํ.\csromanspacer{} เสยฺยถิทํ, สมฺมาทิฏฺฐิ ฯเปฯ สมฺมาสมาธิ.\csromanspacer{} โย โข อาวุโส ราคกฺขโย โทสกฺขโย โมหกฺขโย, อิทํ พฺรหฺมจริยปริโยสานนฺติ.\csromanspacer{}\csromanspacer{}นวมํ.}

\csromanheader{2}{10. ตติยกุกฺกุฏารามสุตฺต}
\proseitem{20}{ปาฏลิปุตฺตนิทานํ. “พฺรหฺมจริยํ พฺรหฺมจริยนฺ”ติ อาวุโส อานนฺท วุจฺจติ, กตมํ นุ โข อาวุโส พฺรหฺมจริยํ, กตโม พฺรหฺมจารี, กตมํ พฺรหฺมจริยปริโยสานนฺติ.\csromanspacer{} สาธุ สาธุ อาวุโส ภทฺท, ภทฺทโก โข เต อาวุโส ภทฺท อุมฺมงฺโค, ภทฺทกํ ปฏิภานํ, กลฺยาณี ปริปุจฺฉา, เอวํ หิ ตฺวํ อาวุโส ภทฺท ปุจฺฉสิ “พฺรหฺมจริยํ พฺรหฺมจริยนฺติ อาวุโส อานนฺท วุจฺจติ, กตมํ นุ โข อาวุโส พฺรหฺมจริยํ, กตโม พฺรหฺมจารี, กตมํ พฺรหฺมจริยปริโยสานนฺ”ติ.\csromanspacer{} เอวมาวุโสติ.\csromanspacer{} อยเมว โข อาวุโส อริโย อฏฺฐงฺคิโก มคฺโค พฺรหฺมจริยํ.\csromanspacer{} เสยฺยถิทํ, สมฺมาทิฏฺฐิ ฯเปฯ สมฺมาสมาธิ.\csromanspacer{} โย โข อาวุโส อิมินา อริเยน อฏฺฐงฺคิเกน มคฺเคน สมนฺนาคโต, อยํ วุจฺจติ พฺรหฺมจารี.\csromanspacer{} โย โข อาวุโส ราคกฺขโย โทสกฺขโย โมหกฺขโย, อิทํ พฺรหฺมจริยปริโยสานนฺติ.\csromanspacer{}\csromanspacer{}ทสมํ.\csromanspacer{} ตีณิ สุตฺตนฺตานิ เอกนิทานานิ.\csromanspacer{} วิหารวคฺโค ทุติโย.}

\titlehead{\textbf{ตสฺสุทฺทานํ}}
\csromangathameasure{เทฺว วิหารา จ เสกฺโข จ, อุปฺปาทา อปเร ทุเว. \\ ปริสุทฺเธน เทฺว วุตฺตา, กุกฺกุฏาราเมน ตโยติ.}
\csromangathagroup{\csromangathastanza{เทฺว วิหารา จ เสกฺโข จ, อุปฺปาทา อปเร ทุเว. \\ ปริสุทฺเธน เทฺว วุตฺตา, กุกฺกุฏาราเมน ตโยติ.}}

\chapterhead{3. มิจฺฉตฺตวคฺค}
\csromanheader{2}{1. มิจฺฉตฺตสุตฺต}
\proseitem{21}{สาวตฺถินิทานํ.\csromanspacer{} มิจฺฉตฺตญฺจ โว ภิกฺขเว เทเสสฺสามิ สมฺมตฺตญฺจ, ตํ สุณาถ.\csromanspacer{} กตมญฺจ ภิกฺขเว มิจฺฉตฺตํ, เสยฺยถิทํ, มิจฺฉทิฏฺฐิ ฯเปฯ \csromanfolio{15}มิจฺฉาสมาธิ.\csromanspacer{} อิทํ วุจฺจติ ภิกฺขเว มิจฺฉตฺตํ.\csromanspacer{} กตมญฺจ ภิกฺขเว สมฺมตฺตํ, เสยฺยถิทํ, สมฺมาทิฏฺฐิ (ป) สมฺมาสมาธิ.\csromanspacer{} อิทํ วุจฺจติ ภิกฺขเว สมฺมตฺตนฺติ.\csromanspacer{}\csromanspacer{}ปฐมํ.}

\csromanheader{2}{2. อกุสลธมฺมสุตฺต}
\proseitem{22}{สาวตฺถินิทานํ.\csromanspacer{} อกุสเล จ โข ภิกฺขเว ธมฺเม เทเสสฺสามิ กุสเล จ ธมฺเม, ตํ สุณาถ.\csromanspacer{} กตเม จ ภิกฺขเว อกุสลา ธมฺมา, เสยฺยถิทํ, มิจฺฉาทิฏฺฐิ ฯเปฯ มิจฺฉาสมาธิ.\csromanspacer{} อิเม วุจฺจนฺติ ภิกฺขเว อกุสลา ธมฺมา.\csromanspacer{} กตเม จ ภิกฺขเว กุสลา ธมฺมา, เสยฺยถิทํ, สมฺมาทิฏฺฐิ ฯเปฯ สมฺมาสมาธิ.\csromanspacer{} อิเม วุจฺจนฺติ ภิกฺขเว กุสลา ธมฺมาติ.\csromanspacer{}\csromanspacer{}ทุติยํ.}

\csromanheader{2}{3. ปฐมปฏิปทาสุตฺต}
\proseitem{23}{สาวตฺถินิทานํ.\csromanspacer{} มิจฺฉาปฏิปทญฺจ โว ภิกฺขเว เทเสสฺสามิ สมฺมาปฏิปทญฺจ, ตํ สุณาถ.\csromanspacer{} กตมา จ ภิกฺขเว มิจฺฉาปฏิปทา, เสยฺยถิทํ, มิจฺฉาทิฏฺฐิ ฯเปฯ มิจฺฉาสมาธิ.\csromanspacer{} อยํ วุจฺจติ ภิกฺขเว มิจฺฉาปฏิปทา.\csromanspacer{} กตมา จ ภิกฺขเว สมฺมาปฏิปทา, เสยฺยถิทํ, สมฺมาทิฏฺฐิ ฯเปฯ สมฺมาสมาธิ.\csromanspacer{} อยํ วุจฺจติ ภิกฺขเว สมฺมาปฏิปทาติ.\csromanspacer{}\csromanspacer{}ตติยํ.}

\csromanheader{2}{4. ทุติยปฏิปทาสุตฺต}
\proseitem{24}{สาวตฺถินิทานํ.\csromanspacer{} คิหิโน วาหํ ภิกฺขเว ปพฺพชิตสฺส วา มิจฺฉาปฏิปทํ น วณฺเณมิ, คิหิ วา ภิกฺขเว ปพฺพชิโต วา มิจฺฉาปฏิปนฺโน มิจฺฉาปฏิปตฺตาธิกรณเหตุ นาราธโก โหติ ญายํ ธมฺมํ กุสลํ.\csromanspacer{} กตมา จ ภิกฺขเว มิจฺฉาปฏิปทา, เสยฺยถิทํ, มิจฺฉาทิฏฺฐิ (ป) มิจฺฉาสมาธิ.\csromanspacer{} อยํ วุจฺจติ ภิกฺขเว มิจฺฉาปฏิปทา.\csromanspacer{} คิหิโน วาหํ ภิกฺขเว ปพฺพชิตสฺส วา มิจฺฉาปฏิปทํ น วณฺเณมิ, คิหิ วา ภิกฺขเว ปพฺพชิโต วา มิจฺฉาปฏิปนฺโน มิจฺฉาปฏิปตฺตาธิกรณเหตุ นาราธโก โหติ ญายํ ธมฺมํ กุสลํ.}

\prose{คิหิโน วาหํ ภิกฺขเว ปพฺพชิตสฺส วา สมฺมาปฏิปทํ วณฺเณมิ, คิหิ วา ภิกฺขเว ปพฺพชิโต วา สมฺมาปฏิปนฺโน สมฺมาปฏิปตฺตาธิกรณเหตุ อาราธโก โหติ ญายํ ธมฺมํ กุสลํ.\csromanspacer{} กตมา จ ภิกฺขเว สมฺมาปฏิปทา, เสยฺยถิทํ, สมฺมาทิฏฺฐิ (ป) สมฺมาสมาธิ.\csromanspacer{} อยํ วุจฺจติ ภิกฺขเว สมฺมาปฏิปทา.\csromanspacer{} คิหิโน วาหํ ภิกฺขเว ปพฺพชิตสฺส วา สมฺมาปฏิปทํ วณฺเณมิ, คิหิ วา \csromanfolio{16}ภิกฺขเว ปพฺพชิโต วา สมฺมาปฏิปนฺโน สมฺมาปฏิปตฺตาธิกรณเหตุ อาราธโก โหติ ญายํ ธมฺมํ กุสลนฺติ.\csromanspacer{}\csromanspacer{}จตุตฺถํ.}

\csromanheader{2}{5. ปฐม-อสปฺปุริสสุตฺต}
\proseitem{25}{สาวตฺถินิทานํ.\csromanspacer{} อสปฺปุริสญฺจ โว ภิกฺขเว เทเสสฺสามิ สปฺปุริสญฺจ, ตํ สุณาถ.\csromanspacer{} กตโม จ ภิกฺขเว อสปฺปุริโส, อิธ ภิกฺขเว เอกจฺโจ มิจฺฉาทิฏฺฐิโก โหติ มิจฺฉาสงฺกปฺโป มิจฺฉาวาโจ มิจฺฉากมฺมนฺโต มิจฺฉา-อาชีโว มิจฺฉาวายาโม มิจฺฉาสติ มิจฺฉาสมาธิ.\csromanspacer{} อยํ วุจฺจติ ภิกฺขเว อสปฺปุริโส.}

\prose{กตโม จ ภิกฺขเว สปฺปุริโส, อิธ ภิกฺขเว เอกจฺโจ สมฺมาทิฏฺฐิโก โหติ สมฺมาสงฺกปฺโป สมฺมาวาโจ สมฺมากมฺมนฺโต สมฺมา-อาชีโว สมฺมาวายาโม สมฺมาสติ สมฺมาสมาธิ.\csromanspacer{} อยํ วุจฺจติ ภิกฺขเว สปฺปุริโสติ.\csromanspacer{}\csromanspacer{}ปญฺจมํ.}

\csromanheader{2}{6. ทุติย-อสปฺปุริสสุตฺต}
\proseitem{26}{สาวตฺถินิทานํ.\csromanspacer{} อสปฺปุริสญฺจ โว ภิกฺขเว เทเสสฺสามิ อสปฺปุริเสน อสปฺปุริสตรญฺจ, สปฺปุริสญฺจ โว ภิกฺขเว เทเสสฺสามิ สปฺปุริเสน สปฺปุริสตรญฺจ, ตํ สุณาถ.\csromanspacer{} กตโม จ ภิกฺขเว อสปฺปุริโส, อิธ ภิกฺขเว เอกจฺโจ มิจฺฉาทิฏฺฐิโก โหติ ฯเปฯ มิจฺฉาสมาธิ.\csromanspacer{} อยํ วุจฺจติ ภิกฺขเว อสปฺปุริโส.}

\prose{กตโม จ ภิกฺขเว อสปฺปุริเสน อสปฺปุริสตโร, อิธ ภิกฺขเว เอกจฺโจ มิจฺฉาทิฏฺฐิโก โหติ ฯเปฯ มิจฺฉาสมาธิ มิจฺฉาญาณี มิจฺฉาวิมุตฺติ.\csromanspacer{} อยํ วุจฺจติ ภิกฺขเว อสปฺปุริเสน อสปฺปุริสตโร.}

\prose{กตโม จ ภิกฺขเว สปฺปุริโส, อิธ ภิกฺขเว เอกจฺโจ สมฺมาทิฏฺฐิโก โหติ ฯเปฯ สมฺมาสมาธิ.\csromanspacer{} อยํ วุจฺจติ ภิกฺขเว สปฺปุริโส.}

\prose{กตโม จ ภิกฺขเว สปฺปุริเสน สปฺปุริสตโร, อิธ ภิกฺขเว เอกจฺโจ สมฺมาทิฏฺฐิโก โหติ ฯเปฯ สมฺมาสมาธิ สมฺมาญาณี สมฺมาวิมุตฺติ.\csromanspacer{} อยํ วุจฺจติ ภิกฺขเว สปฺปุริเสน สปฺปุริสตโรติ.\csromanspacer{}\csromanspacer{}ฉฏฺฐํ.}

\csromanheader{2}{7. กุมฺภสุตฺต}
\proseitem{27}{\csromanfolio{17}สาวตฺถินิทานํ.\csromanspacer{} เสยฺยถาปิ ภิกฺขเว กุมฺโภ อนาธาโร สุปฺปวตฺติโย โหติ, สาธาโร ทุปฺปวตฺติโย โหติ.\csromanspacer{} เอวเมว โข ภิกฺขเว จิตฺตํ อนาธารํ สุปฺปวตฺติยํ โหติ, สาธารํ ทุปฺปวตฺติยํ โหติ.\csromanspacer{} โก จ ภิกฺขเว จิตฺตสฺส อาธาโร, อยเมว อริโย อฏฺฐงฺคิโก มคฺโค.\csromanspacer{} เสยฺยถิทํ, สมฺมาทิฏฺฐิ ฯเปฯ สมฺมาสมาธิ.\csromanspacer{} อยํ จิตฺตสฺส อาธาโร.\csromanspacer{} เสยฺยถาปิ ภิกฺขเว กุมฺโภ อนาธาโร สุปฺปวตฺติโย โหติ, สาธาโร ทุปฺปวตฺติโย โหติ.\csromanspacer{} เอวเมว โข ภิกฺขเว จิตฺตํ อนาธารํ สุปฺปวตฺติยํ โหติ, สาธารํ ทุปฺปวตฺติยํ โหตีติ.\csromanspacer{}\csromanspacer{}สตฺตมํ.}

\csromanheader{2}{8. สมาธิสุตฺต}
\proseitem{28}{สาวตฺถินิทานํ.\csromanspacer{} อริยํ โว ภิกฺขเว สมฺมาสมาธิํ เทเสสฺสามิ ส-อุปนิสํ สปริกฺขารํ, ตํ สุณาถ.\csromanspacer{} กตโม จ ภิกฺขเว อริโย สมฺมาสมาธิ ส-อุปนิโส สปริกฺขาโร.\csromanspacer{} เสยฺยถิทํ, สมฺมาทิฏฺฐิ ฯเปฯ สมฺมาสติ\footnote{สมฺมาสมาธิ (สี, สฺยา, กํ, ก)}.\csromanspacer{} ยา โข ภิกฺขเว อิเมหิ สตฺตหงฺเคหิ จิตฺตสฺส เอกคฺคตา สปริกฺขารตา\footnote{สปริกฺขตา (สี, อิ)}.\csromanspacer{} อยํ วุจฺจติ ภิกฺขเว อริโย สมฺมาสมาธิ “ส-อุปนิโส” อิติปิ “สปริกฺขาโร” อิติปีติ.\csromanspacer{}\csromanspacer{}อฏฺฐมํ.}

\csromanheader{2}{9. เวทนาสุตฺต}
\proseitem{29}{สาวตฺถินิทานํ.\csromanspacer{} ติสฺโส อิมา ภิกฺขเว เวทนา.\csromanspacer{} กตมา ติสฺโส, สุขา เวทนา ทุกฺขา เวทนา อทุกฺขมสุขา เวทนา.\csromanspacer{} อิมา โข ภิกฺขเว ติสฺโส เวทนา.\csromanspacer{} อิมาสํ โข ภิกฺขเว ติสฺสนฺนํ เวทนานํ ปริญฺญาย อริโย อฏฺฐงฺคิโก มคฺโค ภาเวตพฺโพ.\csromanspacer{} กตโม อริโย อฏฺฐงฺคิโก มคฺโค, เสยฺยถิทํ, สมฺมาทิฏฺฐิ ฯเปฯ สมฺมาสมาธิ.\csromanspacer{} อิมาสํ โข ภิกฺขเว ติสฺสนฺนํ เวทนานํ ปริญฺญาย อริโย อฏฺฐงฺคิโก มคฺโค ภาเวตพฺโพติ.\csromanspacer{}\csromanspacer{}นวมํ.}

\csromanheader{2}{10. อุตฺติยสุตฺต}
\proseitem{30}{สาวตฺถินิทานํ.\csromanspacer{} อถ โข อายสฺมา อุตฺติโย เยน ภควา เตนุปสงฺกมิ ฯเปฯ เอกมนฺตํ นิสินฺโน โข อายสฺมา อุตฺติโย ภควนฺตํ เอตทโวจ “อิธ มยฺหํ ภนฺเต รโหคตสฺส ปฏิสลฺลีนสฺส เอวํ เจตโส \csromanfolio{18}ปริวิตกฺโก อุทปาทิ ‘ปญฺจ กามคุณา วุตฺตา ภควตา, กตเม นุ โข ปญฺจ กามคุณา วุตฺตา ภควตา’ติ”.\csromanspacer{} สาธุ สาธุ อุตฺติย, ปญฺจิเม โข อุตฺติย กามคุณา วุตฺตา มยา.\csromanspacer{} กตเม ปญฺจ, จกฺขุวิญฺเญยฺยา รูปา อิฏฺฐา กนฺตา มนาปา ปิยรูปา กามูปสํหิตา รชนียา.\csromanspacer{} โสตวิญฺเญยฺยา สทฺทา ฯเปฯ.\csromanspacer{} ฆานวิญฺเญยฺยา คนฺธา ฯเปฯ.\csromanspacer{} ชิวฺหาวิญฺเญยฺยา รสา ฯเปฯ.\csromanspacer{} กายวิญฺเญยฺยา โผฏฺฐพฺพา อิฏฺฐา กนฺตา มนาปา ปิยรูปา กามูปสํหิตา รชนียา.\csromanspacer{} อิเม โข อุตฺติย ปญฺจ กามคุณา วุตฺตา มยา.\csromanspacer{} อิเมสํ โข อุตฺติย ปญฺจนฺนํ กามคุณานํ ปหานาย อริโย อฏฺฐงฺคิโก มคฺโค ภาเวตพฺโพ.\csromanspacer{} กตโม อริโย อฏฺฐงฺคิโก มคฺโค, เสยฺยถิทํ, สมฺมาทิฏฺฐิ ฯเปฯ สมฺมาสมาธิ.\csromanspacer{} อิเมสํ โข อุตฺติย ปญฺจนฺนํ กามคุณานํ ปหานาย อยํ อริโย อฏฺฐงฺคิโก มคฺโค ภาเวตพฺโพติ.\csromanspacer{}\csromanspacer{}ทสมํ.\csromanspacer{} มิจฺฉตฺตวคฺโค ตติโย.}

\titlehead{\textbf{ตสฺสุทฺทานํ}}
\csromangathameasure{มิจฺฉตฺตํ อกุสลํ ธมฺมํ, ทุเว ปฏิปทาปิ จ. \\ อสปฺปุริเสน เทฺว กุมฺโภ, สมาธิ เวทนุตฺติเยนาติ.}
\csromangathagroup{\csromangathastanza{มิจฺฉตฺตํ อกุสลํ ธมฺมํ, ทุเว ปฏิปทาปิ จ. \\ อสปฺปุริเสน เทฺว กุมฺโภ, สมาธิ เวทนุตฺติเยนาติ.}}

\chapterhead{4. ปฏิปตฺติวคฺค}
\csromanheader{2}{1. ปฐมปฏิปตฺติสุตฺต}
\proseitem{31}{สาวตฺถินิทานํ.\csromanspacer{} มิจฺฉาปฏิปตฺติญฺจ โว ภิกฺขเว เทเสสฺสามิ สมฺมาปฏิปตฺติญฺจ, ตํ สุณาถ.\csromanspacer{} กตมา จ ภิกฺขเว มิจฺฉาปฏิปตฺติ, เสยฺยถิทํ, มิจฺฉาทิฏฺฐิ ฯเปฯ มิจฺฉาสมาธิ.\csromanspacer{} อยํ วุจฺจติ ภิกฺขเว มิจฺฉาปฏิปตฺติ.\csromanspacer{} กตมา จ ภิกฺขเว สมฺมาปฏิปตฺติ, เสยฺยถิทํ, สมฺมาทิฏฺฐิ ฯเปฯ สมฺมาสมาธิ.\csromanspacer{} อยํ วุจฺจติ ภิกฺขเว สมฺมาปฏิปตฺตีติ.\csromanspacer{}\csromanspacer{}ปฐมํ.}

\csromanheader{2}{2. ทุติยปฏิปตฺติสุตฺต}
\proseitem{32}{สาวตฺถินิทานํ.\csromanspacer{} มิจฺฉาปฏิปนฺนญฺจ โว ภิกฺขเว เทเสสฺสามิ สมฺมาปฏิปนฺนญฺจ, ตํ สุณาถ.\csromanspacer{} กตโม จ ภิกฺขเว มิจฺฉาปฏิปนฺโน, อิธ ภิกฺขเว เอกจฺโจ มิจฺฉาทิฏฺฐิโก โหติ ฯเปฯ มิจฺฉาสมาธิ.\csromanspacer{} อยํ วุจฺจติ \csromanfolio{19}ภิกฺขเว มิจฺฉาปฏิปนฺโน.\csromanspacer{} กตโม จ ภิกฺขเว สมฺมาปฏิปนฺโน, อิธ ภิกฺขเว เอกจฺโจ สมฺมาทิฏฺฐิโก โหติ ฯเปฯ สมฺมาสมาธิ.\csromanspacer{} อยํ วุจฺจติ ภิกฺขเว สมฺมาปฏิปนฺโนติ.\csromanspacer{}\csromanspacer{}ทุติยํ.}

\csromanheader{2}{3. วิรทฺธสุตฺต}
\proseitem{33}{สาวตฺถินิทานํ.\csromanspacer{} เยสํ เกสญฺจิ ภิกฺขเว อริโย อฏฺฐงฺคิโก มคฺโค วิรทฺโธ, วิรทฺโธ เตสํ อริโย อฏฺฐงฺคิโก มคฺโค สมฺมา ทุกฺขกฺขยคามี.\csromanspacer{} เยสํ เกสญฺจิ ภิกฺขเว อริโย อฏฺฐงฺคิโก มคฺโค อารทฺโธ, อารทฺโธ เตสํ อริโย อฏฺฐงฺคิโก มคฺโค สมฺมา ทุกฺขกฺขยคามี.\csromanspacer{} กตโม จ ภิกฺขเว อริโย อฏฺฐงฺคิโก มคฺโค, เสยฺยถิทํ, สมฺมาทิฏฺฐิ ฯเปฯ สมฺมาสมาธิ.\csromanspacer{} เยสํ เกสญฺจิ ภิกฺขเว อยํ อริโย อฏฺฐงฺคิโก มคฺโค วิรทฺโธ, วิรทฺโธ เตสํ อริโย อฏฺฐงฺคิโก มคฺโค สมฺมา ทุกฺขกฺขยคามี.\csromanspacer{} เยสํ เกสญฺจิ ภิกฺขเว อยํ อริโย อฏฺฐงฺคิโก มคฺโค อารทฺโธ, อารทฺโธ เตสํ อริโย อฏฺฐงฺคิโก มคฺโค สมฺมา ทุกฺขกฺขยคามีติ.\csromanspacer{}\csromanspacer{}ตติยํ.}

\csromanheader{2}{4. ปารํคมสุตฺต}
\proseitem{34}{สาวตฺถินิทานํ.\csromanspacer{} อฏฺฐิเม ภิกฺขเว ธมฺมา ภาวิตา พหุลีกตา อปารา ปารํ คมนาย สํวตฺตนฺติ.\csromanspacer{} กตเม อฏฺฐ? เสยฺยถิทํ, สมฺมาทิฏฺฐิ ฯเปฯ สมฺมาสมาธิ.\csromanspacer{} อิเม โข ภิกฺขเว อฏฺฐ ธมฺมา ภาวิตา พหุลีกตา อปารา ปารํ คมนาย สํวตฺตนฺตีติ.}

\prose{อิทมโวจ ภควา, อิทํ วตฺวาน สุคโต อถาปรํ เอตทโวจ สตฺถา\csromandash{}}

\csromangathameasure{“อปฺปกา เต มนุสฺเสสุ, เย ชนา ปารคามิโน. \\ อถายํ อิตรา ปชา, ตีรเมวานุธาวติ. \\ เย จ โข สมฺมทกฺขาเต, ธมฺเม ธมฺมานุวตฺติโน. \\ เต ชนา ปารเมสฺสนฺติ, มจฺจุเธยฺยํ สุทุตฺตรํ. \\ กณฺหํ ธมฺมํ วิปฺปหาย, สุกฺกํ ภาเวถ ปณฺฑิโต. \\ โอกา อโนกมาคมฺม, วิเวเก ยตฺถ ทูรมํ. \\ ตตฺราภิรติมิจฺเฉยฺย, หิตฺวา กาเม อกิญฺจโน. \\ ปริโยทเปยฺย อตฺตานํ, จิตฺตเกฺลเสหิ ปณฺฑิโต. \\ เยสํ สมฺโพธิยงฺเคสุ, สมฺมา จิตฺตํ สุภาวิตํ. \\ อาทานปฏินิสฺสคฺเค, อนุปาทาย เย รตา. \\ ขีณาสวา ชุติมนฺโต, เต โลเก ปรินิพฺพุตา”ติ. จตุตฺถํ.}
\csromangathagroup{\csromangathastanza{“อปฺปกา เต มนุสฺเสสุ, เย ชนา ปารคามิโน. \\ อถายํ อิตรา ปชา, ตีรเมวานุธาวติ.}\csromangathastanza{เย จ โข สมฺมทกฺขาเต, ธมฺเม ธมฺมานุวตฺติโน. \\ เต ชนา ปารเมสฺสนฺติ, มจฺจุเธยฺยํ สุทุตฺตรํ.}\csromangathastanza{กณฺหํ ธมฺมํ วิปฺปหาย, สุกฺกํ ภาเวถ ปณฺฑิโต. \\ โอกา อโนกมาคมฺม, วิเวเก ยตฺถ ทูรมํ.}\csromangathastanza{ตตฺราภิรติมิจฺเฉยฺย, หิตฺวา กาเม อกิญฺจโน. \\ ปริโยทเปยฺย อตฺตานํ, จิตฺตเกฺลเสหิ ปณฺฑิโต.}\csromangathastanza{\csromanfolio{20}เยสํ สมฺโพธิยงฺเคสุ, สมฺมา จิตฺตํ สุภาวิตํ. \\ อาทานปฏินิสฺสคฺเค, อนุปาทาย เย รตา. \\ ขีณาสวา ชุติมนฺโต, เต โลเก ปรินิพฺพุตา”ติ.\csromanspacer{} จตุตฺถํ.}}

\csromanheader{2}{5. ปฐมสามญฺญสุตฺต}
\proseitem{35}{สาวตฺถินิทานํ.\csromanspacer{} สามญฺญญฺจ โว ภิกฺขเว เทเสสฺสามิ สามญฺญผลานิ จ, ตํ สุณาถ.\csromanspacer{} กตมญฺจ ภิกฺขเว สามญฺญํ, อยเมว อริโย อฏฺฐงฺคิโก มคฺโค, เสยฺยถิทํ, สมฺมาทิฏฺฐิ ฯเปฯ สมฺมาสมาธิ.\csromanspacer{} อิทํ วุจฺจติ ภิกฺขเว สามญฺญํ.\csromanspacer{} กตมานิ จ ภิกฺขเว สามญฺญผลานิ, โสตาปตฺติผลํ สกทาคามิผลํ อนาคามิผลํ อรหตฺตผลํ.\csromanspacer{} อิมานิ วุจฺจนฺติ ภิกฺขเว สามญฺญผลานีติ.\csromanspacer{}\csromanspacer{}ปญฺจมํ.}

\csromanheader{2}{6. ทุติยสามญฺญสุตฺต}
\proseitem{36}{สาวตฺถินิทานํ.\csromanspacer{} สามญฺญญฺจ โว ภิกฺขเว เทเสสฺสามิ สามญฺญตฺถญฺจ, ตํ สุณาถ.\csromanspacer{} กตมญฺจ โข ภิกฺขเว สามญฺญํ, อยเมว อริโย อฏฺฐงฺคิโก มคฺโค, เสยฺยถิทํ, สมฺมาทิฏฺฐิ ฯเปฯ สมฺมาสมาธิ.\csromanspacer{} อิทํ วุจฺจติ ภิกฺขเว สามญฺญํ.\csromanspacer{} กตโม จ ภิกฺขเว สามญฺญตฺโถ, โย โข ภิกฺขเว ราคกฺขโย โทสกฺขโย โมหกฺขโย.\csromanspacer{} อยํ วุจฺจติ ภิกฺขเว สามญฺญตฺโถติ.\csromanspacer{}\csromanspacer{}ฉฏฺฐํ.}

\csromanheader{2}{7. ปฐมพฺรหฺมญฺญสุตฺต}
\proseitem{37}{สาวตฺถินิทานํ.\csromanspacer{} พฺรหฺมญฺญญฺจ โว ภิกฺขเว เทเสสฺสามิ พฺรหฺมญฺญผลานิ จ, ตํ สุณาถ.\csromanspacer{} กตมญฺจ โข ภิกฺขเว พฺรหฺมญฺญํ, อยเมว อริโย อฏฺฐงฺคิโก มคฺโค, เสยฺยถิทํ, สมฺมาทิฏฺฐิ ฯเปฯ สมฺมาสมาธิ.\csromanspacer{} อิทํ วุจฺจติ ภิกฺขเว พฺรหฺมญฺญํ.\csromanspacer{} กตมานิ จ ภิกฺขเว พฺรหฺมญฺญผลานิ, โสตาปตฺติผลํ สกทาคามิผลํ อนาคามิผลํ อรหตฺตผลํ.\csromanspacer{} อิมานิ วุจฺจนฺติ ภิกฺขเว พฺรหฺมญฺญผลานีติ.\csromanspacer{}\csromanspacer{}สตฺตมํ.}

\csromanheader{2}{8. ทุติยพฺรหฺมญฺญสุตฺต}
\proseitem{38}{สาวตฺถินิทานํ.\csromanspacer{} พฺรหฺมญฺญญฺจ โว ภิกฺขเว เทเสสฺสามิ พฺรหฺมญฺญตฺถญฺจ, ตํ สุณาถ.\csromanspacer{} กตมญฺจ ภิกฺขเว พฺรหฺมญฺญํ, อยเมว อริโย อฏฺฐงฺคิโก มคฺโค, เสยฺยถิทํ, สมฺมาทิฏฺฐิ ฯเปฯ สมฺมาสมาธิ.\csromanspacer{} อิทํ วุจฺจติ ภิกฺขเว พฺรหฺมญฺญํ. \csromanfolio{21}กตโม จ ภิกฺขเว พฺรหฺมญฺญตฺโถ, โย โข ภิกฺขเว ราคกฺขโย โทสกฺขโย โมหกฺขโย.\csromanspacer{} อยํ วุจฺจติ ภิกฺขเว พฺรหฺมญฺญตฺโถติ.\csromanspacer{}\csromanspacer{}อฏฺฐมํ.}

\csromanheader{2}{9. ปฐมพฺรหฺมจริยสุตฺต}
\proseitem{39}{สาวตฺถินิทานํ.\csromanspacer{} พฺรหฺมจริยญฺจ โว ภิกฺขเว เทเสสฺสามิ พฺรหฺมจริยผลานิ จ, ตํ สุณาถ.\csromanspacer{} กตมญฺจ ภิกฺขเว พฺรหฺมจริยํ.\csromanspacer{} อยเมว อริโย อฏฺฐงฺคิโก มคฺโค, เสยฺยถิทํ, สมฺมาทิฏฺฐิ ฯเปฯ สมฺมาสมาธิ.\csromanspacer{} อิทํ วุจฺจติ ภิกฺขเว พฺรหฺมจริยํ.\csromanspacer{} กตมานิ จ ภิกฺขเว พฺรหฺมจริยผลานิ, โสตาปตฺติผลํ สกทาคามิผลํ อนาคามิผลํ อรหตฺตผลํ.\csromanspacer{} อิมานิ วุจฺจนฺติ ภิกฺขเว พฺรหฺมจริยผลานีติ.\csromanspacer{}\csromanspacer{}นวมํ.}

\csromanheader{2}{10. ทุติยพฺรหฺมจริยสุตฺต}
\proseitem{40}{สาวตฺถินิทานํ.\csromanspacer{} พฺรหฺมจริยญฺจ โว ภิกฺขเว เทเสสฺสามิ พฺรหฺมจริยตฺถญฺจ, ตํ สุณาถ.\csromanspacer{} กตมญฺจ ภิกฺขเว พฺรหฺมจริยํ, อยเมว อริโย อฏฺฐงฺคิโก มคฺโค, เสยฺยถิทํ, สมฺมาทิฏฺฐิ ฯเปฯ สมฺมาสมาธิ.\csromanspacer{} อิทํ วุจฺจติ ภิกฺขเว พฺรหฺมจริยํ.\csromanspacer{} กตโม จ ภิกฺขเว พฺรหฺมจริยตฺโถ, โย โข ภิกฺขเว ราคกฺขโย โทสกฺขโย โมหกฺขโย.\csromanspacer{} อยํ วุจฺจติ ภิกฺขเว พฺรหฺมจริยตฺโถติ.\csromanspacer{}\csromanspacer{}ทสมํ.\csromanspacer{} ปฏิปตฺติวคฺโค จตุตฺโถ.}

\titlehead{\textbf{ตสฺสุทฺทานํ}}
\csromangathameasure{ปฏิปตฺติ ปฏิปนฺโน จ, วิรทฺธญฺจ ปารํคมา. \\ สามญฺเญน จ เทฺว วุตฺตา, พฺรหฺมญฺญา อปเร ทุเว. \\ พฺรหฺมจริเยน เทฺว วุตฺตา, วคฺโค เตน ปวุจฺจตีติ.}
\csromangathagroup{\csromangathastanza{ปฏิปตฺติ ปฏิปนฺโน จ, วิรทฺธญฺจ ปารํคมา. \\ สามญฺเญน จ เทฺว วุตฺตา, พฺรหฺมญฺญา อปเร ทุเว. \\ พฺรหฺมจริเยน เทฺว วุตฺตา, วคฺโค เตน ปวุจฺจตีติ.}}

\chapterhead{5. อญฺญติตฺถิยเปยฺยาลวคฺค}
\csromanheader{2}{1. ราควิราคสุตฺต}
\proseitem{41}{สาวตฺถินิทานํ.\csromanspacer{} สเจ โว ภิกฺขเว อญฺญติตฺถิยา ปริพฺพาชกา เอวํ ปุจฺเฉยฺยุํ “กิมตฺถิยํ อาวุโส สมเณ โคตเม พฺรหฺมจริยํ วุสฺสตี”ติ. \csromanfolio{22}เอวํ ปุฏฺฐา ตุเมฺห ภิกฺขเว เตสํ อญฺญติตฺถิยานํ ปริพฺพาชกานํ เอวํ พฺยากเรยฺยาถ “ราควิราคตฺถํ โข อาวุโส ภควติ พฺรหฺมจริยํ วุสฺสตี”ติ.\csromanspacer{} สเจ ปน โว ภิกฺขเว อญฺญติตฺถิยา ปริพฺพาชกา เอวํ ปุจฺเฉยฺยุํ “อตฺถิ ปนาวุโส มคฺโค อตฺถิ ปฏิปทา ราควิราคายา”ติ.\csromanspacer{} เอวํ ปุฏฺฐา ตุเมฺห ภิกฺขเว เตสํ อญฺญติตฺถิยานํ ปริพฺพาชกานํ เอวํ พฺยากเรยฺยาถ “อตฺถิ โข อาวุโส มคฺโค อตฺถิ ปฏิปทา ราควิราคายา”ติ.\csromanspacer{} กตโม จ ภิกฺขเว มคฺโค กตมา จ ปฏิปทา ราควิราคาย.\csromanspacer{} อยเมว อริโย อฏฺฐงฺคิโก มคฺโค, เสยฺยถิทํ, สมฺมาทิฏฺฐิ ฯเปฯ สมฺมาสมาธิ.\csromanspacer{} อยํ ภิกฺขเว มคฺโค อยํ ปฏิปทา ราควิราคายาติ.\csromanspacer{} เอวํ ปุฏฺฐา ตุเมฺห ภิกฺขเว เตสํ อญฺญติตฺถิยานํ ปริพฺพาชกานํ เอวํ พฺยากเรยฺยาถาติ.\csromanspacer{}\csromanspacer{}ปฐมํ.}

\csromanheader{2}{2-7. สํโยชนปฺปหานาทิสุตฺตฉกฺก}
\proseitem{42-47}{สเจ โว ภิกฺขเว อญฺญติตฺถิยา ปริพฺพาชกา เอวํ ปุจฺเฉยฺยุํ “กิมตฺถิยํ อาวุโส สมเณ โคตเม พฺรหฺมจริยํ วุสฺสตี”ติ.\csromanspacer{} เอวํ ปุฏฺฐา ตุเมฺห ภิกฺขเว เตสํ อญฺญติตฺถิยานํ ปริพฺพาชกานํ เอวํ พฺยากเรยฺยาถ “สํโยชนปฺปหานตฺถํ โข อาวุโส ภควติ พฺรหฺมจริยํ วุสฺสตี”ติ ฯเปฯ “อนุสยสมุคฺฆาตนตฺถํ โข อาวุโส ภควติ พฺรหฺมจริยํ วุสฺสตี”ติ ฯเปฯ “อทฺธานปริญฺญตฺถํ โข อาวุโส ภควติ พฺรหฺมจริยํ วุสฺสตี”ติ ฯเปฯ “อาสวานํ ขยตฺถํ โข อาวุโส ภควติ พฺรหฺมจริยํ วุสฺสตี”ติ ฯเปฯ “วิชฺชาวิมุตฺติผล- สจฺฉิกิริยตฺถํ โข อาวุโส ภควติ พฺรหฺมจริยํ วุสฺสตี”ติ ฯเปฯ “ญาณทสฺสนตฺถํ โข อาวุโส ภควติ พฺรหฺมจริยํ วุสฺสตี”ติ ฯเปฯ.\csromanspacer{}\csromanspacer{}สตฺตมํ.}

\csromanheader{2}{8. อนุปาทาปรินิพฺพานสุตฺต}
\proseitem{48}{สาวตฺถินิทานํ.\csromanspacer{} สเจ โว ภิกฺขเว อญฺญติตฺถิยา ปริพฺพาชกา เอวํ ปุจฺเฉยฺยุํ “กิมตฺถิยํ อาวุโส สมเณ โคตเม พฺรหฺมจริยํ วุสฺสตี”ติ.\csromanspacer{} เอวํ ปุฏฺฐา ตุเมฺห ภิกฺขเว เตสํ อญฺญติตฺถิยานํ ปริพฺพาชกานํ เอวํ พฺยากเรยฺยาถ “อนุปาทาปรินิพฺพานตฺถํ โข อาวุโส ภควติ พฺรหฺมจริยํ วุสฺสตี”ติ.\csromanspacer{} สเจ ปน โว ภิกฺขเว อญฺญติตฺถิยา ปริพฺพาชกา เอวํ ปุจฺเฉยฺยุํ “อตฺถิ ปนาวุโส มคฺโค อตฺถิ ปฏิปทา อนุปาทาปรินิพฺพานายา”ติ.\csromanspacer{} เอวํ ปุฏฺฐา ตุเมฺห ภิกฺขเว เตสํ อญฺญติตฺถิยานํ ปริพฺพาชกานํ เอวํ \csromanfolio{23}พฺยากเรยฺยาถ “อตฺถิ โข อาวุโส มคฺโค อตฺถิ ปฏิปทา อนุปาทาปรินิพฺพานายา”ติ.\csromanspacer{} กตโม จ ภิกฺขเว มคฺโค กตมา จ ปฏิปทา อนุปาทาปรินิพฺพานาย, อยเมว อริโย อฏฺฐงฺคิโก มคฺโค, เสยฺยถิทํ, สมฺมาทิฏฺฐิ ฯเปฯ สมฺมาสมาธิ.\csromanspacer{} อยํ ภิกฺขเว มคฺโค อยํ ปฏิปทา อนุปาทาปรินิพฺพานายาติ.\csromanspacer{} เอวํ ปุฏฺฐา ตุเมฺห ภิกฺขเว เตสํ อญฺญติตฺถิยานํ ปริพฺพาชกานํ เอวํ พฺยากเรยฺยาถาติ.\csromanspacer{}\csromanspacer{}อฏฺฐมํ.\csromanspacer{} อญฺญติตฺถิยเปยฺยาลํ.}

\titlehead{\textbf{ตสฺสุทฺทานํ}}
\csromangathameasure{วิราคสํโยชนํ อนุสยํ, อทฺธานํ อาสวา ขยา. \\ วิชฺชา วิมุตฺติ ญาณญฺจ, อนุปาทาย อฏฺฐมี.}
\csromangathagroup{\csromangathastanza{วิราคสํโยชนํ อนุสยํ, อทฺธานํ อาสวา ขยา. \\ วิชฺชา วิมุตฺติ ญาณญฺจ, อนุปาทาย อฏฺฐมี.}}

\chapterhead{6. สูริยเปยฺยาลวคฺค}
\csromanheader{2}{1. กลฺยาณมิตฺตสุตฺต}
\proseitem{49}{สาวตฺถินิทานํ.\csromanspacer{} สูริยสฺส ภิกฺขเว อุทยโต เอตํ ปุพฺพงฺคมํ เอตํ ปุพฺพนิมิตฺตํ, ยทิทํ อรุณุคฺคํ.\csromanspacer{} เอวเมว โข ภิกฺขเว ภิกฺขุโน อริยสฺส อฏฺฐงฺคิกสฺส มคฺคสฺส อุปฺปาทาย เอตํ ปุพฺพงฺคมํ เอตํ ปุพฺพนิมิตฺตํ, ยทิทํ กลฺยาณมิตฺตตา.\csromanspacer{} กลฺยาณมิตฺตสฺเสตํ ภิกฺขเว ภิกฺขุโน ปาฏิกงฺขํ อริยํ อฏฺฐงฺคิกํ มคฺคํ ภาเวสฺสติ, อริยํ อฏฺฐงฺคิกํ มคฺคํ พหุลีกริสฺสติ.\csromanspacer{} กถญฺจ ภิกฺขเว ภิกฺขุ กลฺยาณมิตฺโต อริยํ อฏฺฐงฺคิกํ มคฺคํ ภาเวติ, อริยํ อฏฺฐงฺคิกํ มคฺคํ พหุลีกโรติ, อิธ ภิกฺขเว ภิกฺขุ สมฺมาทิฏฺฐิํ ภาเวติ วิเวกนิสฺสิตํ วิราคนิสฺสิตํ นิโรธนิสฺสิตํ โวสฺสคฺคปริณามิํ ฯเปฯ สมฺมาสมาธิํ ภาเวติ วิเวกนิสฺสิตํ วิราคนิสฺสิตํ นิโรธนิสฺสิตํ โวสฺสคฺคปริณามิํ.\csromanspacer{} เอวํ โข ภิกฺขเว ภิกฺขุ กลฺยาณมิตฺโต อริยํ อฏฺฐงฺคิกํ มคฺคํ ภาเวติ, อริยํ อฏฺฐงฺคิกํ มคฺคํ พหุลีกโรตีติ.\csromanspacer{}\csromanspacer{}ปฐมํ.}

\csromanheader{2}{2-6. สีลสมฺปทาทิสุตฺตปญฺจก}
\proseitem{50-54}{สูริยสฺส ภิกฺขเว อุทยโต เอตํ ปุพฺพงฺคมํ เอตํ ปุพฺพนิมิตฺตํ, ยทิทํ อรุณุคฺคํ.\csromanspacer{} เอวเมว โข ภิกฺขเว ภิกฺขุโน อริยสฺส อฏฺฐงฺคิกสฺส มคฺคสฺส \csromanfolio{24}อุปฺปาทาย เอตํ ปุพฺพงฺคมํ เอตํ ปุพฺพนิมิตฺตํ, ยทิทํ สีลสมฺปทา.\csromanspacer{} สีลสมฺปนฺนสฺเสตํ ภิกฺขเว ภิกฺขุโน ปาฏิกงฺขํ ฯเปฯ ยทิทํ ฉนฺทสมฺปทา ฯเปฯ ยทิทํ อตฺตสมฺปทา ฯเปฯ ยทิทํ ทิฏฺฐิสมฺปทา ฯเปฯ ยทิทํ อปฺปมาทสมฺปทา ฯเปฯ.\csromanspacer{}\csromanspacer{}ฉฏฺฐํ.}

\csromanheader{2}{7. โยนิโสมนสิการสมฺปทาสุตฺต}
\proseitem{55}{สูริยสฺส ภิกฺขเว อุทยโต เอตํ ปุพฺพงฺคมํ เอตํ ปุพฺพนิมิตฺตํ, ยทิทํ อรุณุคฺคํ.\csromanspacer{} เอวเมว โข ภิกฺขเว ภิกฺขุโน อริยสฺส อฏฺฐงฺคิกสฺส มคฺคสฺส อุปฺปาทาย เอตํ ปุพฺพงฺคมํ เอตํ ปุพฺพนิมิตฺตํ, ยทิทํ โยนิโสมนสิการสมฺปทา.\csromanspacer{} โยนิโสมนสิการสมฺปนฺนสฺเสตํ ภิกฺขเว ภิกฺขุโน ปาฏิกงฺขํ อริยํ อฏฺฐงฺคิกํ มคฺคํ ภาเวสฺสติ, อริยํ อฏฺฐงฺคิกํ มคฺคํ พหุลีกริสฺสติ.\csromanspacer{} กถญฺจ ภิกฺขเว ภิกฺขุ โยนิโสมนสิการสมฺปนฺโน อริยํ อฏฺฐงฺคิกํ มคฺคํ ภาเวติ, อริยํ อฏฺฐงฺคิกํ มคฺคํ พหุลีกโรติ, อิธ ภิกฺขเว ภิกฺขุ สมฺมาทิฏฺฐิํ ภาเวติ วิเวกนิสฺสิตํ วิราคนิสฺสิตํ นิโรธนิสฺสิตํ โวสฺสคฺคปริณามิํ ฯเปฯ สมฺมาสมาธิํ ภาเวติ วิเวกนิสฺสิตํ วิราคนิสฺสิตํ นิโรธนิสฺสิตํ โวสฺสคฺคปริณามิํ.\csromanspacer{} เอวํ โข ภิกฺขเว ภิกฺขุ โยนิโสมนสิการสมฺปนฺโน อริยํ อฏฺฐงฺคิกํ มคฺคํ ภาเวติ, อริยํ อฏฺฐงฺคิกํ มคฺคํ พหุลีกโรตีติ.\csromanspacer{}\csromanspacer{}สตฺตมํ.}

\csromanheader{2}{1. กลฺยาณมิตฺตสุตฺต}
\proseitem{56}{สูริยสฺส ภิกฺขเว อุทยโต เอตํ ปุพฺพงฺคมํ เอตํ ปุพฺพนิมิตฺตํ, ยทิทํ อรุณุคฺคํ.\csromanspacer{} เอวเมว โข ภิกฺขเว ภิกฺขุโน อริยสฺส อฏฺฐงฺคิกสฺส มคฺคสฺส อุปฺปาทาย เอตํ ปุพฺพงฺคมํ เอตํ ปุพฺพนิมิตฺตํ, ยทิทํ กลฺยาณมิตฺตตา.\csromanspacer{} กลฺยาณมิตฺตสฺเสตํ ภิกฺขเว ภิกฺขุโน ปาฏิกงฺขํ อริยํ อฏฺฐงฺคิกํ มคฺคํ ภาเวสฺสติ, อริยํ อฏฺฐงฺคิกํ มคฺคํ พหุลีกริสฺสติ.\csromanspacer{} กถญฺจ ภิกฺขเว ภิกฺขุ กลฺยาณมิตฺโต อริยํ อฏฺฐงฺคิกํ มคฺคํ ภาเวติ, อริยํ อฏฺฐงฺคิกํ มคฺคํ พหุลีกโรติ, อิธ หฺหิกฺขเว ภิกฺขุ สมฺมาทิฏฺฐิํ ภาเวติ ราควินยปริโยสานํ โทสวินยปริโยสานํ โมหวินยปริโยสานํ ฯเปฯ สมฺมาสมาธิํ ภาเวติ ราควินยปริโยสานํ โทสวินยปริโยสานํ โมหวินยปริโยสานํ.\csromanspacer{} เอวํ โข ภิกฺขเว หฺหิกฺขุ กลฺยาณมิตฺโต อริยํ อฏฺฐงฺคิกํ มคฺคํ ภาเวติ, อริยํ อฏฺฐงฺคิกํ มคฺคํ พหุลีกโรตีติ.\csromanspacer{}\csromanspacer{}ปฐมํ.}

\csromanheader{2}{2-6. สีลสมฺปทาทิสุตฺตปญฺจก}
\proseitem{57-61}{\csromanfolio{25}สูริยสฺส ภิกฺขเว อุทยโต เอตํ ปุพฺพงฺคมํ เอตํ ปุพฺพนิมิตฺตํ, ยทิทํ อรุณุคฺคํ.\csromanspacer{} เอวเมว โข ภิกฺขเว ภิกฺขุโน อริยสฺส อฏฺฐงฺคิกสฺส มคฺคสฺส อุปฺปาทาย เอตํ ปุพฺพงฺคมํ เอตํ ปุพฺพนิมิตฺตํ, ยทิทํ สีลสมฺปทา ฯเปฯ ยทิทํ ฉนฺทสมฺปทา ฯเปฯ ยทิทํ อตฺตสมฺปทา ฯเปฯ ยทิทํ ทิฏฺฐิสมฺปทา ฯเปฯ ยทิทํ อปฺปมาทสมฺปทา ฯเปฯ.\csromanspacer{}\csromanspacer{}ฉฏฺฐํ.}

\csromanheader{2}{7. โยนิโสมนสิการสมฺปทาสุตฺต}
\proseitem{62}{ยทิทํ โยนิโสมนสิการสมฺปทา.\csromanspacer{} โยนิโสมนสิการสมฺปนฺนสฺเสตํ ภิกฺขเว ภิกฺขุโน ปาฏิกงฺขํ อริยํ อฏฺฐงฺคิกํ มคฺคํ ภาเวสฺสติ, อริยํ อฏฺฐงฺคิกํ มคฺคํ พหุลีกริสฺสติ.\csromanspacer{} กถญฺจ ภิกฺขเว ภิกฺขุ โยนิโสมนสิการสมฺปนฺโน อริยํ อฏฺฐงฺคิกํ มคฺคํ ภาเวติ, อริยํ อฏฺฐงฺคิกํ มคฺคํ พหุลีกโรติ, อิธ ภิกฺขเว ภิกฺขุ สมฺมาทิฏฺฐิํ ภาเวติ ราควินยปริโยสานํ โทสวินยปริโยสานํ โมหวินยปริโยสานํ ฯเปฯ สมฺมาสมาธิํ ภาเวติ ราควินยปริโยสานํ โทสวินยปริโยสานํ โมหวินยปริโยสานํ.\csromanspacer{} เอวํ โข ภิกฺขเว ภิกฺขุ โยนิโสมนสิการสมฺปนฺโน อริยํ อฏฺฐงฺคิกํ มคฺคํ ภาเวติ, อริยํ อฏฺฐงฺคิกํ มคฺคํ พหุลีกโรตีติ.\csromanspacer{}\csromanspacer{}สตฺตมํ.\csromanspacer{} สูริยเปยฺยาลํ.}

\titlehead{\textbf{ตสฺสุทฺทานํ}}
\csromangathameasure{กลฺยาณมิตฺตํ สีลญฺจ, ฉนฺโท จ อตฺตสมฺปทา. \\ ทิฏฺฐิ จ อปฺปมาโท จ, โยนิโส ภวติ สตฺตมํ.}
\csromangathagroup{\csromangathastanza{กลฺยาณมิตฺตํ สีลญฺจ, ฉนฺโท จ อตฺตสมฺปทา. \\ ทิฏฺฐิ จ อปฺปมาโท จ, โยนิโส ภวติ สตฺตมํ.}}

\chapterhead{7. เอกธมฺมเปยฺยาลวคฺค}
\csromanheader{2}{1. กลฺยาณมิตฺตสุตฺต}
\proseitem{63}{สาวตฺถินิทานํ.\csromanspacer{} เอกธมฺโม ภิกฺขเว พหูปกาโร อริยสฺส อฏฺฐงฺคิกสฺส มคฺคสฺส อุปฺปาทาย.\csromanspacer{} กตโม เอกธมฺโม, ยทิทํ กลฺยาณมิตฺตตา. \csromanfolio{26}กลฺยาณมิตฺตสฺเสตํ ภิกฺขเว ภิกฺขุโน ปาฏิกงฺขํ อริยํ อฏฺฐงฺคิกํ มคฺคํ ภาเวสฺสติ, อริยํ อฏฺฐงฺคิกํ มคฺคํ พหุลีกริสฺสติ.\csromanspacer{} กถญฺจ ภิกฺขเว ภิกฺขุ กลฺยาณมิตฺโต อริยํ อฏฺฐงฺคิกํ มคฺคํ ภาเวติ, อริยํ อฏฺฐงฺคิกํ มคฺคํ พหุลีกโรติ, อิธ ภิกฺขเว ภิกฺขุ สมฺมาทิฏฺฐิํ ภาเวติ วิเวกนิสฺสิตํ วิราคนิสฺสิตํ นิโรธนิสฺสิตํ โวสฺสคฺคปริณามิํ ฯเปฯ สมฺมาสมาธิํ ภาเวติ วิเวกนิสฺสิตํ วิราคนิสฺสิตํ นิโรธนิสฺสิตํ โวสฺสคฺคปริณามิํ.\csromanspacer{} เอวํ โข ภิกฺขเว ภิกฺขุ กลฺยาณมิตฺโต อริยํ อฏฺฐงฺคิกํ มคฺคํ ภาเวติ, อริยํ อฏฺฐงฺคิกํ มคฺคํ พหุลีกโรตีติ.\csromanspacer{}\csromanspacer{}ปฐมํ.}

\csromanheader{2}{2-6. สีลสมฺปทาทิสุตฺตปญฺจก}
\proseitem{64-68}{เอกธมฺโม ภิกฺขเว พหูปกาโร อริยสฺส อฏฺฐงฺคิกสฺส มคฺคสฺส อุปฺปาทาย.\csromanspacer{} กตโม เอกธมฺโม, ยทิทํ สีลสมฺปทา ฯเปฯ ยทิทํ ฉนฺทสมฺปทา ฯเปฯ ยทิทํ อตฺตสมฺปทา ฯเปฯ ยทิทํ ทิฏฺฐิสมฺปทา ฯเปฯ ยทิทํ อปฺปมาทสมฺปทา ฯเปฯ.\csromanspacer{}\csromanspacer{}ฉฏฺฐํ.}

\csromanheader{2}{7. โยนิโสมนสิการสมฺปทาสุตฺต}
\proseitem{69}{ยทิทํ โยนิโสมนสิการสมฺปทา.\csromanspacer{} โยนิโสมนสิการสมฺปนฺนสฺเสตํ ภิกฺขเว ภิกฺขุโน ปาฏิกงฺขํ อริยํ อฏฺฐงฺคิกํ มคฺคํ ภาเวสฺสติ, อริยํ อฏฺฐงฺคิกํ มคฺคํ พหุลีกริสฺสติ.\csromanspacer{} กถญฺจ ภิกฺขเว ภิกฺขุ โยนิโสมนสิการสมฺปนฺโน อริยํ อฏฺฐงฺคิกํ มคฺคํ ภาเวติ, อริยํ อฏฺฐงฺคิกํ มคฺคํ พหุลีกโรติ, อิธ ภิกฺขเว ภิกฺขุ สมฺมาทิฏฺฐิํ ภาเวติ วิเวกนิสฺสิตํ วิราคนิสฺสิตํ นิโรธนิสฺสิตํ โวสฺสคฺคปริณามิํ ฯเปฯ สมฺมาสมาธิํ ภาเวติ วิเวกนิสฺสิตํ วิราคนิสฺสิตํ นิโรธนิสฺสิตํ โวสฺสคฺคปริณามิํ.\csromanspacer{} เอวํ โข ภิกฺขเว ภิกฺขุ โยนิโสมนสิการสมฺปนฺโน อริยํ อฏฺฐงฺคิกํ มคฺคํ ภาเวติ, อริยํ อฏฺฐงฺคิกํ มคฺคํ พหุลีกโรตีติ.\csromanspacer{}\csromanspacer{}สตฺตมํ.}

\csromanheader{2}{1. กลฺยาณมิตฺตสุตฺต}
\proseitem{70}{สาวตฺถินิทานํ.\csromanspacer{} เอกธมฺโม ภิกฺขเว พหูปกาโร อริยสฺส อฏฺฐงฺคิกสฺส มคฺคสฺส อุปฺปาทาย.\csromanspacer{} กตโม เอกธมฺโม, ยทิทํ กลฺยาณมิตฺตตา. \csromanfolio{27}กลฺยาณมิตฺตสฺเสตํ ภิกฺขเว ภิกฺขุโน ปาฏิกงฺขํ อริยํ อฏฺฐงฺคิกํ มคฺคํ ภาเวสฺสติ, อริยํ อฏฺฐงฺคิกํ มคฺคํ พหุลีกริสฺสติ.\csromanspacer{} กถญฺจ ภิกฺขเว ภิกฺขุ กลฺยาณมิตฺโต อริยํ อฏฺฐงฺคิกํ มคฺคํ ภาเวติ, อริยํ อฏฺฐงฺคิกํ มคฺคํ พหุลีกโรติ, อิธ ภิกฺขเว ภิกฺขุ สมฺมาทิฏฺฐิํ ภาเวติ ราควินยปริโยสานํ โทสวินยปริโยสานํ โมหวินยปริโยสานํ ฯเปฯ สมฺมาสมาธิํ ภาเวติ ราควินยปริโยสานํ โทสวินยปริโยสานํ โมหวินยปริโยสานํ.\csromanspacer{} เอวํ โข ภิกฺขเว ภิกฺขุ กลฺยาณมิตฺโต อริยํ อฏฺฐงฺคิกํ มคฺคํ ภาเวติ, อริยํ อฏฺฐงฺคิกํ มคฺคํ พหุลีกโรตีติ.\csromanspacer{}\csromanspacer{}ปฐมํ.}

\csromanheader{2}{2-6. สีลสมฺปทาทิสุตฺตปญฺจก}
\proseitem{71-75}{สาวตฺถินิทานํ.\csromanspacer{} เอกธมฺโม ภิกฺขเว พหูปกาโร อริยสฺส อฏฺฐงฺคิกสฺส มคฺคสฺส อุปฺปาทาย.\csromanspacer{} กตโม เอกธมฺโม, ยทิทํ สีลสมฺปทา ฯเปฯ ยทิทํ ฉนฺทสมฺปทา ฯเปฯ ยทิทํ อตฺตสมฺปทา ฯเปฯ ยทิทํ ทิฏฺฐิสมฺปทา ฯเปฯ ยทิทํ อปฺปมาทสมฺปทา ฯเปฯ.\csromanspacer{}\csromanspacer{}ฉฏฺฐํ.}

\csromanheader{2}{7. โยนิโสมนสิการสมฺปทาสุตฺต}
\proseitem{76}{ยทิทํ โยนิโสมนสิการสมฺปทา.\csromanspacer{} โยนิโสมนสิการสมฺปนฺนสฺเสตํ ภิกฺขเว ภิกฺขุโน ปาฏิกงฺขํ อริยํ อฏฺฐงฺคิกํ มคฺคํ ภาเวสฺสติ, อริยํ อฏฺฐงฺคิกํ มคฺคํ พหุลีกริสฺสติ.\csromanspacer{} กถญฺจ ภิกฺขเว ภิกฺขุ โยนิโสมนสิการสมฺปนฺโน อริยํ อฏฺฐงฺคิกํ มคฺคํ ภาเวติ, อริยํ อฏฺฐงฺคิกํ มคฺคํ พหุลีกโรติ, อิธ ภิกฺขเว ภิกฺขุ สมฺมาทิฏฺฐิํ ภาเวติ ฯเปฯ สมฺมาสมาธิํ ภาเวติ ราควินยปริโยสานํ โทสวินยปริโยสานํ โมหวินยปริโยสานํ.\csromanspacer{} เอวํ โข ภิกฺขเว ภิกฺขุ โยนิโสมนสิการสมฺปนฺโน อริยํ อฏฺฐงฺคิกํ มคฺคํ ภาเวติ, อริยํ อฏฺฐงฺคิกํ มคฺคํ พหุลีกโรตีติ.\csromanspacer{}\csromanspacer{}สตฺตมํ.\csromanspacer{} เอกธมฺมเปยฺยาลํ.}

\titlehead{\textbf{ตสฺสุทฺทานํ}}
\csromangathameasure{กลฺยาณมิตฺตํ สีลญฺจ, ฉนฺโท จ อตฺตสมฺปทา. \\ ทิฏฺฐิ จ อปฺปมาโท จ, โยนิโส ภวติ สตฺตมํ.}
\csromangathagroup{\csromangathastanza{กลฺยาณมิตฺตํ สีลญฺจ, ฉนฺโท จ อตฺตสมฺปทา. \\ ทิฏฺฐิ จ อปฺปมาโท จ, โยนิโส ภวติ สตฺตมํ.}}

\chapterheadpageread{8. ทุติย-เอกธมฺมเปยฺยาลวคฺค}
\csromanheader{2}{1. กลฺยาณมิตฺตสุตฺต}
\proseitem{77}{\csromanfolio{28}สาวตฺถินิทานํ.\csromanspacer{} นาหํ ภิกฺขเว อญฺญํ เอกธมฺมมฺปิ สมนุปสฺสามิ, เยน อนุปฺปนฺโน วา อริโย อฏฺฐงฺคิโก มคฺโค อุปฺปชฺชติ, อุปฺปนฺโน วา อริโย อฏฺฐงฺคิโก มคฺโค ภาวนาปาริปูริํ คจฺฉติ, ยถยิทํ ภิกฺขเว กลฺยาณมิตฺตตา.\csromanspacer{} กลฺยาณมิตฺตสฺเสตํ ภิกฺขเว ภิกฺขุโน ปาฏิกงฺขํ อริยํ อฏฺฐงฺคิกํ มคฺคํ ภาเวสฺสติ, อริยํ อฏฺฐงฺคิกํ มคฺคํ พหุลีกริสฺสติ.\csromanspacer{} กถญฺจ ภิกฺขเว ภิกฺขุ กลฺยาณมิตฺโต อริยํ อฏฺฐงฺคิกํ มคฺคํ ภาเวติ, อริยํ อฏฺฐงฺคิกํ มคฺคํ พหุลีกโรติ.\csromanspacer{} อิธ ภิกฺขเว ภิกฺขุ สมฺมาทิฏฺฐิํ ภาเวติ วิเวกนิสฺสิตํ ฯเปฯ สมฺมาสมาธิํ ภาเวติ วิเวกนิสฺสิตํ วิราคนิสฺสิตํ นิโรธนิสฺสิตํ โวสฺสคฺคปริณามิํ.\csromanspacer{} เอวํ โข ภิกฺขเว ภิกฺขุ กลฺยาณมิตฺโต อริยํ อฏฺฐงฺคิกํ มคฺคํ ภาเวติ, อริยํ อฏฺฐงฺคิกํ มคฺคํ พหุลีกโรตีติ.\csromanspacer{}\csromanspacer{}ปฐมํ.}

\csromanheader{2}{2-6. สีลสมฺปทาทิสุตฺตปญฺจก}
\proseitem{78-82}{นาหํ ภิกฺขเว อญฺญํ เอกธมฺมมฺปิ สมนุปสฺสามิ, เยน อนุปฺปนฺโน วา อริโย อฏฺฐงฺคิโก มคฺโค อุปฺปชฺชติ, อุปฺปนฺโน วา อริโย อฏฺฐงฺคิโก มคฺโค ภาวนาปาริปูริํ คจฺฉติ, ยถยิทํ ภิกฺขเว สีลสมฺปทา ฯเปฯ ยถยิทํ ภิกฺขเว ฉนฺทสมฺปทา ฯเปฯ ยถยิทํ ภิกฺขเว อตฺตสมฺปทา ฯเปฯ ยถยิทํ ภิกฺขเว ทิฏฺฐิสมฺปทา ฯเปฯ ยถยิทํ ภิกฺขเว อปฺปมาทสมฺปทา ฯเปฯ.\csromanspacer{}\csromanspacer{}ฉฏฺฐํ.}

\csromanheader{2}{7. โยนิโสมนสิการสมฺปทาสุตฺต}
\proseitem{83}{ยถยิทํ ภิกฺขเว โยนิโสมนสิการสมฺปทา.\csromanspacer{} โยนิโสมนสิการสมฺปนฺนสฺเสตํ ภิกฺขเว ภิกฺขุโน ปาฏิกงฺขํ อริยํ อฏฺฐงฺคิกํ มคฺคํ ภาเวสฺสติ, อริยํ อฏฺฐงฺคิกํ มคฺคํ พหุลีกริสฺสติ.\csromanspacer{} กถญฺจ ภิกฺขเว ภิกฺขุ โยนิโสมนสิการสมฺปนฺโน อริยํ อฏฺฐงฺคิกํ มคฺคํ ภาเวติ, อริยํ อฏฺฐงฺคิกํ มคฺคํ พหุลีกโรติ, อิธ ภิกฺขเว ภิกฺขุ สมฺมาทิฏฺฐิํ ภาเวติ วิเวกนิสฺสิตํ ฯเปฯ สมฺมาสมาธิํ ภาเวติ วิเวกนิสฺสิตํ วิราคนิสฺสิตํ นิโรธนิสฺสิตํ โวสฺสคฺคปริณามิํ.\csromanspacer{} เอวํ โข ภิกฺขเว \csromanfolio{29}ภิกฺขุ โยนิโสมนสิการสมฺปนฺโน อริยํ อฏฺฐงฺคิกํ มคฺคํ ภาเวติ, อริยํ อฏฺฐงฺคิกํ มคฺคํ พหุลีกโรตีติ.\csromanspacer{}\csromanspacer{}สตฺตมํ.}

\csromanheader{2}{1. กลฺยาณมิตฺตสุตฺต}
\proseitem{84}{นาหํ ภิกฺขเว อญฺญํ เอกธมฺมมฺปิ สมนุปสฺสามิ, เยน อนุปฺปนฺโน วา อริโย อฏฺฐงฺคิโก มคฺโค อุปฺปชฺชติ, อุปฺปนฺโน วา อริโย อฏฺฐงฺคิโก มคฺโค ภาวนาปาริปูริํ คจฺฉติ, ยถยิทํ ภิกฺขเว กลฺยาณมิตฺตตา.\csromanspacer{} กลฺยาณมิตฺตสฺเสตํ ภิกฺขเว ภิกฺขุโน ปาฏิกงฺขํ อริยํ อฏฺฐงฺคิกํ มคฺคํ ภาเวสฺสติ, อริยํ อฏฺฐงฺคิกํ มคฺคํ พหุลีกริสฺสติ.\csromanspacer{} กถญฺจ ภิกฺขเว ภิกฺขุ กลฺยาณมิตฺโต อริยํ อฏฺฐงฺคิกํ มคฺคํ ภาเวติ, อริยํ อฏฺฐงฺคิกํ มคฺคํ พหุลีกโรติ, อิธ ภิกฺขเว ภิกฺขุ สมฺมาทิฏฺฐิํ ภาเวติ ราควินยปริโยสานํ โทสวินยปริโยสานํ โมหวินยปริโยสานํ ฯเปฯ สมฺมาสมาธิํ ภาเวติ ราควินยปริโยสานํ โทสวินยปริโยสานํ โมหวินยปริโยสานํ.\csromanspacer{} เอวํ โข ภิกฺขเว ภิกฺขุ กลฺยาณมิตฺโต อริยํ อฏฺฐงฺคิกํ มคฺคํ ภาเวติ, อริยํ อฏฺฐงฺคิกํ มคฺคํ พหุลีกโรตีติ.\csromanspacer{}\csromanspacer{}ปฐมํ.}

\csromanheader{2}{2-6. สีลสมฺปทาทิสุตฺตปญฺจก}
\proseitem{85-89}{นาหํ ภิกฺขเว อญฺญํ เอกธมฺมมฺปิ สมนุปสฺสามิ, เยน อนุปฺปนฺโน วา อริโย อฏฺฐงฺคิโก มคฺโค อุปฺปชฺชติ, อุปฺปนฺโน วา อริโย อฏฺฐงฺคิโก มคฺโค ภาวนาปาริปูริํ คจฺฉติ, ยถยิทํ ภิกฺขเว สีลสมฺปทา ฯเปฯ ยถยิทํ ภิกฺขเว ฉนฺทสมฺปทา ฯเปฯ ยถยิทํ ภิกฺขเว อตฺตสมฺปทา ฯเปฯ ยถยิทํ ภิกฺขเว ทิฏฺฐิสมฺปทา ฯเปฯ ยถยิทํ ภิกฺขเว อปฺปมาทสมฺปทา ฯเปฯ.\csromanspacer{}\csromanspacer{}ฉฏฺฐํ.}

\csromanheader{2}{7. โยนิโสมนสิการสมฺปทาสุตฺต}
\proseitem{90}{ยถยิทํ ภิกฺขเว โยนิโสมนสิการสมฺปทา.\csromanspacer{} โยนิโสมนสิการสมฺปนฺนสฺเสตํ ภิกฺขเว ภิกฺขุโน ปาฏิกงฺขํ อริยํ อฏฺฐงฺคิกํ มคฺคํ ภาเวสฺสติ, อริยํ อฏฺฐงฺคิกํ มคฺคํ พหุลีกริสฺสติ.\csromanspacer{} กถญฺจ ภิกฺขเว ภิกฺขุ โยนิโสมนสิการสมฺปนฺโน อริยํ อฏฺฐงฺคิกํ มคฺคํ ภาเวติ, อริยํ อฏฺฐงฺคิกํ มคฺคํ พหุลีกโรติ, อิธ ภิกฺขเว ภิกฺขุ สมฺมาทิฏฺฐิํ ภาเวติ ราควินยปริโยสานํ โทสวินยปริโยสานํ โมหวินยปริโยสานํ ฯเปฯ สมฺมาสมาธิํ ภาเวติ ราควินยปริโยสานํ โทสวินยปริโยสานํ โมหวินยปริโยสานํ.\csromanspacer{} เอวํ \csromanfolio{30}โข ภิกฺขเว ภิกฺขุ โยนิโสมนสิการสมฺปนฺโน อริยํ อฏฺฐงฺคิกํ มคฺคํ ภาเวติ, อริยํ อฏฺฐงฺคิกํ มคฺคํ พหุลีกโรตีติ.\csromanspacer{}\csromanspacer{}สตฺตมํ.\csromanspacer{} เอกธมฺมเปยฺยาลํ.}

\titlehead{\textbf{ตสฺสุทฺทานํ}}
\csromangathameasure{กลฺยาณมิตฺตํ สีลญฺจ, ฉนฺโท จ อตฺตสมฺปทา. \\ ทิฏฺฐิ จ อปฺปมาโท จ, โยนิโส ภวติ สตฺตมํ.}
\csromangathagroup{\csromangathastanza{กลฺยาณมิตฺตํ สีลญฺจ, ฉนฺโท จ อตฺตสมฺปทา. \\ ทิฏฺฐิ จ อปฺปมาโท จ, โยนิโส ภวติ สตฺตมํ.}}

\csromanmatikaanchor{mtk.1}
\csromanmatikaanchor{mtk.23}
\csromantocmarknopage{chapter}{4. นีวรณวคฺค}{mtk.1}
\bookmark[dest={mtk.1},level=0]{4. นีวรณวคฺค}
\csromanheader{4}{1. คงฺคาเปยฺยาลวคฺค}
\csromanheader{2}{1. ปฐมปาจีนนินฺนสุตฺต}
\proseitem{91}{สาวตฺถินิทานํ.\csromanspacer{} เสยฺยถาปิ ภิกฺขเว คงฺคา นที ปาจีนนินฺนา ปาจีนโปณา ปาจีนปพฺภารา.\csromanspacer{} เอวเมว โข ภิกฺขเว ภิกฺขุ อริยํ อฏฺฐงฺคิกํ มคฺคํ ภาเวนฺโต อริยํ อฏฺฐงฺคิกํ มคฺคํ พหุลีกโรนฺโต นิพฺพานนินฺโน โหติ นิพฺพานโปโณ นิพฺพานปพฺภาโร.\csromanspacer{} กถญฺจ ภิกฺขเว ภิกฺขุ อริยํ อฏฺฐงฺคิกํ มคฺคํ ภาเวนฺโต อริยํ อฏฺฐงฺคิกํ มคฺคํ พหุลีกโรนฺโต นิพฺพานนินฺโน โหติ นิพฺพานโปโณ นิพฺพานปพฺภาโร, อิธ ภิกฺขเว ภิกฺขุ สมฺมาทิฏฺฐิํ ภาเวติ วิเวกนิสฺสิตํ วิราคนิสฺสิตํ นิโรธนิสฺสิตํ โวสฺสคฺคปริณามิํ ฯเปฯ สมฺมาสมาธิํ ภาเวติ วิเวกนิสฺสิตํ วิราคนิสฺสิตํ นิโรธนิสฺสิตํ โวสฺสคฺคปริณามิํ.\csromanspacer{} เอวํ โข ภิกฺขเว ภิกฺขุ อริยํ อฏฺฐงฺคิกํ มคฺคํ ภาเวนฺโต อริยํ อฏฺฐงฺคิกํ มคฺคํ พหุลีกโรนฺโต นิพฺพานนินฺโน โหติ นิพฺพานโปโณ นิพฺพานปพฺภาโรติ.\csromanspacer{}\csromanspacer{}อฐมํ.}

\csromanheader{2}{2-5. ทุติยาทิปาจีนนินฺนสุตฺตจตุกฺก}
\proseitem{92-95}{เสยฺยถาปิ ภิกฺขเว ยมุนา นที ปาจีนนินฺนา ปาจีนโปณา ปาจีนปพฺภารา.\csromanspacer{} เอวเมว โข ภิกฺขเว ฯเปฯ.\csromanspacer{} เสยฺยถาปิ ภิกฺขเว อจิรวตี นที ปาจีนนินฺนา ปาจีนโปณา ปาจีนปพฺภารา.\csromanspacer{} เอวเมว โข ภิกฺขเว ฯเปฯ.\csromanspacer{} เสยฺยถาปิ ภิกฺขเว สรภู นที ปาจีนนินฺนา ปาจีนโปณา ปาจีนปพฺภารา.\csromanspacer{} เอวเมว โข ภิกฺขเว ฯเปฯ.\csromanspacer{} เสยฺยถาปิ ภิกฺขเว มหี นที ปาจีนนินฺนา ปาจีนโปณา ปาจีนปพฺภารา.\csromanspacer{} เอวเมว โข ภิกฺขเว ฯเปฯ.\csromanspacer{}\csromanspacer{}ปญฺจมํ.}

\csromanheader{2}{6. ฉฏฺฐปาจีนนินฺนสุตฺต}
\proseitem{96}{\csromanfolio{31}เสยฺยถาปิ ภิกฺขเว ยา กาจิมา มหานทิโย, เสยฺยถิทํ.\csromanspacer{} คงฺคา ยมุนา อจิรวตี สรภู มหี, สพฺพา ตา ปาจีนนินฺนา ปาจีนโปณา ปาจีนปพฺภารา.\csromanspacer{} เอวเมว โข ภิกฺขเว ภิกฺขุ อริยํ อฏฺฐงฺคิกํ มคฺคํ ภาเวนฺโต อริยํ อฏฺฐงฺคิกํ มคฺคํ พหุลีกโรนฺโต นิพฺพานนินฺโน โหติ นิพฺพานโปโณ นิพฺพานปพฺภาโร.\csromanspacer{} กถญฺจ ภิกฺขเว ภิกฺขุ อริยํ อฏฺฐงฺคิกํ มคฺคํ ภาเวนฺโต อริยํ อฏฺฐงฺคิกํ มคฺคํ พหุลีกโรนฺโต นิพฺพานนินฺโน โหติ นิพฺพานโปโณ นิพฺพานปพฺภาโร, อิธ ภิกฺขเว ภิกฺขุ สมฺมาทิฏฺฐิํ ภาเวติ วิเวกนิสฺสิตํ ฯเปฯ สมฺมาสมาธิํ ภาเวติ วิเวกนิสฺสิตํ วิราคนิสฺสิตํ นิโรธนิสฺสิตํ โวสฺสคฺคปริณามิํ.\csromanspacer{} เอวํ โข ภิกฺขเว ภิกฺขุ อริยํ อฏฺฐงฺคิกํ มคฺคํ ภาเวนฺโต อริยํ อฏฺฐงฺคิกํ มคฺคํ พหุลีกโรนฺโต นิพฺพานนินฺโน โหติ นิพฺพานโปโณ นิพฺพานปพฺภาโรติ.\csromanspacer{}\csromanspacer{}ฉฏฺฐํ.}

\csromanheader{2}{1. ปฐมสมุทฺทนินฺนสุตฺต}
\proseitem{97}{เสยฺยถาปิ ภิกฺขเว คงฺคา นที สมุทฺทนินฺนา สมุทฺทโปณา สมุทฺทปพฺภารา.\csromanspacer{} เอวเมว โข ภิกฺขเว ภิกฺขุ อริยํ อฏฺฐงฺคิกํ มคฺคํ ภาเวนฺโต อริยํ อฏฺฐงฺคิกํ มคฺคํ พหุลีกโรนฺโต นิพฺพานนินฺโน โหติ นิพฺพานโปโณ นิพฺพานปพฺภาโร.\csromanspacer{} กถญฺจ ภิกฺขเว ภิกฺขุ อริยํ อฏฺฐงฺคิกํ มคฺคํ ภาเวนฺโต อริยํ อฏฺฐงฺคิกํ มคฺคํ พหุลีกโรนฺโต นิพฺพานนินฺโน โหติ นิพฺพานโปโณ นิพฺพานปพฺภาโร, อิธ ภิกฺขเว ภิกฺขุ สมฺมาทิฏฺฐิํ ภาเวติ วิเวกนิสฺสิตํ ฯเปฯ สมฺมาสมาธิํ ภาเวติ วิเวกนิสฺสิตํ วิราคนิสฺสิตํ นิโรธนิสฺสิตํ โวสฺสคฺคปริณามิํ.\csromanspacer{} เอวํ โข ภิกฺขเว ภิกฺขุ อริยํ อฏฺฐงฺคิกํ มคฺคํ ภาเวนฺโต อริยํ อฏฺฐงฺคิกํ มคฺคํ พหุลีกโรนฺโต นิพฺพานนินฺโน โหติ นิพฺพานโปโณ นิพฺพานปพฺภาโรติ.\csromanspacer{}\csromanspacer{}ปฐมํ.}

\csromanheader{2}{2-6. ทุติยาทิสมุทฺทนินฺนสุตฺตปญฺจก}
\proseitem{98-102}{เสยฺยถาปิ ภิกฺขเว ยมุนา นที สมุทฺทนินฺนา สมุทฺทโปณา สมุทฺทปพฺภารา.\csromanspacer{} เอวเมว โข ภิกฺขเว ภิกฺขุ ฯเปฯ.\csromanspacer{} เสยฺยถาปิ ภิกฺขเว อจิรวตี นที สมุทฺทนินฺนา สมุทฺทโปณา สมุทฺทปพฺภารา.\csromanspacer{} เอวเมว โข ภิกฺขเว ภิกฺขุ ฯเปฯ.\csromanspacer{} เสยฺยถาปิ ภิกฺขเว สรภู นที สมุทฺทนินฺนา สมุทฺทโปณา สมุทฺทปพฺภรา.\csromanspacer{} เอวเมว โข ภิกฺขเว ภิกฺขุ ฯเปฯ.\csromanspacer{} เสยฺยถาปิ ภิกฺขเว \csromanfolio{32}มหี นที สมุทฺทนินฺนา สมุทฺทโปณา สมุทฺทปพฺภารา.\csromanspacer{} เอวเมว โข ภิกฺขเว ภิกฺขุ ฯเปฯ.\csromanspacer{} เสยฺยถาปิ ภิกฺขเว ยา กาจิมา มหานทิโย, เสยฺยถิทํ, คงฺคา ยมุนา อจิรวตี สรภู มหี, สพฺพา ตา สมุทฺทนินฺนา สมุทฺทโปณา สมุทฺทปพฺภารา.\csromanspacer{} เอวเมว โข ภิกฺขเว ภิกฺขุ อริยํ อฏฺฐงฺคิกํ มคฺคํ ภาเวนฺโต อริยํ อฏฺฐงฺคิกํ มคฺคํ พหุลีกโรนฺโต นิพฺพานนินฺโน โหติ นิพฺพานโปโณ นิพฺพานปพฺภาโร.\csromanspacer{} กถญฺจ ภิกฺขเว ภิกฺขุ อริยํ อฏฺฐงฺคิกํ มคฺคํ ภาเวนฺโต อริยํ อฏฺฐงฺคิกํ มคฺคํ พหุลีกโรนฺโต นิพฺพานนินฺโน โหติ นิพฺพานโปโณ นิพฺพานปพฺภาโร, อิธ ภิกฺขเว ภิกฺขุ สมฺมาทิฏฺฐิํ ภาเวติ วิเวกนิสฺสิตํ วิราคนิสฺสิตํ นิโรธนิสฺสิตํ โวสฺสคฺคปริณามิํ ฯเปฯ สมฺมาสมาธิํ ภาเวติ วิเวกนิสฺสตํ วิราคนิสฺสิตํ นิโรธนิสฺสิตํ โวสฺสคฺคปริณามิํ.\csromanspacer{} เอวํ โข ภิกฺขเว ภิกฺขุ อริยํ อฏฺฐงฺคิกํ มคฺคํ ภาเวนฺโต อริยํ อฏฺฐงฺคิกํ มคฺคํ พหุลีกโรนฺโต นิพฺพานนินฺโน โหติ นิพฺพานโปโณ นิพฺพานปพฺภาโรติ.\csromanspacer{}\csromanspacer{}ฉฏฺฐํ.\csromanspacer{} คงฺคาเปยฺยาลํ.}

\titlehead{\textbf{ตสฺสุทฺทานํ}}
\csromangathameasure{ฉ ปาจีนโต นินฺนา, ฉ นินฺนา จ สมุทฺทโต. \\ เอเต เทฺว ฉ ทฺวาทส โหนฺติ, วคฺโค เตน ปวุจฺจตีติ. \\ คงฺคาเปยฺยาลี ปาจินนินฺนวาจนมคฺคี,}
\csromangathagroup{\csromangathastanza{ฉ ปาจีนโต นินฺนา, ฉ นินฺนา จ สมุทฺทโต. \\ เอเต เทฺว ฉ ทฺวาทส โหนฺติ, วคฺโค เตน ปวุจฺจตีติ. \\ คงฺคาเปยฺยาลี ปาจินนินฺนวาจนมคฺคี,}}

\csromanheader{2}{วิเวกนิสฺสิตํ ทฺวาทสกี ปฐมกี.}
\csromansectionrule
\chapterhead{2. ทุติยคงฺคาเปยฺยาลวคฺค}
\csromanheader{2}{1. ปฐมปาจีนนินฺนสุตฺต}
\proseitem{103}{เสยฺยถาปิ ภิกฺขเว คงฺคา นที ปาจีนนินฺนา ปาจีนโปณา ปาจีนปพฺภารา.\csromanspacer{} เอวเมว โข ภิกฺขเว ภิกฺขุ อริยํ อฏฺฐงฺคิกํ มคฺคํ ภาเวนฺโต อริยํ อฏฺฐงฺคิกํ มคฺคํ พหุลีกโรนฺโต นิพฺพานนินฺโน โหติ นิพฺพานโปโณ นิพฺพานปพฺภาโร.\csromanspacer{} กถญฺจ ภิกฺขเว ภิกฺขุ อริยํ อฏฺฐงฺคิกํ มคฺคํ ภาเวนฺโต อริยํ อฏฺฐงฺคิกํ มคฺคํ พหุลีกโรนฺโต นิพฺพานนินฺโน โหติ นิพฺพานโปโณ นิพฺพานปพฺภาโร, อิธ ภิกฺขเว ภิกฺขุ สมฺมาทิฏฺฐิํ ภาเวติ \csromanfolio{33}ราควินยปริโยสานํ โทสวินยปริโยสานํ โมหวินยปริโยสานํ ฯเปฯ สมฺมาสมาธิํ ภาเวติ ราควินยปริโยสานํ โทสวินยปริโยสานํ โมหวินยปริโยสานํ.\csromanspacer{} เอวํ โข ภิกฺขเว ภิกฺขุ อริยํ อฏฺฐงฺคิกํ มคฺคํ ภาเวนฺโต อริยํ อฏฺฐงฺคิกํ มคฺคํ พหุลีกโรนฺโต นิพฺพานนินฺโน โหติ นิพฺพานโปโณ นิพฺพานปพฺภาโรติ.\csromanspacer{}\csromanspacer{}ปฐมํ.}

\csromanheader{2}{2-6. ทุติยาทิปาจีนนินฺนสุตฺตปญฺจก}
\proseitem{104-108}{เสยฺยถาปิ ภิกฺขเว ยมุนา นที ปาจีนนินฺนา ปาจีนโปณา ปาจีนปพฺภารา.\csromanspacer{} เอวเมว โข ภิกฺขเว ภิกฺขุ ฯเปฯ.\csromanspacer{} เสยฺยถาปิ ภิกฺขเว อจิรวตี นที ปาจีนนินฺนา ปาจีนโปณา ปาจีนปพฺภารา.\csromanspacer{} เอวเมว โข ภิกฺขเว ภิกฺขุ ฯเปฯ.\csromanspacer{} เสยฺยถาปิ ภิกฺขเว สรภู นที ปาจีนนินฺนา ปาจีนโปณา ปาจีนปพฺภารา.\csromanspacer{} เอวเมว โข ภิกฺขเว ภิกฺขุ ฯเปฯ.\csromanspacer{} เสยฺยถาปิ ภิกฺขเว มหี นที ปาจีนนินฺนา ปาจีนโปณา ปาจีนปพฺภารา.\csromanspacer{} เอวเมว โข ภิกฺขเว ภิกฺขุ ฯเปฯ.\csromanspacer{} เสยฺยถาปิ ภิกฺขเว ยา กาจิมา มหานทิโย.\csromanspacer{} เสยฺยถิทํ, คงฺคา ยมุนา อจิรวตี สรภู มหี, สพฺพา ตา ปาจีนนินฺนา ปาจีนโปณา ปาจีนปพฺภารา.\csromanspacer{} เอวเมว โข ภิกฺขเว ภิกฺขุ ฯเปฯ.\csromanspacer{} ฉฏฺฐํ.}

\csromanheader{2}{7. ปฐมสมุทฺทนินฺนสุตฺต}
\proseitem{109}{เสยฺยถาปิ ภิกฺขเว คงฺคา นที สมุทฺทนินฺนา สมุทฺทโปณา สมุทฺทปพฺภารา.\csromanspacer{} เอวเมว โข ภิกฺขเว ภิกฺขุ อริยํ อฏฺฐงฺคิกํ มคฺคํ ภาเวนฺโต อริยํ อฏฺฐงฺคิกํ มคฺคํ พหุลีกโรนฺโต นิพฺพานนินฺโน โหติ นิพฺพานโปโณ นิพฺพานปพฺภาโร.\csromanspacer{} กถญฺจ ภิกฺขเว ภิกฺขุ อริยํ อฏฺฐงฺคิกํ มคฺคํ ภาเวนฺโต อริยํ อฏฺฐงฺคิกํ มคฺคํ พหุลีกโรนฺโต นิพฺพานนินฺโน โหติ นิพฺพานโปโณ นิพฺพานปพฺภาโร, อิธ ภิกฺขเว ภิกฺขุ สมฺมาทิฏฺฐิํ ภาเวติ ราควินยปริโยสานํ โทสวินยปริโยสานํ โมหวินยปริโยสานํ ฯเปฯ สมฺมาสมาธิํ ภาเวติ ราควินยปริโยสานํ โทสวินยปริโยสานํ โมหวินยปริโยสานํ.\csromanspacer{} เอวํ โข ภิกฺขเว ภิกฺขุ อริยํ อฏฺฐงฺคิกํ มคฺคํ ภาเวนฺโต อริยํ อฏฺฐงฺคิกํ มคฺคํ พหุลีกโรนฺโต นิพฺพานนินฺโน โหติ นิพฺพานโปโณ นิพฺพานปพฺภาโรติ.\csromanspacer{}\csromanspacer{}ปฐมํ.}

\csromanheader{2}{8-12. ทุติยาทิสมุทฺทนินฺนสุตฺต}
\proseitem{110-114}{เสยฺยถาปิ ภิกฺขเว ยมุนา นที สมุทฺทนินฺนา สมุทฺทโปณา สมุทฺทปพฺภรา.\csromanspacer{} เอวเมว โข ภิกฺขเว ภิกฺขุ ฯเปฯ.\csromanspacer{} เสยฺยถาปิ ภิกฺขเว \csromanfolio{34}อจิรวตี นที สมุทฺทนินฺนา สมุทฺทโปณา สมุทฺทปพฺภารา.\csromanspacer{} เอวเมว โข ภิกฺขเว ภิกฺขุ ฯเปฯ.\csromanspacer{} เสยฺยถาปิ ภิกฺขเว สรภู นที สมุทฺทนินฺนา สมุทฺทโปณา สมุทฺทปพฺภารา.\csromanspacer{} เอวเมว โข ภิกฺขเว ภิกฺขุ ฯเปฯ.\csromanspacer{} เสยฺยถาปิ ภิกฺขเว มหี นที สมุทฺทนินฺนา สมุทฺทโปณา สมุทฺทปพฺภารา.\csromanspacer{} เอวเมว โข ภิกฺขเว ภิกฺขุ ฯเปฯ.\csromanspacer{} เสยฺยถาปิ ภิกฺขเว ยา กาจิมา มหานทิโย.\csromanspacer{} เสยฺยถิทํ, คงฺคา ยมุนา อจิรวตี สรภู มหี, สพฺพา ตา สมุทฺทนินฺนา สมุทฺทโปณา สมุทฺทปพฺภารา.\csromanspacer{} เอวเมว โข ภิกฺขเว ภิกฺขุ อริยํ อฏฺฐงฺคิกํ มคฺคํ ภาเวนฺโต อริยํ อฏฺฐงฺคิกํ มคฺคํ พหุลีกโรนฺโต นิพฺพานนินฺโน โหติ นิพฺพานโปโณ นิพฺพานปพฺภาโร.\csromanspacer{} กถญฺจ ภิกฺขเว ภิกฺขุ อริยํ อฏฺฐงฺคิกํ มคฺคํ ภาเวนฺโต อริยํ อฏฺฐงฺคิกํ มคฺคํ พหุลีกโรนฺโต นิพฺพานนินฺโน โหติ นิพฺพานโปโณ นิพฺพานปพฺภาโร.\csromanspacer{} อิธ ภิกฺขเว ภิกฺขุ สมฺมาทิฏฺฐิํ ภาเวติ ราควินยปริโยสานํ โทสวินยปริโยสานํ โมหวินยปริโยสานํ ฯเปฯ สมฺมาสมาธิํ ภาเวติ ราควินยปริโยสานํ โทสวินยปริโยสานํ โมหวินยปริโยสานํ.\csromanspacer{} เอวํ โข ภิกฺขเว ภิกฺขุ อริยํ อฏฺฐงฺคิกํ มคฺคํ ภาเวนฺโต อริยํ อฏฺฐงฺคิกํ มคฺคํ พหุลีกโรนฺโต นิพฺพานนินฺโน โหติ นิพฺพานโปโณ นิพฺพานปพฺภาโรติ.\csromanspacer{}\csromanspacer{}ฉฏฺฐํ. (ราควินยทฺวาทสกี ทุติยกี สมุทฺทนินฺนนฺติ.)}

\csromanheader{2}{13. ปฐมปาจีนนินฺนสุตฺต}
\proseitem{115}{เสยฺยถาปิ ภิกฺขเว คงฺคา นที ปาจีนนินฺนา ปาจีนโปณา ปาจีนปพฺภารา.\csromanspacer{} เอวเมว โข ภิกฺขเว ภิกฺขุ อริยํ อฏฺฐงฺคิกํ มคฺคํ ภาเวนฺโต อริยํ อฏฺฐงฺคิกํ มคฺคํ พหุลีกโรนฺโต นิพฺพานนินฺโน โหติ นิพฺพานโปโณ นิพฺพานปพฺภาโร.\csromanspacer{} กถญฺจ ภิกฺขเว ภิกฺขุ อริยํ อฏฺฐงฺคิกํ มคฺคํ ภาเวนฺโต อริยํ อฏฺฐงฺคิกํ มคฺคํ พหุลีกโรนฺโต นิพฺพานนินฺโน โหติ นิพฺพานโปโณ นิพฺพานปพฺภาโร, อิธ ภิกฺขเว ภิกฺขุ สมฺมาทิฏฺฐิํ ภาเวติ อมโตคธํ อมตปรายนํ อมตปริโยสานํ ฯเปฯ สมฺมาสมาธิํ ภาเวติ อมโตคธํ อมตปรายนํ อมตปริโยสานํ.\csromanspacer{} เอวํ โข ภิกฺขเว ภิกฺขุ อริยํ อฏฺฐงฺคิกํ มคฺคํ ภาเวนฺโต อริยํ อฏฺฐงฺคิกํ มคฺคํ พหุลีกโรนฺโต นิพฺพานนินฺโน โหติ นิพฺพานโปโณ นิพฺพานปพฺภาโรติ.\csromanspacer{}\csromanspacer{}ปฐมํ.}

\csromanheader{2}{14-18. ทุติยาทิปาจีนนินฺนสุตฺต}
\proseitem{116-120}{\csromanfolio{35}เสยฺยถาปิ ภิกฺขเว ยมุนา นที ปาจีนนินฺนา ปาจีนโปณา ปาจีนปพฺภารา.\csromanspacer{} เอวเมว โข ภิกฺขเว ภิกฺขุ ฯเปฯ.\csromanspacer{} เสยฺยถาปิ ภิกฺขเว อจิรวตี นที ปาจีนนินฺนา ปาจีนโปณา ปาจีนปพฺภารา.\csromanspacer{} เอวเมว โข ภิกฺขเว ภิกฺขุ ฯเปฯ.\csromanspacer{} เสยฺยถาปิ ภิกฺขเว สรภู นที ปาจีนนินฺนา ปาจีนโปณา ปาจีนปพฺภารา.\csromanspacer{} เอวเมว โข ภิกฺขเว ภิกฺขุ ฯเปฯ.\csromanspacer{} เสยฺยถาปิ ภิกฺขเว มหี นที ปาจีนนินฺนา ปาจีนโปณา ปาจีนปพฺภารา.\csromanspacer{} เอวเมว โข ภิกฺขเว ภิกฺขุ ฯเปฯ.\csromanspacer{} เสยฺยถาปิ ภิกฺขเว ยา กาจิมา มหานทิโย.\csromanspacer{} เสยฺยถิทํ, คงฺคา ยมุนา อจิรวตี สรภู มหี, สพฺพา ตา ปาจีนนินฺนา ปาจีนโปณา ปาจีนปพฺภารา.\csromanspacer{} เอวเมว โข ภิกฺขเว ภิกฺขุ ฯเปฯ.\csromanspacer{}\csromanspacer{}ฉฏฺฐํ.}

\csromanheader{2}{19. ปฐมสมุทฺทนินฺนสุตฺต}
\proseitem{121}{เสยฺยถาปิ ภิกฺขเว คงฺคา นที สมุทฺทนินฺนา สมุทฺทโปณา สมุทฺทปพฺภารา.\csromanspacer{} เอวเมว โข ภิกฺขเว ภิกฺขุ อริยํ อฏฺฐงฺคิกํ มคฺคํ ภาเวนฺโต อริยํ อฏฺฐงฺคิกํ มคฺคํ พหุลีกโรนฺโต นิพฺพานนินฺโน โหติ นิพฺพานโปโณ นิพฺพานปพฺภาโร.\csromanspacer{} กถญฺจ ภิกฺขเว ภิกฺขุ อริยํ อฏฺฐงฺคิกํ มคฺคํ ภาเวนฺโต อริยํ อฏฺฐงฺคิกํ มคฺคํ พหุลีกโรนฺโต นิพฺพานนินฺโน โหติ นิพฺพานโปโณ นิพฺพานปพฺภาโร, อิธ ภิกฺขเว ภิกฺขุ สมฺมาทิฏฺฐิํ ภาเวติ อมโตคธํ อมตปรายนํ อมตปริโยสานํ ฯเปฯ สมฺมาสมาธิํ ภาเวติ อมโตคธํ อมตปรายนํ อมตปริโยสานํ.\csromanspacer{} เอวํ โข ภิกฺขเว ภิกฺขุ อริยํ อฏฺฐงฺคิกํ มคฺคํ ภาเวนฺโต อริยํ อฏฺฐงฺคิกํ มคฺคํ พหุลีกโรนฺโต นิพฺพานนินฺโน โหติ นิพฺพานโปโณ นิพฺพานปพฺภาโรติ.\csromanspacer{}\csromanspacer{}ปฐมํ.}

\csromanheader{2}{20-24. ทุติยาทิสมุทฺทนินฺนสุตฺต}
\proseitem{122-126}{เสยฺยถาปิ ภิกฺขเว ยมุนา นที สมุทฺทนินฺนา สมุทฺทโปณา สมุทฺทาปพฺภารา.\csromanspacer{} เอวเมว โข ภิกฺขเว ภิกฺขุ ฯเปฯ.\csromanspacer{} เสยฺยถาปิ ภิกฺขเว อจิรวตี นที สมุทฺทนินฺนา สมุทฺทโปณา สมุทฺทปพฺภารา.\csromanspacer{} เอวเมว โข ภิกฺขเว ภิกฺขุ ฯเปฯ.\csromanspacer{} เสยฺยถาปิ ภิกฺขเว สรภู นที สมุทฺทนินฺนา สมุทฺทโปณา สมุทฺทปพฺภารา.\csromanspacer{} เอวเมว โข ภิกฺขเว ภิกฺขุ ฯเปฯ.\csromanspacer{} เสยฺยถาปิ ภิกฺขเว มหี นที สมุทฺทนินฺนา สมุทฺทโปณา สมุทฺทปพฺภารา.\csromanspacer{} เอวเมว โข ภิกฺขเว ภิกฺขุ ฯเปฯ.\csromanspacer{} เสยฺยถาปิ ภิกฺขเว ยา กาจิมา มหานทิโย.\csromanspacer{} เสยฺยถิทํ, \csromanfolio{36}คงฺคา ยมุนา อจิรวตี สรภู มหี, สพฺพา ตา สมุทฺทนินฺนา สมุทฺทโปณา สมุทฺทปพฺภารา.\csromanspacer{} เอวเมว โข ภิกฺขเว ภิกฺขุ อริยํ อฏฺฐงฺคิกํ มคฺคํ ภาเวนฺโต อริยํ อฏฺฐงฺคิกํ มคฺคํ พหุลีกโรนฺโต นิพฺพานนินฺโน โหติ นิพฺพานโปโณ นิพฺพานปพฺภาโร.\csromanspacer{} กถญฺจ ภิกฺขเว ภิกฺขุ อริยํ อฏฺฐงฺคิกํ มคฺคํ ภาเวนฺโต อริยํ อฏฺฐงฺคิกํ มคฺคํ พหุลีกโรนฺโต นิพฺพานนินฺโน โหติ นิพฺพานโปโณ นิพฺพานปพฺภาโร, อิธ ภิกฺขเว ภิกฺขุ สมฺมาทิฏฺฐิํ ภาเวติ อมโตคธํ อมตปรายนํ อมตปริโยสานํ ฯเปฯ สมฺมาสมาธิํ ภาเวติ อมโตคธํ อมตปรายนํ อมตปริโยสานํ.\csromanspacer{} เอวํ โข ภิกฺขเว ภิกฺขุ อริยํ อฏฺฐงฺคิกํ มคฺคํ ภาเวนฺโต อริยํ อฏฺฐงฺคิกํ มคฺคํ พหุลีกโรนฺโต นิพฺพานนินฺโน โหติ นิพฺพานโปโณ นิพฺพานปพฺภาโรติ.\csromanspacer{}\csromanspacer{}ฉฏฺฐํ. (อมโตคธทฺวาทสกี ตติยกี.)}

\csromanheader{2}{25. ปฐมปาจีนนินฺนสุตฺต}
\proseitem{127}{เสยฺยถาปิ ภิกฺขเว คงฺคา นที ปาจีนนินฺนา ปาจีนโปณา ปาจีนปพฺภารา.\csromanspacer{} เอวเมว โข ภิกฺขเว ภิกฺขุ อริยํ อฏฺฐงฺคิกํ มคฺคํ ภาเวนฺโต อริยํ อฏฺฐงฺคิกํ มคฺคํ พหุลีกโรนฺโต นิพฺพานนินฺโน โหติ นิพฺพานโปโณ นิพฺพานปพฺภาโร.\csromanspacer{} กถญฺจ ภิกฺขเว ภิกฺขุ อริยํ อฏฺฐงฺคิกํ มคฺคํ ภาเวนฺโต อริยํ อฏฺฐงฺคิกํ มคฺคํ พหุลีกโรนฺโต นิพฺพานนินฺโน โหติ นิพฺพานโปโณ นิพฺพานปพฺภาโร, อิธ ภิกฺขเว ภิกฺขุ สมฺมาทิฏฺฐิํ ภาเวติ นิพฺพานนินฺนํ นิพฺพานโปณํ นิพฺพานปพฺภารํ ฯเปฯ สมฺมาสมาธิํ ภาเวติ นิพฺพานนินฺนํ นิพฺพานโปณํ นิพฺพานปพฺภารํ.\csromanspacer{} เอวํ โข ภิกฺขเว ภิกฺขุ อริยํ อฏฺฐงฺคิกํ มคฺคํ ภาเวนฺโต อริยํ อฏฺฐงฺคิกํ มคฺคํ พหุลีกโรนฺโต นิพฺพานนินฺโน โหติ นิพฺพานโปโณ นิพฺพานปพฺภาโรติ.\csromanspacer{}\csromanspacer{}ปฐมํ.}

\csromanheader{2}{26-30. ทุติยาทิปาจีนนินฺนสุตฺต}
\proseitem{128-132}{เสยฺยถาปิ ภิกฺขเว ยมุนา นที ปาจีนนินฺนา ปาจีนโปณา ปาจีนปพฺภารา.\csromanspacer{} เอวเมว โข ภิกฺขเว ภิกฺขุ ฯเปฯ.\csromanspacer{} เสยฺยถาปิ ภิกฺขเว อจิรวตี นที ปาจีนนินฺนา ปาจีนโปณา ปาจีนปพฺภารา.\csromanspacer{} เอวเมว โข ภิกฺขเว ภิกฺขุ ฯเปฯ.\csromanspacer{} เสยฺยถาปิ ภิกฺขเว สรภู นที ปาจีนนินฺนา ปาจีนโปณา ปาจีนปพฺภารา.\csromanspacer{} เอวเมว โข ภิกฺขเว ภิกฺขุ ฯเปฯ.\csromanspacer{} เสยฺยถาปิ ภิกฺขเว มหี นที ปาจีนนินฺนา \csromanfolio{37}ปาจีนโปณา ปาจีนปพฺพารา.\csromanspacer{} เอวเมว โข ภิกฺขเว ภิกฺขุ ฯเปฯ.\csromanspacer{} เสยฺยถาปิ ภิกฺขเว ยา กาจิมา มหานทิโย.\csromanspacer{} เสยฺยถิทํ, คงฺคา ยมุนา อจิรวตี สรภู มหี, สพฺพา ตา ปาจีนนินฺนา ปาจีนโปณา ปาจีนปพฺภารา.\csromanspacer{} เอวเมว โข ภิกฺขเว ภิกฺขุ อริยํ อฏฺฐงฺคิกํ มคฺคํ ภาเวนฺโต อริยํ อฏฺฐงฺคิกํ มคฺคํ พหุลีกโรนฺโต นิพฺพานนินฺโน โหติ นิพฺพานโปโณ นิพฺพานปพฺภาโร.\csromanspacer{} กถญฺจ ภิกฺขเว ภิกฺขุ อริยํ อฏฺฐงฺคิกํ มคฺคํ ภาเวนฺโต อริยํ อฏฺฐงฺคิกํ มคฺคํ พหุลีกโรนฺโต นิพฺพานนินฺโน โหติ นิพฺพานโปโณ นิพฺพานปพฺภาโร, อิธ ภิกฺขเว ภิกฺขุ สมฺมาทิฏฺฐิํ ภาเวติ นิพฺพานนินฺนํ นิพฺพานโปณํ นิพฺพานปพฺภารํ ฯเปฯ สมฺมาสมาธิํ ภาเวติ นิพฺพานนินฺนํ นิพฺพานโปณํ นิพฺพานปพฺภารํ.\csromanspacer{} เอวํ โข ภิกฺขเว ภิกฺขุ อริยํ อฏฺฐงฺคิกํ มคฺคํ ภาเวนฺโต อริยํ อฏฺฐงฺคิกํ มคฺคํ พหุลีกโรนฺโต นิพฺพานนินฺโน โหติ นิพฺพานโปโณ นิพฺพานปพฺภาโรติ.\csromanspacer{}\csromanspacer{}ฉฏฺฐํ.}

\csromanheader{2}{31. ปฐมสมุทฺทนินฺนสุตฺต}
\proseitem{133}{เสยฺยถาปิ ภิกฺขเว คงฺคา นที สมุทฺทนินฺนา สมุทฺทโปณา สมุทฺทปพฺภารา.\csromanspacer{} เอวเมว โข ภิกฺขเว ภิกฺขุ อริยํ อฏฺฐงฺคิกํ มคฺคํ ภาเวนฺโต อริยํ อฏฺฐงฺคิกํ มคฺคํ พหุลีกโรนฺโต นิพฺพานนินฺโน โหติ นิพฺพานโปโณ นิพฺพานปพฺภาโร.\csromanspacer{} กถญฺจ ภิกฺขเว ภิกฺขุ อริยํ อฏฺฐงฺคิกํ มคฺคํ ภาเวนฺโต อริยํ อฏฺฐงฺคิกํ มคฺคํ พหุลีกโรนฺโต นิพฺพานนินฺโน โหติ นิพฺพานโปโณ นิพฺพานปพฺภาโร, อิธ ภิกฺขเว ภิกฺขุ สมฺมาทิฏฺฐิํ ภาเวติ นิพฺพานนินฺนํ นิพฺพานโปณํ นิพฺพานปพฺภารํ ฯเปฯ สมฺมาสมาธิํ ภาเวติ นิพฺพานนินฺนํ นิพฺพานโปณํ นิพฺพานปพฺภารํ.\csromanspacer{} เอวํ โข ภิกฺขเว ภิกฺขุ อริยํ อฏฺฐงฺคิกํ มคฺคํ ภาเวนฺโต อริยํ อฏฺฐงฺคิกํ มคฺคํ พหุลีกโรนฺโต นิพฺพานนินฺโน โหติ นิพฺพานโปโณ นิพฺพานปพฺภาโรติ.\csromanspacer{}\csromanspacer{}ปฐมํ.}

\csromanheader{2}{32-36. ทุติยาทิสมุทฺทนินฺนสุตฺต}
\proseitem{134-138}{เสยฺยถาปิ ภิกฺขเว ยมุนา นที สมุทฺทนินฺนา สมุทฺทโปณา สมุทฺทปพฺภารา.\csromanspacer{} เอวเมว โข ภิกฺขเว ภิกฺขุ ฯเปฯ.\csromanspacer{} เสยฺยถาปิ ภิกฺขเว อจิรวตี นที สมุทฺทนินฺนา สมุทฺทโปณา สมุทฺทปพฺภารา.\csromanspacer{} เอวเมว โข ภิกฺขเว ภิกฺขุ ฯเปฯ.\csromanspacer{} เสยฺยถาปิ ภิกฺขเว สรภู นที สมุทฺทนินฺนา สมุทฺทโปณา สมุทฺทปพฺภารา.\csromanspacer{} เอวเมว โข ภิกฺขเว ภิกฺขุ ฯเปฯ.\csromanspacer{} เสยฺยถาปิ ภิกฺขเว มหี นที \csromanfolio{38}สมุทฺทนินฺนา สมุทฺทโปณา สมุทฺทปพฺภารา.\csromanspacer{} เอวเมว โข ภิกฺขเว ภิกฺขุ ฯเปฯ.\csromanspacer{} เสยฺยถาปิ ภิกฺขเว ยา กาจิมา มหานทิโย.\csromanspacer{} เสยฺยถิทํ, คงฺคา ยมุนา อจิรวตี สรภู มหี, สพฺพา ตา สมุทฺทนินฺนา สมุทฺทโปณา สมุทฺทปพฺภารา.\csromanspacer{} เอวเมว โข ภิกฺขเว ภิกฺขุ อริยํ อฏฺฐงฺคิกํ มคฺคํ ภาเวนฺโต อริยํ อฏฺฐงฺคิกํ มคฺคํ พหุลีกโรนฺโต นิพฺพานนินฺโน โหติ นิพฺพานโปโณ นิพฺพานปพฺภาโร.\csromanspacer{} กถญฺจ ภิกฺขเว ภิกฺขุ อริยํ อฏฺฐงฺคิกํ มคฺคํ ภาเวนฺโต อริยํ อฏฺฐงฺคิกํ มคฺคํ พหุลีกโรนฺโต นิพฺพานนินฺโน โหติ นิพฺพานโปโณ นิพฺพานปพฺภาโร, อิธ ภิกฺขเว ภิกฺขุ สมฺมาทิฏฺฐิํ ภาเวติ นิพฺพานนินฺนํ นิพฺพานโปณํ นิพฺพานปพฺภารํ ฯเปฯ สมฺมาสมาธิํ ภาเวติ นิพฺพานนินฺนํ นิพฺพานโปณํ นิพฺพานปพฺภารํ.\csromanspacer{} เอวํ โข ภิกฺขเว ภิกฺขุ อริยํ อฏฺฐงฺคิกํ มคฺคํ ภาเวนฺโต อริยํ อฏฺฐงฺคิกํ มคฺคํ พหุลีกโรนฺโต นิพฺพานนินฺโน โหติ นิพฺพานโปโณ นิพฺพานปพฺภาโรติ.\csromanspacer{}\csromanspacer{}ฉฏฺฐํ. (คงฺคาเปยฺยาลี.)}

\titlehead{\textbf{ตสฺสุทฺทานํ}}
\csromangathameasure{ฉ ปาจีนโต นินฺนา, ฉ นินฺนา จ สมุทฺทโต. \\ เอเต เทฺว ฉ ทฺวาทส โหนฺติ, วคฺโค เตน ปวุจฺจตีติ. \\ นิพฺพานนินฺโน ทฺวาทสกี, จตุตฺถกี ฉฏฺฐา นวกี.}
\csromangathagroup{\csromangathastanza{ฉ ปาจีนโต นินฺนา, ฉ นินฺนา จ สมุทฺทโต. \\ เอเต เทฺว ฉ ทฺวาทส โหนฺติ, วคฺโค เตน ปวุจฺจตีติ. \\ นิพฺพานนินฺโน ทฺวาทสกี, จตุตฺถกี ฉฏฺฐา นวกี.}}

\chapterhead{5. อปฺปมาทเปยฺยาลวคฺค}
\csromanheader{2}{1. ตถาคตสุตฺต}
\proseitem{139}{สาวตฺถินิทานํ.\csromanspacer{} ยาวตา ภิกฺขเว สตฺตา อปทา วา ทฺวิปทา\footnote{ทิปทา (สี)} วา จตุปฺปทา วา พหุปฺปทา\footnote{พหุปทา (?)} วา รูปิโน วา อรูปิโน วา สญฺญิโน วา อสญฺญิโน วา เนวสญฺญีนาสญฺญิโน วา, ตถาคโต เตสํ อคฺคมกฺขายติ อรหํ สมฺมาสมฺพุทฺโธ.\csromanspacer{} เอวเมว โข ภิกฺขเว เย เกจิ กุสลา ธมฺมา สพฺเพ เต อปฺปมาทมูลกา อปฺปมาทสโมสรณา, อปฺปมาโท เตสํ ธมฺมานํ อคฺคมกฺขายติ.\csromanspacer{} อปฺปมตฺตสฺเสตํ ภิกฺขเว ภิกฺขุโน ปาฏิกงฺขํ อริยํ อฏฺฐงฺคิกํ มคฺคํ ภาเวสฺสติ อริยํ อฏฺฐงฺคิกํ \csromanfolio{39}มคฺคํ พหุลีกริสฺสติ.\csromanspacer{} กถญฺจ ภิกฺขเว ภิกฺขุ อปฺปมตฺโต อริยํ อฏฺฐงฺคิกํ มคฺคํ ภาเวติ อริยํ อฏฺฐงฺคิกํ มคฺคํ พหุลีกโรติ, อิธ ภิกฺขเว ภิกฺขุ สมฺมาทิฏฺฐิํ ภาเวติ วิเวกนิสฺสิตํ วิราคนิสฺสิตํ นิโรธนิสฺสิตํ โวสฺสคฺคปริณามิํ ฯเปฯ สมฺมาสมาธิํ ภาเวติ วิเวกนิสฺสิตํ วิราคนิสฺสิตํ นิโรธนิสฺสิตํ โวสฺสคฺคปริณามิํ.\csromanspacer{} เอวํ โข ภิกฺขเว ภิกฺขุ อปฺปมตฺโต อริยํ อฏฺฐงฺคิกํ มคฺคํ ภาเวติ อริยํ อฏฺฐงฺคิกํ มคฺคํ พหุลีกโรตีติ.}

\prose{ยาวตา ภิกฺขเว สตฺตา อปทา วา ทฺวิปทา วา จตุปฺปทา วา พหุปฺปทา วา รูปิโน วา อรูปิโน วา สญฺญิโน วา อสญฺญิโน วา เนวสญฺญีนาสญฺญิโน วา, ตถาคโต เตสํ อคฺคมกฺขายติ อรหํ สมฺมาสมฺพุทฺโธ.\csromanspacer{} เอวเมว โข ภิกฺขเว เย เกจิ กุสลา ธมฺมา สพฺเพ เต อปฺปมาทมูลกา อปฺปมาทสโมสรณา, อปฺปมาโท เตสํ ธมฺมานํ อคฺคมกฺขายติ.\csromanspacer{} อปฺปมตฺตสฺเสตํ ภิกฺขเว ภิกฺขุโน ปาฏิกงฺขํ อริยํ อฏฺฐงฺคิกํ มคฺคํ ภาเวสฺสติ อริยํ อฏฺฐงฺคิกํ มคฺคํ พหุลีกริสฺสติ.\csromanspacer{} กถญฺจ ภิกฺขเว ภิกฺขุ อปฺปมตฺโต อริยํ อฏฺฐงฺคิกํ มคฺคํ ภาเวติ, อริยํ อฏฺฐงฺคิกํ มคฺคํ พหุลีกโรติ, อิธ ภิกฺขเว ภิกฺขุ สมฺมาทิฏฺฐิํ ภาเวติ ราควินยปริโยสานํ โทสวินยปริโยสานํ โมหวินยปริโยสานํ ฯเปฯ สมฺมาสมาธิํ ภาเวติ ราควินยปริโยสานํ โทสวินยปริโยสานํ โมหวินยปริโยสานํ.\csromanspacer{} เอวํ โข ภิกฺขเว ภิกฺขุ อปฺปมตฺโต อริยํ อฏฺฐงฺคิกํ มคฺคํ ภาเวติ อริยํ อฏฺฐงฺคิกํ มคฺคํ พหุลีกโรตีติ.}

\prose{ยาวตา ภิกฺขเว สตฺตา อปทา วา ทฺวิปทา วา จตุปฺปทา วา พหุปฺปทา วา รูปิโน วา อรูปิโน วา สญฺญิโน วา อสญฺญิโน วา เนวสญฺญีนาสญฺญิโน วา, ตถาคโต เตสํ อคฺคมกฺขายติ อรหํ สมฺมาสมฺพุทฺโธ.\csromanspacer{} เอวเมว โข ภิกฺขเว เย เกจิ กุสลา ธมฺมา สพฺเพ เต อปฺปมาทมูลกา อปฺปมาทสโมสรณา, อปฺปมาโท เตสํ ธมฺมานํ อคฺคมกฺขายติ.\csromanspacer{} อปฺปมตฺตสฺเสตํ ภิกฺขเว ภิกฺขุโน ปาฏิกงฺขํ อริยํ อฏฺฐงฺคิกํ มคฺคํ ภาเวสฺสติ อริยํ อฏฺฐงฺคิกํ มคฺคํ พหุลีกริสฺสติ.\csromanspacer{} กถญฺจ ภิกฺขเว ภิกฺขุ อปฺปมตฺโต อริยํ อฏฺฐงฺคิกํ มคฺคํ ภาเวติ อริยํ อฏฺฐงฺคิกํ มคฺคํ พหุลีกโรติ, อิธ ภิกฺขเว ภิกฺขุ สมฺมาทิฏฺฐิํ ภาเวติ อมโตคธํ อมตปรายนํ อมตปริโยสานํ ฯเปฯ สมฺมาสมาธิํ ภาเวติ อมโตคธํ อมตปรายนํ อมตปริโยสานํ.\csromanspacer{} เอวํ โข ภิกฺขเว \csromanfolio{40}ภิกฺขุ อปฺปมตฺโต อริยํ อฏฺฐงฺคิกํ มคฺคํ ภาเวติ อริยํ อฏฺฐงฺคิกํ มคฺคํ พหุลีกโรตีติ.}

\prose{ยาวตา ภิกฺขเว สตฺตา อปทา วา ทฺวิปทา วา จตุปฺปทา วา พหุปฺปทา วา รูปิโน วา อรูปิโน วา สญฺญิโน วา อสญฺญิโน วา เนวสญฺญีนาสญฺญิโน วา, ตถาคโต เตสํ อคฺคมกฺขายติ อรหํ สมฺมาสมฺพุทฺโธ.\csromanspacer{} เอวเมว โข ภิกฺขเว เย เกจิ กุสลา ธมฺมา สพฺเพ เต อปฺปมาทมูลกา อปฺปมาทสโมสรณา, อปฺปมาโท เตสํ ธมฺมานํ อคฺคมกฺขายติ.\csromanspacer{} อปฺปมตฺตสฺเสตํ ภิกฺขเว ภิกฺขุโน ปาฏิกงฺขํ อริยํ อฏฺฐงฺคิกํ มคฺคํ ภาเวสฺสติ อริยํ อฏฺฐงฺคิกํ มคฺคํ พหุลีกริสฺสติ.\csromanspacer{} กถญฺจ ภิกฺขเว ภิกฺขุ อปฺปมตฺโต อริยํ อฏฺฐงฺคิกํ มคฺคํ ภาเวติ อริยํ อฏฺฐงฺคิกํ มคฺคํ พหุลีกโรติ, อิธ ภิกฺขเว ภิกฺขุ สมฺมาทิฏฺฐิํ ภาเวติ นิพฺพานนินฺนํ นิพฺพานโปณํ นิพฺพานปพฺภารํ ฯเปฯ สมฺมาสมาธิํ ภาเวติ นิพฺพานนินฺนํ นิพฺพานโปณํ นิพฺพานปพฺภารํ.\csromanspacer{} เอวํ โข ภิกฺขเว ภิกฺขุ อปฺปมตฺโต อริยํ อฏฺฐงฺคิกํ มคฺคํ ภาเวติ อริยํ อฏฺฐงฺคิกํ มคฺคํ พหุลีกโรตีติ.\csromanspacer{}\csromanspacer{}ปฐมํ.}

\csromanheader{2}{2. ปทสุตฺต}
\proseitem{140}{เสยฺยถาปิ ภิกฺขเว ยานิ กานิจิ ชงฺคลานํ ปาณานํ ปทชาตานิ สพฺพานิ ตานิ หตฺถิปเท สโมธานํ คจฺฉนฺติ, หตฺถิปทํ เตสํ อคฺคมกฺขายติ ยทิทํ มหนฺตตฺเตน.\csromanspacer{} เอวเมว โข ภิกฺขเว เย เกจิ กุสลา ธมฺมา สพฺเพ เต อปฺปมาทมูลกา อปฺปมาทสโมสรณา, อปฺปมาโท เตสํ ธมฺมานํ อคฺคมกฺขายติ.\csromanspacer{} อปฺปมตฺตสฺเสตํ ภิกฺขเว ภิกฺขุโน ปาฏิกงฺขํ อริยํ อฏฺฐงฺคิกํ มคฺคํ ภาเวสฺสติ อริยํ อฏฺฐงฺคิกํ มคฺคํ พหุลีกริสฺสติ.\csromanspacer{} กถญฺจ ภิกฺขเว ภิกฺขุ อปฺปมตฺโต อริยํ อฏฺฐงฺคิกํ มคฺคํ ภาเวติ อริยํ อฏฺฐงฺคิกํ มคฺคํ พหุลีกโรติ, อิธ ภิกฺขเว ภิกฺขุ สมฺมาทิฏฺฐิํ ภาเวติ วิเวกนิสฺสิตํ วิราคนิสฺสิตํ นิโรธนิสฺสิตํ โวสฺสคฺคปริณามิํ ฯเปฯ สมฺมาสมาธิํ ภาเวติ วิเวกนิสฺสิตํ วิราคนิสฺสิตํ นิโรธนิสฺสิตํ โวสฺสคฺคปริณามิํ.\csromanspacer{} เอวํ โข ภิกฺขเว ภิกฺขุ อปฺปมตฺโต อริยํ อฏฺฐงฺคิกํ มคฺคํ ภาเวติ อริยํ อฏฺฐงฺคิกํ มคฺคํ พหุลีกโรตีติ.\csromanspacer{}\csromanspacer{}ทุติยํ.}

\csromanheader{2}{3-7. กูฏาทิสุตฺต}
\proseitem{141-145}{เสยฺยถาปิ ภิกฺขเว กูฏาคารสฺส ยา กาจิ โคปานสิโย สพฺพา ตา กูฏงฺคมา กูฏนินฺนา กูฏสโมสรณา, กูฏํ ตาสํ \csromanfolio{41}อคฺคมกฺขายติ.\csromanspacer{} เอวเมว โข ภิกฺขเว. (ยถา เหฏฺฐิมสุตฺตนฺตํ, เอวํ วิตฺถาเรตพฺพํ.) เสยฺยถาปิ ภิกฺขเว เย เกจิ มูลคนฺธา, กาฬานุสาริยํ เตสํ อคฺคมกฺขายติ.\csromanspacer{} เอวเมว โข ภิกฺขเว ฯเปฯ.\csromanspacer{} เสยฺยถาปิ ภิกฺขเว เย เกจิ สารคนฺธา, โลหิตจนฺทนํ เตสํ อคฺคมกฺขายติ.\csromanspacer{} เอวเมว โข ภิกฺขเว ฯเปฯ.\csromanspacer{} เสยฺยถาปิ ภิกฺขเว เย เกจิ ปุปฺผคนฺธา, วสฺสิกํ เตสํ อคฺคมกฺขายติ.\csromanspacer{} เอวเมว โข ภิกฺขเว ฯเปฯ.\csromanspacer{} เสยฺยถาปิ ภิกฺขเว เย เกจิ กุฏฺฏราชาโน สพฺเพ เต รญฺโญ จกฺกวตฺติสฺส อนุยนฺตา ภวนฺติ, ราชา เตสํ จกฺกวตฺติ อคฺคมกฺขายติ.\csromanspacer{} เอวเมว โข ภิกฺขเว ฯเปฯ.\csromanspacer{}\csromanspacer{}สตฺตมํ.}

\csromanheader{2}{8-10. จนฺทิมาทิสุตฺต}
\proseitem{146-148}{เสยฺยถาปิ ภิกฺขเว ยา กาจิ ตารกรูปานํ ปภา, สพฺพา ตา จนฺทิมปฺปภาย\footnote{จนฺทิมาปภาย (สฺยา, ก)} กลํ นาคฺฆนฺติ โสฬสิํ, จนฺทปฺปภา ตาสํ อคฺคมกฺขายติ.\csromanspacer{} เอวเมว โข ภิกฺขเว ฯเปฯ.\csromanspacer{} เสยฺยถาปิ ภิกฺขเว สรทสมเย วิทฺเธ วิคตวลาหเก เทเว อาทิจฺโจ นภํ อพฺภุสฺสกฺกมาโน สพฺพํ อากาสคตํ ตมคตํ อภิวิหจฺจ ภาสเต จ ตปเต จ วิโรจติ จ.\csromanspacer{} เอวเมว โข ภิกฺขเว ฯเปฯ.\csromanspacer{} เสยฺยถาปิ ภิกฺขเว ยานิ กานิจิ ตนฺตาวุตานํ วตฺถานํ, กาสิกวตฺถํ เตสํ อคฺคมกฺขายติ.\csromanspacer{} เอวเมว โข ภิกฺขเว เย เกจิ กุสลา ธมฺมา สพฺเพ เต อปฺปมาทมูลกา อปฺปมาทสโมสรณา, อปฺปมาโท เตสํ ธมฺมานํ อคฺคมกฺขายติ.\csromanspacer{} อปฺปมตฺตสฺเสตํ ภิกฺขเว ภิกฺขุโน ปาฏิกงฺขํ อริยํ อฏฺฐงฺคิกํ มคฺคํ ภาเวสฺสติ อริยํ อฏฺฐงฺคิกํ มคฺคํ พหุลีกริสฺสติ.\csromanspacer{} กถญฺจ ภิกฺขเว ภิกฺขุ อปฺปมตฺโต อริยํ อฏฺฐงฺคิกํ มคฺคํ ภาเวติ อริยํ อฏฺฐงฺคิกํ มคฺคํ พหุลีกโรติ, อิธ ภิกฺขเว ภิกฺขุ สมฺมาทิฏฺฐิํ ภาเวติ วิเวกนิสฺสิตํ วิราคนิสฺสิตํ นิโรธนิสฺสิตํ โวสฺสคฺคปริณามิํ ฯเปฯ สมฺมาสมาธิํ ภาเวติ วิเวกนิสฺสิตํ วิราคนิสฺสิตํ นิโรธนิสฺสิตํ โวสฺสคฺคปริณามิํ.\csromanspacer{} เอวํ โข ภิกฺขเว ภิกฺขุ อปฺปมตฺโต อริยํ อฏฺฐงฺคิกํ มคฺคํ ภาเวติ อริยํ อฏฺฐงฺคิกํ มคฺคํ พหุลีกโรตีติ.\csromanspacer{}\csromanspacer{}ทสมํ. (ยทปิ ตถาคตํ, ตทปิ วิตฺถาเรตพฺพํ.) อปฺปมาทวคฺโค ปญฺจโม.}

\titlehead{\textbf{ตสฺสุทฺทานํ}}
\csromangathameasure{ตถาคตํ ปทํ กูฏํ, มูลํ สาโร จ วสฺสิกํ. \\ ราชา จนฺทิมสูริยา จ, วตฺเถน ทสมํ ปทํ.}
\csromangathagroup{\csromangathastanza{\csromanfolio{42}ตถาคตํ ปทํ กูฏํ, มูลํ สาโร จ วสฺสิกํ. \\ ราชา จนฺทิมสูริยา จ, วตฺเถน ทสมํ ปทํ.}}

\csromanmatikaanchor{mtk.2}
\csromanmatikaanchor{mtk.27}
\csromantocmarkref{subsection}{1-2. กุสลสุตฺตานิ}{mtk.2}
\bookmark[dest={mtk.2},level=2]{1-2. กุสลสุตฺตานิ}
\csromanheader{4}{6. พลกรณียวคฺค}
\csromanheader{2}{1. พลสุตฺต}
\proseitem{149}{สาวตฺถินิทานํ.\csromanspacer{} เสยฺยถาปิ ภิกฺขเว เย เกจิ พลกรณียา กมฺมนฺตา กรียนฺติ, สพฺเพ เต ปถวิํ นิสฺสาย ปถวิยํ ปติฏฺฐาย เอวเมเต พลกรณียา กมฺมนฺตา กรียนฺติ.\csromanspacer{} เอวเมว โข ภิกฺขเว ภิกฺขุ สีลํ นิสฺสาย สีเล ปติฏฺฐาย อริยํ อฏฺฐงฺคิกํ มคฺคํ ภาเวติ อริยํ อฏฺฐงฺคิกํ มคฺคํ พหุลีกโรติ.\csromanspacer{} กถญฺจ ภิกฺขเว ภิกฺขุ สีลํ นิสฺสาย สีเล ปติฏฺฐาย อริยํ อฏฺฐงฺคิกํ มคฺคํ ภาเวติ อริยํ อฏฺฐงฺคิกํ มคฺคํ พหุลีกโรติ, อิธ ภิกฺขเว ภิกฺขุ สมฺมาทิฏฺฐิํ ภาเวติ วิเวกนิสฺสิตํ วิราคนิสฺสิตํ นิโรธนิสฺสิตํ โวสฺสคฺคปริณามิํ.\csromanspacer{} เอวํ โข ภิกฺขเว ภิกฺขุ สีลํ นิสฺสาย สีเล ปติฏฺฐาย อริยํ อฏฺฐงฺคิกํ มคฺคํ ภาเวติ อริยํ อฏฺฐงฺคิกํ มคฺคํ พหุลีกโรตีติ. (ปรคงฺคาเปยฺยาลีวณฺณิยโต ปริปุณฺณสุตฺตนฺติ วิตฺถารมคฺคี.)}

\prose{เสยฺยถาปิ ภิกฺขเว เย เกจิ พลกรณียา กมฺมนฺตา กรียนฺติ, สพฺเพ เต ปถวิํ นิสฺสาย ปถวิยํ ปติฏฺฐาย เอวเมเต พลกรณียา กมฺมนฺตา กรียนฺติ.\csromanspacer{} เอวเมว โข ภิกฺขเว ภิกฺขุ สีลํ นิสฺสาย สีเล ปติฏฺฐาย อริยํ อฏฺฐงฺคิกํ มคฺคํ ภาเวติ อริยํ อฏฺฐงฺคิกํ มคฺคํ พหุลีกโรติ.\csromanspacer{} กถญฺจ ภิกฺขเว ภิกฺขุ สีลํ นิสฺสาย สีเล ปติฏฺฐาย อริยํ อฏฺฐงฺคิกํ มคฺคํ ภาเวติ อริยํ อฏฺฐงฺคิกํ มคฺคํ พหุลีกโรติ, อิธ ภิกฺขเว ภิกฺขุ สมฺมาทิฏฺฐิํ ภาเวติ ราควินยปริโยสานํ โทสวินยปริโยสานํ โมหวินยปริโยสานํ ฯเปฯ สมฺมาสมาธิํ ภาเวติ ราควินยปริโยสานํ โทสวินยปริโยสานํ โมหวินยปริโยสานํ.\csromanspacer{} เอวํ โข ภิกฺขเว ภิกฺขุ สีลํ นิสฺสาย สีเล ปติฏฺฐาย อริยํ อฏฺฐงฺคิกํ มคฺคํ ภาเวติ อริยํ อฏฺฐงฺคิกํ มคฺคํ พหุลีกโรตีติ.}

\prose{\csromanfolio{43}เสยฺยถาปิ ภิกฺขเว เย เกจิ พลกรณียา กมฺมนฺตา กรียนฺติ, สพฺเพ เต ปถวิํ นิสฺสาย ปถวิยํ ปติฏฺฐาย เอวเมเต พลกรณียา กมฺมนฺตา กรียนฺติ.\csromanspacer{} เอวเมว โข ภิกฺขเว ภิกฺขุ สีลํ นิสฺสาย สีเล ปติฏฺฐาย อริยํ อฏฺฐงฺคิกํ มคฺคํ ภาเวติ อริยํ อฏฺฐงฺคิกํ มคฺคํ พหุลีกโรติ.\csromanspacer{} กถญฺจ ภิกฺขเว ภิกฺขุ สีลํ นิสฺสาย สีเล ปติฏฺฐาย อริยํ อฏฺฐงฺคิกํ มคฺคํ ภาเวติ อริยํ อฏฺฐงฺคิกํ มคฺคํ พหุลีกโรติ.\csromanspacer{} อิธ ภิกฺขเว ภิกฺขุ สมฺมาทิฏฺฐิํ ภาเวติ อมโตคธํ อมตปรายนํ อมตปริโยสานํ ฯเปฯ สมฺมาสมาธิํ ภาเวติ อมโตคธํ อมตปรายนํ อมตปริโยสานํ.\csromanspacer{} เอวํ โข ภิกฺขเว ภิกฺขุ สีลํ นิสฺสาย สีเล ปติฏฺฐาย อริยํ อฏฺฐงฺคิกํ มคฺคํ ภาเวติ อริยํ อฏฺฐงฺคิกํ มคฺคํ พหุลีกโรตีติ.}

\prose{เสยฺยถาปิ ภิกฺขเว เย เกจิ พลกรณียา กมฺมนฺตา กรียนฺติ, สพฺเพ เต ปถวิํ นิสฺสาย ปถวิยํ ปติฏฺฐาย เอวเมเต พลกรณียา กมฺมนฺตา กรียนฺติ.\csromanspacer{} เอวเมว โข ภิกฺขเว ภิกฺขุ สีลํ นิสฺสาย สีเล ปติฏฺฐาย อริยํ อฏฺฐงฺคิกํ มคฺคํ ภาเวติ อริยํ อฏฺฐงฺคิกํ มคฺคํ พหุลีกโรติ.\csromanspacer{} กถญฺจ ภิกฺขเว ภิกฺขุ สีลํ นิสฺสาย สีเล ปติฏฺฐาย อริยํ อฏฺฐงฺคิกํ มคฺคํ ภาเวติ อริยํ อฏฺฐงฺคิกํ มคฺคํ พหุลีกโรติ, อิธ ภิกฺขเว ภิกฺขุ สมฺมาทิฏฺฐิํ ภาเวติ นิพฺพานนินฺนํ นิพฺพานโปณํ นิพฺพานปพฺภารํ ฯเปฯ สมฺมาสมาธิํ ภาเวติ นิพฺพานนินฺนํ นิพฺพานโปณํ นิพฺพานปพฺภารํ.\csromanspacer{} เอวํ โข ภิกฺขเว ภิกฺขุ สีลํ นิสฺสาย สีเล ปติฏฺฐาย อริยํ อฏฺฐงฺคิกํ มคฺคํ ภาเวติ อริยํ อฏฺฐงฺคิกํ มคฺคํ พหุลีกโรตีติ.\csromanspacer{} ปฐมํ.}

\csromanheader{2}{2. พีชสุตฺต}
\proseitem{150}{เสยฺยถาปิ ภิกฺขเว เย เกจิเม พีชคามภูตคามา วุฑฺฒิํ วิรูฬฺหิํ เวปุลฺลํ อาปชฺชนฺติ, สพฺเพ เต ปถวิํ นิสฺสาย ปถวิยํ ปติฏฺฐาย เอวเมเต พีชคามภูตคามา วุฑฺฒิํ วิรูฬฺหิํ เวปุลฺลํ อาปชฺชนฺติ.\csromanspacer{} เอวเมว โข ภิกฺขเว ภิกฺขุ สีลํ นิสฺสาย สีเล ปติฏฺฐาย อริยํ อฏฺฐงฺคิกํ มคฺคํ ภาเวนฺโต อริยํ อฏฺฐงฺคิกํ มคฺคํ พหุลีกโรนฺโต วุฑฺฒิํ วิรูฬฺหิํ เวปุลฺลํ ปาปุณาติ ธมฺเมสุ.\csromanspacer{} กถญฺจ ภิกฺขเว ภิกฺขุ สีลํ นิสฺสาย สีเล ปติฏฺฐาย อริยํ อฏฺฐงฺคิกํ มคฺคํ ภาเวนฺโต อริยํ อฏฺฐงฺคิกํ มคฺคํ พหุลีกโรนฺโต \csromanfolio{44}วุฑฺฒิํ วิรูฬฺหิํ เวปุลฺลํ ปาปุณาติ ธมฺเมสุ, อิธ ภิกฺขเว ภิกฺขุ สมฺมาทิฏฺฐิํ ภาเวติ วิเวกนิสฺสิตํ ฯเปฯ สมฺมาสมาธิํ ภาเวติ วิเวกนิสฺสิตํ วิราคนิสฺสิตํ นิโรธนิสฺสิตํ โวสฺสคฺคปริณามิํ.\csromanspacer{} เอวํ โข ภิกฺขเว ภิกฺขุ สีลํ นิสฺสาย สีเล ปติฏฺฐาย อริยํ อฏฺฐงฺคิกํ มคฺคํ ภาเวนฺโต อริยํ อฏฺฐงฺคิกํ มคฺคํ พหุลีกโรนฺโต วุฑฺฒิํ วิรูฬฺหิํ เวปุลฺลํ ปาปุณาติ ธมฺเมสูติ.\csromanspacer{}\csromanspacer{}ทุติยํ.}

\csromanheader{2}{3. นาคสุตฺต}
\proseitem{151}{เสยฺยถาปิ ภิกฺขเว หิมวนฺตํ ปพฺพตราชํ นิสฺสาย นาคา กายํ วฑฺเฒนฺติ, พลํ คาเหนฺติ.\csromanspacer{} เต ตตฺถ กายํ วฑฺเฒตฺวา พลํ คาเหตฺวา กุโสพฺเภ โอตรนฺติ, กุโสพฺเภ\footnote{กุสฺสุพฺเภ (สี, สฺยา), กุสุพฺเภ (อิ, ก)} โอตริตฺวา มหาโสพฺเภ โอตรนฺติ, มหาโสพฺเภ โอตริตฺวา กุนฺนทิโย โอตรนฺติ, กุนฺนทิโย โอตริตฺวา มหานทิโย โอตรนฺติ, มหานทิโย โอตริตฺวา มหาสมุทฺทํ\footnote{มหาสมุทฺทสาครํ (สพฺพตฺถ) สํ 1. 269 ปิฏฺเฐปิ.} โอตรนฺติ.\csromanspacer{} เต ตตฺถ มหนฺตตฺตํ เวปุลฺลตฺตํ อาปชฺชนฺติ กาเยน.\csromanspacer{} เอวเมว โข ภิกฺขเว ภิกฺขุ สีลํ นิสฺสาย สีเล ปติฏฺฐาย อริยํ อฏฺฐงฺคิกํ มคฺคํ ภาเวนฺโต อริยํ อฏฺฐงฺคิกํ มคฺคํ พหุลีกโรนฺโต มหนฺตตฺตํ เวปุลฺลตฺตํ ปาปุณาติ ธมฺเมสุ.\csromanspacer{} กถญฺจ ภิกฺขเว ภิกฺขุ สีลํ นิสฺสาย สีเล ปติฏฺฐาย อริยํ อฏฺฐงฺคิกํ มคฺคํ ภาเวนฺโต อริยํ อฏฺฐงฺคิกํ มคฺคํ พหุลีกโรนฺโต มหนฺตตฺตํ เวปุลฺลตฺตํ ปาปุณาติ ธมฺเมสุ, อิธ ภิกฺขเว ภิกฺขุ สมฺมาทิฏฺฐิํ ภาเวติ วิเวกนิสฺสิตํ วิราคนิสฺสิตํ นิโรธนิสฺสิตํ โวสฺสคฺคปริณามิํ ฯเปฯ สมฺมาสมาธิํ ภาเวติ วิเวกนิสฺสิตํ วิราคนิสฺสิตํ นิโรธนิสฺสิตํ โวสฺสคฺคปริณามิํ.\csromanspacer{} เอวํ โข ภิกฺขเว ภิกฺขุ สีลํ นิสฺสาย สีเล ปติฏฺฐาย อริยํ อฏฺฐงฺคิกํ มคฺคํ ภาเวนฺโต อริยํ อฏฺฐงฺคิกํ มคฺคํ พหุลีกโรนฺโต มหนฺตตฺตํ เวปุลฺลตฺตํ ปาปุณาติ ธมฺเมสูติ.\csromanspacer{}\csromanspacer{}ตติยํ.}

\csromanheader{2}{4. รุกฺขสุตฺต}
\proseitem{152}{เสยฺยถาปิ ภิกฺขเว รุกฺโข ปาจีนนินฺโน ปาจีนโปโณ ปาจีนปพฺภาโร, โส มูลจฺฉินฺโน\footnote{มูลจฺฉินฺเท กเต (สฺยา)} กตเมน ปปเตยฺยาติ.\csromanspacer{} เยน ภนฺเต นินฺโน เยน โปโณ เยน ปพฺภาโรติ.\csromanspacer{} เอวเมว โข ภิกฺขเว ภิกฺขุ อริยํ อฏฺฐงฺคิกํ มคฺคํ ภาเวนฺโต อริยํ อฏฺฐงฺคิกํ มคฺคํ พหุลีกโรนฺโต \csromanfolio{45}นิพฺพานนินฺโน โหติ นิพฺพานโปโณ นิพฺพานปพฺภาโร.\csromanspacer{} กถญฺจ ภิกฺขเว ภิกฺขุ อริยํ อฏฺฐงฺคิกํ มคฺคํ ภาเวนฺโต อริยํ อฏฺฐงฺคิกํ มคฺคํ พหุลีกโรนฺโต นิพฺพานนินฺโน โหติ นิพฺพานโปโณ นิพฺพานปพฺภาโร, อิธ ภิกฺขเว ภิกฺขุ สมฺมาทิฏฺฐิํ ภาเวติ วิเวกนิสฺสิตํ วิราคนิสฺสิตํ นิโรธนิสฺสิตํ โวสฺสคฺคปริณามิํ ฯเปฯ สมฺมาสมาธิํ ภาเวติ วิเวกนิสฺสิตํ วิราคนิสฺสิตํ นิโรธนิสฺสิตํ โวสฺสคฺคปริณามิํ.\csromanspacer{} เอวํ โข ภิกฺขเว ภิกฺขุ อริยํ อฏฺฐงฺคิกํ มคฺคํ ภาเวนฺโต อริยํ อฏฺฐงฺคิกํ มคฺคํ พหุลีกโรนฺโต นิพฺพานนินฺโน โหติ นิพฺพานโปโณ นิพฺพานปพฺภาโรติ.\csromanspacer{}\csromanspacer{}จตุตฺถํ.}

\csromanheader{2}{5. กุมฺภสุตฺต}
\proseitem{153}{เสยฺยถาปิ ภิกฺขเว กุมฺโภ นิกฺกุชฺโช วมเตว อุทกํ โน ปจฺจาวมติ.\csromanspacer{} เอวเมว โข ภิกฺขเว ภิกฺขุ อริยํ อฏฺฐงฺคิกํ มคฺคํ ภาเวนฺโต อริยํ อฏฺฐงฺคิกํ มคฺคํ พหุลีกโรนฺโต วมเตว ปาปเก อกุสเล ธมฺเม โน ปจฺจาวมติ.\csromanspacer{} กถญฺจ ภิกฺขเว ภิกฺขุ อริยํ อฏฺฐงฺคิกํ มคฺคํ ภาเวนฺโต อริยํ อฏฺฐงฺคิกํ มคฺคํ พหุลีกโรนฺโต วมเตว ปาปเก อกุสเล ธมฺเม โน ปจฺจาวมติ, อิธ ภิกฺขเว ภิกฺขุ สมฺมาทิฏฺฐิํ ภาเวติ วิเวกนิสฺสิตํ วิราคนิสฺสิตํ นิโรธนิสฺสิตํ โวสฺสคฺคปริณามิํ ฯเปฯ สมฺมาสมาธิํ ภาเวติ วิเวกนิสฺสิตํ วิราคนิสฺสิตํ นิโรธนิสฺสิตํ โวสฺสคฺคปริณามิํ.\csromanspacer{} เอวํ โข ภิกฺขเว ภิกฺขุ อริยํ อฏฺฐงฺคิกํ มคฺคํ ภาเวนฺโต อริยํ อฏฺฐงฺคิกํ มคฺคํ พหุลีกโรนฺโต วมเตว ปาปเก อกุสเล ธมฺเม โน ปจฺจาวมตีติ.\csromanspacer{}\csromanspacer{}ปญฺจมํ.}

\csromanheader{2}{6. สูกสุตฺต}
\proseitem{154}{เสยฺยถาปิ ภิกฺขเว สาลิสูกํ วา ยวสูกํ วา สมฺมาปณิหิตํ หตฺเถน วา ปาเทน วา อกฺกนฺตํ หตฺถํ วา ปาทํ วา ภินฺทิสฺสติ โลหิตํ วา อุปฺปาเทสฺสตีติ ฐานเมตํ วิชฺชติ.\csromanspacer{} ตํ กิสฺส เหตุ, สมฺมาปณิหิตตฺตา ภิกฺขเว สูกสฺส.\csromanspacer{} เอวเมว โข ภิกฺขเว ภิกฺขุ สมฺมาปณิหิตาย ทิฏฺฐิยา สมฺมาปณิหิตาย มคฺคภาวนาย อวิชฺชํ ภินฺทิสฺสติ วิชฺชํ อุปฺปาเทสฺสติ นิพฺพานํ สจฺฉิกริสฺสตีติ ฐานเมตํ วิชฺชติ.\csromanspacer{} ตํ กิสฺส เหตุ, สมฺมาปณิหิตตฺตา ภิกฺขเว ทิฏฺฐิยา.\csromanspacer{} กถญฺจ ภิกฺขเว ภิกฺขุ สมฺมาปณิหิตาย ทิฏฺฐิยา สมฺมาปณิหิตาย มคฺคภาวนาย อวิชฺชํ ภินฺทติ วิชฺชํ อุปฺปาเทติ นิพฺพานํ สจฺฉิกโรติ, อิธ ภิกฺขเว ภิกฺขุ สมฺมาทิฏฺฐิํ ภาเวติ วิเวกนิสฺสิตํ ฯเปฯ \csromanfolio{46}สมฺมาสมาธิํ ภาเวติ วิเวกนิสฺสิตํ วิราคนิสฺสิตํ นิโรธนิสฺสิตํ โวสฺสคฺคปริณามิํ.\csromanspacer{} เอวํ โข ภิกฺขเว ภิกฺขุ สมฺมาปณิหิตาย ทิฏฺฐิยา สมฺมาปณิหิตาย มคฺคภาวนาย อวิชฺชํ ภินฺทติ วิชฺชํ อุปฺปาเทติ นิพฺพานํ สจฺฉิกโรตีติ.\csromanspacer{}\csromanspacer{}ฉฏฺฐํ.}

\csromanheader{2}{7. อากาสสุตฺต}
\proseitem{155}{เสยฺยถาปิ ภิกฺขเว อากาเส วิวิธา วาตา วายนฺติ, ปุรตฺถิมาปิ วาตา วายนฺติ, ปจฺฉิมาปิ วาตา วายนฺติ, อุตฺตราปิ วาตา วายนฺติ, ทกฺขิณาปิ วาตา วายนฺติ, สรชาปิ วาตา วายนฺติ, อรชาปิ วาตา วายนฺติ, สีตาปิ วาตา วายนฺติ, อุณฺหาปิ วาตา วายนฺติ, ปริตฺตาปิ วาตา วายนฺติ, อธิมตฺตาปิ วาตา วายนฺติ.\csromanspacer{} เอวเมว โข ภิกฺขเว ภิกฺขุโน อริยํ อฏฺฐงฺคิกํ มคฺคํ ภาวยโต อริยํ อฏฺฐงฺคิกํ มคฺคํ พหุลีกโรโต จตฺตาโรปิ สติปฏฺฐานา ภาวนาปาริปูริํ คจฺฉนฺติ, จตฺตาโรปิ สมฺมปฺปธานา ภาวนาปาริปูริํ คจฺฉนฺติ, จตฺตาโรปิ อิทฺธิปาทา ภาวนาปาริปูริํ คจฺฉนฺติ, ปญฺจปิ อินฺทฺริยานิ ภาวนาปาริปูริํ คจฺฉนฺติ, ปญฺจปิ พลานิ ภาวนาปาริปูริํ คจฺฉนฺติ, สตฺตปิ โพชฺฌงฺคา ภาวนาปาริปูริํ คจฺฉนฺติ.\csromanspacer{} กถญฺจ ภิกฺขเว ภิกฺขุโน อริยํ อฏฺฐงฺคิกํ มคฺคํ ภาวยโต อริยํ อฏฺฐงฺคิกํ มคฺคํ พหุลีกโรโต จตฺตาโรปิ สติปฏฺฐานา ภาวนาปาริปูริํ คจฺฉนฺติ, จตฺตาโรปิ สมฺมปฺปธานา ภาวนาปาริปูริํ คจฺฉนฺติ, จตฺตาโรปิ อิทฺธิปาทา ภาวนาปาริปูริํ คจฺฉนฺติ, ปญฺจปิ อินฺทฺริยานิ ภาวนาปาริปูริํ คจฺฉนฺติ, ปญฺจปิ พลานิ ภาวนาปาริปูริํ คจฺฉนฺติ, สตฺตปิ โพชฺฌงฺคา ภาวนาปาริปูริํ คจฺฉนฺติ, อิธ ภิกฺขเว ภิกฺขุ สมฺมาทิฏฺฐิํ ภาเวติ ฯเปฯ สมฺมาสมาธิํ ภาเวติ วิเวกนิสฺสิตํ วิราคนิสฺสิตํ นิโรธนิสฺสิตํ โวสฺสคฺคปริณามิํ.\csromanspacer{} เอวํ โข ภิกฺขเว ภิกฺขุโน อริยํ อฏฺฐงฺคิกํ มคฺคํ ภาวยโต อริยํ อฏฺฐงฺคิกํ มคฺคํ พหุลีกโรโต จตฺตาโรปิ สติปฏฺฐานา ภาวนาปาริปูริํ คจฺฉนฺติ, จตฺตาโรปิ สมฺมปฺปธานา ภาวนาปาริปูริํ คจฺฉนฺติ, จตฺตาโรปิ อิทฺธิปาทา ภาวนาปาริปูริํ คจฺฉนฺติ, ปญฺจปิ อินฺทฺริยานิ ภาวนาปาริปูริํ คจฺฉนฺติ, ปญฺจปิ พลานิ ภาวนาปาริปูริํ คจฺฉนฺติ, สตฺตปิ โพชฺฌงฺคา ภาวนาปาริปูริํ คจฺฉนฺตีติ.\csromanspacer{}\csromanspacer{}สตฺตมํ.}

\csromanheader{2}{8. ปฐมเมฆสุตฺต}
\proseitem{156}{เสยฺยถาปิ ภิกฺขเว คิมฺหานํ ปจฺฉิเม มาเส อูหตํ รโชชลฺลํ, ตเมนํ มหา-อกาลเมโฆ ฐานโส อนฺตรธาเปติ วูปสเมติ. \csromanfolio{47}เอวเมว โข ภิกฺขเว ภิกฺขุ อริยํ อฏฺฐงฺคิกํ มคฺคํ ภาเวนฺโต อริยํ อฏฺฐงฺคิกํ มคฺคํ พหุลีกโรนฺโต อุปฺปนฺนุปฺปนฺเน ปาปเก อกุสเล ธมฺเม ฐานโส อนฺตรธาเปติ วูปสเมติ.\csromanspacer{} กถญฺจ ภิกฺขเว ภิกฺขุ อริยํ อฏฺฐงฺคิกํ มคฺคํ ภาเวนฺโต อริยํ อฏฺฐงฺคิกํ มคฺคํ พหุลีกโรนฺโต อุปฺปนฺนุปฺปนฺเน ปาปเก อกุสเล ธมฺเม ฐานโส อนฺตรธาเปติ วูปสเมติ, อิธ ภิกฺขเว ภิกฺขุ สมฺมาทิฏฺฐิํ ภาเวติ ฯเปฯ สมฺมาสมาธิํ ภาเวติ วิเวกนิสฺสิตํ วิราคนิสฺสิตํ นิโรธนิสฺสิตํ โวสฺสคฺคปริณามิํ.\csromanspacer{} เอวํ โข ภิกฺขเว ภิกฺขุ อริยํ อฏฺฐงฺคิกํ มคฺคํ ภาเวนฺโต อริยํ อฏฺฐงฺคิกํ มคฺคํ พหุลีกโรนฺโต อุปฺปนฺนุปฺปนฺเน ปาปเก อกุสเล ธมฺเม ฐานโส อนฺตรธาเปติ วูปสเมตีติ.\csromanspacer{}\csromanspacer{}อฏฺฐมํ.}

\csromanheader{2}{9. ทุติยเมฆสุตฺต}
\proseitem{157}{เสยฺยถาปิ ภิกฺขเว อุปฺปนฺนํ มหาเมฆํ, ตเมนํ มหาวาโต อนฺตราเยว อนฺตรธาเปติ วูปสเมติ.\csromanspacer{} เอวเมว โข ภิกฺขเว ภิกฺขุ อริยํ อฏฺฐงฺคิกํ มคฺคํ ภาเวนฺโต อริยํ อฏฺฐงฺคิกํ มคฺคํ พหุลีกโรนฺโต อุปฺปนฺนุปฺปนฺเน ปาปเก อกุสเล ธมฺเม อนฺตราเยว อนฺตรธาเปติ วูปสเมติ.\csromanspacer{} กถญฺจ ภิกฺขเว ภิกฺขุ อริยํ อฏฺฐงฺคิกํ มคฺคํ ภาเวนฺโต อริยํ อฏฺฐงฺคิกํ มคฺคํ พหุลีกโรนฺโต อุปฺปนฺนุปฺปนฺเน ปาปเก อกุสเล ธมฺเม อนฺตราเยว อนฺตรธาเปติ วูปสเมติ, อิธ ภิกฺขเว ภิกฺขุ สมฺมาทิฏฺฐิํ ภาเวติ ฯเปฯ สมฺมาสมาธิํ ภาเวติ วิเวกนิสฺสิตํ วิราคนิสฺสิตํ นิโรธนิสฺสิตํ โวสฺสคฺคปริณามิํ.\csromanspacer{} เอวํ โข ภิกฺขเว ภิกฺขุ อริยํ อฏฺฐงฺคิกํ มคฺคํ ภาเวนฺโต อริยํ อฏฺฐงฺคิกํ มคฺคํ พหุลีกโรนฺโต อุปฺปนฺนุปฺปนฺเน ปาปเก อกุสเล ธมฺเม อนฺตราเยว อนฺตรธาเปติ วูปสเมตีติ.\csromanspacer{}\csromanspacer{}นวมํ.}

\csromanheader{2}{10. นาวาสุตฺต}
\proseitem{158}{เสยฺยถาปิ ภิกฺขเว สามุทฺทิกาย นาวาย เวตฺตพนฺธนพนฺธาย ฉ มาสานิ อุทเก ปริยาทาย\footnote{ปริยาตาย (ก), ปริยาหตาย (?)} เหมนฺติเกน ถลํ อุกฺขิตฺตาย วาตาตปปเรตานิ พนฺธนานิ ตานิ ปาวุสฺสเกน เมเฆน อภิปฺปวุฏฺฐานิ อปฺปกสิเรเนว ปฏิปฺปสฺสมฺภนฺติ ปูติกานิ ภวนฺติ.\csromanspacer{} เอวเมว โข ภิกฺขเว ภิกฺขุโน อริยํ อฏฺฐงฺคิกํ มคฺคํ ภาวยโต อริยํ อฏฺฐงฺคิกํ มคฺคํ พหุลีกโรโต อปฺปกสิเรเนว สํโยชนานิ ปฏิปฺปสฺสมฺภนฺติ ปูติกานิ \csromanfolio{48}ภวนฺติ.\csromanspacer{} กถญฺจ ภิกฺขเว ภิกฺขุโน อริยํ อฏฺฐงฺคิกํ มคฺคํ ภาวยโต อริยํ อฏฺฐงฺคิกํ มคฺคํ พหุลีกโรโต อปฺปกสิเรเนว สํโยชนานิ ปฏิปฺปสฺสมฺภนฺติ ปูติกานิ ภวนฺติ, อิธ ภิกฺขเว ภิกฺขุ สมฺมาทิฏฺฐิํ ภาเวติ ฯเปฯ สมฺมาสมาธิํ ภาเวติ วิเวกนิสฺสิตํ วิราคนิสฺสิตํ นิโรธนิสฺสิตํ โวสฺสคฺคปริณามิํ.\csromanspacer{} เอวํ โข ภิกฺขเว ภิกฺขุโน อริยํ อฏฺฐงฺคิกํ มคฺคํ ภาวยโต อริยํ อฏฺฐงฺคิกํ มคฺคํ พหุลีกโรโต อปฺปกสิเรเนว สํโยชนานิ ปฏิปฺปสฺสมฺภนฺติ ปูติกานิ ภวนฺตีติ.\csromanspacer{}\csromanspacer{}ทสมํ.}

\csromanheader{2}{11. อาคนฺตุกสุตฺต}
\proseitem{159}{เสยฺยถาปิ ภิกฺขเว อาคนฺตุกาคารํ, ตตฺถ ปุรตฺถิมายปิ ทิสาย อาคนฺตฺวา วาสํ กปฺเปนฺติ, ปจฺฉิมายปิ ทิสาย อาคนฺตฺวา วาสํ กปฺเปนฺติ, อุตฺตรายปิ ทิสาย อาคนฺตฺวา วาสํ กปฺเปนฺติ, ทกฺขิณายปิ ทิสาย อาคนฺตฺวา วาสํ กปฺเปนฺติ, ขตฺติยาปิ อาคนฺตฺวา วาสํ กปฺเปนฺติ, พฺราหฺมณาปิ อาคนฺตฺวา วาสํ กปฺเปนฺติ, เวสฺสาปิ อาคนฺตฺวา วาสํ กปฺเปนฺติ, สุทฺทาปิ อาคนฺตฺวา วาสํ กปฺเปนฺติ.\csromanspacer{} เอวเมว โข ภิกฺขเว ภิกฺขุ อริยํ อฏฺฐงฺคิกํ มคฺคํ ภาเวนฺโต อริยํ อฏฺฐงฺคิกํ มคฺคํ พหุลีกโรนฺโต เย ธมฺมา อภิญฺญา ปริญฺเญยฺยา, เต ธมฺเม อภิญฺญา ปริชานาติ ฯเปฯ เย ธมฺมา อภิญฺญา ปหาตพฺพา, เต ธมฺเม อภิญฺญา ปชหติ.\csromanspacer{} เย ธมฺมา อภิญฺญา สจฺฉิกาตพฺพา, เต ธมฺเม อภิญฺญา สจฺฉิกโรติ.\csromanspacer{} เย ธมฺมา อภิญฺญา ภาเวตพฺพา, เต ธมฺเม อภิญฺญา ภาเวติ.}

\prose{กตเม จ ภิกฺขเว ธมฺมา อภิญฺญา ปริญฺเญยฺยา, ปญฺจุปาทานกฺขนฺธาติสฺส วจนียํ.\csromanspacer{} กตเม ปญฺจ, เสยฺยถิทํ, รูปุปาทานกฺขนฺโธ ฯเปฯ วิญฺญาณุปาทานกฺขนฺโธ, อิเม ภิกฺขเว ธมฺมา อภิญฺญา ปริญฺเญยฺยา.\csromanspacer{} กตเม จ ภิกฺขเว ธมฺมา อภิญฺญา ปหาตพฺพา, อวิชฺชา จ ภวตณฺหา จ, อิเม ภิกฺขเว ธมฺมา อภิญฺญา ปหาตพฺพา.\csromanspacer{} กตเม จ ภิกฺขเว ธมฺมา อภิญฺญา สจฺฉิกาตพฺพา, วิชฺชา จ วิมุตฺติ จ, อิเม ภิกฺขเว ธมฺมา อภิญฺญา สจฺฉิกาตพฺพา.\csromanspacer{} กตเม จ ภิกฺขเว ธมฺมา อภิญฺญา ภาเวตพฺพา, สมโถ จ วิปสฺสนา จ, อิเม ภิกฺขเว ธมฺมา อภิญฺญา ภาเวตพฺพา.\csromanspacer{} กถญฺจ ภิกฺขเว ภิกฺขุ อริยํ อฏฺฐงฺคิกํ มคฺคํ ภาเวนฺโต อริยํ อฏฺฐงฺคิกํ มคฺคํ พหุลีกโรนฺโต เย ธมฺมา อภิญฺญา ปริญฺเญยฺยา, เต ธมฺเม อภิญฺญา ปริชานาติ. \csromanfolio{49}เย ธมฺมา อภิญฺญา ภาเวตพฺพา, เต ธมฺเม อภิญฺญา ภาเวติ, อิธ ภิกฺขเว ภิกฺขุ สมฺมาทิฏฺฐิํ ภาเวติ ฯเปฯ สมฺมาสมาธิํ ภาเวติ วิเวกนิสฺสิตํ วิราคนิสฺสิตํ นิโรธนิสฺสิตํ โวสฺสคฺคปริณามิํ.\csromanspacer{} เอวํ โข ภิกฺขเว ภิกฺขุ อริยํ อฏฺฐงฺคิกํ มคฺคํ ภาเวนฺโต อริยํ อฏฺฐงฺคิกํ มคฺคํ พหุลีกโรนฺโต เย ธมฺมา อภิญฺญา ปริญฺเญยฺยา, เต ธมฺเม อภิญฺญา ปริชานาติ.\csromanspacer{} เย ธมฺมา อภิญฺญา ปหาตพฺพา, เต ธมฺเม อภิญฺญา ปชหติ.\csromanspacer{} เย ธมฺมา อภิญฺญา สจฺฉิกาตพฺพา, เต ธมฺเม อภิญฺญา สจฺฉิกโรติ.\csromanspacer{} เย ธมฺมา อภิญฺญา ภาเวตพฺพา, เต ธมฺเม อภิญฺญา ภาเวตีติ.\csromanspacer{}\csromanspacer{}เอกาทสมํ.}

\csromanheader{2}{12. นทีสุตฺต}
\proseitem{160}{เสยฺยถาปิ ภิกฺขเว คงฺคา นที ปาจีนนินฺนา ปาจีนโปณา ปาจีนปพฺภารา, อถ มหาชนกาโย อาคจฺเฉยฺย กุทาล ปิฏกํ อาทาย มยํ อิมํ คงฺคํ นทิํ ปจฺฉานินฺนํ กริสฺสาม ปจฺฉาโปณํ ปจฺฉาปพฺภารนฺติ.\csromanspacer{} ตํ กิํ มญฺญถ ภิกฺขเว, อปิ นุ โส มหาชนกาโย คงฺคํ นทิํ ปจฺฉานินฺนํ กเรยฺย ปจฺฉาโปณํ ปจฺฉาปพฺภารนฺติ.\csromanspacer{} โน เหตํ ภนฺเต.\csromanspacer{} ตํ กิสฺส เหตุ, คงฺคา ภนฺเต นที ปาจีนนินฺนา ปาจีนโปณา ปาจีนปพฺภารา.\csromanspacer{} สา น สุกรา ปจฺฉานินฺนํ กาตุํ ปจฺฉาโปณํ ปจฺฉาปพฺภารํ, ยาวเทว ปน โส มหาชนกาโย กิลมถสฺส วิฆาตสฺส ภาคี อสฺสาติ.\csromanspacer{} เอวเมว โข ภิกฺขเว ภิกฺขุํ อริยํ อฏฺฐงฺคิกํ มคฺคํ ภาเวนฺตํ อริยํ อฏฺฐงฺคิกํ มคฺคํ พหุลีกโรนฺตํ ราชาโน วา ราชมหามตฺตา วา มิตฺตา วา อมจฺจา วา ญาตี วา ญาติสาโลหิตา วา โภเคหิ อภิหฏฺฐุํ ปวาเรยฺยุํ “เอหมฺโภ ปุริส กิํ เต อิเม กาสาวา อนุทหนฺติ, กิํ มุณฺโฑ กปาลมนุสํจรสิ, เอหิ หีนายาวตฺติตฺวา โภเค จ ภุญฺชสฺสุ, ปุญฺญานิ จ กโรหี”ติ.\csromanspacer{} โส วต ภิกฺขเว ภิกฺขุ อริยํ อฏฺฐงฺคิกํ มคฺคํ ภาเวนฺโต อริยํ อฏฺฐงฺคิกํ มคฺคํ พหุลีกโรนฺโต สิกฺขํ ปจฺจกฺขาย หีนายาวตฺติสฺสตีติ เนตํ ฐานํ วิชฺชติ.\csromanspacer{} ตํ กิสฺส เหตุ, ยํ หิ ตํ ภิกฺขเว จิตฺตํ ทีฆรตฺตํ วิเวกนินฺนํ วิเวกโปณํ วิเวกปพฺภารํ, ตํ วต หีนายาวตฺติสฺสตีติ เนตํ ฐานํ วิชฺชติ.\csromanspacer{} กถญฺจ ภิกฺขเว ภิกฺขุ อริยํ อฏฺฐงฺคิกํ มคฺคํ ภาเวติ อริยํ อฏฺฐงฺคิกํ มคฺคํ พหุลีกโรติ, อิธ ภิกฺขเว ภิกฺขุ สมฺมาทิฏฺฐิํ ภาเวติ วิเวกนิสฺสิตํ ฯเปฯ สมฺมาสมาธิํ ภาเวติ วิเวกนิสฺสิตํ วิราคนิสฺสิตํ นิโรธนิสฺสิตํ โวสฺสคฺคปริณามิํ.\csromanspacer{} เอวํ โข ภิกฺขเว \csromanfolio{50}ภิกฺขุ อริยํ อฏฺฐงฺคิกํ มคฺคํ ภาเวติ อริยํ อฏฺฐงฺคิกํ มคฺคํ พหุลีกโรตีติ. (ยทปิ พลกรณียํ, ตทปิ วิตฺถาเรตพฺพํ.).\csromanspacer{} ทฺวาทสมํ.\csromanspacer{} พลกรณียวคฺโค ฉฏฺโฐ.}

\titlehead{\textbf{ตสฺสุทฺทานํ}}
\csromangathameasure{พลํ พีชญฺจ นาโค จ, รุกฺโข กุมฺเภน สูกิยา. \\ อากาเสน จ เทฺว เมฆา, นาวา อาคนฺตุกา นทีติ.}
\csromangathagroup{\csromangathastanza{พลํ พีชญฺจ นาโค จ, รุกฺโข กุมฺเภน สูกิยา. \\ อากาเสน จ เทฺว เมฆา, นาวา อาคนฺตุกา นทีติ.}}

\csromanmatikaanchor{mtk.3}
\csromanmatikaanchor{mtk.29}
\csromantocmarkref{subsection}{3-7. อุปกฺกิเลสสุตฺตาทีนิ}{mtk.3}
\bookmark[dest={mtk.3},level=2]{3-7. อุปกฺกิเลสสุตฺตาทีนิ}
\csromanheader{4}{7. เอสนาวคฺค}
\csromanheader{2}{1. เอสนาสุตฺต}
\proseitem{161}{สาวตฺถินิทานํ.\csromanspacer{} ติสฺโส อิมา ภิกฺขเว เอสนา.\csromanspacer{} กตมา ติสฺโส, กาเมสนา ภเวสนา พฺรหฺมจริเยสนา, อิมา โข ภิกฺขเว ติสฺโส เอสนา.\csromanspacer{} อิมาสํ โข ภิกฺขเว ติสฺสนฺนํ เอสนานํ อภิญฺญาย อริโย อฏฺฐงฺคิโก มคฺโค ภาเวตพฺโพ.\csromanspacer{} กตโม อริโย อฏฺฐงฺคิโก มคฺโค, อิธ ภิกฺขเว ภิกฺขุ สมฺมาทิฏฺฐิํ ภาเวติ วิเวกนิสฺสิตํ ฯเปฯ สมฺมาสมาธิํ ภาเวติ วิเวกนิสฺสิตํ วิราคนิสฺสิตํ นิโรธนิสฺสิตํ โวสฺสคฺคปริณามิํ.\csromanspacer{} อิมาสํ โข ภิกฺขเว ติสฺสนฺนํ เอสนานํ อภิญฺญาย อยํ อริโย อฏฺฐงฺคิโก มคฺโค ภาเวตพฺโพติ.}

\prose{ติสฺโส อิมา โข ภิกฺขเว เอสนา.\csromanspacer{} กตมา ติสฺโส, กาเมสนา ภเวสนา พฺรหฺมจริเยสนา, อิมา โข ภิกฺขเว ติสฺโส เอสนา.\csromanspacer{} อิมาสํ โข ภิกฺขเว ติสฺสนฺนํ เอสนานํ อภิญฺญาย อริโย อฏฺฐงฺคิโก มคฺโค ภาเวตพฺโพ.\csromanspacer{} กตโม อริโย อฏฺฐงฺคิโก มคฺโค, อิธ ภิกฺขเว ภิกฺขุ สมฺมาทิฏฺฐิํ ภาเวติ ฯเปฯ สมฺมาสมาธิํ ภาเวติ ราควินยปริโยสานํ โทสวินยปริโยสานํ โมหวินยปริโยสานํ.\csromanspacer{} อิมาสํ โข ภิกฺขเว ติสฺสนฺนํ เอสนานํ อภิญฺญาย อยํ อริโย อฏฺฐงฺคิโก มคฺโค ภาเวตพฺโพติ.}

\prose{ติสฺโส อิมา โข ภิกฺขเว เอสนา.\csromanspacer{} กตมา ติสฺโส, กาเมสนา ภเวสนา พฺรหฺมจริเยสนา, อิมา โข ภิกฺขเว ติสฺโส เอสนา.\csromanspacer{} อิมาสํ \csromanfolio{51}โข ภิกฺขเว ติสฺสนฺนํ เอสนานํ อภิญฺญาย อริโย อฏฺฐงฺคิโก มคฺโค ภาเวตพฺโพ.\csromanspacer{} กตโม อริโย อฏฺฐงฺคิโก มคฺโค, อิธ ภิกฺขเว ภิกฺขุ สมฺมาทิฏฺฐิํ ภาเวติ ฯเปฯ สมฺมาสมาธิํ ภาเวติ อมโตคธํ อมตปรายนํ อมตปริโยสานํ.\csromanspacer{} อิมาสํ โข ภิกฺขเว ติสฺสนฺนํ เอสนานํ อภิญฺญาย อยํ อริโย อฏฺฐงฺคิโก มคฺโค ภาเวตพฺโพติ.}

\prose{ติสฺโส อิมา โข ภิกฺขเว เอสนา.\csromanspacer{} กตมา ติสฺโส, กาเมสนา ภเวสนา พฺรหฺมจริเยสนา, อิมา โข ภิกฺขเว ติสฺโส เอสนา.\csromanspacer{} อิมาสํ โข ภิกฺขเว ติสฺสนฺนํ เอสนานํ อภิญฺญาย อริโย อฏฺฐงฺคิโก มคฺโค ภาเวตพฺโพ.\csromanspacer{} กตโม อริโย อฏฺฐงฺคิโก มคฺโค, อิธ ภิกฺขเว ภิกฺขุ สมฺมาทิฏฺฐิํ ภาเวติ ฯเปฯ สมฺมาสมาธิํ ภาเวติ นิพฺพานนินฺนํ นิพฺพานโปณํ นิพฺพานปพฺภารํ.\csromanspacer{} อิมาสํ โข ภิกฺขเว ติสฺสนฺนํ เอสนานํ อภิญฺญาย อยํ อริโย อฏฺฐงฺคิโก มคฺโค ภาเวตพฺโพติ.}

\prose{ติสฺโส อิมา ภิกฺขเว เอสนา.\csromanspacer{} กตมา ติสฺโส, กาเมสนา ภเวสนา พฺรหฺมจริเยสนา, อิมา โข ภิกฺขเว ติสฺโส เอสนา.\csromanspacer{} อิมาสํ โข ภิกฺขเว ติสฺสนฺนํ เอสนานํ ปริญฺญาย ฯเปฯ อยํ อริโย อฏฺฐงฺคิโก มคฺโค ภาเวตพฺโพติ. (ยทปิ อภิญฺญา, ตทปิ ปริญฺญาย วิตฺถาเรตพฺพํ.)}

\prose{ติสฺโส อิมา ภิกฺขเว เอสนา.\csromanspacer{} กตมา ติสฺโส, กาเมสนา ภเวสนา พฺรหฺมจริเยสนา, อิมา โข ภิกฺขเว ติสฺโส เอสนา.\csromanspacer{} อิมาสํ โข ภิกฺขเว ติสฺสนฺนํ เอสนานํ ปริกฺขยาย ฯเปฯ อยํ อริโย อฏฺฐงฺคิโก มคฺโค ภาเวตพฺโพติ. (ยทปิ อภิญฺญา, ตทปิ ปริกฺขยาย วิตฺถาเรตพฺพํ.)}

\prose{ติสฺโส อิมา ภิกฺขเว เอสนา.\csromanspacer{} กตมา ติสฺโส, กาเมสนา ภเวสนา พฺรหฺมจริเยสนา, อิมา โข ภิกฺขเว ติสฺโส เอสนา.\csromanspacer{} อิมาสํ โข ภิกฺขเว ติสฺสนฺนํ เอสนานํ ปหานาย อริโย อฏฺฐงฺคิโก มคฺโค ภาเวตพฺโพ.\csromanspacer{} กตโม อริโย อฏฺฐงฺคิโก มคฺโค, อิธ ภิกฺขเว ภิกฺขุ สมฺมาทิฏฺฐิํ ภาเวติ ฯเปฯ สมฺมาสมาธิํ ภาเวติ วิเวกนิสฺสิตํ วิราคนิสฺสิตํ นิโรธนิสฺสิตํ โวสฺสคฺคปริณามิํ.\csromanspacer{} อิมาสํ โข ภิกฺขเว ติสฺสนฺนํ เอสนานํ ปหานาย อยํ อริโย อฏฺฐงฺคิโก มคฺโค ภาเวตพฺโพติ. (ยทปิ อภิญฺญา, ตทปิ ปหานาย วิตฺถาเรตพฺพํ.).\csromanspacer{} ปฐมํ.}

\csromanheader{2}{2. วิธาสุตฺต}
\proseitem{162}{\csromanfolio{52}ติสฺโส อิมา ภิกฺขเว วิธา.\csromanspacer{} กตมา ติสฺโส, เสยฺโยหมสฺมีติวิธา สทิโสหมสฺมีติวิธา หีโนหมสฺมีติวิธา, อิมา โข ภิกฺขเว ติสฺโส วิธา.\csromanspacer{} อิมาสํ โข ภิกฺขเว ติสฺสนฺนํ วิธานํ อภิญฺญาย ปริญฺญาย ปริกฺขยาย ปหานาย อริโย อฏฺฐงฺคิโก มคฺโค ภาเวตพฺโพ.\csromanspacer{} กตโม อริโย อฏฺฐงฺคิโก มคฺโค, อิธ ภิกฺขเว ภิกฺขุ สมฺมาทิฏฺฐิํ ภาเวติ ฯเปฯ สมฺมาสมาธิํ ภาเวติ วิเวกนิสฺสิตํ วิราคนิสฺสิตํ นิโรธนิสฺสิตํ โวสฺสคฺคปริณามิํ.\csromanspacer{} อิมาสํ โข ภิกฺขเว ติสฺสนฺนํ วิธานํ อภิญฺญาย ปริญฺญาย ปริกฺขยาย ปหานาย อยํ อริโย อฏฺฐงฺคิโก มคฺโค ภาเวตพฺโพติ. (ยถา เอสนา, เอวํ วิตฺถาเรตพฺพํ.).\csromanspacer{} ทุติยํ.}

\csromanheader{2}{3. อาสวสุตฺต}
\proseitem{163}{ตโยเม ภิกฺขเว อาสวา.\csromanspacer{} กตเม ตโย, กามาสโว ภวาสโว อวิชฺชาสโว, อิเม โข ภิกฺขเว ตโย อาสวา.\csromanspacer{} อิเมสํ โข ภิกฺขเว ติณฺณนฺนํ อาสวานํ อภิญฺญาย ปริญฺญาย ปริกฺขยาย ปหานาย ฯเปฯ อยํ อริโย อฏฺฐงฺคิโก มคฺโค ภาเวตพฺโพติ.\csromanspacer{}\csromanspacer{}ตติยํ.}

\csromanheader{2}{4. ภวสุตฺต}
\proseitem{164}{ตโยเม ภิกฺขเว ภวา.\csromanspacer{} กตเม ตโย, กามภโว รูปภโว อรูปภโว, อิเม โข ภิกฺขเว ตโย ภวา.\csromanspacer{} อิเมสํ โข ภิกฺขเว ติณฺณนฺนํ ภวานํ อภิญฺญาย ปริญฺญาย ปริกฺขยาย ปหานาย ฯเปฯ อยํ อริโย อฏฺฐงฺคิโก มคฺโค ภาเวตพฺโพติ.\csromanspacer{}\csromanspacer{}จตุตฺถํ.}

\csromanheader{2}{5. ทุกฺขตาสุตฺต}
\proseitem{165}{ติสฺโส อิมา ภิกฺขเว ทุกฺขตา.\csromanspacer{} กตมา ติสฺโส, ทุกฺขทุกฺขตา สงฺขารทุกฺขตา วิปริณามทุกฺขตา, อิมา โข ภิกฺขเว ติสฺโส ทุกฺขตา.\csromanspacer{} อิมาสํ โข ภิกฺขเว ติสฺสนฺนํ ทุกฺขตานํ อภิญฺญาย ปริญฺญาย ปริกฺขยาย ปหานาย ฯเปฯ อยํ อริโย อฏฺฐงฺคิโก มคฺโค ภาเวตพฺโพติ.\csromanspacer{}\csromanspacer{}ปญฺจมํ.}

\csromanheader{2}{6. ขิลสุตฺต}
\proseitem{166}{ตโยเม ภิกฺขเว ขิลา.\csromanspacer{} กตเม ตโย, ราโค ขิโล โทโส ขิโล โมโห ขิโล, อิเม โข ภิกฺขเว ตโย ขิลา. \csromanfolio{53}อิเมสํ โข ภิกฺขเว ติณฺณนฺนํ ขิลานํ อภิญฺญาย ปริญฺญาย ปริกฺขยาย ปหานาย ฯเปฯ อยํ อริโย อฏฺฐงฺคิโก มคฺโค ภาเวตพฺโพติ.\csromanspacer{}\csromanspacer{}ฉฏฺฐํ.}

\csromanheader{2}{7. มลสุตฺต}
\proseitem{167}{ตีณิมานิ ภิกฺขเว มลานิ.\csromanspacer{} กตมานิ ตีณิ, ราโค มลํ โทโส มลํ โมโห มลํ, อิมานิ โข ภิกฺขเว ตีณิ มลานิ.\csromanspacer{} อิเมสํ โข ภิกฺขเว ติณฺณนฺนํ มลานํ อภิญฺญาย ปริญฺญาย ปริกฺขยาย ปหานาย ฯเปฯ อยํ อริโย อฏฺฐงฺคิโก มคฺโค ภาเวตพฺโพติ.\csromanspacer{}\csromanspacer{}สตฺตมํ.}

\csromanheader{2}{8. นีฆสุตฺต}
\proseitem{168}{ตโยเม ภิกฺขเว นีฆา.\csromanspacer{} กตเม ตโย, ราโค นีโฆ โทโส นีโฆ โมโห นีโฆ, อิเม โข ภิกฺขเว ตโย นีฆา.\csromanspacer{} อิเมสํ โข ภิกฺขเว ติณฺณนฺนํ นีฆานํ อภิญฺญาย ปริญฺญาย ปริกฺขยาย ปหานาย ฯเปฯ อยํ อริโย อฏฺฐงฺคิโก มคฺโค ภาเวตพฺโพติ.\csromanspacer{}\csromanspacer{}อฏฺฐมํ.}

\csromanheader{2}{9. เวทนาสุตฺต}
\proseitem{169}{ติสฺโส อิมา ภิกฺขเว เวทนา.\csromanspacer{} กตมา ติสฺโส, สุขา เวทนา ทุกฺขา เวทนา อทุกฺขมสุขา เวทนา, อิมา โข ภิกฺขเว ติสฺโส เวทนา.\csromanspacer{} อิมาสํ โข ภิกฺขเว ติสฺสนฺนํ เวทนานํ อภิญฺญาย ปริญฺญาย ปริกฺขยาย ปหานาย ฯเปฯ อยํ อริโย อฏฺฐงฺคิโก มคฺโค ภาเวตพฺโพติ.\csromanspacer{}\csromanspacer{}นวมํ.}

\csromanheader{2}{10. ตณฺหาสุตฺต}
\proseitem{170}{ติสฺโส อิมา ภิกฺขเว ตณฺหา.\csromanspacer{} กตมา ติสฺโส, กามตณฺหา ภวตณฺหา วิภวตณฺหา, อิมา โข ภิกฺขเว ติสฺโส ตณฺหา.\csromanspacer{} อิมาสํ โข ภิกฺขเว ติสฺสนฺนํ ตณฺหานํ อภิญฺญาย ปริญฺญาย ปริกฺขยาย ปหานาย ฯเปฯ อยํ อริโย อฏฺฐงฺคิโก มคฺโค ภาเวตพฺโพ.\csromanspacer{} กตโม อริโย อฏฺฐงฺคิโก มคฺโค, อิธ ภิกฺขเว ภิกฺขุ สมฺมาทิฏฺฐิํ ภาเวติ วิเวกนิสฺสิตํ วิราคนิสฺสิตํ นิโรธนิสฺสิตํ โวสฺสคฺคปริณามิํ ฯเปฯ สมฺมาสมาธิํ ภาเวติ วิเวกนิสฺสิตํ วิราคนิสฺสิตํ นิโรธนิสฺสิตํ โวสฺสคฺคปริณามิํ.\csromanspacer{} อิมาสํ โข ภิกฺขเว ติสฺสนฺนํ ตณฺหานํ อภิญฺญาย ปริญฺญาย ปริกฺขยาย ปหานาย ฯเปฯ อยํ อริโย อฏฺฐงฺคิโก มคฺโค ภาเวตพฺโพติ.\csromanspacer{}\csromanspacer{}ทสมํ.}

\csromanheader{2}{11. ตสินาสุตฺต}
\proseitem{171}{\csromanfolio{54}ติสฺโส อิมา ภิกฺขเว ตสินา.\csromanspacer{} กตมา ติสฺโส, กามตสินา ภวตสินา วิภวตสินา.\csromanspacer{} อิมาสํ โข ภิกฺขเว ติสฺสนฺนํ ตสินานํ อภิญฺญาย ปริญฺญาย ปริกฺขยาย ปหานาย ฯเปฯ ราควินยปริโยสานํ โทสวินยปริโยสานํ โมหวินยปริโยสานํ ฯเปฯ อมโตคธํ อมตปรายนํ อมตปริโยสานํ ฯเปฯ นิพฺพานนินฺนํ นิพฺพานโปณํ นิพฺพานปพฺภารํ.\csromanspacer{} อิมาสํ โข ภิกฺขเว ติสฺสนฺนํ ตสินานํ อภิญฺญาย ปริญฺญาย ปริกฺขายาย ปหานาย ฯเปฯ อยํ อริโย อฏฺฐงฺคิโก มคฺโค ภาเวตพฺโพติ.\csromanspacer{}\csromanspacer{}เอกาทสมํ.\csromanspacer{} เอสนาวคฺโค สตฺตโม.}

\titlehead{\textbf{ตสฺสุทฺทานํ}}
\csromangathameasure{เอสนา วิธา อาสโว, ภโว จ ทุกฺขตา ขิลา. \\ มลํ นีโฆ จ เวทนา, เทฺว ตณฺหา ตสินาย จาติ.}
\csromangathagroup{\csromangathastanza{เอสนา วิธา อาสโว, ภโว จ ทุกฺขตา ขิลา. \\ มลํ นีโฆ จ เวทนา, เทฺว ตณฺหา ตสินาย จาติ.}}

\csromanmatikaanchor{mtk.4}
\csromanmatikaanchor{mtk.31}
\csromantocmarkref{section}{8. อาวรณนีวรณสุตฺต}{mtk.4}
\bookmark[dest={mtk.4},level=1]{8. อาวรณนีวรณสุตฺต}
\csromanheader{4}{8. โอฆวคฺค}
\csromanheader{2}{1. โอฆสุตฺต}
\proseitem{172}{สาวตฺถินิทานํ.\csromanspacer{} จตฺตาโรเม ภิกฺขเว โอฆา.\csromanspacer{} กตเม จตฺตาโร, กาโมโฆ ภโวโฆ ทิฏฺโฐโฆ อวิชฺโชโฆ, อิเม โข ภิกฺขเว จตฺตาโร โอฆา.\csromanspacer{} อิเมสํ โข ภิกฺขเว จตุนฺนํ โอฆานํ อภิญฺญาย ปริญฺญาย ปริกฺขยาย ปหานาย ฯเปฯ อยํ อริโย อฏฺฐงฺคิโก มคฺโค ภาเวตพฺโพติ. (ยถา เอสนา, เอวํ วิตฺถาเรตพฺพํ.).\csromanspacer{} ปฐมํ.}

\csromanheader{2}{2. โยคสุตฺต}
\proseitem{173}{จตฺตาโรเม ภิกฺขเว โยคา.\csromanspacer{} กตเม จตฺตาโร, กามโยโค ภวโยโค ทิฏฺฐิโยโค อวิชฺชาโยโค, อิเม โข ภิกฺขเว จตฺตาโร โยคา.\csromanspacer{} อิเมสํ โข ภิกฺขเว จตุนฺนํ โยคานํ อภิญฺญาย ปริญฺญาย ปริกฺขยาย ปหานาย ฯเปฯ อยํ อริโย อฏฺฐงฺคิโก มคฺโค ภาเวตพฺโพติ.\csromanspacer{}\csromanspacer{}ทุติยํ.}

\csromanheader{2}{3. อุปาทานสุตฺต}
\proseitem{174}{\csromanfolio{55}จตฺตาริมานิ ภิกฺขเว อุปาทานานิ.\csromanspacer{} กตมานิ จตฺตาริ, กามุปาทานํ ทิฏฺฐุปาทานํ สีลพฺพตุปาทานํ อตฺตวาทุปาทานํ, อิมานิ โข ภิกฺขเว จตฺตาริ อุปาทานานิ.\csromanspacer{} อิเมสํ โข ภิกฺขเว จตุนฺนํ อุปาทานานํ อภิญฺญาย ปริญฺญาย ปริกฺขยาย ปหานาย ฯเปฯ อยํ อริโย อฏฺฐงฺคิโก มคฺโค ภาเวตพฺโพติ.\csromanspacer{}\csromanspacer{}ตติยํ.}

\csromanheader{2}{4. คนฺถสุตฺต}
\proseitem{175}{จตฺตาโรเม ภิกฺขเว คนฺถา.\csromanspacer{} กตเม จตฺตาโร, อภิชฺฌา กายคนฺโถ พฺยาปาโท กายคนฺโถ สีลพฺพตปรามาโส กายคนฺโถ อิทํสจฺจาภินิเวโส กายคนฺโถ, อิเม โข ภิกฺขเว จตฺตาโร คนฺถา.\csromanspacer{} อิเมสํ โข ภิกฺขเว จตุนฺนํ คนฺถานํ อภิญฺญาย ปริญฺญาย ปริกฺขยาย ปหานาย ฯเปฯ อยํ อริโย อฏฺฐงฺคิโก มคฺโค ภาเวตพฺโพติ.\csromanspacer{}\csromanspacer{}จตุตฺถํ.}

\csromanheader{2}{5. อนุสยสุตฺต}
\proseitem{176}{สตฺติเม ภิกฺขเว อนุสยา.\csromanspacer{} กตเม สตฺต, กามราคานุสโย ปฏิฆานุสโย ทิฏฺฐานุสโย วิจิกิจฺฉานุสโย มานานุสโย ภวราคานุสโย อวิชฺชานุสโย, อิเม โข ภิกฺขเว สตฺตานุสยา.\csromanspacer{} อิเมสํ โข ภิกฺขเว สตฺตนฺนํ อนุสยานํ อภิญฺญาย ปริญฺญาย ปริกฺขยาย ปหานาย ฯเปฯ อยํ อริโย อฏฺฐงฺคิโก มคฺโค ภาเวตพฺโพติ.\csromanspacer{}\csromanspacer{}ปญฺจมํ.}

\csromanheader{2}{6. กามคุณสุตฺต}
\proseitem{177}{ปญฺจิเม ภิกฺขเว กามคุณา.\csromanspacer{} กตเม ปญฺจ, จกฺขุวิญฺเญยฺยา รูปา อิฏฺฐา กนฺตา มนาปา ปิยรูปา กามูปสํหิตา รชนียา, โสตวิญฺเญยฺยา สทฺทา ฯเปฯ ฆานวิญฺเญยฺยา คนฺธา ฯเปฯ ชิวฺหาวิญฺเญยฺยา รสา ฯเปฯ กายวิญฺเญยฺยา โผฏฺฐพฺพา อิฏฺฐา กนฺตา มนาปา ปิยรูปา กามูปสํหิตา รชนียา, อิเม โข ภิกฺขเว ปญฺจ กามคุณา.\csromanspacer{} อิเมสํ โข ภิกฺขเว ปญฺจนฺนํ กามคุณานํ อภิญฺญาย ปริญฺญาย ปริกฺขยาย ปหานาย ฯเปฯ อยํ อริโย อฏฺฐงฺคิโก มคฺโค ภาเวตพฺโพติ.\csromanspacer{}\csromanspacer{}ฉฏฺฐํ.}

\csromanheader{2}{7. นีวรณสุตฺต}
\proseitem{178}{\csromanfolio{56}ปญฺจิมานิ ภิกฺขเว นีวรณานิ.\csromanspacer{} กตมานิ ปญฺจ, กามจฺฉนฺทนีวรณํ พฺยาปาทนีวรณํ ถินมิทฺธนีวรณํ อุทฺธจฺจกุกฺกุจฺจนีวรณํ วิจิกิจฺฉานีวรณํ, อิมานิ โข ภิกฺขเว ปญฺจ นีวรณานิ.\csromanspacer{} อิเมสํ โข ภิกฺขเว ปญฺจนฺนํ นีวรณานํ อภิญฺญาย ปริญฺญาย ปริกฺขยาย ปหานาย ฯเปฯ อยํ อริโย อฏฺฐงฺคิโก มคฺโค ภาเวตพฺโพติ.\csromanspacer{}\csromanspacer{}สตฺตมํ.}

\csromanheader{2}{8. อุปาทานกฺขนฺธสุตฺต}
\proseitem{179}{ปญฺจิเม ภิกฺขเว อุปาทานกฺขนฺธา.\csromanspacer{} กตเม ปญฺจ, เสยฺยถิทํ, รูปุปาทานกฺขนฺโธ เวทนุปาทานกฺขนฺโธ สญฺญุปาทานกฺขนฺโธ สงฺขารุปาทานกฺขนฺโธ วิญฺญาณุปาทานกฺขนฺโธ, อิเม โข ภิกฺขเว ปญฺจุปาทานกฺขนฺธา.\csromanspacer{} อิเมสํ โข ภิกฺขเว ปญฺจนฺนํ อุปาทานกฺขนฺธานํ อภิญฺญาย ปริญฺญาย ปริกฺขยาย ปหานาย ฯเปฯ อยํ อริโย อฏฺฐงฺคิโก มคฺโค ภาเวตพฺโพติ.\csromanspacer{}\csromanspacer{}อฏฺฐมํ.}

\csromanheader{2}{9. โอรมฺภาคิยสุตฺต}
\proseitem{180}{ปญฺจิมานิ ภิกฺขเว โอรมฺภาคิยานิ สํโยชนานิ.\csromanspacer{} กตมานิ ปญฺจ, สกฺกายทิฏฺฐิ วิจิกิจฺฉา สีลพฺพตปรามาโส กามจฺฉนฺโท พฺยาปาโท, อิมานิ โข ภิกฺขเว ปญฺโจรมฺภาคิยานิ สํโยชนานิ.\csromanspacer{} อิมาสํ โข ภิกฺขเว ปญฺจนฺนํ โอรมฺภาคิยานํ สํโยชนานํ อภิญฺญาย ปริญฺญาย ปริกฺขยาย ปหานาย ฯเปฯ อยํ อริโย อฏฺฐงฺคิโก มคฺโค ภาเวตพฺโพติ.\csromanspacer{}\csromanspacer{}นวมํ.}

\csromanheader{2}{10. อุทฺธมฺภาคิยสุตฺต}
\proseitem{181}{ปญฺจิมานิ ภิกฺขเว อุทฺธมฺภาคิยานิ สํโยชนานิ.\csromanspacer{} กตมานิ ปญฺจ, รูปราโค อรูปราโค มาโน อุทฺธจฺจํ อวิชฺชา, อิมานิ โข ภิกฺขเว ปญฺจุทฺธมฺภาคิยานิ สํโยชนานิ.\csromanspacer{} อิเมสํ โข ภิกฺขเว ปญฺจนฺนํ อุทฺธมฺภาคิยานํ สํโยชนานํ อภิญฺญาย ปริญฺญาย ปริกฺขยาย ปหานาย อริโย อฏฺฐงฺคิโก มคฺโค ภาเวตพฺโพ.\csromanspacer{} กตโม อริโย อฏฺฐงฺคิโก มคฺโค, อิธ ภิกฺขเว ภิกฺขุ สมฺมาทิฏฺฐิํ ภาเวติ วิเวกนิสฺสิตํ ฯเปฯ สมฺมาสมาธิํ ภาเวติ วิเวกนิสฺสิตํ วิราคนิสฺสิตํ นิโรธนิสฺสิตํ โวสฺสคฺคปริณามิํ.\csromanspacer{} อิเมสํ โข ภิกฺขเว ปญฺจนฺนํ อุทฺธมฺภาคิยานํ สํโยชนานํ อภิญฺญาย \csromanfolio{57}ปริญฺญาย ปริกฺขยาย ปหานาย อยํ อริโย อฏฺฐงฺคิโก มคฺโค ภาเวตพฺโพติ.}

\prose{ปญฺจิมานิ ภิกฺขเว อุทฺธมฺภาคิยานิ สํโยชนานิ.\csromanspacer{} กตมานิ ปญฺจ, รูปราโค อรูปราโค มาโน อุทฺธจฺจํ อวิชฺชา, อิมานิ โข ภิกฺขเว ปญฺจุทฺธมฺภาคิยานิ สํโยชนานิ.\csromanspacer{} อิเมสํ โข ภิกฺขเว ปญฺจนฺนํ อุทฺธมฺภาคิยานํ สํโยชนานํ อภิญฺญาย ปริญฺญาย ปริกฺขยาย ปหานาย อริโย อฏฺฐงฺคิโก มคฺโค ภาเวตพฺโพ.\csromanspacer{} กตโม อริโย อฏฺฐงฺคิโก มคฺโค, อิธ ภิกฺขเว ภิกฺขุ สมฺมาทิฏฺฐิํ ภาเวติ ฯเปฯ สมฺมาสมาธิํ ภาเวติ ราควินยปริโยสานํ โทสวินยปริโยสานํ โมหวินยปริโยสานํ.\csromanspacer{} อมโตคธํ อมตปรายนํ อมตปริโยสานํ.\csromanspacer{} นิพฺพานนินฺนํ นิพฺพานโปณํ นิพฺพานปพฺภารํ.\csromanspacer{} อิเมสํ โข ภิกฺขเว ปญฺจนฺนํ อุทฺธมฺภาคิยานํ สํโยชนานํ อภิญฺญาย ปริญฺญาย ปริกฺขยาย ปหานาย อยํ อริโย อฏฺฐงฺคิโก มคฺโค ภาเวตพฺโพติ.\csromanspacer{}\csromanspacer{}ทสมํ.\csromanspacer{} โอฆวคฺโค อฏฺฐโม.}

\titlehead{\textbf{ตสฺสุทฺทานํ}}
\csromangathameasure{โอโฆ โยโค อุปาทานํ, คนฺถํ อนุสเยน จ. \\ กามคุณา นีวรณํ, ขนฺธา โอรุทฺธมฺภาคิยาติ.}
\csromangathagroup{\csromangathastanza{โอโฆ โยโค อุปาทานํ, คนฺถํ อนุสเยน จ. \\ กามคุณา นีวรณํ, ขนฺธา โอรุทฺธมฺภาคิยาติ.}}

\csromanheader{2}{\textbf{วคฺคุทฺทานํ}}
\csromangathameasure{อวิชฺชาวคฺโค ปฐโม, ทุติยํ วิหารํ วุจฺจติ. \\ มิจฺฉตฺตํ ตติโย วคฺโค, จตุตฺถํ ปฏิปนฺเนเนว. \\ ติตฺถิยํ ปญฺจโม วคฺโค, จตสฺโส สูริเยน จ. \\ พหุกเต สตฺตโม วคฺโค, อุปฺปาโท อฏฺฐเมน จ. \\ ทิวสวคฺโค นวโม, ทสโม อปฺปมาเทน จ. \\ เอกาทสพลวคฺโค, ทฺวาทส เอสนา ปาฬิยํ. \\ โอฆวคฺโค ภวติ เตรสาติ. มคฺคสํยุตฺตํ ปฐมํ.}
\csromangathagroup{\csromangathastanza{อวิชฺชาวคฺโค ปฐโม, ทุติยํ วิหารํ วุจฺจติ. \\ มิจฺฉตฺตํ ตติโย วคฺโค, จตุตฺถํ ปฏิปนฺเนเนว.}\csromangathastanza{ติตฺถิยํ ปญฺจโม วคฺโค, จตสฺโส สูริเยน จ. \\ พหุกเต สตฺตโม วคฺโค, อุปฺปาโท อฏฺฐเมน จ.}\csromangathastanza{ทิวสวคฺโค นวโม, ทสโม อปฺปมาเทน จ. \\ เอกาทสพลวคฺโค, ทฺวาทส เอสนา ปาฬิยํ. \\ โอฆวคฺโค ภวติ เตรสาติ. มคฺคสํยุตฺตํ ปฐมํ.}}

\csromanheader{2}{2. โพชฺฌงฺคสํยุตฺต}
\chapterhead{1. ปพฺพตวคฺค}
\csromanheader{2}{1. หิมวนฺตสุตฺต}
\proseitem{182}{\csromanfolio{58}สาวตฺถินิทานํ.\csromanspacer{} เสยฺยถาปิ ภิกฺขเว หิมวนฺตํ ปพฺพตราชานํ นิสฺสาย นาคา กายํ วฑฺเฒนฺติ พลํ คาเหนฺติ, เต ตตฺถ กายํ วฑฺเฒตฺวา พลํ คาเหตฺวา กุโสพฺเภ โอตรนฺติ, กุโสพฺเภ โอตริตฺวา มหาโสพฺเภ โอตรนฺติ, มหาโสพฺเภ โอตริตฺวา กุนฺนทิโย โอตรนฺติ, กุนฺนทิโย โอตริตฺวา มหานทิโย โอตรนฺติ, มหานทิโย โอตริตฺวา มหาสมุทฺทสาครํ โอตรนฺติ, เต ตตฺถ มหนฺตตฺตํ เวปุลฺลตฺตํ อาปชฺชนฺติ กาเยน.\csromanspacer{} เอวเมว โข ภิกฺขเว ภิกฺขุ สีลํ นิสฺสาย สีเล ปติฏฺฐาย สตฺต โพชฺฌงฺเค ภาเวนฺโต สตฺต โพชฺฌงฺเค พหุลีกโรนฺโต มหนฺตตฺตํ เวปุลฺลตฺตํ ปาปุณาติ ธมฺเมสุ.\csromanspacer{} กถญฺจ ภิกฺขเว ภิกฺขุ สีลํ นิสฺสาย สีเล ปติฏฺฐาย สตฺต โพชฺฌงฺเค ภาเวนฺโต สตฺต โพชฺฌงฺเค พหุลีกโรนฺโต มหนฺตตฺตํ เวปุลฺลตฺตํ ปาปุณาติ ธมฺเมสูติ.\csromanspacer{} อิธ ภิกฺขเว ภิกฺขุ สติสมฺโพชฺฌงฺคํ ภาเวติ วิเวกนิสฺสิตํ วิราคนิสฺสิตํ นิโรธนิสฺสิตํ โวสฺสคฺคปริณามิํ.\csromanspacer{} ธมฺมวิจยสมฺโพชฺฌงฺคํ ภาเวติ ฯเปฯ.\csromanspacer{} วีริยสมฺโพชฺฌงฺคํ ภาเวติ ฯเปฯ.\csromanspacer{} ปีติสมฺโพชฺฌงฺคํ ภาเวติ ฯเปฯ.\csromanspacer{} ปสฺสทฺธิสมฺโพชฺฌงฺคํ ภาเวติ ฯเปฯ.\csromanspacer{} สมาธิสมฺโพชฺฌงฺคํ ภาเวติ ฯเปฯ.\csromanspacer{} อุเปกฺขาสมฺโพชฺฌงฺคํ ภาเวติ วิเวกนิสฺสิตํ วิราคนิสฺสิตํ นิโรธนิสฺสิตํ โวสฺสคฺคปริณามิํ.\csromanspacer{} เอวํ โข ภิกฺขเว ภิกฺขุ สีลํ นิสฺสาย สีเล ปติฏฺฐาย สตฺต โพชฺฌงฺเค ภาเวนฺโต สตฺต โพชฺฌางฺเค พหุลีกโรนฺโต มหนฺตตฺตํ เวปุลฺลตฺตํ ปาปุณาติ ธมฺเมสูติ.\csromanspacer{}\csromanspacer{}ปฐมํ.}

\csromanheader{2}{2. กายสุตฺต}
\proseitem{183}{สาวตฺถินิทานํ.\csromanspacer{} เสยฺยถาปิ ภิกฺขเว อยํ กาโย อาหารฏฺฐิติโก อาหารํ ปฏิจฺจ ติฏฺฐติ, อนาหาโร โน ติฏฺฐติ.\csromanspacer{} เอวเมว โข ภิกฺขเว ปญฺจ นีวรณา อาหารฏฺฐิติกา อาหารํ ปฏิจฺจ ติฏฺฐนฺติ, อนาหารา โน ติฏฺฐนฺติ.}

\csromanheader{2}{2. โพชฺฌงฺคสํยุตฺต}
\prose{\csromanfolio{59}โก จ ภิกฺขเว อาหาโร อนุปฺปนฺนสฺส วา กามจฺฉนฺทสฺส อุปฺปาทาย อุปฺปนฺนสฺส วา กามจฺฉนฺทสฺส ภิยฺโยภาวาย เวปุลฺลาย, อตฺถิ ภิกฺขเว สุภนิมิตฺตํ, ตตฺถ อโยนิโสมนสิการพหุลีกาโร, อยมาหาโร อนุปฺปนฺนสฺส วา กามจฺฉนฺทสฺส อุปฺปาทาย อุปฺปนฺนสฺส วา กามจฺฉนฺทสฺส ภิยฺโยภาวาย เวปุลฺลาย.}

\prose{โก จ ภิกฺขเว อาหาโร อนุปฺปนฺนสฺส วา พฺยาปาทสฺส อุปฺปาทาย อุปฺปนฺนสฺส วา พฺยาปาทสฺส ภิยฺโยภาวาย เวปุลฺลาย, อตฺถิ ภิกฺขเว ปฏิฆนิมิตฺตํ, ตตฺถ อโยนิโสมนสิการพหุลีกาโร, อยมาหาโร อนุปฺปนฺนสฺส วา พฺยาปาทสฺส อุปฺปาทาย อุปฺปนฺนสฺส วา พฺยาปาทสฺส ภิยฺโยภาวาย เวปุลฺลาย.}

\prose{โก จ ภิกฺขเว อาหาโร อนุปฺปนฺนสฺส วา ถินมิทฺธสฺส อุปฺปาทาย อุปฺปนฺนสฺส วา ถินมิทฺธสฺส ภิยฺโยภาวาย เวปุลฺลาย, อตฺถิ ภิกฺขเว อรติ ตนฺทิ วิชมฺภิตา ภตฺตสมฺมโท เจตโส จ ลีนตฺตํ, ตตฺถ อโยนิโสมนสิการพหุลีกาโร, อยมาหาโร อนุปฺปนฺนสฺส วา ถินมิทฺธสฺส อุปฺปาทาย อุปฺปนฺนสฺส วา ถินมิทฺธสฺส ภิยฺโยภาวาย เวปุลฺลาย.}

\prose{โก จ ภิกฺขเว อาหาโร อนุปฺปนฺนสฺส วา อุทฺธจฺจกุกฺกุจฺจสฺส อุปฺปาทาย อุปฺปนฺนสฺส วา อุทฺธจฺจกุกฺกุจฺจสฺส ภิยฺโยภาวาย เวปุลฺลาย, อตฺถิ ภิกฺขเว เจตโส อวูปสโม, ตตฺถ อโยนิโสมนสิการพหุลีกาโร, อยมาหาโร อนุปฺปนฺนสฺส วา อุทฺธจฺจกุกฺกุจฺจสฺส อุปฺปาทาย อุปฺปนฺนสฺส วา อุทฺธจฺจกุกฺกุจฺจสฺส ภิยฺโยภาวาย เวปุลฺลาย.}

\prose{โก จ ภิกฺขเว อาหาโร อนุปฺปนฺนาย วา วิจิกิจฺฉาย อุปฺปาทาย อุปฺปนฺนาย วา วิจิกิจฺฉาย ภิยฺโยภาวาย เวปุลฺลาย, อตฺถิ ภิกฺขเว วิจิกิจฺฉาฏฺฐานียา ธมฺมา, ตตฺถ อโยนิโสมนสิการพหุลีกาโร, อยมาหาโร อนุปฺปนฺนาย วา วิจิกิจฺฉาย อุปฺปาทาย อุปฺปนฺนาย วา วิจิกิจฺฉาย ภิยฺโยภาวาย เวปุลฺลาย.}

\prose{เสยฺยถาปิ ภิกฺขเว อยํ กาโย อาหารฏฺฐิติโก อาหารํ ปฏิจฺจ ติฏฺฐติ, อนาหาโร โน ติฏฺฐติ.\csromanspacer{} เอวเมว โข ภิกฺขเว อิเม ปญฺจ นีวรณา อาหารฏฺฐิติกา อาหารํ ปฏิจฺจ ติฏฺฐนฺติ, อนาหารา โน ติฏฺฐนฺติ.}

\prose{\csromanfolio{60}เสยฺยถาปิ ภิกฺขเว อยํ กาโย อาหารฏฺฐิติโก อาหารํ ปฏิจฺจ ติฏฺฐติ, อนาหาโร โน ติฏฺฐติ.\csromanspacer{} เอวเมว โข ภิกฺขเว สตฺต โพชฺฌงฺคา อาหารฏฺฐิติกา อาหารํ ปฏิจฺจ ติฏฺฐนฺติ, อนาหารา โน ติฏฺฐนฺติ.}

\prose{โก จ ภิกฺขเว อาหาโร อนุปฺปนฺนสฺส วา สติสมฺโพชฺฌงฺคสฺส อุปฺปาทาย อุปฺปนฺนสฺส วา สติสมฺโพชฺฌงฺคสฺส ภาวนาย ปาริปูริยา, อตฺถิ ภิกฺขเว สติสมฺโพชฺฌงฺคฏฺฐานียา ธมฺมา, ตตฺถ โยนิโสมนสิการพหุลีกาโร, อยมาหาโร อนุปฺปนฺนสฺส วา สติสมฺโพชฺฌงฺคสฺส อุปฺปาทาย อุปฺปนฺนสฺส วา สติสมฺโพชฺฌงฺคสฺส ภาวนาย ปาริปูริยา.}

\prose{โก จ ภิกฺขเว อาหาโร อนุปฺปนฺนสฺส วา ธมฺมวิจยสมฺโพชฺฌงฺคสฺส อุปฺปาทาย อุปฺปนฺนสฺส วา ธมฺมวิจยสมฺโพชฺฌงฺคสฺส ภาวนาย ปาริปูริยา, อตฺถิ ภิกฺขเว กุสลากุสลา ธมฺมา สาวชฺชานวชฺชา ธมฺมา หีนปณีตา ธมฺมา กณฺหสุกฺกสปฺปฏิภาคา ธมฺมา, ตตฺถ โยนิโสมนสิการพหุลีกาโร, อยมาหาโร อนุปฺปนฺนสฺส วา ธมฺมวิจยสมฺโพชฺฌงฺคสฺส อุปฺปาทาย อุปฺปนฺนสฺส วา ธมฺมวิจยสมฺโพชฺฌงฺคสฺส ภาวนาย ปาริปูริยา.}

\prose{โก จ ภิกฺขเว อาหาโร อนุปฺปนฺนสฺส วา วีริยสมฺโพชฺฌงฺคสฺส อุปฺปาทาย อุปฺปนฺนสฺส วา วีริยสมฺโพชฺฌงฺคสฺส ภาวนาย ปาริปูริยา, อตฺถิ ภิกฺขเว อารมฺภธาตุ\footnote{อารพฺภธาตุ (สฺยา, ก)} นิกฺกมธาตุ ปรกฺกมธาตุ, ตตฺถ โยนิโสมนสิการพหุลีกาโร, อยมาหาโร อนุปฺปนฺนสฺส วา วีริยสมฺโพชฺฌงฺคสฺส อุปฺปาทาย อุปฺปนฺนสฺส วา วีริยสมฺโพชฺฌงฺคสฺส ภาวนาย ปาริปูริยา.}

\prose{โก จ ภิกฺขเว อาหาโร อนุปฺปนฺนสฺส วา ปีติสมฺโพชฺฌงฺคสฺส อุปฺปาทาย อุปฺปนฺนสฺส วา ปีติสมฺโพชฺฌงฺคสฺส ภาวนาย ปาริปูริยา, อตฺถิ ภิกฺขเว ปีติสมฺโพชฺฌงฺคฏฺฐานียา ธมฺมา, ตตฺถ โยนิโสมนสิการพหุลีกาโร, อยมาหาโร อนุปฺปนฺนสฺส วา ปีติสมฺโพชฺฌงฺคสฺส อุปฺปาทาย อุปฺปนฺนสฺส วา ปีติสมฺโพชฺฌงฺคาสฺส ภาวนาย ปาริปูริยา.}

\prose{โก จ ภิกฺขเว อาหาโร อนุปฺปนฺนสฺส วา ปสฺสทฺธิสมฺโพชฺฌงฺคสฺส อุปฺปาทาย อุปฺปนฺนสฺส วา ปสฺสทฺธิสมฺโพชฺฌงฺคสฺส ภาวนาย ปาริปูริยา, อตฺถิ}

\csromanheader{2}{2. โพชฺฌงฺคสํยุตฺต}
\prose{\csromanfolio{61}ภิกฺขเว กายปสฺสทฺธิ จิตฺตปสฺสทฺธิ, ตตฺถ โยนิโสมนสิการพหุลีกาโร, อยมาหาโร อนุปฺปนฺนสฺส วา ปสฺสทฺธิสมฺโพชฺฌงฺคสฺส อุปฺปาทาย อุปฺปนฺนสฺส วา ปสฺสทฺธิสมฺโพชฺฌงฺคสฺส ภาวนาย ปาริปูริยา.}

\prose{โก จ ภิกฺขเว อาหาโร อนุปฺปนฺนสฺส วา สมาธิสมฺโพชฺฌงฺคสฺส อุปฺปาทาย อุปฺปนฺนสฺส วา สมาธิสมฺโพชฺฌงฺคสฺส ภาวนาย ปาริปูริยา, อตฺถิ ภิกฺขเว สมถนิมิตฺตํ\footnote{สมาธินิมิตฺตํ (สฺยา)} อพฺยคฺคนิมิตฺตํ, ตตฺถ โยนิโสมนสิการพหุลีกาโร, อยมาหาโร อนุปฺปนฺนสฺส วา สมาธิสมฺโพชฺฌงฺคสฺส อุปฺปาทาย อุปฺปนฺนสฺส วา สมาธิสมฺโพชฺฌงฺคสฺส ภาวนาย ปาริปูริยา.}

\prose{โก จ ภิกฺขเว อาหาโร อนุปฺปนฺนสฺส วา อุเปกฺขาสมฺโพชฺฌงฺคสฺส อุปฺปาทาย อุปฺปนฺนสฺส วา อุเปกฺขาสมฺโพชฺฌงฺคสฺส ภาวนาย ปาริปูริยา, อตฺถิ ภิกฺขเว อุเปกฺขาสมฺโพชฺฌงฺคฏฺฐานียา ธมฺมา, ตตฺถ โยนิโสมนสิการพหุลีกาโร, อยมาหาโร อนุปฺปนฺนสฺส วา อุเปกฺขาสมฺโพชฺชงฺคสฺส อุปฺปาทาย อุปฺปนฺนสฺส วา อุเปกฺขาสมฺโพชฺฌงฺคสฺส ภาวนาย ปาริปูริยา.}

\prose{เสยฺยถาปิ ภิกฺขเว อยํ กาโย อาหารฏฺฐิติโก อาหารํ ปฏิจฺจ ติฏฺฐติ, อนาหาโร โน ติฏฺฐติ.\csromanspacer{} เอวเมว โข ภิกฺขเว อิเม สตฺต โพชฺฌงฺคา อาหารฏฺฐิติกา อาหารํ ปฏิจฺจ ติฏฺฐนฺติ, อนาหารา โน ติฏฺฐนฺตีติ.\csromanspacer{}\csromanspacer{}ทุติยํ.}

\csromanheader{2}{3. สีลสุตฺต}
\proseitem{184}{เย เต ภิกฺขเว ภิกฺขู สีลสมฺปนฺนา สมาธิสมฺปนฺนา ญาณสมฺปนฺนา วิมุตฺติสมฺปนฺนา วิมุตฺติญาณทสฺสนสมฺปนฺนา, ทสฺสนมฺปาหํ ภิกฺขเว เตสํ ภิกฺขูนํ พหุการํ\footnote{พหูปการํ (สฺยา)} วทามิ, สวนมฺปาหํ ภิกฺขเว เตสํ ภิกฺขูนํ พหุการํ วทามิ, อุปสงฺกมนมฺปาหํ ภิกฺขเว เตสํ ภิกฺขูนํ พหุการํ วทามิ, ปยิรุปาสนมฺปาหํ ภิกฺขเว เตสํ ภิกฺขูนํ พหุการํ วทามิ, อนุสฺสติมฺปาหํ ภิกฺขเว เตสํ ภิกฺขูนํ พหุการํ วทามิ, อนุปพฺพชฺชมฺปาหํ ภิกฺขเว เตสํ ภิกฺขูนํ พหุการํ วทามิ.\csromanspacer{} ตํ กิสฺส เหตุ, ตถารูปานํ ภิกฺขเว ภิกฺขูนํ ธมฺมํ สุตฺวา ทฺวเยน วูปกาเสน วูปกฏฺโฐ\footnote{ทฺวเยน วูปกฏฺโฐ (สี, สฺยา)} วิหรติ กายวูปกาเสน จ จิตฺตวูปกาเสน จ, โส ตถา วูปกฏฺโฐ วิหรนฺโต ตํ ธมฺมํ อนุสฺสรติ อนุวิตกฺเกติ.}

\prose{\csromanfolio{62}ยสฺมิํ สมเย ภิกฺขเว ภิกฺขุ ตถา วูปกฏฺโฐ วิหรนฺโต ตํ ธมฺมํ อนุสฺสรติ อนุวิตกฺเกติ, สติสมฺโพชฺฌงฺโค ตสฺมิํ สมเย ภิกฺขุโน อารทฺโธ โหติ, สติสมฺโพชฺฌงฺคํ ตสฺมิํ สมเย ภิกฺขุ ภาเวติ, สติสมฺโพชฺฌงฺโค ตสฺมิํ สมเย ภิกฺขุโน ภาวนาปาริปูริํ คจฺฉติ, โส ตถา สโต วิหรนฺโต ตํ ธมฺมํ ปญฺญาย ปวิจินติ ปวิจรติ ปริวีมํสมาปชฺชติ.}

\prose{ยสฺมิํ สมเย ภิกฺขเว ภิกฺขุ ตถา สโต วิหรนฺโต ตํ ธมฺมํ ปญฺญาย ปวิจินติ ปวิจรติ ปริวีมํสมาปชฺชติ, ธมฺมวิจยสมฺโพชฺฌงฺโค ตสฺมิํ สมเย ภิกฺขุโน อารทฺโธ โหติ, ธมฺมวิจยสมฺโพชฺฌงฺคํ ตสฺมิํ สมเย ภิกฺขุ ภาเวติ, ธมฺมวิจยสมฺโพชฺฌงฺโค ตสฺมิํ สมเย ภิกฺขุโน ภาวนาปาริปูริํ คจฺฉติ.\csromanspacer{} ตสฺส ตํ ธมฺมํ ปญฺญาย ปวิจินโต ปวิจรโต ปริวีมํสมาปชฺชโต อารทฺธํ โหติ วีริยํ อสลฺลีนํ.}

\prose{ยสฺมิํ สมเย ภิกฺขเว ภิกฺขุโน ตํ ธมฺมํ ปญฺญาย ปวิจินโต ปวิจรโต ปริวีมํสมาปชฺชโต อารทฺธํ โหติ วีริยํ อสลฺลีนํ, วีริยสมฺโพชฺฌงฺโค ตสฺมิํ สมเย ภิกฺขุโน อารทฺโธ โหติ.\csromanspacer{} วีริยสมฺโพชฺฌงฺคํ ตสฺมิํ สมเย ภิกฺขุ ภาเวติ, วีริยสมฺโพชฺฌงฺโค ตสฺมิํ สมเย ภิกฺขุโน ภาวนาปาริปูริํ คจฺฉติ, อารทฺธวีริยสฺส อุปฺปชฺชติ ปีติ นิรามิสา.}

\prose{ยสฺมิํ สมเย ภิกฺขเว ภิกฺขุโน อารทฺธวีริยสฺส อุปฺปชฺชติ ปีติ นิรามิสา, ปีติสมฺโพชฺฌงฺโค ตสฺมิํ สมเย ภิกฺขุโน อารทฺโธ โหติ, ปีติสมฺโพชฺฌงฺคํ ตสฺมิํ สมเย ภิกฺขุ ภาเวติ, ปีติสมฺโพชฺฌงฺโค ตสฺมิํ สมเย ภิกฺขุโน ภาวนาปาริปูริํ คจฺฉติ, ปีติมนสฺส กาโยปิ ปสฺสมฺภติ จิตฺตมฺปิ ปสฺสมฺภติ, ยสฺมิํ สมเย ภิกฺขเว ภิกฺขุโน ปีติมนสฺส กาโยปิ ปสฺสมฺภติ จิตฺตมฺปิ ปสฺสมฺภติ.\csromanspacer{} ปสฺสทฺธิสมฺโพชฺฌงฺโค ตสฺมิํ สมเย ภิกฺขุโน อารทฺโธ โหติ, ปสฺสทฺธิสมฺโพชฺฌงฺคํ ตสฺมิํ สมเย ภิกฺขุ ภาเวติ, ปสฺสทฺธิสมฺโพชฺฌงฺโค ตสฺมิํ สมเย ภิกฺขุโน ภาวนาปาริปูริํ คจฺฉติ, ปสฺสทฺธกายสฺส สุขิโน จิตฺตํ สมาธิยติ.}

\prose{ยสฺมิํ สมเย ภิกฺขเว ภิกฺขุโน ปสฺสทฺธกายสฺส สุขิโน จิตฺตํ สมาธิยติ, สมาธิสมฺโพชฺฌงฺโค ตสฺมิํ สมเย ภิกฺขุโน อารทฺโธ}

\csromanheader{2}{2. โพชฺฌงฺคสํยุตฺต}
\prose{\csromanfolio{63}โหติ, สมาธิสมฺโพชฺฌงฺคํ ตสฺมิํ สมเย ภิกฺขุ ภาเวติ, สมาธิสมฺโพชฺฌงฺโค ตสฺมิํ สมเย ภิกฺขุโน ภาวนาปาริปูริํ คจฺฉติ, โส ตถาสมาหิตํ จิตฺตํ สาธุกํ อชฺฌุเปกฺขิตา โหติ.}

\prose{ยสฺมิํ สมเย ภิกฺขเว ภิกฺขุ ตถาสมาหิตํ จิตฺตํ สาธุกํ อชฺฌุเปกฺขิตา โหติ, อุเปกฺขาสมฺโพชฺฌงฺโค ตสฺมิํ สมเย ภิกฺขุโน อารทฺโธ โหติ, อุเปกฺขาสมฺโพชฺฌงฺคํ ตสฺมิํ สมเย ภิกฺขุ ภาเวติ, อุเปกฺขาสมฺโพชฺฌงฺโค ตสฺมิํ สมเย ภิกฺขุโน ภาวนาปาริปูริํ คจฺฉติ.}

\prose{เอวํ ภาวิเตสุ โข ภิกฺขเว สตฺตสุ สมฺโพชฺฌงฺเคสุ เอวํ พหุลีกเตสุ สตฺต ผลา สตฺตานิสํสา ปาฏิกงฺขา.\csromanspacer{} กตเม สตฺต ผลา สตฺตานิสํสา, ทิฏฺเฐว ธมฺเม ปฏิกจฺจ อญฺญํ อาราเธติ.\csromanspacer{} โน เจ ทิฏฺเฐว ธมฺเม ปฏิกจฺจ อญฺญํ อาราเธติ, อถ มรณกาเล อญฺญํ อาราเธติ.\csromanspacer{} โน เจ ทิฏฺเฐว ธมฺเม ปฏิกจฺจ อญฺญํ อาราเธติ, โน เจ มรณกาเล อญฺญํ อาราเธติ, อถ ปญฺจนฺนํ โอรมฺภาคิยานํ สํโยชนานํ ปริกฺขยา อนฺตราปรินิพฺพายี โหติ.\csromanspacer{} โน เจ ทิฏฺเฐว ธมฺเม ปฏิกจฺจ อญฺญํ อาราเธติ, โน เจ มรณกาเล อญฺญํ อาราเธติ, โน เจ ปญฺจนฺนํ โอรมฺภาคิยานํ สํโยชนานํ ปริกฺขยา อนฺตราปรินิพฺพายี โหติ, อถ ปญฺจนฺนํ โอรมฺภาคิยานํ สํโยชนานํ ปริกฺขยา อุปหจฺจปรินิพฺพายี โหติ.\csromanspacer{} โน เจ ทิฏฺเฐว ธมฺเม ปฏิกจฺจ อญฺญํ อาราเธติ, โน เจ มรณกาเล อญฺญํ อาราเธติ, โน เจ ปญฺจนฺนํ โอรมฺภาคิยานํ สํโยชนานํ ปริกฺขยา อนฺตราปรินิพฺพายี โหติ, โน เจ ปญฺจนฺนํ โอรมฺภาคิยานํ สํโยชนานํ ปริกฺขยา อุปหจฺจปรินิพฺพายี โหติ, อถ ปญฺจนฺนํ โอรมฺภาคิยานํ สํโยชนานํ ปริกฺขยา อสงฺขารปรินิพฺพายี โหติ.\csromanspacer{} โน เจ ทิฏฺเฐว ธมฺเม ปฏิกจฺจ อญฺญํ อารเธติ, โน เจ มรณกาเล อญฺญํ อาราเธติ, โน เจ ปญฺจนฺนํ โอรมฺภาคิยานํ สํโยชนานํ ปริกฺขยา อนฺตราปรินิพฺพายี โหติ, โน เจ ปญฺจนฺนํ โอรมฺภาคิยานํ สํโยชนานํ ปริกฺขยา อุปหจฺจปรินิพฺพายี โหติ, โน เจ ปญฺจนฺนํ โอรมฺภาคิยานํ สํโยชนานํ ปริกฺขยา อสงฺขารปรินิพฺพายี โหติ, อถ ปญฺจนฺนํ โอรมฺภาคิยานํ สํโยชนานํ ปริกฺขยา สสงฺขารปรินิพฺพายี โหติ.\csromanspacer{} โน เจ ทิฏฺเฐว ธมฺเม ปฏิกจฺจ อญฺญํ อาราเธติ, โน เจ มรณกาเล อญฺญํ อาราเธติ, โน เจ ปญฺจนฺนํ โอรมฺภาคิยานํ สํโยชนานํ ปริกฺขยา อนฺตราปรินิพฺพายี}

\prose{\csromanfolio{64}โหติ, โน เจ ปญฺจนฺนํ โอรมฺภาคิยานํ สํโยชนานํ ปริกฺขยา อุปหจฺจปรินิพฺพายี โหติ, โน เจ ปญฺจนฺนํ โอรมฺภาคิยานํ สํโยชนานํ ปริกฺขยา อสงฺขารปรินิพฺพายี โหติ, โน เจ ปญฺจนฺนํ โอรมฺภาคิยานํ สํโยชนานํ ปริกฺขยา สสงฺขารปรินิพฺพายี โหติ, อถ ปญฺจนฺนํ โอรมฺภาคิยานํ สํโยชนานํ ปริกฺขยา อุทฺธํโสโต โหติ อกนิฏฺฐคามี, เอวํ ภาวิเตสุ โข ภิกฺขเว สตฺตสุ โพชฺฌงฺเคสุ เอวํ พหุลีกเตสุ อิเม สตฺต ผลา สตฺตานิสํสา ปาฏิกงฺขาติ.\csromanspacer{}\csromanspacer{}ตติยํ.}

\csromanheader{2}{4. วตฺถสุตฺต}
\proseitem{185}{เอกํ สมยํ อายสฺมา สาริปุตฺโต สาวตฺถิยํ วิหรติ เชตวเน อนาถปิณฺฑิกสฺส อาราเม.\csromanspacer{} ตตฺร โข อายสฺมา สาริปุตฺโต ภิกฺขู อามนฺเตสิ “อาวุโส ภิกฺขโว”ติ. “อาวุโส”ติ โข เต ภิกฺขู อายสฺมโต สาริปุตฺตสฺส ปจฺจสฺโสสุํ.\csromanspacer{} อายสฺมา สาริปุตฺโต เอตทโวจ\csromandash{}}

\prose{สตฺติเม อาวุโส โพชฺฌงฺคา.\csromanspacer{} กตเม สตฺต, สติสมฺโพชฺฌงฺโค ธมฺมวิจยสมฺโพชฺฌงฺโค วีริยสมฺโพชฺฌงฺโค ปีติสมฺโพชฺฌงฺโค ปสฺสทฺธิสมฺโพชฺฌงฺโค สมาธิสมฺโพชฺฌงฺโค อุเปกฺขาสมฺโพชฺฌงฺโค, อิเม โข อาวุโส สตฺต โพชฺฌงฺคา.\csromanspacer{} อิเมสํ ขฺวาหํ อาวุโส สตฺตนฺนํ โพชฺฌงฺคานํ เยน เยน โพชฺฌงฺเคน อากงฺขามิ ปุพฺพณฺหสมยํ วิหริตุํ, เตน เตน โพชฺฌงฺเคน ปุพฺพณฺหสมยํ วิหรามิ.\csromanspacer{} เยน เยน โพชฺฌงฺเคน อากงฺขามิ มชฺฌนฺหิกํ สมยํ วิหริตุํ, เตน เตน โพชฺฌงฺเคน มชฺฌนฺหิกํ สมยํ วิหรามิ.\csromanspacer{} เยน เยน โพชฺฌงฺเคน อากงฺขามิ สายนฺหสมยํ วิหริตุํ, เตน เตน โพชฺฌงฺเคน สายนฺหสมยํ วิหรามิ.\csromanspacer{} สติสมฺโพชฺฌงฺโค อิติ เจ เม อาวุโส โหติ, อปฺปมาโณติ เม โหติ, สุสมารทฺโธติ เม โหติ, ติฏฺฐนฺตํ จ นํ ติฏฺฐตีติ ปชานามิ, สเจปิ เม จวติ, อิทปฺปจฺจยา เม จวตีติ ปชานามิ ฯเปฯ.\csromanspacer{} อุเปกฺขาสมฺโพชฺฌงฺโค อิติ เจ เม อาวุโส โหติ, อปฺปมาโณติ เม โหติ, สุสมารทฺโธติ เม โหติ, ติฏฺฐนฺตํ จ นํ ติฏฺฐตีติ ปชานามิ, สเจปิ เม จวติ, อิทปฺปจฺจยา เม จวตีติ ปชานามิ.}

\prose{เสยฺยถาปิ อาวุโส รญฺโญ วา ราชมหามตฺตสฺส วา นานารตฺตานํ ทุสฺสานํ ทุสฺสกรณฺฑโก ปูโร อสฺส, โส ยญฺญเทว ทุสฺสยุคํ}

\csromanheader{2}{2. โพชฺฌงฺคสํยุตฺต}
\prose{\csromanfolio{65}อากงฺเขยฺย ปุพฺพณฺหสมยํ ปารุปิตุํ, ตํ ตเทว ทุสฺสยุคํ ปุพฺพณฺหสมยํ ปารุเปยฺย.\csromanspacer{} ยญฺญเทว ทุสฺสยุคํ อากงฺเขยฺย มชฺฌนฺหิกํ สมยํ ปารุปิตุํ, ตํ ตเทว ทุสฺสยุคํ มชฺฌนฺหิกํ สมยํ ปารุเปยฺย.\csromanspacer{} ยญฺญเทว ทุสฺสยุคํ อากงฺเขยฺย สายนฺหสมยํ ปารุปิตุํ, ตํ ตเทว ทุสฺสยุคํ สายนฺหสมยํ ปารุเปยฺย.\csromanspacer{} เอวเมว ขฺวาหํ อาวุโส อิเมสํ สตฺตนฺนํ โพชฺฌงฺคานํ เยน เยน โพชฺฌงฺเคน อากงฺขามิ ปุพฺพณฺหสมยํ วิหริตุํ, เตน เตน โพชฺฌงฺเคน ปุพฺพณฺหสมยํ วิหรามิ.\csromanspacer{} เยน เยน โพชฺฌงฺเคน อากงฺขามิ มชฺฌนฺหิกํ สมยํ วิหริตุํ, เตน เตน โพชฺฌงฺเคน มชฺฌนฺหิกํ สมยํ วิหรามิ.\csromanspacer{} เยน เยน โพชฺฌงฺเคน อากงฺขามิ สายนฺหสมยํ วิหริตุํ, เตน เตน โพชฺฌงฺเคน สายนฺหสมยํ วิหรามิ.\csromanspacer{} สติสมฺโพชฺฌงฺโค อิติ เจ เม อาวุโส โหติ, อปฺปมาโณติ เม โหติ, สุสมารทฺโธติ เม โหติ.\csromanspacer{} ติฏฺฐนฺตํ จ นํ ติฏฺฐตีติ ปชานามิ, สเจปิ เม จวติ, อิทปฺปจฺจยา เม จวตีติ ปชานามิ ฯเปฯ.\csromanspacer{} อุเปกฺขาสมฺโพชฺฌงฺโค อิติ เจ เม อาวุโส โหติ, อปฺปมาโณติ เม โหติ, สุสมารทฺโธติ เม โหติ, ติฏฺฐนฺตํ จ นํ ติฏฺฐตีติ ปชานามิ, สเจปิ เม จวติ, อิทปฺปจฺจยา เม จวตีติ ปชานามี”ติ.\csromanspacer{}\csromanspacer{}จตุตฺถํ.}

\csromanheader{2}{5. ภิกฺขุสุตฺต}
\proseitem{186}{สาวตฺถินิทานํ.\csromanspacer{} อถ โข อญฺญตโร ภิกฺขุ เยน ภควา เตนุปสงฺกมิ ฯเปฯ เอกมนฺตํ นิสินฺโน โข โส ภิกฺขุ ภควนฺตํ เอตทโวจ “โพชฺฌงฺคา โพชฺฌงฺคาติ ภนฺเต วุจฺจนฺติ, กิตฺตาวตา นุ โข ภนฺเต โพชฺฌงฺคาติ วุจฺจนฺตี”ติ.\csromanspacer{} โพธาย สํวตฺตนฺตีติ โข ภิกฺขุ ตสฺมา โพชฺฌงฺคาติ วุจฺจนฺติ.\csromanspacer{} อิธ ภิกฺขุ สติสมฺโพชฺฌงฺคํ ภาเวติ วิเวกนิสฺสิตํ วิราคนิสฺสิตํ นิโรธนิสฺสิตํ โวสฺสคฺคปริณามิํ ฯเปฯ.\csromanspacer{} อุเปกฺขาสมฺโพชฺฌงฺคํ ภาเวติ วิเวกนิสฺสิตํ วิราคนิสฺสิตํ นิโรธนิสฺสิตํ โวสฺสคฺคปริณามิํ.\csromanspacer{} ตสฺสิเม สตฺต โพชฺฌงฺเค ภาวยโต กามาสวาปิ จิตฺตํ วิมุจฺจติ, ภวาสวาปิ จิตฺตํ วิมุจฺจติ, อวิชฺชาสวาปิ จิตฺตํ วิมุจฺจติ, วิมุตฺตสฺมิํ “วิมุตฺตมฺ”อิติ ญาณํ โหติ, “ขีณา ชาติ, วุสิตํ พฺรหฺมจริยํ, กตํ กรณียํ, นาปรํ อิตฺถตฺตายา”ติ ปชานาติ.\csromanspacer{} โพธาย สํวตฺตนฺตีติ ภิกฺขุ ตสฺมา โพชฺฌงฺคาติ วุจฺจนฺตีติ.\csromanspacer{}\csromanspacer{}ปญฺจมํ.}

\csromanheader{2}{6. กุณฺฑลิยสุตฺต}
\proseitem{187}{\csromanfolio{66}เอกํ สมยํ ภควา สาเกเต วิหรติ อญฺจนวเน มิคทาเย.\csromanspacer{} อถ โข กุณฺฑลิโย ปริพฺพาชโก เยน ภควา เตนุปสงฺกมิ, อุปสงฺกมิตฺวา ภควตา สทฺธิํ สมฺโมทิ, สมฺโมทนียํ กถํ สารณียํ วีติสาเรตฺวา เอกมนฺตํ นิสีทิ, เอกมนฺตํ นิสินฺโน โข กุณฺฑลิโย ปริพฺพาชโก ภควนฺตํ เอตทโวจ “อหมสฺมิ โภ โคตม อารามนิสฺสิยี\footnote{อารามนิสาที (สี), อารามนิยาที (สฺยา)} ปริสาวจโร, ตสฺส มยฺหํ โภ โคตม ปจฺฉาภตฺตํ ภุตฺตปาตราสสฺส อยมาจาโร\footnote{อยมาหาโร (สฺยา, ก)} โหติ ‘อาราเมน อารามํ อุยฺยาเนน อุยฺยานํ อนุจงฺกมามิ อนุวิจรามิ, โส ตตฺถ ปสฺสามิ เอเก สมณพฺราหฺมเณ อิติวาทปฺปโมกฺขานิสํสญฺเจว กถํ กเถนฺเต อุปารมฺภานิสํสญฺจ, ภวํ ปน โคตโม กิมานิสํโส วิหรตี”ติ.\csromanspacer{} วิชฺชาวิมุตฺติผลานิสํโส โข กุณฺฑลิย ตถาคโต วิหรตีติ.}

\prose{กตเม ปน โภ โคตม ธมฺมา ภาวิตา พหุลีกตา วิชฺชาวิมุตฺติํ ปริปูเรนฺตีติ.\csromanspacer{} สตฺต โข กุณฺฑลิย โพชฺฌงฺคา ภาวิตา พหุลีกตา วิชฺชาวิมุตฺติํ ปริปูเรนฺตีติ.\csromanspacer{} กตเม ปน โภ โคตม ธมฺมา ภาวิตา พหุลีกตา สตฺต โพชฺฌงฺเค ปริปูเรนฺตีติ, จตฺตาโร โข กุณฺฑลิย สติปฏฺฐานา ภาวิตา พหุลีกตา สตฺต โพชฺฌงฺเค ปริปูเรนฺตีติ.\csromanspacer{} กตเม ปน โภ โคตม ธมฺมา ภาวิตา พหุลีกตา จตฺตาโร สติปฏฺฐาเน ปริปูเรนฺตีติ, ตีณิ โข กุณฺฑลิย สุจริตานิ ภาวิตานิ พหุลีกตานิ จตฺตาโร สติปฏฺฐาเน ปริปูเรนฺตีติ.\csromanspacer{} กตเม ปน โภ โคตม ธมฺมา ภาวิตา พหุลีกตา ตีณิ สุจริตานิ ปริปูเรนฺตีติ, อินฺทฺริยสํวโร โข กุณฺฑลิย ภาวิโต พหุลีกโต ตีณิ สุจริตานิ ปริปูเรตีติ.}

\prose{กถํ ภาวิโต จ กุณฺฑลิย อินฺทฺริยสํวโร กถํ พหุลีกโต ตีณิ สุจริตานิ ปริปูเรนฺตีติ, อิธ กุณฺฑลิย ภิกฺขุ จกฺขุนา รูปํ ทิสฺวา มนาปํ นาภิชฺฌติ นาภิหํสติ น ราคํ ชเนติ, ตสฺส ฐิโต จ กาโย โหติ, ฐิตํ จิตฺตํ อชฺฌตฺตํ สุสณฺฐิตํ สุวิมุตฺตํ.\csromanspacer{} จกฺขุนา โข ปเนว รูปํ ทิสฺวา อมนาปํ น มงฺกุ โหติ อปฺปติฏฺฐิตจิตฺโต อทีนมานโส อพฺยาปนฺนเจตโส.\csromanspacer{} ตสฺส ฐิโต จ กาโย โหติ, ฐิตํ จิตฺตํ อชฺฌตฺตํ สุสณฺฐิตํ สุวิมุตฺตํ.}

\csromanheader{2}{2. โพชฺฌงฺคสํยุตฺต}
\prose{\csromanfolio{67}ปุน จปรํ กุณฺฑลิย ภิกฺขุ โสเตน สทฺทํ สุตฺวา ฯเปฯ ฆาเนน คนฺธํ ฆายิตฺวา ฯเปฯ ชิวฺหาย รสํ สายิตฺวา ฯเปฯ กาเยน โผฏฺฐพฺพํ ผุสิตฺวา ฯเปฯ มนสา ธมฺมํ วิญฺญาย มนาปํ นาภิชฺฌติ นาภิหํสติ น ราคํ ชเนติ.\csromanspacer{} ตสฺส ฐิโต จ กาโย โหติ ฐิตํ จิตฺตํ อชฺฌตฺตํ สุสณฺฐิตํ สุวิมุตฺตํ.\csromanspacer{} มนสา โข ปเนว ธมฺมํ วิญฺญาย อมนาปํ น มงฺกุ โหติ อปฺปติฏฺฐิตจิตฺโต อทีนมานโส อพฺยาปนฺนเจตโส.\csromanspacer{} ตสฺส ฐิโต จ กาโย โหติ, ฐิตํ จิตฺตํ อชฺฌตฺตํ สุสณฺฐิตํ สุวิมุตฺตํ.}

\prose{ยโต โข กุณฺฑลิย ภิกฺขุโน จกฺขุนา รูปํ ทิสฺวา มนาปามนาเปสุ รูเปสุ ฐิโต จ กาโย โหติ, ฐิตํ จิตฺตํ อชฺฌตฺตํ สุสณฺฐิตํ สุวิมุตฺตํ.\csromanspacer{} โสเตน สทฺทํ สุตฺวา ฯเปฯ.\csromanspacer{} ฆาเนน คนฺธํ ฆายิตฺวา ฯเปฯ.\csromanspacer{} ชิวฺหาย รสํ สายิตฺวา ฯเปฯ.\csromanspacer{} กาเยน โผฏฺฐพฺพํ ผุสิตฺวา ฯเปฯ.\csromanspacer{} มนสา ธมฺมํ วิญฺญาย มนาปามนาเปสุ ธมฺเมสุ ฐิโต จ กาโย โหติ, ฐิตํ จิตฺตํ อชฺฌตฺตํ สุสณฺฐิตํ สุวิมุตฺตํ.\csromanspacer{} เอวํ ภาวิโต โข กุณฺฑลิย อินฺทฺริยสํวโร เอวํ พหุลีกโต ตีณิ สุจริตานิ ปริปูเรติ.}

\prose{กถํ ภาวิตานิ จ กุณฺฑลิย ตีณิ สุจริตานิ กถํ พหุลีกตานิ จตฺตาโร สติปฏฺฐาเน ปริปูเรนฺติ, อิธ กุณฺฑลิย ภิกฺขุ กายทุจฺจริตํ ปหาย กายสุจริตํ ภาเวติ, วจีทุจฺจริตํ ปหาย วจีสุจริตํ ภาเวติ, มโนทุจฺจริตํ ปหาย มโนสุจริตํ ภาเวติ.\csromanspacer{} เอวํ ภาวิตานิ โข กุณฺฑลิย ตีณิ สุจริตานิ เอวํ พหุลีกตานิ จตฺตาโร สติปฏฺฐาเน ปริปูเรนฺติ.}

\prose{กถํ ภาวิตา จ กุณฺฑลิย จตฺตาโร สติปฏฺฐานา กถํ พหุลีกตา สตฺต โพชฺฌงฺเค ปริปูเรนฺติ, อิธ กุณฺฑลิย ภิกฺขุ กาเย กายานุปสฺสี วิหรติ อาตาปี สมฺปชาโน สติมา วิเนยฺย โลเก อภิชฺฌาโทมนสฺสํ.\csromanspacer{} เวทนาสุ ฯเปฯ.\csromanspacer{} จิตฺเต ฯเปฯ.\csromanspacer{} ธมฺเมสุ ธมฺมานุปสฺสี วิหรติ อาตาปี สมฺปชาโน สติมา วิเนยฺย โลเก อภิชฺฌาโทมนสฺสํ.\csromanspacer{} เอวํ ภาวิตา โข กุณฺฑลิย จตฺตาโร สติปฏฺฐานา เอวํ พหุลีกตา สตฺต โพชฺฌงฺเค ปริปูเรนฺติ.}

\prose{กถํ ภาวิตา จ กุณฺฑลิย สตฺต โพชฺฌงฺคา กถํ พหุลีกตา วิชฺชาวิมุตฺติํ ปริปูเรนฺติ, อิธ กุณฺฑลิย ภิกฺขุ สติสมฺโพชฺฌงฺคํ ภาเวติ วิเวกนิสฺสิตํ วิราคนิสฺสิตํ นิโรธนิสฺสิตํ โวสฺสคฺคปริณามิํ ฯเปฯ.}

\prose{\csromanfolio{68}อุเปกฺขาสมฺโพชฺฌางฺคํ ภาเวติ วิเวกนิสฺสิตํ วิราคนิสฺสิตํ นิโรธนิสฺสิตํ โวสฺสคฺคปริณามิํ.\csromanspacer{} เอวํ ภาวิตา โข กุณฺฑลิย สตฺต โพชฺฌงฺคา เอวํ พหุลีกตา วิชฺชาวิมุตฺติํ ปริปูเรนฺตีติ.}

\prose{เอวํ วุตฺเต กุณฺฑลิโย ปริพฺพาชโก ภควนฺตํ เอตทโวจ “อภิกฺกนฺตํ โภ โคตม, อภิกฺกนฺตํ โภ โคตม, เสยฺยถาปิ โภ โคตม นิกฺกุชฺชิตํ วา อุกฺกุชฺเชยฺย, ปฏิจฺฉนฺนํ วา วิจเรยฺย, มูฬฺหสฺส วา มคฺคํ อาจิกฺเขยฺย, อนฺธกาเร วา เตลปชฺโชตํ ธาเรยฺย ‘จกฺขุมนฺโต รูปานิ ทกฺขนฺตี’ติ, เอวเมว โภตา โคตเมน อเนกปริยาเยน ธมฺโม ปกาสิโต, เอสาหํ ภวนฺตํ โคตมํ สรณํ คจฺฉามิ ธมฺมญฺจ ภิกฺขุสํฆญฺจ, อุปาสกํ มํ ภวํ โคตโม ธาเรตุ อชฺชตคฺเค ปาณุเปตํ สรณํ คตนฺ”ติ.\csromanspacer{}\csromanspacer{}ฉฏฺฐํ.}

\csromanheader{2}{7. กูฏาคารสุตฺต}
\proseitem{188}{เสยฺยถาปิ ภิกฺขเว กูฏาคารสฺส ยา กาจิ โคปานสิโย, สพฺพา ตา กูฏนินฺนา กูฏโปณา กูฏปพฺภารา.\csromanspacer{} เอวเมว โข ภิกฺขเว ภิกฺขุ สตฺต โพชฺฌงฺเค ภาเวนฺโต สตฺต โพชฺฌงฺเค พหุลีกโรนฺโต นิพฺพานนินฺโน โหติ นิพฺพานโปโณ นิพฺพานปพฺภาโร.}

\prose{กถญฺจ ภิกฺขเว ภิกฺขุ สตฺต โพชฺฌงฺเค ภาเวนฺโต สตฺต โพชฺฌงฺเค พหุลีกโรนฺโต นิพฺพานนินฺโน โหติ นิพฺพานโปโณ นิพฺพานปพฺภาโร, อิธ ภิกฺขเว ภิกฺขุ สติสมฺโพชฺฌงฺคํ ภาเวติ วิเวกนิสฺสิตํ วิราคนิสฺสิตํ นิโรธนิสฺสิตํ โวสฺสคฺคปริณามิํ ฯเปฯ.\csromanspacer{} อุเปกฺขาสมฺโพชฺฌงฺคํ ภาเวติ วิเวกนิสฺสิตํ วิราคนิสฺสิตํ นิโรธนิสฺสิตํ โวสฺสคฺคปริณามิํ.\csromanspacer{} เอวํ โข ภิกฺขเว ภิกฺขุ สตฺต โพชฺฌงฺเค ภาเวนฺโต สตฺต โพชฺฌงฺเค พหุลีกโรนฺโต นิพฺพานนินฺโน โหติ นิพฺพานโปโณ นิพฺพานปพฺภาโรติ.\csromanspacer{}\csromanspacer{}สตฺตมํ.}

\csromanheader{2}{8. อุปวานสุตฺต}
\proseitem{189}{เอกํ สมยํ อายสฺมา จ อุปวาโน อายสฺมา จ สาริปุตฺโต โกสมฺพิยํ วิหรนฺติ โฆสิตาราเม.\csromanspacer{} อถ โข อายสฺมา สาริปุตฺโต สายนฺหสมยํ ปฏิสลฺลานา วุฏฺฐิโต เยนายสฺมา อุปวาโน เตนุปสงฺกมิ, อุปสงฺกมิตฺวา อายสฺมตา อุปวาเนน สทฺธิํ สมฺโมทิ, สมฺโมทนียํ กถํ สารณียํ วีติสาเรตฺวา เอกมนฺตํ นิสีทิ, เอกมนฺตํ นิสินฺโน โข อายสฺมา สาริปุตฺโต อายสฺมนฺตํ อุปวานํ เอตทโวจ\csromandash{}}

\csromanheader{2}{2. โพชฺฌงฺคสํยุตฺต}
\prose{\csromanfolio{69}“ชาเนยฺย นุ โข อาวุโส อุปวาน ภิกฺขุ ปจฺจตฺตํ โยนิโสมนสิการา เอวํ สุสมารทฺธา เม สตฺต โพชฺฌงฺคา ผาสุวิหาราย สํวตฺตนฺตี”ติ.\csromanspacer{} ชาเนยฺย โข อาวุโส สาริปุตฺต ภิกฺขุ ปจฺจตฺตํ โยนิโสมนสิการา เอวํ สุสมารทฺธา เม สตฺต โพชฺฌงฺคา ผาสุวิหาราย สํวตฺตนฺตีติ.}

\prose{สติสมฺโพชฺฌงฺคํ โข อาวุโส ภิกฺขุ อารพฺภมาโน ปชานาติ “จิตฺตญฺจ เม สุวิมุตฺตํ, ถินมิทฺธญฺจ เม สุสมูหตํ, อุทฺธจฺจกุกฺกุจฺจญฺจ เม สุปฺปฏิวินีตํ, อารทฺธญฺจ เม วีริยํ, อฏฺฐิํกตฺวา มนสิ กโรมิ โน จ ลีนนฺ”ติ ฯเปฯ.\csromanspacer{} อุเปกฺขาสมฺโพชฺฌงฺคํ อาวุโส ภิกฺขุ อารพฺภมาโน ปชานาติ “จิตฺตญฺจ เม สุวิมุตฺตํ, ถินมิทฺธญฺจ เม สุสมูหตํ, อุทฺธจฺจกุกฺกุจฺจญฺจ เม สุปฺปฏิวินีตํ, อารทฺธญฺจ เม วีริยํ, อฏฺฐิํกตฺวา มนสิ กโรมิ โน จ ลีนนฺ”ติ.\csromanspacer{} เอวํ โข อาวุโส สาริปุตฺต ภิกฺขุ ชาเนยฺย ปจฺจตฺตํ โยนิโสมนสิการา เอวํ สุสมารทฺธา เม สตฺต โพชฺฌงฺคา ผาสุวิหาราย สํวตฺตนฺตีติ.\csromanspacer{}\csromanspacer{}อฏฺฐมํ.}

\csromanheader{2}{9. ปฐม-อุปฺปนฺนสุตฺต}
\proseitem{190}{สตฺติเม ภิกฺขเว โพชฺฌงฺคา ภาวิตา พหุลีกตา อนุปฺปนฺนา อุปฺปชฺชนฺติ นาญฺญตฺร ตถาคตสฺส ปาตุภาวา อรหโต สมฺมาสมฺพุทฺธสฺส.\csromanspacer{} กตเม สตฺต, สติสมฺโพชฺฌงฺโค ฯเปฯ อุเปกฺขาสมฺโพชฺฌงฺโค.\csromanspacer{} อิเม โข ภิกฺขเว สตฺต โพชฺฌงฺคา ภาวิตา พหุลีกตา อนุปฺปนฺนา อุปฺปชฺชนฺติ นาญฺญตฺร ตถาคตสฺส ปาตุภาวา อรหโต สมฺมาสมฺพุทฺธสฺสาติ.\csromanspacer{}\csromanspacer{}นวมํ.}

\csromanheader{2}{10. ทุติย-อุปฺปนฺนสุตฺต}
\proseitem{191}{สตฺติเม ภิกฺขเว โพชฺฌงฺคา ภาวิตา พหุลีกตา อนุปฺปนฺนา อุปฺปชฺชนฺติ นาญฺญตฺร สุคตวินยา.\csromanspacer{} กตเม สตฺต, สติสมฺโพชฺฌงฺโค ฯเปฯ อุเปกฺขาสมฺโพชฺฌงฺโค.\csromanspacer{} อิเม โข ภิกฺขเว สตฺต โพชฺฌงฺคา ภาวิตา พหุลีกตา อนุปฺปนฺนา อุปฺปชฺชนฺติ นาญฺญตฺร สุคตวินยาติ.\csromanspacer{}\csromanspacer{}ทสมํ.\csromanspacer{} ปพฺพตวคฺโค ปฐโม.}

\titlehead{\textbf{ตสฺสุทฺทานํ}}
\csromangathameasure{หิมวนฺตํ กายํ สีลํ, วตฺถํ ภิกฺขุ จ กุณฺฑลิ. \\ กูฏญฺจ อุปวานญฺจ, อุปฺปนฺนา อปเร ทุเวติ.}
\csromangathagroup{\csromangathastanza{หิมวนฺตํ กายํ สีลํ, วตฺถํ ภิกฺขุ จ กุณฺฑลิ. \\ กูฏญฺจ อุปวานญฺจ, อุปฺปนฺนา อปเร ทุเวติ.}}

\chapterheadpageread{2. คิลานวคฺค}
\csromanheader{2}{1. ปาณสุตฺต}
\proseitem{192}{\csromanfolio{70}เสยฺยถาปิ ภิกฺขเว เย เกจิ ปาณา จตฺตาโร อิริยาปเถ กปฺเปนฺติ กาเลน คมนํ กาเลน ฐานํ กาเลน นิสชฺชํ กาเลน เสยฺยํ, สพฺเพ เต ปถวิํ นิสฺสาย, ปถวิยํ ปติฏฺฐาย, เอวเมเต จตฺตาโร อิริยาปเถ กปฺเปนฺติ.\csromanspacer{} เอวเมว โข ภิกฺขเว ภิกฺขุ สีลํ นิสฺสาย สีเล ปติฏฺฐาย สตฺต โพชฺฌงฺเค ภาเวติ สตฺต โพชฺฌงฺเค พหุลีกโรติ.}

\prose{กถญฺจ ภิกฺขเว ภิกฺขุ สีลํ นิสฺสาย สีเล ปติฏฺฐาย สตฺต โพชฺฌงฺเค ภาเวติ สตฺต โพชฺฌงฺเค พหุลีกโรติ อิธ ภิกฺขเว ภิกฺขุ สติสมฺโพชฺฌงฺคํ ภาเวติ วิเวกนิสฺสิตํ วิราคนิสฺสิตํ นิโรธนิสฺสิตํ โวสฺสคฺคปริณามิํ ฯเปฯ.\csromanspacer{} อุเปกฺขาสมฺโพชฺฌงฺคํ ภาเวติ วิเวกนิสฺสิตํ วิราคนิสฺสิตํ นิโรธนิสฺสิตํ โวสฺสคฺคปริณามิํ.\csromanspacer{} เอวํ โข ภิกฺขเว ภิกฺขุ สีลํ นิสฺสาย สีเล ปติฏฺฐาย สตฺต โพชฺฌงฺเค ภาเวติ สตฺต โพชฺฌงฺเค พหุลีกโรตีติ.\csromanspacer{}\csromanspacer{}ปฐมํ.}

\csromanheader{2}{2. ปฐมสูริยูปมสุตฺต}
\proseitem{193}{สูริยสฺส ภิกฺขเว อุทยโต เอตํ ปุพฺพงฺคมํ เอตํ ปุพฺพนิมิตฺตํ, ยทิทํ อรุณุคฺคํ.\csromanspacer{} เอวเมว โข ภิกฺขเว ภิกฺขุโน สตฺตนฺนํ โพชฺฌงฺคานํ อุปฺปาทาย เอตํ ปุพฺพงฺคมํ เอตํ นิมิตฺตํ, ยทิทํ กลฺยาณมิตฺตตา.\csromanspacer{} กลฺยาณมิตฺตสฺเสตํ ภิกฺขเว ภิกฺขุโน ปาฏิกงฺขํ สตฺต โพชฺฌงฺเค ภาเวสฺสติ สตฺต โพชฺฌงฺเค พหุลีกริสฺสติ.}

\prose{กถญฺจ ภิกฺขเว ภิกฺขุ กลฺยาณมิตฺโต สตฺต โพชฺฌงฺเค ภาเวติ สตฺต โพชฺฌงฺเค พหุลีกโรติ, อิธ ภิกฺขเว ภิกฺขุ สติสมฺโพชฺฌงฺคํ ภาเวติ วิเวกนิสฺสิตํ ฯเปฯ.\csromanspacer{} อุเปกฺขาสมฺโพชฺฌงฺคํ ภาเวติ วิเวกนิสฺสิตํ วิราคนิสฺสิตํ นิโรธนิสฺสิตํ โวสฺสคฺคปริณามิํ.\csromanspacer{} เอวํ โข ภิกฺขเว ภิกฺขุ กลฺยาณมิตฺโต สตฺต โพชฺฌงฺเค ภาเวติ สตฺต โพชฺฌงฺเค พหุลีกโรตีติ.\csromanspacer{}\csromanspacer{}ทุติยํ.}

\csromanheader{2}{2. โพชฺฌงฺคสํยุตฺต}
\csromanheader{2}{3. ทุติยสูริยูปมสุตฺต}
\proseitem{194}{\csromanfolio{71}สูริยสฺส ภิกฺขเว อุทยโต เอตํ ปุพฺพงฺคมํ เอตํ ปุพฺพนิมิตฺตํ, ยทิทํ อรุณุคฺคํ.\csromanspacer{} เอวเมว โข ภิกฺขเว ภิกฺขุโน สตฺตนฺนํ โพชฺฌงฺคานํ อุปฺปาทาย เอตํ ปุพฺพงฺคมํ เอตํ ปุพฺพนิมิตฺตํ, ยทิทํ โยนิโสมนสิกาโร.\csromanspacer{} โยนิโสมนสิการสมฺปนฺนสฺเสตํ ภิกฺขเว ภิกฺขุโน ปาฏิกงฺขํ สตฺต โพชฺฌงฺเค ภาเวสฺสติ, สตฺต โพชฺฌงฺเค พหุลีกริสฺสติ.}

\prose{กถญฺจ ภิกฺขเว ภิกฺขุ โยนิโสมนสิการสมฺปนฺโน สตฺต โพชฺฌงฺเค ภาเวติ สตฺต โพชฺฌงฺเค พหุลีกโรติ, อิธ ภิกฺขเว ภิกฺขุ สติสมฺโพชฺฌงฺคํ ภาเวติ วิเวกนิสฺสิตํ ฯเปฯ.\csromanspacer{} อุเปกฺขาสมฺโพชฺฌางฺคํ ภาเวติ วิเวกนิสฺสิตํ วิราคนิสฺสิตํ นิโรธนิสฺสิตํ โวสฺสคฺคปริณามิํ.\csromanspacer{} เอวํ โข ภิกฺขเว ภิกฺขุ โยนิโสมนสิการสมฺปนฺโน สตฺต โพชฺฌงฺเค ภาเวติ สตฺต โพชฺฌงฺเค พหุลีกโรตีติ.\csromanspacer{}\csromanspacer{}ตติยํ.}

\csromanheader{2}{4. ปฐมคิลานสุตฺต}
\proseitem{195}{เอกํ สมยํ ภควา ราชคเห วิหรติ เวฬุวเน กลนฺทกนิวาเป.\csromanspacer{} เตน โข ปน สมเยน อายสฺมา มหากสฺสโป ปิปฺปลิคุหายํ\footnote{ปิปฺผลิคุหายํ (สี)} วิหรติ อาพาธิโก ทุกฺขิโต พาฬฺหคิลาโน.\csromanspacer{} อถ โข ภควา สายนฺหสมยํ ปฏิสลฺลานา วุฏฺฐิโต เยนายสฺมา มหากสฺสโป เตนุปสงฺกมิ, อุปสงฺกมิตฺวา ปญฺญตฺเต อาสเน นิสีทิ, นิสชฺช โข ภควา อายสฺมนฺตํ มหากสฺสปํ เอตทโวจ\csromandash{}}

\prose{“กจฺจิ เต กสฺสป ขมนียํ, กจฺจิ ยาปนียํ, กจฺจิ ทุกฺขา เวทนา ปฏิกฺกมนฺติ โน อภิกฺกมนฺติ, ปฏิกฺกโมสานํ ปญฺญายติ โน อภิกฺกโม”ติ.\csromanspacer{} น เม ภนฺเต ขมนียํ, น ยาปนียํ, พฬฺหา เม ทุกฺขา เวทนา อภิกฺกมนฺติ โน ปฏิกฺกมนฺติ, อภิกฺกโมสานํ ปญฺญายติ โน ปฏิกฺกโมติ.}

\prose{สตฺติเม กสฺสป โพชฺฌงฺคา มยา สมฺมทกฺขาตา ภาวิตา พหุลีกตา อภิญฺญาย สมฺโพธาย นิพฺพานาย สํวตฺตนฺติ.\csromanspacer{} กตเม สตฺต, สติสมฺโพชฺฌงฺโค โข กสฺสป มยา สมฺมทกฺขาโต ภาวิโต พหุลีกโต อภิญฺญาย สมฺโพธาย นิพฺพานาย สํวตฺตติ ฯเปฯ.\csromanspacer{} อุเปกฺขาสมฺโพชฺฌงฺโค โข กสฺสป มยา สมฺมทกฺขาโต ภาวิโต พหุลีกโต อภิญฺญาย \csromanfolio{72}สมฺโพธาย นิพฺพานาย สํวตฺตติ.\csromanspacer{} อิเม โข กสฺสป สตฺต โพชฺฌงฺคา มยา สมฺมทกฺขาตา ภาวิตา พหุลีกตา อภิญฺญาย สมฺโพธาย นิพฺพานาย สํวตฺตนฺตีติ.\csromanspacer{} ตคฺฆ ภควา โพชฺฌงฺคา, ตคฺฆ สุคต โพชฺฌงฺคาติ.}

\prose{อิทมโวจ ภควา.\csromanspacer{} อตฺตมโน อายสฺมา มหากสฺสโป ภควโต ภาสิตํ อภินนฺทิ.\csromanspacer{} วุฏฺฐหิ จายสฺมา มหากสฺสโป ตมฺหา อาพาธา, ตถาปหีโน จายสฺมโต มหากสฺสปสฺส โส อาพาโธ อโหสีติ.\csromanspacer{}\csromanspacer{}จตุตฺถํ.}

\csromanheader{2}{5. ทุติยคิลานสุตฺต}
\proseitem{196}{เอกํ สมยํ ภควา ราชคเห วิหรติ เวฬุวเน กลนฺทกนิวาเป.\csromanspacer{} เตน โข ปน สมเยน อายสฺมา มหาโมคฺคลฺลาโน คิชฺฌกูเฏ ปพฺพเต วิหรติ อาพาธิโก ทุกฺขิโต พาฬฺหคิลาโน.\csromanspacer{} อถ โข ภควา สายนฺหสมยํ ปฏิสลฺลานา วุฏฺฐิโต เยนายสฺมา มหาโมคฺคลฺลาโน เตนุปสงฺกมิ, อุปสงฺกมิตฺวา ปญฺญตฺเต อาสเน นิสีทิ, นิสชฺช โข ภควา อายสฺมนฺตํ มหาโมคฺคลฺลานํ เอตทโวจ\csromandash{}}

\prose{“กจฺจิ เต โมคฺคลฺลาน ขมนียํ, กจฺจิ ยาปนียํ, กจฺจิ ทุกฺขา เวทนา ปฏิกฺกมนฺติ โน อภิกฺกมนฺติ, ปฏิกฺกโมสานํ ปญฺญายติ โน อภิกฺกโม”ติ.\csromanspacer{} น เม ภนฺเต ขมนียํ, น ยาปนียํ, พาฬฺหา เม ทุกฺขา เวทนา อภิกฺกมนฺติ โน ปฏิกฺกมนฺติ, อภิกฺกโมสานํ ปญฺญายติ โน ปฏิกฺกโมติ.}

\prose{สตฺติเม โมคฺคลฺลาน โพชฺฌงฺคา มยา สมฺมทกฺขาตา ภาวิตา พหุลีกตา อภิญฺญาย สมฺโพธาย นิพฺพานาย สํวตฺตนฺติ.\csromanspacer{} กตเม สตฺต, สติสมฺโพชฺฌงฺโค โข โมคฺคลฺลาน มยา สมฺมทกฺขาโต ภาวิโต พหุลีกโต อภิญฺญาย สมฺโพธาย นิพฺพานาย สํวตฺตติ ฯเปฯ.\csromanspacer{} อุเปกฺขาสมฺโพชฺฌงฺโค โข โมคฺคลฺลาน มยา สมฺมทกฺขาโต ภาวิโต พหุลีกโต อภิญฺญาย สมฺโพธาย นิพฺพานาย สํวตฺตติ.\csromanspacer{} อิเม โข โมคฺคลฺลาน สตฺต โพชฺฌงฺคา มยา สมฺมทกฺขาตา ภาวิตา พหุลีกตา อภิญฺญาย สมฺโพธาย นิพฺพานาย สํวตฺตนฺตีติ.\csromanspacer{} ตคฺฆ ภควา โพชฺฌงฺคา, ตคฺฆ สุคต โพชฺฌงฺคาติ.}

\csromanheader{2}{2. โพชฺฌงฺคสํยุตฺต}
\prose{\csromanfolio{73}อิทมโวจ ภควา.\csromanspacer{} อตฺตมโน อายสฺมา มหาโมคฺคลฺลาโน ภควโต ภาสิตํ อภินนฺทิ.\csromanspacer{} วุฏฺฐหิ จายสฺมา มหาโมคฺคลฺลาโน ตมฺหา อาพาธา, ตถาปหีโน จายสฺมโต มหาโมคฺคลฺลานสฺส โส อาพาโธ อโหสีติ.\csromanspacer{}\csromanspacer{}ปญฺจมํ.}

\csromanheader{2}{6. ตติยคิลานสุตฺต}
\proseitem{197}{เอกํ สมยํ ภควา ราชคเห วิหรติ เวฬุวเน กลนฺทกนิวาเป.\csromanspacer{} เตน โข ปน สมเยน ภควา อาพาธิโก โหติ ทุกฺขิโต พาฬฺหคิลาโน.\csromanspacer{} อถ โข อายสฺมา มหาจุนฺโท เยน ภควา เตนุปสงฺกมิ, อุปสงฺกมิตฺวา ภควนฺตํ อภิวาเทตฺวา เอกมนฺตํ นิสีทิ, เอกมนฺตํ นิสินฺนํ โข อายสฺมนฺตํ มหาจุนฺทํ ภควา เอตทโวจ “ปฏิภนฺตุ ตํ จุนฺท โพชฺฌงฺคา”ติ.}

\prose{สตฺติเม ภนฺเต โพชฺฌงฺคา ภควตา สมฺมทกฺขาตา ภาวิตา พหุลีกตา อภิญฺญาย สมฺโพธาย นิพฺพานาย สํวตฺตนฺติ.\csromanspacer{} กตเม สตฺต, สติสมฺโพชฺฌงฺโค โข ภนฺเต ภควตา สมฺมทกฺขาโต ภาวิโต พหุลีกโต อภิญฺญาย สมฺโพธาย นิพฺพานาย สํวตฺตติ ฯเปฯ.\csromanspacer{} อุเปกฺขาสมฺโพชฺฌงฺโค โข ภนฺเต ภควตา สมฺมทกฺขาโต ภาวิโต พหุลีกโต อภิญฺญาย สมฺโพธาย นิพฺพานาย สํวตฺตติ.\csromanspacer{} อิเม โข ภนฺเต สตฺต โพชฺฌงฺคา ภควตา สมฺมทกฺขาตา ภาวิตา พหุลีกตา อภิญฺญาย สมฺโพธาย นิพฺพานาย สํวตฺตนฺตีติ.\csromanspacer{} ตคฺฆ จุนฺท โพชฺฌงฺคา, ตคฺฆ จุนฺท โพชฺฌงฺคาติ.}

\prose{อิทมโวจายสฺมา จุนฺโท.\csromanspacer{} สมนุญฺโญ สตฺถา อโหสิ.\csromanspacer{} วุฏฺฐหิ จ ภควา ตมฺหา อาพาธา, ตถาปหีโน จ ภควโต โส อาพาโธ อโหสีติ.\csromanspacer{}\csromanspacer{}ฉฏฺฐํ.}

\csromanheader{2}{7. ปารงฺคมสุตฺต}
\proseitem{198}{สตฺติเม ภิกฺขเว โพชฺฌงฺคา ภาวิตา พหุลีกตา อปารา ปารํ คมนาย สํวตฺตนฺติ.\csromanspacer{} กตเม สตฺต, สติสมฺโพชฺฌงฺโค ฯเปฯ อุเปกฺขาสมฺโพชฺฌงฺโค.\csromanspacer{} อิเม โข ภิกฺขเว สตฺต โพชฺฌงฺคา ภาวิตา พหุลีกตา อปารา ปารํ คมนาย สํวตฺตนฺตีติ.}

\csromangathameasure{อปฺปกา เต มนุสฺเสสุ, เย ชนา ปารคามิโน. \\ อถายํ อิตรา ปชา, ตีรเมวานุธาวติ. \\ เย จ โข สมฺมทกฺขาเต, ธมฺเม ธมฺมานุวตฺติโน. \\ เต ชนา ปารเมสฺสนฺติ, มจฺจุเธยฺยํ สุทุตฺตรํ. \\ กณฺหํ ธมฺมํ วิปฺปหาย, สุกฺกํ ภาเวถ ปณฺฑิโต. \\ โอกา อโนกมาคมฺม, วิเวเก ยตฺถ ทูรมํ. \\ ตตฺราภิรติมิจฺเฉยฺย, หิตฺวา กาเม อกิญฺจโน. \\ ปริโยทเปยฺย อตฺตานํ, จิตฺตเกฺลเสหิ ปณฺฑิโต. \\ เยสํ สมฺโพธิยงฺเคสุ, สมฺมา จิตฺตํ สุภาวิตํ. \\ อาทานปฺปฏินิสฺสคฺเค, อนุปาทาย เย รตา. \\ ขีณาสวา ชุติมนฺโต, เต โลเก ปรินิพฺพุตาติ. สตฺตมํ.}
\csromangathagroup{\csromangathastanza{\csromanfolio{74}อปฺปกา เต มนุสฺเสสุ, เย ชนา ปารคามิโน. \\ อถายํ อิตรา ปชา, ตีรเมวานุธาวติ.}\csromangathastanza{เย จ โข สมฺมทกฺขาเต, ธมฺเม ธมฺมานุวตฺติโน. \\ เต ชนา ปารเมสฺสนฺติ, มจฺจุเธยฺยํ สุทุตฺตรํ.}\csromangathastanza{กณฺหํ ธมฺมํ วิปฺปหาย, สุกฺกํ ภาเวถ ปณฺฑิโต. \\ โอกา อโนกมาคมฺม, วิเวเก ยตฺถ ทูรมํ.}\csromangathastanza{ตตฺราภิรติมิจฺเฉยฺย, หิตฺวา กาเม อกิญฺจโน. \\ ปริโยทเปยฺย อตฺตานํ, จิตฺตเกฺลเสหิ ปณฺฑิโต.}\csromangathastanza{เยสํ สมฺโพธิยงฺเคสุ, สมฺมา จิตฺตํ สุภาวิตํ. \\ อาทานปฺปฏินิสฺสคฺเค, อนุปาทาย เย รตา. \\ ขีณาสวา ชุติมนฺโต, เต โลเก ปรินิพฺพุตาติ.\csromanspacer{} สตฺตมํ.}}

\csromanheader{2}{8. วิรทฺธสุตฺต}
\proseitem{199}{เยสํ เกสญฺจิ ภิกฺขเว สตฺต โพชฺฌงฺคา วิรทฺธา, วิรทฺโธ เตสํ อริโย มคฺโค สมฺมา ทุกฺขกฺขยคามี.\csromanspacer{} เยสํ เกสญฺจิ ภิกฺขเว สตฺต โพชฺฌงฺคา อารทฺธา, อารทฺโธ เตสํ อริโย มคฺโค สมฺมา ทุกฺขกฺขยคามี.\csromanspacer{} กตเม สตฺต, สติสมฺโพชฺฌงฺโค ฯเปฯ อุเปกฺขาสมฺโพชฺฌงฺโค.\csromanspacer{} เยสํ เกสญฺจิ ภิกฺขเว อิเม สตฺต โพชฺฌงฺคา วิรทฺธา, วิรทฺโธ เตสํ อริโย มคฺโค สมฺมา ทุกฺขกฺขยคามี.\csromanspacer{} เยสํ เกสญฺจิ ภิกฺขเว อิเม สตฺต โพชฺฌงฺคา อารทฺธา, อารทฺโธ เตสํ อริโย มคฺโค สมฺมา ทุกฺขกฺขยคามีติ.\csromanspacer{}\csromanspacer{}อฏฺฐมํ.}

\csromanheader{2}{9. อริยสุตฺต}
\proseitem{200}{สตฺติเม ภิกฺขเว โพชฺฌงฺคา ภาวิตา พหุลีกตา อริยา นิยฺยานิกา นียนฺติ ตกฺกรสฺส สมฺมา ทุกฺขกฺขยาย.\csromanspacer{} กตเม สตฺต, สติสมฺโพชฺฌงฺโค ฯเปฯ อุเปกฺขาสมฺโพชฺฌงฺโค.\csromanspacer{} อิเม โข ภิกฺขเว สตฺต โพชฺฌงฺคา ภาวิตา พหุลีกตา อริยา นิยฺยานิกา นียนฺติ ตกฺกรสฺส สมฺมา ทุกฺขกฺขยายาติ.\csromanspacer{}\csromanspacer{}นวมํ.}

\csromanheader{2}{2. โพชฺฌงฺคสํยุตฺต}
\csromanheader{2}{10. นิพฺพิทาสุตฺต}
\proseitem{201}{\csromanfolio{75}สตฺติเม ภิกฺขเว โพชฺฌงฺคา ภาวิตา พหุลีกตา เอกนฺตนิพฺพิทาย วิราคาย นิโรธาย อุปสมาย อภิญฺญาย สมฺโพธาย นิพฺพานาย สํวตฺตนฺติ.\csromanspacer{} กตเม สตฺต, สติสมฺโพชฺฌงฺโค ฯเปฯ อุเปกฺขาสมฺโพชฺฌงฺโค.\csromanspacer{} อิเม โข ภิกฺขเว สตฺต โพชฺฌงฺคา ภาวิตา พหุลีกตา เอกนฺตนิพฺพิทาย วิราคาย นิโรธาย อุปสมาย อภิญฺญาย สมฺโพธาย นิพฺพานาย สํวตฺตนฺตีติ.\csromanspacer{}\csromanspacer{}ทสมํ.\csromanspacer{} คิลานวคฺโค ทุติโย.}

\titlehead{\textbf{ตสฺสุทฺทานํ}}
\csromangathameasure{ปาณา สูริยูปมา เทฺว, คิลานา อปเร ตโย. \\ ปารํคามี วิรทฺโธ จ, อริโย นิพฺพิทาย จาติ.}
\csromangathagroup{\csromangathastanza{ปาณา สูริยูปมา เทฺว, คิลานา อปเร ตโย. \\ ปารํคามี วิรทฺโธ จ, อริโย นิพฺพิทาย จาติ.}}

\chapterhead{3. อุทายิวคฺค}
\csromanheader{2}{1. โพธายสุตฺต}
\proseitem{202}{อถ โข อญฺญตโร ภิกฺขุ เยน ภควา เตนุปสงฺกมิ ฯเปฯ เอกมนฺตํ นิสินฺโน โข โส ภิกฺขุ ภควนฺตํ เอตทโวจ\csromandash{}}

\prose{“โพชฺฌงฺคา โพชฺฌงฺคา”ติ ภนฺเต วุจฺจนฺติ, กิตฺตาวตา นุ โข ภนฺเต “โพชฺฌงฺคา”ติ วุจฺจนฺตีติ.\csromanspacer{} โพธาย สํวตฺตนฺตีติ โข ภิกฺขุ, ตสฺมา “โพชฺฌงฺคา”ติ วุจฺจนฺติ, อิธ ภิกฺขุ สติสมฺโพชฺฌงฺคํ ภาเวติ วิเวกนิสฺสิตํ ฯเปฯ.\csromanspacer{} อุเปกฺขาสมฺโพชฺฌงฺคํ ภาเวติ วิเวกนิสฺสิตํ วิราคนิสฺสิตํ นิโรธนิสฺสิตํ โวสฺสคฺคปริณามิํ.\csromanspacer{} โพธาย สํวตฺตนฺตีติ โข ภิกฺขุ, ตสฺมา “โพชฺฌงฺคา”ติ วุจฺจนฺตีติ.\csromanspacer{}\csromanspacer{}ปฐมํ.}

\csromanheader{2}{2. โพชฺฌงฺคเทสนาสุตฺต}
\proseitem{203}{สตฺต โว ภิกฺขเว โพชฺฌงฺเค เทเสสฺสามิ ตํ สุณาถ.\csromanspacer{} กตเม จ ภิกฺขเว สตฺต โพชฺฌงฺคา, สติสมฺโพชฺฌงฺโค ฯเปฯ อุเปกฺขาสมฺโพชฺฌงฺโค.\csromanspacer{} อิเม โข ภิกฺขเว สตฺต โพชฺฌงฺคาติ.\csromanspacer{}\csromanspacer{}ทุติยํ.}

\csromanheader{2}{3. ฐานิยสุตฺต}
\proseitem{204}{\csromanfolio{76}กามราคฏฺฐานิยานํ\footnote{กามราคฏฺฐานียานํ (สี)} ภิกฺขเว ธมฺมานํ มนสิการพหุลีการา อนุปฺปนฺโน เจว กามจฺฉนฺโท อุปฺปชฺชติ, อุปฺปนฺโน จ กามจฺฉนฺโท ภิยฺโยภาวาย เวปุลฺลาย สํวตฺตติ.\csromanspacer{} พฺยาปาทฏฺฐานิยานํ ภิกฺขเว ธมฺมานํ มนสิการพหุลีการา อนุปฺปนฺโน เจว พฺยาปาโท อุปฺปชฺชติ, อุปฺปนฺโน จ พฺยาปาโท ภิยฺโยภาวาย เวปุลฺลาย สํวตฺตติ.\csromanspacer{} ถินมิทฺธฏฺฐานิยานํ ภิกฺขเว ธมฺมานํ มนสิการพหุลีการา อนุปฺปนฺนญฺเจว ถินมิทฺธํ อุปฺปชฺชติ, อุปฺปนฺนญฺจ ถินมิทฺธํ ภิยฺโยภาวาย เวปุลฺลาย สํวตฺตติ.\csromanspacer{} อุทฺธจฺจกุกฺกุจฺจฏฺฐานิยานํ ภิกฺขเว ธมฺมานํ มนสิการพหุลีการา อนุปฺปนฺนญฺเจว อุทฺธจฺจกุกฺกุจฺจํ อุปฺปชฺชติ, อุปฺปนฺนญฺจ อุทฺธจฺจกุกฺกุจฺจํ ภิยฺโยภาวาย เวปุลฺลาย สํวตฺตติ.\csromanspacer{} วิจิกิจฺฉาฏฺฐานิยานํ ภิกฺขเว ธมฺมานํ มนสิการพหุลีการา อนุปฺปนฺนา เจว วิจิกิจฺฉา อุปฺปชฺชติ, อุปฺปนฺนา จ วิจิกิจฺฉา ภิยฺโยภาวาย เวปุลฺลาย สํวตฺตติ.}

\prose{สติสมฺโพชฺฌงฺคฏฺฐานิยานํ ภิกฺขเว ธมฺมานํ มนสิการพหุลีการา อนุปฺปนฺโน เจว สติสมฺโพชฺฌงฺโค อุปฺปชฺชติ, อุปฺปนฺโน จ สติสมฺโพชฺฌงฺโค ภาวนาปาริปูริํ คจฺฉติ ฯเปฯ.\csromanspacer{} อุเปกฺขาสมฺโพชฺฌงฺคฏฺฐานิยานํ ภิกฺขเว ธมฺมานํ มนสิการพหุลีการา อนุปฺปนฺโน เจว อุเปกฺขาสมฺโพชฺฌงฺโค อุปฺปชฺชติ, อุปฺปนฺโน จ อุเปกฺขาสมฺโพชฺฌงฺโค ภาวนาปาริปูริํ คจฺฉตีติ.\csromanspacer{}\csromanspacer{}ตติยํ.}

\csromanheader{2}{4. อโยนิโสมนสิการสุตฺต}
\proseitem{205}{อโยนิโส ภิกฺขเว มนสิกโรโต อนุปฺปนฺโน เจว กามจฺฉนฺโท อุปฺปชฺชติ, อุปฺปนฺโน จ กามจฺฉนฺโท ภิยฺโยภาวาย เวปุลฺลาย สํวตฺตติ.\csromanspacer{} อนุปฺปนฺโน เจว พฺยาปาโท อุปฺปชฺชติ, อุปฺปนฺโน จ พฺยาปาโท ภิยฺโยภาวาย เวปุลฺลาย สํวตฺตติ.\csromanspacer{} อนุปฺปนฺนญฺเจว ถินมิทฺธํ อุปฺปชฺชติ, อุปฺปนฺนญฺจ ถินมิทฺธํ ภิยฺโยภาวาย เวปุลฺลาย สํวตฺตติ.\csromanspacer{} อนุปฺปนฺนญฺเจว อุทฺธจฺจกุกฺกุจฺจํ อุปฺปชฺชติ, อุปฺปนฺนญฺจ อุทฺธจฺจกุกฺกุจฺจํ ภิยฺโยภาวาย เวปุลฺลาย สํวตฺตติ.\csromanspacer{} อนุปฺปนฺนา เจว วิจิกิจฺฉา อุปฺปชฺชติ, อุปฺปนฺนา จ วิจิกิจฺฉา ภิยฺโยภาวาย เวปุลฺลาย สํวตฺตติ.\csromanspacer{} อนุปฺปนฺโน เจว สติสมฺโพชฺฌงฺโค นุปฺปชฺชติ, อุปฺปนฺโน จ สติสมฺโพชฺฌงฺโค นิรุชฺฌติ ฯเปฯ.\csromanspacer{} อนุปฺปนฺโน เจว อุเปกฺขาสมฺโพชฺฌงฺโค นุปฺปชฺชติ, อุปฺปนฺโน จ อุเปกฺขาสมฺโพชฺฌงฺโค นิรุชฺฌติ.}

\csromanheader{2}{2. โพชฺฌงฺคสํยุตฺต}
\prose{\csromanfolio{77}โยนิโส จ โข ภิกฺขเว มนสิกโรโต อนุปฺปนฺโน เจว กามจฺฉนฺโท นุปฺปชฺชติ, อุปฺปนฺโน จ กามจฺฉนฺโท ปหียติ.\csromanspacer{} อนุปฺปนฺโน เจว พฺยาปาโท นุปฺปชฺชติ, อุปฺปนฺโน จ พฺยาปาโท ปหียติ.\csromanspacer{} อนุปฺปนฺนญฺเจว ถินมิทฺธํ นุปฺปชฺชติ, อุปฺปนฺนญฺจ ถินมิทฺธํ ปหียติ.\csromanspacer{} อนุปฺปนฺนญฺเจว อุทฺธจฺจกุกฺกุจฺจํ นุปฺปชฺชติ, อุปฺปนฺนญฺจ อุทฺธจฺจกุกฺกุจฺจํ ปหียติ.\csromanspacer{} อนุปฺปนฺนา เจว วิจิกิจฺฉา นุปฺปชฺชติ, อุปฺปนฺนา จ วิจิกิจฺฉา ปหียติ.}

\prose{อนุปฺปนฺโน เจว สติสมฺโพชฺฌงฺโค อุปฺปชฺชติ, อุปฺปนฺโน จ สติสมฺโพชฺฌงฺโค ภาวนาปาริปูริํ คจฺฉติ ฯเปฯ.\csromanspacer{} อนุปฺปนฺโน เจว อุเปกฺขาสมฺโพชฺฌงฺโค อุปฺปชฺชติ, อุปฺปนฺโน จ อุเปกฺขาสมฺโพชฺฌงฺโค ภาวนาปาริปูริํ คจฺฉตีติ.\csromanspacer{}\csromanspacer{}จตุตฺถํ.}

\csromanheader{2}{5. อปริหานิยสุตฺต}
\proseitem{206}{สตฺต โว ภิกฺขเว อปริหานิเย ธมฺเม เทเสสฺสามิ ตํ สุณาถ.\csromanspacer{} กตเม จ ภิกฺขเว สตฺต อปริหานิยา ธมฺมา, ยทิทํ สตฺต โพชฺฌงฺคา.\csromanspacer{} กตเม สตฺต, สติสมฺโพชฺฌงฺโค ฯเปฯ อุเปกฺขาสมฺโพชฺฌงฺโค.\csromanspacer{} อิเม โข ภิกฺขเว สตฺต อปริหานิยา ธมฺมาติ.\csromanspacer{}\csromanspacer{}ปญฺจมํ.}

\csromanheader{2}{6. ตณฺหกฺขยสุตฺต}
\proseitem{207}{โย ภิกฺขเว มคฺโค.\csromanspacer{} ยา ปฏิปทา ตณฺหกฺขยาย สํวตฺตติ, ตํ มคฺคํ ตํ ปฏิปทํ ภาเวถ.\csromanspacer{} กตโม จ ภิกฺขเว มคฺโค กตมา จ ปฏิปทา ตณฺหกฺขยาย สํวตฺตติ, ยทิทํ สตฺต โพชฺฌงฺคา.\csromanspacer{} กตเม สตฺต, สติสมฺโพชฺฌงฺโค ฯเปฯ อุเปกฺขาสมฺโพชฺฌงฺโคติ.\csromanspacer{} เอวํ วุตฺเต อายสฺมา อุทายี ภควนฺตํ เอตทโวจ “กถํ ภาวิตา นุ โข ภนฺเต สตฺต โพชฺฌงฺคา กถํ พหุลีกตา ตณฺหกฺขยาย สํวตฺตนฺตี”ติ.}

\prose{อิธ อุทายิ ภิกฺขุ สติสมฺโพชฺฌงฺคํ ภาเวติ วิเวกนิสฺสิตํ วิราคนิสฺสิตํ นิโรธนิสฺสิตํ โวสฺสคฺคปริณามิํ วิปุลํ มหคฺคตํ อปฺปมาณํ อพฺยาปชฺชํ\footnote{อพฺยาปชฺฌํ (สี, สฺยา, อิ)}.\csromanspacer{} ตสฺส สติสมฺโพชฺฌงฺคํ ภาวยโต วิเวกนิสฺสิตํ วิราคนิสฺสิตํ นิโรธนิสฺสิตํ โวสฺสคฺคปริณามิํ วิปุลํ มหคฺคตํ อปฺปมาณํ อพฺยาปชฺชํ ตณฺหา ปหียติ, ตณฺหาย ปหานา กมฺมํ ปหียติ, กมฺมสฺส ปหานา ทุกฺขํ ปหียติ ฯเปฯ.\csromanspacer{} อุเปกฺขาสมฺโพชฺฌงฺคํ ภาเวติ วิเวกนิสฺสิตํ \csromanfolio{78}วิราคนิสฺสิตํ นิโรธนิสฺสิตํ โวสฺสคฺคปริณามิํ วิปุลํ มหคฺคตํ อปฺปมาณํ อพฺยาปชฺชํ.\csromanspacer{} ตสฺส อุเปกฺขาสมฺโพชฺฌงฺคํ ภาวยโต วิเวกนิสฺสิตํ วิราคนิสฺสิตํ นิโรธนิสฺสิตํ โวสฺสคฺคปริณามิํ วิปุลํ มหคฺคตํ อปฺปมาณํ อพฺยาปชฺชํ ตณฺหา ปหียติ, ตณฺหาย ปหานา กมฺมํ ปหียติ, กมฺมสฺส ปหานา ทุกฺขํ ปหียติ.\csromanspacer{} อิติ โข อุทายิ ตณฺหกฺขยา กมฺมกฺขโย, กมฺมกฺขยา ทุกฺขกฺขโยติ.\csromanspacer{}\csromanspacer{}ฉฏฺฐํ.}

\csromanheader{2}{7. ตณฺหานิโรธสุตฺต}
\proseitem{208}{โย ภิกฺขเว มคฺโค ยา ปฏิปทา ตณฺหานิโรธาย สํวตฺตติ, ตํ มคฺคํ ตํ ปฏิปทํ ภาเวถ.\csromanspacer{} กตโม จ ภิกฺขเว มคฺโค กตมา จ ปฏิปทา ตณฺหานิโรธาย สํวตฺตติ, ยทิทํ สตฺต โพชฺฌงฺคา.\csromanspacer{} กตเม สตฺต, สติสมฺโพชฺฌงฺโค ฯเปฯ อุเปกฺขาสมฺโพชฺฌงฺโค.\csromanspacer{} กถํ ภาวิตา จ ภิกฺขเว สตฺต โพชฺฌงฺคา กถํ พหุลีกตา ตณฺหานิโรธาย สํวตฺตนฺติ.}

\prose{อิธ ภิกฺขเว ภิกฺขุ สติสมฺโพชฺฌงฺคํ ภาเวติ วิเวกนิสฺสิตํ วิราคนิสฺสิตํ นิโรธนิสฺสิตํ โวสฺสคฺคปริณามิํ ฯเปฯ.\csromanspacer{} อุเปกฺขาสมฺโพชฺฌงฺคํ ภาเวติ วิเวกนิสฺสิตํ วิราคนิสฺสิตํ นิโรธนิสฺสิตํ โวสฺสคฺคปริณามิํ.\csromanspacer{} เอวํ ภาวิตา โข ภิกฺขเว สตฺต โพชฺฌงฺคา เอวํ พหุลีกตา ตณฺหานิโรธาย สํวตฺตนฺตีติ.\csromanspacer{}\csromanspacer{}สตฺตมํ.}

\csromanheader{2}{8. นิพฺเพธภาคิยสุตฺต}
\proseitem{209}{นิพฺเพธภาคิยํ โว ภิกฺขเว มคฺคํ เทเสสฺสามิ ตํ สุณาถ.\csromanspacer{} กตโม จ ภิกฺขเว นิพฺเพธภาคิโย มคฺโค, ยทิทํ สตฺต โพชฺฌงฺคา.\csromanspacer{} กตเม สตฺต, สติสมฺโพชฺฌงฺโค ฯเปฯ อุเปกฺขาสมฺโพชฺฌงฺโคติ.\csromanspacer{} เอวํ วุตฺเต อายสฺมา อุทายี ภควนฺตํ เอตทโวจ “กถํ ภาวิตา นุ โข ภนฺเต สตฺต โพชฺฌงฺคา กถํ พหุลีกตา นิพฺเพธาย สํวตฺตนฺตี”ติ.}

\prose{อิธ อุทายิ ภิกฺขุ สติสมฺโพชฺฌงฺคํ ภาเวติ วิเวกนิสฺสิตํ วิราคนิสฺสิตํ นิโรธนิสฺสิตํ โวสฺสคฺคปริณามิํ วิปุลํ มหคฺคตํ อปฺปมาณํ อพฺยาปชฺชํ.\csromanspacer{} โส สติสมฺโพชฺฌงฺคํ ภาวิเตน จิตฺเตน อนิพฺพิทฺธปุพฺพํ อปฺปทาลิตปุพฺพํ โลภกฺขนฺธํ นิพฺพิชฺฌติ ปทาเลติ, อนิพฺพิทฺธปุพฺพํ อปฺปทาลิตปุพฺพํ โทสกฺขนฺธํ นิพฺพิชฺฌติ ปทาเลติ, อนิพฺพิทฺธปุพฺพํ อปฺปทาลิตปุพฺพํ โมหกฺขนฺธํ นิพฺพิชฺฌติ ปทาเลติ ฯเปฯ.\csromanspacer{} อุเปกฺขาสมฺโพชฺฌงฺคํ ภาเวติ วิเวกนิสฺสิตํ วิราคนิสฺสิตํ}

\csromanheader{2}{2. โพชฺฌงฺคสํยุตฺต}
\prose{\csromanfolio{79}นิโรธนิสฺสิตํ โวสฺสคฺคปริณามิํ วิปุลํ มหคฺคตํ อปฺปมาณํ อพฺยาปชฺชํ.\csromanspacer{} โส อุเปกฺขาสมฺโพชฺฌงฺคํ ภาวิเตน จิตฺเตน อนิพฺพิทฺธปุพฺพํ อปฺปทาลิตปุพฺพํ โลภกฺขนฺธํ นิพฺพิชฺฌติ ปทาเลติ, อนิพฺพิทฺธปุพฺพํ อปฺปทาลิตปุพฺพํ โทสกฺขนฺธํ นิพฺพิชฺฌติ ปทาเลติ, อนิพฺพิทฺธปุพฺพํ อปฺปทาลิตปุพฺพํ โมหกฺขนฺธํ นิพฺพิชฺฌติ ปทาเลติ.\csromanspacer{} เอวํ ภาวิตา โข อุทายิ สตฺต โพชฺฌงฺคา เอวํ พหุลีกตา นิพฺเพธาย สํวตฺตนฺตีติ.\csromanspacer{}\csromanspacer{}อฏฺฐมํ.}

\csromanheader{2}{9. เอกธมฺมสุตฺต}
\proseitem{210}{นาหํ ภิกฺขเว อญฺญํ เอกธมฺมมฺปิ สมนุปสฺสามิ.\csromanspacer{} โย เอวํ ภาวิโต พหุลีกโต สํโยชนียานํ ธมฺมานํ ปหานาย สํวตฺตติ, ยถยิทํ ภิกฺขเว สตฺต โพชฺฌงฺคา.\csromanspacer{} กตเม สตฺต, สติสมฺโพชฺฌงฺโค ฯเปฯ อุเปกฺขาสมฺโพชฺฌงฺโค.\csromanspacer{} กถํ ภาวิตา จ ภิกฺขเว สตฺต โพชฺฌงฺคา กถํ พหุลีกตา สํโยชนียานํ ธมฺมานํ ปหานาย สํวตฺตนฺติ.}

\prose{อิธ ภิกฺขเว ภิกฺขุ สติสมฺโพชฺฌงฺคํ ภาเวติ วิเวกนิสฺสิตํ ฯเปฯ.\csromanspacer{} อุเปกฺขาสมฺโพชฺฌงฺคํ ภาเวติ วิเวกนิสฺสิตํ วิราคนิสฺสิตํ นิโรธนิสฺสิตํ โวสฺสคฺคปริณามิํ.\csromanspacer{} เอวํ ภาวิตา โข ภิกฺขเว สตฺต โพชฺฌงฺคา เอวํ พหุลีกตา สํโยชนียานํ ธมฺมานํ ปหานาย สํวตฺตนฺติ.}

\prose{กตเม จ ภิกฺขเว สํโยชนียา ธมฺมา, จกฺขุ ภิกฺขเว สํโยชนีโย ธมฺโม, เอตฺเถเต อุปฺปชฺชนฺติ สํโยชนวินิพนฺธา อชฺโฌสานา ฯเปฯ ชิวฺหา สํโยชนียา ธมฺมา, เอตฺเถเต อุปฺปชฺชนฺติ สํโยชนวินิพนฺธา อชฺโฌสานา ฯเปฯ มโน สํโยชนีโย ธมฺโม, เอตฺเถเต อุปฺปชฺชนฺติ สํโยชนวินิพนฺธา อชฺโฌสานา.\csromanspacer{} อิเม วุจฺจนฺติ ภิกฺขเว สํโยชนียา ธมฺมาติ.\csromanspacer{}\csromanspacer{}นวมํ.}

\csromanheader{2}{10. อุทายิสุตฺต}
\proseitem{211}{เอกํ สมยํ ภควา สุมฺเภสุ วิหรติ เสตกํ นาม สุมฺภานํ นิคโม.\csromanspacer{} อถ โข อายสฺมา อุทายี เยน ภควา เตนุปสงฺกมิ ฯเปฯ เอกมนฺตํ นิสินฺโน โข อายสฺมา อุทายี ภควนฺตํ เอตทโวจ\csromandash{}}

\prose{อจฺฉริยํ ภนฺเต, อพฺภุตํ ภนฺเต, ยาว พหุกตญฺจ เม ภนฺเต ภควติ เปมญฺจ คารโว จ หิรี จ โอตฺตปฺปญฺจ.\csromanspacer{} อหํ หิ ภนฺเต ปุพฺเพ อคาริกภูโต \csromanfolio{80}สมาโน อพหุกโต อโหสิํ ธมฺเมน\footnote{ธมฺเม (?)}, อพหุกโต สํเฆน.\csromanspacer{} โส ขฺวาหํ ภควติ เปมญฺจ คารวญฺจ หิริํ จ โอตฺตปฺปญฺจ สมฺปสฺสมาโน อคารสฺมา อนคาริยํ ปพฺพชิโต, ตสฺส เม ภควา ธมฺมํ เทเสสิ “อิติ รูปํ อิติ รูปสฺส สมุทโย อิติ รูปสฺส อตฺถงฺคโม.\csromanspacer{} อิติ เวทนา ฯเปฯ.\csromanspacer{} อิติ สญฺญา.\csromanspacer{} อิติ สงฺขารา.\csromanspacer{} อิติ วิญฺญาณํ.\csromanspacer{} อิติ วิญฺญาณสฺส สมุทโย อิติ วิญฺญาณสฺส อตฺถงฺคโม”ติ.}

\prose{โส ขฺวาหํ ภนฺเต สุญฺญาคารคโต อิเมสํ ปญฺจุปาทานกฺขนฺธานํ อุกฺกุชฺชาวกุชฺชํ สมฺปริวตฺเตนฺโต “อิทํ ทุกฺขนฺ”ติ ยถาภูตํ อพฺภญฺญาสิํ, “อยํ ทุกฺขสมุทโย”ติ ยถาภูตํ อพฺภญฺญาสิํ, “อยํ ทุกฺขนิโรโธ”ติ ยถาภูตํ อพฺภญฺญาสิํ, “อยํ ทุกฺขนิโรธคามินี ปฏิปทา”ติ ยถาภูตํ อพฺภญฺญาสิํ.\csromanspacer{} ธมฺโม จ เม ภนฺเต อภิสมิโต, มคฺโค จ เม ปฏิลทฺโธ.\csromanspacer{} โย เม ภาวิโต พหุลีกโต, ตถา ตถา วิหรนฺตํ ตถตฺตาย อุปเนสฺสติ, ยถาหํ “ขีณา ชาติ, วุสิตํ พฺรหฺมจริยํ, กตํ กรณียํ, นาปรํ อิตฺถตฺตายา”ติ ปชานิสฺสามิ.}

\prose{สติสมฺโพชฺฌงฺโค เม ภนฺเต ปฏิลทฺโธ, โย เม ภาวิโต พหุลีกโต, ตถา ตถา วิหรนฺตํ ตถตฺตาย อุปเนสฺสติ, ยถาหํ “ขีณา ชาติ, วุสิตํ พฺรหฺมจริยํ, กตํ กรณียํ, นาปรํ อิตฺถตฺตายา”ติ ปชานิสฺสามิ ฯเปฯ.\csromanspacer{} อุเปกฺขาสมฺโพชฺฌงฺโค เม ภนฺเต ปฏิลทฺโธ.\csromanspacer{} โย เม ภาวิโต พหุลีกโต, ตถา ตถา วิหรนฺตํ ตถตฺตาย อุปเนสฺสติ, ยถาหํ “ขีณา ชาติ, วุสิตํ พฺรหฺมจริยํ, กตํ กรณียํ, นาปรํ อิตฺถตฺตายา”ติ ปชานิสฺสามิ.\csromanspacer{} อยํ โข เม ภนฺเต มคฺโค ปฏิลทฺโธ.\csromanspacer{} โย เม ภาวิโต พหุลีกโต, ตถา ตถา วิหรนฺตํ ตถตฺตาย อุปเนสฺสติ, ยถาหํ “ขีณา ชาติ, วุสิตํ พฺรหฺมจริยํ, กตํ กรณียํ, นาปรํ อิตฺถตฺตายา”ติ ปชานิสฺสามีติ.}

\prose{สาธุ สาธุ อุทายิ, เอโส หิ เต อุทายิ มคฺโค ปฏิลทฺโธ.\csromanspacer{} โย เต ภาวิโต พหุลีกโต, ตถา ตถา วิหรนฺตํ ตถตฺตาย อุปเนสฺสติ, ยถา ตฺวํ “ขีณา ชาติ, วุสิตํ พฺรหฺมจริยํ, กตํ กรณียํ, นาปรํ อิตฺถตฺตายา”ติ ปชานิสฺสสีติ.\csromanspacer{}\csromanspacer{}ทสมํ.\csromanspacer{} อุทายิวคฺโค ตติโย.}

\csromanheader{2}{2. โพชฺฌงฺคสํยุตฺต}
\titlehead{\textbf{ตสฺสุทฺทานํ}}
\csromangathameasure{โพธาย เทสนา ฐานา, อโยนิโส จาปริหานี. \\ ขโย นิโรโธ นิพฺเพโธ, เอกธมฺโม อุทายินาติ.}
\csromangathagroup{\csromangathastanza{\csromanfolio{81}โพธาย เทสนา ฐานา, อโยนิโส จาปริหานี. \\ ขโย นิโรโธ นิพฺเพโธ, เอกธมฺโม อุทายินาติ.}}

\chapterhead{4. นีวรณวคฺค}
\csromanheader{2}{1. ปฐมกุสลสุตฺต}
\proseitem{212}{เย เกจิ ภิกฺขเว ธมฺมา กุสลา กุสลภาคิยา กุสลปกฺขิกา, สพฺเพ เต อปฺปมาทมูลกา อปฺปมาทสโมสรณา.\csromanspacer{} อปฺปมาโท เตสํ ธมฺมานํ อคฺคมกฺขายติ.\csromanspacer{} อปฺปมตฺตสฺเสตํ ภิกฺขเว ภิกฺขุโน ปาฏิกงฺขํ สตฺต โพชฺฌงฺเค ภาเวสฺสติ สตฺต โพชฺฌงฺเค พหุลีกริสฺสติ.}

\prose{กถญฺจ ภิกฺขเว ภิกฺขุ อปฺปมตฺโต สตฺต โพชฺฌงฺเค ภาเวติ, สตฺตโพชฺฌงฺเค พหุลีกโรติ, อิธ ภิกฺขเว ภิกฺขุ สติสมฺโพชฺฌงฺคํ ภาเวติ วิเวกนิสฺสิตํ ฯเปฯ.\csromanspacer{} อุเปกฺขาสมฺโพชฺฌงฺคํ ภาเวติ วิเวกนิสฺสิตํ วิราคนิสฺสิตํ นิโรธนิสฺสิตํ โวสฺสคฺคปริณามิํ.\csromanspacer{} เอวํ โข ภิกฺขเว ภิกฺขุ อปฺปมตฺโต สตฺต โพชฺฌงฺเค ภาเวติ สตฺต โพชฺฌงฺเค พหุลีกโรตีติ.\csromanspacer{}\csromanspacer{}ปฐมํ.}

\csromanheader{2}{2. ทุติยกุสลสุตฺต}
\proseitem{213}{เย เกจิ ภิกฺขเว ธมฺมา กุสลา กุสลภาคิยา กุสลปกฺขิกา, สพฺเพ เต โยนิโสมนสิการมูลกา โยนิโสมนสิการสโมสรณา.\csromanspacer{} โยนิโสมนสิกาโร เตสํ ธมฺมานํ อคฺคมกฺขายติ.\csromanspacer{} โยนิโสมนสิการสมฺปนฺนสฺเสตํ ภิกฺขเว ภิกฺขุโน ปาฏิกงฺขํ สตฺต โพชฺฌงฺเค ภาเวสฺสติ สตฺต โพชฺฌงฺเค พหุลีกริสฺสติ.}

\prose{กถญฺจ ภิกฺขเว ภิกฺขุ โยนิโสมนสิการสมฺปนฺโน สตฺต โพชฺฌงฺเค ภาเวติ สตฺต โพชฺฌงฺเค พหุลีกโรติ, อิธ ภิกฺขเว ภิกฺขุ สติสมฺโพชฺฌงฺคํ ภาเวติ วิเวกนิสฺสิตํ ฯเปฯ.\csromanspacer{} อุเปกฺขาสมฺโพชฺฌงฺคํ ภาเวติ วิเวกนิสฺสิตํ วิราคนิสฺสิตํ นิโรธนิสฺสิตํ โวสฺสคฺคปริณามิํ.\csromanspacer{} เอวํ โข \csromanfolio{82}ภิกฺขเว ภิกฺขุ โยนิโสมนสิการสมฺปนฺโน สตฺต โพชฺฌงฺเค ภาเวติ สตฺต โพชฺฌงฺเค พหุลีกโรตีติ.\csromanspacer{}\csromanspacer{}ทุติยํ.}

\csromanheader{2}{3. อุปกฺกิเลสสุตฺต}
\proseitem{214}{ปญฺจิเม ภิกฺขเว ชาตรูปสฺส อุปกฺกิเลสา, เยหิ อุปกฺกิเลเสหิ อุปกฺกิลิฏฺฐํ ชาตรูปํ น เจว มุทุ โหติ น จ กมฺมนิยํ น จ ปภสฺสรํ ปภงฺคุ จ, น จ สมฺมา อุเปติ กมฺมาย.\csromanspacer{} กตเม ปญฺจ, อโย ภิกฺขเว ชาตรูปสฺส อุปกฺกิเลโส, เยน อุปกฺกิเลเสน อุปกฺกิลิฏฺฐํ ชาตรูปํ น เจว มุทุ โหติ น จ กมฺมนิยํ น จ ปภสฺสรํ ปภงฺคุ จ, น จ สมฺมา อุเปติ กมฺมาย.\csromanspacer{} โลหํ ภิกฺขเว ชาตรูปสฺส อุปกฺกิเลโส, เยน อุปกฺกิเลเสน อุปกฺกิลิฏฺฐํ ชาตรูปํ ฯเปฯ.\csromanspacer{} ติปุ ภิกฺขเว ชาตรูปสฺส อุปกฺกิเลโส ฯเปฯ.\csromanspacer{} สีสํ ภิกฺขเว ชาตรูปสฺส อุปกฺกิเลโส ฯเปฯ.\csromanspacer{} สชฺฌุ ภิกฺขเว ชาตรูปสฺส อุปกฺกิเลโส, เยน อุปกฺกิเลเสน อุปกฺกิลิฏฺฐํ ชาตรูปํ น เจว มุทุ โหติ น จ กมฺมนิยํ น จ ปภสฺสรํ ปภงฺคุ จ, น จ สมฺมา อุเปติ กมฺมาย.\csromanspacer{} อิเม โข ภิกฺขเว ปญฺจ ชาตรูปสฺส อุปกฺกิเลสา, เยหิ อุปกฺกิเลเสหิ อุปกฺกิลิฏฺฐํ ชาตรูปํ น เจว มุทุ โหติ น จ กมฺมนิยํ น จ ปภสฺสรํ ปภงฺคุ จ, น จ สมฺมา อุเปติ กมฺมาย.}

\prose{เอวเมว โข ภิกฺขเว ปญฺจิเม จิตฺตสฺส อุปกฺกิเลสา, เยหิ อุปกฺกิเลเสหิ อุปกฺกิลิฏฺฐํ จิตฺตํ น เจว มุทุ โหติ น จ กมฺมนิยํ น จ ปภสฺสรํ ปภงฺคุ จ, น จ สมฺมา สมาธิยติ อาสวานํ ขยาย.\csromanspacer{} กตเม ปญฺจ, กามจฺฉนฺโท ภิกฺขเว จิตฺตสฺส อุปกฺกิเลโส, เยน อุปกฺกิเลเสน อุปกฺกิลิฏฺฐํ จิตฺตํ น เจว มุทุ โหติ น จ กมฺมนิยํ น จ ปภสฺสรํ ปภงฺคุ จ, น จ สมฺมา สมาธิยติ อาสวานํ ขยาย ฯเปฯ.\csromanspacer{} อิเม โข ภิกฺขเว ปญฺจ จิตฺตสฺส อุปกฺกิเลสา, เยหิ อุปกฺกิเลเสหิ อุปกฺกิลิฏฺฐํ จิตฺตํ น เจว มุทุ โหติ น จ กมฺมนิยํ น จ ปภสฺสรํ ปภงฺคุ จ, น จ สมฺมา สมาธิยติ อาสวานํ ขยายาติ.\csromanspacer{}\csromanspacer{}ตติยํ.}

\csromanheader{2}{4. อนุปกฺกิเลสสุตฺต}
\proseitem{215}{สตฺติเม ภิกฺขเว โพชฺฌงฺคา อนาวรณา อนีวรณา เจตโส อนุปกฺกิเลสา ภาวิตา พหุลีกตา วิชฺชาวิมุตฺติผลสจฺฉิกิริยาย สํวตฺตนฺติ.\csromanspacer{} กตเม สตฺต, สติสมฺโพชฺฌงฺโค ภิกฺขเว อนาวรโณ}

\csromanheader{2}{2. โพชฺฌงฺคสํยุตฺต}
\prose{\csromanfolio{83}อนีวรโณ เจตโส อนุปกฺกิเลโส ภาวิโต พหุลีกโต วิชฺชาวิมุตฺติผลสจฺฉิกิริยาย สํวตฺตติ ฯเปฯ.\csromanspacer{} อุเปกฺขาสมฺโพชฺฌงฺโค ภิกฺขเว อนาวรโณ อนีวรโณ เจตโส อนุปกฺกิเลโส ภาวิโต พหุลีกโต วิชฺชาวิมุตฺติผลสจฺฉิกิริยาย สํวตฺตติ.\csromanspacer{} อิเม โข ภิกฺขเว สตฺต โพชฺฌงฺคา อนาวรณา อนีวรณา เจตโส อนุปกฺกิเลสา ภาวิตา พหุลีกตา วิชฺชาวิมุตฺติผลสจฺฉิกิริยาย สํวตฺตนฺตีติ.\csromanspacer{}\csromanspacer{}จตุตฺถํ.}

\csromanheader{2}{5. อโยนิโสมนสิการสุตฺต}
\proseitem{216}{อโยนิโส ภิกฺขเว มนสิกโรโต อนุปฺปนฺโน เจว กามจฺฉนฺโท อุปฺปชฺชติ, อุปฺปนฺโน จ กามจฺฉนฺโท ภิยฺโยภาวาย เวปุลฺลาย สํวตฺตติ.\csromanspacer{} อนุปฺปนฺโน เจว พฺยาปาโท อุปฺปชฺชติ, อุปฺปนฺโน จ พฺยาปาโท ภิยฺโยภาวาย เวปุลฺลาย สํวตฺตติ.\csromanspacer{} อนุปฺปนฺนญฺเจว ถินมิทฺธํ อุปฺปชฺชติ, อุปฺปนฺนญฺจ ถินมิทฺธํ ภิยฺโยภาวาย เวปุลฺลาย สํวตฺตติ.\csromanspacer{} อนุปฺปนฺนญฺเจว อุทฺธจฺจกุกฺกุจฺจํ อุปฺปชฺชติ, อุปฺปนฺนญฺจ อุทฺธจฺจกุกฺกุจฺจํ ภิยฺโยภาวาย เวปุลฺลาย สํวตฺตติ.\csromanspacer{} อนุปฺปนฺนา เจว วิจิกิจฺฉา อุปฺปชฺชติ, อุปฺปนฺนา จ วิจิกิจฺฉา ภิยฺโยภาวาย เวปุลฺลาย สํวตฺตตีติ.\csromanspacer{}\csromanspacer{}ปญฺจมํ.}

\csromanheader{2}{6. โยนิโสมนสิการสุตฺต}
\proseitem{217}{โยนิโส จ โข ภิกฺขเว มนสิกโรโต อนุปฺปนฺโน เจว สติสมฺโพชฺฌงฺโค อุปฺปชฺชติ, อุปฺปนฺโน จ สติสมฺโพชฺฌงฺโค ภาวนาปาริปูริํ คจฺฉติ ฯเปฯ อนุปฺปนฺโน เจว อุเปกฺขาสมฺโพชฺฌงฺโค อุปฺปชฺชติ, อุปฺปนฺโน จ อุเปกฺขาสมฺโพชฺฌงฺโค ภาวนาปาริปูริํ คจฺฉตีติ.\csromanspacer{}\csromanspacer{}ฉฏฺฐํ.}

\csromanheader{2}{7. พุทฺธิสุตฺต}
\proseitem{218}{สตฺติเม ภิกฺขเว โพชฺฌงฺคา ภาวิตา พหุลีกตา พุทฺธิยา อปริหานาย สํวตฺตนฺติ.\csromanspacer{} กตเม สตฺต, สติสมฺโพชฺฌงฺโค ฯเปฯ อุเปกฺขาสมฺโพชฺฌงฺโค.\csromanspacer{} อิเม โข ภิกฺขเว สตฺต โพชฺฌงฺคา ภาวิตา พหุลีกตา พุทฺธิยา อปริหานาย สํวตฺตนฺตีติ.\csromanspacer{}\csromanspacer{}สตฺตมํ.}

\csromanheader{1}{8. อาวรณนีวรณสุตฺต}
\proseitem{219}{ปญฺจิเม ภิกฺขเว อาวรณา นีวรณา เจตโส อุปกฺกิเลสา ปญฺญาย ทุพฺพลีกรณา.\csromanspacer{} กตเม ปญฺจ, กามจฺฉนฺโท ภิกฺขเว อาวรโณ \csromanfolio{84}นีวรโณ เจตโส อุปกฺกิเลโส ปญฺญาย ทุพฺพลีกรโณ, พฺยาปาโท ภิกฺขเว อาวรโณ นีวรโณ เจตโส อุปกฺกิเลโส ปญฺญาย ทุพฺพลีกรโณ, ถินมิทฺธํ ภิกฺขเว อาวรณํ นีวรณํ เจตโส อุปกฺกิเลสํ ปญฺญาย ทุพฺพลีกรณํ, อุทฺธจฺจกุกฺกุจฺจํ ภิกฺขเว อาวรณํ นีวรณํ เจตโส อุปกฺกิเลสํ ปญฺญาย ทุพฺพลีกรณํ, วิจิกิจฺฉา ภิกฺขเว อาวรณา นีวรณา เจตโส อุปกฺกิเลสา ปญฺญาย ทุพฺพลีกรณา.\csromanspacer{} อิเม โข ภิกฺขเว ปญฺจ อาวรณา นีวรณา เจตโส อุปกฺกิเลสา ปญฺญาย ทุพฺพลีกรณา.}

\prose{สตฺติเม ภิกฺขเว โพชฺฌงฺคา อนาวรณา อนีวรณา เจตโส อนุปกฺกิเลสา ภาวิตา พหุลีกตา วิชฺชาวิมุตฺติผลสจฺฉิกิริยาย สํวตฺตนฺติ.\csromanspacer{} กตเม สตฺต, สติสมฺโพชฺฌงฺโค ภิกฺขเว อนาวรโณ อนีวรโณ เจตโส อนุปกฺกิเลโส ภาวิโต พหุลีกโต วิชฺชาวิมุตฺติผลสจฺฉิกิริยาย สํวตฺตติ ฯเปฯ.\csromanspacer{} อุเปกฺขาสมฺโพชฺฌงฺโค ภิกฺขเว อนาวรโณ อนีวรโณ เจตโส อนุปกฺกิเลโส ภาวิโต พหุลีกโต วิชฺชาวิมุตฺติผลสจฺฉิ กิริยาย สํวตฺตติ.\csromanspacer{} อิเม โข ภิกฺขเว สตฺต โพชฺฌงฺคา อนาวรณา อนีวรณา เจตโส อนุปกฺกิเลสา ภาวิตา พหุลีกตา วิชฺชาวิมุตฺติผลสจฺฉิกิริยาย สํวตฺตนฺตีติ.}

\prose{ยสฺมิํ ภิกฺขเว สมเย อริยสาวโก อฏฺฐิํ กตฺวา มนสิ กตฺวา สพฺพํ เจตโส สมนฺนาหริตฺวา โอหิตโสโต ธมฺมํ สุณาติ, อิมสฺส ปญฺจ นีวรณา ตสฺมิํ สมเย น โหนฺติ, สตฺต โพชฺฌงฺคา ตสฺมิํ สมเย ภาวนาปาริปูริํ คจฺฉนฺติ.}

\prose{กตเม ปญฺจ นีวรณา ตสฺมิํ สมเย น โหนฺติ, กามจฺฉนฺทนีวรณํ ตสฺมิํ สมเย น โหติ, พฺยาปาทนีวรณํ ตสฺมิํ สมเย น โหติ, ถินมิทฺธนีวรณํ ตสฺมิํ สมเย น โหติ, อุทฺธจฺจกุกฺกุจฺจนีวรณํ ตสฺมิํ สมเย น โหติ, วิจิกิจฺฉานีวรณํ ตสฺมิํ สมเย น โหติ.\csromanspacer{} อิมสฺส ปญฺจ นีวรณา ตสฺมิํ สมเย น โหนฺติ.}

\prose{กตเม สตฺต โพชฺฌงฺคา ตสฺมิํ สมเย ภาวนาปาริปูริํ คจฺฉนฺติ, สติสมฺโพชฺฌงฺโค ตสฺมิํ สมเย ภาวนาปาริปูริํ คจฺฉติ ฯเปฯ อุเปกฺขาสมฺโพชฺฌงฺโค ตสฺมิํ สมเย ภาวนาปาริปูริํ คจฺฉติ.\csromanspacer{} อิเม สตฺต โพชฺฌงฺคา ตสฺมิํ สมเย ภาวนาปาริปูริํ คจฺฉนฺติ.\csromanspacer{} ยสฺมิํ ภิกฺขเว สมเย อริยสาวโก อฏฺฐิํ กตฺวา มนสิ กตฺวา สพฺพํ เจตโส สมนฺนาหริตฺวา}

\csromanheader{2}{2. โพชฺฌงฺคสํยุตฺต}
\prose{\csromanfolio{85}โอหิตโสโต ธมฺมํ สุณาติ, อิมสฺส ปญฺจ นีวรณา ตสฺมิํ สมเย น โหนฺติ.\csromanspacer{} อิเม สตฺต โพชฺฌงฺคา ตสฺมิํ สมเย ภาวนาปาริปูริํ คจฺฉนฺตีติ.\csromanspacer{}\csromanspacer{}อฏฺฐมํ.}

\csromanmatikaanchor{mtk.5}
\csromantocmarkref{section}{9. รุกฺขสุตฺต}{mtk.5}
\bookmark[dest={mtk.5},level=1]{9. รุกฺขสุตฺต}
\csromanheader{1}{9. รุกฺขสุตฺต}
\proseitem{220}{สนฺติ ภิกฺขเว มหารุกฺขา อณุพีชา มหากายา รุกฺขานํ อชฺฌารุหา, เยหิ รุกฺขา อชฺฌารูฬฺหา โอภคฺควิภคฺคา วิปติตา เสนฺติ.\csromanspacer{} กตเม จ เต ภิกฺขเว มหารุกฺขา อณุพีชา มหากายา รุกฺขานํ อชฺฌารุหา, เยหิ รุกฺขา อชฺฌารูฬฺหา โอภคฺควิภคฺคา วิปติตา เสนฺติ\footnote{เสนฺติ. เสยฺยถิทํ (กตฺถจิ)}, อสฺสตฺโถ นิโคฺรโธ ปิลกฺโข อุทุมฺพโร กจฺฉโก กปิตฺถโน.\csromanspacer{} อิเม โข เต ภิกฺขเว มหารุกฺขา อณุพีชา มหากายา รุกฺขานํ อชฺฌารุหา, เยหิ รุกฺขา อชฺฌารูฬฺหา โอภคฺควิภคฺคา วิปติตา เสนฺติ.\csromanspacer{} เอวเมว โข ภิกฺขเว อิเธกจฺโจ กุลปุตฺโต ยาทิสเก กาเม โอหาย อคารสฺมา อนคาริยํ ปพฺพชิโต โหติ, โส ตาทิสเกหิ กาเมหิ ตโต วา ปาปิฏฺฐตเรหิ โอภคฺควิภคฺโค วิปติโต เสติ.}

\prose{ปญฺจิเม ภิกฺขเว อาวรณา นีวรณา เจตโส อชฺฌารุหา ปญฺญาย ทุพฺพลีกรณา.\csromanspacer{} กตเม ปญฺจ, กามจฺฉนฺโท ภิกฺขเว อาวรโณ นีวรโณ เจตโส อชฺฌารุโห ปญฺญาย ทุพฺพลีกรโณ, พฺยาปาโท ภิกฺขเว อาวรโณ นีวรโณ เจตโส อชฺฌารุโห ปญฺญาย ทุพฺพลีกรโณ, ถินมิทฺธํ ภิกฺขเว อาวรณํ นีวรณํ เจตโส อชฺฌารุหํ ปญฺญาย ทุพฺพลีกรณํ, อุทฺธจฺจกุกฺกุจฺจํ ภิกฺขเว อาวรณํ นีวรณํ เจตโส อชฺฌารุหํ ปญฺญาย ทุพฺพลีกรณํ, วิจิกิจฺฉา ภิกฺขเว อาวรณา นีวรณา เจตโส อชฺฌารุหา ปญฺญาย ทุพฺพลีกรณา.\csromanspacer{} อิเม โข ภิกฺขเว ปญฺจ อาวรณา นีวรณา เจตโส อชฺฌารุหา ปญฺญาย ทุพฺพลีกรณา.}

\prose{สตฺติเม ภิกฺขเว โพชฺฌงฺคา อนาวรณา อนีวรณา เจตโส อนชฺฌารุหา ภาวิตา พหุลีกตา วิชฺชาวิมุตฺติผลสจฺฉิกิริยาย สํวตฺตนฺติ.\csromanspacer{} กตเม สตฺต, สติสมฺโพชฺฌงฺโค ภิกฺขเว อนาวรโณ อนีวรโณ เจตโส อนชฺฌารุโห ภาวิโต พหุลีกโต วิชฺชาวิมุตฺติผลสจฺฉิกิริยาย สํวตฺตติ ฯเปฯ.\csromanspacer{} อุเปกฺขาสมฺโพชฺฌงฺโค ภิกฺขเว \csromanfolio{86}อนาวรโณ อนีวรโณ เจตโส อนชฺฌารุโห ภาวิโต พหุลีกโต วิชฺชาวิมุตฺติผลสจฺฉิ กิริยาย สํวตฺตติ.\csromanspacer{} อิเม โข ภิกฺขเว สตฺต โพชฺฌงฺคา อนาวรณา อนีวรณา เจตโส อนชฺฌารุหา ภาวิตา พหุลีกตา วิชฺชาวิมุตฺติผลสจฺฉิกิริยาย สํวตฺตนฺตีติ.\csromanspacer{}\csromanspacer{}นวมํ.}

\csromanmatikaanchor{mtk.6}
\csromantocmarkref{section}{10. นีวรณสุตฺต}{mtk.6}
\bookmark[dest={mtk.6},level=1]{10. นีวรณสุตฺต}
\csromanheader{1}{10. นีวรณสุตฺต}
\proseitem{221}{ปญฺจิเม ภิกฺขเว นีวรณา อนฺธกรณา อจกฺขุกรณา อญฺญาณกรณา ปญฺญานิโรธิกา วิฆาตปกฺขิยา อนิพฺพานสํวตฺตนิกา.\csromanspacer{} กตเม ปญฺจ, กามจฺฉนฺทนีวรณํ ภิกฺขเว อนฺธกรณํ อจกฺขุกรณํ อญฺญาณกรณํ ปญฺญานิโรธิกํ วิฆาตปกฺขิยํ อนิพฺพานสํวตฺตนิกํ, พฺยาปาทนีวรณํ ภิกฺขเว ฯเปฯ ถินมิทฺธนีวรณํ ภิกฺขเว ฯเปฯ อุทฺธจฺจกุกฺกุจฺจนีวรณํ ภิกฺขเว ฯเปฯ วิจิกิจฺฉานีวรณํ ภิกฺขเว อนฺธกรณํ อจกฺขุกรณํ อญฺญาณกรณํ ปญฺญานิโรธิกํ วิฆาตปกฺขิยํ อนิพฺพานสํวตฺตนิกํ.\csromanspacer{} อิเม โข ภิกฺขเว ปญฺจ นีวรณา อนฺธกรณา อจกฺขุกรณา อญฺญาณกรณา ปญฺญานิโรธิกา วิฆาตปกฺขิยา อนิพฺพานสํวตฺตนิกา.}

\prose{สตฺติเม ภิกฺขเว โพชฺฌงฺคา จกฺขุกรณา ญาณกรณา ปญฺญาพุทฺธิยา อวิฆาตปกฺขิยา นิพฺพานสํวตฺตนิกา.\csromanspacer{} กตเม สตฺต, สติสมฺโพชฺฌงฺโค ภิกฺขเว จกฺขุกรโณ ญาณกรโณ ปญฺญาพุทฺธิโย อวิฆาตปกฺขิโย นิพฺพานสํวตฺตนิโก ฯเปฯ.\csromanspacer{} อุเปกฺขาสมฺโพชฺฌงฺโค ภิกฺขเว จกฺขุกรโณ ญาณกรโณ ปญฺญาพุทฺธิโย อวิฆาตปกฺขิโย นิพฺพานสํวตฺตนิโก.\csromanspacer{} อิเม โข ภิกฺขเว สตฺต โพชฺฌงฺคา จกฺขุกรณา ญาณกรณา ปญฺญาพุทฺธิยา อวิฆาตปกฺขิยา นิพฺพานสํวตฺตนิกาติ.\csromanspacer{}\csromanspacer{}ทสมํ.\csromanspacer{} นีวรณวคฺโค จตุตฺโถ.}

\titlehead{\textbf{ตสฺสุทฺทานํ}}
\csromangathameasure{เทฺว กุสลา กิเลสา จ, เทฺว โยนิโส จ พุทฺธิ จ. \\ อาวรณา นีวรณา รุกฺขํ, นีวรณญฺจ เต ทสาติ.}
\csromangathagroup{\csromangathastanza{เทฺว กุสลา กิเลสา จ, เทฺว โยนิโส จ พุทฺธิ จ. \\ อาวรณา นีวรณา รุกฺขํ, นีวรณญฺจ เต ทสาติ.}}

\csromanheader{2}{2. โพชฺฌงฺคสํยุตฺต}
\csromanmatikaanchor{mtk.7}
\csromantocmarknopage{chapter}{5. จกฺกวตฺติวคฺค}{mtk.7}
\bookmark[dest={mtk.7},level=0]{5. จกฺกวตฺติวคฺค}
\chapterhead{5. จกฺกวตฺติวคฺค}
\csromanmatikaanchor{mtk.8}
\csromantocmarkref{section}{1. วิธาสุตฺต}{mtk.8}
\bookmark[dest={mtk.8},level=1]{1. วิธาสุตฺต}
\csromanheader{1}{1. วิธาสุตฺต}
\csromanmatikaanchor{mtk.9}
\csromantocmarkref{subsection}{2-8. จกฺกวตฺติสุตฺตาทีนิ}{mtk.9}
\bookmark[dest={mtk.9},level=2]{2-8. จกฺกวตฺติสุตฺตาทีนิ}
\proseitem{222}{\csromanfolio{87}สาวตฺถินิทานํ.\csromanspacer{} เย หิ เกจิ ภิกฺขเว อตีตมทฺธานํ สมณา วา พฺราหฺมณา วา ติสฺโส วิธา ปชหิํสุ, สพฺเพ เต สตฺตนฺนํ โพชฺฌงฺคานํ ภาวิตตฺตา พหุลีกตตฺตา.\csromanspacer{} เย หิ เกจิ ภิกฺขเว อนาคตมทฺธานํ สมณา วา พฺราหฺมณา วา ติสฺโส วิธา ปชหิสฺสนฺติ, สพฺเพ เต สตฺตนฺนํ โพชฺฌงฺคานํ ภาวิตตฺตา พหุลีกตตฺตา.\csromanspacer{} เย หิ เกจิ ภิกฺขเว เอตรหิ สมณา วา พฺราหฺมณา วา ติสฺโส วิธา ปชหนฺติ, สพฺเพ เต สตฺตนฺนํ โพชฺฌงฺคานํ ภาวิตตฺตา พหุลีกตตฺตา.\csromanspacer{} กตเมสํ สตฺตนฺนํ โพชฺฌงฺคานํ, สติสมฺโพชฺฌงฺคสฺส ฯเปฯ อุเปกฺขาสมฺโพชฺฌงฺคสฺส.\csromanspacer{} เย หิ เกจิ ภิกฺขเว อตีตมทฺธานํ สมณา วา พฺราหฺมณา วา ติสฺโส วิธา ปชหิํสุ ฯเปฯ ปชหิสฺสนฺติ ฯเปฯ ปชหนฺติ, สพฺเพ เต อิเมสํเยว สตฺตนฺนํ โพชฺฌงฺคานํ ภาวิตตฺตา พหุลีกตตฺตาติ.\csromanspacer{}\csromanspacer{}ปฐมํ.}

\csromanheader{2}{2. จกฺกวตฺติสุตฺต}
\proseitem{223}{รญฺโญ ภิกฺขเว จกฺกวตฺติสฺส ปาตุภาวา สตฺตนฺนํ รตนานํ ปาตุภาโว โหติ.\csromanspacer{} กตเมสํ สตฺตนฺนํ, จกฺกรตนสฺส ปาตุภาโว โหติ, หตฺถิรตนสฺส ปาตุภาโว โหติ, อสฺสรตนสฺส ปาตุภาโว โหติ, มณิรตนสฺส ปาตุภาโว โหติ, อิตฺถิรตนสฺส ปาตุภาโว โหติ, คหปติรตนสฺส ปาตุภาโว โหติ, ปริณายกรตนสฺส ปาตุภาโว โหติ.\csromanspacer{} รญฺโญ ภิกฺขเว จกฺกวตฺติสฺส ปาตุภาวา อิเมสํ สตฺตนฺนํ รตนานํ ปาตุภาโว โหติ.}

\prose{ตถาคตสฺส ภิกฺขเว ปาตุภาวา อรหโต สมฺมาสมฺพุทฺธสฺส สตฺตนฺนํ โพชฺฌงฺครตนานํ ปาตุภาโว โหติ.\csromanspacer{} กตเมสํ สตฺตนฺนํ, สติสมฺโพชฺฌงฺคสฺส รตนสฺส ปาตุภาโว โหติ ฯเปฯ อุเปกฺขาสมฺโพชฺฌงฺคสฺส รตนสฺส ปาตุภาโว โหติ.\csromanspacer{} ตถาคตสฺส ภิกฺขเว ปาตุภาวา อรหโต สมฺมาสมฺพุทฺธสฺส อิเมสํ สตฺตนฺนํ โพชฺฌงฺครตนานํ ปาตุภาโว โหตีติ.\csromanspacer{}\csromanspacer{}ทุติยํ.}

\csromanheader{2}{3. มารสุตฺต}
\proseitem{224}{\csromanfolio{88}มารเสนปฺปมทฺทนํ โว ภิกฺขเว มคฺคํ เทเสสฺสามิ, ตํ สุณาถ.\csromanspacer{} กตโม จ ภิกฺขเว มารเสนปฺปมทฺทโน มคฺโค, ยทิทํ สตฺต โพชฺฌงฺคา.\csromanspacer{} กตเม สตฺต, สติสมฺโพชฺฌงฺโค ฯเปฯ อุเปกฺขาสมฺโพชฺฌงฺโค.\csromanspacer{} อยํ โข ภิกฺขเว มารเสนปฺปมทฺทโน มคฺโคติ.\csromanspacer{}\csromanspacer{}ตติยํ.}

\csromanheader{2}{4. ทุปฺปญฺญสุตฺต}
\proseitem{225}{อถ โข อญฺญตโร ภิกฺขุ เยน ภควา เตนุปสงฺกมิ ฯเปฯ เอกมนฺตํ นิสินฺโน โข โส ภิกฺขุ ภควนฺตํ เอตทโวจ “ทุปฺปญฺโญ เอฬมูโค ทุปฺปญฺโญ เอฬมูโค”ติ ภนฺเต วุจฺจติ, กิตฺตาวตา นุ โข ภนฺเต “ทุปฺปญฺโญ เอฬมูโค”ติ วุจฺจตีติ.\csromanspacer{} สตฺตนฺนํ โข ภิกฺขุ โพชฺฌงฺคานํ อภาวิตตฺตา อพหุลีกตตฺตา “ทุปฺปญฺโญ เอฬมูโค”ติ วุจฺจติ.\csromanspacer{} กตเมสํ สตฺตนฺนํ, สติสมฺโพชฺฌงฺคสฺส ฯเปฯ อุเปกฺขาสมฺโพชฺฌงฺคสฺส.\csromanspacer{} อิเมสํ โข ภิกฺขุ สตฺตนฺนํ โพชฺฌงฺคานํ อภาวิตตฺตา อพหุลีกตตฺตา “ทุปฺปญฺโญ เอฬมูโค”ติ วุจฺจตีติ.\csromanspacer{}\csromanspacer{}จตุตฺถํ.}

\csromanheader{2}{5. ปญฺญวนฺตสุตฺต}
\proseitem{226}{“ปญฺญวา อเนฬมูโค ปญฺญวา อเนฬมูโค”ติ ภนฺเต วุจฺจติ, กิตฺตาวตา นุ โข ภนฺเต “ปญฺญวา อเนฬมูโค”ติ วุจฺจตีติ.\csromanspacer{} สตฺตนฺนํ โข ภิกฺขุ โพชฺฌงฺคานํ ภาวิตตฺตา พหุลีกตตฺตา “ปญฺญวา อเนฬมูโค”ติ วุจฺจติ.\csromanspacer{} กตเมสํ สตฺตนฺนํ, สติสมฺโพชฺฌงฺคสฺส ฯเปฯ อุเปกฺขาสมฺโพชฺฌงฺคสฺส.\csromanspacer{} อิเมสํ โข ภิกฺขุ สตฺตนฺนํ โพชฺฌงฺคานํ ภาวิตตฺตา พหุลีกตตฺตา “ปญฺญวา อเนฬมูโค”ติ วุจฺจตีติ.\csromanspacer{}\csromanspacer{}ปญฺจมํ.}

\csromanheader{2}{6. ทลิทฺทสุตฺต}
\proseitem{227}{“ทลิทฺโท ทลิทฺโท”ติ ภนฺเต วุจฺจติ, กิตฺตาวตา นุ โข ภนฺเต “ทลิทฺโท”ติ วุจฺจตีติ.\csromanspacer{} สตฺตนฺนํ โข ภิกฺขุ โพชฺฌงฺคานํ อภาวิตตฺตา อพหุลีกตตฺตา “ทลิทฺโท”ติ วุจฺจติ.\csromanspacer{} กตเมสํ สตฺตนฺนํ, สติสมฺโพชฺฌงฺคสฺส ฯเปฯ อุเปกฺขาสมฺโพชฺฌงฺคสฺส.\csromanspacer{} อิเมสํ โข ภิกฺขุ สตฺตนฺนํ โพชฺฌงฺคานํ อภาวิตตฺตา อพหุลีกตตฺตา “ทลิทฺโท”ติ วุจฺจตีติ.\csromanspacer{}\csromanspacer{}ฉฏฺฐํ.}

\csromanheader{2}{2. โพชฺฌงฺคสํยุตฺต}
\csromanheader{2}{7. อทลิทฺทสุตฺต}
\proseitem{228}{\csromanfolio{89}“อทลิทฺโท อทลิทฺโท”ติ ภนฺเต วุจฺจติ, กิตฺตาวตา นุ โข ภนฺเต “อทลิทฺโท”ติ วุจฺจตีติ.\csromanspacer{} สตฺตนฺนํ โข ภิกฺขุ โพชฺฌงฺคานํ ภาวิตตฺตา พหุลีกตตฺตา “อทลิทฺโท”ติ วุจฺจติ.\csromanspacer{} กตเมสํ สตฺตนฺนํ, สติสมฺโพชฺฌงฺคสฺส ฯเปฯ อุเปกฺขาสมฺโพชฺฌงฺคสฺส.\csromanspacer{} อิเมสํ โข ภิกฺขุ สตฺตนฺนํ โพชฺฌงฺคานํ ภาวิตตฺตา พหุลีกตตฺตา “อทลิทฺโท”ติ วุจฺจตีติ.\csromanspacer{}\csromanspacer{}สตฺตมํ.}

\csromanheader{2}{8. อาทิจฺจสุตฺต}
\proseitem{229}{อาทิจฺจสฺส ภิกฺขเว อุทยโต เอตํ ปุพฺพงฺคมํ เอตํ ปุพฺพนิมิตฺตํ, ยทิทํ อรุณุคฺคํ.\csromanspacer{} เอวเมว โข ภิกฺขเว ภิกฺขุโน สตฺตนฺนํ โพชฺฌงฺคานํ อุปฺปาทาย เอตํ ปุพฺพงฺคมํ เอตํ ปุพฺพนิมิตฺตํ, ยทิทํ กลฺยาณมิตฺตตา.\csromanspacer{} กลฺยาณมิตฺตสฺเสตํ ภิกฺขเว ภิกฺขุโน ปาฏิกงฺขํ สตฺต โพชฺฌงฺเค ภาเวสฺสติ สตฺต โพชฺฌงฺเค พหุลีกริสฺสติ.\csromanspacer{} กถญฺจ ภิกฺขเว ภิกฺขุ กลฺยาณมิตฺโต สตฺต โพชฺฌงฺเค ภาเวติ สตฺต โพชฺฌงฺเค พหุลีกโรติ, อิธ ภิกฺขเว ภิกฺขุ สติสมฺโพชฺฌงฺคํ ภาเวติ วิเวกนิสฺสิตํ ฯเปฯ อุเปกฺขาสมฺโพชฺฌงฺคํ ภาเวติ วิเวกนิสฺสิตํ วิราคนิสฺสิตํ นิโรธนิสฺสิตํ โวสฺสคฺคปริณามิํ.\csromanspacer{} เอวํ โข ภิกฺขเว ภิกฺขุ กลฺยาณมิตฺโต สตฺต โพชฺฌงฺเค ภาเวติ สตฺต โพชฺฌงฺเค พหุลีกโรตีติ.\csromanspacer{}\csromanspacer{}อฏฺฐมํ.}

\csromanmatikaanchor{mtk.10}
\csromantocmarkref{subsection}{9. อชฺฌตฺติกงฺคสุตฺต}{mtk.10}
\bookmark[dest={mtk.10},level=2]{9. อชฺฌตฺติกงฺคสุตฺต}
\csromanheader{2}{9. อชฺฌตฺติกงฺคสุตฺต}
\proseitem{230}{“อชฺฌตฺติกํ ภิกฺขเว องฺคนฺ”ติ กริตฺวา นาญฺญํ เอกงฺคมฺปิ สมนุปสฺสามิ สตฺตนฺนํ โพชฺฌงฺคานํ อุปฺปาทาย, ยถยิทํ ภิกฺขเว โยนิโสมนสิกาโร.\csromanspacer{} โยนิโสมนสิการสมฺปนฺนสฺเสตํ ภิกฺขเว ภิกฺขุโน ปาฏิกงฺขํ สตฺต โพชฺฌงฺเค ภาเวสฺสติ สตฺต โพชฺฌงฺเค พหุลีกริสฺสติ.\csromanspacer{} กถญฺจ ภิกฺขเว ภิกฺขุ โยนิโสมนสิการสมฺปนฺโน สตฺต โพชฺฌงฺเค ภาเวติ สตฺต โพชฺฌงฺเค พหุลีกโรติ, อิธ ภิกฺขเว ภิกฺขุ สติสมฺโพชฺฌงฺคํ ภาเวติ วิเวกนิสฺสิตํ ฯเปฯ อุเปกฺขาสมฺโพชฺฌงฺคํ ภาเวติ วิเวกนิสฺสิตํ วิราคนิสฺสิตํ นิโรธนิสฺสิตํ โวสฺสคฺคปริณามิํ.\csromanspacer{} เอวํ โข ภิกฺขเว ภิกฺขุ โยนิโสมนสิการสมฺปนฺโน สตฺต โพชฺฌงฺเค ภาเวติ สตฺต โพชฺฌงฺเค พหุลีกโรตีติ.\csromanspacer{}\csromanspacer{}นวมํ.}

\csromanmatikaanchor{mtk.11}
\csromantocmarkref{subsection}{10. พาหิรงฺคสุตฺต}{mtk.11}
\bookmark[dest={mtk.11},level=2]{10. พาหิรงฺคสุตฺต}
\csromanheader{2}{10. พาหิรงฺคสุตฺต}
\proseitem{231}{\csromanfolio{90}“พาหิรํ ภิกฺขเว องฺคนฺ”ติ กริตฺวา นาญฺญํ เอกงฺคมฺปิ สมนุปสฺสามิ สตฺตนฺนํ โพชฺฌงฺคานํ อุปฺปาทาย, ยถยิทํ ภิกฺขเว กลฺยาณมิตฺตตา.\csromanspacer{} กลฺยาณมิตฺตสฺเสตํ ภิกฺขเว ภิกฺขุโน ปาฏิกงฺขํ สตฺต โพชฺฌงฺเค ภาเวสฺสติ สตฺต โพชฺฌงฺเค พหุลีกริสฺสติ.\csromanspacer{} กถญฺจ ภิกฺขเว ภิกฺขุ กลฺยาณมิตฺโต สตฺต โพชฺฌงฺเค ภาเวติ สตฺต โพชฺฌงฺเค พหุลีกโรติ, อิธ ภิกฺขเว ภิกฺขุ สติสมฺโพชฺฌงฺคํ ภาเวติ วิเวกนิสฺสิตํ ฯเปฯ อุเปกฺขาสมฺโพชฺฌงฺคํ ภาเวติ วิเวกนิสฺสิตํ วิราคนิสฺสิตํ นิโรธนิสฺสิตํ โวสฺสคฺคปริณามิํ.\csromanspacer{} เอวํ โข ภิกฺขเว ภิกฺขุ กลฺยาณมิตฺโต สตฺต โพชฺฌงฺเค ภาเวติ สตฺต โพชฺฌงฺเค พหุลีกโรตีติ.\csromanspacer{} จกฺกวตฺติวคฺโค ปญฺจโม.}

\titlehead{\textbf{ตสฺสุทฺทานํ}}
\csromangathameasure{วิธา จกฺกวตฺติ มาโร, ทุปฺปญฺโญ ปญฺญเวน จ. \\ ทลิทฺโท อทลิทฺโท จ, อาทิจฺจงฺเคน เต ทสาติ.}
\csromangathagroup{\csromangathastanza{วิธา จกฺกวตฺติ มาโร, ทุปฺปญฺโญ ปญฺญเวน จ. \\ ทลิทฺโท อทลิทฺโท จ, อาทิจฺจงฺเคน เต ทสาติ.}}

\csromanmatikaanchor{mtk.12}
\csromantocmarknopage{subsection}{6. สากจฺฉวคฺค}{mtk.12}
\bookmark[dest={mtk.12},level=2]{6. สากจฺฉวคฺค}
\csromanheader{2}{6. สากจฺฉวคฺค}
\csromanmatikaanchor{mtk.13}
\csromantocmarkref{paragraph}{1. อาหารสุตฺต}{mtk.13}
\bookmark[dest={mtk.13},level=4]{1. อาหารสุตฺต}
\csromanheader{4}{1. อาหารสุตฺต}
\proseitem{232}{สาวตฺถินิทานํ.\csromanspacer{} ปญฺจนฺนญฺจ ภิกฺขเว นีวรณานํ สตฺตนฺนญฺจ โพชฺฌงฺคานํ อาหารญฺจ อนาหารญฺจ เทเสสฺสามิ, ตํ สุณาถ.\csromanspacer{} โก จ ภิกฺขเว อาหาโร อนุปฺปนฺนสฺส วา กามจฺฉนฺทสฺส อุปฺปาทาย อุปฺปนฺนสฺส วา กามจฺฉนฺทสฺส ภิยฺโยภาวาย เวปุลฺลาย, อตฺถิ ภิกฺขเว สุภนิมิตฺตํ, ตตฺถ อโยนิโสมนสิการพหุลีกาโร, อยมาหาโร อนุปฺปนฺนสฺส วา กามจฺฉนฺทสฺส อุปฺปาทาย อุปฺปนฺนสฺส วา กามจฺฉนฺทสฺส ภิยฺโยภาวาย เวปุลฺลาย.}

\prose{โก จ ภิกฺขเว อาหาโร อนุปฺปนฺนสฺส วา พฺยาปาทสฺส อุปฺปาทาย อุปฺปนฺนสฺส วา พฺยาปาทสฺส ภิยฺโยภาวาย เวปุลฺลาย, อตฺถิ ภิกฺขเว}

\csromanheader{2}{2. โพชฺฌงฺคสํยุตฺต}
\prose{\csromanfolio{91}ปฏิฆนิมิตฺตํ, ตตฺถ อโยนิโสมนสิการพหุลีกาโร, อยมาหาโร อนุปฺปนฺนสฺส วา พฺยาปาทสฺส อุปฺปาทาย อุปฺปนฺนสฺส วา พฺยาปาทสฺส ภิยฺโยภาวาย เวปุลฺลาย.}

\prose{โก จ ภิกฺขเว อาหาโร อนุปฺปนฺนสฺส วา ถินมิทฺธสฺส อุปฺปาทาย อุปฺปนฺนสฺส วา ถินมิทฺธสฺส ภิยฺโยภาวาย เวปุลฺลาย, อตฺถิ ภิกฺขเว อรติ ตนฺทิ วิชมฺภิตา ภตฺตสมฺมโท เจตโส จ ลีนตฺตํ, ตตฺถ อโยนิโสมนสิการพหุลีกาโร, อยมาหาโร อนุปฺปนฺนสฺส วา ถินมิทฺธสฺส อุปฺปาทาย อุปฺปนฺนสฺส วา ถินมิทฺธสฺส ภิยฺโยภาวาย เวปุลฺลาย.}

\prose{โก จ ภิกฺขเว อาหาโร อนุปฺปนฺนสฺส วา อุทฺธจฺจกุกฺกุจฺจสฺส อุปฺปาทาย อุปฺปนฺนสฺส วา อุทฺธจฺจกุกฺกุจฺจสฺส ภิยฺโยภาวาย เวปุลฺลาย, อตฺถิ ภิกฺขเว เจตโส อวูปสโม, ตตฺถ อโยนิโสมนสิการพหุลีกาโร, อยมาหาโร อนุปฺปนฺนสฺส วา อุทฺธจฺจกุกฺกุจฺจสฺส อุปฺปาทาย อุปฺปนฺนสฺส วา อุทฺธจฺจกุกฺกุจฺจสฺส ภิยฺโยภาวาย เวปุลฺลาย.}

\prose{โก จ ภิกฺขเว อาหาโร อนุปฺปนฺนาย วา วิจิกิจฺฉาย อุปฺปาทาย อุปฺปนฺนาย วา วิจิกิจฺฉาย ภิยฺโยภาวาย เวปุลฺลาย, อตฺถิ ภิกฺขเว วิจิกิจฺฉาฏฺฐานียา ธมฺมา, ตตฺถ อโยนิโสมนสิการพหุลีกาโร, อยมาหาโร อนุปฺปนฺนาย วา วิจิกิจฺฉาย อุปฺปาทาย อุปฺปนฺนาย วา วิจิกิจฺฉาย ภิยฺโยภาวาย เวปุลฺลาย.}

\prose{โก จ ภิกฺขเว อาหาโร อนุปฺปนฺนสฺส วา สติสมฺโพชฺฌงฺคสฺส อุปฺปาทาย อุปฺปนฺนสฺส วา สติสมฺโพชฺฌงฺคสฺส ภาวนาย ปาริปูริยา, อตฺถิ ภิกฺขเว สติสมฺโพชฺฌงฺคฏฺฐานียา ธมฺมา, ตตฺถ โยนิโสมนสิการพหุลีกาโร, อยมาหาโร อนุปฺปนฺนสฺส วา สติสมฺโพชฺฌงฺคสฺส อุปฺปาทาย อุปฺปนฺนสฺส วา สติสมฺโพชฺฌงฺคสฺส ภาวนาย ปาริปูริยา.}

\prose{โก จ ภิกฺขเว อาหาโร อนุปฺปนฺนสฺส วา ธมฺมวิจยสมฺโพชฺฌงฺคสฺส อุปฺปาทาย อุปฺปนฺนสฺส วา ธมฺมวิจยสมฺโพชฺฌงฺคสฺส ภาวนาย ปาริปูริยา, อตฺถิ ภิกฺขเว กุสลากุสลา ธมฺมา สาวชฺชานวชฺชา ธมฺมา หีนปณีตา ธมฺมา กณฺหสุกฺกสปฺปฏิภาคา ธมฺมา, ตตฺถ โยนิโสมนสิการพหุลีกาโร, อยมาหาโร อนุปฺปนฺนสฺส วา ธมฺมวิจยสมฺโพชฺฌงฺคสฺส อุปฺปาทาย อุปฺปนฺนสฺส วา ธมฺมวิจยสมฺโพชฺฌงฺคสฺส ภาวนาย ปาริปูริยา.}

\prose{\csromanfolio{92}โก จ ภิกฺขเว อาหาโร อนุปฺปนฺนสฺส วา วีริยสมฺโพชฺฌงฺคสฺส อุปฺปาทาย อุปฺปนฺนสฺส วา วีริยสมฺโพชฺฌงฺคสฺส ภาวนาย ปาริปูริยา, อตฺถิ ภิกฺขเว อารมฺภธาตุ นิกฺกมธาตุ ปรกฺกมธาตุ, ตตฺถ โยนิโสมนสิการพหุลีกาโร, อยมาหาโร อนุปฺปนฺนสฺส วา วีริยสมฺโพชฺฌงฺคสฺส อุปฺปาทาย อุปฺปนฺนสฺส วา วีริยสมฺโพชฺฌงฺคสฺส ภาวนาย ปาริปูริยา.}

\prose{โก จ ภิกฺขเว อาหาโร อนุปฺปนฺนสฺส วา ปีติสมฺโพชฺฌงฺคสฺส อุปฺปาทาย อุปฺปนฺนสฺส วา ปีติสมฺโพชฺฌงฺคสฺส ภาวนาย ปาริปูริยา, อตฺถิ ภิกฺขเว ปีติสมฺโพชฺฌงฺคฏฺฐานียา ธมฺมา, ตตฺถ โยนิโสมนสิการพหุลีกาโร, อยมาหาโร อนุปฺปนฺนสฺส วา ปีติสมฺโพชฺฌงฺคสฺส อุปฺปาทาย อุปฺปนฺนสฺส วา ปีติสมฺโพชฺฌงฺคสฺส ภาวนาย ปาริปูริยา.}

\prose{โก จ ภิกฺขเว อาหาโร อนุปฺปนฺนสฺส วา ปสฺสทฺธิสมฺโพชฺฌงฺคสฺส อุปฺปาทาย อุปฺปนฺนสฺส วา ปสฺสทฺธิสมฺโพชฺฌงฺคสฺส ภาวนาย ปาริปูริยา, อตฺถิ ภิกฺขเว กายปฺปสฺสทฺธิ จิตฺตปฺปสฺสทฺธิ, ตตฺถ โยนิโสมนสิการพหุลีกาโร, อยมาหาโร อนุปฺปนฺนสฺส วา ปสฺสทฺธิสมฺโพชฺฌงฺคสฺส อุปฺปาทาย อุปฺปนฺนสฺส วา ปสฺสทฺธิสมฺโพชฺฌงฺคสฺส ภาวนาย ปาริปูริยา.}

\prose{โก จ ภิกฺขเว อาหาโร อนุปฺปนฺนสฺส วา สมาธิสมฺโพชฺฌงฺคสฺส อุปฺปาทาย อุปฺปนฺนสฺส วา สมาธิสมฺโพชฺฌงฺคสฺส ภาวนาย ปาริปูริยา, อตฺถิ ภิกฺขเว สมถนิมิตฺตํ อพฺยคฺคนิมิตฺตํ, ตตฺถ โยนิโสมนสิการพหุลีกาโร, อยมาหาโร อนุปฺปนฺนสฺส วา สมาธิสมฺโพชฺฌงฺคสฺส อุปฺปาทาย อุปฺปนฺนสฺส วา สมาธิสมฺโพชฺฌงฺคสฺส ภาวนาย ปาริปูริยา.}

\prose{โก จ ภิกฺขเว อาหาโร อนุปฺปนฺนสฺส วา อุเปกฺขาสมฺโพชฺฌงฺคสฺส อุปฺปาทาย อุปฺปนฺนสฺส วา อุเปกฺขาสมฺโพชฺฌงฺคสฺส ภาวนาย ปาริปูริยา, อตฺถิ ภิกฺขเว อุเปกฺขาสมฺโพชฺฌงฺคฏฺฐานียา ธมฺมา, ตตฺถ โยนิโสมนสิการพหุลีกาโร, อยมาหาโร อนุปฺปนฺนสฺส วา อุเปกฺขาสมฺโพชฺฌงฺคสฺส อุปฺปาทาย อุปฺปนฺนสฺส วา อุเปกฺขาสมฺโพชฺฌงฺคสฺส ภาวนาย ปาริปูริยา.}

\prose{โก จ ภิกฺขเว อนาหาโร อนุปฺปนฺนสฺส วา กามจฺฉนฺทสฺส อุปฺปาทาย อุปฺปนฺนสฺส วา กามจฺฉนฺทสฺส ภิยฺโยภาวาย เวปุลฺลาย, อตฺถิ ภิกฺขเว}

\csromanheader{2}{2. โพชฺฌงฺคสํยุตฺต}
\prose{\csromanfolio{93}อสุภนิมิตฺตํ, ตตฺถ โยนิโสมนสิการพหุลีกาโร, อยมนาหาโร อนุปฺปนฺนสฺส วา กามจฺฉนฺทสฺส อุปฺปาทาย อุปฺปนฺนสฺส วา กามจฺฉนฺทสฺส ภิยฺโยภาวาย เวปุลฺลาย.}

\prose{โก จ ภิกฺขเว อนาหาโร อนุปฺปนฺนสฺส วา พฺยาปาทสฺส อุปฺปาทาย อุปฺปนฺนสฺส วา พฺยาปาทสฺส ภิยฺโยภาวาย เวปุลฺลาย, อตฺถิ ภิกฺขเว เมตฺตาเจโตวิมุตฺติ, ตตฺถ โยนิโสมนสิการพหุลีกาโร, อยมนาหาโร อนุปฺปนฺนสฺส วา พฺยาปาทสฺส อุปฺปาทาย อุปฺปนฺนสฺส วา พฺยาปาทสฺส ภิยฺโยภาวาย เวปุลฺลาย.}

\prose{โก จ ภิกฺขเว อนาหาโร อนุปฺปนฺนสฺส วา ถินมิทฺธสฺส อุปฺปาทาย อุปฺปนฺนสฺส วา ถินมิทฺธสฺส ภิยฺโยภาวาย เวปุลฺลาย, อตฺถิ ภิกฺขเว อารมฺภธาตุ นิกฺกมธาตุ ปรกฺกมธาตุ, ตตฺถ โยนิโสมนสิการพหุลีกาโร, อยมนาหาโร อนุปฺปนฺนสฺส วา ถินมิทฺธสฺส อุปฺปาทาย อุปฺปนฺนสฺส วา ถินมิทฺธสฺส ภิยฺโยภาวาย เวปุลฺลาย.}

\prose{โก จ ภิกฺขเว อนาหาโร อนุปฺปนฺนสฺส วา อุทฺธจฺจกุกฺกุจฺจสฺส อุปฺปาทาย อุปฺปนฺนสฺส วา อุทฺธจฺจกุกฺกุจฺจสฺส ภิยฺโยภาวาย เวปุลฺลาย, อตฺถิ ภิกฺขเว เจตโส วูปสโม, ตตฺถ โยนิโสมนสิการพหุลีกาโร, อยมนาหาโร อนุปฺปนฺนสฺส วา อุทฺธจฺจกุกฺกุจฺจสฺส อุปฺปาทาย อุปฺปนฺนสฺส วา อุทฺธจฺจกุกฺกุจฺจสฺส ภิยฺโยภาวาย เวปุลฺลาย.}

\prose{โก จ ภิกฺขเว อนาหาโร อนุปฺปนฺนาย วา วิจิกิจฺฉาย อุปฺปาทาย อุปฺปนฺนาย วา วิจิกิจฺฉาย ภิยฺโยภาวาย เวปุลฺลาย, อตฺถิ ภิกฺขเว กุสลากุสลา ธมฺมา สาวชฺชานวชฺชา ธมฺมา หีนปณีตา ธมฺมา กณฺหสุกฺกสปฺปฏิภาคา ธมฺมา, ตตฺถ โยนิโสมนสิการพหุลีกาโร, อยมนาหาโร อนุปฺปนฺนาย วา วิจิกิจฺฉาย อุปฺปาทาย อุปฺปนฺนาย วา วิจิกิจฺฉาย ภิยฺโยภาวาย เวปุลฺลาย.}

\prose{โก จ ภิกฺขเว อนาหาโร อนุปฺปนฺนสฺส วา สติสมฺโพชฺฌงฺคสฺส อุปฺปาทาย อุปฺปนฺนสฺส วา สติสมฺโพชฺฌงฺคสฺส ภาวนาย ปาริปูริยา, อตฺถิ ภิกฺขเว สติสมฺโพชฺฌงฺคฏฺฐานียา ธมฺมา, ตตฺถ อมนสิการพหุลีกาโร, อยมนาหาโร อนุปฺปนฺนสฺส วา สติสมฺโพชฺฌงฺคสฺส อุปฺปาทาย อุปฺปนฺนสฺส วา สติสมฺโพชฺฌงฺคสฺส ภาวนาย ปาริปูริยา.}

\prose{\csromanfolio{94}โก จ ภิกฺขเว อนาหาโร อนุปฺปนฺนสฺส วา ธมฺมวิจยสมฺโพชฺฌงฺคสฺส อุปฺปาทาย อุปฺปนฺนสฺส วา ธมฺมวิจยสมฺโพชฺฌงฺคสฺส ภาวนาย ปาริปูริยา, อตฺถิ ภิกฺขเว กุสลากุสลา ธมฺมา สาวชฺชานวชฺชา ธมฺมา หีนปณีตา ธมฺมา กณฺหสุกฺกสปฺปฏิภาคา ธมฺมา, ตตฺถ อมนสิการพหุลีกาโร, อยมนาหาโร อนุปฺปนฺนสฺส วา ธมฺมวิจยสมฺโพชฺฌงฺคสฺส อุปฺปาทาย อุปฺปนฺนสฺส วา ธมฺมวิจยสมฺโพชฺฌงฺคสฺส ภาวนาย ปาริปูริยา.}

\prose{โก จ ภิกฺขเว อนาหาโร อนุปฺปนฺนสฺส วา วีริยสมฺโพชฺฌงฺคสฺส อุปฺปาทาย อุปฺปนฺนสฺส วา วีริยสมฺโพชฺฌงฺคสฺส ภาวนาย ปาริปูริยา, อตฺถิ ภิกฺขเว อารมฺภธาตุ นิกฺกมธาตุ ปรกฺกมธาตุ, ตตฺถ อมนสิการพหุลีกาโร, อยมนาหาโร อนุปฺปนฺนสฺส วา วีริยสมฺโพชฺฌงฺคสฺส อุปฺปาทาย อุปฺปนฺนสฺส วา วีริยสมฺโพชฺฌงฺคสฺส ภาวนาย ปาริปูริยา.}

\prose{โก จ ภิกฺขเว อนาหาโร อนุปฺปนฺนสฺส วา ปีติสมฺโพชฺฌงฺคสฺส อุปฺปาทาย อุปฺปนฺนสฺส วา ปีติสมฺโพชฺฌงฺคสฺส ภาวนาย ปาริปูริยา, อตฺถิ ภิกฺขเว ปีติสมฺโพชฺฌงฺคฏฺฐานียา ธมฺมา, ตตฺถ อมนสิการพหุลีกาโร, อยมนาหาโร อนุปฺปนฺนสฺส วา ปีติสมฺโพชฺฌงฺคสฺส อุปฺปาทาย อุปฺปนฺนสฺส วา ปีติสมฺโพชฺฌงฺคสฺส ภาวนาย ปาริปูริยา.}

\prose{โก จ ภิกฺขเว อนาหาโร อนุปฺปนฺนสฺส วา ปสฺสทฺธิสมฺโพชฺฌงฺคสฺส อุปฺปาทาย อุปฺปนฺนสฺส วา ปสฺสทฺธิสมฺโพชฺฌงฺคสฺส ภาวนาย ปาริปูริยา, อตฺถิ ภิกฺขเว กายปฺปสฺสทฺธิ จิตฺตปฺปสฺสทฺธิ, ตตฺถ อมนสิการพหุลีกาโร อยมนาหาโร อนุปฺปนฺนสฺส วา ปสฺสทฺธิสมฺโพชฺฌงฺคสฺส อุปฺปาทาย อุปฺปนฺนสฺส วา ปสฺสทฺธิสมฺโพชฺฌงฺคสฺส ภาวนาย ปาริปูริยา.}

\prose{โก จ ภิกฺขเว อนาหาโร อนุปฺปนฺนสฺส วา สมาธิสมฺโพชฺฌงฺคสฺส อุปฺปาทาย อุปฺปนฺนสฺส วา สมาธิสมฺโพชฺฌงฺคสฺส ภาวนาย ปาริปูริยา, อตฺถิ ภิกฺขเว สมถนิมิตฺตํ อพฺยคฺคนิมิตฺตํ, ตตฺถ อมนสิการพหุลีกาโร อยมนาหาโร อนุปฺปนฺนสฺส วา สมาธิสมฺโพชฺฌงฺคสฺส อุปฺปาทาย อุปฺปนฺนสฺส วา สมาธิสมฺโพงฺคสฺส ภาวนาย ปาริปูริยา.}

\prose{โก จ ภิกฺขเว อนาหาโร อนุปฺปนฺนสฺส วา อุเปกฺขาสมฺโพชฺฌงฺคสฺส อุปฺปาทาย อุปฺปนฺนสฺส วา อุเปกฺขาสมฺโพชฺฌงฺคสฺส ภาวนาย ปาริปูริยา, อตฺถิ ภิกฺขเว อุเปกฺขาสมฺโพชฺฌงฺคฏฺฐานียา ธมฺมา, ตตฺถ อมนสิการพหุลีกาโร, อยมนาหาโร อนุปฺปนฺนสฺส วา อุเปกฺขาสมฺโพชฺฌงฺคสฺส}

\csromanheader{2}{2. โพชฺฌงฺคสํยุตฺต}
\prose{\csromanfolio{95}อุปฺปาทาย อุปฺปนฺนสฺส วา อุเปกฺขาสมฺโพชฺฌงฺคสฺส ภาวนาย ปาริปูริยาติ.\csromanspacer{}\csromanspacer{}ปฐมํ.}

\csromanmatikaanchor{mtk.14}
\csromantocmarkref{paragraph}{2. ปริยายสุตฺต}{mtk.14}
\bookmark[dest={mtk.14},level=4]{2. ปริยายสุตฺต}
\csromanheader{4}{2. ปริยายสุตฺต}
\proseitem{233}{อถ โข สมฺพหุลา ภิกฺขู ปุพฺพณฺหสมยํ นิวาเสตฺวา ปตฺตจีวรมาทาย สวตฺถิํ ปิณฺฑาย ปวิสิํสุ.\csromanspacer{} อถ โข เตสํ ภิกฺขูนํ เอตทโหสิ “อติปฺปโค โข ตาว สาวตฺถิยํ ปิณฺฑาย จริตุํ, ยํนูน มยํ เยน อญฺญติตฺถิยานํ ปริพฺพาชกานํ อาราโม เตนุปสงฺกเมยฺยามา”ติ.}

\prose{อถ โข เต ภิกฺขู เยน อญฺญติตฺถิยานํ ปริพฺพาชกานํ อาราโม เตนุปสงฺกมิํสุ, อุปสงฺกมิตฺวา เตหิ อญฺญติตฺถิเยหิ ปริพฺพาชเกหิ สทฺธิํ สมฺโมทิํสุ, สมฺโมทนียํ กถํ สารณียํ วีติสาเรตฺวา เอกมนฺตํ นิสีทิํสุ, เอกมนฺตํ นิสินฺเน โข เต ภิกฺขู อญฺญติตฺถิยา ปริพฺพาชกา เอตทโวจุํ\csromandash{}}

\prose{สมโณ อาวุโส โคตโม สาวกานํ เอวํ ธมฺมํ เทเสติ “เอถ ตุเมฺห ภิกฺขเว ปญฺจ นีวรเณ ปหาย เจตโส อุปกฺกิเลเส ปญฺญาย ทุพฺพลีกรเณ สตฺต โพชฺฌงฺเค ยถาภูตํ ภาเวถา”ติ.\csromanspacer{} มยมฺปิ โข อาวุโส สาวกานํ เอวํ ธมฺมํ เทเสม “เอถ ตุเมฺห อาวุโส ปญฺจ นีวรเณ ปหาย เจตโส อุปกฺกิเลเส ปญฺญาย ทุพฺพลีกรเณ สตฺต โพชฺฌงฺเค ยถาภูตํ ภาเวถา”ติ.\csromanspacer{} อิธ โน อาวุโส โก วิเสโส, โก อธิปฺปยาโส, กิํ นานากรณํ สมณสฺส วา โคตมสฺส อมฺหากํ วา, ยทิทํ ธมฺมเทสนาย วา ธมฺมเทสนํ อนุสาสนิยา วา อนุสาสนินฺติ.}

\prose{อถ โข เต ภิกฺขู เตสํ อญฺญติตฺถิยานํ ปริพฺพาชกานํ ภาสิตํ เนว อภินนฺทิํสุ นปฺปฏิกฺโกสิํสุ, อนภินนฺทิตฺวา อปฺปฏิกฺโกสิตฺวา อุฏฺฐายาสนา ปกฺกมิํสุ “ภควโต สนฺติเก เอตสฺส ภาสิตสฺส อตฺถํ อาชานิสฺสามา”ติ.\csromanspacer{} อถ โข เต ภิกฺขู สาวตฺถิํ ปิณฺฑาย จริตฺวา ปจฺฉาภตฺตํ ปิณฺฑปาตปฏิกฺกนฺตา เยน ภควา เตนุปสงฺกมิํสุ, อุปสงฺกมิตฺวา ภควนฺตํ อภิวาเทตฺวา เอกมนฺตํ นิสีทิํสุ, เอกมนฺตํ นิสินฺนา โข เต ภิกฺขู ภควนฺตํ เอตทโวจุํ\csromandash{}}

\prose{\csromanfolio{96}อิธ มยํ ภนฺเต ปุพฺพณฺหสมยํ นิวาเสตฺวา ปตฺตจีวรมาทาย สาวตฺถิํ ปิณฺฑาย ปวิสิมฺห, เตสํ โน ภนฺเต อมฺหากํ เอตทโหสิ “อติปฺปโค โข ตาว สาวตฺถิยํ ปิณฺฑาย จริตุํ, ยํนูน มยํ เยน อญฺญติตฺถิยานํ ปริพฺพาชกานํ อาราโม เตนุปสงฺกเมยฺยามา”ติ.\csromanspacer{} อถ โข มยํ ภนฺเต เยน อญฺญติตฺถิยานํ ปริพฺพาชกานํ อาราโม เตนุปสงฺกมิมฺห, อุปสงฺกมิตฺวา เตหิ อญฺญติตฺถิเยหิ ปริพฺพาชเกหิ สทฺธิํ สมฺโมทิมฺห, สมฺโมทนียํ กถํ สารณียํ วีติสาเรตฺวา เอกมนฺตํ นิสีทิมฺห, เอกมนฺตํ นิสินฺเน โข อเมฺห ภนฺเต อญฺญติตฺถิยา ปริพฺพาชกา เอตทโวจุํ\csromandash{}}

\prose{สมโณ อาวุโส โคตโม สาวกานํ เอวํ ธมฺมํ เทเสติ “เอถ ตุเมฺห ภิกฺขเว ปญฺจ นีวรเณ ปหาย เจตโส อุปกฺกิเลเส ปญฺญาย ทุพฺพลีกรเณ สตฺต โพชฺฌงฺเค ยถาภูตํ ภาเวถา”ติ.\csromanspacer{} มยมฺปิ โข อาวุโส สาวกานํ เอวํ ธมฺมํ เทเสม “เอถ ตุเมฺห อาวุโส ปญฺจ นีวรเณ ปหาย เจตโส อุปกฺกิเลเส ปญฺญาย ทุพฺพลีกรเณ สตฺต โพชฺฌงฺเค ยถาภูตํ ภาเวถา”ติ, อิธ โน อาวุโส โก วิเสโส, โก อธิปฺปยาโส, กิํ นานากรณํ สมณสฺส วา โคตมสฺส อมฺหากํ วา, ยทิทํ ธมฺมเทสนาย วา ธมฺมเทสนํ อนุสาสนิยา วา อนุสาสนินฺ”ติ.}

\prose{อถ โข มยํ ภนฺเต เตสํ อญฺญติตฺถิยานํ ปริพฺพาชกานํ ภาสิตํ เนว อภินนฺทิมฺห นปฺปฏิกฺโกสิมฺห, อนภินนฺทิตฺวา อปฺปฏิกฺโกสิตฺวา อุฏฺฐายาสนา ปกฺกมิมฺห “ภควโต สนฺติเก เอตสฺส ภาสิตสฺส อตฺถํ อาชานิสฺสามา”ติ.}

\prose{เอวํวาทิโน ภิกฺขเว อญฺญติตฺถิยา ปริพฺพาชกา เอวมสฺสุวจนียา “อตฺถิ ปนาวุโส ปริยาโย, ยํ ปริยายํ อาคมฺม ปญฺจ นีวรณา ทส โหนฺติ, สตฺต โพชฺฌงฺคา จตุทฺทสา”ติ.\csromanspacer{} เอวํ ปุฏฺฐา ภิกฺขเว อญฺญติตฺถิยา ปริพฺพาชกา น เจว สมฺปายิสฺสนฺติ, อุตฺตริญฺจ วิฆาตํ อาปชฺชิสฺสนฺติ.\csromanspacer{} ตํ กิสฺส เหตุ, ยถา ตํ ภิกฺขเว อวิสยสฺมิํ.\csromanspacer{} นาหํ ตํ ภิกฺขเว ปสฺสามิ สเทวเก โลเก สมารเก สพฺรหฺมเก สสฺสมณพฺราหฺมณิยา ปชาย สเทวมนุสฺสาย, โย อิเมสํ ปญฺหานํ เวยฺยากรเณน จิตฺตํ อาราเธยฺย อญฺญตฺร ตถาคเตน วา ตถาคตสาวเกน วา อิโต วา ปน สุตฺวา.}

\csromanheader{2}{2. โพชฺฌงฺคสํยุตฺต}
\prose{\csromanfolio{97}กตโม จ ภิกฺขเว ปริยาโย, ยํ ปริยายํ อาคมฺม ปญฺจ นีวรณา ทส โหนฺติ.\csromanspacer{} ยทปิ ภิกฺขเว อชฺฌตฺตํ กามจฺฉนฺโท, ตทปิ นีวรณํ, ยทปิ พหิทฺธา กามจฺฉนฺโท, ตทปิ นีวรณํ กามจฺฉนฺทนีวรณนฺติ อิติ หิทํ อุทฺเทสํ คจฺฉติ, ตทมินาเปตํ ปริยาเยน ทฺวยํ โหติ.\csromanspacer{} ยทปิ ภิกฺขเว อชฺฌตฺตํ พฺยาปาโท, ตทปิ นีวรณํ, ยทปิ พหิทฺธา พฺยาปาโท, ตทปิ นีวรณํ, พฺยาปาทนีวรณนฺติ อิติ หิทํ อุทฺเทสํ คจฺฉติ, ตทมินาเปตํ ปริยาเยน ทฺวยํ โหติ.\csromanspacer{} ยทปิ ภิกฺขเว ถินํ, ตทปิ นีวรณํ, ยทปิ มิทฺธํ, ตทปิ นีวรณํ, ถินมิทฺธนีวรณนฺติ อิติ หิทํ อุทฺเทสํ คจฺฉติ, ตทมินาเปตํ ปริยาเยน ทฺวยํ โหติ.\csromanspacer{} ยทปิ ภิกฺขเว อุทฺธจฺจํ, ตทปิ นีวรณํ, ยทปิ กุกฺกุจฺจํ, ตทปิ นีวรณํ, อุทฺธจฺจกุกฺกุจฺจนีวรณนฺติ อิติ หิทํ อุทฺเทสํ คจฺฉติ, ตทมินาเปตํ ปริยาเยน ทฺวยํ โหติ.\csromanspacer{} ยทปิ ภิกฺขเว อชฺฌตฺตํ ธมฺเมสุ วิจิกิจฺฉา, ตทปิ นีวรณํ, ยทปิ พหิทฺธา ธมฺเมสุ วิจิกิจฺฉา, ตทปิ นีวรณํ, วิจิกิจฺฉานีวรณนฺติ อิติ หิทํ อุทฺเทสํ คจฺฉติ, ตทมินาเปตํ ปริยาเยน ทฺวยํ โหติ.\csromanspacer{} อยํ โข ภิกฺขเว ปริยาโย, ยํ ปริยายํ อาคมฺม ปญฺจ นีวรณา ทส โหนฺติ.}

\prose{กตโม จ ภิกฺขเว ปริยาโย, ยํ ปริยายํ อาคมฺม สตฺต โพชฺฌงฺคา จตุทฺทส โหนฺติ.\csromanspacer{} ยทปิ ภิกฺขเว อชฺฌตฺตํ ธมฺเมสุ สติ, ตทปิ สติสมฺโพชฺฌงฺโค, ยทปิ พหิทฺธา ธมฺเมสุ สติ, ตทปิ สติสมฺโพชฺฌงฺโค, สติสมฺโพชฺฌงฺโคติ อิติ หิทํ อุทฺเทสํ คจฺฉติ, ตทมินาเปตํ ปริยาเยน ทฺวยํ โหติ.}

\prose{ยทปิ ภิกฺขเว อชฺฌตฺตํ ธมฺเมสุ ปญฺญาย ปวิจินติ\footnote{ปวิจินาติ (ก)} ปวิจรติ ปริวีมํสมาปชฺชติ, ตทปิ ธมฺมวิจยสมฺโพชฺฌงฺโค.\csromanspacer{} ยทปิ พหิทฺธา ธมฺเมสุ ปญฺญาย ปวิจินติ ปวิจรติ ปริวีมํสมาปชฺชติ, ตทปิ ธมฺมวิจยสมฺโพชฺฌงฺโค, ธมฺมวิจยสมฺโพชฺฌงฺโคติ อิติ หิทํ อุทฺเทสํ คจฺฉติ, ตทมินาเปตํ ปริยาเยน ทฺวยํ โหติ.}

\prose{ยทปิ ภิกฺขเว กายิกํ วีริยํ, ตทปิ วีริยสมฺโพชฺฌงฺโค, ยทปิ เจตสิกํ วีริยํ, ตทปิ วีริยสมฺโพชฺฌงฺโค, วีริยสมฺโพชฺฌงฺโคติ อิติ หิทํ อุทฺเทสํ คจฺฉติ, ตทมินาเปตํ ปริยาเยน ทฺวยํ โหติ.}

\prose{\csromanfolio{98}ยทปิ ภิกฺขเว สวิตกฺกสวิจารา ปีติ, ตทปิ ปีติสมฺโพชฺฌงฺโค, ยทปิ อวิตกฺก-อวิจารา ปีติ, ตทปิ ปีติสมฺโพชฺฌงฺโค, ปีติสมฺโพชฺฌงฺโคติ อิติ หิทํ อุทฺเทสํ คจฺฉติ, ตทมินาเปตํ ปริยาเยน ทฺวยํ โหติ.}

\prose{ยทปิ ภิกฺขเว กายปฺปสฺสทฺธิ, ตทปิ ปสฺสทฺธิสมฺโพชฺฌงฺโค, ยทปิ จิตฺตปฺปสฺสทฺธิ, ตทปิ ปสฺสทฺธิสมฺโพชฺฌงฺโค, ปสฺสทฺธิสมฺโพชฺฌงฺโคติ อิติ หิทํ อุทฺเทสํ คจฺฉติ, ตทมินาเปตํ ปริยาเยน ทฺวยํ โหติ.}

\prose{ยทปิ ภิกฺขเว สวิตกฺโก สวิจาโร สมาธิ, ตทปิ สมาธิสมฺโพชฺฌงฺโค, ยทปิ อวิตกฺก-อวิจาโร สมาธิ, ตทปิ สมาธิสมฺโพชฺฌงฺโค, สมาธิสมฺโพชฺฌงฺโคติ อิติ หิทํ อุทฺเทสํ คจฺฉติ, ตทมินาเปตํ ปริยาเยน ทฺวยํ โหติ.}

\prose{ยทปิ ภิกฺขเว อชฺฌตฺตํ ธมฺเมสุ อุเปกฺขา, ตทปิ อุเปกฺขาสมฺโพชฺฌงฺโค, ยทปิ พหิทฺธา ธมฺเมสุ อุเปกฺขา, ตทปิ อุเปกฺขาสมฺโพชฺฌงฺโค, อุเปกฺขาสมฺโพชฺฌงฺโคติ อิติ หิทํ อุทฺเทสํ คจฺฉติ, ตทมินาเปตํ ปริยาเยน ทฺวยํ โหติ.\csromanspacer{} อยํ โข ภิกฺขเว ปริยาโย, ยํ ปริยายํ อาคมฺม สตฺต โพชฺฌงฺคา จตุทฺทสาติ.\csromanspacer{}\csromanspacer{}ทุติยํ.}

\csromanmatikaanchor{mtk.15}
\csromantocmarkref{paragraph}{3. อคฺคิสุตฺต}{mtk.15}
\bookmark[dest={mtk.15},level=4]{3. อคฺคิสุตฺต}
\csromanheader{4}{3. อคฺคิสุตฺต}
\proseitem{234}{อถ โข สมฺพหุลา ภิกฺขู ปุพฺพณฺหสมยํ นิวาเสตฺวา ปตฺตจีวรมาทาย สาวตฺถิํ ปิณฺฑาย ปวิสิํสุ. (ปริยายสุตฺตสทิสํ.)}

\prose{เอวํวาทิโน ภิกฺขเว อญฺญติตฺถิยา ปริพฺพาชกา เอวมสฺสุ วจนียา “ยสฺมิํ อาวุโส สมเย ลีนํ จิตฺตํ โหติ, กตเมสํ ตสฺมิํ สมเย โพชฺฌงฺคานํ อกาโล ภาวนาย, กตเมสํ ตสฺมิํ สมเย โพชฺฌงฺคานํ กาโล ภาวนาย.\csromanspacer{} ยสฺมิํ ปนาวุโส สมเย อุทฺธตํ จิตฺตํ โหติ, กตเมสํ ตสฺมิํ สมเย โพชฺฌงฺคานํ อกาโล ภาวนาย, กตเมสํ ตสฺมิํ สมเย โพชฺฌงฺคานํ กาโล ภาวนายา”ติ.\csromanspacer{} เอวํ ปุฏฺฐา ภิกฺขเว อญฺญติตฺถิยา ปริพฺพาชกา น เจว สมฺปายิสฺสนฺติ, อุตฺตริญฺจ วิฆาตํ อาปชฺชิสฺสนฺติ.\csromanspacer{} ตํ กิสฺส เหตุ, ยถา ตํ ภิกฺขเว อวิสยสฺมิํ.}

\prose{นาหํ ตํ ภิกฺขเว ปสฺสามิ สเทวเก โลเก สมารเก สพฺรหฺมเก สสฺสมณพฺราหฺมณิยา ปชาย สเทวมนุสฺสาย, โย อิเมสํ ปญฺหานํ}

\csromanheader{2}{2. โพชฺฌงฺคสํยุตฺต}
\prose{\csromanfolio{99}เวยฺยากรเณน จิตฺตํ อาราเธยฺย อญฺญตฺร ตถาคเตน วา ตถาคตสาวเกน วา อิโต วา ปน สุตฺวา.}

\prose{ยสฺมิํ ภิกฺขเว สมเย ลีนํ จิตฺตํ โหติ, อกาโล ตสฺมิํ สมเย ปสฺสทฺธิสมฺโพชฺฌงฺคสฺส ภาวนาย, อกาโล สมาธิสมฺโพชฺฌงฺคสฺส ภาวนาย, อกาโล อุเปกฺขาสมฺโพชฺฌงฺคสฺส ภาวนาย.\csromanspacer{} ตํ กิสฺส เหตุ, ลีนํ ภิกฺขเว จิตฺตํ, ตํ เอเตหิ ธมฺเมหิ ทุสฺสมุฏฺฐาปยํ โหติ.}

\prose{เสยฺยถาปิ ภิกฺขเว ปุริโส ปริตฺตํ อคฺคิํ อุชฺชาเลตุกาโม อสฺส, โส ตตฺถ อลฺลานิ เจว ติณานิ ปกฺขิเปยฺย, อลฺลานิ จ โคมยานิ ปกฺขิเปยฺย, อลฺลานิ จ กฏฺฐานิ ปกฺขิเปยฺย, อุทกวาตญฺจ ทเทยฺย, ปํสุเกน จ โอกิเรยฺย, ภพฺโพ นุ โข โส ปุริโส ปริตฺตํ อคฺคิํ อุชฺชาเลตุนฺติ.\csromanspacer{} โน เหตํ ภนฺเต.}

\prose{เอวํเมว โข ภิกฺขเว ยสฺมิํ สมเย ลีนํ จิตฺตํ โหติ, อกาโล ตสฺมิํ สมเย ปสฺสทฺธิสมฺโพชฺฌงฺคสฺส ภาวนาย, อกาโล สมาธิสมฺโพชฺฌงฺคสฺส ภาวนาย, อกาโล อุเปกฺขาสมฺโพชฺฌงฺคสฺส ภาวนาย.\csromanspacer{} ตํ กิสฺส เหตุ, ลีนํ ภิกฺขเว จิตฺตํ, ตํ เอเตหิ ธมฺเมหิ ทุสฺสมุฏฺฐาปยํ โหติ.}

\prose{ยสฺมิํ จ โข ภิกฺขเว สมเย ลีนํ จิตฺตํ โหติ, กาโล ตสฺมิํ สมเย ธมฺมวิจยสมฺโพชฺฌงฺคสฺส ภาวนาย, กาโล วีริยสมฺโพชฺฌงฺคสฺส ภาวนาย, กาโล ปีติสมฺโพชฺฌงฺคสฺส ภาวนาย.\csromanspacer{} ตํ กิสฺส เหตุ, ลีนํ ภิกฺขเว จิตฺตํ, ตํ เอเตหิ ธมฺเมหิ สุสมุฏฺฐาปยํ โหติ.}

\prose{เสยฺยถาปิ ภิกฺขเว ปุริโส ปริตฺตํ อคฺคิํ อุชฺชาเลตุกาโม อสฺส, โส ตตฺถ สุกฺขานิ เจว ติณานิ ปกฺขิเปยฺย, สุกฺขานิ โคมยานิ ปกฺขิเปยฺย, สุกฺขานิ กฏฺฐานิ ปกฺขิเปยฺย, มุขวาตญฺจ ทเทยฺย, น จ ปํสุเกน โอกิเรยฺย, ภพฺโพ นุ โข โส ปุริโส ปริตฺตํ อคฺคิํ อุชฺชาเลตุนฺติ.\csromanspacer{} เอวํ ภนฺเต.}

\prose{เอวเมว โข ภิกฺขเว ยสฺมิํ สมเย ลีนํ จิตฺตํ โหติ, กาโล ตสฺมิํ สมเย ธมฺมวิจยสมฺโพชฺฌงฺคสฺส ภาวนาย, กาโล วิริยสมฺโพชฺฌงฺคสฺส ภาวนาย, กาโล ปีติสมฺโพชฺฌงฺคสฺส ภาวนาย.\csromanspacer{} ตํ กิสฺส เหตุ, ลีนํ ภิกฺขเว จิตฺตํ, ตํ เอเตหิ ธมฺเมหิ สุสมุฏฺฐาปยํ โหติ.}

\prose{\csromanfolio{100}ยสฺมิํ ภิกฺขเว สมเย อุทฺธตํ จิตฺตํ โหติ, อกาโล ตสฺมิํ สมเย ธมฺมวิจยสมฺโพชฺฌงฺคสฺส ภาวนาย, อกาโล วีริยสมฺโพชฺฌงฺคสฺส ภาวนาย, อกาโล ปีติสมฺโพชฺฌงฺคสฺส ภาวนาย.\csromanspacer{} ตํ กิสฺส เหตุ, อุทฺธตํ ภิกฺขเว จิตฺตํ, ตํ เอเตหิ ธมฺเมหิ ทุวูปสมยํ โหติ.}

\prose{เสยฺยถาปิ ภิกฺขเว ปุริโส มหนฺตํ อคฺคิกฺขนฺธํ นิพฺพาเปตุกาโม อสฺส, โส ตตฺถ สุกฺขานิ เจว ติณานิ ปกฺขิเปยฺย, สุกฺขานิ จ โคมยานิ ปกฺขิเปยฺย, สุกฺขานิ จ กฏฺฐานิ ปกฺขิเปยฺย, มุขวาตญฺจ ทเทยฺย, น จ ปํสุเกน โอกิเรยฺย, ภพฺโพ นุ โข โส ปุริโส มหนฺตํ อคฺคิกฺขนฺธํ นิพฺพาเปตุนฺติ.\csromanspacer{} โน เหตํ ภนฺเต.}

\prose{เอวเมว โข ภิกฺขเว ยสฺมิํ สมเย อุทฺธตํ จิตฺตํ โหติ, อกาโล ตสฺมิํ สมเย ธมฺมวิจยสมฺโพชฺฌงฺคสฺส ภาวนาย, อกาโล วีริยสมฺโพชฺฌงฺคสฺส ภาวนาย, อกาโล ปีติสมฺโพชฺฌงฺคสฺส ภาวนาย.\csromanspacer{} ตํ กิสฺส เหตุ, อุทฺธตํ ภิกฺขเว จิตฺตํ, ตํ เอเตหิ ธมฺเมหิ ทุวูปสมยํ โหติ.}

\prose{ยสฺมิํ จ โข ภิกฺขเว สมเย อุทฺธตํ จิตฺตํ โหติ, กาโล ตสฺมิํ สมเย ปสฺสทฺธิสมฺโพชฺฌงฺคสฺส ภาวนาย, กาโล สมาธิสมฺโพชฺฌงฺคสฺส ภาวนาย, กาโล อุเปกฺขาสมฺโพชฺฌงฺคสฺส ภาวนาย.\csromanspacer{} ตํ กิสฺส เหตุ, อุทฺธตํ ภิกฺขเว จิตฺตํ, ตํ เอเตหิ ธมฺเมหิ สุวูปสมยํ โหติ.}

\prose{เสยฺยถาปิ ภิกฺขเว ปุริโส มหนฺตํ อคฺคิกฺขนฺธํ นิพฺพาเปตุกาโม อสฺส, โส ตตฺถ อลฺลานิ เจว ติณานิ ปกฺขิเปยฺย, อลฺลานิ จ โคมยานิ ปกฺขิเปยฺย, อลฺลานิ จ กฏฺฐานิ ปกฺขิเปยฺย, อุทกวาตญฺจ ทเทยฺย, ปํสุเกน จ โอกิเรยฺย, ภพฺโพ นุ โข ปุริโส มหนฺตํ อคฺคิกฺขนฺธํ นิพฺพาเปตุนฺติ.\csromanspacer{} เอวํ ภนฺเต.}

\prose{เอวเมว โข ภิกฺขเว ยสฺมิํ สมเย อุทฺธตํ จิตฺตํ โหติ, กาโล ตสฺมิํ สมเย ปสฺสทฺธิสมฺโพชฺฌงฺคสฺส ภาวนาย, กาโล สมาธิสมฺโพชฺฌงฺคสฺส ภาวนาย, กาโล อุเปกฺขาสมฺโพชฺฌงฺคสฺส ภาวนาย.\csromanspacer{} ตํ กิสฺส เหตุ, อุทฺธตํ ภิกฺขเว จิตฺตํ, ตํ เอเตหิ ธมฺเมหิ สุวูปสมยํ โหติ.\csromanspacer{} สติํ จ ขฺวาหํ ภิกฺขเว สพฺพตฺถิกํ วทามีติ.\csromanspacer{}\csromanspacer{}ตติยํ.}

\csromanheader{2}{2. โพชฺฌงฺคสํยุตฺต}
\csromanmatikaanchor{mtk.16}
\csromantocmarkref{paragraph}{4. เมตฺตาสหคตสุตฺต (หลิทฺทวสนสุตฺต)}{mtk.16}
\bookmark[dest={mtk.16},level=4]{4. เมตฺตาสหคตสุตฺต (หลิทฺทวสนสุตฺต)}
\csromanheader{4}{4. เมตฺตาสหคตสุตฺต}
\proseitem{235}{\csromanfolio{101}เอกํ สมยํ ภควา โกลิเยสุ วิหรติ หลิทฺทวสนํ นาม โกลิยานํ นิคโม.\csromanspacer{} อถ โข สมฺพหุลา ภิกฺขู ปุพฺพณฺหสมยํ นิวาเสตฺวา ปตฺตจีวรมาทาย หลิทฺทวสนํ ปิณฺฑาย ปวิสิํสุ.\csromanspacer{} อถ โข เตสํ ภิกฺขูนํ เอตทโหสิ “อติปฺปโค โข ตาว หลิทฺทวสเน ปิณฺฑาย จริตุํ, ยํนูน มยํ เยน อญฺญติตฺถิยานํ ปริพฺพาชกานํ อาราโม เตนุปสงฺกเมยฺยามา”ติ.}

\prose{อถ โข เต ภิกฺขู เยน อญฺญติตฺถิยานํ ปริพฺพาชกานํ อาราโม เตนุปสงฺกมิํสุ, อุปสงฺกมิตฺวา เตหิ อญฺญติตฺถิเยหิ ปริพฺพาชเกหิ สทฺธิํ สมฺโมทิํสุ, สมฺโมทนียํ กถํ สารณียํ วีติสาเรตฺวา เอกมนฺตํ นิสีทิํสุ, เอกมนฺตํ นิสินฺเน โข เต ภิกฺขู อญฺญติตฺถิยา ปริพฺพาชกา เอตทโวจุํ\csromandash{}}

\prose{“สมโณ อาวุโส โคตโม สาวกานํ เอวํ ธมฺมํ เทเสติ ‘เอถ ตุเมฺห ภิกฺขเว ปญฺจ นีวรเณ ปหาย เจตโส อุปกฺกิเลเส ปญฺญาย ทุพฺพลีกรเณ เมตฺตาสหคเตน เจตสา เอกํ ทิสํ ผริตฺวา วิหรถ.\csromanspacer{} ตถา ทุติยํ.\csromanspacer{} ตถา ตติยํ.\csromanspacer{} ตถา จตุตฺถํ.\csromanspacer{} อิติ อุทฺธมโธ ติริยํ สพฺพธิ สพฺพตฺตตาย สพฺพาวนฺตํ โลกํ เมตฺตาสหคเตน เจตสา วิปุเลน มหคฺคเตน อปฺปมาเณน อเวเรน อพฺยาปชฺเชน ผริตฺวา วิหรถ.\csromanspacer{} กรุณาสหคเตน เจตสา เอกํ ทิสํ ผริตฺวา วิหรถ.\csromanspacer{} ตถา ทุติยํ.\csromanspacer{} ตถา ตติยํ.\csromanspacer{} ตถา จตุตฺถํ.\csromanspacer{} อิติ อุทฺธมโธ ติริยํ สพฺพธิ สพฺพตฺตตาย สพฺพาวนฺตํ โลกํ กรุณาสหคเตน เจตสา วิปุเลน มหคฺคเตน อปฺปมาเณน อเวเรน อพฺยาปชฺเชน ผริตฺวา วิหรถ.\csromanspacer{} มุทิตาสหคเตน เจตสา เอกํ ทิสํ ผริตฺวา วิหรถ.\csromanspacer{} ตถา ทุติยํ.\csromanspacer{} ตถา ตติยํ.\csromanspacer{} ตถา จตุตฺถํ.\csromanspacer{} อิติ อุทฺธมโธ ติริยํ สพฺพธิ สพฺพตฺตตาย สพฺพาวนฺตํ โลกํ มุทิตาสหคเตน เจตสา วิปุเลน มหคฺคเตน อปฺปมาเณน อเวเรน อพฺยาปชฺเชน ผริตฺวา วิหรถ.\csromanspacer{} อุเปกฺขาสหคเตน เจตสา เอกํ ทิสํ ผริตฺวา วิหรถ.\csromanspacer{} ตถา ทุติยํ.\csromanspacer{} ตถา ตติยํ.\csromanspacer{} ตถา จตุตฺถํ.\csromanspacer{} อิติ อุทฺธมโธ ติริยํ สพฺพธิ สพฺพตฺตตาย สพฺพาวนฺตํ โลกํ อุเปกฺขาสหคเตน เจตสา วิปุเลน มหคฺคเตน อปฺปมาเณน อเวเรน อพฺยาปชฺเชน ผริตฺวา วิหรถา’ติ.}

\prose{\csromanfolio{102}มยมฺปิ โข อาวุโส สาวกานํ เอวํ ธมฺมํ เทเสม ‘เอถ ตุเมฺห อาวุโส ปญฺจ นีวรเณ ปหาย เจตโส อุปกฺกิเลเส ปญฺญาย ทุพฺพลีกรเณ เมตฺตาสหคเตน เจตสา เอกํ ทิสํ ผริตฺวา วิหรถ ฯเปฯ กรุณาสหคเตน เจตสา.\csromanspacer{} มุทิตาสหคเตน เจตสา.\csromanspacer{} อุเปกฺขาสหคเตน เจตสา เอกํ ทิสํ ผริตฺวา วิหรถ.\csromanspacer{} ตถา ทุติยํ.\csromanspacer{} ตถา ตติยํ.\csromanspacer{} ตถา จตุตฺถํ.\csromanspacer{} อิติ อุทฺธมโธ ติริยํ สพฺพธิ สพฺพตฺตตาย สพฺพาวนฺตํ โลกํ อุเปกฺขาสหคเตน เจตสา วิปุเลน มหคฺคเตน อปฺปมาเณน อเวเรน อพฺยาปชฺเชน ผริตฺวา วิหรถา’ติ.\csromanspacer{} อิธ โน อาวุโส โก วิเสโส, โก อธิปฺปยาโส, กิํ นานากรณํ สมณสฺส วา โคตมสฺส อมฺหากํ วา, ยทิทํ ธมฺมเทสนาย วา ธมฺมเทสนํ, อนุสาสนิยา วา อนุสาสนินฺ”ติ.}

\prose{อถ โข เต ภิกฺขู เตสํ อญฺญติตฺถิยานํ ปริพฺพาชกานํ ภาสิตํ เนว อภินนฺทิํสุ นปฺปฏิกฺโกสิํสุ, อนภินนฺทิตฺวา อปฺปฏิกฺโกสิตฺวา อุฏฺฐายาสนา ปกฺกมิํสุ “ภควโต สนฺติเก เอตสฺส ภาสิตสฺส อตฺถํ อาชานิสฺสามา”ติ.\csromanspacer{} อถ โข เต ภิกฺขู หลิทฺทวสเน ปิณฺฑาย จริตฺวา ปจฺฉาภตฺตํ ปิณฺฑปาตปฏิกฺกนฺตา เยน ภควา เตนุปสงฺกมิํสุ, อุปสงฺกมิตฺวา ภควนฺตํ อภิวาเทตฺวา เอกมนฺตํ นิสินฺนา โข เต ภิกฺขู ภควนฺตํ เอตทโวจุํ\csromandash{}}

\prose{“อิธ มยํ ภนฺเต ปุพฺพณฺหสมยํ นิวาเสตฺวา ปตฺตจีวรมาทาย หลิทฺทวสเน ปิณฺฑาย ปวิสิมฺห, เตสํ โน ภนฺเต อมฺหากํ เอตทโหสิ ‘อติปฺปโค โข ตาว หลิทฺทวสเน ปิณฺฑาย จริตุํ, ยํนูน มยํ เยน อญฺญติตฺถิยานํ ปริพฺพาชกานํ อาราโม เตนุปสงฺกเมยฺยามา’ติ.}

\prose{อถ โข มยํ ภนฺเต เยน อญฺญติตฺถิยานํ ปริพฺพาชกานํ อาราโม เตนุปสงฺกมิมฺห, อุปสงฺกมิตฺวา เตหิ อญฺญติตฺถิเยหิ ปริพฺพาชเกหิ สทฺธิํ สมฺโมทิมฺห, สมฺโมทนียํ กถํ สารณียํ วีติสาเรตฺวา เอกมนฺตํ นิสีทิมฺห, เอกมนฺตํ นิสินฺเน โข อเมฺห ภนฺเต เต อญฺญติตฺถิยา ปริพฺพาชกา เอตทโวจุํ\csromandash{}}

\prose{สมโณ อาวุโส โคตโม สาวกานํ เอวํ ธมฺมํ เทเสติ ‘เอถ ตุเมฺห ภิกฺขเว ปญฺจ นีวรเณ ปหาย เจตโส อุปกฺกิเลเส ปญฺญาย ทุพฺพลีกรเณ เมตฺตาสหคเตน เจตสา เอกํ ทิสํ}

\csromanheader{2}{2. โพชฺฌงฺคสํยุตฺต}
\prose{\csromanfolio{103}ผริตฺวา วิหรถ ฯเปฯ กรุณาสหคเตน เจตสา ฯเปฯ มุทิตาสหคเตน เจตสา ฯเปฯ อุเปกฺขาสหคเตน เจตสา เอกํ ทิสํ ผริตฺวา วิหรถ.\csromanspacer{} ตถา ทุติยํ.\csromanspacer{} ตถา ตติยํ.\csromanspacer{} ตถา จตุตฺถํ.\csromanspacer{} อิติ อุทฺธมโธ ติริยํ สพฺพธิ สพฺพตฺตตาย สพฺพาวนฺตํ โลกํ อุเปกฺขาสหคเตน เจตสา วิปุเลน มหคฺคเตน อปฺปมาเณน อเวเรน อพฺยาปชฺเชน ผริตฺวา วิหรถา’ติ.}

\prose{มยมฺปิ โข อาวุโส สาวกานํ เอวํ ธมฺมํ เทเสม ‘เอถ ตุเมฺห อาวุโส ปญฺจ นีวรเณ ปหาย เจตโส อุปกฺกิเลเส ปญฺญาย ทุพฺพลีกรเณ เมตฺตาสหคเตน เจตสา เอกํ ทิสํ ผริตฺวา วิหรถ ฯเปฯ กรุณาสหคเตน เจตสา ฯเปฯ มุทิตาสหคเตน เจตสา ฯเปฯ อุเปกฺขาสหคเตน เจตสา เอกํ ทิสํ ผริตฺวา วิหรถ.\csromanspacer{} ตถา ทุติยํ.\csromanspacer{} ตถ ตติยํ.\csromanspacer{} ตถา จตุตฺถํ.\csromanspacer{} อิติ อุทฺธมโธ ติริยํ สพฺพธิ สพฺพตฺตตาย สพฺพาวนฺตํ โลกํ อุเปกฺขาสหคเตน เจตสา วิปุเลน มหคฺคเตน อปฺปมาเณน อเวเรน อพฺยาปชฺเชน ผริตฺวา วิหรถา’ติ.\csromanspacer{} อิธ โน อาวุโส โก วิเสโส, โก อธิปฺปยาโส, กิํ นานากรณํ สมณสฺส วา โคตมสฺส อมฺหากํ วา, ยทิทํ ธมฺมเทสนาย วา ธมฺมเทสนํ อนุสาสนิยา วา อนุสาสนินฺ’ติ.}

\prose{อถ โข มยํ ภนฺเต เตสํ อญฺญติตฺถิยานํ ปริพฺพาชกานํ ภาสิตํ เนว อภินนฺทิมฺห นปฺปฏิกฺโกสิมฺห, อนภินนฺทิตฺวา อปฺปฏิกฺโกสิตฺวา อุฏฺฐายาสนา ปกฺกมิมฺห ‘ภควโต สนฺติเก เอตสฺส ภาสิตสฺส อตฺถํ อาชานิสฺสามา’ติ”.}

\prose{เอวํวาทิโน ภิกฺขเว อญฺญติตฺถิยา ปริพฺพาชกา เอวมสฺสุ วจนียา “กถํ ภาวิตา ปนาวุโส เมตฺตาเจโตวิมุตฺติ กิํคติกา โหติ กิํปรมา กิํผลา กิํปริโยสานา, กถํ ภาวิตา ปนาวุโส กรุณาเจโตวิมุตฺติ กิํคติกา โหติ กิํปรมา กิํผลา กิํปริโยสานา, กถํ ภาวิตา ปนาวุโส มุทิตาเจโตวิมุตฺติ กิํคติกา โหติ กิํปรมา กิํผลา กิํปริโยสานา, กถํ ภาวิตา ปนาวุโส อุเปกฺขาเจโตวิมุตฺติ กิํคติกา โหติ กิํปรมา กิํผลา กิํปริโยสานา”ติ.\csromanspacer{} เอวํ ปุฏฺฐา ภิกฺขเว อญฺญติตฺถิยา ปริพฺพาชกา น เจว สมฺปายิสฺสนฺติ, อุตฺตริญฺจ วิฆาตํ อาปชฺชิสฺสนฺติ.\csromanspacer{} ตํ กิสฺส เหตุ, ยถา ตํ ภิกฺขเว อวิสยสฺมิํ.\csromanspacer{} นาหํ ตํ ภิกฺขเว ปสฺสามิ}

\prose{\csromanfolio{104}สเทวเก โลเก สมารเก สพฺรหฺมเก สสฺสมณพฺราหฺมณิยา ปชาย สเทวมนุสฺสาย, โย อิเมสํ ปญฺหานํ เวยฺยากรเณน จิตฺตํ อาราเธยฺย อญฺญตฺร ตถาคเตน วา ตถาคตสาวเกน วา อิโต วา ปน สุตฺวา.}

\prose{กถํ ภาวิตา จ ภิกฺขเว เมตฺตาเจโตวิมุตฺติ กิํคติกา โหติ กิํปรมา กิํผลา กิํปริโยสานา, อิธ ภิกฺขเว ภิกฺขุ เมตฺตาสหคตํ สติสมฺโพชฺฌงฺคํ ภาเวติ ฯเปฯ เมตฺตาสหคตํ อุเปกฺขาสมฺโพชฺฌงฺคํ ภาเวติ วิเวกนิสฺสิตํ วิราคนิสฺสิตํ นิโรธนิสฺสิตํ โวสฺสคฺคปริณามิํ.\csromanspacer{} โส สเจ อากงฺขติ “อปฺปฏิกูเล ปฏิกูลสญฺญี วิหเรยฺยนฺ”ติ, ปฏิกูลสญฺญี ตตฺถ วิหรติ.\csromanspacer{} สเจ อากงฺขติ “ปฏิกูเล อปฺปฏิกูลสญฺญี วิหเรยฺยนฺ”ติ, อปฺปฏิกูลสญฺญี ตตฺถ วิหรติ.\csromanspacer{} สเจ อากงฺขติ “อปฺปฏิกูเล จ ปฏิกูเล จ ปฏิกูลสญฺญี วิหเรยฺยนฺ”ติ, ปฏิกูลสญฺญี ตตฺถ วิหรติ.\csromanspacer{} สเจ อากงฺขติ “ปฏิกูเล จ อปฺปฏิกูเล จ อปฺปฏิกูลสญฺญี วิหเรยฺยนฺ”ติ, อปฺปฏิกูลสญฺญี ตตฺถ วิหรติ.\csromanspacer{} สเจ อากงฺขติ “อปฺปฏิกูลญฺจ ปฏิกูลญฺจ ตทุภยํ อภินิวชฺเชตฺวา อุเปกฺขโก วิหเรยฺยํ สโต สมฺปชาโน”ติ, อุเปกฺขโก จ ตตฺถ วิหรติ สโต สมฺปชาโน.\csromanspacer{} สุภํ วา โข ปน วิโมกฺขํ อุปสมฺปชฺช วิหรติ, สุภปรมาหํ ภิกฺขเว เมตฺตาเจโตวิมุตฺติํ วทามิ อิธปญฺญสฺส ภิกฺขุโน อุตฺตริวิมุตฺติํ อปฺปฏิวิชฺฌโต.}

\prose{กถํ ภาวิตา จ ภิกฺขเว กรุณาเจโตวิมุตฺติ กิํคติกา โหติ กิํปรมา กิํผลา กิํปริโยสานา, อิธ ภิกฺขเว ภิกฺขุ กรุณาสหคตํ สติสมฺโพชฺฌงฺคํ ภาเวติ ฯเปฯ กรุณาสหคตํ อุเปกฺขาสมฺโพชฺฌงฺคํ ภาเวติ วิเวกนิสฺสิตํ วิราคนิสฺสิตํ นิโรธนิสฺสิตํ โวสฺสคฺคปริณามิํ.\csromanspacer{} โส สเจ อากงฺขติ “อปฺปฏิกูเล ปฏิกูลสญฺญี วิหเรยฺยนฺ”ติ, ปฏิกูลสญฺญี ตตฺถ วิหรติ ฯเปฯ.\csromanspacer{} สเจ อากงฺขติ “อปฺปฏิกูลญฺจ ปฏิกูลญฺจ ตทุภยํ อภินิวชฺเชตฺวา อุเปกฺขโก วิหเรยฺยํ สโต สมฺปชาโน”ติ, อุเปกฺขโก ตตฺถ วิหรติ สโต สมฺปชาโน.\csromanspacer{} สพฺพโส วา ปน รูปสญฺญานํ สมติกฺกมา ปฏิฆสญฺญานํ อตฺถงฺคมา นานตฺตสญฺญานํ อมนสิการา “อนนฺโต อากาโส”ติ อากาสานญฺจายตนํ อุปสมฺปชฺช วิหรติ, อากาสานญฺจายตนปรมาหํ ภิกฺขเว}

\csromanheader{2}{2. โพชฺฌงฺคสํยุตฺต}
\prose{\csromanfolio{105}กรุณาเจโตวิมุตฺติํ วทามิ อิธปญฺญสฺส ภิกฺขุโน อุตฺตริวิมุตฺติํ อปฺปฏิวิชฺฌโต.}

\prose{กถํ ภาวิตา จ ภิกฺขเว มุทิตาเจโตวิมุตฺติ กิํคติกา โหติ กิํปรมา กิํผลา กิํปริโยสานา, อิธ ภิกฺขเว ภิกฺขุ มุทิตาสหคตํ สติสมฺโพชฺฌงฺคํ ภาเวติ ฯเปฯ มุทิตาสหคตํ อุเปกฺขาสมฺโพชฺฌงฺคํ ภาเวติ วิเวกนิสฺสิตํ วิราคนิสฺสิตํ นิโรธนิสฺสิตํ โวสฺสคฺคปริณามิํ.\csromanspacer{} โส สเจ อากงฺขติ “อปฺปฏิกูเล ปฏิกูลสญฺญี วิหเรยฺยนฺ”ติ, ปฏิกูลสญฺญี ตตฺถ วิหรติ ฯเปฯ.\csromanspacer{} สเจ อากงฺขติ “อปฺปฏิกูลญฺจ ปฏิกูลญฺจ ตทุภยํ อภินิวชฺเชตฺวา อุเปกฺขโก วิหเรยฺยํ สโต สมฺปชาโน”ติ, อุเปกฺขโก ตตฺถ วิหรติ สโต สมฺปชาโน.\csromanspacer{} สพฺพโส วา ปน อากาสานญฺจายตนํ สมติกฺกมฺม “อนนฺตํ วิญฺญาณนฺ”ติ วิญฺญาณญฺจายตนํ อุปสมฺปชฺช วิหรติ, วิญฺญาณญฺจายตนปรมาหํ ภิกฺขเว มุทิตาเจโตวิมุตฺติํ วทามิ อิธปญฺญสฺส ภิกฺขุโน อุตฺตริวิมุตฺติํ อปฺปฏิวิชฺฌโต.}

\prose{กถํ ภาวิตา จ ภิกฺขเว อุเปกฺขาเจโตวิมุตฺติ กิํคติกา โหติ กิํปรมา กิํผลา กิํปริโยสานา, อิธ ภิกฺขเว ภิกฺขุ อุเปกฺขาสหคตํ สติสมฺโพชฺฌงฺคํ ภาเวติ วิเวกนิสฺสิตํ วิราคนิสฺสิตํ นิโรธนิสฺสิตํ โวสฺสคฺคปริณามิํ ฯเปฯ อุเปกฺขาสหคตํ อุเปกฺขาสมฺโพชฺฌงฺคํ ภาเวติ วิเวกนิสฺสิตํ วิราคนิสฺสิตํ นิโรธนิสฺสิตํ โวสฺสคฺคปริณามิํ.\csromanspacer{} โส สเจ อากงฺขติ “อปฺปฏิกูเล ปฏิกูลสญฺญี วิหเรยฺยนฺ”ติ, ปฏิกูลสญฺญี ตตฺถ วิหรติ.\csromanspacer{} สเจ อากงฺขติ “ปฏิกูเล อปฺปฏิกูลสญฺญี วิหเรยฺยนฺ”ติ, อปฺปฏิกูลสญฺญี ตตฺถ วิหรติ.\csromanspacer{} สเจ อากงฺขติ “อปฺปฏิกูเล จ ปฏิกูเล จ ปฏิกูลสญฺญี วิหเรยฺยนฺ”ติ, ปฏิกูลสญฺญี ตตฺถ วิหรติ.\csromanspacer{} สเจ อากงฺขติ “ปฏิกูเล จ อปฺปฏิกูเล จ อปฺปฏิกูลสญฺญี วิหเรยฺยนฺ”ติ, อปฺปฏิกูลสญฺญี ตตฺถ วิหรติ.\csromanspacer{} สเจ อากงฺขติ “อปฺปฏิกูลญฺจ ปฏิกูลญฺจ ตทุภยํ อภินิวชฺเชตฺวา อุเปกฺขโก วิหเรยฺยํ สโต สมฺปชาโน”ติ, อุเปกฺขโก ตตฺถ วิหรติ สโต สมฺผชาโน.\csromanspacer{} สพฺพโส วา ปน วิญฺญาณญฺจายตนํ สมติกฺกมฺม “นตฺถิ กิญฺจี”ติ อากิญฺจญฺญายตนํ อุปสมฺปชฺช วิหรติ, อากิญฺจญฺญายตนปรมาหํ ภิกฺขเว อุเปกฺขาเจโตวิมุตฺติํ วทามิ อิธปญฺญสฺส ภิกฺขุโน อุตฺตริวิมุตฺตํ อปฺปฏิวิชฺฌโตติ.\csromanspacer{}\csromanspacer{}จตุตฺถํ.}

\csromanmatikaanchor{mtk.17}
\csromantocmarkref{paragraph}{5. สงฺคารวสุตฺต}{mtk.17}
\bookmark[dest={mtk.17},level=4]{5. สงฺคารวสุตฺต}
\csromanheader{4}{5. สงฺคารวสุตฺต}
\proseitem{236}{\csromanfolio{106}สาวตฺถินิทานํ.\csromanspacer{} อถ โข สงฺคารโว พฺราหฺมโณ เยน ภควา เตนุปสงฺกมิ, อุปสงฺกมิตฺวา ภควตา สทฺธิํ สมฺโมทิ, สมฺโมทนียํ กถํ สารณียํ วีติสาเรตฺวา เอกมนฺตํ นิสีทิ, เอกมนฺตํ นิสินฺโน โข สงฺคารโว พฺราหฺมโณ ภควนฺตํ เอตทโวจ\csromandash{}}

\prose{“โก นุ โข โภ โคตม เหตุ โก ปจฺจโย, เยเนกทา ทีฆรตฺตํ สชฺฌายกตาปิ มนฺตา นปฺปฏิภนฺติ, ปเคว อสชฺฌายกตา.\csromanspacer{} โก ปน โภ โคตม เหตุ โก ปจฺจโย, เยเนกทา ทีฆรตฺตํ อสชฺฌายกตาปิ มนฺตา ปฏิภนฺติ, ปเคว สชฺฌายกตา”ติ.}

\prose{ยสฺมิํ โข พฺราหฺมณ สมเย กามราคปริยุฏฺฐิเตน เจตสา วิหรติ กามราคปเรเตน, อุปฺปนฺนสฺส จ กามราคสฺส นิสฺสรณํ ยถาภูตํ นปฺปชานาติ.\csromanspacer{} อตฺตตฺถมฺปิ ตสฺมิํ สมเย ยถาภูตํ น ชานาติ น ปสฺสติ, ปรตฺถมฺปิ ตสฺมิํ สมเย ยถาภูตํ น ชานาติ น ปสฺสติ, อุภยตฺถมฺปิ ตสฺมิํ สมเย ยถาภูตํ น ชานาติ น ปสฺสติ, ทีฆรตฺตํ สชฺฌายกตาปิ มนฺตา นปฺปฏิภนฺติ, ปเคว อสชฺฌายกตา.\csromanspacer{} เสยฺยถาปิ พฺราหฺมณ อุทปตฺโต สํสฏฺโฐ ลาขาย วา หลิทฺทิยา วา นีลิยา วา มญฺชิฏฺฐาย วา ตตฺถ จกฺขุมา ปุริโส สกํ มุขนิมิตฺตํ ปจฺจเวกฺขมาโน ยถาภูตํ น ชาเนยฺย น ปสฺเสยฺย.\csromanspacer{} เอวเมว โข พฺราหฺมณ ยสฺมิํ สมเย กามราคปริยุฏฺฐิเตน เจตสา วิหรติ กามราคปเรเตน, อุปฺปนฺนสฺส จ กามราคสฺส นิสฺสรณํ ยถาภูตํ นปฺปชานาติ.\csromanspacer{} อตฺตตฺถมฺปิ ตสฺมิํ สมเย ยถาภูตํ น ชานาติ น ปสฺสติ, ปรตฺถมฺปิ ฯเปฯ อุภยตฺถมฺปิ ตสฺมิํ สมเย ยถาภูตํ น ชานาติ น ปสฺสติ, ทีฆรตฺตํ สชฺฌายกตาปิ มนฺตา นปฺปฏิภนฺติ, ปเคว อสชฺฌายกตา.}

\prose{ปุน จปรํ พฺราหฺมณ ยสฺมิํ สมเย พฺยาปาทปริยุฏฺฐิเตน เจตสา วิหรติ พฺยาปาทปเรเตน, อุปฺปนฺนสฺส จ พฺยาปาทสฺส นิสฺสรณํ ยถาภูตํ นปฺปชานาติ.\csromanspacer{} อตฺตตฺถมฺปิ ตสฺมิํ สมเย ยถาภูตํ น ชานาติ น ปสฺสติ, ปรตฺถมฺปิ ฯเปฯ อุภยตฺถมฺปิ ตสฺมิํ สมเย ยถาภูตํ น ชานาติ น ปสฺสติ, ทีฆรตฺตํ สชฺฌายกตาปิ มนฺตา นปฺปฏิภนฺติ, ปเคว อสชฺฌายกตา.\csromanspacer{} เสยฺยถาปิ พฺรหฺมณ อุทปตฺโต อคฺคินา สนฺตตฺโต ปกฺกุถิโต\footnote{ปกฺกุธิโต (ก), อุกฺกฏฺฐิโต (สี), อุกฺกุฏฺฐิโต (สฺยา)}}

\csromanheader{2}{2. โพชฺฌงฺคสํยุตฺต}
\prose{\csromanfolio{107}อุสฺมุทกชาโต\footnote{อุสฺสทกชาโต (สี), อุสฺมาทกชาโต (สฺยา)}, ตตฺถ จกฺขุมา ปุริโส สกํ มุขนิมิตฺตํ ปจฺจเวกฺขมาโน ยถาภูตํ น ชาเนยฺย น ปสฺเสยฺย.\csromanspacer{} เอวเมว โข พฺราหฺมณ ยสฺมิํ สมเย พฺยาปาทปริยุฏฺฐิเตน เจตสา วิหรติ พฺยาปาทปเรเตน.\csromanspacer{} อุปฺปนฺนสฺส จ พฺยาปาทสฺส นิสฺสรณํ ยถาภูตํ นปฺปชานาติ.\csromanspacer{} อตฺตตฺถมฺปิ ตสฺมิํ สมเย ยถาภูตํ น ชานาติ น ปสฺสติ, ปรตฺถมฺปิ ตสฺมิํ สมเย ฯเปฯ อุภยตฺถมฺปิ ตสฺมิํ สมเย ยถาภูตํ น ชานาติ น ปสฺสติ, ทีฆรตฺตํ สชฺฌายกตาปิ มนฺตา นปฺปฏิภนฺติ, ปเคว อสชฺฌายกตา.}

\prose{ปุน จปรํ พฺราหฺมณ ยสฺมิํ สมเย ถินมิทฺธปริยุฏฺฐิเตน เจตสา วิหรติ ถินมิทฺธปเรเตน, อุปฺปนฺนสฺส จ ถินมิทฺธสฺส นิสฺสรณํ ยถาภูตํ นปฺปชานาติ.\csromanspacer{} อตฺตตฺถมฺปิ ตสฺมิํ สมเย ยถาภูตํ น ชานาติ น ปสฺสติ, ปรตฺถมฺปิ ฯเปฯ อุภยตฺถมฺปิ ตสฺมิํ สมเย ยถาภูตํ น ชานาติ น ปสฺสติ, ทีฆรตฺตํ สชฺฌายกตาปิ มนฺตา นปฺปฏิภนฺติ, ปเคว อสชฺฌายกตา.}

\prose{เสยฺยถาปิ พฺราหฺมณ อุทปตฺโต เสวาลปณกปริโยนทฺโธ, ตตฺถ จกฺขุมา ปุริโส สกํ มุขนิมิตฺตํ ปจฺจเวกฺขมาโน ยถาภูตํ น ชาเนยฺย น ปสฺเสยฺย.\csromanspacer{} เอวเมว โข พฺราหฺมณ ยสฺมิํ สมเย ถินมิทฺธปริยุฏฺฐิเตน เจตสา วิหรติ ถินมิทฺธปเรเตน, อุปฺปนฺนสฺส จ ถินมิทฺธสฺส นิสฺสรณํ ยถาภูตํ นปฺปชานาติ.\csromanspacer{} อตฺตตฺถมฺปิ ตสฺมิํ สมเย ยถาภูตํ น ชานาติ น ปสฺสติ, ปรตฺถมฺปิ ฯเปฯ อุภยตฺถมฺปิ ตสฺมิํ สมเย ยถาภูตํ น ชานาติ น ปสฺสติ, ทีฆรตฺตํ สชฺฌายกตาปิ มนฺตา นปฺปฏิภนฺติ, ปเคว อสชฺฌายกตา.}

\prose{ปุน จปรํ พฺราหฺมณ ยสฺมิํ สมเย อุทฺธจฺจกุกฺกุจฺจปริยุฏฺฐิเตน เจตสา วิหรติ อุทฺธจฺจกุกฺกุจฺจปเรเตน, อุปฺปนฺนสฺส จ อุทฺธจฺจกุกฺกุจฺจสฺส นิสฺสรณํ ยถาภูตํ นปฺปชานาติ.\csromanspacer{} อตฺตตฺถมฺปิ ตสฺมิํ สมเย ยถาภูตํ น ชานาติ น ปสฺสติ, ปรตฺถมฺปิ ฯเปฯ อุภยตฺถมฺปิ ตสฺมิํ สเมเย ยถาภูตํ น ชานาติ น ปสฺสติ, ทีฆรตฺตํ สชฺฌายกตาปิ มนฺตา นปฺปฏิภนฺติ, ปเคว อสชฺฌายกตา.}

\prose{เสยฺยถาปิ พฺราหฺมณ อุทปตฺโต วาเตริโต จลิโต ภนฺโต อูมิชาโต, ตตฺถ จกฺขุมา ปุริโส สกํ มุขนิมิตฺตํ ปจฺจเวกฺขมาโน}

\prose{\csromanfolio{108}ยถาภูตํ น ชาเนยฺย น ปสฺเสยฺย.\csromanspacer{} เอวเมว โข พฺราหฺมณ ยสฺมิํ สมเย อุทฺธจฺจกุกฺกุจฺจปริยุฏฺฐิเตน เจตสา วิหรติ อุทฺธจฺจกุกฺกุจฺจปเรเตน, อุปฺปนฺนสฺส จ อุทฺธจฺจกุกฺกุจฺจสฺส นิสฺสรณํ ยถาภูตํ นปฺปชานาติ.\csromanspacer{} อตฺตตฺถมฺปิ ตสฺมิํ สมเย ยถาภูตํ น ชานาติ น ปสฺสติ, ปรตฺถมฺปิ ฯเปฯ อุภยตฺถมฺปิ ตสฺมิํ สมเย ยถาภูตํ น ชานาติ น ปสฺสติ, ทีฆรตฺตํ สชฺฌายกตาปิ มนฺตา นปฺปฏิภนฺติ, ปเคว อสชฺฌายกตา.}

\prose{ปุน จปรํ พฺราหฺมณ ยสฺมิํ สมเย วิจิกิจฺฉาปริยุฏฺฐิเตน เจตสา วิหรติ วิจิกิจฺฉาปเรเตน, อุปฺปนฺนาย จ วิจิกิจฺฉาย นิสฺสรณํ ยถาภูตํ นปฺปชานาติ.\csromanspacer{} อตฺตตฺถมฺปิ ตสฺมิํ สมเย ยถาภูตํ น ชานาติ น ปสฺสติ, ปรตฺถมฺปิ ฯเปฯ อุภยตฺถมฺปิ ฯเปฯ ทีฆรตฺตํ สชฺฌายกตาปิ มนฺตา นปฺปฏิภนฺติ, ปเคว อสชฺฌยกตา.}

\prose{เสยฺยถาปิ พฺราหฺมณ อุทปตฺโต อาวิโล ลุฬิโต กลลีภูโต อนฺธกาเร นิกฺขิตฺโต, ตตฺถ จกฺขุมา ปุริโส สกํ มุขนิมิตฺตํ ปจฺจเวกฺขมาโน ยถาภูตํ น ชาเนยฺย น ปสฺเสยฺย.\csromanspacer{} เอวเมว โข พฺราหฺมณ ยสฺมิํ สมเย วิจิกิจฺฉาปริยุฏฺฐิเตน เจตสา วิหรติ วิจิกิจฺฉาปเรเตน, อุปฺปนฺนาย จ วิจิกิจฺฉาย นิสฺสรณํ ยถาภูตํ นปฺปชานาติ.\csromanspacer{} อตฺตตฺถมฺปิ ตสฺมิํ สมเย ยถาภูตํ น ชานาติ น ปสฺสติ, ปรตฺถมฺปิ ตสฺมิํ สมเย ยถาภูตํ น ชานาติ น ปสฺสติ, อุภยตฺถมฺปิ ตสฺมิํ สมเย ยถาภูตํ น ชานาติ น ปสฺสติ, ทีฆรตฺตํ สชฺฌายกตาปิ มนฺตา นปฺปฏิภนฺติ, ปเคน อสชฺฌายกตา.\csromanspacer{} อยํ โข พฺราหฺมณ เหตุ อยํ ปจฺจโย, เยเนกทา ทีฆรตฺตํ สชฺฌายกตาปิ มนฺตา นปฺปฏิภนฺติ, ปเคว อสชฺฌายกตา.}

\prose{ยสฺมิํ จ โข พฺราหฺมณ สมเย น กามราคปริยุฏฺฐิเตน เจตสา วิหรติ น กามราคปเรเตน, อุปฺปนฺนสฺส จ กามราคสฺส นิสฺสรณํ ยถาภูตํ ปชานาติ.\csromanspacer{} อตฺตตฺถมฺปิ ตสฺมิํ สมเย ยถาภูตํ ชานาติ ปสฺสติ, ปรตฺถมฺปิ ตสฺมิํ สมเย ยถาภูตํ ชานาติ ปสฺสติ, อุภยตฺถมฺปิ ตสฺมิํ สมเย ยถาภูตํ ชานาติ ปสฺสติ, ทีฆรตฺตํ อสชฺฌายกตาปิ มนฺตา ปฏิภนฺติ, ปเคว สชฺฌายกตา.}

\csromanheader{2}{2. โพชฺฌงฺคสํยุตฺต}
\prose{\csromanfolio{109}เสยฺยถาปิ พฺราหฺมณ อุทปตฺโต อสํสฏฺโฐ ลาขาย วา หลิทฺทิยา วา นีลิยา วา มญฺชิฏฺฐาย วา, ตตฺถ จกฺขุมา ปุริโส สกํ มุขนิมิตฺตํ ปจฺจเวกฺขมาโน ยถาภูตํ ชาเนยฺย ปสฺเสยฺย.\csromanspacer{} เอวเมว โข พฺราหฺมณ ยสฺมิํ สมเย น กามราคปริยุฏฺฐิเตน เจตสา วิหรติ น กามราคปเรเตน, อุปฺปนฺนสฺส จ กามราคสฺส นิสฺสรณํ ยถาภูตํ ปชานาติ ฯเปฯ.}

\prose{ปุน จปรํ พฺราหฺมณ ยสฺมิํ สมเย น พฺยาปาทปริยุฏฺฐิเตน เจตสา วิหรติ น พฺยาปาทปเรเตน, อุปฺปนฺนสฺส จ พฺยาปาทสฺส นิสฺสรณํ ยถาภูตํ ปชานาติ.\csromanspacer{} อตฺตตฺถมฺปิ ตสฺมิํ สมเย ยถาภูตํ ชานาติ ปสฺสติ, ปรตฺถมฺปิ ฯเปฯ อุภยตฺถมฺปิ ฯเปฯ ทีฆรตฺตํ อสชฺฌายกตาปิ มนฺตา ปฏิภนฺติ, ปเคว สชฺฌายกตา.}

\prose{เสยฺยถาปิ พฺราหฺมณ อุทปตฺโต น อคฺคินา สนฺตตฺโต น ปกฺกุถิโต น อุสฺมุทกชาโต, ตตฺถ จกฺขุมา ปุริโส สกํ มุขนิมิตฺตํ ปจฺจเวกฺขมาโน ยถาภูตํ ชาเนยฺย ปสฺเสยฺย.\csromanspacer{} เอวเมว โข พฺราหฺมณ ยสฺมิํ สมเย น พฺยาปาทปริยุฏฺฐิเตน เจตสา วิหรติ น พฺยาปาทปเรเตน, อุปฺปนฺนสฺส จ พฺยาปาทสฺส นิสฺสรณํ ยถาภูตํ ปชานาติ.\csromanspacer{} อตฺตตฺถมฺปิ ตสฺมิํ สมเย ยถาภูตํ ชานาติ ปสฺสติ, ปรตฺถมฺปิ ฯเปฯ อุภยตฺถมฺปิ ฯเปฯ ทีฆรตฺตํ อสชฺฌายกตาปิ มนฺตา ปฏิภนฺติ, ปเคว สชฺฌายกตา.}

\prose{ปุน จปรํ พฺราหฺมณ ยสฺมิํ สมเย น ถินมิทฺธปริยุฏฺฐิเตน เจตสา วิหรติ น ถินมิทฺธปเรเตน, อุปฺปนฺนสฺส จ ถินมิทฺธสฺส นิสฺสรณํ ยถาภูตํ ปชานาติ.\csromanspacer{} อตฺตตฺถมฺปิ ตสฺมิํ สมเย ยถาภูตํ ชานาติ ปสฺสติ, ปรตฺถมฺปิ ฯเปฯ อุภยตฺถมฺปิ ฯเปฯ ทีฆรตฺตํ อสชฺฌายกตาปิ มนฺตา ปฏิภนฺติ, ปเคว สชฺฌายกตา.}

\prose{เสยฺยถาปิ พฺราหฺมณ อุทปตฺโต น เสวาลปณกปริโยนทฺโธ, ตตฺถ จกฺขุมา ปุริโส สกํ มุขนิมิตฺตํ ปจฺจเวกฺขมาโน ยถาภูตํ ชาเนยฺย ปสฺเสยฺย.\csromanspacer{} เอวเมว โข พฺราหฺมณ ยสฺมิํ สมเย น ถินมิทฺธปริยุฏฺฐิเตน เจตสา วิหรติ น ถินมิทฺธปเรเตน, อุปฺปนฺนสฺส จ ถินมิทฺธสฺส นิสฺสรณํ ยถาภูตํ ปชานาติ.\csromanspacer{} อตฺตตฺถมฺปิ ตสฺมิํ สมเย ยถาภูตํ ชานาติ ปสฺสติ, ปรตฺถมฺปิ ฯเปฯ อุภยตฺถมฺปิ ฯเปฯ ทีฆรตฺตํ อสชฺฌายกตาปิ มนฺตา ปฏิภนฺติ, ปเคว สชฺฌายกตา.}

\prose{\csromanfolio{110}ปุน จปรํ พฺราหฺมณ ยสฺมิํ สมเย น อุทฺธจฺจกุกฺกุจฺจปริยุฏฺฐิเตน เจตสา วิหรติ น อุทฺธจฺจกุกฺกุจฺจปเรเตน, อุปฺปนฺนสฺส จ อุทฺธจฺจกุกฺกุจฺจสฺส นิสฺสรณํ ยถาภูตํ ปชานาติ.\csromanspacer{} อตฺตตฺถมฺปิ ตสฺมิํ สมเย ยถาภูตํ ชานาติ ปสฺสติ, ปรตฺถมฺปิ ฯเปฯ อุภยตฺถมฺปิ ฯเปฯ ทีฆรตฺตํ อสชฺฌายกตาปิ มนฺตา ปฏิภนฺติ, ปเคว สชฺฌายกตา.}

\prose{เสยฺยถาปิ พฺราหฺมณ อุทปตฺโต น วาเตริโต น จลิโต น ภนฺโต น อูมิชาโต, ตตฺถ จกฺขุมา ปุริโส สกํ มุขนิมิตฺตํ ปจฺจเวกฺขมาโน ยถาภูตํ ชาเนยฺย ปสฺเสยฺย.\csromanspacer{} เอวเมว โข พฺราหฺมณ ยสฺมิํ สมเย น อุทฺธจฺจกุกฺกุจฺจปริยุฏฺฐิเตน เจตสา วิหรติ น อุทฺธจฺจกุกฺกุจฺจปเรเตน, อุปฺปนฺนสฺส จ อุทฺธจฺจกุกฺกุจฺจสฺส นิสฺสรณํ ยถาภูตํ ปชานาติ.\csromanspacer{} อตฺตตฺถมฺปิ ตสฺมิํ สมเย ยถาภูตํ ชานาติ ปสฺสติ, ปรตฺถมฺปิ ฯเปฯ อุภยตฺถมฺปิ ฯเปฯ ทีฆรตฺตํ อสชฺฌายกตาปิ มนฺตา ปฏิภนฺติ, ปเคว สชฺฌายกตา.}

\prose{ปุน จปรํ พฺราหฺมณ ยสฺมิํ สมเย น วิจิกิจฺฉาปริยุฏฺฐิเตน เจตสา วิหรติ น วิจิกิจฺฉาปเรเตน, อุปฺปนฺนาย จ วิจิกิจฺฉาย นิสฺสรณํ ยถาภูตํ ปชานาติ\footnote{ปชานาติ ปสฺสติ (สฺยา)}.\csromanspacer{} อตฺตตฺถมฺปิ ตสฺมิํ สมเย ยถาภูตํ ชานาติ ปสฺสติ, ปรตฺถมฺปิ ตสฺมิํ สมเย ยถาภูตํ ชานาติ ปสฺสติ, อุภยตฺถมฺปิ ตสฺมิํ สมเย ยถาภูตํ ชานาติ ปสฺสติ, ทีฆรตฺตํ อสชฺฌายกตาปิ มนฺตา ปฏิภนฺติ, ปเคว สชฺฌายกตา.}

\prose{เสยฺยถาปิ พฺราหฺมณ อุทปตฺโต อจฺโฉ วิปฺปสนฺโน อนาวิโล อาโลเก นิกฺขิตฺโต, ตตฺถ จกฺขุมา ปุริโส สกํ มุขนิมิตฺตํ ปจฺจเวกฺขมาโน ยถาภูตํ ชาเนยฺย ปสฺเสยฺย.\csromanspacer{} เอวเมว โข พฺราหฺมณ ยสฺมิํ สมเย น วิจิกิจฺฉาปริยุฏฺฐิเตน เจตสา วิหรติ น วิจิกิจฺฉาปเรเตน, อุปฺปนฺนาย จ วิจิกิจฺฉาย นิสฺสรณํ ยถาภูตํ ปชานาติ, อตฺตตฺถมฺปิ ตสฺมิํ สมเย ยถาภูตํ ชานาติ ปสฺสติ, ปรตฺถมฺปิ ตสฺมิํ สมเย ยถาภูตํ ชานาติ ปสฺสติ, อุภยตฺถมฺปิ ตสฺมิํ สมเย ยถาภูตํ ชานาติ ปสฺสติ, ทีฆรตฺตํ อสชฺฌายกตาปิ มนฺตา ปฏิภนฺติ, ปเคว สชฺฌายกตา.\csromanspacer{} อยํ โข พฺราหฺมณ เหตุ อยํ ปจฺจโย, เยเนกทา ทีฆรตฺตํ อสชฺฌายกตาปิ มนฺตา ปฏิภนฺติ, ปเคว สชฺฌายกตา.}

\csromanheader{2}{2. โพชฺฌงฺคสํยุตฺต}
\prose{\csromanfolio{111}สตฺติเม พฺราหฺมณ โพชฺฌงฺคา อนาวรณา อนีวรณา เจตโส อนุปกฺกิเลสา ภาวิตา พหุลีกตา วิชฺชาวิมุตฺติผลสจฺฉิกิริยาย สํวตฺตนฺติ.\csromanspacer{} กตเม สตฺต, สติสมฺโพชฺฌงฺโค โข พฺราหฺมณ อนาวรโณ อนีวรโณ เจตโส อนุปกฺกิเลโส ภาวิโต พหุลีกโต วิชฺชาวิมุตฺติผลสจฺฉิกิริยาย สํวตฺตติ ฯเปฯ อุเปกฺขาสมฺโพชฺฌงฺโค โข พฺราหฺมณ อนาวรโณ อนีวรโณ เจตโส อนุปกฺกิเลโส ภาวิโต พหุลีกโต วิชฺชาวิมุตฺติผลสจฺฉิ กิริยาย สํวตฺตติ.\csromanspacer{} อิเม โข พฺราหฺมณ สตฺต โพชฺฌงฺคา อนาวรณา อนีวรณา เจตโส อนุปกฺกิเลสา ภาวิตา พหุลีกตา วิชฺชาวิมุตฺติผลสจฺฉิกิริยาย สํวตฺตนฺตีติ.\csromanspacer{} เอวํ วุตฺเต สงฺคารโว พฺราหฺมโณ ภควนฺตํ เอตทโวจ “อภิกฺกนฺตํ โภ โคตม ฯเปฯ อุปาสกํ มํ ภวํ โคตโม ธาเรตุ อชฺชตคฺเค ปาณุเปตํ สรณํ คตนฺ”ติ.\csromanspacer{} ปญฺจมํ.}

\csromanmatikaanchor{mtk.18}
\csromantocmarkref{paragraph}{6. อภยสุตฺต}{mtk.18}
\bookmark[dest={mtk.18},level=4]{6. อภยสุตฺต}
\csromanheader{4}{6. อภยสุตฺต}
\proseitem{237}{เอวํ เม สุตํ\csromandash{}เอกํ สมยํ ภควา ราชคเห วิหรติ คิชฺฌกูเฏ ปพฺพเต.\csromanspacer{} อถ โข อภโย ราชกุมาโร เยน ภควา เตนุปสงฺกมิ, อุปสงฺกมิตฺวา ภควนฺตํ อภิวาเทตฺวา เอกมนฺตํ นิสีทิ, เอกมนฺตํ นิสินฺโน โข อภโย ราชกุมาโร ภควนฺตํ เอตทโวจ “ปูรโณ ภนฺเต กสฺสโป เอวมาห นตฺถิ เหตุ นตฺถิ ปจฺจโย อญฺญาณาย อทสฺสนาย, อเหตุ อปฺปจฺจโย\footnote{อปฺปจฺจยา (สี), อปฺปจฺจยํ (?)} อญฺญาณํ อทสฺสนํ โหติ.\csromanspacer{} นตฺถิ เหตุ นตฺถิ ปจฺจโย ญาณาย ทสฺสนาย, อเหตุ อปฺปจฺจโย ญาณํ ทสฺสนํ โหตี”ติ.\csromanspacer{} อิธ ภควา กิมาหาติ.\csromanspacer{} อตฺถิ ราชกุมาร เหตุ อตฺถิ ปจฺจโย อญฺญาณาย อทสฺสนาย, สเหตุ สปฺปจฺจโย\footnote{สปฺปจฺจยา (สี), สปฺปจฺจยํ (?)} อญฺญาณํ อทสฺสนํ โหติ.\csromanspacer{} อตฺถิ ราชกุมาร เหตุ อตฺถิ ปจฺจโย ญาณาย ทสฺสนาย, สเหตุ สปฺปจฺจโย ญาณํ ทสฺสนํ โหตีติ.}

\prose{กตโม ปน ภนฺเต เหตุ กตโม ปจฺจโย อญฺญาณาย อทสฺสนาย, กถํ สเหตุ สปฺปจฺจโย อญฺญาณํ อทสฺสนํ โหตีติ.\csromanspacer{} ยสฺมิํ โข ราชกุมาร สมเย กามราคปริยุฏฺฐิเตน เจตสา วิหรติ กามราคปเรเตน, อุปฺปนฺนสฺส จ กามราคสฺส นิสฺสรณํ ยถาภูตํ น ชานาติ น ปสฺสติ.\csromanspacer{} อยมฺปิ โข ราชกุมาร เหตุ \csromanfolio{112}อยํ ปจฺจโย อญฺญาณาย อทสฺสนาย, เอวมฺปิ สเหตุ สปฺปจฺจโย อญฺญาณํ อทสฺสนํ โหติ.}

\prose{ปุน จปรํ ราชกุมาร ยสฺมิํ สมเย พฺยาปาทปริยุฏฺฐิเตน เจตสา วิหรติ พฺยาปาทปเรเตน ฯเปฯ ถินมิทฺธปริยุฏฺฐิเตน ฯเปฯ อุทฺธจฺจกุกฺกุจฺจปริยุฏฺฐิเตน ฯเปฯ วิจิกิจฺฉาปริยุฏฺฐิเตน เจตสา วิหรติ วิจิกิจฺฉาปเรเตน, อุปฺปนฺนาย จ วิจิกิจฺฉาย นิสฺสรณํ ยถาภูตํ น ชานาติ น ปสฺสติ.\csromanspacer{} อยมฺปิ โข ราชกุมาร เหตุ อยํ ปจฺจโย อญฺญาณาย อทสฺสนาย, เอวมฺปิ สเหตุ สปฺปจฺจโย อญฺญาณํ อทสฺสนํ โหตีติ.}

\prose{โก นามายํ ภนฺเต ธมฺมปริยาโยติ.\csromanspacer{} นีวรณา นาเมเต ราชกุมาราติ.\csromanspacer{} ตคฺฆ ภควา นีวรณา, ตคฺฆ สุคต นีวรณา, เอกเมเกนปิ โข ภนฺเต นีวรเณน อภิภูโต ยถาภูตํ น ชาเนยฺย น ปสฺเสยฺย, โก ปน วาโท ปญฺจหิ นีวรเณหิ.}

\prose{กตโม ปน ภนฺเต เหตุ กตโม ปจฺจโย ญาณาย ทสฺสนาย, กถํ สเหตุ สปฺปจฺจโย ญาณํ ทสฺสนํ โหตีติ.\csromanspacer{} อิธ ราชกุมาร ภิกฺขุ สติสมฺโพชฺฌงฺคํ ภาเวติ วิเวกนิสฺสิตํ วิราคนิสฺสิตํ นิโรธนิสฺสิตํ โวสฺสคฺคปริณามิํ, โส สติสมฺโพชฺฌงฺคํ ภาวิเตน จิตฺเตน ยถาภูตํ ชานาติ ปสฺสติ.\csromanspacer{} อยมฺปิ โข ราชกุมาร เหตุ อยํ ปจฺจโย ญาณาย ทสฺสนาย, เอวมฺปิ สเหตุ สปฺปจฺจโย ญาณํ ทสฺสนํ โหติ.}

\prose{ปุน จปรํ ราชกุมาร ภิกฺขุ ฯเปฯ อุเปกฺขาสมฺโพชฺฌงฺคํ ภาเวติ วิเวกนิสฺสิตํ วิราคนิสฺสิตํ นิโรธนิสฺสิตํ โวสฺสคฺคปริณามิํ, โส อุเปกฺขาสมฺโพชฺฌงฺคํ ภาวิเตน จิตฺเตน ยถาภูตํ ชานาติ ปสฺสติ.\csromanspacer{} อยมฺปิ โข ราชกุมาร เหตุ อยํ ปจฺจโย ญาณาย ทสฺสนาย, เอวํ สเหตุ สปฺปจฺจโย ญาณํ ทสฺสนํ โหตีติ.}

\prose{โก นามายํ ภนฺเต ธมฺมปริยาโยติ.\csromanspacer{} โพชฺฌงฺคา นาเมเต ราชกุมาราติ.\csromanspacer{} ตคฺฆ ภควา โพชฺฌงฺคา, ตคฺฆ สุคต โพชฺฌงฺคา, เอกเมเกนปิ โข ภนฺเต โพชฺฌงฺเคน สมนฺนาคโต ยถาภูตํ ชาเนยฺย ปสฺเสยฺย, โก ปน วาโท สตฺตหิ โพชฺฌงฺเคหิ.\csromanspacer{} โยปิ เม ภนฺเต คิชฺฌกูฏํ ปพฺพตํ}

\csromanheader{2}{2. โพชฺฌงฺคสํยุตฺต}
\csromanmatikaanchor{mtk.19}
\csromanmatikaanchor{mtk.20}
\csromantocmarknopage{paragraph}{7. อานาปานวคฺค}{mtk.19}
\bookmark[dest={mtk.19},level=4]{7. อานาปานวคฺค}
\csromantocmarkref{subparagraph}{1-10. อฏฺฐิกมหปฺผลสุตฺตาทีนิ}{mtk.20}
\bookmark[dest={mtk.20},level=5]{1-10. อฏฺฐิกมหปฺผลสุตฺตาทีนิ}
\prose{\csromanfolio{113}อาโรหนฺตสฺส กายกิลมโถ จิตฺตกิลมโถ, โสปิ เม ปฏิปฺปสฺสทฺโธ, ธมฺโม จ เม อภิสมิโตติ.\csromanspacer{}\csromanspacer{}ฉฏฺฐํ.\csromanspacer{} สากจฺฉวคฺโค ฉฏฺโฐ.}

\titlehead{\textbf{ตสฺสุทฺทานํ}}
\csromangathameasure{อาหารา ปริยายมคฺคิ, เมตฺตํ สงฺคารเวน จ. \\ อภโย ปุจฺฉิโต ปญฺหํ, คิชฺฌกูฏมฺหิ ปพฺพเตติ.}
\csromangathagroup{\csromangathastanza{อาหารา ปริยายมคฺคิ, เมตฺตํ สงฺคารเวน จ. \\ อภโย ปุจฺฉิโต ปญฺหํ, คิชฺฌกูฏมฺหิ ปพฺพเตติ.}}

\csromanheader{4}{7. อานาปานวคฺค}
\csromanheader{2}{1. อฏฺฐิกมหปฺผลสุตฺต}
\proseitem{238}{สาวตฺถินิทานํ.\csromanspacer{} อฏฺฐิกสญฺญา ภิกฺขเว ภาวิตา พหุลีกตา มหปฺผลา โหติ มหานิสํสา.\csromanspacer{} กถํ ภาวิตา จ ภิกฺขเว อฏฺฐิกสญฺญา กถํ พหุลีกตา มหปฺผลา โหติ มหานิสํสา, อิธ ภิกฺขเว ภิกฺขุ อฏฺฐิกสญฺญาสหคตํ สติสมฺโพชฺฌงฺคํ ภาเวติ วิเวกนิสฺสิตํ วิราคนิสฺสิตํ นิโรธนิสฺสิตํ โวสฺสคฺคปริณามิํ ฯเปฯ อฏฺฐิกสญฺญาสหคตํ อุเปกฺขาสมฺโพชฺฌงฺคํ ภาเวติ วิเวกนิสฺสิตํ วิราคนิสฺสิตํ นิโรธนิสฺสิตํ โวสฺสคฺคปริณามิํ.\csromanspacer{} เอวํ ภาวิตา โข ภิกฺขเว อฏฺฐิกสญฺญา เอวํ พหุลีกตา มหปฺผลา โหติ มหานิสํสาติ.}

\csromanheader{2}{\textbf{อญฺญตรผลสุตฺต}}
\prose{อฏฺฐิกสญฺญาย ภิกฺขเว ภาวิตาย พหุลีกตาย ทฺวินฺนํ ผลานํ อญฺญตรํ ผลํ ปาฏิกงฺขํ, ทิฏฺเฐว ธมฺเม อญฺญา, สติ วา อุปาทิเสเส อนาคามิตา.\csromanspacer{} กถํ ภาวิตาย จ โข ภิกฺขเว อฏฺฐิกสญฺญาย กถํ พหุลีกตาย ทฺวินฺนํ ผลานํ อญฺญตรํ ผลํ ปาฏิกงฺขํ ทิฏฺเฐว ธมฺเม อญฺญา, สติ วา อุปาทิเสเส อนาคามิตา, อิธ ภิกฺขเว ภิกฺขุ อฏฺฐิกสญฺญาสหคตํ สติสมฺโพชฺฌงฺคํ ภาเวติ ฯเปฯ อฏฺฐิกสญฺญาสหคตํ อุเปกฺขาสมฺโพชฺฌงฺคํ ภาเวติ วิเวกนิสฺสิตํ วิราคนิสฺสิตํ นิโรธนิสฺสิตํ โวสฺสคฺคปริณามิํ.\csromanspacer{} เอวํ ภาวิตาย โข ภิกฺขเว อฏฺฐิกสญฺญาย เอวํ}

\prose{\csromanfolio{114}พหุลีกตาย ทฺวินฺนํ ผลานํ อญฺญตรํ ผลํ ปาฏิกงฺขํ ทิฏฺเฐว ธมฺเม อญฺญา, สติ วา อุปาทิเสเส อนาคามิตาติ.}

\csromanheader{2}{\textbf{มหตฺถสุตฺต}}
\prose{อฏฺฐิกสญฺญา ภิกฺขเว ภาวิตา พหุลีกตา มหโต อตฺถาย สํวตฺตติ.\csromanspacer{} กถํ ภาวิตา จ ภิกฺขเว อฏฺฐิกสญฺญา กถํ พหุลีกตา มหโต อตฺถาย สํวตฺตติ, อิธ ภิกฺขเว ภิกฺขุ อฏฺฐิกสญฺญาสหคตํ สติสมฺโพชฺฌงฺคํ ภาเวติ ฯเปฯ อฏฺฐิกสญฺญาสหคตํ อุเปกฺขาสมฺโพชฺฌงฺคํ ภาเวติ วิเวกนิสฺสิตํ วิราคนิสฺสิตํ นิโรธนิสฺสิตํ โวสฺสคฺคปริณามิํ.\csromanspacer{} เอวํ ภาวิตา โข ภิกฺขเว อฏฺฐิกสญฺญา เอวํ พหุลีกตา มหโต อตฺถาย สํวตฺตตีติ.}

\csromanheader{2}{\textbf{โยคกฺเขมสุตฺต}}
\prose{อฏฺฐิกสญฺญา ภิกฺขเว ภาวิตา พหุลีกตา มหโต โยคกฺเขมาย สํวตฺตติ.\csromanspacer{} กถํ ภาวิตา จ ภิกฺขเว อฏฺฐิกสญฺญา กถํ พหุลีกตา มหโต โยคกฺเขมาย สํวตฺตติ, อิธ ภิกฺขเว ภิกฺขุ อฏฺฐิกสญฺญาสหคตํ สติสมฺโพชฺฌงฺคํ ภาเวติ ฯเปฯ อฏฺฐิกสญฺญาสหคตํ อุเปกฺขาสมฺโพชฺฌงฺคํ ภาเวติ วิเวกนิสฺสิตํ วิราคนิสฺสิตํ นิโรธนิสฺสิตํ โวสฺสคฺคปริณามิํ.\csromanspacer{} เอวํ ภาวิตา โข ภิกฺขเว อฏฺฐิกสญฺญา เอวํ พหุลีกตา มหโต โยคกฺเขมาย สํวตฺตตีติ.}

\csromanheader{2}{\textbf{สํเวคสุตฺต}}
\prose{อฏฺฐิกสญฺญา ภิกฺขเว ภาวิตา พหุลีกตา มหโต สํเวคาย สํวตฺตติ.\csromanspacer{} กถํ ภาวิตา จ ภิกฺขเว อฏฺฐิกสญฺญา กถํ พหุลีกตา มหโต สํเวคาย สํวตฺตติ, อิธ ภิกฺขเว ภิกฺขุ อฏฺฐิกสญฺญาสหคตํ สติสมฺโพชฺฌงฺคํ ภาเวติ ฯเปฯ อฏฺฐิกสญฺญาสหคตํ อุเปกฺขาสมฺโพชฺฌงฺคํ ภาเวติ วิเวกนิสฺสิตํ วิราคนิสฺสิตํ นิโรธนิสฺสิตํ โวสฺสคฺคปริณามิํ.\csromanspacer{} เอวํ ภาวิตา โข ภิกฺขเว อฏฺฐิกสญฺญา เอวํ พหุลีกตา มหโต สํเวคาย สํวตฺตตีติ.}

\csromanheader{2}{2. โพชฺฌงฺคสํยุตฺต}
\csromanheader{2}{\textbf{ผาสุวิหารสุตฺต}}
\prose{\csromanfolio{115}อฏฺฐิกสญฺญา ภิกฺขเว ภาวิตา พหุลีกตา มหโต ผาสุวิหาราย สํวตฺตติ.\csromanspacer{} กถํ ภาวิตา จ ภิกฺขเว อฏฺฐิกสญฺญา กถํ พหุลีกตา มหโต ผาสุวิหาราย สํวตฺตติ, อิธ ภิกฺขเว ภิกฺขุ อฏฺฐิกสญฺญาสหคตํ สติสมฺโพชฺฌงฺคํ ภาเวติ ฯเปฯ อฏฺฐิกสญฺญาสหคตํ อุเปกฺขาสมฺโพชฺฌงฺคํ ภาเวติ วิเวกนิสฺสิตํ วิราคนิสฺสิตํ นิโรธนิสฺสิตํ โวสฺสคฺคปริณามิํ.\csromanspacer{} เอวํ ภาวิตา โข ภิกฺขเว อฏฺฐิกสญฺญา เอวํ พหุลีกตา มหโต ผาสุวิหาราย สํวตฺตตีติ.\csromanspacer{}\csromanspacer{}ปฐมํ.}

\csromanheader{2}{2. ปุฬวกสุตฺต}
\vspace{\tipitakaparskip}
\csromancenter{ปุฬวกสญฺญา\footnote{ปุฬุวกสญฺญา (ก)} ภิกฺขเว ภาวิตา ฯเปฯ.\csromanspacer{}\csromanspacer{}ทุติยํ.}
\csromanheader{2}{3. วินีลกสุตฺต}
\vspace{\tipitakaparskip}
\csromancenter{วินีลกสญฺญา ภิกฺขเว ฯเปฯ.\csromanspacer{} ตติยํ.}
\csromanheader{2}{4. วิจฺฉิทฺทกสุตฺต}
\vspace{\tipitakaparskip}
\csromancenter{วิจฺฉิทฺทกสญฺญา ภิกฺขเว ฯเปฯ.\csromanspacer{}\csromanspacer{}จตุตฺถํ.}
\csromanheader{2}{5. อุทฺธุมาตกสุตฺต}
\vspace{\tipitakaparskip}
\csromancenter{อุทฺธุมาตกสญฺญา ภิกฺขเว ฯเปฯ.\csromanspacer{}\csromanspacer{}ปญฺจมํ.}
\csromanheader{2}{6. เมตฺตาสุตฺต}
\vspace{\tipitakaparskip}
\csromancenter{เมตฺตา ภิกฺขเว ภาวิตา ฯเปฯ.\csromanspacer{}\csromanspacer{}ฉฏฺฐํ.}
\csromanheader{2}{7. กรุณาสุตฺต}
\vspace{\tipitakaparskip}
\csromancenter{กรุณา ภิกฺขเว ภาวิตา ฯเปฯ.\csromanspacer{}\csromanspacer{}สตฺตมํ.}
\csromanheader{2}{8. มุทิตาสุตฺต}
\vspace{\tipitakaparskip}
\csromancenter{มุทิตา ภิกฺขเว ภาวิตา ฯเปฯ.\csromanspacer{}\csromanspacer{}อฏฺฐมํ.}
\csromanheader{2}{9. อุเปกฺขาสุตฺต}
\vspace{\tipitakaparskip}
\csromancenter{อุเปกฺขา ภิกฺขเว ภาวิตา ฯเปฯ.\csromanspacer{}\csromanspacer{}นวมํ.}
\csromanheader{2}{10. อานาปานสุตฺต}
\csromanmatikaanchor{mtk.21}
\csromanmatikaanchor{mtk.22}
\csromantocmarknopage{paragraph}{8. นิโรธวคฺค}{mtk.21}
\bookmark[dest={mtk.21},level=4]{8. นิโรธวคฺค}
\csromantocmarkref{subparagraph}{1-10. อสุภสุตฺตาทีนิ}{mtk.22}
\bookmark[dest={mtk.22},level=5]{1-10. อสุภสุตฺตาทีนิ}
\csromantocmarknopage{paragraph}{9. คงฺคาเปยฺยาลวคฺค}{mtk.23}
\bookmark[dest={mtk.23},level=4]{9. คงฺคาเปยฺยาลวคฺค}
\proseitem{247}{\csromanfolio{116}อานาปานสฺสติ ภิกฺขเว ภาวิตา ฯเปฯ.\csromanspacer{}\csromanspacer{}ทสมํ.\csromanspacer{} อานาปานวคฺโค สตฺตโม.}

\titlehead{\textbf{ตสฺสุทฺทานํ}}
\csromangathameasure{อฏฺฐิก ปุฬวกํ วินีลกํ, วิจฺฉิทฺทกํ อุทฺธุมาเตน ปญฺจมํ. \\ เมตฺตา กรุณา มุทิตา, อุเปกฺขา อานาปาเนน เต ทสาติ.}
\csromangathagroup{\csromangathastanza{อฏฺฐิก ปุฬวกํ วินีลกํ, วิจฺฉิทฺทกํ อุทฺธุมาเตน ปญฺจมํ. \\ เมตฺตา กรุณา มุทิตา, อุเปกฺขา อานาปาเนน เต ทสาติ.}}

\csromanheader{4}{8. นิโรธวคฺค}
\csromanheader{2}{1. อสุภสุตฺต}
\vspace{\tipitakaparskip}
\csromancenter{อสุภสญฺญา ภิกฺขเว ฯเปฯ.\csromanspacer{}\csromanspacer{}ปฐมํ.}
\csromanheader{2}{2. มรณสุตฺต}
\vspace{\tipitakaparskip}
\csromancenter{มรณสญฺญา ภิกฺขเว ฯเปฯ.\csromanspacer{}\csromanspacer{}ทุติยํ.}
\csromanheader{2}{3. อาหาเรปฏิกูลสุตฺต}
\vspace{\tipitakaparskip}
\csromancenter{อาหาเร ปฏิกูลสญฺญา ภิกฺขเว ฯเปฯ.\csromanspacer{} ตติยํ.}
\csromanheader{2}{4. อนภิรติสุตฺต}
\vspace{\tipitakaparskip}
\csromancenter{สพฺพโลเก อนภิรติสญฺญา ภิกฺขเว ฯเปฯ.\csromanspacer{}\csromanspacer{}จตุตฺถํ.}
\csromanheader{2}{5. อนิจฺจสุตฺต}
\vspace{\tipitakaparskip}
\csromancenter{อนิจฺจสญฺญา ภิกฺขเว ฯเปฯ.\csromanspacer{}\csromanspacer{}ปญฺจมํ.}
\csromanheader{2}{6. ทุกฺขสุตฺต}
\vspace{\tipitakaparskip}
\csromancenter{อนิจฺเจ ทุกฺขสญฺญา ภิกฺขเว ฯเปฯ.\csromanspacer{}\csromanspacer{}ฉฏฺฐํ.}
\csromanheader{2}{2. โพชฺฌงฺคสํยุตฺต}
\csromanheader{2}{7. อนตฺตสุตฺต}
\vspace{\tipitakaparskip}
\csromancenter{ทุกฺเข อนตฺตสญฺญา ภิกฺขเว ฯเปฯ.\csromanspacer{}\csromanspacer{}สตฺตมํ.}
\csromanheader{2}{8. ปหานสุตฺต}
\vspace{\tipitakaparskip}
\csromancenter{ปหานสญฺญา ภิกฺขเว ฯเปฯ.\csromanspacer{}\csromanspacer{}อฏฺฐมํ.}
\csromanheader{2}{9. วิราคสุตฺต}
\vspace{\tipitakaparskip}
\csromancenter{วิราคสญฺญา ภิกฺขเว ฯเปฯ.\csromanspacer{}\csromanspacer{}นวมํ.}
\csromanheader{2}{10. นิโรธสุตฺต}
\proseitem{257}{\csromanfolio{117}นิโรธสญฺญา ภิกฺขเว ภาวิตา พหุลีกตา มหปฺผลา โหติ มหานิสํสา.\csromanspacer{} กถํ ภาวิตา จ ภิกฺขเว นิโรธสญฺญา กถํ พหุลีกตา มหปฺผลา โหติ มหานิสํสา, อิธ ภิกฺขเว ภิกฺขุ นิโรธสญฺญาสหคตํ สติสมฺโพชฺฌงฺคํ ภาเวติ ฯเปฯ นิโรธสญฺญาสหคตํ อุเปกฺขาสมฺโพชฺฌงฺคํ ภาเวติ วิเวกนิสฺสิตํ วิราคนิสฺสิตํ นิโรธนิสฺสิตํ โวสฺสคฺคปริณามิํ.\csromanspacer{} เอวํ ภาวิตา โข ภิกฺขเว นิโรธสญฺญา เอวํ พหุลีกตา มหปฺผลา โหติ มหานิสํสาติ.}

\prose{นิโรธสญฺญาย ภิกฺขเว ภาวิตาย พหุลีกตาย ทฺวินฺนํ ผลานํ อญฺญตรํ ผลํ ปาฏิกงฺขํ ทิฏฺเฐว ธมฺเม อญฺญา, สติ วา อุปาทิเสเส อนาคามิตา.\csromanspacer{} กถํ ภาวิตาย ภิกฺขเว นิโรธสญฺญาย กถํ พหุลีกตาย ทฺวินฺนํ ผลานํ อญฺญตรํ ผลํ ปาฏิกงฺขํ ทิฏฺเฐว ธมฺเม อญฺญา, สติ วา อุปาทิเสเส อนาคามิตา, อิธ ภิกฺขเว ภิกฺขุ นิโรธสญฺญาสหคตํ สติสมฺโพชฺฌงฺคํ ภาเวติ ฯเปฯ นิโรธสญฺญาสหคตํ อุเปกฺขาสมฺโพชฺฌงฺคํ ภาเวติ วิเวกนิสฺสิตํ วิราคนิสฺสิตํ นิโรธนิสฺสิตํ โวสฺสคฺคปริณามิํ.\csromanspacer{} เอวํ ภาวิตาย โข ภิกฺขเว นิโรธสญฺญาย เอวํ พหุลีกตาย ทฺวินฺนํ ผลานํ อญฺญตรํ ผลํ ปาฏิกงฺขํ ทิฏฺเฐว ธมฺเม อญฺญา, สติ วา อุปาทิเสเส อนาคามิตาติ.}

\prose{นิโรธสญฺญา ภิกฺขเว ภาวิตา พหุลีกตา มหโต อตฺถาย สํวตฺตติ, มหโต โยคกฺเขมาย สํวตฺตติ, มหโต สํเวคาย สํวตฺตติ, มหโต ผาสุวิหาราย สํวตฺตติ.\csromanspacer{} กถํ ภาวิตา จ ภิกฺขเว นิโรธสญฺญา กถํ พหุลีกตา มหโต อตฺถาย สํวตฺตติ, มหโต \csromanfolio{118}โยคกฺเขมาย สํวตฺตติ, มหโต สํเวคาย สํวตฺตติ, มหโต ผาสุวิหาราย สํวตฺตติ.\csromanspacer{} อิธ ภิกฺขเว ภิกฺขุ นิโรธสญฺญาสหคตํ สติสมฺโพชฺฌงฺคํ ภาเวติ ฯเปฯ นิโรธสญฺญาสหคตํ อุเปกฺขาสมฺโพชฺฌงฺคํ ภาเวติ วิเวกนิสฺสิตํ วิราคนิสฺสิตํ นิโรธนิสฺสิตํ โวสฺสคฺคปริณามิํ.\csromanspacer{} เอวํ ภาวิตา โข ภิกฺขเว นิโรธสญฺญา เอวํ พหุลีกตา มหโต อตฺถาย สํวตฺตติ, มหโต โยคกฺเขมาย สํวตฺตติ, มหโต สํเวคาย สํวตฺตติ, มหโต ผาสุวิหาราย สํวตฺตตีติ.\csromanspacer{}\csromanspacer{}ทสมํ.\csromanspacer{} นิโรธวคฺโค อฏฺฐโม.}

\titlehead{\textbf{ตสฺสุทฺทานํ}}
\csromangathameasure{อสุภมรณ-อาหาเร, ปฏิกูล-อนภิรเตน. \\ อนิจฺจทุกฺข-อนตฺตปหานํ, วิราคนิโรเธน เต ทสาติ.}
\csromangathagroup{\csromangathastanza{อสุภมรณ-อาหาเร, ปฏิกูล-อนภิรเตน\footnote{ปฏิกูเลน จ สพฺพโลเก (สฺยา)}. \\ อนิจฺจทุกฺข-อนตฺตปหานํ, วิราคนิโรเธน เต ทสาติ.}}

\chapterhead{9. คงฺคาเปยฺยาลวคฺค}
\csromanmatikaanchor{mtk.24}
\csromantocmarkref{subparagraph}{1-12. คงฺคานที-อาทิสุตฺต}{mtk.24}
\bookmark[dest={mtk.24},level=5]{1-12. คงฺคานที-อาทิสุตฺต}
\csromanheader{5}{1-12. คงฺคานที-อาทิสุตฺต}
\proseitem{258-269}{เสยฺยถาปิ ภิกฺขเว คงฺคา นที ปาจีนนินฺนา ปาจีนโปณา ปาจีนปพฺภารา.\csromanspacer{} เอวเมว โข ภิกฺขเว ภิกฺขุ สตฺต โพชฺฌงฺเค ภาเวนฺโต สตฺต โพชฺฌงฺเค พหุลีกโรนฺโต นิพฺพานนินฺโน โหติ นิพฺพานโปโณ นิพฺพานปพฺภาโร.\csromanspacer{} กถญฺจ ภิกฺขเว ภิกฺขุ สตฺต โพชฺฌงฺเค ภาเวนฺโต สตฺต โพชฺฌงฺเค พหุลีกโรนฺโต นิพฺพานนินฺโน โหติ นิพฺพานโปโณ นิพฺพานปพฺภาโร.\csromanspacer{} อิธ ภิกฺขเว ภิกฺขุ สติสมฺโพชฺฌงฺคํ ภาเวติ ฯเปฯ อุเปกฺขาสมฺโพชฺฌงฺคํ ภาเวติ วิเวกนิสฺสิตํ วิราคนิสฺสิตํ นิโรธนิสฺสิตํ โวสฺสคฺคปริณามิํ.\csromanspacer{} เอวํ โข ภิกฺขเว ภิกฺขุ สตฺต โพชฺฌงฺเค ภาเวนฺโต สตฺต โพชฺฌงฺเค พหุลีกโรนฺโต นิพฺพานนินฺโน โหติ นิพฺพานโปโณ นิพฺพานปพฺภาโรติ. (ยาว เอสนา ปาฬิ วิตฺถาเรตพฺพา.) คงฺคาเปยฺยาลวคฺโค นวโม.}

\csromanheader{2}{2. โพชฺฌงฺคสํยุตฺต}
\titlehead{\textbf{ตสฺสุทฺทานํ}}
\csromangathameasure{ฉ ปาจีนโต นินฺนา, ฉ นินฺนา จ สมุทฺทโต. \\ เทฺว เต ฉ ทฺวาทส โหนฺติ, วคฺโค เตน ปวุจฺจตีติ.}
\csromangathagroup{\csromangathastanza{\csromanfolio{119}ฉ ปาจีนโต นินฺนา, ฉ นินฺนา จ สมุทฺทโต. \\ เทฺว เต ฉ ทฺวาทส โหนฺติ, วคฺโค เตน ปวุจฺจตีติ.}}

\csromanmatikaanchor{mtk.25}
\csromantocmarknopage{paragraph}{10. อปฺปมาทวคฺค}{mtk.25}
\bookmark[dest={mtk.25},level=4]{10. อปฺปมาทวคฺค}
\csromanheader{4}{10. อปฺปมาทวคฺค}
\csromanmatikaanchor{mtk.26}
\csromantocmarkref{subparagraph}{1-10. ตถาคตาทิสุตฺต}{mtk.26}
\bookmark[dest={mtk.26},level=5]{1-10. ตถาคตาทิสุตฺต}
\csromantocmarknopage{paragraph}{11. พลกรณียวคฺค}{mtk.27}
\bookmark[dest={mtk.27},level=4]{11. พลกรณียวคฺค}
\csromanheader{5}{1-10. ตถาคตาทิสุตฺต}
\proseitem{270-279}{ยาวตา ภิกฺขเว สตฺตา อปทา วา ทฺวิปทา วา จตุปฺปทา วา พหุปฺปทา วาติ วิตฺถาเรตพฺพํ.\csromanspacer{} อปฺปมาทวคฺโค ทสโม.}

\titlehead{\textbf{ตสฺสุทฺทานํ}}
\csromangathameasure{ตถาคตํ ปทํ กูฏํ, มูลํ สาเรน วสฺสิกํ. \\ ราชา จนฺทิมสูริยา จ, วตฺเถน ทสมํ ปทนฺติ. (อปฺปมาทวคฺโค โพชฺฌงฺคสํยุตฺตสฺส โพชฺฌงฺควเสน วิตฺถาเรตพฺโพ.)}
\csromangathagroup{\csromangathastanza{ตถาคตํ ปทํ กูฏํ, มูลํ สาเรน วสฺสิกํ. \\ ราชา จนฺทิมสูริยา จ, วตฺเถน ทสมํ ปทนฺติ. (อปฺปมาทวคฺโค โพชฺฌงฺคสํยุตฺตสฺส โพชฺฌงฺควเสน วิตฺถาเรตพฺโพ.)}}

\chapterhead{11. พลกรณียวคฺค}
\csromanmatikaanchor{mtk.28}
\csromantocmarkref{subparagraph}{1-12. พลาทิสุตฺต}{mtk.28}
\bookmark[dest={mtk.28},level=5]{1-12. พลาทิสุตฺต}
\csromantocmarknopage{paragraph}{12. เอสนาวคฺค}{mtk.29}
\bookmark[dest={mtk.29},level=4]{12. เอสนาวคฺค}
\csromanheader{5}{1-12. พลาทิสุตฺต}
\proseitem{280-291}{เสยฺยถาปิ ภิกฺขเว เย เกจิ พลกรณียา กมฺมนฺตา กรียนฺตีติ วิตฺถาเรตพฺพํ.\csromanspacer{} พลกรณียวคฺโค เอกาทสโม.}

\titlehead{\textbf{ตสฺสุทฺทานํ}}
\csromangathameasure{พลํ พีชญฺจ นาโค จ, รุกฺโข กุมฺเภน สูกิยา. \\ อากาเสน จ เทฺว เมฆา, นาวา อาคนฺตุกา นทีติ. (พลกรณียวคฺโค โพชฺฌงฺคสํยุตฺตสฺส โพชฺฌงฺควเสน วิตฺถาเรตพฺโพ.)}
\csromangathagroup{\csromangathastanza{พลํ พีชญฺจ นาโค จ, รุกฺโข กุมฺเภน สูกิยา. \\ อากาเสน จ เทฺว เมฆา, นาวา อาคนฺตุกา นทีติ. (พลกรณียวคฺโค โพชฺฌงฺคสํยุตฺตสฺส โพชฺฌงฺควเสน วิตฺถาเรตพฺโพ.)}}

\chapterheadpageread{12. เอสนาวคฺค}
\csromanmatikaanchor{mtk.30}
\csromantocmarkref{subparagraph}{1-10. เอสนาทิสุตฺต}{mtk.30}
\bookmark[dest={mtk.30},level=5]{1-10. เอสนาทิสุตฺต}
\csromantocmarknopage{paragraph}{13. โอฆวคฺค}{mtk.31}
\bookmark[dest={mtk.31},level=4]{13. โอฆวคฺค}
\csromanheader{5}{1-10. เอสนาทิสุตฺต}
\proseitem{292-301}{\csromanfolio{120}ติสฺโส อิมา ภิกฺขเว เอสนา.\csromanspacer{} กตมา ติสฺโส, กาเมสนา ภเวสนา พฺรหฺมจริเยสนาติ วิตฺถาเรตพฺพํ.\csromanspacer{} เอสนาวคฺโค ทฺวาทสโม.}

\titlehead{\textbf{ตสฺสุทฺทานํ}}
\csromangathameasure{เอสนา วิธา อาสโว, ภโว จ ทุกฺขตา ติสฺโส. \\ ขิลํ มลญฺจ นีโฆ จ, เวทนา ตณฺหา ตสินาย จาติ. (โพชฺฌงฺคสํยุตฺตสฺส เอสนาเปยฺยาลํ วิเวกนิสฺสิตโต วิตฺถาเรตพฺพํ.)}
\csromangathagroup{\csromangathastanza{เอสนา วิธา อาสโว, ภโว จ ทุกฺขตา ติสฺโส. \\ ขิลํ มลญฺจ นีโฆ จ, เวทนา ตณฺหา ตสินาย จาติ. (โพชฺฌงฺคสํยุตฺตสฺส เอสนาเปยฺยาลํ วิเวกนิสฺสิตโต วิตฺถาเรตพฺพํ.)}}

\chapterhead{13. โอฆวคฺค}
\csromanmatikaanchor{mtk.32}
\csromantocmarkref{subparagraph}{1-9. โอฆาทิสุตฺต}{mtk.32}
\bookmark[dest={mtk.32},level=5]{1-9. โอฆาทิสุตฺต}
\csromanheader{5}{1-9. โอฆาทิสุตฺต}
\proseitem{302-310}{จตฺตาโรเม ภิกฺขเว โอฆา.\csromanspacer{} กตเม จตฺตาโร, กาโมโฆ ภโวโฆ ทิฏฺโฐโฆ อวิชฺโชโฆติ วิตฺถาเรตพฺพํ.}

\csromanmatikaanchor{mtk.33}
\csromantocmarkref{paragraph}{10. อุทฺธมฺภาคิยสุตฺต}{mtk.33}
\bookmark[dest={mtk.33},level=4]{10. อุทฺธมฺภาคิยสุตฺต}
\csromanheader{4}{10. อุทฺธมฺภาคิยสุตฺต}
\proseitem{311}{สาวตฺถินิทานํ.\csromanspacer{} ปญฺจิมานิ ภิกฺขเว อุทฺธมฺภาคิยานิ สํโยชนานิ.\csromanspacer{} กตมานิ ปญฺจ, รูปราโค อรูปราโค มาโน อุทฺธจฺจํ อวิชฺชา.\csromanspacer{} อิมานิ โข ภิกฺขเว ปญฺจุทฺธมฺภาคิยานิ สํโยชนานิ.\csromanspacer{} อิเมสํ โข ภิกฺขเว ปญฺจนฺนํ อุทฺธมฺภาคิยานํ สํโยชนานํ อภิญฺญาย ปริญฺญาย ปริกฺขยาย ปหานาย สตฺต โพชฺฌงฺคา ภาเวตพฺพา.\csromanspacer{} กตเม สตฺต, อิธ ภิกฺขเว ภิกฺขุ สติสมฺโพชฺฌงฺคํ ภาเวติ วิเวกนิสฺสิตํ วิราคนิสฺสิตํ นิโรธนิสฺสิตํ โวสฺสคฺคปริณามิํ ฯเปฯ อุเปกฺขาสมฺโพชฺฌงฺคํ ภาเวติ ราควินยปริโยสานํ โทสวินยปริโยสานํ โมหวินยปริโยสานํ.\csromanspacer{} อมโตคธํ อมตปรายนํ อมตปริโยสานํ.\csromanspacer{} นิพฺพานนินฺนํ นิพฺพานโปณํ นิพฺพานปพฺภารํ.\csromanspacer{} อิเมสํ โข ภิกฺขเว ภิกฺขุ ปญฺจนฺนํ อุทฺธมฺภาคิยานํ สํโยชนานํ อภิญฺญาย ปริญฺญาย ปริกฺขยาย ปหานาย อิเม สตฺต โพชฺฌงฺคา ภาเวตพฺพาติ.\csromanspacer{}\csromanspacer{}ทสมํ.\csromanspacer{} โอฆวคฺโค เตรสโม.}

\csromanheader{2}{2. โพชฺฌงฺคสํยุตฺต}
\titlehead{\textbf{ตสฺสุทฺทานํ}}
\csromanmatikaanchor{mtk.34}
\csromanmatikaanchor{mtk.35}
\csromanmatikaanchor{mtk.37}
\csromanmatikaanchor{mtk.39}
\csromantocmarknopage{subparagraph}{14. ปุนคงฺคาเปยฺยาลวคฺค}{mtk.34}
\bookmark[dest={mtk.34},level=5]{14. ปุนคงฺคาเปยฺยาลวคฺค}
\csromantocmarkref{subparagraph}{312-323. ปุนคงฺคานที-อาทิสุตฺต}{mtk.35}
\bookmark[dest={mtk.35},level=5]{312-323. ปุนคงฺคานที-อาทิสุตฺต}
\csromangathameasure{โอโฆ โยโค อุปาทานํ, คนฺถา อนุสเยน จ. \\ กามคุณา นีวรณา, ขนฺธา โอรุทฺธมฺภาคิยานีติ.}
\csromangathagroup{\csromangathastanza{\csromanfolio{121}โอโฆ โยโค อุปาทานํ, คนฺถา อนุสเยน จ. \\ กามคุณา นีวรณา, ขนฺธา โอรุทฺธมฺภาคิยานีติ.}}

\csromanheader{5}{14. ปุนคงฺคาเปยฺยาลวคฺค}
\vspace{\tipitakaparskip}
\csromancenter{ปุนคงฺคานที-อาทิสุตฺต วคฺโค จุทฺทสโม.}
\csromanheader{2}{\textbf{อุทฺทานํ}}
\csromangathameasure{ฉ ปาจีนโต นินฺนา, ฉ นินฺนา จ สมุทฺทโต. \\ เทฺว เต ฉ ทฺวาทส โหนฺติ, วคฺโค เตน ปวุจฺจตีติ. (โพชฺฌงฺคสํยุตฺตสฺส คงฺคาเปยฺยาลํ ราควเสน วิตฺถาเรตพฺพํ.)}
\csromangathagroup{\csromangathastanza{ฉ ปาจีนโต นินฺนา, ฉ นินฺนา จ สมุทฺทโต. \\ เทฺว เต ฉ ทฺวาทส โหนฺติ, วคฺโค เตน ปวุจฺจตีติ. (โพชฺฌงฺคสํยุตฺตสฺส คงฺคาเปยฺยาลํ ราควเสน วิตฺถาเรตพฺพํ.)}}

\csromanmatikaanchor{mtk.36}
\csromantocmarknopage{subparagraph}{15. ปุน-อปฺปมาทวคฺค}{mtk.36}
\bookmark[dest={mtk.36},level=5]{15. ปุน-อปฺปมาทวคฺค}
\csromantocmarkref{subparagraph}{324-333. ตถาคตาทิสุตฺต}{mtk.37}
\bookmark[dest={mtk.37},level=5]{324-333. ตถาคตาทิสุตฺต}
\csromanheader{5}{15. ปุน-อปฺปมาทวคฺค}
\vspace{\tipitakaparskip}
\csromancenter{ตถาคตาทิสุตฺต ปนฺนรสโม.}
\csromanheader{2}{\textbf{อุทฺทานํ}}
\csromangathameasure{ตถาคตํ ปทํ กูฏํ, มูลํ สาเรน วสฺสิกํ. \\ ราชา จนฺทิมสูริยา จ, วตฺเถน ทสมํ ปทนฺติ. (อปฺปมาทวคฺโค ราควเสน วิตฺถาเรตพฺโพ.)}
\csromangathagroup{\csromangathastanza{ตถาคตํ ปทํ กูฏํ, มูลํ สาเรน วสฺสิกํ. \\ ราชา จนฺทิมสูริยา จ, วตฺเถน ทสมํ ปทนฺติ. (อปฺปมาทวคฺโค ราควเสน วิตฺถาเรตพฺโพ.)}}

\csromanmatikaanchor{mtk.38}
\csromantocmarknopage{subparagraph}{16. ปุนพลกรณียวคฺค}{mtk.38}
\bookmark[dest={mtk.38},level=5]{16. ปุนพลกรณียวคฺค}
\csromantocmarkref{subparagraph}{334-356. ปุนพลาทิสุตฺต}{mtk.39}
\bookmark[dest={mtk.39},level=5]{334-356. ปุนพลาทิสุตฺต}
\csromanheader{5}{16. ปุนพลกรณียวคฺค}
\vspace{\tipitakaparskip}
\csromancenter{ปุนพลาทิสุตฺต โสฬสโม.}
\csromanheader{2}{\textbf{อุทฺทานํ}}
\csromangathameasure{พลํ พีชญฺจ นาโค จ, รุกฺโข กุมฺเภน สูกิยา. \\ อากาเสน จ เทฺว เมฆา, นาวา อาคนฺตุกา นทีติ. (โพชฺฌงฺคสํยุตฺตสฺส พลกรณียวคฺโค ราควเสน วิตฺถาเรตพฺโพ.)}
\csromangathagroup{\csromangathastanza{\csromanfolio{122}พลํ พีชญฺจ นาโค จ, รุกฺโข กุมฺเภน สูกิยา. \\ อากาเสน จ เทฺว เมฆา, นาวา อาคนฺตุกา นทีติ. (โพชฺฌงฺคสํยุตฺตสฺส พลกรณียวคฺโค ราควเสน วิตฺถาเรตพฺโพ.)}}

\csromanmatikaanchor{mtk.40}
\csromantocmarknopage{subparagraph}{17. ปุน-เอสนาวคฺค}{mtk.40}
\bookmark[dest={mtk.40},level=5]{17. ปุน-เอสนาวคฺค}
\csromanheader{5}{17. ปุน-เอสนาวคฺค}
\vspace{\tipitakaparskip}
\csromancenter{ปุน-เอสนาทิสุตฺต ปุน-เอสนาวคฺโค สตฺตรสโม.}
\csromanheader{2}{\textbf{อุทฺทานํ}}
\csromangathameasure{เอสนา วิธา อาสโว, ภโว จ ทุกฺขตา ติสฺโส. \\ ขิลํ มลญฺจ นีโฆ จ, เวทนา ตณฺหา ตสินาย จาติ.}
\csromangathagroup{\csromangathastanza{เอสนา วิธา อาสโว, ภโว จ ทุกฺขตา ติสฺโส. \\ ขิลํ มลญฺจ นีโฆ จ, เวทนา ตณฺหา ตสินาย จาติ.}}

\csromanmatikaanchor{mtk.41}
\csromanmatikaanchor{mtk.42}
\csromantocmarkref{subparagraph}{346-356. ปุน-เอสนาทิสุตฺต}{mtk.41}
\bookmark[dest={mtk.41},level=5]{346-356. ปุน-เอสนาทิสุตฺต}
\csromantocmarknopage{subparagraph}{18. ปุน-โอฆวคฺค}{mtk.42}
\bookmark[dest={mtk.42},level=5]{18. ปุน-โอฆวคฺค}
\csromanheader{5}{18. ปุน-โอฆวคฺค}
\csromanmatikaanchor{mtk.43}
\csromantocmarkref{subparagraph}{357-366. ปุน-โอฆาทิสุตฺต}{mtk.43}
\bookmark[dest={mtk.43},level=5]{357-366. ปุน-โอฆาทิสุตฺต}
\csromanheader{6}{357-366. ปุน-โอฆาทิสุตฺต}
\vspace{\tipitakaparskip}
\csromancenter{โพชฺฌงฺคสํยุตฺตสฺส ปุน-โอฆวคฺโค อฏฺฐารสโม.}
\csromanheader{2}{\textbf{อุทฺทานํ}}
\csromangathameasure{โอโฆ โยโค อุปาทานํ, คนฺถา อนุสเยน จ. \\ กามคุณา นีวรณา, ขนฺธา โอรุทฺธมฺภาคิยานีติ.}
\csromangathagroup{\csromangathastanza{โอโฆ โยโค อุปาทานํ, คนฺถา อนุสเยน จ. \\ กามคุณา นีวรณา, ขนฺธา โอรุทฺธมฺภาคิยานีติ.}}

\prose{(ราควินยปริโยสานโทสวินยปริโยสานโมหวินยปริโยสาน- วคฺโค วิตฺถาเรตพฺโพ.) (ยทปิ มคฺคสํยุตฺตํ วิตฺถาเรตพฺพํ, ตทปิ โพชฺฌงฺคสํยุตฺตํ วิตฺถาเรตพฺพํ.)}

\vspace{\tipitakaparskip}
\csromancenter{\textbf{โพชฺฌงฺคสํยุตฺตํ ทุติยํ.}}
\csromanmatikaanchor{mtk.44}
\csromantocmarknopage{subsection}{3. สติปฏฺฐานสํยุตฺต}{mtk.44}
\bookmark[dest={mtk.44},level=2]{3. สติปฏฺฐานสํยุตฺต}
\csromanheader{2}{3. สติปฏฺฐานสํยุตฺต}
\csromanmatikaanchor{mtk.45}
\csromantocmarknopage{paragraph}{1. อมฺพปาลิวคฺค}{mtk.45}
\bookmark[dest={mtk.45},level=4]{1. อมฺพปาลิวคฺค}
\csromanheader{4}{1. อมฺพปาลิวคฺค}
\csromanmatikaanchor{mtk.46}
\csromantocmarkref{subparagraph}{1. อมฺพปาลิสุตฺต}{mtk.46}
\bookmark[dest={mtk.46},level=5]{1. อมฺพปาลิสุตฺต}
\csromanheader{5}{1. อมฺพปาลิสุตฺต}
\proseitem{367}{\csromanfolio{123}เอวํ เม สุตํ\csromandash{}เอกํ สมยํ ภควา เวสาลิยํ วิหรติ อมฺพปาลิวเน.\csromanspacer{} ตตฺร โข ภควา ภิกฺขู อามนฺเตสิ “ภิกฺขโว”ติ. “ภทนฺเต”ติ เต ภิกฺขู ภควโต ปจฺจสฺโสสุํ.\csromanspacer{} ภควา เอตทโวจ\csromandash{}}

\prose{เอกายนฺวายํ ภิกฺขเว มคฺโค สตฺตานํ วิสุทฺธิยา โสกปริเทวานํ สมติกฺกมาย ทุกฺขโทมนสฺสานํ อตฺถงฺคมาย ญายสฺส อธิคมาย นิพฺพานสฺส สจฺฉิกิริยาย, ยทิทํ จตฺตาโร สติปฏฺฐานา.\csromanspacer{} กตเม จตฺตาโร, อิธ ภิกฺขเว ภิกฺขุ กาเย กายานุปสฺสี วิหรติ อาตาปี สมฺปชาโน สติมา วิเนยฺย โลเก อภิชฺฌาโทมนสฺสํ, เวทนาสุ เวทนานุปสฺสี วิหรติ อาตาปี สมฺปชาโน สติมา วิเนยฺย โลเก อภิชฺฌาโทมนสฺสํ, จิตฺเต จิตฺตานุปสฺสี วิหรติ อาตาปี สมฺปชาโน สติมา วิเนยฺย โลเก อภิชฺฌาโทมนสฺสํ, ธมฺเมสุ ธมฺมานุปสฺสี วิหรติ อาตาปี สมฺปชาโน สติมา วิเนยฺย โลเก อภิชฺฌาโทมนสฺสํ.\csromanspacer{} เอกายนฺวายํ ภิกฺขเว มคฺโค สตฺตานํ วิสุทฺธิยา โสกปริเทวานํ สมติกฺกมาย ทุกฺขโทมนสฺสานํ อตฺถงฺคมาย ญายสฺส อธิคมาย นิพฺพานสฺส สจฺฉิกิริยาย, ยทิทํ จตฺตาโร สติปฏฺฐานาติ.}

\prose{อิทมโวจ ภควา.\csromanspacer{} อตฺตมนา เต ภิกฺขู ภควโต ภาสิตํ อภินนฺทุนฺติ.\csromanspacer{}\csromanspacer{}ปฐมํ.}

\csromanmatikaanchor{mtk.47}
\csromantocmarkref{subparagraph}{2. สติสุตฺต}{mtk.47}
\bookmark[dest={mtk.47},level=5]{2. สติสุตฺต}
\csromanheader{5}{2. สติสุตฺต}
\proseitem{368}{เอกํ สมยํ ภควา เวสาลิยํ วิหรติ อมฺพปาลิวเน.\csromanspacer{} ตตฺร โข ภควา ภิกฺขู อามนฺเตสิ “ภิกฺขโว”ติ. “ภทนฺเต”ติ เต ภิกฺขู ภควโต ปจฺจสฺโสสุํ.\csromanspacer{} ภควา เอตทโวจ\csromandash{}}

\prose{สโต ภิกฺขเว ภิกฺขุ วิหเรยฺย สมฺปชาโน อยํ โว อมฺหากํ อนุสาสนี.\csromanspacer{} กถญฺจ ภิกฺขเว ภิกฺขุ สโต โหติ, อิธ ภิกฺขเว ภิกฺขุ กาเย กายานุปสฺสี วิหรติ อาตาปี สมฺปชาโน สติมา \csromanfolio{124}วิเนยฺย โลเก อภิชฺฌาโทมนสฺสํ, เวทนาสุ ฯเปฯ จิตฺเต ฯเปฯ ธมฺเมสุ ธมฺมานุปสฺสี วิหรติ อาตาปี สมฺปชาโน สติมา วิเนยฺย โลเก อภิชฺฌาโทมนสฺสํ.\csromanspacer{} เอวํ โข ภิกฺขเว ภิกฺขุ สโต โหติ.}

\prose{กถญฺจ ภิกฺขเว ภิกฺขุ สมฺปชาโน โหติ, อิธ ภิกฺขเว ภิกฺขุ อภิกฺกนฺเต ปฏิกฺกนฺเต สมฺปชานการี โหติ, อาโลกิเต วิโลกิเต สมฺปชานการี โหติ, สมิญฺชิเต ปสาริเต สมฺปชานการี โหติ, สงฺฆาฏิปตฺตจีวรธารเณ สมฺปชานการี โหติ, อสิเต ปีเต ขายิเต สายิเต สมฺปชานการี โหติ, อุจฺจารปสฺสาวกมฺเม สมฺปชานการี โหติ, คเต ฐิเต นิสินฺเน สุตฺเต ชาคริเต ภาสิเต ตุณฺหีภาเว สมฺปชานการี โหติ.\csromanspacer{} เอวํ โข ภิกฺขเว ภิกฺขุ สมฺปชานการี โหติ, สโต ภิกฺขเว ภิกฺขุ วิหเรยฺย สมฺปชาโน.\csromanspacer{} อยํ โว อมฺหากํ อนุสาสนีติ.\csromanspacer{}\csromanspacer{}ทุติยํ.}

\csromanmatikaanchor{mtk.48}
\csromantocmarkref{subparagraph}{3. ภิกฺขุสุตฺต}{mtk.48}
\bookmark[dest={mtk.48},level=5]{3. ภิกฺขุสุตฺต}
\csromanheader{5}{3. ภิกฺขุสุตฺต}
\proseitem{369}{เอกํ สมยํ ภควา สาวตฺถิยํ วิหรติ เชตวเน อนาถปิณฺฑิกสฺส อาราเม.\csromanspacer{} อถ โข อญฺญตโร ภิกฺขุ เยน ภควา เตนุปสงฺกมิ, อุปสงฺกมิตฺวา ภควนฺตํ อภิวาเทตฺวา เอกมนฺตํ นิสีทิ, เอกมนฺตํ นิสินฺโน โข โส ภิกฺขุ ภควนฺตํ เอตทโวจ “สาธุ เม ภนฺเต ภควา สํขิตฺเตน ธมฺมํ เทเสตุ, ยมหํ ภควโต ธมฺมํ สุตฺวา เอโก วูปกฏฺโฐ อปฺปมตฺโต อาตาปี ปหิตตฺโต วิหเรยฺยนฺ”ติ, เอวเมว ปนิเธกจฺเจ โมฆปุริสา มญฺเจว\footnote{มเมว (สี)} อชฺเฌสนฺติ, ธมฺเม จ ภาสิเต มเมว อนุพนฺธิตพฺพํ มญฺญนฺตีติ.\csromanspacer{} เทเสตุ เม ภนฺเต ภควา สํขิตฺเตน ธมฺมํ, เทเสตุ สุคโต สํขิตฺเตน ธมฺมํ, อปฺเปว นามาหํ ภควโต ภาสิตสฺส อตฺถํ ชาเนยฺยํ.\csromanspacer{} อปฺเปว นามาหํ ภควโต ภาสิตสฺส ทายาโท อสฺสนฺติ.\csromanspacer{} ตสฺมาติห ตฺวํ ภิกฺขุ อาทิเมว วิโสเธหิ กุสเลสุ ธมฺเมสุ.\csromanspacer{} โก จาทิ กุสลานํ ธมฺมานํ, สีลญฺจ สุวิสุทฺธํ ทิฏฺฐิ จ อุชุกา.\csromanspacer{} ยโต โข เต ภิกฺขุ สีลญฺจ สุวิสุทฺธํ ภวิสฺสติ, ทิฏฺฐิ จ อุชุกา.\csromanspacer{} ตโต ตฺวํ ภิกฺขุ สีลํ นิสฺสาย สีเล ปติฏฺฐาย จตฺตาโร สติปฏฺฐาเน ติวิเธน ภาเวยฺยาสิ.}

\csromanheader{2}{3. สติปฏฺฐานสํยุตฺต}
\prose{\csromanfolio{125}กตเม จตฺตาโร, อิธ ตฺวํ ภิกฺขุ อชฺฌตฺตํ วา กาเย กายานุปสฺสี วิหราหิ อาตาปี สมฺปชาโน สติมา วิเนยฺย โลเก อภิชฺฌาโทมนสฺสํ.\csromanspacer{} พหิทฺธา วา กาเย กายานุปสฺสี วิหราหิ อาตาปี สมฺปชาโน สติมา วิเนยฺย โลเก อภิชฺฌาโทมนสฺสํ.\csromanspacer{} อชฺฌตฺตพหิทฺธา วา กาเย กายานุปสฺสี วิหราหิ อาตาปี สมฺปชาโน สติมา วิเนยฺย โลเก อภิชฺฌาโทมนสฺสํ.\csromanspacer{} อชฺฌตฺตํ วา เวทนาสุ.\csromanspacer{} พหิทฺธา วา เวทนาสุ.\csromanspacer{} อชฺฌตฺตพหิทฺธา วา เวทนาสุ เวทนานุปสฺสี วิหราหิ อาตาปี สมฺปชาโน สติมา วิเนยฺย โลเก อภิชฺฌาโทมนสฺสํ.\csromanspacer{} อชฺฌตฺตํ วา จิตฺเต.\csromanspacer{} พหิทฺธา วา จิตฺเต.\csromanspacer{} อชฺฌตฺตพหิทฺธา วา จิตฺเต จิตฺตานุปสฺสี วิหราหิ อาตาปี สมฺปชาโน สติมา วิเนยฺย โลเก อภิชฺฌาโทมนสฺสํ.\csromanspacer{} อชฺฌตฺตํ วา ธมฺเมสุ.\csromanspacer{} พหิทฺธา วา ธมฺเมสุ.\csromanspacer{} อชฺฌตฺตพหิทฺธา วา ธมฺเมสุ ธมฺมานุปสฺสี วิหราหิ อาตาปี สมฺปชาโน สติมา วิเนยฺย โลเก อภิชฺฌาโทมนสฺสํ.\csromanspacer{} ยโต โข ตฺวํ ภิกฺขุ สีลํ นิสฺสาย สีเล ปติฏฺฐาย อิเม จตฺตาโร สติปฏฺฐาเน เอวํ ติวิเธน ภาเวสฺสสิ, ตโต ตุยฺหํ ภิกฺขุ ยา รตฺติ วา ทิวโส วา อาคมิสฺสติ, วุทฺธิเยว ปาฏิกงฺขา กุสเลสุ ธมฺเมสุ โน ปริหานีติ.}

\prose{อถ โข โส ภิกฺขเว ภควโต ภาสิตํ อภินนฺทิตฺวา อนุโมทิตฺวา อุฏฺฐายาสนา ภควนฺตํ อภิวาเทตฺวา ปทกฺขิณํ กตฺวา ปกฺกามิ.\csromanspacer{} อถ โข โส ภิกฺขุ เอโก วูปกฏฺโฐ อปฺปมตฺโต อาตาปี ปหิตตฺโต วิหรนฺโต นจิรสฺเสว ยสฺสตฺถาย กุลปุตฺตา สมฺมเทว อคารสฺมา อนคาริยํ ปพฺพชนฺติ, ตทนุตฺตรํ พฺรหฺมจริยปริโยสานํ ทิฏฺเฐว ธมฺเม สยํ อภิญฺญา สจฺฉิกตฺวา อุปสมฺปชฺช วิหรติ, “ขีณา ชาติ, วุสิตํ พฺรหฺมจริยํ, กตํ กรณียํ, นาปรํ อิตฺถตฺตายา”ติ อพฺภญฺญาสิ.\csromanspacer{} อญฺญตโร จ ปน โส ภิกฺขุ อรหตํ อโหสีติ.\csromanspacer{}\csromanspacer{}ตติยํ.}

\csromanmatikaanchor{mtk.49}
\csromantocmarkref{subparagraph}{4. สาลสุตฺต}{mtk.49}
\bookmark[dest={mtk.49},level=5]{4. สาลสุตฺต}
\csromanheader{5}{4. สาลสุตฺต}
\proseitem{370}{เอกํ สมยํ ภควา โกสเลสุ วิหรติ สาลาย พฺราหฺมณคาเม.\csromanspacer{} ตตฺร โข ภควา ภิกฺขู อามนฺเตสิ ฯเปฯ เอตทโวจ\csromandash{}}

\prose{เย เต ภิกฺขเว ภิกฺขู นวา อจิรปพฺพชิตา อธุนาคตา อิมํ ธมฺมวินยํ, เต โว ภิกฺขเว ภิกฺขู จตุนฺนํ สติปฏฺฐานานํ ภาวนาย สมาทเปตพฺพา นิเวเสตพฺพา ปติฏฺฐาเปตพฺพา.\csromanspacer{} กตเมสํ จตุนฺนํ, เอถ ตุเมฺห อาวุโส \csromanfolio{126}กาเย กายานุปสฺสิโน วิหรถ อาตาปิโน สมฺปชานา เอโกทิภูตา วิปฺปสนฺนจิตฺตา สมาหิตา เอกคฺคจิตฺตา กายสฺส ยถาภูตํ ญาณาย.\csromanspacer{} เวทนาสุ เวทนานุปสฺสิโน วิหรถ อาตาปิโน สมฺปชานา เอโกทิภูตา วิปฺปสนฺนจิตฺตา สมาหิตา เอกคฺคจิตฺตา เวทนานํ ยถาภูตํ ญาณาย.\csromanspacer{} จิตฺเต จิตฺตานุปสฺสิโน วิหรถ อาตาปิโน สมฺปชานา เอโกทิภูตา วิปฺปสนฺนจิตฺตา สมาหิตา เอกคฺคจิตฺตา จิตฺตสฺส ยถาภูตํ ญาณาย.\csromanspacer{} ธมฺเมสุ ธมฺมานุปสฺสิโน วิหรถ อาตาปิโน สมฺปชานา เอโกทิภูตา วิปฺปสนฺนจิตฺตา สมาหิตา เอกคฺคจิตฺตา ธมฺมานํ ยถาภูตํ ญาณาย.\csromanspacer{} เยปิ เต ภิกฺขเว ภิกฺขู เสขา อปฺปตฺตมานสา อนุตฺตรํ โยคกฺเขมํ ปตฺถยมานา วิหรนฺติ, เตปิ กาเย กายานุปสฺสิโน วิหรนฺติ อาตาปิโน สมฺปชานา เอโกทิภูตา วิปฺปสนฺนจิตฺตา สมาหิตา เอกคฺคจิตฺตา กายสฺส ปริญฺญาย.\csromanspacer{} เวทนาสุ เวทนานุปสฺสิโน วิหรนฺติ อาตาปิโน สมฺปชานา เอโกทิภูตา วิปฺปสนฺนจิตฺตา สมาหิตา เอกคฺคจิตฺตา เวทนานํ ปริญฺญาย.\csromanspacer{} จิตฺเต จิตฺตานุปสฺสิโน วิหรนฺติ อาตาปิโน สมฺปชานา เอโกทิภูตา วิปฺปสนฺนจิตฺตา สมาหิตา เอกคฺคจิตฺตา จิตฺตสฺส ปริญฺญาย.\csromanspacer{} ธมฺเมสุ ธมฺมานุปสฺสิโน วิหรนฺติ อาตาปิโน สมฺปชานา เอโกทิภูตา วิปฺปสนฺนจิตฺตา สมาหิตา เอกคฺคจิตฺตา ธมฺมานํ ปริญฺญาย.}

\prose{เยปิ เต ภิกฺขเว ภิกฺขู อรหนฺโต ขีณาสวา วุสิตวนฺโต กตกรณียา โอหิตภารา อนุปฺปตฺตสทตฺถา ปริกฺขีณภวสํโยชนา สมฺมทญฺญา วิมุตฺตา, เตปิ กาเย กายานุปสฺสิโน วิหรนฺติ อาตาวิโน สมฺปชานา เอโกทิภูตา วิปฺปสนฺนจิตฺตา สมาหิตา เอกคฺคจิตฺตา กาเยน วิสํยุตฺตา.\csromanspacer{} เวทนาสุ เวทนานุปสฺสิโน วิหรนฺติ อาตาปิโน สมฺปชานา เอโกทิภูตา วิปฺปสนฺนจิตฺตา สมาหิตา เอกคฺคจิตฺตา เวทนาหิ วิสํยุตฺตา.\csromanspacer{} จิตฺเต จิตฺตานุปสฺสิโน วิหรนฺติ อาตาปิโน สมฺปชานา เอโกทิภูตา วิปฺปสนฺนจิตฺตา สมาหิตา เอกคฺคจิตฺตา จิตฺเตน วิสํยุตฺตา.\csromanspacer{} ธมฺเมสุ ธมฺมานุปสฺสิโน วิหรนฺติ อาตาปิโน สมฺปชานา เอโกทิภูตา วิปฺปสนฺนจิตฺตา สมาหิตา เอกคฺคจิตฺตา ธมฺเมหิ วิสํยุตฺตา.}

\prose{เยปิ เต ภิกฺขเว ภิกฺขู นวา อจิรปพฺพชิตา อธุนาคตา อิมํ ธมฺมวินยํ, เต โว ภิกฺขเว ภิกฺขู อิเมสํ จตุนฺนํ สติปฏฺฐานานํ ภาวนาย สมาทเปตพฺพา นิเวเสตพฺพา ปติฏฺฐาเปตพฺพาติ.\csromanspacer{}\csromanspacer{}จตุตฺถํ.}

\csromanheader{2}{3. สติปฏฺฐานสํยุตฺต}
\csromanmatikaanchor{mtk.50}
\csromantocmarkref{subparagraph}{5. อกุสลราสิสุตฺต}{mtk.50}
\bookmark[dest={mtk.50},level=5]{5. อกุสลราสิสุตฺต}
\csromanheader{5}{5. อกุสลราสิสุตฺต}
\proseitem{371}{\csromanfolio{127}สาวตฺถินิทานํ.\csromanspacer{} ตตฺร โข ภควา เอตทโวจ\csromandash{}“อกุสลราสี”ติ ภิกฺขเว วทมาโน ปญฺจ นีวรเณ สมฺมา วทมาโน วเทยฺย.\csromanspacer{} เกวโลหายํ ภิกฺขเว อกุสลราสิ ยทิทํ ปญฺจ นีวรณา.\csromanspacer{} กตเม ปญฺจ, กามจฺฉนฺทนีวรณํ พฺยาปาทนีวรณํ ถินมิทฺธนีวรณํ อุทฺธจฺจกุกฺกุจฺจนีวรณํ วิจิกิจฺฉานีวรณํ. “อกุสลราสี”ติ ภิกฺขเว วทมาโน อิเม ปญฺจ นีวรเณ สมฺมา วทมาโน วเทยฺย.\csromanspacer{} เกวโล หายํ ภิกฺขเว อกุสลราสิ ยทิทํ ปญฺจ นีวรณา.}

\prose{“กุสลราสี”ติ ภิกฺขเว วทมาโน จตฺตาโร สติปฏฺฐาเน สมฺมา วทมาโน วเทยฺย.\csromanspacer{} เกวโล หายํ ภิกฺขเว กุสลราสิ ยทิทํ จตฺตาโร สติปฏฺฐานา.\csromanspacer{} กตเม จตฺตาโร, อิธ ภิกฺขเว ภิกฺขุ กาเย กายานุปสฺสี วิหรติ อาตาปี สมฺปชาโน สติมา วิเนยฺย โลเก อภิชฺฌาโทมนสฺสํ.\csromanspacer{} เวทนาสุ ฯเปฯ.\csromanspacer{} จิตฺเต ฯเปฯ.\csromanspacer{} ธมฺเมสุ ธมฺมานุปสฺสี วิหรติ อาตาปี สมฺปชาโน สติมา วิเนยฺย โลเก อภิชฺฌาโทมนสฺสํ. “กุสลราสี”ติ ภิกฺขเว วทมาโน อิเม จตฺตาโร สติปฏฺฐาเน สมฺมา วทมาโน วเทยฺย.\csromanspacer{} เกวโล หายํ ภิกฺขเว กุสลราสิ ยทิทํ จตฺตาโร สติปฏฺฐานาติ.\csromanspacer{}\csromanspacer{}ปญฺจมํ.}

\csromanmatikaanchor{mtk.51}
\csromantocmarkref{subparagraph}{6. สกุณคฺฆิสุตฺต}{mtk.51}
\bookmark[dest={mtk.51},level=5]{6. สกุณคฺฆิสุตฺต}
\csromanheader{5}{6. สกุณคฺฆิสุตฺต}
\proseitem{372}{ภูตปุพฺพํ ภิกฺขเว สกุณคฺฆิ ลาปํ สกุณํ สหสา อชฺฌปฺปตฺตา อคฺคเหสิ.\csromanspacer{} อถ โข ภิกฺขเว ลาโป สกุโณ สกุณคฺฆิยา หริยมาโน เอวํ ปริเทวสิ “มยเมวมฺห\footnote{มยเมวามฺห (ก)} อลกฺขิกา, มยํ อปฺปปุญฺญา, เย มยํ อโคจเร จริมฺห ปรวิสเย, สเจชฺช มยํ โคจเร จเรยฺยาม สเก เปตฺติเก วิสเย, น มฺยายํ\footnote{น จายํ (สี)} สกุณคฺฆิ อลํ อภวิสฺส ยทิทํ ยุทฺธายา”ติ.\csromanspacer{} โก ปน เต ลาป โคจโร สโก เปตฺติโก วิสโยติ.\csromanspacer{} ยทิทํ นงฺคลกฏฺฐกรณํ เลฑฺฑุฏฺฐานนฺติ.\csromanspacer{} อถ โข ภิกฺขเว สกุณคฺฆิ สเก พเล อปตฺถทฺธา สเก พเล อสํวทมานา\footnote{อวจมานา (สี)} ลาปํ สกุณํ ปมุญฺจิ “คจฺฉ โข ตฺวํ ลาป ตตฺรปิ เม คนฺตฺวา น โมกฺขสี”ติ.}

\prose{\csromanfolio{128}อถ โข ภิกฺขเว ลาโป สกุโณ นงฺคลกฏฺฐกรณํ โลฑฺฑุฏฺฐานํ คนฺตฺวา มหนฺตํ เลฑฺฑุํ อภิรุหิตฺวา สกุณคฺฆิํ วทมาโน อฏฺฐาสิ “เอหิ โข ทานิ เม สกุณคฺฆิ, เอหิ โข ทานิ เม สกุณคฺฆี”ติ.\csromanspacer{} อถ โข สา ภิกฺขเว สกุณคฺฆิ สเก พเล อปตฺถทฺธา สเก พเล อสํวทมานา อุโภ ปกฺเข สนฺนยฺห\footnote{สนฺธาย (สี, สฺยา)} ลาปํ สกุณํ สหสา อชฺฌปฺปตฺตา.\csromanspacer{} ยทา โข ภิกฺขเว อญฺญาสิ ลาโป สกุโณ พหุ-อาคโต โข มฺยายํ สกุณคฺฆีติ, อถ ตสฺเสว เลฑฺฑุสฺส อนฺตรํ ปจฺจุปาทิ.\csromanspacer{} อถ โข ภิกฺขเว สกุณคฺฆิ ตตฺเถว อุรํ ปจฺจตาเฬสิ.\csromanspacer{} เอวํ หิ ตํ\footnote{เอวํ เหตํ (สี)} ภิกฺขเว โหติ โย อโคจเร จรติ ปรวิสเย.}

\prose{ตสฺมาติห ภิกฺขเว มา อโคจเร จริตฺถ ปรวิสเย, อโคจเร ภิกฺขเว จรตํ ปรวิสเย ลจฺฉติ มาโร โอตารํ, ลจฺฉติ มาโร อารมฺมณํ.\csromanspacer{} โก จ ภิกฺขเว ภิกฺขุโน อโคจโร ปรวิสโย, ยทิทํ ปญฺจ กามคุณา.\csromanspacer{} กตเม ปญฺจ, จกฺขุวิญฺเญยฺยา รูปา อิฏฺฐา กนฺตา มนาปา ปิยรูปา กามูปสํหิตา รชนียา.\csromanspacer{} โสตวิญฺเญยฺยา สทฺทา ฯเปฯ.\csromanspacer{} ฆานวิญฺเญยฺยา คนฺธา ฯเปฯ.\csromanspacer{} ชิวฺหาวิญฺเญยฺยา รสา ฯเปฯ.\csromanspacer{} กายวิญฺเญยฺยา โผฏฺฐพฺพา อิฏฺฐา กนฺตา มนาปา ปิยรูปา กามูปสํหิตา รชนียา.\csromanspacer{} อยํ ภิกฺขเว ภิกฺขุโน อโคจโร ปรวิสโย.}

\prose{โคจเร ภิกฺขเว จรถ สเก เปตฺติเก วิสเย, โคจเร ภิกฺขเว จรตํ สเก เปตฺติเก วิสเย น ลจฺฉติ มาโร โอตารํ, น ลจฺฉติ มาโร อารมฺมณํ.\csromanspacer{} โก จ ภิกฺขเว ภิกฺขุโน โคจโร สโก เปตฺติโก วิสโย, ยทิทํ จตฺตาโร สติปฏฺฐานา.\csromanspacer{} กตเม จตฺตาโร, อิธ ภิกฺขเว ภิกฺขุ กาเย กายานุปสฺสี วิหรติ อาตาปี สมฺปชาโน สติมา วิเนยฺย โลเก อภิชฺฌาโทมนสฺสํ.\csromanspacer{} เวทนาสุ ฯเปฯ.\csromanspacer{} จิตฺเต ฯเปฯ.\csromanspacer{} ธมฺเมสุ ธมฺมานุปสฺสี วิหรติ อาตาปี สมฺปชาโน สติมา วิเนยฺย โลเก อภิชฺฌาโทมนสฺสํ.\csromanspacer{} อยํ ภิกฺขเว ภิกฺขุโน โคจโร สโก เปตฺติโก วิสโยติ.\csromanspacer{}\csromanspacer{}ฉฏฺฐํ.}

\csromanmatikaanchor{mtk.52}
\csromantocmarkref{subparagraph}{7. มกฺกฏสุตฺต}{mtk.52}
\bookmark[dest={mtk.52},level=5]{7. มกฺกฏสุตฺต}
\csromanheader{5}{7. มกฺกฏสุตฺต}
\proseitem{373}{อตฺถิ ภิกฺขเว หิมวโต ปพฺพตราชสฺส ทุคฺคา วิสมา เทสา, ยตฺถ เนว มกฺกฏานํ จารี น มนุสฺสานํ.\csromanspacer{} อตฺถิ ภิกฺขเว หิมวโต}

\csromanheader{2}{3. สติปฏฺฐานสํยุตฺต}
\prose{\csromanfolio{129}ปพฺพตราชสฺส ทุคฺคา วิสมา เทสา, ยตฺถ มกฺกฏานํ หิ โข จารี น มนุสฺสานํ.\csromanspacer{} อตฺถิ ภิกฺขเว หิมวโต ปพฺพตราชสฺส สมา ภูมิภาคา รมณียา, ยตฺถ มกฺกฏานญฺเจว จารี มนุสฺสานญฺจ.\csromanspacer{} ตตฺร ภิกฺขเว ลุทฺทา มกฺกฏวีถีสุ เลปํ โอฑฺเฑนฺติ มกฺกฏานํ พาธนาย.}

\prose{ตตฺร ภิกฺขเว เย เต มกฺกฏา อพาลชาติกา อโลลชาติกา, เต ตํ เลปํ ทิสฺวา อารกา ปริวชฺชนฺติ.\csromanspacer{} โย ปน โส โหติ มกฺกโฏ พาลชาติโก โลลชาติโก, โส ตํ เลปํ อุปสงฺกมิตฺวา หตฺเถน คณฺหาติ, โส ตตฺถ พชฺฌติ. “หตฺถํ โมเจสฺสามี”ติ ทุติเยน หตฺเถน คณฺหาติ, โส ตตฺถ พชฺฌติ. “อุโภ หตฺเถ โมเจสฺสามี”ติ ปาเทน คณฺหาติ, โส ตตฺถ พชฺฌติ. “อุโภ หตฺเถ โมเจสฺสามิ ปาทญฺจา”ติ ทุติเยน ปาเทน คณฺหาติ, โส ตตฺถ พชฺฌติ. “อุโภ หตฺเถ โมเจสฺสามิ ปาเท จา”ติ ตุณฺเฑน คณฺหาติ, โส ตตฺถ พชฺฌติ.\csromanspacer{} เอวํ หิ โส ภิกฺขเว มกฺกโฏ ปญฺโจฑฺฑิโต ถุนํ เสติ อนยํ อาปนฺโน, พฺยสนํ อาปนฺโน, ยถากามกรณีโย ลุทฺทสฺส.\csromanspacer{} ตเมนํ ภิกฺขเว ลุทฺโท วิชฺฌิตฺวา ตสฺมิํ เยว กฏฺฐกตงฺคาเร\footnote{ตสฺมิํ เยว มกฺกฏํ อุทฺธริตฺวา (สี, สฺยา)} อวสฺสชฺเชตฺวา เยน กามํ ปกฺกมติ.\csromanspacer{} เอวํ โส ตํ ภิกฺขเว โหติ, โย อโคจเร จรติ ปรวิสเย.}

\prose{ตสฺมาติห ภิกฺขเว มา อโคจเร จริตฺถ ปรวิสเย, อโคจเร ภิกฺขเว จรตํ ปรวิสเย ลจฺฉติ มาโร โอตารํ, ลจฺฉติ มาโร อารมฺมณํ.\csromanspacer{} โก จ ภิกฺขเว ภิกฺขุโน อโคจโร ปรวิสโย, ยทิทํ ปญฺจ กามคุณา.\csromanspacer{} กตเม ปญฺจ, จกฺขุวิญฺเญยฺยา รูปา อิฏฺฐา กนฺตา มนาปา ปิยรูปา กามูปสํหิตา รชนียา.\csromanspacer{} โสตวิญฺเญยฺยา สทฺทา ฯเปฯ.\csromanspacer{} ฆานวิญฺเญยฺยา คนฺธา ฯเปฯ.\csromanspacer{} ชิวฺหาวิญฺเญยฺยา รสา ฯเปฯ.\csromanspacer{} กายวิญฺเญยฺยา โผฏฺฐพฺพา อิฏฺฐา กนฺตา มนาปา ปิยรูปา กามูปสํหิตา รชนียา.\csromanspacer{} อยํ ภิกฺขเว ภิกฺขุโน อโคจโร ปรวิสโย.}

\prose{โคจเร ภิกฺขเว จรถ สเก เปตฺติเก วิสเย, โคจเร ภิกฺขเว จรตํ สเก เปตฺติเก วิสเย น ลจฺฉติ มาโร โอตารํ, น ลจฺฉติ มาโร อารมฺมณํ.\csromanspacer{} โก จ ภิกฺขเว ภิกฺขุโน โคจโร สโก เปตฺติโก วิสโย, ยทิทํ จตฺตาโร สติปฏฺฐานา.\csromanspacer{} กตเม}

\prose{\csromanfolio{130}จตฺตาโร, อิธ ภิกฺขเว ภิกฺขุ กาเย กายานุปสฺสี วิหรติ อาตาปี สมฺปชาโน สติมา วิเนยฺย โลเก อภิชฺฌาโทมนสฺสํ.\csromanspacer{} เวทนาสุ ฯเปฯ.\csromanspacer{} จิตฺเต ฯเปฯ.\csromanspacer{} ธมฺเมสุ ธมฺมานุปสฺสี วิหรติ อาตาปี สมฺปชาโน สติมา วิเนยฺย โลเก อภิชฺฌาโทมนสฺสํ.\csromanspacer{} อยํ ภิกฺขเว ภิกฺขุโน โคจโร สโก เปตฺติโก วิสโยติ.\csromanspacer{}\csromanspacer{}สตฺตมํ.}

\csromanmatikaanchor{mtk.53}
\csromantocmarkref{subparagraph}{8. สูทสุตฺต}{mtk.53}
\bookmark[dest={mtk.53},level=5]{8. สูทสุตฺต}
\csromanheader{5}{8. สูทสุตฺต}
\proseitem{374}{เสยฺยถาปิ ภิกฺขเว พาโล อพฺยตฺโต อกุสโล สูโท ราชานํ วา ราชมหามตฺตํ วา\footnote{ราชมหามตฺตานํ วา (สี)} นานจฺจเยหิ สูเปหิ ปจฺจุปฏฺฐิโต อสฺส อมฺพิลคฺเคหิปิ ติตฺตกคฺเคหิปิ กฏุกคฺเคหิปิ มธุรคฺเคหิปิ ขาริเกหิปิ อขาริเกหิปิ โลณิเกหิปิ อโลณิเกหิปิ.}

\prose{ส โข โส ภิกฺขเว พาโล อพฺยตฺโต อกุสโล สูโท สกสฺส ภตฺตุ นิมิตฺตํ น อุคฺคณฺหาติ “อิทํ วา เม อชฺช ภตฺตุ สูเปยฺยํ รุจฺจติ, อิมสฺส วา อภิหรติ.\csromanspacer{} อิมสฺส วา พหุํ คณฺหาติ, อิมสฺส วา วณฺณํ ภาสติ.\csromanspacer{} อมฺพิลคฺคํ วา เม อชฺช ภตฺตุ สูเปยฺยํ รุจฺจติ, อมฺพิลคฺคสฺส วา อภิหรติ, อมฺพิลคฺคสฺส วา พหุํ คณฺหาติ, อมฺภิลคฺคสฺส วา วณฺณํ ภาสติ.\csromanspacer{} ติตฺตกคฺคํ วา เม อชฺช.\csromanspacer{} กฏุกคฺคํ วา เม อชฺช.\csromanspacer{} มธุรคฺคํ วา เม อชฺช.\csromanspacer{} ขาริกํ วา เม อชฺช.\csromanspacer{} อขาริกํ วา เม อชฺช.\csromanspacer{} โลณิกํ วา เม อชฺช.\csromanspacer{} อโลณิกํ วา เม อชฺช ภตฺตุ สูเปยฺยํ รุจฺจติ, อโลณิกสฺส วา อภิหรติ, อโลณิกสฺส วา พหุํ คณฺหาติ, อโลณิกสฺส วา วณฺณํ ภาสตี”ติ.}

\prose{ส โข โส ภิกฺขเว พาโล อพฺยตฺโต อกุสโล สูโท น เจว ลาภี โหติ อจฺฉาทนสฺส, น ลาภี เวตนสฺส, น ลาภี อภิหารานํ.\csromanspacer{} ตํ กิสฺส เหตุ, ตถา หิ โส ภิกฺขเว พาโล อพฺยตฺโต อกุสโล สูโท สกสฺส ภตฺตุ นิมิตฺตํ น อุคฺคณฺหาติ.\csromanspacer{} เอวเมว โข ภิกฺขเว อิเธกจฺโจ พาโล อพฺยตฺโต อกุสโล ภิกฺขุ กาเย กายานุปสฺสี วิหรติ อาตาปี สมฺปชาโน สติมา วิเนยฺย โลเก อภิชฺฌาโทมนสฺสํ, ตสฺส กาเย กายานุปสฺสิโน วิหรโต จิตฺตํ น สมาธิยติ, อุปกฺกิเลสา น ปหียนฺติ, โส ตํ นิมิตฺตํ น อุคฺคณฺหาติ.\csromanspacer{} เวทนาสุ เวทนานุปสฺสี วิหรติ ฯเปฯ.\csromanspacer{} จิตฺเต จิตฺตานุปสฺสี}

\csromanheader{2}{3. สติปฏฺฐานสํยุตฺต}
\prose{\csromanfolio{131}วิหรติ ฯเปฯ.\csromanspacer{} ธมฺเมสุ ธมฺมานุปสฺสี วิหรติ อาตาปี สมฺปชาโน สติมา วิเนยฺย โลเก อภิชฺฌาโทมนสฺสํ, ตสฺส ธมฺเมสุ ธมฺมานุปสฺสิโน วิหรโต จิตฺตํ น สมาธิยติ, อุปกฺกิเลสา น ปหียนฺติ, โส ตํ นิมิตฺตํ น อุคฺคณฺหาติ.}

\prose{ส โข โส ภิกฺขเว พาโล อพฺยตฺโต อกุสโล ภิกฺขุ น เจว ลาภี โหติ ทิฏฺเฐว ธมฺเม สุขวิหารานํ, น ลาภี สติสมฺปชญฺญสฺส.\csromanspacer{} ตํ กิสฺส เหตุ, ตถา หิ โส ภิกฺขเว พาโล อพฺยตฺโต อกุสโล ภิกฺขุ สกสฺส จิตฺตสฺส นิมิตฺตํ น อุคฺคณฺหาติ.}

\prose{เสยฺยถาปิ ภิกฺขเว ปณฺฑิโต พฺยตฺโต กุสโล สูโท ราชานํ วา ราชมหามตฺตํ วา นานจฺจเยหิ สูเปหิ ปจฺจุปฏฺฐิโต อสฺส อมฺพิลคฺเคหิปิ ติตฺตกคฺเคหิปิ กฏุกคฺเคหิปิ มธุรคฺเคหิปิ ขาริเกหิปิ อขาริเกหิปิ โลณิเกหิปิ อโลณิเกหิปิ.}

\prose{ส โข โส ภิกฺขเว ปณฺฑิโต พฺยตฺโต กุสโล สูโท สกสฺส ภตฺตุ นิมิตฺตํ อุคฺคณฺหาติ “อิทํ วา เม อชฺช ภตฺตุ สูเปยฺยํ รุจฺจติ, อิมสฺส วา อภิหรติ, อิมสฺส วา พหุํ คณฺหาติ, อิมสฺส วา วณฺณํ ภาสติ.\csromanspacer{} อมฺพิลคฺคํ วา เม อชฺช ภตฺตุ สูเปยฺยํ รุจฺจติ, อมฺพิลคฺคสฺส วา อภิหรติ, อมฺพิลคฺคสฺส วา พหุํ คณฺหาติ, อมฺพิลคฺคสฺส วา วณฺณํ ภาสติ.\csromanspacer{} ติตฺตกคฺคํ วา เม อชฺช.\csromanspacer{} กฏุกคฺคํ วา เม อชฺช.\csromanspacer{} มธุรคฺคํ วา เม อชฺช.\csromanspacer{} ขาริกํ วา เม อชฺช.\csromanspacer{} อขาริกํ วา เม อชฺช.\csromanspacer{} โลณิกํ วา เม อชฺช.\csromanspacer{} อโลณิกํ วา เม อชฺช ภตฺตุ สูเปยฺยํ รุจฺจติ.\csromanspacer{} อโลณิกสฺส วา อภิหรติ, อโลณิกสฺส วา พหุํ คณฺหาติ.\csromanspacer{} อโลณิกสฺส วา วณฺณํ ภาสตี”ติ.}

\prose{ส โข โส ภิกฺขเว ปณฺฑิโต พฺยตฺโต กุสโล สูโท ลาภี เจว โหติ อจฺฉาทนสฺส, ลาภี เวตนสฺส, ลาภี อภิหารานํ.\csromanspacer{} ตํ กิสฺส เหตุ, ตถา หิ โส ภิกฺขเว ปณฺฑิโต พฺยตฺโต กุสโล สูโท สกสฺส ภตฺตุ นิมิตฺตํ อุคฺคณฺหาติ.\csromanspacer{} เอวเมว โข ภิกฺขเว อิเธกจฺโจ ปณฺฑิโต พฺยตฺโต กุสโล ภิกฺขุ กาเย กายานุปสฺสี วิหรติ อาตาปี สมฺปชาโน สติมา วิเนยฺย โลเก อภิชฺฌาโทมนสฺสํ, ตสฺส กาเย กายานุปสฺสิโน วิหรโต จิตฺตํ สมาธิยติ, อุปกฺกิเลสา ปหียนฺติ, โส ตํ นิมิตฺตํ อุคฺคณฺหาติ.\csromanspacer{} เวทนาสุ เวทนานุปสฺสี วิหรติ ฯเปฯ.\csromanspacer{} จิตฺเต จิตฺตานุปสฺสี วิหรติ ฯเปฯ.\csromanspacer{} ธมฺเมสุ ธมฺมานุปสฺสี}

\prose{\csromanfolio{132}วิหรติ อาตาปี สมฺปชาโน สติมา วิเนยฺย โลเก อภิชฺฌาโทมนสฺสํ, ตสฺส ธมฺเมสุ ธมฺมานุปสฺสิโน วิหรโต จิตฺตํ สมาธิยติ, อุปกฺกิเลสา ปหียนฺติ, โส ตํ นิมิตฺตํ อุคฺคณฺหาติ.}

\prose{ส โข โส ภิกฺขเว ปณฺฑิโต พฺยตฺโต กุสโล ภิกฺขุ ลาภี เจว โหติ ทิฏฺเฐว ธมฺเม สุขวิหารานํ, ลาภี โหติ สติสมฺปชญฺญสฺส.\csromanspacer{} ตํ กิสฺส เหตุ, ตถา หิ โส ภิกฺขเว ปณฺฑิโต พฺยตฺโต กุสโล ภิกฺขุ สกสฺส จิตฺตสฺส นิมิตฺตํ อุคฺคณฺหาตีติ.\csromanspacer{}\csromanspacer{}อฏฺฐมํ.}

\csromanmatikaanchor{mtk.54}
\csromantocmarkref{subparagraph}{9. คิลานสุตฺต}{mtk.54}
\bookmark[dest={mtk.54},level=5]{9. คิลานสุตฺต}
\csromanheader{5}{9. คิลานสุตฺต}
\proseitem{375}{เอวํ เม สุตํ\csromandash{}เอกํ สมยํ ภควา เวสาลิยํ วิหรติ เวฬุวคามเก\footnote{เพลุว คามเก (สี, สฺยา, กํ, อิ)}.\csromanspacer{} ตตฺร โข ภควา ภิกฺขู อมนฺเตสิ “เอถ ตุเมฺห ภิกฺขเว สมนฺตา เวสาลิยา ยถามิตฺตํ ยถาสนฺทิฏฺฐํ ยถาสมฺภตฺตํ วสฺสํ อุเปถ, อิเธวาหํ เวฬุวคามเก วสฺสํ อุปคจฺฉามี”ติ. “เอวํ ภนฺเต”ติ โข เต ภิกฺขู ภควโต ปฏิสฺสุตฺวา สมนฺตา เวสาลิยา ยถามิตฺตํ ยถาสนฺทิฏฺฐํ ยถาสมฺภตฺตํ วสฺสํ อุปคจฺฉุํ, ภควา ปน ตตฺเถว เวฬุวคามเก วสฺสํ อุปคจฺฉิ\footnote{อุปคญฺฉิ (สี, อิ)}.}

\prose{อถ โข ภควโต วสฺสูปคตสฺส ขโร อาพาโธ อุปฺปชฺชิ, พาฬฺหา เวทนา วตฺตนฺติ มารณนฺติกา, ตตฺร สุทํ ภควา สโต สมฺปชาโน อธิวาเสสิ อวิหญฺญมาโน.\csromanspacer{} อถ โข ภควโต เอตทโหสิ “น โข เม ตํ ปติรูปํ, โยหํ อนามนฺเตตฺวา อุปฏฺฐาเก อนปโลเกตฺวา ภิกฺขุสํฆํ ปรินิพฺพาเยยฺยํ, ยํนูนาหํ อิมํ อาพาธํ วีริเยน ปฏิปณาเมตฺวา ชีวิตสงฺขารํ อธิฏฺฐาย วิหเรยฺยนฺ”ติ.\csromanspacer{} อถ โข ภควา ตํ อาพาธํ วีริเยน ปฏิปณาเมตฺวา ชีวิตสงฺขารํ อธิฏฺฐาย วิหาสิ. (อถ โข ภควโต โส อาพาโธ ปฏิปฺปสฺสมฺภิ.)\footnote{( ) ที 2. 84 ปิฏฺเฐ ทิสฺสติ.}}

\prose{อถ โข ภควา คิลานา วุฏฺฐิโต\footnote{คิลานวุฏฺฐิโต (สทฺทนีติ)} อจิรวุฏฺฐิโต เคลญฺญา วิหารา นิกฺขมิตฺวา วิหารปจฺฉายายํ\footnote{วิหารปจฺฉาฉายายํ (พหูสุ)} ปญฺญตฺเต อาสเน นิสีทิ.\csromanspacer{} อถ โข อายสฺมา อานนฺโท เยน ภควา เตนุปสงฺกมิ, อุปสงฺกมิตฺวา ภควนฺตํ อภิวาเทตฺวา เอกมนฺตํ นิสีทิ, เอกมนฺตํ นิสินฺโน โข อายสฺมา อานนฺโท}

\csromanheader{2}{3. สติปฏฺฐานสํยุตฺต}
\prose{\csromanfolio{133}ภควนฺตํ เอตทโวจ “ทิฏฺโฐ เม ภนฺเต ภควโต ผาสุ, ทิฏฺฐํ ภนฺเต ภควโต ขมนียํ, ทิฏฺฐํ ภนฺเต ภควโต ยาปนียํ, อปิ จ เม ภนฺเต มธุรกชาโต วิย กาโย, ทิสาปิ เม น ปกฺขายนฺติ, ธมฺมาปิ มํ นปฺปฏิภนฺติ ภควโต เคลญฺเญน, อปิ จ เม ภนฺเต อโหสิ กาจิเทว อสฺสาสมตฺตา, น ตาว ภควา ปรินิพฺพายิสฺสติ, น ยาว ภควา ภิกฺขุสํฆํ อารพฺภ กิญฺจิเทว อุทาหรตี”ติ.}

\prose{กิํ ปน ทานิ อานนฺท ภิกฺขุสํโฆ มยิ ปจฺจาสีสติ\footnote{ปจฺจาสิํ สติ (สี, สฺยา, กํ, อิ)}, เทสิโต อานนฺท มยา ธมฺโม อนนฺตรํ อพาหิรํ กริตฺวา, นตฺถานนฺท ตถาคตสฺส ธมฺเมสุ อาจริยมุฏฺฐิ, ยสฺส นูน อานนฺท เอวมสฺส “อหํ ภิกฺขุสํฆํ ปริหริสฺสามี”ติ วา “มมุทฺเทสิโก ภิกฺขุสํโฆ”ติ วา, โส นูน อานนฺท ภิกฺขุสํฆํ อารพฺภ กิญฺจิเทว อุทาหเรยฺย.\csromanspacer{} ตถาคตสฺส โข อานนฺท น เอวํ โหติ “อหํ ภิกฺขุสํฆํ ปริหริสฺสามี”ติ วา “มมุทฺเทสิโก ภิกฺขุสํโฆ”ติ วา, ส กิํ\footnote{โส นูน (สี, อิ)} อานนฺท ตถาคโต ภิกฺขุสํฆํ อารพฺภ กิญฺจิเทว อุทาหริสฺสติ, เอตรหิ โข ปนาหํ อานนฺท ชิณฺโณ วุทฺโธ มหลฺลโก อทฺธคโต วโย-อนุปฺปตฺโต, อาสีติโก เม วโย วตฺตติ.\csromanspacer{} เสยฺยถาปิ อานนฺท ชชฺชรสกฏํ\footnote{ชรสกฏํ (สพฺพตฺถ) 4. เวคมิสฺสเกน (สี), เวฬุมิสฺสเกน (สฺยา, กํ), เวธมิสฺสเกน (อิ, ก), เวขมิสฺสเกน (ก) 5. ผาสุตรํ (สพฺพตฺถ)} เวฐมิสฺสเกน4 ยาเปติ.\csromanspacer{} เอวเมว โข อานนฺท เวขมิสฺสเกน มญฺเญ ตถาคตสฺส กาโย ยาเปติ.}

\prose{ยสฺมิํ อานนฺท สมเย ตถาคโต สพฺพนิมิตฺตานํ อมนสิการา เอกจฺจานํ เวทนานํ นิโรธา อนิมิตฺตํ เจโตสมาธิํ อุปสมฺปชฺช วิหรติ.\csromanspacer{} ผาสุตโร5 อานนฺท ตสฺมิํ สมเย ตถาคตสฺส กาโย โหติ\footnote{ตถาคตสฺส โหติ (พหูสุ)}.\csromanspacer{} ตสฺมาติหานนฺท อตฺตทีปา วิหรถ อตฺตสรณา อนญฺญสรณา ธมฺมทีปา ธมฺมสรณา อนญฺญสรณา.}

\prose{กถญฺจานนฺท ภิกฺขุ อตฺตทีโป วิหรติ อตฺตสรโณ อนญฺญสรโณ ธมฺมทีโป ธมฺมสรโณ อนญฺญสรโณ, อิธานนฺท ภิกฺขุ กาเย กายานุปสฺสี วิหรติ อาตาปี สมฺปชาโน สติมา วิเนยฺย โลเก อภิชฺฌาโทมนสฺสํ.\csromanspacer{} เวทนาสุ ฯเปฯ.\csromanspacer{} จิตฺเต ฯเปฯ.\csromanspacer{} ธมฺเมสุ ธมฺมานุปสฺสี วิหรติ อาตาปี}

\prose{\csromanfolio{134}สมฺปชาโน สติมา วิเนยฺย โลเก อภิชฺฌาโทมนสฺสํ.\csromanspacer{} เอวํ โข อานนฺท ภิกฺขุ อตฺตทีโป วิหรติ อตฺตสรโณ อนญฺญสรโณ ธมฺมทีโป ธมฺมสรโณ อนญฺญสรโณ.\csromanspacer{} เย หิ เกจิ อานนฺท เอตรหิ วา มมจฺจเย วา อตฺตทีปา วิหริสฺสนฺติ อตฺตสรณา อนญฺญสรณา ธมฺมทีปา ธมฺมสรณา อนญฺญสรณา.\csromanspacer{} ตมตคฺเค เมเต อานนฺท ภิกฺขู ภวิสฺสนฺติ เย เกจิ สิกฺขากามาติ.\csromanspacer{}\csromanspacer{}นวมํ.}

\csromanmatikaanchor{mtk.55}
\csromantocmarkref{subparagraph}{10. ภิกฺขุนุปสฺสยสุตฺต}{mtk.55}
\bookmark[dest={mtk.55},level=5]{10. ภิกฺขุนุปสฺสยสุตฺต}
\csromanheader{5}{10. ภิกฺขุนุปสฺสยสุตฺต}
\proseitem{376}{อถ โข อายสฺมา อานนฺโท ปุพฺพณฺหสมยํ นิวาเสตฺวา ปตฺตจีวรมาทาย เยน อญฺญตโร ภิกฺขุนุปสฺสโย เตนุปสงฺกมิ, อุปสงฺกมิตฺวา ปญฺญตฺเต อาสเน นิสีทิ.\csromanspacer{} อถ โข สมฺพหุลา ภิกฺขุนิโย เยนายสฺมา อานนฺโท เตนุปสงฺกมิํสุ, อุปสงฺกมิตฺวา อายสฺมนฺตํ อานนฺทํ อภิวาเทตฺวา เอกมนฺตํ นิสีทิํสุ, เอกมนฺตํ นิสินฺโน โข ตา ภิกฺขุนิโย อายสฺมนฺตํ อานนฺทํ เอตทโวจุํ\csromandash{}}

\prose{อิธ ภนฺเต อานนฺท สมฺพหุลา ภิกฺขุนิโย จตูสุ สติปฏฺฐาเนสุ สุปฺปติฏฺฐิตจิตฺตา\footnote{สุปฏฺฐิตจิตฺตา (สี, อิ, ก)} วิหรนฺติโย อุฬารํ ปุพฺเพนาปรํ วิเสสํ สญฺชานนฺตีติ\footnote{สมฺปชานนฺตีติ (ก)}.\csromanspacer{} เอวเมตํ ภคินิโย, เอวเมตํ ภคินิโย, โย หิ โกจิ ภคินิโย ภิกฺขุ วา ภิกฺขุนี วา จตูสุ สติปฏฺฐาเนสุ สุปฺปติฏฺฐิตจิตฺโต วิหรติ, ตสฺเสตํ ปาฏิกงฺขํ “อุฬารํ ปุพฺเพนาปรํ วิเสสํ สญฺชานิสฺสตี”ติ.}

\prose{อถ โข อายสฺมา อานนฺโท ตา ภิกฺขุนิโย ธมฺมิยา กถาย สนฺทสฺเสตฺวา สมาทเปตฺวา สมุตฺเตเชตฺวา สมฺปหํเสตฺวา อุฏฺฐายาสนา ปกฺกามิ.\csromanspacer{} อถ โข อายสฺมา อานนฺโท สาวตฺถิยํ ปิณฺฑาย จริตฺวา ปจฺฉาภตฺตํ ปิณฺฑปาตปฏิกฺกนฺโต เยน ภควา เตนุปสงฺกมิ, อุปสงฺกมิตฺวา ภควนฺตํ อภิวาเทตฺวา เอกมนฺตํ นิสีทิ, เอกมนฺตํ นิสินฺโน โข อายสฺมา อานนฺโท ภควนฺตํ เอตทโวจ\csromandash{}}

\prose{อิธาหํ ภนฺเต ปุพฺพณฺหสมยํ นิวาเสตฺวา ปตฺตจีวรมาทาย เยน อญฺญตโร ภิกฺขุนุปสฺสโย เตนุปสงฺกมิํ, อุปสงฺกมิตฺวา ปญฺญตฺเต อาสเน นิสีทิ.\csromanspacer{} อถ โข ภนฺเต สมฺพหุลา ภิกฺขุนิโย เยนาหํ เตนุปสงฺกมิํสุ, อุปสงฺกมิตฺวา มํ อภิวาเทตฺวา เอกมนฺตํ นิสีทิํสุ,}

\csromanheader{2}{3. สติปฏฺฐานสํยุตฺต}
\prose{\csromanfolio{135}เอกมนฺตํ นิสินฺนา โข ภนฺเต ตา ภิกฺขุนิโย มํ เอตทโวจุํ “อิธ ภนฺเต อานนฺท สมฺพหุลา ภิกฺขุนิโย จตูสุ สติปฏฺฐาเนสุ สุปฺปติฏฺฐิตจิตฺตา วิหรนฺติโย อุฬารํ ปุพฺเพนาปรํ วิเสสํ สญฺชานนฺตี”ติ.\csromanspacer{} เอวํ วุตฺตาหํ ภนฺเต ตา ภิกฺขุนิโย เอตทโวจํ “เอวเมตํ ภคินิโย, เอวเมตํ ภคินิโย, โย หิ โกจิ ภคินิโย ภิกฺขุ วา ภิกฺขุนี วา จตูสุ สติปฏฺฐาเนสุ สุปฺปติฏฺฐิตจิตฺโต วิหรติ, ตสฺเสตํ ปาฏิกงฺขํ อุฬารํ ปุพฺเพนาปรํ วิเสสํ สญฺชานิสฺสตี”ติ.}

\prose{เอวเมตํ อานนฺท, เอวเมตํ อานนฺท, โย หิ โกจิ อานนฺท ภิกฺขุ วา ภิกฺขุนี วา จตูสุ สติปฏฺฐาเนสุ สุปฺปติฏฺฐิตจิตฺโต วิหรติ, ตสฺเสตํ ปาฏิกงฺขํ อุฬารํ ปุพฺเพนาปรํ วิเสสํ สญฺชานิสฺสติ\footnote{สญฺชานิสฺสตีติ (พหูสุ)}.}

\prose{กตเมสุ จตูสุ, อิธานนฺท ภิกฺขุ กาเย กายานุปสฺสี วิหรติ อาตาปี สมฺปชาโน สติมา วิเนยฺย โลเก อภิชฺฌาโทมนสฺสํ, ตสฺส กาเย กายานุปสฺสิโน วิหรโต กายารมฺมโณ วา อุปฺปชฺชติ กายสฺมิํ ปริฬาโห เจตโส วา ลีนตฺตํ, พหิทฺธา วา จิตฺตํ วิกฺขิปติ.\csromanspacer{} เตนานนฺท\footnote{เตนหานนฺท (สี)} ภิกฺขุนา กิสฺมิญฺจิเทว ปสาทนีเย นิมิตฺเต จิตฺตํ ปณิทหิตพฺพํ, ตสฺส กิสฺมิญฺจิเทว ปสาทนีเย นิมิตฺเต จิตฺตํ ปณิทหโต ปาโมชฺชํ ชายติ, ปมุทิตสฺส ปีติ ชายติ, ปีติมนสฺส กาโย ปสฺสมฺภติ, ปสฺสทฺธกาโย สุขํ เวทยติ\footnote{เวทิยติ (สี)}, สุขิโน จิตฺตํ สมาธิยติ.\csromanspacer{} โส อิติ ปฏิสญฺจิกฺขติ “ยสฺส ขฺวาหํ อตฺถาย จิตฺตํ ปณิทหิํ, โส เม อตฺโถ อภินิปฺผนฺโน, หนฺท ทานิ ปฏิสํหรามี”ติ.\csromanspacer{} โส ปฏิสํหรติ เจว น จ วิตกฺเกติ น จ วิจาเรติ, “อวิตกฺโกมฺหิ อวิจาโร อชฺฌตฺตํ สติมา สุขมสฺมี”ติ ปชานาติ.}

\prose{ปุน จปรํ อานนฺท ภิกฺขุ เวทนาสุ ฯเปฯ.\csromanspacer{} จิตฺเต ฯเปฯ.\csromanspacer{} ธมฺเมสุ ธมฺมานุปสฺสี วิหรติ อาตาปี สมฺปชาโน สติมา วิเนยฺย โลเก อภิชฺฌาโทมนสฺสํ, ตสฺส ธมฺเมสุ ธมฺมานุปสฺสิโน วิหรโต ธมฺมารมฺมโณ วา อุปฺปชฺชติ กายสฺมิํ ปริฬาโห เจตโส วา ลีนตฺตํ, พหิทฺธา วา จิตฺตํ วิกฺขิปติ.\csromanspacer{} เตนานนฺท ภิกฺขุนา กิสฺมิญฺจิเทว ปสาทนีเย นิมิตฺเต จิตฺตํ ปณิทหิตพฺพํ, ตสฺส กิสฺมิญฺจิ เทว ปสาทนีเย นิมิตฺเต จิตฺตํ ปณิทหโต ปาโมชฺชํ}

\prose{\csromanfolio{136}ชายติ, ปมุทิตสฺส ปีติ ชายติ, ปีติมนสฺส กาโย ปสฺสมฺภติ, ปสฺสทฺธกาโย สุขํ เวทยติ, สุขิโน จิตฺตํ สมาธิยติ.\csromanspacer{} โส อิติ ปฏิสญฺจิกฺขติ “ยสฺส ขฺวาหํ อตฺถาย จิตฺตํ ปณิทหิํ, โส เม อตฺโถ อภินิปฺผนฺโน, หนฺท ทานิ ปฏิสํหรามี”ติ.\csromanspacer{} โส ปฏิสํหรติ เจว น จ วิตกฺเกติ น จ วิจาเรติ, “อวิตกฺโกมฺหิ อวิจาโร อชฺฌตฺตํ สติมา สุขมสฺมี”ติ ปชานาติ.\csromanspacer{} เอวํ โข อานนฺท ปณิธาย ภาวนา โหติ.}

\prose{กถญฺจานนฺท อปฺปณิธาย ภาวนา โหติ, พหิทฺธา อานนฺท ภิกฺขุ จิตฺตํ อปฺปณิธาย “อปฺปณิหิตํ เม พหิทฺธา จิตฺตนฺ”ติ ปชานาติ.\csromanspacer{} อถ “ปจฺฉาปุเร อสํขิตฺตํ วิมุตฺตํ อปฺปณิหิตนฺ”ติ ปชานาติ.\csromanspacer{} อถ จ ปน “กาเย กายานุปสฺสี วิหรามิ, อาตาปี สมฺปชาโน สติมา สุขมสฺมี”ติ ปชานาติ.\csromanspacer{} พหิทฺธา อานนฺท ภิกฺขุ จิตฺตํ อปฺปณิธาย “อปฺปณิหิตํ เม พหิทฺธา จิตฺตนฺ”ติ ปชานาติ.\csromanspacer{} อถ “ปจฺฉาปุเร อสํขิตฺตํ วิมุตฺตํ อปฺปณิหิตนฺ”ติ ปชานาติ.\csromanspacer{} อถ จ ปน “เวทนาสุ เวทนานุปสฺสี วิหรามิ, อาตาปี สมฺปชาโน สติมา สุขมสฺมี”ติ ปชานาติ.\csromanspacer{} พหิทฺธา อานนฺท ภิกฺขุ จิตฺตํ อปฺปณิธาย “อปฺปณิหิตํ เม พหิทฺธา จิตฺตนฺ”ติ ปชานาติ.\csromanspacer{} อถ “ปจฺฉา ปุเร อสํขิตฺตํ วิมุตฺตํ อปฺปณิหิตนฺ”ติ ปชานาติ.\csromanspacer{} อถ จ ปน “จิตฺเต จิตฺตานุปสฺสี วิหรามิ, อาตาปี สมฺปชาโน สติมา สุขมสฺมี”ติ ปชานาติ.\csromanspacer{} พหิทฺธา อานนฺท ภิกฺขุ จิตฺตํ อปฺปณิธาย “อปฺปณิหิตํ เม พหิทฺธา จิตฺตนฺ”ติ ปชานาติ.\csromanspacer{} อถ “ปจฺฉา ปุเร อสํขิตฺตํ วิมุตฺตํ อปฺปณิหิตนฺ”ติ ปชานาติ.\csromanspacer{} อถ จ ปน “ธมฺเมสุ ธมฺมานุปสฺสี วิหรามิ, อาตาปี สมฺปชาโน สติมา สุขมสฺมี”ติ ปชานาติ.\csromanspacer{} เอวํ โข อานนฺท อปฺปณิธาย ภาวนา โหติ.}

\prose{อิติ โข อานนฺท เทสิตา มยา ปณิธาย ภาวนา, เทสิตา อปฺปณิธาย ภาวนา.\csromanspacer{} ยํ อานนฺท สตฺถารา กรณียํ สาวกานํ หิเตสินา อนุกมฺปเกน อนุกมฺปํ อุปาทาย, กตํ โว ตํ มยา.\csromanspacer{} เอตานิ อานนฺท รุกฺขมูลานิ, เอตานิ สุญฺญาคารานิ.\csromanspacer{} ฌายถานนฺท มา ปมาทตฺถ, มา ปจฺฉา วิปฺปฏิสาริโน อหุวตฺถ, อยํ โว อมฺหากํ อนุสาสนีติ.}

\prose{อิทมโวจ ภควา.\csromanspacer{} อตฺตมโน อายสฺมา อานนฺโท ภควโต ภาสิตํ อภินนฺทีติ.\csromanspacer{}\csromanspacer{}ทสมํ.\csromanspacer{} อมฺพปาลิวคฺโค ปฐโม.}

\csromanheader{2}{3. สติปฏฺฐานสํยุตฺต}
\titlehead{\textbf{ตสฺสุทฺทานํ}}
\csromangathameasure{อมฺพปาลิ สโต ภิกฺขุ, สาลา’กุสลราสิ จ. \\ สกุณคฺฆิ มกฺกโฏ สูโท, คิลาโน ภิกฺขุนุปสฺสโยติ.}
\csromangathagroup{\csromangathastanza{\csromanfolio{137}อมฺพปาลิ สโต ภิกฺขุ, สาลา’กุสลราสิ จ. \\ สกุณคฺฆิ มกฺกโฏ สูโท, คิลาโน ภิกฺขุนุปสฺสโยติ.}}

\csromanmatikaanchor{mtk.56}
\csromantocmarknopage{paragraph}{2. นาลนฺทวคฺค}{mtk.56}
\bookmark[dest={mtk.56},level=4]{2. นาลนฺทวคฺค}
\csromanheader{4}{2. นาลนฺทวคฺค}
\csromanmatikaanchor{mtk.57}
\csromantocmarkref{subparagraph}{1. มหาปุริสสุตฺต}{mtk.57}
\bookmark[dest={mtk.57},level=5]{1. มหาปุริสสุตฺต}
\csromanheader{5}{1. มหาปุริสสุตฺต}
\proseitem{377}{สาวตฺถินิทานํ.\csromanspacer{} อถ โข อายสฺมา สาริปุตฺโต เยน ภควา เตนุปสงฺกมิ, อุปสงฺกมิตฺวา ภควนฺตํ อภิวาเทตฺวา เอกมนฺตํ นิสีทิ, เอกมนฺตํ นิสินฺโน โข อายสฺมา สาริปุตฺโต ภควนฺตํ เอตทโวจ “มหาปุริโส มหาปุริโสติ ภนฺเต วุจฺจติ, กิตฺตาวตา นุ โข ภนฺเต มหาปุริโส โหตี”ติ.\csromanspacer{} วิมุตฺตจิตฺตตฺตา ขฺวาหํ สาริปุตฺต “มหาปุริโส”ติ วทามิ, อวิมุตฺตจิตฺตตฺตา “โน มหาปุริโส”ติ วทามิ.}

\prose{กถญฺจ สาริปุตฺต วิมุตฺตจิตฺโต โหติ, อิธ สาริปุตฺต ภิกฺขุ กาเย กายานุปสฺสี วิหรติ อาตาปี สมฺปชาโน สติมา วิเนยฺย โลเก อภิชฺฌาโทมนสฺสํ, ตสฺส กาเย กายานุปสฺสิโน วิหรโต จิตฺตํ วิรชฺชติ วิมุจฺจติ อนุปาทาย อาสเวหิ.\csromanspacer{} เวทนาสุ ฯเปฯ.\csromanspacer{} จิตฺเต ฯเปฯ.\csromanspacer{} ธมฺเมสุ ธมฺมานุปสฺสี วิหรติ อาตาปี สมฺปชาโน สติมา วิเนยฺย โลเก อภิชฺฌาโทมนสฺสํ, ตสฺส ธมฺเมสุ ธมฺมานุปสฺสิโน วิหรโต จิตฺตํ วิรชฺชติ วิมุจฺจติ อนุปาทาย อาสเวหิ.\csromanspacer{} เอวํ โข สาริปุตฺต วิมุตฺตจิตฺโต โหติ, วิมุตฺตจิตฺตตฺตา ขฺวาหํ สาริปุตฺต “มหาปุริโส”ติ วทามิ, อวิมุตฺตจิตฺตตฺตา “โน มหาปุริโส”ติ วทามีติ.\csromanspacer{}\csromanspacer{}ปฐมํ.}

\csromanmatikaanchor{mtk.58}
\csromantocmarkref{subparagraph}{2. นาลนฺทสุตฺต}{mtk.58}
\bookmark[dest={mtk.58},level=5]{2. นาลนฺทสุตฺต}
\csromanheader{5}{2. นาลนฺทสุตฺต}
\proseitem{378}{เอกํ สมยํ ภควา นาลนฺทายํ วิหรติ ปาวาริกมฺพวเน.\csromanspacer{} อถ โข อายสฺมา สาริปุตฺโต เยน ภควา เตนุปสงฺกมิ, อุปสงฺกมิตฺวา ภควนฺตํ อภิวาเทตฺวา เอกมนฺตํ นิสีทิ, เอกมนฺตํ นิสินฺโน โข อายสฺมา สาริปุตฺโต ภควนฺตํ เอตทโวจ “เอวํปสนฺโน อหํ ภนฺเต ภควติ, น จาหุ \csromanfolio{138}น จ ภวิสฺสติ น เจตรหิ วิชฺชติ อญฺโญ สมโณ วา พฺราหฺมโณ วา ภควตา ภิยฺโยภิญฺญตโร ยทิทํ สมฺโพธิยนฺ”ติ.\csromanspacer{} อุฬารา โข ตฺยายํ สาริปุตฺต อาสภี วาจา ภาสิตา, เอกํโส คหิโต, สีหนาโท นทิโต “เอวํปสนฺโน อหํ ภนฺเต ภควติ, น จาหุ น จ ภวิสฺสติ น เจตรหิ วิชฺชติ อญฺโญ สมโณ วา พฺราหฺมโณ วา ภควตา ภิยฺโยภิญฺญตโร ยทิทํ สมฺโพธิยนฺ”ติ.}

\prose{กิํ นุ เต สาริปุตฺต เยเต อเหสุํ อตีตมทฺธานํ อรหนฺโต สมฺมาสมฺพุทฺธา, สพฺเพเต ภควนฺโต เจตสา เจโต ปริจฺจ วิทิตา “เอวํสีลา เต ภควนฺโต อเหสุํ อิติ วา, เอวํธมฺมา เต ภควนฺโต อเหสุํ อิติ วา, เอวํปญฺญา เต ภควนฺโต อเหสุํ อิติ วา, เอวํวิหาริโน เต ภควนฺโต อเหสุํ อิติ วา, เอวํวิมุตฺตา เต ภควนฺโต อเหสุํ อิติ วา”ติ.\csromanspacer{} โน เหตํ ภนฺเต.}

\prose{กิํ ปน เต สาริปุตฺต เย เต ภวิสฺสนฺติ อนาคตมทฺธานํ อรหนฺโต สมฺมาสมฺพุทฺธา, สพฺเพเต ภควนฺโต เจตสา เจโต ปริจฺจ วิทิตา “เอวํสีลา เต ภควนฺโต ภวิสฺสนฺติ อิติ วา, เอวํธมฺมา เต ภควนฺโต ภวิสฺสนฺติ อิติ วา, เอวํปญฺญา เต ภควนฺโต ภวิสฺสนฺติ อิติ วา, เอวํวิหาริโน เต ภควนฺโต ภวิสฺสนฺติ อิติ วา, เอวํวิมุตฺตา เต ภควนฺโต ภวิสฺสนฺติ อิติ วา”ติ.\csromanspacer{} โน เหตํ ภนฺเต.}

\prose{กิํ ปน ตฺยาหํ\footnote{กิํ ปน เต (สี)} สาริปุตฺต เอตรหิ อรหํ สมฺมาสมฺพุทฺโธ เจตสา เจโต ปริจฺจ วิทิโต “เอวํสีโล ภควา อิติ วา, เอวํธมฺโม ภควา อิติ วา, เอวํปญฺโญ ภควา อิติ วา, เอวํวิหารี ภควา อิติ วา, เอวํวิมุตฺโต ภควา อิติ วา”ติ.\csromanspacer{} โน เหตํ ภนฺเต.}

\prose{เอตฺถ จ เต สาริปุตฺต อตีตานาคตปจฺจุปฺปนฺเนสุ อรหนฺเตสุ สมฺมาสมฺพุทฺเธสุ เจโตปริยญาณํ\footnote{เจโตปริยายญาณํ (พหูสุ)} นตฺถิ, อถ กิญฺจรหิ ตฺยายํ สาริปุตฺต อุฬารา อาสภี วาจา ภาสิตา, เอกํโส คหิโต, สีหนาโท นทิโต “เอวํปสนฺโน อหํ ภนฺเต ภควติ, น จาหุ น จ ภวิสฺสติ น เจตรหิ วิชฺชติ อญฺโญ สมโณ วา พฺราหฺมโณ วา ภควตา ภิยฺโยภิญฺญตโร ยทิทํ สมฺโพธิยนฺ”ติ.}

\csromanheader{2}{3. สติปฏฺฐานสํยุตฺต}
\prose{\csromanfolio{139}น โข เม\footnote{น โข เม ตํ (สฺยา, กํ, ก)} ภนฺเต อตีตานาคตปจฺจุปฺปนฺเนสุ อรหนฺเตสุ สมฺมาสมฺพุทฺเธสุ เจโตปริยญาณํ อตฺถิ, อปิ จ เม ธมฺมนฺวโย วิทิโต.\csromanspacer{} เสยฺยถาปิ ภนฺเต รญฺโญ ปจฺจนฺติมํ นครํ ทฬฺหุทฺธาปํ\footnote{ทฬฺหุทฺทาปํ (สี, อิ, ก), ทฬฺหทฺธาปํ (สฺยา, กํ)} ทฬฺหปาการโตรณํ เอกทฺวารํ, ตตฺรสฺส โทวาริโก ปณฺฑิโต พฺยตฺโต เมธาวี อญฺญาตานํ นิวาเรตา ญาตานํ ปเวเสตา.\csromanspacer{} โส ตสฺส นครสฺส สมนฺตา อนุปริยายปถํ อนุกฺกมมาโน น ปสฺเสยฺย ปาการสนฺธิํ วา ปาการวิวรํ วา อนฺตมโส พิฬารนิกฺขมนมตฺตมฺปิ.\csromanspacer{} ตสฺส เอวมสฺส “เย โข เกจิ โอฬาริกา ปาณา อิมํ นครํ ปวิสนฺติ วา นิกฺขมนฺติ วา, สพฺเพเต อิมินาว ทฺวาเรน ปวิสนฺติ วา นิกฺขมนฺติ วา’ติ.\csromanspacer{} เอวเมว โข เม ภนฺเต ธมฺมนฺวโย วิทิโต, เยปิ เต ภนฺเต อเหสุํ อตีตมทฺธานํ อรหนฺโต สมฺมาสมฺพุทฺธา, สพฺเพเต ภควนฺโต ปญฺจ นีวรเณ ปหาย เจตโส อุปกฺกิเลเส ปญฺญาย ทุพฺพลีกรเณ จตูสุ สติปฏฺฐาเนสุ สุปฺปติฏฺฐิตจิตฺตา สตฺต โพชฺฌงฺเค ยถาภูตํ ภาเวตฺวา อนุตฺตรํ สมฺมาสมฺโพธิํ อภิสมฺพุชฺฌิํสุ.\csromanspacer{} เยปิ เต ภนฺเต ภวิสฺสนฺติ อนาคตมทฺธานํ อรหนฺโต สมฺมาสมฺพุทฺธา, สพฺเพเต ภควนฺโต ปญฺจ นีวรเณ ปหาย เจตโส อุปกฺกิเลเส ปญฺญาย ทุพฺพลีกรเณ จตูสุ สติปฏฺฐาเนสุ สุปฺปติฏฺฐิตจิตฺตา สตฺต โพชฺฌงฺเค ยถาภูตํ ภาเวตฺวา อนุตฺตรํ สมฺมาสมฺโพธิํ อภิสมฺพุชฺฌิสฺสนฺติ.\csromanspacer{} ภควาปิ ภนฺเต เอตรหิ อรหํ สมฺมาสมฺพุทฺโธ ปญฺจ นีวรเณ ปหาย เจตโส อุปกฺกิเลเส ปญฺญาย ทุพฺพลีกรเณ จตูสุ สติปฏฺฐาเนสุ สุปฺปติฏฺฐิตจิตฺโต สตฺต โพชฺฌงฺเค ยถาภูตํ ภาเวตฺวา อนุตฺตรํ สมฺมาสมฺโพธิํ อภิสมฺพุทฺโธติ.}

\prose{สาธุ สาธุ สาริปุตฺต, ตสฺมาติห ตฺวํ สาริปุตฺต อิมํ ธมฺมปริยายํ อภิกฺขณํ ภาเสยฺยาสิ ภิกฺขูนํ ภิกฺขุนีนํ อุปาสกานํ อุปาสิกานํ.\csromanspacer{} เยสมฺปิ หิ สาริปุตฺต โมฆปุริสานํ ภวิสฺสติ ตถาคเต กงฺขา วา วิมติ วา, เตสมฺปิมํ ธมฺมปริยายํ สุตฺวา ยา ตถาคเต กงฺขา วา วิมติ วา, สา ปหียิสฺสตีติ.\csromanspacer{}\csromanspacer{}ทุติยํ.}

\csromanmatikaanchor{mtk.59}
\csromantocmarkref{subparagraph}{3. จุนฺทสุตฺต}{mtk.59}
\bookmark[dest={mtk.59},level=5]{3. จุนฺทสุตฺต}
\csromanheader{5}{3. จุนฺทสุตฺต}
\proseitem{379}{เอกํ สมยํ ภควา สาวตฺถิยํ วิหรติ เชตวเน อนาถปิณฺฑิกสฺส อาราเม.\csromanspacer{} เตน โข ปน สมเยน อายสฺมา สาริปุตฺโต \csromanfolio{140}มคเธสุ วิหรติ นาลกคามเก อาพาธิโก ทุกฺขิโต พาฬฺหคิลาโน.\csromanspacer{} จุนฺโท จ สมณุทฺเทโส อายสฺมโต สาริปุตฺตสฺส อุปฏฺฐาโก โหติ.}

\prose{อถ โข อายสฺมา สาริปุตฺโต เตเนว อาพาเธน ปรินิพฺพายิ.\csromanspacer{} อถ โข จุนฺโท สมณุทฺเทโส อายสฺมโต สาริปุตฺตสฺส ปตฺตจีวรมาทาย เยน สาวตฺถิ เชตวนํ อนาถปิณฺฑิกสฺส อาราโม, เยนายสฺมา อานนฺโท เตนุปสงฺกมิ, อุปสงฺกมิตฺวา อายสฺมนฺตํ อานนฺทํ อภิวาเทตฺวา เอกมนฺตํ นิสีทิ, เอกมนฺตํ นิสินฺโน โข จุนฺโท สมณุทฺเทโส อายสฺมนฺตํ อานนฺทํ เอตทโวจ “อายสฺมา ภนฺเต สาริปุตฺโต ปรินิพฺพุโต, อิทมสฺส ปตฺตจีวรนฺ”ติ.\csromanspacer{} อตฺถิ โข อิทํ อาวุโส จุนฺท กถาปาภตํ ภควนฺตํ ทสฺสนาย, อายามาวุโส จุนฺท, เยน ภควา เตนุปสงฺกมิสฺสาม, อุปสงฺกมิตฺวา ภควโต เอตมตฺถํ อาโรเจสฺสามาติ. “เอวํ ภนฺเต”ติ โข จุนฺโท สมณุทฺเทโส อายสฺมโต อานนฺทสฺส ปจฺจสฺโสสิ.}

\prose{อถ โข อายสฺมา จ อานนฺโท จุนฺโท จ สมณุทฺเทโส เยน ภควา เตนุปสงฺกมิํสุ, อุปสงฺกมิตฺวา ภควนฺตํ อภิวาเทตฺวา เอกมนฺตํ นิสีทิํสุ, เอกมนฺตํ นิสินฺโน โข อายสฺมา อานนฺโท ภควนฺตํ เอตทโวจ “อยํ ภนฺเต จุนฺโท สมณุทฺเทโส เอวมาห ‘อยสฺมา ภนฺเต สาริปุตฺโต ปรินิพฺพุโต, อิทมสฺส ปตฺตจีวรนฺ’ติ.\csromanspacer{} อปิ จ เม ภนฺเต มธุรกชาโต วิย กาโย, ทิสาปิ เม น ปกฺขายนฺติ, ธมฺมาปิ มํ นปฺปฏิภนฺติ ‘อายสฺมา สาริปุตฺโต ปรินิพฺพุโต’ติ สุตฺวา”.}

\prose{กิํ นุ โข เต อานนฺท สาริปุตฺโต สีลกฺขนฺธํ วา อาทาย ปรินิพฺพุโต, สมาธิกฺขนฺธํ วา อาทาย ปรินิพฺพุโต, ปญฺญากฺขนฺธํ วา อาทาย ปรินิพฺพุโต, วิมุตฺติกฺขนฺธํ วา อาทาย ปรินิพฺพุโต, วิมุตฺติญาณทสฺสนกฺขนฺธํ วา อาทาย ปรินิพฺพุโตติ.\csromanspacer{} น จ โข เม ภนฺเต อายสฺมา สาริปุตฺโต สีลกฺขนฺธํ วา อาทาย ปรินิพฺพุโต, สมาธิกฺขนฺธํ วา ฯเปฯ ปญฺญากฺขนฺธํ วา ฯเปฯ วิมุตฺติกฺขนฺธํ วา ฯเปฯ วิมุตฺติญาณทสฺสนกฺขนฺธํ วา อาทาย ปรินิพฺพุโต, อปิ จ เม ภนฺเต อายสฺมา สาริปุตฺโต โอวาทโก อโหสิ โอติณฺโณ วิญฺญาปโก สนฺทสฺสโก สมาทปโก สมุตฺเตชโก สมฺปหํสโก, อกิลาสุ ธมฺมเทสนาย, อนุคฺคาหโก สพฺรหฺมจารีนํ, ตํ มยํ อายสฺมโต สาริปุตฺตสฺส ธมฺโมชํ ธมฺมโภคํ ธมฺมานุคฺคหํ อนุสฺสรามาติ.}

\csromanheader{2}{3. สติปฏฺฐานสํยุตฺต}
\prose{\csromanfolio{141}นนุ ตํ อานนฺท มยา ปฏิกจฺเจว\footnote{ปฏิคจฺเจว (สี, อิ)} อกฺขาตํ, สพฺเพหิ ปิเยหิ มนาเปหิ นานาภาโว วินาภาโว อญฺญถาภาโว, ตํ กุเตตฺถ อานนฺท ลพฺภา “ยํ ตํ ชาตํ ภูตํ สงฺขตํ ปโลกธมฺมํ, ตํ วต มา ปลุชฺชี”ติ เนตํ ฐานํ วิชฺชติ.\csromanspacer{} เสยฺยถาปิ อานนฺท มหโต รุกฺขสฺส ติฏฺฐโต สารวโต โย มหนฺตตโร ขนฺโธ โส ปลุชฺเชยฺย.\csromanspacer{} เอวเมว โข อานนฺท มหโต ภิกฺขุสํฆสฺส ติฏฺฐโต สารวโต สาริปุตฺโต ปรินิพฺพุโต, ตํ กุเตตฺถ อานนฺท ลพฺภา “ยํ ตํ ชาตํ ภูตํ สงฺขตํ ปโลกธมฺมํ, ตํ วต มา ปลุชฺชี”ติ เนตํ ฐานํ วิชฺชติ.\csromanspacer{} ตสฺมาติหานนฺท อตฺตทีปา วิหรถ อตฺตสรณา อนญฺญสรณา ธมฺมทีปา ธมฺมสรณา อนญฺญสรณา.}

\prose{กถญฺจานนฺท ภิกฺขุ อตฺตทีโป วิหรติ อตฺตสรโณ อนญฺญสรโณ ธมฺมทีโป ธมฺมสรโณ อนญฺญสรโณ, อิธานนฺท ภิกฺขุ กาเย กายานุปสฺสี วิหรติ อาตาปี สมฺปชาโน สติมา วิเนยฺย โลเก อภิชฺฌาโทมนสฺสํ.\csromanspacer{} เวทนาสุ ฯเปฯ.\csromanspacer{} จิตฺเต ฯเปฯ.\csromanspacer{} ธมฺเมสุ ธมฺมานุปสฺสี วิหรติ อาตาปี สมฺปชาโน สติมา วิเนยฺย โลเก อภิชฺฌาโทมนสฺสํ.\csromanspacer{} เอวํ โข อานนฺท ภิกฺขุ อตฺตทีโป วิหรติ อตฺตสรโณ อนญฺญสรโณ ธมฺมทีโป ธมฺมสรโณ อนญฺญสรโณ.\csromanspacer{} เย หิ เกจิ อานนฺท เอตรหิ วา มมจฺจเย วา อตฺตทีปา วิหริสฺสนฺติ อตฺตสรณา อนญฺญสรณา ธมฺมทีปา ธมฺมสรณา อนญฺญสรณา, ตมตคฺเค เมเต อานนฺท ภิกฺขู ภวิสฺสนฺติ, เย เกจิ สิกฺขากามาติ.\csromanspacer{}\csromanspacer{}ตติยํ.}

\csromanmatikaanchor{mtk.60}
\csromantocmarkref{subparagraph}{4. อุกฺกเจลสุตฺต}{mtk.60}
\bookmark[dest={mtk.60},level=5]{4. อุกฺกเจลสุตฺต}
\csromanheader{5}{4. อุกฺกเจลสุตฺต}
\proseitem{380}{เอกํ สมยํ ภควา วชฺชีสุ วิหรติ อุกฺกเจลายํ คงฺคาย นทิยา ตีเร มหตา ภิกฺขุสํเฆน สทฺธิํ อจิรปรินิพฺพุเตสุ สาริปุตฺตโมคฺคลฺลาเนสุ.\csromanspacer{} เตน โข ปน สมเยน ภควา ภิกฺขุสํฆปริวุโต อชฺโฌกาเส นิสินฺโน โหติ.}

\prose{อถ โข ภควา ตุณฺหีภูตํ ภิกฺขุสํฆํ อนุวิโลเกตฺวา ภิกฺขู อามนฺเตสิ “อปิ มฺยายํ ภิกฺขเว ปริสา สุญฺญา วิย ขายติ ปรินิพฺพุเตสุ สาริปุตฺตโมคฺคลฺลาเนสุ.\csromanspacer{} อสุญฺญา เม ภิกฺขเว ปริสา \csromanfolio{142}โหติ, อนเปกฺขา ตสฺสํ ทิสายํ โหติ, ยสฺสํ ทิสายํ สาริปุตฺตโมคฺคลฺลานา วิหรนฺติ.\csromanspacer{} เย หิ เต ภิกฺขเว อเหสุํ อตีตมทฺธานํ อรหนฺโต สมฺมาสมฺพุทฺธา, เตสมฺปิ ภควนฺตานํ เอตปฺปรมํเยว\footnote{เอตปรมํเยว (สี, สฺยา, กํ, อิ)} สาวกฺยุคํ อโหสิ, เสยฺยถาปิ มยฺหํ สาริปุตฺตโมคฺคลฺลานา.\csromanspacer{} เยปิ เต ภิกฺขเว ภวิสฺสนฺติ อนาคตมทฺธานํ อรหนฺโต สมฺมาสมฺพุทฺธา, เตสมฺปิ ภควนฺตานํ เอตปฺปรมํเยว สาวกฺยุคํ ภวิสฺสติ, เสยฺยถาปิ มยฺหํ สาริปุตฺตโมคฺคลฺลานา.\csromanspacer{} อจฺฉริยํ ภิกฺขเว สาวกานํ, อพฺภุตํ ภิกฺขเว สาวกานํ, สตฺถุ จ นาม สาสนกรา ภวิสฺสนฺติ โอวาทปฺปฏิกรา, จตุนฺนญฺจ ปริสานํ ปิยา ภวิสฺสนฺติ มนาปา ครุภาวนียา จ.\csromanspacer{} อจฺฉริยํ ภิกฺขเว ตถาคตสฺส, อพฺภุตํ ภิกฺขเว ตถคตสฺส, เอวรูเปปิ นาม สาวกฺยุเค ปรินิพฺพุเต นตฺถิ ตถาคตสฺส โสโก วา ปริเทโว วา, ตํ กุเตตฺถ ภิกฺขเว ลพฺภา “ยํ ตํ ชาตํ ภูตํ สงฺขตํ ปโลกธมฺมํ, ตํ วต มา ปลุชฺชี”ติ เนตํ ฐานํ วิชฺชติ.\csromanspacer{} เสยฺยถาปิ ภิกฺขเว มหโต รุกฺขสฺส ติฏฺฐโต สารวโต เย มหนฺตตรา ขนฺธา, เต ปลุชฺเชยฺยุํ.\csromanspacer{} เอวเมว โข ภิกฺขเว มหโต ภิกฺขุสํฆสฺส ติฏฺฐโต สารวโต สาริปุตฺตโมคฺคลฺลานา ปรินิพฺพุตา, ตํ กุเตตฺถ ภิกฺขเว ลพฺภา “ยํ ตํ ชาตํ ภูตํ สงฺขตํ ปโลกธมฺมํ, ตํ วต มา ปลุชฺชี”ติ เนตํ ฐานํ วิชฺชติ.\csromanspacer{} ตสฺมาติห ภิกฺขเว อตฺตทีปา วิหรถ อตฺตสรณา อนญฺญสรณา ธมฺมทีปา ธมฺมสรณา อนญฺญสรณา.}

\prose{กถญฺจ ภิกฺขเว ภิกฺขุ อตฺตทีโป วิหรติ อตฺตสรโณ อนญฺญสรโณ ธมฺมทีโป ธมฺมสรโณ อนญฺญสรโณ, อิธ ภิกฺขเว ภิกฺขุ กาเย กายานุปสฺสี วิหรติ อาตาปี สมฺปชาโน สติมา วิเนยฺย โลเก อภิชฺฌาโทมนสฺสํ.\csromanspacer{} เวทนาสุ ฯเปฯ.\csromanspacer{} จิตฺเต ฯเปฯ.\csromanspacer{} ธมฺเมสุ ธมฺมานุปสฺสี วิหรติ อาตาปี สมฺปชาโน สติมา วิเนยฺย โลเก อภิชฺฌาโทมนสฺสํ.\csromanspacer{} เอวํ โข ภิกฺขเว ภิกฺขุ อตฺตทีโป วิหรติ อตฺตสรโณ อนญฺญสรโณ ธมฺมทีโป ธมฺมสรโณ อนญฺญสรโณ.\csromanspacer{} เย หิ เกจิ ภิกฺขเว เอตรหิ วา มมจฺจเย วา อตฺตทีปา วิหริสฺสนฺติ อตฺตสรณา อนญฺญสรณา ธมฺมทีปา ธมฺมสรณา อนญฺญสรณา.\csromanspacer{} ตมตคฺเค เมเต ภิกฺขเว ภิกฺขู ภวิสฺสนฺติ, เย เกจิ สิกฺขากามาติ.\csromanspacer{}\csromanspacer{}จตุตฺถํ.}

\csromanheader{2}{3. สติปฏฺฐานสํยุตฺต}
\csromanheader{2}{5. พาหิยสุตฺต}
\csromanmatikaanchor{mtk.61}
\csromantocmarkref{subparagraph}{5. ภาหิยสุตฺต}{mtk.61}
\bookmark[dest={mtk.61},level=5]{5. ภาหิยสุตฺต}
\proseitem{381}{\csromanfolio{143}สาวตฺถินิทานํ.\csromanspacer{} อถ โข อายสฺมา พาหิโย เยน ภควา เตนุปสงฺกมิ, อุปสงฺกมิตฺวา ภควนฺตํ อภิวาเทตฺวา เอกมนฺตํ นิสีทิ, เอกมนฺตํ นิสินฺโน โข อายสฺมา พาหิโย ภควนฺตํ เอตทโวจ “สาธุ เม ภนฺเต ภควา สํขิตฺเตน ธมฺมํ เทเสตุ, ยมหํ ภควโต ธมฺมํ สุตฺวา เอโก วูปกฏฺโฐ อปฺปมตฺโต อาตาปี ปหิตตฺโต วิหเรยฺยนฺ”ติ.\csromanspacer{} ตสฺมาติห ตฺวํ พาหิย อาทิเมว วิโสเธหิ กุสเลสุ ธมฺเมสุ.\csromanspacer{} โก จาทิ กุสลานํ ธมฺมานํ, สีลญฺจ สุวิสุทฺธํ ทิฏฺฐิ จ อุชุกา.\csromanspacer{} ยโต จ โข เต พาหิย สีลญฺจ สุวิสุทฺธํ ภวิสฺสติ ทิฏฺฐิ จ อุชุกา, ตโต ตฺวํ พาหิย สีลํ นิสฺสาย สีเล ปติฏฺฐาย จตฺตาโร สติปฏฺฐาเน ภาเวยฺยาสิ.}

\prose{กตเม จตฺตาโร, อิธ ตฺวํ พาหิย กาเย กายานุปสฺสี วิหราหิ อาตาปี สมฺปชาโน สติมา วิเนยฺย โลเก อภิชฺฌาโทมนสฺสํ.\csromanspacer{} เวทนาสุ ฯเปฯ.\csromanspacer{} จิตฺเต ฯเปฯ.\csromanspacer{} ธมฺเมสุ ธมฺมานุปสฺสี วิหราหิ อาตาปี สมฺปชาโน สติมา วิเนยฺย โลเก อภิชฺฌาโทมนสฺสํ.\csromanspacer{} ยโต โข ตฺวํ พาหิย สีลํ นิสฺสาย สีเล ปติฏฺฐาย อิเม จตฺตาโร สติปฏฺฐาเน เอวํ ภาเวสฺสสิ, ตโต ตุยฺหํ พาหิย ยา รตฺติ วา ทิวโส วา อาคมิสฺสติ, วุทฺธิเยว ปาฏิกงฺขา กุสเลสฺสุ ทมฺเมสุ, โน ปริหานีติ.}

\prose{อถ โข อายสฺมา พาหิโย ภควโต ภาสิตํ อภินนฺทิตฺวา อนุโมทิตฺวา อุฏฺฐายาสนา ภควนฺตํ อภิวาเทตฺวา ปทกฺขิณํ กตฺวา ปกฺกามิ.\csromanspacer{} อถ โข อายสฺมา พาหิโย เอโก วูปกฏฺโฐ อปฺปมตฺโต อาตาปี ปหิตตฺโต วิหรนฺโต นจิรสฺเสว ยสฺสตฺถาย กุลปุตฺตา สมฺมเทว อคารสฺมา อนคาริยํ ปพฺพชนฺติ, ตทนุตฺตรํ พฺรหฺมจริยปริโยสานํ ทิฏฺเฐว ธมฺเม สยํ อภิญฺญา สจฺฉิกตฺวา อุปสมฺปชฺช วิหาสิ, “ขีณา ชาติ, วุสิตํ พฺรหฺมจริยํ, กตํ กรณียํ, นาปรํ อิตฺถตฺตายา”ติ อพฺภญฺญาสิ.\csromanspacer{} อญฺญตโร จ ปนายสฺมา พาหิโย อรหตํ อโหสีติ.\csromanspacer{}\csromanspacer{}ปญฺจมํ.}

\csromanmatikaanchor{mtk.62}
\csromantocmarkref{subparagraph}{6. อุตฺติยสุตฺต}{mtk.62}
\bookmark[dest={mtk.62},level=5]{6. อุตฺติยสุตฺต}
\csromanheader{5}{6. อุตฺติยสุตฺต}
\proseitem{382}{สาวตฺถินิทานํ.\csromanspacer{} อถ โข อายสฺมา อุตฺติโย เยน ภควา เตนุปสงฺกมิ ฯเปฯ เอกมนฺตํ นิสินฺโน โข อายสฺมา อุตฺติโย ภควนฺตํ เอตทโวจ “สาธุ เม ภนฺเต ภควา สํขิตฺเตน ธมฺมํ เทเสตุ, ยมหํ \csromanfolio{144}ภควโต ธมฺมํ สุตฺวา เอโก วูปกฏฺโฐ อปฺปมตฺโต อาตาปี ปหิตตฺโต วิหเรยฺยนฺ”ติ.\csromanspacer{} ตสฺมาติห ตฺวํ อุตฺติย อาทิเมว วิโสเธหิ กุสเลสุ ธมฺเมสุ.\csromanspacer{} โก จาทิ กุสลานํ ธมฺมานํ, สีลญฺจ สุวิสุทฺธํ ทิฏฺฐิ จ อุชุกา.\csromanspacer{} ยโต จ โข เต อุตฺติย สีลญฺจ สุวิสุทฺธํ ภวิสฺสติ ทิฏฺฐิ จ อุชุกา, ตโต ตฺวํ อุตฺติย สีลํ นิสฺสาย สีเล ปติฏฺฐาย จตฺตาโร สติปฏฺฐาเน ภาเวยฺยาสิ.}

\prose{กตเม จตฺตาโร, อิธ ตฺวํ อุตฺติย กาเย กายานุปสฺสี วิหราหิ อาตาปี สมฺปชาโน สติมา วิเนยฺย โลเก อภิชฺฌาโทมนสฺสํ.\csromanspacer{} เวทนาสุ ฯเปฯ.\csromanspacer{} จิตฺเต ฯเปฯ.\csromanspacer{} ธมฺเมสุ ธมฺมานุปสฺสี วิหราหิ อาตาปี สมฺปชาโน สติมา วิเนยฺย โลเก อภิชฺฌาโทมนสฺสํ.\csromanspacer{} ยโต โข ตฺวํ อุตฺติย สีลํ นิสฺสาย สีเล ปติฏฺฐาย อิเม จตฺตาโร สติปฏฺฐาเน เอวํ ภาเวสฺสสิ, ตโต ตฺวํ อุตฺติย คมิสฺสสิ มจฺจุเธยฺยสฺส ปารนฺติ.}

\prose{อถ โข อายสฺมา อุตฺติโย ภควโต ภาสิตํ อภินนฺทิตฺวา อนุโมทิตฺวา อุฏฺฐายาสนา ภควนฺตํ อภิวาเทตฺวา ปทกฺขิณํ กตฺวา ปกฺกามิ.\csromanspacer{} อถ โข อายสฺมา อุตฺติโย เอโก วูปกฏฺโฐ อปฺปมตฺโต อาตาปี ปหิตตฺโต วิหรนฺโต นจิรสฺเสว ยสฺสตฺถาย กุลปุตฺตา สมฺมเทว อคารสฺมา อนคาริยํ ปพฺพชนฺติ, ตทนุตฺตรํ พฺรหฺมจริยปริโยสานํ ทิฏฺเฐว ธมฺเม สยํ อภิญฺญา สจฺฉิกตฺวา อุปสมฺปชฺช วิหาสิ, “ขีณา ชาติ, วุสิตํ พฺรหฺมจริยํ, กตํ กรณียํ, นาปรํ อิตฺถตฺตายา”ติ อพฺภญฺญาสิ.\csromanspacer{} อญฺญตโร จ ปนายสฺมา อุตฺติโย อรหตํ อโหสีติ.\csromanspacer{}\csromanspacer{}ฉฏฺฐํ.}

\csromanmatikaanchor{mtk.63}
\csromantocmarkref{subparagraph}{7. อริยสุตฺต}{mtk.63}
\bookmark[dest={mtk.63},level=5]{7. อริยสุตฺต}
\csromanheader{5}{7. อริยสุตฺต}
\proseitem{383}{จตฺตาโรเม ภิกฺขเว สติปฏฺฐานา ภาวิตา พหุลีกตา อริยา นิยฺยานิกา นิยฺยนฺติ ตกฺกรสฺส สมฺมา ทุกฺขกฺขยาย.\csromanspacer{} กตเม จตฺตาโร, อิธ ภิกฺขเว ภิกฺขุ กาเย กายานุปสฺสี วิหรติ อาตาปี สมฺปชาโน สติมา วิเนยฺย โลเก อภิชฺฌาโทมนสฺสํ.\csromanspacer{} เวทนาสุ ฯเปฯ.\csromanspacer{} จิตฺเต ฯเปฯ.\csromanspacer{} ธมฺเมสุ ธมฺมานุปสฺสี วิหรติ อาตาปี สมฺปชาโน สติมา วิเนยฺย โลเก อภิชฺฌาโทมนสฺสํ.\csromanspacer{} อิเม โข ภิกฺขเว จตฺตาโร สติปฏฺฐานา ภาวิตา พหุลีกตา อริยา นิยฺยานิกา นิยฺยนฺติ ตกฺกรสฺส สมฺมา ทุกฺขกฺขยายาติ.\csromanspacer{}\csromanspacer{}สตฺตมํ.}

\csromanheader{2}{3. สติปฏฺฐานสํยุตฺต}
\csromanmatikaanchor{mtk.64}
\csromantocmarkref{subparagraph}{8. พฺรหฺมสุตฺต}{mtk.64}
\bookmark[dest={mtk.64},level=5]{8. พฺรหฺมสุตฺต}
\csromanheader{5}{8. พฺรหฺมสุตฺต}
\proseitem{384}{\csromanfolio{145}เอกํ สมยํ ภควา อุรุเวลายํ วิหรติ นชฺชา เนรญฺชราย ตีเร อชปาลนิโคฺรเธ ปฐมาภิสมฺพุทฺโธ.\csromanspacer{} อถ โข ภควโต รโหคตสฺส ปฏิสลฺลีนสฺส เอวํ เจตโส ปริวิตกฺโก อุทปาทิ\csromandash{}เอกายนฺวายํ มคฺโค สตฺตานํ วิสุทฺธิยา โสกปริเทวานํ สมติกฺกมาย ทุกฺขโทมนสฺสานํ อตฺถงฺคมาย ญายสฺส อธิคมาย นิพฺพานสฺส สจฺฉิกิริยาย, ยทิทํ จตฺตาโร สติปฏฺฐานา.}

\prose{กตเม จตฺตาโร, กาเย วา ภิกฺขุ กายานุปสฺสี วิหเรยฺย อาตาปี สมฺปชาโน สติมา วิเนยฺย โลเก อภิชฺฌาโทมนสฺสํ.\csromanspacer{} เวทนาสุ วา ภิกฺขุ ฯเปฯ.\csromanspacer{} จิตฺเต วา ภิกฺขุ ฯเปฯ.\csromanspacer{} ธมฺเมสุ วา ภิกฺขุ ธมฺมานุปสฺสี วิหเรยฺย อาตาปี สมฺปชาโน สติมา วิเนยฺย โลเก อภิชฺฌาโทมนสฺสํ.\csromanspacer{} เอกายนฺวายํ มคฺโค สตฺตานํ วิสุทฺธิยา โสกปริเทวานํ สมติกฺกมาย ทุกฺขโทมนสฺสานํ อตฺถงฺคมาย ญายสฺส อธิคมาย นิพฺพานสฺส สจฺฉิกิริยาย, ยทิทํ จตฺตาโร สติปฏฺฐานาติ.}

\prose{อถ โข พฺรหฺมา สหมฺปติ ภควโต เจตสา เจโตปริวิตกฺกมญฺญาย เสยฺยถาปิ นาม พลวา ปุริโส สมิญฺชิตํ วา พาหํ ปสาเรยฺย, ปสาริตํ วา พาหํ สมิญฺเชยฺย.\csromanspacer{} เอวเมว โข พฺรหฺมโลเก อนฺตรหิโต ภควโต ปุรโต ปาตุรโหสิ.\csromanspacer{} อถ โข พฺรหฺมา สหมฺปติ เอกํสํ อุตฺตราสงฺคํ กริตฺวา เยน ภควา เตนญฺชลิํ ปณาเมตฺวา ภควนฺตํ เอตทโวจ\csromandash{} เอวเมตํ ภควา, เอวเมตํ สุคต, เอกายนฺวายํ ภนฺเต มคฺโค สตฺตานํ วิสุทฺธิยา โสกปริเทวานํ สมติกฺกมาย ทุกฺขโทมนสฺสานํ อตฺถงฺคมาย ญายสฺส อธิคมาย นิพฺพานสฺส สจฺฉิกิริยาย, ยทิทํ จตฺตาโร สติปฏฺฐานา.}

\prose{กตเม จตฺตาโร, กาเย วา ภนฺเต ภิกฺขุ กายานุปสฺสี วิหเรยฺย อาตาปี สมฺปชาโน สติมา วิเนยฺย โลเก อภิชฺฌาโทมนสฺสํ.\csromanspacer{} เวทนาสุ วา ภนฺเต ภิกฺขุ ฯเปฯ.\csromanspacer{} จิตฺเต วา ภนฺเต ภิกฺขุ ฯเปฯ.\csromanspacer{} ธมฺเมสุ วา ภนฺเต ภิกฺขุ ธมฺมานุปสฺสี วิหเรยฺย อาตาปี สมฺปชาโน สติมา วิเนยฺย โลเก อภิชฺฌาโทมนสฺสํ.\csromanspacer{} เอกายนฺวายํ ภนฺเต มคฺโค สตฺตานํ วิสุทฺธิยา โสกปริเทวานํ สมติกฺกมาย ทุกฺขโทมนสฺสานํ \csromanfolio{146}อตฺถงฺคมาย ญายสฺส อธิคมาย นิพฺพานสฺส สจฺฉิกิริยาย, ยทิทํ จตฺตาโร สติปฏฺฐานาติ.}

\prose{อิทมโวจ พฺรหฺมา สหมฺปติ, อิทํ วตฺวา อถาปรํ เอตทโวจ\csromandash{}}

\csromangathameasure{“เอกายนํ ชาติขยนฺตทสฺสี, มคฺคํ ปชานาติ หิตานุกมฺปี. \\ เอเตน มคฺเคน ตริํสุ ปุพฺเพ, ตริสฺสนฺติ เย จ ตรนฺติ โอฆนฺ”ติ. อฏฺฐมํ.}
\csromangathagroup{\csromangathastanza{“เอกายนํ ชาติขยนฺตทสฺสี, มคฺคํ ปชานาติ หิตานุกมฺปี. \\ เอเตน มคฺเคน ตริํสุ ปุพฺเพ, ตริสฺสนฺติ เย จ ตรนฺติ โอฆนฺ”ติ.\csromanspacer{} อฏฺฐมํ.}}

\csromanmatikaanchor{mtk.65}
\csromantocmarkref{subparagraph}{9. เสทกสุตฺต}{mtk.65}
\bookmark[dest={mtk.65},level=5]{9. เสทกสุตฺต}
\csromanheader{5}{9. เสทกสุตฺต}
\proseitem{385}{เอกํ สมยํ ภควา สุมฺเภสุ วิหรติ เสทกํ นาม สุมฺภานํ นิคโม.\csromanspacer{} ตตฺร โข ภควา ภิกฺขู อามนฺเตสิ “ภูตปุพฺพํ ภิกฺขเว จณฺฑาลวํสิโก จณฺฑาลวํสํ อุสฺสาเปตฺวา เมทกถาลิกํ อนฺเตวาสิํ อามนฺเตสิ ‘เอหิ ตฺวํ สมฺม เมทกถาลิเก จณฺฑาลวํสํ อภิรุหิตฺวา มม อุปริขนฺเธ ติฏฺฐาหี’ติ”. “เอวํ อาจริยา”ติ โข ภิกฺขเว เมทกถาลิกา อนฺเตวาสี จณฺฑาลวํสิกสฺส ปฏิสฺสุตฺวา จณฺฑาลวํสํ อภิรุหิตฺวา อาจริยสฺส อุปริขนฺเธ อฏฺฐาสิ.\csromanspacer{} อถ โข ภิกฺขเว จณฺฑาลวํสิโก เมทกถาลิกํ อนฺเตวาสิํ เอตทโวจ “ตฺวํ สมฺม เมทกถาลิเก มมํ รกฺข, อหํ ตํ รกฺขิสฺสามิ, เอวํ มยํ อญฺญมญฺญํ คุตฺตา อญฺญมญฺญํ รกฺขิตา สิปฺปานิ เจว ทสฺเสสฺสาม, ลาภญฺจ\footnote{ลาเภ จ (สี)} ลจฺฉาม, โสตฺถินา จ จณฺฑาลวํสา โอโรหิสฺสามา”ติ.\csromanspacer{} เอวํ วุตฺเต ภิกฺขเว เมทกถาลิกา อนฺเตวาสี จณฺฑาลวํสิกํ เอตทโวจ “น โข ปเนตํ อาจริย เอวํ ภวิสฺสติ, ตฺวํ อาจริย อตฺตานํ รกฺข, อหํ อตฺตานํ รกฺขิสฺสามิ, เอวํ มยํ อตฺตคุตฺตา อตฺตรกฺขิตา สิปฺปานิ เจว ทสฺเสสฺสาม, ลาภญฺจ ลจฺฉาม, โสตฺถินา จ จณฺฑาลวํสา โอโรหิสฺสามา”ติ.}

\prose{“โส ตตฺถ ญาโย”ติ ภควา เอตทโวจ, ยถา เมทกถาลิกา อนฺเตวาสี อาจริยํ อโวจ “อตฺตานํ ภิกฺขเว รกฺขิสฺสามี”ติ สติปฏฺฐานํ เสวิตพฺพํ, “ปรํ รกฺขิสฺสามี”ติ สติปฏฺฐานํ เสวิตพฺพํ.\csromanspacer{} อตฺตานํ ภิกฺขเว รกฺขนฺโต ปรํ รกฺขติ, ปรํ รกฺขนฺโต อตฺตานํ รกฺขติ.}

\csromanheader{2}{3. สติปฏฺฐานสํยุตฺต}
\prose{\csromanfolio{147}กถญฺจ ภิกฺขเว อตฺตานํ รกฺขนฺโต ปรํ รกฺขติ, อาเสวนาย ภาวนาย พหุลีกมฺเมน.\csromanspacer{} เอวํ โข ภิกฺขเว อตฺตานํ รกฺขนฺโต ปรํ รกฺขติ.\csromanspacer{} กถญฺจ ภิกฺขเว ปรํ รกฺขนฺโต อตฺตานํ รกฺขติ, ขนฺติยา อวิหิํสาย เมตฺตจิตฺตตาย อนุทยตาย.\csromanspacer{} เอวํ โข ภิกฺขเว ปรํ รกฺขนฺโต อตฺตานํ รกฺขติ. “อตฺตานํ ภิกฺขเว รกฺขิสฺสามี”ติ สติปฏฺฐานํ เสวิตพฺพํ, “ปรํ รกฺขิสฺสามี”ติ สติปฏฺฐานํ เสวิตพฺพํ, อตฺตานํ ภิกฺขเว รกฺขนฺโต ปรํ รกฺขติ, ปรํ รกฺขนฺโต อตฺตานํ รกฺขตีติ.\csromanspacer{}\csromanspacer{}นวมํ.}

\csromanmatikaanchor{mtk.66}
\csromantocmarkref{subparagraph}{10. ชนปทกลฺยาณีสุตฺต}{mtk.66}
\bookmark[dest={mtk.66},level=5]{10. ชนปทกลฺยาณีสุตฺต}
\csromanheader{5}{10. ชนปทกลฺยาณีสุตฺต}
\proseitem{386}{เอวํ เม สุตํ\csromandash{}เอกํ สมยํ ภควา สุมฺเภสุ วิหรติ เสทกํ นาม สุมฺภานํ นิคโม.\csromanspacer{} ตตฺร โข ภควา ภิกฺขู อามนฺเตสิ “ภิกฺขโว”ติ. “ภทนฺเต”ติ เต ภิกฺขู ภควโต ปจฺจสฺโสสุํ.\csromanspacer{} ภควา เอตทโวจ\csromandash{}}

\prose{เสยฺยถาปิ ภิกฺขเว “ชนปทกลฺยาณี ชนปทกลฺยาณี”ติ โข ภิกฺขเว มหาชนกาโย สนฺนิปเตยฺย, สา โข ปนสฺส ชนปทกลฺยาณี ปรมปาสาวินี นจฺเจ, ปรมปาสาวินี คีเต. “ชนปทกลฺยาณี นจฺจติ คายตี”ติ โข ภิกฺขเว ภิยฺโยโส มตฺตาย มหาชนกาโย สนฺนิปเตยฺย, อถ ปุริโส อาคจฺเฉยฺย ชีวิตุกาโม อมริตุกาโม สุขกาโม ทุกฺขปฺปฏิกูโล.\csromanspacer{} ตเมนํ เอวํ วเทยฺย “อยํ เต อมฺโภ ปุริส สมติตฺติโก เตลปตฺโต อนฺตเรน จ มหาสมชฺชํ อนฺตเรน จ ชนปทกลฺยาณิยา ปริหริตพฺโพ, ปุริโส จ เต อุกฺขิตฺตาสิโก ปิฏฺฐิโต ปิฏฺฐิโต อนุพนฺธิสฺสติ “ยตฺเถว นํ โถกมฺปิ ฉฑฺเฑสฺสติ, ตตฺเถว เต สิโร ปาเตสฺสตี”ติ.\csromanspacer{} ตํ กิํ มญฺญถ ภิกฺขเว, อปิ นุ โส ปุริโส อมุํ เตลปตฺตํ อมนสิกริตฺวา พหิทฺธา ปมาทํ อาหเรยฺยาติ.\csromanspacer{} โน เหตํ ภนฺเต.}

\prose{อุปมา โข มฺยายํ ภิกฺขเว กตา อตฺถสฺส วิญฺญาปนาย, อยํ เจเวตฺต อตฺโถ, “สมติตฺติโก เตลปตฺโต”ติ โข ภิกฺขเว กายคตาย เอตํ สติยา อธิวจนํ.\csromanspacer{} ตสฺมาติห ภิกฺขเว เอวํ สิกฺขิตพฺพํ “กายคตา สติ โน ภาวิตา ภวิสฺสติ พหุลีกตา ยานีกตา \csromanfolio{148}วตฺถุกตา อนุฏฺฐิตา ปริจิตา สุสมารทฺธา”ติ.\csromanspacer{} เอวํ หิ โข ภิกฺขเว สิกฺขิตพฺพนฺติ.\csromanspacer{}\csromanspacer{}ทสมํ.\csromanspacer{} นาลนฺทวคฺโค ทุติโย.}

\titlehead{\textbf{ตสฺสุทฺทานํ}}
\csromangathameasure{มหาปุริโส นาลนฺทํ, จุนฺโท เจลญฺจ พาหิโย. \\ อุตฺติโย อริโย พฺรหฺมา, เสทกํ ชนปเทน จาติ.}
\csromangathagroup{\csromangathastanza{มหาปุริโส นาลนฺทํ, จุนฺโท เจลญฺจ พาหิโย. \\ อุตฺติโย อริโย พฺรหฺมา, เสทกํ ชนปเทน จาติ.}}

\csromanmatikaanchor{mtk.67}
\csromantocmarknopage{paragraph}{3. สีลฏฺฐิติวคฺค}{mtk.67}
\bookmark[dest={mtk.67},level=4]{3. สีลฏฺฐิติวคฺค}
\csromanheader{4}{3. สีลฏฺฐิติวคฺค}
\csromanmatikaanchor{mtk.68}
\csromantocmarkref{subparagraph}{1. สีลสุตฺต}{mtk.68}
\bookmark[dest={mtk.68},level=5]{1. สีลสุตฺต}
\csromanheader{5}{1. สีลสุตฺต}
\proseitem{387}{เอวํ เม สุตํ\csromandash{}เอกํ สมยํ อายสฺมา จ อานนฺโท อายสฺมา จ ภทฺโท ปาฏลิปุตฺเต วิหรนฺติ กุกฺกุฏาราเม.\csromanspacer{} อถ โข อายสฺมา ภทฺโท สายนฺหสมยํ ปฏิสลฺลานา วุฏฺฐิโต เยนายสฺมา อานนฺโท เตนุปสงฺกมิ, อุปสงฺกมิตฺวา อายสฺมตา อานนฺเทน สทฺธิํ สมฺโมทิ, สมฺโมทนียํ กถํ สารณียํ วีติสาเรตฺวา เอกมนฺตํ นิสีทิ, เอกมนฺตํ นิสินฺโน โข อายสฺมา ภทฺโท อายสฺมนฺตํ อานนฺทํ เอตทโวจ “ยานิมานิ อาวุโส อานนฺท กุสลานิ สีลานิ วุตฺตานิ ภควตา, อิมานิ กุสลานิ สีลานิ กิมตฺถิยานิ วุตฺตานิ ภควตา”ติ.}

\prose{สาธุ สาธุ อาวุโส ภทฺท, ภทฺทโก โข เต อาวุโส ภทฺท อุมฺมงฺโค\footnote{อุมฺมคฺโค (สี, สฺยา, กํ)}, ภทฺทกํ ปฏิภานํ, กลฺยาณี ปริปุจฺฉา, เอวํ หิ ตฺวํ อาวุโส ภทฺท ปุจฺฉสิ “ยานิมานิ อาวุโส อานนฺท กุสลานิ สีลานิ วุตฺตานิ ภควตา, อิมานิ กุสลานิ สีลานิ กิมตฺถิยานิ วุตฺตานิ ภควตา”ติ.\csromanspacer{} เอวมาวุโสติ.\csromanspacer{} ยานิมานิ อาวุโส ภทฺท กุสลานิ สีลานิ วุตฺตานิ ภควตา, อิมานิ กุสลานิ สีลานิ ยาวเทว จตุนฺนํ สติปฏฺฐานานํ ภาวนาย วุตฺตานิ ภควตา.}

\prose{กตเมสํ จตุนฺนํ, อิธาวุโส ภิกฺขุ กาเย กายานุปสฺสี วิหรติ อาตาปี สมฺปชาโน สติมา วิเนยฺย โลเก อภิชฺฌาโทมนสฺสํ.}

\csromanheader{2}{3. สติปฏฺฐานสํยุตฺต}
\prose{\csromanfolio{149}เวทนาสุ ฯเปฯ.\csromanspacer{} จิตฺเต ฯเปฯ.\csromanspacer{} ธมฺเมสุ ธมฺมานุปสฺสี วิหรติ อาตาปี สมฺปชาโน สติมา วิเนยฺย โลเก อภิชฺฌาโทมนสฺสํ.\csromanspacer{} ยานิมานิ อาวุโส ภทฺท กุสลานิ สีลานิ วุตฺตานิ ภควตา, อิมานิ กุสลานิ สีลานิ ยาวเทว อิเมสํ จตุนฺนํ สติปฏฺฐานานํ ภาวนาย วุตฺตานิ ภควตาติ.\csromanspacer{}\csromanspacer{}ปฐมํ.}

\csromanmatikaanchor{mtk.69}
\csromantocmarkref{subparagraph}{2. จิรฏฺฐิติสุตฺต}{mtk.69}
\bookmark[dest={mtk.69},level=5]{2. จิรฏฺฐิติสุตฺต}
\csromanheader{5}{2. จิรฏฺฐิติสุตฺต}
\proseitem{388}{ตํเยว นิทานํ.\csromanspacer{} เอกมนฺตํ นิสินฺโน โข อายสฺมา ภทฺโท อายสฺมนฺตํ อานนฺทํ เอตทโวจ “โก นุ โข อาวุโส อานนฺท เหตุ โก ปจฺจโย, เยน ตถาคเต ปรินิพฺพุเต สทฺธมฺโม น จิรฏฺฐิติโก โหติ, โก ปนาวุโส อานนฺท เหตุ โก ปจฺจโย, เยน ตถาคเต ปรินิพฺพุเต สทฺธมฺโม จิรฏฺฐิติโก โหตี”ติ.}

\prose{สาธุ สาธุ อาวุโส ภทฺท, ภทฺทโก โข เต อาวุโส ภทฺท อุมฺมงฺโค, ภทฺทกํ ปฏิภานํ, กลฺยาณี ปริปุจฺฉา, เอวํ หิ ตฺวํ อาวุโส ภทฺท ปุจฺฉสิ “โก นุ โข อาวุโส อานนฺท เหตุ โก ปจฺจโย, เยน ตถาคเต ปรินิพฺพุเต สทฺธมฺโม น จิรฏฺฐิติโก โหติ, โก ปนาวุโส อานนฺท เหตุ โก ปจฺจโย, เยน ตถาคเต ปรินิพฺพุเต สทฺธมฺโม จิรฏฺฐิติโก โหตี”ติ.\csromanspacer{} เอวมาวุโสติ.\csromanspacer{} จตุนฺนํ โข อาวุโส สติปฏฺฐานานํ อภาวิตตฺตา อพหุลีกตตฺตา ตถาคเต ปรินิพฺพุเต สทฺธมฺโม น จิรฏฺฐิติโก โหติ.\csromanspacer{} จตุนฺนญฺจ โข อาวุโส สติปฏฺฐานานํ ภาวิตตฺตา พหุลีกตตฺตา ตถาคเต ปรินิพฺพุเต สทฺธมฺโม จิรฏฺฐิติโก โหติ.}

\prose{กตเมสํ จตุนฺนํ, อิธาวุโส ภิกฺขุ กาเย กายานุปสฺสี วิหรติ อาตาปี สมฺปชาโน สติมา วิเนยฺย โลเก อภิชฺฌาโทมนสฺสํ.\csromanspacer{} เวทนาสุ ฯเปฯ.\csromanspacer{} จิตฺเต ฯเปฯ.\csromanspacer{} ธมฺเมสุ ธมฺมานุปสฺสี วิหรติ อาตาปี สมฺปชาโน สติมา วิเนยฺย โลเก อภิชฺฌาโทมนสฺสํ.\csromanspacer{} อิเมสํ โข อาวุโส จตุนฺนํ สติปฏฺฐานานํ อภาวิตตฺตา อพหุลีกตตฺตา ตถาคเต ปรินิพฺพุเต สทฺธมฺโม น จิรฏฺฐิติโก โหติ.\csromanspacer{} อิเมสญฺจ โข อาวุโส จตุนฺนํ สติปฏฺฐานานํ ภาวิตตฺตา พหุลีกตตฺตา ตถาคเต ปรินิพฺพุเต สทฺธมฺโม จิรฏฺฐิติโก โหตีติ.\csromanspacer{}\csromanspacer{}ทุติยํ.}

\csromanheader{2}{3. ปริหานสุตฺต}
\csromanmatikaanchor{mtk.70}
\csromantocmarkref{subparagraph}{3-4. ปริหานสุตฺตาทีนิ}{mtk.70}
\bookmark[dest={mtk.70},level=5]{3-4. ปริหานสุตฺตาทีนิ}
\proseitem{389}{\csromanfolio{150}เอกํ สมยํ อายสฺมา จ อานนฺโท อายสฺมา จ ภทฺโท ปาฏลิปุตฺเต วิหรนฺติ กุกฺกุฏาราเม.\csromanspacer{} อถ โข อายสฺมา ภทฺโท สายนฺหสมยํ ปฏิสลฺลานา วุฏฺฐิโต เยนายสฺมา อานนฺโท เตนุปสงฺกมิ, อุปสงฺกมิตฺวา อายสฺมตา อานนฺเทน สทฺธิํ สมฺโมทิ, สมฺโมทนียํ กถํ สารณียํ วีติสาเรตฺวา เอกมนฺตํ นิสีทิ, เอกมนฺตํ นิสินฺโน โข อายสฺมา ภทฺโท อายสฺมนฺตํ อานนฺทํ เอตทโวจ “โก นุ โข อาวุโส อานนฺท เหตุ โก ปจฺจโย, เยน สทฺธมฺมปริหานํ โหติ, โก นุ โข อาวุโส อานนฺท เหตุ โก ปจฺจโย, เยน สทฺธมฺม-อปริหานํ โหตี”ติ.}

\prose{สาธุ สาธุ อาวุโส ภทฺท, ภทฺทโก โข เต อาวุโส ภทฺท อุมฺมงฺโค, ภทฺทกํ ปฏิภานํ, กลฺยาณี ปริปุจฺฉา, เอวํ หิ ตฺวํ อาวุโส ภทฺท ปุจฺฉสิ “โก นุ โข อาวุโส อานนฺท เหตุ โก ปจฺจโย, เยน สทฺธมฺมปริหานํ โหติ, โก ปนาวุโส อานนฺท เหตุ โก ปจฺจโย, เยน สทฺธมฺม-อปริหานํ โหตี”ติ.\csromanspacer{} เอวมาวุโสติ.\csromanspacer{} จตุนฺนํ โข อาวุโส สติปฏฺฐานานํ อภาวิตตฺตา อพหุลีกตตฺตา สทฺธมฺมปริหานํ โหติ.\csromanspacer{} จตุนฺนญฺจ โข อาวุโส สติปฏฺฐานานํ ภาวิตตฺตา พหุลีกตตฺตา สทฺธมฺม- อปริหานํ โหติ.}

\prose{กตเมสํ จตุนฺนํ, อิธาวุโส ภิกฺขุ กาเย กายานุปสฺสี วิหรติ อาตาปี สมฺปชาโน สติมา วิเนยฺย โลเก อภิชฺฌาโทมนสฺสํ.\csromanspacer{} เวทนาสุ ฯเปฯ.\csromanspacer{} จิตฺเต ฯเปฯ.\csromanspacer{} ธมฺเมสุ ธมฺมานุปสฺสี วิหรติ อาตาปี สมฺปชาโน สติมา วิเนยฺย โลเก อภิชฺฌาโทมนสฺสํ.\csromanspacer{} อิเมสํ โข อาวุโส จตุนฺนํ สติปฏฺฐานานํ อภาวิตตฺตา อพหุลีกตตฺตา สทฺธมฺมปริหานํ โหติ.\csromanspacer{} อิเมสญฺจ โข อาวุโส จตุนฺนํ สติปฏฺฐานานํ ภาวิตตฺตา พหุลีกตตฺตา สทฺธมฺม-อปริหานํ โหตีติ.\csromanspacer{}\csromanspacer{}ตติยํ.}

\csromanheader{2}{4. สุทฺธสุตฺต}
\proseitem{390}{สาวตฺถินิทานํ.\csromanspacer{} จตฺตาโรเม ภิกฺขเว สติปฏฺฐานา.\csromanspacer{} กตเม จตฺตาโร, อิธ ภิกฺขเว ภิกฺขุ กาเย กายานุปสฺสี วิหรติ อาตาปี สมฺปชาโน สติมา วิเนยฺย โลเก อภิชฺฌาโทมนสฺสํ.\csromanspacer{} เวทนาสุ ฯเปฯ.}

\csromanheader{2}{3. สติปฏฺฐานสํยุตฺต}
\prose{\csromanfolio{151}จิตฺเต ฯเปฯ.\csromanspacer{} ธมฺเมสุ ธมฺมานุปสฺสี วิหรติ อาตาปี สมฺปชาโน สติมา วิเนยฺย โลเก อภิชฺฌาโทมนสฺสํ.\csromanspacer{} อิเม โข ภิกฺขเว จตฺตาโร สติปฏฺฐานาติ.\csromanspacer{}\csromanspacer{}จตุตฺถํ.}

\csromanmatikaanchor{mtk.71}
\csromantocmarkref{subparagraph}{5. อญฺญตรพฺราหฺมณสุตฺต}{mtk.71}
\bookmark[dest={mtk.71},level=5]{5. อญฺญตรพฺราหฺมณสุตฺต}
\csromanheader{6}{5. อญฺญตรพฺราหฺมณสุตฺต}
\proseitem{391}{เอวํ เม สุตํ\csromandash{}เอกํ สมยํ ภควา สาวตฺถิยํ วิหรติ เชตวเน อนาถปิณฺฑิกสฺส อาราเม.\csromanspacer{} อถ โข อญฺญตโร พฺราหฺมโณ เยน ภควา เตนุปสงฺกมิ, อุปสงฺกมิตฺวา ภควตา สทฺธิํ สมฺโมทิ, สมฺโมทนียํ กถํ สารณียํ วีติสาเรตฺวา เอกมนฺตํ นิสีทิ, เอกมนฺตํ นิสินฺโน โข โส พฺราหฺมโณ ภควนฺตํ เอตทโวจ “โก นุ โข โภ โคตม เหตุ โก ปจฺจโย, เยน ตถาคเต ปรินิพฺพุเต สทฺธมฺโม น จิรฏฺฐิติโก โหติ, โก ปน โภ โคตม เหตุ โก ปจฺจโย, เยน ตถาคเต ปรินิพฺพุเต สทฺธมฺโม จิรฏฺฐิติโก โหตี”ติ.}

\prose{จตุนฺนํ โข พฺราหฺมณ สติปฏฺฐานานํ อภาวิตตฺตา อพหุลีกตตฺตา ตถาคเต ปรินิพฺพุเต สทฺธมฺโม น จิรฏฺฐิติโก โหติ, จตุนฺนญฺจ โข พฺราหฺมณ สติปฏฺฐานานํ ภาวิตตฺตา พหุลีกตตฺตา ตถาคเต ปรินิพฺพุเต สทฺธมฺโม จิรฏฺฐิติโก โหติ.}

\prose{กตเมสํ จตุนฺนํ, อิธ พฺราหฺมณ ภิกฺขุ กาเย กายานุปสฺสี วิหรติ อาตาปี สมฺปชาโน สติมา วิเนยฺย โลเก อภิชฺฌาโทมนสฺสํ.\csromanspacer{} เวทนาสุ ฯเปฯ.\csromanspacer{} จิตฺเต ฯเปฯ.\csromanspacer{} ธมฺเมสุ ธมฺมานุปสฺสี วิหรติ อาตาปี สมฺปชาโน สติมา วิเนยฺย โลเก อภิชฺฌาโทมนสฺสํ.\csromanspacer{} อิเมสํ โข พฺราหฺมณ จตุนฺนํ สติปฏฺฐานานํ อภาวิตตฺตา อพหุลีกตตฺตา ตถาคเต ปรินิพฺพุเต สทฺธมฺโม น จิรฏฺฐิติโก โหติ.\csromanspacer{} อิเมสญฺจ โข พฺราหฺมณ จตุนฺนํ สติปฏฺฐานานํ ภาวิตตฺตา พหุลีกตตฺตา ตถาคเต ปรินิพฺพุเต สทฺธมฺโม จิรฏฺฐิติโก โหตีติ.}

\prose{เอวํ วุตฺเต โส พฺราหฺมโณ ภควนฺตํ เอตทโวจ “อภิกฺกนฺตํ โภ โคตม ฯเปฯ อุปาสกํ มํ ภวํ โคตโม ธาเรตุ อชฺชตคฺเค ปาณุเปตํ สรณํ คตนฺ”ติ.\csromanspacer{}\csromanspacer{}ปญฺจมํ.}

\csromanmatikaanchor{mtk.72}
\csromantocmarkref{subparagraph}{6. ปเทสสุตฺต}{mtk.72}
\bookmark[dest={mtk.72},level=5]{6. ปเทสสุตฺต}
\csromanheader{6}{6. ปเทสสุตฺต}
\proseitem{392}{\csromanfolio{152}เอกํ สมยํ อายสฺมา จ สาริปุตฺโต อายสฺมา จ มหาโมคฺคลฺลาโน อายสฺมา จ อนุรุทฺโธ สาเกเต วิหรนฺติ กณฺฑกีวเน\footnote{กณฺฑกีวเน (สี, สฺยา, กํ, อิ)}.\csromanspacer{} อถ โข อายสฺมา จ สาริปุตฺโต อายสฺมา จ มหาโมคฺคลฺลาโน สายนฺหสมยํ ปฏิสลฺลานา วุฏฺฐิตา เยนายสฺมา อนุรุทฺโธ เตนุปสงฺกมิํสุ, อุปสงฺกมิตฺวา อายสฺมตา อนุรุทฺเธน สทฺธิํ สมฺโมทิํสุ, สมฺโมทนียํ กถํ สารณียํ วีติสาเรตฺวา เอกมนฺตํ นิสีทิํสุ, เอกมนฺตํ นิสินฺโน โข อายสฺมา สาริปุตฺโต อายสฺมนฺตํ อนุรุทฺธํ เอตทโวจ “เสโข เสโขติ\footnote{เสกฺโข เสกฺโขติ (สฺยา, กํ)} อาวุโส อนุรุทฺธ วุจฺจติ, กิตฺตาวตา นุ โข อาวุโส เสโข โหตี”ติ, จตุนฺนํ โข อาวุโส สติปฏฺฐานานํ ปเทสํ ภาวิตตฺตา เสโข โหติ.}

\prose{กตเมสํ จตุนฺนํ, อิธาวุโส ภิกฺขุ กาเย กายานุปสฺสี วิหรติ อาตาปี สมฺปชาโน สติมา วิเนยฺย โลเก อภิชฺฌาโทมนสฺสํ.\csromanspacer{} เวทนาสุ ฯเปฯ.\csromanspacer{} จิตฺเต ฯเปฯ.\csromanspacer{} ธมฺเมสุ ธมฺมานุปสฺสี วิหรติ อาตาปี สมฺปชาโน สติมา วิเนยฺย โลเก อภิชฺฌาโทมนสฺสํ.\csromanspacer{} อิเมสํ โข อาวุโส จตุนฺนํ สติปฏฺฐานานํ ปเทสํ ภาวิตตฺตา เสโข โหตีติ.\csromanspacer{}\csromanspacer{}ฉฏฺฐํ.}

\csromanmatikaanchor{mtk.73}
\csromantocmarkref{subparagraph}{7. สมตฺตสุตฺต}{mtk.73}
\bookmark[dest={mtk.73},level=5]{7. สมตฺตสุตฺต}
\csromanheader{6}{7. สมตฺตสุตฺต}
\proseitem{393}{ตํเยว นิทานํ.\csromanspacer{} เอกมนฺตํ นิสินฺโน โข อายสฺมา สาริปุตฺโต อายสฺมนฺตํ อนุรุทฺธํ เอตทโวจ “อเสโข อเสโขติ อาวุโส อนุรุทฺธ วุจฺจติ, กิตฺตาวตา นุ โข อาวุโส อเสโข โหตี”ติ.\csromanspacer{} จตุนฺนํ โข อาวุโส สติปฏฺฐานานํ สมตฺตํ ภาวิตตฺตา อเสโข โหติ.}

\prose{กตเมสํ จตุนฺนํ, อิธาวุโส ภิกฺขุ กาเย กายานุปสฺสี วิหรติ อาตาปี สมฺปชาโน สติมา วิเนยฺย โลเก อภิชฺฌาโทมนสฺสํ.\csromanspacer{} เวทนาสุ ฯเปฯ.\csromanspacer{} จิตฺเต ฯเปฯ.\csromanspacer{} ธมฺเมสุ ธมฺมานุปสฺสี วิหรติ อาตาปี สมฺปชาโน สติมา วิเนยฺย โลเก อภิชฺฌาโทมนสฺสํ.\csromanspacer{} อิเมสํ}

\csromanheader{2}{3. สติปฏฺฐานสํยุตฺต}
\prose{\csromanfolio{153}โข อาวุโส จตุนฺนํ สติปฏฺฐานานํ สมตฺตํ ภาวิตตฺตา อเสโข โหตีติ.\csromanspacer{}\csromanspacer{}สตฺตมํ.}

\csromanmatikaanchor{mtk.74}
\csromantocmarkref{subparagraph}{8. โลกสุตฺต}{mtk.74}
\bookmark[dest={mtk.74},level=5]{8. โลกสุตฺต}
\csromanheader{6}{8. โลกสุตฺต}
\proseitem{394}{ตํเยว นิทานํ.\csromanspacer{} เอกมนฺตํ นิสินฺโน โข อายสฺมา สาริปุตฺโต อายสฺมนฺตํ อนุรุทฺธํ เอตทโวจ “กตเมสํ อาวุโส อนุรุทฺธ ธมฺมานํ ภาวิตตฺตา พหุลีกตตฺตา มหาภิญฺญตํ\footnote{มหาภิญฺญาตํ (อิ)} ปตฺโต”ติ.\csromanspacer{} จตุนฺนํ อาวุโส สติปฏฺฐานานํ ภาวิตตฺตา พหุลีกตตฺตา มหาภิญฺญตํ ปตฺโต.}

\prose{กตเมสํ จตุนฺนํ, อิธาหํ อาวุโส กาเย กายานุปสฺสี วิหรามิ อาตาปี สมฺปชาโน สติมา วิเนยฺย โลเก อภิชฺฌาโทมนสฺสํ.\csromanspacer{} เวทนาสุ ฯเปฯ.\csromanspacer{} จิตฺเต ฯเปฯ.\csromanspacer{} ธมฺเมสุ ธมฺมานุปสฺสี วิหรามิ อาตาปี สมฺปชาโน สติมา วิเนยฺย โลเก อภิชฺฌาโทมนสฺสํ.\csromanspacer{} อิเมสํ ขฺวาหํ อาวุโส จตุนฺนํ สติปฏฺฐานานํ ภาวิตตฺตา พหุลีกตตฺตา มหาภิญฺญตํ ปตฺโต.\csromanspacer{} อิเมสญฺจ ปนาหํ อาวุโส จตุนฺนํ สติปฏฺฐานานํ ภาวิตตฺตา พหุลีกตตฺตา สหสฺสํ โลกํ อภิชานามีติ.\csromanspacer{}\csromanspacer{}อฏฺฐมํ.}

\csromanmatikaanchor{mtk.75}
\csromantocmarkref{subparagraph}{9. สิริวฑฺฒสุตฺต}{mtk.75}
\bookmark[dest={mtk.75},level=5]{9. สิริวฑฺฒสุตฺต}
\csromanheader{6}{9. สิริวฑฺฒสุตฺต}
\proseitem{395}{เอกํ สมยํ อายสฺมา อานนฺโท ราชคเห วิหรติ เวฬุวเน กลนฺทกนิวาเป.\csromanspacer{} เตน โข ปน สมเยน สิริวฑฺโฒ\footnote{สิริวฑฺโฒ (ก)} คหปติ อาพาธิโก โหติ ทุกฺขิโต พาฬฺหคิลาโน.\csromanspacer{} อถ โข สิริวฑฺโฒ คหปติ อญฺญตรํ ปุริสํ อามนฺเตสิ “เอหิ ตฺวํ อมฺโภ ปุริส เยนายสฺมา อานนฺโท เตนุปสงฺกมิ, อุปสงฺกมิตฺวา มม วจเนน อายสฺมโต อานนฺทสฺส ปาเท สิรสา วนฺท ‘สิริวฑฺโฒ ภนฺเต คหปติ อาพาธิโก ทุกฺขิโต พาฬฺหคิลาโน, โส อายสฺมโต อานนฺทสฺส ปาเท สิรสา วนฺทตี’ติ, เอวญฺจ วเทหิ ‘สาธุ กิร ภนฺเต อายสฺมา อานนฺโท เยน สิริวฑฺฒสฺส คหปติสฺส นิเวสนํ เตนุปสงฺกมตุ อนุกมฺปํ อุปาทายา’ติ”. “เอวํ ภนฺเต”ติ โข โส ปุริโส สิริวฑฺฒสฺส คหปติสฺส ปฏิสฺสุตฺวา เยนายสฺมา อานนฺโท เตนุปสงฺกมิ, อุปสงฺกมิตฺวา อายสฺมนฺตํ อานนฺทํ อภิวาเทตฺวา เอกมนฺตํ นิสีทิ, เอกมนฺตํ นิสินฺโน \csromanfolio{154}โข โส ปุริโส อายสฺมนฺตํ อานนฺทํ เอตทโวจ “สิริวฑฺโฒ ภนฺเต คหปติ อาพาธิโก ทุกฺขิโต พาฬฺหคิลาโน, โส อายสฺมโต อานนฺทสฺส ปาเท สิรสา วนฺทติ, เอวญฺจ วเทติ ‘สาธุ กิร ภนฺเต อายสฺมา อานนฺโท เยน สิริวฑฺฒสฺส คหปติสฺส นิเวสนํ เตนุปสงฺกมตุ อนุกมฺปํ อุปาทายา’ติ”.\csromanspacer{} อธิวาเสสิ โข อายสฺมา อานนฺโท ตุณฺหีภาเวน.}

\prose{อถ โข อายสฺมา อานนฺโท ปุพฺพณฺหสมยํ นิวาเสตฺวา ปตฺตจีวรมาทาย เยน สิริวฑฺฒสฺส คหปติสฺส นิเวสนํ เตนุปสงฺกมิ, อุปสงฺกมิตฺวา ปญฺญตฺเต อาสเน นิสีทิ, นิสชฺช โข อายสฺมา อานนฺโท สิริวฑฺฒํ คหปติํ เอตทโวจ “กจฺจิ เต คหปติ ขมนียํ, กจฺจิ ยาปนียํ, กจฺจิ ทุกฺขา เวทนา ปฏิกฺกมนฺติ โน อภิกฺกมนฺติ, ปฏิกฺกโมสานํ ปญฺญายติ โน อภิกฺกโม”ติ.\csromanspacer{} น เม ภนฺเต ขมนียํ, น ยาปนียํ, พาฬฺหา เม ทุกฺขา เวทนา อภิกฺกมนฺติ โน ปฏิกฺกมนฺติ, อภิกฺกโมสานํ ปญฺญายติ โน ปฏิกฺกโมติ.}

\prose{ตสฺมาติห เต คหปติ เอวํ สิกฺขิตพฺพํ “กาเย กายานุปสฺสี วิหริสฺสามิ อาตาปี สมฺปชาโน สติมา วิเนยฺย โลเก อภิชฺฌาโทมนสฺสํ.\csromanspacer{} เวทนาสุ ฯเปฯ.\csromanspacer{} จิตฺเต ฯเปฯ.\csromanspacer{} ธมฺเมสุ ธมฺมานุปสฺสี วิหริสฺสามิ อาตาปี สมฺปชาโน สติมา วิเนยฺย โลเก อภิชฺฌาโทมนสฺสนฺ”ติ.\csromanspacer{} เอวํ หิ เต คหปติ สิกฺขิตพฺพนฺติ.}

\prose{เยเม ภนฺเต ภควตา จตฺตาโร สติปฏฺฐานา เทสิตา, สํวิชฺชนฺติ เต ธมฺมา\footnote{สํวิชฺชนฺเต รตนธมฺมา (สี)} มยิ, อหญฺจ เตสุ ธมฺเมสุ สนฺทิสฺสามิ.\csromanspacer{} อหํ หิ ภนฺเต กาเย กายานุปสฺสี วิหรามิ อาตาปี สมฺปชาโน สติมา วิเนยฺย โลเก อภิชฺฌาโทมนสฺสํ.\csromanspacer{} เวทนาสุ ฯเปฯ.\csromanspacer{} จิตฺเต ฯเปฯ.\csromanspacer{} ธมฺเมสุ ธมฺมานุปสฺสี วิหรามิ อาตาปี สมฺปชาโน สติมา วิเนยฺย โลเก อภิชฺฌาโทมนสฺสํ.\csromanspacer{} ยานิ จิมานิ ภนฺเต ภควตา ปญฺโจรมฺภาคิยานิ สํโยชนานิ เทสิตานิ, นาหํ ภนฺเต เตสํ กิญฺจิ อตฺตนิ อปฺปหีนํ สมนุปสฺสามีติ.\csromanspacer{} ลาภา เต คหปติ, สุลทฺธํ เต คหปติ, อนาคามิผลํ ตยา คหปติ พฺยากตนฺติ.\csromanspacer{}\csromanspacer{}นวมํ.}

\csromanheader{2}{3. สติปฏฺฐานสํยุตฺต}
\csromanmatikaanchor{mtk.76}
\csromantocmarkref{subparagraph}{10. มานทินฺนสุตฺต}{mtk.76}
\bookmark[dest={mtk.76},level=5]{10. มานทินฺนสุตฺต}
\csromanheader{6}{10. มานทินฺนสุตฺต}
\proseitem{396}{\csromanfolio{155}ตํเยว นิทานํ.\csromanspacer{} เตน โข ปน สมเยน มานทินฺโน คหปติ อาพาธิโก โหติ ทุกฺขิโต พาฬฺหคิลาโน.\csromanspacer{} อถ โข มานทินฺโน คหปติ อญฺญตรํ ปุริสํ อามนฺเตสิ “เอหิ ตฺวํ อมฺโภ ปุริส ฯเปฯ” น เม ภนฺเต ขมนียํ, น ยาปนียํ, พาฬฺหา เม ทุกฺขา เวทนา อภิกฺกมนฺติ โน ปฏิกฺกมนฺติ, อภิกฺกโมสานํ ปญฺญายติ โน ปฏิกฺกโมติ.\csromanspacer{} เอวรูปาย จาหํ ภนฺเต ทุกฺขาย เวทนาย ผุฏฺโฐ สมาโน กาเย กายานุปสฺสี วิหรามิ อาตาปี สมฺปชาโน สติมา วิเนยฺย โลเก อภิชฺฌาโทมนสฺสํ.\csromanspacer{} เวทนาสุ ฯเปฯ.\csromanspacer{} จิตฺเต ฯเปฯ.\csromanspacer{} ธมฺเมสุ ธมฺมานุปสฺสี วิหรามิ อาตาปี สมฺปชาโน สติมา วิเนยฺย โลเก อภิชฺฌาโทมนสฺสํ.\csromanspacer{} ยานิ จิมานิ ภนฺเต ภควตา ปญฺโจรมฺภาคิยานิ สํโยชนานิ เทสิตานิ, นาหํ ภนฺเต เตสํ กิญฺจิ อตฺตนิ อปฺปหีนํ สมนุปสฺสามีติ.\csromanspacer{} ลาภา เต คหปติ, สุลทฺธํ เต คหปติ, อนาคามิผลํ ตยา คหปติ พฺยากตนฺติ.\csromanspacer{}\csromanspacer{}ทสมํ.\csromanspacer{} สีลฏฺฐิติวคฺโค ตติโย.}

\titlehead{\textbf{ตสฺสุทฺทานํ}}
\csromangathameasure{สีลํ ฐิติ ปริหานํ, สุทฺธํ พฺราหฺมณปเทสํ. \\ สมตฺตํ โลโก สิริวฑฺโฒ, มานทินฺเนน เต ทสาติ.}
\csromangathagroup{\csromangathastanza{สีลํ ฐิติ ปริหานํ, สุทฺธํ พฺราหฺมณปเทสํ. \\ สมตฺตํ โลโก สิริวฑฺโฒ, มานทินฺเนน เต ทสาติ.}}

\csromanmatikaanchor{mtk.77}
\csromantocmarknopage{subparagraph}{4. อนนุสฺสุตวคฺค}{mtk.77}
\bookmark[dest={mtk.77},level=5]{4. อนนุสฺสุตวคฺค}
\csromanheader{6}{4. อนนุสฺสุตวคฺค}
\csromanmatikaanchor{mtk.78}
\csromantocmarkref{subparagraph}{1. อนนุสฺสุตสุตฺต}{mtk.78}
\bookmark[dest={mtk.78},level=5]{1. อนนุสฺสุตสุตฺต}
\csromanheader{6}{1. อนนุสฺสุตสุตฺต}
\proseitem{397}{สาวตฺถินิทานํ.\csromanspacer{} อยํ กาเย กายานุปสฺสนาติ เม ภิกฺขเว ปุพฺเพ อนนุสฺสุเตสุ ธมฺเมสุ จกฺขุํ อุทปาทิ, ญาณํ อุทปาทิ, ปญฺญา อุทปาทิ, วิชฺชา อุทปาทิ, อาโลโก อุทปาทิ.\csromanspacer{} สา โข ปนายํ กาเย กายานุปสฺสนา ภาเวตพฺพาติ เม ภิกฺขเว ฯเปฯ ภาวิตาติ เม ภิกฺขเว ปุพฺเพ อนนุสฺสุเตสุ ธมฺเมสุ จกฺขุํ อุทปาทิ, ญาณํ อุทปาทิ, ปญฺญา อุทปาทิ, วิชฺชา อุทปาทิ, อาโลโก อุทปาทิ.}

\prose{\csromanfolio{156}อยํ เวทนาสุ เวทนานุปสฺสนาติ เม ภิกฺขเว ปุพฺเพ อนนุสฺสุเตสุ ธมฺเมสุ จกฺขุํ อุทปาทิ, ญาณํ อุทปาทิ, ปญฺญา อุทปาทิ, วิชฺชา อุทปาทิ, อาโลโก อุทปาทิ.\csromanspacer{} สา โข ปนายํ เวทนาสุ เวทนานุปสฺสนา ภาเวตพฺพาติ เม ภิกฺขเว ฯเปฯ ภาวิตาติ เม ภิกฺขเว ปุพฺเพ อนนุสฺสุเตสุ ธมฺเมสุ จกฺขุํ อุทปาทิ, ญาณํ อุทปาทิ, ปญฺญา อุทปาทิ, วิชฺชา อุทปาทิ, อาโลโก อุทปาทิ.}

\prose{อยํ จิตฺเต จิตฺตานุปสฺสนาติ เม ภิกฺขเว ปุพฺเพ อนนุสฺสุเตสุ ธมฺเมสุ จกฺขุํ อุทปาทิ, ญาณํ อุทปาทิ, ปญฺญา อุทปาทิ, วิชฺชา อุทปาทิ, อาโลโก อุทปาทิ.\csromanspacer{} สา โข ปนายํ จิตฺเต จิตฺตานุปสฺสนา ภาเวตพฺพาติ เม ภิกฺขเว ฯเปฯ ภาวิตาติ เม ภิกฺขเว ปุพฺเพ อนนุสฺสุเตสุ ธมฺเมสุ จกฺขุํ อุทปาทิ, ญาณํ อุทปาทิ, ปญฺญา อุทปาทิ, วิชฺชา อุทปาทิ, อาโลโก อุทปาทิ.}

\prose{อยํ ธมฺเมสุ ธมฺมานุปสฺสนาติ เม ภิกฺขเว ปุพฺเพ อนนุสฺสุเตสุ ธมฺเมสุ จกฺขุํ อุทปาทิ, ญาณํ อุทปาทิ, ปญฺญา อุทปาทิ, วิชฺชา อุทปาทิ, อาโลโก อุทปาทิ.\csromanspacer{} สา โข ปนายํ ธมฺเมสุ ธมฺมานุปสฺสนา ภาเวตพฺพาติ เม ภิกฺขเว ฯเปฯ ภาวิตาติ เม ภิกฺขเว ปุพฺเพ อนนุสฺสุเตสุ ธมฺเมสุ จกฺขุํ อุทปาทิ, ญาณํ อุทปาทิ, ปญฺญา อุทปาทิ, วิชฺชา อุทปาทิ, อาโลโก อุทปาทีติ.\csromanspacer{}\csromanspacer{}ปฐมํ.}

\csromanmatikaanchor{mtk.79}
\csromantocmarkref{subparagraph}{2. วิราคสุตฺต}{mtk.79}
\bookmark[dest={mtk.79},level=5]{2. วิราคสุตฺต}
\csromanheader{6}{2. วิราคสุตฺต}
\proseitem{398}{จตฺตาโรเม ภิกฺขเว สติปฏฺฐานา ภาวิตา พหุลีกตา เอกนฺตนิพฺพิทาย วิราคาย นิโรธาย อุปสมาย อภิญฺญาย สมฺโพธาย นิพฺพานาย สํวตฺตนฺติ.}

\prose{กตเม จตฺตาโร, อิธ ภิกฺขเว ภิกฺขุ กาเย กายานุปสฺสี วิหรติ อาตาปี สมฺปชาโน สติมา วิเนยฺย โลเก อภิชฺฌาโทมนสฺสํ.\csromanspacer{} เวทนาสุ ฯเปฯ.\csromanspacer{} จิตฺเต ฯเปฯ.\csromanspacer{} ธมฺเมสุ ธมฺมานุปสฺสี วิหรติ อาตาปี สมฺปชาโน สติมา วิเนยฺย โลเก อภิชฺฌาโทมนสฺสํ.\csromanspacer{} อิเม โข ภิกฺขเว จตฺตาโร สติปฏฺฐานา ภาวิตา พหุลีกตา เอกนฺตนิพฺพิทาย วิราคาย นิโรธาย อุปสมาย อภิญฺญาย สมฺโพธาย นิพฺพานาย สํวตฺตนฺตีติ.\csromanspacer{}\csromanspacer{}ทุติยํ.}

\csromanheader{2}{3. สติปฏฺฐานสํยุตฺต}
\csromanheader{2}{3. วิรทฺธสุตฺต}
\csromanmatikaanchor{mtk.80}
\csromanmatikaanchor{mtk.81}
\csromantocmarkref{subparagraph}{3-4. วิรทฺธสุตฺตาทีนิ}{mtk.80}
\bookmark[dest={mtk.80},level=5]{3-4. วิรทฺธสุตฺตาทีนิ}
\csromantocmarkref{subparagraph}{5-9. สติสุตฺตาทีนิ}{mtk.81}
\bookmark[dest={mtk.81},level=5]{5-9. สติสุตฺตาทีนิ}
\proseitem{399}{\csromanfolio{157}เยสํ เกสญฺจิ ภิกฺขเว จตฺตาโร สติปฏฺฐานา วิรทฺธา, วิรทฺโธ เตสํ อริโย มคฺโค\footnote{อริโย อฏฺฐงฺคิโก มคฺโค (ก) อิมสฺมิํ เยว สุตฺเต ทิสฺสติ อฏฺฐงฺคิโกติปทํ, น ปนาญฺญตฺถ อิทฺธิปาท อนุรุทฺธาทีสุ.} สมฺมา ทุกฺขกฺขยคามี.\csromanspacer{} เยสํ เกสญฺจิ ภิกฺขเว จตฺตาโร สติปฏฺฐานา อารทฺธา, อารทฺโธ เตสํ อริโย มคฺโค สมฺมา ทุกฺขกฺขยคามี.}

\prose{กตเม จตฺตาโร, อิธ ภิกฺขเว ภิกฺขุ กาเย กายานุปสฺสี วิหรติ อาตาปี สมฺปชาโน สติมา วิเนยฺย โลเก อภิชฺฌาโทมนสฺสํ.\csromanspacer{} เวทนาสุ ฯเปฯ.\csromanspacer{} จิตฺเต ฯเปฯ.\csromanspacer{} ธมฺเมสุ ธมฺมานุปสฺสี วิหรติ อาตาปี สมฺปชาโน สติมา วิเนยฺย โลเก อภิชฺฌาโทมนสฺสํ.\csromanspacer{} เยสํ เกสญฺจิ ภิกฺขเว อิเม จตฺตาโร สติปฏฺฐานา วิรทฺธา, วิรทฺโธ เตสํ อริโย มคฺโค สมฺมา ทุกฺขกฺขยคามี.\csromanspacer{} เยสํ เกสญฺจิ ภิกฺขเว อิเม จตฺตาโร สติปฏฺฐานา อารทฺธา, อารทฺโธ เตสํ อริโย มคฺโค สมฺมา ทุกฺขกฺขยคามีติ.\csromanspacer{}\csromanspacer{}ตติยํ.}

\csromanheader{2}{4. ภาวิตสุตฺต}
\proseitem{400}{จตฺตาโรเม ภิกฺขเว สติปฏฺฐานา ภาวิตา พหุลีกตา อปารา ปารํ คมนาย สํวตฺตนฺติ.}

\prose{กตเม จตฺตาโร, อิธ ภิกฺขเว ภิกฺขุ กาเย กายานุปสฺสี วิหรติ อาตาปี สมฺปชาโน สติมา วิเนยฺย โลเก อภิชฺฌาโทมนสฺสํ.\csromanspacer{} เวทนาสุ ฯเปฯ.\csromanspacer{} จิตฺเต ฯเปฯ.\csromanspacer{} ธมฺเมสุ ธมฺมานุปสฺสี วิหรติ อาตาปี สมฺปชาโน สติมา วิเนยฺย โลเก อภิชฺฌาโทมนสฺสํ.\csromanspacer{} อิเม โข ภิกฺขเว จตฺตาโร สติปฏฺฐานา ภาวิตา พหุลีกตา อปารา ปารํ คมนาย สํวตฺตนฺตีติ.\csromanspacer{}\csromanspacer{}จตุตฺถํ.}

\csromanheader{2}{5. สติสุตฺต}
\proseitem{401}{สาวตฺถินิทานํ.\csromanspacer{} สโต ภิกฺขเว ภิกฺขุ วิหเรยฺย สมฺปชาโน, อยํ โว อมฺหากํ อนุสาสนี.}

\prose{กถญฺจ ภิกฺขเว ภิกฺขุ สโต โหติ, อิธ ภิกฺขเว ภิกฺขุ กาเย กายานุปสฺสี วิหรติ อาตาปี สมฺปชาโน สติมา วิเนยฺย โลเก \csromanfolio{158}อภิชฺฌาโทมนสฺสํ.\csromanspacer{} เวทนาสุ ฯเปฯ.\csromanspacer{} จิตฺเต ฯเปฯ.\csromanspacer{} ธมฺเมสุ ธมฺมานุปสฺสี วิหรติ อาตาปี สมฺปชาโน สติมา วิเนยฺย โลเก อภิชฺฌาโทมนสฺสํ.\csromanspacer{} เอวํ โข ภิกฺขเว ภิกฺขุ สโต โหติ.}

\prose{กถญฺจ ภิกฺขเว ภิกฺขุ สมฺปชาโน โหติ, อิธ ภิกฺขเว ภิกฺขุโน วิทิตา เวทนา อุปฺปชฺชนฺติ, วิทิตา อุปฏฺฐหนฺติ, วิทิตา อพฺภตฺถํ คจฺฉนฺติ.\csromanspacer{} วิทิตา วิตกฺกา อุปฺปชฺชนฺติ, วิทิตา อุปฏฺฐหนฺติ, วิทิตา อพฺภตฺถํ คจฺฉนฺติ.\csromanspacer{} วิทิตา สญฺญา อุปฺปชฺชนฺติ, วิทิตา อุปฏฺฐหนฺติ, วิทิตา อพฺภตฺถํ คจฺฉนฺติ.\csromanspacer{} เอวํ โข ภิกฺขเว ภิกฺขุ สมฺปชาโน โหติ, สโต ภิกฺขเว ภิกฺขุ วิหเรยฺย สมฺปชาโน, อยํ โว อมฺหากํ อนุสาสนีติ.\csromanspacer{}\csromanspacer{}ปญฺจมํ.}

\csromanheader{2}{6. อญฺญาสุตฺต}
\proseitem{402}{จตฺตาโรเม ภิกฺขเว สติปฏฺฐานา.\csromanspacer{} กตเม จตฺตาโร, อิธ ภิกฺขเว ภิกฺขุ กาเย กายานุปสฺสี วิหรติ อาตาปี สมฺปชาโน สติมา วิเนยฺย โลเก อภิชฺฌาโทมนสฺสํ.\csromanspacer{} เวทนาสุ ฯเปฯ.\csromanspacer{} จิตฺเต ฯเปฯ.\csromanspacer{} ธมฺเมสุ ธมฺมานุปสฺสี วิหรติ อาตาปี สมฺปชาโน สติมา วิเนยฺย โลเก อภิชฺฌาโทมนสฺสํ.\csromanspacer{} อิเม โข ภิกฺขเว จตฺตาโร สติปฏฺฐานา.\csromanspacer{} อิเมสํ โข ภิกฺขเว จตุนฺนํ สติปฏฺฐานานํ ภาวิตตฺตา พหุลีกตตฺตา ทฺวินฺนํ ผลานํ อญฺญตรํ ผลํ ปาฏิกงฺขํ, ทิฏฺเฐว ธมฺเม อญฺญา, สติ วา อุปาทิเสเส อนาคามิตาติ.\csromanspacer{}\csromanspacer{}ฉฏฺฐํ.}

\csromanheader{2}{7. ฉนฺทสุตฺต}
\proseitem{403}{จตฺตาโรเม ภิกฺขเว สติปฏฺฐานา.\csromanspacer{} กตเม จตฺตาโร, อิธ ภิกฺขเว ภิกฺขุ กาเย กายานุปสฺสี วิหรติ อาตาปี สมฺปชาโน สติมา วิเนยฺย โลเก อภิชฺฌาโทมนสฺสํ, ตสฺส กาเย กายานุปสฺสิโน วิหรโต โย กายสฺมิํ ฉนฺโท, โส ปหียติ, ฉนฺทสฺส ปหานา อมตํ สจฺฉิกตํ โหติ.}

\prose{เวทนาสุ เวทนานุปสฺสี วิหรติ อาตาปี สมฺปชาโน สติมา วิเนยฺย โลเก อภิชฺฌาโทมนสฺสํ, ตสฺส เวทนาสุ เวทนานุปสฺสิโน วิหรโต โย เวทนาสุ ฉนฺโท, โส ปหียติ, ฉนฺทสฺส ปหานา อมตํ สจฺฉิกตํ โหติ.}

\csromanheader{2}{3. สติปฏฺฐานสํยุตฺต}
\prose{\csromanfolio{159}จิตฺเต จิตฺตานุปสฺสี วิหรติ อาตาปี สมฺปชาโน สติมา วิเนยฺย โลเก อภิชฺฌาโทมนสฺสํ, ตสฺส จิตฺเต จิตฺตานุปสฺสิโน วิหรโต โย จิตฺตมฺหิ ฉนฺโท, โส ปหียติ, ฉนฺทสฺส ปหานา อมตํ สจฺฉิกตํ โหติ.}

\prose{ธมฺเมสุ ธมฺมานุปสฺสี วิหรติ อาตาปี สมฺปชาโน สติมา วิเนยฺย โลเก อภิชฺฌาโทมนสฺสํ, ตสฺส ธมฺเมสุ ธมฺมานุปสฺสิโน วิหรโต โย ธมฺเมสุ ฉนฺโท, โส ปหียติ, ฉนฺทสฺส ปหานา อมตํ สจฺฉิกตํ โหตีติ.\csromanspacer{}\csromanspacer{}สตฺตมํ.}

\csromanheader{2}{8. ปริญฺญาตสุตฺต}
\proseitem{404}{จตฺตาโรเม ภิกฺขเว สติปฏฺฐานา.\csromanspacer{} กตเม จตฺตาโร, อิธ ภิกฺขเว ภิกฺขุ กาเย กายานุปสฺสี วิหรติ อาตาปี สมฺปชาโน สติมา วิเนยฺย โลเก อภิชฺฌาโทมนสฺสํ, ตสฺส กาเย กายานุปสฺสิโน วิหรโต กาโย ปริญฺญาโต โหติ, กายสฺส ปริญฺญาตตฺตา อมตํ สจฺฉิกตํ โหติ.}

\prose{เวทนาสุ เวทนานุปสฺสี วิหรติ อาตาปี สมฺปชาโน สติมา วิเนยฺย โลเก อภิชฺฌาโทมนสฺสํ, ตสฺส เวทนาสุ เวทนานุปสฺสิโน วิหรโต เวทนา ปริญฺญาตา โหนฺติ, เวทนานํ ปริญฺญาตตฺตา อมตํ สจฺฉิกตํ โหติ.}

\prose{จิตฺเต จิตฺตานุปสฺสี วิหรติ อาตาปี สมฺปชาโน สติมา วิเนยฺย โลเก อภิชฺฌาโทมนสฺสํ, ตสฺส จิตฺเต จิตฺตานุปสฺสิโน วิหรโต จิตฺตํ ปริญฺญาตํ โหติ, จิตฺตสฺส ปริญฺญาตตฺตา อมตํ สจฺฉิกตํ โหติ.}

\prose{ธมฺเมสุ ธมฺมานุปสฺสี วิหรติ อาตาปี สมฺปชาโน สติมา วิเนยฺย โลเก อภิชฺฌาโทมนสฺสํ, ตสฺส ธมฺเมสุ ธมฺมานุปสฺสีโน วิหรโต ธมฺมา ปริญฺญาตา โหนฺติ, ธมฺมานํ ปริญฺญาตตฺตา อมตํ สจฺฉิกตํ โหตีติ.\csromanspacer{}\csromanspacer{}อฏฺฐมํ.}

\csromanheader{2}{9. ภาวนาสุตฺต}
\proseitem{405}{\csromanfolio{160}จตุนฺนํ ภิกฺขเว สติปฏฺฐานานํ ภาวนํ เทเสสฺสามิ, ตํ สุณาถ.\csromanspacer{} กตมา ภิกฺขเว จตุนฺนํ สติปฏฺฐานานํ ภาวนา, อิธ ภิกฺขเว ภิกฺขุ กาเย กายานุปสฺสี วิหรติ อาตาปี สมฺปชาโน สติมา วิเนยฺย โลเก อภิชฺฌาโทมนสฺสํ.\csromanspacer{} เวทนาสุ ฯเปฯ.\csromanspacer{} จิตฺเต ฯเปฯ.\csromanspacer{} ธมฺเมสุ ธมฺมานุปสฺสี วิหรติ อาตาปี สมฺปชาโน สติมา วิเนยฺย โลเก อภิชฺฌาโทมนสฺสํ.\csromanspacer{} อยํ โข ภิกฺขเว จตุนฺนํ สติปฏฺฐานานํ ภาวนาติ.\csromanspacer{}\csromanspacer{}นวมํ.}

\csromanmatikaanchor{mtk.82}
\csromantocmarkref{subparagraph}{10. วิภงฺคสุตฺต}{mtk.82}
\bookmark[dest={mtk.82},level=5]{10. วิภงฺคสุตฺต}
\csromanheader{6}{10. วิภงฺคสุตฺต}
\proseitem{406}{สติปฏฺฐานญฺจ โว ภิกฺขเว เทเสสฺสามิ สติปฏฺฐานภาวนญฺจ สติปฏฺฐานภาวนาคามินิญฺจ ปฏิปทํ, ตํ สุณาถ.\csromanspacer{} กตมญฺจ ภิกฺขเว สติปฏฺฐานํ, อิธ ภิกฺขเว ภิกฺขุ กาเย กายานุปสฺสี วิหรติ อาตาปี สมฺปชาโน สติมา วิเนยฺย โลเก อภิชฺฌาโทมนสฺสํ.\csromanspacer{} เวทนาสุ เวทนานุปสฺสี วิหรติ ฯเปฯ.\csromanspacer{} จิตฺเต จิตฺตานุปสฺสี วิหรติ ฯเปฯ.\csromanspacer{} ธมฺเมสุ ธมฺมานุปสฺสี วิหรติ อาตาปี สมฺปชาโน สติมา วิเนยฺย โลเก อภิชฺฌาโทมนสฺสํ.\csromanspacer{} อิทํ วุจฺจติ ภิกฺขเว สติปฏฺฐานํ.}

\prose{กตมา จ ภิกฺขเว สติปฏฺฐานภาวนา, อิธ ภิกฺขเว ภิกฺขุ สมุทยธมฺมานุปสฺสี กายสฺมิํ วิหรติ, วยธมฺมานุปสฺสี กายสฺมิํ วิหรติ, สมุทยวยธมฺมานุปสฺสี กายสฺมิํ วิหรติ อาตาปี สมฺปชาโน สติมา วิเนยฺย โลเก อภิชฺฌาโทมนสฺสํ.\csromanspacer{} สมุทยธมฺมานุปสฺสี เวทนาสุ วิหรติ ฯเปฯ.\csromanspacer{} สมุทยธมฺมานุปสฺสี จิตฺเต วิหรติ ฯเปฯ.\csromanspacer{} สมุทยธมฺมานุปสฺสี ธมฺเมสุ วิหรติ, วยธมฺมานุปสฺสี ธมฺเมสุ วิหรติ, สมุทยวยธมฺมานุปสฺสี ธมฺเมสุ วิหรติ อาตาปี สมฺปชาโน สติมา วิเนยฺย โลเก อภิชฺฌาโทมนสฺสํ.\csromanspacer{} อยํ วุจฺจติ ภิกฺขเว สติปฏฺฐานภาวนา.}

\prose{กตมา จ ภิกฺขเว สติปฏฺฐานภาวนาคามินี ปฏิปทา, อยเมว อริโย อฏฺฐงฺคิโก มคฺโค.\csromanspacer{} เสยฺยถิทํ, สมฺมาทิฏฺฐิ สมฺมาสงฺกปฺโป สมฺมาวาจา สมฺมากมฺมนฺโต สมฺมา-อาชีโว สมฺมาวายาโม สมฺมาสติ สมฺมาสมาธิ.\csromanspacer{} อยํ วุจฺจติ ภิกฺขเว สติปฏฺฐานภาวนาคามินี ปฏิปทาติ.\csromanspacer{}\csromanspacer{}ทสมํ.\csromanspacer{} อนนุสฺสุตวคฺโค จตุตฺโถ.}

\csromanheader{2}{3. สติปฏฺฐานสํยุตฺต}
\titlehead{\textbf{ตสฺสุทฺทานํ}}
\csromanmatikaanchor{mtk.83}
\csromanmatikaanchor{mtk.84}
\csromanmatikaanchor{mtk.85}
\csromantocmarknopage{subparagraph}{5. อมตวคฺค}{mtk.83}
\bookmark[dest={mtk.83},level=5]{5. อมตวคฺค}
\csromantocmarkref{subparagraph}{1-2. อมตสุตฺตาทีนิ}{mtk.84}
\bookmark[dest={mtk.84},level=5]{1-2. อมตสุตฺตาทีนิ}
\csromantocmarkref{subparagraph}{3-4. มคฺคสุตฺตาทีนิ}{mtk.85}
\bookmark[dest={mtk.85},level=5]{3-4. มคฺคสุตฺตาทีนิ}
\csromangathameasure{อนนุสฺสุตํ วิราโค, วิรทฺโธ ภาวนา สติ. \\ อญฺญา ฉนฺทํ ปริญฺญาย, ภาวนา วิภงฺเคน จาติ.}
\csromangathagroup{\csromangathastanza{\csromanfolio{161}อนนุสฺสุตํ วิราโค, วิรทฺโธ ภาวนา สติ. \\ อญฺญา ฉนฺทํ ปริญฺญาย, ภาวนา วิภงฺเคน จาติ.}}

\csromanheader{6}{5. อมตวคฺค}
\csromanheader{2}{1. อมตสุตฺต}
\proseitem{407}{สาวตฺถินิทานํ.\csromanspacer{} จตูสุ ภิกฺขเว สติปฏฺฐาเนสุ สุปฺปติฏฺฐิตจิตฺตา วิหรถ, มา โว อมตํ ปนสฺส.\csromanspacer{} กตเมสุ จตูสุ, อิธ ภิกฺขเว ภิกฺขุ กาเย กายานุปสฺสี วิหรติ อาตาปี สมฺปชาโน สติมา วิเนยฺย โลเก อภิชฺฌาโทมนสฺสํ.\csromanspacer{} เวทนาสุ ฯเปฯ.\csromanspacer{} จิตฺเต ฯเปฯ.\csromanspacer{} ธมฺเมสุ ธมฺมานุปสฺสี วิหรติ อาตาปี สมฺปชาโน สติมา วิเนยฺย โลเก อภิชฺฌาโทมนสฺสํ.\csromanspacer{} อิเมสุ ภิกฺขเว จตูสุ สติปฏฺฐาเนสุ สุปฺปติฏฺฐิตจิตฺตา วิหรถ, มา โว อมตํ ปนสฺสาติ.\csromanspacer{}\csromanspacer{}ปฐมํ.}

\csromanheader{2}{2. สมุทยสุตฺต}
\proseitem{408}{จตุนฺนํ ภิกฺขเว สติปฏฺฐานานํ สมุทยญฺจ อตฺถงฺคมญฺจ เทเสสฺสามิ, ตํ สุณาถ.\csromanspacer{} โก จ ภิกฺขเว กายสฺส สมุทโย, อาหารสมุทยา กายสฺส สมุทโย, อาหารนิโรธา กายสฺส อตฺถงฺคโม.\csromanspacer{} ผสฺสสมุทยา เวทนานํ สมุทโย, ผสฺสนิโรธา เวทนานํ อตฺถงฺคโม.\csromanspacer{} นามรูปสมุทยา จิตฺตสฺส สมุทโย, นามรูปนิโรธา จิตฺตสฺส อตฺถงฺคโม.\csromanspacer{} มนสิการสมุทยา ธมฺมานํ สมุทโย, มนสิการนิโรธา ธมฺมานํ อตฺถงฺคโมติ.\csromanspacer{}\csromanspacer{}ทุติยํ.}

\csromanheader{2}{3. มคฺคสุตฺต}
\proseitem{409}{สาวตฺถินิทานํ.\csromanspacer{} ตตฺร โข ภควา ภิกฺขู อามนฺเตสิ\csromandash{} เอกมิทาหํ ภิกฺขเว สมยํ อุรุเวลายํ วิหรามิ นชฺชา เนรญฺชราย ตีเร อชปาลนิโคฺรเธ ปฐมาภิสมฺพุทฺโธ.\csromanspacer{} ตสฺส มยฺหํ ภิกฺขเว รโหคตสฺส ปฏิสลฺลีนสฺส เอวํ เจตโส ปริวิตกฺโก อุทปาทิ เอกายนฺวายํ มคฺโค \csromanfolio{162}สตฺตานํ วิสุทฺธิยา โสกปริเทวานํ สมติกฺกมาย ทุกฺขโทมนสฺสานํ อตฺถงฺคมาย ญายสฺส อธิคมาย นิพฺพานสฺส สจฺฉิกิริยาย ยทิทํ จตฺตาโร สติปฏฺฐานา.}

\prose{กตเม จตฺตาโร, กาเย วา ภิกฺขุ กายานุปสฺสี วิหเรยฺย อาตาปี สมฺปชาโน สติมา วิเนยฺย โลเก อภิชฺฌาโทมนสฺสํ.\csromanspacer{} เวทนาสุ วา ภิกฺขุ เวทนานุปสฺสี วิหเรยฺย ฯเปฯ.\csromanspacer{} จิตฺเต วา ภิกฺขุ จิตฺตานุปสฺสี วิหเรยฺย ฯเปฯ.\csromanspacer{} ธมฺเมสุ วา ภิกฺขุ ธมฺมานุปสฺสี วิหเรยฺย อาตาปี สมฺปชาโน สติมา วิเนยฺย โลเก อภิชฺฌาโทมนสฺสํ.\csromanspacer{} เอกายนฺวายํ มคฺโค สตฺตานํ วิสุทฺธิยา โสกปริเทวานํ สมติกฺกมาย ทุกฺขโทมนสฺสานํ อตฺถงฺคมาย ญายสฺส อธิคมาย นิพฺพานสฺส สจฺฉิกิริยาย ยทิทํ จตฺตาโร สติปฏฺฐานาติ.}

\prose{อถ โข ภิกฺขเว พฺรหฺมา สหมฺปติ มม เจตสา เจโตปริวิตกฺกมญฺญาย, เสยฺยถาปิ นาม พลวา ปุริโส สมิญฺชิตํ วา พาหํ ปสาเรยฺย, ปสาริตํ วา พาหํ สมิญฺเชยฺย, เอวเมว พฺรหฺมโลเก อนฺตรหิโต มม ปุรโต ปาตุรโหสิ.\csromanspacer{} อถ โข ภิกฺขเว พฺรหฺมา สหมฺปติ เอกํสํ อุตฺตราสงฺคํ กริตฺวา เยนาหํ เตนญฺชลิํ ปณาเมตฺวา มํ เอตทโวจ\csromandash{}เอวเมตํ ภควา, เอวเมตํ สุคต, เอกายนฺวายํ ภนฺเต มคฺโค สตฺตานํ วิสุทฺธิยา โสกปริเทวานํ สมติกฺกมาย ทุกฺขโทมนสฺสานํ อตฺถงฺคมาย ญายสฺส อธิคมาย นิพฺพานสฺส สจฺฉิกิริยาย ยทิทํ จตฺตาโร สติปฏฺฐานา.}

\prose{กตเม จตฺตาโร, กาเย วา ภนฺเต ภิกฺขุ กายานุปสฺสี วิหเรยฺย อาตาปี สมฺปชาโน สติมา วิเนยฺย โลเก อภิชฺฌาโทมนสฺสํ.\csromanspacer{} เวทนาสุ วา ฯเปฯ.\csromanspacer{} จิตฺเต วา ฯเปฯ.\csromanspacer{} ธมฺเมสุ วา ภนฺเต ภิกฺขุ ธมฺมานุปสฺสี วิหเรยฺย อาตาปี สมฺปชาโน สติมา วิเนยฺย โลเก อภิชฺฌาโทมนสฺสํ.\csromanspacer{} เอกายนฺวายํ ภนฺเต มคฺโค สตฺตานํ วิสุทฺธิยา โสกปริเทวานํ สมติกฺกมาย ทุกฺขโทมนสฺสานํ อตฺถงฺคมาย ญายสฺส อธิคมาย นิพฺพานสฺส สจฺฉิกิริยาย ยทิทํ จตฺตาโร สติปฏฺฐานาติ.}

\prose{อิทมโวจ ภิกฺขเว พฺรหฺมา สหมฺปติ, อิทํ วตฺวา อถาปรํ เอตทโวจ\csromandash{}}

\csromangathameasure{“เอกายนํ ชาติขยนฺตทสฺสี, มคฺคํ ปชานาติ หิตานุกมฺปี. \\ เอเตน มคฺเคน ตริํสุ ปุพฺเพ, ตริสฺสนฺติ เย จ ตรนฺติ โอฆนฺ”ติ. ตติยํ.}
\csromangathagroup{\csromangathastanza{“เอกายนํ ชาติขยนฺตทสฺสี, มคฺคํ ปชานาติ หิตานุกมฺปี. \\ เอเตน มคฺเคน ตริํสุ ปุพฺเพ, ตริสฺสนฺติ เย จ ตรนฺติ โอฆนฺ”ติ.\csromanspacer{} ตติยํ.}}

\csromanheader{2}{3. สติปฏฺฐานสํยุตฺต}
\csromanheader{2}{4. สติสุตฺต}
\proseitem{410}{\csromanfolio{163}สโต ภิกฺขเว ภิกฺขุ วิหเรยฺย, อยํ โว อมฺหากํ อนุสาสนี.\csromanspacer{} กถญฺจ ภิกฺขเว ภิกฺขุ สโต โหติ, อิธ ภิกฺขเว ภิกฺขุ กาเย กายานุปสฺสี วิหรติ อาตาปี สมฺปชาโน สติมา วิเนยฺย โลเก อภิชฺฌาโทมนสฺสํ.\csromanspacer{} เวทนาสุ ฯเปฯ.\csromanspacer{} จิตฺเต ฯเปฯ.\csromanspacer{} ธมฺเมสุ ธมฺมานุปสฺสี วิหรติ อาตาปี สมฺปชาโน สติมา วิเนยฺย โลเก อภิชฺฌาโทมนสฺสํ.\csromanspacer{} เอวํ โข ภิกฺขเว ภิกฺขุ สโต โหติ, สโต ภิกฺขเว ภิกฺขุ วิหเรยฺย, อยํ โว อมฺหากํ อนุสาสนีติ.\csromanspacer{}\csromanspacer{}จตุตฺถํ.}

\csromanmatikaanchor{mtk.86}
\csromantocmarkref{subparagraph}{5. กุสลราสิสุตฺต}{mtk.86}
\bookmark[dest={mtk.86},level=5]{5. กุสลราสิสุตฺต}
\csromanheader{6}{5. กุสลราสิสุตฺต}
\proseitem{411}{“กุสลราสี”ติ ภิกฺขเว วทมาโน จตฺตาโร สติปฏฺฐาเน สมฺมา วทมาโน วเทยฺย.\csromanspacer{} เกวโล หายํ ภิกฺขเว กุสลราสิ ยทิทํ จตฺตาโร สติปฏฺฐานา.}

\prose{กตเม จตฺตาโร, อิธ ภิกฺขเว ภิกฺขุ กาเย กายานุปสฺสี วิหรติ อาตาปี สมฺปชาโน สติมา วิเนยฺย โลเก อภิชฺฌาโทมนสฺสํ.\csromanspacer{} เวทนาสุ ฯเปฯ.\csromanspacer{} จิตฺเต ฯเปฯ.\csromanspacer{} ธมฺเมสุ ธมฺมานุปสฺสี วิหรติ อาตาปี สมฺปชาโน สติมา วิเนยฺย โลเก อภิชฺฌาโทมนสฺสํ.\csromanspacer{} เวทนาสุ ฯเปฯ.\csromanspacer{} จิตฺเต ฯเปฯ.\csromanspacer{} ธมฺเมสุ ธมฺมานุปสฺสี วิหรติ อาตาปิ สมฺปชาโน สติมา วิเนยฺย โลเก อภิชฺฌาโทมนสฺสํ. “กุสลราสี”ติ ภิกฺขเว วทมาโน อิเม จตฺตาโร สติปฏฺฐาเน สมฺมา วทมาโน วเทยฺย.\csromanspacer{} เกวโล หายํ ภิกฺขเว กุสลราสิ ยทิทํ จตฺตาโร สติปฏฺฐานาติ.\csromanspacer{}\csromanspacer{}ปญฺจมํ.}

\csromanmatikaanchor{mtk.87}
\csromantocmarkref{subparagraph}{6. ปาติโมกฺขสํวรสุตฺต}{mtk.87}
\bookmark[dest={mtk.87},level=5]{6. ปาติโมกฺขสํวรสุตฺต}
\csromanheader{6}{6. ปาติโมกฺขสํวรสุตฺต}
\proseitem{412}{อถ โข อญฺญตโร ภิกฺขุ เยน ภควา เตนุปสงฺกมิ ฯเปฯ เอกมนฺตํ นิสินฺโน โข โส ภิกฺขุ ภควนฺตํ เอตทโวจ\csromandash{}}

\prose{สาธุ เม ภนฺเต ภควา สํขิตฺเตน ธมฺมํ เทเสตุ, ยมหํ ภควโต ธมฺมํ สุตฺวา เอโก วูปกฏฺโฐ อปฺปมตฺโต อาตาปี ปหิตตฺโต วิหเรยฺยนฺติ.\csromanspacer{} ตสฺมาติห ตฺวํ ภิกฺขุ อาทิเมว วิโสเธหิ กุสเลสุ ธมฺเมสุ.\csromanspacer{} โก จาทิ กุสลานํ ธมฺมานํ, อิธ ตฺวํ ภิกฺขุ ปาติโมกฺขสํวรสํวุโต วิหราหิ อาจารโคจรสมฺปนฺโน อณุมตฺเตสุ วชฺเชสุ ภยทสฺสาวี, สมาทาย สิกฺขสฺสุ สิกฺขาปเทสุ, ยโต โข ตฺวํ ภิกฺขุ ปาติโมกฺขสํวรสํวุโต วิหริสฺสสิ อาจารโคจรสมฺปนฺโน อณุมตฺเตสุ \csromanfolio{164}วชฺเชสุ ภยทสฺสาวี, สมาทาย สิกฺขสฺสุ สิกฺขาปเทสุ, ตโต ตฺวํ ภิกฺขุ สีลํ นิสฺสาย สีเล ปติฏฺฐาย จตฺตาโร สติปฏฺฐาเน ภาเวยฺยาสิ.}

\prose{กตเม จตฺตาโร, อิธ ตฺวํ ภิกฺขุ กาเย กายานุปสฺสี วิหราหิ อาตาปี สมฺปชาโน สติมา วิเนยฺย โลเก อภิชฺฌาโทมนสฺสํ.\csromanspacer{} เวทนาสุ ฯเปฯ.\csromanspacer{} จิตฺเต ฯเปฯ.\csromanspacer{} ธมฺเมสุ ธมฺมานุปสฺสี วิหราหิ อาตาปี สมฺปชาโน สติมา วิเนยฺย โลเก อภิชฺฌาโทมนสฺสํ.\csromanspacer{} ยโต โข ตฺวํ ภิกฺขุ สีลํ นิสฺสาย สีเล ปติฏฺฐาย อิเม จตฺตาโร สติปฏฺฐาเน เอวํ ภาเวสฺสสิ, ตโต ตุยฺหํ ภิกฺขุ ยา รตฺติ วา ทิวโส วา อาคมิสฺสติ, วุทฺธิเยว ปาฏิกงฺขา กุสเลสุ ธมฺเมสุ, โน ปริหานีติ.\csromanspacer{} อถ โข โส ภิกฺขุ ภควโต ภาสิตํ อภินนฺทิตฺวา อนุโมทิตฺวา อุฏฺฐายาสนา ภควนฺตํ อภิวาเทตฺวา ปทกฺขิณํ กตฺวา ปกฺกามิ.}

\prose{อถ โข โส ภิกฺขุ เอโก วูปกฏฺโฐ อปฺปมตฺโต อาตาปี ปหิตตฺโต วิหรนฺโต นจิรสฺเสว ยสฺสตฺถาย กุลปุตฺตา สมฺมเทว อคารสฺมา อนคาริยํ ปพฺพชนฺติ, ตทนุตฺตรํ พฺรหฺมจริยปริโยสานํ ทิฏฺเฐว ธมฺเม สยํ อภิญฺญา สจฺฉิกตฺวา อุปสมฺปชฺช วิหรติ, “ขีณา ชาติ, วุสิตํ พฺรหฺมจริยํ, กตํ กรณียํ, นาปรํ อิตฺถตฺตายา”ติ อพฺภญฺญาสิ, อญฺญตโร จ ปน โส ภิกฺขุ อรหตํ อโหสีติ.\csromanspacer{}\csromanspacer{}ฉฏฺฐํ.}

\csromanmatikaanchor{mtk.88}
\csromantocmarkref{subparagraph}{7. ทุจฺจริตสุตฺต}{mtk.88}
\bookmark[dest={mtk.88},level=5]{7. ทุจฺจริตสุตฺต}
\csromanheader{6}{7. ทุจฺจริตสุตฺต}
\proseitem{413}{อถ โข อญฺญตโร ภิกฺขุ เยน ภควา เตนุปสงฺกมิ ฯเปฯ สาธุ เม ภนฺเต ภควา สํขิตฺเตน ธมฺมํ เทเสตุ, ยมหํ ภควโต ธมฺมํ สุตฺวา เอโก วูปกฏฺโฐ อปฺปมตฺโต อาตาปี ปหิตตฺโต วิหเรยฺยนฺติ.\csromanspacer{} ตสฺมาติห ตฺวํ ภิกฺขุ อาทิเมว วิโสเธหิ กุสเลสุ ธมฺเมสุ.\csromanspacer{} โก จาทิ กุสลานํ ธมฺมานํ, อิธ ตฺวํ ภิกฺขุ กายทุจฺจริตํ ปหาย กายสุจริตํ ภาเวสฺสสิ, วจีทุจฺจริตํ ปหาย วจีสุจริตํ ภาเวสฺสสิ, มโนทุจฺจริตํ ปหาย มโนสุจริตํ ภาเวสฺสสิ, ยโต โข ตฺวํ ภิกฺขุ กายทุจฺจริตํ ปหาย กายสุจริตํ ภาเวสฺสสิ, วจีทุจฺจริตํ ปหาย วจีสุจริตํ ภาเวสฺสสิ, มโนทุจฺจริตํ ปหาย มโนสุจริตํ ภาเวสฺสสิ, ตโต ตฺวํ ภิกฺขุ สีลํ นิสฺสาย สีเล ปติฏฺฐาย จตฺตาโร สติปฏฺฐาเน ภาเวยฺยาสิ.}

\csromanheader{2}{3. สติปฏฺฐานสํยุตฺต}
\csromanmatikaanchor{mtk.89}
\csromantocmarkref{subparagraph}{8-10. มิตฺตสุตฺตาทีนิ}{mtk.89}
\bookmark[dest={mtk.89},level=5]{8-10. มิตฺตสุตฺตาทีนิ}
\prose{\csromanfolio{165}กตเม จตฺตาโร, อิธ ตฺวํ ภิกฺขุ กาเย กายานุปสฺสี วิหราหิ อาตาปี สมฺปชาโน สติมา วิเนยฺย โลเก อภิชฺฌาโทมนสฺสํ.\csromanspacer{} เวทนาสุ ฯเปฯ.\csromanspacer{} จิตฺเต ฯเปฯ.\csromanspacer{} ธมฺเมสุ ธมฺมานุปสฺสี วิหราหิ อาตาปี สมฺปชาโน สติมา วิเนยฺย โลเก อภิชฺฌาโทมนสฺสํ.\csromanspacer{} ยโต โข ตฺวํ ภิกฺขุ สีลํ นิสฺสาย สีเล ปติฏฺฐาย อิเม จตฺตาโร สติปฏฺฐาเน เอวํ ภาเวสฺสสิ, ตโต ตุยฺหํ ภิกฺขุ ยา รตฺติ วา ทิวโส วา อาคมิสฺสติ, วุทฺธิเยว ปาฏิกงฺขา กุสเลสุ ธมฺเมสุ, โน ปริหานีติ ฯเปฯ อญฺญตโร จ ปน โส ภิกฺขุ อรหตํ อโหสีติ.\csromanspacer{}\csromanspacer{}สตฺตมํ.}

\csromanheader{2}{8. มิตฺตสุตฺต}
\proseitem{414}{เย ภิกฺขเว อนุกมฺเปยฺยาถ, เย จ โข โสตพฺพํ มญฺเญยฺยุํ มิตฺตา วา อมจฺจา วา ญาตี วา สาโลหิตา วา, เต โว ภิกฺขเว จตุนฺนํ สติปฏฺฐานานํ ภาวนาย สมาทเปตพฺพา นิเวเสตพฺพา ปติฏฺฐาเปตพฺพา.}

\prose{กตเมสํ จตุนฺนํ, อิธ ภิกฺขเว ภิกฺขุ กาเย กายานุปสฺสี วิหรติ อาตาปี สมฺปชาโน สติมา วิเนยฺย โลเก อภิชฺฌาโทมนสฺสํ.\csromanspacer{} เวทนาสุ ฯเปฯ.\csromanspacer{} จิตฺเต ฯเปฯ.\csromanspacer{} ธมฺเมสุ ธมฺมานุปสฺสี วิหรติ อาตาปี สมฺปชาโน สติมา วิเนยฺย โลเก อภิชฺฌาโทมนสฺสํ.\csromanspacer{} เย ภิกฺขเว อนุกมฺเปยฺยาถ, เย จ โสตพฺพํ มญฺเญยฺยุํ มิตฺตา วา อมจฺจา วา ญาตี วา สาโลหิตา วา, เต โว ภิกฺขเว อิเมสํ จตุนฺนํ สติปฏฺฐานานํ ภาวนาย สมาทเปตพฺพา นิเวเสตพฺพา ปติฏฺฐาเปตพฺพาติ.\csromanspacer{}\csromanspacer{}อฏฺฐมํ.}

\csromanheader{2}{9. เวทนาสุตฺต}
\proseitem{415}{ติสฺโส อิมา ภิกฺขเว เวทนา.\csromanspacer{} กตมา ติสฺโส, สุขา เวทนา ทุกฺขา เวทนา อทุกฺขมสุขา เวทนา, อิมา โข ภิกฺขเว ติสฺโส เวทนา.\csromanspacer{} อิมาสํ โข ภิกฺขเว ติสฺสนฺนํ เวทนานํ ปริญฺญาย จตฺตาโร สติปฏฺฐานา ภาเวตพฺพา.}

\prose{กตเม จตฺตาโร, อิธ ภิกฺขเว ภิกฺขุ กาเย กายานุปสฺสี วิหรติ อาตาปี สมฺปชาโน สติมา วิเนยฺย โลเก อภิชฺฌาโทมนสฺสํ.\csromanspacer{} เวทนาสุ ฯเปฯ.\csromanspacer{} จิตฺเต ฯเปฯ.\csromanspacer{} ธมฺเมสุ ธมฺมานุปสฺสี วิหรติ อาตาปี สมฺปชาโน สติมา วิเนยฺย โลเก อภิชฺฌาโทมนสฺสํ.\csromanspacer{} อิมาสํ โข ภิกฺขเว ติสฺสนฺนํ เวทนานํ ปริญฺญาย อิเม จตฺตาโร สติปฏฺฐานา ภาเวตพฺพาติ.\csromanspacer{}\csromanspacer{}นวมํ.}

\csromanheader{2}{10. อาสวสุตฺต}
\proseitem{416}{\csromanfolio{166}ตโยเม ภิกฺขเว อาสวา.\csromanspacer{} กตเม ตโย, กามาสโว ภวาสโว อวิชฺชาสโว, อิเม โข ภิกฺขเว ตโย อาสวา.\csromanspacer{} อิเมสํ โข ภิกฺขเว ติณฺณนฺนํ อาสวานํ ปหานาย จตฺตาโร สติปฏฺฐานา ภาเวตพฺพา.\csromanspacer{} กตเม จตฺตาโร, อิธ ภิกฺขเว ภิกฺขุ กาเย กายานุปสฺสี วิหรติ อาตาปี สมฺปชาโน สติมา วิเนยฺย โลเก อภิชฺฌาโทมนสฺสํ.\csromanspacer{} เวทนาสุ ฯเปฯ.\csromanspacer{} จิตฺเต ฯเปฯ.\csromanspacer{} ธมฺเมสุ ธมฺมานุปสฺสี วิหรติ อาตาปี สมฺปชาโน สติมา วิเนยฺย โลเก อภิชฺฌาโทมนสฺสํ.\csromanspacer{} อิเมสํ โข ภิกฺขเว ติณฺณนฺนํ อาสวานํ ปหานาย อิเม จตฺตาโร สติปฏฺฐานา ภาเวตพฺพาติ.\csromanspacer{}\csromanspacer{}ทสมํ.\csromanspacer{} อมตวคฺโค ปญฺจโม.}

\titlehead{\textbf{ตสฺสุทฺทานํ}}
\csromangathameasure{อมตํ สมุทโย มคฺโค, สติ กุสลราสิ จ. \\ ปาติโมกฺขํ ทุจฺจริตํ, มิตฺตเวทนา อาสเวน จาติ.}
\csromangathagroup{\csromangathastanza{อมตํ สมุทโย มคฺโค, สติ กุสลราสิ จ. \\ ปาติโมกฺขํ ทุจฺจริตํ, มิตฺตเวทนา อาสเวน จาติ.}}

\chapterhead{6. คงฺคาเปยฺยาลวคฺค}
\csromanmatikaanchor{mtk.90}
\csromanmatikaanchor{mtk.91}
\csromantocmarknopage{subparagraph}{6. คงฺคาเปยฺยาลวคฺค}{mtk.90}
\bookmark[dest={mtk.90},level=5]{6. คงฺคาเปยฺยาลวคฺค}
\csromantocmarkref{subparagraph}{1-12. คงฺคานที-อาทิสุตฺตทฺวาทสก}{mtk.91}
\bookmark[dest={mtk.91},level=5]{1-12. คงฺคานที-อาทิสุตฺตทฺวาทสก}
\csromanheader{6}{1-12. คงฺคานที-อาทิสุตฺตทฺวาทสก}
\proseitem{417-428}{เสยฺยถาปิ ภิกฺขเว คงฺคา นที ปาจีนนินฺนา ปาจีนโปณา ปาจีนปพฺภารา, เอวเมว โข ภิกฺขเว ภิกฺขุ จตฺตาโร สติปฏฺฐาเน ภาเวนฺโต จตฺตาโร สติปฏฺฐาเน พหุลีกโรนฺโต นิพฺพานนินฺโน โหติ นิพฺพานโปโณ นิพฺพานปพฺภาโร.}

\prose{กถญฺจ ภิกฺขเว ภิกฺขุ จตฺตาโร สติปฏฺฐาเน ภาเวนฺโต จตฺตาโร สติปฏฺฐาเน พหุลีกโรนฺโต นิพฺพานนินฺโน โหติ นิพฺพานโปโณ นิพฺพานปพฺภาโร, อิธ ภิกฺขเว ภิกฺขุ กาเย กายานุปสฺสี วิหรติ อาตาปี สมฺปชาโน สติมา วิเนยฺย โลเก อภิชฺฌาโทมนสฺสํ.}

\csromanheader{2}{3. สติปฏฺฐานสํยุตฺต}
\prose{\csromanfolio{167}เวทนาสุ ฯเปฯ.\csromanspacer{} จิตฺเต ฯเปฯ.\csromanspacer{} ธมฺเมสุ ธมฺมานุปสฺสี วิหรติ อาตาปี สมฺปชาโน สติมา วิเนยฺย โลเก อภิชฺฌาโทมนสฺสํ.\csromanspacer{} เอวํ โข ภิกฺขเว ภิกฺขุ จตฺตาโร สติปฏฺฐาเน ภาเวนฺโต จตฺตาโร สติปฏฺฐาเน พหุลีกโรนฺโต นิพฺพานนินฺโน โหติ นิพฺพานโปโณ นิพฺพานปพฺภาโรติ วิตฺถาเรตพฺพํ.\csromanspacer{} คงฺคาเปยฺยาลวคฺโค ฉฏฺโฐ.}

\titlehead{\textbf{ตสฺสุทฺทานํ}}
\csromangathameasure{ฉ ปาจีนโต นินฺนา, ฉ นินฺนา จ สมุทฺทโต. \\ เอเต เทฺว ฉ ทฺวาทส โหนฺติ, วคฺโค เตน ปวุจฺจตีติ.}
\csromangathagroup{\csromangathastanza{ฉ ปาจีนโต นินฺนา, ฉ นินฺนา จ สมุทฺทโต. \\ เอเต เทฺว ฉ ทฺวาทส โหนฺติ, วคฺโค เตน ปวุจฺจตีติ.}}

\chapterhead{7. อปฺปมาทวคฺค}
\csromanmatikaanchor{mtk.92}
\csromanmatikaanchor{mtk.93}
\csromantocmarknopage{subparagraph}{7. อปฺปมาทวคฺค}{mtk.92}
\bookmark[dest={mtk.92},level=5]{7. อปฺปมาทวคฺค}
\csromantocmarkref{subparagraph}{1-10. ตถาคตาทิสุตฺต}{mtk.93}
\bookmark[dest={mtk.93},level=5]{1-10. ตถาคตาทิสุตฺต}
\csromanheader{6}{1-10. ตถาคตาทิสุตฺต}
\proseitem{429-438}{ยาวตา ภิกฺขเว สตฺตา อปทา วา ทฺวิปทา วา จตุปฺปทา วา พหุปฺปทา วาติ วิตฺถาเรตพฺพํ.\csromanspacer{} อปฺปมาทวคฺโค สตฺตโม.}

\titlehead{\textbf{ตสฺสุทฺทานํ}}
\csromangathameasure{ตถาคต ปทํ กูฏํ, มูลํ สาโร จ วสฺสิกํ. \\ ราชา จนฺทิมสูริยา, วตฺเถน ทสมํ ปทนฺติ.}
\csromangathagroup{\csromangathastanza{ตถาคต ปทํ กูฏํ, มูลํ สาโร จ วสฺสิกํ. \\ ราชา จนฺทิมสูริยา, วตฺเถน ทสมํ ปทนฺติ.}}

\chapterhead{8. พลกรณียวคฺค}
\csromanmatikaanchor{mtk.94}
\csromanmatikaanchor{mtk.95}
\csromantocmarknopage{subparagraph}{8. พลกรณียวคฺค}{mtk.94}
\bookmark[dest={mtk.94},level=5]{8. พลกรณียวคฺค}
\csromantocmarkref{subparagraph}{1-12. พลาทิสุตฺต}{mtk.95}
\bookmark[dest={mtk.95},level=5]{1-12. พลาทิสุตฺต}
\csromanheader{6}{1-12. พลาทิสุตฺต}
\proseitem{439-450}{เสยฺยถาปิ ภิกฺขเว เย เกจิ พลกรณียา กมฺมนฺตา กรียนฺตีติ วิตฺถาเรตพฺพํ.\csromanspacer{} พลกรณียวคฺโค อฏฺฐโม.}

\titlehead{\textbf{ตสฺสุทฺทานํ}}
\csromangathameasure{พลํ พีชญฺจ นาโค จ, รุกฺโข กุมฺเภน สูกิยา. \\ อากาเสน จ เทฺว เมฆา, นาวา อาคนฺตุกา นทีติ.}
\csromangathagroup{\csromangathastanza{\csromanfolio{168}พลํ พีชญฺจ นาโค จ, รุกฺโข กุมฺเภน สูกิยา. \\ อากาเสน จ เทฺว เมฆา, นาวา อาคนฺตุกา นทีติ.}}

\chapterhead{9. เอสนาวคฺค}
\csromanmatikaanchor{mtk.96}
\csromanmatikaanchor{mtk.97}
\csromantocmarknopage{subparagraph}{9. เอสนาวคฺค}{mtk.96}
\bookmark[dest={mtk.96},level=5]{9. เอสนาวคฺค}
\csromantocmarkref{subparagraph}{1-10. เอสนาทิสุตฺต}{mtk.97}
\bookmark[dest={mtk.97},level=5]{1-10. เอสนาทิสุตฺต}
\csromanheader{6}{1-10. เอสนาทิสุตฺต}
\proseitem{451-460}{ติสฺโส อิมา ภิกฺขเว เอสนา.\csromanspacer{} กตมา ติสฺโส, กาเมสนา ภเวสนา พฺรหฺมจริเยสนาติ วิตฺถาเรตพฺพํ.\csromanspacer{} เอสนาวคฺโค นวโม.}

\titlehead{\textbf{ตสฺสุทฺทานํ}}
\csromangathameasure{เอสนา วิธา อาสโว, ภโว จ ทุกฺขตา ติสฺโส. \\ ขิลํ มลญฺจ นีโฆ จ, เวทนา ตณฺหา ตสินาย จาติ.}
\csromangathagroup{\csromangathastanza{เอสนา วิธา อาสโว, ภโว จ ทุกฺขตา ติสฺโส. \\ ขิลํ มลญฺจ นีโฆ จ, เวทนา ตณฺหา ตสินาย จาติ.}}

\chapterhead{10. โอฆวคฺค}
\csromanmatikaanchor{mtk.98}
\csromanmatikaanchor{mtk.99}
\csromantocmarknopage{subparagraph}{10. โอฆวคฺค}{mtk.98}
\bookmark[dest={mtk.98},level=5]{10. โอฆวคฺค}
\csromantocmarkref{subparagraph}{1-10. อุทฺธมฺภาคิยาทิสุตฺต}{mtk.99}
\bookmark[dest={mtk.99},level=5]{1-10. อุทฺธมฺภาคิยาทิสุตฺต}
\csromanheader{6}{1-10. อุทฺธมฺภาคิยาทิสุตฺต}
\proseitem{461-470}{ปญฺจิมานิ ภิกฺขเว อุทฺธมฺภาคิยานิ สํโยชนานิ.\csromanspacer{} กตมานิ ปญฺจ, รูปราโค อรูปราโค มาโน อุทฺธจฺจํ อวิชฺชา, อิมานิ โข ภิกฺขเว ปญฺจุทฺธมฺภาคิยานิ สํโยชนานิ.\csromanspacer{} อิเมสํ โข ภิกฺขเว ปญฺจนฺนํ อุทฺธมฺภาคิยานํ สํโยชนานํ อภิญฺญาย ปริญฺญาย ปริกฺขยาย ปหานาย จตฺตาโร สติปฏฺฐานา ภาเวตพฺพา.}

\csromanheader{2}{3. สติปฏฺฐานสํยุตฺต}
\prose{\csromanfolio{169}กตเม จตฺตาโร, อิธ ภิกฺขเว ภิกฺขุ กาเย กายานุปสฺสี วิหรติ อาตาปี สมฺปชาโน สติมา วิเนยฺย โลเก อภิชฺฌาโทมนสฺสํ.\csromanspacer{} เวทนาสุ ฯเปฯ.\csromanspacer{} จิตฺเต ฯเปฯ.\csromanspacer{} ธมฺเมสุ ธมฺมานุปสฺสี วิหรติ อาตาปี สมฺปชาโน สติมา วิเนยฺย โลเก อภิชฺฌาโทมนสฺสํ.\csromanspacer{} อิเมสํ โข ภิกฺขเว ปญฺจนฺนํ อุทฺธมฺภาคิยานํ สํโยชนานํ อภิญฺญาย ปริญฺญาย ปริกฺขยาย ปหานาย อิเม จตฺตาโร สติปฏฺฐานา ภาเวตพฺพาติ.\csromanspacer{}\csromanspacer{}ทสมํ. (ยถา มคฺคสํยุตฺตํ, ตถา สติปฏฺฐานสํยุตฺตํ วิตฺถาเรตพฺพํ.) โอฆวคฺโค ทสโม.}

\titlehead{\textbf{ตสฺสุทฺทานํ}}
\csromangathameasure{โอโฆ โยโค อุปาทานํ, คนฺถา อนุสเยน จ. \\ กามคุณา นีวรณา, ขนฺธา โอรุทฺธมฺภาคิยาติ. \\ \textbf{สติปฏฺฐานสํยุตฺตํ ตติยํ.}}
\csromangathagroup{\csromangathastanza{โอโฆ โยโค อุปาทานํ, คนฺถา อนุสเยน จ. \\ กามคุณา นีวรณา, ขนฺธา โอรุทฺธมฺภาคิยาติ. \\ \textbf{สติปฏฺฐานสํยุตฺตํ ตติยํ.}}}

\csromanheader{2}{4. อินฺทฺริยสํยุตฺต}
\csromanmatikaanchor{mtk.100}
\csromanmatikaanchor{mtk.101}
\csromantocmarknopage{subparagraph}{4. อินฺทฺริยสุํยุตฺต}{mtk.100}
\bookmark[dest={mtk.100},level=5]{4. อินฺทฺริยสุํยุตฺต}
\csromantocmarknopage{subparagraph}{1. สุทฺธิกวคฺค}{mtk.101}
\bookmark[dest={mtk.101},level=5]{1. สุทฺธิกวคฺค}
\csromanheader{6}{1. สุทฺธิกวคฺค}
\csromanheader{2}{1. สุทฺธิกสุตฺต}
\csromanmatikaanchor{mtk.102}
\csromantocmarkref{subparagraph}{1-7. สุทฺธิกสุตฺตาทีนิ}{mtk.102}
\bookmark[dest={mtk.102},level=5]{1-7. สุทฺธิกสุตฺตาทีนิ}
\proseitem{471}{\csromanfolio{170}สาวตฺถินิทานํ.\csromanspacer{} ตตฺร โข ภควา เอตทโวจ\csromandash{}ปญฺจิมานิ ภิกฺขเว อินฺทฺริยานิ.\csromanspacer{} กตมานิ ปญฺจ, สทฺธินฺทฺริยํ วีริยินฺทฺริยํ สตินฺทฺริยํ สมาธินฺทฺริยํ ปญฺญินฺทฺริยํ.\csromanspacer{} อิมานิ โข ภิกฺขเว ปญฺจินฺทฺริยานีติ.\csromanspacer{}\csromanspacer{}ปฐมํ.}

\csromanheader{2}{2. ปฐมโสตาปนฺนสุตฺต}
\proseitem{472}{ปญฺจิมานิ ภิกฺขเว อินฺทฺริยานิ.\csromanspacer{} กตมานิ ปญฺจ, สทฺธินฺทฺริยํ วีริยินฺทฺริยํ สตินฺทฺริยํ สมาธินฺทฺริยํ ปญฺญินฺทฺริยํ.\csromanspacer{} ยโต โข ภิกฺขเว อริยสาวโก อิเมสํ ปญฺจนฺนํ อินฺทฺริยานํ อสฺสาทญฺจ\footnote{สมุทยญฺจ อตฺถงฺคมญฺจ อสฺสาทญฺจ (สฺยา, กํ, อิ, ก) สํ 1. 433 ปิฏฺเฐปิ.} อาทีนวญฺจ นิสฺสรณญฺจ ยถาภูตํ ปชานาติ.\csromanspacer{} อยํ วุจฺจติ ภิกฺขเว อริยสาวโก โสตาปนฺโน อวินิปาตธมฺโม นิยโต สมฺโพธิปรายโณติ.\csromanspacer{}\csromanspacer{}ทุติยํ.}

\csromanheader{2}{3. ทุติยโสตาปนฺนสุตฺต}
\proseitem{473}{ปญฺจิมานิ ภิกฺขเว อินฺทฺริยานิ.\csromanspacer{} กตมานิ ปญฺจ, สทฺธินฺทฺริยํ วีริยินฺทฺริยํ สตินฺทฺริยํ สมาธินฺทฺริยํ ปญฺญินฺทฺริยํ.\csromanspacer{} ยโต โข ภิกฺขเว อริยสาวโก อิเมสํ ปญฺจนฺนํ อินฺทฺริยานํ สมุทยญฺจ อตฺถงฺคมญฺจ อสฺสาทญฺจ อาทีนวญฺจ นิสฺสรณญฺจ ยถาภูตํ ปชานาติ.\csromanspacer{} อยํ วุจฺจติ ภิกฺขเว อริยสาวโก โสตาปนฺโน อวินิปาตธมฺโม นิยโต สมฺโพธิปรายโณติ.\csromanspacer{}\csromanspacer{}ตติยํ.}

\csromanheader{2}{4. ปฐม-อรหนฺตสุตฺต}
\proseitem{474}{ปญฺจิมานิ ภิกฺขเว อินฺทฺริยานิ.\csromanspacer{} กตมานิ ปญฺจ, สทฺธินฺทฺริยํ วีริยินฺทฺริยํ สตินฺทฺริยํ สมาธินฺทฺริยํ ปญฺญินฺทฺริยํ.\csromanspacer{} ยโต โข ภิกฺขเว อริยสาวโก อิเมสํ ปญฺจนฺนํ อินฺทฺริยานํ อสฺสาทญฺจ1 อาทีนวญฺจ นิสฺสรณญฺจ ยถาภูตํ}

\csromanheader{2}{4. อินฺทฺริยสํยุตฺต}
\prose{\csromanfolio{171}วิทิตฺวา อนุปาทาวิมุตฺโต โหติ.\csromanspacer{} อยํ วุจฺจติ ภิกฺขเว ภิกฺขุ อรหํ ขีณาสโว วุสิตวา กตกรณีโย โอหิตภาโร อนุปฺปตฺตสทตฺโถ ปริกฺขีณภวสํโยชโน สมฺมทญฺญาวิมุตฺโตติ.\csromanspacer{}\csromanspacer{}จตุตฺถํ.}

\csromanheader{2}{5. ทุติย-อรหนฺตสุตฺต}
\proseitem{475}{ปญฺจิมานิ ภิกฺขเว อินฺทฺริยานิ.\csromanspacer{} กตมานิ ปญฺจ, สทฺธินฺทฺริยํ วีริยินฺทฺริยํ สตินฺทฺริยํ สมาธินฺทฺริยํ ปญฺญินฺทฺริยํ.\csromanspacer{} ยโต โข ภิกฺขเว ภิกฺขุ อิเมสํ ปญฺจนฺนํ อินฺทฺริยานํ สมุทยญฺจ อตฺถงฺคมญฺจ อสฺสาทญฺจ อาทีนวญฺจ นิสฺสรณญฺจ ยถาภูตํ วิทิตฺวา อนุปาทาวิมุตฺโต โหติ, อยํ วุจฺจติ ภิกฺขเว ภิกฺขุ อรหํ ขีณาสโว วุสิตวา กตกรณีโย โอหิตภาโร อนุปฺปตฺตสทตฺโถ ปริกฺขีณภวสํโยชโน สมฺมทญฺญาวิมุตฺโตติ.\csromanspacer{}\csromanspacer{}ปญฺจมํ.}

\csromanheader{2}{6. ปฐมสมณพฺราหฺมณสุตฺต}
\proseitem{476}{ปญฺจิมานิ ภิกฺขเว อินฺทฺริยานิ.\csromanspacer{} กตมานิ ปญฺจ สทฺธินฺทฺริยํ ฯเปฯ ปญฺญินฺทฺริยํ.\csromanspacer{} เย หิ เกจิ ภิกฺขเว สมณา วา พฺราหฺมณา วา อิเมสํ ปญฺจนฺนํ อินฺทฺริยานํ อสฺสาทญฺจ อาทีนวญฺจ นิสฺสรณญฺจ ยถาภูตํ นปฺปชานนฺติ, น เม เต ภิกฺขเว สมณา วา พฺราหฺมณา วา สมเณสุ วา สมณสมฺมตา, พฺราหฺมเณสุ วา พฺราหฺมณสมฺมตา, น จ ปเนเต อายสฺมนฺโต สามญฺญตฺถํ วา พฺรหฺมญฺญตฺถํ วา ทิฏฺเฐว ธมฺเม สยํ อภิญฺญา สจฺฉิกตฺวา อุปสมฺปชฺช วิหรนฺติ.}

\prose{เย จ โข เกจิ\footnote{เย จ โข เต (สฺยา, กํ, ก) สํ 1. 432 ปิฏฺเฐปิ.} ภิกฺขเว สมณา วา พฺราหฺมณา วา อิเมสํ ปญฺจนฺนํ อินฺทฺริยานํ สมุทยญฺจ อตฺถงฺคมญฺจ อสฺสาทญฺจ อาทีนวญฺจ นิสฺสรณญฺจ ยถาภูตํ ปชานนฺติ, เต โข เม ภิกฺขเว สมณา วา พฺราหฺมณา วา สมเณสุ เจว สมณสมฺมตา, พฺราหฺมเณสุ จ พฺราหฺมณสมฺมตา, เต จ ปนายสฺมนฺโต สามญฺญตฺถญฺจ พฺรหฺมญฺญตฺถญฺจ ทิฏฺเฐว ธมฺเม สยํ อภิญฺญา สจฺฉิกตฺวา อุปสมฺปชฺช วิหรนฺตีติ.\csromanspacer{}\csromanspacer{}ฉฏฺฐํ.}

\csromanheader{2}{7. ทุติยสมณพฺราหฺมณสุตฺต}
\csromanmatikaanchor{mtk.103}
\csromanmatikaanchor{mtk.104}
\csromantocmarkref{subparagraph}{8. ทฏฺฐพฺพสุตฺต}{mtk.103}
\bookmark[dest={mtk.103},level=5]{8. ทฏฺฐพฺพสุตฺต}
\csromantocmarkref{subparagraph}{9-10. วิภงฺคสุตฺตานิ}{mtk.104}
\bookmark[dest={mtk.104},level=5]{9-10. วิภงฺคสุตฺตานิ}
\proseitem{477}{เย หิ เกจิ ภิกฺขเว สมณา วา พฺราหฺมณา วา สทฺธินฺทฺริยํ นปฺปชานนฺติ, สทฺธินฺทฺริยสมุทยํ นปฺปชานนฺติ, สทฺธินฺทฺริยนิโรธํ นปฺปชานนฺติ, สทฺธินฺทฺริยนิโรธคามินิํ ปฏิปทํ นปฺปชานนฺติ.\csromanspacer{} วีริยินฺทฺริยํ นปฺปชานนฺติ ฯเปฯ.\csromanspacer{} สตินฺทฺริยํ \csromanfolio{172}นปฺปชานนฺติ ฯเปฯ.\csromanspacer{} สมาธินฺทฺริยํ นปฺปชานนฺติ ฯเปฯ.\csromanspacer{} ปญฺญินฺทฺริยํ นปฺปชานนฺติ, ปญฺญินฺทฺริยสมุทยํ นปฺปชานนฺติ, ปญฺญินฺทฺริยนิโรธํ นปฺปชานนฺติ, ปญฺญินฺทฺริยนิโรธคามินิํ ปฏิปทํ นปฺปชานนฺติ.\csromanspacer{} น เม เต ภิกฺขเว สมณา วา พฺราหฺมณา วา สมเณสุ วา สมณสมฺมตา, พฺราหฺมเณสุ วา พฺราหฺมณสมฺมตา, น จ ปเนเต อายสฺมนฺโต สามญฺญตฺถํ วา พฺราหฺมญฺญตฺถํ วา ทิฏฺเฐว ธมฺเม สยํ อภิญฺญา สจฺฉิกตฺวา อุปสมฺปชฺช วิหรนฺติ.}

\prose{เย จ โข เกจิ ภิกฺขเว สมณา วา พฺราหฺมณา วา สทฺธินฺทฺริยํ ปชานนฺติ, สทฺธินฺทฺริยสมุทยํ ปชานนฺติ, สทฺธินฺทฺริยนิโรธํ ปชานนฺติ, สทฺธินฺทฺริยนิโรธคามินิํ ปฏิปทํ ปชานนฺติ.\csromanspacer{} วีริยินฺทฺริยํ ปชานนฺติ, วีริยินฺทฺริยสมุทยํ ปชานนฺติ, วีริยินฺทฺริยนิโรธํ ปชานนฺติ, วีริยินฺทฺริยนิโรธคามินิํ ปฏิปทํ ปชานนฺติ.\csromanspacer{} สตินฺทฺริยํ ปชานนฺติ ฯเปฯ.\csromanspacer{} สมาธินฺทฺริยํ ปชานนฺติ ฯเปฯ.\csromanspacer{} ปญฺญินฺทฺริยํ ปชานนฺติ, ปญฺญินฺทฺริยสมุทยํ ปชานนฺติ, ปญฺญินฺทฺริยนิโรธํ ปชานนฺติ, ปญฺญินฺทฺริยนิโรธคามินิํ ปฏิปทํ ปชานนฺติ.\csromanspacer{} เต โข เม ภิกฺขเว สมณา วา พฺราหฺมณา วา สมเณสุ เจว สมณสมฺมตา, พฺราหฺมเณสุ จ พฺราหฺมณสมฺมตา, เต จ ปนายสฺมนฺโต สามญฺญตฺถญฺจ พฺราหฺมญฺญตฺถญฺจ ทิฏฺเฐว ธมฺเม สยํ อภิญฺญา สจฺฉิกตฺวา อุปสมฺปชฺช วิหรนฺตีติ.\csromanspacer{}\csromanspacer{}สตฺตมํ.}

\csromanheader{6}{8. ทฏฺฐพฺพสุตฺต}
\proseitem{478}{ปญฺจิมานิ ภิกฺขเว อินฺทฺริยานิ.\csromanspacer{} กตมานิ ปญฺจ, สทฺธินฺทฺริยํ ฯเปฯ ปญฺญินฺทฺริยํ.\csromanspacer{} กตฺถ จ ภิกฺขเว สทฺธินฺทฺริยํ ทฏฺฐพฺพํ, จตูสุ โสตาปตฺติยงฺเคสุ เอตฺถ สทฺธินฺทฺริยํ ทฏฺฐพฺพํ.\csromanspacer{} กตฺถ จ ภิกฺขเว วีริยินฺทฺริยํ ทฏฺฐพฺพํ, จตูสุ สมฺมปฺปธาเนสุ เอตฺถ วีริยินฺทฺริยํ ทฏฺฐพฺพํ.\csromanspacer{} กตฺถ จ ภิกฺขเว สตินฺทฺริยํ ทฏฺฐพฺพํ, จตูสุ สติปฏฺฐาเนสุ เอตฺถ สตินฺทฺริยํ ทฏฺฐพฺพํ.\csromanspacer{} กตฺถ จ ภิกฺขเว สมาธินฺทฺริยํ ทฏฺฐพฺพํ, จตูสุ ฌาเนสุ เอตฺถ สมาธินฺทฺริยํ ทฏฺฐพฺพํ.\csromanspacer{} กตฺถ จ ภิกฺขเว ปญฺญินฺทฺริยํ ทฏฺฐพฺพํ, จตูสุ อริยสจฺเจสุ เอตฺถ ปญฺญินฺทฺริยํ ทฏฺฐพฺพํ.\csromanspacer{} อิมานิ โข ภิกฺขเว ปญฺจินฺทฺริยานีติ.\csromanspacer{}\csromanspacer{}อฏฺฐมํ.}

\csromanheader{2}{9. ปฐมวิภงฺคสุตฺต}
\proseitem{479}{ปญฺจิมานิ ภิกฺขเว อินฺทฺริยานิ.\csromanspacer{} กตมานิ ปญฺจ, สทฺธินฺทฺริยํ ฯเปฯ ปญฺญินฺทฺริยํ.\csromanspacer{} กตมญฺจ ภิกฺขเว สทฺธินฺทฺริยํ, อิธ ภิกฺขเว อริยสาวโก สทฺโธ โหติ, สทฺทหติ ตถาคตสฺส โพธิํ “อิติปิ โส ภควา อรหํ สมฺมาสมฺพุทฺโธ}

\csromanheader{2}{4. อินฺทฺริยสํยุตฺต}
\prose{\csromanfolio{173}วิชฺชาจรณสมฺปนฺโน สุคโต โลกวิทู อนุตฺตโร ปุริสทมฺมสารถิ สตฺถา เทวมนุสฺสานํ พุทฺโธ ภควา”ติ.\csromanspacer{} อิทํ วุจฺจติ ภิกฺขเว สทฺธินฺทฺริยํ.}

\prose{กตมญฺจ ภิกฺขเว วีริยินฺทฺริยํ, อิธ ภิกฺขเว อริยสาวโก อารทฺธวีริโย วิหรติ อกุสลานํ ธมฺมานํ ปหานาย กุสลานํ ธมฺมานํ อุปสมฺปทาย ถามวา ทฬฺหปรกฺกโม อนิกฺขิตฺตธุโร กุสเลสุ ธมฺเมสุ.\csromanspacer{} อิทํ วุจฺจติ ภิกฺขเว วีริยินฺทฺริยํ.}

\prose{กตมญฺจ ภิกฺขเว สตินฺทฺริยํ, อิธ ภิกฺขเว อริยสาวโก สติมา โหติ ปรเมน สติเนปกฺเกน สมนฺนาคโต, จิรกตมฺปิ จิรภาสิตมฺปิ สริตา อนุสฺสริตา.\csromanspacer{} อิทํ วุจฺจติ ภิกฺขเว สตินฺทฺริยํ.}

\prose{กตมญฺจ ภิกฺขเว สมาธินฺทฺริยํ, อิธ ภิกฺขเว อริยสาวโก โวสฺสคฺคารมฺมณํ กริตฺวา ลภติ สมาธิํ, ลภติ จิตฺตสฺส เอกคฺคตํ.\csromanspacer{} อิทํ วุจฺจติ ภิกฺขเว สมาธินฺทฺริยํ.}

\prose{กตมญฺจ ภิกฺขเว ปญฺญินฺทฺริยํ, อิธ ภิกฺขเว อริยสาวโก ปญฺญวา โหติ อุทยตฺถคามินิยา ปญฺญาย สมนฺนาคโต อริยาย นิพฺเพธิกาย สมฺมา ทุกฺขกฺขยคามินิยา.\csromanspacer{} อิทํ วุจฺจติ ภิกฺขเว ปญฺญินฺทฺริยํ.\csromanspacer{} อิมานิ โข ภิกฺขเว ปญฺจินฺทฺริยานีติ.\csromanspacer{}\csromanspacer{}นวมํ.}

\csromanheader{2}{10. ทุติยวิภงฺคสุตฺต}
\proseitem{480}{ปญฺจิมานิ ภิกฺขเว อินฺทฺริยานิ.\csromanspacer{} กตมานิ ปญฺจ, สทฺธินฺทฺริยํ ฯเปฯ ปญฺญินฺทฺริยํ.\csromanspacer{} กตมญฺจ ภิกฺขเว สทฺธินฺทฺริยํ, อิธ ภิกฺขเว อริยสาวโก สทฺโธ โหติ, สทฺทหติ ตถาคตสฺส โพธิํ “อิติปิ โส ภควา อรหํ สมฺมาสมฺพุทฺโธ วิชฺชาจรณสมฺปนฺโน สุคโต โลกวิทู อนุตฺตโร ปุริสทมฺมสารถิ สตฺถา เทวมนุสฺสานํ พุทฺโธ ภควา”ติ.\csromanspacer{} อิทํ วุจฺจติ ภิกฺขเว สทฺธินฺทฺริยํ.}

\prose{กตมญฺจ ภิกฺขเว วีริยินฺทฺริยํ, อิธ ภิกฺขเว อริยสาวโก อารทฺธวีริโย วิหรติ อกุสลานํ ธมฺมานํ ปหานาย กุสลานํ ธมฺมานํ อุปสมฺปทาย ถามวา ทฬฺหปรกฺกโม อนิกฺขิตฺตธุโร กุสเลสุ ธมฺเมสุ, โส อนุปฺปนฺนานํ ปาปกานํ อกุสลานํ ธมฺมานํ อนุปฺปาทาย ฉนฺทํ ชเนติ วายมติ วีริยํ อารภติ จิตฺตํ ปคฺคณฺหาติ ปทหติ, อุปฺปนฺนานํ ปาปกานํ อกุสลานํ ธมฺมานํ ปหานาย ฉนฺทํ ชเนติ วายมติ \csromanfolio{174}วีริยํ อารภติ จิตฺตํ ปคฺคณฺหาติ ปทหติ, อนุปฺปนฺนานํ กุสลานํ ธมฺมานํ อุปฺปาทาย ฉนฺทํ ชเนติ วายมติ วีริยํ อารภติ จิตฺตํ ปคฺคณฺหาติ ปทหติ, อุปฺปนฺนานํ กุสลานํ ธมฺมานํ ฐิติยา\footnote{สมาปตฺติยา (สฺยา, กํ, ก)} อสมฺโมสาย ภิยฺโยภาวาย เวปุลฺลาย ภาวนาย ปาริปูริยา ฉนฺทํ ชเนติ วายมติ วีริยํ อารภติ จิตฺตํ ปคฺคณฺหาติ ปทหติ.\csromanspacer{} อิทํ วุจฺจติ ภิกฺขเว วีริยินฺทฺริยํ.}

\prose{กตมญฺจ ภิกฺขเว สตินฺทฺริยํ, อิธ ภิกฺขเว อริยสาวโก สติมา โหติ ปรเมน สติเนปกฺเกน สมนฺนาคโต, จิรกตมฺปิ จิรภาสิตมฺปิ สริตา อนุสฺสริตา, โส กาเย กายานุปสฺสี วิหรติ อาตาปี สมฺปชาโน สติมา วิเนยฺย โลเก อภิชฺฌาโทมนสฺสํ.\csromanspacer{} เวทนาสุ ฯเปฯ.\csromanspacer{} จิตฺเต ฯเปฯ.\csromanspacer{} ธมฺเมสุ ธมฺมานุปสฺสี วิหรติ อาตาปี สมฺปชาโน สติมา วิเนยฺย โลเก อภิชฺฌาโทมนสฺสํ.\csromanspacer{} อิทํ วุจฺจติ ภิกฺขเว สตินฺทฺริยํ.}

\prose{กตมญฺจ ภิกฺขเว สมาธินฺทฺริยํ, อิธ ภิกฺขเว อริยสาวโก โวสฺสคฺคารมฺมณํ กริตฺวา ลภติ สมาธิํ, ลภติ จิตฺตสฺส เอกคฺคตํ, โส วิวิจฺเจว กาเมหิ วิวิจฺจ อกุสเลหิ ธมฺเมหิ สวิตกฺกํ สวิจารํ วิเวกชํ ปีติสุขํ ปฐมํ ฌานํ อุปสมฺปชฺช วิหรติ, วิตกฺกวิจารานํ วูปสมา อชฺฌตฺตํ สมฺปสาทนํ เจตโส เอโกทิภาวํ อวิตกฺกํ อวิจารํ สมาธิชํ ปีติสุขํ ทุติยํ ฌานํ อุปสมฺปชฺช วิหรติ, ปีติยา จ วิราคา อุเปกฺขโก จ วิหรติ สโต จ สมฺปชาโน สุขญฺจ กาเยน ปฏิสํเวเทติ, ยํ ตํ อริยา อาจิกฺขนฺติ “อุเปกฺขโก สติมา สุขวิหารี”ติ, ตติยํ ฌานํ อุปสมฺปชฺช วิหรติ.\csromanspacer{} สุขสฺส จ ปหานา ทุกฺขสฺส จ ปหานา ปุพฺเพว โสมนสฺสโทมนสฺสานํ อตฺถงฺคมา อทุกฺขมสุขํ อุเปกฺขาสติปาริสุทฺธิํ จตุตฺถํ ฌานํ อุปสมฺปชฺช วิหรติ.\csromanspacer{} อิทํ วุจฺจติ ภิกฺขเว สมาธินฺทฺริยํ.}

\prose{กตมญฺจ ภิกฺขเว ปญฺญินฺทฺริยํ, อิธ ภิกฺขเว อริยสาวโก ปญฺญวา โหติ อุทยตฺถคามินิยา ปญฺญาย สมนฺนาคโต อริยาย นิพฺเพธิกาย สมฺมา ทุกฺขกฺขยคามินิยา, โส “อิทํ ทุกฺขนฺ”ติ ยถาภูตํ ปชานาติ, “อยํ ทุกฺขสมุทโย”ติ ยถาภูตํ ปชานาติ, “อยํ ทุกฺขนิโรโธ”ติ ยถาภูตํ ปชานาติ. “อยํ ทุกฺขนิโรธคามินี ปฏิปทา”ติ}

\csromanheader{2}{4. อินฺทฺริยสํยุตฺต}
\prose{\csromanfolio{175}ยถาภูตํ ปชานาติ.\csromanspacer{} อิทํ วุจฺจติ ภิกฺขเว ปญฺญินฺทฺริยํ.\csromanspacer{} อิมานิ โข ภิกฺขเว ปญฺจินฺทฺริยานีติ.\csromanspacer{}\csromanspacer{}ทสมํ.\csromanspacer{} สุทฺธิกวคฺโค ปฐโม.}

\titlehead{\textbf{ตสฺสุทฺทานํ}}
\csromangathameasure{สุทฺธิกญฺเจว เทฺว โสตา, อรหนฺตา อปเร ทุเว. \\ สมณพฺราหฺมณา ทฏฺฐพฺพํ, วิภงฺคา อปเร ทุเวติ.}
\csromangathagroup{\csromangathastanza{สุทฺธิกญฺเจว เทฺว โสตา, อรหนฺตา อปเร ทุเว. \\ สมณพฺราหฺมณา ทฏฺฐพฺพํ, วิภงฺคา อปเร ทุเวติ.}}

\csromanmatikaanchor{mtk.105}
\csromantocmarknopage{subparagraph}{2. มุทุตรวคฺค}{mtk.105}
\bookmark[dest={mtk.105},level=5]{2. มุทุตรวคฺค}
\csromanheader{6}{2. มุทุตรวคฺค}
\csromanmatikaanchor{mtk.106}
\csromantocmarkref{subparagraph}{1. ปฏิลาภสุตฺต}{mtk.106}
\bookmark[dest={mtk.106},level=5]{1. ปฏิลาภสุตฺต}
\csromanheader{6}{1. ปฏิลาภสุตฺต}
\proseitem{481}{ปญฺจิมานิ ภิกฺขเว อินฺทฺริยานิ.\csromanspacer{} กตมานิ ปญฺจ, สทฺธินฺทฺริยํ ฯเปฯ ปญฺญินฺทฺริยํ.\csromanspacer{} กตมญฺจ ภิกฺขเว สทฺธินฺทฺริยํ, อิธ ภิกฺขเว อริยสาวโก สทฺโธ โหติ, สทฺทหติ ตถาคตสฺส โพธิํ “อิติปิ โส ภควา อรหํ สมฺมาสมฺพุทฺโธ วิชฺชาจรณสมฺปนฺโน สุคโต โลกวิทู อนุตฺตโร ปุริสทมฺมสารถิ สตฺถา เทวมนุสฺสานํ พุทฺโธ ภควา”ติ.\csromanspacer{} อิทํ วุจฺจติ ภิกฺขเว สทฺธินฺทฺริยํ.}

\prose{กตมญฺจ ภิกฺขเว วีริยินฺทฺริยํ, ยํ โข ภิกฺขเว จตฺตาโร สมฺมปฺปธาเน อารพฺภ วีริยํ ปฏิลภติ.\csromanspacer{} อิทํ วุจฺจติ ภิกฺขเว วีริยินฺทฺริยํ.}

\prose{กตมญฺจ ภิกฺขเว สตินฺทฺริยํ, ยํ โข ภิกฺขเว จตฺตาโร สติปฏฺฐาเน อารพฺภ สติํ ปฏิลภติ.\csromanspacer{} อิทํ วุจฺจติ ภิกฺขเว สตินฺทฺริยํ.}

\prose{กตมญฺจ ภิกฺขเว สมาธินฺทฺริยํ, อิธ ภิกฺขเว อริยสาวโก โวสฺสคฺคารมฺมณํ กริตฺวา ลภติ สมาธิํ, ลภติ จิตฺตสฺส เอกคฺคตํ.\csromanspacer{} อิทํ วุจฺจติ ภิกฺขเว สมาธินฺทฺริยํ.}

\prose{กตมญฺจ ภิกฺขเว ปญฺญินฺทฺริยํ, อิธ ภิกฺขเว อริยสาวโก ปญฺญวา โหติ อุทยตฺถคามินิยา ปญฺญาย สมนฺนาคโต อริยาย นิพฺเพธิกาย สมฺมา ทุกฺขกฺขยคามินิยา.\csromanspacer{} อิทํ วุจฺจติ ภิกฺขเว ปญฺญินฺทฺริยํ.\csromanspacer{} อิมานิ โข ภิกฺขเว ปญฺจินฺทฺริยานีติ.\csromanspacer{}\csromanspacer{}ปฐมํ.}

\csromanheader{2}{2. ปฐมสํขิตฺตสุตฺต}
\csromanmatikaanchor{mtk.107}
\csromanmatikaanchor{mtk.108}
\csromantocmarkref{subparagraph}{2-4. สํขิตฺตสุตฺตานิ}{mtk.107}
\bookmark[dest={mtk.107},level=5]{2-4. สํขิตฺตสุตฺตานิ}
\csromantocmarkref{subparagraph}{5-7. วิตฺถารสุตฺตานิ}{mtk.108}
\bookmark[dest={mtk.108},level=5]{5-7. วิตฺถารสุตฺตานิ}
\proseitem{482}{\csromanfolio{176}ปญฺจิมานิ ภิกฺขเว อินฺทฺริยานิ.\csromanspacer{} กตมานิ ปญฺจ, สทฺธินฺทฺริยํ ฯเปฯ ปญฺญินฺทฺริยํ.\csromanspacer{} อิมานิ โข ภิกฺขเว ปญฺจินฺทฺริยานิ.\csromanspacer{} อิเมสํ โข ภิกฺขเว ปญฺจนฺนํ อินฺทฺริยานํ สมตฺตา ปริปูรตฺตา อรหํ โหติ, ตโต มุทุตเรหิ อนาคามี โหติ, ตโต มุทุตเรหิ สกทาคามี โหติ, ตโต มุทุตเรหิ โสตาปนฺโน โหติ, ตโต มุทุตเรหิ ธมฺมานุสารี โหติ, ตโต มุทุตเรหิ สทฺธานุสารี โหตีติ.\csromanspacer{}\csromanspacer{}ทุติยํ.}

\csromanheader{2}{3. ทุติยสํขิตฺตสุตฺต}
\proseitem{483}{ปญฺจิมานิ ภิกฺขเว อินฺทฺริยานิ.\csromanspacer{} กตมานิ ปญฺจ, สทฺธินฺทฺริยํ ฯเปฯ ปญฺญินฺทฺริยํ.\csromanspacer{} อิมานิ โข ภิกฺขเว ปญฺจินฺทฺริยานิ.\csromanspacer{} อิเมสํ โข ภิกฺขเว ปญฺจนฺนํ อินฺทฺริยานํ สมตฺตา ปริปูรตฺตา อรหํ โหติ, ตโต มุทุตเรหิ อนาคามี โหติ, ตโต มุทุตเรหิ สกทาคามี โหติ, ตโต มุทุตเรหิ โสตาปนฺโน โหติ, ตโต มุทุตเรหิ ธมฺมานุสารี โหติ.\csromanspacer{} ตโต มุทุตเรหิ สทฺธานุสารี โหติ.\csromanspacer{} อิติ โข ภิกฺขเว อินฺทฺริยเวมตฺตตา ผลเวมตฺตตา โหติ, ผลเวมตฺตตา ปุคฺคลเวมตฺตตาติ.\csromanspacer{}\csromanspacer{}ตติยํ.}

\csromanheader{2}{4. ตติยสํขิตฺตสุตฺต}
\proseitem{484}{ปญฺจิมานิ ภิกฺขเว อินฺทฺริยานิ.\csromanspacer{} กตมานิ ปญฺจ, สทฺธินฺทฺริยํ ฯเปฯ ปญฺญินฺทฺริยํ.\csromanspacer{} อิมานิ โข ภิกฺขเว ปญฺจินฺทฺริยานิ.\csromanspacer{} อิเมสํ โข ภิกฺขเว ปญฺจนฺนํ อินฺทฺริยานํ สมตฺตา ปริปูรตฺตา อรหํ โหติ, ตโต มุทุตเรหิ อนาคามี โหติ, ตโต มุทุตเรหิ สกทาคามี โหติ, ตโต มุทุตเรหิ โสตาปนฺโน โหติ, ตโต มุทุตเรหิ ธมฺมานุสารี โหติ, ตโต มุทุตเรหิ สทฺธานุสารี โหติ.\csromanspacer{} อิติ โข ภิกฺขเว ปริปูรํ ปริปูรการี อาราเธติ, ปเทสํ ปเทสการี อาราเธติ, อวญฺฌานิเตฺววาหํ ภิกฺขเว ปญฺจนฺทฺริยานีติ วทามีติ.\csromanspacer{}\csromanspacer{}จตุตฺถํ.}

\csromanheader{2}{5. ปฐมวิตฺถารสุตฺต}
\proseitem{485}{ปญฺจิมานิ ภิกฺขเว อินฺทฺริยานิ.\csromanspacer{} กตมานิ ปญฺจ, สทฺธินฺทฺริยํ ฯเปฯ ปญฺญินฺทฺริยํ.\csromanspacer{} อิมานิ โข ภิกฺขเว ปญฺจินฺทฺริยานิ.\csromanspacer{} อิเมสํ โข ภิกฺขเว ปญฺจนฺนํ อินฺทฺริยานํ}

\csromanheader{2}{4. อินฺทฺริยสํยุตฺต}
\prose{\csromanfolio{177}สมตฺตา ปริปูรตฺตา อรหํ โหติ, ตโต มุทุตเรหิ อนฺตราปรินิพฺพายี โหติ, ตโต มุทุตเรหิ อุปหจฺจปรินิพฺพายี โหติ, ตโต มุทุตเรหิ อสงฺขารปรินิพฺพายี โหติ, ตโต มุทุตเรหิ สสงฺขารปรินิพฺพายี โหติ, ตโต มุทุตเรหิ อุทฺธํโสโต โหติ อกนิฏฺฐคามี, ตโต มุทุตเรหิ สกทาคามี โหติ, ตโต มุทุตเรหิ โสตาปนฺโน โหติ, ตโต มุทุตเรหิ ธมฺมานุสารี โหติ, ตโต มุทุตเรหิ สทฺธานุสารี โหตีติ.\csromanspacer{}\csromanspacer{}ปญฺจมํ.}

\csromanheader{2}{6. ทุติยวิตฺถารสุตฺต}
\proseitem{486}{ปญฺจิมานิ ภิกฺขเว อินฺทฺริยานิ.\csromanspacer{} กตมานิ ปญฺจ, สทฺธินฺทฺริยํ ฯเปฯ ปญฺญินฺทฺริยํ.\csromanspacer{} อิมานิ โข ภิกฺขเว ปญฺจินฺทฺริยานิ.\csromanspacer{} อิเมสํ โข ภิกฺขเว ปญฺจนฺนํ อินฺทฺริยานํ สมตฺตา ปริปูรตฺตา อรหํ โหติ, ตโต มุทุตเรหิ อนฺตราปรินิพฺพายี โหติ, ตโต มุทุตเรหิ อุปหจฺจปรินิพฺพายี โหติ, ตโต มุทุตเรหิ อสงฺขารปรินิพฺพายี โหติ, ตโต มุทุตเรหิ สสงฺขารปรินิพฺพายี โหติ, ตโต มุทุตเรหิ อุทฺธํโสโต โหติ อกนิฏฺฐคามี, ตโต มุทุตเรหิ สกทาคามี โหติ, ตโต มุทุตเรหิ โสตาปนฺโน โหติ, ตโต มุทุตเรหิ ธมฺมานุสารี โหติ, ตโต มุทุตเรหิ สทฺธานุสารี โหติ.\csromanspacer{} อิติ โข ภิกฺขเว อินฺทฺริยเวมตฺตตา ผลเวมตฺตตา โหติ, ผลเวมตฺตตา ปุคฺคลเวมตฺตตา โหตีติ.\csromanspacer{}\csromanspacer{}ฉฏฺฐํ.}

\csromanheader{2}{7. ตติยวิตฺถารสุตฺต}
\csromanmatikaanchor{mtk.109}
\csromantocmarkref{subparagraph}{8-10. ปฏิปนฺนสุตฺตาทีนิ}{mtk.109}
\bookmark[dest={mtk.109},level=5]{8-10. ปฏิปนฺนสุตฺตาทีนิ}
\proseitem{487}{ปญฺจิมานิ ภิกฺขเว อินฺทฺริยานิ.\csromanspacer{} กตมานิ ปญฺจ, สทฺธินฺทฺริยํ ฯเปฯ ปญฺญินฺทฺริยํ.\csromanspacer{} อิมานิ โข ภิกฺขเว ปญฺจินฺทฺริยานิ.\csromanspacer{} อิเมสํ โข ภิกฺขเว ปญฺจนฺนํ อินฺทฺริยานํ สมตฺตา ปริปูรตฺตา อรหํ โหติ, ตโต มุทุตเรหิ อนฺตราปรินิพฺพายี โหติ, ตโต มุทุตเรหิ อุปหจฺจปรินิพฺพายี โหติ, ตโต มุทุตเรหิ อสงฺขารปรินิพฺพายี โหติ, ตโต มุทุตเรหิ สสงฺขารปรินิพฺพายี โหติ, ตโต มุทุตเรหิ อุทฺธํโสโต โหติ อกนิฏฺฐคามี, ตโต มุทุตเรหิ สกทาคามี โหติ, ตโต มุทุตเรหิ โสตาปนฺโน โหติ, ตโต มุทุตเรหิ ธมฺมานุสารี โหติ, ตโต มุทุตเรหิ สทฺธานุสารี โหติ.\csromanspacer{} อิติ โข ภิกฺขเว \csromanfolio{178}ปริปูรํ ปริปูรการี อาราเธติ, ปเทสํ ปเทสการี อาราเธติ, อวญฺฌานิ เตฺววาหํ ภิกฺขเว ปญฺจินฺทฺริยานีติ วทามีติ.\csromanspacer{}\csromanspacer{}สตฺตมํ.}

\csromanheader{2}{8. ปฏิปนฺนสุตฺต}
\proseitem{488}{ปญฺจิมานิ ภิกฺขเว อินฺทฺริยานิ.\csromanspacer{} กตมานิ ปญฺจ, สทฺธินฺทฺริยํ ฯเปฯ ปญฺญินฺทฺริยํ.\csromanspacer{} อิมานิ โข ภิกฺขเว ปญฺจินฺทฺริยานิ.\csromanspacer{} อิเมสํ โข ภิกฺขเว ปญฺจนฺนํ อินฺทฺริยานํ สมตฺตา ปริปูรตฺตา อรหํ โหติ, ตโต มุทุตเรหิ อรหตฺตผลสจฺฉิ กิริยาย ปฏิปนฺโน โหติ, ตโต มุทุตเรหิ อนาคามี โหติ, ตโต มุทุตเรหิ อนาคามิผลสจฺฉิกิริยาย ปฏิปนฺโน โหติ, ตโต มุทุตเรหิ สกทาคามี โหติ, ตโต มุทุตเรหิ สกทาคามิผลสจฺฉิกิริยาย ปฏิปนฺโน โหติ, ตโต มุทุตเรหิ โสตาปนฺโน โหติ, ตโต มุทุตเรหิ โสตาปตฺติผลสจฺฉิกิริยาย ปฏิปนฺโน โหติ.\csromanspacer{} ยสฺส โข ภิกฺขเว อิมานิ ปญฺจินฺทฺริยานิ สพฺเพน สพฺพํ สพฺพถา สพฺพํ นตฺถิ, ตมหํ “พาหิโร ปุถุชฺชนปกฺเข ฐิโต”ติ วทามีติ.\csromanspacer{}\csromanspacer{}อฏฺฐมํ.}

\csromanheader{2}{9. สมฺปนฺนสุตฺต}
\proseitem{489}{อถ โข อญฺญตโร ภิกฺขุ เยน ภควา เตนุปสงฺกมิ, อุปสงฺกมิตฺวา ภควนฺตํ อภิวาเทตฺวา เอกมนฺตํ นิสีทิ, เอกมนฺตํ นิสินฺโน โข โส ภิกฺขุ ภควนฺตํ เอตทโวจ\csromandash{}}

\prose{“อินฺทฺริยสมฺปนฺโน อินฺทฺริยสมฺปนฺโนติ ภนฺเต วุจฺจติ, กิตฺตาวตา นุ โข ภนฺเต อินฺทฺริยสมฺปนฺโน โหตี”ติ.\csromanspacer{} อิธ ภิกฺขุ ภิกฺขุ สทฺธินฺทฺริยํ ภาเวติ อุปสมคามิํ สมฺโพธคามิํ, วีริยินฺทฺริยํ ภาเวติ อุปสมคามิํ สมฺโพธคามิํ, สตินฺทฺริยํ ภาเวติ อุปสมคามิํ สมฺโพธคามิํ, สมาธินฺทฺริยํ ภาเวติ อุปสมคามิํ สมฺโพธคามิํ, ปญฺญินฺทฺริยํ ภาเวติ อุปสมคามิํ สมฺโพธคามิํ.\csromanspacer{} เอตฺตาวตา โข ภิกฺขุ ภิกฺขุ อินฺทฺริยสมฺปนฺโน โหตีติ.\csromanspacer{}\csromanspacer{}นวมํ.}

\csromanheader{2}{10. อาสวกฺขยสุตฺต}
\proseitem{490}{ปญฺจิมานิ ภิกฺขเว อินฺทฺริยานิ.\csromanspacer{} กตมานิ ปญฺจ, สทฺธินฺทฺริยํ ฯเปฯ ปญฺญินฺทฺริยํ.\csromanspacer{} อิมานิ โข ภิกฺขเว ปญฺจินฺทฺริยานิ.\csromanspacer{} อิเมสํ โข ภิกฺขเว ปญฺจนฺนํ อินฺทฺริยานํ}

\csromanheader{2}{4. อินฺทฺริยสํยุตฺต}
\csromanmatikaanchor{mtk.110}
\csromanmatikaanchor{mtk.111}
\csromantocmarknopage{subparagraph}{3. ฉฬินฺทฺริยวคฺค}{mtk.110}
\bookmark[dest={mtk.110},level=5]{3. ฉฬินฺทฺริยวคฺค}
\csromantocmarkref{subparagraph}{1-10. ปุนพฺภวสุตฺตาทีนิ}{mtk.111}
\bookmark[dest={mtk.111},level=5]{1-10. ปุนพฺภวสุตฺตาทีนิ}
\prose{\csromanfolio{179}ภาวิตตฺตา พหุลีกตตฺตา ภิกฺขุ อาสวานํ ขยา อนาสวํ เจโตวิมุตฺติํ ปญฺญาวิมุตฺติํ ทิฏฺเฐว ธมฺเม สยํ อภิญฺญา สจฺฉิกตฺวา อุปสมฺปชฺช วิหรตีติ.\csromanspacer{}\csromanspacer{}ทสมํ.\csromanspacer{} มุทุตรวคฺโค ทุติโย.}

\titlehead{\textbf{ตสฺสุทฺทานํ}}
\csromangathameasure{ปฏิลาโภ ตโย สํขิตฺตา, วิตฺถารา อปเร ตโย. \\ ปฏิปนฺโน จ สมฺปนฺโน, ทสมํ อาสวกฺขยนฺติ.}
\csromangathagroup{\csromangathastanza{ปฏิลาโภ ตโย สํขิตฺตา, วิตฺถารา อปเร ตโย. \\ ปฏิปนฺโน จ สมฺปนฺโน\footnote{ปฏิปนฺโน จูปสโม (สฺยา, กํ, อิ, ก)}, ทสมํ อาสวกฺขยนฺติ.}}

\csromanheader{6}{3. ฉฬินฺทฺริยวคฺค}
\csromanheader{2}{1. ปุนพฺภวสุตฺต}
\proseitem{491}{ปญฺจิมานิ ภิกฺขเว อินฺทฺริยานิ.\csromanspacer{} กตมานิ ปญฺจ, สทฺธินฺทฺริยํ ฯเปฯ ปญฺญินฺทฺริยํ.\csromanspacer{} ยาวกีวญฺจาหํ ภิกฺขเว อิเมสํ ปญฺจนฺนํ อินฺทฺริยานํ สมุทยญฺจ อตฺถงฺคมญฺจ อสฺสาทญฺจ อาทีนวญฺจ นิสฺสรณญฺจ ยถาภูตํ นาพฺภญฺญาสิํ, เนว ตาวาหํ ภิกฺขเว สเทวเก โลเก สมารเก สพฺรหฺมเก สสฺสมณพฺราหฺมณิยา ปชาย สเทวมนุสฺสาย อนุตฺตรํ สมฺมาสมฺโพธิํ อภิสมฺพุทฺโธติ ปจฺจญฺญาสิํ\footnote{อภิสมฺพุทฺโธ ปจฺจญฺญาสิํ (สี, สฺยา, กํ)}.\csromanspacer{} ยโต จ ขฺวาหํ ภิกฺขเว อิเมสํ ปญฺจนฺนํ อินฺทฺริยานํ สมุทยญฺจ อตฺถงฺคมญฺจ อสฺสาทญฺจ อาทีนวญฺจ นิสฺสรณญฺจ ยถาภูตํ อพฺภญฺญาสิํ, อถาหํ ภิกฺขเว สเทวเก โลเก สมารเก สพฺรหฺมเก สสฺสมณพฺราหฺมณิยา ปชาย สเทวมนุสฺสาย อนุตฺตรํ สมฺมาสมฺโพธิํ อภิสมฺพุทฺโธติ ปจฺจญฺญาสิํ, ญาณญฺจ ปน เม ทสฺสนํ อุทปาทิ, อกุปฺปา เม วิมุตฺติ\footnote{เจโตวิมุตฺติ (สี, อิ, ก)}, อยมนฺติมา ชาติ, นตฺถิทานิ ปุนพฺภโวติ.\csromanspacer{}\csromanspacer{}ปฐมํ.}

\csromanheader{2}{2. ชีวิตินฺทฺริยสุตฺต}
\proseitem{492}{ตีณิมานิ ภิกฺขเว อินฺทฺริยานิ.\csromanspacer{} กตมานิ ตีณิ, อิตฺถินฺทฺริยํ ปุริสินฺทฺริยํ ชีวิตินฺทฺริยํ.\csromanspacer{} อิมานิ โข ภิกฺขเว ตีณิ อินฺทฺริยานีติ.\csromanspacer{}\csromanspacer{}ทุติยํ.}

\csromanheader{2}{3. อญฺญินฺทฺริยสุตฺต}
\proseitem{493}{\csromanfolio{180}ตีณิมานิ ภิกฺขเว อินฺทฺริยานิ.\csromanspacer{} กตมานิ ตีณิ, อนญฺญาตญฺญสฺสามีตินฺทฺริยํ อญฺญินฺทฺริยํ อญฺญาตาวินฺทฺริยํ.\csromanspacer{} อิมานิ โข ภิกฺขเว ตีณิ อินฺทฺริยานีติ.\csromanspacer{}\csromanspacer{}ตติยํ.}

\csromanheader{2}{4. เอกพีชีสุตฺต}
\proseitem{494}{ปญฺจิมานิ ภิกฺขเว อินฺทฺริยานิ.\csromanspacer{} กตมานิ ปญฺจ, สทฺธินฺทฺริยํ ฯเปฯ ปญฺญินฺทฺริยํ.\csromanspacer{} อิมานิ โข ภิกฺขเว ปญฺจินฺทฺริยานิ.\csromanspacer{} อิเมสํ โข ภิกฺขเว ปญฺจนฺนํ อินฺทฺริยานํ สมตฺตา ปริปูรตฺตา อรหํ โหติ, ตโต มุทุตเรหิ อนฺตราปรินิพฺพายี โหติ, ตโต มุทุตเรหิ อุปหจฺจปรินิพฺพายี โหติ, ตโต มุทุตเรหิ อสงฺขารปรินิพฺพายี โหติ, ตโต มุทุตเรหิ สสงฺขารปรินิพฺพายี โหติ, ตโต มุทุตเรหิ อุทฺธํโสโต โหติ อกนิฏฺฐคามี, ตโต มุทุตเรหิ สกทาคามี โหติ, ตโต มุทุตเรหิ เอกพีชี\footnote{เอกพีชิ (ก)} โหติ, ตโต มุทุตเรหิ โกลํโกโล โหติ, ตโต มุทุตเรหิ สตฺตกฺขตฺตุปรโม โหติ, ตโต มุทุตเรหิ ธมฺมานุสารี โหติ, ตโต มุทุตเรหิ สทฺธานุสารี โหตีติ.\csromanspacer{}\csromanspacer{}จตุตฺถํ.}

\csromanheader{2}{5. สุทฺธกสุตฺต}
\proseitem{495}{ฉยิมานิ ภิกฺขเว อินฺทฺริยานิ.\csromanspacer{} กตมานิ ฉ, จกฺขุนฺทฺริยํ โสตินฺทฺริยํ ฆานินฺทฺริยํ ชิวฺหินฺทฺริยํ กายินฺทฺริยํ มนินฺทฺริยํ.\csromanspacer{} อิมานิ โข ภิกฺขเว ฉ อินฺทฺริยานีติ.\csromanspacer{}\csromanspacer{}ปญฺจมํ.}

\csromanheader{2}{6. โสตาปนฺนสุตฺต}
\proseitem{496}{ฉยิมานิ ภิกฺขเว อินฺทฺริยานิ.\csromanspacer{} กตมานิ ฉ, จกฺขุนฺทฺริยํ ฯเปฯ มนินฺทฺริยํ.\csromanspacer{} ยโต โข ภิกฺขเว อริยสาวโก อิเมสํ ฉนฺนํ อินฺทฺริยานํ สมุทยญฺจ อตฺถงฺคมญฺจ อสฺสาทญฺจ อาทีนวญฺจ นิสฺสรณญฺจ ยถาภูตํ ปชานาติ.\csromanspacer{} อยํ วุจฺจติ ภิกฺขเว อริยสาวโก โสตาปนฺโน อวินิปาตธมฺโม นิยโต สมฺโพธิปรายโณติ.\csromanspacer{}\csromanspacer{}ฉฏฺฐํ.}

\csromanheader{2}{4. อินฺทฺริยสํยุตฺต}
\csromanheader{2}{7. อรหนฺตสุตฺต}
\proseitem{497}{\csromanfolio{181}ฉยิมานิ ภิกฺขเว อินฺทฺริยานิ.\csromanspacer{} กตมานิ ฉ, จกฺขุนฺทฺริยํ โสตินฺทฺริยํ ฆานินฺทฺริยํ ชิวฺหินฺทฺริยํ กายินฺทฺริยํ มนินฺทฺริยํ.\csromanspacer{} ยโต โข ภิกฺขเว ภิกฺขุ อิเมสํ ฉนฺนํ อินฺทฺริยานํ สมุทยญฺจ อตฺถงฺคมญฺจ อสฺสาทญฺจ อาทีนวญฺจ นิสฺสรณญฺจ ยถาภูตํ วิทิตฺวา อนุปาทาวิมุตฺโต โหติ.\csromanspacer{} อยํ วุจฺจติ ภิกฺขเว ภิกฺขุ อรหํ ขีณาสโว วุสิตวา กตกรณีโย โอหิตภาโร อนุปฺปตฺตสทตฺโถ ปริกฺขีณภวสํโยชโน สมฺมทญฺญาวิมุตฺโตติ.\csromanspacer{}\csromanspacer{}สตฺตมํ.}

\csromanheader{2}{8. สมฺพุทฺธสุตฺต}
\proseitem{498}{ฉยิมานิ ภิกฺขเว อินฺทฺริยานิ.\csromanspacer{} กตมานิ ฉ, จกฺขุนฺทฺริยํ โสตินฺทฺริยํ ฆานินฺทฺริยํ ชิวฺหินฺทฺริยํ กายินฺทฺริยํ มนินฺทฺริยํ.\csromanspacer{} ยาวกีวญฺจาหํ ภิกฺขเว อิเมสํ ฉนฺนํ อินฺทฺริยานํ สมุทยญฺจ อตฺถงฺคมญฺจ อสฺสาทญฺจ อาทีนวญฺจ นิสฺสรณญฺจ ยถาภูตํ นาพฺภญฺญาสิํ, เนว ตาวาหํ ภิกฺขเว สเทวเก โลเก สมารเก สพฺรหฺมเก สสฺสมณพฺราหฺมณิยา ปชาย สเทวมนุสฺสาย อนุตฺตรํ สมฺมาสมฺโพธิํ อภิสมฺพุทฺโธติ ปจฺจญฺญาสิํ.\csromanspacer{} ยโต จ ขฺวาหํ ภิกฺขเว อิเมสํ ฉนฺนํ อินฺทฺริยานํ สมุทยญฺจ อตฺถงฺคมญฺจ อสฺสาทญฺจ อาทีนวญฺจ นิสฺสรณญฺจ ยถาภูตํ อพฺภญฺญาสิํ, อถาหํ ภิกฺขเว สเทวเก โลเก สมารเก สพฺรหฺมเก สสฺสมณพฺราหฺมณิยา ปชาย สเทวมนุสฺสาย อนุตฺตรํ สมฺมาสมฺโพธิํ อภิสมฺพุทฺโธติ ปจฺจญฺญาสิํ, ญาณญฺจ ปน เม ทสฺสนํ อุทปาทิ, อกุปฺปา เม วิมุตฺติ, อยมนฺติมา ชาติ, นตฺถิทานิ ปุนพฺภโวติ.\csromanspacer{}\csromanspacer{}อฏฺฐมํ.}

\csromanheader{2}{9. ปฐมสมณพฺราหฺมณสุตฺต}
\proseitem{499}{ฉยิมานิ ภิกฺขเว อินฺทฺริยานิ.\csromanspacer{} กตมานิ ฉ, จกฺขุนฺทฺริยํ โสตินฺทฺริยํ ฆานินฺทฺริยํ ชิวฺหินฺทฺริยํ กายินฺทฺริยํ มนินฺทฺริยํ.\csromanspacer{} เย หิ เกจิ ภิกฺขเว สมณา วา พฺราหฺมณา วา อิเมสํ ฉนฺนํ อินฺทฺริยานํ สมุทยญฺจ อตฺถงฺคมญฺจ อสฺสาทญฺจ อาทีนวญฺจ นิสฺสรณญฺจ ยถาภูตํ นปฺปชานนฺติ, น เม เต ภิกฺขเว สมณา วา พฺราหฺมณา วา สมเณสุ วา สมณสมฺมตา, พฺราหฺมเณสุ วา พฺราหฺมณสมฺมตา, น จ ปเนเต อายสฺมนฺโต สามญฺญตฺถํ วา พฺรหฺมญฺญตฺถํ \csromanfolio{182}วา ทิฏฺเฐว ธมฺเม สยํ อภิญฺญา สจฺฉิกตฺวา อุปสมฺปชฺช วิหรนฺติ.\csromanspacer{} เย จ โข เกจิ ภิกฺขเว สมณา วา พฺราหฺมณา วา อิเมสํ ฉนฺนํ อินฺทฺริยานํ สมุทยญฺจ อตฺถงฺคมญฺจ อสฺสาทญฺจ อาทีนวญฺจ นิสฺสรณญฺจ ยถาภูตํ ปชานนฺติ, เต โข เม ภิกฺขเว สมณา วา พฺราหฺมณา วา สมเณสุ เจว สมณสมฺมตา, พฺราหฺมเณสุ จ พฺราหฺมณสมฺมตา, เต จ ปนายสฺมนฺโต สามญฺญตฺถญฺจ พฺรหฺมญฺญตฺถญฺจ ทิฏฺเฐว ธมฺเม สยํ อภิญฺญา สจฺฉิกตฺวา อุปสมฺปชฺช วิหรนฺตีติ.\csromanspacer{}\csromanspacer{}นวมํ.}

\csromanheader{2}{10. ทุติยสมณพฺราหฺมณสุตฺต}
\proseitem{500}{เย หิ เกจิ ภิกฺขเว สมณา วา พฺราหฺมณา วา จกฺขุนฺทฺริยํ นปฺปชานนฺติ, จกฺขุนฺทฺริยสมุทยํ นปฺปชานนฺติ, จกฺขุนฺทฺริยนิโรธํ นปฺปชานนฺติ, จกฺขุนฺทฺริยนิโรธคามินิํ ปฏิปทํ นปฺปชานนฺติ.\csromanspacer{} โสตินฺทฺริยํ ฯเปฯ.\csromanspacer{} ฆานินฺทฺริยํ ฯเปฯ.\csromanspacer{} ชิวฺหานฺทฺริยํ ฯเปฯ.\csromanspacer{} กายินฺทฺริยํ ฯเปฯ.\csromanspacer{} มนินฺทฺริยํ นปฺปชานนฺติ, มนินฺทฺริยสมุทยํ นปฺปชานนฺติ, มนินฺทฺริยนิโรธํ นปฺปชานนฺติ.\csromanspacer{} มนินฺทฺริยนิโรธคามินิํ ปฏิปทํ นปฺปชานนฺติ.\csromanspacer{} น เม เต ภิกฺขเว ฯเปฯ สยํ อภิญฺญา สจฺฉิกตฺวา อุปสมฺปชฺช วิหรนฺติ.}

\prose{เย จ โข เกจิ ภิกฺขเว สมณา วา พฺราหฺมณา วา จกฺขุนฺทฺริยํ ปชานนฺติ, จกฺขุนฺทฺริยสมุทยํ ปชานนฺติ, จกฺขุนฺทฺริยนิโรธํ ปชานนฺติ, จกฺขุนฺทฺริยนิโรธคามินิํ ปฏิปทํ ปชานนฺติ.\csromanspacer{} โสตินฺทฺริยํ ฯเปฯ.\csromanspacer{} ฆานินฺทฺริยํ ฯเปฯ.\csromanspacer{} ชิวฺหินฺทฺริยํ ฯเปฯ.\csromanspacer{} กายินฺทฺริยํ ฯเปฯ.\csromanspacer{} มนินฺทฺริยํ ปชานนฺติ, มนินฺทฺริยสมุทยํ ปชานนฺติ, มนินฺทฺริยนิโรธํ ปชานนฺติ, มนินฺทฺริยนิโรธคามินิํ ปฏิปทํ ปชานนฺติ, เต โข เม ภิกฺขเว สมณา วา พฺราหฺมณา วา สมเณสุ เจว สมณสมฺมตา, พฺราหฺมเณสุ จ พฺราหฺมณสมฺมตา, เต จ ปนายสฺมนฺโต สามญฺญตฺถญฺจ พฺรหฺมญฺญตฺถญฺจ ทิฏฺเฐว ธมฺเม สยํ อภิญฺญา สจฺฉิกตฺวา อุปสมฺปชฺช วิหรนฺตีติ.\csromanspacer{}\csromanspacer{}ทสมํ.}

\vspace{\tipitakaparskip}
\csromancenter{ฉฬินฺทิยวคฺโค ตติโย.}
\titlehead{\textbf{ตสฺสุทฺทานํ}}
\csromangathameasure{ปุนพฺภโว ชีวิตญฺญาย, เอกพีชี จ สุทฺธกํ. \\ โสโต อรหสมฺพุทฺโธ, เทฺว จ สมณพฺราหฺมณาติ.}
\csromangathagroup{\csromangathastanza{ปุนพฺภโว ชีวิตญฺญาย, เอกพีชี จ สุทฺธกํ. \\ โสโต อรหสมฺพุทฺโธ, เทฺว จ สมณพฺราหฺมณาติ.}}

\csromanheader{2}{4. อินฺทฺริยสํยุตฺต}
\csromanmatikaanchor{mtk.112}
\csromantocmarknopage{subparagraph}{4. สุขินฺทฺริยวคฺค}{mtk.112}
\bookmark[dest={mtk.112},level=5]{4. สุขินฺทฺริยวคฺค}
\csromanheader{6}{4. สุขินฺทฺริยวคฺค}
\csromanheader{2}{1. สุทฺธิกสุตฺต}
\csromanmatikaanchor{mtk.113}
\csromanmatikaanchor{mtk.114}
\csromantocmarkref{subparagraph}{1-3. สุทฺธิกสุตฺตาทีนิ}{mtk.113}
\bookmark[dest={mtk.113},level=5]{1-3. สุทฺธิกสุตฺตาทีนิ}
\csromantocmarkref{subparagraph}{4-8. ปฐมสมณพฺราหฺมณสุตฺตาทีนิ}{mtk.114}
\bookmark[dest={mtk.114},level=5]{4-8. ปฐมสมณพฺราหฺมณสุตฺตาทีนิ}
\proseitem{501}{\csromanfolio{183}ปญฺจิมานิ ภิกฺขเว อินฺทฺริยานิ.\csromanspacer{} กตมานิ ปญฺจ, สุขินฺทฺริยํ ทุกฺขินฺทฺริยํ โสมนสฺสินฺทฺริยํ โทมนสฺสินฺทฺริยํ อุเปกฺขินฺทฺริยํ.\csromanspacer{} อิมานิ โข ภิกฺขเว ปญฺจินฺทฺริยานีติ.\csromanspacer{}\csromanspacer{}ปฐมํ.}

\csromanheader{2}{2. โสตาปนฺนสุตฺต}
\proseitem{502}{ปญฺจิมานิ ภิกฺขเว อินฺทฺริยานิ.\csromanspacer{} กตมานิ ปญฺจ, สุขินฺทฺริยํ ทุกฺขินฺทฺริยํ โสมนสฺสินฺทฺริยํ โทมนสฺสินฺทฺริยํ อุเปกฺขินฺทฺริยํ.\csromanspacer{} ยโต โข ภิกฺขเว อริยสาวโก อิเมสํ ปญฺจนฺนํ อินฺทฺริยานํ สมุทยญฺจ อตฺถงฺคมญฺจ อสฺสาทญฺจ อาทีนวญฺจ นิสฺสรณญฺจ ยถาภูตํ ปชานาติ.\csromanspacer{} อยํ วุจฺจติ ภิกฺขเว อริยสาวโก โสตาปนฺโน อวินิปาตธมฺโม นิยโต สมฺโพธิปรายโณติ.\csromanspacer{}\csromanspacer{}ทุติยํ.}

\csromanheader{2}{3. อรหนฺตสุตฺต}
\proseitem{503}{ปญฺจิมานิ ภิกฺขเว อินฺทฺริยานิ.\csromanspacer{} กตมานิ ปญฺจ, สุขินฺทฺริยํ ทุกฺขินฺทฺริยํ โสมนสฺสินฺทฺริยํ โทมนสฺสินฺทฺริยํ อุเปกฺขินฺทฺริยํ.\csromanspacer{} ยโต โข ภิกฺขเว ภิกฺขุ อิเมสํ ปญฺจนฺนํ อินฺทฺริยานํ สมุทยญฺจ อตฺถงฺคมญฺจ อสฺสาทญฺจ อาทีนวญฺจ นิสฺสรณญฺจ ยถาภูตํ วิทิตฺวา อนุปาทาวิมุตฺโต โหติ.\csromanspacer{} อยํ วุจฺจติ ภิกฺขเว ภิกฺขุ อรหํ ขีณาสโว วุสิตวา กตกรณีโย โอหิตภาโร อนุปฺปตฺตสทตฺโถ ปริกฺขีณภวสํโยชโน สมฺมทญฺญาวิมุตฺโตติ.\csromanspacer{}\csromanspacer{}ตติยํ.}

\csromanheader{2}{4. ปฐมสมณพฺราหฺมณสุตฺต}
\proseitem{504}{ปญฺจิมานิ ภิกฺขเว อินฺทฺริยานิ.\csromanspacer{} กตมานิ ปญฺจ, สุขินฺทฺริยํ ทุกฺขินฺทฺริยํ โสมนสฺสินฺทฺริยํ โทมนสฺสินฺทฺริยํ อุเปกฺขินฺทฺริยํ.\csromanspacer{} เย หิ เกจิ ภิกฺขเว สมณา วา พฺราหฺมณา วา อิเมสํ ปญฺจนฺนํ อินฺทฺริยานํ สมุทยญฺจ อตฺถงฺคมญฺจ อสฺสาทญฺจ อาทีนวญฺจ นิสฺสรณญฺจ ยถาภูตํ นปฺปชานนฺติ, น เม เต ภิกฺขเว สมณา วา พฺราหฺมณา วา สมเณสุ วา สมณสมฺมตา, พฺราหฺมเณสุ วา พฺราหฺมณสมฺมตา, น จ ปเนเต อายสฺมนฺโต สามญฺญตฺถํ วา พฺรหฺมญฺญตฺถํ วา ทิฏฺเฐว ธมฺเม สยํ อภิญฺญา สจฺฉิกตฺวา อุปสมฺปชฺช วิหรนฺติ.}

\prose{\csromanfolio{184}เย จ โข เกจิ ภิกฺขเว สมณา วา พฺราหฺมณา วา อิเมสํ ปญฺจนฺนํ อินฺทฺริยานํ สมุทยญฺจ อตฺถงฺคมญฺจ อสฺสาทญฺจ อาทีนวญฺจ นิสฺสรณญฺจ ยถาภูตํ ปชานนฺติ, เต โข เม ภิกฺขเว สมณา วา พฺราหฺมณา วา สมเณสุ เจว สมณสมฺมตา, พฺราหฺมเณสุ จ พฺราหฺมณสมฺมตา, เต จ ปนายสฺมนฺโต สามญฺญตฺถญฺจ พฺรหฺมญฺญตฺถญฺจ ทิฏฺเฐว ธมฺเม สยํ อภิญฺญา สจฺฉิกตฺวา อุปสมฺปชฺช วิหรนฺตีติ.\csromanspacer{}\csromanspacer{}จตุตฺถํ.}

\csromanheader{2}{5. ทุติยสมณพฺราหฺมณสุตฺต}
\proseitem{505}{ปญฺจิมานิ ภิกฺขเว อินฺทฺริยานิ.\csromanspacer{} กตมานิ ปญฺจ, สุขินฺทฺริยํ ทุกฺขินฺทฺริยํ โสมนสฺสินฺทฺริยํ โทมนสฺสินฺทฺริยํ อุเปกฺขินฺทฺริยํ.\csromanspacer{} เย หิ เกจิ ภิกฺขเว สมณา วา พฺราหฺมณา วา สุขินฺทฺริยํ นปฺปชานนฺติ, สุขินฺทฺริยสมุทยํ นปฺปชานนฺติ, สุขินฺทฺริยนิโรธํ นปฺปชานนฺติ, สุขินฺทฺริยนิโรธคามินิํ ปฏิปทํ นปฺปชานนฺติ.\csromanspacer{} ทุกฺขินฺทฺริยํ นปฺปชานนฺติ ฯเปฯ.\csromanspacer{} โสมนสฺสินฺทฺริยํ นปฺปชานนฺติ ฯเปฯ.\csromanspacer{} โทมนสฺสินฺทฺริยํ นปฺปชานนฺติ ฯเปฯ.\csromanspacer{} อุเปกฺขินฺทฺริยํ นปฺปชานนฺติ, อุเปกฺขินฺทฺริยสมุทยํ นปฺปชานนฺติ, อุเปกฺขินฺทฺริยนิโรธํ นปฺปชานนฺติ, อุเปกฺขินฺทฺริยนิโรธคามินิํ ปฏิปทํ นปฺปชานนฺติ.\csromanspacer{} น เม เต ภิกฺขเว สมณา วา พฺราหฺมณา วา สมเณสุ วา สมณสมฺมตา, พฺราหฺมเณสุ วา พฺราหฺมณสมฺมตา, น จ ปเนเต อายสฺมนฺโต สามญฺญตฺถํ วา พฺราหฺมญฺญตฺถํ วา ทิฏฺเฐว ธมฺเม สยํ อภิญฺญา สจฺฉิกตฺวา อุปสมฺปชฺช วิหรนฺติ.}

\prose{เย จ โข เกจิ ภิกฺขเว สมณา วา พฺราหฺมณา วา สุขินฺทฺริยํ ปชานนฺติ, สุขินฺทฺริยสมุทยํ ปชานนฺติ, สุขินฺทฺริยนิโรธํ ปชานนฺติ, สุขินฺทฺริยนิโรธคามินิํ ปฏิปทํ ปชานนฺติ.\csromanspacer{} ทุกฺขินฺทฺริยํ ปชานนฺติ ฯเปฯ.\csromanspacer{} โสมนสฺสินฺทฺริยํ ปชานนฺติ ฯเปฯ.\csromanspacer{} โทมนสฺสินฺทฺริยํ ปชานนฺติ ฯเปฯ.\csromanspacer{} อุเปกฺขินฺทฺริยํ ปชานนฺติ, อุเปกฺขินฺทฺริยสมุทยํ ปชานนฺติ, อุเปกฺขินฺทฺริยนิโรธํ ปชานนฺติ, อุเปกฺขินฺทฺริยนิโรธคามินิํ ปฏิปทํ ปชานนฺติ.\csromanspacer{} เต จ โข เม ภิกฺขเว สมณา วา พฺราหฺมณา วา สมเณสุ เจว สมณสมฺมตา, พฺราหฺมเณสุ จ พฺราหฺมณสมฺมตา, เต จ ปนายสฺมนฺโต สามญฺญตฺถญฺจ พฺรหฺมญฺญตฺถญฺจ ทิฏฺเฐว ธมฺเม สยํ อภิญฺญา สจฺฉิกตฺวา อุปสมฺปชฺช วิหรนฺตีติ.\csromanspacer{}\csromanspacer{}ปญฺจมํ.}

\csromanheader{2}{6. ปฐมวิภงฺคสุตฺต}
\proseitem{506}{ปญฺจิมานิ ภิกฺขเว อินฺทฺริยานิ.\csromanspacer{} กตมานิ ปญฺจ, สุขินฺทฺริยํ ทุกฺขินฺทฺริยํ โสมนสฺสินฺทฺริยํ โทมนสฺสินฺทฺริยํ อุเปกฺขินฺทฺริยํ.}

\csromanheader{2}{4. อินฺทฺริยสํยุตฺต}
\prose{\csromanfolio{185}กตมญฺจ ภิกฺขเว สุขินฺทฺริยํ, ยํ โข ภิกฺขเว กายิกํ สุขํ กายิกํ สาตํ กายสมฺผสฺสชํ สุขํ สาตํ เวทยิตํ.\csromanspacer{} อิทํ วุจฺจติ ภิกฺขเว สุขินฺทฺริยํ.}

\prose{กตมญฺจ ภิกฺขเว ทุกฺขินฺทฺริยํ, ยํ โข ภิกฺขเว กายิกํ ทุกฺขํ กายิกํ อสาตํ กายสมฺผสฺสชํ ทุกฺขํ อสาตํ เวทยิตํ.\csromanspacer{} อิทํ วุจฺจติ ภิกฺขเว ทุกฺขินฺทฺริยํ.}

\prose{กตมญฺจ ภิกฺขเว โสมนสฺสินฺทฺริยํ, ยํ โข ภิกฺขเว เจตสิกํ สุขํ เจตสิกํ สาตํ มโนสมฺผสฺสชํ สุขํ สาตํ เวทยิตํ.\csromanspacer{} อิทํ วุจฺจติ ภิกฺขเว โสมนสฺสินฺทฺริยํ.}

\prose{กตมญฺจ ภิกฺขเว โทมนสฺสินฺทฺริยํ, ยํ โข ภิกฺขเว เจตสิกํ ทุกฺขํ เจตสิกํ อสาตํ มโนสมฺผสฺสชํ ทุกฺขํ อสาตํ เวทยิตํ.\csromanspacer{} อิทํ วุจฺจติ ภิกฺขเว โทมนสฺสินฺทฺริยํ.}

\prose{กตมญฺจ ภิกฺขเว อุเปกฺขินฺทฺริยํ, ยํ โข ภิกฺขเว กายิกํ วา เจตสิกํ วา เนวสาตํ นาสาตํ เวทยิตํ.\csromanspacer{} อิทํ วุจฺจติ ภิกฺขเว อุเปกฺขินฺทฺริยํ.\csromanspacer{} อิมานิ โข ภิกฺขเว ปญฺจินฺทฺริยานีติ.\csromanspacer{}\csromanspacer{}ฉฏฺฐํ.}

\csromanheader{2}{7. ทุติยวิภงฺคสุตฺต}
\proseitem{507}{ปญฺจิมานิ ภิกฺขเว อินฺทฺริยานิ.\csromanspacer{} กตมานิ ปญฺจ, สุขินฺทฺริยํ ทุกฺขินฺทฺริยํ โสมนสฺสินฺทฺริยํ โทมนสฺสินฺทฺริยํ อุเปกฺขินฺทฺริยํ.}

\prose{กตมญฺจ ภิกฺขเว สุขินฺทฺริยํ, ยํ โข ภิกฺขเว กายิกํ สุขํ กายิกํ สาตํ กายสมฺผสฺสชํ สุขํ สาตํ เวทยิตํ.\csromanspacer{} อิทํ วุจฺจติ ภิกฺขเว สุขินฺทฺริยํ.}

\prose{กตมญฺจ ภิกฺขเว ทุกฺขินฺทฺริยํ, ยํ โข ภิกฺขเว กายิกํ ทุกฺขํ กายิกํ อสาตํ กายสมฺผสฺสชํ ทุกฺขํ อสาตํ เวทยิตํ.\csromanspacer{} อิทํ วุจฺจติ ภิกฺขเว ทุกฺขินฺทฺริยํ.}

\prose{กตมญฺจ ภิกฺขเว โสมนสฺสินฺทฺริยํ, ยํ โข ภิกฺขเว เจตสิกํ สุขํ เจตสิกํ สาตํ มโนสมฺผสฺสชํ สุขํ สาตํ เวทยิตํ.\csromanspacer{} อิทํ วุจฺจติ ภิกฺขเว โสมนสฺสินฺทฺริยํ.}

\prose{\csromanfolio{186}กตมญฺจ ภิกฺขเว โทมนสฺสินฺทฺริยํ, ยํ โข ภิกฺขเว เจตสิกํ ทุกฺขํ เจตสิกํ อสาตํ มโนสมฺผสฺสชํ ทุกฺขํ อสาตํ เวทยิตํ.\csromanspacer{} อิทํ วุจฺจติ ภิกฺขเว โทมนสฺสินฺทฺริยํ.}

\prose{กตมญฺจ ภิกฺขเว อุเปกฺขินฺทฺริยํ, ยํ โข ภิกฺขเว กายิกํ วา เจตสิกํ วา เนวสาตํ นาสาตํ เวทยิตํ.\csromanspacer{} อิทํ วุจฺจติ ภิกฺขเว อุเปกฺขินฺทฺริยํ.}

\prose{ตตฺร ภิกฺขเว ยญฺจ สุขินฺทฺริยํ ยญฺจ โสมนสฺสินฺทฺริยํ, สุขา สา เวทนา ทฏฺฐพฺพา.\csromanspacer{} ตตฺร ภิกฺขเว ยญฺจ ทุกฺขินฺทฺริยํ ยญฺจ โทมนสฺสินฺทฺริยํ, ทุกฺขา สา เวทนา ทฏฺฐพฺพา.\csromanspacer{} ตตฺร ภิกฺขเว ยทิทํ อุเปกฺขินฺทฺริยํ, อทุกฺขมสุขา สา เวทนา ทฏฺฐพฺพา.\csromanspacer{} อิมานิ โข ภิกฺขเว ปญฺจินฺทฺริยานีติ.\csromanspacer{}\csromanspacer{}สตฺตมํ.}

\csromanheader{2}{8. ตติยวิภงฺคสุตฺต}
\proseitem{508}{ปญฺจิมานิ ภิกฺขเว อินฺทฺริยานิ.\csromanspacer{} กตมานิ ปญฺจ, สุขินฺทฺริยํ ทุกฺขินฺทฺริยํ โสมนสฺสินฺทฺริยํ โทมนสฺสินฺทฺริยํ อุเปกฺขินฺทฺริยํ.}

\prose{กตมญฺจ ภิกฺขเว สุขินฺทฺริยํ, ยํ โข ภิกฺขเว กายิกํ สุขํ กายิกํ สาตํ กายสมฺผสฺสชํ สุขํ สาตํ เวทยิตํ.\csromanspacer{} อิทํ วุจฺจติ ภิกฺขเว สุขินฺทฺริยํ.}

\prose{กตมญฺจ ภิกฺขเว ทุกฺขินฺทฺริยํ, ยํ โข ภิกฺขเว กายิกํ ทุกฺขํ กายิกํ อสาตํ กายสมฺผสฺสชํ ทุกฺขํ อสาตํ เวทยิตํ.\csromanspacer{} อิทํ วุจฺจติ ภิกฺขเว ทุกฺขินฺทฺริยํ.}

\prose{กตมญฺจ ภิกฺขเว โสมนสฺสินฺทฺริยํ, ยํ โข ภิกฺขเว เจตสิกํ สุขํ เจตสิกํ สาตํ มโนสมฺผสฺสชํ สุขํ สาตํ เวทยิตํ.\csromanspacer{} อิทํ วุจฺจติ ภิกฺขเว โสมนสฺสินฺทฺริยํ.}

\prose{กตมญฺจ ภิกฺขเว โทมนสฺสินฺทฺริยํ, ยํ โข ภิกฺขเว เจตสิกํ ทุกฺขํ เจตสิกํ อสาตํ มโนสมฺผสฺสชํ ทุกฺขํ อสาตํ เวทยิตํ.\csromanspacer{} อิทํ วุจฺจติ ภิกฺขเว โทมนสฺสินฺทฺริยํ.}

\prose{กตมญฺจ ภิกฺขเว อุเปกฺขินฺทฺริยํ, ยํ โข ภิกฺขเว กายิกํ วา เจตสิกํ วา เนวสาตํ นาสาตํ เวทยิตํ.\csromanspacer{} อิทํ วุจฺจติ ภิกฺขเว อุเปกฺขินฺทฺริยํ.}

\prose{ตตฺร ภิกฺขเว ยญฺจ สุขินฺทฺริยํ ยญฺจ โสมนสฺสินฺทฺริยํ, สุขา สา เวทนา ทฏฺฐพฺพา.\csromanspacer{} ตตฺร ภิกฺขเว ยญฺจ ทุกฺขินฺทฺริยํ ยญฺจ โทมนสฺสินฺทฺริยํ, ทุกฺขา สา เวทนา ทฏฺฐพฺพา.\csromanspacer{} ตตฺร ภิกฺขเว ยทิทํ อุเปกฺขินฺทฺริยํ.\csromanspacer{} อทุกฺขมสุขา สา เวทนา}

\csromanheader{2}{4. อินฺทฺริยสํยุตฺต}
\prose{\csromanfolio{187}ทฏฺฐพฺพา.\csromanspacer{} อิติ โข ภิกฺขเว อิมานิ ปญฺจินฺทฺริยานิ ปญฺจ หุตฺวา ตีณิ โหนฺติ, ตีณิ หุตฺวา ปญฺจ โหนฺติ ปริยาเยนาติ.\csromanspacer{}\csromanspacer{}อฏฺฐมํ.}

\csromanmatikaanchor{mtk.115}
\csromantocmarkref{subparagraph}{9. กฏฺโฐปมสุตฺต}{mtk.115}
\bookmark[dest={mtk.115},level=5]{9. กฏฺโฐปมสุตฺต}
\csromanheader{6}{9. กฏฺโฐปมสุตฺต}
\proseitem{509}{ปญฺจิมานิ ภิกฺขเว อินฺทฺริยานิ.\csromanspacer{} กตมานิ ปญฺจ, สุขินฺทฺริยํ ทุกฺขินฺทฺริยํ โสมนสฺสินฺทฺริยํ โทมนสฺสินฺทฺริยํ อุเปกฺขินฺทฺริยํ.\csromanspacer{} สุขเวทนิยํ ภิกฺขเว ผสฺสํ ปฏิจฺจ อุปฺปชฺชติ สุขินฺทฺริยํ, โส สุขิโตว สมาโน “สุขิโตสฺมี”ติ ปชานาติ, ตสฺเสว สุขเวทนิยสฺส ผสฺสสฺส นิโรธา ยํ ตชฺชํ เวทยิตํ สุขเวทนิยํ ผสฺสํ ปฏิจฺจ อุปฺปนฺนํ สุขินฺทฺริยํ, ตํ นิรุชฺฌติ, ตํ วูปสมฺมตีติ ปชานาติ.}

\prose{ทุกฺขเวทนิยํ ภิกฺขเว ผสฺสํ ปฏิจฺจ อุปฺปชฺชติ ทุกฺขินฺทฺริยํ, โส ทุกฺขิโตว สมาโน “ทุกฺขิโตสฺมี”ติ ปชานาติ, ตสฺเสว ทุกฺขเวทนิยสฺส ผสฺสสฺส นิโรธา ยํ ตชฺชํ เวทยิตํ ทุกฺขเวทนิยํ ผสฺสํ ปฏิจฺจ อุปฺปนฺนํ ทุกฺขินฺทฺริยํ, ตํ นิรุชฺฌติ, ตํ วูปสมฺมตีติ ปชานาติ.}

\prose{โสมนสฺสเวทนิยํ ภิกฺขเว ผสฺสํ ปฏิจฺจ อุปฺปชฺชติ โสมนสฺสินฺทฺริยํ, โส สุมโนว สมาโน “สุมโนสฺมี”ติ ปชานาติ, ตสฺเสว โสมนสฺสเวทนิยสฺส ผสฺสสฺส นิโรธา ยํ ตชฺชํ เวทยิตํ โสมนสฺสเวทนิยํ ผสฺสํ ปฏิจฺจ อุปฺปนฺนํ โสมนสฺสินฺทฺริยํ, ตํ นิรุชฺฌติ, ตํ วูปสมฺมตีติ ปชานาติ.}

\prose{โทมนสฺสเวทนิยํ ภิกฺขเว ผสฺสํ ปฏิจฺจ อุปฺปชฺชติ โทมนสฺสินฺทฺริยํ, โส ทุมฺมโนว สมาโน “ทุมฺมโนสฺมี”ติ ปชานาติ, ตสฺเสว โทมนสฺสเวทนิยสฺส ผสฺสสฺสฺส นิโรธา ยํ ตชฺชํ เวทยิตํ โทมนสฺสเวทนิยํ ผสฺสํ ปฏิจฺจ อุปฺปนฺนํ โทมนสฺสินฺทฺริยํ, ตํ นิรุชฺฌติ, ตํ วูปสมฺมตีติ ปชานาติ.}

\prose{อุเปกฺขาเวทนิยํ ภิกฺขเว ผสฺสํ ปฏิจฺจ อุปฺปชฺชติ อุเปกฺขินฺทฺริยํ, โส อุเปกฺขโกว สมาโน “อุเปกฺขโกสฺมี”ติ ปชานาติ, ตสฺเสว อุเปกฺขาเวทนิยสฺส ผสฺสสฺส นิโรธา ยํ ตชฺชํ เวทยิตํ อุเปกฺขาเวทนิยํ ผสฺสํ ปฏิจฺจ อุปฺปนฺนํ อุเปกฺขินฺทฺริยํ, ตํ นิรุชฺฌติ, ตํ วูปสมฺมตีติ ปชานาติ.}

\prose{เสยฺยถาปิ ภิกฺขเว ทฺวินฺนํ กฏฺฐานํ สํฆฏฺฏนสโมธานา\footnote{สํฆฏฺฏนา สโมธานา (อิ, ก), สํฆฏนสโมธานา (สฺยา, กํ)} อุสฺมา ชายติ, เตโช อภินิพฺพตฺตติ, เตสํเยว กฏฺฐานํ นานาภาวาวินิกฺเขปา\footnote{นานาภาวนิกฺเขปา (สฺยา, กํ, อิ, ก)} \csromanfolio{188}ยา ตชฺชา อุสฺมา, สา นิรุชฺฌติ, สา วูปสมฺมติ.\csromanspacer{} เอวเมว โข ภิกฺขเว สุขเวทนิยํ ผสฺสํ ปฏิจฺจ อุปฺปชฺชติ สุขินฺทฺริยํ, โส สุขิโตว สมาโน “สุขิโตสฺมี”ติ ปชานาติ, ตสฺเสว สุขเวทนิยสฺส ผสฺสสฺส นิโรธา ยํ ตชฺชํ เวทยิตํ สุขเวทนิยํ ผสฺสํ ปฏิจฺจ อุปฺปชฺชติ สุขินฺทฺริยํ, ตํ นิรุชฺฌติ, ตํ วูปสมฺมตีติ ปชานาติ.\csromanspacer{} ทุกฺขเวทนิยํ ภิกฺขเว ผสฺสํ ปฏิจฺจ ฯเปฯ.\csromanspacer{} โสมนสฺสเวทนิยํ ภิกฺขเว ผสฺสํ ปฏิจฺจ ฯเปฯ.\csromanspacer{} โทมนสฺสเวทนิยํ ภิกฺขเว ผสฺสํ ปฏิจฺจ ฯเปฯ.\csromanspacer{} อุเปกฺขาเวทนิยํ ภิกฺขเว ผสฺสํ ปฏิจฺจ อุปฺปชฺชติ อุเปกฺขินฺทฺริยํ, โส อุเปกฺขโกว สมาโน “อุเปกฺขโกสฺมี”ติ ปชานาติ, ตสฺเสว อุเปกฺขาเวทนิยสฺส ผสฺสสฺส นิโรธา ยํ ตชฺชํ เวทยิตํ อุเปกฺขาเวทนิยํ ผสฺสํ ปฏิจฺจ อุปฺปชฺชติ อุเปกฺขินฺทฺริยํ.\csromanspacer{} ตํ นิรุชฺฌติ, ตํ วูปสมฺมตีติ ปชานาติ.\csromanspacer{}\csromanspacer{}นวมํ.}

\csromanmatikaanchor{mtk.116}
\csromantocmarkref{subparagraph}{10. อุปฺปฏิปาฏิกสุตฺต}{mtk.116}
\bookmark[dest={mtk.116},level=5]{10. อุปฺปฏิปาฏิกสุตฺต}
\csromanheader{6}{10. อุปฺปฏิปาฏิกสุตฺต}
\proseitem{510}{ปญฺจิมานิ ภิกฺขเว อินฺทฺริยานิ.\csromanspacer{} กตมานิ ปญฺจ, ทุกฺขินฺทฺริยํ โทมนสฺสินฺทฺริยํ สุขินฺทฺริยํ โสมนสฺสินฺทฺริยํ อุเปกฺขินฺทฺริยํ.\csromanspacer{} อิธ ภิกฺขเว ภิกฺขุโน อปฺปมตฺตสฺส อาตาปิโน ปหิตตฺตสฺส วิหรโต อุปฺปชฺชติ ทุกฺขินฺทฺริยํ, โส เอวํ ปชานาติ “อุปฺปนฺนํ โข เม อิทํ ทุกฺขินฺทฺริยํ, ตญฺจ โข สนิมิตฺตํ สนิทานํ สสงฺขารํ สปฺปจฺจยํ, ตญฺจ อนิมิตฺตํ อนิทานํ อสงฺขารํ อปฺปจฺจยํ ทุกฺขินฺทฺริยํ อุปฺปชฺชิสฺสตี”ติ เนตํ ฐานํ วิชฺชติ.\csromanspacer{} โส ทุกฺขินฺทฺริยญฺจ ปชานาติ, ทุกฺขินฺทฺริยสมุทยญฺจ ปชานาติ, ทุกฺขินฺทฺริยนิโรธญฺจ ปชานาติ.\csromanspacer{} ยตฺถ จุปฺปนฺนํ ทุกฺขินฺทฺริยํ อปริเสสํ นิรุชฺฌติ, ตญฺจ ปชานาติ.\csromanspacer{} กตฺถ จุปฺปนฺนํ ทุกฺขินฺทฺริยํ อปริเสสํ นิรุชฺฌติ, อิธ ภิกฺขเว ภิกฺขุ วิวิจฺเจว กาเมหิ วิวิจฺจ อกุสเลหิ ธมฺเมหิ สวิตกฺกํ สวิจารํ วิเวกชํ ปีติสุขํ ปฐมํ ฌานํ อุปสมฺปชฺช วิหรติ, เอตฺถ จุปฺปนฺนํ ทุกฺขินฺทฺริยํ อปริเสสํ นิรุชฺฌติ.\csromanspacer{} อยํ วุจฺจติ ภิกฺขเว ภิกฺขุ อญฺญาสิ ทุกฺขินฺทฺริยสฺส นิโรธํ, ตทตฺถาย จิตฺตํ อุปสํหรติ.}

\prose{อิธ ปน ภิกฺขเว ภิกฺขุโน อปฺปมตฺตสฺส อาตาปิโน ปหิตตฺตสฺส วิหรโต อุปฺปชฺชติ โทมนสฺสินฺทฺริยํ, โส เอวํ ปชานาติ “อุปฺปนฺนํ โข เม อิทํ โทมนสฺสินฺทฺริยํ, ตญฺจ โข สนิมิตฺตํ สนิทานํ สสงฺขารํ สปฺปจฺจยํ, ตญฺจ อนิมิตฺตํ อนิทานํ อสงฺขารํ อปฺปจฺจยํ โทมนสฺสินฺทฺริยํ อุปฺปชฺชิสฺสตี”ติ เนตํ ฐานํ วิชฺชติ.\csromanspacer{} โส โทมนสฺสินฺทฺริยญฺจ ปชานาติ, โทมนสฺสินฺทฺริยสมุทยญฺจ ปชานาติ, โทมนสฺสินฺทฺริยนิโรธญฺจ ปชานาติ.\csromanspacer{} ยตฺถ จุปฺปนฺนํ โทมนสฺสินฺทฺริยํ}

\csromanheader{2}{4. อินฺทฺริยสํยุตฺต}
\prose{\csromanfolio{189}อปริเสสํ นิรุชฺฌติ, ตญฺจ ปชานาติ.\csromanspacer{} กตฺถ จุปฺปนฺนํ โทมนสฺสินฺทฺริยํ อปริเสสํ นิรุชฺฌติ, อิธ ภิกฺขเว ภิกฺขุ วิตกฺกวิจารานํ วูปสมา อชฺฌตฺตํ สมฺปสาทนํ เจตโส เอโกทิภาวํ อวิตกฺกํ อวิจารํ สมาธิชํ ปีติสุขํ ทุติยํ ฌานํ อุปสมฺปชฺช วิหรติ, เอตฺถ จุปฺปนฺนํ โทมนสฺสินฺทฺริยํ อปริเสสํ นิรุชฺฌติ.\csromanspacer{} อยํ วุจฺจติ ภิกฺขเว ภิกฺขุ อญฺญาสิ โทมนสฺสินฺทฺริยสฺส นิโรธํ, ตทตฺถาย จิตฺตํ อุปสํหรติ.}

\prose{อิธ ปน ภิกฺขเว ภิกฺขุโน อปฺปมตฺตสฺส อาตาปิโน ปหิตตฺตสฺส วิหรโต อุปฺปชฺชติ สุขินฺทฺริยํ, โส เอวํ ปชานาติ “อุปฺปนฺนํ โข เม อิทํ สุขินฺทฺริยํ, ตญฺจ โข สนิมิตฺตํ สนิทานํ สสงฺขารํ สปฺปจฺจยํ, ตญฺจ อนิมิตฺตํ อนิทานํ อสงฺขารํ อปฺปจฺจยํ สุขินฺทฺริยํ อุปฺปชฺชิสฺสตี”ติ เนตํ ฐานํ วิชฺชติ.\csromanspacer{} โส สุขินฺทฺริยญฺจ ปชานาติ สุขินฺทฺริยสมุทยญฺจ ปชานาติ, สุขินฺทฺริยนิโรธญฺจ ปชานาติ.\csromanspacer{} ยตฺถ จุปฺปนฺนํ สุขินฺทฺริยํ อปริเสสํ นิรุชฺฌติ, ตญฺจ ปชานาติ.\csromanspacer{} กตฺถ จุปฺปนฺนํ สุขินฺทฺริยํ อปริเสสํ นิรุชฺฌติ, อิธ ภิกฺขเว ภิกฺขุ ปีติยา จ วิราคา อุเปกฺขโก จ วิหรติ สโต จ สมฺปชาโน สุขญฺจ กาเยน ปฏิสํเวเทติ.\csromanspacer{} ยํ ตํ อริยา อาจิกฺขนฺติ “อุเปกฺขโก สติมา สุขวิหารี”ติ, ตติยํ ฌานํ อุปสมฺปชฺช วิหรติ, เอตฺถ จุปฺปนฺนํ สุขินฺทฺริยํ อปริเสสํ นิรุชฺฌติ.\csromanspacer{} อยํ วุจฺจติ ภิกฺขเว ภิกฺขุ อญฺญาสิ สุขินฺทฺริยสฺส นิโรธํ, ตทตฺถาย จิตฺตํ อุปสํหรติ.}

\prose{อิธ ปน ภิกฺขเว ภิกฺขุโน อปฺปมตฺตสฺส อาตาปิโน ปหิตตฺตสฺส วิหรโต อุปฺปชฺชติ โสมนสฺสินฺทฺริยํ, โส เอวํ ปชานาติ “อุปฺปนฺนํ โข เม อิทํ โสมนสฺสินฺทฺริยํ, ตญฺจ โข สนิมิตฺตํ สนิทานํ สสงฺขารํ สปฺปจฺจยํ, ตญฺจ อนิมิตฺตํ อนิทานํ อสงฺขารํ อปฺปจฺจยํ โสมนสฺสินฺทฺริยํ อุปฺปชฺชิสฺสตี”ติ เนตํ ฐานํ วิชฺชติ.\csromanspacer{} โส โสมนสฺสินฺทฺริยญฺจ ปชานาติ, โสมนสฺสินฺทฺริยสมุทยญฺจ ปชานาติ, โสมนสฺสินฺทิยนิโรธญฺจ ปชานาติ.\csromanspacer{} ยตฺถ จุปฺปนฺนํ โสมนสฺสินฺทฺริยํ อปริเสสํ นิรุชฺฌติ, ตญฺจ ปชานาติ.\csromanspacer{} กตฺถ จุปฺปนฺนํ โสมนสฺสินฺทฺริยํ อปริเสสํ นิรุชฺฌติ, อิธ ภิกฺขเว ภิกฺขุ สุขสฺส จ ปหานา ทุกฺขสฺส จ ปหานา ปุพฺเพว โสมนสฺสโทมนสฺสานํ อตฺถงฺคมา อทุกฺขมสุขํ อุเปกฺขาสติปาริสุทฺธิํ จตุตฺถํ ฌานํ อุปสมฺปชฺช วิหรติ, เอตฺถ จุปฺปนฺนํ โสมนสฺสินฺทฺริยํ อปริเสสํ นิรุชฺฌติ.\csromanspacer{} อยํ วุจฺจติ ภิกฺขเว ภิกฺขุ อญฺญาสิ โสมนสฺสินฺทฺริยสฺส นิโรธํ, ตทตฺถาย จิตฺตํ อุปสํหรติ.}

\prose{\csromanfolio{190}อิธ ปน ภิกฺขเว ภิกฺขุโน อปฺปมตฺตสฺส อาตาปิโน ปหิตตฺตสฺส วิหรโต อุปฺปชฺชติ อุเปกฺขินฺทฺริยํ, โส เอวํ ปชานาติ “อุปฺปนฺนํ โข เม อิทํ อุเปกฺขินฺทฺริยํ, ตญฺจ โข สนิมิตฺตํ สนิทานํ สสงฺขารํ สปฺปจฺจยํ, ตญฺจ อนิมิตฺตํ อนิทานํ อสงฺขารํ อปฺปจฺจยํ อุเปกฺขินฺทฺริยํ อุปฺปชฺชิสฺสตี”ติ เนตํ ฐานํ วิชฺชติ.\csromanspacer{} โส อุเปกฺขินฺทฺริยญฺจ ปชานาติ, อุเปกฺขินฺทฺริยสมุทยญฺจ ปชานาติ, อุเปกฺขินฺทฺริยนิโรธญฺจ ปชานาติ.\csromanspacer{} ยตฺถ จุปฺปนฺนํ อุเปกฺขินฺทฺริยํ อปริเสสํ นิรุชฺฌติ, ตญฺจ ปชานาติ.\csromanspacer{} กตฺถ จุปฺปนฺนํ อุเปกฺขินฺทฺริยํ อปริเสสํ นิรุชฺฌติ, อิธ ภิกฺขเว ภิกฺขุ สพฺพโส เนวสญฺญานาสญฺญายตนํ สมติกฺกมฺม สญฺญาเวทยิตนิโรธํ อุปสมฺปชฺช วิหรติ, เอตฺถ จุปฺปนฺนํ อุเปกฺขินฺทฺริยํ อปริเสสํ นิรุชฺฌติ.\csromanspacer{} อยํ วุจฺจติ ภิกฺขเว ภิกฺขุ อญฺญาสิ อุเปกฺขินฺทฺริยสฺส นิโรธํ, ตทตฺถาย จิตฺตํ อุปสํหรตีติ.\csromanspacer{}\csromanspacer{}ทสมํ.\csromanspacer{} สุขินฺทฺริยวคฺโค จตุตฺโถ.}

\titlehead{\textbf{ตสฺสุทฺทานํ}}
\csromangathameasure{สุทฺธิกญฺจ โสโต อรหา, ทุเว สมณพฺราหฺมณา. \\ วิภงฺเคน ตโย วุตฺตา, กฏฺโฐ อุปฺปฏิปาฏิกนฺติ.}
\csromangathagroup{\csromangathastanza{สุทฺธิกญฺจ โสโต อรหา, ทุเว สมณพฺราหฺมณา. \\ วิภงฺเคน ตโย วุตฺตา, กฏฺโฐ อุปฺปฏิปาฏิกนฺติ.}}

\csromanmatikaanchor{mtk.117}
\csromantocmarknopage{subparagraph}{5. ชราวคฺค}{mtk.117}
\bookmark[dest={mtk.117},level=5]{5. ชราวคฺค}
\csromanheader{6}{5. ชราวคฺค}
\csromanmatikaanchor{mtk.118}
\csromantocmarkref{subparagraph}{1. ชราธมฺมสุตฺต}{mtk.118}
\bookmark[dest={mtk.118},level=5]{1. ชราธมฺมสุตฺต}
\csromanheader{6}{1. ชราธมฺมสุตฺต}
\proseitem{511}{เอวํ เม สุตํ\csromandash{}เอกํ สมยํ ภควา สาวตฺถิยํ วิหรติ ปุพฺพาราเม มิคารมาตุปาสาเท.\csromanspacer{} เตน โข ปน สมเยน ภควา สายนฺหสมยํ ปฏิสลฺลานา วุฏฺฐิโต ปจฺฉาตเป นิสินฺโน โหติ ปิฏฺฐิํ โอตาปยมาโน.}

\prose{อถ โข อายสฺมา อานนฺโท เยน ภควา เตนุปสงฺกมิ, อุปสงฺกมิตฺวา ภควนฺตํ อภิวาเทตฺวา ภควโต คตฺตานิ ปาณินา อโนมชฺชนฺโต ภควนฺตํ เอตทโวจ “อจฺฉริยํ ภนฺเต, อพฺภุตํ ภนฺเต, น เจวํ ทานิ ภนฺเต ภควโต ตาว ปริสุทฺโธ ฉวิวณฺโณ ปริโยทาโต, สิถิลานิ จ คตฺตานิ สพฺพานิ วลิยชาตานิ, ปุรโต}

\csromanheader{2}{4. อินฺทฺริยสํยุตฺต}
\prose{\csromanfolio{191}ปพฺภาโร จ กาโย, ทิสฺสติ จ อินฺทฺริยานํ อญฺญถตฺตํ จกฺขุนฺทฺริยสฺส โสตินฺทฺริยสฺส ฆานินฺทฺริยสฺส ชิวฺหินฺทฺริยสฺส กายินฺทฺริยสฺสา”ติ.}

\prose{เอวํ เหตํ อานนฺท โหติ ชราธมฺโม โยพฺพญฺเญ, พฺยาธิธมฺโม อาโรเคฺย, มรณธมฺโม ชีวิเต, น เจว ตาว ปริสุทฺโธ โหติ ฉวิวณฺโณ ปริโยทาโต, สิถิลานิ จ โหนฺติ คตฺตานิ สพฺพานิ วลิยชาตานิ, ปุรโต ปพฺภาโร จ กาโย, ทิสฺสติ จ อินฺทฺริยานํ อญฺญถตฺตํ จกฺขุนฺทฺริยสฺส โสตินฺทฺริยสฺส ฆานินฺทฺริยสฺส ชิวฺหินฺทฺริยสฺส กายินฺทฺริยสฺสาติ.}

\prose{อิทมโวจ ภควา, อิทํ วตฺวา จ พุคโต อถาปรํ เอตทโวจ สตฺถา\csromandash{}}

\csromangathameasure{“ธี ตํ ชมฺมิ ชเร อตฺถุ, ทุพฺพณฺณกรณี ชเร. \\ ตาว มโนรมํ พิมฺพํ, ชราย อภิมทฺทิตํ. \\ โยปิ วสฺสสตํ ชีเว, โสปิ มจฺจุปรายโณ. \\ น กิญฺจิ ปริวชฺเชติ, สพฺพเมวาภิมทฺทตี”ติ. ปฐมํ.}
\csromangathagroup{\csromangathastanza{“ธี ตํ ชมฺมิ ชเร อตฺถุ, ทุพฺพณฺณกรณี ชเร. \\ ตาว มโนรมํ พิมฺพํ, ชราย อภิมทฺทิตํ.}\csromangathastanza{โยปิ วสฺสสตํ ชีเว, โสปิ มจฺจุปรายโณ\footnote{สพฺเพ มจฺจุปรายนา (สฺยา, กํ, ก)}. \\ น กิญฺจิ ปริวชฺเชติ, สพฺพเมวาภิมทฺทตี”ติ.\csromanspacer{} ปฐมํ.}}

\csromanmatikaanchor{mtk.119}
\csromantocmarkref{subparagraph}{2. อุณฺณาภพฺราหฺมณสุตฺต}{mtk.119}
\bookmark[dest={mtk.119},level=5]{2. อุณฺณาภพฺราหฺมณสุตฺต}
\csromanheader{6}{2. อุณฺณาภพฺราหฺมณสุตฺต}
\proseitem{512}{สาวตฺถินิทานํ.\csromanspacer{} อถ โข อุณฺณาโภ พฺราหฺมโณ เยน ภควา เตนุปสงฺกมิ, อุปสงฺกมิตฺวา ภควตา สทฺธิํ สมฺโมทิ, สมฺโมทนียํ กถํ สารณียํ วีติสาเรตฺวา เอกมนฺตํ นิสีทิ, เอกมนฺตํ นิสินฺโน โข อุณฺณาโภ พฺราหฺมโณ ภควนฺตํ เอตทโวจ\csromandash{}}

\prose{ปญฺจิมานิ โภ โคตม อินฺทฺริยานิ นานาวิสยานิ นานาโคจรานิ น อญฺญมญฺญสฺส โคจรวิสยํ ปจฺจนุโภนฺติ.\csromanspacer{} กตมานิ ปญฺจ, จกฺขุนฺทฺริยํ โสตินฺทฺริยํ ฆานินฺทฺริยํ ชิวฺหินฺทฺริยํ กายินฺทฺริยํ.\csromanspacer{} อิเมสํ นุ โข โภ โคตม ปญฺจนฺนํ อินฺทฺริยานํ นานาวิสยานํ นานาโคจรานํ น อญฺญมญฺญสฺส โคจรวิสยํ ปจฺจนุโภนฺตานํ กิํ ปฏิสรณํ, โก จ เนสํ โคจรวิสยํ ปจฺจนุโภตีติ.}

\prose{ปญฺจิมานิ พฺราหฺมณ อินฺทฺริยานิ นานาวิสยานิ นานาโคจรานิ น อญฺญมญฺญสฺส โคจรวิสยํ ปจฺจนุโภนฺติ.\csromanspacer{} กตมานิ ปญฺจ, จกฺขุนฺทฺริยํ \csromanfolio{192}โสตินฺทฺริยํ ฆานินฺทฺริยํ ชิวฺหินฺทฺริยํ กายินฺทฺริยํ.\csromanspacer{} อิเมสํ โข พฺราหฺมณ ปญฺจนฺนํ อินฺทฺริยานํ นานาวิสยานํ นานาโคจรานํ น อญฺญมญฺญสฺส โคจรวิสยํ ปจฺจนุโภนฺตานํ มโน ปฏิสรณํ, มโนว เนสํ โคจรวิสยํ ปจฺจนุโภตีติ.}

\prose{มนสฺส ปน โภ โคตม กิํ ปฏิสรณนฺติ, มนสฺส โข พฺราหฺมณ สติ ปฏิสรณนฺติ.\csromanspacer{} สติยา ปน โภ โคตม กิํ ปฏิสรณนฺติ, สติยา โข พฺราหฺมณ วิมุตฺติ ปฏิสรณนฺติ.\csromanspacer{} วิมุตฺติยา ปน โภ โคตม กิํ ปฏิสรณนฺติ, วิมุตฺติยา โข พฺราหฺมณ นิพฺพานํ ปฏิสรณนฺติ.\csromanspacer{} นิพฺพานสฺส ปน โภ โคตม กิํ ปฏิสรณนฺติ, อจฺจยาสิ\footnote{อจฺจสรา (สี, สฺยา, กํ), อชฺฌปรํ (อิ, ก)} พฺราหฺมณ ปญฺหํ นาสกฺขิ ปญฺหสฺส ปริยนฺตํ คเหตุํ, นิพฺพาโนคธํ หิ พฺราหฺมณ พฺรหฺมจริยํ วุสฺสติ นิพฺพานปรายณํ นิพฺพานปริโยสานนฺติ.}

\prose{อถ โข อุณฺณาโภ พฺราหฺมโณ ภควโต ภาสิตํ อภินนฺทิตฺวา อนุโมทิตฺวา อุฏฺฐายาสนา ภควนฺตํ อภิวาเทตฺวา ปทกฺขิณํ กตฺวา ปกฺกามิ.}

\prose{อถ โข ภควา อจิรปกฺกนฺเต อุณฺณาเภ พฺราหฺมเณ ภิกฺขู อามนฺเตสิ “เสยฺยถาปิ ภิกฺขเว กูฏาคาเร วา กูฏาคารสาลายํ วา\footnote{กูฏาคารํ วา กูฏาคารสาลํ วา อุตฺตราย (ก-สี)} ปาจีนวาตปานา สูริเย อุคฺคจฺฉนฺเต วาตปาเนน รสฺมิ\footnote{รสฺมิโย (สฺยา, ก)} ปวิสิตฺวา กฺวาสฺส\footnote{กาย (สฺยา, ก)} ปติฏฺฐิตา”ติ.\csromanspacer{} ปจฺฉิมายํ ภนฺเต ภิตฺติยนฺติ.\csromanspacer{} เอวเมว โข ภิกฺขเว อุณฺณาภสฺส พฺราหฺมณสฺส ตถาคเต สทฺธา นิวิฏฺฐา, มูลชาตา ปติฏฺฐิตา, ทฬฺหา อสํหาริยา สมเณน วา พฺราหฺมเณน วา เทเวน วา มาเรน วา พฺรหฺมุนา วา เกนจิ วา โลกสฺมิํ.\csromanspacer{} อิมมฺหิ เจ ภิกฺขเว สมเย อุณฺณาโภ พฺราหฺมโณ กาลํ กเรยฺย, นตฺถิ ตํ สํโยชนํ, เยน สํโยชเนน สํยุตฺโต อุณฺณาโภ พฺราหฺมโณ ปุน อิมํ โลกํ อาคจฺเฉยฺยาติ.\csromanspacer{}\csromanspacer{}ทุติยํ.}

\csromanmatikaanchor{mtk.120}
\csromantocmarkref{subparagraph}{3. สาเกตสุตฺต}{mtk.120}
\bookmark[dest={mtk.120},level=5]{3. สาเกตสุตฺต}
\csromanheader{6}{3. สาเกตสุตฺต}
\proseitem{513}{เอวํ เม สุตํ\csromandash{}เอกํ สมยํ ภควา สาเกเต วิหรติ อญฺชนวเน มิคทาเย.\csromanspacer{} ตตฺร โข ภควา ภิกฺขู อามนฺเตสิ “อตฺถิ นุ โข}

\csromanheader{2}{4. อินฺทฺริยสํยุตฺต}
\prose{\csromanfolio{193}ภิกฺขเว ปริยาโย, ยํ ปริยายํ อาคมฺม ยานิ ปญฺจินฺทฺริยานิ, ตานิ ปญฺจ พลานิ โหนฺติ.\csromanspacer{} ยานิ ปญฺจ พลานิ, ตานิ ปญฺจินฺทฺริยานิ โหนฺตี”ติ.}

\prose{ภควํมูลกา โน ภนฺเต ธมฺมา ภควํเนตฺติกา ภควํปฏิสรณา, สาธุ วต ภนฺเต, ภควนฺตํเยว ปฏิภาตุ เอตสฺส ภาสิตสฺส อตฺโถ, ภควโต สุตฺวา ภิกฺขู ธาเรสฺสนฺตีติ.\csromanspacer{} อตฺถิ ภิกฺขเว ปริยาโย, ยํ ปริยายํ อาคมฺม ยานิ ปญฺจินฺทฺริยานิ, ตานิ ปญฺจ พลานิ, โหนฺติ.\csromanspacer{} ยานิ ปญฺจ พลานิ, ตานิ ปญฺจินฺทฺริยานิ โหนฺติ.}

\prose{กตโม จ ภิกฺขเว ปริยาโย, ยํ ปริยายํ อาคมฺม ยานิ ปญฺจินฺทฺริยานิ, ตานิ ปญฺจ พลานิ โหนฺติ.\csromanspacer{} ยานิ ปญฺจ พลานิ, ตานิ ปญฺจินฺทฺริยานิ โหนฺติ.\csromanspacer{} ยํ ภิกฺขเว สทฺธินฺทฺริยํ, ตํ สทฺธาพลํ, ยํ สทฺธาพลํ, ตํ สทฺธินฺทฺริยํ, ยํ วีริยินฺทฺริยํ, ตํ วีริยพลํ, ยํ วีริยพลํ, ตํ วีริยินฺทฺริยํ, ยํ สตินฺทฺริยํ, ตํ สติพลํ, ยํ สติพลํ, ตํ สตินฺทฺริยํ, ยํ สมาธินฺทฺริยํ, ตํ สมาธิพลํ, ยํ สมาธิพลํ, ตํ สมาธินฺทฺริยํ, ยํ ปญฺญินฺทฺริยํ, ตํ ปญฺญาพลํ, ยํ ปญฺญาพลํ, ตํ ปญฺญินฺทฺริยํ.\csromanspacer{} เสยฺยถาปิ ภิกฺขเว นที ปาจีนนินฺนา ปาจีนโปณา ปาจีนปพฺภารา, ตสฺส มชฺเฌ ทีโป, อตฺถิ ภิกฺขเว ปริยาโย, ยํ ปริยายํ อาคมฺม ตสฺสา นทิยา เอโก โสโตเตฺวว สงฺขฺยํ\footnote{สงฺขํ (สี, สฺยา, กํ)} คจฺฉติ, อตฺถิ ปน ภิกฺขเว ปริยาโย, ยํ ปริยายํ อาคมฺม ตสฺสา นทิยา เทฺว โสตานิเตฺวว สงฺขฺยํ คจฺฉติ.}

\prose{กตโม จ ภิกฺขเว ปริยาโย, ยํ ปริยายํ อาคมฺม ตสฺสา นทิยา เอโก โสโตเตฺวว สงฺขฺยํ คจฺฉติ.\csromanspacer{} ยญฺจ ภิกฺขเว ตสฺส ทีปสฺส ปุริมนฺเต\footnote{ปุรตฺถิมนฺเต (สี, สฺยา, กํ, อิ)} อุทกํ, ยญฺจ ปจฺฉิมนฺเต อุทกํ, อยํ โข ภิกฺขเว ปริยาโย, ยํ ปริยายํ อาคมฺม ตสฺสา นทิยา เอโก โสโตเตฺวว สงฺขฺยํ คจฺฉติ.}

\prose{กตโม จ ภิกฺขเว ปริยาโย, ยํ ปริยายํ อาคมฺม ตสฺสา นทิยา เทฺว โสตานิเตฺวว สงฺขฺยํ คจฺฉนฺติ.\csromanspacer{} ยญฺจ ภิกฺขเว ตสฺส ทีปสฺส อุตฺตรนฺเต อุทกํ, ยญฺจ ทกฺขิณนฺเต อุทกํ, อยํ โข ภิกฺขเว ปริยาโย, ยํ ปริยายํ อาคมฺม ตสฺสา นทิยา เทฺว โสตานิเตฺวว สงฺขฺยํ คจฺฉนฺติ.\csromanspacer{} เอวเมว โข ภิกฺขเว ยํ สทฺธินฺทฺริยํ, ตํ สทฺธาพลํ, ยํ สทฺธาพลํ, ตํ สทฺธินฺทฺริยํ, ยํ วีริยินฺทฺริยํ, ตํ วีริยพลํ, ยํ วีริยพลํ, ตํ วีริยินฺทฺริยํ, ยํ สตินฺทฺริยํ, ตํ}

\prose{\csromanfolio{194}สติพลํ, ยํ สติพลํ, ตํ สตินฺทฺริยํ, ยํ สมาธินฺทฺริยํ, ตํ สมาธิพลํ, ยํ สมาธิพลํ, ตํ สมาธินฺทฺริยํ, ยํ ปญฺญินฺทฺริยํ, ตํ ปญฺญาพลํ, ยํ ปญฺญาพลํ, ตํ ปญฺญินฺทฺริยํ.\csromanspacer{} ปญฺจนฺนํ ภิกฺขเว อินฺทฺริยานํ ภาวิตตฺตา พหุลีกตตฺตา ภิกฺขุ อาสวานํ ขยา อนาสวํ เจโตวิมุตฺติํ ปญฺญาวิมุตฺติํ ทิฏฺเฐว ธมฺเม สยํ อภิญฺญา สจฺฉิกตฺวา อุปสมฺปชฺช วิหรตีติ.\csromanspacer{}\csromanspacer{}ตติยํ.}

\csromanmatikaanchor{mtk.121}
\csromantocmarkref{subparagraph}{4. ปุพฺพโกฏฺฐกสุตฺต}{mtk.121}
\bookmark[dest={mtk.121},level=5]{4. ปุพฺพโกฏฺฐกสุตฺต}
\csromanheader{6}{4. ปุพฺพโกฏฺฐกสุตฺต}
\proseitem{514}{เอวํ เม สุตํ\csromandash{}เอกํ สมยํ ภควา สาวตฺถิยํ วิหรติ ปุพฺพโกฏฺฐเก.\csromanspacer{} ตตฺร โข ภควา อายสฺมนฺตํ สาริปุตฺตํ อามนฺเตสิ “สทฺทหสิ\footnote{สทฺทหาสิ (สี, อิ)} ตฺวํ สาริปุตฺต สทฺธินฺทฺริยํ ภาวิตํ พหุลีกตํ อมโตคธํ โหติ อมตปรายณํ อมตปริโยสานํ ฯเปฯ ปญฺญินฺทฺริยํ ภาวิตํ พหุลีกตํ อมโตคธํ โหติ อมตปรายณํ อมตปริโยสานนฺ”ติ.}

\prose{น ขฺวาหํ เอตฺถ ภนฺเต ภควโต สทฺธาย คจฺฉามิ สทฺธินฺทฺริยํ ฯเปฯ ปญฺญินฺทฺริยํ ภาวิตํ พหุลีกตํ อมโตคธํ โหติ อมตปรายณํ อมตปริโยสานํ.\csromanspacer{} เยสํ เหตํ ภนฺเต อญฺญาตํ อสฺส อทิฏฺฐํ อวิทิตํ อสจฺฉิกตํ อผสฺสิตํ\footnote{อปสฺสิตํ (สี, สฺยา, กํ, ก), อผุสิตํ (อิ)} ปญฺญาย, เต ตตฺถ ปเรสํ สทฺธาย คจฺเฉยฺยุํ.\csromanspacer{} สทฺธินฺทฺริยํ ฯเปฯ ปญฺญินฺทฺริยํ ภาวิตํ พหุลีกตํ อมโตคธํ โหติ อมตปรายณํ อมตปริโยสานํ.\csromanspacer{} เยสญฺจ โข เอตํ ภนฺเต ญาตํ ทิฏฺฐํ วิทิตํ สจฺฉิกตํ ผสฺสิตํ ปญฺญาย, นิกฺกงฺขา เต ตตฺถ นิพฺพิจิกิจฺฉา สทฺธินฺทฺริยํ ฯเปฯ ปญฺญินฺทฺริยํ ภาวิตํ พหุลีกตํ อมโตคธํ โหติ อมตปรายณํ อมตปริโยสานํ.\csromanspacer{} มยฺหญฺจ โข เอตํ ภนฺเต ญาตํ ทิฏฺฐํ วิทิตํ สจฺฉิกตํ ผสฺสิตํ ปญฺญาย, นิกฺกงฺขฺวาหํ ตตฺถ นิพฺพิจิกิจฺโฉ สทฺธินฺทฺริยํ ฯเปฯ ปญฺญินฺทฺริยํ ภาวิตํ พหุลีกตํ อมโตคธํ โหติ อมตปรายณํ อมตปริโยสานนฺติ.}

\prose{สาธุ สาธุ สาริปุตฺต, เยสํ เหตํ สาริปุตฺต อญฺญาตํ อสฺส อทิฏฺฐํ อวิทิตํ อสจฺฉิกตํ อผสฺสิตํ ปญฺญาย, เต ตตฺถ ปเรสํ สทฺธาย คจฺเฉยฺยุํ.\csromanspacer{} สทฺธินฺทฺริยํ ภาวิตํ พหุลีกตํ อมโตคธํ โหติ อมตปรายณํ อมตปริโยสานํ ฯเปฯ ปญฺญินฺทฺริยํ ภาวิตํ พหุลีกตํ อมโตคธํ โหติ อมตปรายณํ อมตปริโยสานํ.\csromanspacer{} เยสญฺจ โข เอตํ สาริปุตฺต ญาตํ ทิฏฺฐํ วิทิตํ สจฺฉิกตํ ผสฺสิตํ ปญฺญาย, นิกฺกงฺขา เต ตตฺถ}

\csromanheader{2}{4. อินฺทฺริยสํยุตฺต}
\csromanmatikaanchor{mtk.122}
\csromantocmarkref{subparagraph}{5-8. ปุพฺพารามสุตฺตานิ}{mtk.122}
\bookmark[dest={mtk.122},level=5]{5-8. ปุพฺพารามสุตฺตานิ}
\prose{\csromanfolio{195}นิพฺพิจิกิจฺฉา สทฺธินฺทฺริยํ ภาวิตํ พหุลีกตํ อมโตคธํ โหติ อมตปรายณํ อมตปริโยสานํ ฯเปฯ ปญฺญินฺทฺริยํ ภาวิตํ พหุลีกตํ อมโตคธํ โหติ อมตปรายณํ อมตปริโยสานนฺติ.\csromanspacer{}\csromanspacer{}จตุตฺถํ.}

\csromanheader{2}{5. ปฐมปุพฺพารามสุตฺต}
\proseitem{515}{เอวํ เม สุตํ\csromandash{}เอกํ สมยํ ภควา สาวตฺถิยํ วิหรติ ปุพฺพาราเม มิคารมาตุปาสาเท.\csromanspacer{} ตตฺร โข ภวคา ภิกฺขู อมนฺเตสิ “กตินํ นุ โข ภิกฺขเว อินฺทฺริยานํ ภาวิตตฺตา พหุลีกตตฺตา ขีณาสโว ภิกฺขุ อญฺญํ พฺยากโรติ ‘ขีณา ชาติ, วุสิตํ พฺรหฺมจริยํ, กตํ กรณียํ, นาปรํ อิตฺถตฺตายาติ ปชานามี’ติ”.}

\prose{ภควํมูลกา โน ภนฺเต ธมฺมา ฯเปฯ.\csromanspacer{} เอกสฺส โข ภิกฺขเว อินฺทฺริยสฺส ภาวิตตฺตา พหุลีกตตฺตา ขีณาสโว ภิกฺขุ อญฺญํ พฺยากโรติ “ขีณา ชาติ, วุสิตํ พฺรหฺมจริยํ, กตํ กรณียํ, นาปรํ อิตฺถตฺตายาติ ปชานามี”ติ.\csromanspacer{} กตมสฺส เอกสฺส, ปญฺญินฺทฺริยสฺส.\csromanspacer{} ปญฺญวโต ภิกฺขเว อริยสาวกสฺส ตทนฺวยา สทฺธา สณฺฐาติ, ตทนฺวยํ วีริยํ สณฺฐาติ, ตทนฺวยา สติ สณฺฐาติ, ตทนฺวโย สมาธิ สณฺฐาติ.\csromanspacer{} อิมสฺส โข ภิกฺขเว เอกสฺส อินฺทฺริยสฺส ภาวิตตฺตา พหุลีกตตฺตา ขีณาสโว ภิกฺขุ อญฺญํ พฺยากโรติ “ขีณา ชาติ, วุสิตํ พฺรหฺมจริยํ, กตํ กรณียํ, นาปรํ อิตฺถตฺตายาติ ปชานามี”ติ.\csromanspacer{}\csromanspacer{}ปญฺจมํ.}

\csromanheader{2}{6. ทุติยปุพฺพารามสุตฺต}
\proseitem{516}{ตํเยว นิทานํ.\csromanspacer{} กตินํ นุ โข ภิกฺขเว อินฺทฺริยานํ ภาวิตตฺตา พหุลีกตตฺตา ขีณาสโว ภิกฺขุ อญฺญํ พฺยากโรติ “ขีณา ชาติ, วุสิตํ พฺรหฺมจริยํ, กตํ กรณียํ, นาปรํ อิตฺถตฺตายาติ ปชานามี”ติ.\csromanspacer{} ภควํมูลกา โน ภนฺเต ธมฺมา ฯเปฯ.\csromanspacer{} ทฺวินฺนํ โข ภิกฺขเว อินฺทฺริยานํ ภาวิตตฺตา พหุลีกตตฺตา ขีณาสโว ภิกฺขุ อญฺญํ พฺยากโรติ “ขีณา ชาติ, วุสิตํ พฺรหฺมจริยํ, กตํ กรณียํ, นาปรํ อิตฺถตฺตายาติ ปชานามี”ติ.\csromanspacer{} กตเมสํ ทฺวินฺนํ, อริยาย จ ปญฺญาย อริยาย จ วิมุตฺติยา.\csromanspacer{} ยา หิสฺส ภิกฺขเว อริยา ปญฺญา, ตทสฺส ปญฺญินฺทฺริยํ.\csromanspacer{} ยา หิสฺส ภิกฺขเว อริยา วิมุตฺติ, ตทสฺส สมาธินฺทฺริยํ.\csromanspacer{} อิเมสํ โข ภิกฺขเว ทฺวินฺนํ อินฺทฺริยานํ ภาวิตตฺตา พหุลีกตตฺตา ขีณาสโว ภิกฺขุ อญฺญํ พฺยากโรติ “ขีณา ชาติ, \csromanfolio{196}วุสิตํ พฺรหฺมจริยํ, กตํ กรณียํ, นาปรํ อิตฺถตฺตายาติ ปชานามี”ติ.\csromanspacer{}\csromanspacer{}ฉฏฺฐํ.}

\csromanheader{2}{7. ตติยปุพฺพารามสุตฺต}
\proseitem{517}{ตํเยว นิทานํ.\csromanspacer{} กตินํ นุ โข ภิกฺขเว อินฺทฺริยานํ ภาวิตตฺตา พหุลีกตตฺตา ขีณาสโว ภิกฺขุ อญฺญํ พฺยากโรติ “ขีณา ชาติ, วุสิตํ พฺรหฺมจริยํ, กตํ กรณียํ, นาปรํ อิตฺถตฺตายาติ ปชานามี”ติ.\csromanspacer{} ภควํมูลกา โน ภนฺเต ธมฺมา ฯเปฯ.\csromanspacer{} จตุนฺนํ โข ภิกฺขเว อินฺทฺริยานํ ภาวิตตฺตา พหุลีกตตฺตา ขีณาสโว ภิกฺขุ อญฺญํ พฺยากโรติ “ขีณา ชาติ, วุสิตํ พฺรหฺมจริยํ, กตํ กรณียํ, นาปรํ อิตฺถตฺตายาติ ปชานามี”ติ.\csromanspacer{} กตเมสํ จตุนฺนํ, วีริยินฺทฺริยสฺส สตินฺทฺริยสฺส สมาธินฺทฺริยสฺส ปญฺญินฺทฺริยสฺส.\csromanspacer{} อิเมสํ โข ภิกฺขเว จตุนฺนํ อินฺทฺริยานํ ภาวิตตฺตา พหุลีกตตฺตา ขีณาสโว ภิกฺขุ อญฺญํ พฺยากโรติ “ขีณา ชาติ, วุสิตํ พฺรหฺมจริยํ, กตํ กรณียํ, นาปรํ อิตฺถตฺตายาติ ปชานามี”ติ.\csromanspacer{}\csromanspacer{}สตฺตมํ.}

\csromanheader{2}{8. จตุตฺถปุพฺพารามสุตฺต}
\proseitem{518}{ตํเยว นิทานํ.\csromanspacer{} กตินํ นุ โข ภิกฺขเว อินฺทฺริยานํ ภาวิตตฺตา พหุลีกตตฺตา ขีณาสโว ภิกฺขุ อญฺญํ พฺยากโรติ “ขีณา ชาติ, วุสิตํ พฺรหฺมจริยํ, กตํ กรณียํ, นาปรํ อิตฺถตฺตายาติ ปชานามี”ติ.\csromanspacer{} ภควํมูลกา โน ภนฺเต ธมฺมา ฯเปฯ.\csromanspacer{} ปญฺจนฺนํ โข ภิกฺขเว อินฺทฺริยานํ ภาวิตตฺตา พหุลีกตตฺตา ขีณาสโว ภิกฺขุ อญฺญํ พฺยากโรติ “ขีณา ชาติ, วุสิตํ พฺรหฺมจริยํ, กตํ กรณียํ, นาปรํ อิตฺถตฺตายาติ ปชานามี”ติ.\csromanspacer{} กตเมสํ ปญฺจนฺนํ, สทฺธินฺทฺริยสฺส วีริยินฺทฺริยสฺส สตินฺทฺริยสฺส สมาธินฺทฺริยสฺส ปญฺญินฺทฺริยสฺส.\csromanspacer{} อิเมสํ โข ภิกฺขเว ปญฺจนฺนํ อินฺทฺริยานํ ภาวิตตฺตา พหุลีกตตฺตา ขีณาสโว ภิกฺขุ อญฺญํ พฺยากโรติ “ขีณา ชาติ, วุสิตํ พฺรหฺมจริยํ, กตํ กรณียํ, นาปรํ อิตฺถตฺตายาติ ปชานามี”ติ.\csromanspacer{}\csromanspacer{}อฏฺฐมํ.}

\csromanmatikaanchor{mtk.123}
\csromantocmarkref{subparagraph}{9. ปิณฺโฑลภารทฺวาชสุตฺต}{mtk.123}
\bookmark[dest={mtk.123},level=5]{9. ปิณฺโฑลภารทฺวาชสุตฺต}
\csromanheader{6}{9. ปิณฺโฑลภารทฺวาชสุตฺต}
\proseitem{519}{เอวํ เม สุตํ\csromandash{}เอกํ สมยํ ภควา โกสมฺพิยํ วิหรติ โฆสิตาราเม.\csromanspacer{} เตน โข ปน สมเยน อายสฺมตา ปิณฺโฑลภารทฺวาเชน อญฺญา พฺยากตา โหติ “ขีณา ชาติ, วุสิตํ พฺรหฺมจริยํ, กตํ}

\csromanheader{2}{4. อินฺทฺริยสํยุตฺต}
\prose{\csromanfolio{197}กรณียํ, นาปรํ อิตฺถตฺตายาติ ปชานามี”ติ.\csromanspacer{} อถ โข สมฺพหุลา ภิกฺขู เยน ภควา เตนุปสงฺกมิํสุ, อุปสงฺกมิตฺวา ภควนฺตํ อภิวาเทตฺวา เอกมนฺตํ นิสีทิํสุ, เอกมนฺตํ นิสินฺนา โข เต ภิกฺขู ภควนฺตํ เอตทโวจุํ\csromandash{}}

\prose{อายสฺมตา ภนฺเต ปิณฺโฑลภารทฺวาเชน อญฺญา พฺยากตา “ขีณา ชาติ, วุสิตํ พฺรหฺมจริยํ, กตํ กรณียํ, นาปรํ อิตฺถตฺตายาติ ปชานามี”ติ.\csromanspacer{} กิํ นุ โข ภนฺเต อตฺถวสํ สมฺปสฺสมาเนน อายสฺมตา ปิณฺโฑลภารทฺวาเชน อญฺญา พฺยากตา “ขีณา ชาติ, วุสิตํ พฺรหฺมจริยํ, กตํ กรณียํ, นาปรํ อิตฺถตฺตายาติ ปชานามี”ติ.}

\prose{ติณฺณนฺนํ โข ภิกฺขเว อินฺทฺริยานํ ภาวิตตฺตา พหุลีกตตฺตา ปิณฺโฑลภารทฺวาเชน ภิกฺขุนา อญฺญา พฺยากตา “ขีณา ชาติ, วุสิตํ พฺรหฺมจริยํ, กตํ กรณียํ, นาปรํ อิตฺถตฺตายาติ ปชานามี”ติ.\csromanspacer{} กตเมสํ ติณฺณนฺนํ, สตินฺทฺริยสฺส สมาธินฺทฺริยสฺส ปญฺญินฺทฺริยสฺส.\csromanspacer{} อิเมสํ โข ภิกฺขเว ติณฺณนฺนํ อินฺทฺริยานํ ภาวิตตฺตา พหุลีกตตฺตา ปิณฺโฑลภารทฺวาเชน ภิกฺขุนา อญฺญา พฺยากตา “ขีณา ชาติ, วุสิตํ พฺรหฺมจริยํ, กตํ กรณียํ, นาปรํ อิตฺถตฺตายาติ ปชานามี”ติ.\csromanspacer{} อิมานิ จ ภิกฺขเว ตีณินฺทฺริยานิ กิมนฺตานิ, ขยนฺตานิ.\csromanspacer{} กิสฺส ขยนฺตานิ, ชาติชรามรณสฺส, “ชาติชรามรณํ ขยนฺ”ติ โข ภิกฺขเว สมฺปสฺสมาเนน ปิณฺโฑลภารทฺวาเชน ภิกฺขุนา อญฺญา พฺยากตา “ขีณา ชาติ, วุสิตํ พฺรหฺมจริยํ, กตํ กรณียํ, นาปรํ อิตฺถตฺตายาติ ปชานามี”ติ.\csromanspacer{}\csromanspacer{}นวมํ.}

\csromanmatikaanchor{mtk.124}
\csromantocmarkref{subparagraph}{10. อาปณสุตฺต}{mtk.124}
\bookmark[dest={mtk.124},level=5]{10. อาปณสุตฺต}
\csromanheader{6}{10. อาปณสุตฺต}
\proseitem{520}{เอวํ เม สุตํ\csromandash{}เอกํ สมยํ ภควา องฺเคสุ วิหรติ อาปณํ นาม องฺคานํ นิคโม.\csromanspacer{} ตตฺร โข ภควา อายสฺมนฺตํ สาริปุตฺตํ อามนฺเตสิ “โย โส สาริปุตฺต อริยสาวโก ตถาคเต เอกนฺตคโต\footnote{เอกนฺติคโต (สี)} อภิปฺปสนฺโน, น โส ตถาคเต วา ตถาคตสาสเน วา กงฺเขยฺย วา วิจิกิจฺเฉยฺย วา”ติ.}

\prose{โย โส ภนฺเต อริยสาวโก ตถาคเต เอกนฺตคโต อภิปฺปสนฺโน, น โส ตถาคเต วา ตถาคตสาสเน วา กงฺเขยฺย วา วิจิกิจฺเฉยฺย วา.\csromanspacer{} สทฺธสฺส หิ ภนฺเต อริยสาวกสฺส เอตํ ปาฏิกงฺขํ \csromanfolio{198}ยํ อารทฺธวีริโย วิหริสฺสติ อกุสลานํ ธมฺมานํ ปหานาย กุสลานํ ธมฺมานํ อุปสมฺปทาย ถามวา ทฬฺหปรกฺกโม อนิกฺขิตฺตธุโร กุสเลสุ ธมฺเมสุ.\csromanspacer{} ยํ หิสฺส ภนฺเต วีริยํ, ตทสฺส วีริยินฺทฺริยํ.}

\prose{สทฺธสฺส หิ ภนฺเต อริยสาวกสฺส อารทฺธวีริยสฺส เอตํ ปาฏิกงฺขํ ยํ สติมา ภวิสฺสติ ปรเมน สติเนปกฺเกน สมนฺนาคโต จิรกตมฺปิ จิรภาสิตมฺปิ สริตา อนุสฺสริตา.\csromanspacer{} ยา หิสฺส ภนฺเต สติ, ตทสฺส สตินฺทฺริยํ.}

\prose{สทฺธสฺส หิ ภนฺเต อริยสาวกสฺส อารทฺธวีริยสฺส อุปฏฺฐิตสฺสติโน เอตํ ปาฏิกงฺขํ ยํ โวสฺสคฺคารมฺมณํ กริตฺวา ลภิสฺสติ สมาธิํ, ลภิสฺสติ จิตฺตสฺส เอกคฺคตํ.\csromanspacer{} โย หิสฺส ภนฺเต สมาธิ, ตทสฺส สมาธินฺทฺริยํ.}

\prose{สทฺธสฺส หิ ภนฺเต อริยสาวกสฺส อารทฺธวีริยสฺส อุปฏฺฐิตสฺสติโน สมาหิตจิตฺตสฺส เอตํ ปาฏิกงฺขํ ยํ เอวํ ปชานิสฺสติ “อนมตคฺโค โข สํสาโร, ปุพฺพโกฏิ น ปญฺญายติ อวิชฺชานีวรณานํ สตฺตานํ ตณฺหาสํโยชนานํ สนฺธาวตํ สํสรตํ, อวิชฺชาย เตฺวว ตโมกายสฺส อเสสวิราคนิโรโธ สนฺตเมตํ ปทํ, ปณีตเมตํ ปทํ, ยทิทํ สพฺพสงฺขารสมโถ สพฺพูปธิปฏินิสฺสคฺโค ตณฺหากฺขโย วิราโค นิโรโธ นิพฺพานํ”\footnote{นิพฺพานนฺติ (?)}.\csromanspacer{} ยา หิสฺส ภนฺเต ปญฺญา, ตทสฺส ปญฺญินฺทฺริยํ.}

\prose{สทฺโธ โส\footnote{ส โข โส (สี, สฺยา, กํ)} ภนฺเต อริยสาวโก เอวํ ปทหิตฺวา ปทหิตฺวา เอวํ สริตฺวา สริตฺวา เอวํ สมาทหิตฺวา สมาทหิตฺวา เอวํ ปชานิตฺวา ปชานิตฺวา เอวํ อภิสทฺทหติ “อิเม โข เต ธมฺมา, เย เม ปุพฺเพ สุตวา อเหสุํ, เตนาหํ เอตรหิ กาเยน จ ผุสิตฺวา วิหรามิ, ปญฺญาย จ อติวิชฺฌ\footnote{ปฏิวิชฺฌ (สี, ก) ตทฏฺฐกถาสุ ปน อติวิชฺฌิตฺวาติ วณฺณิตํ.} ปสฺสามี”ติ.\csromanspacer{} ยา หิสฺส ภนฺเต สทฺธา, ตทสฺส สทฺธินฺทฺริยนฺติ.}

\prose{สาธุ สาธุ สาริปุตฺต, โย โส สาริปุตฺต อริยสาวโก ตถาคเต เอกนฺตคโต อภิปฺปสนฺโน, น โส ตถาคเต วา ตถาคตสาสเน วา กงฺเขยฺย วา วิจิกิจฺเฉยฺย วา.\csromanspacer{} สทฺธสฺส หิ สาริปุตฺต อริยสาวกสฺส เอตํ ปาฏิกงฺขํ ยํ อารทฺธวีริโย วิหริสฺสติ}

\csromanheader{2}{4. อินฺทฺริยสํยุตฺต}
\prose{\csromanfolio{199}อกุสลานํ ธมฺมานํ ปหานาย กุสลานํ ธมฺมานํ อุปสมฺปทาย ถามวา ทฬฺหปรกฺกโม อนิกฺขิตฺตธุโร กุสเลสุ ธมฺเมสุ.\csromanspacer{} ยํ หิสฺส สาริปุตฺต วีริยํ, ตทสฺส วีริยินฺทฺริยํ.}

\prose{สทฺธสฺส หิ สาริปุตฺต อริยสาวกสฺส อารทฺธวีริยสฺส เอตํ ปาฏิกงฺขํ ยํ สติมา ภวิสฺสติ ปรเมน สติเนปกฺเกน สมนฺนาคโต จิรกตมฺปิ จิรภาสิตมฺปิ สริตา อนุสฺสริตา.\csromanspacer{} ยา หิสฺส สาริปุตฺต สติ, ตทสฺส สตินฺทฺริยํ.}

\prose{สทฺธสฺส หิ สาริปุตฺต อริยสาวกสฺส อารทฺธวีริยสฺส อุปฏฺฐิตสฺสติโน เอตํ ปาฏิกงฺขํ ยํ โวสฺสคฺคารมฺมณํ กริตฺวา ลภิสฺสติ สมาธิํ, ลภิสฺสติ จิตฺตสฺส เอกคฺคตํ.\csromanspacer{} โย หิสฺส สาริปุตฺต สมาธิ, ตทสฺส สมาธินฺทฺริยํ.}

\prose{สทฺธสฺส หิ สาริปุตฺต อริยสาวกสฺส อารทฺธวีริยสฺส อุปฏฺฐิตสฺสติโน สมาหิตจิตฺตสฺส เอตํ ปาฏิกงฺขํ ยํ เอวํ ปชานิสฺสติ “อนมตคฺโค โข สํสาโร, ปุพฺพโกฏิ น ปญฺญายติ อวิชฺชานีวรณานํ สตฺตานํ ตณฺหาสํโยชนานํ สนฺธาวตํ สํสรตํ, อวิชฺชาย เตฺวว ตโมกายสฺส อเสสวิราคนิโรโธ สนฺตเมตํ ปทํ, ปณีตเมตํ ปทํ, ยทิทํ สพฺพสงฺขารสมโถ สพฺพูปธิปฏินิสฺสคฺโค ตณฺหากฺขโย วิราโค นิโรโธ นิพฺพานํ”.\csromanspacer{} ยา หิสฺส สาริปุตฺต ปญฺญา, ตทสฺส ปญฺญินฺทฺริยํ.}

\prose{สทฺโธ โส\footnote{ส โข โส (สี, สฺยา, กํ, อิ)} สาริปุตฺต อริยสาวโก เอวํ ปทหิตฺวา ปทหิตฺวา เอวํ สริตฺวา สริตฺวา เอวํ สมาทหิตฺวา สมาทหิตฺวา เอวํ ปชานิตฺวา ปชานิตฺวา เอวํ อภิสทฺทหติ “อิเม โข เต ธมฺมา, เย เม ปุพฺเพ สุตวา อเหสุํ, เตนาหํ เอตรหิ กาเยน จ ผุสิตฺวา วิหรามิ, ปญฺญาย จ อติวิชฺฌ\footnote{ปฏิวิชฺฌ (ก-สี, ก)} ปสฺสามี”ติ.\csromanspacer{} ยา หิสฺส สาริปุตฺต สทฺธา, ตทสฺส สทฺธินฺทฺริยนฺติ.\csromanspacer{}\csromanspacer{}ทสมํ.\csromanspacer{} ชราวคฺโค ปญฺจโม.}

\titlehead{\textbf{ตสฺสุทฺทานํ}}
\csromangathameasure{ชรา อุณฺณาโภ พฺราหฺมโณ, สาเกโต ปุพฺพโกฏฺฐโก. \\ ปุพฺพาราเม จ จตฺตาริ, ปิณฺโฑโล อาปเณน จาติ.}
\csromangathagroup{\csromangathastanza{ชรา อุณฺณาโภ พฺราหฺมโณ, สาเกโต ปุพฺพโกฏฺฐโก. \\ ปุพฺพาราเม จ จตฺตาริ, ปิณฺโฑโล อาปเณน จาติ\footnote{สทฺเธน เต ทสาติ (สฺยา, กํ, ก)}.}}

\csromanmatikaanchor{mtk.125}
\csromantocmarknopage{subparagraph}{6. สูกรขตวคฺค}{mtk.125}
\bookmark[dest={mtk.125},level=5]{6. สูกรขตวคฺค}
\csromanheader{6}{6. สูกรขตวคฺค}
\csromanmatikaanchor{mtk.126}
\csromantocmarkref{subparagraph}{1. สาลสุตฺต}{mtk.126}
\bookmark[dest={mtk.126},level=5]{1. สาลสุตฺต}
\csromanheader{6}{1. สาลสุตฺต}
\proseitem{521}{\csromanfolio{200}เอวํ เม สุตํ\csromandash{}เอกํ สมยํ ภควา โกสเลสุ วิหรติ สาลาย พฺราหฺมณคาเม.\csromanspacer{} ตตฺร โข ภควา ภิกฺขู อามนฺเตสิ\csromandash{}เสยฺยถาปิ ภิกฺขเว เย เกจิ ติรจฺฉานคตา ปาณา, สีโห มิคราชา เตสํ อคฺคมกฺขายติ ยทิทํ ถาเมน ชเวน สูเรน\footnote{สูริเยน (สี, สฺยา, กํ)}.\csromanspacer{} เอวเมว โข ภิกฺขเว เย เกจิ โพธิปกฺขิยา ธมฺมา, ปญฺญินฺทฺริยํ เตสํ อคฺคมกฺขายติ ยทิทํ โพธาย.}

\prose{กตเม จ ภิกฺขเว โพธิปกฺขิยา ธมฺมา, สทฺธินฺทฺริยํ ภิกฺขเว โพธิปกฺขิโย ธมฺโม, ตํ โพธาย สํวตฺตติ.\csromanspacer{} วีริยินฺทฺริยํ โพธิปกฺขิโย ธมฺโม, ตํ โพธาย สํวตฺตติ.\csromanspacer{} สตินฺทฺริยํ โพธิปกฺขิโย ธมฺโม, ตํ โพธาย สํวตฺตติ.\csromanspacer{} สมาธินฺทฺริยํ โพธิปกฺขิโย ธมฺโม, ตํ โพธาย สํวตฺตติ, ปญฺญินฺทฺริยํ โพธิปกฺขิโย ธมฺโม, ตํ โพธาย สํวตฺตติ.\csromanspacer{} เสยฺยถาปิ ภิกฺขเว เย เกจิ ติรจฺฉานคตา ปาณา, สีโห มิคราชา เตสํ อคฺคมกฺขายติ ยทิทํ ถาเมน ชเวน สูเรน.\csromanspacer{} เอวเมว โข ภิกฺขเว เย เกจิ โพธิปกฺขิยา ธมฺมา, ปญฺญินฺทฺริยํ เตสํ อคฺคมกฺขายติ ยทิทํ โพธายาติ.\csromanspacer{}\csromanspacer{}ปฐมํ.}

\csromanmatikaanchor{mtk.127}
\csromantocmarkref{subparagraph}{2. มลฺลิกสุตฺต}{mtk.127}
\bookmark[dest={mtk.127},level=5]{2. มลฺลิกสุตฺต}
\csromanheader{6}{2. มลฺลิกสุตฺต}
\proseitem{522}{เอวํ เม สุตํ\csromandash{}เอกํ สมยํ ภควา มลฺเลสุ\footnote{มลฺลเกสุ (สี, สฺยา, กํ), มลฺลิเกสุ (ก)} วิหรติ อุรุเวลกปฺปํ นาม มลฺลานํ นิคโม.\csromanspacer{} ตตฺร โข ภควา ภิกฺขู อามนฺเตสิ\csromandash{}ยาวกีวญฺจ ภิกฺขเว อริยสาวกสฺส อริยญาณํ น อุปฺปนฺนํ โหติ, เนว ตาว จตุนฺนํ อินฺทฺริยานํ สณฺฐิติ โหติ, เนว ตาว จตุนฺนํ อินฺทฺริยานํ อวฏฺฐิติ โหติ.\csromanspacer{} ยโต จ โข ภิกฺขเว อริยสาวกสฺส อริยญาณํ อุปฺปนฺนํ โหติ, อถ จตุนฺนํ อินฺทฺริยานํ สณฺฐิติ โหติ, อถ จตุนฺนํ อินฺทฺริยานํ อวฏฺฐิติ โหติ.}

\prose{เสยฺยถาปิ ภิกฺขเว ยาวกีวญฺจ กูฏาคารสฺส กูฏํ น อุสฺสิตํ โหติ, เนว ตาว โคปานสีนํ สณฺฐิติ โหติ, เนว ตาว โคปานสีนํ อวฏฺฐิติ โหติ.\csromanspacer{} ยโต จ โข ภิกฺขเว กูฏาคารสฺส กูฏํ}

\csromanheader{2}{4. อินฺทฺริยสํยุตฺต}
\prose{\csromanfolio{201}อุสฺสิตํ โหติ, อถ โคปานสีนํ สณฺฐิติ โหติ, อถ โคปานสีนํ อวฏฺฐิติ โหติ.\csromanspacer{} เอวเมว โข ภิกฺขเว ยาวกีวญฺจ อริยสาวกสฺส อริยญาณํ น อุปฺปนฺนํ โหติ, เนว ตาว จตุนฺนํ อินฺทฺริยานํ สณฺฐิติ โหติ, เนว ตาว จตุนฺนํ อินฺทฺริยานํ อวฏฺฐิติ โหติ.\csromanspacer{} ยโต จ โข ภิกฺขเว อริยสาวกสฺส อริยญาณํ อุปฺปนฺนํ โหติ, อถ จตุนฺนํ อินฺทฺริยานํ ฯเปฯ อวฏฺฐิติ โหติ.}

\prose{กตเมสํ จตุนฺนํ, สทฺธินฺทฺริยสฺส วีริยินฺทฺริยสฺส สตินฺทฺริยสฺส สมาธินฺทฺริยสฺส.\csromanspacer{} ปญฺญวโต ภิกฺขเว อริยสาวกสฺส ตทนฺวยา สทฺธา สณฺฐาติ, ตทนฺวยํ วีริยํ สณฺฐาติ, ตทนฺวยา สติ สณฺฐาติ, ตทนฺวโย สมาธิ สณฺฐาตีติ.\csromanspacer{}\csromanspacer{}ทุติยํ.}

\csromanmatikaanchor{mtk.128}
\csromantocmarkref{subparagraph}{3. เสขสุตฺต}{mtk.128}
\bookmark[dest={mtk.128},level=5]{3. เสขสุตฺต}
\csromanheader{6}{3. เสขสุตฺต}
\proseitem{523}{เอวํ เม สุตํ\csromandash{}เอกํ สมยํ ภควา โกสมฺพิยํ วิหรติ โฆสิตาราเม.\csromanspacer{} ตตฺร โข ภควา ภิกฺขู อามนฺเตสิ\csromandash{}อตฺถิ นุ โข ภิกฺขเว ปริยาโย, ยํ ปริยายํ อาคมฺม เสโข ภิกฺขุ เสขภูมิยํ ฐิโต “เสโขสฺมี”ติ ปชาเนยฺย, อเสโข ภิกฺขุ อเสขภูมิยํ ฐิโต “อเสโขสฺมี”ติ ปชาเนยฺยาติ.}

\prose{ภควํมูลกา โน ภนฺเต ธมฺมา ฯเปฯ.\csromanspacer{} อตฺถิ ภิกฺขเว ปริยาโย, ยํ ปริยายํ อาคมฺม เสโข ภิกฺขุ เสขภูมิยํ ฐิโต “เสโขสฺมี”ติ ปชาเนยฺย, อเสโข ภิกฺขุ อเสขภูมิยํ ฐิโต “อเสโขสฺมี”ติ ปชาเนยฺย.}

\prose{กตโม จ ภิกฺขเว ปริยาโย, ยํ ปริยายํ อาคมฺม เสโข ภิกฺขุ เสขภูมิยํ ฐิโต “เสโขสฺมี”ติ ปชานาติ, อิธ ภิกฺขเว เสโข ภิกฺขุ “อิทํ ทุกฺขนฺ”ติ ยถาภูตํ ปชานาติ, “อยํ ทุกฺขสมุทโย”ติ ยถาภูตํ ปชานาติ, “อยํ ทุกฺขนิโรโธ”ติ ยถาภูตํ ปชานาติ, “อยํ ทุกฺขนิโรธคามินี ปฏิปทา”ติ ยถาภูตํ ปชานาติ.\csromanspacer{} อยมฺปิ โข ภิกฺขเว ปริยาโย, ยํ ปริยายํ อาคมฺม เสโข ภิกฺขุ เสขภูมิยํ ฐิโต “เสโขสฺมี”ติ ปชานาติ.}

\csromanmatikaanchor{mtk.129}
\csromantocmarkref{subparagraph}{4-6. ปทสุตฺตาทีนิ}{mtk.129}
\bookmark[dest={mtk.129},level=5]{4-6. ปทสุตฺตาทีนิ}
\prose{ปุน จปรํ ภิกฺขเว เสโข ภิกฺขุ อิติ ปฏิสญฺจิกฺขติ “อตฺถิ นุ โข อิโต พหิทฺธา อญฺโญ สมโณ วา พฺราหฺมโณ วา, โย เอวํ ภูตํ ตจฺฉํ ตถํ \csromanfolio{202}ธมฺมํ เทเสติ ยถา ภควา”ติ.\csromanspacer{} โส เอวํ ปชานาติ “นตฺถิ โข อิโต พหิทฺธา อญฺโญ สมโณ วา พฺราหฺมโณ วา, โย เอวํ ภูตํ ตจฺฉํ ตถํ ธมฺมํ เทเสติ ยถา ภควา”ติ.\csromanspacer{} อยมฺปิ โข ภิกฺขเว ปริยาโย, ยํ ปริยายํ อาคมฺม เสโข ภิกฺขุ เสขภูมิยํ ฐิโต “เสโขสฺมี”ติ ปชานาติ.}

\prose{ปุน จปรํ ภิกฺขเว เสโข ภิกฺขุ ปญฺจินฺทฺริยานิ ปชานาติ, สทฺธินฺทฺริยํ วีริยินฺทฺริยํ สตินฺทฺริยํ สมาธินฺทฺริยํ ปญฺญินฺทฺริยํ.\csromanspacer{} ยํคติกานิ ยํปรมานิ ยํผลานิ ยํปริโยสานานิ น เหว โข กาเยน ผุสิตฺวา วิหรติ, ปญฺญาย จ อติวิชฺฌ ปสฺสติ.\csromanspacer{} อยมฺปิ โข ภิกฺขเว ปริยาโย, ยํ ปริยายํ อาคมฺม เสโข ภิกฺขุ เสขภูมิยํ ฐิโต “เสโขสฺมี”ติ ปชานาติ.}

\prose{กตโม จ ภิกฺขเว ปริยาโย, ยํ ปริยายํ อาคมฺม อเสโข ภิกฺขุ อเสขภูมิยํ “ฐิโต อเสโขสฺมี”ติ ปชานาติ, อิธ ภิกฺขเว อเสโข ภิกฺขุ ปญฺจินฺทฺริยานิ ปชานาติ สทฺธินฺทฺริยํ วีริยินฺทฺริยํ สตินฺทฺริยํ สมาธินฺทฺริยํ ปญฺญินฺทฺริยํ.\csromanspacer{} ยํคติกานิ ยํปรมานิ ยํผลานิ ยํปริโยสานานิ กาเยน จ ผุสิตฺวา วิหรติ, ปญฺญาย จ อติวิชฺฌ ปสฺสติ.\csromanspacer{} อยมฺปิ โข ภิกฺขเว ปริยาโย, ยํ ปริยายํ อาคมฺม อเสโข ภิกฺขุ อเสขภูมิยํ ฐิโต “อเสโขสฺมี”ติ ปชานาติ.}

\prose{ปุน จปรํ ภิกฺขเว อเสโข ภิกฺขุ ฉ อินฺทฺริยานิ ปชานาติ, จกฺขุนฺทฺริยํ โสตินฺทฺริยํ ฆานินฺทฺริยํ ชิวฺหินฺทฺริยํ กายินฺทฺริยํ มนินฺทฺริยํ.\csromanspacer{} อิมานิ โข ฉ อินฺทฺริยานิ สพฺเพน สพฺพํ สพฺพถา สพฺพํ อปริเสสํ นิรุชฺฌิสฺสนฺติ, อญฺญานิ จ ฉ อินฺทฺริยานิ “น กุหิญฺจิ กิสฺมิญฺจิ อุปฺปชฺชิสฺสนฺตี”ติ ปชานาติ.\csromanspacer{} อยมฺปิ โข ภิกฺขเว ปริยาโย, ยํ ปริยายํ อาคมฺม อเสโข ภิกฺขุ อเสขภูมิยํ ฐิโต “อเสโขสฺมี”ติ ปชานาตีติ.\csromanspacer{}\csromanspacer{}ตติยํ.}

\csromanheader{2}{4. ปทสุตฺต}
\proseitem{524}{เสยฺยถาปิ ภิกฺขเว ยานิ กานิจิ ชงฺคลานํ\footnote{ชงฺคมานํ (สี, อิ)} ปาณานํ ปทชาตานิ, สพฺพานิ ตานิ หตฺถิปเท สโมธานํ คจฺฉนฺติ, หตฺถิปทํ เตสํ อคฺคมกฺขายติ ยทิทํ มหนฺตตฺเตน.\csromanspacer{} เอวเมว โข ภิกฺขเว ยานิ กานิจิ ปทานิ โพธาย}

\csromanheader{2}{4. อินฺทฺริยสํยุตฺต}
\prose{\csromanfolio{203}สํวตฺตนฺติ, ปญฺญินฺทฺริยํ ปทํ เตสํ อคฺคมกฺขายติ ยทิทํ โพธาย.\csromanspacer{} กตมานิ จ ภิกฺขเว ปทานิ โพธาย สํวตฺตนฺติ, สทฺธินฺทฺริยํ ภิกฺขเว ปทํ, ตํ โพธาย สํวตฺตติ.\csromanspacer{} วีริยินฺทฺริยํ ปทํ, ตํ โพธาย สํวตฺตติ.\csromanspacer{} สตินฺทฺริยํ ปทํ, ตํ โพธาย สํวตฺตติ.\csromanspacer{} สมาธินฺทฺริยํ ปทํ, ตํ โพธาย สํวตฺตติ.\csromanspacer{} ปญฺญินฺทฺริยํ ปทํ, ตํ โพธาย สํวตฺตติ.\csromanspacer{} เสยฺยถาปิ ภิกฺขเว ยานิ กานิจิ ชงฺคลานํ ปาณานํ ปทชาตานิ, สพฺพานิ ตานิ หตฺถิปเท สโมธานํ คจฺฉนฺติ, หตฺถิปทํ เตสํ อคฺคมกฺขายติ ยทิทํ มหนฺตตฺเตน.\csromanspacer{} เอวเมว โข ภิกฺขเว ยานิ กานิจิ ปทานิ โพธาย สํวตฺตนฺติ, ปญฺญินฺทฺริยํ ปทํ เตสํ อคฺคมกฺขายติ ยทิทํ โพธายาติ.\csromanspacer{}\csromanspacer{}จตุตฺถํ.}

\csromanheader{2}{5. สารสุตฺต}
\proseitem{525}{เสยฺยถาปิ ภิกฺขเว เย เกจิ สารคนฺธา, โลหิตจนฺทนํ เตสํ อคฺคมกฺขายติ.\csromanspacer{} เอวเมว โข ภิกฺขเว เย เกจิ โพธิปกฺขิยา ธมฺมา, ปญฺญินฺทฺริยํ เตสํ อคฺคมกฺขายติ ยทิทํ โพธาย.\csromanspacer{} กตเม จ ภิกฺขเว โพธิปกฺขิยา ธมฺมา, สทฺธินฺทฺริยํ ภิกฺขเว โพธิปกฺขิโย ธมฺโม, ตํ โพธาย สํวตฺตติ.\csromanspacer{} วีริยินฺทฺริยํ ฯเปฯ.\csromanspacer{} สตินฺทฺริยํ ฯเปฯ.\csromanspacer{} สมาธินฺทฺริยํ ฯเปฯ.\csromanspacer{} ปญฺญินฺทฺริยํ โพธิปกฺขิโย ธมฺโม, ตํ โพธาย สํวตฺตติ.\csromanspacer{} เสยฺยถาปิ ภิกฺขเว เย เกจิ สารคนฺธา, โลหิตจนฺทนํ เตสํ อคฺคมกฺขายติ.\csromanspacer{} เอวเมว โข ภิกฺขเว เย เกจิ โพธิปกฺขิยา ธมฺมา, ปญฺญินฺทฺริยํ เตสํ อคฺคมกฺขายติ ยทิทํ โพธายาติ.\csromanspacer{}\csromanspacer{}ปญฺจมํ.}

\csromanheader{2}{6. ปติฏฺฐิตสุตฺต}
\proseitem{526}{เอกธมฺเม ปติฏฺฐิตสฺส ภิกฺขเว ภิกฺขุโน ปญฺจินฺทฺริยานิ ภาวิตานิ โหนฺติ สุภาวิตานิ.\csromanspacer{} กตมสฺมิํ เอกธมฺเม อปฺปมาเท.\csromanspacer{} กตโม จ ภิกฺขเว อปฺปมาโท, อิธ ภิกฺขเว ภิกฺขุ จิตฺตํ รกฺขติ อาสเวสุ จ สาสเวสุ จ ธมฺเมสุ, ตสฺส จิตฺตํ รกฺขโต อาสเวสุ จ สาสเวสุ จ ธมฺเมสุ สทฺธินฺทฺริยมฺปิ ภาวนาปาริปูริํ คจฺฉติ, วีริยินฺทฺริยมฺปิ ภาวนาปาริปูริํ คจฺฉติ, สตินฺทฺริยมฺปิ ภาวนาปาริปูริํ คจฺฉติ, สมาธินฺทฺริยมฺปิ ภาวนาปาริปูริํ คจฺฉติ, ปญฺญินฺทฺริยมฺปิ ภาวนาปาริปูริํ คจฺฉติ.\csromanspacer{} เอวมฺปิ โข ภิกฺขเว เอกธมฺเม ปติฏฺฐิตสฺส ภิกฺขุโน ปญฺจินฺทฺริยานิ ภาวิตานิ โหนฺติ สุภาวิตานีติ.\csromanspacer{}\csromanspacer{}ฉฏฺฐํ.}

\csromanmatikaanchor{mtk.130}
\csromantocmarkref{subparagraph}{7. สหมฺปติพฺรหฺมสุตฺต}{mtk.130}
\bookmark[dest={mtk.130},level=5]{7. สหมฺปติพฺรหฺมสุตฺต}
\csromanheader{6}{7. สหมฺปติพฺรหฺมสุตฺต}
\proseitem{527}{\csromanfolio{204}เอกํ สมยํ ภควา อุรุเวลายํ วิหรติ นชฺชา เนรญฺชราย ตีเร อชปาลนิโคฺรเธ ปฐมาภิสมฺพุทฺโธ.\csromanspacer{} อถ โข ภควโต รโหคตสฺส ปฏิสลฺลีนสฺส เอวํ เจตโส ปริวิตกฺโก อุทปาทิ “ปญฺจินฺทฺริยานิ ภาวิตานิ พหุลีกตานิ อมโตคธานิ โหนฺติ อมตปรายณานิ อมตปริโยสานานิ.\csromanspacer{} กตมานิ ปญฺจ, สทฺธินฺทฺริยํ ภาวิตํ พหุลีกตํ อมโตคธํ โหติ อมตปรายณํ อมตปริโยสานํ.\csromanspacer{} วีริยินฺทฺริยํ ฯเปฯ.\csromanspacer{} สตินฺทฺริยํ ฯเปฯ.\csromanspacer{} สมาธินฺทฺริยํ ฯเปฯ.\csromanspacer{} ปญฺญินฺทฺริยํ ภาวิตํ พหุลีกตํ อมโตคธํ โหติ อมตปรายณํ อมตปริโยสานํ.\csromanspacer{} อิมานิ ปญฺจินฺทฺริยานิ ภาวิตานิ พหุลีกตานิ อมโตคธานิ โหนฺติ อมตปรายณานิ อมตปริโยสานานี”ติ.}

\prose{อถ โข พฺรหฺมา สหมฺปติ ภควโต เจตสา เจโตปริวิตกฺกมญฺญาย เสยฺยถาปิ นาม พลวา ปุริโส สมิญฺชิตํ วา พาหํ ปสาเรยฺย, ปสาริตํ วา พาหํ สมิญฺเชยฺย, เอวเมว พฺรหฺมโลเก อนฺตรหิโต ภควโต ปุรโต ปาตุรโหสิ.\csromanspacer{} อถ โข พฺรหฺมา สหมฺปติ เอกํสํ อุตฺตราสงฺคํ กริตฺวา เยน ภควา เตนญฺชลิํ ปณาเมตฺวา ภควนฺตํ เอตทโวจ\csromandash{}เอวเมตํ ภควา, เอวเมตํ สุคต, ปญฺจินฺทฺริยานิ ภาวิตานิ พหุลีกตานิ อมโตคธานิ โหนฺติ อมตปรายณานิ อมตปริโยสานานิ.\csromanspacer{} กตมานิ ปญฺจ, สทฺธินฺทฺริยํ ภาวิตํ พหุลีกตํ อมโตคธํ โหติ อมตปรายณํ อมตปริโยสานํ ฯเปฯ ปญฺญินฺทฺริยํ ภาวิตํ พหุลีกตํ อมโตคธํ โหติ อมตปรายณํ อมตปริโยสานํ.\csromanspacer{} อิมานิ ปญฺจินฺทฺริยานิ ภาวิตานิ พหุลีกตานิ อมโตคธานิ โหนฺติ อมตปรายณานิ อมตปริโยสานานิ.}

\prose{ภูตปุพฺพาหํ ภนฺเต กสฺสเป สมฺมาสมฺพุทฺเธ พฺรหฺมจริยํ อจริํ, ตตฺรปิ มํ เอวํ ชานนฺติ “สหโก ภิกฺขุ สหโก ภิกฺขู”ติ.\csromanspacer{} โส ขฺวาหํ ภนฺเต อิเมสํเยว ปญฺจนฺนํ อินฺทฺริยานํ ภาวิตตฺตา พหุลีกตตฺตา กาเมสุ กามจฺฉนฺทํ วิราเชตฺวา กายสฺส เภทา ปรํ มรณา สุคติํ พฺรหฺมโลกํ อุปปนฺโน, ตตฺรปิ มํ เอวํ ชานนฺติ “พฺรหฺมา สหมฺปติ พฺรหฺมา สหมฺปตี”ติ.\csromanspacer{} เอวเมตํ ภควา, เอวเมตํ สุคต, อหเมตํ ชานามิ, อหเมตํ ปสฺสามิ, ยถา อิมานิ ปญฺจินฺทฺริยานิ ภาวิตานิ พหุลีกตานิ อมโตคธานิ โหนฺติ อมตปรายณานิ อมตปริโยสานานีติ.\csromanspacer{}\csromanspacer{}สตฺตมํ.}

\csromanheader{2}{4. อินฺทฺริยสํยุตฺต}
\csromanheader{2}{8. สูกรขตสุตฺต}
\csromanmatikaanchor{mtk.131}
\csromantocmarkref{subparagraph}{8-10. สูกรขตสุตฺตาทีนิ}{mtk.131}
\bookmark[dest={mtk.131},level=5]{8-10. สูกรขตสุตฺตาทีนิ}
\proseitem{528}{\csromanfolio{205}เอกํ สมยํ ภควา ราชคเห วิหรติ คิชฺฌกูเฏ ปพฺพเต สูกรขตายํ.\csromanspacer{} ตตฺร โข ภควา อายสฺมนฺตํ สาริปุตฺตํ อามนฺเตสิ “กิํ นุ โข สาริปุตฺต อตฺถวสํ สมฺปสฺสมาโน ขีณาสโว ภิกฺขุ ตถาคเต วา ตถาคตสาสเน วา ปรมนิปจฺจการํ ปวตฺตมาโน ปวตฺตตี”ติ\footnote{ปวตฺเตตีติ (สี)}.\csromanspacer{} อนุตฺตรํ หิ ภนฺเต โยคกฺเขมํ สมฺปสฺสมาโน ขีณาสโว ภิกฺขุ ตถาคเต วา ตถาคตสาสเน วา ปรมนิปจฺจการํ ปวตฺตมาโน ปวตฺตตีติ.\csromanspacer{} สาธุ สาธุ สาริปุตฺต, อนุตฺตรํ หิ สาริปุตฺต โยคกฺเขมํ สมฺปสฺสมาโน ขีณาสโว ภิกฺขุ ตถาคเต วา ตถาคตสาสเน วา ปรมนิปจฺจการํ ปวตฺตมาโน ปวตฺตติ.}

\prose{กตโม จ สาริปุตฺต อนุตฺตโร โยคกฺเขโม, ยํ สมฺปสฺสมาโน ขีณาสโว ภิกฺขุ ตถาคเต วา ตถาคตสาสเน วา ปรมนิปจฺจการํ ปวตฺตมาโน ปวตฺตตีติ.\csromanspacer{} อิธ ภนฺเต ขีณาสโว ภิกฺขุ สทฺธินฺทฺริยํ ภาเวติ อุปสมคามิํ สมฺโพธคามิํ.\csromanspacer{} วีริยินฺทฺริยํ ภาเวติ ฯเปฯ.\csromanspacer{} สตินฺทฺริยํ ภาเวติ.\csromanspacer{} สมาธินฺทฺริยํ ภาเวติ.\csromanspacer{} ปญฺญินฺทฺริยํ ภาเวติ อุปสมคามิํ สมฺโพธคามิํ.\csromanspacer{} อยํ โข ภนฺเต อนุตฺตโร โยคกฺเขโม, ยํ สมฺปสฺสมาโน ขีณาสโว ภิกฺขุ ตถาคเต วา ตถาคตสาสเน วา ปรมนิปจฺจการํ ปวตฺตมาโน ปวตฺตตีติ.\csromanspacer{} สาธุ สาธุ สาริปุตฺต, เอโส หิ สาริปุตฺต อนุตฺตโร โยคกฺเขโม, ยํ สมฺปสฺสมาโน ขีณาสโว ภิกฺขุ ตถาคเต วา ตถาคตสาสเน วา ปรมนิปจฺจการํ ปวตฺตมาโน ปวตฺตตีติ.}

\csromanmatikaanchor{mtk.132}
\csromanmatikaanchor{mtk.133}
\csromantocmarknopage{subparagraph}{7. โพธิปกฺขิยวคฺค}{mtk.132}
\bookmark[dest={mtk.132},level=5]{7. โพธิปกฺขิยวคฺค}
\csromantocmarkref{subparagraph}{1-6. สํโยชนสุตฺตาทีนิ}{mtk.133}
\bookmark[dest={mtk.133},level=5]{1-6. สํโยชนสุตฺตาทีนิ}
\prose{กตโม จ สาริปุตฺต ปรมนิปจฺจกาโร, ยํ ขีณาสโว ภิกฺขุ ตถาคเต วา ตถาคตสาสเน วา ปรมนิปจฺจการํ ปวตฺตมาโน ปวตฺตตีติ.\csromanspacer{} อิธ ภนฺเต ขีณาสโว ภิกฺขุ สตฺถริ สคารโว วิหรติ สปฺปติสฺโส\footnote{สปฺปฏิสฺโส (สฺยา, กํ, ก)}, ธมฺเม สคารโว วิหรติ สปฺปติสฺโส, สํเฆ สคารโว วิหรติ สปฺปติสฺโส, สิกฺขาย สคารโว วิหรติ สปฺปติสฺโส, สมาธิสฺมิํ สคารโว วิหรติ สปฺปติสฺโส.\csromanspacer{} อยํ โข ภนฺเต ปรมนิปจฺจกาโร, ยํ ขีณาสโว ภิกฺขุ ตถาคเต วา ตถาคตสาสเน วา ปรมนิปจฺจการํ ปวตฺตมาโน ปวตฺตตีติ.\csromanspacer{} สาธุ สาธุ สาริปุตฺต, เอโส \csromanfolio{206}หิ สาริปุตฺต ปรมนิปจฺจกาโร, ยํ ขีณาสโว ภิกฺขุ ตถาคเต วา ตถาคตสาสเน วา ปรมนิปจฺจการํ ปวตฺตมาโน ปวตฺตตีติ.\csromanspacer{}\csromanspacer{}อฏฺฐมํ.}

\csromanheader{2}{9. ปฐม-อุปฺปาทสุตฺต}
\proseitem{529}{สาวตฺถินิทานํ.\csromanspacer{} ปญฺจิมานิ ภิกฺขเว อินฺทฺริยานิ ภาวิตานิ พหุลีกตานิ อนุปฺปนฺนานิ อุปฺปชฺชนฺติ นาญฺญตฺร ตถาคตสฺส ปาตุภาวา อรหโต สมฺมาสมฺพุทฺธสฺส.\csromanspacer{} กตมานิ ปญฺจ, สทฺธินฺทฺริยํ วีริยินฺทฺริยํ สตินฺทฺริยํ สมาธินฺทฺริยํ ปญฺญินฺทฺริยํ.\csromanspacer{} อิมานิ โข ภิกฺขเว ปญฺจินฺทฺริยานิ ภาวิตานิ พหุลีกตานิ อนุปฺปนฺนานิ อุปฺปชฺชนฺติ นาญฺญตฺร ตถาคตสฺส ปาตุภาวา อรหโต สมฺมาสมฺพุทฺธสฺสาติ.\csromanspacer{}\csromanspacer{}นวมํ.}

\csromanheader{2}{10. ทุติย-อุปฺปาทสุตฺต}
\proseitem{530}{ปญฺจิมานิ ภิกฺขเว อินฺทฺริยานิ ภาวิตานิ พหุลีกตานิ อนุปฺปนฺนานิ อุปฺปชฺชนฺติ นาญฺญตฺร สุคตวินยา.\csromanspacer{} กตมานิ ปญฺจ, สทฺธินฺทฺริยํ วีริยินฺทฺริยํ สตินฺทฺริยํ สมาธินฺทฺริยํ ปญฺญินฺทฺริยํ.\csromanspacer{} อิมานิ โข ภิกฺขเว ปญฺจินฺทฺริยานิ ภาวิตานิ พหุลีกตานิ อนุปฺปนฺนานิ อุปฺปชฺชนฺติ นาญฺญตฺร สุคตวินยาติ.\csromanspacer{}\csromanspacer{}ทสมํ.\csromanspacer{} สูกรขตวคฺโค ฉฏฺโฐ.}

\titlehead{\textbf{ตสฺสุทฺทานํ}}
\csromangathameasure{สาลํ มลฺลิกํ เสโข จ, ปทํ สารํ ปติฏฺฐิตํ. \\ พฺรหฺมสูกรขตาโย, อุปฺปาทา อปเร ทุเวติ.}
\csromangathagroup{\csromangathastanza{สาลํ มลฺลิกํ เสโข จ, ปทํ สารํ ปติฏฺฐิตํ. \\ พฺรหฺมสูกรขตาโย, อุปฺปาทา อปเร ทุเวติ.}}

\csromanheader{6}{7. โพธิปกฺขิยวคฺค}
\csromanheader{2}{1. สํโยชนสุตฺต}
\proseitem{531}{สาวตฺถินิทานํ.\csromanspacer{} ปญฺจิมานิ ภิกฺขเว อินฺทฺริยานิ ภาวิตานิ พหุลีกตานิ สํโยชนปฺปหานาย\footnote{สํโยชนานํ ปหานาย (สฺยา, ก)} สํวตฺตนฺติ.\csromanspacer{} กตมานิ ปญฺจ, สทฺธินฺทฺริยํ ฯเปฯ ปญฺญินฺทฺริยํ.\csromanspacer{} อิมานิ โข ภิกฺขเว ปญฺจินฺทฺริยานิ ภาวิตานิ พหุลีกตานิ สํโยชนปฺปหานาย สํวตฺตนฺตีติ.\csromanspacer{}\csromanspacer{}ปฐมํ.}

\csromanheader{2}{4. อินฺทฺริยสํยุตฺต}
\csromanheader{2}{2. อนุสยสุตฺต}
\proseitem{532}{\csromanfolio{207}ปญฺจิมานิ ภิกฺขเว อินฺทฺริยานิ ภาวิตานิ พหุลีกตานิ อนุสยสมุคฺฆาตาย สํวตฺตนฺติ.\csromanspacer{} กตมานิ ปญฺจ, สทฺธินฺทฺริยํ ฯเปฯ ปญฺญินฺทฺริยํ.\csromanspacer{} อิมานิ โข ภิกฺขเว ปญฺจินฺทฺริยานิ ภาวิตานิ พหุลีกตานิ อนุสยสมุคฺฆาตาย สํวตฺตนฺตีติ.\csromanspacer{}\csromanspacer{}ทุติยํ.}

\csromanheader{2}{3. ปริญฺญาสุตฺต}
\proseitem{533}{ปญฺจิมานิ ภิกฺขเว อินฺทฺริยานิ ภาวิตานิ พหุลีกตานิ อทฺธานปริญฺญาย สํวตฺตนฺติ.\csromanspacer{} กตมานิ ปญฺจ, สทฺธินฺทฺริยํ ฯเปฯ ปญฺญินฺทฺริยํ.\csromanspacer{} อิมานิ โข ภิกฺขเว ปญฺจินฺทฺริยานิ ภาวิตานิ พหุลีกตานิ อทฺธานปริญฺญาย สํวตฺตนฺตีติ.\csromanspacer{}\csromanspacer{}ตติยํ.}

\csromanheader{2}{4. อาสวกฺขยสุตฺต}
\proseitem{534}{ปญฺจิมานิ ภิกฺขเว อินฺทฺริยานิ ภาวิตานิ พหุลีกตานิ อาสวานํ ขยาย สํวตฺตนฺติ.\csromanspacer{} กตมานิ ปญฺจ, สทฺธินฺทฺริยํ ฯเปฯ ปญฺญินฺทฺริยํ.\csromanspacer{} อิมานิ โข ภิกฺขเว ปญฺจินฺทฺริยานิ ภาวิตานิ พหุลีกตานิ อาสวานํ ขยาย สํวตฺตนฺตีติ.}

\prose{ปญฺจิมานิ ภิกฺขเว อินฺทฺริยานิ ภาวิตานิ พหุลีกตานิ สํโยชนปฺปหานาย สํวตฺตนฺติ, อนุสยสมุคฺฆาตาย สํวตฺตนฺติ, อทฺธานปริญฺญาย สํวตฺตนฺติ, อาสวานํ ขยาย สํวตฺตนฺติ.\csromanspacer{} กตมานิ ปญฺจ, สทฺธินฺทฺริยํ ฯเปฯ ปญฺญินฺทฺริยํ.\csromanspacer{} อิมานิ โข ภิกฺขเว ปญฺจินฺทฺริยานิ ภาวิตานิ พหุลีกตานิ สํโยชนปฺปหานาย สํวตฺตนฺติ, อนุสยสมุคฺฆาตาย สํวตฺตนฺติ, อทฺธานปริญฺญาย สํวตฺตนฺติ, อาสวานํ ขยาย สํวตฺตนฺตีติ.\csromanspacer{}\csromanspacer{}จตุตฺถํ.}

\csromanheader{2}{5. ปฐมผลสุตฺต}
\proseitem{535}{ปญฺจิมานิ ภิกฺขเว อินฺทฺริยานิ.\csromanspacer{} กตมานิ ปญฺจ, สทฺธินฺทฺริยํ ฯเปฯ ปญฺญินฺทฺริยํ.\csromanspacer{} อิมานิ โข ภิกฺขเว ปญฺจินฺทฺริยานิ.\csromanspacer{} อิเมสํ โข ภิกฺขเว ปญฺจนฺนํ อินฺทฺริยานํ ภาวิตตฺตา พหุลีกตตฺตา ทฺวินฺนํ ผลานํ อญฺญตรํ ผลํ ปาฏิกงฺขํ, ทิฏฺเฐว ธมฺเม อญฺญา, สติ วา อุปาทิเสเส อนาคามิตาติ.\csromanspacer{}\csromanspacer{}ปญฺจมํ.}

\csromanheader{2}{6. ทุติยผลสุตฺต}
\csromanmatikaanchor{mtk.134}
\csromantocmarkref{subparagraph}{7-10. รุกฺขสุตฺต}{mtk.134}
\bookmark[dest={mtk.134},level=5]{7-10. รุกฺขสุตฺต}
\proseitem{536}{\csromanfolio{208}ปญฺจิมานิ ภิกฺขเว อินฺทฺริยานิ.\csromanspacer{} กตมานิ ปญฺจ, สทฺธินฺทฺริยํ ฯเปฯ ปญฺญินฺทฺริยํ.\csromanspacer{} อิมานิ โข ภิกฺขเว ปญฺจินฺทฺริยานิ.\csromanspacer{} อิเมสํ โข ภิกฺขเว ปญฺจนฺนํ อินฺทฺริยานํ ภาวิตตฺตา พหุลีกตตฺตา สตฺต ผลา สตฺตานิสํสา ปาฏิกงฺขา.\csromanspacer{} กตเม สตฺต ผลา สตฺตานิสํสา, ทิฏฺเฐว ธมฺเม ปฏิกจฺจ อญฺญํ อาราเธติ.\csromanspacer{} โน เจ ทิฏฺเฐว ธมฺเม ปฏิกจฺจ อญฺญํ อาราเธติ, อถ มรณกาเล อญฺญํ อาราเธติ.\csromanspacer{} โน เจ ทิฏฺเฐว ธมฺเม อญฺญํ อาราเธติ, โน เจ มรณกาเล อญฺญํ อาราเธติ, อถ ปญฺจนฺนํ โอรมฺภาคิยานํ สํโยชนานํ ปริกฺขยา อนฺตราปรินิพฺพายี โหติ.\csromanspacer{} อุปหจฺจปรินิพฺพายี โหติ.\csromanspacer{} อสงฺขารปรินิพฺพายี โหติ.\csromanspacer{} สสงฺขารปรินิพฺพายี โหติ.\csromanspacer{} อุทฺธํโสโต โหติ อกนิฏฺฐคามี.\csromanspacer{} อิเมสํ โข ภิกฺขเว ปญฺจนฺนํ อินฺทฺริยานํ ภาวิตตฺตา พหุลีกตตฺตา อิเม สตฺต ผลา สตฺตานิสํสา ปาฏิกงฺขาติ.\csromanspacer{}\csromanspacer{}ฉฏฺฐํ.}

\csromanheader{2}{7. ปฐมรุกฺขสุตฺต}
\proseitem{537}{เสยฺยถาปิ ภิกฺขเว เย เกจิ ชมฺพุทีปกา รุกฺขา, ชมฺพู เตสํ อคฺคมกฺขายติ.\csromanspacer{} เอวเมว โข ภิกฺขเว เย เกจิ โพธิปกฺขิยา ธมฺมา, ปญฺญินฺทฺริยํ เตสํ อคฺคมกฺขายติ ยทิทํ โพธาย.\csromanspacer{} กตเม จ ภิกฺขเว โพธิปกฺขิยา ธมฺมา, สทฺธินฺทฺริยํ ภิกฺขเว โพธิปกฺขิโย ธมฺโม, ตํ โพธาย สํวตฺตติ.\csromanspacer{} วีริยินฺทฺริยํ ฯเปฯ.\csromanspacer{} สตินฺทฺริยํ.\csromanspacer{} สมาธินฺทฺริยํ.\csromanspacer{} ปญฺญินฺทฺริยํ โพธิปกฺขิโย ธมฺโม, ตํ โพธาย สํวตฺตติ.\csromanspacer{} เสยฺยถาปิ ภิกฺขเว เย เกจิ ชมฺพุทีปกา รุกฺขา, ชมฺพู เตสํ อคฺคมกฺขายติ.\csromanspacer{} เอวเมว โข ภิกฺขเว เย เกจิ โพธิปกฺขิยา ธมฺมา, ปญฺญินฺทฺริยํ เตสํ อคฺคมกฺขายติ ยทิทํ โพธายาติ.\csromanspacer{}\csromanspacer{}สตฺตมํ.}

\csromanheader{2}{8. ทุติยรุกฺขสุตฺต}
\proseitem{538}{เสยฺยถาปิ ภิกฺขเว เย เกจิ เทวานํ ตาวติํสานํ รุกฺขา, ปาริฉตฺตโก เตสํ อคฺคมกฺขายติ.\csromanspacer{} เอวเมว โข ภิกฺขเว เย เกจิ โพธิปกฺขิยา ธมฺมา, ปญฺญินฺทฺริยํ เตสํ อคฺคมกฺขายติ ยทิทํ โพธาย.\csromanspacer{} กตเม จ ภิกฺขเว โพธิปกฺขิยา ธมฺมา, สทฺธินฺทฺริยํ ภิกฺขเว โพธิปกฺขิโย ธมฺโม, ตํ โพธาย สํวตฺตติ.\csromanspacer{} วีริยินฺทฺริยํ ฯเปฯ.\csromanspacer{} สตินฺทฺริยํ.\csromanspacer{} สมาธินฺทฺริยํ.\csromanspacer{} ปญฺญินฺทฺริยํ โพธิปกฺขิโย ธมฺโม, ตํ โพธาย สํวตฺตติ.\csromanspacer{} เสยฺยถาปิ}

\csromanheader{2}{4. อินฺทฺริยสํยุตฺต}
\prose{\csromanfolio{209}ภิกฺขเว เย เกจิ เทวานํ ตาวติํสานํ รุกฺขา, ปาริฉตฺตโก เตสํ อคฺคมกฺขายติ.\csromanspacer{} เอวเมว โข ภิกฺขเว เย เกจิ โพธิปกฺขิยา ธมฺมา, ปญฺญินฺทฺริยํ เตสํ อคฺคมกฺขายติ ยทิทํ โพธายาติ.\csromanspacer{}\csromanspacer{}อฏฺฐมํ.}

\csromanheader{2}{9. ตติยรุกฺขสุตฺต}
\proseitem{539}{เสยฺยถาปิ ภิกฺขเว เย เกจิ อสุรานํ รุกฺขา, จิตฺตปาฏลิ เตสํ อคฺคมกฺขายติ.\csromanspacer{} เอวเมว โข ภิกฺขเว เย เกจิ โพธิปกฺขิยา ธมฺมา, ปญฺญินฺทฺริยํ เตสํ อคฺคมกฺขายติ ยทิทํ โพธาย.\csromanspacer{} กตเม จ ภิกฺขเว โพธิปกฺขิยา ธมฺมา, สทฺธินฺทฺริยํ ภิกฺขเว โพธิปกฺขิโย ธมฺโม, ตํ โพธาย สํวตฺตติ ฯเปฯ.\csromanspacer{} ปญฺญินฺทฺริยํ โพธิปกฺขิโย ธมฺโม, ตํ โพธาย สํวตฺตติ.\csromanspacer{} เสยฺยถาปิ ภิกฺขเว เย เกจิ อสุรานํ รุกฺขา, จิตฺตปาฏลิ เตสํ อคฺคมกฺขายติ.\csromanspacer{} เอวเมว โข ภิกฺขเว เย เกจิ โพธิปกฺขิยา ธมฺมา, ปญฺญินฺทฺริยํ เตสํ อคฺคมกฺขายติ ยทิทํ โพธายาติ.\csromanspacer{}\csromanspacer{}นวมํ.}

\csromanheader{2}{10. จตุตฺถรุกฺขสุตฺต}
\proseitem{540}{เสยฺยถาปิ ภิกฺขเว เย เกจิ สุปณฺณานํ รุกฺขา, กูฏสิมฺพลี\footnote{โกฏสิมฺพลิ (สฺยา, กํ)} เตสํ อคฺคมกฺขายติ.\csromanspacer{} เอวเมว โข ภิกฺขเว เย เกจิ โพธิปกฺขิยา ธมฺมา, ปญฺญินฺทฺริยํ เตสํ อคฺคมกฺขายติ ยทิทํ โพธาย.\csromanspacer{} กตเม จ ภิกฺขเว โพธิปกฺขิยา ธมฺมา, สทฺธินฺทฺริยํ ภิกฺขเว โพธิปกฺขิโย ธมฺโม, ตํ โพธาย สํวตฺตติ ฯเปฯ.\csromanspacer{} ปญฺญินฺทฺริยํ โพธิปกฺขิโย ธมฺโม, ตํ โพธาย สํวตฺตติ.\csromanspacer{} เสยฺยถาปิ ภิกฺขเว เย เกจิ สุปณฺณานํ รุกฺขา, กูฏสิมฺพลี เตสํ อคฺคมกฺขายติ.\csromanspacer{} เอวเมว โข ภิกฺขเว เย เกจิ โพธิปกฺขิยา ธมฺมา, ปญฺญินฺทฺริยํ เตสํ อคฺคมกฺขายติ ยทิทํ โพธายาติ.\csromanspacer{}\csromanspacer{}ทสมํ.\csromanspacer{} โพธิปกฺขิยวคฺโค สตฺตโม.}

\titlehead{\textbf{ตสฺสุทฺทานํ}}
\csromangathameasure{สํโยชนา อนุสยา, ปริญฺญา อาสวกฺขยา. \\ เทฺว ผลา จตุโร รุกฺขา, วคฺโค เตน ปวุจฺจตีติ.}
\csromangathagroup{\csromangathastanza{สํโยชนา อนุสยา, ปริญฺญา อาสวกฺขยา. \\ เทฺว ผลา จตุโร รุกฺขา, วคฺโค เตน ปวุจฺจตีติ.}}

\chapterheadpageread{8. คงฺคาเปยฺยาลวคฺค}
\csromanmatikaanchor{mtk.135}
\csromanmatikaanchor{mtk.136}
\csromantocmarknopage{subparagraph}{8. คงฺคาเปยฺยาลวคฺค}{mtk.135}
\bookmark[dest={mtk.135},level=5]{8. คงฺคาเปยฺยาลวคฺค}
\csromantocmarkref{subparagraph}{1-12. ปาจีนาทิสุตฺต}{mtk.136}
\bookmark[dest={mtk.136},level=5]{1-12. ปาจีนาทิสุตฺต}
\csromanheader{6}{1-12. ปาจีนาทิสุตฺต}
\proseitem{541-552}{\csromanfolio{210}เสยฺยถาปิ ภิกฺขเว คงฺคา นที ปาจีนนินฺนา ปาจีนโปณา ปาจีนปพฺภารา.\csromanspacer{} เอวเมว โข ภิกฺขเว ภิกฺขุ ปญฺจินฺทฺริยานิ ภาเวนฺโต ปญฺจินฺทฺริยานิ พหุลีกโรนฺโต นิพฺพานนินฺโน โหติ นิพฺพานโปโณ นิพฺพานปพฺภาโร.\csromanspacer{} กถญฺจ ภิกฺขเว ภิกฺขุ ปญฺจินฺทฺริยานิ ภาเวนฺโต ปญฺจินฺทฺริยานิ พหุลีกโรนฺโต นิพฺพานนินฺโน โหติ นิพฺพานโปโณ นิพฺพานปพฺภาโร, อิธ ภิกฺขเว ภิกฺขุ สทฺธินฺทฺริยํ ภาเวติ วิเวกนิสฺสิตํ วิราคนิสฺสิตํ นิโรธนิสฺสิตํ โวสฺสคฺคปริณามิํ.\csromanspacer{} วีริยินฺทฺริยํ ฯเปฯ.\csromanspacer{} สตินฺทฺริยํ.\csromanspacer{} สมาธินฺทฺริยํ.\csromanspacer{} ปญฺญินฺทฺริยํ ภาเวติ วิเวกนิสฺสิตํ วิราคนิสฺสิตํ นิโรธนิสฺสิตํ โวสฺสคฺคปริณามิํ.\csromanspacer{} เอวํ โข ภิกฺขเว ภิกฺขุ ปญฺจินฺทฺริยานิ ภาเวนฺโต ปญฺจินฺทฺริยานิ พหุลีกโรนฺโต นิพฺพานนินฺโน โหติ นิพฺพานโปโณ นิพฺพานปพฺภาโรติ.\csromanspacer{}\csromanspacer{}ทฺวาทสมํ.\csromanspacer{} คงฺคาเปยฺยาลวคฺโค อฏฺฐโม.}

\titlehead{\textbf{ตสฺสุทฺทานํ}}
\csromangathameasure{ฉ ปาจีนโต นินฺนา, ฉ นินฺนา จ สมุทฺทโต. \\ เทฺวเต ฉ ทฺวาทส โหนฺติ, วคฺโค เตน ปวุจฺจตีติ. อปฺปมาทวคฺโค วิตฺถาเรตพฺโพ.}
\csromangathagroup{\csromangathastanza{ฉ ปาจีนโต นินฺนา, ฉ นินฺนา จ สมุทฺทโต. \\ เทฺวเต ฉ ทฺวาทส โหนฺติ, วคฺโค เตน ปวุจฺจตีติ. อปฺปมาทวคฺโค วิตฺถาเรตพฺโพ.}}

\titlehead{\textbf{ตสฺสุทฺทานํ}}
\csromangathameasure{ตถาคตํ ปทํ กูฏํ, มูลํ สาเรน วสฺสิกํ. \\ ราชา จนฺทิมสูริยา, วตฺเถน ทสมํ ปทนฺติ. พลกรณียวคฺโค วิตฺถาเรตพฺโพ.}
\csromangathagroup{\csromangathastanza{ตถาคตํ ปทํ กูฏํ, มูลํ สาเรน วสฺสิกํ. \\ ราชา จนฺทิมสูริยา, วตฺเถน ทสมํ ปทนฺติ. พลกรณียวคฺโค วิตฺถาเรตพฺโพ.}}

\titlehead{\textbf{ตสฺสุทฺทานํ}}
\csromangathameasure{พลํ พีชญฺจ นาโค จ, รุกฺโข กุมฺเภน สูกิยา. \\ อากาเสน จ เทฺว เมฆา, นาวา อาคนฺตุกา นทีติ.}
\csromangathagroup{\csromangathastanza{พลํ พีชญฺจ นาโค จ, รุกฺโข กุมฺเภน สูกิยา. \\ อากาเสน จ เทฺว เมฆา, นาวา อาคนฺตุกา นทีติ.}}

\csromanheader{2}{4. อินฺทฺริยสํยุตฺต}
\vspace{\tipitakaparskip}
\csromancenter{เอสนาวคฺโค วิตฺถาเรตพฺโพ.}
\titlehead{\textbf{ตสฺสุทฺทานํ}}
\csromangathameasure{เอสนา วิธา อาสโว, ภโว จ ทุกฺขตา ติสฺโส. \\ ขิลํ มลญฺจ นีโฆ จ, เวทนา ตณฺหา ตสินา จาติ.}
\csromangathagroup{\csromangathastanza{\csromanfolio{211}เอสนา วิธา อาสโว, ภโว จ ทุกฺขตา ติสฺโส. \\ ขิลํ มลญฺจ นีโฆ จ, เวทนา ตณฺหา ตสินา จาติ.}}

\chapterhead{12. โอฆวคฺค}
\csromanmatikaanchor{mtk.137}
\csromanmatikaanchor{mtk.138}
\csromantocmarknopage{subparagraph}{12. โอฆวคฺค}{mtk.137}
\bookmark[dest={mtk.137},level=5]{12. โอฆวคฺค}
\csromantocmarkref{subparagraph}{1-10. โอฆาทิสุตฺต}{mtk.138}
\bookmark[dest={mtk.138},level=5]{1-10. โอฆาทิสุตฺต}
\csromanheader{6}{1-10. โอฆาทิสุตฺต}
\proseitem{587-596}{ปญฺจิมานิ ภิกฺขเว อุทฺธมฺภาคิยานิ สํโยชนานิ.\csromanspacer{} กตมานิ ปญฺจ, รูปราโค อรูปราโค มาโน อุทฺธจฺจํ อวิชฺชา, อิมานิ โข ภิกฺขเว ปญฺจุทฺธมฺภาคิยานิ สํโยชนานิ.\csromanspacer{} อิเมสํ โข ภิกฺขเว ปญฺจนฺนํ อุทฺธมฺภาคิยานํ สํโยชนานํ อภิญฺญาย ปริญฺญาย ปริกฺขยาย ปหานาย ปญฺจินฺทฺริยานิ ภาเวตพฺพานิ.\csromanspacer{} กตมานิ ปญฺจ, อิธ ภิกฺขเว ภิกฺขุ สทฺธินฺทฺริยํ ภาเวติ วิเวกนิสฺสิตํ วิราคนิสฺสิตํ นิโรธนิสฺสิตํ โวสฺสคฺคปริณามิํ ฯเปฯ.\csromanspacer{} ปญฺญินฺทฺริยํ ภาเวติ วิเวกนิสฺสิตํ วิราคนิสฺสิตํ นิโรธนิสฺสิตํ โวสฺสคฺคปริณามิํ.\csromanspacer{} อิเมสํ โข ภิกฺขเว ปญฺจนฺนํ อุทฺธมฺภาคิยานํ สํโยชนานํ อภิญฺญาย ปริญฺญาย ปริกฺขยาย ปหานาย อิมานิ ปญฺจินฺทฺริยานิ ภาเวตพฺพานีติ.\csromanspacer{}\csromanspacer{}ทสมํ.\csromanspacer{}\csromanspacer{}(ยถา มคฺคสํยุตฺตํ, ตถา วิตฺถาเรตพฺพํ.) โอฆวคฺโค ทฺวาทสโม.}

\titlehead{\textbf{ตสฺสุทฺทานํ}}
\csromangathameasure{โอโฆ โยโค อุปาทานํ, คนฺถา อนุสเยน จ. \\ กามคุณา นีวรณา, ขนฺธา โอรุทฺธมฺภาคิยาติ.}
\csromangathagroup{\csromangathastanza{โอโฆ โยโค อุปาทานํ, คนฺถา อนุสเยน จ. \\ กามคุณา นีวรณา, ขนฺธา โอรุทฺธมฺภาคิยาติ.}}

\chapterheadpageread{13. คงฺคาเปยฺยาลวคฺค}
\csromanmatikaanchor{mtk.139}
\csromanmatikaanchor{mtk.140}
\csromantocmarknopage{subparagraph}{13. คงฺคาเปยฺยาลวคฺค}{mtk.139}
\bookmark[dest={mtk.139},level=5]{13. คงฺคาเปยฺยาลวคฺค}
\csromantocmarkref{subparagraph}{1-12. ปาจีนาทิสุตฺต}{mtk.140}
\bookmark[dest={mtk.140},level=5]{1-12. ปาจีนาทิสุตฺต}
\csromanheader{6}{1-12. ปาจีนาทิสุตฺต}
\proseitem{597-608}{\csromanfolio{212}เสยฺยถาปิ ภิกฺขเว คงฺคา นที ปาจีนนินฺนา ปาจีนโปณา ปาจีนปพฺภารา.\csromanspacer{} เอวเมว โข ภิกฺขเว ภิกฺขุ ปญฺจินฺทฺริยานิ ภาเวนฺโต ปญฺจินฺทฺริยานิ พหุลีกโรนฺโต นิพฺพานนินฺโน โหติ นิพฺพานโปโณ นิพฺพานปพฺภาโร.\csromanspacer{} กถญฺจ ภิกฺขเว ภิกฺขุ ปญฺจินฺทฺริยานิ ภาเวนฺโต ปญฺจินฺทฺริยานิ พหุลีกโรนฺโต นิพฺพานนินฺโน โหติ นิพฺพานโปโณ นิพฺพานปพฺภาโร, อิธ ภิกฺขเว ภิกฺขุ สทฺธินฺทฺริยํ ภาเวติ ราควินยปริโยสานํ โทสวินยปริโยสานํ โมหวินยปริโยสานํ ฯเปฯ.\csromanspacer{} ปญฺญินฺทฺริยํ ภาเวติ ราควินยปริโยสานํ โทสวินยปริโยสานํ โมหวินยปริโยสานํ.\csromanspacer{} เอวํ โข ภิกฺขเว ภิกฺขุ ปญฺจินฺทฺริยานิ ภาเวนฺโต ปญฺจินฺทฺริยานิ พหุลีกโรนฺโต นิพฺพานนินฺโน โหติ นิพฺพานโปโณ นิพฺพานปพฺภาโรติ.\csromanspacer{}\csromanspacer{}ทฺวาทสมํ.\csromanspacer{} คงฺคาเปยฺยาลวคฺโค เตรสโม.}

\titlehead{\textbf{ตสฺสุทฺทานํ}}
\csromangathameasure{ฉ ปาจีนโต นินฺนา, ฉ นินฺนา จ สมุทฺทโต. \\ เทฺวเต ฉ ทฺวาทส โหนฺติ, วคฺโค เตน ปวุจฺจตีติ. (อปฺปมาทวคฺค พลกรณียวคฺค เอสนาวคฺคา วิตฺถาเรตพฺพา.)}
\csromangathagroup{\csromangathastanza{ฉ ปาจีนโต นินฺนา, ฉ นินฺนา จ สมุทฺทโต. \\ เทฺวเต ฉ ทฺวาทส โหนฺติ, วคฺโค เตน ปวุจฺจตีติ. (อปฺปมาทวคฺค พลกรณียวคฺค เอสนาวคฺคา วิตฺถาเรตพฺพา.)}}

\chapterhead{17. โอฆวคฺค}
\csromanmatikaanchor{mtk.141}
\csromanmatikaanchor{mtk.142}
\csromantocmarknopage{subparagraph}{17. โอฆวคฺค}{mtk.141}
\bookmark[dest={mtk.141},level=5]{17. โอฆวคฺค}
\csromantocmarkref{subparagraph}{1-10. โอฆาทิสุตฺต}{mtk.142}
\bookmark[dest={mtk.142},level=5]{1-10. โอฆาทิสุตฺต}
\csromanheader{6}{1-10. โอฆาทิสุตฺต}
\proseitem{641-650}{ปญฺจิมานิ ภิกฺขเว อุทฺธมฺภาคิยานิ สํโยชนานิ.\csromanspacer{} กตมานิ ปญฺจ, รูปราโค อรูปราโค มาโน อุทฺธจฺจํ อวิชฺชา, อิมานิ โข ภิกฺขเว ปญฺจุทฺธมฺภาคิยานิ สํโยชนานิ, อิเมสํ โข ภิกฺขเว ปญฺจนฺนํ อุทฺธมฺภาคิยานํ สํโยชนานํ อภิญฺญาย ปริญฺญาย ปริกฺขยาย ปหานาย ปญฺจินฺทฺริยานิ ภาเวตพฺพานิ.\csromanspacer{} กตมานิ ปญฺจ, อิธ ภิกฺขเว ภิกฺขุ สทฺธินฺทฺริยํ ภาเวติ ราควินยปริโยสานํ โทสวินยปริโยสานํ โมหวินยปริโยสานํ.}

\csromanheader{2}{4. อินฺทฺริยสํยุตฺต}
\prose{\csromanfolio{213}วีริยินฺทฺริยํ ฯเปฯ.\csromanspacer{} สตินฺทฺริยํ.\csromanspacer{} สมาธินฺทฺริยํ.\csromanspacer{} ปญฺญินฺทฺริยํ ภาเวติ ราควินยปริโยสานํ โทสวินยปริโยสานํ โมหวินยปริโยสานํ.\csromanspacer{} อิเมสํ โข ภิกฺขเว ปญฺจนฺนํ อุทฺธมฺภาคิยานํ สํโยชนานํ อภิญฺญาย ปริญฺญาย ปริกฺขยาย ปหานาย อิมานิ ปญฺจินฺทฺริยานิ ภาเวตพฺพานีติ.\csromanspacer{} โอฆวคฺโค สตฺตรสโม.}

\titlehead{\textbf{ตสฺสุทฺทานํ}}
\csromangathameasure{โอโฆ โยโค อุปาทานํ, คนฺถา อนุสเยน จ. \\ กามคุณา นีวรณา, ขนฺธา โอรุทฺธมฺภาคิยาติ. \\ \textbf{อินฺทฺริยสํยุตฺตํ จตุตฺถํ.}}
\csromangathagroup{\csromangathastanza{โอโฆ โยโค อุปาทานํ, คนฺถา อนุสเยน จ. \\ กามคุณา นีวรณา, ขนฺธา โอรุทฺธมฺภาคิยาติ. \\ \textbf{อินฺทฺริยสํยุตฺตํ จตุตฺถํ.}}}

\csromanmatikaanchor{mtk.143}
\csromantocmarknopage{subparagraph}{5. สมฺมปฺปธานสํยุตฺต}{mtk.143}
\bookmark[dest={mtk.143},level=5]{5. สมฺมปฺปธานสํยุตฺต}
\csromanheader{6}{5. สมฺมปฺปธานสํยุตฺต}
\chapterhead{1. คงฺคาเปยฺยาลวคฺค}
\csromanmatikaanchor{mtk.144}
\csromanmatikaanchor{mtk.145}
\csromantocmarknopage{subparagraph}{1. คงฺคาเปยฺยาลวคฺค}{mtk.144}
\bookmark[dest={mtk.144},level=5]{1. คงฺคาเปยฺยาลวคฺค}
\csromantocmarkref{subparagraph}{1-12. ปาจีนาทิสุตฺต}{mtk.145}
\bookmark[dest={mtk.145},level=5]{1-12. ปาจีนาทิสุตฺต}
\csromanheader{6}{1-12. ปาจีนาทิสุตฺต}
\proseitem{651-662}{\csromanfolio{214}สาวตฺถินิทานํ.\csromanspacer{} ตตฺร โข ภควา เอตทโวจ\csromandash{}จตฺตาโรเม ภิกฺขเว สมฺมปฺปธานา.\csromanspacer{} กตเม จตฺตาโร, อิธ ภิกฺขเว ภิกฺขุ อนุปฺปนฺนานํ ปาปกานํ อกุสลานํ ธมฺมานํ อนุปฺปาทาย ฉนฺทํ ชเนติ วายมติ, วีริยํ อารภติ, จิตฺตํ ปคฺคณฺหาติ ปทหติ, อุปฺปนฺนานํ ปาปกานํ อกุสลานํ ธมฺมานํ ปหานาย ฉนฺทํ ชเนติ วายมติ, วีริยํ อารภติ, จิตฺตํ ปคฺคณฺหาติ ปทหติ, อนุปฺปนฺนานํ กุสลานํ ธมฺมานํ อุปฺปาทาย ฉนฺทํ ชเนติ วายมติ, วีริยํ อารภติ, จิตฺตํ ปคฺคณฺหาติ ปทหติ, อุปฺปนฺนานํ กุสลานํ ธมฺมานํ ฐิติยา อสมฺโมสาย ภิยฺโยภาวาย เวปุลฺลาย ภาวนาย ปาริปูริยา ฉนฺทํ ชเนติ วายมติ, วีริยํ อารภติ, จิตฺตํ ปคฺคณฺหาติ ปทหติ.\csromanspacer{} อิเม โข ภิกฺขเว จตฺตาโร สมฺมปฺปธานาติ.}

\prose{เสยฺยถาปิ ภิกฺขเว คงฺคา นที ปาจีนนินฺนา ปาจีนโปณา ปาจีนปพฺภารา.\csromanspacer{} เอวเมว โข ภิกฺขเว ภิกฺขุ จตฺตาโร สมฺมปฺปธาเน ภาเวนฺโต จตฺตาโร สมฺมปฺปธาเน พหุลีกโรนฺโต นิพฺพานนินฺโน โหติ นิพฺพานโปโณ นิพฺพานปพฺภาโร.\csromanspacer{} กถญฺจ ภิกฺขเว ภิกฺขุ จตฺตาโร สมฺมปฺปธาเน ภาเวนฺโต จตฺตาโร สมฺมปฺปธาเน พหุลีกโรนฺโต นิพฺพานนินฺโน โหติ นิพฺพานโปโณ นิพฺพานปพฺภาโร, อิธ ภิกฺขเว ภิกฺขุ อนุปฺปนฺนานํ ปาปกานํ อกุสลานํ ธมฺมานํ อนุปฺปาทาย ฉนฺทํ ชเนติ, วายมติ วีริยํ อารภติ, จิตฺตํ ปคฺคณฺหาติ ปทหติ, อุปฺปนฺนานํ ปาปกานํ อกุสลานํ ธมฺมานํ ปหานาย ฉนฺทํ ชเนติ วายมติ, วีริยํ อารภติ, จิตฺตํ ปคฺคณฺหาติ, ปทหติ, อนุปฺปนฺนานํ กุสลานํ ธมฺมานํ อุปฺปาทาย ฉนฺทํ ชเนติ วายมติ, วีริยํ อารภติ, จิตฺตํ ปคฺคณฺหาติ ปทหติ, อุปฺปนฺนานํ กุสลานํ ธมฺมานํ ฐิติยา อสมฺโมสาย ภิยฺโยภาวาย เวปุลฺลาย ภาวนาย ปาริปูริยา ฉนฺทํ ชเนติ วายมติ, วีริยํ อารภติ, จิตฺตํ ปคฺคณฺหาติ ปทหติ.\csromanspacer{} เอวํ โข ภิกฺขเว ภิกฺขุ จตฺตาโร สมฺมปฺปธาเน ภาเวนฺโต จตฺตาโร สมฺมปฺปธาเน พหุลีกโรนฺโต นิพฺพานนินฺโน โหติ นิพฺพานโปโณ นิพฺพานปพฺภาโรติ.\csromanspacer{}\csromanspacer{}ทฺวาทสมํ.}

\csromanheader{2}{5. สมฺมปฺปธานสํยุตฺต}
\vspace{\tipitakaparskip}
\csromancenter{คงฺคาเปยฺยาลวคฺโค ปฐโม.}
\titlehead{\textbf{ตสฺสุทฺทานํ}}
\csromanmatikaanchor{mtk.146}
\csromanmatikaanchor{mtk.147}
\csromantocmarknopage{subparagraph}{2. อปฺปมาทวคฺค}{mtk.146}
\bookmark[dest={mtk.146},level=5]{2. อปฺปมาทวคฺค}
\csromantocmarkref{section}{อุทฺทานคาถา}{mtk.147}
\bookmark[dest={mtk.147},level=1]{อุทฺทานคาถา}
\csromangathameasure{ฉ ปาจีนโต นินฺนา, ฉ นินฺนา จ สมุทฺทโต. \\ เทฺวเต ฉ ทฺวาทส โหนฺติ, วคฺโค เตน ปวุจฺจตีติ. (สมฺมปฺปธานสํยุตฺตสฺส คงฺคาเปยฺยาลี สมฺมปฺปธานวเสน วิตฺถาเรตพฺพา.)}
\csromangathagroup{\csromangathastanza{\csromanfolio{215}ฉ ปาจีนโต นินฺนา, ฉ นินฺนา จ สมุทฺทโต. \\ เทฺวเต ฉ ทฺวาทส โหนฺติ, วคฺโค เตน ปวุจฺจตีติ. (สมฺมปฺปธานสํยุตฺตสฺส คงฺคาเปยฺยาลี สมฺมปฺปธานวเสน วิตฺถาเรตพฺพา.)}}

\chapterhead{2. อปฺปมาทวคฺค}
\titlehead{\textbf{ตสฺสุทฺทานํ}}
\csromangathameasure{ตถาคตํ ปทํ กูฏํ, มูลํ สาเรน วสฺสิกํ. \\ ราชา จนฺทิมสูริยา, วตฺเถน ทสมํ ปทนฺติ. (อปฺปมาทวคฺโค สมฺมปฺปธานวเสน วิตฺถาเรตพฺโพ.)}
\csromangathagroup{\csromangathastanza{ตถาคตํ ปทํ กูฏํ, มูลํ สาเรน วสฺสิกํ. \\ ราชา จนฺทิมสูริยา, วตฺเถน ทสมํ ปทนฺติ. (อปฺปมาทวคฺโค สมฺมปฺปธานวเสน วิตฺถาเรตพฺโพ.)}}

\chapterhead{3. พลกรณียวคฺค}
\csromanmatikaanchor{mtk.148}
\csromanmatikaanchor{mtk.149}
\csromantocmarknopage{subparagraph}{3. พลกรณียวคฺค}{mtk.148}
\bookmark[dest={mtk.148},level=5]{3. พลกรณียวคฺค}
\csromantocmarkref{subparagraph}{1-12. พลกรณียาทิสุตฺต}{mtk.149}
\bookmark[dest={mtk.149},level=5]{1-12. พลกรณียาทิสุตฺต}
\csromanheader{6}{1-12. พลกรณียาทิสุตฺต}
\proseitem{673-684}{เสยฺยถาปิ ภิกฺขเว เย เกจิ พลกรณียา กมฺมนฺตา กยิรนฺติ, สพฺเพเต ปถวิํ นิสฺสาย ปถวิยํ ปติฏฺฐาย, เอวเมเต พลกรณียา กมฺมนฺตา กยิรนฺติ.\csromanspacer{} เอวเมว โข ภิกฺขเว ภิกฺขุ สีลํ นิสฺสาย สีเล ปติฏฺฐาย จตฺตาโร สมฺมปฺปธาเน ภาเวติ จตฺตาโร สมฺมปฺปธาเน พหุลีกโรติ.\csromanspacer{} กถญฺจ ภิกฺขเว ภิกฺขุ สีลํ นิสฺสาย สีเล ปติฏฺฐาย จตฺตาโร สมฺมปฺปธาเน ภาเวติ จตฺตาโร สมฺมปฺปธาเน พหุลีกโรติ, อิธ ภิกฺขเว ภิกฺขุ อนุปฺปนฺนานํ ปาปกานํ อกุสลานํ ธมฺมานํ อนุปฺปาทาย ฉนฺทํ ชเนติ วายมติ, วีริยํ อารภติ, จิตฺตํ ปคฺคณฺหาติ ปทหติ ฯเปฯ อุปฺปนฺนานํ กุสลานํ ธมฺมานํ ฐิติยา อสมฺโมสาย ภิยฺโยภาวาย เวปุลฺลาย ภาวนาย ปาริปูริยา ฉนฺทํ ชเนติ วายมติ, วีริยํ อารภติ, จิตฺตํ ปคฺคณฺหาติ ปทหติ.\csromanspacer{} เอวํ โข ภิกฺขเว ภิกฺขุ สีลํ \csromanfolio{216}นิสฺสาย สีเล ปติฏฺฐาย จตฺตาโร สมฺมปฺปธาเน ภาเวติ จตฺตาโร สมฺมปฺปธาเน พหุลีกโรตีติ. (เอวํ พลกรณียวคฺโค สมฺมปฺปธานวเสน วิตฺถาเรตพฺโพ.).\csromanspacer{} ทฺวาทสมํ.\csromanspacer{} พลกรณียวคฺโค ตติโย.}

\csromanheader{2}{\textbf{ตสฺสุทานํ}}
\csromangathameasure{พลํ พีชญฺจ นาโค จ, รุกฺโข กุมฺเภน สูกิยา. \\ อากาเสน จ เทฺว เมฆา, นาวา อาคนฺตุกา นทีติ.}
\csromangathagroup{\csromangathastanza{พลํ พีชญฺจ นาโค จ, รุกฺโข กุมฺเภน สูกิยา.}\csromangathastanza{อากาเสน จ เทฺว เมฆา, นาวา อาคนฺตุกา นทีติ.}}

\chapterhead{4. เอสนาวคฺค}
\csromanmatikaanchor{mtk.150}
\csromanmatikaanchor{mtk.151}
\csromantocmarknopage{subparagraph}{4. เอสนาวคฺค}{mtk.150}
\bookmark[dest={mtk.150},level=5]{4. เอสนาวคฺค}
\csromantocmarkref{subparagraph}{1-10. เอสนาทิสุตฺตทสก}{mtk.151}
\bookmark[dest={mtk.151},level=5]{1-10. เอสนาทิสุตฺตทสก}
\csromanheader{6}{1-10. เอสนาทิสุตฺตทสก}
\proseitem{685-694}{ติสฺโส อิมา ภิกฺขเว เอสนา.\csromanspacer{} กตมา ติสฺโส, กาเมสนา ภเวสนา พฺรหฺมจริเยสนา.\csromanspacer{} อิมา โข ภิกฺขเว ติสฺโส เอสนา.\csromanspacer{} อิมาสํ โข ภิกฺขเว ติสฺสนฺนํ เอสนานํ อภิญฺญาย ปริญฺญาย ปริกฺขยาย ปหานาย จตฺตาโร สมฺมปฺปธานา ภาเวตพฺพา.\csromanspacer{} กตเม จตฺตาโร, อิธ ภิกฺขเว ภิกฺขุ อนุปฺปนฺนานํ ฯเปฯ อุปฺปนฺนานํ กุสลานํ ธมฺมานํ ฐิติยา อสมฺโมสาย ภิยฺโยภาวาย เวปุลฺลาย ภาวนาย ปาริปูริยา ฉนฺทํ ชเนติ วายมติ, วีริยํ อารภติ, จิตฺตํ ปคฺคณฺหาติ ปทหติ.\csromanspacer{} อิมาสํ โข ภิกฺขเว ติสฺสนฺนํ เอสนานํ อภิญฺญาย ปริญฺญาย ปริกฺขยาย ปหานาย อิเม จตฺตาโร สมฺมปฺปธานา ภาเวตพฺพาติ.\csromanspacer{} วิตฺถาเรตพฺพํ.\csromanspacer{}\csromanspacer{}ทสมํ.\csromanspacer{} เอสนาวคฺโค จตุตฺโถ.}

\titlehead{\textbf{ตสฺสุทฺทานํ}}
\csromangathameasure{เอสนา วิธา อาสโว, ภโว จ ทุกฺขตา ติสฺโส. \\ ขิลํ มลญฺจ นีโฆ จ, เวทนา ตณฺหา ตสินา จาติ.}
\csromangathagroup{\csromangathastanza{เอสนา วิธา อาสโว, ภโว จ ทุกฺขตา ติสฺโส. \\ ขิลํ มลญฺจ นีโฆ จ, เวทนา ตณฺหา ตสินา จาติ.}}

\csromanheader{2}{5. สมฺมปฺปธานสํยุตฺต}
\chapterhead{5. โอฆวคฺค}
\csromanmatikaanchor{mtk.152}
\csromanmatikaanchor{mtk.153}
\csromantocmarknopage{subparagraph}{5. โอฆวคฺค}{mtk.152}
\bookmark[dest={mtk.152},level=5]{5. โอฆวคฺค}
\csromantocmarkref{subparagraph}{1-10. โอฆาทิสุตฺต}{mtk.153}
\bookmark[dest={mtk.153},level=5]{1-10. โอฆาทิสุตฺต}
\csromanheader{6}{1-10. โอฆาทิสุตฺต}
\proseitem{695-704}{\csromanfolio{217}ปญฺจิมานิ ภิกฺขเว อุทฺธมฺภาคิยานิ สํโยชนานิ.\csromanspacer{} กตมานิ ปญฺจ, รูปราโค อรูปราโค มาโน อุทฺธจฺจํ อวิชฺชา, อิมานิ โข ภิกฺขเว ปญฺจุทฺธมฺภาคิยานิ สํโยชนานิ.\csromanspacer{} อิเมสํ โข ภิกฺขเว ปญฺจนฺนํ อุทฺธมฺภาคิยานํ สํโยชนานํ อภิญฺญาย ปริญฺญาย ปริกฺขยาย ปหานาย จตฺตาโร สมฺมปฺปธานา ภาเวตพฺพา.\csromanspacer{} กตเม จตฺตาโร, อิธ ภิกฺขเว ภิกฺขุ อนุปฺปนฺนานํ ฯเปฯ อุปฺปนฺนานํ กุสลานํ ธมฺมานํ ฐิติยา อสมฺโมสาย ภิยฺโยภาวาย เวปุลฺลาย ภาวนาย ปาริปูริยา ฉนฺทํ ชเนติ วายมติ, วีริยํ อารภติ, จิตฺตํ ปคฺคณฺหาติ ปทหติ.\csromanspacer{} อิเมสํ โข ภิกฺขเว ปญฺจนฺนํ อุทฺธมฺภาคิยานํ สํโยชนานํ อภิญฺญาย ปริญฺญาย ปริกฺขยาย ปหานาย อิเม จตฺตาโร สมฺมปฺปธานา ภาเวตพฺพาติ.\csromanspacer{} วิตฺถาเรตพฺพา.\csromanspacer{} ทสมํ.\csromanspacer{} โอฆวคฺโค ปญฺจโม.}

\titlehead{\textbf{ตสฺสุทฺทานํ}}
\csromangathameasure{โอโฆ โยโค อุปาทานํ, คนฺถา อนุสเยน จ. \\ กามคุณา นีวรณา, ขนฺธา โอรุทฺธมฺภาคิยาติ. \\ \textbf{สมฺมปฺปธานสํยุตฺตํ ปญฺจมํ.}}
\csromangathagroup{\csromangathastanza{โอโฆ โยโค อุปาทานํ, คนฺถา อนุสเยน จ. \\ กามคุณา นีวรณา, ขนฺธา โอรุทฺธมฺภาคิยาติ. \\ \textbf{สมฺมปฺปธานสํยุตฺตํ ปญฺจมํ.}}}

\csromanmatikaanchor{mtk.154}
\csromantocmarknopage{subparagraph}{6. พลสํยุตฺต}{mtk.154}
\bookmark[dest={mtk.154},level=5]{6. พลสํยุตฺต}
\csromanheader{6}{6. พลสํยุตฺต}
\chapterhead{1. คงฺคาเปยฺยาลวคฺค}
\csromanmatikaanchor{mtk.155}
\csromanmatikaanchor{mtk.156}
\csromantocmarknopage{subparagraph}{1. คงฺคาเปยฺยาลวคฺค}{mtk.155}
\bookmark[dest={mtk.155},level=5]{1. คงฺคาเปยฺยาลวคฺค}
\csromantocmarkref{subparagraph}{1-12. พลาทิสุตฺต}{mtk.156}
\bookmark[dest={mtk.156},level=5]{1-12. พลาทิสุตฺต}
\csromanheader{6}{1-12. พลาทิสุตฺต}
\proseitem{705-716}{\csromanfolio{218}ปญฺจิมานิ ภิกฺขเว พลานิ.\csromanspacer{} กตมานิ ปญฺจ, สทฺธาพลํ วีริยพลํ สติพลํ สมาธิพลํ ปญฺญาพลํ.\csromanspacer{} อิมานิ โข ภิกฺขเว ปญฺจ พลานีติ.\csromanspacer{} เสยฺยถาปิ ภิกฺขเว คงฺคา นที ปาจีนนินฺนา ปาจีนโปณา ปาจีนปพฺภารา.\csromanspacer{} เอวเมว โข ภิกฺขเว ภิกฺขุ ปญฺจ พลานิ ภาเวนฺโต ปญฺจพลานิ พหุลีกโรนฺโต นิพฺพานนินฺโน โหติ นิพฺพานโปโณ นิพฺพานปพฺภาโร.\csromanspacer{} กถญฺจ ภิกฺขเว ภิกฺขุ ปญฺจ พลานิ ภาเวนฺโต ปญฺจ พลานิ พหุลีกโรนฺโต นิพฺพานนินฺโน โหติ นิพฺพานโปโณ นิพฺพานปพฺภาโร, อิธ ภิกฺขเว ภิกฺขุ สทฺธาพลํ ภาเวติ วิเวกนิสฺสิตํ วิราคนิสฺสิตํ นิโรธนิสฺสิตํ โวสฺสคฺคปริณามิํ.\csromanspacer{} วีริยพลํ ฯเปฯ.\csromanspacer{} สติพลํ.\csromanspacer{} สมาธิพลํ.\csromanspacer{} ปญฺญาพลํ ภาเวติ วิเวกนิสฺสิตํ วิราคนิสฺสิตํ นิโรธนิสฺสิตํ โวสฺสคฺคปริณามิํ.\csromanspacer{} เอวํ โข ภิกฺขเว ภิกฺขุ ปญฺจ พลานิ ภาเวนฺโต ปญฺจ พลานิ พหุลีกโรนฺโต นิพฺพานนินฺโน โหติ นิพฺพานโปโณ นิพฺพานปพฺภาโรติ.\csromanspacer{}\csromanspacer{}ทฺวาทสมํ.\csromanspacer{} คงฺคาเปยฺยาลวคฺโค ปฐโม.}

\titlehead{\textbf{ตสฺสุทฺทานํ}}
\csromangathameasure{ฉ ปาจีนโต นินฺนา, ฉ นินฺนา จ สมุทฺทโต. \\ เทฺวเต ฉ ทฺวาทส โหนฺติ, วคฺโค เตน ปวุจฺจตีติ. อปฺปมาทวคฺโค วิตฺถาเรตพฺโพ.}
\csromangathagroup{\csromangathastanza{ฉ ปาจีนโต นินฺนา, ฉ นินฺนา จ สมุทฺทโต. \\ เทฺวเต ฉ ทฺวาทส โหนฺติ, วคฺโค เตน ปวุจฺจตีติ. อปฺปมาทวคฺโค วิตฺถาเรตพฺโพ.}}

\titlehead{\textbf{ตสฺสุทฺทานํ}}
\csromangathameasure{ตถาคตํ ปทํ กูฏํ, มูลํ สาเรน วสฺสิกํ. \\ ราชา จนฺทิมสูริยา, วตฺเถน ทสมํ ปทนฺติ.}
\csromangathagroup{\csromangathastanza{ตถาคตํ ปทํ กูฏํ, มูลํ สาเรน วสฺสิกํ. \\ ราชา จนฺทิมสูริยา, วตฺเถน ทสมํ ปทนฺติ.}}

\csromanheader{2}{6. พลสํยุตฺต}
\vspace{\tipitakaparskip}
\csromancenter{พลกรณียวคฺโค วิตฺถาเรตพฺโพ.}
\titlehead{\textbf{ตสฺสุทฺทานํ}}
\csromangathameasure{พลํ พีชญฺจ นาโค จ, รุกฺโข กุมฺเภน สูกิยา. \\ อากาเสน จ เทฺว เมฆา, นาวา อาคนฺตุกา นทีติ. เอสนาวคฺโค วิตฺถาเรตพฺโพ.}
\csromangathagroup{\csromangathastanza{\csromanfolio{219}พลํ พีชญฺจ นาโค จ, รุกฺโข กุมฺเภน สูกิยา. \\ อากาเสน จ เทฺว เมฆา, นาวา อาคนฺตุกา นทีติ. เอสนาวคฺโค วิตฺถาเรตพฺโพ.}}

\titlehead{\textbf{ตสฺสุทฺทานํ}}
\csromangathameasure{เอสนา วิธา อาสโว, ภโว จ ทุกฺขตา ติสฺโส. \\ ขิลํ มลญฺจ นีโฆ จ, เวทนา ตณฺหา ตสินา จาติ.}
\csromangathagroup{\csromangathastanza{เอสนา วิธา อาสโว, ภโว จ ทุกฺขตา ติสฺโส. \\ ขิลํ มลญฺจ นีโฆ จ, เวทนา ตณฺหา ตสินา จาติ.}}

\chapterhead{5. โอฆวคฺค}
\csromanmatikaanchor{mtk.157}
\csromanmatikaanchor{mtk.158}
\csromantocmarknopage{subparagraph}{5. โอฆวคฺค}{mtk.157}
\bookmark[dest={mtk.157},level=5]{5. โอฆวคฺค}
\csromantocmarkref{subparagraph}{1-10. โอฆาทิสุตฺต}{mtk.158}
\bookmark[dest={mtk.158},level=5]{1-10. โอฆาทิสุตฺต}
\csromanheader{6}{1-10. โอฆาทิสุตฺต}
\proseitem{749-758}{ปญฺจิมานิ ภิกฺขเว อุทฺธมฺภาคิยานิ สํโยชนานิ.\csromanspacer{} กตมานิ ปญฺจ, รูปราโค อรูปราโค มาโน อุทฺธจฺจํ อวิชฺชา, อิมานิ โข ภิกฺขเว ปญฺจุทฺธมฺภาคิยานิ สํโยชนานิ.\csromanspacer{} อิเมสํ โข ภิกฺขเว ปญฺจนฺนํ อุทฺธมฺภาคิยานํ สํโยชนานํ อภิญฺญาย ปริญฺญาย ปริกฺขยาย ปหานาย ปญฺจ พลานิ ภาเวตพฺพานิ.\csromanspacer{} กตมานิ ปญฺจ, อิธ ภิกฺขเว ภิกฺขุ สทฺธาพลํ ภาเวติ วิเวกนิสฺสิตํ วิราคนิสฺสิตํ นิโรธนิสฺสิตํ โวสฺสคฺคปริณามิํ.\csromanspacer{} วีริยพลํ ฯเปฯ.\csromanspacer{} สติพลํ สมาธิพลํ ปญฺญาพลํ ภาเวติ วิเวกนิสฺสิตํ วิราคนิสฺสิตํ นิโรธนิสฺสิตํ โวสฺสคฺคปริณามิํ.\csromanspacer{} อิเมสํ โข ภิกฺขเว ปญฺจนฺนํ อุทฺธมฺภาคิยานํ สํโยชนานํ อภิญฺญาย ปริญฺญาย ปริกฺขยาย ปหานาย อิมานิ ปญฺจ พลานิ ภาเวตพฺพานีติ. (เอวํ วิตฺถาเรตพฺพา.)}

\chapterheadpageread{6. คงฺคาเปยฺยาลวคฺค}
\csromanmatikaanchor{mtk.159}
\csromanmatikaanchor{mtk.160}
\csromantocmarknopage{subparagraph}{6. คงฺคาเปยฺยาลวคฺค}{mtk.159}
\bookmark[dest={mtk.159},level=5]{6. คงฺคาเปยฺยาลวคฺค}
\csromantocmarkref{subparagraph}{1-12. ปาจีนาทิสุตฺต}{mtk.160}
\bookmark[dest={mtk.160},level=5]{1-12. ปาจีนาทิสุตฺต}
\csromanheader{6}{1-12. ปาจีนาทิสุตฺต}
\proseitem{759-770}{\csromanfolio{220}เสยฺยถาปิ ภิกฺขเว คงฺคา นที ปาจีนนินฺนา ปาจีนโปณา ปาจีนปพฺภารา.\csromanspacer{} เอวเมว โข ภิกฺขเว ภิกฺขุ ปญฺจ พลานิ ภาเวนฺโต ปญฺจ พลานิ พหุลีกโรนฺโต นิพฺพานนินฺโน โหติ นิพฺพานโปโณ นิพฺพานปพฺภาโร.\csromanspacer{} กถญฺจ ภิกฺขเว ภิกฺขุ ปญฺจ พลานิ ภาเวนฺโต ปญฺจ พลานิ พหุลีกโรนฺโต นิพฺพานนินฺโน โหติ นิพฺพานโปโณ นิพฺพานปพฺภาโร, อิธ ภิกฺขเว ภิกฺขุ สทฺธาพลํ ภาเวติ ราควินยปริโยสานํ โทสวินยปริโยสานํ โมหวินยปริโยสานํ.\csromanspacer{} เอวํ โข ภิกฺขเว ภิกฺขุ ปญฺจ พลานิ ภาเวนฺโต ปญฺจ พลานิ พหุลีกโรนฺโต นิพฺพานนินฺโน โหติ นิพฺพานโปโณ นิพฺพานปพฺภาโรติ.\csromanspacer{} วิตฺถาเรตพฺพา.\csromanspacer{}\csromanspacer{}ทฺวาทสมํ.\csromanspacer{} คงฺคาเปยฺยาลวคฺโค ฉฏฺโฐ.}

\titlehead{\textbf{ตสฺสุทฺทานํ}}
\csromangathameasure{ฉ ปาจีนโต นินฺนา, ฉ นินฺนา จ สมุทฺทโต. \\ เทฺวเต ฉ ทฺวาทส โหนฺติ, วคฺโค เตน ปวุจฺจตีติ. (อปฺปมาท พลกรณียวคฺคา วิตฺถาเรตพฺพา.)}
\csromangathagroup{\csromangathastanza{ฉ ปาจีนโต นินฺนา, ฉ นินฺนา จ สมุทฺทโต. \\ เทฺวเต ฉ ทฺวาทส โหนฺติ, วคฺโค เตน ปวุจฺจตีติ. (อปฺปมาท พลกรณียวคฺคา วิตฺถาเรตพฺพา.)}}

\chapterhead{9. เอสนาวคฺค}
\csromanmatikaanchor{mtk.161}
\csromanmatikaanchor{mtk.162}
\csromantocmarknopage{subparagraph}{9. เอสนาวคฺค}{mtk.161}
\bookmark[dest={mtk.161},level=5]{9. เอสนาวคฺค}
\csromantocmarkref{subparagraph}{1-12. เอสนาทิสุตฺต}{mtk.162}
\bookmark[dest={mtk.162},level=5]{1-12. เอสนาทิสุตฺต}
\csromanheader{6}{1-12. เอสนาทิสุตฺต}
\proseitem{792-802}{เอวํ เอสนา ปาฬิ วิตฺถาเรตพฺพา ราควินยปริโยสานํ โทสวินยปริโยสานํ โมหวินยปริโยสานํ.\csromanspacer{} เอสนาวคฺโค นวโม.}

\titlehead{\textbf{ตสฺสุทฺทานํ}}
\csromangathameasure{เอสนา วิธา อาสโว, ภโว จ ทุกฺขตา ติสฺโส. \\ ขิลํ มลญฺจ นีโฆ จ, เวทนา ตณฺหา ตสินา จาติ.}
\csromangathagroup{\csromangathastanza{เอสนา วิธา อาสโว, ภโว จ ทุกฺขตา ติสฺโส. \\ ขิลํ มลญฺจ นีโฆ จ, เวทนา ตณฺหา ตสินา จาติ.}}

\csromanheader{2}{6. พลสํยุตฺต}
\chapterhead{10. โอฆวคฺค}
\csromanmatikaanchor{mtk.163}
\csromanmatikaanchor{mtk.164}
\csromantocmarknopage{subparagraph}{10. โอฆวคฺค}{mtk.163}
\bookmark[dest={mtk.163},level=5]{10. โอฆวคฺค}
\csromantocmarkref{subparagraph}{1-10. โอฆาทิสุตฺต}{mtk.164}
\bookmark[dest={mtk.164},level=5]{1-10. โอฆาทิสุตฺต}
\csromanheader{6}{1-10. โอฆาทิสุตฺต}
\proseitem{803-821}{\csromanfolio{221}ปญฺจิมานิ ภิกฺขเว อุทฺธมฺภาคิยานิ สํโยชนานิ.\csromanspacer{} กตมานิ ปญฺจ, รูปราโค อรูปราโค มาโน อุทฺธจฺจํ อวิชฺชา, อิมานิ โข ภิกฺขเว ปญฺจุทฺธมฺภาคิยานิ สํโยชนานิ.\csromanspacer{} อิเมสํ โข ภิกฺขเว ปญฺจนฺนํ อุทฺธมฺภาคิยานํ สํโยชนานํ อภิญฺญาย ปริญฺญาย ปริกฺขยาย ปหานาย ปญฺจ พลานิ ภาเวตพฺพานิ.\csromanspacer{} กตมานิ ปญฺจ, อิธ ภิกฺขเว ภิกฺขุ สทฺธาพลํ ภาเวติ ฯเปฯ ปญฺญาพลํ ภาเวติ ราควินยปริโยสานํ โทสวินยปริโยสานํ โมหวินยปริโยสานํ.\csromanspacer{} อิเมสํ โข ภิกฺขเว ปญฺจนฺนํ อุทฺธมฺภาคิยานํ สํโยชนานํ อภิญฺญาย ปริญฺญาย ปริกฺขยาย ปหานาย อิมานิ ปญฺจ พลานิ ภาเวตพฺพานีติ.\csromanspacer{}\csromanspacer{}ทสมํ.\csromanspacer{} โอฆวคฺโค ทสโม.}

\titlehead{\textbf{ตสฺสุทฺทานํ}}
\csromangathameasure{โอโฆ โยโค อุปาทานํ, คนฺถา อนุสเยน จ. \\ กามคุณา นีวรณา, ขนฺธา โอรุทฺธมฺภาคิยาติ. \\ \textbf{พลสํยุตฺตํ ฉฏฺฐํ.}}
\csromangathagroup{\csromangathastanza{โอโฆ โยโค อุปาทานํ, คนฺถา อนุสเยน จ. \\ กามคุณา นีวรณา, ขนฺธา โอรุทฺธมฺภาคิยาติ. \\ \textbf{พลสํยุตฺตํ ฉฏฺฐํ.}}}

\csromanmatikaanchor{mtk.165}
\csromantocmarknopage{subparagraph}{7. อิทฺธิปาทสํยุตฺต}{mtk.165}
\bookmark[dest={mtk.165},level=5]{7. อิทฺธิปาทสํยุตฺต}
\csromanheader{6}{7. อิทฺธิปาทสํยุตฺต}
\csromanmatikaanchor{mtk.166}
\csromantocmarknopage{subparagraph}{1. จาปาลวคฺค}{mtk.166}
\bookmark[dest={mtk.166},level=5]{1. จาปาลวคฺค}
\csromanheader{6}{1. จาปาลวคฺค}
\csromanmatikaanchor{mtk.167}
\csromantocmarkref{subparagraph}{1. อปารสุตฺต}{mtk.167}
\bookmark[dest={mtk.167},level=5]{1. อปารสุตฺต}
\csromanheader{6}{1. อปารสุตฺต}
\csromanmatikaanchor{mtk.168}
\csromantocmarkref{subparagraph}{2-4. วิรทฺธสุตฺตาทีนิ}{mtk.168}
\bookmark[dest={mtk.168},level=5]{2-4. วิรทฺธสุตฺตาทีนิ}
\proseitem{813}{\csromanfolio{222}จตฺตาโรเม ภิกฺขเว อิทฺธิปาทา ภาวิตา พหุลีกตา อปารา ปารํ คมนาย สํวตฺตนฺติ.\csromanspacer{} กตเม จตฺตาโร, อิธ ภิกฺขเว ภิกฺขุ ฉนฺทสมาธิปฺปธานสงฺขารสมนฺนาคตํ อิทฺธิปาทํ ภาเวติ, วีริยสมาธิปฺปธานสงฺขารสมนฺนาคตํ อิทฺธิปาทํ ภาเวติ, จิตฺตสมาธิปฺปธานสงฺขารสมนฺนาคตํ อิทฺธิปาทํ ภาเวติ, วีมํสาสมาธิปฺปธานสงฺขารสมนฺนาคตํ อิทฺธิปาทํ ภาเวติ.\csromanspacer{} อิเม โข ภิกฺขเว จตฺตาโร อิทฺธิปาทา ภาวิตา พหุลีกตา อปารา ปารํ คมนาย สํวตฺตนฺตีติ.\csromanspacer{}\csromanspacer{}ปฐมํ.}

\csromanheader{2}{2. วิรทฺธสุตฺต}
\proseitem{814}{เยสํ เกสญฺจิ ภิกฺขเว จตฺตาโร อิทฺธิปาทา วิรทฺธา, วิรทฺโธ เตสํ อริโย มคฺโค สมฺมา ทุกฺขกฺขยคามี.\csromanspacer{} เยสํ เกสญฺจิ ภิกฺขเว จตฺตาโร อิทฺธิปาทา อารทฺธา, อารทฺโธ เตสํ อริโย มคฺโค สมฺมา ทุกฺขกฺขยคามี.\csromanspacer{} กตเม จตฺตาโร, อิธ ภิกฺขเว ภิกฺขุ ฉนฺทสมาธิปฺปธานสงฺขารสมนฺนาคตํ อิทฺธิปาทํ ภาเวติ, วีริยสมาธิ ฯเปฯ จิตฺตสมาธิ ฯเปฯ วีมํสาสมาธิปฺปธานสงฺขารสมนฺนาคตํ อิทฺธิปาทํ ภาเวติ.\csromanspacer{} เยสํ เกสญฺจิ ภิกฺขเว อิเม จตฺตาโร อิทฺธิปาทา วิรทฺธา, วิรทฺโธ เตสํ อริโย มคฺโค สมฺมา ทุกฺขกฺขยคามี.\csromanspacer{} เยสํ เกสญฺจิ ภิกฺขเว อิเม จตฺตาโร อิทฺธิปาทา อารทฺธา, อารทฺโธ เตสํ อริโย มคฺโค สมฺมา ทุกฺขกฺขยคามีติ.\csromanspacer{}\csromanspacer{}ทุติยํ.}

\csromanheader{2}{3. อริยสุตฺต}
\proseitem{815}{จตฺตาโรเม ภิกฺขเว อิทฺธิปาทา ภาวิตา พหุลีกตา อริยา นิยฺยานิกา นิยฺยนฺติ ตกฺกรสฺส สมฺมา ทุกฺขกฺขยาย.\csromanspacer{} กตเม จตฺตาโร, อิธ ภิกฺขเว ภิกฺขุ ฉนฺทสมาธิปฺปธานสงฺขารสมนฺนาคตํ อิทฺธิปาทํ ภาเวติ, วีริยสมาธิ ฯเปฯ จิตฺตสมาธิ ฯเปฯ วีมํสาสมาธิปฺปธานสงฺขารสมนฺนาคตํ อิทฺธิปาทํ ภาเวติ.\csromanspacer{} อิเม โข ภิกฺขเว จตฺตาโร อิทฺธิปาทา}

\csromanheader{2}{7. อิทฺธิปาทสํยุตฺต}
\prose{\csromanfolio{223}ภาวิตา พหุลีกตา อริยา นิยฺยานิกา นิยฺยนฺติ ตกฺกรสฺส สมฺมา ทุกฺขกฺขยายาติ.\csromanspacer{}\csromanspacer{}ตติยํ.}

\csromanheader{2}{4. นิพฺพิทาสุตฺต}
\proseitem{816}{จตฺตาโรเม ภิกฺขเว อิทฺธิปาทา ภาวิตา พหุลีกตา เอกนฺตนิพฺพิทาย วิราคาย นิโรธาย อุปสมาย อภิญฺญาย สมฺโพธาย นิพฺพานาย สํวตฺตนฺติ.\csromanspacer{} กตเม จตฺตาโร, อิธ ภิกฺขเว ภิกฺขุ ฉนฺทสมาธิปฺปธานสงฺขารสมนฺนาคตํ อิทฺธิปาทํ ภาเวติ, วีริยสมาธิ ฯเปฯ จิตฺตสมาธิ ฯเปฯ วีมํสาสมาธิปฺปธานสงฺขารสมนฺนาคตํ อิทฺธิปาทํ ภาเวติ.\csromanspacer{} อิเม โข ภิกฺขเว จตฺตาโร อิทฺธิปาทา ภาวิตา พหุลีกตา เอกนฺตนิพฺพิทาย วิราคาย นิโรธาย อุปสมาย อภิญฺญาย สมฺโพธาย นิพฺพานาย สํวตฺตนฺตีติ.\csromanspacer{}\csromanspacer{}จตุตฺถํ.}

\csromanmatikaanchor{mtk.169}
\csromantocmarkref{subparagraph}{5. อิทฺธิปเทสสุตฺต}{mtk.169}
\bookmark[dest={mtk.169},level=5]{5. อิทฺธิปเทสสุตฺต}
\csromanheader{6}{5. อิทฺธิปเทสสุตฺต}
\proseitem{817}{เย หิ เกจิ ภิกฺขเว อตีตมทฺธานํ สมณา วา พฺราหฺมณา วา อิทฺธิปเทสํ อภินิปฺผาเทสุํ, สพฺเพ เต จตุนฺนํ อิทฺธิปาทานํ ภาวิตตฺตา พหุลีกตตฺตา.\csromanspacer{} เย หิ เกจิ ภิกฺขเว อนาคตมทฺธานํ สมณา วา พฺราหฺมณา วา อิทฺธิปเทสํ อภินิปฺผาเทสฺสนฺติ, สพฺเพ เต จตุนฺนํ อิทฺธิปาทานํ ภาวิตตฺตา พหุลีกตตฺตา.\csromanspacer{} เย หิ เกจิ ภิกฺขเว เอตรหิ สมณา วา พฺราหฺมณา วา อิทฺธิปเทสํ อภินิปฺผาเทนฺติ, สพฺเพ เต จตุนฺนํ อิทฺธิปาทานํ ภาวิตตฺตา พหุลีกตตฺตา.}

\prose{กตเมสํ จตุนฺนํ, อิธ ภิกฺขเว ภิกฺขุ ฉนฺทสมาธิปฺปธานสงฺขารสมนฺนาคตํ อิทฺธิปาทํ ภาเวติ, วีริยสมาธิ ฯเปฯ จิตฺตสมาธิ ฯเปฯ วีมํสาสมาธิปฺปธานสงฺขารสมนฺนาคตํ อิทฺธิปาทํ ภาเวติ.\csromanspacer{} เย หิ เกจิ ภิกฺขเว อตีตมทฺธานํ สมณา วา พฺราหฺมณา วา อิทฺธิปเทสํ อภินิปฺผาเทสุํ, สพฺเพ เต อิเมสํเยว จตุนฺนํ อิทฺธิปาทานํ ภาวิตตฺตา พหุลีกตตฺตา.\csromanspacer{} เย หิ เกจิ ภิกฺขเว อนาคตมทฺธานํ สมณา วา พฺราหฺมณา วา อิทฺธิปเทสํ อภินิปฺผาเทสฺสนฺติ, สพฺเพ เต อิเมสํเยว จตุนฺนํ อิทฺธิปาทานํ ภาวิตตฺตา พหุลีกตตฺตา.\csromanspacer{} เย หิ เกจิ ภิกฺขเว เอตรหิ สมณา วา พฺราหฺมณา วา อิทฺธิปเทสํ อภินิปฺผาเทนฺติ, สพฺเพ เต อิเมสํเยว จตุนฺนํ อิทฺธิปาทานํ ภาวิตตฺตา พหุลีกตตฺตาติ.\csromanspacer{}\csromanspacer{}ปญฺจมํ.}

\csromanmatikaanchor{mtk.170}
\csromantocmarkref{subparagraph}{6. สมตฺตสุตฺต}{mtk.170}
\bookmark[dest={mtk.170},level=5]{6. สมตฺตสุตฺต}
\csromanheader{6}{6. สมตฺตสุตฺต}
\csromanmatikaanchor{mtk.171}
\csromantocmarkref{subparagraph}{7-9. ภิกฺขุสุตฺตาทีนิ}{mtk.171}
\bookmark[dest={mtk.171},level=5]{7-9. ภิกฺขุสุตฺตาทีนิ}
\proseitem{818}{\csromanfolio{224}เย หิ เกจิ ภิกฺขเว อตีตมทฺธานํ สมณา วา พฺราหฺมณา วา สมตฺตํ อิทฺธิํ อภินิปฺผาเทสุํ, สพฺเพ เต จตุนฺนํ อิทฺธิปาทานํ ภาวิตตฺตา พหุลีกตตฺตา.\csromanspacer{} เย หิ เกจิ ภิกฺขเว อนาคตมทฺธานํ สมณา วา พฺราหฺมณา วา สมตฺตํ อิทฺธิํ อภินิปฺผาเทสฺสนฺติ, สพฺเพ เต จตุนฺนํ อิทฺธิปาทานํ ภาวิตตฺตา พหุลีกตตฺตา.\csromanspacer{} เย หิ เกจิ ภิกฺขเว เอตรหิ สมณา วา พฺราหฺมณา วา สมตฺตํ อิทฺธิํ อภินิปฺผาเทนฺติ, สพฺเพ เต จตุนฺนํ อิทฺธิปาทานํ ภาวิตตฺตา พหุลีกตตฺตา.}

\prose{กตเมสํ จตุนฺนํ, อิธ ภิกฺขเว ภิกฺขุ ฉนฺทสมาธิปฺปธานสงฺขารสมนฺนาคตํ อิทฺธิปาทํ ภาเวติ, วีริยสมาธิ ฯเปฯ จิตฺตสมาธิ ฯเปฯ วีมํสาสมาธิปฺปธานสงฺขารสมนฺนาคตํ อิทฺธิปาทํ ภาเวติ.\csromanspacer{} เย หิ เกจิ ภิกฺขเว อตีตมทฺธานํ สมณา วา พฺราหฺมณา วา สมตฺตํ อิทฺธิํ อภินิปฺผาเทสุํ, สพฺเพ เต อิเมสํเยว จตุนฺนํ อิทฺธิปาทานํ ภาวิตตฺตา พหุลีกตตฺตา.\csromanspacer{} เย หิ เกจิ ภิกฺขเว อนาคตมทฺธานํ สมณา วา พฺราหฺมณา วา สมตฺตํ อิทฺธิํ อภินิปฺผาเทสฺสนฺติ, สพฺเพ เต อิเมสํเยว จตุนฺนํ อิทฺธิปาทานํ ภาวิตตฺตา พหุลีกตตฺตา.\csromanspacer{} เย หิ เกจิ ภิกฺขเว เอตรหิ สมณา วา พฺราหฺมณา วา สมตฺตํ อิทฺธิํ อภินิปฺผาเทนฺติ, สพฺเพ เต อิเมสํเยว จตุนฺนํ อิทฺธิปาทานํ ภาวิตตฺตา พหุลีกตตฺตาติ.\csromanspacer{}\csromanspacer{}ฉฏฺฐํ.}

\csromanheader{2}{7. ภิกฺขุสุตฺต}
\proseitem{819}{เย หิ เกจิ ภิกฺขเว อตีตมทฺธานํ ภิกฺขู อาสวานํ ขยา อนาสวํ เจโตวิมุตฺติํ ปญฺญาวิมุตฺติํ ทิฏฺเฐว ธมฺเม สยํ อภิญฺญา สจฺฉิกตฺวา อุปสมฺปชฺช วิหริํสุ, สพฺเพ เต จตุนฺนํ อิทฺธิปาทานํ ภาวิตตฺตา พหุลีกตตฺตา.\csromanspacer{} เย หิ เกจิ ภิกฺขเว อนาคตมทฺธานํ ภิกฺขู อาสวานํ ขยา อนาสวํ เจโตวิมุตฺติํ ปญฺญาวิมุตฺติํ ทิฏฺเฐว ธมฺเม สยํ อภิญฺญา สจฺฉิกตฺวา อุปสมฺปชฺช วิหริสฺสนฺติ, สพฺเพ เต จตุนฺนํ อิทฺธิปาทานํ ภาวิตตฺตา พหุลีกตตฺตา.\csromanspacer{} เย หิ เกจิ ภิกฺขเว เอตรหิ ภิกฺขู อาสวานํ ขยา อนาสวํ เจโตวิมุตฺติํ ปญฺญาวิมุตฺติํ ทิฏฺเฐว ธมฺเม สยํ อภิญฺญา สจฺฉิกตฺวา อุปสมฺปชฺช วิหรนฺติ, สพฺเพ เต จตุนฺนํ อิทฺธิปาทานํ ภาวิตตฺตา พหุลีกตตฺตา.}

\csromanheader{2}{7. อิทฺธิปาทสํยุตฺต}
\prose{\csromanfolio{225}กตเมสํ จตุนฺนํ, อิธ ภิกฺขเว ภิกฺขุ ฉนฺทสมาธิปฺปธานสงฺขารสมนฺนาคตํ อิทฺธิปาทํ ภาเวติ, วีริยสมาธิ ฯเปฯ จิตฺตสมาธิ.\csromanspacer{} วีมํสาสมาธิปฺปธานสงฺขารสมนฺนาคตํ อิทฺธิปาทํ ภาเวติ.\csromanspacer{} เย หิ เกจิ ภิกฺขเว อตีตมทฺธานํ ภิกฺขู อาสวานํ ขยา อนาสวํ เจโตวิมุตฺติํ ปญฺญาวิมุตฺติํ ทิฏฺเฐว ธมฺเม สยํ อภิญฺญา สจฺฉิกตฺวา อุปสมฺปชฺช วิหริํสุ, สพฺเพ เต อิเมสํเยว จตุนฺนํ อิทฺธิปาทานํ ภาวิตตฺตา พหุลีกตตฺตา.\csromanspacer{} เย หิ เกจิ ภิกฺขเว อนาคตมทฺธานํ ภิกฺขู อาสวานํ ขยา อนาสวํ เจโตวิมุตฺติํ ปญฺญาวิมุตฺติํ ทิฏฺเฐว ธมฺเม สยํ อภิญฺญา สจฺฉิกตฺวา อุปสมฺปชฺช วิหริสฺสนฺติ, สพฺเพ เต อิเมสํเยว จตุนฺนํ อิทฺธิปาทานํ ภาวิตตฺตา พหุลีกตตฺตา.\csromanspacer{} เย หิ เกจิ ภิกฺขเว เอตรหิ ภิกฺขู อาสวานํ ขยา อนาสวํ เจโตวิมุตฺติํ ปญฺญาวิมุตฺติํ ทิฏฺเฐว ธมฺเม สยํ อภิญฺญา สจฺฉิกตฺวา อุปสมฺปชฺช วิหรนฺติ, สพฺเพ เต อิเมสํเยว จตุนฺนํ อิทฺธิปาทานํ ภาวิตตฺตา พหุลีกตตฺตาติ.\csromanspacer{}\csromanspacer{}สตฺตมํ.}

\csromanheader{2}{8. พุทฺธสุตฺต}
\proseitem{820}{จตฺตาโรเม ภิกฺขเว อิทฺธิปาทา.\csromanspacer{} กตเม จตฺตาโร, อิธ ภิกฺขเว ภิกฺขุ ฉนฺทสมาธิปฺปธานสงฺขารสมนฺนาคตํ อิทฺธิปาทํ ภาเวติ, วีริยสมาธิ ฯเปฯ จิตฺตสมาธิ.\csromanspacer{} วีมํสาสมาธิปฺปธานสงฺขารสมนฺนาคตํ อิทฺธิปาทํ ภาเวติ.\csromanspacer{} อิเม โข ภิกฺขเว จตฺตาโร อิทฺธิปาทา, อิเมสํ โข ภิกฺขเว จตุนฺนํ อิทฺธิปาทานํ ภาวิตตฺตา พหุลีกตตฺตา ตถาคโต อรหํ สมฺมาสมฺพุทฺโธติ วุจฺจตีติ.\csromanspacer{}\csromanspacer{}อฏฺฐมํ.}

\csromanheader{2}{9. ญาณสุตฺต}
\proseitem{821}{“อยํ ฉนฺทสมาธิปฺปธานสงฺขารสมนฺนาคโต อิทฺธิปาโท”ติ เม ภิกฺขเว ปุพฺเพ อนนุสฺสุเตสุ ธมฺเมสุ จกฺขุํ อุทปาทิ, ญาณํ อุทปาทิ, ปญฺญา อุทปาทิ, วิชฺชา อุทปาทิ, อาโลโก อุทปาทิ. “โส โข ปนายํ ฉนฺทสมาธิปฺปธานสงฺขารสมนฺนาคโต อิทฺธิปาโท ภาเวตพฺโพ”ติ เม ภิกฺขเว. “ภาวิโต”ติ เม ภิกฺขเว ปุพฺเพ อนนุสฺสุเตสุ ธมฺเมสุ จกฺขุํ อุทปาทิ, ญาณํ อุทปาทิ, ปญฺญา อุทปาทิ, วิชฺชา อุทปาทิ, อาโลโก อุทปาทิ.}

\prose{\csromanfolio{226}“อยํ วีริยสมาธิปฺปธานสงฺขารสมนฺนาคโต อิทฺธิปาโท”ติ เม ภิกฺขเว ปุพฺเพ อนนุสฺสุเตสุ ธมฺเมสุ จกฺขุํ อุทปาทิ, ญาณํ อุทปาทิ, ปญฺญา อุทปาทิ, วิชฺชา อุทปาทิ, อาโลโก อุทปาทิ. “โส โข ปนายํ วีริยสมาธิปฺปธานสงฺขารสมนฺนาคโต อิทฺธิปาโท ภาเวตพฺโพ”ติ เม ภิกฺขเว. “ภาวิโต”ติ เม ภิกฺขเว ปุพฺเพ อนนุสฺสุเตสุ ธมฺเมสุ จกฺขุํ อุทปาทิ, ญาณํ อุทปาทิ, ปญฺญา อุทปาทิ, วิชฺชา อุทปาทิ, อาโลโก อุทปาทิ.}

\prose{“อยํ จิตฺตสมาธิปฺปธานสงฺขารสมนฺนาคโต อิทฺธิปาโท”ติ เม ภิกฺขเว ปุพฺเพ อนนุสฺสุเตสุ ธมฺเมสุ จกฺขุํ อุทปาทิ, ญาณํ อุทปาทิ, ปญฺญา อุทปาทิ, วิชฺชา อุทปาทิ, อาโลโก อุทปาทิ. “โส โข ปนายํ จิตฺตสมาธิปฺปธานสงฺขารสมนฺนาคโต อิทฺธิปาโท ภาเวตพฺโพ”ติ เม ภิกฺขเว. “ภาวิโต”ติ เม ภิกฺขเว ปุพฺเพ อนนุสฺสุเตสุ ธมฺเมสุ จกฺขุํ อุทปาทิ, ญาณํ อุทปาทิ, ปญฺญา อุทปาทิ, วิชฺชา อุทปาทิ, อาโลโก อุทปาทิ.}

\prose{“อยํ วีมํสาสมาธิปฺปธานสงฺขารสมนฺนาคโต อิทฺธิปาโท”ติ เม ภิกฺขเว ปุพฺเพ อนนุสฺสุเตสุ ธมฺเมสุ จกฺขุํ อุทปาทิ, ญาณํ อุทปาทิ, ปญฺญา อุทปาทิ, วิชฺชา อุทปาทิ, อาโลโก อุทปาทิ. “โส โข ปนายํ วีมํสาสมาธิปฺปธานสงฺขารสมนฺนาคโต อิทฺธิปาโท ภาเวตพฺโพ”ติ เม ภิกฺขเว. “ภาวิโต”ติ เม ภิกฺขเว ปุพฺเพ อนนุสฺสุเตสุ ธมฺเมสุ จกฺขุํ อุทปาทิ, ญาณํ อุทปาทิ, ปญฺญา อุทปาทิ, วิชฺชา อุทปาทิ, อาโลโก อุทปาทีติ.\csromanspacer{}\csromanspacer{}นวมํ.}

\csromanmatikaanchor{mtk.172}
\csromantocmarkref{subparagraph}{10. เจติยสุตฺต}{mtk.172}
\bookmark[dest={mtk.172},level=5]{10. เจติยสุตฺต}
\csromanheader{6}{10. เจติยสุตฺต}
\proseitem{822}{เอวํ เม สุตํ\csromandash{}เอกํ สมยํ ภควา เวสาลิยํ วิหรติ มหาวเน กูฏาคารสาลายํ.\csromanspacer{} อถ โข ภควา ปุพฺพณฺหสมยํ นิวาเสตฺวา ปตฺตจีวรมาทาย เวสาลิํ ปิณฺฑาย ปาวิสิ, เวสาลิยํ ปิณฺฑาย จริตฺวา ปจฺฉาภตฺตํ ปิณฺฑปาตปฏิกฺกนฺโต อายสฺมนฺตํ อานนฺทํ อามนฺเตสิ “คณฺหาหิ อานนฺท นิสีทนํ, เยน จาปาลํ เจติยํ\footnote{ปาวาลเจติยํ (สฺยา, กํ)} เตนุปสงฺกมิสฺสาม ทิวาวิหารายา”ติ. “เอวํ ภนฺเต”ติ โข อายสฺมา อานนฺโท ภควโต ปฏิสฺสุตฺวา นิสีทนํ อาทาย ภควนฺตํ ปิฏฺฐิโต ปิฏฺฐิโต อนุพนฺธิ.\csromanspacer{} อถ โข ภควา เยน จาปาลํ เจติยํ เตนุปสงฺกมิ,}

\csromanheader{2}{7. อิทฺธิปาทสํยุตฺต}
\prose{\csromanfolio{227}อุปสงฺกมิตฺวา ปญฺญตฺเต อาสเน นิสีทิ.\csromanspacer{} อายสฺมาปิ โข อานนฺโท ภควนฺตํ อภิวาเทตฺวา เอกมนฺตํ นิสีทิ, เอกมนฺตํ นิสินฺนํ โข อายสฺมนฺตํ อานนฺทํ ภควา เอตทโวจ\csromandash{}}

\prose{รมณียา อานนฺท เวสาลี, รมณียํ อุเทนํ เจติยํ, รมณียํ โคตมกํ เจติยํ, รมณียํ สตฺตมฺพํ เจติยํ, รมณียํ พหุปุตฺตํ เจติยํ\footnote{พหุปุตฺตกเจติยํ (สฺยา, กํ, อิ, ก)}, รมณียํ สารนฺททํ เจติยํ\footnote{อานนฺทเจติยํ (ก), สานนฺทรํ (ก)}, รมณียํ จาปาลํ เจติยํ.\csromanspacer{} ยสฺส กสฺสจิ อานนฺท จตฺตาโร อิทฺธิปาทา ภาวิตา พหุลีกตา ยานีกตา วตฺถุกตา อนุฏฺฐิตา ปริจิตา สุสมารทฺธา, โส อากงฺขมาโน กปฺปํ วา ติฏฺเฐยฺย กปฺปาวเสสํ วา, ตถาคตสฺส โข อานนฺท จตฺตาโร อิทฺธิปาทา ภาวิตา พหุลีกตา ยานีกตา วตฺถุกตา อนุฏฺฐิตา ปริจิตา สุสมารทฺธา, อากงฺขมาโน อานนฺท ตถาคโต กปฺปํ วา ติฏฺเฐยฺย กปฺปาวเสสํ วาติ.}

\prose{เอวมฺปิ โข อายสฺมา อานนฺโท ภควตา โอฬาริเก นิมิตฺเต กยิรมาเน โอฬาริเก โอภาเส กยิรมาเน นาสกฺขิ ปฏิวิชฺฌิตุํ, น ภควนฺตํ ยาจิ “ติฏฺฐตุ ภนฺเต ภควา กปฺปํ, ติฏฺฐตุ สุคโต กปฺปํ\footnote{สุคโต กปฺปาวเสสํ (อิ, ก) ที 2. 86 ปิฏฺเฐปิ.} พหุชนหิตาย พหุชนสุขาย โลกานุกมฺปาย อตฺถาย หิตาย สุขาย เทวมนุสฺสานนฺ”ติ, ยถา ตํ มาเรน ปริยุฏฺฐิตจิตฺโต.}

\prose{ทุติยมฺปิ โข ภควา ฯเปฯ.\csromanspacer{} ตติยมฺปิ โข ภควา อายสฺมนฺตํ อานนฺทํ อามนฺเตสิ “รมณียา อานนฺท เวสาลี, รมณียํ อุเทนํ เจติยํ, รมณียํ โคตมกํ เจติยํ, รมณียํ สตฺตพฺพํ เจติยํ, รมณียํ พหุปุตฺตํ เจติยํ, รมณียํ สารนฺททํ เจติยํ, รมณียํ จาปาลํ เจติยํ.\csromanspacer{} ยสฺส กสฺสจิ อานนฺท จตฺตาโร อิทฺธิปาทา ภาวิตา พหุลีกตา ยานีกตา วตฺถุกตา อนุฏฺฐิตา ปริจิตา สุสมารทฺธา, โส อากงฺขมาโน กปฺปํ วา ติฏฺเฐยฺย กปฺปาวเสสํ วา, ตถาคตสฺส โข อานนฺท จตฺตาโร อิทฺธิปาทา ภาวิตา พหุลีกตา ยานีกตา วตฺถุกตา อนุฏฺฐิตา ปริจิตา สุสมารทฺธา, อากงฺขมาโน อานนฺท ตถาคโต กปฺปํ วา ติฏฺเฐยฺย กปฺปาวเสสํ วา”ติ.}

\prose{\csromanfolio{228}เอวมฺปิ โข อายสฺมา อานนฺโท ภควตา โอฬาริเก นิมิตฺเต กยิรมาเน โอฬาริเก โอภาเส กยิรมาเน นาสกฺขิ ปฏิวิชฺฌิตุํ, น ภควนฺตํ ยาจิ “ติฏฺฐตุ ภนฺเต ภควา กปฺปํ, ติฏฺฐตุ สุคโต กปฺปํ พหุชนหิตาย พหุชนสุขาย โลกานุกมฺปาย อตฺถาย หิตาย สุขาย เทวมนุสฺสานนฺ”ติ, ยถา ตํ มาเรน ปริยุฏฺฐิตจิตฺโต.}

\prose{อถ โข ภควา อายสฺมนฺตํ อานนฺทํ อามนฺเตสิ “คจฺฉ โข ตฺวํ อานนฺท, ยสฺสทานิ กาลํ มญฺญสี”ติ. “เอวํ ภนฺเต”ติ โข อายสฺมา อานนฺโท ภควโต ปฏิสฺสุตฺวา อุฏฺฐายาสนา ภควนฺตํ อภิวาเทตฺวา ปทกฺขิณํ กตฺวา อวิทูเร อญฺญตรสฺมิํ รุกฺขมูเล นิสีทิ.\csromanspacer{} อถ โข มาโร ปาปิมา อจิรปกฺกนฺเต อายสฺมนฺเต อานนฺเท เยน ภควา เตนุปสงฺกมิ, อุปสงฺกมิตฺวา ภควนฺตํ เอตทโวจ “ปรินิพฺพาตุทานิ ภนฺเต ภควา, ปรินิพฺพาตุทานิ สุคโต\footnote{ปรินิพฺพาตุ สุคโต (สี, สฺยา, กํ) เอวมุปริปิ.}, ปรินิพฺพานกาโลทานิ ภนฺเต ภควโต.\csromanspacer{} ภาสิตา โข ปเนสา ภนฺเต ภควตา วาจา “น ตาวาหํ ปาปิม ปรินิพฺพายิสฺสามิ, ยาว เม ภิกฺขู น สาวกา ภวิสฺสนฺติ วิยตฺตา วินีตา วิสารทา พหุสฺสุตา ธมฺมธรา ธมฺมานุธมฺมปฺปฏิปนฺนา สามีจิปฺปฏิปนฺนา อนุธมฺมจาริโน สกํ อาจริยกํ อุคฺคเหตฺวา อาจิกฺขิสฺสนฺติ เทเสสฺสนฺติ ปญฺญเปสฺสนฺติ ปฏฺฐเปสฺสนฺติ วิวริสฺสนฺติ วิภชิสฺสนฺติ อุตฺตานีกริสฺสนฺติ\footnote{อุตฺตานิํกริสฺสนฺติ (ก), อุตฺตานิกริสฺสนฺติ (อิ)}, อุปฺปนฺนํ ปรปฺปวาทํ สหธมฺเมน สุนิคฺคหิตํ นิคฺคเหตฺวา สปฺปาฏิหาริยํ ธมฺมํ เทเสสฺสนฺตี”ติ.}

\prose{สนฺติ โข ปน ภนฺเต เอตรหิ ภิกฺขู ภควโต สาวกา วิยตฺตา วินีตา วิสารทา พหุสฺสุตา ธมฺมธรา ธมฺมานุธมฺมปฺปฏิปนฺนา สามีจิปฺปฏิปนฺนา อนุธมฺมจาริโน สกํ อาจริยกํ อุคฺคเหตฺวา อาจิกฺขนฺติ เทเสนฺติ ปญฺญเปนฺติ ปฏฺฐเปนฺติ วิวรนฺติ วิภชนฺติ อุตฺตานีกโรนฺติ, อุปฺปนฺนํ ปรปฺปวาทํ สหธมฺเมน สุนิคฺคหิตํ นิคฺคเหตฺวา สปฺปาฏิหาริยํ ธมฺมํ เทเสนฺติ.\csromanspacer{} ปรินิพฺพาตุทานิ ภนฺเต ภควา, ปรินิพฺพาตุทานิ สุคโต, ปรินิพฺพานกาโลทานิ ภนฺเต ภควโต.}

\prose{ภาสิตา โข ปเนสา ภนฺเต ภควตา วาจา “น ตาวาหํ ปาปิม ปรินิพฺพายิสฺสามิ, ยาว เม ภิกฺขุนิโย น สาวิกา ภวิสฺสนฺติ วิยตฺตา}

\csromanheader{2}{7. อิทฺธิปาทสํยุตฺต}
\prose{\csromanfolio{229}วินีตา วิสารทา พหุสฺสุตา ธมฺมธรา ธมฺมานุธมฺมปฺปฏิปนฺนา สามีจิปฺปฏิปนฺนา อนุธมฺมจารินิโย สกํ อาจริยกํ อุคฺคเหตฺวา อาจิกฺขิสฺสนฺติ เทเสสฺสนฺติ ปญฺญเปสฺสนฺติ ปฏฺฐเปสฺสนฺติ วิวริสฺสนฺติ วิภชิสฺสนฺติ อุตฺตานีกริสฺสนฺติ, อุปฺปนฺนํ ปรปฺปวาทํ สหธมฺเมน สุนิคฺคหิตํ นิคฺคเหตฺวา สปฺปาฏิหาริยํ ธมฺมํ เทเสสฺสนฺตี”ติ.}

\prose{สนฺติ โข ปน ภนฺเต เอตรหิ ภิกฺขุนิโย ภควโต สาวิกา วิยตฺตา วินีตา วิสารทา พหุสฺสุตา ธมฺมธรา ธมฺมานุธมฺมปฺปฏิปนฺนา สามีจิปฺปฏิปนฺนา อนุธมฺมจารินิโย สกํ อาจริยกํ อุคฺคเหตฺวา อาจิกฺขนฺติ เทเสนฺติ ปญฺญเปนฺติ ปฏฺฐเปนฺติ วิวรนฺติ วิภชนฺติ อุตฺตานีกโรนฺติ, อุปฺปนฺนํ ปรปฺปวาทํ สหธมฺเมน สุนิคฺคหิตํ นิคฺคเหตฺวา สปฺปาฏิหาริยํ ธมฺมํ เทเสนฺติ, ปรินิพฺพาตุทานิ ภนฺเต ภควา, ปรินิพฺพาตุทานิ สุคโต, ปรินิพฺพานกาโลทานิ ภนฺเต ภควโต.}

\prose{ภาสิตา โข ปเนสา ภนฺเต ภควตา วาจา “น ตาวาหํ ปาปิม ปรินิพฺพายิสฺสามิ, ยาว เม อุปาสกา ฯเปฯ ยาว เม อุปาสิกา น สาวิกา ภวิสฺสนฺติ วิยตฺตา วินีตา วิสารทา พหุสฺสุตา ธมฺมธรา ธมฺมานุธมฺมปฺปฏิปนฺนา สามีจิปฺปฏิปนฺนา อนุธมฺมจาริโย สกํ อาจริยกํ อุคฺคเหตฺวา อาจิกฺขิสฺสนฺติ เทเสสฺสนฺติ ปญฺญเปสฺสนฺติ ปฏฺฐเปสฺสนฺติ วิวริสฺสนฺติ วิภชิสฺสนฺติ อุตฺตานีกริสฺสนฺติ, อุปฺปนฺนํ ปรปฺปวาทํ สหธมฺเมน สุนิคฺคหิตํ นิคฺคเหตฺวา สปฺปาฏิหาริยํ ธมฺมํ เทเสสฺสนฺตี”ติ.}

\prose{สนฺติ โข ปน ภนฺเต เอตรหิ อุปาสกา อุปาสิกา ภควโต สาวิกา วิยตฺตา วินีตา วิสารทา พหุสฺสุตา ธมฺมธรา ธมฺมานุธมฺมปฺปฏิปนฺนา สามีจิปฺปฏิปนฺนา อนุธมฺมจารินิโย สกํ อาจริยกํ อุคฺคเหตฺวา อาจิกฺขนฺติ เทเสนฺติ ปญฺญเปนฺติ ปฏฺฐเปนฺติ วิวรนฺติ วิภชนฺติ อุตฺตานีกโรนฺติ, อุปฺปนฺนํ ปรปฺปวาทํ สหธมฺเมน สุนิคฺคหิตํ นิคฺคเหตฺวา สปฺปาฏิหาริยํ ธมฺมํ เทเสนฺติ.\csromanspacer{} ปรินิพฺพาตุทานิ ภนฺเต ภควา, ปรินิพฺพาตุทานิ สุคโต, ปรินิพฺพานกาโลทานิ ภนฺเต ภควโต.}

\prose{ภาสิตา โข ปเนสา ภนฺเต ภควตา วาจา “น ตาวาหํ ปาปิม ปรินิพฺพายิสฺสามิ, ยาว เม อิทํ พฺรหฺมจริยํ น อิทฺธญฺเจว ภวิสฺสติ ผีตญฺจ, วิตฺถาริตํ พาหุชญฺญํ ปุถุภูตํ ยาว เทวมนุสฺเสหิ สุปฺปกาสิตนฺ”ติ.\csromanspacer{} ตยิทํ ภนฺเต ภควโต พฺรหฺมจริยํ อิทฺธญฺเจว ผีตญฺจ, วิตฺถาริตํ พาหุชญฺญํ}

\prose{\csromanfolio{230}ปุถุภูตํ ยาว เทวมนุสฺเสหิ สุปฺปกาสิตํ.\csromanspacer{} ปรินิพฺพาตุทานิ ภนฺเต ภควา, ปรินิพฺพาตุทานิ สุคโต, ปรินิพฺพานกาโลทานิ ภนฺเต ภควโตติ.}

\prose{เอวํ วุตฺเต ภควา มารํ ปาปิมนฺตํ เอตทโวจ “อปฺโปสฺสุกฺโก ตฺวํ ปาปิม โหหิ, น จิรํ\footnote{นจิรสฺเสว (ก)} ตถาคตสฺส ปรินิพฺพานํ ภวิสฺสติ, อิโต ติณฺณํ มาสานํ อจฺจเยน ตถาคโต ปรินิพฺพายิสฺสตี”ติ.\csromanspacer{} อถ โข ภควา จาปาเล เจติเย สโต สมฺปชาโน อายุสงฺขารํ โอสฺสชิ, โอสฺสฏฺเฐ จ\footnote{โอสชฺเช ปน (ก)} ภควตา อายุสงฺขาเร มหาภูมิจาโล อโหสิ ภิํสนโก โลมหํโส, เทวทุนฺทุภิโย\footnote{เทวทุทฺรภิโย (ก)} จ ผลิํสุ.\csromanspacer{} อถ โข ภควา เอตมตฺถํ วิทิตฺวา ตายํ เวลายํ อิมํ อุทานํ อุทาเนสิ\csromandash{}}

\csromangathameasure{“ตุลมตุลญฺจ สมฺภวํ, ภวสงฺขารมวสฺสชิ มุนิ. \\ อชฺฌตฺตรโต สมาหิโต, อภินฺทิ กวจมิวตฺตสมฺภวนฺ”ติ. ทสมํ. จาปาลวคฺโค ปฐโม.}
\csromangathagroup{\csromangathastanza{“ตุลมตุลญฺจ สมฺภวํ, ภวสงฺขารมวสฺสชิ มุนิ. \\ อชฺฌตฺตรโต สมาหิโต, อภินฺทิ กวจมิวตฺตสมฺภวนฺ”ติ.\csromanspacer{} ทสมํ.\csromanspacer{} จาปาลวคฺโค ปฐโม.}}

\titlehead{\textbf{ตสฺสุทฺทานํ}}
\csromangathameasure{อปาราปิ วิรทฺโธ จ, อริยา นิพฺพิทาปิ จ. \\ ปเทสํ สมตฺตํ ภิกฺขุ, พุทฺธํ ญาณญฺจ เจติยนฺติ.}
\csromangathagroup{\csromangathastanza{อปาราปิ วิรทฺโธ จ, อริยา นิพฺพิทาปิ จ. \\ ปเทสํ สมตฺตํ ภิกฺขุ, พุทฺธํ ญาณญฺจ เจติยนฺติ.}}

\csromanmatikaanchor{mtk.173}
\csromantocmarknopage{subparagraph}{2. ปาสาทกมฺปนวคฺค}{mtk.173}
\bookmark[dest={mtk.173},level=5]{2. ปาสาทกมฺปนวคฺค}
\csromanheader{6}{2. ปาสาทกมฺปนวคฺค}
\csromanmatikaanchor{mtk.174}
\csromantocmarkref{subparagraph}{1. ปุพฺพสุตฺต}{mtk.174}
\bookmark[dest={mtk.174},level=5]{1. ปุพฺพสุตฺต}
\csromanheader{6}{1. ปุพฺพสุตฺต}
\proseitem{823}{สาวตฺถินิทานํ.\csromanspacer{} ปุพฺเพว เม ภิกฺขเว สมฺโพธา อนภิสมฺพุทฺธสฺส โพธิสตฺตสฺเสว สโต เอตทโหสิ “โก นุ โข เหตุ โก ปจฺจโย อิทฺธิปาทภาวนายา”ติ.\csromanspacer{} ตสฺส มยฺหํ ภิกฺขเว เอตทโหสิ “อิธ ภิกฺขุ ฉนฺทสมาธิปฺปธานสงฺขารสมนฺนาคตํ อิทฺธิปาทํ ภาเวติ, อิติ เม ฉนฺโท น จ อติลีโน ภวิสฺสติ, น จ อติปฺปคฺคหิโต ภวิสฺสติ, น จ อชฺฌตฺตํ สํขิตฺโต ภวิสฺสติ.\csromanspacer{} น จ พหิทฺธา วิกฺขิตฺโต ภวิสฺสติ, ปจฺฉาปุเรสญฺญี จ วิหรติ, ยถา ปุเร ตถา ปจฺฉา, ยถา ปจฺฉา ตถา ปุเร,}

\csromanheader{2}{7. อิทฺธิปาทสํยุตฺต}
\prose{\csromanfolio{231}ยถา อโธ ตถา อุทฺธํ, ยถา อุทฺธํ ตถา อโธ, ยถา ทิวา ตถา รตฺติํ, ยถา รตฺติํ ตถา ทิวา, อิติ วิวเฏน เจตสา อปริโยนทฺเธน สปฺปภาสํ จิตฺตํ ภาเวติ.}

\prose{วีริยสมาธิปฺปธานสงฺขารสมนฺนาคตํ อิทฺธิปาทํ ภาเวติ, อิติ เม วีริยํ น จ อติลีนํ ภวิสฺสติ, น จ อติปฺปคฺคหิตํ ภวิสฺสติ, น จ อชฺฌตฺตํ สํขิตฺตํ ภวิสฺสติ, น จ พหิทฺธา วิกฺขิตฺตํ ภวิสฺสติ, ปจฺฉาปุเรสญฺญี จ วิหรติ, ยถา ปุเร ตถา ปจฺฉา, ยถา ปจฺฉา ตถา ปุเร, ยถา อโธ ตถา อุทฺธํ, ยถา อุทฺธํ ตถา อโธ, ยถา ทิวา ตถา รตฺติํ, ยถา รตฺติํ ตถา ทิวา, อิติ วิวเฏน เจตสา อปริโยนทฺเธน สปฺปภาสํ จิตฺตํ ภาเวติ.}

\prose{จิตฺตสมาธิปฺปธานสงฺขารสมนฺนาคตํ อิทฺธิปาทํ ภาเวติ, อิติ เม จิตฺตํ น จ อติลีนํ ภวิสฺสติ, น จ อติปฺปคฺคหิตํ ภวิสฺสติ, น จ อชฺฌตฺตํ สํขิตฺตํ ภวิสฺสติ, น จ พหิทฺธา วิกฺขิตฺตํ ภวิสฺสติ, ปจฺฉาปุเรสญฺญี จ วิหรติ, ยถา ปุเร ตถา ปจฺฉา, ยถา ปจฺฉา ตถา ปุเร, ยถา อโธ ตถา อุทฺธํ, ยถา อุทฺธํ ตถา อโธ, ยถา ทิวา ตถา รตฺติํ, ยถา รตฺติํ ตถา ทิวา, อิติ วิวเฏน เจตสา อปริโยนทฺเธน สปฺปภาสํ จิตฺตํ ภาเวติ.}

\prose{วีมํสาสมาธิปฺปธานสงฺขารสมนฺนาคตํ อิทฺธิปาทํ ภาเวติ, อิติ เม วีมํสา น จ อติลีนา ภวิสฺสติ, น จ อติปฺปคฺคหิตา ภวิสฺสติ, น จ อชฺฌตฺตํ สํขิตฺตา ภวิสฺสติ, น จ พหิทฺธา วิกฺขิตฺตา ภวิสฺสติ, ปจฺฉาปุเรสญฺญี จ วิหรติ, ยถา ปุเร ตถา ปจฺฉา, ยถา ปจฺฉา ตถา ปุเร, ยถา อโธ ตถา อุทฺธํ, ยถา อุทฺธํ ตถา อโธ, ยถา ทิวา ตถา รตฺติํ, ยถา รตฺติํ ตถา ทิวา, อิติ วิวเฏน เจตสา อปริโยนทฺเธน สปฺปภาสํ จิตฺตํ ภาเวติ.}

\prose{เอวํ ภาวิเตสุ โข ภิกฺขุ จตูสุ อิทฺธิปาเทสุ เอวํ พหุลีกเตสุ อเนกวิหิตํ อิทฺธิวิธํ ปจฺจนุโภติ, เอโกปิ หุตฺวา พหุธา โหติ, พหุธาปิ หุตฺวา เอโก โหติ, อาวิภาวํ ติโรภาวํ ติโรกุฏฺฏํ ติโรปาการํ ติโรปพฺพตํ อสชฺชมาโน คจฺฉติ เสยฺยถาปิ อากาเส.\csromanspacer{} ปถวิยาปิ อุมฺมุชฺชนิมุชฺชํ กโรติ เสยฺยถาปิ อุทเก.\csromanspacer{} อุทเกปิ}

\prose{\csromanfolio{232}อภิชฺชมาเน\footnote{อภิชฺชมาโน (สี, อิ, ก)} คจฺฉติ เสยฺยถาปิ ปถวิยํ.\csromanspacer{} อากาเสปิ ปลฺลงฺเกน กมติ เสยฺยถาปิ ปกฺขี สกุโณ.\csromanspacer{} อิเมปิ จนฺทิมสูริเย เอวํมหิทฺธิเก เอวํมหานุภาเว ปาณินา ปริมสติ\footnote{ปรามสติ (สี, สฺยา, กํ)} ปริมชฺชติ, ยาว พฺรหฺมโลกาปิ กาเยน วสํ วตฺเตติ.}

\prose{เอวํ ภาวิเตสุ โข ภิกฺขุ จตูสุ อิทฺธิปาเทสุ เอวํ พหุลีกเตสุ ทิพฺพาย โสตธาตุยา วิสุทฺธาย อติกฺกนฺตมานุสิกาย อุโภ สทฺเท สุณาติ ทิพฺเพ จ มานุเส จ ทูเร สนฺติเก จาติ.}

\prose{เอวํ ภาวิเตสุ โข ภิกฺขุ จตูสุ อิทฺธิปาเทสุ เอวํ พหุลีกเตสุ ปรสตฺตานํ ปรปุคฺคลานํ เจตสา เจโต ปริจฺจ ปชานาติ, สราคํ วา จิตฺตํ “สราคํ จิตฺตนฺ”ติ ปชานาติ, วีตราคํ วา จิตฺตํ “วีตราคํ จิตฺตนฺ”ติ ปชานาติ, สโทสํ วา จิตฺตํ “สโทสํ จิตฺตนฺ”ติ ปชานาติ, วีตโทสํ วา จิตฺตํ “วีตโทสํ จิตฺตนฺ”ติ ปชานาติ, สโมหํ วา จิตฺตํ “สโมหํ จิตฺตนฺ”ติ ปชานาติ, วีตโมหํ วา จิตฺตํ “วีตโมหํ จิตฺตนฺ”ติ ปชานาติ, สํขิตฺตํ วา จิตฺตํ “สํขิตฺตํ จิตฺตนฺ”ติ ปชานาติ, วิกฺขิตฺตํ วา จิตฺตํ “วิกฺขิตฺตํ จิตฺตนฺ”ติ ปชานาติ, มหคฺคตํ วา จิตฺตํ “มหคฺคตํ จิตฺตํนฺ”ติ ปชานาติ, อมหคฺคตํ วา จิตฺตํ “อมหคฺคตํ จิตฺตนฺ”ติ ปชานาติ, ส-อุตฺตรํ วา จิตฺตํ “ส-อุตฺตรํ จิตฺตนฺ”ติ ปชานาติ, อนุตฺตรํ วา จิตฺตํ “อนุตฺตรํ จิตฺตนฺ”ติ ปชานาติ, สมาหิตํ วา จิตฺตํ “สมาหิตํ จิตฺตนฺ”ติ ปชานาติ, อสมาหิตํ วา จิตฺตํ “อสมาหิตํ จิตฺตนฺ”ติ ปชานาติ, วิมุตฺตํ วา จิตฺตํ “วิมุตฺตํ จิตฺตนฺ”ติ ปชานาติ, อวิมุตฺตํ วา จิตฺตํ “อวิมุตฺตํ จิตฺตนฺ”ติ ปชานาติ.}

\prose{เอวํ ภาวิเตสุ โข ภิกฺขุ จตูสุ อิทฺธิปาเทสุ เอวํ พหุลีกเตสุ อเนกวิหิตํ ปุพฺเพนิวาสํ อนุสฺสรติ.\csromanspacer{} เสยฺยถิทํ, เอกมฺปิ ชาติํ เทฺวปิ ชาติโย ติสฺโสปิ ชาติโย จตสฺโสปิ ชาติโย ปญฺจปิ ชาติโย ทสปิ ชาติโย วีสมฺปิ ชาติโย ติํสมฺปิ ชาติโย จตฺตาลีสมฺปิ ชาติโย ปญฺญาสมฺปิ ชาติโย ชาติสตมฺปิ ชาติสหสฺสมฺปิ ชาติสตสหสฺสมฺปิ อเนเกปิ สํวฏฺฏกปฺเป อเนเกปิ วิวฏฺฏกปฺเป อเนเกปิ สํวฏฺฏวิวฏฺฏกปฺเป, “อมุตฺราสิํ เอวํนาโม เอวํโคตฺโต เอวํวณฺโณ เอวมาหาโร เอวํสุขทุกฺขปฺปฏิสํเวที เอวมายุปริยนฺโต, โส ตโต จุโต อมุตฺร อุทปาทิํ\footnote{อุปฺปาทิํ (สี)}, ตตฺราปาสิํ}

\csromanheader{2}{7. อิทฺธิปาทสํยุตฺต}
\prose{\csromanfolio{233}“เอวํนาโม เอวํโคตฺโต เอวํวณฺโณ เอวมาหาโร เอวํสุขทุกฺขปฺปฏิสํเวที เอวมายุปริยนฺโต, โส ตโต จุโต อิธูปปนฺโน”ติ.\csromanspacer{} อิติ สาการํ ส-อุทฺเทสํ อเนกวิหิตํ ปุพฺเพนิวาสํ อนุสฺสรติ.}

\prose{เอวํ ภาวิเตสุ โข ภิกฺขุ จตูสุ อิทฺธิปาเทสุ เอวํ พหุลีกเตสุ ทิพฺเพน จกฺขุนา วิสุทฺเธน อติกฺกนฺตมานุสเกน สตฺเต ปสฺสติ จวมาเน อุปปชฺชมาเน หีเน ปณีเต สุวณฺเณ ทุพฺพณฺเณ สุคเต ทุคฺคเต, ยถากมฺมูปเค สตฺเต ปชานาติ “อิเม วต โภนฺโต สตฺตา กายทุจฺจริเตน สมนฺนาคตา วจีทุจฺจริเตน สมนฺนาคตา มโนทุจฺจริเตน สมนฺนาคตา, อริยานํ อุปวาทกา มิจฺฉาทิฏฺฐิกา มิจฺฉาทิฏฺฐิกมฺมสมาทานา, เต กายสฺส เภทา ปรํ มรณา อปายํ ทุคฺคติํ วินิปาตํ นิรยํ อุปปนฺนา.\csromanspacer{} อิเม วา ปน โภนฺโต สตฺตา กายสุจริเตน สมนฺนาคตา วจีสุจริเตน สมนฺนาคตา, มโนสุจริเตน สมนฺนาคตา อริยานํ อนุปวาทกา สมฺมาทิฏฺฐิกา สมฺมาทิฏฺฐิกมฺมสมาทานา, เต กายสฺส เภทา ปรํ มรณา สุคติํ สคฺคํ โลกํ อุปปนฺนา”ติ.\csromanspacer{} อิติ ทิพฺเพน จกฺขุนา วิสุทฺเธน อติกฺกนฺตมานุสเกน สตฺเต ปสฺสติ จวมาเน อุปปชฺชมาเน หีเน ปณีเต สุวณฺเณ ทุพฺพณฺเณ สุคเต ทุคฺคเต, ยถากมฺมูปเค สตฺเต ปชานาติ.}

\prose{เอวํ ภาวิเตสุ โข ภิกฺขุ จตูสุ อิทฺธิปาเทสุ เอวํ พหุลีกเตสุ อาสวานํ ขยา อนาสวํ เจโตวิมุตฺติํ ปญฺญาวิมุตฺติํ ทิฏฺเฐว ธมฺเม สยํ อภิญฺญา สจฺฉิกตฺวา อุปสมฺปชฺช วิหรตีติ.\csromanspacer{}\csromanspacer{}ปฐมํ.}

\csromanmatikaanchor{mtk.175}
\csromantocmarkref{subparagraph}{2. มหปฺผลสุตฺต}{mtk.175}
\bookmark[dest={mtk.175},level=5]{2. มหปฺผลสุตฺต}
\csromanheader{6}{2. มหปฺผลสุตฺต}
\csromanmatikaanchor{mtk.176}
\csromantocmarkref{subparagraph}{3. ฉนฺทสมาทิสุตฺต}{mtk.176}
\bookmark[dest={mtk.176},level=5]{3. ฉนฺทสมาทิสุตฺต}
\proseitem{824}{จตฺตาโรเม ภิกฺขเว อิทฺธิปาทา ภาวิตา พหุลีกตา มหปฺผลา โหนฺติ มหานิสํสา.\csromanspacer{} กถํ ภาวิตา จ ภิกฺขเว จตฺตาโร อิทฺธิปาทา กถํ พหุลีกตา มหปฺผลา โหนฺติ มหานิสํสา.\csromanspacer{} อิธ ภิกฺขเว ภิกฺขุ ฉนฺทสมาธิปฺปธานสงฺขารสมนฺนาคตํ อิทฺธิปาทํ ภาเวติ, อิติ เม ฉนฺโท น จ อติลีโน ภวิสฺสติ, น จ อติปฺปคฺคหิโต ภวิสฺสติ, น จ อชฺฌตฺตํ สํขิตฺโต ภวิสฺสติ, น จ พหิทฺธา วิกฺขิตฺโต ภวิสฺสติ, ปจฺฉาปุเรสญฺญี จ วิหรติ, ยถา ปุเร ตถา ปจฺฉา, ยถา ปจฺฉา ตถา ปุเร, ยถา อโธ ตถา อุทฺธํ, ยถา อุทฺธํ ตถา อโธ, ยถา ทิวา ตถา รตฺติํ, ยถา รตฺติํ \csromanfolio{234}ตถา ทิวา, อิติ วิวเฏน เจตสา อปริโยนทฺเธน สปฺปภาสํ จิตฺตํ ภาเวติ.}

\prose{วีริยสมาธิ ฯเปฯ.\csromanspacer{} จิตฺตสมาธิ.\csromanspacer{} วีมํสาสมาธิปฺปธานสงฺขารสมนฺนาคตํ อิทฺธิปาทํ ภาเวติ, อิติ เม วีมํสา น จ อติลีนา ภวิสฺสติ, น จ อติปฺปคฺคหิตา ภวิสฺสติ, น จ อชฺฌตฺตํ สํขิตฺตา ภวิสฺสติ, น จ พหิทฺธา วิกฺขิตฺตา ภวิสฺสติ, ปจฺฉาปุเรสญฺญี จ วิหรติ, ยถา ปุเร ตถา ปจฺฉา, ยถา ปจฺฉา ตถา ปุเร, ยถา อโธ ตถา อุทฺธํ, ยถา อุทฺธํ ตถา อโธ, ยถา ทิวา ตถา รตฺติํ, ยถา รตฺติํ ตถา ทิวา, อิติ วิวเฏน เจตสา อปริโยนทฺเธน สปฺปภาสํ จิตฺตํ ภาเวติ.\csromanspacer{} เอวํ ภาวิตา โข ภิกฺขเว จตฺตาโร อิทฺธิปาทา เอวํ พหุลีกตา มหปฺผลา โหนฺติ มหานิสํสา.}

\prose{เอวํ ภาวิเตสุ โข ภิกฺขเว ภิกฺขุ จตูสุ อิทฺธิปาเทสุ เอวํ พหุลีกเตสุ อเนกวิหิตํ อิทฺธิวิธํ ปจฺจนุโภติ, เอโกปิ หุตฺวา พหุธา โหติ ฯเปฯ ยาว พฺรหฺมโลกาปิ กาเยน วสํ วตฺเตติ.}

\prose{เอวํ ภาวิเตสุ โข ภิกฺขเว ภิกฺขุ จตูสุ อิทฺธิปาเทสุ เอวํ พหุลีกเตสุ อาสวานํ ขยา อนาสวํ เจโตวิมุตฺติํ ปญฺญาวิมุตฺติํ ทิฏฺเฐว ธมฺเม สยํ อภิญฺญา สจฺฉิกตฺวา อุปสมฺปชฺช วิหรตีติ.\csromanspacer{}\csromanspacer{}ทุติยํ.}

\csromanheader{2}{3. ฉนฺทสมาธิสุตฺต}
\proseitem{825}{ฉนฺทํ เจ ภิกฺขเว ภิกฺขุ นิสฺสาย ลภติ สมาธิํ, ลภติ จิตฺตสฺส เอกคฺคตํ, อยํ วุจฺจติ ฉนฺทสมาธิ.\csromanspacer{} โส อนุปฺปนฺนานํ ปาปกานํ อกุสลานํ ธมฺมานํ อนุปฺปาทาย ฉนฺทํ ชเนติ วายมติ, วีริยํ อารภติ, จิตฺตํ ปคฺคณฺหาติ ปทหติ.\csromanspacer{} อุปฺปนฺนานํ ปาปกานํ อกุสลานํ ธมฺมานํ ปหานาย ฉนฺทํ ชเนติ วายมติ, วีริยํ อารภติ, จิตฺตํ ปคฺคณฺหาติ ปทหติ.\csromanspacer{} อนุปฺปนฺนานํ กุสลานํ ธมฺมานํ อุปฺปาทาย ฉนฺทํ ชเนติ วายมติ, วีริยํ อารภติ, จิตฺตํ ปคฺคณฺหาติ ปทหติ.\csromanspacer{} อุปฺปนฺนานํ กุสลานํ ธมฺมานํ ฐิติยา อสมฺโมสาย ภิยฺโยภาวาย เวปุลฺลาย ภาวนาย ปาริปูริยา ฉนฺทํ ชเนติ วายมติ, วีริยํ อารภติ, จิตฺตํ ปคฺคณฺหาติ ปทหติ.\csromanspacer{} อิเม วุจฺจนฺติ ปธานสงฺขาราติ.\csromanspacer{} อิติ อยญฺจ ฉนฺโท อยญฺจ ฉนฺทสมาธิ อิเม จ ปธานสงฺขารา, อยํ วุจฺจติ ภิกฺขเว ฉนฺทสมาธิปฺปธานสงฺขารสมนฺนาคโต อิทฺธิปาโท.}

\csromanheader{2}{7. อิทฺธิปาทสํยุตฺต}
\prose{\csromanfolio{235}วีริยํ เจ ภิกฺขเว ภิกฺขุ นิสฺสาย ลภติ สมาธิํ, ลภติ จิตฺตสฺส เอกคฺคตํ, อยํ วุจฺจติ วีริยสมาธิ.\csromanspacer{} โส อนุปฺปนฺนานํ ฯเปฯ.\csromanspacer{} อุปฺปนฺนานํ กุสลานํ ธมฺมานํ ฐิติยา อสมฺโมสาย ภิยฺโยภาวาย เวปุลฺลาย ภาวนาย ปาริปูริยา ฉนฺทํ ชเนติ วายมติ, วีริยํ อารภติ, จิตฺตํ ปคฺคณฺหาติ ปทหติ.\csromanspacer{} อิเม วุจฺจนฺติ ปธานสงฺขาราติ.\csromanspacer{} อิติ อิทญฺจ วีริยํ อยญฺจ วีริยสมาธิ อิเม จ ปธานสงฺขารา, อยํ วุจฺจติ ภิกฺขเว วีริยสมาธิปฺปธานสงฺขารสมนฺนาคโต อิทฺธิปาโท.}

\prose{จิตฺตํ เจ ภิกฺขเว ภิกฺขุ นิสฺสาย ลภติ สมาธิํ, ลภติ จิตฺตสฺส เอกคฺคตํ, อยํ วุจฺจติ จิตฺตสมาธิ.\csromanspacer{} โส อนุปฺปนฺนานํ ปาปกานํ ฯเปฯ.\csromanspacer{} อุปฺปนฺนานํ กุสลานํ ธมฺมานํ ฐิติยา อสมฺโมสาย ภิยฺโยภาวาย เวปุลฺลาย ภาวนาย ปาริปูริยา ฉนฺทํ ชเนติ วายมติ, วีริยํ อารภติ, จิตฺตํ ปคฺคณฺหาติ ปทหติ.\csromanspacer{} อิเม วุจฺจนฺติ ปธานสงฺขาราติ.\csromanspacer{} อิติ อิทญฺจ จิตฺตํ อยญฺจ จิตฺตสมาธิ อิเม จ ปธานสงฺขารา, อยํ วุจฺจติ ภิกฺขเว จิตฺตสมาธิปฺปธานสงฺขารสมนฺนาคโต อิทฺธิปาโท.}

\prose{วีมํสํ เจ ภิกฺขเว ภิกฺขุ นิสฺสาย ลภติ สมาธิํ, ลภติ จิตฺตสฺส เอกคฺคตํ, อยํ วุจฺจติ วีมํสาสมาธิ.\csromanspacer{} โส อนุปฺปนฺนานํ ปาปกานํ อกุสลานํ ธมฺมานํ อนุปฺปาทาย ฉนฺทํ ชเนติ วายมติ, วีริยํ อารภติ, จิตฺตํ ปคฺคณฺหาติ ปทหติ ฯเปฯ.\csromanspacer{} อุปฺปนฺนานํ กุสลานํ ธมฺมานํ ฐิติยา อสมฺโมสาย ภิยฺโยภาวาย เวปุลฺลาย ภาวนาย ปาริปูริยา ฉนฺทํ ชเนติ วายมติ, วีริยํ อารภติ, จิตฺตํ ปคฺคณฺหาติ ปทหติ.\csromanspacer{} อิเม วุจฺจนฺติ ปธานสงฺขารา.\csromanspacer{} อิติ อยญฺจ วีมํสา อยญฺจ วีมํสาสมาธิ อิเม จ ปธานสงฺขารา, อยํ วุจฺจติ ภิกฺขเว วีมํสาสมาธิปฺปธานสงฺขารสมนฺนาคโต อิทฺธิปาโทติ.\csromanspacer{}\csromanspacer{}ตติยํ.}

\csromanmatikaanchor{mtk.177}
\csromantocmarkref{subparagraph}{4. โมคฺคลฺลานสุตฺต}{mtk.177}
\bookmark[dest={mtk.177},level=5]{4. โมคฺคลฺลานสุตฺต}
\csromanheader{6}{4. โมคฺคลฺลานสุตฺต}
\proseitem{826}{เอวํ เม สุตํ\csromandash{}เอกํ สมยํ ภควา สาวตฺถิยํ วิหรติ ปุพฺพาราเม มิคารมาตุปาสาเท.\csromanspacer{} เตน โข ปน สมเยน สมฺพหุลา ภิกฺขู เหฏฺฐามิคารมาตุปาสาเท วิหรนฺติ อุทฺธตา อุนฺนฬา จปลา มุขรา วิกิณฺณวาจา มุฏฺฐสฺสติโน อสมฺปชานา อสมาหิตา ภนฺตจิตฺตา\footnote{วิพฺภนฺตจิตฺตา (สี, สฺยา, กํ)} ปากตินฺทฺริยา.}

\prose{\csromanfolio{236}อถ โข ภควา อายสฺมนฺตํ มหาโมคฺคลฺลานํ อามนฺเตสิ “เอเต โข โมคฺคลฺลาน สพฺรหฺมจาริโน เหฏฺฐามิคารมาตุปาสาเท วิหรนฺติ อุทฺธตา อุนฺนฬา จปลา มุขรา วิกิณฺณวาจา มุฏฺฐสฺสติโน อสมฺปชานา อสมาหิตา ภนฺตจิตฺตา ปากตินฺทฺริยา, คจฺฉ โมคฺคลฺลาน, เต ภิกฺขู สํเวเชหี”ติ.}

\prose{“เอวํ ภนฺเต”ติ โข อายสฺมา มหาโมคฺคลฺลาโน ภควโต ปฏิสฺสุตฺวา ตถารูปํ อิทฺธาภิสงฺขารํ อภิสงฺขาสิ\footnote{อภิสงฺขเรสิ (สฺยา, อิ, ก)}.\csromanspacer{} ยถา ปาทงฺคุฏฺฐเกน มิคารมาตุปาสาทํ สงฺกมฺเปสิ สมฺปกมฺเปสิ สมฺปจาเลสิ.\csromanspacer{} อถ โข เต ภิกฺขู สํวิคฺคา โลมหฏฺฐชาตา เอกมนฺตํ อฏฺฐํสุ, อจฺฉริยํ วต โภ, อพฺภุตํ วต โภ, นิวาตญฺจ วต, อยญฺจ มิคารมาตุปาสาโท คมฺภีรเนโม สุนิขาโต อจโล อสมฺปกมฺปี, อถ จ ปน สงฺกมฺปิโต สมฺปกมฺปิโต สมฺปจาลิโตติ.}

\prose{อถ โข ภควา เยน เต ภิกฺขู เตนุปสงฺกมิ, อุปสงฺกมิตฺวา เต ภิกฺขู ภควา เอตทโวจ “กิํ นุ ตุเมฺห ภิกฺขเว สํวคฺคา โลมหฏฺฐชาตา เอกมนฺตํ ฐิตา”ติ.\csromanspacer{} อจฺฉริยํ ภนฺเต, อพฺภุตํ ภนฺเต, นิวาตญฺจ วต, อยญฺจ มิคารมาตุปาสาโท คมฺภีรเนโม สุนิขาโต อจโล อสมฺปกมฺปี, อถ จ ปน สงฺกมฺปิโต สมฺปกมฺปิโต สมฺปจาลิโตติ.\csromanspacer{} ตุเมฺหว โข ภิกฺขเว สํเวเชตุกาเมน โมคฺคลฺลาเนน ภิกฺขุนา ปาทงฺคุฏฺฐเกน มิคารมาตุปาสาโท สงฺกมฺปิโต สมฺปกมฺปิโต สมฺปจาลิโต.\csromanspacer{} ตํ กิํ มญฺญถ ภิกฺขเว, กตเมสํ ธมฺมานํ ภาวิตตฺตา พหุลีกตตฺตา โมคฺคลฺลาโน ภิกฺขุ เอวํมหิทฺธิโก เอวํมหานุภาโวติ.\csromanspacer{} ภควํมูลกา โน ภนฺเต ธมฺมา ภควํเนตฺติกา ภควํปฏิสรณา, สาธุ วต ภนฺเต ภควนฺตํเยว ปฏิภาตุ เอตสฺส ภาสิตสฺส อตฺโถ, ภควโต สุตฺวา ภิกฺขู ธาเรสฺสนฺตีติ.}

\prose{เตน หิ ภิกฺขเว สุณาถ, จตุนฺนํ โข ภิกฺขเว อิทฺธิปาทานํ ภาวิตตฺตา พหุลีกตตฺตา โมคฺคลฺลาโน ภิกฺขุ เอวํมหิทฺธิโก เอวํมหานุภาโว.\csromanspacer{} กตเมสํ จตุนฺนํ, อิธ ภิกฺขเว โมคฺคลฺลาโน ภิกฺขุ ฉนฺทสมาธิปฺปธานสงฺขารสมนฺนาคตํ อิทฺธิปาทํ ภาเวติ, วีริยสมาธิ ฯเปฯ จิตฺตสมาธิ, วีมํสาสมาธิปฺปธานสงฺขารสมนฺนาคตํ อิทฺธิปาทํ ภาเวติ, อิติ เม วีมํสา น}

\csromanheader{2}{7. อิทฺธิปาทสํยุตฺต}
\prose{\csromanfolio{237}จ อติลีนา ภวิสฺสติ, น จ อติปฺปคฺคหิตา ภวิสฺสติ น จ อชฺฌตฺตํ สํขิตฺตา ภวิสฺสติ, น จ พหิทฺธา วิกฺขิตฺตา ภวิสฺสติ.\csromanspacer{} ปจฺฉาปุเรสญฺญี จ วิหรติ, ยถา ปุเร ตถา ปจฺฉา, ยถา ปจฺฉา ตถา ปุเร, ยถา อโธ ตถา อุทฺธํ, ยถา อุทฺธํ ตถา อโธ, ยถา ทิวา ตถา รตฺติํ, ยถา รตฺติํ ตถา ทิวา, อิติ วิวเฏน เจตสา อปริโยนทฺเธน สปฺปภาสํ จิตฺตํ ภาเวติ.\csromanspacer{} อิเมสํ โข ภิกฺขเว จตุนฺนํ อิทฺธิปาทานํ ภาวิตตฺตา พหุลีกตตฺตา โมคฺคลฺลาโน ภิกฺขุ เอวํมหิทฺธิโก เอวํมหานุภาโว, อิเมสญฺจ ปน ภิกฺขเว จตุนฺนํ อิทฺธิปาทานํ ภาวิตตฺตา พหุลีกตตฺตา โมคฺคลฺลาโน ภิกฺขุ อเนกวิหิตํ อิทฺธิวิธํ ปจฺจนุโภติ ฯเปฯ ยาว พฺรหฺมโลกาปิ กาเยน วสํ วตฺเตติ.\csromanspacer{} อิเมสญฺจ ปน ภิกฺขเว จตุนฺนํ อิทฺธิปาทานํ ภาวิตตฺตา พหุลีกตตฺตา โมคฺคลฺลาโน ภิกฺขุ อาสวานํ ขยา อนาสวํ เจโตวิมุตฺติํ ปญฺญาวิมุตฺติํ ทิฏฺเฐว ธมฺเม สยํ อภิญฺญา สจฺฉิกตฺวา อุปสมฺปชฺช วิหรตีติ.\csromanspacer{}\csromanspacer{}จตุตฺถํ.}

\csromanmatikaanchor{mtk.178}
\csromantocmarkref{subparagraph}{5. อุณฺณาภพฺราหฺมณสุตฺต}{mtk.178}
\bookmark[dest={mtk.178},level=5]{5. อุณฺณาภพฺราหฺมณสุตฺต}
\csromanheader{6}{5. อุณฺณาภพฺราหฺมณสุตฺต}
\proseitem{827}{เอวํ เม สุตํ\csromandash{}เอกํ สมยํ อายสฺมา อานนฺโท โกสมฺพิยํ วิหรติ โฆสิตาราเม.\csromanspacer{} อถ โข อุณฺณาโภ พฺราหฺมโณ เยนายสฺมา อานนฺโท เตนุปสงฺกมิ, อุปสงฺกมิตฺวา อายสฺมตา อานนฺเทน สทฺธิํ สมฺโมทิ, สมฺโมทนียํ กถํ สารณียํ วีติสาเรตฺวา เอกมนฺตํ นิสีทิ, เอกมนฺตํ นิสินฺโน โข อุณฺณาโภ พฺราหฺมโณ อายสฺมนฺตํ อานนฺทํ เอตทโวจ “กิมตฺถิยํ นุ โข โภ อานนฺท สมเณ โคตเม พฺรหฺมจริยํ วุสฺสตี”ติ.\csromanspacer{} ฉนฺทปฺปหานตฺถํ โข พฺราหฺมณ ภควติ พฺรหฺมจริยํ วุสฺสตีติ.}

\prose{อตฺถิ ปน โภ อานนฺท มคฺโค, อตฺถิ ปฏิปทา เอตสฺส ฉนฺทสฺส ปหานายาติ.\csromanspacer{} อตฺถิ โข พฺราหฺมณ มคฺโค, อตฺถิ ปฏิปทา เอตสฺส ฉนฺทสฺส ปหานายาติ.}

\prose{กตโม ปน โภ อานนฺท มคฺโค กตมา ปฏิปทา เอตสฺส ฉนฺทสฺส ปหานายาติ, อิธ พฺราหฺมณ ภิกฺขุ, ฉนฺทสมาธิปฺปธานสงฺขารสมนฺนาคตํ อิทฺธิปาทํ ภาเวติ, วีริยสมาธิ ฯเปฯ.\csromanspacer{} จิตฺตสมาธิ.\csromanspacer{} วีมํสาสมาธิปฺปธานสงฺขารสมนฺนาคตํ อิทฺธิปาทํ ภาเวติ.\csromanspacer{} อยํ โข พฺราหฺมณ มคฺโค อยํ ปฏิปทา เอตสฺส ฉนฺทสฺส ปหานายาติ.}

\prose{\csromanfolio{238}เอวํ สนฺเต โภ อานนฺท สนฺตกํ โหติ โน อสนฺตกํ.\csromanspacer{} ฉนฺเทเนว ฉนฺทํ ปชหิสฺสตีติ เนตํ ฐานํ วิชฺชติ.\csromanspacer{} เตน หิ พฺราหฺมณ ตญฺเญเวตฺถ ปฏิปุจฺฉิสฺสามิ.\csromanspacer{} ยถา เต ขเมยฺย, ตถา ตํ พฺยากเรยฺยาสิ.\csromanspacer{} ตํ กิํ มญฺญสิ พฺราหฺมณ, อโหสิ เต ปุพฺเพ ฉนฺโท “อารามํ คมิสฺสามี”ติ, ตสฺส เต อารามคตสฺส โย ตชฺโช ฉนฺโท, โส ปฏิปฺปสฺสทฺโธติ.\csromanspacer{} เอวํ โภ.\csromanspacer{} อโหสิ เต ปุพฺเพ วีริยํ “อารามํ คมิสฺสามี”ติ, ตสฺส เต อารามคตสฺส ยํ ตชฺชํ วีริยํ, ตํ ปฏิปฺปสฺสทฺธนฺติ.\csromanspacer{} เอวํ โภ.\csromanspacer{} อโหสิ เต ปุพฺเพ จิตฺตํ “อารามํ คมิสฺสามี”ติ, ตสฺส เต อารามคตสฺส ยํ ตชฺชํ จิตฺตํ, ตํ ปฏิปฺปสฺสทฺธนฺติ.\csromanspacer{} เอวํ โภ.\csromanspacer{} อโหสิ เต ปุพฺเพ วีมํสา “อารามํ คมิสฺสามี”ติ, ตสฺส เต อารามคตสฺส ยา ตชฺชา วีมํสา, สา ปฏิปฺปสฺสทฺธาติ.\csromanspacer{} เอวํ โภ.}

\prose{เอวเมว โข พฺราหฺมณ โย โส ภิกฺขุ อรหํ ขีณาสโว วุสิตวา กตกรณีโย โอหิตภาโร อนุปฺปตฺตสทตฺโถ ปริกฺขีณภวสํโยชโน สมฺมทญฺญาวิมุตฺโต, ตสฺส โย ปุพฺเพ ฉนฺโท อโหสิ อรหตฺตปฺปตฺติยา, อรหตฺตปฺปตฺเต\footnote{อรหตฺเต ปตฺเต (สี, สฺยา, กํ)} โย ตชฺโช ฉนฺโท, โส ปฏิปฺปสฺสทฺโธ.\csromanspacer{} ยํ ปุพฺเพ วีริยํ อโหสิ อรหตฺตปฺปตฺติยา, อรหตฺตปฺปตฺเต ยํ ตชฺชํ วีริยํ, ตํ ปฏิปฺปสฺสทฺธํ.\csromanspacer{} ยํ ปุพฺเพ จิตฺตํ อโหสิ อรหตฺตปฺปตฺติยา, อรหตฺตปฺปตฺเต ยํ ตชฺชํ จิตฺตํ, ตํ ปฏิปฺปสฺสทฺธํ.\csromanspacer{} ยา ปุพฺเพ วีมํสา อโหสิ อรหตฺตปฺปตฺติยา, อรหตฺตปฺปตฺเต ยา ตชฺชา วีมํสา, สา ปฏิปฺปสฺสทฺธา.\csromanspacer{} ตํ กิํ มญฺญสิ พฺราหฺมณ, อิติ เอวํ สนฺเต สนฺตกํ วา โหติ โน อสนฺตกํ วาติ.}

\prose{อทฺธา โภ อานนฺท เอวํ สนฺเต สนฺตกํ โหติ โน อสนฺตกํ.\csromanspacer{} อภิกฺกนฺตํ โภ อานนฺท, อภิกฺกนฺตํ โภ อานนฺท, เสยฺยถาปิ โภ อานนฺท นิกฺกุชฺชิตํ วา อุกฺกุชฺเชยฺย, ปฏิจฺฉนฺนํ วา วิวเรยฺย, มูฬฺหสฺส วา มคฺคํ อาจิกฺเขยฺย, อนฺธกาเร วา เตลปชฺโชตํ ธาเรยฺย “จกฺขุมนฺโต รูปานิ ทกฺขนฺตี”ติ, เอวเมวํ โภตา อานนฺเทน อเนกปริยาเยน ธมฺโม ปกาสิโต, เอสาหํ โภ อานนฺท ตํ ภวนฺตํ โคตมํ สรณํ คจฺฉามิ ธมฺมญฺจ ภิกฺขุสํฆญฺจ, อุปาสกํ มํ ภวํ อานนฺโท ธาเรตุ อชฺชตคฺเค ปาณุเปตํ สรณํ คตนฺติ.\csromanspacer{}\csromanspacer{}ปญฺจมํ.}

\csromanheader{2}{7. อิทฺธิปาทสํยุตฺต}
\csromanheader{2}{6. ปฐมสมณพฺราหฺมณสุตฺต}
\csromanmatikaanchor{mtk.179}
\csromantocmarkref{subparagraph}{6-9. ปฐมสมณพฺราหฺมณสุตฺตาทีนิ}{mtk.179}
\bookmark[dest={mtk.179},level=5]{6-9. ปฐมสมณพฺราหฺมณสุตฺตาทีนิ}
\proseitem{828}{\csromanfolio{239}เย หิ เกจิ ภิกฺขเว อตีตมทฺธานํ สมณา วา พฺราหฺมณา วา มหิทฺธิกา อเหสุํ มหานุภาวา, สพฺเพ เต จตุนฺนํ อิทฺธิปาทานํ ภาวิตตฺตา พหุลีกตตฺตา.\csromanspacer{} เย หิ เกจิ ภิกฺขเว อนาคตมทฺธานํ สมณา วา พฺราหฺมณา วา มหิทฺธิกา ภวิสฺสนฺติ มหานุภาวา, สพฺเพ เต จตุนฺนํ อิทฺธิปาทานํ ภาวิตตฺตา พหุลีกตตฺตา.\csromanspacer{} เย หิ เกจิ ภิกฺขเว เอตรหิ สมณา วา พฺราหฺมณา วา มหิทฺธิกา มหานุภาวา, สพฺเพ เต จตุนฺนํ อิทฺธิปาทานํ ภาวิตตฺตา พหุลีกตตฺตา.}

\prose{กตเมสํ จตุนฺนํ, อิธ ภิกฺขเว ภิกฺขุ ฉนฺทสมาธิปฺปธานสงฺขารสมนฺนาคตํ อิทฺธิปาทํ ภาเวติ.\csromanspacer{} วีริยสมาธิ ฯเปฯ.\csromanspacer{} จิตฺตสมาธิ.\csromanspacer{} วีมํสาสมาธิปฺปธานสงฺขารสมนฺนาคตํ อิทฺธิปาทํ ภาเวติ.\csromanspacer{} เย หิ เกจิ ภิกฺขเว อตีตมทฺธานํ สมณา วา พฺราหฺมณา วา มหิทฺธิกา อเหสุํ มหานุภาวา, สพฺเพ เต อิเมสํเยว จตุนฺนํ อิทฺธิปาทานํ ภาวิตตฺตา พหุลีกตตฺตา.\csromanspacer{} เย หิ เกจิ ภิกฺขเว อนาคตมทฺธานํ สมณา วา พฺราหฺมณา วา มหิทฺธิกา ภวิสฺสนฺติ มหานุภาวา, สพฺเพ เต อิเมสํเยว จตุนฺนํ อิทฺธิปาทานํ ภาวิตตฺตา พหุลีกตตฺตา.\csromanspacer{} เย หิ เกจิ ภิกฺขเว เอตรหิ สมณา วา พฺราหฺมณา วา มหิทฺธิกา มหานุภาวา, สพฺเพ เต อิเมสํเยว จตุนฺนํ อิทฺธิปาทานํ ภาวิตตฺตา พหุลีกตตฺตาติ.\csromanspacer{}\csromanspacer{}ฉฏฺฐํ.}

\csromanheader{2}{7. ทุติยสมณพฺราหฺมณสุตฺต}
\proseitem{829}{เย หิ เกจิ ภิกฺขเว อตีตมทฺธานํ สมณา วา พฺราหฺมณา วา อเนกวิหิตํ อิทฺธิวิธํ ปจฺจนุโภสุํ, เอโกปิ หุตฺวา พหุธา อเหสุํ, พหุธาปิ หุตฺวา เอโก อเหสุํ, อาวิภาวํ ติโรภาวํ ติโรกุฏฺฏํ ติโรปาการํ ติโรปพฺพตํ อสชฺชมานา อคมํสุ เสยฺยถาปิ อากาเส.\csromanspacer{} ปถวิยาปิ อุมฺมุชฺชนิมุชฺชํ อกํสุ เสยฺยถาปิ อุทเก.\csromanspacer{} อุทเกปิ อภิชฺชมาเน\footnote{อภิชฺชมานา (สี, อิ, ก)} อคมํสุ เสยฺยถาปิ ปถวิยํ.\csromanspacer{} อากาเสปิ ปลฺลงฺเกน กมิํสุ เสยฺยถาปิ ปกฺขี สกุโณ.\csromanspacer{} อิเมปิ จนฺทิมสูริเย เอวํมหิทฺธิเก เอวํมหานุภาเว ปาณินา ปริมสิํสุ\footnote{ปรามสิํสุ (สฺยา, กํ, ก)} ปริมชฺชิํสุ, ยาว พฺรหฺมโลกาปิ กาเยน วสํ วตฺเตสุํ, สพฺเพ เต จตุนฺนํ อิทฺธิปาทานํ ภาวิตตฺตา พหุลีกตตฺตา.}

\prose{\csromanfolio{240}เย หิ เกจิ ภิกฺขเว อนาคตมทฺธานํ สมณา วา พฺราหฺมณา วา อเนกวิหิตํ อิทฺธิวิธํ ปจฺจนุโภสฺสนฺติ, เอโกปิ หุตฺวา พหุธา ภวิสฺสนฺติ, พหุธาปิ หุตฺวา เอโก ภวิสฺสนฺติ, อาวิภาวํ ติโรภาวํ ติโรกุฏฺฏํ ติโรปาการํ ติโรปพฺพตํ อสชฺชมานา คมิสฺสนฺติ เสยฺยถาปิ อากาเส.\csromanspacer{} ปถวิยาปิ อุมฺมุชฺชนิมุชฺชํ กริสฺสนฺติ เสยฺยถาปิ อุทเก.\csromanspacer{} อุทเกปิ อภิชฺชมาเน คมิสฺสนฺติ เสยฺยถาปิ ปถวิยํ.\csromanspacer{} อากาเสปิ ปลฺลงฺเกน กมิสฺสนฺติ เสยฺยถาปิ ปกฺขี สกุโณ.\csromanspacer{} อิเมปิ จนฺทิมสูริเย เอวํมหิทฺธิเก เอวํมหานุภาเว ปาณินา ปริมสิสฺสนฺติ ปริมชฺชิสฺสนฺติ, ยาว พฺรหฺมโลกาปิ กาเยน วสํ วตฺติสฺสนฺติ, สพฺเพ เต จตุนฺนํ อิทฺธิปาทานํ ภาวิตตฺตา พหุลีกตตฺตา.}

\prose{เย หิ เกจิ ภิกฺขเว เอตรหิ สมณา วา พฺราหฺมณา วา อเนกวิหิตํ อิทฺธิวิธํ ปจฺจนุโภนฺติ, เอโกปิ หุตฺวา พหุธา โหนฺติ, พหุธาปิ หุตฺวา เอโก โหนฺติ, อาวิภาวํ ติโรภาวํ ติโรกุฏฺฏํ ติโรปาการํ ติโรปพฺพตํ อสชฺชมานา คจฺฉนฺติ เสยฺยถาปิ อากาเส.\csromanspacer{} ปถวิยาปิ อุมฺมุชฺชนิมุชฺชํ กโรนฺติ เสยฺยถาปิ อุทเก.\csromanspacer{} อุทเกปิ อภิชฺชมาเน คจฺฉนฺติ เสยฺยถาปิ ปถวิยํ.\csromanspacer{} อากาเสปิ ปลฺลงฺเกน กมนฺติ เสยฺยถาปิ ปกฺขี สกุโณ.\csromanspacer{} อิเมปิ จนฺทิมสูริเย เอวํมหิทฺธิเก เอวํมหานุภาเว ปาณินา ปริมสนฺติ\footnote{ปรามสนฺติ (สฺยา, กํ)} ปริมชฺชนฺติ, ยาว พฺรหฺมโลกาปิ กาเยน วสํ วตฺเตนฺติ, สพฺเพ เต จตุนฺนํ อิทฺธิปาทานํ ภาวิตตฺตา พหุลีกตตฺตาติ.}

\prose{กตเมสํ จตุนฺนํ, อิธ ภิกฺขเว ภิกฺขุ ฉนฺทสมาธิปฺปธานสงฺขารสมนฺนาคตํ อิทฺธิปาทํ ภาเวติ.\csromanspacer{} วีริยสมาธิ ฯเปฯ.\csromanspacer{} จิตฺตสมาธิ.\csromanspacer{} วีมํสาสมาธิปฺปธานสงฺขารสมนฺนาคตํ อิทฺธิปาทํ ภาเวติ.\csromanspacer{} เย หิ เกจิ ภิกฺขเว อตีตมทฺธานํ สมณา วา พฺราหฺมณา วา อเนกวิหิตํ อิทฺธิวิธํ ปจฺจนุโภสุํ, เอโกปิ หุตฺวา พหุธา อเหสุํ ฯเปฯ ยาว พฺรหฺมโลกาปิ กาเยน วสํ วตฺเตสุํ, สพฺเพ เต อิเมสํเยว จตุนฺนํ อิทฺธิปาทานํ ภาวิตตฺตา พหุลีกตตฺตา.}

\prose{เย หิ เกจิ ภิกฺขเว อนาคตมทฺธานํ สมณา วา พฺราหฺมณา วา อเนกวิหิตํ อิทฺธิวิธํ ปจฺจนุโภสฺสนฺติ, เอโกปิ หุตฺวา พหุธา ภวิสฺสนฺติ ฯเปฯ ยาว พฺรหฺมโลกาปิ กาเยน วสํ วตฺติสฺสนฺติ, สพฺเพ เต อิเมสํเยว จตุนฺนํ อิทฺธิปาทานํ ภาวิตตฺตา พหุลีกตตฺตา.}

\csromanheader{2}{7. อิทฺธิปาทสํยุตฺต}
\prose{\csromanfolio{241}เย หิ เกจิ ภิกฺขเว เอตรหิ สมณา วา พฺราหฺมณา วา อเนกวิหิตํ อิทฺธิวิธํ ปจฺจนุโภนฺติ, เอโกปิ หุตฺวา พหุธา โหนฺติ ฯเปฯ ยาว พฺรหฺมโลกาปิ กาเยน วสํ วตฺเตนฺติ, สพฺเพ เต อิเมสํเยว จตุนฺนํ อิทฺธิปาทานํ ภาวิตตฺตา พหุลีกตตฺตาติ.\csromanspacer{}\csromanspacer{}สตฺตมํ.}

\csromanheader{2}{8. ภิกฺขุสุตฺต}
\proseitem{830}{จตุนฺนํ ภิกฺขเว อิทฺธิปาทานํ ภาวิตตฺตา พหุลีกตตฺตา ภิกฺขุ อาสวานํ ขยา อนาสวํ เจโตวิมุตฺติํ ปญฺญาวิมุตฺติํ ทิฏฺเฐว ธมฺเม สยํ อภิญฺญา สจฺฉิกตฺวา อุปสมฺปชฺช วิหรติ.}

\prose{กตเมสํ จตุนฺนํ, อิธ ภิกฺขเว ภิกฺขุ ฉนฺทสมาธิปฺปธานสงฺขารสมนฺนาคตํ อิทฺธิปาทํ ภาเวติ.\csromanspacer{} วีริยสมาธิ ฯเปฯ.\csromanspacer{} จิตฺตสมาธิ.\csromanspacer{} วีมํสาสมาธิปฺปธานสงฺขารสมนฺนาคตํ อิทฺธิปาทํ ภาเวติ.\csromanspacer{} อิเมสํ โข ภิกฺขเว จตุนฺนํ อิทฺธิปาทานํ ภาวิตตฺตา พหุลีกตตฺตา ภิกฺขุ อาสวานํ ขยา อนาสวํ เจโตวิมุตฺติํ ปญฺญาวิมุตฺติํ ทิฏฺเฐว ธมฺเม สยํ อภิญฺญา สจฺฉิกตฺวา อุปสมฺปชฺช วิหรตีติ.\csromanspacer{}\csromanspacer{}อฏฺฐมํ.}

\csromanheader{2}{9. อิทฺธาทิเทสนาสุตฺต}
\proseitem{831}{อิทฺธิํ โว ภิกฺขเว เทเสสฺสามิ อิทฺธิปาทญฺจ อิทฺธิปาทภาวนญฺจ อิทฺธิปาทภาวนาคามินิญฺจ ปฏิปทํ, ตํ สุณาถ.}

\prose{กตมา จ ภิกฺขเว อิทฺธิ, อิธ ภิกฺขเว ภิกฺขุ อเนกวิหิตํ อิทฺธิวิธํ ปจฺจนุโภติ, เอโกปิ หุตฺวา พหุธา โหติ, พหุธาปิ หุตฺวา เอโก โหติ ฯเปฯ ยาว พฺรหฺมโลกาปิ กาเยน วสํ วตฺเตติ, อยํ วุจฺจติ ภิกฺขเว อิทฺธิ.}

\prose{กตโม จ ภิกฺขเว อิทฺธิปาโท, โย โส ภิกฺขเว มคฺโค ยา ปฏิปทา อิทฺธิลาภาย อิทฺธิปฏิลาภาย สํวตฺตติ, อยํ วุจฺจติ ภิกฺขเว อิทฺธิปาโท.}

\prose{กตมา จ ภิกฺขเว อิทฺธิปาทภาวนา, อิธ ภิกฺขเว ภิกฺขุ ฉนฺทสมาธิปฺปธานสงฺขารสมนฺนาคตํ อิทฺธิปาทํ ภาเวติ.\csromanspacer{} วีริยสมาธิ ฯเปฯ.\csromanspacer{} จิตฺตสมาธิ.\csromanspacer{} วีมํสาสมาธิปฺปธานสงฺขารสมนฺนาคตํ อิทฺธิปาทํ ภาเวติ, อยํ วุจฺจติ ภิกฺขเว อิทฺธิปาทภาวนา.}

\prose{\csromanfolio{242}กตมา จ ภิกฺขเว อิทฺธิปาทภาวนาคามินี ปฏิปทา, อยเมว อริโย อฏฺฐงฺคิโก มคฺโค.\csromanspacer{} เสยฺยถิทํ, สมฺมาทิฏฺฐิ สมฺมาสงฺกปฺโป สมฺมาวาจา สมฺมากมฺมนฺโต สมฺมา-อาชีโว สมฺมาวายาโม สมฺมาสติ สมฺมาสมาธิ.\csromanspacer{} อยํ วุจฺจติ ภิกฺขเว อิทฺธิปาทภาวนาคามินี ปฏิปทาติ.\csromanspacer{}\csromanspacer{}นวมํ.}

\csromanmatikaanchor{mtk.180}
\csromantocmarkref{subparagraph}{10. วิภงฺคสุตฺต}{mtk.180}
\bookmark[dest={mtk.180},level=5]{10. วิภงฺคสุตฺต}
\csromanheader{6}{10. วิภงฺคสุตฺต}
\proseitem{832}{จตฺตาโรเม ภิกฺขเว อิทฺธิปาทา ภาวิตา พหุลีกตา มหปฺผลา โหนฺติ มหานิสํสา.}

\prose{กถํ ภาวิตา จ ภิกฺขเว จตฺตาโร อิทฺธิปาทา กถํ พหุลีกตา มหปฺผลา โหนฺติ มหานิสํสา, อิธ ภิกฺขเว ภิกฺขุ ฉนฺทสมาธิปฺปธานสงฺขารสมนฺนาคตํ อิทฺธิปาทํ ภาเวติ, อิติ เม ฉนฺโท น จ อติลีโน ภวิสฺสติ, น จ อติปฺปคฺคหิโต ภวิสฺสติ, น จ อชฺฌตฺตํ สํขิตฺโต ภวิสฺสติ, น จ พหิทฺธา วิกฺขิตฺโต ภวิสฺสติ.\csromanspacer{} ปจฺฉาปุเรสญฺญี จ วิหรติ, ยถา ปุเร ตถา ปจฺฉา, ยถา ปจฺฉา ตถา ปุเร.\csromanspacer{} ยถา อโธ ตถา อุทฺธํ, ยถา อุทฺธํ ตถา อโธ.\csromanspacer{} ยถา ทิวา ตถา รตฺติํ, ยถา รตฺติํ ตถา ทิวา.\csromanspacer{} อิติ วิวเฏน เจตสา อปริโยนทฺเธน สปฺปภาสํ จิตฺตํ ภาเวติ.\csromanspacer{} วีริยสมาธิ ฯเปฯ.\csromanspacer{} จิตฺตสมาธิ.\csromanspacer{} วีมํสาสมาธิปฺปธานสงฺขารสมนฺนาคตํ อิทฺธิปาทํ ภาเวติ, อิติ เม วีมํสา น จ อติลีนา ภวิสฺสติ, น จ อติปฺปคฺคหิตา ภวิสฺสติ, น จ อชฺฌตฺตํ สํขิตฺตา ภวิสฺสติ, น จ พหิทฺธา วิกฺขิตฺตา ภวิสฺสติ.\csromanspacer{} ปจฺฉาปุเรสญฺญี จ วิหรติ, ยถา ปุเร ตถา ปจฺฉา, ยถา ปจฺฉา ตถา ปุเร.\csromanspacer{} ยถา อโธ ตถา อุทฺธํ, ยถา อุทฺธํ ตถา อโธ.\csromanspacer{} ยถา ทิวา ตถา รตฺติํ, ยถา รตฺติํ ตถา ทิวา.\csromanspacer{} อิติ วิวเฏน เจตสา อปริโยนทฺเธน สปฺปภาสํ จิตฺตํ ภาเวติ.}

\prose{กตโม จ ภิกฺขเว อติลีโน ฉนฺโท, โย ภิกฺขเว ฉนฺโท โกสชฺชสหคโต โกสชฺชสมฺปยุตฺโต, อยํ วุจฺจติ ภิกฺขเว อติลีโน ฉนฺโท.}

\prose{กตโม จ ภิกฺขเว อติปฺปคฺคหิโต ฉนฺโท, โย ภิกฺขเว ฉนฺโท อุทฺธจฺจสหคโต อุทฺธจฺจสมฺปยุตฺโต, อยํ วุจฺจติ ภิกฺขเว อติปฺปคฺคหิโต ฉนฺโท.}

\csromanheader{2}{7. อิทฺธิปาทสํยุตฺต}
\prose{\csromanfolio{243}กตโม จ ภิกฺขเว อชฺฌตฺตํ สํขิตฺโต ฉนฺโท, โย ภิกฺขเว ฉนฺโท ถินมิทฺธสหคโต ถินมิทฺธสมฺปยุตฺโต, อยํ วุจฺจติ ภิกฺขเว อชฺฌตฺตํ สํขิตฺโต ฉนฺโท.}

\prose{กตโม จ ภิกฺขเว พหิทฺธา วิกฺขิตฺโต ฉนฺโท, โย ภิกฺขเว ฉนฺโท พหิทฺธา ปญฺจ กามคุเณ อารพฺภ อนุวิกฺขิตฺโต อนุวิสโฏ, อยํ วุจฺจติ ภิกฺขเว พหิทฺธา วิกฺขิตฺโต ฉนฺโท.}

\prose{กถญฺจ ภิกฺขเว ภิกฺขุ ปจฺฉาปุเรสญฺญี จ วิหรติ, ยถา ปุเร ตถา ปจฺฉา, ยถา ปจฺฉา ตถา ปุเร, อิธ ภิกฺขเว ภิกฺขุโน ปจฺฉาปุเรสญฺญา สุคฺคหิตา โหติ สุมนสิกตา สูปธาริตา สุปฺปฏิวิทฺธา ปญฺญาย.\csromanspacer{} เอวํ โข ภิกฺขเว ภิกฺขุ ปจฺฉาปุเรสญฺญี จ วิหรติ, ยถา ปุเร ตถา ปจฺฉา, ยถา ปจฺฉา ตถา ปุเร.}

\prose{กถญฺจ ภิกฺขเว ภิกฺขุ ยถา อโธ ตถา อุทฺธํ, ยถา อุทฺธํ ตถา อโธ วิหรติ, อิธ ภิกฺขเว ภิกฺขุ อิมเมว กายํ อุทฺธํ ปาทตลา อโธ เกสมตฺถกา ตจปริยนฺตํ ปูรํ นานปฺปการสฺส อสุจิโน ปจฺจเวกฺขติ “อตฺถิ อิมสฺมิํ กาเย เกสา โลมา นขา ทนฺตา ตโจ, มํสํ นฺหารุ อฏฺฐิ อฏฺฐิมิญฺชํ วกฺกํ, หทยํ ยกนํ กิโลมกํ ปิหกํ ปปฺผาสํ, อนฺตํ อนฺตคุณํ อุทริยํ กรีสํ, ปิตฺตํ เสมฺหํ ปุพฺโพ โลหิตํ เสโท เมโท, อสฺสุ วสา เขโฬ สิงฺฆาณิกา ลสิกา มุตฺตนฺ”ติ.\csromanspacer{} เอวํ โข ภิกฺขเว ภิกฺขุ ยถา อโธ ตถา อุทฺธํ, ยถา อุทฺธํ ตถา อโธ วิหรติ.}

\prose{กถญฺจ ภิกฺขเว ภิกฺขุ ยถา ทิวา ตถา รตฺติํ, ยถา รตฺติํ ตถา ทิวา วิหรติ, อิธ ภิกฺขเว ภิกฺขุ เยหิ อากาเรหิ เยหิ ลิงฺเคหิ เยหิ นิมิตฺเตหิ ทิวา ฉนฺทสมาธิปฺปธานสงฺขารสมนฺนาคตํ อิทฺธิปาทํ ภาเวติ, โส เตหิ อากาเรหิ เตหิ ลิงฺเคหิ เตหิ นิมิตฺเตหิ รตฺติํ ฉนฺทสมาธิปฺปธานสงฺขารสมนฺนาคตํ อิทฺธิปาทํ ภาเวติ.\csromanspacer{} เยหิ วา ปน อากาเรหิ เยหิ ลิงฺเคหิ เยหิ นิมิตฺเตหิ รตฺติํ ฉนฺทสมาธิปฺปธานสงฺขารสมนฺนาคตํ อิทฺธิปาทํ ภาเวติ, โส เตหิ อากาเรหิ เตหิ ลิงฺเคหิ เตหิ นิมิตฺเตหิ ทิวา ฉนฺทสมาธิปฺปธานสงฺขารสมนฺนาคตํ อิทฺธิปาทํ ภาเวติ.\csromanspacer{} เอวํ โข ภิกฺขเว ภิกฺขุ ยถา ทิวา ตถา รตฺติํ, ยถา รตฺติํ ตถา ทิวา วิหรติ.}

\prose{\csromanfolio{244}กถญฺจ ภิกฺขเว ภิกฺขุ วิวเฏน เจตสา อปริโยนทฺเธน สปฺปภาสํ จิตฺตํ ภาเวติ, อิธ ภิกฺขเว ภิกฺขุโน อาโลกสญฺญา สุคฺคหิตา โหติ, ทิวาสญฺญา สฺวาธิฏฺฐิตา.\csromanspacer{} เอวํ โข ภิกฺขเว ภิกฺขุ วิวเฏน เจตสา อปริโยนทฺเธน สปฺปภาสํ จิตฺตํ ภาเวติ.}

\prose{กถมญฺจ ภิกฺขเว อติลีนํ วีริยํ, ยํ ภิกฺขเว วีริยํ โกสชฺชสหคตํ โกสชฺชสมฺปยุตฺตํ, อิทํ วุจฺจติ ภิกฺขเว อติลีนํ วีริยํ.}

\prose{กตมญฺจ ภิกฺขเว อติปฺปคฺคหิตํ วีริยํ, ยํ ภิกฺขเว วีริยํ อุทฺธจฺจสหคตํ อุทฺธจฺจสมฺปยุตฺตํ, อิทํ วุจฺจติ ภิกฺขเว อติปฺปคฺคหิตํ วีริยํ.}

\prose{กตมญฺจ ภิกฺขเว อชฺฌตฺตํ สํขิตฺตํ วีริยํ, ยํ ภิกฺขเว วีริยํ ถินมิทฺธสหคตํ ถินมิทฺธสมฺปยุตฺตํ, อิทํ วุจฺจติ ภิกฺขเว อชฺฌตฺตํ สํขิตฺตํ วีริยํ.}

\prose{กตมญฺจ ภิกฺขเว พหิทฺธา วิกฺขิตฺตํ วีริยํ, ยํ ภิกฺขเว วีริยํ พหิทฺธา ปญฺจ กามคุเณ อารพฺภ อนุวิกฺขิตฺตํ อนุวิสฏํ, อิทํ วุจฺจติ ภิกฺขเว พหิทฺธา วิกฺขิตฺตํ วีริยํ ฯเปฯ.}

\prose{กถญฺจ ภิกฺขเว ภิกฺขุ วิวเฏน เจตสา อปริโยนทฺเธน สปฺปภาสํ จิตฺตํ ภาเวติ, อิธ ภิกฺขเว ภิกฺขุโน อาโลกสญฺญา สุคฺคหิตา โหติ, ทิวาสญฺญา สฺวาธิฏฺฐิตา.\csromanspacer{} เอวํ โข ภิกฺขเว ภิกฺขุ วิวเฏน เจตสา อปริโยนทฺเธน สปฺปภาสํ จิตฺตํ ภาเวติ.}

\prose{กถมญฺจ ภิกฺขเว อติลีนํ จิตฺตํ, ยํ ภิกฺขเว จิตฺตํ โกสชฺชสหคตํ โกสชฺชสมฺปยุตฺตํ, อิทํ วุจฺจติ ภิกฺขเว อติลีนํ จิตฺตํ.}

\prose{กตมญฺจ ภิกฺขเว อติปฺปคฺคหิตํ จิตฺตํ, ยํ ภิกฺขเว จิตฺตํ อุทฺธจฺจสหคตํ อุทฺธจฺจสมฺปยุตฺตํ, อิทํ วุจฺจติ ภิกฺขเว อติปฺปคฺคหิตํ จิตฺตํ.}

\prose{กตมญฺจ ภิกฺขเว อชฺฌตฺตํ สํขิตฺตํ จิตฺตํ, ยํ ภิกฺขเว จิตฺตํ ถินมิทฺธสหคตํ ถินมิทฺธสมฺปยุตฺตํ, อิทํ วุจฺจติ ภิกฺขเว อชฺฌตฺตํ สํขิตฺตํ จิตฺตํ.}

\prose{กตมญฺจ ภิกฺขเว พหิทฺธา วิกฺขิตฺตํ จิตฺตํ, ยํ ภิกฺขเว จิตฺตํ พหิทฺธา ปญฺจ กามคุเณ อารพฺภ อนุวิกฺขิตฺตํ อนุวิสฏํ, อิทํ วุจฺจติ ภิกฺขเว พหิทฺธา วิกฺขิตฺตํ จิตฺตํ ฯเปฯ.\csromanspacer{} เอวํ โข ภิกฺขเว ภิกฺขุ วิวเฏน เจตสา อปริโยนทฺเธน สปฺปภาสํ จิตฺตํ ภาเวติ.}

\prose{กตมา จ ภิกฺขเว อติลีนา วีมํสา, ยา ภิกฺขเว วีมํสา โกสชฺชสหคตา โกสชฺชสมฺปยุตฺตา, อยํ วุจฺจติ ภิกฺขเว อติลีนา วิมํสา.}

\csromanheader{2}{7. อิทฺธิปาทสํยุตฺต}
\prose{\csromanfolio{245}กตมา จ ภิกฺขเว อติปฺปคฺคหิตา วีมํสา, ยา ภิกฺขเว วีมํสา อุทฺธจฺจสหคตา อุทฺธจฺจสมฺปยุตฺตา, อยํ วุจฺจติ ภิกฺขเว อติปฺปคฺคหิตา วีมํสา.}

\prose{กตมา จ ภิกฺขเว อชฺฌตฺตํ สํขิตฺตา วีมํสา, ยา ภิกฺขเว วีมํสา ถินมิทฺธสหคตา ถินมิทฺธสมฺปยุตฺตา, อยํ วุจฺจติ ภิกฺขเว อชฺฌตฺตํ สํขิตฺตา วีมํสา.}

\prose{กตมา จ ภิกฺขเว พหิทฺธา วิกฺขิตฺตา วีมํสา, ยา ภิกฺขเว วีมํสา พหิทฺธา ปญฺจ กามคุเณ อารพฺภ อนุวิกฺขิตฺตา อนุวิสฏา, อยํ วุจฺจติ ภิกฺขเว พหิทฺธา วิกฺขิตฺตา วีมํสา ฯเปฯ.\csromanspacer{} เอวํ โข ภิกฺขเว ภิกฺขุ วิวเฏน เจตสา อปริโยนทฺเธน สปฺปภาสํ จิตฺตํ ภาเวติ, เอวํ ภาวิตา โข ภิกฺขเว จตฺตาโร อิทฺธิปาทา เอวํ พหุลีกตา มหปฺผลา โหนฺติ มหานิสํสา.}

\prose{เอวํ ภาวิเตสุ โข ภิกฺขเว ภิกฺขุ จตูสุ อิทฺธิปาเทสุ เอวํ พหุลีกเตสุ อเนกวิหิตํ อิทฺธิวิธํ ปจฺจนุโภติ, เอโกปิ หุตฺวา พหุธา โหติ, พหุธาปิ หุตฺวา เอโก โหติ ฯเปฯ ยาว พฺรหฺมโลกาปิ กาเยน วสํ วตฺเตติ.\csromanspacer{} เอวํ ภาวิเตสุ โข ภิกฺขเว ภิกฺขุ จตูสุ อิทฺธิปาเทสุ เอวํ พหุลีกเตสุ อาสวานํ ขยา อนาสวํ เจโตวิมุตฺติํ ปญฺญาวิมุตฺติํ ทิฏฺเฐว ธมฺเม สยํ อภิญฺญา สจฺฉิกตฺวา อุปสมฺปชฺช วิหรตีติ.\csromanspacer{}\csromanspacer{}ทสมํ.\csromanspacer{} ปาสาทกมฺปนวคฺโค ทุติโย.}

\titlehead{\textbf{ตสฺสุทฺทานํ}}
\csromangathameasure{ปุพฺพํ มหปฺผลํ ฉนฺทํ, โมคฺคลฺลานญฺจ อุณฺณาภํ. \\ เทฺว สมณพฺราหฺมณา ภิกฺขุ, เทสนา วิภงฺเคน จาติ.}
\csromangathagroup{\csromangathastanza{ปุพฺพํ มหปฺผลํ ฉนฺทํ, โมคฺคลฺลานญฺจ อุณฺณาภํ. \\ เทฺว สมณพฺราหฺมณา ภิกฺขุ, เทสนา วิภงฺเคน จาติ.}}

\csromanmatikaanchor{mtk.181}
\csromantocmarknopage{subparagraph}{3. อโยคุฬวคฺค}{mtk.181}
\bookmark[dest={mtk.181},level=5]{3. อโยคุฬวคฺค}
\csromanheader{6}{3. อโยคุฬวคฺค}
\csromanmatikaanchor{mtk.182}
\csromantocmarkref{subparagraph}{1. มคฺคสุตฺต}{mtk.182}
\bookmark[dest={mtk.182},level=5]{1. มคฺคสุตฺต}
\csromanheader{6}{1. มคฺคสุตฺต}
\csromanmatikaanchor{mtk.183}
\csromantocmarkref{subparagraph}{2-10. อโยคุฬสุตฺตาทีนิ}{mtk.183}
\bookmark[dest={mtk.183},level=5]{2-10. อโยคุฬสุตฺตาทีนิ}
\proseitem{833}{สาวตฺถินิทานํ.\csromanspacer{} ปุพฺเพว เม ภิกฺขเว สมฺโพธา อนภิสมฺพุทฺธสฺส โพธิสตฺตสฺเสว สโต เอตทโหสิ “โก นุ โข มคฺโค กา ปฏิปทา \csromanfolio{246}อิทฺธิปาทภาวนายา”ติ.\csromanspacer{} ตสฺส มยฺหํ ภิกฺขเว เอตทโหสิ\csromandash{}อิธ ภิกฺขุ ฉนฺทสมาธิปฺปธานสงฺขารสมนฺนาคตํ อิทฺธิปาทํ ภาเวติ, อิติ เม ฉนฺโท น จ อติลีโน ภวิสฺสติ, น จ อติปฺปคฺคหิโต ภวิสฺสติ, น จ อชฺฌตฺตํ สํขิตฺโต ภวิสฺสติ, น จ พหิทฺธา วิกฺขิตฺโต ภวิสฺสติ.\csromanspacer{} ปจฺฉาปุเรสญฺญี จ วิหรติ, ยถา ปุเร ตถา ปจฺฉา, ยถา ปจฺฉา ตถา ปุเร.\csromanspacer{} ยถา อโธ ตถา อุทฺธํ, ยถา อุทฺธํ ตถา อโธ.\csromanspacer{} ยถา ทิวา ตถา รตฺติํ ยถา รตฺติํ ตถา ทิวา.\csromanspacer{} อิติ วิวเฏน เจตสา อปริโยนทฺเธน สปฺปภาสํ จิตฺตํ ภาเวติ.\csromanspacer{} วีริยสมาธิ ฯเปฯ.\csromanspacer{} จิตฺตสมาธิ.\csromanspacer{} วีมํสาสมาธิปฺปธานสงฺขารสมนฺนาคตํ อิทฺธิปาทํ ภาเวติ, อิติ เม วีมํสา น จ อติลีนา ภวิสฺสติ, น จ อติปฺปคฺคหิตา ภวิสฺสติ, น จ อชฺฌตฺตํ สํขิตฺตา ภวิสฺสติ, น จ พหิทฺธา วิกฺขิตฺตา ภวิสฺสติ.\csromanspacer{} ปจฺฉาปุเรสญฺญี จ วิหรติ, ยถา ปุเร ตถา ปจฺฉา, ยถา ปจฺฉา ตถา ปุเร.\csromanspacer{} ยถา อโธ ตถา อุทฺธํ, ยถา อุทฺธํ ตถา อโธ.\csromanspacer{} ยถา ทิวา ตถา รตฺติํ, ยถา รตฺติํ ตถา ทิวา.\csromanspacer{} อิติ วิวเฏน เจตสา อปริโยนทฺเธน สปฺปภาสํ จิตฺตํ ภาเวติ.}

\prose{เอวํ ภาวิเตสุ โข ภิกฺขเว ภิกฺขุ จตูสุ อิทฺธิปาเทสุ เอวํ พหุลีกเตสุ อเนกวิหิตํ อิทฺธิวิธํ ปจฺจนุโภติ, เอโกปิ หุตฺวา พหุธา โหติ, พหุธาปิ หุตฺวา เอโก โหติ ฯเปฯ ยาว พฺรหฺมโลกาปิ กาเยน วสํ วตฺเตติ.\csromanspacer{} เอวํ ภาวิเตสุ โข ภิกฺขเว ภิกฺขุ จตูสุ อิทฺธิปาเทสุ เอวํ พหุลีกเตสุ อาสวานํ ขยา อนาสวํ เจโตวิมุตฺติํ ปญฺญาวิมุตฺติํ ทิฏฺเฐว ธมฺเม สยํ อภิญฺญา สจฺฉิกตฺวา อุปสมฺปชฺช วิหรตีติ.\csromanspacer{}\csromanspacer{}ปฐมํ. (ฉปิ อภิญฺญาโย วิตฺถาเรตพฺพา.)}

\csromanheader{2}{2. อโยคุฬสุตฺต}
\proseitem{834}{สาวตฺถินิทานํ.\csromanspacer{} อถ โข อายสฺมา อานนฺโท เยน ภควา เตนุปสงฺกมิ, อุปสงฺกมิตฺวา ภควนฺตํ อภิวาเทตฺวา เอกมนฺตํ นิสีทิ, เอกมนฺตํ นิสินฺโน โข อายสฺมา อานนฺโท ภควนฺตํ เอตทโวจ “อภิชานาติ นุ โข ภนฺเต ภควา อิทฺธิยา มโนมเยน กาเยน พฺรหฺมโลกํ อุปสงฺกมิตา”ติ.\csromanspacer{} อภิชานามิ ขฺวาหํ อานนฺท อิทฺธิยา มโนมเยน กาเยน พฺรหฺมโลกํ อุปสงฺกมิตาติ.\csromanspacer{} อภิชานาติ ปน ภนฺเต ภควา อิมินา จาตุมหาภูติเกน กาเยน อิทฺธิยา พฺรหฺมโลกํ}

\csromanheader{2}{7. อิทฺธิปาทสํยุตฺต}
\prose{\csromanfolio{247}อุปสงฺกมิตาติ.\csromanspacer{} อภิชานามิ ขฺวาหํ อานนฺท อิมินา จาตุมหาภูติเกน\footnote{จาตุมฺมหาภูติเกน (สี, สฺยา, กํ)} กาเยน อิทฺธิยา พฺรหฺมโลกํ อุปสงฺกมิตาติ.}

\prose{ยญฺจ โข โอมาติ ภนฺเต ภควา อิทฺธิยา มโนมเยน กาเยน พฺรหฺมโลกํ อุปสงฺกมิตุํ, ยญฺจ โข อภิชานาติ ภนฺเต ภควา อิมินา จาตุมหาภูติเกน กาเยน อิทฺธิยา พฺรหฺมโลกํ อุปสงฺกมิตา, ตยิทํ ภนฺเต ภควโต อจฺฉริยํ เจว อพฺภุตํ จาติ.\csromanspacer{} อจฺฉริยา เจว อานนฺท ตถาคตา อจฺฉริยธมฺมสมนฺนาคตา จ, อพฺภุตา เจว อานนฺท ตถาคตา อพฺภุตธมฺมสมนฺนาคตา จ.}

\prose{ยสฺมิํ อานนฺท สมเย ตถาคโต กายมฺปิ จิตฺเต สโมทหติ\footnote{สมาทหติ (สี, สฺยา, อิ)}, จิตฺตมฺปิ กาเย สโมทหติ, สุขสญฺญญฺจ ลหุสญฺญญฺจ กาเย โอกฺกมิตฺวา วิหรติ.\csromanspacer{} ตสฺมิํ อานนฺท สมเย ตถาคตสฺส กาโย ลหุตโร เจว โหติ มุทุตโร จ กมฺมนิยตโร จ ปภสฺสรตโร จ.}

\prose{เสยฺยถาปิ อานนฺท อโยคุโฬ ทิวสํ สนฺตตฺโต ลหุตโร เจว โหติ มุทุตโร จ กมฺมนิยตโร จ ปภสฺสรตโร จ.\csromanspacer{} เอวเมว โข อานนฺท ยสฺมิํ สมเย ตถาคโต กายมฺปิ จิตฺเต สโมทหติ, จิตฺตมฺปิ กาเย สโมทหติ, สุขสญฺญญฺจ ลหุสญฺญญฺจ กาเย โอกฺกมิตฺวา วิหรติ, ตสฺมิํ อานนฺท สมเย ตถาคตสฺส กาโย ลหุตโร เจว โหติ มุทุตโร จ กมฺมนิยตโร จ ปภสฺสรตโร จ.}

\prose{ยสฺมิํ อานนฺท สมเย ตถาคโต กายมฺปิ จิตฺเต สโมทหติ, จิตฺตมฺปิ กาเย สโมทหติ, สุขสญฺญญฺจ ลหุสญฺญญฺจ กาเย โอกฺกมิตฺวา วิหรติ, ตสฺมิํ อานนฺท สมเย ตถาคตสฺส กาโย อปฺปกสิเรเนว ปถวิยา เวหาสํ อพฺภุคฺคจฺฉติ.\csromanspacer{} โส อเนกวิหิตํ อิทฺธิวิธํ ปจฺจนุโภติ, เอโกปิ หุตฺวา พหุธา โหติ, พหุธาปิ หุตฺวา เอโก โหติ ฯเปฯ ยาว พฺรหฺมโลกาปิ กาเยน วสํ วตฺเตติ.}

\prose{เสยฺยถาปิ อานนฺท ตูลปิจุ วา กปฺปาสปิจุ วา ลหุโก วาตูปาทาโน อปฺปกสิเรเนว ปถวิยา เวหาสํ อพฺภุคฺคจฺฉติ.\csromanspacer{} เอวเมว โข อานนฺท ยสฺมิํ สมเย ตถาคโต กายมฺปิ จิตฺเต สโมทหติ,}

\prose{\csromanfolio{248}จิตฺตมฺปิ กาเย สโมทหติ, สุขสญฺญญฺจ ลหุสญฺญญฺจ กาเย โอกฺกมิตฺวา วิหรติ, ตสฺมิํ อานนฺท สมเย ตถาคตสฺส กาโย อปฺปกสิเรเนว ปถวิยา เวหาสํ อพฺภุคฺคจฺฉติ.\csromanspacer{} โส อเนกวิหิตํ อิทฺธิวิธํ ปจฺจนุโภติ.\csromanspacer{} เอโกปิ หุตฺวา พหุธา โหติ ฯเปฯ ยาว พฺรหฺมโลกาปิ กาเยน วสํ วตฺเตตีติ.\csromanspacer{}\csromanspacer{}ทุติยํ.}

\csromanheader{2}{3. ภิกฺขุสุตฺต}
\proseitem{835}{จตฺตาโรเม ภิกฺขเว อิทฺธิปาทา.\csromanspacer{} กตเม จตฺตาโร, อิธ ภิกฺขเว ภิกฺขุ ฉนฺทสมาธิปฺปธานสงฺขารสมนฺนาคตํ อิทฺธิปาทํ ภาเวติ.\csromanspacer{} วีริยสมาธิ ฯเปฯ.\csromanspacer{} จิตฺตสมาธิ.\csromanspacer{} วีมํสาสมาธิปฺปธานสงฺขารสมนฺนาคตํ อิทฺธิปาทํ ภาเวติ.\csromanspacer{} อิเม โข ภิกฺขเว จตฺตาโร อิทฺธิปาทา, อิเมสํ โข ภิกฺขเว จตุนฺนํ อิทฺธิปาทานํ ภาวิตตฺตา พหุลีกตตฺตา ภิกฺขุ อาสวานํ ขยา อนาสวํ เจโตวิมุตฺติํ ปญฺญาวิมุตฺติํ ทิฏฺเฐว ธมฺเม สยํ อภิญฺญา สจฺฉิกตฺวา อุปสมฺปชฺช วิหรตีติ.\csromanspacer{}\csromanspacer{}ตติยํ.}

\csromanheader{2}{4. สุทฺธิกสุตฺต}
\proseitem{836}{จตฺตาโรเม ภิกฺขเว อิทฺธิปาทา.\csromanspacer{} กตเม จตฺตาโร, อิธ ภิกฺขเว ภิกฺขุ ฉนฺทสมาธิปฺปธานสงฺขารสมนฺนาคตํ อิทฺธิปาทํ ภาเวติ.\csromanspacer{} วีริยสมาธิ ฯเปฯ.\csromanspacer{} จิตฺตสมาธิ.\csromanspacer{} วีมํสาสมาธิปฺปธานสงฺขารสมนฺนาคตํ อิทฺธิปาทํ ภาเวติ.\csromanspacer{} อิเม โข ภิกฺขเว จตฺตาโร อิทฺธิปาทาติ.\csromanspacer{}\csromanspacer{}จตุตฺถํ.}

\csromanheader{2}{5. ปฐมผลสุตฺต}
\proseitem{837}{จตฺตาโรเม ภิกฺขเว อิทฺธิปาทา.\csromanspacer{} กตเม จตฺตาโร, อิธ ภิกฺขเว ภิกฺขุ ฉนฺทสมาธิปฺปธานสงฺขารสมนฺนาคตํ อิทฺธิปาทํ ภาเวติ.\csromanspacer{} วีริยสมาธิ ฯเปฯ.\csromanspacer{} จิตฺตสมาธิ.\csromanspacer{} วีมํสาสมาธิปฺปธานสงฺขารสมนฺนาคตํ อิทฺธิปาทํ ภาเวติ.\csromanspacer{} อิเม โข ภิกฺขเว จตฺตาโร อิทฺธิปาทา, อิเมสํ โข ภิกฺขเว จตุนฺนํ อิทฺธิปาทานํ ภาวิตตฺตา พหุลีกตตฺตา ภิกฺขุนา ทฺวินฺนํ ผลานํ อญฺญตรํ ผลํ ปาฏิกงฺขํ, ทิฏฺเฐว ธมฺเม อญฺญา, สติ วา อุปาทิเสเส อนาคามิตาติ.\csromanspacer{}\csromanspacer{}ปญฺจมํ.}

\csromanheader{2}{7. อิทฺธิปาทสํยุตฺต}
\csromanheader{2}{6. ทุติยผลสุตฺต}
\proseitem{838}{\csromanfolio{249}จตฺตาโรเม ภิกฺขเว อิทฺธิปาทา.\csromanspacer{} กตเม จตฺตาโร, อิธ ภิกฺขเว ภิกฺขุ ฉนฺทสมาธิปฺปธานสงฺขารสมนฺนาคตํ อิทฺธิปาทํ ภาเวติ.\csromanspacer{} วีริยสมาธิ ฯเปฯ.\csromanspacer{} จิตฺตสมาธิ.\csromanspacer{} วีมํสาสมาธิปฺปธานสงฺขารสมนฺนาคตํ อิทฺธิปาทํ ภาเวติ.\csromanspacer{} อิเม โข ภิกฺขเว จตฺตาโร อิทฺธิปาทา, อิเมสํ โข ภิกฺขเว จตุนฺนํ อิทฺธิปาทานํ ภาวิตตฺตา พหุลีกตตฺตา สตฺต ผลา สตฺตานิสํสา ปาฏิกงฺขา.}

\prose{กตเม สตฺต ผลา สตฺตานิสํสา, ทิฏฺเฐว ธมฺเม ปฏิกจฺจ\footnote{ปฏิคจฺจ (สี), ปฏิหจฺจ (อิ)} อญฺญํ อาราเธติ.\csromanspacer{} โน เจ ทิฏฺเฐว ธมฺเม ปฏิกจฺจ อญฺญํ อาราเธติ, อถ มรณกาเล อญฺญํ อาราเธติ.\csromanspacer{} โน เจ ทิฏฺเฐว ธมฺเม ปฏิกจฺจ อญฺญํ อาราเธติ, โน เจ มรณกาเล อญฺญํ อาราเธติ, อถ ปญฺจนฺนํ โอรมฺภาคิยานํ สํโยชนานํ ปริกฺขยา อนฺตราปรินิพฺพายี โหติ, อุปหจฺจปรินิพฺพายี โหติ, อสงฺขารปรินิพฺพายี โหติ, สสงฺขารปรินิพฺพายี โหติ, อุทฺธํโสโต โหติ อกนิฏฺฐคามี.\csromanspacer{} อิเมสํ โข ภิกฺขเว จตุนฺนํ อิทฺธิปาทานํ ภาวิตตฺตา พหุลีกตตฺตา อิเม สตฺต ผลา สตฺตานิสํสา ปาฏิกงฺขาติ.\csromanspacer{}\csromanspacer{}ฉฏฺฐํ.}

\csromanheader{2}{7. ปฐม-อานนฺทสุตฺต}
\proseitem{839}{สาวตฺถินิทานํ.\csromanspacer{} อถ โข อายสฺมา อานนฺโท เยน ภควา เตนุปสงฺกมิ, อุปสงฺกมิตฺวา ภควนฺตํ อภิวาเทตฺวา เอกมนฺตํ นิสีทิ, เอกมนฺตํ นิสินฺโน โข อายสฺมา อานนฺโท ภควนฺตํ เอตทโวจ\csromandash{}}

\prose{กตมา นุ โข ภนฺเต อิทฺธิ, กตโม อิทฺธิปาโท, กตมา อิทฺธิปาทภาวนา, กตมา อิทฺธิปาทภาวนาคามินี ปฏิปทาติ.\csromanspacer{} อิธานนฺท ภิกฺขุ อเนกวิหิตํ อิทฺธิวิธํ ปจฺจนุโภติ, เอโกปิ หุตฺวา พหุธา โหติ, พหุธาปิ หุตฺวา เอโก โหติ ฯเปฯ ยาว พฺรหฺมโลกาปิ กาเยน วสํ วตฺเตติ, อยํ วุจฺจตานนฺท อิทฺธิ.}

\prose{กตโม จานนฺท อิทฺธิปาโท, โย อานนฺท มคฺโค ยา ปฏิปทา อิทฺธิลาภาย อิทฺธิปฏิลาภาย สํวตฺตติ, อยํ วุจฺจตานนฺท อิทฺธิปาโท.}

\prose{\csromanfolio{250}กตมา จานนฺท อิทฺธิปาทภาวนา, อิธานนฺท ภิกฺขุ ฉนฺทสมาธิปฺปธานสงฺขารสมนฺนาคตํ อิทฺธิปาทํ ภาเวติ.\csromanspacer{} วีริยสมาธิ ฯเปฯ.\csromanspacer{} จิตฺตสมาธิ.\csromanspacer{} วีมํสาสมาธิปฺปธานสงฺขารสมนฺนาคตํ อิทฺธิปาทํ ภาเวติ, อยํ วุจฺจตานนฺท อิทฺธิปาทภาวนา.}

\prose{กตมา จานนฺท อิทฺธิปาทภาวนาคามินี ปฏิปทา, อยเมว อริโย อฏฺฐงฺคิโก มคฺโค, เสยฺยถิทํ, สมฺมาทิฏฺฐิ ฯเปฯ สมฺมาสมาธิ, อยํ วุจฺจตานนฺท อิทฺธิปาทภาวนาคามินี ปฏิปทาติ.\csromanspacer{}\csromanspacer{}สตฺตมํ.}

\csromanheader{2}{8. ทุติย-อานนฺทสุตฺต}
\proseitem{840}{เอกมนฺตํ นิสินฺนํ โข อายสฺมนฺตํ อานนฺทํ ภควา เอตทโวจ “กตมา นุ โข อานนฺท อิทฺธิ, กตโม อิทฺธิปาโท, กตมา อิทฺธิปาทภาวนา, กตมา อิทฺธิปาทภาวนาคามินี ปฏิปทา”ติ.\csromanspacer{} ภควํมูลกา โน ภนฺเต ธมฺมา, ภควํเนตฺติกา ฯเปฯ.}

\prose{อิธานนฺท ภิกฺขุ อเนกวิหิตํ อิทฺธิวิธํ ปจฺจนุโภติ, เอโกปิ หุตฺวา พหุธา โหติ ฯเปฯ ยาว พฺรหฺมโลกาปิ กาเยน วสํ วตฺเตติ, อยํ วุจฺจตานนฺท อิทฺธิ.}

\prose{กตโม จานนฺท อิทฺธิปาโท, โย อานนฺท มคฺโค ยา ปฏิปทา อิทฺธิลาภาย อิทฺธิปฏิลาภาย สํวตฺตติ, อยํ วุจฺจตานนฺท อิทฺธิปาโท.}

\prose{กตมา จานนฺท อิทฺธิปาทภาวนา, อิธานนฺท ภิกฺขุ ฉนฺทสมาธิปฺปธานสงฺขารสมนฺนาคตํ อิทฺธิปาทํ ภาเวติ.\csromanspacer{} วีริยสมาธิ ฯเปฯ.\csromanspacer{} จิตฺตสมาธิ.\csromanspacer{} วีมํสาสมาธิปฺปธานสงฺขารสมนฺนาคตํ อิทฺธิปาทํ ภาเวติ, อยํ วุจฺจตานนฺท อิทฺธิปาทภาวนา.}

\prose{กตมา จานนฺท อิทฺธิปาทภาวนาคามินี ปฏิปทา, อยเมว อริโย อฏฺฐงฺคิโก มคฺโค.\csromanspacer{} เสยฺยถิทํ, สมฺมาทิฏฺฐิ ฯเปฯ สมฺมาสมาธิ, อยํ วุจฺจตานนฺท อิทฺธิปาทภาวนาคามินี ปฏิปทาติ.\csromanspacer{}\csromanspacer{}อฏฺฐมํ.}

\csromanheader{2}{9. ปฐมภิกฺขุสุตฺต}
\proseitem{841}{อถ โข สมฺพหุลา ภิกฺขู เยน ภควา เตนุปสงฺกมิํสุ, อุปสงฺกมิตฺวา ภควนฺตํ อภิวาเทตฺวา เอกมนฺตํ นิสีทิํสุ, เอกมนฺตํ นิสินฺนา}

\csromanheader{2}{7. อิทฺธิปาทสํยุตฺต}
\prose{\csromanfolio{251}โข เต ภิกฺขู ภควนฺตํ เอตทโวจุํ “กตมา นุ โข ภนฺเต อิทฺธิ, กตโม อิทฺธิปาโท, กตมา อิทฺธิปาทภาวนา, กตมา อิทฺธิปาทภาวนาคามินี ปฏิปทา”ติ.}

\prose{อิธ ภิกฺขเว ภิกฺขุ อเนกวิหิตํ อิทฺธิวิธํ ปจฺจนุโภติ, เอโกปิ หุตฺวา พหุธา โหติ ฯเปฯ ยาว พฺรหฺมโลกาปิ กาเยน วสํ วตฺเตติ, อยํ วุจฺจติ ภิกฺขเว อิทฺธิ.}

\prose{กตโม จ ภิกฺขเว อิทฺธิปาโท, โย ภิกฺขเว มคฺโค ยา ปฏิปทา อิทฺธิลาภาย อิทฺธิปฏิลาภาย สํวตฺตติ, อยํ วุจฺจติ ภิกฺขเว อิทฺธิปาโท.}

\prose{กตมา จ ภิกฺขเว อิทฺธิปาทภาวนา, อิธ ภิกฺขเว ภิกฺขุ ฉนฺทสมาธิปฺปธานสงฺขารสมนฺนาคตํ อิทฺธิปาทํ ภาเวติ.\csromanspacer{} วีริยสมาธิ ฯเปฯ.\csromanspacer{} จิตฺตสมาธิ.\csromanspacer{} วีมํสาสมาธิปฺปธานสงฺขารสมนฺนาคตํ อิทฺธิปาทํ ภาเวติ, อยํ วุจฺจติ ภิกฺขเว อิทฺธิปาทภาวนา.}

\prose{กตมา จ ภิกฺขเว อิทฺธิปาทภาวนาคามินี ปฏิปทา, อยเมว อริโย อฏฺฐงฺคิโก มคฺโค.\csromanspacer{} เสยฺยถิทํ, สมฺมาทิฏฺฐิ ฯเปฯ สมฺมาสมาธิ, อยํ วุจฺจติ ภิกฺขเว อิทฺธิปาทภาวนาคามินี ปฏิปทาติ.\csromanspacer{}\csromanspacer{}นวมํ.}

\csromanheader{2}{10. ทุติยภิกฺขุสุตฺต}
\proseitem{842}{อถ โข สมฺพหุลา ภิกฺขู เยน ภควา เตนุปสงฺกมิํสุ ฯเปฯ เอกมนฺตํ นิสินฺเน โข เต ภิกฺขู ภควา เอตทโวจ “กตมา นุ โข ภิกฺขเว อิทฺธิ, กตโม อิทฺธิปาโท, กตมา อิทฺธิปาทภาวนา, กตมา อิทฺธิปาทภาวนาคามินี ปฏิปทา”ติ.\csromanspacer{} ภควํมูลกา โน ภนฺเต ธมฺมา, ภควํเนตฺติกา ฯเปฯ.}

\prose{กตมา จ ภิกฺขเว อิทฺธิ, อิธ ภิกฺขเว ภิกฺขุ อเนกวิหิตํ อิทฺธิวิธํ ปจฺจนุโภติ, เอโกปิ หุตฺวา พหุธา โหติ ฯเปฯ ยาว พฺรหฺมโลกาปิ กาเยน วสํ วตฺเตติ, อยํ วุจฺจติ ภิกฺขเว อิทฺธิ.}

\prose{กตโม จ ภิกฺขเว อิทฺธิปาโท, โย ภิกฺขเว มคฺโค ยา ปฏิปทา อิทฺธิลาภาย อิทฺธิปฏิลาภาย สํวตฺตติ, อยํ วุจฺจติ ภิกฺขเว อิทฺธิปาโท.}

\prose{กตมา จ ภิกฺขเว อิทฺธิปาทภาวนา, อิธ ภิกฺขเว ภิกฺขุ ฉนฺทสมาธิปฺปธานสงฺขารสมนฺนาคตํ อิทฺธิปาทํ ภาเวติ.\csromanspacer{} วีริยสมาธิ ฯเปฯ.\csromanspacer{} จิตฺตสมาธิ. \csromanfolio{252}วีมํสาสมาธิปฺปธานสงฺขารสมนฺนาคตํ อิทฺธิปาทํ ภาเวติ, อยํ วุจฺจติ ภิกฺขเว อิทฺธิปาทภาวนา.}

\prose{กตมา จ ภิกฺขเว อิทฺธิปาทภาวนาคามินี ปฏิปทา, อยเมว อริโย อฏฺฐงฺคิโก มคฺโค.\csromanspacer{} เสยฺยถิทํ, สมฺมาทิฏฺฐิ ฯเปฯ สมฺมาสมาธิ, อยํ วุจฺจติ ภิกฺขเว อิทฺธิปาทภาวนาคามินี ปฏิปทาติ.\csromanspacer{}\csromanspacer{}ทสมํ.}

\csromanmatikaanchor{mtk.184}
\csromantocmarkref{subparagraph}{11. โมคฺคลฺลานสุตฺต}{mtk.184}
\bookmark[dest={mtk.184},level=5]{11. โมคฺคลฺลานสุตฺต}
\csromanheader{6}{11. โมคฺคลฺลานสุตฺต}
\proseitem{843}{ตตฺร โข ภควา ภิกฺขู อามนฺเตสิ “ตํ กิํ มญฺญถ ภิกฺขเว, กตเมสํ ธมฺมานํ ภาวิตตฺตา พหุลีกตตฺตา โมคฺคลฺลาโน ภิกฺขุ เอวํมหิทฺธิโก เอวํมหานุภาโว”ติ.\csromanspacer{} ภควํมูลกา โน ภนฺเต ธมฺมา, ภควํเนตฺติกา ฯเปฯ.\csromanspacer{} จตุนฺนํ โข ภิกฺขเว อิทฺธิปาทานํ ภาวิตตฺตา พหุลีกตตฺตา โมคฺคลฺลาโน ภิกฺขุ เอวํมหิทฺธิโก เอวํมหานุภาโว.}

\prose{กตเมสํ จตุนฺนํ, อิธ ภิกฺขเว โมคฺคลฺลาโน ภิกฺขุ ฉนฺทสมาธิปฺปธานสงฺขารสมนฺนาคตํ อิทฺธิปาทํ ภาเวติ, อิติ เม ฉนฺโท น จ อติลีโน ภวิสฺสติ, น จ อติปฺปคฺคหิโต ภวิสฺสติ, น จ อชฺฌตฺตํ สํขิตฺโต ภวิสฺสติ, น จ พหิทฺธา วิกฺขิตฺโต ภวิสฺสติ.\csromanspacer{} ปจฺฉาปุเรสญฺญี จ วิหรติ, ยถา ปุเร ตถา ปจฺฉา, ยถา ปจฺฉา ตถา ปุเร.\csromanspacer{} ยถา อโธ ตถา อุทฺธํ, ยถา อุทฺธํ ตถา อโธ.\csromanspacer{} ยถา ทิวา ตถา รตฺติํ, ยถา รตฺติํ ตถา ทิวา.\csromanspacer{} อิติ วิวเฏน เจตสา อปริโยนทฺเธน สปฺปภาสํ จิตฺตํ ภาเวติ.\csromanspacer{} วีริยสมาธิ ฯเปฯ.\csromanspacer{} จิตฺตสมาธิ.\csromanspacer{} วีมํสาสมาธิปฺปธานสงฺขารสมนฺนาคตํ อิทฺธิปาทํ ภาเวติ, อิติ เม วีมํสา น จ อติลีนา ภวิสฺสติ, น จ อติปฺปคฺคหิตา ภวิสฺสติ, น จ อชฺฌตฺตํ สํขิตฺตา ภวิสฺสติ, น จ พหิทฺธา วิกฺขิตฺตา ภวิสฺสติ ฯเปฯ.\csromanspacer{} อิติ วิวเฏน เจตสา อปริโยนทฺเธน สปฺปภาสํ จิตฺตํ ภาเวติ.\csromanspacer{} อิเมสํ โข ภิกฺขเว จตุนฺนํ อิทฺธิปาทานํ ภาวิตตฺตา พหุลีกตตฺตา โมคฺคลฺลาโน ภิกฺขุ เอวํมหิทฺธิโก เอวํมหานุภาโว.}

\prose{อิเมสญฺจ ปน ภิกฺขเว จตุนฺนํ อิทฺธิปาทานํ ภาวิตตฺตา พหุลีกตตฺตา โมคฺคลฺลาโน ภิกฺขุ เอวํ อเนกวิหิตํ อิทฺธิวิธํ ปจฺจนุโภติ, เอโกปิ หุตฺวา พหุธา โหติ, พหุธาปิ หุตฺวา เอโก โหติ ฯเปฯ.}

\csromanheader{2}{7. อิทฺธิปาทสํยุตฺต}
\prose{\csromanfolio{253}ยาว พฺรหฺมโลกาปิ กาเยน วสํ วตฺเตติ.\csromanspacer{} อิเมสญฺจ ปน ภิกฺขเว จตุนฺนํ อิทฺธิปาทานํ ภาวิตตฺตา พหุลีกตตฺตา โมคฺคลฺลาโน ภิกฺขุ อาสวานํ ขยา อนาสวํ เจโตวิมุตฺติํ ปญฺญาวิมุตฺติํ ทิฏฺเฐว ธมฺเม สยํ อภิญฺญา สจฺฉิกตฺวา อุปสมฺปชฺช วิหรตีติ.\csromanspacer{}\csromanspacer{}เอกาทสมํ.}

\csromanmatikaanchor{mtk.185}
\csromantocmarkref{subparagraph}{12. ตถาคตสุตฺต}{mtk.185}
\bookmark[dest={mtk.185},level=5]{12. ตถาคตสุตฺต}
\csromanheader{6}{12. ตถาคตสุตฺต}
\proseitem{844}{ตตฺร โข ภควา ภิกฺขู อามนฺเตสิ “ตํ กิํ มญฺญถ ภิกฺขเว, กตเมสํ ธมฺมานํ ภาวิตตฺตา พหุลีกตตฺตา ตถาคโต เอวํมหิทฺธิโก เอวํมหานุภาโว”ติ.\csromanspacer{} ภควํมูลกา โน ภนฺเต ธมฺมา ฯเปฯ.\csromanspacer{} จตุนฺนํ โข ภิกฺขเว อิทฺธิปาทานํ ภาวิตตฺตา พหุลีกตตฺตา ตถาคโต เอวํมหิทฺธิโก เอวํมหานุภาโว.}

\prose{กตเมสํ จตุนฺนํ, อิธ ภิกฺขเว ตถาคโต ฉนฺทสมาธิปฺปธานสงฺขารสมนฺนาคตํ อิทฺธิปาทํ ภาเวติ, อิติ เม ฉนฺโท น จ อติลีโน ภวิสฺสติ, น จ อติปฺปคฺคหิโต ภวิสฺสติ, น จ อชฺฌตฺตํ สํขิตฺโต ภวิสฺสติ, น จ พหิทฺธา วิกฺขิตฺโต ภวิสฺสติ.\csromanspacer{} ปจฺฉาปุเรสญฺญี จ วิหรติ, ยถา ปุเร ตถา ปจฺฉา, ยถา ปจฺฉา ตถา ปุเร.\csromanspacer{} ยถา อโธ ตถา อุทฺธํ, ยถา อุทฺธํ ตถา อโธ.\csromanspacer{} ยถา ทิวา ตถา รตฺติํ, ยถา รตฺติํ ตถา ทิวา.\csromanspacer{} อิติ วิวเฏน เจตสา อปริโยนทฺเธน สปฺปภาสํ จิตฺตํ ภาเวติ.\csromanspacer{} วีริยสมาธิ ฯเปฯ.\csromanspacer{} จิตฺตสมาธิ.\csromanspacer{} วีมํสาสมาธิปฺปธานสงฺขารสมนฺนาคตํ อิทฺธิปาทํ ภาเวติ, อิติ เม วีมํสา น จ อติลีนา ภวิสฺสติ, น จ อติปฺปคฺคหิตา ภวิสฺสติ, น จ อชฺฌตฺตํ สํขิตฺตา ภวิสฺสติ, น จ พหิทฺธา วิกฺขิตฺตา ภวิสฺสติ.\csromanspacer{} ปจฺฉาปุเรสญฺญี จ วิหรติ, ยถา ปุเร ตถา ปจฺฉา, ยถา ปจฺฉา ตถา ปุเร.\csromanspacer{} ยถา อโธ ตถา อุทฺธํ, ยถา อุทฺธํ ตถา อโธ.\csromanspacer{} ยถา ทิวา ตถา รตฺติํ, ยถา รตฺติํ ตถา ทิวา.\csromanspacer{} อิติ วิวเฏน เจตสา อปริโยนทฺเธน สปฺปภาสํ จิตฺตํ ภาเวติ.\csromanspacer{} อิเมสํ โข ภิกฺขเว จตุนฺนํ อิทฺธิปาทานํ ภาวิตตฺตา พหุลีกตตฺตา ตถาคโต เอวํมหิทฺธิโก เอวํมหานุภาโว.}

\prose{อิเมสญฺจ ปน ภิกฺขเว จตุนฺนํ อิทฺธิปาทานํ ภาวิตตฺตา พหุลีกตตฺตา ตถาคโต อเนกวิหิตํ อิทฺธิวิธํ ปจฺจนุโภติ, เอโกปิ หุตฺวา พหุธา โหติ ฯเปฯ ยาว พฺรหฺมโลกาปิ กาเยน วสํ วตฺเตติ.\csromanspacer{} อิเมสญฺจ ปน ภิกฺขเว จตุนฺนํ อิทฺธิปาทานํ ภาวิตตฺตา พหุลีกตตฺตา ตถาคโต \csromanfolio{254}อาสวานํ ขยา อนาสวํ เจโตวิมุตฺติํ ปญฺญาวิมุตฺติํ ทิฏฺเฐว ธมฺเม สยํ อภิญฺญา สจฺฉิกตฺวา อุปสมฺปชฺช วิหรตีติ.\csromanspacer{}\csromanspacer{}ทฺวาทสมํ. (ฉปิ อภิญฺญาโย วิตฺถาเรตพฺพา.) อโยคุฬวคฺโค ตติโย.}

\titlehead{\textbf{ตสฺสุทฺทานํ}}
\csromangathameasure{มคฺโค อโยคุโฬ ภิกฺขุ, สุทฺธิกญฺจาปิ เทฺว ผลา. \\ เทฺว จานนฺทา ทุเว ภิกฺขู, โมคฺคลฺลาโน ตถาคโตติ.}
\csromangathagroup{\csromangathastanza{มคฺโค อโยคุโฬ ภิกฺขุ, สุทฺธิกญฺจาปิ เทฺว ผลา. \\ เทฺว จานนฺทา ทุเว ภิกฺขู, โมคฺคลฺลาโน ตถาคโตติ.}}

\chapterhead{4. คงฺคาเปยฺยาลวคฺค}
\csromanmatikaanchor{mtk.186}
\csromanmatikaanchor{mtk.187}
\csromantocmarknopage{subparagraph}{4. คงฺคาเปยฺยาลวคฺค}{mtk.186}
\bookmark[dest={mtk.186},level=5]{4. คงฺคาเปยฺยาลวคฺค}
\csromantocmarkref{subparagraph}{1-12. คงฺคานที-อาทิสุตฺต}{mtk.187}
\bookmark[dest={mtk.187},level=5]{1-12. คงฺคานที-อาทิสุตฺต}
\csromanheader{6}{1-12. คงฺคานที-อาทิสุตฺต}
\proseitem{845-856}{เสยฺยถาปิ ภิกฺขเว คงฺคา นที ปาจีนนินฺนา ปาจีนโปณา ปาจีนปพฺภารา, เอวเมว โข ภิกฺขเว ภิกฺขุ จตฺตาโร อิทฺธิปาเท ภาเวนฺโต จตฺตาโร อิทฺธิปาเท พหุลีกโรนฺโต นิพฺพานนินฺโน โหติ นิพฺพานโปโณ นิพฺพานปพฺภาโร.\csromanspacer{} กถญฺจ ภิกฺขเว ภิกฺขุ จตฺตาโร อิทฺธิปาเท ภาเวนฺโต จตฺตาโร อิทฺธิปาเท พหุลีกโรนฺโต นิพฺพานนินฺโน โหติ นิพฺพานโปโณ นิพฺพานปพฺภาโร, อิธ ภิกฺขเว ภิกฺขุ ฉนฺทสมาธิปฺปธานสงฺขารสมนฺนาคตํ อิทฺธิปาทํ ภาเวติ.\csromanspacer{} วีริยสมาธิ ฯเปฯ.\csromanspacer{} จิตฺตสมาธิ.\csromanspacer{} วีมํสาสมาธิปฺปธานสงฺขารสมนฺนาคตํ อิทฺธิปาทํ ภาเวติ.}

\prose{เอวํ โข ภิกฺขเว ภิกฺขุ จตฺตาโร อิทฺธิปาเท ภาเวนฺโต จตฺตาโร อิทฺธิปาเท พหุลีกโรนฺโต นิพฺพานนินฺโน โหติ นิพฺพานโปโณ นิพฺพานปพฺภาโรติ.\csromanspacer{}\csromanspacer{}ทฺวาทสมํ.\csromanspacer{} คงฺคาเปยฺยาลวคฺโค จตุตฺโถ.}

\titlehead{\textbf{ตสฺสุทฺทานํ}}
\csromangathameasure{ฉ ปาจีนโต นินฺนา, ฉ นินฺนา จ สมุทฺทโต. \\ เทฺวเต ฉ ทฺวาทส โหนฺติ, วคฺโค เตน ปวุจฺจตีติ.}
\csromangathagroup{\csromangathastanza{ฉ ปาจีนโต นินฺนา, ฉ นินฺนา จ สมุทฺทโต. \\ เทฺวเต ฉ ทฺวาทส โหนฺติ, วคฺโค เตน ปวุจฺจตีติ.}}

\csromanheader{2}{7. อิทฺธิปาทสํยุตฺต}
\vspace{\tipitakaparskip}
\csromancenter{อปฺปมาทวคฺโค วิตฺถาเรตพฺโพ.}
\titlehead{\textbf{ตสฺสุทฺทานํ}}
\csromangathameasure{ตถาคตํ ปทํ กูฏํ, มูลํ สาโร จ วสฺสิกํ. \\ ราชา จนฺทิมสูริยา, วตฺเถน ทสมํ ปทนฺติ. พลกรณียวคฺโค วิตฺถาเรตพฺโพ.}
\csromangathagroup{\csromangathastanza{\csromanfolio{255}ตถาคตํ ปทํ กูฏํ, มูลํ สาโร จ วสฺสิกํ. \\ ราชา จนฺทิมสูริยา, วตฺเถน ทสมํ ปทนฺติ. พลกรณียวคฺโค วิตฺถาเรตพฺโพ.}}

\titlehead{\textbf{ตสฺสุทฺทานํ}}
\csromangathameasure{พลํ พีชญฺจ นาโค จ, รุกฺโข กุมฺเภน สูกิยา. \\ อากาเสน จ เทฺว เมฆา, นาวา อาคนฺตุกา นทีติ. เอสนาวคฺโค วิตฺถาเรตพฺโพ.}
\csromangathagroup{\csromangathastanza{พลํ พีชญฺจ นาโค จ, รุกฺโข กุมฺเภน สูกิยา. \\ อากาเสน จ เทฺว เมฆา, นาวา อาคนฺตุกา นทีติ. เอสนาวคฺโค วิตฺถาเรตพฺโพ.}}

\titlehead{\textbf{ตสฺสุทฺทานํ}}
\csromangathameasure{เอสนา วิธา อาสโว, ภโว จ ทุกฺขตา ติสฺโส. \\ ขิลํ มลญฺจ นีโฆ จ, เวทนา ตณฺหา ตสินา จาติ.}
\csromangathagroup{\csromangathastanza{เอสนา วิธา อาสโว, ภโว จ ทุกฺขตา ติสฺโส. \\ ขิลํ มลญฺจ นีโฆ จ, เวทนา ตณฺหา ตสินา จาติ.}}

\chapterhead{8. โอฆวคฺค}
\csromanmatikaanchor{mtk.188}
\csromanmatikaanchor{mtk.189}
\csromantocmarknopage{subparagraph}{8. โอฆวคฺค}{mtk.188}
\bookmark[dest={mtk.188},level=5]{8. โอฆวคฺค}
\csromantocmarkref{subparagraph}{1-10. โอฆาทิสุตฺต}{mtk.189}
\bookmark[dest={mtk.189},level=5]{1-10. โอฆาทิสุตฺต}
\csromanheader{6}{1-10. โอฆาทิสุตฺต}
\proseitem{889-898}{ปญฺจิมานิ ภิกฺขเว อุทฺธมฺภาคิยานิ สํโยชนานิ.\csromanspacer{} กตมานิ ปญฺจ, รูปราโค อรูปราโค มาโน อุทฺธจฺจํ อวิชฺชา, อิมานิ โข ภิกฺขเว ปญฺจุทฺธมฺภาคิยานิ สํโยชนานิ.\csromanspacer{} อิเมสํ โข ภิกฺขเว ปญฺจนฺนํ อุทฺธมฺภาคิยานํ สํโยชนานํ อภิญฺญาย ปริญฺญาย ปริกฺขยาย ปหานาย จตฺตาโร อิทฺธิปาทา ภาเวตพฺพา.\csromanspacer{} กตเม จตฺตาโร, อิธ ภิกฺขเว ภิกฺขุ ฉนฺทสมาธิปฺปธานสงฺขารสมนฺนาคตํ อิทฺธิปาทํ ภาเวติ.\csromanspacer{} วีริยสมาธิ ฯเปฯ. \csromanfolio{256}จิตฺตสมาธิ.\csromanspacer{} วีมํสาสมาธิปฺปธานสงฺขารสมนฺนาคตํ อิทฺธิปาทํ ภาเวติ.\csromanspacer{} อิเมสํ โข ภิกฺขเว ปญฺจนฺนํ อุทฺธมฺภาคิยานํ สํโยชนานํ อภิญฺญาย ปริญฺญาย ปริกฺขยาย ปหานาย อิเม จตฺตาโร อิทฺธิปาทา ภาเวตพฺพาติ. (ยถา มคฺคสํยุตฺตํ, ตถา วิตฺถาเรตพฺพํ.) โอฆวคฺโค อฏฺฐโม.}

\titlehead{\textbf{ตสฺสุทฺทานํ}}
\csromangathameasure{โอโฆ โยโค อุปาทานํ, คนฺถา อนุสเยน จ. \\ กามคุณา นีวรณา, ขนฺธา โอรุทฺธมฺภาคิยาติ. \\ \textbf{อิทฺธิปาทสํยุตฺตํ สตฺตมํ.}}
\csromangathagroup{\csromangathastanza{โอโฆ โยโค อุปาทานํ, คนฺถา อนุสเยน จ. \\ กามคุณา นีวรณา, ขนฺธา โอรุทฺธมฺภาคิยาติ. \\ \textbf{อิทฺธิปาทสํยุตฺตํ สตฺตมํ.}}}

\csromanmatikaanchor{mtk.190}
\csromantocmarknopage{subparagraph}{8. อนุรุทฺธสํยุตฺต}{mtk.190}
\bookmark[dest={mtk.190},level=5]{8. อนุรุทฺธสํยุตฺต}
\csromanheader{6}{8. อนุรุทฺธสํยุตฺต}
\csromanmatikaanchor{mtk.191}
\csromantocmarknopage{subparagraph}{1. รโหคตวคฺค}{mtk.191}
\bookmark[dest={mtk.191},level=5]{1. รโหคตวคฺค}
\csromanheader{6}{1. รโหคตวคฺค}
\csromanmatikaanchor{mtk.192}
\csromantocmarkref{subparagraph}{1. ปฐมรโหคตสุตฺต}{mtk.192}
\bookmark[dest={mtk.192},level=5]{1. ปฐมรโหคตสุตฺต}
\csromanheader{6}{1. ปฐมรโหคตสุตฺต}
\proseitem{899}{\csromanfolio{257}เอวํ เม สุตํ\csromandash{}เอกํ สมยํ อายสฺมา อนุรุทฺโธ สาวตฺถิยํ วิหรติ เชตวเน อนาถปิณฺฑิกสฺส อาราเม.\csromanspacer{} อถ โข อายสฺมโต อนุรุทฺธสฺส รโหคตสฺส ปฏิสลฺลีนสฺส เอวํ เจตโส ปริวิตกฺโก อุทปาทิ “เยสํ เกสญฺจิ จตฺถาโร สติปฏฺฐานา วิรทฺธา, วิรทฺโธ เตสํ อริโย มคฺโค สมฺมา ทุกฺขกฺขยคามี.\csromanspacer{} เยสํ เกสญฺจิ จตฺตาโร สติปฏฺฐานา อารทฺธา, อารทฺโธ เตสํ อริโย มคฺโค สมฺมา ทุกฺขกฺขยคามี”ติ.}

\prose{อถ โข อายสฺมา มหาโมคฺคลฺลาโน อายสฺมโต อนุรุทฺธสฺส เจตสา เจโตปริวิตกฺกมญฺญาย เสยฺยถาปิ นาม พลวา ปุริโส สมิญฺชิตํ วา พาหํ ปสาเรยฺย, ปสาริตํ วา พาหํ สมิญฺเชยฺย, เอวเมว อายสฺมโต อนุรุทฺธสฺส สมฺมุเข ปาตุรโหสิ.\csromanspacer{} อถ โข อายสฺมา มหาโมคฺคลฺลาโน อายสฺมนฺตํ อนุรุทฺธํ เอตทโวจ “กิตฺตาวตา นุ โข อาวุโส อนุรุทฺธ ภิกฺขุโน จตฺตาโร สติปฏฺฐานา อารทฺธา โหนฺตี”ติ.}

\prose{อิธาวุโส ภิกฺขุ อชฺฌตฺตํ กาเย สมุทยธมฺมานุปสฺสี วิหรติ.\csromanspacer{} อชฺฌตฺตํ กาเย วยธมฺมานุปสฺสี วิหรติ.\csromanspacer{} อชฺฌตฺตํ กาเย สมุทยวยธมฺมานุปสฺสี วิหรติ อาตาปี สมฺปชาโน สติมา วิเนยฺย โลเก อภิชฺฌาโทมนสฺสํ.\csromanspacer{} พหิทฺธา กาเย สมุทยธมฺมานุปสฺสี วิหรติ.\csromanspacer{} พหิทฺธา กาเย วยธมฺมานุปสฺสี วิหรติ.\csromanspacer{} พหิทฺธา กาเย สมุทยวยธมฺมานุปสฺสี วิหรติ อาตาปี สมฺปชาโน สติมา วิเนยฺย โลเก อภิชฺฌาโทมนสฺสํ.\csromanspacer{} อชฺฌตฺตพหิทฺธา กาเย สมุทยธมฺมานุปสฺสี วิหรติ.\csromanspacer{} อชฺฌตฺตพหิทฺธา กาเย วยธมฺมานุปสฺสี วิหรติ.\csromanspacer{} อชฺฌตฺตพหิทฺธา กาเย สมุทยวยธมฺมานุปสฺสี วิหรติ อาตาปี สมฺปชาโน สติมา วิเนยฺย โลเก อภิชฺฌาโทมนสฺสํ.}

\prose{\csromanfolio{258}โส สเจ อากงฺขติ “อปฺปฏิกูเล ปฏิกูลสญฺญี วิหเรยฺยนฺ”ติ, ปฏิกูลสญฺญี ตตฺถ วิหรติ.\csromanspacer{} สเจ อากงฺขติ “ปฏิกูเล อปฺปฏิกูลสญฺญี วิหเรยฺยนฺ”ติ, อปฺปฏิกูลสญฺญี ตตฺถ วิหรติ.\csromanspacer{} สเจ อากงฺขติ “อปฺปฏิกูเล จ ปฏิกูเล จ ปฏิกูลสญฺญี วิหเรยฺยนฺ”ติ, ปฏิกูลสญฺญี ตตฺถ วิหรติ.\csromanspacer{} สเจ อากงฺขติ “ปฏิกูเล จ อปฺปฏิกูเล จ อปฺปฏิกูลสญฺญี วิหเรยฺยนฺ”ติ, อปฺปฏิกูลสญฺญี ตตฺถ วิหรติ.\csromanspacer{} สเจ อากงฺขติ “อปฺปฏิกูลญฺจ ปฏิกูลญฺจ ตทุภยํ อภินิวชฺเชตฺวา อุเปกฺขโก วิหเรยฺยํ สโต สมฺปชาโน”ติ, อุเปกฺขโก ตตฺถ วิหรติ สโต สมฺปชาโน.}

\prose{อชฺฌตฺตํ เวทนาสุ สมุทยธมฺมานุปสฺสี วิหรติ.\csromanspacer{} อชฺฌตฺตํ เวทนาสุ วยธมฺมานุปสฺสี วิหรติ.\csromanspacer{} อชฺฌตฺตํ เวทนาสุ สมุทยวยธมฺมานุปสฺสี วิหรติ อาตาปี สมฺปชาโน สติมา วิเนยฺย โลเก อภิชฺฌาโทมนสฺสํ.\csromanspacer{} พหิทฺธา เวทนาสุ สมุทยธมฺมานุปสฺสี วิหรติ.\csromanspacer{} พหิทฺธา เวทนาสุ วยธมฺมานุปสฺสี วิหรติ.\csromanspacer{} พหิทฺธา เวทนาสุ สมุทยวยธมฺมานุปสฺสี วิหรติ อาตาปี สมฺปชาโน สติมา วิเนยฺย โลเก อภิชฺฌาโทมนสฺสํ.\csromanspacer{} อชฺฌตฺตพหิทฺธา เวทนาสุ สมุทยธมฺมานุปสฺสี วิหรติ.\csromanspacer{} อชฺฌตฺตพหิทฺธา เวทนาสุ วยธมฺมานุปสฺสี วิหรติ.\csromanspacer{} อชฺฌตฺตพหิทฺธา เวทนาสุ สมุทยวยธมฺมานุปสฺสี วิหรติ อาตาปี สมฺปชาโน สติมา วิเนยฺย โลเก อภิชฺฌาโทมนสฺสํ.}

\prose{โส สเจ อากงฺขติ “อปฺปฏิกูเล ปฏิกูลสญฺญี วิหเรยฺยนฺ”ติ, ปฏิกูลสญฺญี ตตฺถ วิหรติ.\csromanspacer{} สเจ อากงฺขติ “ปฏิกูเล อปฺปฏิกูลสญฺญี วิหเรยฺยนฺ”ติ, อปฺปฏิกูลสญฺญี ตตฺถ วิหรติ.\csromanspacer{} สเจ อากงฺขติ “อปฺปฏิกูเล จ ปฏิกูเล จ ปฏิกูลสญฺญี วิหเรยฺยนฺ”ติ, ปฏิกูลสญฺญี ตตฺถ วิหรติ.\csromanspacer{} สเจ อากงฺขติ “ปฏิกูเล จ อปฺปฏิกูเล จ อปฺปฏิกูลสญฺญี วิหเรยฺยนฺ”ติ, อปฺปฏิกูลสญฺญี ตตฺถ วิหรติ.\csromanspacer{} สเจ อากงฺขติ “อปฺปฏิกูลญฺจ ปฏิกูลญฺจ ตทุภยํ อภินิวชฺเชตฺวา อุเปกฺขโก วิหเรยฺยํ สโต สมฺปชาโน”ติ, อุเปกฺขโก ตตฺถ วิหรติ สโต สมฺปชาโน.}

\prose{อชฺฌตฺตํ จิตฺเต ฯเปฯ.\csromanspacer{} พหิทฺธา จิตฺเต ฯเปฯ.\csromanspacer{} อชฺฌตฺตพหิทฺธา จิตฺเต สมุทยธมฺมานุปสฺสี วิหรติ.\csromanspacer{} อชฺฌตฺตพหิทฺธา จิตฺเต วยธมฺมานุปสฺสี วิหรติ.\csromanspacer{} อชฺฌตฺตพหิทฺธา จิตฺเต สมุทยวยธมฺมานุปสฺสี วิหรติ อาตาปี ฯเปฯ อภิชฺฌาโทมนสฺสํ.\csromanspacer{} โส สเจ อากงฺขติ “อปฺปฏิกูเล ปฏิกูลสญฺญี}

\csromanheader{2}{8. อนุรุทฺธสํยุตฺต}
\prose{\csromanfolio{259}วิหเรยฺยนฺ”ติ, ปฏิกูลสญฺญี ตตฺถ วิหรติ ฯเปฯ อุเปกฺขโก ตตฺถ วิหรติ สโต สมฺปชาโน.}

\prose{อชฺฌตฺตํ ธมฺเมสุ ฯเปฯ.\csromanspacer{} พหิทฺธา ธมฺเมสุ ฯเปฯ.\csromanspacer{} อชฺฌตฺตพหิทฺธา ธมฺเมสุ สมุทยธมฺมานุปสฺสี วิหรติ.\csromanspacer{} อชฺฌตฺตพหิทฺธา ธมฺเมสุ วยธมฺมานุปสฺสี วิหรติ.\csromanspacer{} อชฺฌตฺตพหิทฺธา ธมฺเมสุ สมุทยวยธมฺมานุปสฺสี วิหรติ อาตาปี สมฺปชาโน สติมา วิเนยฺย โลเก อภิชฺฌาโทมนสฺสํ.\csromanspacer{} โส สเจ อากงฺขติ “อปฺปฏิกูเล ปฏิกูลสญฺญี วิหเรยฺยนฺ”ติ, ปฏิกูลสญฺญี ตตฺถ วิหรติ ฯเปฯ อุเปกฺขโก ตตฺถ วิหรติ สโต สมฺปชาโน.\csromanspacer{} เอตฺตาวตา โข อาวุโส ภิกฺขุโน จตฺตาโร สติปฏฺฐานา อารทฺธา โหนฺตีติ.\csromanspacer{}\csromanspacer{}ปฐมํ.}

\csromanmatikaanchor{mtk.193}
\csromantocmarkref{subparagraph}{2. ทุติยรโหคตสุตฺต}{mtk.193}
\bookmark[dest={mtk.193},level=5]{2. ทุติยรโหคตสุตฺต}
\csromanheader{6}{2. ทุติยรโหคตสุตฺต}
\proseitem{900}{สาวตฺถินิทานํ.\csromanspacer{} อถ โข อายสฺมโต อนุรุทฺธสฺส รโหคตสฺส ปฏิสลฺลีนสฺส เอวํ เจตโส ปริวิตกฺโก อุทปาทิ “เยสํ เกสญฺจิ จตฺตาโร สติปฏฺฐานา วิรทฺธา, วิรทฺโธ เตสํ อริโย มคฺโค สมฺมา ทุกฺขกฺขยคามี.\csromanspacer{} เยสํ เกสญฺจิ จตฺตาโร สติปฏฺฐานา อารทฺธา, อารทฺโธ เตสํ อริโย มคฺโค สมฺมา ทุกฺขกฺขยคามี”ติ.}

\prose{อถ โข อายสฺมา มหาโมคฺคลฺลาโน อายสฺมโต อนุรุทฺธสฺส เจตสา เจโตปริวิตกฺกมญฺญาย เสยฺยถาปิ นาม พลวา ปุริโส สมิญฺชิตํ วา พาหํ ปสาเรยฺย, ปสาริตํ วา พาหํ สมิญฺเชยฺย, เอวเมว อายสฺมโต อนุรุทฺธสฺส สมฺมุเข ปาตุรโหสิ.}

\prose{อถ โข อายสฺมา มหาโมคฺคลฺลาโน อายสฺมนฺตํ อนุรุทฺธํ เอตทโวจ “กิตฺตาวตา นุ โข อาวุโส อนุรุทฺธ ภิกฺขุโน จตฺตาโร สติปฏฺฐานา อารทฺธา โหนฺตี”ติ.\csromanspacer{} อิธาวุโส ภิกฺขุ อชฺฌตฺตํ กาเย กายานุปสฺสี วิหรติ อาตาปี สมฺปชาโน สติมา วิเนยฺย โลเก อภิชฺฌาโทมนสฺสํ.\csromanspacer{} พหิทฺธา กาเย กายานุปสฺสี วิหรติ ฯเปฯ.\csromanspacer{} อชฺฌตฺตพหิทฺธา กาเย กายานุปสฺสี วิหรติ อาตาปี สมฺปชาโน สติมา วิเนยฺย โลเก อภิชฺฌาโทมนสฺสํ.}

\csromanmatikaanchor{mtk.194}
\csromanmatikaanchor{mtk.195}
\csromantocmarkref{subparagraph}{3. สุตนุสุตฺต}{mtk.194}
\bookmark[dest={mtk.194},level=5]{3. สุตนุสุตฺต}
\csromantocmarkref{subparagraph}{4-8. ปฐมกณฺฑกีสุตฺตาทีนิ}{mtk.195}
\bookmark[dest={mtk.195},level=5]{4-8. ปฐมกณฺฑกีสุตฺตาทีนิ}
\prose{อชฺฌตฺตํ เวทนาสุ เวทนานุปสฺสี วิหรติ อาตาปี สมฺปชาโน สติมา วิเนยฺย โลเก อภิชฺฌาโทมนสฺสํ.\csromanspacer{} พหิทฺธา เวทนาสุ \csromanfolio{260}เวทนานุปสฺสี วิหรติ อาตาปี สมฺปชาโน สติมา วิเนยฺย โลเก อภิชฺฌาโทมนสฺสํ.\csromanspacer{} อชฺฌตฺตพหิทฺธา เวทนาสุ เวทนานุปสฺสี วิหรติ อาตาปี สมฺปชาโน สติมา วิเนยฺย โลเก อภิชฺฌาโทมนสฺสํ.}

\prose{อชฺฌตฺตํ จิตฺเต ฯเปฯ.\csromanspacer{} พหิทฺธา จิตฺเต ฯเปฯ.\csromanspacer{} อชฺฌตฺตพหิทฺธา จิตฺเต จิตฺตานุปสฺสี วิหรติ อาตาปี สมฺปชาโน สติมา วิเนยฺย โลเก อภิชฺฌาโทมนสฺสํ.}

\prose{อชฺฌตฺตํ ธมฺเมสุ ฯเปฯ.\csromanspacer{} พหิทฺธา ธมฺเมสุ ฯเปฯ.\csromanspacer{} อชฺฌตฺตพหิทฺธา ธมฺเมสุ ธมฺมานุปสฺสี วิหรติ อาตาปี สมฺปชาโน สติมา วิเนยฺย โลเก อภิชฺฌาโทมนสฺสํ.\csromanspacer{} เอตฺตาวตา โข อาวุโส ภิกฺขุโน จตฺตาโร สติปฏฺฐานา อารทฺธา โหนฺตีติ.\csromanspacer{}\csromanspacer{}ทุติยํ.}

\csromanheader{6}{3. สุตนุสุตฺต}
\proseitem{901}{เอกํ สมยํ อายสฺมา อนุรุทฺโธ สาวตฺถิยํ วิหรติ สุตนุตีเร.\csromanspacer{} อถ โข สมฺพหุลา ภิกฺขู เยนายสฺมา อนุรุทฺโธ เตนุปสงฺกมิํสุ, อุปสงฺกมิตฺวา อายสฺมตา อนุรุทฺเธน สทฺธิํ สมฺโมทิํสุ, สมฺโมทนียํ กถํ สารณียํ วีติสาเรตฺวา เอกมนฺตํ นิสีทิํสุ, เอกมนฺตํ นิสินฺนา โข เต ภิกฺขู อายสฺมนฺตํ อนุรุทฺธํ เอตทโวจุํ “กตเมสํ อายสฺมา อนุรุทฺโธ ธมฺมานํ ภาวิตตฺตา พหุลีกตตฺตา มหาภิญฺญตํ ปตฺโต”ติ.}

\prose{จตุนฺนํ ขฺวาหํ อาวุโส สติปฏฺฐานานํ ภาวิตตฺตา พหุลีกตตฺตา มหาภิญฺญตํ ปตฺโต.\csromanspacer{} กตเมสํ จตุนฺนํ, อิธาหํ อาวุโส กาเย กายานุปสฺสี วิหรามิ อาตาปี สมฺปชาโน สติมา วิเนยฺย โลเก อภิชฺฌาโทมนสฺสํ.\csromanspacer{} เวทนาสุ ฯเปฯ.\csromanspacer{} จิตฺเต ฯเปฯ.\csromanspacer{} ธมฺเมสุ ธมฺมานุปสฺสี วิหรามิ อาตาปี สมฺปชาโน สติมา วิเนยฺย โลเก อภิชฺฌาโทมนสฺสํ.\csromanspacer{} อิเมสํ ขฺวาหํ อาวุโส จตุนฺนํ สติปฏฺฐานานํ ภาวิตตฺตา พหุลีกตตฺตา มหาภิญฺญตํ ปตฺโต.\csromanspacer{} อิเมสญฺจ ปนาหํ อาวุโส จตุนฺนํ สติปฏฺฐานานํ ภาวิตตฺตา พหุลีกตตฺตา หีนํ ธมฺมํ หีนโต อพฺภญฺญาสิํ, มชฺฌิมํ ธมฺมํ มชฺฌิมโต อพฺภญฺญาสิํ, ปณีตํ ธมฺมํ ปณีตโต อพฺภญฺญาสินฺติ.\csromanspacer{}\csromanspacer{}ตติยํ.}

\csromanheader{2}{4. ปฐมกณฺฑกีสุตฺต}
\proseitem{902}{เอกํ สมยํ อายสฺมา จ อนุรุทฺโธ อายสฺมา จ สาริปุตฺโต อายสฺมา จ มหาโมคฺคลฺลาโน สาเกเต วิหรนฺติ กณฺฑกีวเน.}

\csromanheader{2}{8. อนุรุทฺธสํยุตฺต}
\prose{\csromanfolio{261}อถ โข อายสฺมา จ สาริปุตฺโต อายสฺมา จ มหาโมคฺคลฺลาโน สายนฺหสมยํ ปฏิสลฺลานา วุฏฺฐิตา เยนายสฺมา อนุรุทฺโธ เตนุปสงฺกมิํสุ, อุปสงฺกมิตฺวา อายสฺมตา อนุรุทฺเธน สทฺธิํ สมฺโมทิํสุ, สมฺโมทนียํ กถํ สารณียํ วีติสาเรตฺวา เอกมนฺตํ นิสีทิํสุ, เอกมนฺตํ นิสินฺโน โข อายสฺมา สาริปุตฺโต อายสฺมนฺตํ อนุรุทฺธํ เอตทโวจ “เสเขนาวุโส อนุรุทฺธ ภิกฺขุนา กตเม ธมฺมา อุปสมฺปชฺช วิหาตพฺพา”ติ.}

\prose{เสเขนาวุโส สาริปุตฺต ภิกฺขุนา จตฺตาโร สติปฏฺฐานา อุปสมฺปชฺช วิหาตพฺพา.\csromanspacer{} กตเม จตฺตาโร, อิธาวุโส ภิกฺขุ กาเย กายานุปสฺสี วิหรติ อาตาปี สมฺปชาโน สติมา วิเนยฺย โลเก อภิชฺฌาโทมนสฺสํ.\csromanspacer{} เวทนาสุ ฯเปฯ.\csromanspacer{} จิตฺเต ฯเปฯ.\csromanspacer{} ธมฺเมสุ ธมฺมานุปสฺสี วิหรติ อาตาปี สมฺปชาโน สติมา วิเนยฺย โลเก อภิชฺฌาโทมนสฺสํ.\csromanspacer{} เสเขนาวุโส สาริปุตฺต ภิกฺขุนา อิเม จตฺตาโร สติปฏฺฐานา อุปสมฺปชฺช วิหาตพฺพาติ.\csromanspacer{}\csromanspacer{}จตุตฺถํ.}

\csromanheader{2}{5. ทุติยกณฺฑกีสุตฺต}
\proseitem{903}{สาเกตนิทานํ.\csromanspacer{} เอกมนฺตํ นิสินฺโน โข อายสฺมา สาริปุตฺโต อายสฺมนฺตํ อนุรุทฺธํ เอตทโวจ “อเสเขนาวุโส อนุรุทฺธ ภิกฺขุนา กตเม ธมฺมา อุปสมฺปชฺช วิหาตพฺพา”ติ.\csromanspacer{} อเสเขนาวุโส สาริปุตฺต ภิกฺขุนา จตฺตาโร สติปฏฺฐานา อุปสมฺปชฺช วิหาตพฺพา.\csromanspacer{} กตเม จตฺตาโร, อิธาวุโส ภิกฺขุ กาเย กายานุปสฺสี วิหรติ อาตาปี สมฺปชาโน สติมา วิเนยฺย โลเก อภิชฺฌาโทมนสฺสํ.\csromanspacer{} เวทนาสุ ฯเปฯ.\csromanspacer{} จิตฺเต ฯเปฯ.\csromanspacer{} ธมฺเมสุ ธมฺมานุปสฺสี วิหรติ อาตาปี สมฺปชาโน สติมา วิเนยฺย โลเก อภิชฺฌาโทมนสฺสํ.\csromanspacer{} อเสเขนาวุโส สาริปุตฺต ภิกฺขุนา อิเม จตฺตาโร สติปฏฺฐานา อุปสมฺปชฺช วิหาตพฺพาติ.\csromanspacer{}\csromanspacer{}ปญฺจมํ.}

\csromanheader{2}{6. ตติยกณฺฑกีสุตฺต}
\proseitem{904}{สาเกตนิทานํ.\csromanspacer{} เอกมนฺตํ นิสินฺโน โข อายสฺมา สาริปุตฺโต อายสฺมนฺตํ อนุรุทฺธํ เอตทโวจ “กตเมสํ อายสฺมา อนุรุทฺโธ ธมฺมานํ ภาวิตตฺตา พหุลีกตตฺตา มหาภิญฺญตํ ปตฺโต”ติ.\csromanspacer{} จตุนฺนํ ขฺวาหํ อาวุโส สติปฏฺฐานานํ ภาวิตตฺตา พหุลีกตตฺตา มหาภิญฺญตํ \csromanfolio{262}ปตฺโต.\csromanspacer{} กตเมสํ จตุนฺนํ, อิธาหํ อาวุโส กาเย กายานุปสฺสี วิหรามิ อาตาปี สมฺปชาโน สติมา วิเนยฺย โลเก อภิชฺฌาโทมนสฺสํ.\csromanspacer{} เวทนาสุ ฯเปฯ.\csromanspacer{} จิตฺเต ฯเปฯ.\csromanspacer{} ธมฺเมสุ ธมฺมานุปสฺสี วิหรามิ อาตาปี สมฺปชาโน สติมา วิเนยฺย โลเก อภิชฺฌาโทมนสฺสํ.\csromanspacer{} อิเมสํ ขฺวาหํ อาวุโส จตุนฺนํ สติปฏฺฐานานํ ภาวิตตฺตา พหุลีกตตฺตา มหาภิญฺญตํ ปตฺโต.\csromanspacer{} อิเมสญฺจ ปนาหํ อาวุโส จตุนฺนํ สติปฏฺฐานานํ ภาวิตตฺตา พหุลีกตตฺตา สหสฺสํ โลกํ อภิชานามีติ.\csromanspacer{}\csromanspacer{}ฉฏฺฐํ.}

\csromanheader{2}{7. ตณฺหากฺขยสุตฺต}
\proseitem{905}{สาวตฺถินิทานํ.\csromanspacer{} ตตฺร โข อายสฺมา อนุรุทฺโธ ภิกฺขู อามนฺเตสิ “อาวุโส ภิกฺขโว”ติ. “อาวุโส”ติ โข เต ภิกฺขู อายสฺมโต อนุรุทฺธสฺส ปจฺจสฺโสสุํ.\csromanspacer{} อายสฺมา อนุรุทฺโธ เอตทโวจ\csromandash{}}

\prose{จตฺตาโรเม อาวุโส สติปฏฺฐานา ภาวิตา พหุลีกตา ตณฺหากฺขยาย สํวตฺตนฺติ.\csromanspacer{} กตเม จตฺตาโร, อิธาวุโส ภิกฺขุ กาเย กายานุปสฺสี วิหรติ ฯเปฯ.\csromanspacer{} เวทนาสุ ฯเปฯ.\csromanspacer{} จิตฺเต ฯเปฯ.\csromanspacer{} ธมฺเมสุ ธมฺมานุปสฺสี วิหรติ อาตาปี สมฺปชาโน สติมา วิเนยฺย โลเก อภิชฺฌาโทมนสฺสํ.\csromanspacer{} อิเม โข อาวุโส จตฺตาโร สติปฏฺฐานา ภาวิตา พหุลีกตา ตณฺหากฺขยาย สํวตฺตนฺตีติ.\csromanspacer{}\csromanspacer{}สตฺตมํ.}

\csromanheader{2}{8. สลฬาคารสุตฺต}
\proseitem{906}{เอกํ สมยํ อายสฺมา อนุรุทฺโธ สาวตฺถิยํ วิหรติ สลฬาคาเร.\csromanspacer{} ตตฺร โข อายสฺมา อนุรุทฺโธ ภิกฺขู อามนฺเตสิ ฯเปฯ เอตทโวจ\csromandash{}เสยฺยถาปิ อาวุโส คงฺคา นที ปาจีนนินฺนา ปาจีนโปณา ปาจีนปพฺภารา, อถ มหาชนกาโย อาคจฺเฉยฺย กุทาลปิฏกํ\footnote{กุทฺทาลปิฏกํ (พหูสุ)} อาทาย “มยํ อิมํ คงฺคานทิํ ปจฺฉานินฺนํ กริสฺสาม ปจฺฉาโปณํ ปจฺฉาปพฺภารนฺ”ติ.\csromanspacer{} ตํ กิํ มญฺญถาวุโส, อปิ นุ โส มหาชนกาโย คงฺคานทิํ ปจฺฉานินฺนํ กเรยฺย ปจฺฉาโปณํ ปจฺฉาปพฺภารนฺติ.\csromanspacer{} โน เหตํ อาวุโส.\csromanspacer{} ตํ กิสฺส เหตุ, คงฺคา อาวุโส นที ปาจีนนินฺนา ปาจีนโปณา ปาจีนปพฺภารา.\csromanspacer{} สา น สุกรา ปจฺฉานินฺนํ กาตุํ ปจฺฉาโปณํ ปจฺฉาปพฺภารํ, ยาวเทว จ ปน โส มหาชนกาโย กิลมถสฺส วิฆาตสฺส ภาคี อสฺสาติ.}

\csromanheader{2}{8. อนุรุทฺธสํยุตฺต}
\prose{\csromanfolio{263}เอวเมว โข อาวุโส ภิกฺขุํ จตฺตาโร สติปฏฺฐาเน ภาเวนฺตํ จตฺตาโร สติปฏฺฐาเน พหุลีกโรนฺตํ ราชาโน วา ราชมหามตฺตา วา มิตฺตา วา อมจฺจา วา ญาตี วา สาโลหิตา วา โภเคหิ อภิหฏฺฐุํ ปวาเรยฺยุํ “เอหมฺโภ ปุริส, กิํ เต อิเม กาสาวา อนุทหนฺติ, กิํ มุณฺโฑ กปาลมนุสญฺจรสิ, เอหิ หีนายาวตฺติตฺวา โภเค จ ภุญฺชสฺสุ, ปุญฺญานิ จ กโรหี”ติ.}

\prose{โส วต อาวุโส ภิกฺขุ จตฺตาโร สติปฏฺฐาเน ภาเวนฺโต จตฺตาโร สติปฏฺฐาเน พหุลีกโรนฺโต สิกฺขํ ปจฺจกฺขาย หีนายาวตฺติสฺสตีติ เนตํ ฐานํ วิชฺชติ.\csromanspacer{} ตํ กิสฺส เหตุ, ยํ หิ ตํ อาวุโส จิตฺตํ ทีฆรตฺตํ วิเวกนินฺนํ วิเวกโปณํ วิเวกปพฺภารํ, ตํ วต หีนายาวตฺติสฺสตีติ เนตํ ฐานํ วิชฺชติ.\csromanspacer{} กถญฺจาวุโส ภิกฺขุ จตฺตาโร สติปฏฺฐาเน ภาเวติ, จตฺตาโร สติปฏฺฐาเน พหุลีกโรตีติ.\csromanspacer{} อิธาวุโส ภิกฺขุ กาเย กายานุปสฺสี วิหรติ ฯเปฯ.\csromanspacer{} เวทนาสุ ฯเปฯ.\csromanspacer{} จิตฺเต ฯเปฯ.\csromanspacer{} ธมฺเมสุ ธมฺมานุปสฺสี วิหรติ อาตาปี สมฺปชาโน สติมา วิเนยฺย โลเก อภิชฺฌาโทมนสฺสํ.\csromanspacer{} เอวํ โข อาวุโส ภิกฺขุ จตฺตาโร สติปฏฺฐาเน ภาเวติ.\csromanspacer{} จตฺตาโร สติปฏฺฐาเน พหุลีกโรตีติ.\csromanspacer{}\csromanspacer{}อฏฺฐมํ.}

\csromanmatikaanchor{mtk.196}
\csromantocmarkref{subparagraph}{9. อมฺพปาลิวนสุตฺต}{mtk.196}
\bookmark[dest={mtk.196},level=5]{9. อมฺพปาลิวนสุตฺต}
\csromanheader{6}{9. อมฺพปาลิวนสุตฺต}
\proseitem{907}{เอกํ สมยํ อายสฺมา จ อนุรุทฺโธ อายสฺมา จ สาริปุตฺโต เวสาลิยํ วิหรนฺติ อมฺพปาลิวเน.\csromanspacer{} อถ โข อายสฺมา สาริปุตฺโต สายนฺหสมยํ ปฏิสลฺลานา วุฏฺฐิโต ฯเปฯ เอกมนฺตํ นิสินฺโน โข อายสฺมา สาริปุตฺโต อายสฺมนฺตํ อนุรุทฺธํ เอตทโวจ\csromandash{}}

\prose{วิปฺปสนฺนานิ โข เต อาวุโส อนุรุทฺธ อินฺทฺริยานิ, ปริสุทฺโธ มุขวณฺโณ ปริโยทาโต, กตเมนายสฺมา อนุรุทฺโธ วิหาเรน เอตรหิ พหุลํ วิหรตีติ.\csromanspacer{} จตูสุ ขฺวาหํ อาวุโส สติปฏฺฐาเนสุ สุปฺปติฏฺฐิตจิตฺโต เอตรหิ พหุลํ วิหรามิ.\csromanspacer{} กตเมสุ จตูสุ, อิธาหํ อาวุโส กาเย กายานุปสฺสี วิหรามิ อาตาปี สมฺปชาโน สติมา วิเนยฺย โลเก อภิชฺฌาโทมนสฺสํ.\csromanspacer{} เวทนาสุ ฯเปฯ.\csromanspacer{} จิตฺเต ฯเปฯ.\csromanspacer{} ธมฺเมสุ ธมฺมานุปสฺสี วิหรามิ \csromanfolio{264}อาตาปี สมฺปชาโน สติมา วิเนยฺย โลเก อภิชฺฌาโทมนสฺสํ.\csromanspacer{} อิเมสุ ขฺวาหํ อาวุโส จตูสุ สติปฏฺฐาเนสุ สุปฺปติฏฺฐิตจิตฺโต เอตรหิ พหุลํ วิหรามิ.\csromanspacer{} โย โส อาวุโส ภิกฺขุ อรหํ ขีณาสโว วุสิตวา กตกรณีโย โอหิตภาโร อนุปฺปตฺตสทตฺโถ ปริกฺขีณภวสํโยชโน สมฺมทญฺญาวิมุตฺโต, โส อิเมสุ จตูสุ สติปฏฺฐาเนสุ สุปฺปติฏฺฐิตจิตฺโต พหุลํ วิหรตีติ.}

\prose{ลาภา วต โน อาวุโส, สุลทฺธํ วต โน อาวุโส, เย มยํ อายสฺมโต อนุรุทฺธสฺส สมฺมุขาว อสฺสุมฺห อาสภิํ วาจํ ภาสมานสฺสาติ.\csromanspacer{}\csromanspacer{}นวมํ.}

\csromanmatikaanchor{mtk.197}
\csromantocmarkref{subparagraph}{10. พาฬฺหคิลานสุตฺต}{mtk.197}
\bookmark[dest={mtk.197},level=5]{10. พาฬฺหคิลานสุตฺต}
\csromanheader{6}{10. พาฬฺหคิลานสุตฺต}
\proseitem{908}{เอกํ สมยํ อายสฺมา อนุรุทฺโธ สาวตฺถิยํ วิหรติ อนฺธวนสฺมิํ อาพาธิโก ทุกฺขิโต พาฬฺหคิลาโน.\csromanspacer{} อถ โข สมฺพหุลา ภิกฺขู เยนายสฺมา อนุรุทฺโธ เตนุปสงฺกมิํสุ, อุปสงฺกมิตฺวา อายสฺมนฺตํ อนุรุทฺธํ เอตทโวจุํ\csromandash{}}

\prose{กตเมนายสฺมโต อนุรุทฺธสฺส วิหาเรน วิหรโต อุปฺปนฺนา สารีริกา ทุกฺขา เวทนา จิตฺตํ น ปริยาทาย ติฏฺฐนฺตีติ.\csromanspacer{} จตูสุ โข เม อาวุโส สติปฏฺฐาเนสุ สุปฺปติฏฺฐิตจิตฺตสฺส วิหรโต อุปฺปนฺนา สารีริกา ทุกฺขา เวทนา จิตฺตํ น ปริยาทาย ติฏฺฐนฺติ.\csromanspacer{} กตเมสุ จตูสุ, อิธาหํ อาวุโส กาเย กายานุปสฺสี วิหรามิ ฯเปฯ.\csromanspacer{} เวทนาสุ ฯเปฯ.\csromanspacer{} จิตฺเต ฯเปฯ.\csromanspacer{} ธมฺเมสุ ธมฺมานุปสฺสี วิหรามิ อาตาปี สมฺปชาโน สติมา วิเนยฺย โลเก อภิชฺฌาโทมนสฺสํ.\csromanspacer{} อิเมสุ โข เม อาวุโส จตูสุ สติปฏฺฐาเนสุ สุปฺปติฏฺฐิตจิตฺตสฺส วิหรโต อุปฺปนฺนา สารีริกา ทุกฺขา เวทนา จิตฺตํ น ปริยาทาย ติฏฺฐนฺตีติ.\csromanspacer{}\csromanspacer{}ทสมํ.\csromanspacer{} รโหคตวคฺโค ปฐโม.}

\titlehead{\textbf{ตสฺสุทฺทานํ}}
\csromangathameasure{รโหคเตน เทฺว วุตฺตา, สุตนุ กณฺฑกี ตโย. \\ ตณฺหากฺขยสลฬาคารํ, อมฺพปาลิ จ คิลานนฺติ.}
\csromangathagroup{\csromangathastanza{รโหคเตน เทฺว วุตฺตา, สุตนุ กณฺฑกี ตโย. \\ ตณฺหากฺขยสลฬาคารํ, อมฺพปาลิ จ คิลานนฺติ.}}

\csromanheader{2}{8. อนุรุทฺธสํยุตฺต}
\csromanmatikaanchor{mtk.198}
\csromantocmarknopage{subparagraph}{2. ทุติยวคฺค}{mtk.198}
\bookmark[dest={mtk.198},level=5]{2. ทุติยวคฺค}
\csromanheader{6}{2. ทุติยวคฺค}
\csromanmatikaanchor{mtk.199}
\csromantocmarkref{subparagraph}{1. กปฺปสหสฺสสุตฺต}{mtk.199}
\bookmark[dest={mtk.199},level=5]{1. กปฺปสหสฺสสุตฺต}
\csromanheader{6}{1. กปฺปสหสฺสสุตฺต}
\csromanmatikaanchor{mtk.200}
\csromantocmarkref{subparagraph}{2-14. อิทฺธิวิธสุตฺตาทีนิ}{mtk.200}
\bookmark[dest={mtk.200},level=5]{2-14. อิทฺธิวิธสุตฺตาทีนิ}
\proseitem{909}{\csromanfolio{265}เอกํ สมยํ อายสฺมา อนุรุทฺโธ สาวตฺถิยํ วิหรติ เชตวเน อนาถปิณฺฑิกสฺส อาราเม.\csromanspacer{} อถ โข สมฺพหุลา ภิกฺขู เยนายสฺมา อนุรุทฺโธ เตนุปสงฺกมิํสุ, อุปสงฺกมิตฺวา อายสฺมตา อนุรุทฺเธน สทฺธิํ ฯเปฯ เอกมนฺตํ นิสินฺนา โข เต ภิกฺขู อายสฺมนฺตํ อนุรุทฺธํ เอตทโวจุํ\csromandash{}}

\prose{กตเมสํ อายสฺมา อนุรุทฺโธ ธมฺมานํ ภาวิตตฺตา พหุลีกตตฺตา มหาภิญฺญตํ ปตฺโตติ.\csromanspacer{} จตุนฺนํ ขฺวาหํ อาวุโส สติปฏฺฐานานํ ภาวิตตฺตา พหุลีกตตฺตา มหาภิญฺญตํ ปตฺโต.\csromanspacer{} กตเมสํ จตุนฺนํ, อิธาหํ อาวุโส กาเย กายานุปสฺสี วิหรามิ ฯเปฯ.\csromanspacer{} เวทนาสุ ฯเปฯ.\csromanspacer{} จิตฺเต ฯเปฯ.\csromanspacer{} ธมฺเมสุ ธมฺมานุปสฺสี วิหรามิ อาตาปี สมฺปชาโน สติมา วิเนยฺย โลเก อภิชฺฌาโทมนสฺสํ.\csromanspacer{} อิเมสํ ขฺวาหํ อาวุโส จตุนฺนํ สติปฏฺฐานานํ ภาวิตตฺตา พหุลีกตตฺตา มหาภิญฺญตํ ปตฺโต.\csromanspacer{} อิเมสญฺจ ปนาหํ อาวุโส จตุนฺนํ สติปฏฺฐานานํ ภาวิตตฺตา พหุลีกตตฺตา กปฺปสหสฺสํ อนุสฺสรามีติ.\csromanspacer{}\csromanspacer{}ปฐมํ.}

\csromanheader{2}{2. อิทฺธิวิธสุตฺต}
\proseitem{910}{อิเมสญฺจ ปนาหํ อาวุโส จตุนฺนํ สติปฏฺฐานานํ ภาวิตตฺตา พหุลีกตตฺตา อเนกวิหิตํ อิทฺธิวิธํ ปจฺจนุโภมิ, เอโกปิ หุตฺวา พหุธา โหมิ ฯเปฯ ยาว พฺรหฺมโลกาปิ กาเยน วสํ วตฺเตมีติ.\csromanspacer{}\csromanspacer{}ทุติยํ.}

\csromanheader{2}{3. ทิพฺพโสตสุตฺต}
\proseitem{911}{อิเมสญฺจ ปนาหํ อาวุโส จตุนฺนํ สติปฏฺฐานานํ ภาวิตตฺตา พหุลีกตตฺตา ทิพฺพาย โสตธาตุยา วิสุทฺธาย อติกฺกนฺตมานุสิกาย อุโภ สทฺเท สุณามิ ทิพฺเพ จ มานุเส จ เย ทูเร สนฺติเก จาติ.\csromanspacer{}\csromanspacer{}ตติยํ.}

\csromanheader{2}{4. เจโตปริยสุตฺต}
\proseitem{912}{อิเมสญฺจ ปนาหํ อาวุโส จตุนฺนํ สติปฏฺฐานานํ ภาวิตตฺตา พหุลีกตตฺตา ปรสตฺตานํ ปรปุคฺคลานํ เจตสา เจโต ปริจฺจ \csromanfolio{266}ปชานามิ, สราคํ วา จิตฺตํ “สราคํ จิตฺตนฺ”ติ ปชานามิ ฯเปฯ อวิมุตฺตํ วา จิตฺตํ “อวิมุตฺตํ จิตฺตนฺ”ติ ปชานามีติ.\csromanspacer{}\csromanspacer{}จตุตฺถํ.}

\csromanheader{2}{5. ฐานสุตฺต}
\proseitem{913}{อิเมสญฺจ ปนาหํ อาวุโส จตุนฺนํ สติปฏฺฐานานํ ภาวิตตฺตา พหุลีกตตฺตา ฐานญฺจ ฐานโต อฏฺฐานญฺจ อฏฺฐานโต ยถาภูตํ ปชานามีติ.\csromanspacer{}\csromanspacer{}ปญฺจมํ.}

\csromanheader{2}{6. กมฺมสมาทานสุตฺต}
\proseitem{914}{อิเมสญฺจ ปนาหํ อาวุโส จตุนฺนํ สติปฏฺฐานานํ ภาวิตตฺตา พหุลีกตตฺตา อตีตานาคตปจฺจุปฺปนฺนานํ กมฺมสมาทานานํ ฐานโส เหตุโส วิปากํ ยถาภูตํ ปชานามีติ.\csromanspacer{}\csromanspacer{}ฉฏฺฐํ.}

\csromanheader{2}{7. สพฺพตฺถคามินิสุตฺต}
\proseitem{915}{อิเมสญฺจ ปนาหํ อาวุโส จตุนฺนํ สติปฏฺฐานานํ ภาวิตตฺตา พหุลีกตตฺตา สพฺพตฺถคามินิปฺปฏิปทํ ยถาภูตํ ปชานามีติ.\csromanspacer{}\csromanspacer{}สตฺตมํ.}

\csromanheader{2}{8. นานาธาตุสุตฺต}
\proseitem{916}{อิเมสญฺจ ปนาหํ อาวุโส จตุนฺนํ สติปฏฺฐานานํ ภาวิตตฺตา พหุลีกตตฺตา อเนกธาตุนานาธาตุโลกํ ยถาภูตํ ปชานามีติ.\csromanspacer{}\csromanspacer{}อฏฺฐมํ.}

\csromanheader{2}{9. นานาธิมุตฺติสุตฺต}
\proseitem{917}{อิเมสญฺจ ปนาหํ อาวุโส จตุนฺนํ สติปฏฺฐานานํ ภาวิตตฺตา พหุลีกตตฺตา สตฺตานํ นานาธิมุตฺติกตํ ยถาภูตํ ปชานามีติ.\csromanspacer{}\csromanspacer{}นวมํ.}

\csromanheader{2}{10. อินฺทฺริยปโรปริยตฺตสุตฺต}
\proseitem{918}{อิเมสญฺจ ปนาหํ อาวุโส จตุนฺนํ สติปฏฺฐานานํ ภาวิตตฺตา พหุลีกตตฺตา ปรสตฺตานํ ปรปุคฺคลานํ อินฺทฺริยปโรปริยตฺตํ ยถาภูตํ ปชานามีติ.\csromanspacer{}\csromanspacer{}ทสมํ.}

\csromanheader{2}{8. อนุรุทฺธสํยุตฺต}
\csromanheader{2}{11. ฌานาทิสุตฺต}
\proseitem{919}{\csromanfolio{267}อิเมสญฺจ ปนาหํ อาวุโส จตุนฺนํ สติปฏฺฐานานํ ภาวิตตฺตา พหุลีกตตฺตา ฌานวิโมกฺขสมาธิสมาปตฺตีนํ สํกิเลสํ โวทานํ วุฏฺฐานํ ยถาภูตํ ปชานามีติ.\csromanspacer{}\csromanspacer{}เอกาทสมํ.}

\csromanheader{2}{12. ปุพฺเพนิวาสสุตฺต}
\proseitem{920}{อิเมสญฺจ ปนาหํ อาวุโส จตุนฺนํ สติปฏฺฐานานํ ภาวิตตฺตา พหุลีกตตฺตา อเนกวิหิตํ ปุพฺเพนิวาสํ อนุสฺสารามิ.\csromanspacer{} เสยฺยถิทํ, เอกมฺปิ ชาติํ เทฺวปิ ชาติโย ฯเปฯ.\csromanspacer{} อิติ สาการํ ส-อุทฺเทสํ อเนกวิหิตํ ปุพฺเพนิวาสํ อนุสฺสรามีติ.\csromanspacer{}\csromanspacer{}ทฺวาทสมํ.}

\csromanheader{2}{13. ทิพฺพจกฺขุสุตฺต}
\proseitem{921}{อิเมสญฺจ ปนาหํ อาวุโส จตุนฺนํ สติปฏฺฐานานํ ภาวิตตฺตา พหุลีกตตฺตา ทิพฺเพน จกฺขุนา วิสุทฺเธน อติกฺกนฺตมานุสเกน สตฺเต ปสฺสามิ จวมาเน อุปปชฺชมาเน ฯเปฯ.\csromanspacer{} อิติ ทิพฺเพน จกฺขุนา วิสุทฺเธน อติกฺกนฺตมานุสเกน ยถากมฺมูปเค สตฺเต ปชานามีติ.\csromanspacer{}\csromanspacer{}เตรสมํ.}

\csromanheader{2}{14. อาสวกฺขยสุตฺต}
\proseitem{922}{อิเมสญฺจ ปนาหํ อาวุโส จตุนฺนํ สติปฏฺฐานานํ ภาวิตตฺตา พหุลีกตฺตา อาสวานํ ขยา อนาสวํ เจโตวิมุตฺติํ ปญฺญาวิมุตฺติํ ทิฏฺเฐว ธมฺเม สยํ อภิญฺญา สจฺฉิกตฺวา อุปสมฺปชฺช วิหรามีติ.\csromanspacer{}\csromanspacer{}จุทฺทสมํ.\csromanspacer{} ทุติโย วคฺโค.}

\titlehead{\textbf{ตสฺสุทฺทานํ}}
\csromangathameasure{มหาภิญฺญํ อิทฺธิ ทิพฺพํ, เจโตปริยํ ฐานํ กมฺมํ. \\ สพฺพตฺถธาตุธิมุตฺติ, อินฺทฺริยํ ฌานํ ติสฺโส วิชฺชาติ. \\ \textbf{อนุรุทฺธสํยุตฺตํ อฏฺฐมํ.}}
\csromangathagroup{\csromangathastanza{มหาภิญฺญํ อิทฺธิ ทิพฺพํ, เจโตปริยํ ฐานํ กมฺมํ. \\ สพฺพตฺถธาตุธิมุตฺติ, อินฺทฺริยํ ฌานํ ติสฺโส วิชฺชาติ. \\ \textbf{อนุรุทฺธสํยุตฺตํ อฏฺฐมํ.}}}

\csromanmatikaanchor{mtk.201}
\csromantocmarknopage{subparagraph}{9. ฌานสํยุตฺต}{mtk.201}
\bookmark[dest={mtk.201},level=5]{9. ฌานสํยุตฺต}
\csromanheader{6}{9. ฌานสํยุตฺต}
\chapterhead{1. คงฺคาเปยฺยาลวคฺค}
\csromanmatikaanchor{mtk.202}
\csromanmatikaanchor{mtk.203}
\csromantocmarknopage{subparagraph}{1. คงฺคาเปยฺยาลวคฺค}{mtk.202}
\bookmark[dest={mtk.202},level=5]{1. คงฺคาเปยฺยาลวคฺค}
\csromantocmarkref{subparagraph}{1-12. ฌานาทิสุตฺต}{mtk.203}
\bookmark[dest={mtk.203},level=5]{1-12. ฌานาทิสุตฺต}
\csromanheader{6}{1-12. ฌานาทิสุตฺต}
\proseitem{923-934}{\csromanfolio{268}สาวตฺถินิทานํ.\csromanspacer{} ตตฺร โข ฯเปฯ จตฺตาโรเม ภิกฺขเว ฌานา.\csromanspacer{} กตเมน จตฺตาโร, อิธ ภิกฺขเว ภิกฺขุ วิวิจฺเจว กาเมหิ วิวิจฺจ อกุสเลหิ ธมฺเมหิ สวิตกฺกํ สวิจารํ วิเวกชํ ปีติสุขํ ปฐมํ ฌานํ อุปสมฺปชฺช วิหรติ.\csromanspacer{} วิตกฺกวิจารานํ วูปสมา อชฺฌตฺตํ สมฺปสาทนํ เจตโส เอโกทิภาวํ อวิตกฺกํ อวิจารํ สมาธิชํ ปีติสุขํ ทุติยํ ฌานํ อุปสมฺปชฺช วิหรติ, ปีติยา จ วิราคา อุเปกฺขโก จ วิหรติ สโต จ สมฺปชาโน สุขญฺจ กาเยน ปฏิสํเวเทติ, ยํ ตํ อริยา อาจิกฺขนฺติ “อุเปกฺขโก สติมา สุขวิหารี”ติ, ตติยํ ฌานํ อุปสมฺปชฺช วิหรติ, สุขสฺส จ ปหานา ทุกฺขสฺส จ ปหานา ปุพฺเพว โสมนสฺสโทมนสฺสานํ อตฺถงฺคมา อทุกฺขมสุขํ อุเปกฺขาสติปาริสุทฺธิํ จตุตฺถํ ฌานํ อุปสมฺปชฺช วิหรติ.\csromanspacer{} อิเม โข ภิกฺขเว จตฺตาโร ฌานาติ.}

\prose{เสยฺยถาปิ ภิกฺขเว คงฺคา นที ปาจีนนินฺนา ปาจีนโปณา ปาจีนปพฺภารา, เอวเมว โข ภิกฺขเว ภิกฺขุ จตฺตาโร ฌาเน ภาเวนฺโต จตฺตาโร ฌาเน พหุลีกโรนฺโต นิพฺพานนินฺโน โหติ นิพฺพานโปโณ นิพฺพานปพฺภาโร.\csromanspacer{} กถญฺจ ภิกฺขเว ภิกฺขุ จตฺตาโร ฌาเน ภาเวนฺโต จตฺตาโร ฌาเน พหุลีกโรนฺโต นิพฺพานนินฺโน โหติ นิพฺพานโปโณ นิพฺพานปพฺภาโร, อิธ ภิกฺขเว ภิกฺขุ วิวิจฺเจว กาเมหิ วิวิจฺจ อกุสเลหิ ธมฺเมหิ สวิตกฺกํ สวิจารํ วิเวกชํ ปีติสุขํ ปฐมํ ฌานํ อุปสมฺปชฺช วิหรติ, วิตกฺกวิจารานํ วูปสมา ฯเปฯ ทุติยํ ฌานํ.\csromanspacer{} ตติยํ ฌานํ.\csromanspacer{} จตุตฺถํ ฌานํ อุปสมฺปชฺช วิหรติ.\csromanspacer{} เอวํ โข ภิกฺขเว ภิกฺขุ จตฺตาโร ฌาเน ภาเวนฺโต จตฺตาโร ฌาเน พหุลีกโรนฺโต นิพฺพานนินฺโน โหติ นิพฺพานโปโณ นิพฺพานปพฺภาโรติ.\csromanspacer{}\csromanspacer{}ทฺวาทสมํ.\csromanspacer{} คงฺคาเปยฺยาลวคฺโค ปฐโม.}

\csromanheader{2}{9. ฌานสํยุตฺต}
\titlehead{\textbf{ตสฺสุทฺทานํ}}
\csromangathameasure{ฉ ปาจีนโต นินฺนา, ฉ นินฺนา จ สมุทฺทโต. \\ เทฺวเต ฉ ทฺวาทส โหนฺติ, วคฺโค เตน ปวุจฺจตีติ. อปฺปมาทวคฺโค วิตฺถาเรตพฺโพ.}
\csromangathagroup{\csromangathastanza{\csromanfolio{269}ฉ ปาจีนโต นินฺนา, ฉ นินฺนา จ สมุทฺทโต. \\ เทฺวเต ฉ ทฺวาทส โหนฺติ, วคฺโค เตน ปวุจฺจตีติ. อปฺปมาทวคฺโค วิตฺถาเรตพฺโพ.}}

\titlehead{\textbf{ตสฺสุทฺทานํ}}
\csromangathameasure{ตถาคตํ ปทํ กูฏํ, มูลํ สาโร จ วสฺสิกํ. \\ ราชา จนฺทิมสูริยา, วตฺเถน ทสมํ ปทนฺติ. พลกรณียวคฺโค วิตฺถาเรตพฺโพ.}
\csromangathagroup{\csromangathastanza{ตถาคตํ ปทํ กูฏํ, มูลํ สาโร จ วสฺสิกํ. \\ ราชา จนฺทิมสูริยา, วตฺเถน ทสมํ ปทนฺติ. พลกรณียวคฺโค วิตฺถาเรตพฺโพ.}}

\titlehead{\textbf{ตสฺสุทฺทานํ}}
\csromangathameasure{พลํ พีชญฺจ นาโค จ, รุกฺโข กุมฺเภน สูกิยา. \\ อากาเสน จ เทฺว เมฆา, นาวา อาคนฺตุกา นทีติ. เอสนาวคฺโค วิตฺถาเรตพฺโพ.}
\csromangathagroup{\csromangathastanza{พลํ พีชญฺจ นาโค จ, รุกฺโข กุมฺเภน สูกิยา. \\ อากาเสน จ เทฺว เมฆา, นาวา อาคนฺตุกา นทีติ. เอสนาวคฺโค วิตฺถาเรตพฺโพ.}}

\titlehead{\textbf{ตสฺสุทฺทานํ}}
\csromangathameasure{เอสนา วิธา อาสโว, ภโว จ ทุกฺขตา ติสฺโส. \\ ขิลํ มลญฺจ นีโฆ จ, เวทนา ตณฺหา ตสินา จาติ.}
\csromangathagroup{\csromangathastanza{เอสนา วิธา อาสโว, ภโว จ ทุกฺขตา ติสฺโส. \\ ขิลํ มลญฺจ นีโฆ จ, เวทนา ตณฺหา ตสินา จาติ.}}

\chapterhead{5. โอฆวคฺค}
\csromanmatikaanchor{mtk.204}
\csromanmatikaanchor{mtk.205}
\csromantocmarknopage{subparagraph}{5. โอฆวคฺค}{mtk.204}
\bookmark[dest={mtk.204},level=5]{5. โอฆวคฺค}
\csromantocmarkref{subparagraph}{1-10. โอฆาทิสุตฺต}{mtk.205}
\bookmark[dest={mtk.205},level=5]{1-10. โอฆาทิสุตฺต}
\csromanheader{6}{1-10. โอฆาทิสุตฺต}
\proseitem{967-976}{ปญฺจิมานิ ภิกฺขเว อุทฺธมฺภาคิยานิ สํโยชนานิ.\csromanspacer{} กตมานิ ปญฺจ, รูปราโค อรูปราโค มาโน อุทฺธจฺจํ อวิชฺชา.\csromanspacer{} อิมานิ โข ภิกฺขเว \csromanfolio{270}ปญฺจุทฺธมฺภาคิยานิ สํโยชนานิ.\csromanspacer{} อิเมสํ โข ภิกฺขเว ปญฺจนฺนํ อุทฺธมฺภาคิยานํ สํโยชนานํ อภิญฺญาย ปริญฺญาย ปริกฺขยาย ปหานาย จตฺตาโร ฌาเน ภาเวตพฺพา.\csromanspacer{} กตเม จตฺตาโร, อิธ ภิกฺขเว ภิกฺขุ วิวิจฺเจว กาเมหิ วิวิจฺจ อกุสเลหิ ธมฺเมหิ สวิตกฺกํ สวิจารํ วิเวกชํ ปีติสุขํ ปฐมํ ฌานํ อุปสมฺปชฺช วิหรติ, วิตกฺกวิจารานํ วูปสมา อชฺฌตฺตํ สมฺปสาทนํ เจตโส เอโกทิภาวํ อวิตกฺกํ อวิจารํ สมาธิชํ ปีติสุขํ ทุติยํ ฌานํ ฯเปฯ ตติยํ ฌานํ ฯเปฯ จตุตฺถํ ฌานํ อุปสมฺปชฺช วิหรติ, อิเมสํ โข ภิกฺขเว ปญฺจนฺนํ อุทฺธมฺภาคิยานํ สํโยชนานํ อภิญฺญาย ปริญฺญาย ปริกฺขยาย ปหานาย อิเม จตฺตาโร ฌานา ภาเวตพฺพาติ วิตฺถาเรตพฺพํ. (ยถา มคฺคสํยุตฺตํ, ตถา วิตฺถาเรตพฺพํ.) โอฆวคฺโค ปญฺจโม.}

\titlehead{\textbf{ตสฺสุทฺทานํ}}
\csromangathameasure{โอโฆ โยโค อุปาทานํ, คนฺถา อนุสเยน จ. \\ กามคุณา นีวรณา, ขนฺธา โอรุทฺธมฺภาคิยาติ. \\ \textbf{ฌานสํยุตฺตํ นวมํ.}}
\csromangathagroup{\csromangathastanza{โอโฆ โยโค อุปาทานํ, คนฺถา อนุสเยน จ. \\ กามคุณา นีวรณา, ขนฺธา โอรุทฺธมฺภาคิยาติ. \\ \textbf{ฌานสํยุตฺตํ นวมํ.}}}

\csromanmatikaanchor{mtk.206}
\csromantocmarknopage{subparagraph}{10. อานาปานสํยุตฺต}{mtk.206}
\bookmark[dest={mtk.206},level=5]{10. อานาปานสํยุตฺต}
\csromanheader{6}{10. อานาปานสํยุตฺต}
\csromanmatikaanchor{mtk.207}
\csromantocmarknopage{subparagraph}{1. เอกธมฺมวคฺค}{mtk.207}
\bookmark[dest={mtk.207},level=5]{1. เอกธมฺมวคฺค}
\csromanheader{6}{1. เอกธมฺมวคฺค}
\csromanmatikaanchor{mtk.208}
\csromantocmarkref{subparagraph}{1. เอกธมฺมสุตฺต}{mtk.208}
\bookmark[dest={mtk.208},level=5]{1. เอกธมฺมสุตฺต}
\csromanheader{6}{1. เอกธมฺมสุตฺต}
\csromanmatikaanchor{mtk.209}
\csromantocmarkref{subparagraph}{2-4. โพชฺฌงฺคสุตฺตาทีนิ}{mtk.209}
\bookmark[dest={mtk.209},level=5]{2-4. โพชฺฌงฺคสุตฺตาทีนิ}
\proseitem{977}{\csromanfolio{271}สาวตฺถินิทานํ.\csromanspacer{} ตตฺร โข ฯเปฯ เอตทโวจ\csromandash{}เอกธมฺโม ภิกฺขเว ภาวิโต พหุลีกโต มหปฺผโล โหติ มหานิสํโส.\csromanspacer{} กตโม เอกธมฺโม, อานาปานสฺสติ\footnote{อานาปานสติ (สี, อิ)}.\csromanspacer{} กถํ ภาวิตา จ ภิกฺขเว อานาปานสฺสติ กถํ พหุลีกตา มหปฺผลา โหติ มหานิสํสา.\csromanspacer{} อิธ ภิกฺขเว ภิกฺขุ อรญฺญคโต วา รุกฺขมูลคโต วา สุญฺญาคารคโต วา นิสีทติ ปลฺลงฺกํ อาภุชิตฺวา อุชุํ กายํ ปณิธาย ปริมุขํ สติํ อุปฏฺฐเปตฺวา, โส สโตว อสฺสสติ, สโตว\footnote{สโต (พหูสุ) ตติยปาราชิเกปิ.} ปสฺสสติ.\csromanspacer{} ทีฆํ วา อสฺสสนฺโต “ทีฆํ อสฺสสามี”ติ ปชานาติ, ทีฆํ วา ปสฺสสนฺโต “ทีฆํ ปสฺสสามี”ติ ปชานาติ, รสฺสํ วา ปสฺสสนฺโต “รสฺสํ ปสฺสสามี”ติ ปชานาติ, รสฺสํ วา ปสฺสสนฺโต “รสฺสํ ปสฺสสามี”ติ ปชานาติ.\csromanspacer{} สพฺพกายปฺปฏิสํเวที “อสฺสสิสฺสามี”ติ สิกฺขติ, สพฺพกายปฺปฏิสํเวที “ปสฺสสิสฺสามี”ติ สิกฺขติ, ปสฺสมฺภยํ กายสงฺขารํ อสฺสสิสฺสามี”ติ สิกฺขติ, ปสฺสมฺภยํ กายสงฺขารํ “ปสฺสสิสฺสามี”ติ สิกฺขติ, ปีติปฺปฏิสํเวที “อสฺสสิสฺสามี”ติ สิกฺขติ, ปีติปฺปฏิสํเวที “ปสฺสสิสฺสามี”ติ สิกฺขติ, สุขปฺปฏิสํเวที “อสฺสสิสฺสามี”ติ สิกฺขติ สุขปฺปฏิสํเวที “ปสฺสสิสฺสามี”ติ สิกฺขติ, จิตฺตสงฺขารปฺปฏิสํเวที “อสฺสสิสฺสามี”ติ สิกฺขติ, จิตฺตสงฺขารปฺปฏิสํเวที “ปสฺสสิสฺสามี”ติ สิกฺขติ, ปสฺสมฺภยํ จิตฺตสงฺขารํ “อสฺสสิสฺสามี”ติ สิกฺขติ, ปสฺสมฺภยํ จิตฺตสงฺขารํ “ปสสิสฺสามี”ติ สิกฺขติ, จิตฺตปฺปฏิสํเวที “อสฺสสิสฺสามี”ติ สิกฺขติ, จิตฺตปฺปฏิสํเวที “ปสฺสสิสฺสามี”ติ สิกฺขติ, อภิปฺปโมทยํ จิตฺตํ “อสฺสสิสฺสามี”ติ สิกฺขติ, อภิปฺปโมทยํ จิตฺตํ “ปสฺสสิสฺสามี”ติ สิกฺขติ, สมาทหํ จิตฺตํ “อสฺสสิสฺสามี”ติ สิกฺขติ, สมาทหํ จิตฺตํ “ปสฺสสิสฺสามี”ติ สิกฺขติ, วิโมจยํ จิตฺตํ “อสฺสสิสฺสามี”ติ สิกฺขติ, \csromanfolio{272}วิโมจยํ จิตฺตํ “ปสฺสสิสฺสามี”ติ สิกฺขติ, อนิจฺจานุปสฺสี “อสฺสสิสฺสามี”ติ สิกฺขติ, อนิจฺจานุปสฺสี “ปสฺสสิสฺสามี”ติ สิกฺขติ, วิราคานุปสฺสี “อสฺสสิสฺสามี”ติ สิกฺขติ, วิราคานุปสฺสี “ปสฺสสิสฺสามี”ติ สิกฺขติ, “นิโรธานุปสฺสี อสฺสสิสฺสามี”ติ สิกฺขติ, “นิโรธานุปสฺสี ปสฺสสิสฺสามี”ติ สิกฺขติ, “ปฏินิสฺสคฺคานุปสฺสี อสฺสสิสฺสามี”ติ สิกฺขติ, “ปฏินิสฺสคฺคานุปสฺสี ปสฺสสิสฺสามี”ติ สิกฺขติ.\csromanspacer{} เอวํ ภาวิตา โข ภิกฺขเว อานาปานสฺสติ เอวํ พหุลีกตา มหปฺผลา โหติ มหานิสํสาติ.\csromanspacer{}\csromanspacer{}ปฐมํ.}

\csromanheader{2}{2. โพชฺฌงฺคสุตฺต}
\proseitem{978}{อานาปานสฺสติ ภิกฺขเว ภาวิตา พหูลีกตา มหปฺผลา โหติ มหานิสํสา.\csromanspacer{} กถํ ภาวิตา จ ภิกฺขเว อานาปานสฺสติ กถํ พหุลีกตา มหปฺผลา โหติ มหานิสํสา, อิธ ภิกฺขเว ภิกฺขุ อานาปานสฺสติสหคตํ สติสมฺโพชฺฌงฺคํ ภาเวติ วิเวกนิสฺสิตํ วิราคนิสฺสิตํ นิโรธนิสฺสิตํ โวสฺสคฺคปริณามิํ.\csromanspacer{} อานาปานสฺสติสหคตํ ธมฺมวิจยสมฺโพชฺฌงฺคํ ภาเวติ ฯเปฯ.\csromanspacer{} อานาปานสฺสติสหคตํ อุเปกฺขาสมฺโพชฺฌงฺคํ ภาเวติ วิเวกนิสฺสิตํ วิราคนิสฺสิตํ นิโรธนิสฺสิตํ โวสฺสคฺคปริณามิํ.\csromanspacer{} เอวํ ภาวิตา โข ภิกฺขเว อานาปานสฺสติ เอวํ พหุลีกตา มหปฺผลา โหติ มหานิสํสาติ.\csromanspacer{}\csromanspacer{}ทุติยํ.}

\csromanheader{2}{3. สุทฺธิกสุตฺต}
\proseitem{979}{อานาปานสฺสติ ภิกฺขเว ภาวิตา พหุลีกตา มหปฺผลา โหติ มหานิสํสา.\csromanspacer{} กถํ ภาวิตา จ ภิกฺขเว อานาปานสฺสติ กถํ พหุลีกตา มหปฺผลา โหติ มหานิสํสา, อิธ ภิกฺขเว ภิกฺขุ อรญฺญคโต วา รุกฺขมูลคโต วา สุญฺญาคารคโต วา นิสีทติ ปลฺลงฺกํ อาภุชิตฺวา อุชุํ กายํ ปณิธาย ปริมุขํ สติํ อุปฏฺฐเปตฺวา, โส สโตว อสฺสสติ, สโตว ปสฺสสติ ฯเปฯ “ปฏินิสฺสคฺคานุปสฺสี อสฺสสิสฺสามี”ติ สิกฺขติ, “ปฏินิสฺสคฺคานุปสฺสี ปสฺสสิสฺสามี”ติ สิกฺขติ.\csromanspacer{} เอวํ ภาวิตา โข ภิกฺขเว อานาปานสฺสติ เอวํ พหุลีกตา มหปฺผลา โหติ มหานิสํสาติ.\csromanspacer{}\csromanspacer{}ตติยํ.}

\csromanheader{2}{10. อานาปานสํยุตฺต}
\csromanheader{2}{4. ปฐมผลสุตฺต}
\csromanmatikaanchor{mtk.210}
\csromantocmarkref{subparagraph}{5-6. ทุติยผลสุตฺตาทีนิ}{mtk.210}
\bookmark[dest={mtk.210},level=5]{5-6. ทุติยผลสุตฺตาทีนิ}
\proseitem{980}{\csromanfolio{273}อานาปานสฺสติ ภิกฺขเว ภาวิตา พหุลีกตา มหปฺผลา โหติ มหานิสํสา.\csromanspacer{} กถํ ภาวิตา จ ภิกฺขเว อานาปานสฺสติ กถํ พหุลีกตา มหปฺผลา โหติ มหานิสํสา, อิธ ภิกฺขเว ภิกฺขุ อรญฺญคโต วา รุกฺขมูลคโต วา สุญฺญาคารคโต วา นิสีทติ ปลฺลงฺกํ อาภุชิตฺวา อุชุํ กายํ ปณิธาย ปริมุขํ สติํ อุปฏฺฐเปตฺวา, โส สโตว อสฺสสติ, สโตว ปสฺสสติ ฯเปฯ “ปฏินิสฺสคฺคานุปสฺสี อสฺสสิสฺสามี”ติ สิกฺขติ, “ปฏินิสฺสคฺคานุปสฺสี ปสฺสสิสฺสามี”ติ สิกฺขติ.\csromanspacer{} เอวํ ภาวิตา โข ภิกฺขเว อานาปานสฺสติ เอวํ พหุลีกตา มหปฺผลา โหติ มหานิสํสา.\csromanspacer{} เอวํ ภาวิตาย โข ภิกฺขเว อานาปานสฺสติยา เอวํ พหุลีกตาย ทฺวินฺนํ ผลานํ อญฺญตรํ ผลํ ปาฏิกงฺขํ ทิฏฺเฐว ธมฺเม อญฺญา, สติ วา อุปาทิเสเส อนาคามิตาติ.\csromanspacer{}\csromanspacer{}จตุตฺถํ.}

\csromanheader{2}{5. ทุติยผลสุตฺต}
\proseitem{981}{อานาปานสฺสติ ภิกฺขเว ภาวิตา พหุลีกตา มหปฺผลา โหติ มหานิสํสา.\csromanspacer{} กถํ ภาวิตา จ ภิกฺขเว อานาปานสฺสติ กถํ พหุลีกตา มหปฺผลา โหติ มหานิสํสา, อิธ ภิกฺขเว ภิกฺขุ อรญฺญคโต วา รุกฺขมูลคโต วา สุญฺญาคารคโต วา นิสีทติ ปลฺลงฺกํ อาภุชิตฺวา อุชุํ กายํ ปณิธาย ปริมุขํ สติํ อุปฏฺฐเปตฺวา, โส สโตว อสฺสสติ, สโตว ปสฺสสติ ฯเปฯ “ปฏินิสฺสคฺคานุปสฺสี อสฺสสิสฺสามี”ติ สิกฺขติ, “ปฏินิสฺสคฺคานุปสฺสี ปสฺสสิสฺสามี”ติ สิกฺขติ.\csromanspacer{} เอวํ ภาวิตา โข ภิกฺขเว อานาปานสฺสติ เอวํ พหุลีกตา มหปฺผลา โหติ มหานิสํสา.}

\prose{เอวํ ภาวิตาย โข ภิกฺขเว อานาปานสฺสติยา เอวํ พหุลีกตาย สตฺต ผลา สตฺตานิสํสา ปาฏิกงฺขา.\csromanspacer{} กตเม สตฺต ผลา สตฺตานิสํสา, ทิฏฺเฐว ธมฺเม ปฏิกจฺจ อญฺญํ อาราเธติ, โน เจ ทิฏฺเฐว ธมฺเม ปฏิกจฺจ อญฺญํ อาราเธติ, อถ มรณกาเล อญฺญํ อาราเธติ, โน เจ ทิฏฺเฐว ธมฺเม ปฏิกจฺจ อญฺญํ อาราเธติ, โน เจ มรณกาเล อญฺญํ อาราเธติ, อถ ปญฺจนฺนํ โอรมฺภาคิยานํ สํโยชนานํ ปริกฺขยา อนฺตราปรินิพฺพายี โหติ, อุปหจฺจปรินิพฺพายี โหติ, อสงฺขารปรินิพฺพายี โหติ, สสงฺขารปรินิพฺพายี โหติ, อุทฺธํโสโต โหติ อกนิฏฺฐคามี.\csromanspacer{} เอวํ \csromanfolio{274}ภาวิตาย โข ภิกฺขเว อานาปานสฺสติยา เอวํ พหุลีกตาย อิเม สตฺต ผลา สตฺตานิสํสา ปาฏิกงฺขาติ.\csromanspacer{}\csromanspacer{}ปญฺจมํ.}

\csromanheader{2}{6. อริฏฺฐสุตฺต}
\proseitem{982}{สาวตฺถินิทานํ.\csromanspacer{} ตตฺร โข ภควา ฯเปฯ เอตทโวจ “ภาเวถ โน ตุเมฺห ภิกฺขเว อานาปานสฺสตินฺ”ติ.\csromanspacer{} เอวํ วุตฺเต อายสฺมา อริฏฺโฐ ภควนฺตํ เอตทโวจ “อหํ โข ภนฺเต ภาเวมิ อานาปานสฺสตินฺ”ติ.\csromanspacer{} ยถา กถํ ปน ตฺวํ อริฏฺฐ ภาเวสิ อานาปานสฺสตินฺติ.\csromanspacer{} อตีเตสุ เม ภนฺเต กาเมสุ กามจฺฉนฺโท ปหีโน, อนาคเตสุ เม กาเมสุ กามจฺฉนฺโท วิคโต, อชฺฌตฺตพหิทฺธา\footnote{อชฺฌตฺตํ พหิทฺธา (สฺยา, กํ, อิ, ก)} จ เม ธมฺเมสุ ปฏิฆสญฺญา สุปฺปฏิวินีตา, โส\footnote{โสหํ (?)} สโตว อสฺสสิสฺสามิ, สโตว ปสฺสสิสฺสามิ.\csromanspacer{} เอวํ ขฺวาหํ ภนฺเต ภาเวมิ อานาปานสฺสตินฺติ.}

\prose{อตฺเถสา อริฏฺฐ อานาปานสฺสติ, เนสา นตฺถีติ วทามิ, อปิ จ อริฏฺฐ ยถา อานาปานสฺสติ วิตฺถาเรน ปริปุณฺณา โหติ, ตํ สุณาหิ สาธุกํ มนสิ กโรหิ, ภาสิสฺสามีติ. “เอวํ ภนฺเต”ติ โข อายสฺมา อริฏฺโฐ ภควโต ปจฺจสฺโสสิ.\csromanspacer{} ภควา เอตทโวจ\csromandash{}}

\prose{กถญฺจ อริฏฺฐ อานาปานสฺสติ วิตฺถาเรน ปริปุณฺณา โหติ, อิธ อริฏฺฐ ภิกฺขุ อรญฺญคโต วา รุกฺขมูลคโต วา สุญฺญาคารคโต วา นิสีทติ ปลฺลงฺกํ อาภุชิตฺวา อุชุํ กายํ ปณิธาย ปริมุขํ สติํ อุปฏฺฐเปตฺวา, โส สโตว อสฺสสติ, สโตว ปสฺสสติ, ทีฆํ วา อสฺสสนฺโต “ทีฆํ อสฺสสามี”ติ ปชานาติ ฯเปฯ “ปฏินิสฺสคฺคานุปสฺสี อสฺสสิสฺสามี”ติ สิกฺขติ, “ปฏินิสฺสคฺคานุปสฺสี ปสฺสสิสฺสามี”ติ สิกฺขติ.\csromanspacer{} เอวํ โข อริฏฺฐ อานาปานสฺสติ วิตฺถาเรน ปริปุณฺณา โหตีติ.\csromanspacer{}\csromanspacer{}ฉฏฺฐํ.}

\csromanmatikaanchor{mtk.211}
\csromantocmarkref{subparagraph}{7. มหากปฺปินสุตฺต}{mtk.211}
\bookmark[dest={mtk.211},level=5]{7. มหากปฺปินสุตฺต}
\csromanheader{6}{7. มหากปฺปินสุตฺต}
\proseitem{983}{สาวตฺถินิทานํ.\csromanspacer{} เตน โข ปน สมเยน อายสฺมา มหากปฺปิโน ภควโต อวิทูเร นิสินฺโน โหติ ปลฺลงฺกํ อาภุชิตฺวา อุชุํ กายํ ปณิธาย ปริมุขํ สติํ อุปฏฺฐเปตฺวา.\csromanspacer{} อทฺทสา โข ภควา อายสฺมนฺตํ มหากปฺปินํ อวิทูเร นิสินฺนํ ปลฺลงฺกํ อาภุชิตฺวา อุชุํ กายํ ปณิธาย ปริมุขํ สติํ อุปฏฺฐเปตฺวา, ทิสฺวาน ภิกฺขู อามนฺเตสิ\csromandash{}}

\csromanheader{2}{10. อานาปานสํยุตฺต}
\prose{\csromanfolio{275}ปสฺสถ โน ตุเมฺห ภิกฺขเว เอตสฺส ภิกฺขุโน กายสฺส อิญฺชิตตฺตํ วา ผนฺทิตตฺตํ วาติ.\csromanspacer{} ยทาปิ มยํ ภนฺเต ตํ อายสฺมนฺตํ ปสฺสาม สํฆมชฺเฌ วา นิสินฺนํ เอกํ วา รโห นิสินฺนํ, ตทาปิ มยํ ตสฺส อายสฺมโต น ปสฺสาม กายสฺส อิญฺชิตตฺตํ วา ผนฺทิตตฺตํ วาติ.}

\prose{ยสฺส ภิกฺขเว สมาธิสฺส ภาวิตตฺตา พหุลีกตตฺตา เนว กายสฺส อิญฺชิตตฺตํ วา โหติ ผนฺทิตตฺตํ วา, น จิตฺตสฺส อิญฺชิตตฺตํ วา โหติ ผนฺทิตตฺตํ วา, ตสฺส โส ภิกฺขเว ภิกฺขุ สมาธิสฺส นิกามลาภี อกิจฺฉลาภี อกสิรลาภี.\csromanspacer{} กตมสฺส จ ภิกฺขเว สมาธิสฺส ภาวิตตฺตา พหุลีกตตฺตา เนว กายสฺส อิญฺชิตตฺตํ วา โหติ ผนฺทิตตฺตํ วา, น จิตฺตสฺส อิญฺชิตตฺตํ วา โหติ ผนฺทิตตฺตํ วา.}

\prose{อานาปานสฺสติสมาธิสฺส ภิกฺขเว ภาวิตตฺตา พหุลีกตตฺตา เนว กายสฺส อิญฺชิตตฺตํ วา โหติ ผนฺทิตตฺตํ วา, น จิตฺตสฺส อิญฺชิตตฺตํ วา โหติ ผนฺทิตตฺตํ วา.\csromanspacer{} กถํ ภาวิเต จ ภิกฺขเว อานาปานสฺสติสมาธิมฺหิ กถํ พหุลีกเต เนว กายสฺส อิญฺชิตตฺตํ วา โหติ ผนฺทิตตฺตํ วา, น จิตฺตสฺส อิญฺชิตตฺตํ วา โหติ ผนฺทิตตฺตํ วา.}

\prose{อิธ ภิกฺขเว ภิกฺขุ อรญฺญคโต วา รุกฺขมูลคโต วา สุญฺญาคารคโต วา นิสีทติ ปลฺลงฺกํ อาภุชิตฺวา อุชุํ กายํ ปณิธาย ปริมุขํ สติํ อุปฏฺฐเปตฺวา, โส สโตว อสฺสสติ, สโตว ปสฺสสติ ฯเปฯ “ปฏินิสฺสคฺคานุปสฺสี อสฺสสิสฺสามี”ติ สิกฺขติ, “ปฏินิสฺสคฺคานุปสฺสี ปสฺสสิสฺสามี”ติ สิกฺขติ.\csromanspacer{} เอวํ ภาวิเต จ โข ภิกฺขเว อานาปานสฺสติสมาธิมฺหิ เอวํ พหุลีกเต เนว กายสฺส อิญฺชิตตฺตํ วา โหติ ผนฺทิตตฺตํ วา, น จิตฺตสฺส อิญฺชิตตฺตํ วา โหติ ผนฺทิตตฺตํ วาติ.\csromanspacer{}\csromanspacer{}สตฺตมํ.}

\csromanmatikaanchor{mtk.212}
\csromantocmarkref{subparagraph}{8. ปทีโปปมสุตฺต}{mtk.212}
\bookmark[dest={mtk.212},level=5]{8. ปทีโปปมสุตฺต}
\csromanheader{6}{8. ปทีโปปมสุตฺต}
\proseitem{984}{อานาปานสฺสติสมาธิ ภิกฺขเว ภาวิโต พหุลีกโต มหปฺผโล โหติ มหานิสํโส.\csromanspacer{} กถํ ภาวิโต จ ภิกฺขเว อานาปานสฺสติสมาธิ กถํ พหุลีกโต มหปฺผโล โหติ มหานิสํโส.}

\prose{อิธ ภิกฺขเว ภิกฺขุ อรญฺญคโต วา รุกฺขมูลคโต วา สุญฺญาคารคโต วา นิสีทติ ปลฺลงฺกํ อาภุชิตฺวา อุชุํ กายํ ปณิธาย \csromanfolio{276}ปริมุขํ สติํ อุปฏฺฐเปตฺวา, โส สโตว อสฺสสติ, สโตว ปสฺสสติ, ทีฆํ วา อสฺสสนฺโต “ทีฆํ อสฺสสามี”ติ ปชานาติ ฯเปฯ “ปฏินิสฺสคฺคานุปสฺสี อสฺสสิสฺสามี”ติ สิกฺขติ, “ปฏินิสฺสคฺคานุปสฺสี ปสฺสสิสฺสามี”ติ สิกฺขติ.\csromanspacer{} เอวํ ภาวิโต โข ภิกฺขเว อานาปานสฺสติสมาธิ เอวํ พหุลีกโต มหปฺผโล โหติ มหานิสํโส.}

\prose{อหมฺปิ สุทํ ภิกฺขเว ปุพฺเพว สมฺโพธา อนภิสมฺพุทฺโธ โพธิสตฺโตว สมาโน อิมินา วิหาเรน พหุลํ วิหรามิ, ตสฺส มยฺหํ ภิกฺขเว อิมินา วิหาเรน พหุลํ วิหรโต เนว กาโย กิลมติ น จกฺขูนิ, อนุปาทาย จ เม อาสเวหิ จิตฺตํ วิมุจฺจิ.}

\prose{ตสฺมาติห ภิกฺขเว ภิกฺขุ เจปิ อากงฺเขยฺย “เนว เม กาโย กิลเมยฺย น จกฺขูนิ, อนุปาทาย จ เม อาสเวหิ จิตฺตํ วิมุจฺเจยฺยา”ติ, อยเมว อานาปานสฺสติสมาธิ สาธุกํ มนสิ กาตพฺโพ.\csromanspacer{} ตสฺมาติห ภิกฺขเว ภิกฺขุ เจปิ อากงฺเขยฺย “เย เม เคหสิตา สรสงฺกปฺปา, เต ปหีเยยฺยุนฺ”ติ, อยเมว อานาปานสฺสติสมาธิ สาธุกํ มนสิ กาตพฺโพ.\csromanspacer{} ตสฺมาติห ภิกฺขเว ภิกฺขุ เจปิ อากงฺเขยฺย “อปฺปฏิกูเล ปฏิกูลสญฺญี วิหเรยฺยนฺ”ติ, อยเมว อานาปานสฺสติสมาธิ สาธุกํ มนสิ กาตพฺโพ.\csromanspacer{} ตสฺมาติห ภิกฺขเว ภิกฺขุ เจปิ อากงฺเขยฺย “ปฏิกูเล อปฺปฏิกูลสญฺญี วิหเรยฺยนฺ”ติ, อยเมว อานาปานสฺสติสมาธิ สาธุกํ มนสิ กาตพฺโพ.\csromanspacer{} ตสฺมาติห ภิกฺขเว ภิกฺขุ เจปิ อากงฺเขยฺย “ปฏิกูเล จ อปฺปฏิกูเล จ ปฏิกูลสญฺญี วิหเรยฺยนฺ”ติ, อยเมว อานาปานสฺสติสมาธิ สาธุกํ มนสิ กาตพฺโพ.\csromanspacer{} ตสฺมาติห ภิกฺขเว ภิกฺขุ เจปิ อากงฺเขยฺย “ปฏิกูเล จ อปฺปฏิกูเล, จ อปฺปฏิกูลสญฺญี วิหเรยฺยนฺ”ติ, อยเมว อานาปานสฺสติสมาธิ สาธุกํ มนสิ กาตพฺโพ.\csromanspacer{} ตสฺมาติห ภิกฺขเว ภิกฺขุ เจปิ อากงฺเขยฺย “อปฺปฏิกูลญฺจ ปฏิกูลญฺจ ตทุภยํ อภินิวชฺเชตฺวา อุเปกฺขโก วิหเรยฺยํ สโต สมฺปชาโน”ติ, อยเมว อานาปานสฺสติสมาธิ สาธุกํ มนสิ กาตพฺโพ.}

\prose{ตสฺมาติห ภิกฺขเว ภิกฺขุ เจปิ อากงฺเขยฺย “วิวิจฺเจว กาเมหิ วิวิจฺจ อกุสเลหิ ธมฺเมหิ สวิตกฺกํ สวิจารํ วิเวกชํ ปีติสุขํ ปฐมํ ฌานํ อุปสมฺปชฺช วิหเรยฺยนฺ”ติ, อยเมว อานาปานสฺสติสมาธิ สาธุกํ มนสิ กาตพฺโพ.\csromanspacer{} ตสฺมาติห ภิกฺขเว ภิกฺขุ เจปิ อากงฺเขยฺย}

\csromanheader{2}{10. อานาปานสํยุตฺต}
\prose{\csromanfolio{277}“วิตกฺกวิจารานํ วูปสมา อชฺฌตฺตํ สมฺปสาทนํ เจตโส เอโกทิภาวํ อวิตกฺกํ อวิจารํ สมาธิชํ ปีติสุขํ ทุติยํ ฌานํ อุปสมฺปชฺช วิหเรยฺยนฺ”ติ, อยเมว อานาปานสฺสติสมาธิ สาธุกํ มนสิ กาตพฺโพ.\csromanspacer{} ตสฺมาติห ภิกฺขเว ภิกฺขุ เจปิ อากงฺเขยฺย “ปีติยา จ วิราคา อุเปกฺขโก จ วิหเรยฺยํ, สโต จ สมฺปชาโน สุขญฺจ กาเยน ปฏิสํเวเทยฺยํ, ยํ ตํ อริยา อาจิกฺกนฺติ ‘อุเปกฺขโก สติมา สุขวิหารี’ติ, ตติยํ ฌานํ อุปสมฺปชฺช วิหเรยฺยนฺ”ติ, อยเมว อานาปานสฺสติสมาธิ สาธุกํ มนสิ กาตพฺโพ.\csromanspacer{} ตสฺมาติห ภิกฺขเว ภิกฺขุ เจปิ อากงฺเขยฺย “สุขสฺส จ ปหานา ทุกฺขสฺส จ ปหานา ปุพฺเพว โสมนสฺสโทมนสฺสานํ อตฺถงฺคมา อทุกฺขมสุขํ อุเปกฺขาสติปาริสุทฺธิํ จตุตฺถํ ฌานํ อุปสมฺปชฺช วิหเรยฺยนฺ”ติ, อยเมว อานาปานสฺสติสมาธิ สาธุกํ มนสิ กาตพฺโพ.}

\prose{ตสฺมาติห ภิกฺขเว ภิกฺขุ เจปิ อากงฺเขยฺย “สพฺพโส รูปสญฺญานํ สมติกฺกมา ปฏิฆสญฺญานํ อตฺถงฺคมา นานตฺตสญฺญานํ อมนสิการา ‘อนนฺโต อากาโส’ติ อากาสานญฺจายตนํ อุปสมฺปชฺช วิหเรยฺยนฺ”ติ, อยเมว อานาปานสฺสติสมาธิ สาธุกํ มนสิ กาตพฺโพ.\csromanspacer{} ตสฺมาติห ภิกฺขเว ภิกฺขุ เจปิ อากงฺเขยฺย “สพฺพโส อากาสานญฺจายตนํ สมติกฺกมฺม ‘อนนฺตํ วิญฺญาณนฺ’ติ วิญฺญาณญฺจายตนํ อุปสมฺปชฺช วิหเรยฺยนฺ”ติ, อยเมว อานาปานสฺสติสมาธิ สาธุกํ มนสิ กาตพฺโพ.\csromanspacer{} ตสฺมาติห ภิกฺขเว ภิกฺขุ เจปิ อากงฺเขยฺย “สพฺพโส วิญฺญาณญฺจายตนํ สมติกฺกมฺม ‘นตฺถิ กิญฺจี’ติ อากิญฺจญฺญายตนํ อุปสมฺปชฺช วิหเรยฺยนฺ”ติ, อยเมว อานาปานสฺสติสมาธิ สาธุกํ มนสิ กาตพฺโพ.\csromanspacer{} ตสฺมาติห ภิกฺขเว ภิกฺขุ เจปิ อากงฺเขยฺย “สพฺพโส อากิญฺจญฺญายตนํ สมติกฺกมฺม เนวสญฺญานาสญฺญายตนํ อุปสมฺปชฺช วิหเรยฺยนฺ”ติ, อยเมว อานาปานสฺสติสมาธิ สาธุกํ มนสิ กาตพฺโพ.\csromanspacer{} ตสฺมาติห ภิกฺขเว ภิกฺขุ เจปิ อากงฺเขยฺย “สพฺพโส เนวสญฺญานาสญฺญายตนํ สมติกฺกมฺม สญฺญาเวทยิตนิโรธํ อุปสมฺปชฺช วิหเรยฺยนฺ”ติ, อยเมว อานาปานสฺสติสมาธิ สาธุกํ มนสิ กาตพฺโพ.}

\prose{เอวํ ภาวิเต โข ภิกฺขเว อานาปานสฺสติสมาธิมฺหิ เอวํ พหุลีกเต สุขํ เจ เวทนํ เวทยติ, สา “อนิจฺจา”ติ ปชานาติ, “อนชฺโฌสิตา”ติ ปชานาติ, “อนภินนฺทิตา”ติ ปชานาติ.\csromanspacer{} ทุกฺขํ เจ เวทนํ เวทยติ, สา “อนิจฺจา”ติ ปชานาติ, “อนชฺโฌสิตา”ติ ปชานาติ,}

\prose{\csromanfolio{278}“อนภินนฺทิตา”ติ ปชานาติ.\csromanspacer{} อทุกฺขมสุขํ เจ เวทนํ เวทยติ, สา “อนิจฺจา”ติ ปชานาติ, “อนชฺโฌสิตา”ติ ปชานาติ, “อนภินนฺทิตา”ติ ปชานาติ.}

\prose{สุขํ\footnote{โส สุขํ (สี, สฺยา, กํ, อิ) ม 3. 288 ปิฏฺเฐปิ. อฏฺฐกถาฏีกา โอโลเกตพฺพา.} เจ เวทนํ เวทยติ, วิสํยุตฺโต นํ เวทยติ.\csromanspacer{} ทุกฺขํ เจ เวทนํ เวทยติ, วิสํยุตฺโต นํ เวทยติ.\csromanspacer{} อทุกฺขมสุขํ เจ เวทนํ เวทยติ, วิสํยุตฺโต นํ เวทยติ.\csromanspacer{} โส กายปริยนฺติกํ เวทนํ เวทยมาโน “กายปริยนฺติกํ เวทนํ เวทยามี”ติ ปชานาติ, ชีวิตปริยนฺติกํ เวทนํ เวทยมาโน “ชีวิตปริยนฺติกํ เวทนํ เวทยามี”ติ ปชานาติ, กายสฺส เภทา อุทฺธํ ชีวิตปริยาทานา “อิเธว สพฺพเวทยิตานิ อนภินนฺทิตานิ สีตีภวิสฺสนฺตี”ติ ปชานาติ.}

\prose{เสยฺยถาปิ ภิกฺขเว เตลญฺจ ปฏิจฺจ วฏฺฏิญฺจ ปฏิจฺจ เตลปฺปทีโป ฌาเยยฺย, ตสฺเสว เตลสฺส จ วฏฺฏิยา จ ปริยาทานา อนาหาโร นิพฺพาเยยฺย.\csromanspacer{} เอวเมว โข ภิกฺขเว ภิกฺขุ กายปริยนฺติกํ เวทนํ เวทยมาโน “กายปริยนฺติกํ เวทนํ เวทยามี”ติ ปชานาติ, ชีวิตปริยนฺติกํ เวทนํ เวทยมาโน “ชีวิตปริยนฺติกํ เวทนํ เวทยามี”ติ ปชานาติ, กายสฺส เภทา อุทฺธํ ชีวิตปริยาทานา “อิเธว สพฺพเวทยิตานิ อนภินนฺทิตานิ สีตีภวิสฺสนฺตี”ติ ปชานาตีติ.\csromanspacer{}\csromanspacer{}อฏฺฐมํ.}

\csromanmatikaanchor{mtk.213}
\csromantocmarkref{subparagraph}{9. เวสาลีสุตฺต}{mtk.213}
\bookmark[dest={mtk.213},level=5]{9. เวสาลีสุตฺต}
\csromanheader{6}{9. เวสาลีสุตฺต}
\proseitem{985}{เอวํ เม สุตํ\csromandash{}เอกํ สมยํ ภควา เวสาลิยํ วิหรติ มหาวเน กูฏาคารสาลายํ.\csromanspacer{} เตน โข ปน สมเยน ภควา ภิกฺขูนํ อเนกปริยาเยน อสุภกถํ กเถติ, อสุภาย วณฺณํ ภาสติ, อสุภภาวนาย วณฺณํ ภาสติ.}

\prose{อถ โข ภควา ภิกฺขู อามนฺเตสิ “อิจฺฉามหํ ภิกฺขเว อฑฺฒมาสํ ปฏิสลฺลียิตุํ, นามฺหิ เกนจิ อุปสงฺกมิตพฺโพ อญฺญตฺร เอเกน ปิณฺฑปาตนีหารเกนา”ติ. “เอวํ ภนฺเต”ติ โข เต ภิกฺขู ภควโต ปฏิสฺสุตฺวา นาสฺสุธ โกจิ ภควนฺตํ อุปสงฺกมติ อญฺญตฺร เอเกน ปิณฺฑปาตนีหารเกน.}

\csromanheader{2}{10. อานาปานสํยุตฺต}
\prose{\csromanfolio{279}อถ โข เต ภิกฺขู “ภควา อเนกปริยาเยน อสุภกถํ กเถติ, อสุภาย วณฺณํ ภาสติ, อสุภภาวนาย วณฺณํ ภาสตี”ติ อเนกาการโวการํ อสุภภาวนานุโยคมนุยุตฺตา วิหรนฺติ.\csromanspacer{} เต อิมินา กาเยน อฏฺฏียมานา\footnote{อฏฺฏิยมานา (สี, สฺยา, กํ, อิ, ก)} หรายมานา ชิคุจฺฉมานา สตฺถหารกํ ปริเยสนฺติ, ทสปิ ภิกฺขู เอกาเหน สตฺถํ อาหรนฺติ, วีสมฺปิ ฯเปฯ ติํสมฺปิ ภิกฺขู เอกาเหน สตฺถํ อาหรนฺติ.}

\prose{อถ โข ภควา ตสฺส อฑฺฒมาสสฺส อจฺจเยน ปฏิสลฺลานา วุฏฺฐิโต อายสฺมนฺตํ อานนฺทํ อามนฺเตสิ “กิํ นุ โข อานนฺท ตนุภูโต วิย ภิกฺขุสํโฆ”ติ.\csromanspacer{} ตถา หิ ปน ภนฺเต “ภควา ภิกฺขูนํ อเนกปริยาเยน อสุภกตํ กเถติ, อสุภาย วณฺณํ ภาสติ, อสุภภาวนาย วณฺณํ ภาสตี”ติ อเนกาการโวการํ อสุภภาวนานุโยคมนุยุตฺตา วิหรนฺติ.\csromanspacer{} เต อิมินา กาเยน อฏฺฏียมานา หรายมานา ชิคุจฺฉมานา สตฺถหารกํ ปริเยสนฺติ, ทสปิ ภิกฺขู เอกาเหน สตฺถํ อาหรนฺติ, วีสมฺปิ ภิกฺขู.\csromanspacer{} ติํสมฺปิ ภิกฺขู เอกาเหน สตฺถํ อาหรนฺติ, สาธุ ภนฺเต ภควา อญฺญํ ปริยายํ อาจิกฺขตุ, ยถายํ ภิกฺขุสํโฆ อญฺญาย สณฺฐเหยฺยาติ.}

\prose{เตนหานนฺท ยาวติกา ภิกฺขู เวสาลิํ อุปนิสฺสาย วิหรนฺติ, เต สพฺเพ อุปฏฺฐานสาลายํ สนฺนิปาเตหีติ. “เอวํ ภนฺเต”ติ โข อายสฺมา อานนฺโท ภควโต ปฏิสฺสุตฺวา ยาวติกา ภิกฺขู เวสาลิํ อุปนิสฺสาย วิหรนฺติ, เต สพฺเพ อุปฏฺฐานสาลายํ สนฺนิปาเตตฺวา เยน ภควา เตนุปสงฺกมิ, อุปสงฺกมิตฺวา ภควนฺตํ เอตทโวจ “สนฺนิปติโต\footnote{สนฺนิปาติโต (สี)} ภนฺเต ภิกฺขุสํโฆ, ยสฺสทานิ ภนฺเต ภควา กาลํ มญฺญตี”ติ.}

\prose{อถ โข ภควา เยน อุปฏฺฐานสาลา เตนุปสงฺกมิ, อุปสงฺกมิตฺวา ปญฺญตฺเต อาสเน นิสีทิ, นิสฺสชฺช โข ภควา ภิกฺขู อามนฺเตสิ\csromandash{}อยมฺปิ โข ภิกฺขเว อานาปานสฺสติสมาธิ ภาวิโต พหุลีกโต สนฺโต เจว ปณีโต จ อเสจนโก จ สุโข จ วิหาโร, อุปฺปนฺนุปฺปนฺเน จ ปาปเก อกุสเล ธมฺเม ฐานโส อนฺตรธาเปติ วูปสเมติ.}

\prose{\csromanfolio{280}เสยฺยถาปิ ภิกฺขเว คิมฺหานํ ปจฺฉิเม มาเส อูหตํ รโชชลฺลํ, ตเมนํ มหา-อกาลเมโฆ ฐานโส อนฺตรธาเปติ วูปสเมติ.\csromanspacer{} เอวเมว โข ภิกฺขเว อานาปานสฺสติสมาธิ ภาวิโต พหุลีกโต สนฺโต เจว ปณีโต จ อเสจนโก จ สุโข จ วิหาโร, อุปฺปนฺนุปฺปนฺเน จ ปาปเก อกุสเล ธมฺเม ฐานโส อนฺตรธาเปติ วูปสเมติ.\csromanspacer{} กถํ ภาวิโต จ ภิกฺขเว อานาปานสฺสติสมาธิ กถํ พหุลีกโต สนฺโต เจว ปณีโต จ อเสจนโก จ สุโข จ วิหาโร, อุปฺปนฺนุปฺปนฺเน จ ปาปเก อกุสเล ธมฺเม ฐานโส อนฺตรธาเปติ วูปสเมติ.}

\prose{อิธ ภิกฺขเว ภิกฺขุ อรญฺญคโต วา รุกฺขมูลคโต วา สุญฺญาคารคโต วา นิสีทติ ปลฺลงฺกํ อาภุชิตฺวา อุชุํ กายํ ปณิธาย ปริมุขํ สติํ อุปฏฺฐเปตฺวา, โส สโตว อสฺสสติ, สโตว ปสฺสสติ ฯเปฯ “ปฏินิสฺสคฺคานุปสฺสี อสฺสสิสฺสามี”ติ สิกฺขติ, “ปฏินิสฺสคฺคานุปสฺสี ปสฺสสิสฺสามี”ติ สิกฺขติ.\csromanspacer{} เอวํ ภาวิโต โข ภิกฺขเว อานาปานสฺสติสมาธิ เอวํ พหุลีกโต สนฺโต เจว ปณีโต จ อเสจนโก จ สุโข จ วิหาโร, อุปฺปนฺนุปฺปนฺเน จ ปาปเก อกุสเล ธมฺเม ฐานโส อนฺตรธาเปติ วูปสเมตีติ.\csromanspacer{}\csromanspacer{}นวมํ.}

\csromanmatikaanchor{mtk.214}
\csromantocmarkref{subparagraph}{10. กิมิลสุตฺต}{mtk.214}
\bookmark[dest={mtk.214},level=5]{10. กิมิลสุตฺต}
\csromanheader{6}{10. กิมิลสุตฺต}
\proseitem{986}{เอวํ เม สุตํ\csromandash{}เอกํ สมยํ ภควา กิมิลายํ\footnote{กิมฺพิลายํ (สี, อิ)} วิหรติ เวฬุวเน.\csromanspacer{} ตตฺร โข ภควา อายสฺมนฺตํ กิมิลํ อามนฺเตสิ “กถํ ภาวิโต นุ โข กิมิล อานาปานสฺสติสมาธิ กถํ พหุลีกโต มหปฺผโล โหติ มหานิสํโส”ติ.}

\prose{เอวํ วุตฺเต อายสฺมา กิมิโล ตุณฺหี อโหสิ.\csromanspacer{} ทุติยมฺปิ โข ภควา ฯเปฯ.\csromanspacer{} ตติยมฺปิ โข ภควา อายสฺมนฺตํ กิมิลํ อามนฺเตสิ “กถํ ภาวิโต นุ โข กิมิล อานาปานสฺสติสมาธิ กถํ พหุลีกโต มหปฺผโล โหติ มหานิสํโส”ติ.\csromanspacer{} ตติยมฺปิ โข อายสฺมา กิมิโล ตุณฺหี อโหสิ.}

\prose{เอวํ วุตฺเต อายสฺมา อานนฺโท ภควนฺตํ เอตทโวจ “เอตสฺส ภควา กาโล, เอตสฺส สุคต กาโล, ยํ ภควา อานาปานสฺสติสมาธิํ}

\csromanheader{2}{10. อานาปานสํยุตฺต}
\prose{\csromanfolio{281}ภาเสยฺย, ภควโต สุตฺวา ภิกฺขู ธาเรสฺสนฺตีติ.\csromanspacer{} เตนหานนฺท สุณาหิ สาธุกํ มนสิ กโรหิ, ภาสิสฺสามีติ. “เอวํ ภนฺเต”ติ โข อายสฺมา อานนฺโท ภควโต ปจฺจสฺโสสิ.\csromanspacer{} ภควา เอตทโวจ\csromandash{}}

\prose{กถํ ภาวิโต จ อานนฺท อานาปานสฺสติสมาธิ กถํ พหุลีกโต มหปฺผโล โหติ มหานิสํโส, อิธานนฺท ภิกฺขุ อรญฺญคโต วา รุกฺขมูลคโต วา สุญฺญาคารคโต ว นิสีทติ ปลฺลงฺกํ อาภุชิตฺวา อุชุํ กายํ ปณิธาย ปริมุขํ สติํ อุปฏฺฐเปตฺวา, โส สโตว อสฺสสติ, สโตว ปสฺสสติ ฯเปฯ “ปฏินิสฺสคฺคานุปสฺสี อสฺสสิสฺสามี”ติ สิกฺขติ, “ปฏินิสฺสคฺคานุปสฺสี ปสฺสสิสฺสามี”ติ สิกฺขติ.\csromanspacer{} เอวํ ภาวิโต โข อานนฺท อานาปานสฺสติสมาธิ เอวํ พหุลีกโต มหปฺผโล โหติ มหานิสํโส.}

\prose{ยสฺมิํ สมเย อานนฺท ภิกฺขุ ทีฆํ วา อสฺสสนฺโต “ทีฆํ อสฺสสามี”ติ ปชานาติ, ทีฆํ วา ปสฺสสนฺโต “ทีฆํ ปสฺสสามี”ติ ปชานาติ, รสฺสํ วา อสฺสสนฺโต “รสฺสํ อสฺสสามี”ติ ปชานาติ, รสฺสํ วา ปสฺสสนฺโต “รสฺสํ ปสฺสสามี”ติ ปชานาติ, “สพฺพกายปฺปฏิสํเวที อสฺสสิสฺสามี”ติ สิกฺขติ, “สพฺพกายปฺปฏิสํเวที ปสฺสสิสฺสามี”ติ สิกฺขติ, “ปสฺสมฺภยํ กายสงฺขารํ อสฺสสิสฺสามี”ติ สิกฺขติ, “ปสฺสมฺภยํ กายสงฺขารํ ปสฺสสิสฺสามี”ติ สิกฺขติ, กาเย กายานุปสฺสี อานนฺท ภิกฺขุ ตสฺมิํ สมเย วิหรติ อาตาปี สมฺปชาโน สติมา วิเนยฺย โลเก อภิชฺฌาโทมนสฺสํ.\csromanspacer{} ตํ กิสฺส เหตุ, กายญฺญตราหํ อานนฺท เอตํ วทามิ ยทิทํ อสฺสาสปสฺสาสํ, ตสฺมาติหานนฺท กาเย กายานุปสฺสี ภิกฺขุ ตสฺมิํ สมเย วิหรติ อาตาปี สมฺปชาโน สติมา วิเนยฺย โลเก อภิชฺฌาโทมนสฺสํ. (1)}

\prose{ยสฺมิํ สมเย อานนฺท ภิกฺขุ “ปีติปฺปฏิสํเวที อสฺสสิสฺสามี”ติ สิกฺขติ, “ปีติปฺปฏิสํเวที ปสฺสสิสฺสามี”ติ สิกฺขติ, “สุขปฺปฏิสํเวที อสฺสสิสฺสามี”ติ สิกฺขติ, “สุขปฺปฏิสํเวที ปสฺสสิสฺสามี”ติ สิกฺขติ, “จิตฺตสงฺขารปฺปฏิสํเวที อสฺสสิสฺสามี”ติ สิกฺขติ, “จิตฺตสงฺขารปฺปฏิสํเวที ปสฺสสิสฺสามี”ติ สิกฺขติ, “ปสฺสมฺภยํ จิตฺตสงฺขารํ ปสฺสสิสฺสามี”ติ สิกฺขติ, เวทนาสุ เวทนานุปสฺสี อานนฺท ภิกฺขุ ตสฺมิํ สมเย วิหรติ อาตาปี}

\prose{\csromanfolio{282}สมฺปชาโน สติมา วิเนยฺย โลเก อภิชฺฌาโทมนสฺสํ.\csromanspacer{} ตํ กิสฺส เหตุ, เวทนาญฺญตราหํ อานนฺท เอตํ วทามิ ยทิทํ อสฺสาสปสฺสาสานํ\footnote{อสฺสาสปสฺสาสํ (อิ, ก) ม 3. 126 ปิฏฺเฐปิ.} สาธุกํ มนสิการํ.\csromanspacer{} ตสฺมาติหานนฺท เวทนาสุ เวทนานุปสฺสี ภิกฺขุ ตสฺมิํ สมเย วิหรติ อาตาปี สมฺปชาโน สติมา วิเนยฺย โลเก อภิชฺฌาโทมนสฺสํ. (2)}

\prose{ยสฺมิํ สมเย อานนฺท ภิกฺขุ “จิตฺตปฺปฏิสํเวที อสฺสสิสฺสามี”ติ สิกฺขติ, “จิตฺตปฺปฏิสํเวที ปสฺสสิสฺสามี”ติ สิกฺขติ, อภิปฺปโมทยํ จิตฺตํ ฯเปฯ สมาทหํ จิตฺตํ. “วิโมจยํ จิตฺตํ อสฺสสิสฺสามี”ติ สิกฺขติ, “วิโมจยํ จิตฺตํ ปสฺสสิสฺสามี”ติ สิกฺขติ.\csromanspacer{} จิตฺเต จิตฺตานุปสฺสี อานนฺท ภิกฺขุ ตสฺมิํ สมเย วิหรติ อาตาปี สมฺปชาโน สติมา วิเนยฺย โลเก อภิชฺฌาโทมนสฺสํ.\csromanspacer{} ตํ กิสฺส เหตุ, นาหํ อานนฺท มุฏฺฐสฺสติสฺส อสมฺปชานสฺส อานาปานสฺสติสมาธิภาวนํ วทามิ.\csromanspacer{} ตสฺมาติหานนฺท จิตฺเต จิตฺตานุปสฺสี ภิกฺขุ ตสฺมิํ สมเย วิหรติ อาตาปี สมฺปชาโน สติมา วิเนยฺย โลเก อภิชฺฌาโทมนสฺสํ. (3)}

\prose{ยสฺมิํ สมเย อานนฺท ภิกฺขุ “อนิจฺจานุปสฺสี อสฺสสิสฺสามี”ติ สิกฺขติ ฯเปฯ วิราคานุปสฺสี.\csromanspacer{} นิโรธานุปสฺสี. “ปฏินิสฺสคฺคานุปสฺสี อสฺสสิสฺสามี”ติ สิกฺขติ. “ปฏินิสฺสคฺคานุปสฺสี ปสฺสสิสฺสามี”ติ สิกฺขติ.\csromanspacer{} ธมฺเมสุ ธมฺมานุปสฺสี อานนฺท ภิกฺขุ ตสฺมิํ สมเย วิหรติ อาตาปี สมฺปชาโน สติมา วิเนยฺย โลเก อภิชฺฌาโทมนสฺสํ.\csromanspacer{} โส ยํ ตํ โหติ อภิชฺฌาโทมนสฺสานํ ปหานํ, ตํ ปญฺญาย ทิสฺวา สาธุกํ อชฺฌุเปกฺขิตา โหติ.\csromanspacer{} ตสฺมาติหานนฺท ธมฺเมสุ ธมฺมานุปสฺสี ภิกฺขุ ตสฺมิํ สมเย วิหรติ อาตาปี สมฺปชาโน สติมา วิเนยฺย โลเก อภิชฺฌาโทมนสฺสํ. (4)}

\prose{เสยฺยถาปิ อานนฺท จตุมหาปเถ\footnote{จาตุมฺมหาปเถ (สี, สฺยา, กํ)} มหาปํสุปุญฺโช ปุรตฺถิมาย เจปิ ทิสาย อาคจฺเฉยฺย สกฏํ วา รโถ วา, อุปหนเตว ตํ ปํสุปุญฺชํ.\csromanspacer{} ปจฺฉิมาย เจปิ ทิสาย อาคจฺเฉยฺย.\csromanspacer{} อุตฺตราย เจปิ ทิสาย.\csromanspacer{} ทกฺขิณาย เจปิ ทิสาย อาคจฺเฉยฺย สกฏํ วา รโถ วา, อุปหนเตว ตํ ปํสุปุญฺชํ.\csromanspacer{} เอวเมว โข อานนฺท ภิกฺขุ กาเย กายานุปสฺสี วิหรนฺโตปิ}

\csromanheader{2}{10. อานาปานสํยุตฺต}
\prose{\csromanfolio{283}อุปหนเตว ปาปเก อกุสเล ธมฺเม.\csromanspacer{} เวทนาสุ ฯเปฯ.\csromanspacer{} จิตฺเต ฯเปฯ.\csromanspacer{} ธมฺเมสุ ธมฺมานุปสฺสี วิหรนฺโตปิ อุปหนเตว ปาปเก อกุสเล ธมฺเมติ.\csromanspacer{}\csromanspacer{}ทสมํ.\csromanspacer{} เอกธมฺมวคฺโค ปฐโม.}

\titlehead{\textbf{ตสฺสุทฺทานํ}}
\csromangathameasure{เอกธมฺโม จ โพชฺฌงฺโค, สุทฺธิกญฺจ ทุเว ผลา. \\ อริฏฺโฐ กปฺปิโน ทีโป, เวสาลี กิมิเลน จาติ.}
\csromangathagroup{\csromangathastanza{เอกธมฺโม จ โพชฺฌงฺโค, สุทฺธิกญฺจ ทุเว ผลา. \\ อริฏฺโฐ กปฺปิโน ทีโป, เวสาลี กิมิเลน จาติ.}}

\chapterhead{2. ทุติยวคฺค}
\csromanmatikaanchor{mtk.215}
\csromanmatikaanchor{mtk.216}
\csromantocmarknopage{subparagraph}{2. ทุติยวคฺค}{mtk.215}
\bookmark[dest={mtk.215},level=5]{2. ทุติยวคฺค}
\csromantocmarkref{subparagraph}{1. อิจฺฉานงฺคลสุตฺต}{mtk.216}
\bookmark[dest={mtk.216},level=5]{1. อิจฺฉานงฺคลสุตฺต}
\csromanheader{6}{1. อิจฺฉานงฺคลสุตฺต}
\proseitem{987}{เอกํ สมยํ ภควา อิจฺฉานงฺคเล วิหรติ อิจฺฉานงฺคลวนสณฺเฑ.\csromanspacer{} ตตฺร โข ภควา ภิกฺขู อามนฺเตสิ “อิจฺฉามหํ ภิกฺขเว เตมาสํ ปฏิสลฺลียิตุํ, นามฺหิ เกนจิ อุปสงฺกมิตพฺโพ อญฺญตฺร เอเกน ปิณฺฑปาตนีหารเกนา”ติ. “เอวํ ภนฺเต”ติ โข เต ภิกฺขู ภควโต ปฏิสฺสุตฺวา นาสฺสุธ โกจิ ภควนฺตํ อุปสงฺกมติ อญฺญตฺร เอเกน ปิณฺฑปาตนีหารเกน.}

\prose{อถ โข ภควา ตสฺส เตมาสสฺส อจฺจเยน ปฏิสลฺลานา วุฏฺฐิโต ภิกฺขู อามนฺเตสิ\csromandash{}สเจ โข ภิกฺขเว อญฺญติตฺถิยา ปริพฺพาชกา เอวํ ปุจฺเฉยฺยุํ “กตเมนาวุโส วิหาเรน สมโณ โคตโม วสฺสาวาสํ พหุลํ วิหาสี”ติ, เอวํ ปุฏฺฐา ตุเมฺห ภิกฺขเว เตสํ อญฺญติตฺถิยานํ ปริพฺพาชกานํ เอวํ พฺยากเรยฺยาถ “อานาปานสฺสติสมาธินา โข อาวุโส ภควา วสฺสาวาสํ พหุลํ วิหาสี”ติ.\csromanspacer{} อิธาหํ ภิกฺขเว สโต อสฺสสามิ, สโต ปสฺสสามิ.\csromanspacer{} ทีฆํ อสฺสสนฺโต “ทีฆํ อสฺสสามี”ติ ปชานามิ, ทีฆํ ปสฺสสนฺโต “ทีฆํ ปสฺสสามี”ติ ปชานามิ, รสฺสํ อสฺสสนฺโต “รสฺสํ อสฺสสามี”ติ ปชานามิ, รสฺสํ ปสฺสสนฺโต “รสฺสํ ปสฺสสามี”ติ ปชานามิ, “สพฺพกายปฺปฏิสํเวที อสฺสสิสฺสามี”ติ ปชานามิ ฯเปฯ “ปฏินิสฺสคฺคานุปสฺสี อสฺสสิสฺสามี”ติ ปชานามิ, “ปฏินิสฺสคฺคานุปสฺสี ปสฺสสิสฺสามี”ติ ปชานามิ.}

\prose{\csromanfolio{284}ยํ หิ ตํ ภิกฺขเว สมฺมา วทมาโน วเทยฺย อริยวิหาโร อิติปิ, พฺรหฺมวิหาโร อิติปิ, ตถาคตวิหาโร อิติปิ.\csromanspacer{} อานาปานสฺสติสมาธิํ สมฺมา วทมาโน วเทยฺย อริยวิหาโร อิติปิ, พฺรหฺมวิหาโร อิติปิ, ตถาคตวิหาโร อิติปิ.\csromanspacer{} เย เต ภิกฺขเว ภิกฺขู เสขา อปฺปตฺตมานสา อนุตฺตรํ โยคกฺเขมํ ปตฺถยมานา วิหรนฺติ, เตสํ อานาปานสฺสติสมาธิ ภาวิโต พหุลีกโต อาสวานํ ขยาย สํวตฺตติ, เย จ โข เต ภิกฺขเว ภิกฺขู อรหนฺโต ขีณาสวา วุสิตวนฺโต กตกรณียา โอหิตภารา อนุปฺปตฺตสทตฺถา ปริกฺขีณภวสํโยชนา สมฺมทญฺญาวิมุตฺตา, เตสํ อานาปานสฺสติสมาธิ ภาวิโต พหุลีกโต ทิฏฺฐธมฺมสุขวิหาราย เจว สํวตฺตติ สติสมฺปชญฺญาย จ.}

\prose{ยํ หิ ตํ ภิกฺขเว สมฺมา วทมาโน วเทยฺย อริยวิหาโร อิติปิ, พฺรหฺมวิหาโร อิติปิ, ตถาคตวิหาโร อิติปิ.\csromanspacer{} อานาปานสฺสติสมาธิํ สมฺมา วทมาโน วเทยฺย อริยวิหาโร อิติปิ, พฺรหฺมวิหาโร อิติปิ, ตถาคตวิหาโร อิติปีติ.\csromanspacer{}\csromanspacer{}ปฐมํ.}

\csromanmatikaanchor{mtk.217}
\csromantocmarkref{subparagraph}{2. กงฺเขยฺยสุตฺต}{mtk.217}
\bookmark[dest={mtk.217},level=5]{2. กงฺเขยฺยสุตฺต}
\csromanheader{6}{2. กงฺเขยฺยสุตฺต}
\proseitem{988}{เอกํ สมยํ อายสฺมา โลมสกํภิโย\footnote{โลมสวงฺคิโส (สี, อิ)} สกฺเกสุ วิหรติ กปิลวตฺถุสฺมิํ นิโคฺรธาราเม.\csromanspacer{} อถ โข มาหานาโม สกฺโก เยนายสฺมา โลมสกํภิโย เตนุปสงฺกมิ, อุปสงฺกมิตฺวา อายสฺมนฺตํ โลมสกํภิยํ อภิวาเทตฺวา เอกมนฺตํ นิสีทิ, เอกมนฺตํ นิสินฺโน โข มหานาโม สกฺโก อายสฺมนฺตํ โลมสกํภิยํ เอตทโวจ “โส เอว นุ โข ภนฺเต เสโข วิหาโร โส ตถาคตวิหาโร, อุทาหุ อญฺโญว\footnote{อญฺโญ (สฺยา, กํ, อิ, ก)} เสโข วิหาโร อญฺโญ ตถาคตวิหาโร”ติ.}

\prose{น โข อาวุโส มหานาม เสฺวว เสโข วิหาโร, โส ตถาคตวิหาโร, อญฺโญ โข อาวุโส มหานาม เสโข วิหาโร, อญฺโญ ตถาคตวิหาโร.\csromanspacer{} เย เต อาวุโส มหานาม ภิกฺขู เสขา อปฺปตฺตมานสา อนุตฺตรํ โยคกฺเขมํ ปตฺถยมานา วิหรนฺติ, เต ปญฺจ นีวรเณ ปหาย วิหรนฺติ.\csromanspacer{} กตเม ปญฺจ, กามจฺฉนฺทนีวรณํ ปหาย วิหรนฺติ, พฺยาปาทนีวรณํ ฯเปฯ ถินมิทฺธนีวรณํ.\csromanspacer{} อุทฺธจฺจกุกฺกุจฺจนีวรณํ.\csromanspacer{} วิจิกิจฺฉานีวรณํ ปหาย วิหรนฺติ.}

\csromanheader{2}{10. อานาปานสํยุตฺต}
\prose{\csromanfolio{285}เยปิ เต อาวุโส มหานาม ภิกฺขู เสขา อปฺปตฺตมานสา อนุตฺตรํ โยคกฺเขมํ ปตฺถยมานา วิหรนฺติ, เต อิเม ปญฺจ นีวรเณ ปหาย วิหรนฺติ.}

\prose{เย จ โข เต อาวุโส มหานาม ภิกฺขู อรหนฺโต ขีณาสวา วุสิตวนฺโต กตกรณียา โอหิตภารา อนุปฺปตฺตสทตฺถา ปริกฺขีณภวสํโยชนา สมฺมทญฺญาวิมุตฺตา, เตสํ ปญฺจ นีวรณา ปหีนา อุจฺฉินฺนมูลา ตาลาวตฺถุกตา อนภาวํกตา\footnote{อนภาวกตา (สี, อิ)} อายติํ อนุปฺปาทธมฺมา.\csromanspacer{} กตเม ปญฺจ, กามจฺฉนฺทนีวรณํ ปหีนํ อุจฺฉินฺนมูลํ ตาลาวตฺถุกตํ อนภาวํกตํ อายติํ อนุปฺปาทธมฺมํ.\csromanspacer{} พฺยาปาทนีวรณํ ปหีนํ ฯเปฯ ถินมิทฺธนีวรณํ.\csromanspacer{} อุทฺธจฺจกุกฺกุจฺจนีวรณํ.\csromanspacer{} วิจิกิจฺฉานีวรณํ ปหีนํ อุจฺฉินฺนมูลํ กาลาวตฺถุกตํ อนภาวํกตํ อายติํ อนุปฺปาทธมฺมํ.}

\prose{เย เต อาวุโส มหานาม ภิกฺขู อรหนฺโต ขีณาสวา วุสิตวนฺโต กตกรณียา โอหิตภารา อนุปฺปตฺตสทตฺถา ปริกฺขีณภวสํโยชนา สมฺมทญฺญาวิมุตฺตา, เตสํ อิเม ปญฺจ นีวรณา ปหีนา อุจฺฉินฺนมูลา ตาลาวตฺถุกตา อนภาวํกตา อายติํ อนุปฺปาทธมฺมา.\csromanspacer{} ตทมินาเปตํ อาวุโส มหานาม ปริยาเยน เวทิตพฺพํ, ยถา อญฺโญว เสโข วิหาโร, อญฺโญ ตถาคตวิหาโร.}

\prose{เอกมิทํ อาวุโส มหานาม สมยํ ภควา อิจฺฉานงฺคเล วิหรติ อิจฺฉานงฺคลวนสณฺเฑ.\csromanspacer{} ตตฺร โข อาวุโส มหานาม ภควา ภิกฺขู อามนฺเตสิ “อิจฺฉามหํ ภิกฺขเว เตมาสํ ปฏิสลฺลียิตุํ, นามฺหิ เกนจิ อุปสงฺกมิตพฺโพ อญฺญตฺร เอเกน ปิณฺฑปาตนีหารเกนา”ติ. “เอวํ ภนฺเต”ติ โข อาวุโส มหานาม เต ภิกฺขู ภควโต ปฏิสฺสุตฺวา นาสฺสุธ โกจิ ภควนฺตํ อุปสงฺกมติ อญฺญตฺร เอเกน ปิณฺฑปาตนีหารเกน.}

\prose{อถ โข อาวุโส ภควา ตสฺส เตมาสสฺส อจฺจเยน ปฏิสลฺลานา วุฏฺฐิโต ภิกฺขู อามนฺเตสิ\csromandash{}สเจ โข ภิกฺขเว อญฺญติตฺถิยา ปริพฺพาชกา เอวํ ปุจฺเฉยฺยุํ “กตเมนาวุโส วิหาเรน สมโณ โคตโม วสฺสาวาสํ พหุลํ วิหาสี”ติ.\csromanspacer{} เอวํ ปุฏฺฐา ตุเมฺห ภิกฺขเว}

\csromanmatikaanchor{mtk.218}
\csromantocmarkref{subparagraph}{3-4. อานนฺทสุตฺตานิ}{mtk.218}
\bookmark[dest={mtk.218},level=5]{3-4. อานนฺทสุตฺตานิ}
\prose{\csromanfolio{286}เตสํ อญฺญติตฺถิยานํ ปริพฺพาชกานํ เอวํ พฺยากเรยฺยาถ “อานาปานสฺสติสมาธินา โข อาวุโส ภควา วสฺสาวาสํ พหุลํ วิหาสี”ติ.\csromanspacer{} อิธาหํ ภิกฺขเว สโต อสฺสสามิ, สโต ปสฺสสามิ, ทีฆํ อสฺสสนฺโต “ทีฆํ อสฺสสามี”ติ ปชานามิ, ทีฆํ ปสฺสสนฺโต “ทีฆํ ปสฺสสามี”ติ ปชานามิ ฯเปฯ “ปฏินิสฺสคฺคานุปสฺสี อสฺสสิสฺสามี”ติ ปชานามิ, “ปฏินิสฺสคฺคานุปสฺสี ปสฺสสิสฺสามี”ติ ปชานามิ.}

\prose{ยํ หิ ตํ ภิกฺขเว สมฺมา วทมาโน วเทยฺย อริยวิหาโร อิติปิ, พฺรหฺมวิหาโร อิติปิ, ตถาคตวิหาโร อิติปิ.\csromanspacer{} อานาปานสฺสติสมาธิํ สมฺมา วทมาโน วเทยฺย อริยวิหาโร อิติปิ, พฺรหฺมวิหาโร อิติปิ, ตถาคตวิหาโร อิติปิ.}

\prose{เย เต ภิกฺขเว ภิกฺขู เสขา อปฺปตฺตมานสา อนุตฺตรํ โยคกฺเขมํ ปตฺถยมานา วิหรนฺติ, เตสํ อานาปานสฺสติสมาธิ ภาวิโต พหุลีกโต อาสวานํ ขยาย สํวตฺตติ.}

\prose{เย จ โข เต ภิกฺขเว ภิกฺขู อรหนฺโต ขีณาสวา วุสิตวนฺโต กตกรณียา โอหิตภารา อนุปฺปตฺตสทตฺถา ปริกฺขีณภวสํโยชนา สมฺมทญฺญาวิมุตฺตา, เตสํ อานาปานสฺสติสมาธิ ภาวิโต พหุลีกโต ทิฏฺเฐว ธมฺเม สุขวิหาราย เจว สํวตฺตติ สติสมฺปชญฺญาย จ.}

\prose{ยํ หิ ตํ ภิกฺขเว สมฺมา วทมาโน วเทยฺย อริยวิหาโร อิติปิ, พฺรหฺมวิหาโร อิติปิ, ตถาคตวิหาโร อิติปิ.\csromanspacer{} อานาปานสฺสติสมาธิํ สมฺมา วทมาโน วเทยฺย อริยวิหาโร อิติปิ, พฺรหฺมวิหาโร อิติปิ, ตถาคตวิหาโร อิติปีติ.\csromanspacer{} อิมินา โข เอตํ อาวุโส มหานาม ปริยาเยน เวทิตพฺพํ, ยถา อญฺโญว เสโข วิหาโร, อญฺโญ ตถาคตวิหาโรติ.}

\csromanheader{2}{3. ปฐม-อานนฺทสุตฺต}
\proseitem{989}{สาวตฺถินิทานํ.\csromanspacer{} อถ โข อายสฺมา อานนฺโท เยน ภควา เตนุปสงฺกมิ, อุปสงฺกมิตฺวา ภควนฺตํ อภิวาเทตฺวา เอกมนฺตํ นิสีทิ, เอกมนฺตํ นิสินฺโน โข อายสฺมา อานนฺโท ภควนฺตํ เอตทโวจ “อตฺถิ นุ โข}

\csromanheader{2}{10. อานาปานสํยุตฺต}
\prose{\csromanfolio{287}ภนฺเต เอกธมฺโม\footnote{เอโก ธมฺโม (สี)} ภาวิโต พหุลีกโต จตฺตาโร ธมฺเม ปริปูเรติ, จตฺตาโร ธมฺมา ภาวิตา พหุลีกตา สตฺต ธมฺเม ปริปูเรนฺติ, สตฺต ธมฺมา ภาวิตา พหุลีกตา เทฺว ธมฺเม ปริปูเรนฺตี”ติ.}

\prose{อตฺถิ โข อานนฺท เอกธมฺโม ภาวิโต พหุลีกโต จตฺตาโร ธมฺเม ปริปูเรติ, จตฺตาโร ธมฺมา ภาวิตา พหุลีกตา สตฺต ธมฺเม ปริปูเรนฺติ.\csromanspacer{} สตฺต ธมฺมา ภาวิตา พหุลีกตา เทฺว ธมฺเม ปริปูเรนฺตีติ.}

\prose{กตโม ปน ภนฺเต เอกธมฺโม1 ภาวิโต พหุลีกโต จตฺตาโร ธมฺเม ปริปูเรติ, จตฺตาโร ธมฺมา ภาวิตา พหุลีกตา สตฺต ธมฺเม ปริปูเรนฺติ, สตฺต ธมฺมา ภาวิตา พหุลีกตา เทฺว ธมฺเม ปริปูเรนฺตีติ.\csromanspacer{} อานาปานสฺสติสมาธิ โข อานนฺท เอกธมฺโม ภาวิโต พหุลีกโต จตฺตาโร สติปฏฺฐาเน ปริปูเรติ, จตฺตาโร สติปฏฺฐานา ภาวิตา พหุลีกตา สตฺต โพชฺฌงฺเค ปริปูเรนฺติ, สตฺต โพชฺฌงฺคา ภาวิตา พหุลีกตา วิชฺชาวิมุตฺติํ ปริปูเรนฺติ.}

\prose{กถํ ภาวิโต อานนฺท อานาปานสฺสติสมาธิ กถํ พหุลีกโต จตฺตาโร สติปฏฺฐาเน ปริปูเรติ.\csromanspacer{} อิธานนฺท ภิกฺขุ อรญฺญคโต วา รุกฺขมูลคโต วา สุญฺญาคารคโต วา นิสีทติ ปลฺลงฺกํ อาภุชิตฺวา อุชุํ กายํ ปณิธาย ปริมุขํ สติํ อุปฏฺฐเปตฺวา, โส สโตว อสฺสสติ, สโตว ปสฺสสติ, ทีฆํ วา อสฺสสนฺโต “ทีฆํ อสฺสสามี”ติ ปชานาติ, ทีฆํ วา ปสฺสสนฺโต “ทีฆํ ปสฺสสามี”ติ ปชานาติ ฯเปฯ “ปฏินิสฺสคฺคานุปสฺสี อสฺสสิสฺสามี”ติ สิกฺขติ, “ปฏินิสฺสคฺคานุปสฺสี ปสฺสสิสฺสามี”ติ สิกฺขติ, ยสฺมิํ สมเย อานนฺท ภิกฺขุ ทีฆํ วา อสฺสสนฺโต “ทีฆํ อสฺสสามี”ติ ปชานาติ, ทีฆํ วา ปสฺสสนฺโต “ทีฆํ ปสฺสสามี”ติ ปชานาติ, รสฺสํ วา ฯเปฯ “ปสฺสมฺภยํ กายสงฺขารํ อสฺสสิสฺสามี”ติ สิกฺขติ, “ปสฺสมฺภยํ กายสงฺขารํ ปสฺสสิสฺสามี”ติ สิกฺขติ.\csromanspacer{} กาเย กายานุปสฺสี อานนฺท ภิกฺขุ ตสฺมิํ สมเย วิหรติ อาตาปี สมฺปชาโน สติมา วิเนยฺย โลเก อภิชฺฌาโทมนสฺสํ.\csromanspacer{} ตํ กิสฺส เหตุ, กายญฺญตราหํ อานนฺท เอตํ วทามิ ยทิทํ อสฺสาสปสฺสาสํ.\csromanspacer{} ตสฺมาติหานนฺท กาเย กายานุปสฺสี ภิกฺขุ ตสฺมิํ สมเย วิหรติ อาตาปี สมฺปชาโน สติมา วิเนยฺย โลเก อภิชฺฌาโทมนสฺสํ. (1)}

\prose{\csromanfolio{288}ยสฺมิํ สมเย อานนฺท ภิกฺขุ “ปีติปฺปฏิสํเวที อสฺสสิสฺสามี”ติ สิกฺขติ, สุขปฺปฏิสํเวที ฯเปฯ จิตฺตสงฺขารปฺปฏิสํเวที. “ปสฺสมฺภยํ จิตฺตสงฺขารํ อสฺสสิสฺสามี”ติ สิกฺขติ, “ปสฺสมฺภยํ จิตฺตสงฺขารํ ปสฺสสิสฺสามี”ติ สิกฺขติ.\csromanspacer{} เวทนาสุ เวทนานุปสฺสี อานนฺท ภิกฺขุ ตสฺมิํ สมเย วิหรติ อาตาปี สมฺปชาโน สติมา วิเนยฺย โลเก อภิชฺฌาโทมนสฺสํ.\csromanspacer{} ตํ กิสฺส เหตุ, เวทนาญฺญตราหํ อานนฺท เอตํ วทามิ, ยทิทํ อสฺสาสปสฺสาสานํ สาธุกํ มนสิการํ.\csromanspacer{} ตสฺมาติหานนฺท เวทนาสุ เวทนานุปสฺสี ภิกฺขุ ตสฺมิํ สมเย วิหรติ อาตาปี สมฺปชาโน สติมา วิเนยฺย โลเก อภิชฺฌาโทมนสฺสํ. (2)}

\prose{ยสฺมิํ สมเย อานนฺท ภิกฺขุ “จิตฺตปฺปฏิสํเวที อสฺสสิสฺสามี”ติ สิกฺขติ, “จิตฺตปฺปฏิสํเวที ปสฺสสิสฺสามี”ติ สิกฺขติ, อภิปฺปโมทยํ จิตฺตํ ฯเปฯ สมาทหํ จิตฺตํ. “วิโมจยํ จิตฺตํ อสฺสสิสฺสามี”ติ สิกฺขติ, “วิโมจยํ จิตฺตํ ปสฺสสิสฺสามี”ติ สิกฺขติ.\csromanspacer{} จิตฺเต จิตฺตานุปสฺสี อานนฺท ภิกฺขุ ตสฺมิํ สมเย วิหรติ อาตาปี สมฺปชาโน สติมา วิเนยฺย โลเก อภิชฺฌาโทมนสฺสํ.\csromanspacer{} ตํ กิสฺส เหตุ, นาหํ อานนฺท มุฏฺฐสฺสติสฺส อสมฺปชานสฺส อานาปานสฺสติสมาธิภาวนํ วทามิ.\csromanspacer{} ตสฺมา ติหานนฺท จิตฺเต จิตฺตานุปสฺสี ภิกฺขุ ตสฺมิํ สมเย วิหรติ อาตาปี สมฺปชาโน สติมา วิเนยฺย โลเก อภิชฺฌาโทมนสฺสํ. (3)}

\prose{ยสฺมิํ สมเย อานนฺท ภิกฺขุ อนิจฺจานุปสฺสี ฯเปฯ วิราคานุปสฺสี.\csromanspacer{} นิโรธานุปสฺสี. “ปฏินิสฺสคฺคานุปสฺสี อสฺสสิสฺสามี”ติ สิกฺขติ, “ปฏินิสฺสคฺคานุปสฺสี ปสฺสสิสฺสามี”ติ สิกฺขติ.\csromanspacer{} ธมฺเมสุ ธมฺมานุปสฺสี อานนฺท ภิกฺขุ ตสฺมิํ สมเย วิหรติ อาตาปี สมฺปชาโน สติมา วิเนยฺย โลเก อภิชฺฌาโทมนสฺสํ.\csromanspacer{} โส ยํ ตํ โหติ อภิชฺฌาโทมนสฺสานํ ปหานํ, ตํ ปญฺญาย ทิสฺวา สาธุกํ อชฺฌุเปกฺขิตา โหติ.\csromanspacer{} ตสฺมาติหานนฺท ธมฺเมสุ ธมฺมานุปสฺสี ภิกฺขุ ตสฺมิํ สมเย วิหรติ อาตาปี สมฺปชาโน สติมา วิเนยฺย โลเก อภิชฺฌาโทมนสฺสํ. (4)}

\prose{เอวํ ภาวิโต โข อานนฺท อานาปานสฺสติสมาธิ เอวํ พหุลีกโต จตฺตาโร สติปฏฺฐาเน ปริปูเรติ.}

\prose{กถํ ภาวิตา จานนฺท จตฺตาโร สติปฏฺฐานา กถํ พหุลีกตา สตฺต โพชฺฌงฺเค ปริปูเรนฺติ, ยสฺมิํ สมเย อานนฺท ภิกฺขุ กาเย}

\csromanheader{2}{10. อานาปานสํยุตฺต}
\prose{\csromanfolio{289}กายานุปสฺสี วิหรติ, อุปฏฺฐิตาสฺส\footnote{อุปฏฺฐิตสฺสติ (อิ, ก)} ตสฺมิํ สมเย ภิกฺขุโน\footnote{ตสฺมิํ สมเย อานนฺท ภิกฺขุโน (อิ, ก), ตสฺมิํ สมเย (ม 3. 127 ปิฏฺเฐ.)} สติ โหติ อสมฺมุฏฺฐา.\csromanspacer{} ยสฺมิํ สมเย อานนฺท ภิกฺขุโน อุปฏฺฐิตา สติ โหติ อสมฺมุฏฺฐา, สติสมฺโพชฺฌงฺโค ตสฺมิํ สมเย ภิกฺขุโน อารทฺโธ โหติ, สติสมฺโพชฺฌงฺคํ ตสฺมิํ สมเย ภิกฺขุ ภาเวติ, สติสมฺโพชฺฌงฺโค ตสฺมิํ สมเย ภิกฺขุโน ภาวนาปาริปูริํ คจฺฉติ. (1)}

\prose{โส ตถาสโต วิหรนฺโต ตํ ธมฺมํ ปญฺญาย ปวิจินติ ปวิจรติ ปริวีมํสมาปชฺชติ.\csromanspacer{} ยสฺมิํ สมเย อานนฺท ภิกฺขุ ตถาสโต วิหรนฺโต ตํ ธมฺมํ ปญฺญาย ปวิจินติ ปวิจรติ ปริวีมํสมาปชฺชติ.\csromanspacer{} ธมฺมวิจยสมฺโพชฺฌงฺโค ตสฺมิํ สมเย ภิกฺขุโน อารทฺโธ โหติ, ธมฺมวิจยสมฺโพชฺฌงฺคํ ตสฺมิํ สมเย ภิกฺขุ ภาเวติ, ธมฺมวิจยสมฺโพชฺฌงฺโค ตสฺมิํ สมเย ภิกฺขุโน ภาวนาปาริปูริํ คจฺฉติ. (2)}

\prose{ตสฺส ตํ ธมฺมํ ปญฺญาย ปวิจินโต ปวิจรโต ปริวีมํสมาปชฺชโต อารทฺธํ โหติ วีริยํ อสลฺลีนํ.\csromanspacer{} ยสฺมิํ สมเย อานนฺท ภิกฺขุโน ตํ ธมฺมํ ปญฺญาย ปวิจินโต ปวิจรโต ปริวีมํสมาปชฺชโต อารทฺธํ โหติ วีริยํ อสลฺลีนํ, วีริยสมฺโพชฺฌงฺโค ตสฺมิํ สมเย ภิกฺขุโน อารทฺโธ โหติ, วีริยสมฺโพชฺฌงฺคํ ตสฺมิํ สมเย ภิกฺขุ ภาเวติ, วีริยสมฺโพชฺฌงฺโค ตสฺมิํ สมเย ภิกฺขุโน ภาวนาปาริปูริํ คจฺฉติ. (3)}

\prose{อารทฺธวีริยสฺส อุปฺปชฺชติ ปีติ นิรามิสา.\csromanspacer{} ยสฺมิํ สมเย อานนฺท ภิกฺขุโน อารทฺธวีริยสฺส อุปฺปชฺชติ ปีติ นิรามิสา, ปีติสมฺโพชฺฌงฺโค ตสฺมิํ สมเย ภิกฺขุโน อารทฺโธ โหติ, ปีติสมฺโพชฺฌงฺคํ ตสฺมิํ สมเย ภิกฺขุ ภาเวติ, ปีติสมฺโพชฺฌงฺโค ตสฺมิํ สมเย ภิกฺขุโน ภาวนาปาริปูริํ คจฺฉติ. (4)}

\prose{ปีติมนสฺส กาโยปิ ปสฺสมฺภติ, จิตฺตมฺปิ ปสฺสมฺภติ.\csromanspacer{} ยสฺมิํ สมเย อานนฺท ภิกฺขุโน ปีติมนสฺส กาโยปิ ปสฺสมฺภติ, จิตฺตมฺปิ ปสฺสมฺภติ, ปสฺสทฺธิสมฺโพชฺฌงฺโค ตสฺมิํ สมเย ภิกฺขุโน อารทฺโธ โหติ, ปสฺสทฺธิสมฺโพชฺฌงฺคํ ตสฺมิํ สมเย ภิกฺขุ ภาเวติ, ปสฺสทฺธิสมฺโพชฺฌงฺโค ตสฺมิํ สมเย ภิกฺขุโน ภาวนาปาริปูริํ คจฺฉติ. (5)}

\prose{\csromanfolio{290}ปสฺสทฺธกายสฺส สุขิโน จิตฺตํ สมาธิยติ.\csromanspacer{} ยสฺมิํ สมเย อานนฺท ภิกฺขุโน ปสฺสทฺธกายสฺส สุขิโน จิตฺตํ สมาธิยติ, สมาธิสมฺโพชฺฌงฺโค ตสฺมิํ สมเย ภิกฺขุโน อารทฺโธ โหติ, สมาธิสมฺโพชฺฌงฺคํ ตสฺมิํ สมเย ภิกฺขุ ภาเวติ, สมาธิสมฺโพชฺฌงฺโค ตสฺมิํ สมเย ภิกฺขุโน ภาวนาปาริปูริํ คจฺฉติ. (6)}

\prose{โส ตถาสมาหิตํ จิตฺตํ สาธุกํ อชฺฌุเปกฺขิตา โหติ.\csromanspacer{} ยสฺมิํ สมเย อานนฺท ภิกฺขุ ตถาสมาหิตํ จิตฺตํ สาธุกํ อชฺฌุเปกฺขิตา โหติ, อุเปกฺขาสมฺโพชฺฌงฺโค ตสฺมิํ สมเย ภิกฺขุโน อารทฺโธ โหติ, อุเปกฺขาสมฺโพชฺฌงฺคํ ตสฺมิํ สมเย ภิกฺขุ ภาเวติ, อุเปกฺขาสมฺโพชฺฌงฺโค ตสฺมิํ สมเย ภิกฺขุโน ภาวนาปาริปูริํ คจฺฉติ. (7)}

\prose{ยสฺมิํ สมเย อานนฺท ภิกฺขุ เวทนาสุ ฯเปฯ.\csromanspacer{} จิตฺเต ฯเปฯ.\csromanspacer{} ธมฺเมสุ ธมฺมานุปสฺสี วิหรติ, อุปฏฺฐิตาสฺส ตสฺมิํ สมเย ภิกฺขุโน สติ โหติ อสมฺมุฏฺฐา.\csromanspacer{} ยสฺมิํ สมเย อานนฺท ภิกฺขุโน อุปฏฺฐิตา สติ โหติ อสมฺมุฏฺฐา, สติสมฺโพชฺฌงฺโค ตสฺมิํ สมเย ภิกฺขุโน อารทฺโธ โหติ, สติสมฺโพชฺฌงฺคํ ตสฺมิํ สมเย ภิกฺขุ ภาเวติ, สติสมฺโพชฺฌงฺโค ตสฺมิํ สมเย ภิกฺขุโน ภาวนาปาริปูริํ คจฺฉติ. (ยถา ปฐมํ สติปฏฺฐานํ, เอวํ วิตฺถาเรตพฺพํ.)}

\prose{โส ตถาสมาหิตํ จิตฺตํ สาธุกํ อชฺฌุเปกฺขิตา โหติ, ยสฺมิํ สมเย อานนฺท ภิกฺขุ ตถาสมาหิตํ จิตฺตํ สาธุกํ อชฺฌุเปกฺขิตา โหติ, อุเปกฺขาสมฺโพชฺฌงฺโค ตสฺมิํ สมเย ภิกฺขุโน อารทฺโธ โหติ, อุเปกฺขาสมฺโพชฺฌงฺคํ ตสฺมิํ สมเย ภิกฺขุ ภาเวติ, อุเปกฺขาสมฺโพชฺฌงฺโค ตสฺมิํ สมเย ภิกฺขุโน ภาวนาปาริปูริํ คจฺฉติ.\csromanspacer{} เอวํ ภาวิตา โข อานนฺท จตฺตาโร สติปฏฺฐานา เอวํ พหุลีกตา สตฺต โพชฺฌงฺเค ปริปูเรนฺติ.}

\prose{กถํ ภาวิตา อานนฺท สตฺต โพชฺฌงฺคา กถํ พหุลีกตา วิชฺชาวิมุตฺติํ ปริปูเรนฺติ.\csromanspacer{} อิธานนฺท ภิกฺขุ สติสมฺโพชฺฌงฺคํ ภาเวติ วิเวกนิสฺสิตํ วิราคนิสฺสิตํ นิโรธนิสฺสิตํ โวสฺสคฺคปริณามิํ.\csromanspacer{} ธมฺมวิจยสมฺโพชฺฌงฺคํ}

\csromanheader{2}{10. อานาปานสํยุตฺต}
\csromanmatikaanchor{mtk.219}
\csromantocmarkref{subparagraph}{5-6. ภิกฺขุสุตฺตานิ}{mtk.219}
\bookmark[dest={mtk.219},level=5]{5-6. ภิกฺขุสุตฺตานิ}
\prose{\csromanfolio{291}ภาเวติ ฯเปฯ.\csromanspacer{} อุเปกฺขาสมฺโพชฺฌงฺคํ ภาเวติ วิเวกนิสฺสิตํ วิราคนิสฺสิตํ นิโรธนิสฺสิตํ โวสฺสคฺคปริณามิํ.\csromanspacer{} เอวํ ภาวิตา โข อานนฺท สตฺต โพชฺฌงฺคา เอวํ พหุลีกตา วิชฺชาวิมุตฺติํ ปริปูเรนฺตีติ.\csromanspacer{}\csromanspacer{}ตติยํ.}

\csromanheader{2}{4. ทุติย-อานนฺทสุตฺต}
\proseitem{990}{อถ โข อายสฺมา อานนฺโท เยน ภควา เตนุปสงฺกมิ, อุปสงฺกมิตฺวา ภควนฺตํ อภิวาเทตฺวา เอกมนฺตํ นิสีทิ, เอกมนฺตํ นิสินฺนํ โข อายสฺมนฺตํ อานนฺทํ ภควา เอตทโวจ “อตฺถิ นุ โข อานนฺท เอกธมฺโม ภาวิโต พหุลีกโต จตฺตาโร ธมฺเม ปริปูเรติ, จตฺตาโร ธมฺมา ภาวิตา พหุลีกตา สตฺต ธมฺเม ปริปูเรนฺติ, สตฺต ธมฺมา ภาวิตา พหุลีกตา เทฺว ธมฺเม ปริปูเรนฺตี”ติ.\csromanspacer{} ภควํมูลกา โน ภนฺเต ธมฺมา ฯเปฯ.\csromanspacer{} อตฺถานนฺท เอกธมฺโม ภาวิโต พหุลีกโต จตฺตาโร ธมฺเม ปริปูเรติ, จตฺตาโร ธมฺมา ภาวิตา พหุลีกตา สตฺต ธมฺเม ปริปูเรนฺติ, สตฺต ธมฺมา ภาวิตา พหุลีกตา เทฺว ธมฺเม ปริปูเรนฺติ.}

\prose{กตโม จานนฺท เอกธมฺโม ภาวิโต พหุลีกโต จตฺตาโร ธมฺเม ปริปูเรติ, จตฺตาโร ธมฺมา ภาวิตา พหุลีกตา สตฺต ธมฺเม ปริปูเรนฺติ, สตฺต ธมฺมา ภาวิตา พหุลีกตา เทฺว ธมฺเม ปริปูเรนฺติ.\csromanspacer{} อานาปานสฺสติสมาธิ อานนฺท เอกธมฺโม ภาวิโต พหุลีกโต จตฺตาโร สติปฏฺฐาเน ปริปูเรติ, จตฺตาโร สติปฏฺฐานา ภาวิตา พหุลีกตา สตฺต โพชฺฌงฺเค ปริปูเรนฺติ, สตฺต โพชฺฌงฺคา ภาวิตา พหุลีกตา วิชฺชาวิมุตฺติํ ปริปูเรนฺตีติ.\csromanspacer{} กถํ ภาวิโต จานนฺท อานาปานสฺสติสมาธิ กถํ พหุลีกโต จตฺตาโร สติปฏฺฐาเน ปริปูเรติ, อิธานนฺท ภิกฺขุ อรญฺญคโต วา ฯเปฯ.\csromanspacer{} เอวํ ภาวิตา โข อานนฺท สตฺต โพชฺฌงฺคา เอวํ พหุลีกตา วิชฺชาวิมุตฺติํ ปริปูเรนฺตีติ.\csromanspacer{}\csromanspacer{}จตุตฺถํ.}

\csromanheader{2}{5. ปฐมภิกฺขุสุตฺต}
\proseitem{991}{อถ โข สมฺพหุลา ภิกฺขู เยน ภควา เตนุปสงฺกมิํสุ, อุปสงฺกมิตฺวา ภควนฺตํ อภิวาเทตฺวา เอกมนฺตํ นิสีทิํสุ, เอกมนฺตํ นิสินฺนา โข เต ภิกฺขู ภควนฺตํ เอตทโวจุํ “อตฺถิ นุ โข ภนฺเต เอกธมฺโม ภาวิโต พหุลีกโต จตฺตาโร ธมฺเม ปริปูเรติ, จตฺตาโร ธมฺมา ภาวิตา พหุลีกตา สตฺต ธมฺเม ปริปูเรนฺติ, สตฺต ธมฺมา \csromanfolio{292}ภาวิตา พหุลีกตา เทฺว ธมฺเม ปริปูเรนฺตี”ติ.\csromanspacer{} อตฺถิ โข ภิกฺขเว เอกธมฺโม ภาวิโต พหุลีกโต จตฺตาโร ธมฺเม ปริปูเรติ, จตฺตาโร ธมฺมา ภาวิตา พหุลีกตา สตฺต ธมฺเม ปริปูเรนฺติ, สตฺต ธมฺมา ภาวิตา พหุลีกตา เทฺว ธมฺเม ปริปูเรนฺตีติ.}

\prose{กตโม ปน ภนฺเต เอกธมฺโม ภาวิโต พหุลีกโต จตฺตาโร ธมฺเม ปริปูเรติ, จตฺตาโร ธมฺมา ภาวิตา พหุลีกตา สตฺต ธมฺเม ปริปูเรนฺติ, สตฺต ธมฺมา ภาวิตา พหุลีกตา เทฺว ธมฺเม ปริปูเรนฺตีติ.\csromanspacer{} อานาปานสฺสติสมาธิ โข ภิกฺขเว เอกธมฺโม ภาวิโต พหุลีกโต จตฺตาโร สติปฏฺฐาเน ปริปูเรติ, จตฺตาโร สติปฏฺฐานา ภาวิตา พหุลีกตา สตฺต โพชฺฌงฺเค ปริปูเรนฺติ, สตฺต โพชฺฌงฺคา ภาวิตา พหุลีกตา วิชฺชาวิมุตฺติํ ปริปูเรนฺตีติ.}

\prose{กถํ ภาวิโต จ ภิกฺขเว อานาปานสฺสติสมาธิ กถํ พหุลีกโต จตฺตาโร สติปฏฺฐาเน ปริปูเรติ.\csromanspacer{} อิธ ภิกฺขเว ภิกฺขุ อรญฺญคโต วา ฯเปฯ.\csromanspacer{} เอวํ ภาวิตา โข ภิกฺขเว สตฺต โพชฺฌงฺคา เอวํ พหุลีกตา วิชฺชาวิมุตฺติํ ปริปูเรนฺตีติ.\csromanspacer{}\csromanspacer{}ปญฺจมํ.}

\csromanheader{2}{6. ทุติยภิกฺขุสุตฺต}
\proseitem{992}{อถ โข สมฺพหุลา ภิกฺขู เยน ภควา เตนุปสงฺกมิํสุ, อุปสงฺกมิตฺวา ภควนฺตํ อภิวาเทตฺวา เอกมนฺตํ นิสีทิํสุ, เอกมนฺตํ นิสินฺเน โข เต ภิกฺขู ภควา เอตทโวจ “อตฺถิ นุ โข ภิกฺขเว เอกธมฺโม ภาวิโต พหุลีกโต จตฺตาโร ธมฺเม ปริปูเรติ, จตฺตาโร ธมฺมา ภาวิตา พหุลีกตา สตฺต ธมฺเม ปริปูเรนฺติ, สตฺต ธมฺมา ภาวิตา พหุลีกตา เทฺว ธมฺเม ปริปูเรนฺตี”ติ.\csromanspacer{} ภควํมูลกา โน ภนฺเต ธมฺมา ฯเปฯ ภควโต สุตฺวา ภิกฺขู ธาเรสฺสนฺตีติ.\csromanspacer{} อตฺถิ ภิกฺขเว เอกธมฺโม ภาวิโต พหุลีกโต จตฺตาโร ธมฺเม ปริปูเรติ, จตฺตาโร ธมฺมา ภาวิตา พหุลีกตา สตฺต ธมฺเม ปริปูเรนฺติ, สตฺต ธมฺมา ภาวิตา พหุลีกตา เทฺว ธมฺเม ปริปูเรนฺติ.}

\prose{กตโม จ ภิกฺขเว เอกธมฺโม ภาวิโต พหุลีกโต จตฺตาโร ธมฺเม ปริปูเรติ, จตฺตาโร ธมฺมา ภาวิตา พหุลีกตา สตฺต ธมฺเม ปริปูเรนฺติ, สตฺต ธมฺมา ภาวิตา พหุลีกตา เทฺว ธมฺเม ปริปูเรนฺติ.\csromanspacer{} อานาปานสฺสติสมาธิ ภิกฺขเว เอกธมฺโม ภาวิโต พหุลีกโต}

\csromanheader{2}{10. อานาปานสํยุตฺต}
\prose{\csromanfolio{293}จตฺตาโร สติปฏฺฐาเน ปริปูเรติ, จตฺตาโร สติปฏฺฐานา ภาวิตา พหุลีกตา สตฺต โพชฺฌงฺเค ปริปูเรนฺติ, สตฺต โพชฺฌงฺคา ภาวิตา พหุลีกตา วิชฺชาวิมุตฺติํ ปริปูเรนฺตีติ.}

\prose{กถํ ภาวิโต จ ภิกฺขเว อานาปานสฺสติสมาธิ กถํ พหุลีกโต จตฺตาโร สติปฏฺฐาเน ปริปูเรติ.\csromanspacer{} อิธ ภิกฺขเว ภิกฺขุ อรญฺญคโต วา รุกฺขมูลคโต วา สุญฺญาคารคโต วา นิสีทติ, ปลฺลงฺกํ อาภุชิตฺวา อุชุํ กายํ ปณิธาย ปริมุขํ สติํ อุปฏฺฐเปตฺวา, โส สโตว อสฺสสติ, สโตว ปสฺสสติ ฯเปฯ “ปฏินิสฺสคฺคานุปสฺสี อสฺสสิสฺสามี”ติ สิกฺขติ, “ปฏินิสฺสคฺคานุปสฺสี ปสฺสสิสฺสามี”ติ สิกฺขติ.}

\prose{ยสฺมิํ สมเย ภิกฺขเว ภิกฺขุ ทีฆํ วา อสฺสสนฺโต “ทีฆํ อสฺสสามี”ติ ปชานาติ, ทีฆํ วา ปสฺสสนฺโต “ทีฆํ ปสฺสสามี”ติ ปชานาติ.\csromanspacer{} รสฺสํ วา อสฺสสนฺโต “รสฺสํ อสฺสสามี”ติ ปชานาติ ฯเปฯ สพฺพกายปฺปฏิสํเวที ฯเปฯ “ปสฺสมฺภยํ กายสงฺขารํ อสฺสสิสฺสามี”ติ สิกฺขติ, “ปสฺสมฺภยํ กายสงฺขารํ ปสฺสสิสฺสามี”ติ สิกฺขติ.\csromanspacer{} กาเย กายานุปสฺสี ภิกฺขเว ภิกฺขุ ตสฺมิํ สมเย วิหรติ อาตาปี สมฺปชาโน สติมา วิเนยฺย โลเก อภิชฺฌาโทมนสฺสํ.\csromanspacer{} ตํ กิสฺส เหตุ, กายญฺญตราหํ ภิกฺขเว เอตํ วทามิ ยทิทํ อสฺสาสปสฺสาสํ.\csromanspacer{} ตสฺมาติห ภิกฺขเว กาเย กายานุปสฺสี ภิกฺขุ ตสฺมิํ สมเย วิหรติ อาตาปี สมฺปชาโน สติมา วิเนยฺย โลเก อภิชฺฌาโทมนสฺสํ. (1)}

\prose{ยสฺมิํ สมเย ภิกฺขเว ภิกฺขุ ปีติปฺปฏิสํเวที ฯเปฯ สุขปฺปฏิสํเวที.\csromanspacer{} จิตฺตสงฺขารปฺปฏิสํเวที. “ปสฺสมฺภยํ จิตฺตสงฺขารํ อสฺสสิสฺสามี”ติ สิกฺขติ, ปสฺสมฺภยํ จิตฺตสงฺขารํ “ปสฺสสิสฺสามี”ติ สิกฺขติ.\csromanspacer{} เวทนาสุ เวทนานุปสฺสี ภิกฺขเว ภิกฺขุ ตสฺมิํ สมเย วิหรติ อาตาปี สมฺปชาโน สติมา วิเนยฺย โลเก อภิชฺฌาโทมนสฺสํ.\csromanspacer{} ตํ กิสฺส เหตุ, เวทนาญฺญตราหํ ภิกฺขเว เอตํ วทามิ ยทิทํ อสฺสาสปสฺสาสานํ สาธุกํ มนสิการํ.\csromanspacer{} ตสฺมาติห ภิกฺขเว เวทนาสุ เวทนานุปสฺสี ภิกฺขุ ตสฺมิํ สมเย วิหรติ อาตาปี สมฺปชาโน สติมา วิเนยฺย โลเก อภิชฺฌาโทมนสฺสํ. (2)}

\prose{ยสฺมิํ สมเย ภิกฺขเว ภิกฺขุ จิตฺตปฺปฏิสํเวที ฯเปฯ อภิปฺปโมทยํ จิตฺตํ ฯเปฯ “สมาทหํ จิตฺตํ อสฺสสิสฺสามี”ติ สิกฺขติ, “สมาทหํ จิตฺตํ}

\prose{\csromanfolio{294}ปสฺสสิสฺสามี”ติ สิกฺขติ, “วิโมจยํ จิตฺตํ อสฺสสิสฺสามี”ติ สิกฺขติ, “วิโมจยํ จิตฺตํ ปสฺสสิสฺสามี”ติ สิกฺขติ.\csromanspacer{} จิตฺเต จิตฺตานุปสฺสี ภิกฺขเว ภิกฺขุ ตสฺมิํ สมเย วิหรติ อาตาปี สมฺปชาโน สติมา วิเนยฺย โลเก อภิชฺฌาโทมนสฺสํ.\csromanspacer{} ตํ กิสฺส เหตุ, นาหํ ภิกฺขเว มุฏฺฐสฺสติสฺส อสมฺปชานสฺส อานาปานสฺสติสมาธิภาวนํ วทามิ.\csromanspacer{} ตสฺมาติห ภิกฺขเว จิตฺเต จิตฺตานุปสฺสี ภิกฺขุ ตสฺมิํ สมเย วิหรติ อาตาปี สมฺปชาโน สติมา วิเนยฺย โลเก อภิชฺฌาโทมนสฺสํ. (3)}

\prose{ยสฺมิํ สมเย ภิกฺขเว ภิกฺขุ อนิจฺจานุปสฺสี ฯเปฯ วิราคานุปสฺสี ฯเปฯ นิโรธานุปสฺสี ฯเปฯ “ปฏินิสฺสคฺคานุปสฺสี อสฺสสิสฺสามี”ติ สิกฺขติ, “ปฏินิสฺสคฺคานุปสฺสี ปสฺสสิสฺสามี”ติ สิกฺขติ.\csromanspacer{} ธมฺเมสุ ธมฺมานุปสฺสี ภิกฺขเว ภิกฺขุ ตสฺมิํ สมเย วิหรติ อาตาปี สมฺปชาโน สติมา วิเนยฺย โลเก อภิชฺฌาโทมนสฺสํ.\csromanspacer{} โส ยํ ตํ โหติ อภิชฺฌาโทมนสฺสานํ ปหานํ, ตํ ปญฺญาย ทิสฺวา สาธุกํ อชฺฌุเปกฺขิตา โหติ.\csromanspacer{} ตสฺมาติห ภิกฺขเว ธมฺเมสุ ธมฺมานุปสฺสี ภิกฺขุ ตสฺมิํ สมเย วิหรติ อาตาปี สมฺปชาโน สติมา วิเนยฺย โลเก อภิชฺฌาโทมนสฺสํ. (4)}

\prose{เอวํ ภาวิโต โข ภิกฺขเว อานาปานสฺสติสมาธิ เอวํ พหุลีกโต จตฺตาโร สติปฏฺฐาเน ปริปูเรติ.}

\prose{กถํ ภาวิตา จ ภิกฺขเว จตฺตาโร สติปฏฺฐานา กถํ พหุลีกตา สตฺต โพชฺฌงฺเค ปริปูเรนฺติ.\csromanspacer{} ยสฺมิํ สมเย ภิกฺขเว ภิกฺขุ กาเย กายานุปสฺสี วิหรติ, อุปฏฺฐิตาสฺส ตสฺมิํ สมเย ภิกฺขุโน สติ โหติ อสมฺมุฏฺฐา.\csromanspacer{} ยสฺมิํ สมเย ภิกฺขเว ภิกฺขุโน อุปฏฺฐิตา สติ อสมฺมุฏฺฐา, สติสมฺโพชฺฌงฺโค ตสฺมิํ สมเย ภิกฺขุโน อารทฺโธ โหติ, สติสมฺโพชฺฌงฺคํ ตสฺมิํ สมเย ภิกฺขุ ภาเวติ, สติสมฺโพชฺฌงฺโค ตสฺมิํ สมเย ภิกฺขุโน ภาวนาปาริปูริํ คจฺฉติ. (1)}

\prose{โส ตถาสโต วิหรนฺโต ตํ ธมฺมํ ปญฺญาย ปวิจินติ ปวิจรติ ปริวีมํสมาปชฺชติ.\csromanspacer{} ยสฺมิํ สมเย ภิกฺขเว ภิกฺขุ ตถาสโต วิหรนฺโต ตํ ธมฺมํ ปญฺญาย ปวิจินติ ปวิจรติ ปริวีมํสมาปชฺชติ, ธมฺมวิจยสมฺโพชฺฌงฺโค ตสฺมิํ สมเย ภิกฺขุโน อารทฺโธ โหติ, ธมฺมวิจยสมฺโพชฺฌงฺคํ ตสฺมิํ สมเย ภิกฺขุ ภาเวติ, ธมฺมวิจยสมฺโพชฺฌงฺโค ตสฺมิํ สมเย ภิกฺขุโน ภาวนาปาริปูริํ คจฺฉติ. (2)}

\csromanheader{2}{10. อานาปานสํยุตฺต}
\prose{\csromanfolio{295}ตสฺส ตํ ธมฺมํ ปญฺญาย ปวิจินโต ปวิจรโต ปริวีมํสมาปชฺชโต อารทฺธํ โหติ วีริยํ อสลฺลีนํ.\csromanspacer{} ยสฺมิํ สมเย ภิกฺขเว ภิกฺขุโน ตํ ธมฺมํ ปญฺญาย ปวิจินโต ปวิจรโต ปริวีมํสมาปชฺชโต อารทฺธํ โหติ วีริยํ อสลฺลีนํ, วีริยสมฺโพชฺฌงฺโค ตสฺมิํ สมเย ภิกฺขุโน อารทฺโธ โหติ, วีริยสมฺโพชฺฌงฺคํ ตสฺมิํ สมเย ภิกฺขุ ภาเวติ, วีริยสมฺโพชฺฌงฺโค ตสฺมิํ สมเย ภิกฺขุโน ภาวนาปาริปูริํ คจฺฉติ. (3)}

\prose{อารทฺธวีริยสฺส อุปฺปชฺชติ ปีติ นิรามิสา.\csromanspacer{} ยสฺมิํ สมเย ภิกฺขเว ภิกฺขุโน อารทฺธวีริยสฺส อุปฺปชฺชติ ปีติ นิรามิสา, ปีติสมฺโพชฺฌงฺโค ตสฺมิํ สมเย ภิกฺขุโน อารทฺโธ โหติ, ปีติสมฺโพชฺฌงฺคํ ตสฺมิํ สมเย ภิกฺขุ ภาเวติ, ปีติสมฺโพชฺฌงฺโค ตสฺมิํ สมเย ภิกฺขุโน ภาวนาปาริปูริํ คจฺฉติ. (4)}

\prose{ปีติมนสฺส กาโยปิ ปสฺสมฺภติ, จิตฺตมฺปิ ปสฺสมฺภติ.\csromanspacer{} ยสฺมิํ สมเย ภิกฺขเว ภิกฺขุโน ปีติมนสฺส กาโยปิ ปสฺสมฺภติ, จิตฺตมฺปิ ปสฺสมฺภติ, ปสฺสทฺธิสมฺโพชฺฌงฺโค ตสฺมิํ สมเย ภิกฺขุโน อารทฺโธ โหติ, ปสฺสทฺธิสมฺโพชฺฌงฺคํ ตสฺมิํ สมเย ภิกฺขุ ภาเวติ, ปสฺสทฺธิสมฺโพชฺฌงฺโค ตสฺมิํ สมเย ภิกฺขุโน ภาวนาปาริปูริํ คจฺฉติ. (5)}

\prose{ปสฺสทฺธกายสฺส สุขิโน จิตฺตํ สมาธิยติ.\csromanspacer{} ยสฺมิํ สมเย ภิกฺขเว ภิกฺขุโน ปสฺสทฺธกายสฺส สุขิโน จิตฺตํ สมาธิยติ, สมาธิสมฺโพชฺฌงฺโค ตสฺมิํ สมเย ภิกฺขุโน อารทฺโธ โหติ, สมาธิสมฺโพชฺฌงฺคํ ตสฺมิํ สมเย ภิกฺขุ ภาเวติ, สมาธิสมฺโพชฺฌงฺโค ตสฺมิํ สมเย ภิกฺขุโน ภาวนาปาริปูริํ คจฺฉติ. (6)}

\prose{โส ตถาสมาหิตํ จิตฺตํ สาธุกํ อชฺฌุเปกฺขิตา โหติ.\csromanspacer{} ยสฺมิํ สมเย ภิกฺขเว ภิกฺขุ ตถาสมาหิตํ จิตฺตํ สาธุกํ อชฺฌุเปกฺขิตา โหติ, อุเปกฺขาสมฺโพชฺฌงฺโค ตสฺมิํ สมเย ภิกฺขุโน อารทฺโธ โหติ, อุเปกฺขาสมฺโพชฺฌงฺคํ ตสฺมิํ สมเย ภิกฺขุ ภาเวติ, อุเปกฺขาสมฺโพชฺฌงฺโค ตสฺมิํ สมเย ภิกฺขุโน ภาวนาปาริปูริํ คจฺฉติ. (7)}

\prose{ยสฺมิํ สมเย ภิกฺขเว ภิกฺขุ เวทนาสุ ฯเปฯ.\csromanspacer{} จิตฺเต ฯเปฯ.\csromanspacer{} ธมฺเมสุ ธมฺมานุปสฺสี วิหรติ, อุปฏฺฐิตาสฺส ตสฺมิํ สมเย ภิกฺขุโน สติ}

\csromanmatikaanchor{mtk.220}
\csromantocmarkref{subparagraph}{7-10. สํโยชนปฺปหานสุตฺตาทีนิ}{mtk.220}
\bookmark[dest={mtk.220},level=5]{7-10. สํโยชนปฺปหานสุตฺตาทีนิ}
\prose{\csromanfolio{296}โหติ อสมฺมุฏฺฐา.\csromanspacer{} ยสฺมิํ สมเย ภิกฺขเว ภิกฺขุโน อุปฏฺฐิตา สติ โหติ อสมฺมุฏฺฐา, สติสมฺโพชฺฌงฺโค ตสฺมิํ สมเย ภิกฺขุโน อารทฺโธ โหติ, สติสมฺโพชฺฌงฺคํ ตสฺมิํ สมเย ภิกฺขุ ภาเวติ, สติสมฺโพชฺฌงฺโค ตสฺมิํ สมเย ภิกฺขุโน ภาวนาปาริปูริํ คจฺฉติ ฯเปฯ.}

\prose{โส ตถาสมาหิตํ จิตฺตํ สาธุกํ อชฺฌุเปกฺขิตา โหติ.\csromanspacer{} ยสฺมิํ สมเย ภิกฺขเว ภิกฺขุ ตถาสมาหิตํ จิตฺตํ สาธุกํ อชฺฌุเปกฺขิตา โหติ, อุเปกฺขาสมฺโพชฺฌงฺโค ตสฺมิํ สมเย ภิกฺขุโน อารทฺโธ โหติ, อุเปกฺขาสมฺโพชฺฌงฺคํ ตสฺมิํ สมเย ภิกฺขุ ภาเวติ, อุเปกฺขาสมฺโพชฺฌงฺโค ตสฺมิํ สมเย ภิกฺขุโน ภาวนาปาริปูริํ คจฺฉติ.\csromanspacer{} เอวํ ภาวิตา โข ภิกฺขเว จตฺตาโร สติปฏฺฐานา เอวํ พหุลีกตา สตฺต โพชฺฌงฺเค ปริปูเรนฺติ.}

\prose{กถํ ภาวิตา จ ภิกฺขเว สตฺต โพชฺฌงฺคา กถํ พหุลีกตา วิชฺชาวิมุตฺติํ ปริปูเรนฺติ, อิธ ภิกฺขเว ภิกฺขุ สติสมฺโพชฺฌงฺคํ ภาเวติ วิเวกนิสฺสิตํ วิราคนิสฺสิตํ นิโรธนิสฺสิตํ โวสฺสคฺคปริณามิํ.\csromanspacer{} ธมฺมวิจยสมฺโพชฺฌงฺคํ ภาเวติ วิเวกนิสฺสิตํ วิราคนิสฺสิตํ นิโรธนิสฺสิตํ โวสฺสคฺคปริณามิํ ฯเปฯ.\csromanspacer{} อุเปกฺขาสมฺโพชฺฌงฺคํ ภาเวติ วิเวกนิสฺสิตํ วิราคนิสฺสิตํ นิโรธนิสฺสิตํ โวสฺสคฺคปริณามิํ.\csromanspacer{} เอวํ ภาวิตา โข ภิกฺขเว สตฺต โพชฺฌงฺคา เอวํ พหุลีกตา วิชฺชาวิมุตฺติํ ปริปูเรนฺตีติ.\csromanspacer{}\csromanspacer{}ฉฏฺฐํ.}

\csromanheader{2}{7. สํโยชนปฺปหานสุตฺต}
\proseitem{993}{อานาปานสฺสติสมาธิ ภิกฺขเว ภาวิโต พหุลีกโต สํโยชนปฺปหานาย สํวตฺตติ ฯเปฯ.\csromanspacer{}\csromanspacer{}สตฺตมํ.}

\csromanheader{2}{8. อนุสยสมุคฺฆาตสุตฺต}
\vspace{\tipitakaparskip}
\csromancenter{อนุสยสมุคฺฆาตาย สํวตฺตติ.\csromanspacer{}\csromanspacer{}อฏฺฐมํ.}
\csromanheader{2}{9. อทฺธานปริญฺญาสุตฺต}
\vspace{\tipitakaparskip}
\csromancenter{อทฺธานปริญฺญาย สํวตฺตติ.\csromanspacer{}\csromanspacer{}นวมํ.}
\csromanheader{2}{10. อานาปานสํยุตฺต}
\csromanheader{2}{10. อาสวกฺขยสุตฺต}
\proseitem{996}{\csromanfolio{297}อาสวานํ ขยาย สํวตฺตติ.\csromanspacer{} กถํ ภาวิโต จ ภิกฺขเว อานาปานสฺสติสมาธิ กถํ พหุลีกโต สํโยชนปฺปหานาย สํวตฺตติ.\csromanspacer{} อนุสยสมุคฺฆาตาย สํวตฺตติ.\csromanspacer{} อทฺธานปริญฺญาย สํวตฺตติ.\csromanspacer{} อาสวานํ ขยาย สํวตฺตติ, อิธ ภิกฺขเว ภิกฺขุ อรญฺญคโต วา รุกฺขมูลคโต วา ฯเปฯ “ปฏินิสฺสคฺคานุปสฺสี อสฺสสิสฺสามี”ติ สิกฺขติ, “ปฏินิสฺสคฺคานุปสฺสี ปสฺสสิสฺสามี”ติ สิกฺขติ.\csromanspacer{} เอวํ ภาวิโต โข ภิกฺขเว อานาปานสฺสติสมาธิ เอวํ พหุลีกโต สํโยชนปฺปหานาย สํวตฺตติ ฯเปฯ.\csromanspacer{} อนุสยสมุคฺฆาตาย สํวตฺตติ ฯเปฯ.\csromanspacer{} อทฺธานปริญฺญาย สํวตฺตติ ฯเปฯ.\csromanspacer{} สาสวานํ ขยาย สํวตฺตตีติ.\csromanspacer{}\csromanspacer{}ทสมํ.\csromanspacer{} ทุติโย วคฺโค.}

\titlehead{\textbf{ตสฺสุทฺทานํ}}
\csromangathameasure{อิจฺฉานงฺคลํ กงฺเขยฺยํ, อานนฺท อปเร ทุเว. \\ ภิกฺขู สํโยชนานุสยา, อทฺธานํ อาสวกฺขยนฺติ. \\ \textbf{อานาปานสํยุตฺตํ ทสมํ.}}
\csromangathagroup{\csromangathastanza{อิจฺฉานงฺคลํ กงฺเขยฺยํ, อานนฺท อปเร ทุเว. \\ ภิกฺขู สํโยชนานุสยา, อทฺธานํ อาสวกฺขยนฺติ. \\ \textbf{อานาปานสํยุตฺตํ ทสมํ.}}}

\csromanmatikaanchor{mtk.221}
\csromantocmarknopage{subparagraph}{11. โสตาปตฺติสํยุตฺต}{mtk.221}
\bookmark[dest={mtk.221},level=5]{11. โสตาปตฺติสํยุตฺต}
\csromanheader{6}{11. โสตาปตฺติสํยุตฺต}
\csromanmatikaanchor{mtk.222}
\csromantocmarknopage{subparagraph}{1. เวฬุทฺวารวคฺค}{mtk.222}
\bookmark[dest={mtk.222},level=5]{1. เวฬุทฺวารวคฺค}
\csromanheader{6}{1. เวฬุทฺวารวคฺค}
\csromanmatikaanchor{mtk.223}
\csromantocmarkref{subparagraph}{1. จกฺกวตฺติราชสุตฺต}{mtk.223}
\bookmark[dest={mtk.223},level=5]{1. จกฺกวตฺติราชสุตฺต}
\csromanheader{6}{1. จกฺกวตฺติราชสุตฺต}
\proseitem{997}{\csromanfolio{298}สาวตฺถินิทานํ.\csromanspacer{} ตตฺร โข ภควา ฯเปฯ เอตทโวจ “กิญฺจาปิ ภิกฺขเว ราชา จกฺกวตฺตี\footnote{จกฺกวตฺติ (สฺยา, กํ, อิ, ก)} จตุนฺนํ ทีปานํ อิสฺสริยาธิปจฺจํ รชฺชํ กาเรตฺวา กายสฺส เภทา ปรํ มรณา สุคติํ สคฺคํ โลกํ อุปปชฺชติ เทวานํ ตาวติํสานํ สหพฺยตํ, โส ตตฺถ นนฺทเน วเน อจฺฉราสงฺฆปริวุโต ทิพฺเพหิ จ ปญฺจหิ กามคุเณหิ สมปฺปิโต สมงฺคีภูโต ปริจาเรติ โส จตูหิ ธมฺเมหิ อสมนฺนาคโต, อถ โข โส อปริมุตฺโตว\footnote{อปริมุตฺโต จ (สฺยา, กํ, ก)} นิรยา, อปริมุตฺโต ติรจฺฉานโยนิยา, อปริมุตฺโต เปตฺติวิสยา, อปริมุตฺโต อปายทุคฺคติวินิปาตา.\csromanspacer{} กิญฺจาปิ ภิกฺขเว อริยสาวโก ปิณฺฑิยาโลเปน ยาเปติ, นนฺตกานิ จ ธาเรติ.\csromanspacer{} โส จตูหิ ธมฺเมหิ สมนฺนาคโต, อถ โข โส ปริมุตฺโต\footnote{ปริมุตฺโต จ (สฺยา, กํ, ก)} นิรยา, ปริมุตฺโต ติรจฺฉานโยนิยา, ปริมุตฺโต เปตฺติวิสยา, ปริมุตฺโต อปายทุคฺคติวินิปาตา.}

\prose{กตเมหิ จตูหิ, อิธ ภิกฺขเว อริยสาวโก พุทฺเธ อเวจฺจปฺปสาเทน สมนฺนาคโต โหติ “อิติปิ โส ภควา อรหํ สมฺมาสมฺพุทฺโธ วิชฺชาจรณสมฺปนฺโน สุคโต โลกวิทู อนุตฺตโร ปุริสทมฺมสารถิ สตฺถา เทวมนุสฺสานํ พุทฺโธ ภควา”ติ.\csromanspacer{} ธมฺเม อเวจฺจปฺปสาเทน สมนฺนาคโต โหติ “สฺวากฺขาโต ภควตา ธมฺโม สนฺทิฏฺฐิโก อกาลิโก เอหิปสฺสิโก โอปเนยฺยิโก ปจฺจตฺตํ เวทิตพฺโพ วิญฺญูหี”ติ.\csromanspacer{} สํเฆ อเวจฺจปฺปสาเทน สมนฺนาคโต โหติ “สุปฺปฏิปนฺโน ภควโต สาวกสํโฆ อุชุปฺปฏิปนฺโน ภควโต สาวกสํโฆ ญายปฺปฏิปนฺโน ภควโต สาวกสํโฆ สามีจิปฺปฏิปนฺโน ภควโต สาวกสํโฆ, ยทิทํ จตฺตาริ ปุริสยุคานิ อฏฺฐ ปุริสปุคฺคลา, เอส ภควโต สาวกสํโฆ อาหุเนยฺโย ปาหุเนยฺโย ทกฺขิเณยฺโย อญฺชลิกรณีโย อนุตฺตรํ ปุญฺญกฺเขตฺตํ โลกสฺสา”ติ.\csromanspacer{} อริยกนฺเตหิ}

\csromanheader{2}{11. โสตาปตฺติสํยุตฺต}
\csromanmatikaanchor{mtk.224}
\csromanmatikaanchor{mtk.225}
\csromantocmarkref{subparagraph}{2. พฺรหฺมจริโยคธสุตฺต}{mtk.224}
\bookmark[dest={mtk.224},level=5]{2. พฺรหฺมจริโยคธสุตฺต}
\csromantocmarkref{subparagraph}{3. ทีฆาวุอุปาสกสุตฺต}{mtk.225}
\bookmark[dest={mtk.225},level=5]{3. ทีฆาวุอุปาสกสุตฺต}
\prose{\csromanfolio{299}สีเลหิ สมนฺนาคโต โหติ อขณฺเฑหิ อจฺฉิทฺเทหิ อสพเลหิ อกมฺมาเสหิ ภุชิสฺเสหิ วิญฺญุปฺปสตฺเถหิ อปรามฏฺเฐหิ สมาธิสํวตฺตนิเกหิ.\csromanspacer{} อิเมหิ จตูหิ ธมฺเมหิ สมนฺนาคโต โหติ, โย จ ภิกฺขเว จตุนฺนํ ทีปานํ ปฏิลาโภ, โย จตุนฺนํ ธมฺมานํ ปฏิลาโภ, จตุนฺนํ ทีปานํ ปฏิลาโภ จตุนฺนํ ธมฺมานํ ปฏิลาภสฺส กลํ นาคฺฆติ โสฬสินฺติ.\csromanspacer{}\csromanspacer{}ปฐมํ.}

\csromanheader{6}{2. พฺรหฺมจริโยคธสุตฺต}
\proseitem{998}{จตูหิ ภิกฺขเว ธมฺเมหิ สมนฺนาคโต อริยสาวโก โสตาปนฺโน โหติ อวินิปาตธมฺโม นิยโต สมฺโพธิปรายโณ.}

\prose{กตเมหิ จตูหิ, อิธ ภิกฺขเว อริยสาวโก พุทฺเธ อเวจฺจปฺปสาเทน สมนฺนาคโต โหติ “อิติปิ โส ภควา อรหํ สมฺมาสมฺพุทฺโธ วิชฺชาจรณสมฺปนฺโน สุคโต โลกวิทู อนุตฺตโร ปุริสทมฺมสารถิ สตฺถา เทวมนุสฺสานํ พุทฺโธ ภควา”ติ.\csromanspacer{} ธมฺเม ฯเปฯ.\csromanspacer{} สํเฆ ฯเปฯ.\csromanspacer{} อริยกนฺเตหิ สีเลหิ สมนฺนาคโต โหติ อขณฺเฑหิ ฯเปฯ สมาธิสํวตฺตนิเกหิ.\csromanspacer{} อิเมหิ โข ภิกฺขเว จตูหิ ธมฺเมหิ สมนฺนาคโต อริยสาวโก โสตาปนฺโน โหติ อวินิปาตธมฺโม นิยโต สมฺโพธิปรายโณติ.}

\prose{อิทมโวจ ภควา, อิทํ วตฺวาน\footnote{วตฺวา (สี, อิ) เอวมีทิเสสุ ฐาเนสุ.} สุคโต, อถาปรํ เอตทโวจ สตฺถา\csromandash{}}

\csromangathameasure{“เยสํ สทฺธา จ สีลญฺจ, ปสาโท ธมฺมทสฺสนํ. \\ เต เว กาเลน ปจฺเจนฺติ, พฺรหฺมจริโยคธํ สุขนฺ”ติ. ทุติยํ.}
\csromangathagroup{\csromangathastanza{“เยสํ สทฺธา จ สีลญฺจ, ปสาโท ธมฺมทสฺสนํ. \\ เต เว กาเลน ปจฺเจนฺติ, พฺรหฺมจริโยคธํ สุขนฺ”ติ.\csromanspacer{} ทุติยํ.}}

\csromanheader{2}{3. ทีฆาวุ-อุปาสกสุตฺต}
\proseitem{999}{เอกํ สมยํ ภควา ราชคเห วิหรติ เวฬุวเน กลนฺทกนิวาเป.\csromanspacer{} เตน โข ปน สมเยน ทีฆาวุ อุปาสโก อาพาธิโก โหติ ทุกฺขิโต พาฬฺหคิลาโน.\csromanspacer{} อถ โข ทีฆาวุ อุปาสโก ปิตรํ โชติกํ คหปติํ อามนฺเตสิ “เอหิ ตฺวํ คหปติ เยน ภควา เตนุปสงฺกม, อุปสงฺกมิตฺวา มม วจเนน ภควโต ปาเท สิรสา วนฺท ‘ทีฆาวุ \csromanfolio{300}ภนฺเต อุปาสโก อาพาธิโก โหติ ทุกฺขิโต พาฬฺหคิลาโน, โส ภควโต ปาเท สิรสา วนฺทตี’ติ.\csromanspacer{} เอวญฺจ วเทหิ ‘สาธุ กิร ภนฺเต ภควา เยน ทีฆาวุสฺส อุปาสกสฺส นิเวสนํ เตนุปสงฺกมตุ อนุกมฺปํ อุปาทายา’ติ”. “เอวํ ตาตา”ติ โข โชติโก คหปติ ทีฆาวุสฺส อุปาสกสฺส ปฏิสฺสุตฺวา เยน ภควา เตนุปสงฺกมิ, อุปสงฺกมิตฺวา ภควนฺตํ อภิวาเทตฺวา เอกมนฺตํ นิสีทิ, เอกมนฺตํ นิสินฺโน โข โชติโก คหปติ ภควนฺตํ เอตทโวจ “ทีฆาวุ ภนฺเต อุปาสโก อาพาธิโก โหติ ทุกฺขิโต พาฬฺหคิลาโน, โส ภควโต ปาเท สิรสา วนฺทติ, เอวญฺจ วเทติ ‘สาธุ กิร ภนฺเต ภควา เยน ทีฆาวุสฺส อุปาสกสฺส นิเวสนํ เตนุปสงฺกมตุ อนุกมฺปํ อุปาทายา’ติ”.\csromanspacer{} อธิวาเสสิ ภควา ตุณฺหีภาเวน.}

\prose{อถ โข ภควา นิวาเสตฺวา ปตฺตจีวรมาทาย เยน ทีฆาวุสฺส อุปาสกสฺส นิเวสนํ เตนุปสงฺกมิ, อุปสงฺกมิตฺวา ปญฺญตฺเต อาสเน นิสีทิ, นิสชฺช โข ภควา ทีฆาวุํ อุปาสกํ เอตทโวจ “กจฺจิ เต ทีฆาวุ ขมนียํ กจฺจิ ยาปนียํ, กจฺจิ ทุกฺขา เวทนา ปฏิกฺกมนฺติ โน อภิกฺกมนฺติ, ปฏิกฺกโมสานํ ปญฺญายติ โน อภิกฺกโม”ติ.\csromanspacer{} น เม ภนฺเต ขมนียํ น ยาปนียํ, พาฬฺหา เม ทุกฺขา เวทนา อภิกฺกมนฺติ โน ปฏิกฺกมนฺติ, อภิกฺกโมสานํ ปญฺญายติ โน ปฏิกฺกโมติ.\csromanspacer{} ตสฺมาติห เต ทีฆาวุ เอวํ สิกฺขิตพฺพํ “พุทฺเธ อเวจฺจปฺปสาเทน สมนฺนาคโต ภวิสฺสามิ, “อิติปิ โส ภควา อรหํ สมฺมาสมฺพุทฺโธ วิชฺชาจรณสมฺปนฺโน สุคโต โลกวิทู อนุตฺตโร ปุริสทมฺมสารถิ สตฺถา เทวมนุสฺสานํ พุทฺโธ ภควา”ติ.\csromanspacer{} ธมฺเม ฯเปฯ.\csromanspacer{} สํเฆ ฯเปฯ.\csromanspacer{} อริยกนฺเตหิ สีเลหิ สมนฺนาคโต ภวิสฺสามิ อขณฺเฑหิ ฯเปฯ สมาธิสํวตฺตนิเกหิ”.\csromanspacer{} เอวํ หิ เต ทีฆาวุ สิกฺขิตพฺพนฺติ.}

\prose{ยานิมานิ ภนฺเต ภควตา จตฺตาริ โสตาปตฺติยงฺคานิ เทสิตานิ, สํวิชฺชนฺเตเต ธมฺมา มยิ, อหญฺจ เตสุ ธมฺเมสุ สนฺทิสฺสามิ.\csromanspacer{} อหํ หิ ภนฺเต พุทฺเธ อเวจฺจปฺปสาเทน สมนฺนาคโต “อิติปิ โส ภควา ฯเปฯ สตฺถา เทวมนุสฺสานํ พุทฺโธ ภควา”ติ.\csromanspacer{} ธมฺเม ฯเปฯ.\csromanspacer{} สํเฆ ฯเปฯ.\csromanspacer{} อริยกนฺเตหิ สีเลหิ สมนฺนาคโต อขณฺเฑหิ ฯเปฯ สมาธิสํวตฺตนิเกหีติ.\csromanspacer{} ตสฺมาติห ตฺวํ ทีฆาวุ อิเมสุ จตูสุ โสตาปตฺติยงฺเคสุ ปติฏฺฐาย ฉ วิชฺชาภาคิเย ธมฺเม อุตฺตริ ภาเวยฺยาสิ “อิธ ตฺวํ ทีฆาวุ สพฺพสงฺขาเรสุ}

\csromanheader{2}{11. โสตาปตฺติสํยุตฺต}
\csromanmatikaanchor{mtk.226}
\csromantocmarkref{subparagraph}{4-5. สาริปุตฺตสุตฺตานิ}{mtk.226}
\bookmark[dest={mtk.226},level=5]{4-5. สาริปุตฺตสุตฺตานิ}
\prose{\csromanfolio{301}อนิจฺจานุปสฺสี วิหราหิ อนิจฺเจ ทุกฺขสญฺญี ทุกฺเข อนตฺตสญฺญี ปหานสญฺญี วิราคสญฺญี นิโรธสญฺญี”ติ.\csromanspacer{} เอวํ หิ เต ทีฆาวุ สิกฺขิตพฺพนฺติ.}

\prose{เยเม ภนฺเต ภควตา ฉ วิชฺชาภาคิยา ธมฺมา เทสิตา, สํวิชฺชนฺเตเต ธมฺมา มยิ, อหญฺจ เตสุ ธมฺเมสุ สนฺทิสฺสามิ.\csromanspacer{} อหํ หิ เภนฺเต สพฺพสงฺขาเรสุ อนิจฺจานุปสฺสี วิหรามิ อนิจฺเจ ทุกฺขสญฺญิ ทุกฺเข อนตฺตสญฺญี ปหานสญฺญี วิราคสญฺญี นิโรธสญฺญี.\csromanspacer{} อปิ จ เม ภนฺเต เอวํ โหติ “มา เหวายํ โชติโก คหปติ มมจฺจเยน วิฆาตํ อาปชฺชีติ\footnote{อาปชฺชติ (ก)}, มา ตฺวํ ตาต ทีฆาวุ เอวํ มนสากาสิ, อิงฺฆ ตฺวํ ตาต ทีฆาวุ ยเทว เต ภควา อาห, ตเทว ตฺวํ สาธุกํ มนสิ กโรหี”ติ.}

\prose{อถ โข ภควา ทีฆาวุํ อุปาสกํ อิมินา โอวาเทน โอวทิตฺวา อุฏฺฐายาสนา ปกฺกามิ.\csromanspacer{} อถ โข ทีฆาวุ อุปาสโก อจิรปกฺกนฺตสฺส ภควโต กาลมกาสิ.\csromanspacer{} อถ โข สมฺพหุลา ภิกฺขู เยน ภควา เตนุปสงฺกมิํสุ, อุปสงฺกมิตฺวา ภควนฺตํ อภิวาเทตฺวา เอกมนฺตํ นิสีทิํสุ, เอกมนฺตํ นิสินฺนา โข เต ภิกฺขู ภควนฺตํ เอตทโวจุํ “โย โส ภนฺเต ทีฆาวุ นาม อุปาสโก ภควตา สํขิตฺเตน โอวาเทน โอวทิโต, โส กาลงฺกโต, ตสฺส กา คติ, โก อภิสมฺปราโย”ติ.\csromanspacer{} ปณฺฑิโต ภิกฺขเว ทีฆาวุ อุปาสโก ปจฺจปาทิ\footnote{อโหสิ สจฺจวาที (สฺยา, กํ, อิ, ก)} ธมฺมสฺสานุธมฺมํ, น จ มํ ธมฺมาธิกรณํ\footnote{น จ ธมฺมาธิกรณํ (สฺยา, กํ, อิ, ก)} วิเหเสสิ\footnote{วิเหเฐสิ (อิติปิ อญฺญตฺถ)}, ทีฆาวุ ภิกฺขเว อุปาสโก ปญฺจนฺนํ โอรมฺภาคิยานํ สํโยชนานํ ปริกฺขยา โอปปาติโก ตตฺถ ปรินิพฺพายี อนาวตฺติธมฺโม ตสฺมา โลกาติ.\csromanspacer{}\csromanspacer{}ตติยํ.}

\csromanheader{2}{4. ปฐมสาริปุตฺตสุตฺต}
\proseitem{1000}{เอกํ สมยํ อายสฺมา จ สาริปุตฺโต อายสฺมา จ อานนฺโท สาวตฺถิยํ วิหรนฺติ เชตวเน อนาถปิณฺฑิกสฺส อาราเม.\csromanspacer{} อถ โข อายสฺมา อานนฺโท สายนฺหสมยํ ปฏิสลฺลานา วุฏฺฐิโต ฯเปฯ เอกมนฺตํ นิสินฺโน โข อายสฺมา อานนฺโท อายสฺมนฺตํ สาริปุตฺตํ เอตทโวจ “กตินํ นุ โข อาวุโส สาริปุตฺต ธมฺมานํ สมนฺนาคมนเหตุ เอวมยํ ปชา ภควตา พฺยากตา ‘โสตาปนฺนา อวินิปาตธมฺมา นิยตา \csromanfolio{302}สมฺโพธิปรายณา’ติ”.\csromanspacer{} จตุนฺนํ โข อาวุโส ธมฺมานํ สมนฺนาคมนเหตุ เอวมยํ ปชา ภควตา พฺยากตา “โสตาปนฺนา อวินิปาตธมฺมา นิยตา สมฺโพธิปรายณา”.}

\prose{กตเมสํ จตุนฺนํ, อิธาวุโส อริยสาวโก พุทฺเธ อเวจฺจปฺปสาเทน สมนฺนาคโต โหติ “อิติปิ โส ภควา ฯเปฯ สตฺถา เทวมนุสฺสานํ พุทฺโธ ภควา”ติ.\csromanspacer{} ธมฺเม ฯเปฯ.\csromanspacer{} สํเฆ ฯเปฯ.\csromanspacer{} อริยกนฺเตหิ สีเลหิ สมนฺนาคโต โหติ อขณฺเฑหิ ฯเปฯ สมาธิสํวตฺตนิเกหิ.\csromanspacer{} อิเมสํ โข อาวุโส จตุนฺนํ ธมฺมานํ สมนฺนาคมนเหตุ เอวมยํ ปชา ภควตา พฺยากตา “โสตาปนฺนา อวินิปาตธมฺมา นิยตา สมฺโพธิปรายณา”ติ.\csromanspacer{}\csromanspacer{}จตุตฺถํ.}

\csromanheader{2}{5. ทุติยสาริปุตฺตสุตฺต}
\proseitem{1001}{อถ โข อายสฺมา สาริปุตฺโต เยน ภควา เตนุปสงฺกมิ, อุปสงฺกมิตฺวา ภควนฺตํ อภิวาเทตฺวา เอกมนฺตํ นิสีทิ, เอกมนฺตํ นิสินฺนํ โข อายสฺมนฺตํ สาริปุตฺตํ ภควา เอตทโวจ “โสตาปตฺติยงฺคํ โสตาปตฺติยงฺคนฺติ หิทํ สาริปุตฺต วุจฺจติ, กตมํ นุ โข สาริปุตฺต โสตาปตฺติยงฺคนฺ”ติ.\csromanspacer{} สปฺปุริสสํเสโว หิ ภนฺเต โสตาปตฺติยงฺคํ, สทฺธมฺมสฺสวนํ โสตาปตฺติยงฺคํ, โยนิโสมนสิกาโร โสตาปตฺติยงฺคํ, ธมฺมานุธมฺมปฺปฏิปตฺติ โสตาปตฺติยงฺคนฺติ.\csromanspacer{} สาธุ สาธุ สาริปุตฺต, สปฺปุริสสํเสโว หิ สาริปุตฺต โสตาปตฺติยงฺคํ, สทฺธมฺมสฺสวนํ โสตาปตฺติยงฺคํ, โยนิโสมนสิกาโร โสตาปตฺติยงฺคํ, ธมฺมานธมฺมปฺปฏิปตฺติ โสตาปตฺติยงฺคํ.}

\prose{“โสโต โสโต”ติ หิทํ สาริปุตฺต วุจฺจติ, กตโม นุ โข สาริปุตฺต โสโตติ.\csromanspacer{} อยเมว หิ ภนฺเต อริโย อฏฺฐงฺคิโก มคฺโค โสโต.\csromanspacer{} เสยฺยถิทํ, สมฺมาทิฏฺฐิ สมฺมาสงฺกปฺโป สมฺมาวาจา สมฺมากมฺมนฺโต สมฺมา-อาชีโว สมฺมาวายาโม สมฺมาสติ สมฺมาสมาธีติ.\csromanspacer{} สาธุ สาธุ สาริปุตฺต, อยเมว หิ สาริปุตฺต อริโย อฏฺฐงฺคิโก มคฺโค โสโต.\csromanspacer{} เสยฺยถิทํ, สมฺมาทิฏฺฐิ ฯเปฯ สมฺมาสมาธิ.}

\prose{“โสตาปนฺโน โสตาปนฺโน”ติ หิทํ สาริปุตฺต วุจฺจติ, กตโม นุ โข สาริปุตฺต โสตาปนฺโนติ.\csromanspacer{} โย หิ ภนฺเต อิมินา อริเยน อฏฺฐงฺคิเกน มคฺคเน สมนฺนาคโต, อยํ วุจฺจติ โสตาปนฺโน, สฺวายํ}

\csromanheader{2}{11. โสตาปตฺติสํยุตฺต}
\prose{\csromanfolio{303}อายสฺมา เอวํนาโม เอวํโคตฺโตติ.\csromanspacer{} สาธุ สาธุ สาริปุตฺต, โย หิ สาริปุตฺต อิมินา อริเยน อฏฺฐงฺคิเกน มคฺเคน สมนฺนาคโต, อยํ วุจฺจติ โสตาปนฺโน, สฺวายํ อายสฺมา เอวํนาโม เอวํโคตฺโตติ.\csromanspacer{}\csromanspacer{}ปญฺจมํ.}

\csromanmatikaanchor{mtk.227}
\csromantocmarkref{subparagraph}{6. ถปติสุตฺต}{mtk.227}
\bookmark[dest={mtk.227},level=5]{6. ถปติสุตฺต}
\csromanheader{6}{6. ถปติสุตฺต}
\proseitem{1002}{สาวตฺถินิทานํ.\csromanspacer{} เตน โข ปน สมเยน สมฺพหุลา ภิกฺขู ภควโต จีวรกมฺมํ กโรนฺติ “นิฏฺฐิตจีวโร ภควา เตมาสจฺจเยน จาริกํ ปกฺกมิสฺสตี”ติ.\csromanspacer{} เตน โข ปน สมเยน อิสิทตฺตปุราณา ถปตโย สาธุเก ปฏิวสนฺติ เกนจิเทว กรณีเยน.\csromanspacer{} อสฺโสสุํ โข อิสิทตฺตปุราณา ถปตโย “สมฺพหุลา กิร ภิกฺขู ภควโต จีวรกมฺมํ กโรนฺติ “นิฏฺฐิตจีวโร ภควา เตมาสจฺจเยน จาริกํ ปกฺกมิสฺสตี”ติ.}

\prose{อถ โข อิสิทตฺตปุราณา ถปตโย มคฺเค ปุริสํ ฐเปสุํ “ยทา ตฺวํ อมฺโภ ปุริส ปสฺเสยฺยาสิ ภควนฺตํ อาคจฺฉนฺตํ อรหนฺตํ สมฺมาสมฺพุทฺธํ, อถ อมฺหากํ อาโรเจยฺยาสี”ติ.\csromanspacer{} ทฺวีหตีหํ ฐิโต โข โส ปุริโส อทฺทส ภควนฺตํ ทูรโตว อาคจฺฉนฺตํ ทิสฺวาน เยน อิสิทตฺตปุราณา ถปตโย เตนุปสงฺกมิ, อุปสงฺกมิตฺวา อิสิทตฺตปุราเณ ถปตโย เอตทโวจ “อยํ โส ภนฺเต ภควา อาคจฺฉติ อรหํ สมฺมาสมฺพุทฺโธ, ยสฺสทานิ กาลํ มญฺญถา”ติ.}

\prose{อถ โข อิสิทตฺตปุราณา ถปตโย เยน ภควา เตนุปสงฺกมิํสุ, อุปสงฺกมิตฺวา ภควนฺตํ อภิวาเทตฺวา ภควนฺตํ ปิฏฺฐิโต ปิฏฺฐิโต อนุพนฺธิํสุ.\csromanspacer{} อถ โข ภควา มคฺคา โอกฺกมฺม เยน อญฺญตรํ รุกฺขมูลํ เตนุปสงฺกมิ, อุปสงฺกมิตฺวา ปญฺญตฺเต อาสเน นิสีทิ, อิสิทตฺตปุราณา ถปตโย ภควนฺตํ อภิวาเทตฺวา เอกมนฺตํ นิสีทิํสุ, เอกมนฺตํ นิสินฺนา โข เต อิสิทตฺตปุราณา ถปตโย ภควนฺตํ เอตทโวจุํ\csromandash{}}

\prose{ยทา มยํ ภนฺเต ภควนฺตํ สุณาม “สาวตฺถิยา โกสเลสุ จาริกํ ปกฺกมิสฺสตี”ติ, โหติ โน ตสฺมิํ สมเย อนตฺตมนตา โหติ โทมนสฺสํ “ทูเร โน ภควา ภวิสฺสตี”ติ.\csromanspacer{} ยทา ปน มยํ ภนฺเต ภควนฺตํ สุณาม “สาวตฺถิยา โกสเลสุ จาริกํ ปกฺกนฺโต”ติ, โหติ โน ตสฺมิํ สมเย อนตฺตมนตา โหติ โทมนสฺสํ, “ทูเร โน ภควา”ติ.}

\prose{\csromanfolio{304}ยทา ปน มยํ ภนฺเต ภควนฺตํ สุณาม “โกสเลหิ มลฺเลสุ จาริกํ ปกฺกมิสฺสตี”ติ, โหติ โน ตสฺมิํ สมเย อนตฺตมนตา โหติ โทมนสฺสํ “ทูเร โน ภควา ภวิสฺสตี”ติ.\csromanspacer{} ยทา ปน มยํ ภนฺเต ภควนฺตํ สุณาม “โกสเลหิ มลฺเลสุ จาริกํ ปกฺกนฺโต”ติ, โหติ โน ตสฺมิํ สมเย อนตฺตมนตา โหติ โทมนสฺสํ “ทูเร โน ภควา”ติ.}

\prose{ยทา ปน มยํ ภนฺเต ภควนฺตํ สุณาม “มลฺเลหิ วชฺชีสุ จาริกํ ปกฺกมิสฺสตี”ติ, โหติ โน ตสฺมิํ สมเย อนตฺตมนตา โหติ โทมนสฺสํ “ทูเร โน ภควา ภวิสฺสตี”ติ.\csromanspacer{} ยทา ปน มยํ ภนฺเต ภควนฺตํ สุณาม “มลฺเลหิ วชฺชีสุ จาริกํ ปกฺกนฺโต”ติ, โหติ โน ตสฺมิํ สมเย อนตฺตมนตา โหติ โทมนสฺสํ “ทูเร โน ภควา”ติ.}

\prose{ยทา ปน มยํ ภนฺเต ภควนฺตํ สุณาม “วชฺชีหิ กาสีสุ จาริกํ ปกฺกมิสฺสตี”ติ, โหติ โน ตสฺมิํ สมเย อนตฺตมนตา โหติ โทมนสฺสํ “ทูเร โน ภควา ภวิสฺสตี”ติ.\csromanspacer{} ยทา ปน มยํ ภนฺเต ภควนฺตํ สุณาม “วชฺชีหิ กาสีสุ จาริกํ ปกฺกนฺโต”ติ, โหติ โน ตสฺมิํ สมเย อนตฺตมนตา โหติ โทมนสฺสํ “ทูเร โน ภควา”ติ.}

\prose{ยทา ปน มยํ ภนฺเต ภควนฺตํ สุณาม “กาสีหิ มาคเธ จาริกํ ปกฺกมิสฺสตี”ติ, โหติ โน ตสฺมิํ สมเย อนตฺตมนตา โหติ โทมนสฺสํ “ทูเร โน ภควา ภวิสฺสตี”ติ.\csromanspacer{} ยทา ปน มยํ ภนฺเต ภควนฺตํ สุณาม “กาสีหิ มาคเธ จาริกํ ปกฺกนฺโต”ติ, โหติ อนปฺปกา โน ตสฺมิํ สมเย อนตฺตมนตา โหติ อนปฺปกํ โทมนสฺสํ “ทูเร โน ภควา”ติ.}

\prose{ยทา ปน มยํ ภนฺเต ภควนฺตํ สุณาม “มาคเธหิ กาสีสุ จาริกํ ปกฺกมิสฺสตี”ติ, โหติ โน ตสฺมิํ สมเย อตฺตมนตา โหติ โสมนสฺสํ “อาสนฺเน โน ภควา ภวิสฺสตี”ติ.\csromanspacer{} ยทา ปน มยํ ภนฺเต ภควนฺตํ สุณาม “มาคเธหิ กาสีสุ จาริกํ ปกฺกนฺโต”ติ, โหติ โน ตสฺมิํ สมเย อตฺตมนตา โหติ โสมนสฺสํ “อาสนฺเน โน ภควา”ติ.}

\prose{ยทา ปน มยํ ภนฺเต ภควนฺตํ สุณาม “กาสีหิ วชฺชีสุ จาริกํ ปกฺกมิสฺสตี”ติ, โหติ โน ตสฺมิํ สมเย อตฺตมนตา, โหติ}

\csromanheader{2}{11. โสตาปตฺติสํยุตฺต}
\prose{\csromanfolio{305}โสมนสฺสํ “อาสนฺเน โน ภควา ภวิสฺสตี”ติ.\csromanspacer{} ยทา ปน มยํ ภนฺเต ภควนฺตํ สุณม “กาสีหิ วชฺชีสุ จาริกํ ปกฺกนฺโต”ติ, โหติ โน ตสฺมิํ สมเย อตฺตมนตา, โหติ โสมนสฺสํ “อาสนฺเน โน ภควา”ติ.}

\prose{ยทา ปน มยํ ภนฺเต ภควนฺตํ สุณาม “วชฺชีหิ มลฺเลสุ จาริกํ ปกฺกมิสฺสตี”ติ, โหติ โน ตสฺมิํ สมเย อตฺตมนตา, โหติ โสมนสฺสํ “อาสนฺเน โน ภควา ภวิสฺสตี”ติ.\csromanspacer{} ยทา ปน มยํ ภนฺเต ภควนฺตํ สุณาม “วชฺชีหิ มลฺเลสุ จาริกํ ปกฺกนฺโต”ติ, โหติ โน ตสฺมิํ สมเย อตฺตมนตา, โหติ โสมนสฺสํ “อาสนฺเน โน ภควา”ติ.}

\prose{ยทา ปน มยํ ภนฺเต ภควนฺตํ สุณาม “มลฺเลหิ โกสเล จาริกํ ปกฺกมิสฺสตี”ติ, โหติ โน ตสฺมิํ สมเย อตฺตมนตา, โหติ โสมนสฺสํ “อาสนฺเน โน ภควา ภวิสฺสตี”ติ.\csromanspacer{} ยทา ปน มยํ ภนฺเต ภควนฺตํ สุณาม “มลฺเลหิ โกสเล จาริกํ ปกฺกนฺโต”ติ, โหติ โน ตสฺมิํ สมเย อตฺตมนตา, โหติ โสมนสฺสํ “อาสนฺเน โน ภควา”ติ.}

\prose{ยทา ปน มยํ ภนฺเต ภควนฺตํ สุณาม “โกสเลหิ สาวตฺถิํ จาริกํ ปกฺกมิสฺสตี”ติ, โหติ โน ตสฺมิํ สมเย อตฺตมนตา, โหติ โสมนสฺสํ “อาสนฺเน โน ภควา ภวิสฺสตี”ติ.\csromanspacer{} ยทา ปน มยํ ภนฺเต ภควนฺตํ สุณาม “สาวตฺถิยํ วิหรติ เชตวเน อนาถปิณฺฑิกสฺส อาราเม”ติ, โหติ อนปฺปกา โน ตสฺมิํ สมเย อตฺตมนตา, โหติ อนปฺปกํ โสมนสฺสํ “อาสนฺเน โน ภควา”ติ.}

\prose{ตสฺมาติห ถปตโย สมฺพาโธ ฆราวาโส รชาปโถ, อพฺโภกาโส ปพฺพชฺชา, อลญฺจ ปน โว ถปตโย อปฺปมาทายาติ.\csromanspacer{} อตฺถิ โข โน ภนฺเต เอตมฺหา สมฺพาธา อญฺโญ สมฺพาโธ สมฺพาธตโร เจว สมฺพาธสงฺขาตตโร จาติ.\csromanspacer{} กตโม ปน โว ถปตโย เอตมฺหา สมฺพาธา อญฺโญ สมฺพาโธ สมฺพาธตโร เจว สมฺพาธสงฺขาตตโร จาติ.}

\prose{อิธ มยํ ภนฺเต ยทา ราชา ปเสนทิ โกสโล อุยฺยานภูมิํ นิยฺยาตุกาโม โหติ, เย เต รญฺโญ ปเสนทิสฺส โกสลสฺส นาคา โอปวยฺหา, เต กปฺเปตฺวา, ยา ตา รญฺโญ ปเสนทิสฺส โกสลสฺส ปชาปติโย ปิยา มนาปา, ตา เอกํ ปุรโต}

\prose{\csromanfolio{306}เอกํ ปจฺฉโต นิสีทาเปม.\csromanspacer{} ตาสํ โข ปน ภนฺเต ภคินีนํ เอวรูโป คนฺโธ โหติ, เสยฺยถาปิ นาม คนฺธกรณฺฑกสฺส ตาวเทว วิวริยมานสฺส ยถา ตํ ราชกญฺญานํ คนฺเธน วิภูสิตานํ.\csromanspacer{} ตาสํ โข ปน ภนฺเต ภคินีนํ เอวรูโป กายสมฺผสฺโส โหติ, เสยฺยถาปิ นาม ตูลปิจุโน วา กปฺปาสปิจุโน วา ยถา ตํ ราชกญฺญานํ สุเขธิตานํ.\csromanspacer{} ตสฺมิํ โข ปน ภนฺเต สมเย นาโคปิ รกฺขิตพฺโพ โหติ, ตาปิ ภคินิโย รกฺขิตพฺพา โหนฺติ, อตฺตาปิ รกฺขิตพฺโพ โหติ, น โข ปน มยํ ภนฺเต อภิชานาม ตาสุ ภคินีสุ ปาปกํ จิตฺตํ อุปฺปาเทตา.\csromanspacer{} อยํ โข โน ภนฺเต เอตมฺหา สมฺพาธา อญฺโญ สมฺพาโธ สมฺพาธตโร เจว สมฺพาธสงฺขาตตโร จาติ.}

\prose{ตสฺมาติห ถปตโย สมฺพาโธ ฆราวาโส รชาปโถ, อพฺโภกาโส ปพฺพชฺชา, อลญฺจ ปน โว ถปตโย อปฺปมาทาย.\csromanspacer{} จตูหิ โข ถปตโย ธมฺเมหิ สมนฺนาคโต อริยสาวโก โสตาปนฺโน โหติ อวินิปาตธมฺโม นิยโต สมฺโพธิปรายโณ.}

\prose{กตเมหิ จตูหิ, อิธ ถปตโย อริยสาวโก พุทฺเธ อเวจฺจปฺปสาเทน สมนฺนาคโต โหติ “อิติปิ โส ภควา ฯเปฯ สตฺถา เทวมนุสฺสานํ พุทฺโธ ภควา”ติ.\csromanspacer{} ธมฺเม ฯเปฯ.\csromanspacer{} สํเฆ ฯเปฯ.\csromanspacer{} วิคตมลมจฺเฉเรน เจตสา อชฺฌาคารํ วสติ มุตฺตจาโค ปยตปาณิ โวสฺสคฺครโต ยาจโยโค ทานสํวิภาครโต, อิเมหิ โข ถปตโย จตูหิ ธมฺเมหิ สมนฺนาคโต อริยสาวโก โสตาปนฺโน โหติ อวินิปาตธมฺโม นิยโต สมฺโพธิปรายโณ.}

\prose{ตุเมฺห โข ถปตโย พุทฺเธ อเวจฺจปฺปสาเทน สมนฺนาคตา “อิติปิ โส ภควา ฯเปฯ สตฺถา เทวมนุสฺสานํ พุทฺโธ ภควา”ติ.\csromanspacer{} ธมฺเม ฯเปฯ.\csromanspacer{} สํเฆ ฯเปฯ.\csromanspacer{} ยํ โข ปน กิญฺจิ กุเล เทยฺยธมฺมํ, สพฺพํ ตํ อปฺปฏิวิภตฺตํ สีลวนฺเตหิ กลฺยาณธมฺเมหิ.\csromanspacer{} ตํ กิํ มญฺญถ ถปตโย, กติวิธา เต โกสเลสุ มนุสฺสา เย ตุมฺหากํ สมสมา ยทิทํ ทานสํวิภาเคติ.\csromanspacer{} ลาภา โน ภนฺเต, สุลทฺธํ โน ภนฺเต, เยสํ โน ภควา เอวํ ปชานาตีติ.\csromanspacer{}\csromanspacer{}ฉฏฺฐํ.}

\csromanmatikaanchor{mtk.228}
\csromantocmarkref{subparagraph}{7. เวฬุทฺวาเรยฺยสุตฺต}{mtk.228}
\bookmark[dest={mtk.228},level=5]{7. เวฬุทฺวาเรยฺยสุตฺต}
\csromanheader{6}{7. เวฬุทฺวาเรยฺยสุตฺต}
\proseitem{1003}{เอวํ เม สุตํ\csromandash{}เอกํ สมยํ ภควา โกสเลสุ จาริกํ จรมาโน มหตา ภิกฺขุสํเฆน สทฺธิํ เยน เวฬุทฺวารํ นาม โกสลานํ}

\csromanheader{2}{11. โสตาปตฺติสํยุตฺต}
\prose{\csromanfolio{307}พฺราหฺมณคาโม ตทวสริ.\csromanspacer{} อสฺโสสุํ โข เต เวฬุทฺวาเรยฺยกา พฺราหฺมณคหปติกา “สมโณ ขลุ โภ โคตโม สกฺยปุตฺโต สกฺยกุลา ปพฺพชิโต โกสเลสุ จาริกํ จรมาโน มหตา ภิกฺขุสํเฆน สทฺธิํ เวฬุทฺวารํ อนุปฺปตฺโต, ตํ โข ปน ภวนฺตํ โคตมํ เอวํ กลฺยาโณ กิตฺติสทฺโท อพฺภุคฺคโต ‘อิติปิ โส ภควา อรหํ สมฺมาสมฺพุทฺโธ วิชฺชาจรณสมฺปนฺโน สุคโต โลกวิทู อนุตฺตโร ปุริสทมฺมสารถิ สตฺถา เทวมนุสฺสานํ พุทฺโธ ภควา’\footnote{ภควาติ (สี, สฺยา, กํ, อิ)} โส อิมํ โลกํ สเทวกํ สมารกํ สพฺรหฺมกํ สสฺสมณพฺราหฺมณิํ ปชํ สเทวมนุสฺสํ สยํ อภิญฺญา สจฺฉิกตฺวา ปเวเทติ, โส ธมฺมํ เทเสติ อาทิกลฺยาณํ มชฺเฌกลฺยาณํ ปริโยสานกลฺยาณํ สาตฺถํ สพฺยญฺชนํ เกวลปริปุณฺณํ ปริสุทฺธํ พฺรหฺมจริยํ ปกาเสติ, สาธุ โข ปน ตถารูปานํ อรหตํ ทสฺสนํ โหตี”ติ.}

\prose{อถ โข เต เวฬุทฺวาเรยฺยกา พฺราหฺมณคหปติกา เยน ภควา เตนุปสงฺกมิํสุ, อุปสงฺกมิตฺวา อปฺเปกจฺเจ ภควนฺตํ อภิวาเทตฺวา เอกมนฺตํ นิสีทิํสุ, อปฺเปกจฺเจ ภควตา สทฺธิํ สมฺโมทิํสุ, สมฺโมทนียํ กถํ สารณียํ วีติสาเรตฺวา เอกมนฺตํ นิสีทิํสุ, อปฺเปกจฺเจ เยน ภควา เตนญฺชลิํ ปณาเมตฺวา เอกมนฺตํ นิสีทิํสุ, อปฺเปกจฺเจ ภควโต สนฺติเก นามโคตฺตํ สาเวตฺวา เอกมนฺตํ นิสีทิํสุ, อปฺเปกจฺเจ ตุณฺหีภูตา เอกมนฺตํ นิสีทิํสุ.\csromanspacer{} เอกมนฺตํ นิสินฺนา โข เต เวฬุทฺวาเรยฺยกา พฺราหฺมณคหปติกา ภควนฺตํ เอตทโวจุํ “มยํ โภ โคตม เอวํกามา เอวํฉนฺทา เอวํอธิปฺปายา ปุตฺตสมฺพาธสยนํ อชฺฌาวเสยฺยาม, กาสิกจนฺทนํ ปจฺจนุภเวยฺยาม, มาลาคนฺธวิเลปนํ ธาเรยฺยาม, ชาตรูปรชตํ สาทิเยยฺยาม, กายสฺส เภทา ปรํ มรณา สุคติํ สคฺคํ โลกํ อุปปชฺเชยฺยาม, เตสํ โน ภวํ โคตโม อมฺหากํ เอวํกามานํ เอวํฉนฺทานํ เอวํอธิปฺปายานํ ตถา ธมฺมํ เทเสตุ, ยถา มยํ ปุตฺตสมฺพาธสยนํ อชฺฌาวเสยฺยาม ฯเปฯ สุคติํ สคฺคํ โลกํ อุปปชฺเชยฺยามา”ติ.}

\prose{อตฺตูปนายิกํ โว คหปตโย ธมฺมปริยายํ เทเสสฺสามิ, ตํ สุณาถ สาธุกํ มนสิ กโรถ, ภาสิสฺสามีติ. “เอวํ โภ”ติ โข เต เวฬุทฺวาเรยฺยกา พฺราหฺมณคหปติกา ภควโต ปจฺจสฺโสสุํ.\csromanspacer{} ภควา เอตทโวจ\csromandash{}}

\prose{\csromanfolio{308}กตโม จ คหปตโย อตฺตุปนายิโก ธมฺมปริยาโย, อิธ คหปตโย อริยสาวโก อิติ ปฏิสญฺจิกฺขติ “อหํ โขสฺมิ ชีวิตุกาโม อมริตุกาโม สุขกาโม ทุกฺขปฺปฏิกูโล.\csromanspacer{} โย โข มํ ชีวิตุกามํ อมริตุกามํ สุขกามํ ทุกฺขปฺปฏิกูลํ ชีวิตา โวโรเปยฺย, น เมตํ อสฺส ปิยํ มนาปํ.\csromanspacer{} อหญฺเจว โข ปน ปรํ ชีวิตุกามํ อมริตุกามํ สุขกามํ ทุกฺขปฺปฏิกูลํ ชีวิตา โวโรเปยฺยํ, ปรสฺสปิ ตํ อสฺส อปฺปิยํ อมนาปํ.\csromanspacer{} โย โข มฺยายํ ธมฺโม อปฺปิโย อมนาโป, ปรสฺสเปโส ธมฺโม อปฺปิโย อมนาโป.\csromanspacer{} โย โข มฺยายํ ธมฺโม อปฺปิโย อมนาโป, กถาหํ ปรํ เตน สํโยเชยฺยนฺ”ติ.\csromanspacer{} โส อิติ ปฏิสงฺขาย อตฺตนา จ ปาณาติปาตา ปฏิวิรโต โหติ, ปรญฺจ ปาณาติปาตา เวรมณิยา สมาทเปติ, ปาณาติปาตา เวรมณิยา จ วณฺณํ ภาสติ, เอวมสฺสายํ กายสมาจาโร ติโกฏิปริสุทฺโธ โหติ.}

\prose{ปุน จปรํ คหปตโย อริยสาวโก อิติ ปฏิสญฺจิกฺขติ “โย โข เม อทินฺนํ เถยฺยสงฺขาตํ อาทิเยยฺย, น เมตํ อสฺส ปิยํ มนาปํ.\csromanspacer{} อหญฺเจว โข ปน ปรสฺส อทินฺนํ เถยฺยสงฺขาตํ อาทิยฺเยยฺยํ, ปรสฺสปิ ตํ อสฺส อปฺปิยํ อมนาปํ.\csromanspacer{} โย โข มฺยายํ ธมฺโม อปฺปิโย อมนาโป, ปรสฺสเปโส ธมฺโม อปฺปิโย อมนาโป.\csromanspacer{} โย โข มฺยายํ ธมฺโม อปฺปิโย อมนาโป, กถาหํ ปรํ เตน สํโยเชยฺยนฺ”ติ.\csromanspacer{} โส อิติ ปฏิสงฺขาย อตฺตนา จ อทินฺนาทานํ ปฏิวิรโต โหติ, ปรญฺจ อทินฺนาทานา เวรมณิยา สมาทเปติ, อทินฺนาทานา เวรมณิยา จ วณฺณํ ภาสติ, เอวมสฺสายํ กายสมาจาโร ติโกฏิปริสุทฺโธ โหติ.}

\prose{ปุน จปรํ คหปตโย อริยสาวโก อิติ ปฏิสญฺจิกฺขติ “โย โข เม ทาเรสุ จาริตฺตํ อาปชฺเชยฺย, น เมตํ อสฺส ปิยํ มนาปํ.\csromanspacer{} อยญฺเจว โข ปน ปรสฺส ทาเรสุ จาริตฺตํ อาปชฺเชยฺยํ, ปรสฺสปิ ตํ อสฺส อปฺปิยํ อมนาปํ.\csromanspacer{} โย โข มฺยายํ ธมฺโม อปฺปิโย อมนาโป, ปรสฺสเปโส ธมฺโม อปฺปิโย อมนาโป.\csromanspacer{} โย โข มฺยายํ ธมฺโม อปฺปิโย อมนาโป, กถาหํ ปรํ เตน สํโยเชยฺยนฺ”ติ.\csromanspacer{} โส อิติ ปฏิสงฺขาย อตฺตนา จ กาเมสุมิจฺฉาจารา ปฏิวิรโต โหติ, ปรญฺจ กาเมสุมิจฺฉาจารา เวรมณิยา สมาทเปติ, กาเมสุมิจฺฉาจารา เวรมณิยา จ}

\csromanheader{2}{11. โสตาปตฺติสํยุตฺต}
\prose{\csromanfolio{309}วณฺณํ ภาสติ, เอวมสฺสายํ กายสมาจาโร ติโกฏิปริสุทฺโธ โหติ.}

\prose{ปุน จปรํ คหปตโย อริยสาวโก อิติ ปฏิสญฺจิกฺขติ “โย โข เม มุสาวาเทน อตฺถํ ภญฺเชยฺย, น เมตํ อสฺส ปิยํ มนาปํ.\csromanspacer{} อหญฺเจว โข ปน ปรสฺส มุสาวาเทน อตฺถํ ภญฺเชยฺยํ, ปรสฺสปิ ตํ อสฺส อปฺปิยํ อมนาปํ.\csromanspacer{} โย โข มฺยายํ ธมฺโม อปฺปิโย อมนาโป, ปรสฺสเปโส ธมฺโม อปฺปิโย อมนาโป.\csromanspacer{} โย โข มฺยายํ ธมฺโม อปฺปิโย อมนาโป, กถาหํ ปรํ เตน สํโยเชยฺยนฺ”ติ.\csromanspacer{} โส อิติ ปฏิสงฺขาย อตฺตนา จ มุสาวาทา ปฏิวิรโต โหติ, ปรญฺจ มุสาวาทา เวรมณิยา สมาทเปติ, มุสาวาทา เวรมณิยา จ วณฺณํ ภาสติ, เอวมสฺสายํ วจีสมาจาโร ติโกฏิปริสุทฺโธ โหติ.}

\prose{ปุน จปรํ คหปตโย อริยสาวโก อิติ ปฏิสญฺจิกฺขติ โย โข มํ ปิสุณาย วาจาย มิตฺเต ภินฺเทยฺย\footnote{มิตฺเตหิ เภเทยฺย (สฺยา, กํ, อิ, ก)}, น เมตํ อสฺส ปิยํ มนาปํ.\csromanspacer{} อหญฺเจว โข ปน ปรํ ปิสุณาย วาจาย มิตฺเต ภินฺเทยฺยํ, ปรสฺสปิ ตํ อสฺส อปฺปิยํ อมนาปํ ฯเปฯ เอวมสฺสายํ วจีสมาจาโร ติโกฏิปริสุทฺโธ โหติ.}

\prose{ปุน จปรํ คหปตโย อริยสาวโก อิติ ปฏิสญฺจิกฺขติ โย โข มํ ผรุสาย วาจาย สมุทาจเรยฺย, น เมตํ อสฺส ปิยํ มนาปํ.\csromanspacer{} อหญฺเจว โข ปน ปรํ ผรุสาย วาจาย สมุทาจเรยฺยํ, ปรสฺสปิ ตํ อสฺส อปฺปิยํ อมนาปํ.\csromanspacer{} โย โข มฺยายํ ธมฺโม ฯเปฯ เอวมสฺสายํ วจีสมาจาโร ติโกฏิปริสุทฺโธ โหติ.}

\prose{ปุน จปรํ คหปตโย อริยสาวโก อิติ ปฏิสญฺจิกฺขติ “โย โข มํ สมฺผภาเสน สมฺผปฺปลาปภาเสน สมุทาจเรยฺย, น เมตํ อสฺส ปิยํ มนาปํ.\csromanspacer{} อหญฺเจว โข ปน ปรํ สมฺผภาเสน สมฺผปฺปลาปภาเสน สมุทาจเรยฺยํ, ปรสฺสปิ ตํ อสฺส อปฺปิยํ อมนาปํ.\csromanspacer{} โย โข มฺยายํ ธมฺโม อปฺปิโย อมนาโป, ปรสฺสเปโส ธมฺโม อปฺปิโย อมนาโป.\csromanspacer{} โย โข มฺยายํ ธมฺโม อปฺปิโย อมนาโป, กถาหํ ปรํ เตน สโยเชยฺยนฺ”ติ.\csromanspacer{} โส อิติ ปฏิสงฺขาย อตฺตนา จ สมฺผปฺปลาปา ปฏิวิรโต}

\csromanmatikaanchor{mtk.229}
\csromantocmarkref{subparagraph}{8-10. คิญฺชกาวสถสุตฺตานิ}{mtk.229}
\bookmark[dest={mtk.229},level=5]{8-10. คิญฺชกาวสถสุตฺตานิ}
\prose{\csromanfolio{310}โหติ, ปรญฺจ สมฺผปฺปลาปา เวรมณิยา สมาทเปติ, สมฺผปฺปลาปา เวรมณิยา จ วณฺณํ ภาสติ, เอวมสฺสายํ วจีสมาจาโร ติโกฏิปริสุทฺโธ โหติ.}

\prose{โส พุทฺเธ อเวจฺจปฺปสาเทน สมนฺนาคโต โหติ “อิติปิ โส ภควา ฯเปฯ สตฺถา เทวมนุสฺสานํ พุทฺโธ ภควา”ติ.\csromanspacer{} ธมฺเม ฯเปฯ.\csromanspacer{} สํเฆ อเวจฺจปฺปสาเทน สมนฺนาคโต โหติ “สุปฺปฏิปนฺโน ภควโต สาวกสํโฆ ฯเปฯ อนุตฺตรํ ปุญฺญกฺเขตฺตํ โลกสฺสา”ติ.\csromanspacer{} อริยกนฺเตหิ สีเลหิ สมนฺนาคโต โหติ อขณฺเฑหิ ฯเปฯ สมาธิสํวตฺตนิเกหิ.\csromanspacer{} ยโต โข คหปตโย อริยสาวโก อิเมหิ สตฺตหิ สทฺธมฺเมหิ\footnote{ธมฺเมหิ (สี)} สมนฺนาคโต โหติ อิเมหิ จตูหิ อากงฺขิเยหิ ฐาเนหิ, โส อากงฺขมาโน อตฺตนาว อตฺตานํ พฺยากเรยฺย “ขีณนิรโยมฺหิ ขีณติรจฺฉานโยนิ\footnote{ขีณติรจฺฉานโยนิโย (สี, สฺยา, กํ, อิ), ขีณติรจฺฉานโยนิโก (ก)} ขีณเปตฺติวิสโย ขีณาปายทุคฺคติวินิปาโต, โสตาปนฺโนหมสฺมิ อวินิปาตธมฺโม นิยโต สมฺโพธิปรายโณ”ติ.}

\prose{เอวํ วุตฺเต เวฬุทฺวาเรยฺยกา พฺราหฺมณคหปติกา ภควนฺตํ เอตทโวจุํ “อภิกฺกนฺตํ โภ โคตม ฯเปฯ เอเต มยํ ภวนฺตํ โคตมํ สรณํ คจฺฉาม ธมฺมญฺจ ภิกฺขุสํฆญฺจ, อุปาสเก โน ภวํ โคตโม ธาเรตุ อชฺชตคฺเค ปาณุเปเต\footnote{ปาณุเปตํ (ก)} สรณํ คเต”ติ.\csromanspacer{} สตฺตมํ.}

\csromanheader{2}{8. ปฐมคิญฺชกาวสถสุตฺต}
\proseitem{1004}{เอวํ เม สุตํ\csromandash{}เอกํ สมยํ ภควา ญาติเก วิหรติ คิญฺชกาวสเถ.\csromanspacer{} อถ โข อายสฺมา อานนฺโท เยน ภควา เตนุปสงฺกมิ, อุปสงฺกมิตฺวา ภควนฺตํ อภิวาเทตฺวา เอกมนฺตํ นิสีทิ, เอกมนฺตํ นิสินฺโน โข อายสฺมา อานนฺโท ภควนฺตํ เอตทโวจ\csromandash{}}

\prose{สาโฬฺห นาม ภนฺเต ภิกฺขุ กาลงฺกโต, ตสฺส กา คติ, โก อภิสมฺปราโย.\csromanspacer{} นนฺทา นาม ภนฺเต ภิกฺขุนี กาลงฺกตา, ตสฺสา กา คติ, โก อภิสมฺปราโย.\csromanspacer{} สุทตฺโต นาม ภนฺเต อุปาสโก กาลงฺกโต, ตสฺส กา คติ, โก อภิสมฺปราโย.\csromanspacer{} สุชาตา นาม ภนฺเต อุปาสิกา กาลงฺกตา, ตสฺสา กา คติ, โก อภิสมฺปราโยติ.}

\csromanheader{2}{11. โสตาปตฺติสํยุตฺต}
\prose{\csromanfolio{311}สาโฬฺห อานนฺท ภิกฺขุ กาลงฺกโต อาสวานํ ขยา อนาสวํ เจโตวิมุตฺติํ ปญฺญาวิมุตฺติํ ทิฏฺเฐว ธมฺเม สยํ อภิญฺญา สจฺฉิกตฺวา อุปสมฺปชฺช วิหาสิ.\csromanspacer{} นนฺทา อานนฺท ภิกฺขุนี กาลงฺกตา ปญฺจนฺนํ โอรมฺภาคิยานํ สํโยชนานํ ปริกฺขยา โอปปาติกา ตตฺถ ปรินิพฺพายินี อนาวตฺติธมฺมา ตสฺมา โลกา.\csromanspacer{} สุทตฺโต อานนฺท อุปาสโก กาลงฺกโต ติณฺณํ สํโยชนานํ ปริกฺขยา ราคโทสโมหานํ ตนุตฺตา สกทาคามี สกิเทว อิมํ โลกํ อาคนฺตฺวา ทุกฺขสฺสนฺตํ กริสฺสติ.\csromanspacer{} สุชาตา อานนฺท อุปาสิกา กาลงฺกตา ติณฺณํ สํโยชนานํ ปริกฺขยา โสตาปนฺนา อวินิปาตธมฺมา นิยตา สมฺโพธิปรายณา.}

\prose{อนจฺฉริยํ โข ปเนตํ อานนฺท ยํ มนุสฺสภูโต กาลํ กเรยฺย, ตสฺมิํ ตสฺมิํ เจ มํ กาลงฺกเต อุปสงฺกมิตฺวา เอตมตฺถํ ปฏิปุจฺฉิสฺสถ, วิเหสาเปสา อานนฺท อสฺส ตถาคตสฺส, ตสฺมาติหานนฺท ธมฺมาทาสํ นาม ธมฺมปริยายํ เทเสสฺสามิ, เยน สมนฺนาคโต อริยสาวโก อากงฺขมาโน อตฺตนาว อตฺตานํ พฺยากเรยฺย “ขีณนิรโยมฺหิ ขีณติรจฺฉานโยนิ ขีณเปตฺติวิสโย ขีณาปายทุคฺคติวินิปาโต, โสตาปนฺโนหมสฺมิ อวินิปาตธมฺโม นิยโต สมฺโพธิปรายโณ”.}

\prose{กตโม จ โส อานนฺท ธมฺมาทาโส ธมฺมปริยาโย, เยน สมนฺนาคโต อริยสาวโก อากงฺขมาโน อตฺตนาว อตฺตานํ พฺยากเรยฺย “ขีณนิรโยมฺหิ ขีณติรจฺฉานโยนิ ขีณเปตฺติวิสโย ขีณาปายทุคฺคติวินิปาโต, โสตาปนฺโนหมสฺมิ อวินิปาตธมฺโม นิยโต สมฺโพธิปรายโณ”.}

\prose{อิธ อานนฺท อริยสาวโก พุทฺเธ อเวจฺจปฺปสาเทน สมนฺนาคโต โหติ “อิติปิ โส ภควา ฯเปฯ สตฺถา เทวมนุสฺสานํ พุทฺโธ ภควา”ติ.\csromanspacer{} ธมฺเม ฯเปฯ.\csromanspacer{} สํเฆ ฯเปฯ.\csromanspacer{} อริยกนฺเตหิ สีเลหิ สมนฺนาคโต โหติ อขณฺเฑหิ ฯเปฯ สมาธิสํวตฺตนิเกหิ.\csromanspacer{} อยํ โข โส อานนฺท ธมฺมฺทาโส ธมฺมปรุยาโย, เยน สมนฺนาคโต อริยสาวโก อากงฺขมาโน อตฺตนาว อตฺตานํ พฺยากเรยฺย “ขีณนิรโยมฺหิ ขีณติรจฺฉานโยนิ ขีณเปตฺติวิสโย ขีณาปายทุคฺคติวินิปาโต, โสตาปนฺโนหมสฺมิ อวินิปาตธมฺโม นิยโต สมฺโพธิปรายโณ”ติ.\csromanspacer{} อฏฺฐมํ. (ตีณิปิ สุตฺตนฺตานิ เอกนิทานานิ.)}

\csromanheader{2}{9. ทุติยคิญฺชกาวสถสุตฺต}
\proseitem{1005}{\csromanfolio{312}เอกมนฺตํ นิสินฺโน โข อายสฺมา อานนฺโท ภควนฺตํ เอตทโวจ “อโสโก นาม ภนฺเต ภิกฺขุ กาลงฺกโต, ตสฺส กา คติ, โก อภิสมฺปราโย.\csromanspacer{} อโสกา นาม ภนฺเต ภิกฺขุนี กาลงฺกตา ฯเปฯ.\csromanspacer{} อโสโก นาม ภนฺเต อุปาสโก กาลงฺกโต ฯเปฯ อโสกา นาม ภนฺเต อุปาสิกา กาลงฺกตา, ตสฺสา กา คติ, โก อภิสมฺปราโย”ติ.}

\prose{อโสโก อานนฺท ภิกฺขุ กาลงฺกโต อาสวานํ ขยา อนาสวํ เจโตวิมุตฺติํ ปญฺญาวิมุตฺติํ ทิฏฺเฐว ธมฺเม สยํ อภิญฺญา สจฺฉิกตฺวา อุปสมฺปชฺช วิหาสิ ฯเปฯ. (ปุริมเวยฺยากรเณน เอกนิทานํ.)}

\prose{อยํ โข โส อานนฺท ธมฺมาทาโส ธมฺมปริยาโย, เยน สมนฺนาคโต อริยสาวโก อากงฺขมาโน อตฺตนาว อตฺตานํ พฺยากเรยฺย “ขีณนิรโยมฺหิ ขีณติรจฺฉานโยนิ ขีณเปตฺติวิสโย ขิณาปายทุคฺคติวินิปาโต, โสตาปนฺโนหมสฺมิ อวินิปาตธมฺโม นิยโต สมฺโพธิปรายโณ”ติ.\csromanspacer{}\csromanspacer{}นวมํ.}

\csromanheader{2}{10. ตติยคิญฺชกาวสถสุตฺต}
\proseitem{1006}{เอกมนฺตํ นิสินฺโน โข อายสฺมา อานนฺโท ภควนฺตํ เอตทโวจ “กกฺกโฏ นาม ภนฺเต ญาติเก อุปาสโก กาลงฺกโต, ตสฺส กา คติ, โก อภิสมฺปราโย.\csromanspacer{} กฬิโภ นาม ภนฺเต ญาติเก อุปาสโก ฯเปฯ นิกโต นาม ภนฺเต ญาติเก อุปาสโก.\csromanspacer{} กฏิสฺสโห นาม ภนฺเต ญาติเก อุปาสโก.\csromanspacer{} ตุฏฺโฐ นาม ภนฺเต ญาติเก อุปาสโก.\csromanspacer{} สนฺตุฏฺโฐ นาม ภนฺเต ญาติเก อุปาสโก.\csromanspacer{} ภทฺโท นาม ภนฺเต ญาติเก อุปาสโก.\csromanspacer{} สุภทฺโท นาม ภนฺเต ญาติเก อุปาสโก กาลงฺกโต, ตสฺส กา คติ, โก อภิสมฺปราโยติ.}

\prose{กกฺกโฏ อานนฺท อุปาสโก กาลงฺกโต ปญฺจนฺนํ โอรมฺภาคิยานํ สํโยชนานํ ปริกฺขยา โอปปาติโก ตตฺถ ปรินิพฺพายี อนาวตฺติธมฺโม ตสฺมา โลกา.\csromanspacer{} กฬิโภ อานนฺท ฯเปฯ.\csromanspacer{} นิกโต อานนฺท.\csromanspacer{} กฏิสฺสโห}

\csromanheader{2}{11. โสตาปตฺติสํยุตฺต}
\prose{\csromanfolio{313}อานนฺท.\csromanspacer{} ตุฏฺโฐ อานนฺท.\csromanspacer{} สนฺตุฏฺโฐ อานนฺท.\csromanspacer{} ภทฺโท อานนฺท.\csromanspacer{} สุภทฺโท อานนฺท อุปาสโก กาลงฺกโต ปญฺจนฺนํ โอรมฺภาคิยานํ สํโยชนานํ ปริกฺขยา โอปปาติโก ตตฺถ ปรินิพฺพายี อนาวตฺติธมฺโม ตสฺมา โลกา. (สพฺเพ เอกคติกา กาตพฺพา.)}

\prose{ปโรปญฺญาส อานนฺท ญาติเก อุปาสกา กาลงฺกตา ปญฺจนฺนํ โอรมฺภาคิยานํ สํโยชนานํ ปริกฺขยา โอปปาติกา ตตฺถ ปรินิพฺพายิโน อนาวตฺติธมฺมา ตสฺมา โลกา.\csromanspacer{} สาธิกนวุติ อานนฺท ญาติเก อุปาสกา กาลงฺกตา ติณฺณํ สํโยชนานํ ปริกฺขยา ราคโทสโมหานํ ตนุตฺตา สกทาคามิโน สกิเทว อิมํ โลกํ อาคนฺตฺวา ทุกฺขสฺสนฺตํ กริสฺสนฺติ.\csromanspacer{} ฉาติเรกานิ โข อานนฺท ปญฺจสตานิ ญาติเก อุปาสกา กาลงฺกตา ติณฺณํ สํโยชนานํ ปริกฺขยา โสตาปนฺนา อวินิปาตธมฺมา นิยตา สมฺโพธิปรายณา.}

\prose{อนจฺฉริยํ โข ปเนตํ อานนฺท ยํ มนุสฺสภูโต กาลํ กเรยฺย, ตสฺมิํ ตสฺมิํ เจ มํ กาลงฺกเต อุปสงฺกมิตฺวา เอตมตฺถํ ปฏิปุจฺฉิสฺสถ, วิเหสาเปสา อาอานนฺท อสฺส ตถาคตสฺส.\csromanspacer{} ตสฺมาติหานนฺท ธมฺมาทาสํ นาม ธมฺมปริยายํ เทเสสฺสามิ, เยน สมนฺนาคโต อริยสาวโก อากงฺขมาโน อตฺตนาว อตฺตานํ พฺยากเรยฺย “ขีณนิรโยมฺหิ ขีณติรจฺฉานโยนิ ขีณเปตฺติวิสโย ขีณาปายทุคฺคติวินิปาโต, โสตาปนฺโนหมสฺมิ อวินิปาตธมฺโม นิยโต สมฺโพธิปรายโณ”.}

\prose{กตโม จ โส อานนฺท ธมฺมาทาโส ธมฺมปริยาโย เยน สมนฺนาคโต อริยสาวโก อากงฺขมาโน อตฺตนาว อตฺตานํ พฺยากเรยฺย “ขีณนิรโยมฺหิ ขีณติรจฺฉานโยนิ ขีณเปตฺติวิสโย ขีณาปายทุคฺคติวินิปาโต, โสตาปนฺโนหมสฺมิ อวินิปาตธมฺโม นิยโต สมฺโพธิปรายโณ”.}

\prose{อิธานนฺท อริยสาวโก พุทฺเธ อเวจฺจปฺปสาเทน สมนฺนาคโต โหติ “อิติปิ โส ภควา ฯเปฯ สตฺถา เทวมนุสฺสานํ พุทฺโธ ภควา”ติ.\csromanspacer{} ธมฺเม ฯเปฯ.\csromanspacer{} สํเฆ ฯเปฯ.\csromanspacer{} อริยกนฺเตหิ สีเลหิ สมนฺนาคโต โหติ อขณฺเฑหิ ฯเปฯ สมาธิสํวตฺตนิเกหิ.\csromanspacer{} อยํ โข โส อานนฺท ธมฺมาทาโส ธมฺมปริยาโย เยน สมนฺนาคโต อริยสาวโก อากงฺขมาโน อตฺตนาว อตฺตานํ พฺยากเรยฺย “ขีณนิรโยมฺหิ ขีณติรจฺฉานโยนิ}

\prose{\csromanfolio{314}ขีณเปตฺติวิสโย ขีณาปายทุคฺคติวินิปาโต, โสตาปนฺโนหมสฺมิ อวินิปาตธมฺโม นิยโต สมฺโพธิปรายโณ”ติ.\csromanspacer{}\csromanspacer{}ทสมํ.\csromanspacer{} เวฬุทฺวารวคฺโค ปฐโม.}

\titlehead{\textbf{ตสฺสุทฺทานํ}}
\csromangathameasure{ราชา โอคธทีฆาวุ, สาริปุตฺตาปเร ทุเว. \\ ถปตี เวฬุทฺวาเรยฺยา, คิญฺชกาวสเถ ตโยติ.}
\csromangathagroup{\csromangathastanza{ราชา โอคธทีฆาวุ, สาริปุตฺตาปเร ทุเว. \\ ถปตี เวฬุทฺวาเรยฺยา, คิญฺชกาวสเถ ตโยติ.}}

\csromanmatikaanchor{mtk.230}
\csromantocmarknopage{subparagraph}{10. ราชการามวคฺค}{mtk.230}
\bookmark[dest={mtk.230},level=5]{10. ราชการามวคฺค}
\csromanheader{6}{2. ราชการามวคฺค}
\csromanmatikaanchor{mtk.231}
\csromantocmarkref{subparagraph}{1. สหสฺสภิกฺขุนิสํฆสุตฺต}{mtk.231}
\bookmark[dest={mtk.231},level=5]{1. สหสฺสภิกฺขุนิสํฆสุตฺต}
\csromanheader{6}{1. สหสฺสภิกฺขุนิสํฆสุตฺต}
\proseitem{1007}{เอกํ สมยํ ภควา สาวตฺถิยํ วิหรติ ราชการาเม.\csromanspacer{} อถ โข สหสฺสภิกฺขุนิสํโฆ\footnote{สหสฺโส ภิกฺขุนิสํโฆ (สี)} เยน ภควา เตนุปสงฺกมิ, อุปสงฺกมิตฺวา ภควนฺตํ อภิวาเทตฺวา เอกมนฺตํ อฏฺฐาสิ, เอกมนฺตํ ฐิตา โข ตา ภิกฺขุนิโย ภควา เอตทโวจ\csromandash{}}

\prose{จตูหิ โข ภิกฺขุนิโย ธมฺเมหิ สมนฺนาคโต อริยสาวโก โสตาปนฺโน โหติ อวินิปาตธมฺโม นิยโต สมฺโพธิปรายโณ.\csromanspacer{} กตเมหิ จตูหิ, อิธ ภิกฺขุนิโย อริยสาวโก พุทฺเธ อเวจฺจปฺปสาเทน สมนฺนาคโต โหติ “อิติปิ โส ภควา ฯเปฯ สตฺถา เทวมนุสฺสานํ พุทฺโธ ภควา”ติ.\csromanspacer{} ธมฺเม ฯเปฯ.\csromanspacer{} สํเฆ ฯเปฯ.\csromanspacer{} อริยกนฺเตหิ สีเลหิ สมนฺนาคโต โหติ อขณฺเฑหิ ฯเปฯ สมาธิสํวตฺตนิเกหิ.\csromanspacer{} อิเมหิ โข ภิกฺขุนิโย จตูหิ ธมฺเมหิ สมนฺนาคโต อริยสาวโก โสตาปนฺโน โหติ อวินิปาตธมฺโม นิยโต สมฺโพธิปรายโณติ.\csromanspacer{}\csromanspacer{}ปฐมํ.}

\csromanmatikaanchor{mtk.232}
\csromantocmarkref{subparagraph}{2. พฺราหฺมณสุตฺต}{mtk.232}
\bookmark[dest={mtk.232},level=5]{2. พฺราหฺมณสุตฺต}
\csromanheader{6}{2. พฺราหฺมณสุตฺต}
\proseitem{1008}{สาวตฺถินิทานํ.\csromanspacer{} พฺราหฺมณา ภิกฺขเว อุทยคามินิํ นาม ปฏิปทํ ปญฺญเปนฺติ, เต สาวกํ เอวํ สมาทเปนฺติ “เอหิ ตฺวํ อมฺโภ ปุริส, กาลสฺเสว อุฏฺฐาย ปาจีนมุโข ยาหิ, โส ตฺวํ มา โสพฺภํ ปริวชฺเชหิ}

\csromanheader{2}{11. โสตาปตฺติสํยุตฺต}
\csromanmatikaanchor{mtk.233}
\csromantocmarkref{subparagraph}{3-5. อานนฺทตฺเถรสุตฺตาทีนิ}{mtk.233}
\bookmark[dest={mtk.233},level=5]{3-5. อานนฺทตฺเถรสุตฺตาทีนิ}
\prose{\csromanfolio{315}มา ปปาตํ มา ขาณุํ มา กณฺฑกฐานํ\footnote{กณฺฑกํ ฐานํ (อิ, ก)} มา จนฺทนิยํ มา โอฬิคลฺลํ, ยตฺถ\footnote{ยตฺเถว (สฺยา, กํ), ยานิ วา (สี)} ปปเตยฺยาสิ, ตตฺเถว มรณํ อาคเมยฺยาสิ.\csromanspacer{} เอวํ ตฺวํ อมฺโภ ปุริส กายสฺส เภทา ปรํ มรณา สุคติํ สคฺคํ โลกํ อุปปชฺชิสฺสสี”ติ.}

\prose{ตํ โข ปเนตํ ภิกฺขเว พฺราหฺมณานํ พาลคมนเมตํ\footnote{พาลานํ คมนเมตํ (สี)} มูฬฺหคมนเมตํ, น นิพฺพิทาย น วิราคาย น นิโรธาย น อุปสมาย น อภิญฺญาย น สมฺโพธาย น นิพฺพานาย สํวตฺตติ.\csromanspacer{} อหญฺจ โข ภิกฺขเว อริยสฺส วินเย อุทยคามินิํ ปฏิปทํ ปญฺญาเปมิ, ยา เอกนฺตนิพฺพิทาย วิราคาย นิโรธาย อุปสมาย อภิญฺญาย สมฺโพธาย นิพฺพานาย สํวตฺตติ.}

\prose{กตมา จ สา ภิกฺขเว อุทยคามินี ปฏิปทา, ยา เอกนฺตนิพฺพิทาย ฯเปฯ นิพฺพานาย สํวตฺตติ, อิธ ภิกฺขเว อริยสาวโก พุทฺเธ อเวจฺจปฺปสาเทน สมนฺนาคโต โหติ “อิติปิ โส ภควา อรหํ สมฺมาสมฺพุทฺโธ ฯเปฯ สตฺถา เทวมนุสฺสานํ พุทฺโธ ภควา”ติ.\csromanspacer{} ธมฺเม ฯเปฯ.\csromanspacer{} สํเฆ ฯเปฯ.\csromanspacer{} อริยกนฺเตหิ สีเลหิ สมนฺนาคโต โหติ อขณฺเฑหิ ฯเปฯ สมาธิสํวตฺตนิเกหิ.\csromanspacer{} อยํ โข สา ภิกฺขเว อุทยคามินี ปฏิปทา เอกนฺตนิพฺพิทาย ฯเปฯ นิพฺพานาย สํวตฺตตีติ.\csromanspacer{}\csromanspacer{}ทุติยํ.}

\csromanheader{2}{3. อานนฺทตฺเถรสุตฺต}
\proseitem{1009}{เอกํ สมยํ อายสฺมา จ อานนฺโท อายสฺมา จ สาริปุตฺโต สาวตฺถิยํ วิหรนฺติ เชตวเน อนาถปิณฺฑิกสฺส อาราเม.\csromanspacer{} อถ โข อายสฺมา สาริปุตฺโต สายนฺหสมยํ ปฏิสลฺลานา วุฏฺฐิโต เยนายสฺมา อานนฺโท เตนุปสงฺกมิ, อุปสงฺกมิตฺวา อายสฺมตา อานนฺเทน สทฺธิํ สมฺโมทิ, สมฺโมทนียํ กถํ สารณียํ วีติสาเรตฺวา เอกมนฺตํ นิสีทิ, เอกมนฺตํ นิสินฺโน โข อายสฺมา สาริปุตฺโต อายสฺมนฺตํ อานนฺทํ เอตทโวจ “กตินํ โข อาวุโส อานนฺท ธมฺมานํ ปหานา, กตินํ ธมฺมานํ สมนฺนาคมนเหตุ เอวมยํ ปชา ภควตา พฺยากตา ‘โสตาปนฺนา อวินิปาตธมฺมา นิยตา สมฺโพธิปรายณา’ติ”.\csromanspacer{} จตุนฺนํ โข อาวุโส ธมฺมานํ ปหานา จตุนฺนํ ธมฺมานํ สมนฺนาคมนเหตุ เอวมยํ ปชา ภควตา พฺยากตา “โสตาปนฺนา อวินิปาตธมฺมา นิยตา สมฺโพธิปรายณา”ติ.}

\prose{\csromanfolio{316}กตเมสํ จตุนฺนํ, ยถารูเปน โข อาวุโส พุทฺเธ อปฺปสาเทน สมนฺนาคโต อสฺสุตวา ปุถุชฺชโน กายสฺส เภทา ปรํ มรณา อปายํ ทุคฺคติํ วินิปาตํ นิรยํ อุปปชฺชติ, ตถารูปสฺส พุทฺเธ อปฺปสาโท น โหติ.\csromanspacer{} ยถารูเปน จ โข อาวุโส พุทฺเธ อเวจฺจปฺปสาเทน สมนฺนาคโต สุตวา อริยสาวโก กายสฺส เภทา ปรํ มรณา สุคติํ สคฺคํ โลกํ อุปปชฺชติ, ตถารูปสฺส พุทฺเธ อเวจฺจปฺปสาโท โหติ “อิติปิ โส ภควา ฯเปฯ สตฺถา เทวมนุสฺสานํ พุทฺโธ ภควา”ติ.}

\prose{ยถารูเปน จ โข อาวุโส ธมฺเม อปฺปสาเทน สมนฺนาคโต อสฺสุตวา ปุถุชฺชโน กายสฺส เภทา ปรํ มรณา อปายํ ทุคฺคติํ วินิปาตํ นิรยํ อุปปชฺชติ, ตถารูปสฺส ธมฺเม อปฺปสาโท น โหติ.\csromanspacer{} ยถารูเปน จ โข อาวุโส ธมฺเม อเวจฺจปฺปสาเทน สมนฺนาคโต สุตวา อริยสาวโก กายสฺส เภทา ปรํ มรณา สุคติํ สคฺคํ โลกํ อุปปชฺชติ, ตถารูปสฺส ธมฺเม อเวจฺจปฺปสาโท โหติ “สฺวากฺขาโต ภควตา ธมฺโม ฯเปฯ วิญฺญูหี”ติ.}

\prose{ยถารูเปน จ โข อาวุโส สํเฆ อปฺปสาเทน สมนฺนาคโต อสฺสุตวา ปุถุชฺชโน กายสฺส เภทา ปรํ มรณา อปายํ ทุคฺคติํ วินิปาตํ นิรยํ อุปปชฺชติ, ตถารูปสฺส สํเฆ อปฺปสาโท น โหติ.\csromanspacer{} ยถารูเปน จ โข อาวุโส สํเฆ อเวจฺจปฺปสาเทน สมนฺนาคโต สุตวา อริยสาวโก กายสฺส เภทา ปรํ มรณา สุคติํ สคฺคํ โลกํ อุปปชฺชติ, ตถารูปสฺส สํเฆ อเวจฺจปฺปสาโท โหติ “สุปฺปฏิปนฺโน ภควโต สาวกสํโฆ ฯเปฯ อนุตฺตรํ ปุญฺญกฺเขตฺตํ โลกสฺสา”ติ.}

\prose{ยถารูเปน จ โข อาวุโส ทุสฺสีเลฺยน สมนฺนาคโต อสฺสุตวา ปุถุชฺชโน กายสฺส เภทา ปรํ มรณา อปายํ ทุคฺคติํ วินิปาตํ นิรยํ อุปปชฺชติ, ตถารูปสฺส ทุสฺสีลฺยํ น โหติ.\csromanspacer{} ยถารูเปหิ จ โข อาวุโส อริยกนฺเตหิ สีเลหิ สมนฺนาคโต สุตวา อริยสาวโก กายสฺส เภทา ปรํ มรณา สุคติํ สคฺคํ โลกํ อุปปชฺชติ, ตถารูปานิ อริยกนฺตานิ สีลานิ โหนฺติ อขณฺฑานิ ฯเปฯ สมาธิสํวตฺตนิกานิ.\csromanspacer{} อิเมสํ โข อาวุโส จตุนฺนํ ธมฺมานํ ปหานา อิเมสํ}

\csromanheader{2}{11. โสตาปตฺติสํยุตฺต}
\csromanmatikaanchor{mtk.234}
\csromantocmarkref{subparagraph}{6-7. มิตฺตามจฺจสุตฺตานิ}{mtk.234}
\bookmark[dest={mtk.234},level=5]{6-7. มิตฺตามจฺจสุตฺตานิ}
\prose{\csromanfolio{317}จตุนฺนํ ธมฺมานํ สมนฺนาคมนเหตุ เอวมยํ ปชา ภควตา พฺยากตา “โสตาปนฺนา อวินิปาตธมฺมา นิยตา สมฺโพธิปรายณา”ติ.\csromanspacer{}\csromanspacer{}ตติยํ.}

\csromanheader{2}{4. ทุคฺคติภยสุตฺต}
\proseitem{1010}{จตูหิ ภิกฺขเว ธมฺเมหิ สมนฺนาคโต อริยสาวโก สพฺพทุคฺคติภยํ สมติกฺกนฺโต โหติ.\csromanspacer{} กตเมหิ จตูหิ, อิธ ภิกฺขเว อริยสาวโก พุทฺเธ อเวจฺจปฺปสาเทน สมนฺนาคโต โหติ “อิติปิ โส ภควา ฯเปฯ สตฺถา เทวมนุสฺสานํ พุทฺโธ ภควา”ติ.\csromanspacer{} ธมฺเม ฯเปฯ.\csromanspacer{} สํเฆ ฯเปฯ.\csromanspacer{} อริยกนฺเตหิ สีเลหิ สมนฺนาคโต โหติ อขณฺเฑหิ ฯเปฯ สมาธิสํวตฺตนิเกหิ.\csromanspacer{} อิเมหิ โข ภิกฺขเว จตูหิ ธมฺเมหิ สมนฺนาคโต อริยสาวโก สพฺพทุคฺคติภยํ สมติกฺกนฺโต โหตีติ.\csromanspacer{}\csromanspacer{}จตุตฺถํ.}

\csromanheader{2}{5. ทุคฺคติวินิปาตภยสุตฺต}
\proseitem{1011}{จตูหิ ภิกฺขเว ธมฺเมหิ สมนฺนาคโต อริยสาวโก สพฺพทุคฺคติวินิปาตภยํ สมติกฺกนฺโต โหติ.\csromanspacer{} กตเมหิ จตูหิ, อิธ ภิกฺขเว อริยสาวโก พุทฺเธ อเวจฺจปฺปสาเทน สมนฺนาคโต โหติ “อิติปิ โส ภควา ฯเปฯ สตฺถา เทวมนุสฺสานํ พุทฺโธ ภควา”ติ.\csromanspacer{} ธมฺเม ฯเปฯ.\csromanspacer{} สํเฆ ฯเปฯ.\csromanspacer{} อริยกนฺเตหิ สีเลหิ สมนฺนาคโต โหติ อขณฺเฑหิ ฯเปฯ สมาธิสํวตฺตนิเกหิ.\csromanspacer{} อิเมหิ โข ภิกฺขเว จตูหิ ธมฺเมหิ สมนฺนาคโต อริยสาวโก สพฺพทุคฺคติวินิปาตภยํ สมติกฺกนฺโต โหตีติ.\csromanspacer{}\csromanspacer{}ปญฺจมํ.}

\csromanheader{2}{6. ปฐมมิตฺตามจฺจสุตฺต}
\csromanmatikaanchor{mtk.235}
\csromantocmarkref{subparagraph}{8-10. เทวจริกสุตฺตานิ}{mtk.235}
\bookmark[dest={mtk.235},level=5]{8-10. เทวจริกสุตฺตานิ}
\proseitem{1012}{เย เต ภิกฺขเว อนุกมฺเปยฺยาถ, เย จ โสตพฺพํ มญฺเญยฺยุํ มิตฺตา วา อมจฺจา วา ญาตี วา สาโลหิตา วา, เต ภิกฺขเว จตูสุ โสตาปตฺติยงฺเคสุ สมาทเปตพฺพา นิเวเสตพฺพา ปติฏฺฐาเปตพฺพา.\csromanspacer{} กตเมสุ จตูสุ, พุทฺเธ อเวจฺจปฺปสาเท สมาทเปตพฺพา นิเวเสตพฺพา ปติฏฺฐาเปตพฺพา “อิติปิ โส ภควา ฯเปฯ สตฺถา เทวมนุสฺสานํ พุทฺโธ ภควา”ติ.\csromanspacer{} ธมฺเม ฯเปฯ.\csromanspacer{} สํเฆ ฯเปฯ.\csromanspacer{} อริยกนฺเตสุ สีเลสุ สมาทเปตพฺพา นิเวเสตพฺพา ปติฏฺฐาเปตพฺพา อขณฺเฑสุ ฯเปฯ สมาธิสํวตฺตนิเกสุ.\csromanspacer{} เย เต ภิกฺขเว อนุกมฺเปยฺยาถ, เย จ โสตพฺพํ มญฺเญยฺยุํ มิตฺตา วา \csromanfolio{318}อมจฺจา วา ญาตี วา สาโลหิตา วา, เต ภิกฺขเว อิเมสุ จตูสุ โสตาปตฺติยงฺเคสุ สมาทเปตพฺพา นิเวเสตพฺพา ปติฏฺฐาเปตพฺพาติ.\csromanspacer{}\csromanspacer{}ฉฏฺฐํ.}

\csromanheader{2}{7. ทุติยมิตฺตามจฺจสุตฺต}
\proseitem{1013}{เย เต ภิกฺขเว อนุกมฺเปยฺยาถ, เย จ โสตพฺพํ มญฺเญยฺยุํ มิตฺตา วา อมจฺจา วา ญาตี วา สาโลหิตา วา, เต ภิกฺขเว จตูสุ โสตาปตฺติยงฺเคสุ สมาทเปตพฺพา นิเวเสตพฺพา ปติฏฺฐาเปตพฺพา.\csromanspacer{} กตเมสุ จตูสุ, พุทฺเธ อเวจฺจปฺปสาเท สมาทเปตพฺพา นิเวเสตพฺพา ปติฏฺฐาเปตพฺพา “อิติปิ โส ภควา ฯเปฯ สตฺถา เทวมนุสฺสานํ พุทฺโธ ภควา”ติ.}

\prose{สิยา ภิกฺขเว จตุนฺนํ มหาภูตานํ อญฺญถตฺตํ ปถวีธาตุยา อาโปธาตุยา เตโชธาตุยา วาโยธาตุยา, น เตฺวว พุทฺเธ อเวจฺจปฺปสาเทน สมนฺนาคตสฺส อริยสาวกสฺส สิยา อญฺญถตฺตํ.\csromanspacer{} ตตฺริทํ อญฺญถตฺตํ, โส วต พุทฺเธ อเวจฺจปฺปสาเทน สมนฺนาคโต อริยสาวโก นิรยํ วา ติรจฺฉานโยนิํ วา เปตฺติวิสยํ วา อุปปชฺชิสฺสตีติ เนตํ ฐานํ วิชฺชติ.\csromanspacer{} ธมฺเม ฯเปฯ.\csromanspacer{} สํเฆ ฯเปฯ.\csromanspacer{} อริยกนฺเตสุ สีเลสุ สมาทเปตพฺพา นิเวเสตพฺพา ปติฏฺฐาเปตพฺพา อขณฺเฑสุ ฯเปฯ สมาธิสํวตฺตนิเกสุ.\csromanspacer{} สิยา ภิกฺขเว จตุนฺนํ มหาภูตานํ อญฺญถตฺตํ ปถวีธาตุยา อาโปธาตุยา เตโชธาตุยา วาโยธาตุยา, น เตฺวว อริยกนฺเตหิ สีเลหิ สมนฺนาคตสฺส อริยสาวกสฺส สิยา อญฺญถตฺตํ.\csromanspacer{} ตตฺริทํ อญฺญถตฺตํ, โส วต อริยกนฺเตหิ สีเลหิ สมนฺนาคโต อริยสาวโก นิรยํ วา ติรจฺฉานโยนิํ วา เปตฺติวิสยํ วา อุปปชฺชิสฺสตีติ เนตํ ฐานํ วิชฺชติ.\csromanspacer{} เย เต ภิกฺขเว อนุกมฺเปยฺยาถ, เย จ โสตพฺพํ มญฺเญยฺยุํ มิตฺตา วา อมจฺจา วา ญาตี วา สาโลหิตา วา, เต ภิกฺขเว อิเมสุ จตูสุ โสตาปตฺติยงฺเคสุ สมาทเปตพฺพา นิเวเสตพฺพา ปติฏฺฐาเปตพฺพาติ.\csromanspacer{}\csromanspacer{}สตฺตมํ.}

\csromanheader{2}{8. ปฐมเทวจาริกสุตฺต}
\proseitem{1014}{สาวตฺถินิทานํ.\csromanspacer{} อถ โข อายสฺมา มหาโมคฺคลฺลาโน เสยฺยถาปิ นาม พลวา ปุริโส สมิญฺชิตํ วา พาหํ ปสาเรยฺย, ปสาริตํ}

\csromanheader{2}{11. โสตาปตฺติสํยุตฺต}
\prose{\csromanfolio{319}วา พาหํ สมิญฺเชยฺย, เอวเมว เชตวเน อนฺตรหิโต เทเวสุ ตาวติํเสสุ ปาตุรโหสิ.\csromanspacer{} อถ โข สมฺพหุลา ตาวติํสกายิกา เทวตาโย เยนายสฺมา มหาโมคฺคลฺลาโน เตนุปสงฺกมิํสุ, อุปสงฺกมิตฺวา อายสฺมนฺตํ มหาโมคฺคลฺลานํ อภิวาเทตฺวา เอกมนฺตํ อฏฺฐํสุ, เอกมนฺตํ ฐิตา โข ตา เทวตาโย อายสฺมา มหาโมคฺคลฺลาโน เอตทโวจ\csromandash{}}

\prose{สาธุ โข อาวุโส พุทฺเธ อเวจฺจปฺปสาเทน สมนฺนาคมนํ โหติ “อิติปิ โส ภควา ฯเปฯ สตฺถา เทวมนุสฺสานํ พุทฺโธ ภควา”ติ.\csromanspacer{} พุทฺเธ อเวจฺจปฺปสาเทน สมนฺนาคมนเหตุ โข อาวุโส เอวมิเธกจฺเจ สตฺตา กายสฺส เภทา ปรํ มรณา สุคติํ สคฺคํ โลกํ อุปปชฺชนฺติ.\csromanspacer{} สาธุ โข อาวุโส ธมฺเม ฯเปฯ.\csromanspacer{} สํเฆ ฯเปฯ.\csromanspacer{} สาธุ โข อาวุโส อริยกนฺเตหิ สีเลหิ สมนฺนาคมนํ โหติ อขณฺเฑหิ ฯเปฯ สมาธิสํวตฺตนิเกหิ.\csromanspacer{} อริยกนฺเตหิ สีเลหิ สมนฺนาคมนเหตุ โข อาวุโส เอวมิเธกจฺเจ สตฺตา กายสฺส เภทา ปรํ มรณา สุคติํ สคฺคํ โลกํ อุปปชฺชนฺตีติ.}

\prose{สาธุ โข มาริส โมคฺคลฺลาน, พุทฺเธ อเวจฺจปฺปสาเทน สมนฺนาคมนํ โหติ “อิติปิ โส ภควา ฯเปฯ สตฺถา เทวมนุสฺสานํ พุทฺโธ ภควา”ติ.\csromanspacer{} พุทฺเธ อเวจฺจปฺปสาเทน สมนฺนาคมนเหตุ โข มาริส โมคฺคลฺลาน เอวมิเธกจฺเจ สตฺตา กายสฺส เภทา ปรํ มรณา สุคติํ สคฺคํ โลกํ อุปปชฺชนฺติ.\csromanspacer{} สาธุ โข มาริส โมคฺคลฺลาน ธมฺเม ฯเปฯ.\csromanspacer{} สํเฆ ฯเปฯ.\csromanspacer{} อริยกนฺเตหิ สีเลหิ สมนฺนาคมนํ โหติ อขณฺเฑหิ ฯเปฯ สมาธิสํวตฺตนิเกหิ.\csromanspacer{} อริยกนฺเตหิ สีเลหิ สมนฺนาคมนเหตุ โข มาริส โมคฺคลฺลาน เอวมิเธกจฺเจ สตฺตา กายสฺส เภทา ปรํ มรณา สุคติํ สคฺคํ โลกํ อุปปชฺชนฺตีติ.\csromanspacer{}\csromanspacer{}อฏฺฐมํ.}

\csromanheader{2}{9. ทุติยเทวจาริกสุตฺต}
\proseitem{1015}{สาวตฺถินิทานํ.\csromanspacer{} อถ โข อายสฺมา มหาโมคฺคลฺลาโน เสยฺยถาปิ นาม พลวา ปุริโส สมิญฺชิตํ วา พาหํ ปสาเรยฺย, ปสาริตํ วา พาหํ สมิญฺเชยฺย, เอวเมว เชตวเน อนฺตรหิโต เทเวสุ ตาวติํเสสุ ปาตุรโหสิ.\csromanspacer{} อถ โข สมฺพหุลา ตาวติํสกายิกา เทวตาโย เยนายสฺมา มหาโมคฺคลฺลาโน เตนุปสงฺกมิํสุ, อุปสงฺกมิตฺวา อายสฺมนฺตํ มหาโมคฺคลฺลานํ อภิวาเทตฺวา เอกมนฺตํ อฏฺฐํสุ, \csromanfolio{320}เอกมนฺตํ ฐิตา โข ตา เทวตาโย ยสฺมา มหาโมคฺคลฺลาโน เอตทโวจ\csromandash{}}

\prose{สาธุ โข อาวุโส พุทฺเธ อเวจฺจปฺปสาเทน สมนฺนาคมนํ โหติ “อิติปิ โส ภควา ฯเปฯ สตฺถา เทวมนุสฺสานํ พุทฺโธ ภควา”ติ.\csromanspacer{} พุทฺเธ อเวจฺจปฺปสาเทน สมนฺนาคมนเหตุ โข อาวุโส เอวมิเธกจฺเจ สตฺตา กายสฺส เภทา ปรํ มรณา สุคติํ สคฺคํ โลกํ อุปปนฺนา.\csromanspacer{} สาธุ โข อาวุโส ธมฺเม ฯเปฯ.\csromanspacer{} สํเฆ ฯเปฯ.\csromanspacer{} อริยกนฺเตหิ สีเลหิ สมนฺนาคมนํ โหติ อขณฺเฑหฺ ฯเปฯ สมาธิสํวตฺตนิเกหิ.\csromanspacer{} อริยกนฺเตหิ สีเลหิ สมนฺนาคมนเหตุ โข อาวุโส เอวมิเธกจฺเจ สตฺตา กายสฺส เภทา ปรํ มรณา สุคติํ สคฺคํ โลกํ อุปปนฺนาติ.}

\prose{สาธุ โข มาริส โมคฺคลฺลาน, พุทฺเธ อเวจฺจปฺปสาเทน สมนฺนาคมนํ โหติ “อิติปิ โส ภควา ฯเปฯ สตฺถา เทวมนุสฺสานํ พุทฺโธ ภควา”ติ.\csromanspacer{} พุทฺเธ อเวจฺจปฺปสาเทน สมนฺนาคมนเหตุ โข มาริส โมคฺคลฺลาน เอวมิเธกจฺเจ สตฺตา กายสฺส เภทา ปรํ มรณา สุคติํ สคฺคํ โลกํ อุปปนฺนา.\csromanspacer{} สาธุ โข มาริส โมคฺคลฺลาน ธมฺเม ฯเปฯ.\csromanspacer{} สํเฆ ฯเปฯ.\csromanspacer{} อริยกนฺเตหิ สีเลหิ สมนฺนาคมนํ โหติ อขณฺเฑหิ ฯเปฯ สมาธิสํวตฺตนิเกหิ.\csromanspacer{} อริยกนฺเตหิ สีเลหิ สมนฺนาคมนเหตุ โข มาริส โมคฺคลฺลาน เอวมิเธกจฺเจ สตฺตา กายสฺส เภทา ปรํ มรณา สุคติํ สคฺคํ โลกํ อุปปนฺนาติ.\csromanspacer{}\csromanspacer{}นวมํ.}

\csromanheader{2}{10. ตติยเทวจาริกสุตฺต}
\proseitem{1016}{อถ โข ภควา เสยฺยถาปิ นาม พลวา ปุริโส สมิญฺชิตํ วา พาหํ ปสาเรยฺย, ปสาริตํ วา พาหํ สมิญฺเชยฺย, เอวเมว เชตวเน อนฺตรหิโต เทเวสุ ตาวติํเสสุ ปาตุรโหสิ.\csromanspacer{} อถ โข สมฺพหุลา ตาวติํสกายิกา เทวตาโย เยน ภควา เตนุปสงฺกมิํสุ, อุปสงฺกมิตฺวา ภควนฺตํ อภิวาเทตฺวา เอกมนฺตํ อฏฺฐํสุ, เอกมนฺตํ ฐิตา โข ตา เทวตาโย ภควา เอตทโวจ\csromandash{}}

\prose{สาธุ โข อาวุโส พุทฺเธ อเวจฺจปฺปสาเทน สมนฺนาคมนํ โหติ “อิติปิ โส ภควา ฯเปฯ สตฺถา เทวมนุสฺสานํ พุทฺโธ ภควา”ติ, พุทฺเธ อเวจฺจปฺปสาเทน สมนฺนาคมนเหตุ โข อาวุโส เอวมิเธกจฺเจ สตฺตา โสตาปนฺนา อวินิปาตธมฺมา นิยตา สมฺโพธิปรายณา.\csromanspacer{} สาธุ โข}

\csromanheader{2}{11. โสตาปตฺติสํยุตฺต}
\csromanmatikaanchor{mtk.236}
\csromanmatikaanchor{mtk.237}
\csromantocmarknopage{subparagraph}{3. สรณานิวคฺค}{mtk.236}
\bookmark[dest={mtk.236},level=5]{3. สรณานิวคฺค}
\csromantocmarkref{subparagraph}{1-2. มหานามสุตฺตานิ}{mtk.237}
\bookmark[dest={mtk.237},level=5]{1-2. มหานามสุตฺตานิ}
\prose{\csromanfolio{321}อาวุโส ธมฺเม ฯเปฯ.\csromanspacer{} สํเฆ ฯเปฯ.\csromanspacer{} อริยกนฺเตหิ สีเลหิ สมนฺนาคมนํ โหติ อขณฺเฑหิ ฯเปฯ สมาธิสํวตฺตนิเกหิ.\csromanspacer{} อริยกนฺเตหิ สีเลหิ สมนฺนาคมนเหตุ โข อาวุโส เอวมิเธกจฺเจ สตฺตา โสตาปนฺนา อวินิปาตธมฺมา นิยตา สมฺโพธิปรายณาติ.}

\prose{สาธุ โข มาริส พุทฺเธ อเวจฺจปฺปสาเทน สมนฺนาคมนํ โหติ “อิติปิ โส ภควา ฯเปฯ สตฺถา เทวมนุสฺสานํ พุทฺโธ ภควา”ติ, พุทฺเธ อเวจฺจปฺปสาเทน สมนฺนาคมนเหตุ โข มาริส เอวมยํ ปชา โสตาปนฺนา อวินิปาตธมฺมา นิยตา สมฺโพธิปรายณา.\csromanspacer{} สาธุ โข มาริส ธมฺเม ฯเปฯ.\csromanspacer{} สํเฆ ฯเปฯ.\csromanspacer{} อริยกนฺเตหิ สีเลหิ สมนฺนาคมนํ โหติ อขณฺเฑหิ ฯเปฯ สมาธิสํวตฺตนิเกหิ.\csromanspacer{} อริยกนฺเตหิ สีเลหิ สมนฺนาคมนเหตุ โข มาริส เอวมยํ ปชา โสตาปนฺนา อวินิปาตธมฺมา นิยตา สมฺโพธิปรายณาติ.\csromanspacer{}\csromanspacer{}ทสมํ.\csromanspacer{} ราชการามวคฺโค ทุติโย.}

\titlehead{\textbf{ตสฺสุทฺทานํ}}
\csromangathameasure{สหสฺสพฺราหฺมณานนฺท, ทุคฺคติ อปเร ทุเว. \\ มิตฺตามจฺจา ทุเว วุตฺตา, ตโย จ เทวจาริกาติ.}
\csromangathagroup{\csromangathastanza{สหสฺสพฺราหฺมณานนฺท, ทุคฺคติ อปเร ทุเว. \\ มิตฺตามจฺจา ทุเว วุตฺตา, ตโย จ เทวจาริกาติ.}}

\csromanheader{6}{3. สรณานิวคฺค}
\csromanheader{2}{1. ปฐมมหานามสุตฺต}
\proseitem{1017}{เอวํ เม สุตํ\csromandash{}เอกํ สมยํ ภควา สกฺเกสุ วิหรติ กปิลวตฺถุสฺมิํ นิโคฺรธาราเม.\csromanspacer{} อถ โข มหานาโม สกฺโก เยน ภควา เตนุปสงฺกมิ, อุปสงฺกมิตฺวา ภควนฺตํ อภิวาเทตฺวา เอกมนฺตํ นิสีทิ, เอกมนฺตํ นิสินฺโน โข มหานาโม สกฺโก ภควนฺตํ เอตทโวจ “อิทํ ภนฺเต กปิลวตฺถุ อิทฺธญฺเจว ผีตญฺจ พาหุชญฺญํ อากิณฺณมนุสฺสํ สมฺพาธพฺยูหํ, โส ขฺวาหํ ภนฺเต ภควนฺตํ วา ปยิรุปาสิตฺวา มโนภาวนีเย วา ภิกฺขู สายนฺหสมยํ กปิลวตฺถุํ ปวิสนฺโต ภนฺเตนปิ\footnote{วิพฺภนฺเตนปิ (สี), ภมนฺเตนปิ (ก)} หตฺถินา \csromanfolio{322}สมาคจฺฉามิ, ภนฺเตนปิ อสฺเสน สมาคจฺฉามิ, ภนฺเตนปิ รเถน สมาคจฺฉามิ, ภนฺเตนปิ สกเฏน สมาคจฺฉามิ, ภนฺเตนปิ ปุริเสน สมาคจฺฉามิ.\csromanspacer{} ตสฺส มยฺหํ ภนฺเต ตสฺมิํ สมเย มุสฺสเตว\footnote{มุสเตว (?)} ภควนฺตํ อารพฺภ สติ, มุสฺสติ\footnote{มุสติ (?)} ธมฺมํ อารพฺภ สติ, มุสฺสติ สํฆํ อารพฺภ สติ.\csromanspacer{} ตสฺส มยฺหํ ภนฺเต เอวํ โหติ “อิมมฺหิ จาหํ สมเย กาลํ กเรยฺยํ, กา มยฺหํ คติ, โก อภิสมฺปราโย”ติ.}

\prose{มา ภายิ มหานาม, มา ภายิ มหานาม, อปาปกํ เต มรณํ ภวิสฺสติ อปาปิกา กาลํกิริยา\footnote{กาลกิริยา (สี, สฺยา, กํ)}.\csromanspacer{} ยสฺส กสฺสจิ มหานาม ทีฆรตฺตํ สทฺธาปริภาวิตํ จิตฺตํ, สีลปริภาวิตํ จิตฺตํ, สุตปริภาวิตํ จิตฺตํ, จาคปริภาวิตํ จิตฺตํ, ปญฺญาปริภาวิตํ จิตฺตํ, ตสฺส โย หิ ขฺวายํ กาโย รูปี จาตุมหาภูติโก\footnote{จาตุมฺมหาภูติโก (สี, สฺยา, กํ) 5. สิคาลา (สี, สฺยา, กํ, อิ)} มาตาเปตฺติกสมฺภโว โอทนกุมฺมาสูปจโย อนิจฺจุจฺฉาทนปริมทฺทนเภทน- วิทฺธํสนธมฺโม.\csromanspacer{} ตํ อิเธว กากา วา ขาทนฺติ, คิชฺฌา วา ขาทนฺติ, กุลลา วา ขาทนฺติ, สุนขา วา ขาทนฺติ, สิงฺคาลา5 วา ขาทนฺติ, วิวิธา วา ปาณกชาตา ขาทนฺติ.\csromanspacer{} ยญฺจ ขฺวสฺส จิตฺตํ ทีฆรตฺตํ สทฺธาปริภาวิตํ ฯเปฯ ปญฺญาปริภาวิตํ, ตํ อุทฺธคามิ โหติ วิเสสคามิ.}

\prose{เสยฺยถาปิ มหานาม ปุริโส สปฺปิกุมฺภํ วา เตลกุมฺภํ วา คมฺภีรํ อุทกรหทํ โอคาหิตฺวา ภินฺเทยฺย.\csromanspacer{} ตตฺร ยา อสฺส สกฺขรา วา กฐลา\footnote{กถลา (อิ, ก)} วา, สา อโธคามี อสฺส.\csromanspacer{} ยญฺจ ขฺวสฺส ตตฺร สปฺปิ วา เตลํ วา, ตํ อุทฺธคามิ อสฺส วิเสสคามิ.\csromanspacer{} เอวเมว โข มหานาม ยสฺส กสฺสจิ ทีฆรตฺตํ สทฺธาปริภาวิตํ จิตฺตํ ฯเปฯ ปญฺญาปริภาวิตํ จิตฺตํ.\csromanspacer{} ตสฺส โย หิ ขฺวายํ กาโย รูปี จาตุมหาภูติโก มาตาเปตฺติกสมฺภโว โอทนกุมฺมาสูปวจโย อนิจฺจุจฺฉาทนปริมทฺทนเภทน- วิทฺธํสนธมฺโม.\csromanspacer{} ตํ อิเธว กากา วา ขาทนฺติ, คิชฺฌา วา ขาทนฺติ, กุลลา วา ขาทนฺติ, สุนขา วา ขาทนฺติ, สิงฺคาลา วา ขาทนฺติ, วิวิธา วา ปาณกชาตา ขาทนฺติ.\csromanspacer{} ยญฺจ ขฺวสฺส จิตฺตํ ทีฆรตฺตํ สทฺธาปริภาวิตํ ฯเปฯ ปญฺญาปริภาวิตํ, ตํ อุทฺธคามิ โหติ วิเสสคามิ.\csromanspacer{} ตุยฺหํ โข ปน มหานาม ทีฆรตฺตํ สทฺธาปริภาวิตํ จิตฺตํ ฯเปฯ ปญฺญาปริภาวิตํ จิตฺตํ, มา ภายิ}

\csromanheader{2}{11. โสตาปตฺติสํยุตฺต}
\prose{\csromanfolio{323}มหานาม, มา ภายิ มหานาม, อปาปกํ เต มรณํ ภวิสฺสติ อปาปิกา กาลํกิริยาติ.\csromanspacer{} ปฐมํ.}

\csromanheader{2}{2. ทุติยมหานามสุตฺต}
\proseitem{1018}{เอวํ เม สุตํ\csromandash{}เอกํ สมยํ ภควา สกฺเกสุ วิหรติ กปิลวตฺถุสฺมิํ นิโคฺรธาราเม.\csromanspacer{} อถ โข มหานาโม สกฺโก เยน ภควา เตนุปสงฺกมิ, อุปสงฺกมิตฺวา ภควนฺตํ อภิวาเทตฺวา เอกมนฺตํ นิสีทิ, เอกมนฺตํ นิสินฺโน โข มหานาโม สกฺโก ภควนฺตํ เอตทโวจ “อิทํ ภนฺเต กปิลวตฺถุ อิทฺธญฺเจว ผีตญฺจ พาหุชญฺญํ อากิณฺณมนุสฺสํ สมฺพาธพฺยูหํ, โส ขฺวาหํ ภนฺเต ภควนฺตํ วา ปยิรุปาสิตฺวา มโนภาวนีเย วา ภิกฺขู สายนฺหสมยํ กปิลวตฺถุํ ปวิสนฺโต ภนฺเตนปิ หตฺถินา สมาคจฺฉามิ, ภนฺเตนปิ อสฺเสน สมาคจฺฉามิ, ภนฺเตนปิ รเถน สมาคจฺฉามิ, ภนฺเตนปิ สกเฏน สมาคจฺฉามิ, ภนฺเตนปิ ปุริเสน สมาคจฺฉามิ, ตสฺส มยฺหํ ภนฺเต ตสฺมิํ สมเย มุสฺสเตว ภควนฺตํ อารพฺภ สติ, มุสฺสติ ธมฺมํ อารพฺภ สติ, มุสฺสติ สํฆํ อารพฺภ สติ.\csromanspacer{} ตสฺส มยฺหํ ภนฺเต เอวํ โหติ “อิมมฺหิ จาหํ สมเย กาลํกเรยฺยํ, กา มยฺหํ คติ, โก อภิสมฺปราโย”ติ.}

\prose{มา ภายิ มหานาม, มา ภายิ มหานาม, อปาปกํ เต มรณํ ภวิสฺสติ อปาปิกา กาลํกิริยา.\csromanspacer{} จตูหิ โข มหานาม ธมฺเมหิ สมนฺนาคโต อริยสาวโก นิพฺพานนินฺโน โหติ นิพฺพานโปโณ นิพฺพานปพฺภาโร.\csromanspacer{} กตเมหิ จตูหิ, อิธ มหานาม อริยสาวโก พุทฺเธ อเวจฺจปฺปสาเทน สมนฺนาคโต โหติ “อิติปิ โส ภควา ฯเปฯ สตฺถา เทวมนุสฺสานํ พุทฺโธ ภควา”ติ.\csromanspacer{} ธมฺเม ฯเปฯ.\csromanspacer{} สํเฆ ฯเปฯ.\csromanspacer{} อริยกนฺเตหิ สีเลหิ สมนฺนาคโต โหติ อขณฺเฑหิ ฯเปฯ สมาธิสํวตฺตนิเกหิ.}

\prose{เสยฺยถาปิ มหานาม รุกฺโข ปาจีนนินฺโน ปาจีนโปโณ ปาจีนปพฺภาโร, โส มูลจฺฉินฺโน กตเมน ปปเตยฺยาติ.\csromanspacer{} เยน ภนฺเต นินฺโน เยน โปโณ เยน ปพฺภาโรติ.\csromanspacer{} เอวเมว โข มหานาม อิเมหิ จตูหิ ธมฺเมหิ สมนฺนาคโต อริยสาวโก นิพฺพานนินฺโน โหติ นิพฺพานโปโณ นิพฺพานปพฺภาโรติ.\csromanspacer{} ทุติยํ.}

\csromanmatikaanchor{mtk.238}
\csromantocmarkref{subparagraph}{3. โคธสกฺกสุตฺต}{mtk.238}
\bookmark[dest={mtk.238},level=5]{3. โคธสกฺกสุตฺต}
\csromanheader{6}{3. โคธสกฺกสุตฺต}
\proseitem{1019}{\csromanfolio{324}กปิลวตฺถุนิทานํ.\csromanspacer{} อถ โข มหานาโม สกฺโก เยน โคธา สกฺโก เตนุปสงฺกมิ, อุปสงฺกมิตฺวา โคธํ สกฺกํ เอตทโวจ “กติหิ\footnote{กตีหิ (อิ, ก) รูปสิทฺธิ โอโลเกตพฺพา.} ตฺวํ โคเธ ธมฺเมหิ สมนฺนาคตํ โสตาปนฺนปุคฺคลํ อาชานาสิ อวินิปาตธมฺมํ นิยตํ สมฺโพธิปรายณนฺติ.}

\prose{ตีหิ ขฺวาหํ มหานาม ธมฺเมหิ สมนฺนาคตํ โสตาปนฺนปุคฺคลํ อาชานามิ อวินิปาตธมฺมํ นิยตํ สมฺโพธิปรายณํ.\csromanspacer{} กตเมหิ ตีหิ, อิธ มหานาม อริยสาวโก พุทฺเธ อเวจฺจปฺปสาเทน สมนฺนาคโต โหติ “อิติปิ โส ภควา ฯเปฯ สตฺถา เทวมนุสฺสานํ พุทฺโธ ภควา”ติ.\csromanspacer{} ธมฺเม ฯเปฯ.\csromanspacer{} สํเฆ อเวจฺจปฺปสาเทน สมนฺนาคโต โหติ “สุปฺปฏิปนฺโน ภควโต สาวกสํโฆ ฯเปฯ อนุตฺตรํ ปุญฺญกฺเขตฺตํ โลกสฺสา”ติ.\csromanspacer{} อิเมหิ ขฺวาหํ มหานาม ตีหิ ธมฺเมหิ สมนฺนาคตํ โสตาปนฺนปุคฺคลํ อาชานามิ อวินิปาตธมฺมํ นิยตํ สมฺโพธิปรายณํ.}

\prose{ตฺวํ ปน มหานาม กติหิ ธมฺเมหิ สมนฺนาคตํ โสตาปนฺนปุคฺคลํ อาชานาสิ อวินิปาตธมฺมํ นิยตํ สมฺโพธิปรายณนฺติ.\csromanspacer{} จตูหิ ขฺวาหํ โคเธ ธมฺเมหิ สมนฺนาคตํ โสตาปนฺนปุคฺคลํ อาชานามิ อวินิปาตธมฺมํ นิยตํ สมฺโพธิปรายณํ.\csromanspacer{} กตเมหิ จตูหิ, อิธ โคเธ อริยสาวโก พุทฺเธ อเวจฺจปฺปสาเทน สมนฺนาคโต โหติ “อิติปิ โส ภควา ฯเปฯ สตฺถา เทวมนุสฺสานํ พุทฺโธ ภควา”ติ.\csromanspacer{} ธมฺเม ฯเปฯ.\csromanspacer{} สํเฆ ฯเปฯ.\csromanspacer{} อริยกนฺเตหิ สีเลหิ สมนฺนาคโต โหติ อขณฺเฑหิ ฯเปฯ สมาธิสํวตฺตนิเกหิ.\csromanspacer{} อิเมหิ ขฺวาหํ โคเธ จตูหิ ธมฺเมหิ สมนฺนาคตํ โสตาปนฺนปุคฺคลํ อาชานามิ อวินิปาตธมฺมํ นิยตํ สมฺโพธิปรายณนฺติ.}

\prose{อาคเมหิ ตฺวํ มหานาม, อาคเมหิ ตฺวํ มหานาม, ภควาว เอตํ ชาเนยฺย เอเตหิ ธมฺเมหิ สมนฺนาคตํ วา อสมนฺนาคตํ วาติ, อายาม โคเธ, เยน ภควา เตนุปสงฺกเมยฺยาม, อุปสงฺกมิตฺวา ภควโต เอตมตฺถํ อาโรเจสฺสามาติ.\csromanspacer{} อถ โข มหานาโม สกฺโก โคธา จ สกฺโก เยน ภควา เตนุปสงฺกมิํสุ, อุปสงฺกมิตฺวา ภควนฺตํ อภิวาเทตฺวา เอกมนฺตํ นิสีทิํสุ, เอกมนฺตํ นิสินฺโน โข มหานาโม สกฺโก ภควนฺตํ เอตทโวจ\csromandash{}}

\csromanheader{2}{11. โสตาปตฺติสํยุตฺต}
\prose{\csromanfolio{325}อิธาหํ ภนฺเต เยน โคธา สกฺโก เตนุปสงฺกมิํ, อุปสงฺกมิตฺวา โคธํ สกฺกํ เอตทโวจํ “กติหิ ตฺวํ โคเธ ธมฺเมหิ สมนฺนาคตํ โสตาปนฺนปุคฺคลํ อาชานาสิ อวินิปาตธมฺมํ นิยตํ สมฺโพธิปรายณํ”.\csromanspacer{} เอวํ วุตฺเต พนฺเต โคธา สกฺโก มํ เอตทโวจ\csromandash{}}

\prose{ตีหิ ขฺวาหํ มหานาม ธมฺเมหิ สมนฺนาคตํ โสตาปนฺนปุคฺคลํ อาชานามิ อวินิปาตธมฺมํ นิยตํ สมฺโพธิปรายณํ.\csromanspacer{} กตเมหิ ตีหิ, อิธ มหานาม อริยสาวโก พุทฺเธ อเวจฺจปฺปสาเทน สมนฺนาคโต โหติ “อิติปิ โส ภควา ฯเปฯ สตฺถา เทวมนุสฺสานํ พุโธ ภควา”ติ.\csromanspacer{} ธมฺเม ฯเปฯ.\csromanspacer{} สํเฆ อเวจฺจปฺปสาเทน สมนฺนาคโต โหติ “สุปฺปฏิปนฺโน ภควโต สาวกสํโฆ ฯเปฯ อนุตฺตรํ ปุญฺญกฺเขตฺตํ โลกสฺสา”ติ.\csromanspacer{} อิเมหิ ขฺวาหํ มหานาม ตีหิ ธมฺเมหิ สมนฺนาคตํ โสตาปนฺนปุคฺคลํ อาชานามิ อวินิปาตธมฺมํ นิยตํ สมฺโพธิปรายณํ.\csromanspacer{} ตฺวํ ปน มหานาม กตเมหิ ธมฺเมหิ สมนฺนาคตํ โสตาปนฺนปุคฺคลํ อาชานาสิ อวินิปาตธมฺมํ นิยตํ สมฺโพธิปรายณนฺติ.}

\prose{เอวํ วุตฺตาหํ ภนฺเต โคธํ สกฺกํ เอตทโวจํ “จตูหิ ขฺวาหํ โคเธ ธมฺเมหิ สมนฺนาคตํ โสตาปนฺนปุคฺคลํ อาชานามิ อวินิปาตธมฺมํ นิยตํ สมฺโพธิปรายณํ.\csromanspacer{} กตเมหิ จตูหิ, อิธ โคเธ อริยสาวโก พุทฺเธ อเวจฺจปฺปสาเทน สมนฺนาคโต โหติ ‘อิติปิ โส ภควา ฯเปฯ สตฺถา เทวมนุสฺสานํ พุทฺโธ ภควา’ติ.\csromanspacer{} ธมฺเม ฯเปฯ.\csromanspacer{} สํเฆ ฯเปฯ.\csromanspacer{} อริยกนฺเตหิ สีเลหิ สมนฺนาคโต โหติ อขณฺเฑหิ ฯเปฯ สมาธิสํวตฺตนิเกหิ.\csromanspacer{} อิเมหิ ขฺวาหํ โคเธ จตูหิ ธมฺเมหิ สมนฺนาคตํ โสตาปนฺนปุคฺคลํ อาชานามิ อวินิปาตธมฺมํ นิยตํ สมฺโพธิปรายณนฺ”ติ.}

\prose{เอวํ วุตฺเต ภนฺเต โคธา สกฺโก มํ เอตทโวจ “อาคเมหิ ตฺวํ มหานาม, อาคเมหิ ตฺวํ มหานาม, ภควาว เอตํ ชาเนยฺย เอเตหิ ธมฺเมหิ สมนฺนาคตํ วา อสมนฺนาคตํ วา”ติ, อิธ ภนฺเต โกจิเทว ธมฺโม สมุปฺปาโท อุปฺปชฺเชยฺย, เอกโต อสฺส ภควา เอกโต ภิกฺขุสํโฆ จ, เยเนว ภควา, เตเนวาหํ อสฺสํ, เอวํ ปสนฺนํ มํ ภนฺเต ภควา ธาเรตุ, อิธ ภนฺเต โกจิเทว ภมฺโม สมุปฺปาโท อุปฺปชฺเชยฺย, เอกโต อสฺส ภควา, เอกโต ภิกฺขุสํโฆ ภิกฺขุนิสํโฆ จ, เยเนว ภควา, เตเนวาหํ อสฺสํ, เอวํ ปสนฺนํ มํ ภนฺเต ภควา ธาเรตุ.\csromanspacer{} อิธ ภนฺเต โกจิเทว ธมฺโม สมุปฺปาโท อุปฺปชฺเชยฺย, เอกโต อสฺส ภควา, เอกโต}

\csromanmatikaanchor{mtk.239}
\csromantocmarkref{subparagraph}{4-5. สรณานิสกฺกสุตฺตานิ}{mtk.239}
\bookmark[dest={mtk.239},level=5]{4-5. สรณานิสกฺกสุตฺตานิ}
\prose{\csromanfolio{326}ภิกฺขุสํโฆ ภิกฺขุนิสํโฆ จ อุปาสกา จ, เยเนว ภควา, เตเนวาหํ อสฺสํ, เอวํ ปสนฺนํ มํ ภนฺเต ภควา ธาเรตุ, อิธ ภนฺเต โกจิเทว ธมฺโม สมุปฺปาโท อุปฺปชฺเชยฺย, เอกโต อสฺส ภควา, เอกโต ภิกฺขุสํโฆ ภิกฺขุนิสํโฆ อุปาสกา อุปาสิกาโย จ, เยเนว ภควา, เตเนวาหํ อสฺสํ, เอวํ ปสนฺนํ มํ ภนฺเต ภควา ธาเรตุ, อิธ ภนฺเต โกจิเทว ธมฺโม สมุปฺปาโท อุปฺปชฺเชยฺย, เอกโต อสฺส ภควา, เอกโต ภิกฺขุสํโฆ ภิกฺขุนิสํโฆ อุปาสกา อุปาสิกาโย สเทวโก จ โลโก สมารโก สพฺรหฺมโก สสฺสมณพฺราหฺมณี ปชา สเทวมนุสฺสา, เยเนว ภควา, เตเนวาหํ อสฺสํ, เอวํ ปสนฺนํ มํ ภนฺเต ภควา ธาเรตูติ.\csromanspacer{} เอวํวาที ตฺวํ โคเธ มหานามํ สกฺกํ กิํ วเทสีติ.\csromanspacer{} เอวํวาทาหํ ภนฺเต มหานามํ สกฺกํ น กิญฺจิ วทามิ อญฺญตฺร กลฺยาณา อญฺญตฺร กุสลาติ.\csromanspacer{} ตติยํ.}

\csromanheader{2}{4. ปฐมสรณานิสกฺกสุตฺต}
\proseitem{1020}{กปิลวตฺถุนิทานํ.\csromanspacer{} เตน โข ปน สมเยน สรณานิ\footnote{สรกานิ (สี, สฺยา, กํ, อิ)} สกฺโก กาลงฺกโต โหติ, โส ภควตา พฺยากโต “โสตาปนฺโน อวินิปาตธมฺโม นิยโต สมฺโพธิปรายโณ”ติ.\csromanspacer{} ตตฺร สุทํ สมฺพหุลา สกฺกา สงฺคมฺม สมาคมฺม อุชฺฌายนฺติ ขียนฺติ วิปาเจนฺติ “อจฺฉริยํ วต โภ, อพฺภุตํ วต โภ, เอตฺถ ทานิ โก น โสตาปนฺโน ภวิสฺสติ, ยตฺร หิ นาม สรณานิ สกฺโก กาลงฺกโต, โส ภควตา พฺยากโต ‘โสตาปนฺโน อวินิปาตธมฺโม นิยโต สมฺโพธิปรายโณ’ติ.\csromanspacer{} สรณานิ สกฺโก สิกฺขาทุพฺพลฺยมาปาทิ, มชฺชปานํ อปายี”ติ.}

\prose{อถ โข มหานาโม สกฺโก เยน ภควา เตนุปสงฺกมิ, อุปสงฺกมิตฺวา ภควนฺตํ อภิวาเทตฺวา เอกมนฺตํ นิสีทิ, เอกมนฺตํ นิสินฺโน โข มหานาโม สกฺโก ภควนฺตํ เอตทโวจ\csromandash{}อิธ ภนฺเต สรณานิ สกฺโก กาลงฺกโต, โส ภควตา พฺยากโต “โสตาปนฺโน อวินิปาตธมฺโม นิยโต สมฺโพธิปรายโณ”ติ.\csromanspacer{} ตตฺร สุทํ ภนฺเต สมฺพหุลา สกฺกา สงฺคมฺม สมาคมฺม อุชฺฌายนฺติ ขียนฺติ วิปาเจนฺติ “อจฺฉริยํวต โภ, อพฺภุตํ วต โภ, เอตฺถ ทานิ โก น โสตาปนฺโน ภวิสฺสติ, ยตฺร หิ นาม สรณานิ สกฺโก กาลงฺกโต, โส ภควตา}

\csromanheader{2}{11. โสตาปตฺติสํยุตฺต}
\prose{\csromanfolio{327}พฺยากโต ‘โสตาปนฺโน อวินิปาตธมฺโม นิยโต สมฺโพธิปรายโณ’ติ, สรณานิ สกฺโก สิกฺขาทุพฺพลฺยมาปาทิ, มชฺชปานํ อปายี”ติ.}

\prose{โย โส มหานาม ทีฆรตฺตํ อุปาสโก พุทฺธํ สรณํ คโต, ธมฺมํ สรณํ คโต, สํฆํ สรณํ คโต, โส กถํ วินิปาตํ คจฺเฉยฺย.\csromanspacer{} ยํ หิ ตํ มหานาม สมฺมา วทมาโน วเทยฺย “ทีฆรตฺตํ อุปาสโก พุทฺธํ สรณํ คโต, ธมฺมํ สรณํ คโต, สํฆํ สรณํ คโต”ติ.\csromanspacer{} สรณานิ สกฺกํ สมฺมา วทมาโน วเทยฺย, สรณานิ มหานาม สกฺโก ทีฆรตฺตํ อุปาสโก พุทฺธํ สรณํ คโต, ธมฺมํ สรณํ คโต, สํฆํ สรณํ คโต, โส กถํ วินิปาตํ คจฺเฉยฺย.}

\prose{อิธ มหานาม เอกจฺโจ ปุคฺคโล พุทฺเธ อเวจฺจปฺปสาเทน สมนฺนาคโต โหติ “อิติปิ โส ภควา ฯเปฯ สตฺถา เทวมนุสฺสานํ พุทฺโธ ภควา”ติ.\csromanspacer{} ธมฺเม ฯเปฯ.\csromanspacer{} สํเฆ ฯเปฯ หาสปญฺโญ, ชวนปญฺโญ, วิมุตฺติยา จ สมนฺนาคโต, โส อาสวานํ ขยา อนาสวํ เจโตวิมุตฺติํ ปญฺญาวิมุตฺติํ ทิฏฺเฐว ธมฺเม สยํ อภิญฺญา สจฺฉิกตฺวา อุปสมฺปชฺช วิหรติ.\csromanspacer{} อยมฺปิ โข มหานาม ปุคฺคโล ปริมุตฺโต นิรยา, ปริมุตฺโต ติรจฺฉานโยนิยา, ปริมุตฺโต เปตฺติวิสยา, ปริมุตฺโต อปายทุคฺคติวินิปาตา.}

\prose{อิธ ปน มหานาม เอกจฺโจ ปุคฺคโล พุทฺเธ อเวจฺจปฺปสาเทน สมนฺนาคโต โหติ “อิติปิ โส ภควา ฯเปฯ สตฺถา เทวมนุสฺสานํ พุทฺโธ ภควา”ติ.\csromanspacer{} ธมฺเม ฯเปฯ.\csromanspacer{} สํเฆ ฯเปฯ หาสปญฺโญ, ชวนปญฺโญ, น จ วิมุตฺติยา สมนฺนาคโต, โส ปญฺจนฺนํ โอรมฺภาคิยานํ สํโยชนานํ ปริกฺขยา โอปปาติโก โหติ ตตฺถ ปรินิพฺพายี อนาวตฺติธมฺโม ตสฺมา\footnote{อสฺมา (สฺยา, กํ, อิ, ก)} โลกา, อยมฺปิ โข มหานาม ปุคฺคโล ปริมุตฺโต นิรยา, ปริมุตฺโต ติรจฺฉานโยนิยา, ปริมุตฺโต เปตฺติวิสยา, ปริมุตฺโต อปายทุคฺคติวินิปาตา.}

\prose{“อิธ ปน มหานาม เอกจฺโจ ปุคฺคโล พุทฺเธ อเวจฺจปฺปสาเทน สมนฺนาคโต โหติ อิติปิ โส ภควา ฯเปฯ สตฺถา เทวมนุสฺสานํ พุทฺโธ ภควา”ติ.\csromanspacer{} ธมฺเม ฯเปฯ.\csromanspacer{} สํเฆ ฯเปฯ น หาสปญฺโญ, น ชวนปญฺโญ, น จ วิมุตฺติยา สมนฺนาคโต, โส ติณฺณํ สํโยชนานํ ปริกฺขยา ราคโทสโมหานํ ตนุตฺตา สกทาคามี โหติ, สกิเทว อิมํ โลกํ อาคนฺตฺวา ทุกฺขสฺสนฺตํ}

\prose{\csromanfolio{328}กโรติ, อยมฺปิ โข มหานาม ปุคฺคโล ปริมุตฺโต นิรยา, ปริมุตฺโต ติรจฺฉานโยนิยา, ปริมุตฺโต เปตฺติวิสยา, ปริมุตฺโต อปายทุคฺคติวินิปาตา.}

\prose{อิธ ปน มหานาม เอกจฺโจ ปุคฺคโล พุทฺเธ อเวจฺจปฺปสาเทน สมนฺนาคโต โหติ “อิติปิ โส ภควา ฯเปฯ สตฺถา เทวมนุสฺสานํ พุทฺโธ ภควา”ติ.\csromanspacer{} ธมฺเม ฯเปฯ.\csromanspacer{} สํเฆ ฯเปฯ น หาสปญฺโญ, น ชวนปญฺโญ, น จ วิมุตฺติยา สมนฺนาคโต, โส ติณฺณํ สํโยชนานํ ปริกฺขยา โสตาปนฺโน โหติ อวินิปาตธมฺโม นิยโต สมฺโพธิปรายโณติ.\csromanspacer{} อยมฺปิ โข มหานาม ปุคฺคโล ปริมุตฺโต นิรยา, ปริมุตฺโต ติรจฺฉานโยนิยา, ปริมุตฺโต เปตฺติวิสยา, ปริมุตฺโต อปายทุคฺคติวินิปาตา.}

\prose{อิธ ปน มหานาม เอกจฺโจ ปุคฺคโล น เหว โข พุทฺเธ อเวจฺจปฺปสาเทน สมนฺนาคโต โหติ.\csromanspacer{} น ธมฺเม ฯเปฯ.\csromanspacer{} น สํเฆ ฯเปฯ น หาสปญฺโญ, น ชวนปญฺโญ, น จ วิมุตฺติยา สมนฺนาคโต.\csromanspacer{} อปิ จสฺส อิเม ธมฺมา โหนฺติ สทฺธินฺทฺริยํ วีริยินฺทฺริยํ สตินฺทฺริยํ สมาธินฺทฺริยํ ปญฺญินฺทฺริยํ, ตถาคตปฺปเวทิตา จสฺส ธมฺมา ปญฺญาย มตฺตโส นิชฺฌานํ ขมนฺติ, อยมฺปิ โข มหานาม ปุคฺคโล อคนฺตา นิรยํ อคนฺตา ติรจฺฉานโยนิํ อคนฺตา เปตฺติวิสยํ อคนฺตา อปายํ ทุคฺคติํ วินิปาตํ.}

\prose{อิธ ปน มหานาม เอกจฺโจ ปุคฺคโล น เหว โข พุทฺเธ อเวจฺจปฺปสาเทน สมนฺนาคโต โหติ.\csromanspacer{} น ธมฺเม ฯเปฯ.\csromanspacer{} น สํเฆ ฯเปฯ น หาสปญฺโญ, น ชวนปญฺโญ, น จ วิมุตฺติยา สมนฺนาคโต.\csromanspacer{} อปิ จสฺส อิเม ธมฺมา โหนฺติ สทฺธินฺทฺริยํ ฯเปฯ ปญฺญินฺทฺริยํ.\csromanspacer{} ตถาคเต จสฺส สทฺธามตฺตํ โหติ เปมมตฺตํ.\csromanspacer{} อยมฺปิ โข มหานาม ปุคฺคโล อคนฺตา นิรยํ อคนฺตา ติรจฺฉานโยนิํ อคนฺตา เปตฺติวิสยํ อคนฺตา อปายํ ทุคฺคติํ วินิปาตํ.\csromanspacer{} อิเม เจปิ มหานาม มหาสาลา สุภาสิตํ ทุพฺภาสิตํ อาชาเนยฺยุํ, อิเมจาหํ\footnote{อิเมวาหํ (สฺยา, กํ), อิเมสาหํ (ก)} มหาสาเล พฺยากเรยฺยํ “โสตาปนฺนา อวินิปาตธมฺมา นิยตา สมฺโพธิปรายณา”ติ, กิมงฺคํ\footnote{กิมงฺค (สี, สฺยา, กํ, อิ)} ปน สรณานิํ สกฺกํ, สรณานิ มหานาม สกฺโก มรณกาเล สิกฺขํ สมาทิยีติ.\csromanspacer{}\csromanspacer{}จตุตฺถํ.}

\csromanheader{2}{11. โสตาปตฺติสํยุตฺต}
\csromanheader{2}{5. ทุติยสรณานิสกฺกสุตฺต}
\proseitem{1021}{\csromanfolio{329}กปิลวตฺถุนิทานํ.\csromanspacer{} เตน โข ปน สมเยน สรณานิ สกฺโก กาลงฺกโต โหติ.\csromanspacer{} โส ภควตา พฺยากโต “โสตาปนฺโน อวินิปาตธมฺโม นิยโต สมฺโพธิปรายโณ”ติ.\csromanspacer{} ตตฺร สุทํ สมฺพหุลา สกฺกา สงฺคมฺม สมาคมฺม อุชฺฌายนฺติ ขียนฺติ วิปาเจนฺติ “อจฺฉริยํ วต โภ, อพฺภุตํ วต โภ, เอตฺถ ทานิ โก น โสตาปนฺโน ภวิสฺสติ, ยตฺร หิ นาม สรณานิ สกฺโก กาลงฺกโต, โส ภควตา พฺยากโต ‘โสตาปนฺโน อวินิปาตธมฺโม นิยโต สมฺโพธิปรายโณ’ติ.\csromanspacer{} สรณานิ สกฺโก สิกฺขาย อปริปูรการี อโหสี”ติ.\csromanspacer{} อถ โข มหานาโม สกฺโก เยน ภควา เตนุปสงฺกมิ, อุปสงฺกมิตฺวา ภควนฺตํ อภิวาเทตฺวา เอกมนฺตํ นิสีทิ, เอกมนฺตํ นิสินฺโน โข มหานาโม สกฺโก ภควนฺตํ เอตทโวจ\csromandash{}}

\prose{“อิธ ภนฺเต สรณานิ สกฺโก กาลงฺกโต, โส ภควตา พฺยากโต ‘โสตาปนฺโน อวินิปาตธมฺโม นิยโต สมฺโพธิปรายโณ’ติ.\csromanspacer{} ตตฺร สุทํ ภนฺเต สมฺพหุลา สกฺกา สงฺคมฺม สมาคมฺม อุชฺฌายนฺติ ขียนฺติ วิปาเจนฺติ ‘อจฺฉริยํ วต โภ, อพฺภุตํ วต โภ, เอตฺถ ทานิ โก น โสตาปนฺโน ภวิสฺสติ, ยตฺร หิ นาม สรณานิ สกฺโก กาลงฺกโต, โส ภควตา พฺยากโต ‘โสตาปนฺโน อวินิปาตธมฺโม นิยโต สมฺโพธิปรายโณ’ติ, สรณานิ สกฺโก สิกฺขาย อปริปูรการี อโหสี”ติ.}

\prose{โย โส มหานาม ทีฆรตฺตํ อุปาสโก พุทฺธํ สรณํ คโต, ธมฺมํ สรณํ คโต, สํฆํ สรณํ คโต, โส กถํ วินิปาตํ คจฺเฉยฺย.\csromanspacer{} ยํ หิ ตํ มหานาม สมฺมา วทมาโน วเทยฺย “ทีฆรตฺตํ อุปาสโก พุทฺธํ สรณํ คโต, ธมฺมํ สรณํ คโต, สํฆํ สรณํ คโต”.\csromanspacer{} สรณานิํ สกฺกํ สมฺมา วทมาโน วเทยฺย “สรณานิ มหานาม สกฺโก ทีฆรตฺตํ อุปาสโก พุทฺธํ สรณํ คโต, ธมฺมํ สรณํ คโต, สํฆํ สรณํ คโต, โส กถํ วินิปาตํ คจฺเฉยฺย”.}

\prose{อิธ มหานาม เอกจฺโจ ปุคฺคโล พุทฺเธ เอกนฺตคโต โหติ อภิปฺปสนฺโน “อิติปิ โส ภควา ฯเปฯ สตฺถา เทวมนุสฺสานํ พุทฺโธ ภควา”ติ.\csromanspacer{} ธมฺเม ฯเปฯ.\csromanspacer{} สํเฆ ฯเปฯ หาสปญฺโญ, ชวนปญฺโญ, วิมุตฺติยา จ \csromanfolio{330}สมนฺนาคโต, โส อาสวานํ ขยา อนาสวํ เจโตวิมุตฺติํ ปญฺญาวิมุตฺติํ ทิฏฺเฐว ธมฺเม สยํ อภิญฺญา สจฺฉิกตฺวา อุปสมฺปชฺช วิหรติ, อยมฺปิ โข มหานาม ปุคฺคโล ปริมุตฺโต นิรยา, ปริมุตฺโต ติรจฺฉานโยนิยา, ปริมุตฺโต เปตฺติวิสยา, ปริมุตฺโต อปายทุคฺคติวินิปาตา.}

\prose{อิธ ปน มหานาม เอกจฺโจ ปุคฺคโล พุทฺเธ เอกนฺตคโต โหติ อภิปฺปสนฺโน “อิติปิ โส ภควา ฯเปฯ สตฺถา เทวมนุสฺสานํ พุทฺโธ ภควา”ติ.\csromanspacer{} ธมฺเม ฯเปฯ.\csromanspacer{} สํเฆ ฯเปฯ หาสปญฺโญ, ชวนปญฺโญ, น จ วิมุตฺติยา สมนฺนาคโต, โส ปญฺจนฺนํ โอรมฺภาคิยานํ สํโยชนานํ ปริกฺขยา อนฺตราปรินิพฺพายี โหติ, อุปหจฺจปรินิพฺพายี โหติ, อสงฺขารปรินิพฺพายี โหติ, สสงฺขารปรินิพฺพายี โหติ, อุทฺธํโสโต โหติ อกนิฏฺฐคามี.\csromanspacer{} อยมฺปิ โข มหานาม ปุคฺคโล ปริมุตฺโต นิรยา, ปริมุตฺโต ติรจฺฉานโยนิยา, ปริมุตฺโต เปตฺติวิสยา, ปริมุตฺโต อปายทุคฺคติวินิปาตา.}

\prose{อิธ ปน มหานาม เอกจฺโจ ปุคฺคโล พุทฺเธ เอกนฺตคโต โหติ อภิปฺปสนฺโน “อิติปิ โส ภควา ฯเปฯ สตฺถา เทวมนุสฺสานํ พุทฺโธ ภควา”ติ.\csromanspacer{} ธมฺเม ฯเปฯ.\csromanspacer{} สํเฆ ฯเปฯ น หาสปญฺโญ, น ชวนปญฺโญ, น จ วิมุตฺติยา สมนฺนาคโต, โส ติณฺณํ สํโยชนานํ ปริกฺขยา ราคโทสโมหานํ ตนุตฺตา สกทาคามี โหติ, สกิเทว อิมํ โลกํ อาคนฺตฺวา ทุกฺขสฺสนฺตํ กโรติ.\csromanspacer{} อยมฺปิ โข มหานาม ปุคฺคโล ปริมุตฺโต นิรยา, ปริมุตฺโต ติรจฺฉานโยนิยา, ปริมุตฺโต เปตฺติวิสยา, ปริมุตฺโต อปายทุคฺคติวินิปาตา.}

\prose{อิธ ปน มหานาม เอกจฺโจ ปุคฺคโล พุทฺเธ เอกนฺตคโต โหติ อภิปฺปสนฺโน “อิติปิ โส ภควา ฯเปฯ สตฺถา เทวมนุสฺสานํ พุทฺโธ ภควา”ติ.\csromanspacer{} ธมฺเม ฯเปฯ.\csromanspacer{} สํเฆ ฯเปฯ น หาสปญฺโญ, น ชวนปญฺโญ, น จ วิมุตฺติยา สมนฺนาคโต.\csromanspacer{} โส ติณฺณํ สํโยชนานํ ปริกฺขยา โสตาปนฺโน โหติ อวินิปาตธมฺโม นิยโต สมฺโพธิปรายโณ, อยมฺปิ โข มหานาม ปุคฺคโล ปริมุตฺโต นิรยา, ปริมุตฺโต ติรจฺฉานโยนิยา, ปริมุตฺโต เปตฺติวิสยา, ปริมุตฺโต อปายทุคฺคติวินิปาตา.}

\prose{อิธ ปน มหานาม เอกจฺโจ ปุคฺคโล น เหว โข พุทฺเธ เอกนฺตคโต โหติ อภิปฺปสนฺโน ฯเปฯ.\csromanspacer{} น ธมฺเม ฯเปฯ.\csromanspacer{} น สํเฆ ฯเปฯ น หาสปญฺโญ, น ชวนปญฺโญ, น จ วิมุตฺติยา สมนฺนาคโต, อปิ จสฺส อิเม ธมฺมา โหนฺติ}

\csromanheader{2}{11. โสตาปตฺติสํยุตฺต}
\prose{\csromanfolio{331}สทฺธินฺทฺริยํ ฯเปฯ ปญฺญินฺทฺริยํ, ตถาคตปฺปเวทิตา จสฺส ธมฺมา ปญฺญาย มตฺตโส นิชฺฌานํ ขมนฺติ.\csromanspacer{} อยมฺปิ โข มหานาม ปุคฺคโล อคนฺตา นิรยํ อคนฺตา ติรจฺฉานโยนิํ อคนฺตา เปตฺติวิสยํ อคนฺตา อปายํ ทุคฺคติํ วินิปาตํ.}

\prose{อิธ ปน มหานาม เอกจฺโจ ปุคฺคโล น เหว โข พุทฺเธ เอกนฺตคโต โหติ อภิปฺปสนฺโน.\csromanspacer{} น ธมฺเม ฯเปฯ.\csromanspacer{} น สํเฆ ฯเปฯ น หาสปญฺโญ, น ชวนปญฺโญ, น จ วิมุตฺติยา สมนฺนาคโต.\csromanspacer{} อปิ จสฺส อิเม ธมฺมา โหนฺติ สทฺธินฺทฺริยํ ฯเปฯ ปญฺญินฺทฺริยํ, ตถาคเต จสฺส สทฺธามตฺตํ โหติ เปมมตฺตํ.\csromanspacer{} อยมฺปิ โข มหานาม ปุคฺคโล อคนฺตา นิรยํ อคนฺตา ติรจฺฉานโยนิํ อคนฺตา เปตฺติวิสยํ อคนฺตา อปายํ ทุคฺคติํ วินิปาตํ.}

\prose{เสยฺยถาปิ มหานาม ทุกฺเขตฺตํ ทุพฺภูมํ อวิหตขาณุกํ, พีชานิ จสฺสุ ขณฺฑานิ ปูตีนิ วาตาตปหตานิ อสาราทานิ อสุขสยิตานิ\footnote{อสุขาปสฺสยิตานิ (ก)}, เทโว จ น สมฺมา\footnote{เทโว ปน สมฺมา (สฺยา, กํ), เทโว น สมฺมา (ก) ที 2. 279 ปิฏฺเฐปิ.} ธารํ อนุปฺปเวจฺเฉยฺย, อปิ นุ ตานิ พีชานิ วุทฺธิํ วุรูฬฺหิํ เวปุลฺลํ อาปชฺเชยฺยุนฺติ.\csromanspacer{} โน เหตํ ภนฺเต.\csromanspacer{} เอวเมว โข มหานาม อิธ ธมฺโม ทุรกฺขาโต\footnote{ทฺวากฺขาโต (อิ, ก)} โหติ ทุปฺปเวทิโต อนิยฺยานิโก อนุปสมสํวตฺตนิโก อสมฺมาสมฺพุทฺธปฺปเวทิโต.\csromanspacer{} อิทมหํ ทุกฺเขตฺตสฺมิํ วทามิ, ตสฺมิํ จ ธมฺเม สาวโก วิหรติ ธมฺมานุธมฺมปฺปฏิปนฺโน สามีจิปฺปฏิปนฺโน อนุธมฺมจารี, อิทมหํ ทุพฺพีชสฺมิํ วทามิ.}

\prose{เสยฺยถาปิ มหานาม สุเขตฺตํ สุภูมํ สุวิหตขาณุกํ, พีชานิ จสฺสุ อขณฺฑานิ อปูตีนิ อวาตาตปหตานิ สาราทานิ สุขสยิตานิ, เทโว จ\footnote{เทโว จสฺส (สฺยา, กํ, ก)} สมฺมา ธารํ อนุปฺปเวจฺเฉยฺย.\csromanspacer{} อปิ นุ ตานิ พีชานิ วุทฺธิํ วิรูฬฺหิํ เวปุลฺลํ อาปชฺเชยฺยุนฺติ.\csromanspacer{} เอวํ ภนฺเต.\csromanspacer{} เอวเมว โข มหานาม อิธ ธมฺโม สฺวากฺขาโต โหติ สุปฺปเวทิโต นิยฺยานิโก อุปสมสํวตฺตนิโก สมฺมาสมฺพุทฺธปฺปเวทิโต, อิทมหํ สุเขตฺตสฺมิํ วทามิ, ตสฺมิํ จ ธมฺเม สาวโก วิหรติ ธมฺมานุธมฺมปฺปฏิปนฺโน สามีจิปฺปฏิปนฺโน อนุธมฺมจรี, อิทมหํ สุพีชสฺมิํ วทามิ, กิมงฺคํ ปน สรณานิํ สกฺกํ, สรณานิ มหานาม สกฺโก มรณกาเล สิกฺขาย ปริปูรการี อโหสีติ.\csromanspacer{}\csromanspacer{}ปญฺจมํ.}

\csromanheader{2}{6. ปฐม-อนาถปิณฺฑิกสุตฺต}
\csromanmatikaanchor{mtk.240}
\csromantocmarkref{subparagraph}{6-7. อนาถปิณฺฑิกสุตฺตานิ}{mtk.240}
\bookmark[dest={mtk.240},level=5]{6-7. อนาถปิณฺฑิกสุตฺตานิ}
\proseitem{1022}{\csromanfolio{332}สาวตฺถินิทานํ.\csromanspacer{} เตน โข ปน สมเยน อนาถปิณฺฑิโก คหปติ อาพาธิโก โหติ ทุกฺขิโต พาฬฺหคิลาโน.\csromanspacer{} อถ โข อนาถปิณฺฑิโก คหปติ อญฺญตรํ ปุริสํ อามนฺเตสิ “เอหิ ตฺวํ อมฺโภ ปุริส เยนายสฺมา สาริปุตฺโต เตนุปสงฺกม, อุปสงฺกมิตฺวา มม วจเนน อายสฺมโต สาริปุตฺตสฺส ปาเท สิรสา วนฺท, อนาถปิณฺฑิโก ภนฺเต คหปติ อาพาธิโก ทุกฺขิโต พาฬฺหคิลาโน, โส อายสฺมโต สาริปุตฺตสฺส ปาเท สิรสา วนฺทตี”ติ, เอวญฺจ วเทหิ “สาธุ กิร ภนฺเต อายสฺมา สาริปุตฺโต เยน อนาถปิณฺฑิกสฺส คหปติสฺส นิเวสนํ เตนุปสงฺกมตุ อนุกมฺปํ อุปาทายา”ติ.}

\prose{“เอวํ ภนฺเต”ติ โข โส ปุริโส อนาถปิณฺฑิกสฺส คหปติสฺส ปฏิสฺสุตฺวา เยนายสฺมา สาริปุตฺโต เตนุปสงฺกมิ, อุปสงฺกมิตฺวา อายสฺมนฺตํ สาริปุตฺตํ อภิวาเทตฺวา เอกมนฺตํ นิสีทิ, เอกมนฺตํ นิสินฺโน โข โส ปุริโส อายสฺมนฺตํ สาริปุตฺตํ เอตทโวจ\csromandash{}}

\prose{อนาถิปิณฺฑิโก ภนฺเต คหปติ อาพาธิโก ทุกฺขิโต พาฬฺหคิลาโน, โส อายสฺมโต สาริปุตฺตสฺส ปาเท สิรสา วนฺทติ, เอวญฺจ วทติ “สาธุ กิร ภนฺเต อายสฺมา สาริปุตฺโต เยน อนาถปิณฺฑิกสฺส คหปติสฺส นิเวสนํ เตนุปสงฺกมตุ อนุกมฺปํ อุปาทายา”ติ.\csromanspacer{} อธิวาเสสิ โข อายสฺมา สาริปุตฺโต ตุณฺหีภาเวน.}

\prose{อถ โข อายสฺมา สาริปุตฺโต ปุพฺพณฺหสมยํ นิวาเสตฺวา ปตฺตจีวรมาทาย อายสฺมตา อานนฺเทน ปจฺฉาสมเณน เยน อนาถปิณฺฑิกสฺส คหปติสฺส นิเวสนํ เตนุปสงฺกมิ, อุปสงฺกมิตฺวา ปญฺญตฺเต อาสเน นิสีทิ, นิสชฺช โข อายสฺมา สาริปุตฺโต อนาถปิณฺฑิกํ คหปติํ เอตทโวจ “กจฺจิ เต คหปติ ขมนียํ, กจฺจิ ยาปนียํ, กจฺจิ ทุกฺขา เวทนา ปฏิกฺกมนฺติ, โน อภิกฺกมนฺติ, ปฏิกฺกโมสานํ ปญฺญายติ, โน อภิกฺกโม”ติ.\csromanspacer{} น เม ภนฺเต ขมนียํ, น ยาปนียํ, พาฬฺหา เม ทุกฺขา เวทนา อภิกฺกมนฺติ, โน ปฏิกฺกมนฺติ, อภิกฺกโมสานํ ปญฺญายติ, โน ปฏิกฺกโมติ.}

\prose{ยถารูเปน โข คหปติ พุทฺเธ อปฺปสาเทน สมนฺนาคโต อสฺสุตวา ปุถุชฺชโน กายสฺส เภทา ปรํ มรณา อปายํ ทุคฺคติํ วินิปาตํ นิรยํ}

\csromanheader{2}{11. โสตาปตฺติสํยุตฺต}
\prose{\csromanfolio{333}อุปปชฺชติ, ตถารูโป เต พุทฺเธ อปฺปสาโท นตฺถิ, อตฺถิ จ โข เต คหปติ พุทฺเธ อเวจฺจปฺปสาโท “อิติปิ โส ภควา ฯเปฯ สตฺถา เทวมนุสฺสานํ พุทฺโธ ภควา”ติ, ตญฺจ ปน เต พุทฺเธ อเวจฺจปฺปสาทํ อตฺตนิ สมนุปสฺสโต ฐานโส เวทนา ปฏิปฺปสฺสมฺเภยฺย.}

\prose{ยถารูเปน โข คหปติ ธมฺเม อปฺปสาเทน สมนฺนาคโต อสฺสุตวา ปุถุชฺชโน กายสฺส เภทา ปรํ มรณา อปายํ ทุคฺคติํ วินิปาตํ นิรยํ อุปปชฺชติ, ตถารูโป เต ธมฺเม อปฺปสาโท นตฺถิ, อตฺถิ จ โข เต คหปติ ธมฺเม อเวจฺจปฺปสาโท “สฺวากฺขาโต ภควตา ธมฺโม ฯเปฯ ปจฺจตฺตํ เวทิตพฺโพ วิญฺญูหี”ติ, ตญฺจ ปน เต ธมฺเม อเวจฺจปฺปสาทํ อตฺตนิ สมนุปสฺสโต ฐานโส เวทนา ปฏิปฺปสฺสมฺเภยฺย.}

\prose{ยถารูเปน โข คหปติ สํเฆ อปฺปสาเทน สมนฺนาคโต อสฺสุตวา ปุถุชฺชโน กายสฺส เภทา ปรํ มรณา อปายํ ทุคฺคติํ วินิปาตํ นิรยํ อุปปชฺชติ, ตถารูโป เต สํเฆ อปฺปสาโท นตฺถิ, อตฺถิ จ โข เต คหปติ สํเฆ อเวจฺจปฺปสาโท “สุปฺปฏิปนฺโน ภควโต สาวกสํโฆ ฯเปฯ อนุตฺตรํ ปุญฺญกฺเขตฺตํ โลกสฺสา”ติ, ตญฺจ ปน เต สํเฆ อเวจฺจปฺปสาทํ อตฺตนิ สมนุปสฺสโต ฐานโส เวทนา ปฏิปฺปสฺสมฺเภยฺย.}

\prose{ยถารูเปน โข คหปติ ทุสฺสีเลฺยน สมนฺนาคโต อสฺสุตวา ปุถุชฺชโน กายสฺส เภทา ปรํ มรณา อปายํ ทุคฺคติํ วินิปาตํ นิรยํ อุปปชฺชติ, ตถารูปํ เต ทุสฺสีลฺยํ นตฺถิ.\csromanspacer{} อตฺถิ จ โข เต คหปติ อริยกนฺตานิ สีลานิ ฯเปฯ สมาธิสํวตฺตนิกานิ, ตานิ จ ปน เต อริยกนฺตานิ สีลานิ อตฺตนิ สมนุปสฺสโต ฐานโส เวทนา ปฏิปฺปสฺสมฺเภยฺย.}

\prose{ยถารูปาย โข คหปติ มิจฺฉาทิฏฺฐิยา สมนฺนาคโต อสฺสุตวา ปุถุชฺชโน กายสฺส เภทา ปรํ มรณา อปายํ ทุคฺคติํ วินิปาตํ นิรยํ อุปปชฺชติ, ตถารูปา เต มิจฺฉาทิฏฺฐิ นตฺถิ.\csromanspacer{} อตฺถิ จ โข เต คหปติ สมฺมาทิฏฺฐิ, ตญฺจ ปน สมฺมาทิฏฺฐิํ อตฺตนิ สมนุปสฺสโต ฐานโส เวทนา ปฏิปฺปสฺสมฺเภยฺย.}

\prose{ยถารูเปน โข คหปติ มิจฺฉาสงฺกปฺเปน สมนฺนาคโต อสฺสุตวา ปุถุชฺชโน กายสฺส เภทา ปรํ มรณา อปายํ ทุคฺคติํ วินิปาตํ นิรยํ}

\prose{\csromanfolio{334}อุปปชฺชติ, ตถารูโป เต มิจฺฉาสงฺกปฺโป นตฺถิ.\csromanspacer{} อตฺถิ จ โข เต คหปติ สมฺมาสงฺกปฺโป, ตญฺจ ปน เต สมฺมาสงฺกปฺปํ อตฺตนิ สมนุปสฺสโต ฐานโส เวทนา ปฏิปฺปสฺสมฺเภยฺย.}

\prose{ยถารูปาย โข คหปติ มิจฺฉาวาจาย สมนฺนาคโต อสฺสุตวา ปุถุชฺชโน กายสฺส เภทา ปรํ มรณา อปายํ ทุคฺคติํ วินิปาตํ นิรยํ อุปปชฺชติ, ตถารูปา เต มิจฺฉาวาจา นตฺถิ.\csromanspacer{} อตฺถิ จ โข เต คหปติ สมฺมาวาจา, ตญฺจ ปน เต สมฺมาวาจํ อตฺตนิ สมนุปสฺสโต ฐานโส เวทนา ปฏิปฺปสฺสมฺเภยฺย.}

\prose{ยถารูเปน โข คหปติ มิจฺฉากมฺมนฺเตน สมนฺนาคโต อสฺสุตวา ปุถุชฺชโน กายสฺส เภทา ปรํ มรณา อปายํ ทุคฺคติํ วินิปาตํ นิรยํ อุปปชฺชติ, ตถารูโป เต มิจฺฉากมฺมนฺโต นตฺถิ.\csromanspacer{} อตฺถิ จ โข เต คหปติ สมฺมากมฺมนฺโต, ตญฺจ ปน เต สมฺมากมฺมนฺตํ อตฺตนิ สมนุปสฺสโต ฐานโส เวทนา ปฏิปฺปสฺสมฺเภยฺย.}

\prose{ยถารูเปน โข คหปติ มิจฺฉา-อาชีเวน สมนฺนาคโต อสฺสุตวา ปุถุชฺชโน กายสฺส เภทา ปรํ มรณา อปายํ ทุคฺคติํ วินิปาตํ นิรยํ อุปปชฺชติ, ตถารูโป เต มิจฺฉา-อาชีโว นตฺถิ.\csromanspacer{} อตฺถิ จ โข เต คหปติ สมฺมา-อาชีโว, ตญฺจ ปน เต สมฺมา-อาชีวํ อตฺตนิ สมนุปสฺสโต ฐานโส เวทนา ปฏิปฺปสฺสมฺเภยฺย.}

\prose{ยถารูเปน โข คหปติ มิจฺฉาวายาเมน สมนฺนาคโต อสฺสุตวา ปุถุชฺชโน กายสฺส เภทา ปรํ มรณา อปายํ ทุคฺคติํ วินิปาตํ นิรยํ อุปปชฺชติ, ตถารูโป เต มิจฺฉาวายาโม นตฺถิ.\csromanspacer{} อตฺถิ จ โข เต คหปติ สมฺมาวายาโม, ตญฺจ ปน เต สมฺมาวายามํ อตฺตนิ สมนุปสฺสโต ฐานโส เวทนา ปฏิปฺปสฺสมฺเภยฺย.}

\prose{ยถารูเปน โข คหปติ มิจฺฉาสติยา สมนฺนาคโต อสฺสุตวา ปุถุชฺชโน กายสฺส เภทา ปรํ มรณา อปายํ ทุคฺคติํ วินิปาตํ นิรยํ อุปปชฺชติ, ตถารูปา เต มิจฺฉาสติ นตฺถิ.\csromanspacer{} อตฺถิ จ โข เต คหปติ สมฺมาสติ, ตญฺจ ปน เต สมฺมาสติํ อตฺตนิ สมนุปสฺสโต ฐานโส เวทนา ปฏิปฺปสฺสมฺเภยฺย.}

\prose{ยถารูเปน โข คหปติ มิจฺฉาสมาธินา สมนฺนาคโต อสฺสุตวา ปุถุชฺชโน กายสฺส เภทา ปรํ มรณา อปายํ ทุคฺคติํ วินิปาตํ นิรยํ}

\csromanheader{2}{11. โสตาปตฺติสํยุตฺต}
\prose{\csromanfolio{335}อุปปชฺชติ, ตถารูโป เต มิจฺฉาสมาธิ นตฺถิ.\csromanspacer{} อตฺถิ จ โข เต คหปติ สมฺมาสมาธิ, ตญฺจ ปน เต สมฺมาสมาธิํ อตฺตนิ สมนุปสฺสโต ฐานโส เวทนา ปฏิปฺปสฺสมฺเภยฺย.}

\prose{ยถารูเปน โข คหปติ มิจฺฉาญาเณน สมนฺนาคโต อสฺสุตวา ปุถุชฺชโน กายสฺส เภทา ปรํ มรณา อปายํ ทุคฺคติํ วินิปาตํ นิรยํ อุปปชฺชติ, ตถารูปํ เต มิจฺฉาญาณํ นตฺถิ.\csromanspacer{} อตฺถิ จ โข เต คหปติ สมฺมาญาณํ, ตญฺจ ปน เต สมฺมาญาณํ อตฺตนิ สมนุปสฺสโต ฐานโส เวทนา ปฏิปฺปสฺสมฺเภยฺย.}

\prose{ยถารูปาย โข คหปติ มิจฺฉาวิมุตฺติยา สมนฺนาคโต อสฺสุตวา ปุถุชฺชโน กายสฺส เภทา ปรํ มรณา อปายํ ทุคฺคติํ วินิปาตํ นิรยํ อุปปชฺชติ, ตถารูปา เต มิจฺฉาวิมุตฺติ นตฺถิ.\csromanspacer{} อตฺถิ จ โข เต คหปติ สมฺมาวิมุตฺติ, ตญฺจ ปน เต สมฺมาวิมุตฺติํ อตฺตนิ สมนุปสฺสโต ฐานโส เวทนา ปฏิปฺปสฺสมฺเภยฺยาติ.}

\prose{อถ โข อนาถปิณฺฑิกสฺส คหปติสฺส ฐานโส เวทนา ปฏิปฺปสฺสมฺภิํสุ.\csromanspacer{} อถ โข อนาถปิณฺฑิโก คหปติ อายสฺมนฺตญฺจ สาริปุตฺตํ อายสฺมนฺตญฺจ อานนฺทํ สเกเนว ถาลิปาเกน ปริวิสิ.\csromanspacer{} อถ โข อนาถปิณฺฑิโก คหปติ อายสฺมนฺตํ สาริปุตฺตํ ภุตฺตาวิํ โอนีตปตฺตปาณิํ อญฺญตรํ นีจาสนํ คเหตฺวา เอกมนฺตํ นิสีทิ, เอกมนฺตํ นิสินฺนํ โข อนาถปิณฺฑิกํ คหปติํ อายสฺมา สาริปุตฺโต อิมาหิ คาถาหิ อนุโมทิ\csromandash{}}

\csromangathameasure{“ยสฺส สทฺธา ตถาคเต, อจลา สุปฺปติฏฺฐิตา. \\ สีลญฺจ ยสฺส กลฺยาณํ, อริยกนฺตํ ปสํสิตํ. \\ สํเฆ ปสาโท ยสฺสตฺถิ, อุชุภูตญฺจ ทสฺสนํ. \\ อทลิทฺโทติ ตํ อาหุ, อโมฆํ ตสฺส ชีวิตํ. \\ ตสฺมา สทฺธญฺจ สีลญฺจ, ปสาทํ ธมฺมทสฺสนํ. \\ อนุยุญฺเชถ เมธาวี, สรํ พุทฺธานสาสนนฺ”ติ.}
\csromangathagroup{\csromangathastanza{“ยสฺส สทฺธา ตถาคเต, อจลา สุปฺปติฏฺฐิตา. \\ สีลญฺจ ยสฺส กลฺยาณํ, อริยกนฺตํ ปสํสิตํ.}\csromangathastanza{สํเฆ ปสาโท ยสฺสตฺถิ, อุชุภูตญฺจ ทสฺสนํ. \\ อทลิทฺโทติ\footnote{อทฬิทฺโทติ (สี, สฺยา, กํ)} ตํ อาหุ, อโมฆํ ตสฺส ชีวิตํ.}\csromangathastanza{ตสฺมา สทฺธญฺจ สีลญฺจ, ปสาทํ ธมฺมทสฺสนํ. \\ อนุยุญฺเชถ เมธาวี, สรํ พุทฺธานสาสนนฺ”ติ.}}

\prose{อถ โข อายสฺมา สาริปุตฺโต อนาถปิณฺฑิกํ คหปติํ อิมาหิ คาถาหิ อนุโมทิตฺวา อุฏฺฐายาสนา ปกฺกามิ.}

\prose{\csromanfolio{336}อถ โข อายสฺมา อานนฺโท เยน ภควา เตนุปสงฺกมิ, อุปสงฺกมิตฺวา ภควนฺตํ อภิวาเทตฺวา เอกมนฺตํ นิสีทิ, เอกมนฺตํ นิสินฺนํ โข อายสฺมนฺตํ อานนฺทํ ภควา เอตทโวจ “หนฺท กุโต นุ ตฺวํ อานนฺท อาคจฺฉสิ ทิวา ทิวสฺสาติ.\csromanspacer{} อายสฺมตา ภนฺเต สาริปุตฺเตน อนาถปิณฺฑิโก คหปติ อิมินา จ อิมินา จ โอวาเทน โอวทิโตติ.\csromanspacer{} ปณฺฑิโต อานนฺท สาริปุตฺโต, มหาปญฺโญ อานนฺท สาริปุตฺโต, ยตฺร หิ นาม จตฺตาริ โสตาปตฺติยงฺคานิ ทสหากาเรหิ วิภชิสฺสตี”ติ.\csromanspacer{}\csromanspacer{}ฉฏฺฐํ.}

\csromanheader{2}{7. ทุติย-อนาถปิณฺฑิกสุตฺต}
\proseitem{1023}{สาวตฺถินิทานํ.\csromanspacer{} เตน โข ปน สมเยน อนาถปิณฺฑิโก คหปติ อาพาธิโก โหติ ทุกฺขิโต พาฬฺหคิลาโน.\csromanspacer{} อถ โข อนาถปิณฺฑิโก คหปติ อญฺญตรํ ปุริสํ อามนฺเตสิ “เอหิ ตฺวํ อมฺโภ ปุริส, เยนายสฺมา อานนฺโท เตนุปสงฺกม, อุปสงฺกมิตฺวา มม วจเนน อายสฺมโต อานนฺทสฺส ปาเท สิรสา วนฺท, อนาถปิณฺฑิโก ภนฺเต คหปติ อาพาธิโก ทุกฺขิโต พาฬฺหคิลาโน, โส อายสฺมโต อานนฺทสฺส ปาเท สิรสา วนฺทตีติ, เอวญฺจ วเทหิ ‘สาธุ กิร ภนฺเต อายสฺมา อานนฺโท เยน อนาถปิณฺฑิกสฺส คหปติสฺส นิเวสนํ เตนุปสงฺกมตุ, อนุกมฺปํ อุปาทายา’ติ”.}

\prose{“เอวํ ภนฺเต”ติ โข โส ปุริโส อนาถปิณฺฑิกสฺส คหปติสฺส ปฏิสฺสุตฺวา เยนายสฺมา อานนฺโท เตนุปสงฺกมิ, อุปสงฺกมิตฺวา อายสฺมานฺตํ อานนฺทํ อภิวาเทตฺวา เอกมนฺตํ นิสีทิ, เอกมนฺตํ นิสินฺโน โข โส ปุริโส อายสฺมนฺตํ อานนฺทํ เอตทโวจ “อนาถปิณฺฑิโก ภนฺเต คหปติ อาพาธิโก ทุกฺขิโต พาฬฺหคิลาโน, โส อายสฺมโต อานนฺทสฺส ปาเท สิรสา วนฺทติ, เอวญฺจ วทติ ‘สาธุ กิร ภนฺเต อายสฺมา อานนฺโท เยน อนาถปิณฺฑิกสฺส คหปติสฺส นิเวสนํ เตนุปสงฺกมตุ อนุกมฺปํ อุปาทายา’ติ”.\csromanspacer{} อธิวาเสสิ โข อายสฺมา อานนฺโท ตุณฺหีภาเวน.}

\prose{อถ โข อายสฺมา อานนฺโท ปุพฺพณฺหสมยํ นิวาเสตฺวา ปตฺตจีวรมาทาย เยน อนาถปิณฺฑิกสฺส คหปติสฺส นิเวสนํ เตนุปสงฺกมิ, อุปสงฺกมิตฺวา ปญฺญตฺเต อาสเน นิสีทิ, นิสชฺช โข อายสฺมา อานนฺโท}

\csromanheader{2}{11. โสตาปตฺติสํยุตฺต}
\prose{\csromanfolio{337}อนาถปิณฺฑิกํ คหปติํ เอตทโวจ “กจฺจิ เต คหปติ ขมนียํ, กจฺจิ ยาปนียํ, กจฺจิ ทุกฺขา เวทนา ปฏิกฺกมนฺติ, โน อภิกฺกมนฺติ, ปฏิกฺกโมสานํ ปญฺญายติ, โน อภิกฺกโม”ติ.\csromanspacer{} น เม ภนฺเต ขมนียํ, น ยาปนียํ, พาฬฺหา เม ทุกฺขา เวทนา อภิกฺกมนฺติ, โน ปฏิกฺกมนฺติ, อภิกฺกโมสานํ ปญฺญายติ, โน ปฏิกฺกโมติ.}

\prose{จตูหิ โข คหปติ ธมฺเมหิ สมนฺนาคตสฺส อสฺสุตวโต ปุถุชฺชนสฺส โหติ อุตฺตาโส, โหติ ฉมฺภิตตฺตํ, โหติ สมฺปรายิกํ มรณภยํ.\csromanspacer{} กตเมหิ จตูหิ, อิธ คหปติ อสฺสุตวา ปุถุชฺชโน พุทฺเธ อปฺปสาเทน สมนฺนาคโต โหติ, ตญฺจ ปนสฺส พุทฺเธ อปฺปสาทํ อตฺตนิ สมนุปสฺสโต โหติ อุตฺตาโส, โหติ ฉมฺภิตตฺตํ, โหติ สมฺปรายิกํ มรณภยํ.}

\prose{ปุน จปรํ คหปติ อสฺสุตวา ปุถุชฺชโน ธมฺเม อปฺปสาเทน สมนฺนาคโต โหติ, ตญฺจ ปนสฺส ธมฺเม อปฺปสาทํ อตฺตนิ สมนุปสฺสโต โหติ อุตฺตาโส, โหติ ฉมฺภิตตฺตํ, โหติ สมฺปรายิกํ มรณภยํ.}

\prose{ปุน จปรํ คหปติ อสฺสุตวา ปุถุชฺชโน สํเฆ อปฺปสาเทน สมนฺนาคโต โหติ, ตญฺจ ปนสฺส สํเฆ อปฺปสาทํ อตฺตนิ สมนุปสฺสโต โหติ อุตฺตาโส, โหติ ฉมฺภิตตฺตํ, โหติ สมฺปรายิกํ มรณภยํ.}

\prose{ปุน จปรํ คหปติ อสฺสุตวา ปุถุชฺชโน ทุสฺสีเลฺยน สมนฺนาคโต โหติ, ตญฺจ ปนสฺส ทุสฺสีลฺยํ อตฺตนิ สมนุปสฺสโต โหติ อุตฺตาโส, โหติ ฉมฺภิตตฺตํ, โหติ สมฺปรายิกํ มรณภยํ.\csromanspacer{} อิเมหิ โข คหปติ จตูหิ ธมฺเมหิ สมนฺนาคตสฺส อสฺสุตวโต ปุถุชฺชนสฺส โหติ อุตฺตาโส, โหติ ฉมฺภิตตฺตํ, โหติ สมฺปรายิกํ มรณภยํ.}

\prose{จตูหิ โข คหปติ ธมฺเมหิ สมนฺนาคตสฺส สุตวโต อริยสาวกสฺส น โหติ อุตฺตาโส, น โหติ ฉมฺภิตตฺตํ, น โหติ สมฺปรายิกํ มรณภยํ.\csromanspacer{} กตเมหิ จตูหิ, อิธ คหปติ สุตวา อริยสาวโก พุทฺเธ อเวจฺจปฺปสาเทน สมนฺนาคโต โหติ “อิติปิ โส ภควา ฯเปฯ สตฺถา เทวมนุสฺสานํ พุทฺโธ ภควา”ติ, ตญฺจ ปนสฺส พุทฺเธ อเวจฺจปฺปสาทํ อตฺตนิ สมนุปสฺสโต น โหติ อุตฺตาโส, น โหติ ฉมฺภิตตฺตํ, น โหติ สมฺปรายิกํ มรณภยํ.}

\prose{ปุน จปรํ คหปติ สุตวา อริยสาวโก ธมฺเม อเวจฺจปฺปสาเทน สมนฺนาคโต โหติ “สฺวากฺขาโต ภควตา ธมฺโม ฯเปฯ ปจฺจตฺตํ}

\csromanmatikaanchor{mtk.241}
\csromantocmarkref{subparagraph}{8-9. ภยเวรูปสนฺตสุตฺตานิ}{mtk.241}
\bookmark[dest={mtk.241},level=5]{8-9. ภยเวรูปสนฺตสุตฺตานิ}
\prose{\csromanfolio{338}เวทิตพฺโพ วิญฺญูหี”ติ, ตญฺจ ปนสฺส ธมฺเม อเวจฺจปฺปสาทํ อตฺตนิ สมนุปสฺสโต น โหติ อุตฺตาโส, น โหติ ฉมฺภิตตฺตํ, น โหติ สมฺปรายิกํ มรณภยํ.}

\prose{ปุน จปรํ คหปติ สุตวา อริยสาวโก สํเฆ อเวจฺจปฺปสาเทน สมนฺนาคโต โหติ “สุปฺปฏิปนฺโน ภควโต สาวกสํโฆ ฯเปฯ อนุตฺตรํ ปุญฺญกฺเขตฺตํ โลกสฺสา”ติ.\csromanspacer{} ตญฺจ ปนสฺส สํเฆ อเวจฺจปฺปสาทํ อตฺตนิ สมนุปสฺสโต น โหติ อุตฺตาโส, น โหติ ฉมฺภิตตฺตํ, น โหติ สมฺปรายิกํ มรณภยํ.}

\prose{ปุน จปรํ คหปติ สุตวา อริยสาวโก อริยกนฺเตหิ สีเลหิ สมนฺนาคโต โหติ อขณฺเฑหิ ฯเปฯ สมาธิสํวตฺตนิเกหิ, ตานิ จ ปนสฺส อริยกนฺตานิ สีลานิ อตฺตนิ สมนุปสฺสโต น โหติ อุตฺตาโส, น โหติ ฉมฺภิตตฺตํ, น โหติ สมฺปรายิกํ มรณภยํ.\csromanspacer{} อิเมหิ โข คหปติ จตูหิ ธมฺเมหิ สมนฺนาคตสฺส สุตวโต อริยสาวกสฺส น โหติ อุตฺตาโส, น โหติ ฉมฺภิตตฺตํ, น โหติ สมฺปรายิกํ มรณภยนฺติ.}

\prose{นาหํ ภนฺเต อานนฺท ภายามิ, กฺยาหํ ภายิสฺสามิ.\csromanspacer{} อหํ หิ ภนฺเต พุทฺเธ อเวจฺจปฺปสาเทน สมนฺนาคโต โหมิ “อิติปิ โส ภควา ฯเปฯ สตฺถา เทวมนุสฺสานํ พุทฺโธ ภควา”ติ.\csromanspacer{} ธมฺเม ฯเปฯ.\csromanspacer{} สํเฆ อเวจฺจปฺปสาเทน สมนฺนาคโต โหมิ “สุปฺปฏิปนฺโน ภควโต สาวกสํโฆ ฯเปฯ อนุตฺตรํ ปุญฺญกฺเขตฺตํ โลกสฺสา”ติ, ยานิ จิมานิ ภนฺเต ภควตา คิหิสามีจิกานิ สิกฺขาปทานิ เทสิตานิ, นาหํ เตสํ กิญฺจิ อตฺตนิ ขณฺฑํ สมนุปสฺสามีติ.\csromanspacer{} ลาภา เต คหปติ, สุลทฺธํ เต คหปติ, โสตาปตฺติผลํ ตยา คหปติ พฺยากตนฺติ.\csromanspacer{}\csromanspacer{}สตฺตมํ.}

\csromanheader{2}{8. ปฐมภยเวรูปสนฺตสุตฺต}
\proseitem{1024}{สาวตฺถินิทานํ.\csromanspacer{} เอกมนฺตํ นิสินฺนํ โข อนาถปิณฺฑิกํ คหปติํ ภควา เอตทโวจ\csromandash{}ยโต โข คหปติ อริยสาวกสฺส ปญฺจ ภยานิ เวรานิ วูปสนฺตานิ จ โหนฺติ, จตูหิ จ โสตาปตฺติยงฺเคหิ สมนฺนาคโต โหติ, อริโย จสฺส ญาโย ปญฺญาย สุทิฏฺโฐ โหติ สุปฺปฏิวิทฺโธ, โส อากงฺขมาโน อตฺตนาว อตฺตานํ พฺยากเรยฺย, ขีณนิรโยมฺหิ}

\csromanheader{2}{11. โสตาปตฺติสํยุตฺต}
\prose{\csromanfolio{339}ขีณติรจฺฉานโยนิ\footnote{ขีณติรจฺฉานโยนิโย (สพฺพตฺถ)} ขีณาเปตฺติวิสโย ขีณาปายทุคฺคติวินิปาโต, โสตาปนฺโนหมสฺมิ อวินิปาตธมฺโม นิยโต สมฺโพธิปรายโณ.}

\prose{กตมานิ ปญฺจ ภยานิ เวรานิ วูปสนฺตานิ โหนฺติ, ยํ คหปติ ปาณาติปาตี ปาณาติปาตปฺปจฺจยา ทิฏฺฐธมฺมิกมฺปิ ภยํ เวรํ ปสวติ, สมฺปรายิกมฺปิ ภยํ เวรํ ปสวติ, เจตสิกมฺปิ ทุกฺขํ โทมนสฺสํ ปฏิสํเวทิยติ.\csromanspacer{} ปาณาติปาตา ปฏิวิรตสฺส เอวํ ตํ ภยํ เวรํ วูปสนฺตํ โหติ.\csromanspacer{} ยํ คหปติ อทินฺนาทายี ฯเปฯ.\csromanspacer{} ยํ คหปติ กาเมสุมิจฺฉาจารี ฯเปฯ.\csromanspacer{} ยํ คหปติ มุสาวาที ฯเปฯ.\csromanspacer{} ยํ คหปติ สุราเมรยมชฺชปฺปมาทฏฺฐายี สุราเมรยมชฺชปฺปมาทฏฺฐานปฺปจฺจยา ทิฏฺฐธมฺมิกมฺปิ ภยํ เวรํ ปสวติ, สมฺปรายิกมฺปิ ภยํ เวรํ ปสวติ, เจตสิกมฺปิ ทุกฺขํ โทมนสฺสํ ปฏิสํเวทิยติ, สุราเมรยมชฺชปฺปมาทฏฺฐานา ปฏิวิรตสฺส เอวํ ตํ ภยํ เวรํ วูปสนฺตํ โหติ.\csromanspacer{} อิมานิ ปญฺจ ภยานิ เวรานิ วูปสนฺตานิ โหนฺติ.}

\prose{กตเมหิ จตูหิ โสตาปตฺติยงฺเคหิ สมนฺนาคโต โหติ, อิธ คหปติ อริยสาวโก พุทฺเธ อเวจฺจปฺปสาเทน สมนฺนาคโต โหติ “อิติปิ โส ภควา ฯเปฯ สตฺถา เทวมนุสฺสานํ พุทฺโธ ภควา”ติ.\csromanspacer{} ธมฺเม ฯเปฯ.\csromanspacer{} สํเฆ ฯเปฯ.\csromanspacer{} อริยกนฺเตหิ สีเลหิ สมนฺนาคโต โหติ อขณฺเฑหิ ฯเปฯ สมาธิสํวตฺตนิเกหิ.\csromanspacer{} อิเมหิ จตูหิ โสตาปตฺติยงฺเคหิ สมนฺนาคโต โหติ.}

\prose{กตโม จสฺส อริโย ญาโย ปญฺญาย สุทิฏฺโฐ โหติ สุปฺปฏิวิทฺโธ, อิธ คหปติ อริยสาวโก ปฏิจฺจสมุปฺปาทญฺเญว สาธุกํ โยนิโส มนสิ กโรติ “อิติ อิมสฺมิํ สติ อิทํ โหติ, อิมสฺสุปฺปาทา อิทํ อุปฺปชฺชติ.\csromanspacer{} อิติ อิมสฺมิํ อสติ อิทํ น โหติ, อิมสฺส นิโรธา อิทํ นิรุชฺฌติ, ยทิทํ อวิชฺชาปจฺจยา สงฺขารา, สงฺขารปจฺจยา วิญฺญาณํ ฯเปฯ เอวเมตสฺส เกวลสฺส ทุกฺขกฺขนฺธสฺส สมุทโย โหติ.\csromanspacer{} อวิชฺชาย เตฺวว อเสสวิราคนิโรธา สงฺขารนิโรโธ ฯเปฯ เอวเมตสฺส เกวลสฺส ทุกฺขกฺขนฺธสฺส นิโรโธ โหติ.\csromanspacer{} อยมสฺส อริโย ญาโย ปญฺญาย สุทิฏฺโฐ โหติ สุปฺปฏิวิทฺโธ.}

\prose{ยโต โข คหปติ อริยสาวกสฺส อิมานิ ปญฺจ ภยานิ เวรานิ วูปสนฺตานิ โหนฺติ, อิเมหิ จตูหิ โสตาปตฺติยงฺเคหิ สมนฺนาคโต โหติ, อยญฺจสฺส อริโย ญาโย ปญฺญาย สุทิฏฺโฐ โหติ}

\prose{\csromanfolio{340}สุปฺปฏิวิทฺโธ, โส อากงฺขมาโน อตฺตนาว อตฺตานํ พฺยากเรยฺย “ขีณนิรโยมฺหิ ขีณติรจฺฉานโยนิ ขีณเปตฺติวิสโย ขีณาปายทุคฺคติวินิปาโต, โสตาปนฺโนหมสฺมิ อวินิปาตธมฺโม นิยโต สมฺโพธิปรายโณ”ติ.\csromanspacer{}\csromanspacer{}อฏฺฐมํ.}

\csromanheader{2}{9. ทุติยภยเวรูปสนฺตสุตฺต}
\proseitem{1025}{สาวตฺถินิทานํ ฯเปฯ.\csromanspacer{} ยโต โข ภิกฺขเว อริยสาวกสฺส อิมานิ ปญฺจ ภยานิ เวรานิ วูปสนฺตานิ โหนฺติ, อิเมหิ จตูหิ โสตาปตฺติยงฺเคหิ สมนฺนาคโต โหติ, อยญฺจสฺส อริโย ญาโย ปญฺญาย สุทิฏฺโฐ โหติ สุปฺปฏิวิทฺโธ, โส อากงฺขมาโน อตฺตนาว อตฺตานํ พฺยากเรยฺย “ขีณนิรโยมฺหิ ขีณติรจฺฉานโยนิ ขีณเปตฺติวิสโย ขีณาปายทุคฺคติวินิปาโต, โสตาปนฺโนหมสฺมิ อวินิปาตธมฺโม นิยโต สมฺโพธิปรายโณ”ติ.\csromanspacer{}\csromanspacer{}นวมํ.}

\csromanmatikaanchor{mtk.242}
\csromantocmarkref{subparagraph}{10. นนฺทกลิจฺฉวิสุตฺต}{mtk.242}
\bookmark[dest={mtk.242},level=5]{10. นนฺทกลิจฺฉวิสุตฺต}
\csromanheader{6}{10. นนฺทกลิจฺฉวิสุตฺต}
\proseitem{1026}{เอกํ สมยํ ภควา เวสาลิยํ วิหรติ มหาวเน กูฏาคารสาลายํ.\csromanspacer{} อถ โข นนฺทโก ลิจฺฉาวิมหามตฺโต เยน ภควา เตนุปสงฺกมิ, อุปสงฺกมิตฺวา ภควนฺตํ อภิวาเทตฺวา เอกมนฺตํ นิสีทิ, เอกมนฺตํ นิสินฺนํ โข นนฺทกํ ลิจฺฉวิมหามตฺตํ ภควา เอตทโวจ\csromandash{}}

\prose{จตูหิ โข นนฺทก ธมฺเมหิ สมนฺนาคโต อริยสาวโก โสตาปนฺโน โหติ อวินิปาตธมฺโม นิยโต สมฺโพธิปรายโณ.\csromanspacer{} กตเมหิ จตูหิ, อิธ นนฺทก อริยสาวโก พุทฺเธ อเวจฺจปฺปสาเทน สมนฺนาคโต โหติ “อิติปิ โส ภควา ฯเปฯ สตฺถา เทวมนุสฺสานํ พุทฺโธ ภควา”ติ.\csromanspacer{} ธมฺเม ฯเปฯ.\csromanspacer{} สํเฆ ฯเปฯ.\csromanspacer{} อริยกนฺเตหิ สีเลหิ สมนฺนาคโต โหติ อขณฺเฑหิ ฯเปฯ สมาธิสํวตฺตนิเกหิ.\csromanspacer{} อิเมหิ โข นนฺทก จตูหิ ธมฺเมหิ สมนฺนาคโต อริยสาวโก โสตาปนฺโน โหติ อวินิปาตธมฺโม นิยโต สมฺโพธิปรายโณ.}

\prose{อิเมหิ จ ปน นนฺทก จตูหิ ธมฺเมหิ สมนฺนาคโต อริยสาวโก อายุนา สํยุตฺโต โหติ ทิพฺเพนปิ มานุเสนปิ, วณฺเณน สํยุตฺโต โหติ ทิพฺเพนปิ มานุเสนปิ, สุเขน สํยุตฺโต โหติ ทิพฺเพนปิ มานุเสนปิ, ยเสน สํยุตฺโต โหติ ทิพฺเพนปิ มานุเสนปิ, อาธิปเตยฺเยน สํยุตฺโต โหติ ทิพฺเพนปิ มานุเสนปิ.\csromanspacer{} ตํ โข ปนาหํ}

\csromanheader{2}{11. โสตาปตฺติสํยุตฺต}
\csromanmatikaanchor{mtk.243}
\csromanmatikaanchor{mtk.244}
\csromantocmarknopage{subparagraph}{4. ปุญฺญาภิสนฺทวคฺค}{mtk.243}
\bookmark[dest={mtk.243},level=5]{4. ปุญฺญาภิสนฺทวคฺค}
\csromantocmarkref{subparagraph}{1-3. ปุญฺญาภิสนฺทสุตฺตานิ}{mtk.244}
\bookmark[dest={mtk.244},level=5]{1-3. ปุญฺญาภิสนฺทสุตฺตานิ}
\prose{\csromanfolio{341}นนฺทก นาญฺญสฺส สมณสฺส วา พฺราหฺมณสฺส วา สุตฺวา วทามิ, อปิ จ ยเทว มยา สามํ ญาตํ สามํ ทิฏฺฐํ สามํ วิทิตํ, ตเทวาหํ วทามีติ.}

\prose{เอวํ วุตฺเต อญฺญตโร ปุริโส นนฺทกํ ลิจฺฉวิมหามตฺตํ เอตทโวจ “นหานกาโล ภนฺเต”ติ.\csromanspacer{} อลํ ทานิ ภเณ เอเตน พาหิเรน นหาเนน, อลมิทํ อชฺฌตฺตํ นหานํ ภวิสฺสติ ยทิทํ ภควติ ปสาโทติ.\csromanspacer{}\csromanspacer{}ทสมํ.\csromanspacer{} สรณานิวคฺโค ตติโย.}

\titlehead{\textbf{ตสฺสุทฺทานํ}}
\csromangathameasure{มหานาเมน เทฺว วุตฺตา, โคธา จ สรณา ทุเว. \\ ทุเว อนาถปิณฺฑิกา, ทุเว เวรภเยน จ. \\ ลิจฺฉวี ทสโม วุตฺโต, วคฺโค เตน ปวุจฺจตีติ.}
\csromangathagroup{\csromangathastanza{มหานาเมน เทฺว วุตฺตา, โคธา จ สรณา ทุเว. \\ ทุเว อนาถปิณฺฑิกา, ทุเว เวรภเยน จ. \\ ลิจฺฉวี ทสโม วุตฺโต, วคฺโค เตน ปวุจฺจตีติ.}}

\csromanheader{6}{4. ปุญฺญาภิสนฺทวคฺค}
\csromanheader{2}{1. ปฐมปุญฺญาภิสนฺทสุตฺต}
\proseitem{1027}{สาวตฺถินิทานํ.\csromanspacer{} จตฺตาโรเม ภิกฺขเว ปุญฺญาภิสนฺทา กุสลาภิสนฺทา สุขสฺสาหารา.\csromanspacer{} กตเม จตฺตาโร, อิธ ภิกฺขเว อริยสาวโก พุทฺเธ อเวจฺจปฺปสาเทน สมนฺนาคโต โหติ “อิติปิ โส ภควา ฯเปฯ สตฺถา เทวมนุสฺสานํ พุทฺโธ ภควา”ติ, อยํ ปฐโม ปุญฺญาภิสนฺโท กุสลาภิสนฺโท สุขสฺสาหาโร.}

\prose{ปุน จปรํ ภิกฺขเว อริยสาวโก ธมฺเม อเวจฺจปฺปสาเทน สมนฺนาคโต โหติ “สฺวากฺขาโต ภควตา ธมฺโม ฯเปฯ ปจฺจตฺตํ เวทิตพฺโพ วิญฺญูหี”ติ, อยํ ทุติโย ปุญฺญาภิสนฺโท กุสลาภิสนฺโท สุขสฺสาหาโร.}

\prose{ปุน จปรํ ภิกฺขเว อริยสาวโก สํเฆ อเวจฺจปฺปสาเทน สมนฺนาคโต โหติ “สุปฺปฏิปนฺโน ภควโต สาวกสํโฆ ฯเปฯ อนุตฺตรํ ปุญฺญกฺเขตฺตํ โลกสฺสา”ติ, อยํ ตติโย ปุญฺญาภิสนฺโท กุสลาภิสนฺโท สุขสฺสาหาโร.}

\prose{\csromanfolio{342}ปุน จปรํ ภิกฺขเว อริยสาวโก อริยกนฺเตหิ สีเลหิ สมนฺนาคโต โหติ อขณฺเฑหิ ฯเปฯ สมาธิสํวตฺตนิเกหิ, อยํ จตุตฺโถ ปุญฺญาภิสนฺโท กุสลาภิสนฺโท สุขสฺสาหาโร.\csromanspacer{} อิเม โข ภิกฺขเว จตฺตาโร ปุญฺญาภิสนฺทา กุสลาภิสนฺทา สุขสฺสาหาราติ.}

\csromanheader{2}{2. ทุติยปุญฺญาภิสนฺทสุตฺต}
\proseitem{1028}{จตฺตาโรเม ภิกฺขเว ปุญฺญาภิสนฺทา กุสลาภิสนฺทา สุขสฺสาหารา.\csromanspacer{} กตเม จตฺตาโร, อิธ ภิกฺขเว อริยสาวโก พุทฺเธ อเวจฺจปฺปสาเทน สมนฺนาคโต โหติ “อิติปิ โส ภควา ฯเปฯ สตฺถา เทวมนุสฺสานํ พุทฺโธ ภควา”ติ, อยํ ปฐโม ปุญฺญาภิสนฺโท กุสลาภิสนฺโท สุขสฺสาหาโร.}

\prose{ปุน จปรํ ภิกฺขเว อริยสาวโก ธมฺเม ฯเปฯ.\csromanspacer{} สํเฆ ฯเปฯ.}

\prose{ปุน จปรํ ภิกฺขเว อริยสาวโก วิคตมลมจฺเฉเรน เจตสา อคารํ อชฺฌาวสติ มุตฺตจาโค ปยตปาณิ โวสฺสคฺครโต ยาจโยโค ทานสํวิภาครโต, อยํ จตุตฺโถ ปุญฺญาภิสนฺโท กุสลาภิสนฺโท สุขสฺสาหาโร.\csromanspacer{} อิเม โข ภิกฺขเว จตฺตาโร ปุญฺญาภิสนฺทา กุสลาภิสนฺทา สุขสฺสาหาราติ.\csromanspacer{}\csromanspacer{}ทุติยํ.}

\csromanheader{2}{3. ตติยปุญฺญาภิสนฺทสุตฺต}
\proseitem{1029}{จตฺตาโรเม ภิกฺขเว ปุญฺญาภิสนฺทา กุสลาภิสนฺทา สุขสฺสาหารา.\csromanspacer{} กตเม จตฺตาโร, อิธ ภิกฺขเว อริยสาวโก พุทฺเธ อเวจฺจปฺปสาเทน สมนฺนาคโต โหติ “อิติปิ โส ภควา ฯเปฯ สตฺถา เทวมนุสฺสานํ พุทฺโธ ภควา”ติ, อยํ ปฐโม ปุณฺณาภิสนฺโท กุลสาภิสนฺโท สุขสฺสาหาโร.}

\prose{ปุน จปรํ ภิกฺขเว อริยสาวโก ธมฺเม ฯเปฯ.\csromanspacer{} สํเฆ ฯเปฯ.}

\prose{ปุน จปรํ ภิกฺขเว อริยสาวโก ปญฺญวา โหติ อุทยตฺถคามินิยา ปญฺญาย สมนฺนาคโต อริยาย นิพฺเพธิกาย สมฺมา ทุกฺขกฺขยคามินิยา, อยํ จตุตฺโถ ปุญฺญาภิสนฺโท กุสลาภิสนฺโท สุขสฺสาหาโร.\csromanspacer{} อิเม โข ภิกฺขเว จตฺตาโร ปุญฺญาภิสนฺทา กุสลาภิสนฺทา สุขสฺสาหาราติ.\csromanspacer{}\csromanspacer{}ตติยํ.}

\csromanheader{2}{11. โสตาปตฺติสํยุตฺต}
\csromanheader{2}{4. ปฐมเทวปทสุตฺต}
\csromanmatikaanchor{mtk.245}
\csromantocmarkref{subparagraph}{4-6. เทวปทสุตฺตานิ}{mtk.245}
\bookmark[dest={mtk.245},level=5]{4-6. เทวปทสุตฺตานิ}
\proseitem{1030}{\csromanfolio{343}สาวตฺถินิทานํ.\csromanspacer{} จตฺตาริมานิ ภิกฺขเว เทวานํ เทวปทานิ อวิสุทฺธานํ สตฺตานํ วิสุทฺธิยา อปริโยทาตานํ สตฺตานํ ปริโยทปนาย.}

\prose{กตมานิ จตฺตาริ, อิธ ภิกฺขเว อริยสาวโก พุทฺเธ อเวจฺจปฺปสาเทน สมนฺนาคโต โหติ “อิติปิ โส ภควา ฯเปฯ สตฺถา เทวมนุสฺสานํ พุทฺโธ ภควา”ติ, อิทํ ปฐมํ เทวปทํ อวิสุทฺธานํ สตฺตานํ วิสุทฺธิยา อปริโยทาตานํ สตฺตานํ ปริโยทปนาย.}

\prose{ปุน จปรํ ภิกฺขเว อริยสาวโก ธมฺเม ฯเปฯ.\csromanspacer{} สํเฆ ฯเปฯ.}

\prose{ปุน จปรํ ภิกฺขเว อริยสาวโก อริยกนฺเตหิ สีเลหิ สมนฺนาคโต โหติ อขณฺเฑหิ ฯเปฯ สมาธิสํวตฺตนิเกหิ.\csromanspacer{} อิทํ จตุตฺถํ เทวานํ เทวปทํ อวิสุทฺธานํ สตฺตานํ วิสุทฺธิยา อปริโยทาตานํ สตฺตานํ ปริโยทปนาย.\csromanspacer{} อิมานิ โข ภิกฺขเว จตฺตาริ เทวานํ เทวปทานิ อวิสุทฺธานํ สตฺตานํ วิสุทฺธิยา อปริโยทาตานํ สตฺตานํ ปริโยทปนายาติ.\csromanspacer{}\csromanspacer{}จตุตฺถํ.}

\csromanheader{2}{5. ทุติยเทวปทสุตฺต}
\proseitem{1031}{จตฺตาริมานิ ภิกฺขเว เทวานํ เทวปทานิ อวิสุทฺธานํ สตฺตานํ วิสุทฺธิยา อปริโยทาตานํ สตฺตานํ ปริโยทปนาย.}

\prose{กตมานิ จตฺตาริ, อิธ ภิกฺขเว อริยสาวโก พุทฺเธ อเวจฺจปฺปสาเทน สมนฺนาคโต โหติ “อิติปิ โส ภควา ฯเปฯ สตฺถา เทวมนุสฺสานํ พุทฺโธ ภควา”ติ.\csromanspacer{} โส อิติ ปฏิสญฺจิกฺขติ “กิํ นุ โข เทวานํ เทวปทนฺ”ติ.\csromanspacer{} โส เอวํ ปชานาติ “อพฺยาพชฺฌปรเม ขฺวาหํ เอตรหิ เทเว สุณามิ, น จ โข ปนาหํ กิญฺจิ พฺยาพาเธมิ ตสํ วา ถาวรํ วา, อทฺธาหํ เทวปทธมฺมสมนฺนาคโต วิหรามี”ติ, อิทํ ปฐมํ เทวานํ เทวปทํ อวิสุทฺธานํ สตฺตานํ วิสุทฺธิยา อปริโยทาตานํ สตฺตานํ ปริโยทปนาย.}

\prose{ปุน จปรํ ภิกฺขเว อริยสาวโก ธมฺเม ฯเปฯ.\csromanspacer{} สํเฆ ฯเปฯ.}

\prose{ปุน จปรํ ภิกฺขเว อริยสาวโก อริยกนฺเตหิ สีเลหิ สมนฺนาคโต โหติ อขณฺเฑหิ ฯเปฯ สมาธิสํวตฺตนิเกหิ, โส อิติ ปฏิสญฺจิกฺขติ “กิํ นุ โข เทวานํ เทวปทนฺ”ติ.\csromanspacer{} โส เอวํ ปชานาติ \csromanfolio{344}“อพฺยาพชฺฌปรเม ขฺวาหํ เอตรหิ เทเว สุณามิ, น โข ปนาหํ กิญฺจิ พฺยาพาเธมิ ตสํ วา ถาวรํ วา, อทฺธาหํ เทวปทธมฺมสมนฺนาคโต วิหรามี”ติ, อิทํ จตุตฺถํ เทวานํ เทวปทํ อวิสุทฺธานํ สตฺตานํ วิสุทฺธิยา อปริโยทาตานํ สตฺตานํ ปริโยทปนาย.\csromanspacer{} อิมานิ โข ภิกฺขเว จตฺตาริ เทวานํ เทวปาทานิ อวิสุทฺธานํ สตฺตานํ วิสุทฺธิยา อปริโยทาตานํ สตฺตานํ ปริโยทปนายาติ.\csromanspacer{}\csromanspacer{}ปญฺจมํ.}

\csromanheader{2}{6. เทวสภาคสุตฺต}
\proseitem{1032}{จตูหิ ภิกฺขเว ธมฺเมหิ สมนฺนาคตํ อตฺตมนา เทวา สภาคตํ กเถนฺติ.\csromanspacer{} กตเมหิ จตูหิ, อิธ ภิกฺขเว อริยสาวโก พุทฺเธ อเวจฺจปฺปสาเทน สมนฺนาคโต โหติ “อิติปิ โส ภควา ฯเปฯ สตฺถา เทวมนุสฺสานํ พุทฺโธ ภควา”ติ.\csromanspacer{} ยา ตา เทวตา พุทฺเธ อเวจฺจปฺปสาเทน สมนฺนาคตา อิโต จุตา ตตฺรูปปนฺนา, ตาสํ เอวํ โหติ “ยถารูเปน โข มยํ พุทฺเธ อเวจฺจปฺปสาเทน สมนฺนาคตา ตโต จุตา อิธูปปนฺนา อริยสาวโกปิ ตถารูเปน พุทฺเธ อเวจฺจปฺปสาเทน สมนฺนาคโต เอหีติ เทวานํ สนฺติเก”ติ.}

\prose{ปุน จปรํ ภิกฺขเว อริยสาวโก ธมฺเม ฯเปฯ.\csromanspacer{} สํเฆ ฯเปฯ.\csromanspacer{} อริยกนฺเตหิ สีเลหิ สมนฺนาคโต โหติ อขณฺเฑหิ ฯเปฯ สมาธิสํวตฺตนิเกหิ, ยา ตา เทวตา อริยกนฺเตหิ สีเลหิ สมนฺนาคตา อิโต จุตา ตตฺรูปปนฺนา, ตาสํ เอวํ โหติ “ยถารูเปหิ โข มยํ อริยกนฺเตหิ สีเลหิ สมนฺนาคตา ตโต จุตา อิธูปปนฺนา อริยสาวโกปิ ตถารูเปหิ อริยกนฺเตหิ สีเลหิ สมนฺนาคโต เอหีติ เทวานํ สนฺติเก”ติ.\csromanspacer{} อิเมหิ โข ภิกฺขเว จตูหิ สมนฺนาคตํ อตฺตมนา เทวา สภาคตํ กเถนฺตีติ.\csromanspacer{}\csromanspacer{}ฉฏฺฐํ.}

\csromanmatikaanchor{mtk.246}
\csromantocmarkref{subparagraph}{7. มหานามสุตฺต}{mtk.246}
\bookmark[dest={mtk.246},level=5]{7. มหานามสุตฺต}
\csromanheader{6}{7. มหานามสุตฺต}
\proseitem{1033}{เอกํ สมยํ ภควา สกฺเกสุ วิหรติ กปิลวตฺถุสฺมิํ นิโคฺรธาราเม.\csromanspacer{} อถ โข มหานาโม สกฺโก เยน ภควา เตนุปสงฺกมิ, อุปสงฺกมิตฺวา ภควนฺตํ อภิวาเทตฺวา เอกมนฺตํ นิสีทิ, เอกมนฺตํ นิสินฺโน โข มหานาโม สกฺโก ภควนฺตํ เอตทโวจ\csromandash{}}

\prose{กิตฺตาวตา นุ โข ภนฺเต อุปาสโก โหตีติ.\csromanspacer{} ยโต โข มหานาม พุทฺธํ สรณํ คโต โหติ, ธมฺมํ สรณํ คโต โหติ,}

\csromanheader{2}{11. โสตาปตฺติสํยุตฺต}
\prose{\csromanfolio{345}สํฆํ สรณํ คโต โหติ.\csromanspacer{} เอตฺตาวตา โข มหานาม อุปาสโก โหตีติ.}

\prose{กิตฺตาวตา ปน ภนฺเต อุปาสโก สีลสมฺปนฺโน โหตีติ.\csromanspacer{} ยโต โข มหานาม อุปาสโก ปาณาติปาตา ปฏิวิรโต โหติ, อทินฺนาทานา ปฏิวิรโต โหติ, กาเมสุมิจฺฉาจารา ปฏิวิรโต โหติ, มุสาวาทา ปฏิวิรโต โหติ, สุราเมรยมชฺชปฺปมาทฏฺฐานา ปฏิวิรโต โหติ.\csromanspacer{} เอตฺตาวตา โข มหานาม อุปาสโก สีลสมฺปนฺโน โหตีติ.}

\prose{กิตฺตาวตา ปน ภนฺเต อุปาสโก สทฺธาสมฺปนฺโน โหตีติ, อิธ มหานาม อุปาสโก สทฺโธ โหติ, สทฺทหติ ตถาคตสฺส โพธิํ “อิติปิ โส ภควา ฯเปฯ สตฺถา เทวมนุสฺสานํ พุทฺโธ ภควา”ติ.\csromanspacer{} เอตฺตาวตา โข มหานาม อุปาสโก สทฺธาสมฺปนฺโน โหตีติ.}

\prose{กิตฺตาวตา ปน ภนฺเต อุปาสโก จาคสมฺปนฺโน โหตีติ.\csromanspacer{} อิธ มหานาม อุปาสโก วิคตมลมจฺเฉเรน เจตสา อคารํ อชฺฌาวสติ มุตฺตจาโค ปยตปาณิ โวสฺสคฺครโต ยาจโยโค ทานสํวิภาครโต.\csromanspacer{} เอตฺตาวตา โข มหานาม อุปาสโก จาคสมฺปนฺโน โหตีติ.}

\prose{กิตฺตาวตา ปน ภนฺเต อุปาสโก ปญฺญาสมฺปนฺโน โหตีติ.\csromanspacer{} อิธ มหานาม อุปาสโก ปญฺญวา โหติ, อุทยตฺถคามินิยา ปญฺญาย สมนฺนาคโต อริยาย นิพฺเพธิกาย สมฺมา ทุกฺขกฺขยคามินิยา.\csromanspacer{} เอตฺตาวตา โข มหานาม อุปาสโก ปญฺญาสมฺปนฺโน โหตีติ.\csromanspacer{}\csromanspacer{}สตฺตมํ.}

\csromanmatikaanchor{mtk.247}
\csromantocmarkref{subparagraph}{8. วสฺสสุตฺต}{mtk.247}
\bookmark[dest={mtk.247},level=5]{8. วสฺสสุตฺต}
\csromanheader{6}{8. วสฺสสุตฺต}
\proseitem{1034}{เสยฺยถาปิ ภิกฺขเว อุปริปพฺพเต ถุลฺลผุสิตเก เทเว วสฺสนฺเต ตํ อุทกํ ยถานินฺนํ ปวตฺตมานํ ปพฺพตกนฺทรปทรสาขา ปริปูเรติ, ปพฺพตกนฺทรปทรสาขา ปริปูรา กุโสพฺเภ ปริปูเรนฺติ, กุโสพฺภา ปริปูรา มหาโสพฺเภ ปริปูเรนฺติ, มหาโสพฺภา ปริปูรา กุนฺนทิโย ปริปูเรนฺติ, กุนฺนทิโย ปริปูรา มหานทิโย ปริปูเรนฺติ, มหานทิโย ปริปูรา มหาสมุทฺทํ\footnote{มหาสมุทฺทสาครํ (สพฺพตฺถ) สํ 1. 269 ปิฏฺเฐปิ.} ปริปูเรนฺติ.\csromanspacer{} เอวเมว โข ภิกฺขเว อริยสาวกสฺส โย จ พุทฺเธ อเวจฺจปฺปสาโท โย จ ธมฺเม อเวจฺจปฺปสาโท โย จ \csromanfolio{346}สํเฆ อเวจฺจปฺปสาโท ยานิ จ อริยกนฺตานิ สีลานิ, อิเม ธมฺมา สนฺทมานา ปารํ คนฺตฺวา อาสวานํ ขยาย สํวตฺตนฺตีติ.\csromanspacer{}\csromanspacer{}อฏฺฐมํ.}

\csromanmatikaanchor{mtk.248}
\csromantocmarkref{subparagraph}{9. กาฬิโคธสุตฺต}{mtk.248}
\bookmark[dest={mtk.248},level=5]{9. กาฬิโคธสุตฺต}
\csromanheader{6}{9. กาฬิโคธสุตฺต}
\proseitem{1035}{เอกํ สมยํ ภควา สกฺเกสุ วิหรติ กปิลวตฺถุสฺมิํ นิโคฺรธาราเม.\csromanspacer{} อถ โข ภควา ปุพฺพณฺหสมยํ นิวาเสตฺวา ปตฺตจีวรมาทาย เยน กาฬิโคธาย สากิยานิยา นิเวสนํ เตนุปสงฺกมิ, อุปสงฺกมิตฺวา ปญฺญตฺเต อาสเน นิสีทิ.\csromanspacer{} อถ โข กาฬิโคธา สากิยานี เยน ภควา เตนุปสงฺกมิ, อุปสงฺกมิตฺวา ภควนฺตํ อภิวาเทตฺวา เอกมนฺตํ นิสีทิ, เอกมนฺตํ นิสินฺนํ โข กาฬิโคธํ สากิยานิํ ภควา เอตทโวจ\csromandash{}}

\prose{จตูหิ โข โคเธ ธมฺเมหิ สมนฺนาคตา อริยสาวิกา โสตาปนฺนา โหติ อวินิปาตธมฺมา นิยตา สมฺโพธิปรายณา.\csromanspacer{} กตเมหิ จตูหิ, อิธ โคเธ อริยสาวิกา พุทฺเธ อเวจฺจปฺปสาเทน สมนฺนาคตา โหติ “อิติปิ โส ภควา ฯเปฯ สตฺถา เทวมนุสฺสานํ พุทฺโธ ภควา”ติ.\csromanspacer{} ธมฺเม ฯเปฯ.\csromanspacer{} สํเฆ ฯเปฯ.\csromanspacer{} วิคตมลมจฺเฉเรน เจตสา อคารํ อชฺฌาวสติ มุตฺตจาคา ปยตปาณินี\footnote{ปยตปาณี (สพฺพตฺถ) 3. 30 โมคฺคลฺลานสุตฺตํ โอโลเกตพฺพํ.} โวสฺสคฺครตา ยาจโยคา ทานสํวิภาครตา.\csromanspacer{} อิเมหิ โข โคเธ จตูหิ ธมฺเมหิ สมนฺนาคตา อริยสาวิกา โสตาปนฺนา โหติ อวินิปาตธมฺมา นิยตา สมฺโพธิปรายณาติ.}

\prose{ยานิมานิ ภนฺเต ภควตา จตฺตาริ โสตาปตฺติยงฺคานิ เทสิตานิ, สํวิชฺชนฺเตเต ธมฺมา มยิ, อหญฺจ เตสุ ธมฺเมสุ สนฺทิสฺสามิ, อหํ หิ ภนฺเต พุทฺเธ อเวจฺจปฺปสาเทน สมนฺนาคตา “อิติปิ โส ภควา ฯเปฯ สตฺถา เทวมนุสฺสานํ พุทฺโธ ภควา”ติ.\csromanspacer{} ธมฺเม ฯเปฯ.\csromanspacer{} สํเฆ ฯเปฯ.\csromanspacer{} ยํ โข ปน กิญฺจิ กุเล เทยฺยธมฺมํ, สพฺพํ ตํ อปฺปฏิวิภตฺตํ สีลวนฺเตหิ กลฺยาณธมฺเมหีติ.\csromanspacer{} ลาภา เต โคเธ, สุลทฺธํ เต โคเธ, โสตาปตฺติผลํ ตยา โคเธ พฺยากตนฺติ.\csromanspacer{}\csromanspacer{}นวมํ.}

\csromanmatikaanchor{mtk.249}
\csromantocmarkref{subparagraph}{10. นนฺทิยสกฺกสุตฺต}{mtk.249}
\bookmark[dest={mtk.249},level=5]{10. นนฺทิยสกฺกสุตฺต}
\csromanheader{6}{10. นนฺทิยสกฺกสุตฺต}
\proseitem{1036}{เอกํ สมยํ ภควา สกฺเกสุ วิหรติ กปิลวตฺถุสฺมิํ นิโคฺรธาราเม.\csromanspacer{} อถ โข นนฺทิโย สกฺโก เยน ภควา เตนุปสงฺกมิ, อุปสงฺกมิตฺวา ภควนฺตํ อภิวาเทตฺวา เอกมนฺตํ นิสีทิ, เอกมนฺตํ นิสินฺโน}

\csromanheader{2}{11. โสตาปตฺติสํยุตฺต}
\prose{\csromanfolio{347}โข นนฺทิโย สกฺโก ภควนฺตํ เอตทโวจ “ยสฺเสว นุ โข ภนฺเต อริยสาวกสฺส จตฺตาริ โสตาปตฺติยงฺคานิ สพฺเพน สพฺพํ สพฺพถา สพฺพํ นตฺถิ, เสฺวว นุ โข ภนฺเต อริยสาวโก ปมาทวิหารี”ติ.}

\prose{ยสฺส โข นนฺทิย จตฺตาริ โสตาปตฺติยงฺคานิ สพฺเพน สพฺพํ สพฺพถา สพฺพํ นตฺถิ, ตมหํ พาหิโร ปุถุชฺชนปกฺเข ฐิโตติ วทามิ, อปิ จ นนฺทิย ยถา อริยสาวโก ปมาทวิหารี เจว โหติ อปฺปมาทวิหารี จ, ตํ สุณาหิ สาธุกํ มนสิ กโรหิ, ภาสิสฺสามีติ. “เอวํ ภนฺเต”ติ โข นนฺทิโย สกฺโก ภควโต ปจฺจสฺโสสิ.\csromanspacer{} ภควา เอตทโวจ\csromandash{}}

\prose{กถญฺจ นนฺทิย อริยสาวโก ปมาทวิหารี โหติ, อิธ นนฺทิย อริยสาวโก พุทฺเธ อเวจฺจปฺปสาเทน สมนฺนาคโต โหติ “อิติปิ โส ภควา ฯเปฯ สตฺถา เทวมนุสฺสานํ พุทฺโธ ภควา”ติ.\csromanspacer{} โส เตน พุทฺเธ อเวจฺจปฺปสาเทน สนฺตุฏฺโฐ น อุตฺตริ วายมติ ทิวา ปวิเวกาย รตฺติํ ปฏิสลฺลานาย, ตสฺส เอวํ ปมตฺตสฺส วิหรโต ปาโมชฺชํ น โหติ, ปาโมชฺเช อสติ ปีติ น โหติ, ปีติยา อสติ ปสฺสทฺธิ น โหติ, ปสฺสทฺธิยา อสติ ทุกฺขํ วิหรติ, ทุกฺขิโน จิตฺตํ น สมาธิยติ, อสมาหิเต จิตฺเต ธมฺมา น ปาตุภวนฺติ, ธมฺมานํ อปาตุภาวา ปมาทวิหารีเตฺวว สงฺขฺยํ คจฺฉติ.}

\prose{ปุน จปรํ นนฺทิย อริยสาวโก ธมฺเม ฯเปฯ.\csromanspacer{} สํเฆ ฯเปฯ.\csromanspacer{} อริยกนฺเตหิ สีเลหิ สมนฺนาคโต โหติ อขณฺเฑหิ ฯเปฯ สมาธิสํวตฺตนิเกหิ, โส เตหิ อริยกนฺเตหิ สีเลหิ สนฺตุฏฺโฐ น อุตฺตริ วายมติ ทิวา ปวิเวกาย รตฺติํ ปฏิสลฺลานาย, ตสฺส เอวํ ปมตฺตสฺส วิหรโต ปาโมชฺชํ น โหติ, ปาโมชฺเช อสติ ปีติ น โหติ, ปีติยา อสติ ปสฺสทฺธิ น โหติ, ปสฺสทฺธิยา อสติ ทุกฺขํ วิหรติ, ทุกฺขิโน จิตฺตํ น สมาธิยติ, อสมาหิเต จิตฺเต ธมฺมา น ปาตุภวนฺติ, ธมฺมานํ อปาตุภาวา ปมาทวิหารีเตฺวว สงฺขฺยํ คจฺฉติ.\csromanspacer{} เอวํ โข นนฺทิย อริยสาวโก ปมาทวิหารี โหติ.}

\prose{กถญฺจ นนฺทิย อริยสาวโก อปฺปมาทวิหารี โหติ, อิธ นนฺทิย อริยสาวโก พุทฺเธ อเวจฺจปฺปสาเทน สมนฺนาคโต โหติ “อิติปิ โส ภควา ฯเปฯ สตฺถา เทวมนุสฺสานํ พุทฺโธ ภควา”ติ, โส เตน พุทฺเธ อเวจฺจปฺปสาเทน อสนฺตุฏฺโฐ อุตฺตริ วายมติ ทิวา ปวิเวกาย รตฺติํ ปฏิสลฺลานาย, ตสฺส เอวํ อปฺปมตฺตสฺส วิหรโต ปาโมชฺชํ ชายติ,}

\csromanmatikaanchor{mtk.250}
\csromanmatikaanchor{mtk.251}
\csromantocmarknopage{subparagraph}{5. สคาถกปุญฺญาภิสนฺทวคฺค}{mtk.250}
\bookmark[dest={mtk.250},level=5]{5. สคาถกปุญฺญาภิสนฺทวคฺค}
\csromantocmarkref{subparagraph}{1-3. อภิสนฺทสุตฺตานิ}{mtk.251}
\bookmark[dest={mtk.251},level=5]{1-3. อภิสนฺทสุตฺตานิ}
\prose{\csromanfolio{348}ปมุทิตสฺส ปีติ ชายติ, ปีติมนสฺส กาโย ปสฺสมฺภติ, ปสฺสทฺธกาโย สุขํ เวทิยติ, สุขิโน จิตฺตํ สมาธิยติ, สมาหิเต จิตฺเต ธมฺมา ปาตุภวนฺติ, ธมฺมานํ ปาตุภาวา อปฺปมาทวิหารีเตฺวว สงฺขฺยํ คจฺฉติ.}

\prose{ปุน จปรํ นนฺทิย อริยสาวโก ธมฺเม ฯเปฯ.\csromanspacer{} สํเฆ ฯเปฯ.\csromanspacer{} อริยกนฺเตหิ สีเลหิ สมนฺนาคโต โหติ อขณฺเฑหิ ฯเปฯ สมาธิสํวตฺตนิเกหิ.\csromanspacer{} โส เตหิ อริยกนฺเตหิ สีเลหิ อสนฺตุฏฺโฐ อุตฺตริ วายมติ ทิวา ปวิเวกาย รตฺติํ ปฏิสลฺลานาย, ตสฺส เอวํ อปฺปมตฺตสฺส วิหรโต ปาโมชฺชํ ชายติ, ปมุทิตสฺส ปีติ ชายติ, ปีติมนสฺส กาโย ปสฺสมฺภติ, ปสฺสทฺธกาโย สุขํ เวทิยติ, สุขิโน จิตฺตํ สมาธิยติ, สมาหิเต จิตฺเต ธมฺมา ปาตุภวนฺติ, ธมฺมานํ ปาตุภาวา อปฺปมาทวิหารีเตฺวว สงฺขฺยํ คจฺฉติ.\csromanspacer{} เอวํ โข นนฺทิย อริยสาวโก อปฺปมาทวิหารี โหตีติ.\csromanspacer{}\csromanspacer{}ทสมํ.\csromanspacer{} ปุญฺญาภิสนฺทวคฺโค จตุตฺโถ.}

\titlehead{\textbf{ตสฺสุทฺทานํ}}
\csromangathameasure{อภิสนฺทา ตโย วุตฺตา, ทุเว เทวปทานิ จ. \\ สภาคตํ มหานาโม, วสฺสํ กาฬี จ นนฺทิยาติ.}
\csromangathagroup{\csromangathastanza{อภิสนฺทา ตโย วุตฺตา, ทุเว เทวปทานิ จ. \\ สภาคตํ มหานาโม, วสฺสํ กาฬี จ นนฺทิยาติ.}}

\csromanheader{6}{5. สคาถกปุญฺญาภิสนฺทวคฺค}
\csromanheader{2}{1. ปฐม-อภิสนฺทสุตฺต}
\proseitem{1037}{จตฺตาโรเม ภิกฺขเว ปุญฺญาภิสนฺทา กุสลาภิสนฺทา สุขสฺสาหารา.\csromanspacer{} กตเม จตฺตาโร, อิธ ภิกฺขเว อริยสาวโก พุทฺเธ อเวจฺจปฺปสาเทน สมนฺนาคโต โหติ “อิติปิ โส ภควา ฯเปฯ สตฺถา เทวมนุสฺสานํ พุทฺโธ ภควา”ติ, อยํ ปฐโม ปุญฺญาภิสนฺโท กุสลาภิสนฺโท สุขสฺสาหาโร.}

\prose{ปุน จปรํ ภิกฺขเว อริยสาวโก ธมฺเม ฯเปฯ.\csromanspacer{} สํเฆ ฯเปฯ.}

\prose{ปุน จปรํ ภิกฺขเว อริยสาวโก อริยกนฺเตหิ สีเลหิ สมนฺนาคโต โหติ อขณฺเฑหิ ฯเปฯ สมาธิสํวตฺตนิเกหิ, อยํ}

\csromanheader{2}{11. โสตาปตฺติสํยุตฺต}
\prose{\csromanfolio{349}จตุตฺโถ ปุญฺญาภิสนฺโท กุสลาภิสนฺโท สุขสฺสาหาโร.\csromanspacer{} อิเม โข ภิกฺขเว จตฺตาโร ปุญฺญาภิสนฺทา กุสลาภิสนฺทา สุขสฺสาหารา.}

\prose{อิเมหิ โข ภิกฺขเว จตูหิ ปุญฺญาภิสนฺเทหิ กุสลาภิสนฺเทหิ สมนฺนาคตสฺส อริยสาวกสฺส น สุกรํ ปุญฺญสฺส ปมาณํ คเณตุํ “เอตฺตโก ปุญฺญาภิสนฺโท กุสลาภิสนฺโท สุขสฺสาหาโร”ติ.\csromanspacer{} อถ โข อสงฺเขฺยยฺโย อปฺปเมยฺโย มหาปุญฺญกฺขนฺโธเตฺวว สงฺขฺยํ คจฺฉติ.}

\prose{เสยฺยถาปิ ภิกฺขเว มหาสมุทฺเท น สุกรํ อุทกสฺส ปมาณํ คเณตุํ “เอตฺตกานิ อุทกาฬฺหกานี”ติ วา “เอตฺตกานิ อุทกาฬฺหกสตานี”ติ วา “เอตฺตกานิ อุทกาฬฺหกสหสฺสานี”ติ วา “เอตฺตกานิ อุทกาฬฺหกสตสหสฺสานี”ติ วาติ, อถ โข อสงฺเขฺยยฺโย อปฺปเมยฺโย มหา-อุทกกฺขนฺโธเตฺวว สงฺขฺยํ คจฺฉติ.\csromanspacer{} เอวเมว โข ภิกฺขเว อิเมหิ จตูหิ ปุญฺญาภิสนฺเทหิ กุสลาภิสนฺเทหิ สมนฺนาคตสฺส อริยสาวกสฺส น สุกรํ ปุญฺญสฺส ปมาณํ คเณตุํ “เอตฺตโก ปุญฺญาภิสนฺโท กุสลาภิสนฺโท สุขสฺสาหาโร”ติ.\csromanspacer{} อถ โข อสงฺเขฺยยฺโย อปฺปเมยฺโย มหาปุญฺญกฺขนฺโธเตฺวว สงฺขฺยํ คจฺฉตีติ.}

\prose{อิทมโวจ ภควา, อิทํ วตฺวาน สุคโต อถาปรํ เอตทโวจ สตฺถา\csromandash{}}

\csromangathameasure{“มโหทธิํ อปริมิตํ มหาสรํ, \\ พหุเภรวํ รตนคณานมาลยํ. \\ นชฺโช ยถา นรคณสํฆเสวิตา, \\ ปุถู สวนฺตี อุปยนฺติ สาครํ. \\ เอวํ นรํ อนฺนปานวตฺถททํ, \\ เสยฺยานิ ปจฺจตฺถรณสฺส ทายกํ. \\ ปุญฺญสฺส ธารา อุปยนฺติ ปณฺฑิตํ, \\ นชฺโช ยถา วาริวหาว สาครนฺ”ติ. ปฐมํ.}
\csromangathagroup{\csromangathastanza{“มโหทธิํ อปริมิตํ มหาสรํ, \\ พหุเภรวํ รตนคณานมาลยํ. \\ นชฺโช ยถา นรคณสํฆเสวิตา, \\ ปุถู สวนฺตี อุปยนฺติ สาครํ.}\csromangathastanza{เอวํ นรํ อนฺนปานวตฺถททํ, \\ เสยฺยานิ ปจฺจตฺถรณสฺส\footnote{สชฺชตฺถรณสฺส (สี, สฺยา, กํ, อิ)} ทายกํ. \\ ปุญฺญสฺส ธารา อุปยนฺติ ปณฺฑิตํ, \\ นชฺโช ยถา วาริวหาว สาครนฺ”ติ.\csromanspacer{} ปฐมํ.}}

\csromanheader{2}{2. ทุติย-อภิสนฺทสุตฺต}
\proseitem{1038}{\csromanfolio{350}จตฺตาโรเม ภิกฺขเว ปุญฺญาภิสนฺทา กุสลาภิสนฺทา สุขสฺสาหารา.\csromanspacer{} กตเม จตฺตาโร, อิธ ภิกฺขเว อริยสาวโก พุทฺเธ อเวจฺจปฺปสาเทน สมนฺนาคโต โหติ “อิติปิ โส ภควา ฯเปฯ สตฺถา เทวมนุสฺสานํ พุทฺโธ ภควา”ติ, อยํ ปฐโม ปุญฺญาภิสนฺโท กุสลาภิสนฺโท สุขสฺสาหาโร.}

\prose{ปุน จปรํ ภิกฺขเว อริยสาวโก ธมฺเม ฯเปฯ.\csromanspacer{} สํเฆ ฯเปฯ.}

\prose{ปุน จปรํ ภิกฺขเว อริยสาวโก วิคตมลมจฺเฉเรน เจตสา อคารํ อชฺฌาวสติ มุตฺตจาโค ปยตปาณิ โวสฺสคฺครโต ยาจโยโค ทานสํวิภาครโต, อยํ จตุตฺโถ ปุญฺญาภิสนฺโท กุสลาภิสนฺโท สุขสฺสาหาโร.\csromanspacer{} อิเม โข ภิกฺขเว จตฺตาโร ปุญฺญาภิสนฺทา กุสลาภิสนฺทา สุขสฺสาหารา.}

\prose{อิเมหิ โข ภิกฺขเว จตูหิ ปุญฺญาภิสนฺเทหิ กุสลาภิสนฺเทหิ สมนฺนาคตสฺส อริยสาวกสฺส น สุกรํ ปุญฺญสฺส ปมาณํ คเณตุํ “เอตฺตโก ปุญฺญาภิสนฺโท กุสลาภิสนฺโท สุขสฺสาหาโร”ติ.\csromanspacer{} อถ โข อาสงฺเขฺยยฺโย อปฺปเมยฺโย มหาปุญฺญกฺขนฺโธเตฺวว สงฺขฺยํ คจฺฉติ.}

\prose{เสยฺยถาปิ ภิกฺขเว ยตฺถิมา มหานทิโย สํสนฺทนฺติ สเมนฺติ.\csromanspacer{} เสยฺยถิทํ, คงฺคา ยมุนา อจิรวตี สรภู มหี, ตตฺถ น สุกรํ อุทกสฺส ปมาณํ คเณตุํ “เอตฺตกานิ อุทกาฬฺหกานี”ติ วา “เอตฺตกานิ อุทกาฬฺหกสตานี”ติ วา “เอตฺตกานิ อุทกาฬฺหกสหสฺสานี”ติ วา “เอตฺตกานิ อุทกาฬฺหกสตสหสฺสานี”ติ วาติ.\csromanspacer{} อถ โข อสงฺเขฺยยฺโย อปฺปเมยฺโย มหา-อุทกกฺขนฺโธเตฺวว สงฺขฺยํ คจฺฉติ.\csromanspacer{} เอวเมว โข ภิกฺขเว อิเมหิ จตูหิ ปุญฺญาภิสนฺเทหิ กุสลาภิสนฺเทหิ สมนฺนาคตสฺส อริยสาวกสฺส น สุกรํ ปุญฺญสฺส ปมาณํ คเณตุํ “เอตฺตโก ปุญฺญาภิสนฺโท กุสลาภิสนฺโท สุขสฺสาหาโร”ติ.\csromanspacer{} อถ โข อสงฺเขฺยยฺโย อปฺปเมยฺโย มหาปุญฺญกฺขนฺโธเตฺวว สงฺขฺยํ คจฺฉตีติ.\csromanspacer{} อิทมโวจ ภควา ฯเปฯ สตฺถา\csromandash{}}

\csromangathameasure{“มโหทธิํ อปริมิตํ มหาสรํ, \\ พหุเภรวํ รตนคณานมาลยํ. \\ นชฺโช ยถา นรคณสํฆเสวิตา, \\ ปุถู สวนฺตี อุปยนฺติ สาครํ.}
\csromangathagroup{\csromangathastanza{“มโหทธิํ อปริมิตํ มหาสรํ, \\ พหุเภรวํ รตนคณานมาลยํ. \\ นชฺโช ยถา นรคณสํฆเสวิตา, \\ ปุถู สวนฺตี อุปยนฺติ สาครํ.}}

\csromanheader{2}{11. โสตาปตฺติสํยุตฺต}
\csromangathameasure{เอวํ นรํ อนฺนปานวตฺถททํ, \\ เสยฺยานิ ปจฺจตฺถรณสฺส ทายกํ. \\ ปุญฺญสฺส ธารา อุปยนฺติ ปณฺฑิตํ, \\ นชฺโช ยถา วาริวหาว สาครนฺ”ติ. ทุติยํ.}
\csromangathagroup{\csromangathastanza{\csromanfolio{351}เอวํ นรํ อนฺนปานวตฺถททํ, \\ เสยฺยานิ ปจฺจตฺถรณสฺส ทายกํ. \\ ปุญฺญสฺส ธารา อุปยนฺติ ปณฺฑิตํ, \\ นชฺโช ยถา วาริวหาว สาครนฺ”ติ.\csromanspacer{} ทุติยํ.}}

\csromanheader{2}{3. ตติย-อภิสนฺทสุตฺต}
\proseitem{1039}{จตฺตาโรเม ภิกฺขเว ปุญฺญาภิสนฺทา กุสลาภิสนฺทา สุขสฺสาหารา.\csromanspacer{} กตเม จตฺตาโร, อิธ ภิกฺขเว อริยสาวโก พุทฺเธ อเวจฺจปฺปสาเทน สมนฺนาคโต โหติ “อิติปิ โส ภควา ฯเปฯ สตฺถา เทวมนุสฺสานํ พุทฺโธ ภควา”ติ, อยํ ปฐโม ปุญฺญาภิสนฺโท กุสลาภิสนฺโท สุขสฺสาหาโร.}

\prose{ปุน จปรํ ภิกฺขเว อริยสาวโก ธมฺเม ฯเปฯ.\csromanspacer{} สํเฆ ฯเปฯ.}

\prose{ปุน จปรํ ภิกฺขเว อริยสาวโก ปญฺญวา โหติ, อุทยตฺถคามินิยา ปญฺญาย สมนฺนาคโต อริยาย นิพฺเพธิกาย สมฺมา ทุกฺขกฺขยคามินิยา, อยํ จตุตฺโถ ปุญฺญาภิสนฺโท กุสลาภิสนฺโท สุขสฺสาหาโร.\csromanspacer{} อิเม โข ภิกฺขเว จตฺตาโร ปุญฺญาภิสนฺทา กุสลาภิสนฺทา สุขสฺสาหารา.}

\prose{อิเมหิ โข ภิกฺขเว จตูหิ ปุญฺญาภิสนฺเทหิ กุสลาภิสนฺเทหิ สมนฺนาคตสฺส อริยสาวกสฺส น สุกรํ ปุญฺญสฺส ปมาณํ คเณตุํ “เอตฺตโก ปุญฺญาภิสนฺโท กุสลาภิสนฺโท สุขสฺสาหาโร”ติ.\csromanspacer{} อถ โข อสงฺเขฺยยฺโย อปฺปเมยฺโย มหาปุญฺญกฺขนฺโธเตฺวว สงฺขฺยํ คจฺฉตีติ.\csromanspacer{} อิทมโวจ ภควา ฯเปฯ สตฺถา\csromandash{}}

\csromangathameasure{“โย ปุญฺญกาโม กุสเล ปติฏฺฐิโต, \\ ภาเวติ มคฺคํ อมตสฺส ปตฺติยา. \\ โส ธมฺมสาราธิคโม ขเย รโต, \\ น เวธติ มจฺจุราชาคมนสฺมินฺ”ติ. ตติยํ.}
\csromangathagroup{\csromangathastanza{“โย ปุญฺญกาโม กุสเล ปติฏฺฐิโต, \\ ภาเวติ มคฺคํ อมตสฺส ปตฺติยา. \\ โส ธมฺมสาราธิคโม ขเย รโต, \\ น เวธติ มจฺจุราชาคมนสฺมินฺ”ติ\footnote{มจฺจุราชาคมิสฺสตีติ (สี, อิ), มจฺจุชรากมฺปิสฺมินฺติ (ก)}.\csromanspacer{} ตติยํ.}}

\csromanheader{2}{4. ปฐมมหทฺธนสุตฺต}
\csromanmatikaanchor{mtk.252}
\csromantocmarkref{subparagraph}{4-10. ปฐมมหทฺธนสุตฺตาทีนิ}{mtk.252}
\bookmark[dest={mtk.252},level=5]{4-10. ปฐมมหทฺธนสุตฺตาทีนิ}
\proseitem{1040}{\csromanfolio{352}จตูหิ ภิกฺขเว ธมฺเมหิ สมนฺนาคโต อริยสาวโก อฑฺโฒ มหทฺธโน มหาโภโคติ\footnote{มหาโภโค มหายโสติ (สฺยา, อิ, ก)} วุจฺจติ.}

\prose{กตเมหิ จตูหิ, อิธ ภิกฺขเว อริยสาวโก พุทฺเธ อเวจฺจปฺปสาเทน สมนฺนาคโต โหติ “อิติปิ โส ภควา ฯเปฯ สตฺถา เทวมนุสฺสานํ พุทฺโธ ภควา”ติ.\csromanspacer{} ธมฺเม ฯเปฯ.\csromanspacer{} สํเฆ ฯเปฯ.\csromanspacer{} อริยกนฺเตหิ สีเลหิ สมนฺนาคโต โหติ อขณฺเฑหิ ฯเปฯ สมาธิสํวตฺตนิเกหิ.\csromanspacer{} อิเมหิ โข ภิกฺขเว จตูหิ ธมฺเมหิ สมนฺนาคโต อริยสาวโก อฑฺโฒ มหทฺธโน มหาโภโคติ วุจฺจตีติ.\csromanspacer{}\csromanspacer{}จตุตฺถํ.}

\csromanheader{2}{5. ทุติยมหทฺธนสุตฺต}
\proseitem{1041}{จตูหิ ภิกฺขเว ธมฺเมหิ สมนฺนาคโต อริยสาวโก อฑฺโฒ มหทฺธโน มหาโภโค มหายโสติ วุจฺจติ.}

\prose{กตเมหิ จตูหิ, อิธ ภิกฺขเว อริยสาวโก พุทฺเธ อเวจฺจปฺปสาเทน สมนฺนาคโต โหติ “อิติปิ โส ภควา ฯเปฯ สตฺถา เทวมนุสฺสานํ พุทฺโธ ภควา”ติ.\csromanspacer{} ธมฺเม ฯเปฯ.\csromanspacer{} สํเฆ ฯเปฯ.\csromanspacer{} อริยกนฺเตหิ สีเลหิ สมนฺนาคโต โหติ อขณฺเฑหิ ฯเปฯ.\csromanspacer{} สมาธิสํวตฺตนิเกหิ.\csromanspacer{} อิเมหิ โข ภิกฺขเว จตูหิ ธมฺเมหิ สมนฺนาคโต อริยสาวโก อฑฺโฒ มหทฺธโน มหาโภโค มหายโสติ วุจฺจตีติ.\csromanspacer{}\csromanspacer{}ปญฺจมํ.}

\csromanheader{2}{6. สุทฺธกสุตฺต}
\proseitem{1042}{จตูหิ ภิกฺขเว ธมฺเมหิ สมนฺนาคโต อริยสาวโก โสตาปนฺโน โหติ อวินิปาตธมฺโม นิยโต สมฺโพธิปรายโณ.}

\prose{กตเมหิ จตูหิ, อิธ ภิกฺขเว อริยสาวโก พุทฺเธ อเวจฺจปฺปสาเทน สมนฺนาคโต โหติ “อิติปิ โส ภควา ฯเปฯ สตฺถา เทวมนุสฺสานํ พุทฺโธ ภควา”ติ.\csromanspacer{} ธมฺเม ฯเปฯ.\csromanspacer{} สํเฆ ฯเปฯ.\csromanspacer{} อริยกนฺเตหิ สีเลหิ สมนฺนาคโต โหติ อขณฺเฑหิ ฯเปฯ สมาธิสํวตฺตนิเกหิ.\csromanspacer{} อิเมหิ โข ภิกฺขเว จตูหิ ธมฺเมหิ สมนฺนาคโต อริยสาวโก โสตาปนฺโน โหติ อวินิปาตธมฺโม นิยโต สมฺโพธิปรายโณติ.\csromanspacer{}\csromanspacer{}ฉฏฺฐํ.}

\csromanheader{2}{11. โสตาปตฺติสํยุตฺต}
\csromanheader{2}{7. นนฺทิยสุตฺต}
\proseitem{1043}{\csromanfolio{353}กปิลวตฺถุนิทานํ.\csromanspacer{} เอกมนฺตํ นิสินฺนํ โข นนฺทิยํ สกฺกํ ภควา เอตทโวจ “จตูหิ โข นนฺทิย ธมฺเมหิ สมนฺนาคโต อริยสาวโก โสตาปนฺโน โหติ อวินิปาตธมฺโม นิยโต สมฺโพธิปรายโณ.}

\prose{กตเมหิ จตูหิ, อิธ นนฺทิย อริยสาวโก พุทฺเธ อเวจฺจปฺปสาเทน สมนฺนาคโต โหติ “อิติปิ โส ภควา ฯเปฯ สตฺถา เทวมนุสฺสานํ พุทฺโธ ภควา”ติ.\csromanspacer{} ธมฺเม ฯเปฯ.\csromanspacer{} สํเฆ ฯเปฯ.\csromanspacer{} อริยกนฺเตหิ สีเลหิ สมนฺนาคโต โหติ อขณฺเฑหิ ฯเปฯ สมาธิสํวตฺตนิเกหิ.\csromanspacer{} อิเมหิ โข นนฺทิย จตูหิ ธมฺเมหิ สมนฺนาคโต อริยสาวโก โสตาปนฺโน โหติ อวินิปาตธมฺโม นิยโต สมฺโพธิปรายโณ”ติ.\csromanspacer{}\csromanspacer{}สตฺตมํ.}

\csromanheader{2}{8. ภทฺทิยสุตฺต}
\proseitem{1044}{กปิลวตฺถุนิทานํ.\csromanspacer{} เอกมนฺตํ นิสินฺนํ โข ภทฺทิยํ สกฺกํ ภควา เอตทโวจ “จตูหิ โข ภทฺทิย ธมฺเมหิ สมนฺนาคโต อริยสาวโก โสตาปนฺโน โหติ อวินิปาตธมฺโม นิยโต สมฺโพธิปรายโณ.}

\prose{กตเมหิ จตูหิ, อิธ ภทฺทิย อริยสาวโก พุทฺเธ ฯเปฯ.\csromanspacer{} ธมฺเม ฯเปฯ.\csromanspacer{} สํเฆ ฯเปฯ.\csromanspacer{} อริยกนฺเตหิ สีเลหิ สมนฺนาคโต โหติ อขณฺเฑหิ ฯเปฯ สมาธิสํวตฺตนิเกหิ.\csromanspacer{} อิเมหิ โข ภทฺทิย จตูหิ ธมฺเมหิ สมนฺนาคโต อริยสาวโก โสตาปนฺโน โหติ, อวินิปาตธมฺโม นิยโต สมฺโพธิปรายโณ”ติ.\csromanspacer{}\csromanspacer{}อฏฺฐมํ.}

\csromanheader{2}{9. มหานามสุตฺต}
\proseitem{1045}{กปิลวตฺถุนิทานํ.\csromanspacer{} เอกมนฺตํ นิสินฺนํ โข มหานามํ สกฺกํ ภควา เอตทโวจ “จตูหิ โข มหานาม ธมฺเมหิ สมนฺนาคโต อริยสาวโก โสตาปนฺโน โหติ ฯเปฯ สมฺโพธิปรายโณ.}

\prose{กตเมหิ จตูหิ, อิธ มหานาม อริยสาวโก พุทฺเธ ฯเปฯ.\csromanspacer{} ธมฺเม ฯเปฯ.\csromanspacer{} สํเฆ ฯเปฯ.\csromanspacer{} อริยกนฺเตหิ สีเลหิ สมนฺนาคโต โหติ อขณฺเฑหิ ฯเปฯ สมาธิสํวตฺตนิเกหิ.\csromanspacer{} อิเมหิ โข มหานาม จตูหิ ธมฺเมหิ สมนฺนาคโต อริยสาวโก โสตาปนฺโน โหติ, อวินิปาตธมฺโม นิยโต สมฺโพธิปรายโณ”ติ.\csromanspacer{}\csromanspacer{}นวมํ.}

\csromanheader{2}{10. องฺคสุตฺต}
\proseitem{1046}{\csromanfolio{354}จตฺตาริมานิ ภิกฺขเว โสตาปตฺติยงฺคานิ.\csromanspacer{} กตมานิ จตฺตาริ, สปฺปุริสสํเสโว สทฺธมฺมสฺสวนํ โยนิโสมนสิกาโร ธมฺมานุธมฺมปฺปฏิปตฺติ.\csromanspacer{} อิมานิ โข ภิกฺขเว จตฺตาริ โสตาปตฺติยงฺคานีติ.\csromanspacer{}\csromanspacer{}ทสมํ.\csromanspacer{} สคาถกปุญฺญาภิสนฺทวคฺโค ปญฺจโม.}

\titlehead{\textbf{ตสฺสุทฺทานํ}}
\csromangathameasure{อภิสนฺทา ตโย วุตฺตา, ทุเว มหทฺธเนน จ. \\ สุทฺธํ นนฺทิยํ ภทฺทิยํ, มหานามงฺเคน เต ทสาติ.}
\csromangathagroup{\csromangathastanza{อภิสนฺทา ตโย วุตฺตา, ทุเว มหทฺธเนน จ. \\ สุทฺธํ นนฺทิยํ ภทฺทิยํ, มหานามงฺเคน เต ทสาติ.}}

\csromanmatikaanchor{mtk.253}
\csromantocmarknopage{subparagraph}{6. สปฺปญฺญวคฺค}{mtk.253}
\bookmark[dest={mtk.253},level=5]{6. สปฺปญฺญวคฺค}
\csromanheader{6}{6. สปฺปญฺญวคฺค}
\csromanmatikaanchor{mtk.254}
\csromantocmarkref{subparagraph}{1. สคาถกสุตฺต}{mtk.254}
\bookmark[dest={mtk.254},level=5]{1. สคาถกสุตฺต}
\csromanheader{6}{1. สคาถกสุตฺต}
\proseitem{1047}{จตูหิ ภิกฺขเว ธมฺเมหิ สมนฺนาคโต อริยสาวโก โสตาปนฺโน โหติ อวินิปาตธมฺโม นิยโต สมฺโพธิปรายโณ.}

\prose{กตเมหิ จตูหิ, อิธ ภิกฺขเว อริยสาวโก พุทฺเธ อเวจฺจปฺปสาเทน สมนฺนาคโต โหติ “อิติปิ โส ภควา ฯเปฯ สตฺถา เทวมนุสฺสานํ พุทฺโธ ภควา”ติ.\csromanspacer{} ธมฺเม ฯเปฯ.\csromanspacer{} สํเฆ ฯเปฯ.\csromanspacer{} อริยกนฺเตหิ สีเลหิ สมนฺนาคโต โหติ อขณฺเฑหิ ฯเปฯ สมาธิสํวตฺตนิเกหิ.\csromanspacer{} อิเมหิ โข ภิกฺขเว จตูหิ ธมฺเมหิ สมนฺนาคโต อริยสาวโก โสตาปนฺโน โหติ อวินิปาตธมฺโม นิยโต สมฺโพธิปรายโณติ.\csromanspacer{} อิทมโวจ ภควา, อิทํ วตฺวาน สุคโต อถาปรํ เอตทโวจ สตฺถา\csromandash{}}

\csromangathameasure{“ยสฺส สทฺธา ตถาคเต, อจลา สุปฺปติฏฺฐิตา. \\ สีลญฺจ ยสฺส กลฺยาณํ, อริยกนฺตํ ปสํสิตํ. \\ สํเฆ ปสาโท ยสฺสตฺถิ, อุชุภูตญฺจ ทสฺสนํ. \\ อทลิทฺโทติ ตํ อาหุ, อโมฆํ ตสฺส ชีวิตํ. \\ ตสฺมา สทฺธญฺจ สีลญฺจ, ปสาทํ ธมฺมทสฺสนํ. \\ อนุยุญฺเชถ เมธาวี, สรํ พุทฺธานสาสนนฺ”ติ. ปฐมํ.}
\csromangathagroup{\csromangathastanza{“ยสฺส สทฺธา ตถาคเต, อจลา สุปฺปติฏฺฐิตา. \\ สีลญฺจ ยสฺส กลฺยาณํ, อริยกนฺตํ ปสํสิตํ.}\csromangathastanza{สํเฆ ปสาโท ยสฺสตฺถิ, อุชุภูตญฺจ ทสฺสนํ. \\ อทลิทฺโทติ ตํ อาหุ, อโมฆํ ตสฺส ชีวิตํ.}\csromangathastanza{ตสฺมา สทฺธญฺจ สีลญฺจ, ปสาทํ ธมฺมทสฺสนํ. \\ อนุยุญฺเชถ เมธาวี, สรํ พุทฺธานสาสนนฺ”ติ.\csromanspacer{} ปฐมํ.}}

\csromanheader{2}{11. โสตาปตฺติสํยุตฺต}
\csromanmatikaanchor{mtk.255}
\csromantocmarkref{subparagraph}{2. วสฺสํวุตฺถสุตฺต}{mtk.255}
\bookmark[dest={mtk.255},level=5]{2. วสฺสํวุตฺถสุตฺต}
\csromanheader{6}{2. วสฺสํวุตฺถสุตฺต}
\proseitem{1048}{\csromanfolio{355}เอกํ สมยํ ภควา สาวตฺถิยํ วิหรติ เชตวเน อนาถปิณฺฑิกสฺส อาราเม.\csromanspacer{} เตน โข ปน สมเยน อญฺญตโร ภิกฺขุ สาวตฺถิยํ วสฺสํวุตฺโถ กปิลวตฺถุํ อนุปฺปตฺโต โหติ เกนจิเทว กรณีเยน.\csromanspacer{} อสฺโสสุํ โข กาปิลวตฺถวา สกฺยา “อญฺญตโร กิร ภิกฺขุ สาวตฺถิยํ วสฺสํวุตฺโถ กปิลวตฺถุํ อนุปฺปตฺโต”ติ.}

\prose{อถ โข กาปิลวตฺถวา สกฺยา เยน โส ภิกฺขุ เตนุปสงฺกมิํสุ, อุปสงฺกมิตฺวา ตํ ภิกฺขุํ อภิวาเทตฺวา เอกมนฺตํ นิสีทิํสุ, เอกมนฺตํ นิสินฺนา โข กาปิลวตฺถวา สกฺยา ตํ ภิกฺขุํ เอตทโวจุํ “กจฺจิ ภนฺเต ภควา อโรโค เจว พลวา จาติ.\csromanspacer{} อโรโค จาวุโส ภควา พลวา จาติ.\csromanspacer{} กจฺจิ ปน ภนฺเต สาริปุตฺตโมคฺคลฺลานา อโรคา เจว พลวนฺโต จาติ.\csromanspacer{} สาริปุตฺตโมคฺคลฺลานาปิ โข อาวุโส อโรคา เจว พลวนฺโต จาติ.\csromanspacer{} กจฺจิ ปน ภนฺเต ภิกฺขุสํโฆ อโรโค จ พลวา จาติ.\csromanspacer{} ภิกฺขุสํโฆปิ โข อาวุโส อโรโค จ พลวา จาติ.\csromanspacer{} อตฺถิ ปน เต ภนฺเต กิญฺจิ อิมินา อนฺตรวสฺเสน ภควโต สมฺมุขา สุตํ, สมฺมุขา ปฏิคฺคหิตนฺติ.\csromanspacer{} สมฺมุขา เมตํ อาวุโส ภควโต สุตํ, สมฺมุขา ปฏิคฺคหิตํ “อปฺปกา เต ภิกฺขเว ภิกฺขู, เย อาสวานํ ขยา อนาสวํ เจโตวิมุตฺติํ ปญฺญาวิมุตฺติํ ทิฏฺเฐว ธมฺเม สยํ อภิญฺญา สจฺฉิกตฺวา อุปสมฺปชฺช วิหรนฺติ, อถ โข เอเตว พหุตรา ภิกฺขู, เย ปญฺจนฺนํ โอรมฺภาคิยานํ สํโยชนานํ ปริกฺขยา โอปปาติกา, ตตฺถ ปรินิพฺพายิโน อนาวตฺติธมฺมา ตสฺมา โลกา”ติ.}

\prose{อปรมฺปิ โข เม อาวุโส ภควโต สมฺมุขา สุตํ, สมฺมุขา ปฏิคฺคหิตํ “อปฺปกา เต ภิกฺขเว ภิกฺขู, เย ปญฺจนฺนํ โอรมฺภาคิยานํ สํโยชนานํ ปริกฺขยา โอปปาติกา, ตตฺถ ปรินิพฺพายิโน อนาวตฺติธมฺมา ตสฺมา โลกา, อถ โข เอเตว พหุตรา ภิกฺขู, เย ติณฺณํ สํโยชนานํ ปริกฺขยา ราคโทสโมหานํ ตนุตฺตา สกทาคามิโน สกิเทว อิมํ โลกํ อาคนฺตฺวา ทุกฺขสฺสนฺตํ กริสฺสนฺตี”ติ.}

\prose{อปรมฺปิ โข เม อาวุโส ภควโต สมฺมุขา สุตํ, สมฺมุขา ปฏิคฺคหิตํ “อปฺปกา เต ภิกฺขเว ภิกฺขู, เย ติณฺณํ สํโยชนานํ ปริกฺขยา ราคโทสโมหานํ ตนุตฺตา สกทาคามิโน สกิเทว อิมํ โลกํ \csromanfolio{356}อาคนฺตฺวา ทุกฺขสฺสนฺตํ กริสฺสนฺติ, อถ โข เอเตว พหุตรา ภิกฺขู, เย ติณฺณํ สํโยชนานํ ปริกฺขยา โสตาปนฺนา อวินิปาตธมฺมา นิยตา สมฺโพธิปรายณาติ.\csromanspacer{}\csromanspacer{}ทุติยํ.}

\csromanmatikaanchor{mtk.256}
\csromantocmarkref{subparagraph}{3. ธมฺมทินฺนสุตฺต}{mtk.256}
\bookmark[dest={mtk.256},level=5]{3. ธมฺมทินฺนสุตฺต}
\csromanheader{6}{3. ธมฺมทินฺนสุตฺต}
\proseitem{1049}{เอกํ สมยํ ภควา พาราณสิยํ วิหรติ อิสิปตเน มิคทาเย.\csromanspacer{} อถ โข ธมฺมทินฺโน อุปาสโก ปญฺจหิ อุปาสกสเตหิ สทฺธิํ เยน ภควา เตนุปสงฺกมิ, อุปสงฺกมิตฺวา ภควนฺตํ อภิวาเทตฺวา เอกมนฺตํ นิสีทิ, เอกมนฺตํ นิสินฺโน โข ธมฺมทินฺโน อุปาสโก ภควนฺตํ เอตทโวจ “โอวทตุ โน ภนฺเต ภควา, อนุสาสตุ โน ภนฺเต ภควา, ยํ อมฺหากํ อสฺส ทีฆรตฺตํ หิตาย สุขายา”ติ.}

\prose{ตสฺมาติห โว ธมฺมทินฺน เอวํ สิกฺขิตพฺพํ “เย เต สุตฺตนฺตา ตถาคตภาสิตา คมฺภีรา คมฺภีรตฺถา โลกุตฺตรา สุญฺญตปฏิสํยุตฺตา, เต กาเลน กาลํ อุปสมฺปชฺช วิหริสฺสามา”ติ, เอวํ หิ โว ธมฺมทินฺน สิกฺขิตพฺพนฺติ.\csromanspacer{} น โข เนตํ ภนฺเต สุกรํ อเมฺหหิ ปุตฺตสมฺพาธสยนํ อชฺฌาวสนฺเตหิ กาสิกจนฺทนํ ปจฺจนุโภนฺเตหิ มาลาคนฺธวิเลปนํ ธารยนฺเตหิ ชาตรูปรชตํ สาทิยนฺเตหิ, เย เต สุตฺตนฺตา ตถาคตภาสิตา คมฺภีรา คมฺภีรตฺถา โลกุตฺตรา สุญฺญตปฏิสํยุตฺตา, เต กาเลน กาลํ อุปสมฺปชฺช วิหริตุํ, เตสํ โน ภนฺเต ภควา อมฺหากํ ปญฺจสุ สิกฺขาปเทสุ ฐิตานํ อุตฺตริธมฺมํ เทเสตูติ.}

\prose{ตสฺมาติห โว ธมฺมทินฺน เอวํ สิกฺขิตพฺพํ “พุทฺเธ อเวจฺจปฺปสาเทน สมนฺนาคตา ภวิสฺสาม “อิติปิ โส ภควา ฯเปฯ สตฺถา เทวมนุสฺสานํ พุทฺโธ ภควา”ติ.\csromanspacer{} ธมฺเม ฯเปฯ.\csromanspacer{} สํเฆ ฯเปฯ.\csromanspacer{} อริยกนฺเตหิ สีเลหิ สมนฺนาคตา ภวิสฺสาม อขณฺเฑหิ ฯเปฯ สมาธิสํวตฺตนิเกหีติ.\csromanspacer{} เอวํ หิ โว ธมฺมทินฺน สิกฺขิตพฺพนฺติ.}

\prose{ยานิมานิ ภนฺเต ภควตา จตฺตาริ โสตาปตฺติยงฺคานิ เทสิตานิ, สํวิชฺชนฺเตเต ธมฺมา อเมฺหสุ, มยญฺจ เตสุ ธมฺเมสุ สนฺทิสฺสาม.\csromanspacer{} มยํ หิ ภนฺเต พุทฺเธ อเวจฺจปฺปสาเทน สมนฺนาคตา “อิติปิ โส ภควา ฯเปฯ สตฺถา เทวมนุสฺสานํ พุทฺโธ ภควา”ติ.\csromanspacer{} ธมฺเม ฯเปฯ.\csromanspacer{} สํเฆ ฯเปฯ.\csromanspacer{} อริยกนฺเตหิ สีเลหิ สมนฺนาคตา อขณฺเฑหิ ฯเปฯ สมาธิสํวตฺตนิเกหีติ.\csromanspacer{} ลาภา}

\csromanheader{2}{11. โสตาปตฺติสํยุตฺต}
\prose{\csromanfolio{357}โว ธมฺมทินฺน, สุลทฺธํ โว ธมฺมทินฺน, โสตาปตฺติผลํ ตุเมฺหหิ พฺยากตนฺติ.\csromanspacer{}\csromanspacer{}ตติยํ.}

\csromanmatikaanchor{mtk.257}
\csromantocmarkref{subparagraph}{4. คิลานสุตฺต}{mtk.257}
\bookmark[dest={mtk.257},level=5]{4. คิลานสุตฺต}
\csromanheader{6}{4. คิลานสุตฺต}
\proseitem{1050}{เอกํ สมยํ ภควา สกฺเกสุ วิหรติ กปิลวตฺถุสฺมิํ นิโคฺรธาราเม.\csromanspacer{} เตน โข ปน สมเยน สมฺพหุลา ภิกฺขู ภควโต จีวรกมฺมํ กโรนฺติ “นิฏฺฐิตจีวโร ภควา เตมาสจฺจเยน จาริกํ ปกฺกมิสฺสตี”ติ.\csromanspacer{} อสฺโสสิ โข มหานาโม สกฺโก “สมฺพหุลา กิร ภิกฺขู ภควโต จีวรกมฺมํ กโรนฺติ ‘นิฏฺฐิตจีวโร ภควา เตมาสจฺจเยน จาริกํ ปกฺกมิสฺสตี’ติ”.\csromanspacer{} อถ โข มหานาโม สกฺโก เยน ภควา เตนุปสงฺกมิ, อุปสงฺกมิตฺวา ภควนฺตํ อภิวาเทตฺวา เอกมนฺตํ นิสีทิ, เอกมนฺตํ นิสินฺโน โข มหานาโม สกฺโก ภควนฺตํ เอตทโวจ “สุตเมตํ ภนฺเต สมฺพหุลา กิร ภิกฺขู ภควโต จีวรกมฺมํ กโรนฺติ ‘นิฏฺฐิตจีวโร ภควา เตมาสจฺจเยน จาริกํ ปกฺกมิสฺสตี’ติ.\csromanspacer{} น โข เนตํ\footnote{น โข เต เอตํ (สี, อิ)} ภนฺเต ภควโต สมฺมุขา สุตํ, สมฺมุขา ปฏิคฺคหิตํ ‘สปฺปญฺเญน อุปาสเกน สปฺปญฺโญ อุปาสโก อาพาธิโก ทุกฺขิโต พาฬฺหคิลาโน โอวทิตพฺโพ’ติ”.}

\prose{สปฺปญฺเญน มหานาม อุปาสเกน สปฺปญฺโญ อุปาสโก อาพาธิโก ทุกฺขิโต พาฬฺหคิลาโน จตูหิ อสฺสาสนีเยหิ ธมฺเมหิ อสฺสาเสตพฺโพ “อสฺสาสตายสฺมา, อตฺถายสฺมโต พุทฺเธ อเวจฺจปฺปสาโท ‘อิติปิ โส ภควา ฯเปฯ สตฺถา เทวมนุสฺสานํ พุทฺโธ ภควา’ติ.\csromanspacer{} อสฺสาสตายสฺมา, อตฺถายสฺมโต ธมฺเม ฯเปฯ.\csromanspacer{} สํเฆ ฯเปฯ.\csromanspacer{} อริยกนฺตานิ สีลานิ อขณฺฑานิ ฯเปฯ สมาธิสํวตฺตนิกานี”ติ.}

\prose{สปฺปญฺเญน มหานาม อุปาสเกน สปฺปญฺโญ อุปาสโก อาพาธิโก ทุกฺขิโต พาฬฺหคิลาโน อิเมหิ จตูหิ อสฺสาสนีเยหิ ธมฺเมหิ อสฺสาเสตฺวา เอวมสฺส วจนีโย “อตฺถายสฺมโต มาตาปิตูสุ อเปกฺขา”ติ.\csromanspacer{} โส เจ เอวํ วเทยฺย “อตฺถิ เม มาตาปิตูสุ อเปกฺขา”ติ.\csromanspacer{} โส เอวมสฺส วจนีโย “อายสฺมา โข มาริโส มรณธมฺโม, สเจปายสฺมา มาตาปิตูสุ อเปกฺขํ กริสฺสติ มริสฺสเตว, โน \csromanfolio{358}เจปายสฺมา มาตาปิตูสุ อเปกฺขํ กริสฺสติ มริสฺสเตว.\csromanspacer{} สาธายสฺมา ยา เต มาตาปิตูสุ อเปกฺขา, ตํ ปชหา”ติ.}

\prose{โส เจ เอวํ วเทยฺย “ยา เม มาตาปิตูสุ อเปกฺขา, สา ปหีนา”ติ, โส เอวมสฺส วจนีโย “อตฺถิ ปนายสฺมโต ปุตฺตทาเรสุ อเปกฺขา”ติ.\csromanspacer{} โส เจ เอวํ วเทยฺย “อตฺถิ เม ปุตฺตทาเรสุ อเปกฺขา”ติ.\csromanspacer{} โส เอวมสฺส วจนีโย “อายสฺมา โข มาริโส มรณธมฺโม, สเจปายสฺมา ปุตฺตทาเรสุ อเปกฺขํ กริสฺสติ มริสฺสเตว, โน เจปายสฺมา ปุตฺตทาเรสุ อเปกฺขํ กริสฺสติ มริสฺสเตว.\csromanspacer{} สาธายสฺมา ยา เต ปุตฺตทาเรสุ อเปกฺขา, ตํ ปชหา”ติ.}

\prose{โส เจ เอวํ วเทยฺย “ยา เม ปุตฺตทาเรสุ อเปกฺขา, สา ปหีนาติ, โส เอวมสฺส วจนีโย “อตฺถิ ปนายสฺมโต มานุสเกสุ ปญฺจสุ กามคุเณสุ อเปกฺขา”ติ, โส เจ เอวํ วเทยฺย “อตฺถิ เม มานุสเกสุ ปญฺจสุ กามคุเณสุ อเปกฺขา”ติ, โส เอวมสฺส วจนีโย “มานุสเกหิ โข อาวุโส กาเมหิ ทิพฺพา กามา อภิกฺกนฺตตรา จ ปณีตตรา จ.\csromanspacer{} สาธายสฺมา มานุสเกหิ กาเมหิ จิตฺตํ วุฏฺฐาเปตฺวา จาตุมหาราชิเกสุ\footnote{จาตุมฺมหาราชิเกสุ (สี, สฺยา, กํ, อิ)} เทเวสุ จิตฺตํ อธิโมเจหี”ติ.}

\prose{โส เจ เอวํ วเทยฺย “มานุสเกหิ เม กาเมหิ จิตฺตํ วุฏฺฐิตํ, จาตุมหาราชิเกสุ เทเวสุ จิตฺตํ อธิโมจิตนฺ”ติ, โส เอวมสฺส วจนีโย “จาตุมหาราชิเกหิ โข อาวุโส เทเวหิ ตาวติํสา เทวา อภิกฺกนฺตตรา จ ปณีตตรา จ.\csromanspacer{} สาธายสฺมา จาตุมหาราชิเกหิ เทเวหิ จิตฺตํ วุฏฺฐาเปตฺวา ตาวติํเสสุ เทเวสุ จิตฺตํ อธิโมเจหี”ติ.}

\prose{โส เจ เอวํ วเทยฺย “จาตุมหาราชิเกหิ เม เทเวหิ จิตฺตํ วุฏฺฐิตํ, ตาวติํเสสุ เทเวสุ จิตฺตํ อธิโมจิตนฺ”ติ, โส เอวมสฺส วจนีโย “ตาวติํเสหิ โข อาวุโส เทเวหิ ยามา เทวา ฯเปฯ ตุสิตา เทวา ฯเปฯ นิมฺมานรตี เทวา ฯเปฯ ปรนิมฺมิตวสวตฺตี เทวา ฯเปฯ ปรนิมฺมิตวสวตฺตีหิ โข อาวุโส เทเวหิ พฺรหฺมโลโก อภิกฺกนฺตตโร จ ปณีตตโร จ.\csromanspacer{} สาธายสฺมา ปรนิมฺมิตวสวตฺตีหิ เทเวหิ จิตฺตํ วุฏฺฐาเปตฺวา พฺรหฺมโลเก จิตฺตํ อธิโมเจหี”ติ.}

\csromanheader{2}{11. โสตาปตฺติสํยุตฺต}
\csromanmatikaanchor{mtk.258}
\csromantocmarkref{subparagraph}{5-11. โสตาปตฺติผลสุตฺตาทีนิ}{mtk.258}
\bookmark[dest={mtk.258},level=5]{5-11. โสตาปตฺติผลสุตฺตาทีนิ}
\prose{\csromanfolio{359}โส เจ เอวํ วเทยฺย “ปรนิมฺมิตวสวตฺตีหิ เม เทเวหิ จิตฺตํ วุฏฺฐิตํ, พฺรหฺมโลเก จิตฺตํ อธิโมจิตนฺ”ติ, โส เอวมสฺส วจนีโย “พฺรหฺมโลโกปิ โข อาวุโส อนิจฺโจ อทฺธุโว สกฺกายปริยาปนฺโน.\csromanspacer{} สาธายสฺมา พฺรหฺมโลกา จิตฺตํ วุฏฺฐาเปตฺวา สกฺกายนิโรเธ จิตฺตํ อุปสํหราหี”ติ.}

\prose{โส เจ เอวํ วเทยฺย “พฺรหฺมโลกา เม จิตฺตํ วุฏฺฐิตํ, สกฺกายนิโรเธ จิตฺตํ อุปสํหรามี”ติ.\csromanspacer{} เอวํ วิมุตฺตจิตฺตสฺส โข มหานาม อุปาสกสฺส อาสวา\footnote{วสฺสสต (สี, สฺยา)} วิมุตฺตจิตฺเตน ภิกฺขุนา น กิญฺจิ นานากรณํ วทามิ, ยทิทํ วิมุตฺติยา วิมุตฺตนฺติ.\csromanspacer{}\csromanspacer{}จตุตฺถํ.}

\csromanheader{2}{5. โสตาปตฺติผลสุตฺต}
\proseitem{1051}{จตฺตาโรเม ภิกฺขเว ธมฺมา ภาวิตา พหุลีกตา โสตาปตฺติผลสจฺฉิกิริยาย สํวตฺตนฺติ.\csromanspacer{} กตเม จตฺตาโร, สปฺปุริสสํเสโว สทฺธมฺมสฺสวนํ โยนิโสมนสิกาโร ธมฺมานุธมฺมปฺปฏิปตฺติ.\csromanspacer{} อิเม โข ภิกฺขเว จตฺตาโร ธมฺมา ภาวิตา พหุลีกตา โสตาปตฺติผลสจฺฉิกิริยาย สํวตฺตนฺตีติ.\csromanspacer{}\csromanspacer{}ปญฺจมํ.}

\csromanheader{2}{6. สกทาคามิผลสุตฺต}
\proseitem{1052}{จตฺตาโรเม ภิกฺขเว ธมฺมา ภาวิตา พหุลีกตา สกทาคามิผลสจฺฉิกิริยาย สํวตฺตนฺติ.\csromanspacer{} กตเม จตฺตาโร ฯเปฯ สํวตฺตนฺตีติ.\csromanspacer{}\csromanspacer{}ฉฏฺฐํ.}

\csromanheader{2}{7. อนาคามิผลสุตฺต}
\vspace{\tipitakaparskip}
\csromancenter{อนาคามิผลสจฺฉิกิริยาย สํวตฺตนฺตีติ.\csromanspacer{}\csromanspacer{}สตฺตมํ.}
\csromanheader{2}{8. อรหตฺตผลสุตฺต}
\vspace{\tipitakaparskip}
\csromancenter{อรหตฺตผลสจฺฉิกิริยาย สํวตฺตนฺตีติ.\csromanspacer{}\csromanspacer{}อฏฺฐมํ.}
\csromanheader{2}{9. ปญฺญาปฏิลาภสุตฺต}
\vspace{\tipitakaparskip}
\csromancenter{ปญฺญาปฏิลาภาย สํวตฺตนฺตีติ.\csromanspacer{}\csromanspacer{}นวมํ.}
\csromanheader{2}{10. ปญฺญาวุทฺธิสุตฺต}
\vspace{\tipitakaparskip}
\csromancenter{ปญฺญาวุทฺธิยา สํวตฺตนฺตีติ.\csromanspacer{}\csromanspacer{}ทสมํ.}
\csromanheader{2}{11. ปญฺญาเวปุลฺลสุตฺต}
\csromanmatikaanchor{mtk.259}
\csromanmatikaanchor{mtk.260}
\csromantocmarknopage{subparagraph}{7. มหาปญฺญวคฺค}{mtk.259}
\bookmark[dest={mtk.259},level=5]{7. มหาปญฺญวคฺค}
\csromantocmarkref{subparagraph}{1-13. มหาปญฺญาสุตฺตาทีนิ}{mtk.260}
\bookmark[dest={mtk.260},level=5]{1-13. มหาปญฺญาสุตฺตาทีนิ}
\proseitem{1057}{\csromanfolio{360}ปญฺญาเวปุลฺลาย สํวตฺตนฺตีติ.\csromanspacer{}\csromanspacer{}เอกาทสมํ.\csromanspacer{} สปฺปญฺญวคฺโค ฉฏฺโฐ.}

\titlehead{\textbf{ตสฺสุทฺทานํ}}
\csromangathameasure{สคาถกํ วสฺสํวุตฺถํ, ธมฺมทินฺนญฺจ คิลานํ. \\ จตุโร ผลา ปฏิลาโภ, วุทฺธิ เวปุลฺลตาย จาติ.}
\csromangathagroup{\csromangathastanza{สคาถกํ วสฺสํวุตฺถํ, ธมฺมทินฺนญฺจ คิลานํ. \\ จตุโร ผลา ปฏิลาโภ, วุทฺธิ เวปุลฺลตาย จาติ.}}

\csromanheader{6}{7. มหาปญฺญวคฺค}
\csromanheader{2}{1. มหาปญฺญาสุตฺต}
\proseitem{1058}{จตฺตาโรเม ภิกฺขเว ธมฺมา ภาวิตา พหุลีกตา มหาปญฺญตาย สํวตฺตนฺติ.\csromanspacer{} กตเม จตฺตาโร, สปฺปุริสสํเสโว สทฺธมฺมสฺสวนํ โยนิโสมนสิกาโร ธมฺมานุธมฺมปฺปฏิปตฺติ.\csromanspacer{} อิเม โข ภิกฺขเว จตฺตาโร ธมฺมา ภาวิตา พหุลีกตา มหาปญฺญตาย สํวตฺตนฺตีติ.\csromanspacer{}\csromanspacer{}ปฐมํ.}

\csromanheader{2}{2. ปุถุปญฺญาสุตฺต}
\vspace{\tipitakaparskip}
\csromancenter{ปุถุปญฺญตาย สํวตฺตนฺตีติ.\csromanspacer{}\csromanspacer{}ทุติยํ.}
\csromanheader{2}{3. วิปุลปญฺญาสุตฺต}
\vspace{\tipitakaparskip}
\csromancenter{วิปุลปญฺญตาย สํวตฺตนฺตีติ.\csromanspacer{}\csromanspacer{}ตติยํ.}
\csromanheader{2}{4. คมฺภีรปญฺญาสุตฺต}
\vspace{\tipitakaparskip}
\csromancenter{คมฺภีรปญฺญตาย สํวตฺตนฺตีติ.\csromanspacer{}\csromanspacer{}จตุตฺถํ.}
\csromanheader{2}{11. โสตาปตฺติสํยุตฺต}
\csromanheader{2}{5. อปฺปมตฺตปญฺญาสุตฺต}
\vspace{\tipitakaparskip}
\csromancenter{อปฺปมตฺตปญฺญตาย สํวตฺตนฺตีติ.\csromanspacer{}\csromanspacer{}ปญฺจมํ.}
\csromanheader{2}{6. ภูริปญฺญาสุตฺต}
\vspace{\tipitakaparskip}
\csromancenter{ภูรปญฺญตาย สํวตฺตนฺตีติ.\csromanspacer{}\csromanspacer{}ฉฏฺฐํ.}
\csromanheader{2}{7. ปญฺญาพาหุลฺลสุตฺต}
\vspace{\tipitakaparskip}
\csromancenter{ปญฺญาพาหุลฺลาย สํวตฺตนฺตีติ.\csromanspacer{}\csromanspacer{}สตฺตมํ.}
\csromanheader{2}{8. สีฆปญฺญาสุตฺต}
\vspace{\tipitakaparskip}
\csromancenter{สีฆปญฺญตาย สํวตฺตนฺตีติ.\csromanspacer{}\csromanspacer{}อฏฺฐมํ.}
\csromanheader{2}{9. ลหุปญฺญาสุตฺต}
\vspace{\tipitakaparskip}
\csromancenter{ลหุปญฺญตาย สํวตฺตนฺตีติ.\csromanspacer{}\csromanspacer{}นวมํ.}
\csromanheader{2}{10. หาสปญฺญาสุตฺต}
\vspace{\tipitakaparskip}
\csromancenter{หาสปญฺญตาย สํวตฺตนฺตีติ.\csromanspacer{}\csromanspacer{}ทสมํ.}
\csromanheader{2}{11. ชวนปญฺญาสุตฺต}
\vspace{\tipitakaparskip}
\csromancenter{ชวนปญฺญตาย สํวตฺตนฺตีติ.\csromanspacer{}\csromanspacer{}เอกาทสมํ.}
\csromanheader{2}{12. ติกฺขปญฺญาสุตฺต}
\vspace{\tipitakaparskip}
\csromancenter{ติกฺขปญฺญตาย สํวตฺตนฺตีติ.\csromanspacer{} ทฺวาทสมํ.}
\csromanheader{2}{13. นิพฺเพธิกปญฺญาสุตฺต}
\proseitem{1070}{\csromanfolio{361}นิพฺเพธิกปญฺญตาย สํวตฺตนฺติ.\csromanspacer{} กตเม จตฺตาโร สปฺปุริสํเสโว สทฺธมฺมสฺสวนํ โยนิโสมนสิกาโร ธมฺมานุธมฺมปฺปฏิปตฺติ. \csromanfolio{362}อิเม โข ภิกฺขเว จตฺตาโร ธมฺมา ภาวิตา พหุลีกตา นิพฺเพธิกปญฺญตาย สํวตฺตนฺตีติ.\csromanspacer{}\csromanspacer{}เตเรสมํ.\csromanspacer{} มหาปญฺญวคฺโค สตฺตโม.}

\titlehead{\textbf{ตสฺสุทฺทานํ}}
\csromangathameasure{มหา ปุถุ วิปุล คมฺภีรํ, อปฺปมตฺต ภูริ พาหุลฺลํ. \\ สีฆลหุ หาส ชวน, ติกฺข นิพฺเพธิกาย จาติ. \\ \textbf{โสตาปตฺติสํยุตฺตํ เอกาทสมํ.}}
\csromangathagroup{\csromangathastanza{มหา ปุถุ วิปุล คมฺภีรํ, อปฺปมตฺต ภูริ พาหุลฺลํ. \\ สีฆลหุ หาส ชวน, ติกฺข นิพฺเพธิกาย จาติ. \\ \textbf{โสตาปตฺติสํยุตฺตํ เอกาทสมํ.}}}

\csromanmatikaanchor{mtk.261}
\csromantocmarknopage{subparagraph}{12. สจฺจสํยุตฺต}{mtk.261}
\bookmark[dest={mtk.261},level=5]{12. สจฺจสํยุตฺต}
\csromanheader{6}{12. สจฺจสํยุตฺต}
\csromanmatikaanchor{mtk.262}
\csromantocmarknopage{subparagraph}{1. สมาธิวคฺค}{mtk.262}
\bookmark[dest={mtk.262},level=5]{1. สมาธิวคฺค}
\csromanheader{6}{1. สมาธิวคฺค}
\csromanmatikaanchor{mtk.263}
\csromantocmarkref{subparagraph}{1. สมาธิสุตฺต}{mtk.263}
\bookmark[dest={mtk.263},level=5]{1. สมาธิสุตฺต}
\csromanheader{6}{1. สมาธิสุตฺต}
\proseitem{1071}{\csromanfolio{363}สาวตฺถินิทานํ ฯเปฯ.\csromanspacer{} สมาธิํ ภิกฺขเว ภาเวถ, สมาหิโต ภิกฺขเว ภิกฺขุ ยถาภูตํ ปชานาติ.\csromanspacer{} กิญฺจ ยถาภูตํ ปชานาติ, “อิทํ ทุกฺขนฺ”ติ ยถาภูตํ ปชานาติ, “อยํ ทุกฺขสมุทโย”ติ ยถาภูตํ ปชานาติ, “อยํ ทุกฺขนิโรโธ”ติ ยถาภูตํ ปชานาติ, “อยํ ทุกฺขนิโรธคามินี ปฏิปทา”ติ ยถาภูตํ ปชานาติ.\csromanspacer{} สมาธิํ ภิกฺขเว ภาเวถ, สมาหิโต ภิกฺขเว ภิกฺขุ ยถาภูตํ ปชานาติ.}

\prose{ตสฺมาติห ภิกฺขเว “อิทํ ทุกฺขนฺ”ติ โยโค กรณีโย, “อยํ ทุกฺขสมุทโย”ติ โยโค กรณีโย, “อยํ ทุกฺขนิโรโธ”ติ โยโค กรณีโย, “อยํ ทุกฺขนิโรธคามินี ปฏิปทา”ติ โยโค กรณีโยติ.\csromanspacer{}\csromanspacer{}ปฐมํ.}

\csromanmatikaanchor{mtk.264}
\csromantocmarkref{subparagraph}{2. ปฏิสลฺลานสุตฺต}{mtk.264}
\bookmark[dest={mtk.264},level=5]{2. ปฏิสลฺลานสุตฺต}
\csromanheader{6}{2. ปฏิสลฺลานสุตฺต}
\proseitem{1072}{ปฏิสลฺลาเน ภิกฺขเว โยคมาปชฺชถ, ปฏิสลฺลีโน ภิกฺขเว ภิกฺขุ ยถาภูตํ ปชานาติ.\csromanspacer{} กิญฺจ ยถาภูตํ ปชานาติ, “อิทํ ทุกฺขนฺ”ติ ยถาภูตํ ปชานาติ, “อยํ ทุกฺขสมุทโย”ติ ยถาภูตํ ปชานาติ, “อยํ ทุกฺขนิโรโธ”ติ ยถาภูตํ ปชานาติ, “อยํ ทุกฺขนิโรธคามินี ปฏิปทา”ติ ยถาภูตํ ปชานาติ.\csromanspacer{} ปฏิสลฺลาเน ภิกฺขเว โยคมาปชฺชถ, ปฏิสลฺลีโน ภิกฺขเว ภิกฺขุ ยถาภูตํ ปชานาติ.}

\prose{ตสฺมาติห ภิกฺขเว “อิทํ ทุกฺขนฺ”ติ โยโค กรณีโย, “อยํ ทุกฺขสมุทโย”ติ โยโค กรณีโย, “อยํ ทุกฺขนิโรโธ”ติ โยโค กรณีโย, “อยํ ทุกฺขนิโรธคามินี ปฏิปทา”ติ โยโค กรณีโยติ.\csromanspacer{}\csromanspacer{}ทุติยํ.}

\csromanheader{2}{3. ปฐมกุลปุตฺตสุตฺต}
\csromanmatikaanchor{mtk.265}
\csromantocmarkref{subparagraph}{3-4. กุลปุตฺตสุตฺตานิ}{mtk.265}
\bookmark[dest={mtk.265},level=5]{3-4. กุลปุตฺตสุตฺตานิ}
\proseitem{1073}{\csromanfolio{364}เย หิ เกจิ ภิกฺขเว อตีตมทฺธานํ กุลปุตฺตา สมฺมา อคารสฺมา อนคาริยํ ปพฺพชิํสุ, สพฺเพ เต จตุนฺนํ อริยสจฺจานํ ยถาภูตํ อภิสมยาย.\csromanspacer{} เย หิ เกจิ ภิกฺขเว อนาคตมทฺธานํ กุลปุตฺตา สมฺมา อคารสฺมา อนคาริยํ ปพฺพชิสฺสนฺติ, สพฺเพ เต จตุนฺนํ อริยสจฺจานํ ยถาภูตํ อภิสมยาย.\csromanspacer{} เย หิ เกจิ ภิกฺขเว เอตรหิ กุลปุตฺตา สมฺมา อคารสฺมา อนคาริยํ ปพฺพชนฺติ, สพฺเพ เต จตุนฺนํ อริยสจฺจานํ ยถาภูตํ อภิสมยาย.}

\prose{กตเมสํ จตุนฺนํ, ทุกฺขสฺส อริยสจฺจสฺส ทุกฺขสมุทยสฺส อริยสจฺจสฺส ทุกฺขนิโรธสฺส อริยสจฺจสฺส ทุกฺขนิโรธคามินิยา ปฏิปทาย อริยสจฺจสฺส.\csromanspacer{} เย หิ เกจิ ภิกฺขเว อตีตมทฺธานํ กุลปุตฺตา สมฺมา อคารสฺมา อนคาริยํ ปพฺพชิํสุ ฯเปฯ ปพฺพชิสฺสนฺติ ฯเปฯ ปพฺพชนฺติ, สพฺเพ เต อิเมสํเยว จตุนฺนํ อริยสจฺจานํ ยถาภูตํ อภิสมยาย.}

\prose{ตสฺมาติห ภิกฺขเว “อิทํ ทุกฺขนฺ”ติ โยโค กรณีโย, “อยํ ทุกฺขสมุทโย”ติ โยโค กรณีโย, “อยํ ทุกฺขนิโรโธ”ติ โยโค กรณีโย, “อยํ ทุกฺขนิโรธคามินี ปฏิปทา”ติ โยโค กรณีโยติ.\csromanspacer{}\csromanspacer{}ตติยํ.}

\csromanheader{2}{4. ทุติยกุลปุตฺตสุตฺต}
\proseitem{1074}{เย หิ เกจิ ภิกฺขเว อตีตมทฺธานํ กุลปุตฺตา สมฺมา อคารสฺมา อนคาริยํ ปพฺพชิตา ยถาภูตํ อภิสเมสุํ, สพฺเพ เต จตฺตาริ อริยสจฺจานิ ยถาภูตํ อภิสเมสุํ.\csromanspacer{} เย หิ เกจิ ภิกฺขเว อนาคตมทฺธานํ กุลปุตฺตา สมฺมา อคารสฺมา อนคาริยํ ปพฺพชิตา ยถาภูตํ อภิสเมสฺสนฺติ, สพฺเพ เต จตฺตาริ อริยสจฺจานิ ยถาภูตํ อภิสเมสฺสนฺติ.\csromanspacer{} เย หิ เกจิ ภิกฺขเว เอตรหิ กุลปุตฺตา สมฺมา อคารสฺมา อนคาริยํ ปพฺพชิตา ยถาภูตํ อภิสเมนฺติ, สพฺเพ เต จตฺตาริ อริยสจฺจานิ ยถาภูตํ อภิสเมนฺติ.}

\prose{กตมานิ จตฺตาริ, ทุกฺขํ อริยสจฺจํ ทุกฺขสมุทยํ อริยสจฺจํ ทุกฺขนิโรธํ อริยสจฺจํ ทุกฺขนิโรธคามินี ปฏิปทา อริยสจฺจํ.\csromanspacer{} เย หิ เกจิ ภิกฺขเว อตีตมทฺธานํ กุลปุตฺตา สมฺมา อคารสฺมา อนคาริยํ ปพฺพชิตา ยถาภูตํ}

\csromanheader{2}{12. สจฺจสํยุตฺต}
\csromanmatikaanchor{mtk.266}
\csromantocmarkref{subparagraph}{5-6. สมณพฺราหฺมณสุตฺตานิ}{mtk.266}
\bookmark[dest={mtk.266},level=5]{5-6. สมณพฺราหฺมณสุตฺตานิ}
\prose{\csromanfolio{365}อภิสเมสุํ ฯเปฯ อภิสเมสฺสนฺติ ฯเปฯ อภิสเมนฺติ, สพฺเพ เต อิมานิ จตฺตาริ อริยสจฺจานิ ยถาภูตํ อภิสเมนฺติ.}

\prose{ตสฺมาติห ภิกฺขเว “อิทํ ทุกฺขนฺ”ติ โยโค กรณีโย ฯเปฯ “อยํ ทุกฺขนิโรธคามินี ปฏิปาทา”ติ โยโค กรณีโยติ.\csromanspacer{}\csromanspacer{}จตุตฺถํ.}

\csromanheader{2}{5. ปฐมสมณพฺราหฺมณสุตฺต}
\proseitem{1075}{เย หิ เกจิ ภิกฺขเว อตีตมทฺธานํ สมณา วา พฺราหฺมณา วา ยถาภูตํ อภิสมฺโพชฺฌิํสุ, สพฺเพ เต จตฺตาริ อริยสจฺจานิ ยถาภูตํ อภิสมฺโพชฺฌิํสุ.\csromanspacer{} เย หิ เกจิ ภิกฺขเว อนาคตมทฺธานํ สมณา วา พฺราหฺมณา วา ยถาภูตํ อภิสมฺโพชฺฌิสฺสนฺติ, สพฺเพ เต จตฺตาริ อริยสจฺจานิ ยถาภูตํ อภิสมฺโพชฺฌิสฺสนฺติ.\csromanspacer{} เย หิ เกจิ ภิกฺขเว เอตรหิ สมณา วา พฺราหฺมณา วา ยถาภูตํ อภิสมฺโพชฺฌนฺติ, สพฺเพ เต จตฺตาริ อริยสจฺจานิ ยถาภูตํ อภิสมฺโพชฺฌนฺติ.}

\prose{กตมานิ จตฺตาริ? ทุกฺขํ อริยสจฺจํ ฯเปฯ ทุกฺขนิโรธคามินี ปฏิปทา อริยสจฺจํ.\csromanspacer{} เย หิ เกจิ ภิกฺขเว อตีตมทฺธานํ สมณา วา พฺราหฺมณา วา ยถาภูตํ อภิสมฺโพชฺฌิํสุ ฯเปฯ อภิสมฺโพชฺฌิสฺสนฺติ ฯเปฯ อภิสมฺโพชฺฌนฺติ, สพฺเพ เต อิมานิ จตฺตาริ อริยสจฺจานิ ยถาภูตํ อภิสมฺโพชฺฌนฺติ.}

\prose{ตสฺมาติห ภิกฺขเว “อิทํ ทุกฺขนฺ”ติ โยโค กรณีโย ฯเปฯ “อยํ ทุกฺขนิโรธคามินี ปฏิปทา”ติ โยโค กรณีโยติ.\csromanspacer{}\csromanspacer{}ปญฺจมํ.}

\csromanheader{2}{6. ทุติยสมณพฺราหฺมณสุตฺต}
\proseitem{1076}{เย หิ เกจิ ภิกฺขเว อตีตมทฺธานํ สมณา วา พฺราหฺมณาวา ยถาภูตํ อภิสมฺพุทฺธํ ปกาเสสุํ, สพฺเพ เต จตฺตาริ อริยสจฺจานิ ยถาภูตํ อภิสมฺพุทฺธํ ปกาเสสุํ.\csromanspacer{} เย หิ เกจิ ภิกฺขเว อนาคตมทฺธานํ สมณา วา พฺราหฺมณา วา ยถาภูตํ อภิสมฺพุทฺธํ ปกาเสสฺสนฺติ, สพฺเพ เต จตฺตาริ อริยสจฺจานิ ยถาภูตํ อภิสมฺพุทฺธํ ปกาเสสฺสนฺติ.\csromanspacer{} เย หิ เกจิ ภิกฺขเว เอตรหิ สมณา วา พฺราหฺมณา วา ยถาภูตํ อภิสมฺพุทฺธํ ปกาเสนฺติ, สพฺเพ เต จตฺตาริ อริยสจฺจานิ ยถาภูตํ อภิสมฺพุทฺธํ ปกาเสนฺติ.}

\csromanmatikaanchor{mtk.267}
\csromantocmarkref{subparagraph}{7-9. วิตกฺกสุตฺตาทีนิ}{mtk.267}
\bookmark[dest={mtk.267},level=5]{7-9. วิตกฺกสุตฺตาทีนิ}
\prose{\csromanfolio{366}กตมานิ จตฺตาริ, ทุกฺขํ อริย สจฺจํ ฯเปฯ ทุกฺขนิโรธคามินี ปฏิปทา อริยสจฺจํ.\csromanspacer{} เย หิ เกจิ ภิกฺขเว อตีตมทฺธานํ สมณา วา พฺราหฺมณา วา ยถาภูตํ อภิสมฺพุทฺธํ ปกาเสสุํ ฯเปฯ ปกาเสสฺสนฺติ ฯเปฯ ปกาเสนฺติ, สพฺเพ เต อิมานิ จตฺตาริ อริยสจฺจานิ ยถาภูตํ อภิสมฺพุทฺธํ ปกาเสนฺติ.}

\prose{ตสฺมาติห ภิกฺขเว “อิทํ ทุกฺขนฺ”ติ โยโค กรณีโย ฯเปฯ “อยํ ทุกฺขนิโรธคามินี ปฏิปทา”ติ โยโค กรณีโยติ.\csromanspacer{}\csromanspacer{}ฉฏฺฐํ.}

\csromanheader{2}{7. วิตกฺกสุตฺต}
\proseitem{1077}{มา ภิกฺขเว ปาปเก อกุสเล วิตกฺเก วิตกฺเกยฺยาถ\footnote{วิตกฺเกถ (สี, สฺยา, กํ)}.\csromanspacer{} เสยฺยถิทํ, กามวิตกฺกํ พฺยาปาทวิตกฺกํ วิหิํสาวิตกฺกํ.\csromanspacer{} ตํ กิสฺส เหตุ, เนเต ภิกฺขเว วิตกฺกา อตฺถสํหิตา, นาทิพฺรหฺมจริยกา, น นิพฺพิทาย น วิราคาย น นิโรธาย น อุปสมาย น อภิญฺญาย น สมฺโพธาย น นิพฺพานาย สํวตฺตนฺติ.}

\prose{วิตกฺเกนฺตา จ โข ตุเมฺห ภิกฺขเว “อิทํ ทุกฺขนฺ”ติ วิตกฺเกยฺยาถ, “อยํ ทุกฺขสมุทโย”ติ วิตกฺเกยฺยาถ, “อยํ ทุกฺขนิโรโธ”ติ วิตกฺเกยฺยาถ, “อยํ ทุกฺขนิโรธคามินี ปฏิปทา”ติ วิตกฺเกยฺยาถ.\csromanspacer{} ตํ กิสฺส เหตุ, เอเต ภิกฺขเว วิตกฺกา อตฺถสํหิตา, เอเต อาทิพฺรหฺมจริยกา, เอเต นิพฺพิทาย วิราคาย นิโรธาย อุปสมาย อภิญฺญาย สมฺโพธาย นิพฺพานาย สํวตฺตนฺติ.}

\prose{ตสฺมาติห ภิกฺขเว “อิทํ ทุกฺขนฺ”ติ โยโค กรณีโย ฯเปฯ “อยํ ทุกฺขนิโรธคามินี ปฏิปทา”ติ โยโค กรณีโยติ.\csromanspacer{}\csromanspacer{}สตฺตมํ.}

\csromanheader{2}{8. จินฺตสุตฺต}
\proseitem{1078}{มา ภิกฺขเว ปาปกํ อกุสลํ จิตฺตํ จินฺเตยฺยาถ\footnote{จินฺเตถ (สี, สฺยา, กํ)}, “สสฺสโต โลโก”ติ วา “อสสฺสโต โลโก”ติ วา, “อนฺตวา โลโก”ติ วา, “อนนฺตวา โลโก”ติ วา, “ตํ ชีวํ ตํ สรีรนฺ”ติ วา, อญฺญํ ชีวํ “อญฺญํ สรีรนฺ”ติ วา, “โหติ ตถาคโต ปรํ มรณา”ติ วา, “น โหติ ตถาคโต ปรํ มรณา”ติ วา, “โหติ จ น จ}

\csromanheader{2}{12. สจฺจสํยุตฺต}
\prose{\csromanfolio{367}โหติ ตถาคโต ปรํ มรณา”ติ วา, “เนว โหติ น น โหติ ตถาคโต ปรํ มรณา”ติ วา.\csromanspacer{} ตํ กิสฺส เหตุ, เนสา ภิกฺขเว จินฺตา อตฺถสํหิตา, นาทิพฺรหฺมจริยกา, น นิพฺพิทาย น วิราคาย น นิโรธาย น อุปสมาย น อภิญฺญาย น สมฺโพธาย น นิพฺพานาย สํวตฺตติ.}

\prose{จินฺเตนฺตา จ โข ตุเมฺห ภิกฺขเว “อิทํ ทุกฺขนฺ”ติ จินฺเตยฺยาถ, “อยํ ทุกฺขสมุทโย”ติ จินฺเตยฺยาถ, “อยํ ทุกฺขนิโรโธ”ติ จินฺเตยฺยาถ, “อยํ ทุกฺขนิโรธคามินี ปฏิปทา”ติ จินฺเตยฺยาถ.\csromanspacer{} ตํ กิสฺส เหตุ, เอสา ภิกฺขเว จินฺตา อตฺถสํหิตา, เอสา อาทิพฺรหฺมจริยกา, เอสา นิพฺพิทาย วิราคาย นิโรธาย อุปสมาย อภิญฺญาย สมฺโพธาย นิพฺพานาย สํวตฺตติ.}

\prose{ตสฺมาติห ภิกฺขเว “อิทํ ทุกฺขนฺ”ติ โยโค กรณีโย ฯเปฯ “อยํ ทุกฺขนิโรธคามินี ปฏิปทา”ติ โยโค กรณีโยติ.\csromanspacer{}\csromanspacer{}อฏฺฐมํ.}

\csromanheader{2}{9. วิคฺคาหิกกถาสุตฺต}
\proseitem{1079}{มา ภิกฺขเว วิคฺคาหิกกถํ กเถยฺยาถ\footnote{กเถถ (สี, สฺยา, กํ)} “น ตฺวํ อิมํ ธมฺมวินยํ อาชานาสิ, อหํ อิมํ ธมฺมวินยํ อาชานามิ.\csromanspacer{} กิํ ตฺวํ อิมํ ธมฺมวินยํ อาชานิสฺสสิ, มิจฺฉาปฏิปนฺโน ตฺวมสิ, อหมสฺมิ สมฺมาปฏิปนฺโน.\csromanspacer{} สหิตํ เม, อสหิตํ เต, ปุเรวจนียํ ปจฺฉา อวจ, ปจฺฉาวจนียํ ปุเร อวจ.\csromanspacer{} อธิจิณฺณํ\footnote{อาจิณฺณํ (สฺยา, กํ, อิ)} เต วิปราวตฺตํ, อาโรปิโต เต วาโท, จร วาทปฺปโมกฺขาย นิคฺคหิโตสิ, นิพฺเพเฐหิ วา สเจ ปโหสี”ติ.\csromanspacer{} ตํ กิสฺส เหตุ, เนสา ภิกฺขเว กถา อตฺถสํหิตา, นาทิพฺรหฺมจริยกา, น นิพฺพิทาย น วิราคาย น นิโรธาย น อุปสมาย น อภิญฺญาย น สมฺโพธาย น นิพฺพานาย สํวตฺตติ.}

\prose{กเถนฺตา จ โข ตุเมฺห ภิกฺขเว “อิทํ ทุกฺขนฺ”ติ กเถยฺยาถ, “อยํ ทุกฺขสมุทโย”ติ กเถยฺยาถ, “อยํ ทุกฺขนิโรโธ”ติ กเถยฺยาถ, “อยํ ทุกฺขนิโรธคามินี ปฏิปทา”ติ กเถยฺยาถ ฯเปฯ โยโค กรณีโยติ.\csromanspacer{}\csromanspacer{}นวมํ.}

\csromanmatikaanchor{mtk.268}
\csromantocmarkref{subparagraph}{10. ติรจฺฉานกถาสุตฺต}{mtk.268}
\bookmark[dest={mtk.268},level=5]{10. ติรจฺฉานกถาสุตฺต}
\csromanheader{6}{10. ติรจฺฉานกถาสุตฺต}
\proseitem{1080}{\csromanfolio{368}มา ภิกฺขเว อเนกวิหิตํ ติรจฺฉานกถํ กเถยฺยาถ.\csromanspacer{} เสยฺยถิทํ, ราชกถํ โจรกถํ มหามตฺตกถํ เสนากถํ ภยกถํ ยุทฺธกถํ อนฺนกถํ ปานกถํ วตฺถกถํ สยนกถํ มาลากถํ คนฺธกถํ ญาติกถํ ยานกถํ คามกถํ นิคมกถํ นครกถํ ชนปทกถํ อิตฺถิกถํ\footnote{อิตฺถิกถํ ปุริสกถํ (สฺยา, กํ, อิ, ก)} สูรกถํ วิสิขากถํ กุมฺภฏฺฐานกถํ ปุพฺพเปตกถํ นานตฺตกถํ โลกกฺขายิกํ สมุทฺทกฺขายิกํ อิติภวาภวกถํ อิติ วา.\csromanspacer{} ตํ กิสฺส เหตุ, เนสา ภิกฺขเว กถา อตฺถสํหิตา, นาทิพฺรหฺมจริยกา, น นิพฺพิทาย น วิราคาย น นิโรธาย น อุปสมาย น อภิญฺญาย น สมฺโพธาย น นิพฺพานาย สํวตฺตติ.}

\prose{กเถนฺตา จ โข ตุเมฺห ภิกฺขเว “อิทํ ทุกฺขนฺ”ติ กเถยฺยาถ, “อยํ ทุกฺขสมุทโย”ติ กเถยฺยาถ, “อยํ ทุกฺขนิโรโธ”ติ กเถยฺยาถ, “อยํ ทุกฺขนิโรธคามินี ปฏิปทา”ติ กเถยฺยาถ.\csromanspacer{} ตํ กิสฺส เหตุ, เอสา ภิกฺขเว กถา อตฺถสํหิตา, เอสา อาทิพฺรหฺมจริยกา, เอสา นิพฺพิทาย วิราคาย นิโรธาย อุปสมาย อภิญฺญาย สมฺโพธาย นิพฺพานาย สํวตฺตติ.}

\prose{ตสฺมาติห ภิกฺขเว “อิทํ ทุกฺขนฺ”ติ โยโค กรณีโย ฯเปฯ “อยํ ทุกฺขนิโรธคามินี ปฏิปทา”ติ โยโค กรณีโยติ.\csromanspacer{} สมาธิวคฺโค ปฐโม.}

\titlehead{\textbf{ตสฺสุทฺทานํ}}
\csromangathameasure{สมาธิ ปฏิสลฺลานา, กุลปุตฺตา อปเร ทุเว. \\ สมณพฺราหฺมณา วิตกฺกํ, จินฺตา วิคฺคาหิกา กถาติ.}
\csromangathagroup{\csromangathastanza{สมาธิ ปฏิสลฺลานา, กุลปุตฺตา อปเร ทุเว. \\ สมณพฺราหฺมณา วิตกฺกํ, จินฺตา วิคฺคาหิกา กถาติ.}}

\csromanmatikaanchor{mtk.269}
\csromantocmarknopage{subparagraph}{2. ธมฺมจกฺกปฺปวตฺตนวคฺค}{mtk.269}
\bookmark[dest={mtk.269},level=5]{2. ธมฺมจกฺกปฺปวตฺตนวคฺค}
\csromanheader{6}{2. ธมฺมจกฺกปฺปวตฺตนวคฺค}
\csromanmatikaanchor{mtk.270}
\csromantocmarkref{subparagraph}{1. ธมฺมจกฺกปฺปวตฺตนสุตฺต}{mtk.270}
\bookmark[dest={mtk.270},level=5]{1. ธมฺมจกฺกปฺปวตฺตนสุตฺต}
\csromanheader{6}{1. ธมฺมจกฺกปฺปวตฺตนสุตฺต}
\proseitem{1081}{เอกํ สมยํ ภควา พาราณสิยํ วิหรติ อิสิปตเน มิคทาเย.\csromanspacer{} ตตฺร โข ภควา ปญฺจวคฺคิเย ภิกฺขู อามนฺเตสิ\csromandash{}เทฺวเม ภิกฺขเว อนฺตา}

\csromanheader{2}{12. สจฺจสํยุตฺต}
\prose{\csromanfolio{369}ปพฺพชิเตน น เสวิตพฺพา.\csromanspacer{} กตเม เทฺว, โย จายํ กาเมสุ กามสุขลฺลิกานุโยโค หีโน คมฺโม โปถุชฺชนิโก อนริโย อนตฺถสํหิโต, โย จายํ อตฺตกิลมถานุโยโค ทุกฺโข อนริโย อนตฺถสํหิโต.\csromanspacer{} เอเต โข ภิกฺขเว อุโภ อนฺเต อนุปคมฺม มชฺฌิมา ปฏิปทา ตถาคเตน อภิสมฺพุทฺธา จกฺขุกรณี ญาณกรณี อุปสมาย อภิญฺญาย สมฺโพธาย นิพฺพานาย สํวตฺตติ.}

\prose{กตมา จ สา ภิกฺขเว มชฺฌิมา ปฏิปทา ตถาคเตน อภิสมฺพุทฺธา จกฺขุกรณี ญาณกรณี อุปสมาย อภิญฺญาย สมฺโพธาย นิพฺพานาย สํวตฺตติ, อยเมว อริโย อฏฺฐงฺคิโก มคฺโค.\csromanspacer{} เสยฺยถิทํ, สมฺมาทิฏฺฐิ สมฺมาสงฺกปฺโป สมฺมาวาจา สมฺมากมฺมนฺโต สมฺมา-อาชีโว สมฺมาวายาโม สมฺมาสติ สมฺมาสมาธิ.\csromanspacer{} อยํ โข สา ภิกฺขเว มชฺฌิมา ปฏิปทา ตถาคเตน อภิสมฺพุทฺธา จกฺขุกรณี ญาณกรณี อุปสมาย อภิญฺญาย สมฺโพธาย นิพฺพานาย สํวตฺตติ.}

\prose{อิทํ โข ปน ภิกฺขเว ทุกฺขํ อริยสจฺจํ, ชาติปิ ทุกฺขา, ชราปิ ทุกฺขา, พฺยาธิปิ ทุกฺโข, มรณมฺปิ ทุกฺขํ, อปฺปิเยหิ สมฺปโยโค ทุกฺโข, ปิเยหิ วิปฺปโยโค ทุกฺโข, ยมฺปิจฺฉํ น ลภติ, ตมฺปิ ทุกฺขํ, สํขิตฺเตน ปญฺจุปาทานกฺขนฺธา\footnote{ปญฺจุปาทานกฺขนฺธาปิ (อิ, ก)} ทุกฺขา.\csromanspacer{} อิทํ โข ปน ภิกฺขเว ทุกฺขสมุทยํ อริยสจฺจํ, ยายํ ตณฺหา โปโนพฺภวิกา\footnote{โปโนภวิกา (สี, อิ)} นนฺทิราคสหคตา ตตฺรตตฺราภินนฺทินี.\csromanspacer{} เสยฺยถิทํ\footnote{เสยฺยถีทํ (สี, สฺยา, กํ, อิ)}, กามตณฺหา ภวตณฺหา วิภวตณฺหา.\csromanspacer{} อิทํ โข ปน ภิกฺขเว ทุกฺขนิโรธํ อริยสจฺจํ, โย ตสฺสาเยว ตณฺหาย อเสสวิราคนิโรโธ จาโค ปฏินิสฺสคฺโค มุตฺติ อนาลโย.\csromanspacer{} อิทํ โข ปน ภิกฺขเว ทุกฺขนิโรธคามินี ปฏิปทา อริยสจฺจํ, อยเมว อริโย อฏฺฐงฺคิโก มคฺโค.\csromanspacer{} เสยฺยถิทํ, สมฺมาทิฏฺฐิ ฯเปฯ สมฺมาสมาธิ.}

\prose{“อิทํ ทุกฺขํ อริยสจฺจนฺ”ติ เม ภิกฺขเว ปุพฺเพ อนนุสฺสุเตสุ ธมฺเมสุ จกฺขุํ อุทปาทิ, ญาณํ อุทปาทิ, ปญฺญา อุทปาทิ, วิชฺชา อุทปาทิ, อาโลโก อุทปาทิ. “ตํ โข ปนิทํ ทุกฺขํ อริยสจฺจํ ปริญฺเญยฺยนฺ”ติ เม ภิกฺขเว ปุพฺเพ ฯเปฯ อุทปาทิ. “ตํ โข ปนิทํ ทุกฺขํ อริยสจฺจํ ปริญฺญาตนฺ”ติ เม ภิกฺขเว ปุพฺเพ อนนุสฺสุเตสุ ธมฺเมสุ จกฺขุํ อุทปาทิ, ญาณํ อุทปาทิ, ปญฺญา อุทปาทิ วิชฺชา อุทปาทิ, อาโลโก อุทปาทิ.}

\prose{\csromanfolio{370}“อิทํ ทุกฺขสมุทยํ อริยสจฺจนฺ”ติ เม ภิกฺขเว ปุพฺเพ อนนุสฺสุเตสุ ธมฺเมสุ จกฺขุํ อุทปาทิ, ญาณํ อุทปาทิ, ปญฺญา อุทปาทิ, วิชฺชา อุทปาทิ, อาโลโก อุทปาทิ. “ตํ โข ปนิทํ ทุกฺขสมุทยํ อริยสจฺจํ ปหาตพฺพนฺ”ติ เม ภิกฺขเว ปุพฺเพ ฯเปฯ อุทปาทิ. “ตํ โข ปนิทํ ทุกฺขสมุทยํ อริยสจฺจํ ปหีนนฺ”ติ เม ภิกฺขเว ปุพฺเพ อนนุสฺสุเตสุ ธมฺเมสุ จกฺขุํ อุทปาทิ, ญาณํ อุทปาทิ, ปญฺญา อุทปาทิ, วิชฺชา อุทปาทิ, อาโลโก อุทปาทิ.}

\prose{“อิทํ ทุกฺขนิโรธํ อริยสจฺจนฺ”ติ เม ภิกฺขเว ปุพฺเพ อนนุสฺสุเตสุ ธมฺเมสุ จกฺขุํ อุทปาทิ, ญาณํ อุทปาทิ, ปญฺญา อุทปาทิ, วิชฺชา อุทปาทิ, อาโลโก อุทปาทิ. “ตํ โข ปนิทํ ทุกฺขนิโรธํ อริยสจฺจํ สจฺฉิกาตพฺพนฺ”ติ เม ภิกฺขเว ปุพฺเพ ฯเปฯ อุทปาทิ. “ตํ โข ปนิทํ ทุกฺขนิโรธํ อริยสจฺจํ สจฺฉิกตนฺ”ติ เม ภิกฺขเว ปุพฺเพ อนนุสฺสุเตสุ ธมฺเมสุ จกฺขุํ อุทปาทิ, ญาณํ อุทปาทิ, ปญฺญา อุทปาทิ, วิชฺชา อุทปาทิ, อาโลโก อุทปาทิ.}

\prose{“อิทํ ทุกฺขนิโรธคามินี ปฏิปทา อริยสจฺจนฺ”ติ เม ภิกฺขเว ปุพฺเพ อนนุสฺสุเตสุ ธมฺเมสุ จกฺขุํ อุทปาทิ, ญาณํ อุทปาทิ, ปญฺญา อุทปาทิ, วิชฺชา อุทปาทิ, อาโลโก อุทปาทิ. “ตํ โข ปนิทํ ทุกฺขนิโรธคามินี ปฏิปทา อริยสจฺจํ ภาเวตพฺพนฺ”ติ เม ภิกฺขเว ปุพฺเพ ฯเปฯ อุทปาทิ. “ตํ โข ปนิทํ ทุกฺขนิโรธคามินี ปฏิปทา อริยสจฺจํ ภาวิตนฺ”ติ เม ภิกฺขเว ปุพฺเพ อนนุสฺสุเตสุ ธมฺเมสุ จกฺขุํ อุทปาทิ, ญาณํ อุทปาทิ, ปญฺญา อุทปาทิ, วิชฺชา อุทปาทิ, อาโลโก อุทปาทิ.}

\prose{ยาวกีวญฺจ เม ภิกฺขเว อิเมสุ จตูสุ อริยสจฺเจสุ เอวํ ติปริวฏฺฏํ ทฺวาทสาการํ ยถาภูตํ ญาณทสฺสนํ น สุวิสุทฺธํ อโหสิ, เนว ตาวาหํ ภิกฺขเว สเทวเก โลเก สมารเก สพฺรหฺมเก สสฺสมณพฺราหฺมณิยา ปชาย สเทวมนุสฺสาย อนุตฺตรํ สมฺมาสมฺโพธิํ อภิสมฺพุทฺโธติ ปจฺจญฺญาสิํ\footnote{อภิสมฺพุทฺโธ ปจฺจญฺญาสิํ (สี, สฺยา, กํ)}.}

\prose{ยโต จ โข เม ภิกฺขเว อิเมสุ จตูสุ อริยสจฺเจสุ เอวํ ติปริวฏฺฏํ ทฺวาทสาการํ ยถาภูตํ ญาณทสฺสนํ สุวิสุทฺธํ อโหสิ, อถาหํ ภิกฺขเว สเทวเก โลเก สมารเก สพฺรหฺมเก สสฺสมณพฺราหฺมณิยา ปชาย สเทวมนุสฺสาย อนุตฺตรํ สมฺมาสมฺโพธิํ อภิสมฺพุทฺโธติ ปจฺจญฺญาสิํ, ญาณญฺจ ปน เม ทสฺสนํ อุทปาทิ, อกุปฺปา เม}

\csromanheader{2}{12. สจฺจสํยุตฺต}
\csromanmatikaanchor{mtk.271}
\csromantocmarkref{subparagraph}{2-8. ตถาคตสุตฺตาทีนิ}{mtk.271}
\bookmark[dest={mtk.271},level=5]{2-8. ตถาคตสุตฺตาทีนิ}
\prose{\csromanfolio{371}วิมุตฺติ\footnote{เจโตวิมุตฺติ (สี, อิ)}, อยมนฺติมา ชาติ, นตฺถิทานิ ปุนพฺภโวติ.\csromanspacer{} อิทมโวจ ภควา, อตฺตมนา ปญฺจวคฺคิยา ภิกฺขู ภควโต ภาสิตํ อภินนฺทุนฺติ.}

\prose{อิมสฺมิํ จ ปน เวยฺยากรณสฺมิํ ภญฺญมาเน อายสฺมโต โกณฺฑญฺญสฺส วิรชํ วีตมลํ ธมฺมจกฺขุํ อุทปาทิ “ยํ กิญฺจิ สมุทยธมฺมํ, สพฺพํ ตํ นิโรธธมฺมนฺ”ติ.}

\prose{ปวตฺติเต จ ปน ภควตา ธมฺมจกฺเก ภุมฺมา เทวา สทฺทมนุสฺสาเวสุํ “เอตํ ภควตา พาราณสิยํ อิสิปตเน มิคทาเย อนุตฺตรํ ธมฺมจกฺกํ ปวตฺติตํ อปฺปฏิวตฺติยํ สมเณน วา พฺราหฺมเณน วา เทเวน วา มาเรน วา พฺรหฺมุนา วา เกนจิ วา โลกสฺมินฺ”ติ.\csromanspacer{} ภุมฺมานํ เทวานํ สทฺทํ สุตฺวา จาตุมหาราชิกา เทวา สทฺทมนุสฺสาเวสุํ “เอตํ ภควตา พาราณสิยํ อิสิปตเน มิคทาเย อนุตฺตรํ ธมฺมจกฺกํ ปวตฺติตํ อปฺปฏิวตฺติยํ สมเณน วา พฺราหฺมเณน วา เทเวน วา มาเรน วา พฺรหฺมุนา วา เกนจิ วา โลกสฺมินฺ”ติ.\csromanspacer{} จาตุมหาราชิกานํ เทวานํ สทฺทํ สุตฺวา ตาวติํสา เทวา ฯเปฯ.\csromanspacer{} ยามา เทวา ฯเปฯ.\csromanspacer{} ตุสิตา เทวา ฯเปฯ.\csromanspacer{} นิมฺมานรตี เทวา ฯเปฯ.\csromanspacer{} ปรนิมฺมิตวสวตฺตี เทวา ฯเปฯ.\csromanspacer{} พฺรหฺมกายิกา เทวา สทฺทมนุสฺสาเวสุํ “เอตํ ภควตา พาราณสิยํ อิสิปตเน มิคทาเย อนุตฺตรํ ธมฺมจกฺกํ ปวตฺติตํ อปฺปฏิวตฺติยํ สมเณน วา พฺราหฺมเณน วา เทเวน วา มาเรน วา พฺรหฺมุนา วา เกนจิ วา โลกสฺมินฺ”ติ.}

\prose{อิติห เตน ขเณน (เตน ลเยน)\footnote{( ) นตฺถิ (สี, สฺยา, กํ)} เตน มุหุตฺเตน ยาว พฺรหฺมโลกา สทฺโท อพฺภุคฺคจฺฉิ, อยญฺจ ทสสหสฺสิโลกธาตุ สงฺกมฺปิ สมฺปกมฺปิ สมฺปเวธิ, อปฺปมาโณ จ อุฬาโร โอภาโส โลเก ปาตุรโหสิ อติกฺกมฺม เทวานํ เทวานุภาวนฺติ.}

\prose{อถ โข ภควา อิมํ อุทานํ อุทาเนสิ “อญฺญาสิ วต โภ โกณฺฑญฺโญ, อญฺญาสิ วต โภ โกณฺฑญฺโญ”ติ.\csromanspacer{} อิติ หิทํ อายสฺมโต โกณฺฑญฺญสฺส “อญฺญาสิโกณฺฑญฺโญ”เตฺวว นามํ อโหสีติ.\csromanspacer{}\csromanspacer{}ปฐมํ.}

\csromanheader{2}{2. ตถาคตสุตฺต}
\proseitem{1082}{“อิทํ ทุกฺขํ อริยสจฺจนฺ”ติ ภิกฺขเว ตถาคตานํ ปุพฺเพ อนนุสฺสุเตสุ ธมฺเมสุ จกฺขุํ อุทปาทิ, ญาณํ อุทปาทิ, ปญฺญา อุทปาทิ, วิชฺชา อุทปาทิ, อาโลโก อุทปาทิ. “ตํ โข ปนิทํ ทุกฺขํ อริยสจฺจํ \csromanfolio{372}ปริญฺเญยฺยนฺ”ติ ภิกฺขเว ตถาคตานํ ปุพฺเพ ฯเปฯ อุทปาทิ. “ตํ โข ปนิทํ ทุกฺขํ อริยสจฺจํ ปริญฺญาตนฺ”ติ ภิกฺขเว ตถาคตานํ ปุพฺเพ อนนุสฺสุเตสุ ธมฺเมสุ จกฺขุํ อุทปาทิ, ญาณํ อุทปาทิ, ปญฺญา อุทปาทิ, วิชฺชา อุทปาทิ, อาโลโก อุทปาทิ.}

\prose{“อิทํ ทุกฺขสมุทยํ อริยสจฺจนฺ”ติ ภิกฺขเว ตถาคตานํ ปุพฺเพ อนนุสฺสุเตสุ ธมฺเมสุ จกฺขุํ อุทปาทิ, ญาณํ อุทปาทิ, ปญฺญา อุทปาทิ, วิชฺชา อุทปาทิ, อาโลโก อุทปาทิ. “ตํ โข ปนิทํ ทุกฺขสมุทยํ อริยสจฺจํ ปหาตพฺพนฺ”ติ ภิกฺขเว ตถาคตานํ ปุพฺเพ ฯเปฯ อุทปาทิ. “ตํ โข ปนิทํ ทุกฺขสมุทยํ อริยสจฺจํ ปหีนนฺ”ติ ภิกฺขเว ตถาคตานํ ปุพฺเพ อนนุสฺสุเตสุ ธมฺเมสุ จกฺขุํ อุทปาทิ, ญาณํ อุทปาทิ, ปญฺญา อุทปาทิ, วิชฺชา อุทปาทิ, อาโลโก อุทปาทิ.}

\prose{“อิทํ ทุกฺขนิโรธํ อริยสจฺจนฺ”ติ ภิกฺขเว ตถาคตานํ ปุพฺเพ อนนุสฺสุเตสุ ธมฺเมสุ จกฺขุํ อุทปาทิ, ญาณํ อุทปาทิ, ปญฺญา อุทปาทิ, วิชฺชา อุทปาทิ, อาโลโก อุทปาทิ. “ตํ โข ปนิทํ ทุกฺขนิโรธํ อริยสจฺจํ สจฺฉิกาตพฺพนฺ”ติ ภิกฺขเว ตถาคตานํ ปุพฺเพ ฯเปฯ อุทปาทิ. “ตํ โข ปนิทํ ทุกฺขนิโรธํ อริยสจฺจํ สจฺฉิกตนฺ”ติ ภิกฺขเว ตถาคตานํ ปุพฺเพ อนนุสฺสุเตสุ ธมฺเมสุ จกฺขุํ อุทปาทิ, ญาณํ อุทปาทิ, ปญฺญา อุทปาทิ, วิชฺชา อุทปาทิ, อาโลโก อุทปาทิ.}

\prose{“อิทํ ทุกฺขนิโรธคามินี ปฏิปทา อริยสจฺจนฺ”ติ ภิกฺขเว ตถาคตานํ ปุพฺเพ อนนุสฺสุเตสุ ธมฺเมสุ จกฺขุํ อุทปาทิ, ญาณํ อุทปาทิ, ปญฺญา อุทปาทิ, วิชฺชา อุทปาทิ, อาโลโก อุทปาทิ. “ตํ โข ปนิทํ ทุกฺขนิโรธคามินี ปฏิปทา อริยสจฺจํ ภาเวตพฺพนฺ”ติ ภิกฺขเว ตถาคตานํ ปุพฺเพ ฯเปฯ อุทปาทิ. “ตํ โข ปนิทํ ทุกฺขนิโรธคามินี ปฏิปทา อริยสจฺจํ ภาวิตนฺ”ติ ภิกฺขเว ตถาคตานํ ปุพฺเพ อนนุสฺสุเตสุ ธมฺเมสุ จกฺขุํ อุทปาทิ, ญาณํ อุทปาทิ, ปญฺญา อุทปาทิ, วิชฺชา อุทปาทิ, อาโลโก อุทปาทีติ.\csromanspacer{}\csromanspacer{}ทุติยํ.}

\csromanheader{2}{3. ขนฺธสุตฺต}
\proseitem{1083}{จตฺตาริมานิ ภิกฺขเว อริยสจฺจานิ.\csromanspacer{} กตมานิ จตฺตาริ, ทุกฺขํ อริยสจฺจํ ทุกฺขสมุทยํ อริยสจฺจํ ทุกฺขนิโรธํ อริยสจฺจํ\footnote{ทุกฺขสมุทโย อริยสจฺจํ ทุกฺขนิโรโธ อริยสจฺจํ (สฺยา)} ทุกฺขนิโรธคามินี ปฏิปทา อริยสจฺจํ.}

\csromanheader{2}{12. สจฺจสํยุตฺต}
\prose{\csromanfolio{373}กตมญฺจ ภิกฺขเว ทุกฺขํ อริยสจฺจํ, “ปญฺจุปาทานกฺขนฺธา”ติสฺส วจนียํ.\csromanspacer{} เสยฺยถิทํ\footnote{กตเม ปญฺจ (สี, สฺยา, กํ)}, รูปุปาทานกฺขนฺโท ฯเปฯ วิญฺญาณุปาทานกฺขนฺโธ.\csromanspacer{} อิทํ วุจฺจติ ภิกฺขเว ทุกฺขํ อริยสจฺจํ.}

\prose{กตมญฺจ ภิกฺขเว ทุกฺขสมุทยํ อริยสจฺจํ, ยายํ ตณฺหา โปโนพฺภวิกา นนฺทิราคสหคตา ตตฺรตตฺราภินนฺทินี.\csromanspacer{} เสยฺยถิทํ, กามตณฺหา ภวตณฺหา วิภวตณฺหา.\csromanspacer{} อิทํ วุจฺจติ ภิกฺขเว ทุกฺขสมุทยํ อริยสจฺจํ.}

\prose{กตมญฺจ ภิกฺขเว ทุกฺขนิโรธํ อริยสจฺจํ, โย ตสฺสาเยว ตณฺหาย อเสสวิราคนิโรโธ จาโค ปฏินิสฺสคฺโค มุตฺติ อนาลโย.\csromanspacer{} อิทํ วุจฺจติ ภิกฺขเว ทุกฺขนิโรธํ อริยสจฺจํ.}

\prose{กตมญฺจ ภิกฺขเว ทุกฺขนิโรธคามินี ปฏิปทา อริยสจฺจํ, อยเมว อริโย อฏฺฐงฺคิโก มคฺโค.\csromanspacer{} เสยฺยถิทํ, สมฺมาทิฏฺฐิ ฯเปฯ สมฺมาสมาธิ.\csromanspacer{} อิทํ วุจฺจติ ภิกฺขเว ทุกฺขนิโรธคามินี ปฏิปทา อริยสจฺจํ.\csromanspacer{} อิมานิ โข ภิกฺขเว จตฺตาริ อริยสจฺจานิ.}

\prose{ตสฺมาติห ภิกฺขเว “อิทํ ทุกฺขนฺ”ติ โยโค กรณีโย ฯเปฯ “อยํ ทุกฺขนิโรธคามินี ปฏิปทา”ติ โยโค กรณีโยติ.\csromanspacer{}\csromanspacer{}ตติยํ.}

\csromanheader{2}{4. อชฺฌตฺติกายตนสุตฺต}
\proseitem{1084}{จตฺตาริมานิ ภิกฺขเว อริยสจฺจานิ.\csromanspacer{} กตมานิ จตฺตาริ, ทุกฺขํ อริยสจฺจํ ทุกฺขสมุทยํ อริยสจฺจํ ทุกฺขนิโรธํ อริยสจฺจํ ทุกฺขนิโรธคามินี ปฏิปทา อริยสจฺจํ.}

\prose{กตมญฺจ ภิกฺขเว ทุกฺขํ อริยสจฺจํ, “ฉ อชฺฌตฺติกานิ อายตนานี”ติสฺส วจนียํ.\csromanspacer{} กตมานิ ฉ, จกฺขายตนํ ฯเปฯ มนายตนํ.\csromanspacer{} อิทํ วุจฺจติ ภิกฺขเว ทุกฺขํ อริยสจฺจํ.}

\prose{กตมญฺจ ภิกฺขเว ทุกฺขสมุทยํ อริยสจฺจํ, ยายํ ตณฺหา โปโนพฺภวิกา นนฺทิราคสหคตา ตตฺรตตฺราภินนฺทินี.\csromanspacer{} เสยฺยถิทํ, กามตณฺหา ภวตณฺหา วิภวตณฺหา.\csromanspacer{} อิทํ วุจฺจติ ภิกฺขเว ทุกฺขสมุทยํ อริยสจฺจํ.}

\prose{\csromanfolio{374}กตมญฺจ ภิกฺขเว ทุกฺขนิโรธํ อริยสจฺจํ, โย ตสฺสาเยว ตณฺหาย อเสสวิราคนิโรโธ จาโค ปฏินิสฺสคฺโค มุตฺติ อนาลโย.\csromanspacer{} อิทํ วุจฺจติ ภิกฺขเว ทุกฺขนิโรธํ อริยสจฺจํ.}

\prose{กตมญฺจ ภิกฺขเว ทุกฺขนิโรธคามินี ปฏิปทา อริยสจฺจํ, อยเมว อริโย อฏฺฐงฺคิโก มคฺโค.\csromanspacer{} เสยฺยถิทํ, สมฺมาทิฏฺฐิ ฯเปฯ สมฺมาสมาธิ.\csromanspacer{} อิทํ วุจฺจติ ภิกฺขเว ทุกฺขนิโรธคามินี ปฏิปทา อริยสจฺจํ.\csromanspacer{} อิมานิ โข ภิกฺขเว จตฺตาริ อริยสจฺจานิ.}

\prose{ตสฺมาติห ภิกฺขเว “อิทํ ทุกฺขนฺ”ติ โยโค กรณีโย ฯเปฯ “อยํ ทุกฺขนิโรธคามินี ปฏิปทา”ติ โยโค กรณีโยติ.\csromanspacer{}\csromanspacer{}จตุตฺถํ.}

\csromanheader{2}{5. ปฐมธารณสุตฺต}
\proseitem{1085}{ธาเรถ โน ตุเมฺห ภิกฺขเว “มยา จตฺตาริ อริยสจฺจานิ เทสิตานี”ติ.\csromanspacer{} เอวํ วุตฺเต อญฺญตโร ภิกฺขุ ภควนฺตํ เอตทโวจ “อหํ โข ภนฺเต ธาเรมิ ภควตา จตฺตาริ อริยสจฺจานิ เทสิตานี”ติ.\csromanspacer{} ยถา กถํ ปน ตฺวํ ภิกฺขุ ธาเรสิ “มยา จตฺตาริ อริยสจฺจานิ เทสิตานี”ติ.\csromanspacer{} ทุกฺขํ ขฺวาหํ ภนฺเต ภควตา ปฐมํ อริยสจฺจํ เทสิตํ ธาเรมิ, ทุกฺขสมุทยํ ขฺวาหํ ภนฺเต ภควตา ทุติยํ อริยสจฺจํ เทสิตํ ธาเรมิ, ทุกฺขนิโรธํ ขฺวาหํ ภนฺเต ภควตา ตติยํ อริยสจฺจํ เทสิตํ ธาเรมิ, ทุกฺขนิโรธคามินิํ ปฏิปทํ ขฺวาหํ ภนฺเต ภควตา จตุตฺถํ อริยสจฺจํ เทสิตํ ธาเรมิ.\csromanspacer{} เอวํ ขฺวาหํ ภนฺเต ธาเรมิ “ภควตา จตฺตาริ อริยสจฺจานิ เทสิตานี”ติ.}

\prose{สาธุ สาธุ ภิกฺขุ, สาธุ โข ตฺวํ ภิกฺขุ ธาเรสิ “มยา จตฺตาริ อริยสจฺจานิ เทสิตานี”ติ.\csromanspacer{} ทุกฺขํ โข ภิกฺขุ มยา ปฐมํ อริยสจฺจํ เทสิตํ, ตถา นํ ธาเรหิ, ทุกฺขสมุทยํ\footnote{ทุกฺขสมุทโย (สฺยา, กํ)} โข ภิกฺขุ มยา ทุติยํ อริยสจฺจํ, เทสิตํ, ตถา นํ ธาเรหิ, ทุกฺขนิโรธํ\footnote{ทุกฺขนิโรโธ (สฺยา, กํ)} โข ภิกฺขุ มยา ตติยํ อริยสจฺจํ เทสิตํ, ตถา นํ ธาเรหิ, ทุกฺขนิโรธคามินี ปฏิปทา\footnote{ทุกฺขนิโรธคามินิปฏิปทํ (อิ), ทุกฺขนิโรธคามินิํ ปฏิปทํ (ก)} โข ภิกฺขุ มยา จตุตฺถํ อริยสจฺจํ เทสิตํ, ตถา นํ ธาเรหิ.\csromanspacer{} เอวํ โข ภิกฺขุ ธาเรหิ “มยา จตฺตาริ อริยสจฺจานิ เทสิตานี”ติ.}

\csromanheader{2}{12. สจฺจสํยุตฺต}
\prose{\csromanfolio{375}ตสฺมาติห ภิกฺขุ “อิทํ ทุกฺขนฺ”ติ โยโค กรณีโย ฯเปฯ “อยํ ทุกฺขนิโรธคามินี ปฏิปทา”ติ โยโค กรณีโยติ.\csromanspacer{}\csromanspacer{}ปญฺจมํ.}

\csromanheader{2}{6. ทุติยธารณสุตฺต}
\proseitem{1086}{ธาเรถ โน ตุเมฺห ภิกฺขเว “มยา จตฺตาริ อริยสจฺจานิ เทสิตานี”ติ.\csromanspacer{} เอวํ วุตฺเต อญฺญตโร ภิกฺขุ ภควนฺตํ เอตทโวจ “อหํ โข ภนฺเต ธาเรมิ ภควตา จตฺตาริ อริยสจฺจานิ เทสิตานี”ติ.\csromanspacer{} ยถา กถํ ปน ตฺวํ ภิกฺขุ ธาเรสิ “มยา จตฺตาริ อริยสจฺจานิ เทสิตานี”ติ.}

\prose{ทุกฺขํ ขฺวาหํ ภนฺเต ภควตา ปฐมํ อริยสจฺจํ เทสิตํ ธาเรมิ.\csromanspacer{} โย หิ โกจิ ภนฺเต สมโณ วา พฺราหฺมโณ วา เอวํ วเทยฺย “เนตํ ทุกฺขํ ปฐมํ อริยสจฺจํ, ยํ สมเณน โคตเมน เทสิตํ, อหเมตํ ทุกฺขํ ปฐมํ อริยสจฺจํ ปจฺจกฺขาย อญฺญํ ทุกฺขํ ปฐมํ อริยสจฺจํ ปญฺญเปสฺสามี”ติ เนตํ ฐานํ วิชฺชติ.\csromanspacer{} ทุกฺขสมุทยํ ขฺวาหํ ภนฺเต ภควตา ฯเปฯ ทุกฺขนิโรธคามินิํ ปฏิปทํ ขฺวาหํ ภนฺเต ภควตา จตุตฺถํ อริยสจฺจํ เทสิตํ ธาเรมิ, โย หิ โกจิ ภนฺเต สมโณ วา พฺราหฺมโณ วา เอวํ วเทยฺย “เนตํ ทุกฺขนิโรธคามินี ปฏิปทา จตุตฺถํ อริยสจฺจํ, ยํ สมเณน โคตเมน เทสิตํ, อหเมตํ ทุกฺขนิโรธคามินิํ ปฏิปทํ จตุตฺถํ อริยสจฺจํ ปจฺจกฺขาย อญฺญํ ทุกฺขนิโรธคามินิํ ปฏิปทํ จตุตฺถํ อริยสจฺจํ ปญฺญเปสฺสามี”ติ เนตํ ฐานํ วิชฺชติ.\csromanspacer{} เอวํ ขฺวาหํ ภนฺเต ธาเรมิ “ภควตา จตฺตาริ อริยสจฺจานิ เทสิตานี”ติ.}

\prose{สาธุ สาธุ ภิกฺขุ, สาธุ โข ตฺวํ ภิกฺขุ ธาเรสิ “มยา จตฺตาริ อริยสจฺจานิ เทสิตานี”ติ.\csromanspacer{} ทุกฺขํ โข ภิกฺขุ มยา ปฐมํ อริยสจฺจํ เทสิตํ, ตถา นํ ธาเรหิ.\csromanspacer{} โย หิ โกจิ ภิกฺขุ สมโณ วา พฺราหฺมโณ วา เอวํ วเทยฺย “เนตํ ทุกฺขํ ปฐมํ อริยสจฺจํ, ยํ สมเณน โคตเมน เทสิตํ, อหเมตํ ทุกฺขํ ปฐมํ อริยสจฺจํ ปจฺจกฺขาย อญฺญํ ทุกฺขํ ปฐมํ อริยสจฺจํ ปญฺญเปสฺสามี”ติ เนตํ ฐานํ วิชฺชติ.\csromanspacer{} ทุกฺขสมุทยํ โข ภิกฺขุ ฯเปฯ.\csromanspacer{} ทุกฺขนิโรธํ โข ภิกฺขุ ฯเปฯ.\csromanspacer{} ทุกฺขนิโรธคามินี ปฏิปทา โข ภิกฺขุ มยา จตุตฺถํ อริยสจฺจํ เทสิตํ, ตถา นํ ธาเรหิ.\csromanspacer{} โย หิ โกจิ ภิกฺขุ สมโณ วา พฺราหฺมโณ วา เอวํ วเทยฺย “เนตํ ทุกฺขนิโรธคามินี ปฏิปทา จตุตฺถํ อริยสจฺจํ, ยํ สมเณน โคตเมน \csromanfolio{376}เทสิตํ, อหเมตํ ทุกฺขนิโรธคามินิํ ปฏิปทํ จตุตฺถํ อริยสจฺจํ ปจฺจกฺขาย อญฺญํ ทุกฺขนิโรธคามินิํ ปฏิปทํ จตุตฺถํ อริยสจฺจํ ปญฺญเปสฺสามี”ติ เนตํ ฐานํ วิชฺชติ.\csromanspacer{} เอวํ โข ตฺวํ ภิกฺขุ ธาเรหิ “มยา จตฺตาริ อริยสจฺจานิ เทสิตานี”ติ.}

\prose{ตสฺมาติห ภิกฺขุ “อิทํ ทุกฺขนฺ”ติ โยโค กรณีโย ฯเปฯ “อยํ ทุกฺขนิโรธคามินี ปฏิปทา”ติ โยโค กรณีโยติ.\csromanspacer{}\csromanspacer{}ฉฏฺฐํ.}

\csromanheader{2}{7. อวิชฺชาสุตฺต}
\proseitem{1087}{เอกมนฺตํ นิสินฺโน โข โส ภิกฺขุ ภควนฺตํ เอตทโวจ “อวิชฺชา อวิชฺชาติ ภนฺเต วุจฺจติ, กตมา นุ โข ภนฺเต อวิชฺชา, กิตฺตาวตา จ อวิชฺชาคโต โหตี”ติ.\csromanspacer{} ยํ โข ภิกฺขุ ทุกฺเข อญฺญาณํ, ทุกฺขสมุทเย อญฺญาณํ, ทุกฺขนิโรเธ อญฺญาณํ, ทุกฺขนิโรธคามินิยา ปฏิปทาย อญฺญาณํ.\csromanspacer{} อยํ วุจฺจติ ภิกฺขุ อวิชฺชา, เอตฺตาวตา จ อวิชฺชาคโต โหตีติ.}

\prose{ตสฺมาติห ภิกฺขุ “อิทํ ทุกฺขนฺ”ติ โยโค กรณีโย ฯเปฯ “อยํ ทุกฺขนิโรธคามินี ปฏิปทา”ติ โยโค กรณีโยติ.\csromanspacer{}\csromanspacer{}สตฺตมํ.}

\csromanheader{2}{8. วิชฺชาสุตฺต}
\proseitem{1088}{อถ โข อญฺญตโร ภิกฺขุ เยน ภควา เตนุปสงฺกมิ, อุปสงฺกมิตฺวา ภควนฺตํ อภิวาเทตฺวา เอกมนฺตํ นิสีทิ, เอกมนฺตํ นิสินฺโน โข โส ภิกฺขุ ภควนฺตํ เอตทโวจ “วิชฺชา วิชฺชาติ ภนฺเต วุจฺจติ, กตมา นุ โข ภนฺเต วิชฺชา, กิตฺตาวตา จ วิชฺชาคโต โหตี”ติ.\csromanspacer{} ยํ โข ภิกฺขุ ทุกฺเข ญาณํ, ทุกฺขสมุทเย ญาณํ, ทุกฺขนิโรเธ ญาณํ, ทุกฺขนิโรธคามินิยา ปฏิปทาย ญาณํ.\csromanspacer{} อยํ วุจฺจติ ภิกฺขุ วิชฺชา, เอตฺตาวตา จ วิชฺชาคโต โหตีติ.}

\prose{ตสฺมาติห ภิกฺขุ “อิทํ ทุกฺขนฺ”ติ โยโค กรณีโย ฯเปฯ “อยํ ทุกฺขนิโรธคามินี ปฏิปทา”ติ โยโค กรณีโยติ.\csromanspacer{}\csromanspacer{}อฏฺฐมํ.}

\csromanmatikaanchor{mtk.272}
\csromantocmarkref{subparagraph}{9. สงฺกาสนสุตฺต}{mtk.272}
\bookmark[dest={mtk.272},level=5]{9. สงฺกาสนสุตฺต}
\csromanheader{6}{9. สงฺกาสนสุตฺต}
\proseitem{1089}{“อิทํ ทุกฺขํ อริยสจฺจนฺ”ติ ภิกฺขเว มยา ปญฺญตฺตํ, ตตฺถ อปริมาณา วณฺณา อปริมาณา พฺยญฺชนา อปริมาณา สงฺกาสนา อิติปิทํ ทุกฺขํ}

\csromanheader{2}{12. สจฺจสํยุตฺต}
\csromanmatikaanchor{mtk.273}
\csromanmatikaanchor{mtk.275}
\csromantocmarkref{subparagraph}{10. ตถสุตฺต}{mtk.273}
\bookmark[dest={mtk.273},level=5]{10. ตถสุตฺต}
\prose{\csromanfolio{377}อริยสจฺจนฺติ.\csromanspacer{} อิทํ ทุกฺขสมุทยํ ฯเปฯ.\csromanspacer{} อิทํ ทุกฺขนิโรธํ ฯเปฯ “อิทํ ทุกฺขนิโรธคามินี ปฏิปทา อริยสจฺจนฺ”ติ ภิกฺขเว มยา ปญฺญตฺตํ, ตตฺถ อปริมาณา วณฺณา อปริมาณา พฺยญฺชนา อปริมาณา สงฺกาสนา อิติปิทํ ทุกฺขนิโรธคามินี ปฏิปทา อริยสจฺจนฺติ.}

\prose{ตสฺมาติห ภิกฺขเว “อิทํ ทุกฺขํนฺ”ติ โยโค กรณีโย ฯเปฯ “อยํ ทุกฺขนิโรธคามินี ปฏิปทา”ติ โยโค กรณีโยติ.\csromanspacer{}\csromanspacer{}นวมํ.}

\csromanheader{6}{10. ตถสุตฺต}
\proseitem{1090}{จตฺตาริมานิ ภิกฺขเว ตถานิ อวิตถานิ อนญฺญถานิ.\csromanspacer{} กตมานิ จตฺตาริ, “อิทํ ทุกฺขนฺ”ติ ภิกฺขเว ตถเมตํ อวิตถเมตํ อนญฺญถเมตํ, “อยํ ทุกฺขสมุทโย”ติ ตถเมตํ อวิตถเมตํ อนญฺญถเมตํ, “อยํ ทุกฺขนิโรโธ”ติ ตถเมตํ อวิตถเมตํ อนญฺญถเมตํ, “อยํ ทุกฺขนิโรธคามินี ปฏิปทา”ติ ตถเมตํ อวิตถเมตํ อนญฺญถเมตํ.\csromanspacer{} อิมานิ โข ภิกฺขเว จตฺตาริ ตถานิ อวิตถานิ อนญฺญถานิ.}

\prose{ตสฺมาติห ภิกฺขเว “อิทํ ทุกฺขนฺ”ติ โยโค กรณีโย ฯเปฯ “อยํ ทุกฺขนิโรธคามินี ปฏิปทา”ติ โยโค กรณีโยติ.\csromanspacer{}\csromanspacer{}ทสมํ.\csromanspacer{} ธมฺมจกฺกปฺปวตฺตนวคฺโค ทุติโย.}

\titlehead{\textbf{ตสฺสุทฺทานํ}}
\csromangathameasure{ธมฺมจกฺกํ ตถาคตํ, ขนฺธา อายตเนน จ. \\ ธารณา จ เทฺว อวิชฺชา, วิชฺชา สงฺกาสนา ตถาติ.}
\csromangathagroup{\csromangathastanza{ธมฺมจกฺกํ ตถาคตํ, ขนฺธา อายตเนน จ. \\ ธารณา จ เทฺว อวิชฺชา, วิชฺชา สงฺกาสนา ตถาติ.}}

\csromanmatikaanchor{mtk.274}
\csromantocmarknopage{subparagraph}{3. โกฏิคามวคฺค}{mtk.274}
\bookmark[dest={mtk.274},level=5]{3. โกฏิคามวคฺค}
\csromantocmarkref{subparagraph}{1-2. โกฏิคามสุตฺตานิ}{mtk.275}
\bookmark[dest={mtk.275},level=5]{1-2. โกฏิคามสุตฺตานิ}
\csromanheader{6}{3. โกฏิคามวคฺค}
\csromanheader{2}{1. ปฐมโกฏิคามสุตฺต}
\proseitem{1091}{เอกํ สมยํ ภควา วชฺชีสุ วิหรติ โกฏิคาเม.\csromanspacer{} ตตฺร โข ภควา ภิกฺขู อามนฺเตสิ “จตุนฺนํ ภิกฺขเว อริยสจฺจานํ อนนุโพธา อปฺปฏิเวธา เอวมิทํ ทีฆมทฺธานํ สนฺธาวิตํ สํสริตํ มมญฺเจว ตุมฺหากญฺจ.}

\prose{\csromanfolio{378}กตเมสํ จตุนฺนํ, ทุกฺขสฺส ภิกฺขเว อริยสจฺจสฺส อนนุโพธา อปฺปฏิเวธา เอวมิทํ ทีฆมทฺธานํ สนฺธาวิตํ สํสริตํ มมญฺเจว ตุมฺหากญฺจ.\csromanspacer{} ทุกฺขสมุทยสฺส อริยสจฺจสฺส ฯเปฯ.\csromanspacer{} ทุกฺขนิโรธสฺส อริยสจฺจสฺส ฯเปฯ.\csromanspacer{} ทุกฺขนิโรธคามินิยา ปฏิปทาย อริยสจฺจสฺส อนนุโพธา อปฺปฏิเวธา เอวมิทํ ทีฆมทฺธานํ สนฺธาวิตํ สํสริตํ มมญฺเจว ตุมฺหากญฺจ.\csromanspacer{} ตยิทํ ภิกฺขเว ทุกฺขํ อริยสจฺจํ อนุพุทฺธํ ปฏิวิทฺธํ, ทุกฺขสมุทยํ อริยสจฺจํ อนุพุทฺธํ ปฏิวิทฺธํ, ทุกฺขนิโรธํ อริยสจฺจํ อนุพุทฺธํ ปฏิวิทฺธํ, ทุกฺขนิโรธคามินี ปฏิปทา อริยสจฺจํ อนุพุทฺธํ ปฏิวิทฺธํ.\csromanspacer{} อุจฺฉินฺนา ภวตณฺหา, ขีณา ภวเนตฺติ, นตฺถิทานิ ปุนพฺภโว”ติ.}

\prose{อิทมโวจ ภควา, อิทํ วตฺวาน สุคโต อถาปรํ เอตทโวจ สตฺถา\csromandash{}}

\csromangathameasure{“จตุนฺนํ อริยสจฺจานํ, ยถาภูตํ อทสฺสนา. \\ สํสิตํ ทีฆมทฺธานํ, ตาสุ ตาเสฺวว ชาติสุ. \\ ตานิ เอตานิ ทิฏฺฐานิ, ภวเนตฺติ สมูหตา. \\ อุจฺฉินฺนํ มูลํ ทุกฺขสฺส, นตฺถิทานิ ปุนพฺภโว”ติ. ปฐมํ.}
\csromangathagroup{\csromangathastanza{“จตุนฺนํ อริยสจฺจานํ, ยถาภูตํ อทสฺสนา. \\ สํสิตํ\footnote{สํสริตํ (สฺยา, กํ, ก) ที 2. 77 ปิฏฺเฐปิ.} ทีฆมทฺธานํ, ตาสุ ตาเสฺวว ชาติสุ.}\csromangathastanza{ตานิ\footnote{ยานิ (สฺยา, กํ, อิ, ก)} เอตานิ ทิฏฺฐานิ, ภวเนตฺติ สมูหตา. \\ อุจฺฉินฺนํ มูลํ ทุกฺขสฺส, นตฺถิทานิ ปุนพฺภโว”ติ.\csromanspacer{} ปฐมํ.}}

\csromanheader{2}{2. ทุติยโกฏิคามสุตฺต}
\proseitem{1092}{เย หิ เกจิ ภิกฺขเว สมณา วา พฺราหฺมณา วา “อิทํ ทุกฺขนฺ”ติ ยถาภูตํ นปฺปชานนฺติ, “อยํ ทุกฺขสมุทโย”ติ ยถาภูตํ นปฺปชานนฺติ, “อยํ ทุกฺขนิโรโธ”ติ ยถาภูตํ นปฺปชานนฺติ, “อยํ ทุกฺขนิโรธคามินี ปฏิปทา”ติ ยถาภูตํ นปฺปชานนฺติ, น เม เต ภิกฺขเว สมณา วา พฺราหฺมณา วา สมเณสุ วา สมณสมฺมตา, พฺราหฺมเณสุ วา พฺราหฺมณสมฺมตา, น จ ปน เต อายสฺมนฺโต สามญฺญตฺถํ วา พฺรหฺมญฺญตฺถํ วา ทิฏฺเฐว ธมฺเม สยํ อภิญฺญา สจฺฉิกตฺวา อุปสมฺปชฺช วิหรนฺติ.}

\prose{เย จ โข เกจิ ภิกฺขเว สมณา วา พฺราหฺมณา วา “อิทํ ทุกฺขนฺ”ติ ยถาภูตํ ปชานนฺติ, “อยํ ทุกฺขสมุทโย”ติ ยถาภูตํ ปชานนฺติ, “อยํ ทุกฺขนิโรโธ”ติ ยถาภูตํ ปชานนฺติ, “อยํ ทุกฺขนิโรธคามินี ปฏิปทา”ติ ยถาภูตํ ปชานนฺติ.\csromanspacer{} เต โข เม ภิกฺขเว สมณา วา}

\csromanheader{2}{12. สจฺจสํยุตฺต}
\prose{\csromanfolio{379}พฺราหฺมณา วา สมเณสุ เจว สมณสมฺมตา, พฺราหฺมเณสุ จ พฺราหฺมณสมฺมตา, เต จ ปนายสฺมนฺโต สามญฺญตฺถญฺจ พฺรหฺมญฺญตฺถญฺจ ทิฏฺเฐว ธมฺเม สยํ อภิญฺญา สจฺฉิกตฺวา อุปสมฺปชฺช วิหรนฺตีติ.}

\prose{อิทมโวจ ภควา, อิทํ วตฺวาน สุคโต อถาปรํ เอตทโวจ สตฺถา\csromandash{}}

\csromangathameasure{“เย ทุกฺขํ นปฺปชานนฺติ, อโถ ทุกฺขสฺส สมฺภวํ. \\ ยตฺถ จ สพฺพโส ทุกฺขํ, อเสสํ อุปรุชฺฌติ. \\ ตญฺจ มคฺคํ น ชานนฺติ, ทุกฺขูปสมคามินํ. \\ เจโตวิมุตฺติหีนา เต, อโถ ปญฺญาวิมุตฺติยา. \\ อภพฺพา เต อนฺตกิริยาย, เต เว ชาติชรูปคา. \\ เย จ ทุกฺขํ ปชานนฺติ, อโถ ทุกฺขสฺส สมฺภวํ. \\ ยตฺถ จ สพฺพโส ทุกฺขํ, อเสสํ อุปรุชฺฌติ. \\ ตญฺจ มคฺคํ ปชานนฺติ, ทุกฺขูปสมคามินํ. \\ เจโตวิมุตฺติสมฺปนฺนา, อโถ ปญฺญาวิมุตฺติยา. \\ ภพฺพา เต อนฺตกิริยาย, น เต ชาติชรูปคา”ติ. ทุติยํ.}
\csromangathagroup{\csromangathastanza{“เย ทุกฺขํ นปฺปชานนฺติ, อโถ ทุกฺขสฺส สมฺภวํ. \\ ยตฺถ จ สพฺพโส ทุกฺขํ, อเสสํ อุปรุชฺฌติ.}\csromangathastanza{ตญฺจ มคฺคํ น ชานนฺติ, ทุกฺขูปสมคามินํ. \\ เจโตวิมุตฺติหีนา เต, อโถ ปญฺญาวิมุตฺติยา.}\csromangathastanza{อภพฺพา เต อนฺตกิริยาย, เต เว ชาติชรูปคา. \\ เย จ ทุกฺขํ ปชานนฺติ, อโถ ทุกฺขสฺส สมฺภวํ.}\csromangathastanza{ยตฺถ จ สพฺพโส ทุกฺขํ, อเสสํ อุปรุชฺฌติ. \\ ตญฺจ มคฺคํ ปชานนฺติ, ทุกฺขูปสมคามินํ.}\csromangathastanza{เจโตวิมุตฺติสมฺปนฺนา, อโถ ปญฺญาวิมุตฺติยา. \\ ภพฺพา เต อนฺตกิริยาย, น เต ชาติชรูปคา”ติ.\csromanspacer{} ทุติยํ.}}

\csromanmatikaanchor{mtk.276}
\csromantocmarkref{subparagraph}{3. สมฺมาสมฺพุทฺธสุตฺต}{mtk.276}
\bookmark[dest={mtk.276},level=5]{3. สมฺมาสมฺพุทฺธสุตฺต}
\csromanheader{6}{3. สมฺมาสมฺพุทฺธสุตฺต}
\proseitem{1093}{สาวตฺถินิทานํ.\csromanspacer{} จตฺตาริมานิ ภิกฺขเว อริยสจฺจานิ.\csromanspacer{} กตมานิ จตฺตาริ, ทุกฺขํ อริยสจฺจํ ฯเปฯ.\csromanspacer{} ทุกฺขนิโรธคามินี ปฏิปทา อริยสจฺจํ.\csromanspacer{} อิมานิ โข ภิกฺขเว จตฺตาริ อริยสจฺจานิ.\csromanspacer{} อิเมสํ โข ภิกฺขเว จตุนฺนํ อริยสจฺจานํ ยถาภูตํ อภิสมฺพุทฺธตฺตา ตถาคโต อรหํ “สมฺมาสมฺพุทฺโธ”ติ วุจฺจติ.}

\prose{ตสฺมาติห ภิกฺขเว “อิทํ ทุกฺขนฺ”ติ โยโค กรณีโย ฯเปฯ “อยํ ทุกฺขนิโรธคามินี ปฏิปทา”ติ โยโค กรณีโยติ.\csromanspacer{}\csromanspacer{}ตติยํ.}

\csromanmatikaanchor{mtk.277}
\csromantocmarkref{subparagraph}{4. อรหนฺตสุตฺต}{mtk.277}
\bookmark[dest={mtk.277},level=5]{4. อรหนฺตสุตฺต}
\csromanheader{6}{4. อรหนฺตสุตฺต}
\csromanmatikaanchor{mtk.278}
\csromanmatikaanchor{mtk.279}
\csromantocmarkref{subparagraph}{5. อาสวกฺขยสุตฺต}{mtk.278}
\bookmark[dest={mtk.278},level=5]{5. อาสวกฺขยสุตฺต}
\csromantocmarkref{subparagraph}{6-8. มิตฺตสุตฺตาทีนิ}{mtk.279}
\bookmark[dest={mtk.279},level=5]{6-8. มิตฺตสุตฺตาทีนิ}
\proseitem{1094}{สาวตฺถินิทานํ.\csromanspacer{} เย หิ เกจิ ภิกฺขเว อตีตมทฺธานํ อรหนฺโต สมฺมาสมฺพุทฺธา ยถาภูตํ อภิสมฺพุชฺฌิํสุ, สพฺเพ เต จตฺตาริ อริยสจฺจานิ \csromanfolio{380}ยถาภูตํ อภิสมฺพุชฺฌิํสุ.\csromanspacer{} เย หิ\footnote{เยปิ หิ (พหูสุ)} เกจิ ภิกฺขเว อนาคตมทฺธานํ อรหนฺโต สมฺมาสมฺพุทฺธา ยถาภูตํ อภิสมฺพุชฺฌิสฺสนฺติ, สพฺเพ เต จตฺตาริ อริยสจฺจานิ ยถาภูตํ อภิสมฺพุชฺฌิสฺสนฺติ.\csromanspacer{} เย หิ เกจิ ภิกฺขเว เอตรหิ อรหนฺโต สมฺมาสมฺพุทฺธา ยถาภูตํ อภิสมฺพุชฺฌนฺติ, สพฺเพ เต จตฺตาริ อริยสจฺจานิ ยถาภูตํ อภิสมฺพุชฺฌนฺติ.}

\prose{กตมานิ จตฺตาริ, ทุกฺขํ อริยสจฺจํ ทุกฺขสมุทยํ อริยสจฺจํ ทุกฺขนิโรธํ อริยสจฺจํ ทุกฺขนิโรธคามินี ปฏิปทา อริยสจฺจํ.\csromanspacer{} เย หิ เกจิ ภิกฺขเว อตีตมทฺธานํ อรหนฺโต สมฺมาสมฺพุทฺธา ยถาภูตํ อภิสมฺพุชฺฌิํสุ ฯเปฯ อภิสมฺพุชฺฌิสฺสนฺติ.\csromanspacer{} อภิสมฺพุชฺฌนฺติ, สพฺเพ เต อิมานิ จตฺตาริ อริยสจฺจานิ ยถาภูตํ อภิสมฺพุชฺฌนฺติ.}

\prose{ตสฺมาติห ภิกฺขเว “อิทํ ทุกฺขนฺ”ติ โยโค กรณีโย ฯเปฯ “อยํ ทุกฺขนิโรธคามินี ปฏิปทา”ติ โยโค กรณีโยติ.\csromanspacer{}\csromanspacer{}จตุตฺถํ.}

\csromanheader{6}{5. อาสวกฺขยสุตฺต}
\proseitem{1095}{ชานโตหํ ภิกฺขเว ปสฺสโต อาสวานํ ขยํ วทามิ, โน อชานโต อปสฺสโต.\csromanspacer{} กิญฺจ ภิกฺขเว ชานโต ปสฺสโต อาสวานํ ขโย โหติ, “อิทํ ทุกฺขนฺ”ติ ภิกฺขเว ชานโต ปสฺสโต อาสวานํ ขโย โหติ, “อยํ ทุกฺขสมุทโย”ติ ชานโต ปสฺสโต อาสวานํ ขโย โหติ, “อยํ ทุกฺขนิโรโธ”ติ ชานโต ปสฺสโต อาสวานํ ขโย โหติ, “อยํ ทุกฺขนิโรธคามินี ปฏิปทา”ติ ชานโต ปสฺสโต อาสวานํ ขโย โหติ.\csromanspacer{} เอวํ โข ภิกฺขเว ชานโต เอวํ ปสฺสโต อาสวานํ ขโย โหติ.}

\prose{ตสฺมาติห ภิกฺขเว “อิทํ ทุกฺขนฺ”ติ โยโค กรณีโย ฯเปฯ “อยํ ทุกฺขนิโรธคามินี ปฏิปทา”ติ โยโค กรณีโยติ.\csromanspacer{}\csromanspacer{}ปญฺจมํ.}

\csromanheader{2}{6. มิตฺตสุตฺต}
\proseitem{1096}{เย หิ เกจิ ภิกฺขเว อนุกมฺเปยฺยาถ, เย จ โสตพฺพํ มญฺเญยฺยุํ มิตฺตา วา อมจฺจา วา ญาตี วา สาโลหิตา วา, เต โว ภิกฺขเว จตุนฺนํ อริยสจฺจานํ ยถาภูตํ อภิสมยาย สมาทเปตพฺพา นิเวเสตพฺพา ปติฏฺฐาเปตพฺพา.}

\csromanheader{2}{12. สจฺจสํยุตฺต}
\prose{\csromanfolio{381}กตเมสํ จตุนฺนํ, ทุกฺขสฺส อริยสจฺจสฺส ทุกฺขสมุทยสฺส อริยสจฺจสฺส ทุกฺขนิโรธสฺส อริยสจฺจสฺส ทุกฺขนิโรธคามินิยา ปฏิปทาย อริยสจฺจสฺส.\csromanspacer{} เย หิ เกจิ ภิกฺขเว อนุกมฺเปยฺยาถ, เย จ โสตพฺพํ มญฺเญยฺยุํ มิตฺตา วา อมจฺจา วา ญาตี วา สาโลหิตา วา, เต โว ภิกฺขเว อิเมสํ จตุนฺนํ อริยสจฺจานํ ยถาภูตํ อภิสมยาย สมาทเปตพฺพา นิเวเสตพฺพา ปติฏฺฐาเปตพฺพา.}

\prose{ตสฺมาติห ภิกฺขเว “อิทํ ทุกฺขนฺ”ติ โยโค กรณีโย ฯเปฯ “อยํ ทุกฺขนิโรธคามินี ปฏิปาทา”ติ โยโค กรณีโยติ.\csromanspacer{}\csromanspacer{}ฉฏฺฐํ.}

\csromanheader{2}{7. ตถสุตฺต}
\proseitem{1097}{จตฺตาริมานิ ภิกฺขเว อริยสจฺจานิ.\csromanspacer{} กตมานิ จตฺตาริ, ทุกฺขํ อริยสจฺจํ ทุกฺขสมุทยํ อริยสจฺจํ ทุกฺขนิโรธํ อริยสจฺจํ ทุกฺขนิโรธคามินี ปฏิปทา อริยสจฺจํ.\csromanspacer{} อิมานิ โข ภิกฺขเว จตฺตาริ อริยสจฺจานิ ตถานิ อวิตถานิ อนญฺญถานิ, ตสฺมา “อริยสจฺจานี”ติ วุจฺจนฺติ.}

\prose{ตสฺมาติห ภิกฺขเว “อิทํ ทุกฺขนฺ”ติ โยโค กรณีโย ฯเปฯ “อยํ ทุกฺขนิโรธคามินี ปฏิปทา”ติ โยโค กรณีโยติ.\csromanspacer{}\csromanspacer{}สตฺตมํ.}

\csromanheader{2}{8. โลกสุตฺต}
\proseitem{1098}{จตฺตาริมานิ ภิกฺขเว อริยสจฺจานิ.\csromanspacer{} กตมานิ จตฺตาริ, ทุกฺขํ อริยสจฺจํ ทุกฺขสมุทยํ อริยสจฺจํ ทุกฺขนิโรธํ อริยสจฺจํ ทุกฺขนิโรธคามินี ปฏิปทา อริยสจฺจํ.\csromanspacer{} สเทวเก โลเก สมารเก สพฺรหฺมเก สสฺสมณพฺราหฺมณิยา ปชาย สเทวมนุสฺสาย ตถาคโต อริโย, ตสฺมา “อริยสจฺจานี”ติ วุจฺจนฺติ.}

\prose{ตสฺมาติห ภิกฺขเว “อิทํ ทุกฺขนฺ”ติ โยโค กรณีโย ฯเปฯ “อยํ ทุกฺขนิโรธคามินี ปฏิปทา”ติ โยโค กรณีโยติ.\csromanspacer{}\csromanspacer{}อฏฺฐมํ.}

\csromanmatikaanchor{mtk.280}
\csromantocmarkref{subparagraph}{9. ปริญฺเญยฺยสุตฺต}{mtk.280}
\bookmark[dest={mtk.280},level=5]{9. ปริญฺเญยฺยสุตฺต}
\csromanheader{6}{9. ปริญฺเญยฺยสุตฺต}
\proseitem{1099}{จตฺตาริมานิ ภิกฺขเว อริยสจฺจานิ.\csromanspacer{} กตมานิ จตฺตาริ, ทุกฺขํ อริยสจฺจํ ทุกฺขสมุทยํ อริยสจฺจํ ทุกฺขนิโรธํ อริยสจฺจํ ทุกฺขนิโรธคามินี ปฏิปาทา อริยสจฺจํ.\csromanspacer{} อิมานิ โข ภิกฺขเว จตฺตาริ อริยสจฺจานิ. \csromanfolio{382}อิเมสํ โข ภิกฺขเว จตุนฺนํ อริยสจฺจานํ อตฺถิ อริยสจฺจํ ปริญฺเญยฺยํ, อตฺถิ อริยสจฺจํ ปหาตพฺพํ, อตฺถิ อริยสจฺจํ สจฺฉิกาตพฺพํ, อตฺถิ อริยสจฺจํ ภาเวตพฺพํ.}

\prose{กตมญฺจ ภิกฺขเว อริยสจฺจํ ปริญฺเญยฺยํ, ทุกฺขํ ภิกฺขเว อริยสจฺจํ ปริญฺเญยฺยํ, ทุกฺขสมุทยํ อริยสจฺจํ ปหาตพฺพํ, ทุกฺขนิโรธํ อริยสจฺจํ สจฺฉิกาตพฺพํ, ทุกฺขนิโรธคามินี ปฏิปทา อริยสจฺจํ ภาเวตพฺพํ.}

\prose{ตสฺมาติห ภิกฺขเว “อิทํ ทุกฺขนฺ”ติ โยโค กรณีโย ฯเปฯ “อยํ ทุกฺขนิโรธคามินี ปฏิปทา”ติ โยโค กรณีโยติ.\csromanspacer{}\csromanspacer{}นวมํ.}

\csromanmatikaanchor{mtk.281}
\csromantocmarkref{subparagraph}{10. ควมฺปติสุตฺต}{mtk.281}
\bookmark[dest={mtk.281},level=5]{10. ควมฺปติสุตฺต}
\csromanheader{6}{10. ควมฺปติสุตฺต}
\proseitem{1100}{เอกํ สมยํ สมฺพหุลา เถรา ภิกฺขู เจเตสุ\footnote{เจติเยสุ (สฺยา)} วิหรนฺติ สหญฺจนิเก\footnote{สหชนิเย (สี, สฺยา, กํ)}.\csromanspacer{} เตน โข ปน สมเยน สมฺพหุลานํ เถรานํ ภิกฺขูนํ ปจฺฉาภตฺตํ ปิณฺฑปาตปฏิกฺกนฺตานํ มณฺฑลมาเฬ สนฺนิสินฺนานํ สนฺนิปติตานํ อยมนฺตรากถา อุทปาทิ “โย นุ โข อาวุโส ทุกฺขํ ปสฺสติ, ทุกฺขสมุทยมฺปิ โส ปสฺสติ, ทุกฺขนิโรธมฺปิ ปสฺสติ, ทุกฺขนิโรธคามินิํ ปฏิปทมฺปิ ปสฺสตี”ติ.}

\prose{เอวํ วุตฺเต อายสฺมา ควมฺปติ เถโร\footnote{ควมฺปติตฺเถโร (สฺยา, กํ)} ภิกฺขู เอตทโวจ “สมฺมุขา เมตํ อาวุโส ภควโต สุตํ, สมฺมุขา ปฏิคฺคหิตํ โย ภิกฺขเว ทุกฺขํ ปสฺสติ, ทุกฺขสมุทยมฺปิ โส ปสฺสติ, ทุกฺขนิโรธมฺปิ ปสฺสติ, ทุกฺขนิโรธคามินิํ ปฏิปทมฺปิ ปสฺสติ.\csromanspacer{} โย ทุกฺขสมุทยํ ปสฺสติ, ทุกฺขมฺปิ โส ปสฺสติ, ทุกฺขนิโรธมฺปิ ปสฺสติ, ทุกฺขนิโรธคามินิํ ปฏิปทมฺปิ ปสฺสติ.\csromanspacer{} โย ทุกฺขนิโรธํ ปสฺสติ, ทุกฺขมฺปิ โส ปสฺสติ, ทุกฺขสมุทยมฺปิ ปสฺสติ, ทุกฺขนิโรธคามินิํ ปฏิปทมฺปิ ปสฺสติ.\csromanspacer{} โย ทุกฺขนิโรธคามินิํ ปฏิปทํ ปสฺสติ, ทุกฺขมฺปิ โส ปสฺสติ, ทุกฺขสมุทยมฺปิ ปสฺสติ, ทุกฺขนิโรธมฺปิ ปสฺสตี”ติ.\csromanspacer{}\csromanspacer{}ทสมํ.\csromanspacer{} โกฏิคามวคฺโค ตติโย.}

\titlehead{\textbf{ตสฺสุทฺทานํ}}
\csromangathameasure{เทฺว วชฺชี สมฺมาสมฺพุทฺโธ, อรหํ อาสวกฺขโย. \\ มิตฺตํ ตถา จ โลโก จ, ปริญฺเญยฺยํ ควมฺปตีติ.}
\csromangathagroup{\csromangathastanza{เทฺว วชฺชี สมฺมาสมฺพุทฺโธ, อรหํ อาสวกฺขโย. \\ มิตฺตํ ตถา จ โลโก จ, ปริญฺเญยฺยํ ควมฺปตีติ.}}

\csromanheader{2}{12. สจฺจสํยุตฺต}
\csromanmatikaanchor{mtk.282}
\csromantocmarknopage{subparagraph}{4. สีสปาวนวคฺค}{mtk.282}
\bookmark[dest={mtk.282},level=5]{4. สีสปาวนวคฺค}
\csromanheader{6}{4. สีสปาวนวคฺค}
\csromanmatikaanchor{mtk.283}
\csromantocmarkref{subparagraph}{1. สีสปาวนสุตฺต}{mtk.283}
\bookmark[dest={mtk.283},level=5]{1. สีสปาวนสุตฺต}
\csromanheader{6}{1. สีสปาวนสุตฺต}
\proseitem{1101}{\csromanfolio{383}เอกํ สมยํ ภควา โกสมฺพิยํ วิหรติ สีสปาวเน\footnote{สิํสปาวเน (สี, อิ)}.\csromanspacer{} อถ โข ภควา ปริตฺตานิ สีสปาปณฺณานิ ปาณินา คเหตฺวา ภิกฺขู อามนฺเตสิ “ตํ กิํ มญฺญถ ภิกฺขเว, กตมํ นุ โข พหุตรํ ยานิ วา มยา ปริตฺตานิ สีสปาปณฺณานิ ปาณินา คหิตานิ ยทิทํ อุปริ สีสปาวเน”ติ.\csromanspacer{} อปฺปมตฺตกานิ ภนฺเต ภควตา ปริตฺตานิ สีสปาปณฺณานิ ปาณินา คหิตานิ, อถ โข เอตาเนว พหุตรานิ ยทิทํ อุปริ สีสปาวเนติ.\csromanspacer{} เอวเมว โข ภิกฺขเว เอตเทว พหุตรํ ยํ โว มยา อภิญฺญาย อนกฺขาตํ.\csromanspacer{} กสฺมา เจตํ ภิกฺขเว มยา อนกฺขาตํ, น เหตํ ภิกฺขเว อตฺถสํหิตํ นาทิพฺรหฺมจริยกํ น นิพฺพิทาย น วิราคาย น นิโรธาย น อุปสมาย น อภิญฺญาย น สมฺโพธาย น นิพฺพานาย สํวตฺตติ, ตสฺมา ตํ มยา อนกฺขาตํ.}

\prose{กิญฺจ ภิกฺขเว มยา อกฺขาตํ, “อิทํ ทุกฺขนฺ”ติ ภิกฺขเว มยา อกฺขาตํ, “อยํ ทุกฺขสมุทโย”ติ มยา อกฺขาตํ, “อยํ ทุกฺขนิโรโธ”ติ มยา อกฺขาตํ, “อยํ ทุกฺขนิโรธคามินี ปฏิปทา”ติ มยา อกฺขาตํ.}

\prose{กสฺมา เจตํ ภิกฺขเว มยา อกฺขาตํ, เอตํ หิ ภิกฺขเว อตฺถสํหิตํ เอตํ อาทิพฺรหฺมจริยกํ เอตํ นิพฺพิทาย วิราคาย นิโรธาย อุปสมาย อภิญฺญาย สมฺโพธาย นิพฺพานาย สํวตฺตติ, ตสฺมา ตํ มยา อกฺขาตํ.}

\prose{ตสฺมาติห ภิกฺขเว “อิทํ ทุกฺขนฺ”ติ โยโค กรณีโย ฯเปฯ “อยํ ทุกฺขนิโรธคามินี ปฏิปทา”ติ โยโค กรณีโยติ.\csromanspacer{}\csromanspacer{}ปฐมํ.}

\csromanmatikaanchor{mtk.284}
\csromantocmarkref{subparagraph}{2. ขทิรปตฺตสุตฺต}{mtk.284}
\bookmark[dest={mtk.284},level=5]{2. ขทิรปตฺตสุตฺต}
\csromanheader{6}{2. ขทิรปตฺตสุตฺต}
\proseitem{1102}{โย ภิกฺขเว เอวํ วเทยฺย “อหํ ทุกฺขํ อริยสจฺจํ ยถาภูตํ อนภิสเมจฺจ ทุกฺขสมุทยํ อริยสจฺจํ ยถาภูตํ อนภิสเมจฺจ ทุกฺขนิโรธํ อริยสจฺจํ ยถาภูตํ อนภิสเมจฺจ ทุกฺขนิโรธคามินิํ ปฏิปทํ อริยสจฺจํ ยถาภูตํ อนภิสเมจฺจ สมฺมา ทุกฺขสฺสนฺตํ กริสฺสามี”ติ เนตํ ฐานํ วิชฺชติ.}

\prose{\csromanfolio{384}เสยฺยถาปิ ภิกฺขเว โย เอวํ วเทยฺย “อหํ ขทิรปตฺตานํ วา สรลปตฺตานํ\footnote{ปลาสปตฺตานํ (สี, สฺยา, กํ, อิ)} วา อามลกปตฺตานํ วา ปุฏํ กริตฺวา อุทกํ วา ตาลปตฺตํ วา อาหริสฺสามี”ติ เนตํ ฐานํ วิชฺชติ.\csromanspacer{} เอวเมว โข ภิกฺขเว โย เอวํ วเทยฺย “อหํ ทุกฺขํ อริยสจฺจํ ยถาภูตํ อนภิสเมจฺจ ฯเปฯ ทุกฺขนิโรธคามินิํ ปฏิปทํ อริยสจฺจํ ยถาภูตํ อนภิสเมจฺจ สมฺมา ทุกฺขสฺสนฺตํ กริสฺสามี”ติ เนตํ ฐานํ วิชฺชติ.}

\prose{โย จ โข ภิกฺขเว เอวํ วเทยฺย “อหํ ทุกฺขํ อริยสจฺจํ ยถาภูตํ อภิสเมจฺจ ทุกฺขสมุทยํ อริยสจฺจํ ยถาภูตํ อภิสเมจฺจ ทุกฺขนิโรธํ อริยสจฺจํ ยถาภูตํ อภิสเมจฺจ ทุกฺขนิโรธคามินิํ ปฏิปทํ อริยสจฺจํ ยถาภูตํ อภิสเมจฺจ สมฺมา ทุกฺขสฺสนฺตํ กริสฺสามี”ติ ฐานเมตํ วิชฺชติ.}

\prose{เสยฺยถาปิ ภิกฺขเว โย เอวํ วเทยฺย “อหํ ปทุมปตฺตานํ วา ปลาสปตฺตานํ วา มาลุวปตฺตานํ วา ปุฏํ กริตฺวา อุทกํ วา ตาลปตฺตํ วา อาหริสฺสามี”ติ ฐานเมตํ วิชฺชติ.\csromanspacer{} เอวเมว โข ภิกฺขเว โย เอวํ วเทยฺย “อหํ ทุกฺขํ อริยสจฺจํ ยถาภูตํ อภิสเมจฺจ ฯเปฯ ทุกฺขนิโรธคามินิํ ปฏิปทํ อริยสจฺจํ ยถาภูตํ อภิสเมจฺจ สมฺมา ทุกฺขสฺสนฺตํ กริสฺสามี”ติ ฐานเมตํ วิชฺชติ.}

\prose{ตสฺมาติห ภิกฺขเว “อิทํ ทุกฺขนฺ”ติ โยโค กรณีโย ฯเปฯ “อยํ ทุกฺขนิโรธคามินี ปฏิปทา”ติ โยโค กรณีโยติ.\csromanspacer{}\csromanspacer{}ทุติยํ.}

\csromanmatikaanchor{mtk.285}
\csromantocmarkref{subparagraph}{3. ทณฺฑสุตฺต}{mtk.285}
\bookmark[dest={mtk.285},level=5]{3. ทณฺฑสุตฺต}
\csromanheader{6}{3. ทณฺฑสุตฺต}
\proseitem{1103}{เสยฺยถาปิ ภิกฺขเว ทณฺโฑ อุปริเวหาสํ ขิตฺโต สกิมฺปิ มูเลน นิปตติ, สกิมฺปิ อคฺเคน นิปตติ.\csromanspacer{} เอวเมว โข ภิกฺขเว อวิชฺชานีวรณา สตฺตา ตณฺหาสํโยชนา สนฺธาวนฺตา สํสรนฺตา\footnote{ตณฺหาสํโยชนพนฺธา สนฺธาวตา (ก)} สกิมฺปิ อสฺมา โลกา ปรํ โลกํ คจฺฉนฺติ, สกิมฺปิ ปรสฺมา โลกา อิมํ โลกํ อาคจฺฉนฺติ.\csromanspacer{} ตํ กิสฺส เหตุ, อทิฏฺฐตฺตา ภิกฺขเว จตุนฺนํ อริยสจฺจานํ.\csromanspacer{} กตเมสํ จตุนฺนํ, ทุกฺขสฺส อริยสจฺจสฺส ฯเปฯ ทุกฺขนิโรธคามินิยา ปฏิปทาย อริยสจฺจสฺส.}

\prose{ตสฺมาติห ภิกฺขเว “อิทํ ทุกฺขนฺ”ติ โยโค กรณีโย ฯเปฯ “อยํ ทุกฺขนิโรธคามินี ปฏิปทา”ติ โยโค กรณีโยติ.\csromanspacer{}\csromanspacer{}ตติยํ.}

\csromanheader{2}{12. สจฺจสํยุตฺต}
\csromanmatikaanchor{mtk.286}
\csromantocmarkref{subparagraph}{4. เจลสุตฺต}{mtk.286}
\bookmark[dest={mtk.286},level=5]{4. เจลสุตฺต}
\csromanheader{6}{4. เจลสุตฺต}
\proseitem{1104}{\csromanfolio{385}อาทิตฺเต ภิกฺขเว เจเล วา สีเส วา กิมสฺส กรณียนฺติ.\csromanspacer{} อาทิตฺเต ภนฺเต เจเล วา สีเส วา ตสฺเสว เจลสฺส วา สีสสฺส วา นิพฺพาปนาย อธิมตฺโต ฉนฺโท จ วายาโม จ อุสฺสาโห จ อุสฺโสฬฺหี จ อปฺปฏิวานี จ สติ จ สมฺปชญฺญญฺจ กรณียนฺติ.}

\prose{อาทิตฺตํ ภิกฺขเว เจลํ วา สีสํ วา อชฺฌุเปกฺขิตฺวา อมนสิกริตฺวา อนภิสเมตานํ จตุนฺนํ อริยสจฺจานํ ยถาภูตํ อภิสมยาย อธิมตฺโต ฉนฺโท จ วายาโม จ อุสฺสาโห จ อุสฺโสฬฺหี จ อปฺปฏิวานี จ สติ จ สมฺปชญฺญญฺจ กรณียํ.\csromanspacer{} กตเมสํ จตุนฺนํ, ทุกฺขสฺส อริยสจฺจสฺส ฯเปฯ ทุกฺขนิโรธคามินิยา ปฏิปทาย อริยสจฺจสฺส.}

\prose{ตสฺมาติห ภิกฺขเว “อิทํ ทุกฺขนฺ”ติ โยโค กรณีโย ฯเปฯ “อยํ ทุกฺขนิโรธคามินี ปฏิปทา”ติ โยโค กรณีโยติ.\csromanspacer{}\csromanspacer{}จตุตฺถํ.}

\csromanmatikaanchor{mtk.287}
\csromantocmarkref{subparagraph}{5. สตฺติสตสุตฺต}{mtk.287}
\bookmark[dest={mtk.287},level=5]{5. สตฺติสตสุตฺต}
\csromanheader{6}{5. สตฺติสตสุตฺต}
\proseitem{1105}{เสยฺยถาปิ ภิกฺขเว ปุริโส วสฺสสตายุโก วสฺสสตชีวี, ตเมนํ เอวํ วเทยฺย “เอหมฺโภ ปุริส, ปุพฺพณฺหสมยํ ตํ สตฺติสเตน หนิสฺสนฺติ, มชฺฌนฺหิกสมยํ สตฺติสเตน หนิสฺสนฺติ, สายนฺหสมยํ สตฺติสเตน หนิสฺสนฺติ.\csromanspacer{} โส โข ตฺวํ อมฺโภ ปุริส ทิวเส ทิวเส ตีหิ ตีหิ สตฺติสเตหิ หญฺญมาโน วสฺสสตายุโก วสฺสสตชีวี วสฺสสตสฺส อจฺจเยน อนภิสเมตานิ จตฺตาริ อริยสจฺจานิ อภิสเมสฺสสี”ติ.}

\prose{อตฺถวสิเกน ภิกฺขเว กุลปุตฺเตน อลํ อุปคนฺตุํ.\csromanspacer{} ตํ กิสฺส เหตุ, อนมตคฺโคยํ ภิกฺขเว สํสาโร, ปุพฺพา โกฏิ นปฺปญฺญายติ สตฺติปฺปหารานํ อสิปฺปหารานํ อุสุปฺปหารานํ ผรสุปฺปหารานํ\footnote{อสิปฺปหารานํ ผรสุปฺปหารานํ (ก)}.\csromanspacer{} เอวญฺเจตํ ภิกฺขเว อสฺส.\csromanspacer{} น โข ปนาหํ ภิกฺขเว สห ทุกฺเขน สห โทมนสฺเสน จตุนฺนํ อริยสจฺจานํ อภิสมยํ วทามิ, อปิ จาหํ ภิกฺขเว สหาว สุเขน สหาว โสมนสฺเสน จตุนฺนํ อริยสจฺจานํ อภิสมยํ วทามิ.\csromanspacer{} กตเมสํ จตุนฺนํ, ทุกฺขสฺส อริยสจฺจสฺส ฯเปฯ ทุกฺขนิโรธคามินิยา ปฏิปทาย อริยสจฺจสฺส.}

\csromanmatikaanchor{mtk.288}
\csromanmatikaanchor{mtk.289}
\csromantocmarkref{subparagraph}{6. ปณสุตฺต}{mtk.288}
\bookmark[dest={mtk.288},level=5]{6. ปณสุตฺต}
\csromantocmarkref{subparagraph}{7-8. สูริยสุตฺตานิ}{mtk.289}
\bookmark[dest={mtk.289},level=5]{7-8. สูริยสุตฺตานิ}
\prose{\csromanfolio{386}ตสฺมาติห ภิกฺขเว “อิทํ ทุกฺขนฺ”ติ โยโค กรณีโย ฯเปฯ “อยํ ทุกฺขนิโรธคามินี ปฏิปทา”ติ โยโค กรณีโยติ.\csromanspacer{}\csromanspacer{}ปญฺจมํ.}

\csromanheader{2}{6. ปาณสุตฺต}
\proseitem{1106}{เสยฺยถาปิ ภิกฺขเว ปุริโส ยํ อิมสฺมิํ ชมฺพุทีเป ติณกฏฺฐสาขาปลาสํ ตจฺเฉตฺวา เอกชฺฌํ สํหเรยฺย, เอกชฺฌํ สํหริตฺวา สูลํ กเรยฺย, สูลํ กริตฺวา เย มหาสมุทฺเท มหนฺตกา ปาณา, เต มหนฺตเกสุ สูเลสุ อาวุเนยฺย.\csromanspacer{} เย มหาสมุทฺเท มชฺฌิมกา ปาณา, เต มชฺฌิมเกสุ สูเลสุ อาวุเนยฺย.\csromanspacer{} เย มหาสมุทฺเท สุขุมกา ปาณา, เต สุขุมเกสุ สูเลสุ อาวุเนยฺย, อปริยาทินฺนา จ ภิกฺขเว มหาสมุทฺเท โอฬาริกา ปาณา อสฺสุ.}

\prose{อถ อิมสฺมิํ ชมฺพุทีเป ติณกฏฺฐสาขาปลาสํ ปริกฺขยํ ปริยาทานํ คจฺเฉยฺย.\csromanspacer{} อิโต พหุตรา โข ภิกฺขเว มหาสมุทฺเท สุขุมกา ปาณา, เย น สุกรา สูเลสุ อาวุนิตุํ.\csromanspacer{} ตํ กิสฺส เหตุ, สุขุมตฺตา ภิกฺขเว อตฺตภาวสฺส.\csromanspacer{} เอวํ มหา โข ภิกฺขเว อปาโย, เอวํ มหนฺตสฺมา โข ภิกฺขเว อปายสฺมา ปริมุตฺโต ทิฏฺฐิสมฺปนฺโน ปุคฺคโล “อิทํ ทุกฺขนฺ”ติ ยถาภูตํ ปชานาติ ฯเปฯ “อยํ ทุกฺขนิโรธคามินี ปฏิปทา”ติ ยถาภูตํ ปชานาติ.}

\prose{ตสฺมาติห ภิกฺขเว “อิทํ ทุกฺขนฺ”ติ โยโค กรณีโย ฯเปฯ “อยํ ทุกฺขนิโรธคามินี ปฏิปทา”ติ โยโค กรณีโยติ.\csromanspacer{}\csromanspacer{}ฉฏฺฐํ.}

\csromanheader{2}{7. ปฐมสูริยสุตฺต}
\proseitem{1107}{สูริยสฺส\footnote{สุริยสฺส (สี, สฺยา, กํ, อิ)} ภิกฺขเว อุทยโต เอตํ ปุพฺพงฺคมํ เอตํ ปุพฺพนิมิตฺตํ, ยทิทํ อรุณุคฺคํ.\csromanspacer{} เอวเมว โข ภิกฺขเว ภิกฺขุโน จตุนฺนํ อริยสจฺจานํ ยถาภูตํ อภิสมยาย เอตํ ปุพฺพงฺคมํ เอตํ ปุพฺพนิมิตฺตํ, ยทิทํ สมฺมาทิฏฺฐิ.\csromanspacer{} ตสฺเสตํ ภิกฺขเว ภิกฺขุโน ปาฏิกงฺขํ “อิทํ ทุกฺขนฺ”ติ ยถาภูตํ ปชานิสฺสติ ฯเปฯ “อยํ ทุกฺขนิโรธคามินี ปฏิปทา”ติ ยถาภูตํ ปชานิสฺสติ.}

\prose{ตสฺมาติห ภิกฺขเว “อิทํ ทุกฺขนฺ”ติ โยโค กรณีโย ฯเปฯ “อยํ ทุกฺขนิโรธคามินี ปฏิปทา”ติ โยโค กรณีโยติ.\csromanspacer{}\csromanspacer{}สตฺตมํ.}

\csromanheader{2}{12. สจฺจสํยุตฺต}
\csromanheader{2}{8. ทุติยสูริยสุตฺต}
\proseitem{1108}{\csromanfolio{387}ยาวกีวญฺจ ภิกฺขเว จนฺทิมสูริยา โลเก นุปฺปชฺชนฺติ, เนว ตาว มหโต อาโลกสฺส ปาตุภาโว โหติ มหโต โอภาสสฺส, อนฺธตมํ ตทา โหติ อนฺธการติมิสา, เนว ตาว รตฺตินฺทิวา\footnote{รตฺติทิวา (ก)} ปญฺญายนฺติ, น มาสทฺธมาสา ปญฺญายนฺติ, น อุตุสํวจฺฉรา ปญฺญายนฺติ.}

\prose{ยโต จ โข ภิกฺขเว จนฺทิมสูริยา โลเก อุปฺปชฺชนฺติ, อถ มหโต อาโลกสฺส ปาตุภาโว โหติ มหโต โอภาสสฺส, เนว อนฺธการตมํ ตทา โหติ น อนฺธการติมิสา, อถ รตฺตินฺทิวา ปญฺญายนฺติ, มาสทฺธมาสา ปญฺญายนฺติ, อุตุสํวจฺฉรา ปญฺญายนฺติ.\csromanspacer{} เอวเมว โข ภิกฺขเว ยาวกีวญฺจ ตถาคโต โลเก นุปฺปชฺชติ อรหํ สมฺมาสมฺพุทฺโธ, เนว ตาว มหโต อาโลกสฺส ปาตุภาโว โหติ มหโต โอภาสสฺส, อนฺธตมํ ตทา โหติ อนฺธการติมิสา, เนว ตาว จตุนฺนํ อริยสจฺจานํ อาจิกฺขณา โหติ เทสนา ปญฺญาปนา ปฏฺฐปนา วิวรณา วิภชนา อุตฺตานีกมฺมํ.}

\prose{ยโต จ โข ภิกฺขเว ตถาคโต โลเก อุปฺปชฺชติ อรหํ สมฺมาสมฺพุทฺโธ, อถ มหโต อาโลกสฺส ปาตุภาโว โหติ มหโต โอภาสสฺส, เนว อนฺธตมํ ตทา โหติ น อนฺธการติมิสา, อถ โข จตุนฺนํ อริยสจฺจานํ อาจิกฺขณา โหติ เทสนา ปญฺญาปนา ปฏฺฐปนา วิวรณา วิภชนา อุตฺตานีกมฺมํ.\csromanspacer{} กตเมสํ จตุนฺนํ, ทุกฺขสฺส อริยสจฺจสฺส ฯเปฯ ทุกฺขนิโรธคามินิยา ปฏิปทาย อริยสจฺจสฺส.}

\prose{ตสฺมาติห ภิกฺขเว “อิทํ ทุกฺขนฺ”ติ โยโค กรณีโย ฯเปฯ “อยํ ทุกฺขนิโรธคามินี ปฏิปทา”ติ โยโค กรณีโยติ.\csromanspacer{}\csromanspacer{}ฉฏฺฐํ.}

\csromanmatikaanchor{mtk.290}
\csromantocmarkref{subparagraph}{9. อินฺทขีลสุตฺต}{mtk.290}
\bookmark[dest={mtk.290},level=5]{9. อินฺทขีลสุตฺต}
\csromanheader{6}{9. อินฺทขีลสุตฺต}
\proseitem{1109}{เย หิ เกจิ ภิกฺขเว สมณา วา พฺราหฺมณา วา “อิทํ ทุกฺขนฺ”ติ ยถาภูตํ นปฺปชานนฺติ ฯเปฯ “อยํ ทุกฺขนิโรธคามินี ปฏิปทา”ติ ยถาภูตํ นปฺปชานนฺติ, เต อญฺญสฺส สมณสฺส วา พฺราหฺมณสฺส วา มุขํ อุลฺโลเกนฺติ\footnote{โอโลเกนฺติ (สี, สฺยา)} “อยํ นูน ภวํ ชานํ ชานาติ, ปสฺสํ ปสฺสตี”ติ.}

\prose{\csromanfolio{388}เสยฺยถาปิ ภิกฺขเว ตูลปิจุ วา กปฺปาสปิจุ วา ลหุโก วาตูปาทาโน สเม ภูมิภาเค นิกฺขิตฺโต, ตเมนํ ปุรตฺถิโม วาโต ปจฺฉิเมน สํหเรยฺย, ปจฺฉิโม วาโต ปุรตฺถิเมน สํหเรยฺย, อุตฺตโร วาโต ทกฺขิเณน สํหเรยฺย, ทกฺขิโณ วาโต อุตฺตเรน สํหเรยฺย.\csromanspacer{} ตํ กิสฺส เหตุ, ลหุกตฺตา ภิกฺขเว กปฺปาสปิจุโน.\csromanspacer{} เอวเมว โข ภิกฺขเว เย หิ เกจิ สมณา วา พฺราหฺมณา วา “อิทํ ทุกฺขนฺ”ติ ยถาภูตํ นปฺปชานนฺติ ฯเปฯ “อยํ ทุกฺขนิโรธคามินี ปฏิปทา”ติ ยถาภูตํ นปฺปชานนฺติ, เต อญฺญสฺส สมณสฺส วา พฺราหฺมณสฺส วา มุขํ อุลฺโลเกนฺติ “อยํ นุน ภวํ ชานํ ชานาติ, ปสฺสํ ปสฺสตี”ติ.\csromanspacer{} ตํ กิสฺส เหตุ, อทิฏฺฐตฺตา ภิกฺขเว จตุนฺนํ อริยสจฺจานํ.}

\prose{เย จ โข เกจิ ภิกฺขเว สมณา วา พฺราหฺมณา วา “อิทํ ทุกฺขํนฺ”ติ ยถาภูตํ ปชานนฺติ ฯเปฯ “อยํ ทุกฺขนิโรธคามินี ปฏิปทา”ติ ยถาภูตํ ปชานนฺติ, เต น อญฺญสฺส สมณสฺส วา พฺราหฺมณสฺส วา มุขํ อุลฺโลเกนฺติ “อยํ นุน ภวํ ชานํ ชานาติ, ปสฺสํ ปสฺสตี”ติ.}

\prose{เสยฺยถาปิ ภิกฺขเว อโยขีโล วา อินฺทขีโล วา คมฺภีรเนโม สุนิขาโต อจโล อสมฺปกมฺปี, ปุรตฺถิมาย เจปิ ทิสาย อาคจฺเฉยฺย ภุสา วาตวุฏฺฐิ, เนว สงฺกมฺเปยฺย\footnote{เนว นํ สงฺกมฺเปยฺย (สี, อิ)} น สมฺปกมฺเปยฺย น สมฺปจาเลยฺย.\csromanspacer{} ปจฺฉิมาย เจปิ ทิสาย ฯเปฯ.\csromanspacer{} อุตฺตราย เจปิ ทิสาย ฯเปฯ.\csromanspacer{} ทกฺขิณาย เจปิ ทิสาย อาคจฺเฉยฺย ภุสา วาตวุฏฺฐิ, เนว สงฺกมฺเปยฺย น สมฺปกมฺเปยฺย น สมฺปจาเลยฺย.\csromanspacer{} ตํ กิสฺส เหตุ, คมฺภีรตฺตา ภิกฺขเว เนมสฺส สุนิขาตตฺตา อินฺทขีลสฺส.\csromanspacer{} เอวเมว โข ภิกฺขเว เย จ โข เกจิ สมณา วา พฺราหฺมณา วา “อิทํ ทุกฺขนฺ”ติ ยถาภูตํ ปชานนฺติ ฯเปฯ “อยํ ทุกฺขนิโรธคามินี ปฏิปทา”ติ ยถาภูตํ ปชานนฺติ, เต น อญฺญสฺส สมณสฺส วา พฺราหฺมณสฺส วา มุขํ อุลฺโลเกนฺติ “อยํ นูน ภวํ ชานํ ชานาติ, ปสฺสํ ปสฺสตี”ติ.\csromanspacer{} ตํ กิสฺส เหตุ, สุทิฏฺฐตฺตา ภิกฺขเว จตุนฺนํ อริยสจฺจานํ.\csromanspacer{} กตเมสํ จตุนฺนํ, ทุกฺขสฺส อริยสจฺจสฺส ฯเปฯ ทุกฺขนิโรธคามินิยา ปฏิปทาย อริยสจฺจสฺส.}

\prose{ตสฺมาติห ภิกฺขเว “อิทํ ทุกฺขนฺ”ติ โยโค กรณีโย ฯเปฯ “อยํ ทุกฺขนิโรธคามินี ปฏิปทา”ติ โยโค กรณีโยติ.\csromanspacer{}\csromanspacer{}นวมํ.}

\csromanheader{2}{12. สจฺจสํยุตฺต}
\csromanmatikaanchor{mtk.291}
\csromantocmarkref{subparagraph}{10. วาทตฺถิกสุตฺต}{mtk.291}
\bookmark[dest={mtk.291},level=5]{10. วาทตฺถิกสุตฺต}
\csromanheader{6}{10. วาทตฺถิกสุตฺต}
\proseitem{1110}{\csromanfolio{389}โย หิ โกจิ ภิกฺขเว ภิกฺขุ “อิทํ ทุกฺขนฺ”ติ ยถาภูตํ ปชานาติ ฯเปฯ “อยํ ทุกฺขนิโรธคามินี ปฏิปทา”ติ ยถาภูตํ ปชานาติ.\csromanspacer{} ปุรตฺถิมาย เจปิ ทิสาย อาคจฺเฉยฺย สมโณ วา พฺราหฺมโณ วา วาทตฺถิโก วาทคเวสี “วาทมสฺส อาโรเปสฺสามี”ติ, ตํ วต สหธมฺเมน สงฺกมฺเปสฺสติ วา สมฺปกมฺเปสฺสติ วา สมฺปจาเลสฺสติ วาติ เนตํ ฐานํ วิชฺชติ.\csromanspacer{} ปจฺฉิมาย เจปิ ทิสาย ฯเปฯ.\csromanspacer{} อุตฺตราย เจปิ ทิสาย ฯเปฯ.\csromanspacer{} ทกฺขิณาย เจปิ ทิสาย อาคจฺเฉยฺย สมโณ วา พฺราหฺมโณ วา วาทตฺถิโก วาทคเวสี “วาทมสฺส อาโรเปสฺสามี”ติ, ตํ วต สหธมฺเมน สงฺกมฺเปสฺสติ วา สมฺปกมฺเปสฺสติ วา สมฺปจาเลสฺสติ วาติ เนตํ ฐานํ วิชฺชติ.}

\prose{เสยฺยถาปิ ภิกฺขเว สิลายูโป โสฬสกุกฺกุโก, ตสฺสสฺสุ อฏฺฐ กุกฺกุ เหฏฺฐา เนมงฺคมา, อฏฺฐ กุกฺกุ อุปริเนมสฺส, ปุรตฺถิมาย เจปิ ทิสาย อาคจฺเฉยฺย ภุสา วาตวุฏฺฐิ, เนว สงฺกมฺเปยฺย น สมฺปกมฺเปยฺย น สมฺปจาเลยฺย.\csromanspacer{} ปจฺฉิมาย เจปิ ทิสาย ฯเปฯ.\csromanspacer{} อุตฺตราย เจปิ ทิสาย ฯเปฯ.\csromanspacer{} ทกฺขิณาย เจปิ ทิสาย อาคจฺเฉยฺย ภุสา วาตวุฏฺฐิ, เนว สงฺกมฺเปยฺย น สมฺปกมฺเปยฺย น สมฺปจาเลยฺย.\csromanspacer{} ตํ กิสฺส เหตุ, คมฺภีรตฺตา ภิกฺขเว เนมสฺส สุนิขาตตฺตา สิลายูปสฺส.\csromanspacer{} เอวเมว โข ภิกฺขเว โย หิ โกจิ ภิกฺขุ “อิทํ ทุกฺขนฺ”ติ ยถาภูตํ ปชานาติ ฯเปฯ “อยํ ทุกฺขนิโรธคามินี ปฏิปทา”ติ ยถาภูตํ ปชานาติ.\csromanspacer{} ปุรตฺถิมาย เจปิ ทิสาย อาคจฺเฉยฺย สมโณ วา พฺราหฺมโณ วา วาทตฺถิโก วาทคเวสี “วาทมสฺส อาโรเปสฺสามี”ติ, ตํ วต สหธมฺเมน สงฺกมฺเปสฺสติ วา สมฺปกมฺเปสฺสติ วา สมฺปจาเลสฺสติ วาติ เนตํ ฐานํ วิชฺชติ.\csromanspacer{} ปจฺฉิมาย เจปิ ทิสาย ฯเปฯ.\csromanspacer{} อุตฺตราย เจปิ ทิสาย ฯเปฯ.\csromanspacer{} ทกฺขิณาย เจปิ ทิสาย อาคจฺเฉยฺย สมโณ วา พฺราหฺมโณ วา วาทตฺถิโก วาทคเวสี “วาทมสฺส อาโรเปสฺสามี”ติ, ตํ วต สหธมฺเมน สงฺกมฺเปสฺสติ วา สมฺปกมฺเปสฺสติ วา สมฺปจาเลสฺสติ วาติ เนตํ ฐานํ วิชฺชติ.\csromanspacer{} ตํ กิสฺส เหตุ, สุทิฏฺฐตฺตา ภิกฺขเว จตุนฺนํ อริยสจฺจานํ.\csromanspacer{} กตเมสํ จตุนฺนํ, ทุกฺขสฺส อริยสจฺจสฺส ฯเปฯ ทุกฺขนิโรธคามินิยา ปฏิปทาย อริยสจฺจสฺส.}

\prose{ตสฺมาติห ภิกฺขเว “อิทํ ทุกฺขนฺ”ติ โยโค กรณีโย ฯเปฯ “อยํ ทุกฺขนิโรธคามินี ปฏิปทา”ติ โยโค กรณีโยติ.\csromanspacer{}\csromanspacer{}ทสมํ.\csromanspacer{} สีสปาวนวคฺโค จตุตฺโถ.}

\titlehead{\textbf{ตสฺสุทฺทานํ}}
\csromangathameasure{สีสปา ขทิโร ทณฺโฑ, เจลา สตฺติสเตน จ. \\ ปาณา สูริยูปมา เทฺวธา, อินฺทขีโล จ วาทิโนติ.}
\csromangathagroup{\csromangathastanza{\csromanfolio{390}สีสปา ขทิโร ทณฺโฑ, เจลา สตฺติสเตน จ. \\ ปาณา สูริยูปมา เทฺวธา, อินฺทขีโล จ วาทิโนติ.}}

\csromanmatikaanchor{mtk.292}
\csromantocmarknopage{subparagraph}{5. ปปาตวคฺค}{mtk.292}
\bookmark[dest={mtk.292},level=5]{5. ปปาตวคฺค}
\csromanheader{6}{5. ปปาตวคฺค}
\csromanmatikaanchor{mtk.293}
\csromantocmarkref{subparagraph}{1. โลกจินฺตาสุตฺต}{mtk.293}
\bookmark[dest={mtk.293},level=5]{1. โลกจินฺตาสุตฺต}
\csromanheader{6}{1. โลกจินฺตาสุตฺต}
\proseitem{1111}{เอกํ สมยํ ภควา ราชคเห วิหรติ เวฬุวเน กลนฺทกนิวาเป.\csromanspacer{} ตตฺร โข ภควา ภิกฺขู อามนฺเตสิ “ภูตปุพฺพํ ภิกฺขเว อญฺญตโร ปุริโส ราชคหา นิกฺขมิตฺวา ‘โลกจินฺตํ จินฺเตสฺสามี’ติ เยน สุมาคธา โปกฺขรณี เตนุปสงฺกมิ, อุปสงฺกมิตฺวา สุมาคธาย โปกฺขรณิยา ตีเร นิสีทิ โลกจินฺตํ จินฺเตนฺโต, อทฺทสา โข ภิกฺขเว โส ปุริโส สุมาคธาย โปกฺขรณิยา ตีเร จตุรงฺคินิํ เสนํ\footnote{จตุรงฺคินิเสนํ (ก)} ภิสมุฬาลํ\footnote{ภิสมูลาลํ (อิ, ก)} ปวิสนฺตํ, ทิสฺวานสฺส เอตทโหสิ ‘อุมฺมตฺโตสฺมิ นามาหํ, วิเจโตสฺมิ นามาหํ, ยํ โลเก นตฺถิ ตํ มยา ทิฏฺฐนฺ’ติ.\csromanspacer{} อถ โข โส ภิกฺขเว ปุริโส นครํ ปวิสิตฺวา มหาชนกายสฺส อาโรเจสิ ‘อุมฺมตฺโตสฺมิ นามาหํ ภนฺเต, วิเจโตสฺมิ นามาหํ ภนฺเต, ยํ โลเก นตฺถิ ตํ มยา ทิฏฺฐนฺ’ติ”.}

\prose{กถํ ปน ตฺวํ อมฺโภ ปุริส อุมฺมตฺโต, กถํ วิเจโต, กิญฺจ โลเก นตฺถิ ยํ ตยา ทิฏฺฐนฺติ.\csromanspacer{} อิธาหํ ภนฺเต ราชคหา นิกฺขมิตฺวา “โลกจินฺตํ จินฺเตสฺสามี”ติ เยน สุมาคธา โปกฺขรณี เตนุปสงฺกมิํ, อุปสงฺกมิตฺวา สุมาคธาย โปกฺขรณิยา ตีเร นิสีทิํ โลกจินฺตํ จินฺเตนฺโต, อทฺทสํ ขฺวาหํ ภนฺเต สุมาคธาย โปกฺขรณิยา ตีเร จตุรงฺคินิํ เสนํ ภิสมุฬาลํ ปวิสนฺตํ, เอวํ ขฺวาหํ ภนฺเต อุมฺมตฺโต, เอวํ วิเจโต, อิทญฺจ โลเก นตฺถิ ยํ มยา ทิฏฺฐนฺติ.\csromanspacer{} ตคฺฆ ตฺวํ อมฺโภ ปุริส อุมฺมตฺโต, ตคฺฆ วิเจโต, อิทญฺจ โลเก นตฺถิ ยํ ตยา ทิฏฺฐนฺติ.}

\prose{ตํ โข ปน ภิกฺขเว โส ปุริโส ภูตํเยว อทฺทส, โน อภูตํ.\csromanspacer{} ภูตปุพฺพํ ภิกฺขเว เทวาสุรสงฺคาโม สมุปพฺยูโฬฺห อโหสิ, ตสฺมิํ โข ปน ภิกฺขเว สงฺคาเม เทวา ชินิํสุ, อสุรา ปราชินิํสุ.}

\csromanheader{2}{12. สจฺจสํยุตฺต}
\prose{\csromanfolio{391}ปราชิตา จ โข ภิกฺขเว อสุรา ภีตา ภิสมุฬาเลน อสุรปุรํ ปวิสิํสุ เทวานํเยว โมหยมานา.}

\prose{ตสฺมาติห ภิกฺขเว มา โลกจินฺตํ จินฺเตถ, “สสฺสโต โลโก”ติ วา, “อสสฺสโต โลโก”ติ วา, “อนฺตวา โลโก”ติ วา, “อนนฺตวา โลโก”ติ วา, “ตํ ชีวํ ตํ สรีรนฺ”ติ วา, “อญฺญํ ชีวํ อญฺญํ สรีรนฺ”ติ วา, “โหติ ตถาคโต ปรํ มรณา”ติ วา, “น โหติ ตถาคโต ปรํ มรณา”ติ วา, “โหติ จ น จ โหติ ตถาคโต ปรํ มรณา”ติ วา, “เนว โหติ น น โหติ ตถาคโต ปรํ มรณา”ติ วา.\csromanspacer{} ตํ กิสฺส เหตุ, เนสา ภิกฺขเว จินฺตา อตฺถสํหิตา, นาทิพฺรหฺมจริยกา, น นิพฺพิทาย น วิราคาย น นิโรธาย น อุปสมาย น อภิญฺญาย น สมฺโพธาย น นิพฺพานาย สํวตฺตติ.}

\prose{จินฺเตนฺตา โข ตุเมฺห ภิกฺขเว “อิทํ ทุกฺขนฺ”ติ จินฺเตยฺยาถ ฯเปฯ “อยํ ทุกฺขนิโรธคามินี ปฏิปทา”ติ จินฺเตยฺยาถ.\csromanspacer{} ตํ กิสฺส เหตุ, เอสา ภิกฺขเว จินฺตา อตฺถสํหิตา, เอสา อาทิพฺรหฺมจริยกา, เอสา นิพฺพิทาย วิราคาย นิโรธาย อุปสมาย อภิญฺญาย สมฺโพธาย นิพฺพานาย สํวตฺตติ.}

\prose{ตสฺมาติห ภิกฺขเว “อิทํ ทุกฺขนฺ”ติ โยโค กรณีโย ฯเปฯ “อยํ ทุกฺขนิโรธคามินี ปฏิปทา”ติ โยโค กรณีโยติ.\csromanspacer{}\csromanspacer{}ปฐมํ.}

\csromanmatikaanchor{mtk.294}
\csromantocmarkref{subparagraph}{2. ปปาตสุตฺต}{mtk.294}
\bookmark[dest={mtk.294},level=5]{2. ปปาตสุตฺต}
\csromanheader{6}{2. ปปาตสุตฺต}
\proseitem{1112}{เอกํ สมยํ ภควา ราชคเห วิหรติ คิชฺฌกูเฏ ปพฺพเต.\csromanspacer{} อถ โข ภควา ภิกฺขู อามนฺเตสิ “อายาม ภิกฺขเว เยน ปฏิภานกูโฏ, เตนุปสงฺกมิสฺสาม ทิวาวิหารายา”ติ. “เอวํ ภนฺเต”ติ โข เต ภิกฺขู ภควโต ปจฺจสฺโสสุํ.\csromanspacer{} อถ โข ภควา สมฺพหุเลหิ ภิกฺขูหิ สทฺธิํ เยน ปฏิภานกูโฏ เตนุปสงฺกมิ.\csromanspacer{} อทฺทสา โข อญฺญตโร ภิกฺขุ ปฏิภานกูเฏ มหนฺตํ ปปาตํ, ทิสฺวาน ภควนฺตํ เอตทโวจ “มหา วตายํ ภนฺเต ปปาโต, สุภยานโก ภนฺเต ปปาโต, อตฺถิ นุ โข ภนฺเต อิมมฺหา ปปาตา อญฺโญ ปปาโต มหนฺตตโร จ ภยานกตโร จา”ติ.\csromanspacer{} อตฺถิ โข ภิกฺขุ อิมมฺหา ปปาตา อญฺโญ ปปาโต มหนฺตตโร จ ภยานกตโร จาติ.}

\prose{\csromanfolio{392}กตโม ปน ภนฺเต อิมมฺหา ปปาตา อญฺโญ ปปาโต มหนฺตตโร จ ภยานกตโร จาติ.\csromanspacer{} เย หิ เกจิ ภิกฺขเว สมณา วา พฺราหฺมณา วา “อิทํ ทุกฺขนฺ”ติ ยถาภูตํ นปฺปชานนฺติ, “อยํ ทุกฺขสมุทโย”ติ ยถาภูตํ นปฺปชานนฺติ, “อยํ ทุกฺขนิโรโธ”ติ ยถาภูตํ นปฺปชานนฺติ, “อยํ ทุกฺขนิโรธคามินี ปฏิปทา”ติ ยถาภูตํ นปฺปชานนฺติ.\csromanspacer{} เต ชาติสํวตฺตนิเกสุ สงฺขาเรสุ อภิรมนฺติ, ชราสํวตฺตนิเกสุ สงฺขาเรสุ อภิรมนฺติ, มรณสํวตฺตนิเกสุ สงฺขาเรสุ อภิรมนฺติ, โสกปริเทวทุกฺขโทมนสฺสุปายาสสํวตฺตนิเกสุ สงฺขาเรสุ อภิรมนฺติ.\csromanspacer{} เต ชาติสํวตฺตนิเกสุ สงฺขาเรสุ อภิรตา ชราสํวตฺตนิเกสุ สงฺขาเรสุ อภิรตา มรณสํวตฺตนิเกสุ สงฺขาเรสุ อภิรตา โสกปริเทวทุกฺขโทมนสฺสุปายาสสํวตฺตนิเกสุ สงฺขาเรสุ อภิรตา ชาติสํวตฺตนิเกปิ สงฺขาเร อภิสงฺขโรนฺติ, ชราสํวตฺตนิเกปิ สงฺขาเร อภิสงฺขโรนฺติ, มรณสํวตฺตนิเกปิ สงฺขาเร อภิสงฺขโรนฺติ โสกปริเทวทุกฺขโทนมสฺสุปายาสสํวตฺตนิเกปิ สงฺขาเร อภิสงฺขโรนฺติ.\csromanspacer{} เต ชาติสํวตฺตนิเกปิ สงฺขาเร อภิสงฺขริตฺวา ชราสํวตฺตนิเกปิ สงฺขาเร อภิสงฺขริตฺวา มรณสํวตฺตนิเกปิ สงฺขาเร อภิสงฺขริตฺวา โสกปริเทวทุกฺขโทมนสฺสุปายาสสํวตฺตนิเกปิ สงฺขาเร อภิสงฺขริตฺวา ชาติปปาตมฺปิ ปปตนฺติ, ชราปปาตมฺปิ ปปตนฺติ, มรณปปาตมฺปิ ปปตนฺติ, โสกปริเทวทุกฺขโทมนสฺสุปายาสปปาตมฺปิ ปปตนฺติ.\csromanspacer{} เต น ปริมุจฺจนฺติ ชาติยา ชราย มรเณน โสเกหิ ปริเทเวหิ ทุกฺเขหิ โทมนสฺเสหิ อุปายาเสหิ, น ปริมุจฺจนฺติ ทุกฺขสฺมาติ วทามิ.}

\prose{เย จ โข เกจิ ภิกฺขเว สมณา วา พฺราหฺมณา วา “อิทํ ทุกฺขนฺ”ติ ยถาภูตํ ปชานนฺติ ฯเปฯ “อยํ ทุกฺขนิโรธคามินี ปฏิปทา”ติ ยถาภูตํ ปชานนฺติ.\csromanspacer{} เต ชาติสํวตฺตนิเกสุ สงฺขาเรสุ นาภิรมนฺติ, ชราสํวตฺตนิเกสุ สงฺขาเรสุ นาภิรมนฺติ, มรณสํวตฺตนิเกสุ สงฺขาเรสุ นาภิรมนฺติ, โสกปริเทวทุกฺขโทมนสฺสุปายาสสํวตฺตนิเกสุ สงฺขาเรสุ นาภิรมนฺติ, เต ชาติสํวตฺตนิเกสุ สงฺขาเรสุ อนภิรตา ชราสํวตฺตนิเกสุ สงฺขาเรสุ อนภิรตา มรณสํวตฺตนิเกสุ สงฺขาเรสุ อนภิรตา โสกปริเทวทุกฺขโทมนสฺสุปายาสสํวตฺตนิเกสุ สงฺขาเรสุ อนภิรตา ชาติสํวตฺตนิเกปิ สงฺขาเร นาภิสงฺขาโรนฺติ, ชราสํวตฺตนิเกปิ สงฺขาเร นาภิสงฺขโรนฺติ, มรณสํวตฺตนิเกปิ สงฺขาเร นาภิสงฺขโรนฺติ,}

\csromanheader{2}{12. สจฺจสํยุตฺต}
\prose{\csromanfolio{393}โสกปริเทวทุกฺขโทมนสฺสุปายาสสํวตฺตนิเกปิ สงฺขาเร นาภิสงฺขโรนฺติ, เต ชาติสํวตฺตนิเกปิ สงฺขาเร อนภิสงฺขริตฺวา ชราสํวตฺตนิเกปิ สงฺขาเร อนภิสงฺขริตฺวา มรณสํวตฺตนิเกปิ สงฺขาเร อนภิสงฺขริตฺวา โสกปริเทวทุกฺขโทมนสฺสุปายาสสํวตฺตนิเกปิ สงฺขาเร อนภิสงฺขริตฺวา ชาติปปตมฺปิ นปฺปปตนฺติ, ชราปปาตมฺปิ นปฺปปตนฺติ, มรณปปาตมฺปิ นปฺปปตนฺติ, โสกปริเทวทุกฺขโทมนสฺสุปายาสปปาตมฺปิ นปฺปปตนฺติ, เต ปริมุจฺจนฺติ ชาติยา ชราย มรเณน โสเกหิ ปริเทเวหิ ทุกฺเขหิ โทมนสฺเสหิ อุปายาเสหิ, ปริมุจฺจนฺติ ทุกฺขสฺมาติ วทามิ.}

\prose{ตสฺมาติห ภิกฺขเว “อิทํ ทุกฺขนฺ”ติ โยโค กรณีโย ฯเปฯ “อยํ ทุกฺขนิโรธคามินี ปฏิปทา”ติ โยโค กรณีโยติ.\csromanspacer{}\csromanspacer{}ทุติยํ.}

\csromanmatikaanchor{mtk.295}
\csromantocmarkref{subparagraph}{3. มหาปริฬาหสุตฺต}{mtk.295}
\bookmark[dest={mtk.295},level=5]{3. มหาปริฬาหสุตฺต}
\csromanheader{6}{3. มหาปริฬาหสุตฺต}
\proseitem{1113}{อตฺถิ ภิกฺขเว มหาปริฬาโห นาม นิรโย, ตตฺถ ยํ กิญฺจิ จกฺขุนา รูปํ ปสฺสติ, อนิฏฺฐรูปญฺเญว ปสฺสติ โน อิฏฺฐรูปํ, อกนฺตรูปญฺเญว ปสฺสติ โน กนฺตรูปํ, อมนาปรูปญฺเญว ปสฺสติ โน มนาปรูปํ.\csromanspacer{} ยํ กิญฺจิ โสเตน สทฺทํ สุณาติ ฯเปฯ.\csromanspacer{} ยํ กิญฺจิ กาเยน โผฏฺฐพฺพํ ผุสติ ฯเปฯ.\csromanspacer{} ยํ กิญฺจิ มนสา ธมฺมํ วิชานาติ, อนิฏฺฐรูปญฺเญว วิชานาติ โน อิฏฺฐรูปํ, อกนฺตรูปญฺเญว วิชานาติ, โน กนฺตรูปํ, อมนาปรูปญฺเญว วิชานาติ, โน มนาปรูปนฺติ.}

\prose{เอวํ วุตฺเต อญฺญตโร ภิกฺขุ ภควนฺตํ เอตทโวจ “มหา วต โส ภนฺเต ปริฬาโห, สุมหา วต โส ภนฺเต ปริฬาโห, อตฺถิ นุ โข ภนฺเต เอตมฺหา ปริฬาหา อญฺโญ ปริฬาโห มหนฺตตโร เจว ภยานกตโร จา”ติ.\csromanspacer{} อตฺถิ โข ภิกฺขุ เอตมฺหา ปริฬาหา อญฺโญ ปริฬาโห มหนฺตตโร จ ภยานกตโร จาติ.}

\prose{กตโม ปน ภนฺเต เอตมฺหา ปริฬาหา อญฺโญ ปริฬาโห มหนฺตตโร จ ภยานกตโร จาติ.\csromanspacer{} เยหิ เกจิ ภิกฺขเว สมณา วา พฺราหฺมณา วา “อิทํ ทุกฺขนฺ”ติ ยถาภูตํ นปฺปชานนฺติ ฯเปฯ “อยํ ทุกฺขนิโรธคามินี ปฏิปทา”ติ ยถาภูตํ นปฺปชานนฺติ.\csromanspacer{} เต ชาติสํวตฺตนิเกสุ สงฺขาเรสุ อภิรมนฺติ ฯเปฯ อภิรตา ฯเปฯ อภิสงฺขโรนฺติ ฯเปฯ อภิสงฺขริตฺวา ชาติปริฬาเหนปิ ปริฑยฺหนฺติ, ชราปริฬาเหนปิ ปริฑยฺหนฺติ, มรณปริฬาเหนปิ ปริฑยฺหนฺติ, โสกปริเทวทุกฺขโทมนสฺสุปายาสปริฬาเหนปิ \csromanfolio{394}ปริฑยฺหนฺติ, เต น ปริมุจฺจนฺติ ชาติยา ชราย มรเณน โสเกหิ ปริเทเวหิ ทุกฺเขหิ โทมนสฺเสหิ อุปายาเสหิ, น ปริมุจฺจนฺติ ทุกฺขสฺมาติ วทามิ.}

\prose{เย จ โข เกจิ ภิกฺขเว สมณา วา พฺราหฺมณา วา “อิทํ ทุกฺขนฺ”ติ ยถาภูตํ ปชานนฺติ ฯเปฯ “อยํ ทุกฺขนิโรธคามินี ปฏิปทา”ติ ยถาภูตํ ปชานนฺติ.\csromanspacer{} เต ชาติสํวตฺตนิเกสุ สงฺขาเรสุ นาภิรมนฺติ ฯเปฯ อนภิรตา ฯเปฯ นาภิสงฺขโรนฺติ ฯเปฯ อนภิสงฺขริตฺวา ชาติปริฬาเหนปิ น ปริฑยฺหนฺติ, ชราปริฬาเหนปิ น ปริฑยฺหนฺติ, มรณปริฬาเหนปิ น ปริฑยฺหนฺติ, โสกปริเทว ทุกฺข โทมนสฺสุปายาสปริฬาเหนปิ น ปริฑยฺหนฺติ, เต ปริมุจฺจนฺติ ชาติยา ชราย มรเณน โสเกหิ ปริเทเวหิ ทุกฺเขหิ โทมนสฺเสหิ อุปายาเสหิ, ปริมุจฺจนฺติ ทุกฺขสฺมาติ วทามิ.}

\prose{ตสฺมาติห ภิกฺขเว “อิทํ ทุกฺขนฺ”ติ โยโค กรณีโย ฯเปฯ “อยํ ทุกฺขนิโรธคามินี ปฏิปทา”ติ โยโค กรณีโยติ.\csromanspacer{}\csromanspacer{}ตติยํ.}

\csromanmatikaanchor{mtk.296}
\csromantocmarkref{subparagraph}{4. กูฏาคารสุตฺต}{mtk.296}
\bookmark[dest={mtk.296},level=5]{4. กูฏาคารสุตฺต}
\csromanheader{6}{4. กูฏาคารสุตฺต}
\proseitem{1114}{โย หิ ภิกฺขเว\footnote{โย จ โข ภิกฺขเว (สฺยา, ก)} เอวํ วเทยฺย “อหํ ทุกฺขํ อริยสจฺจํ ยถาภูตํ อนภิสเมจฺจ ฯเปฯ ทุกฺขนิโรธคามินิํ ปฏิปทํ อริยสจฺจํ ยถาภูตํ อนภิสเมจฺจ สมฺมา ทุกฺขสฺสนฺตํ กริสฺสามี”ติ เนตํ ฐานํ วิชฺชติ.}

\prose{เสยฺยถาปิ ภิกฺขเว โย เอวํ วเทยฺย “อหํ กูฏาคารสฺส เหฏฺฐิมํ ฆรํ อกริตฺวา อุปริมํ ฆรํ อาโรเปสฺสามี”ติ เนตํ ฐานํ วิชฺชติ.\csromanspacer{} เอวเมว โข ภิกฺขเว โย เอวํ วเทยฺย “อหํ ทุกฺขํ อริยสจฺจํ ยถาภูตํ อนภิสเมจฺจ ฯเปฯ ทุกฺขนิโรธคามินิํ ปฏิปทํ อริยสจฺจํ ยถาภูตํ อนภิสเมจฺจ สมฺมา ทุกฺขสฺสนฺตํ กริสฺสามี”ติ เนตํ ฐานํ วิชฺชติ.}

\prose{โย จ โข ภิกฺขเว เอวํ วเทยฺย “อหํ ทุกฺขํ อริยสจฺจํ ยถาภูตํ อภิสเมจฺจ ฯเปฯ ทุกฺขนิโรธคามินิํ ปฏิปทํ อริยสจฺจํ ยถาภูตํ อภิสเมจฺจ สมฺมา ทุกฺขสฺสนฺตํ กริสฺสามี”ติ ฐานเมตํ วิชฺชติ.}

\prose{เสยฺยถาปิ ภิกฺขเว โย เอวํ วเทยฺย “อหํ กูฏาคารสฺส เหฏฺฐิมํ ฆรํ กริตฺวา อุปริมํ ฆรํ อาโรเปสฺสามี”ติ ฐานเมตํ วิชฺชติ.\csromanspacer{} เอวเมว}

\csromanheader{2}{12. สจฺจสํยุตฺต}
\prose{\csromanfolio{395}โข ภิกฺขเว โย เอวํ วเทยฺย “อหํ ทุกฺขํ อริยสจฺจํ ยถาภูตํ อภิสเมจฺจ ฯเปฯ ทุกฺขนิโรธคามินิํ ปฏิปทํ อริยสจฺจํ ยถาภูตํ อภิสเมจฺจ สมฺมา ทุกฺขสฺสนฺตํ กริสฺสามี”ติ ฐานเมตํ วิชฺชติ.}

\prose{ตสฺมาติห ภิกฺขเว “อิทํ ทุกฺขนฺ”ติ โยโค กรณีโย ฯเปฯ “อยํ ทุกฺขนิโรธคามินี ปฏิปทา”ติ โยโค กรณีโยติ.\csromanspacer{}\csromanspacer{}จตุตฺถํ.}

\csromanmatikaanchor{mtk.297}
\csromantocmarkref{subparagraph}{5. วาลสุตฺต}{mtk.297}
\bookmark[dest={mtk.297},level=5]{5. วาลสุตฺต}
\csromanheader{6}{5. วาลสุตฺต}
\proseitem{1115}{เอกํ สมยํ ภควา เวสาลิยํ วิหรติ มหาวเน กูฏาคารสาลายํ.\csromanspacer{} อถ โข อายสฺมา อานนฺโท ปุพฺพณฺหสมยํ นิวาเสตฺวา ปตฺตจีวรมาทาย เวสาลิํ ปิณฺฑาย ปาวิสิ.\csromanspacer{} อทฺทสา โข อายสฺมา อานนฺโท สมฺพหุเล ลิจฺฉวิกุมารเก สนฺถาคาเร อุปาสนํ กโรนฺเต ทูรโตว สุขุเมน ตาฬจฺฉิคฺคเฬน อสนํ อติปาเตนฺเต โปงฺขานุโปงฺขํ\footnote{โปขานุโปขํ (สฺยา, กํ)} อวิราธิตํ, ทิสฺวานสฺส เอตทโหสิ “สิกฺขิตา วติเม ลิจฺฉวิกุมารกา, สุสิกฺขิตา วติเม ลิจฺฉวิกุมารกา, ยตฺร หิ นาม ทูรโตว สุขุเมน ตาฬจฺฉิคฺคเฬน อสนํ อติปาเตสฺสนฺติ โปงฺขานุโปงฺขํ อวิราธิตนฺ”ติ.}

\prose{อถ โข อายสฺมา อานนฺโท เวสาลิํ ปิณฺฑาย จริตฺวา ปจฺฉาภตฺตํ ปิณฺฑปาตปฏิกฺกนฺโต เยน ภควา เตนุปสงฺกมิ, อุปสงฺกมิตฺวา ภควนฺตํ อภิวาเทตฺวา เอกมนฺตํ นิสีทิ, เอกมนฺตํ นิสินฺโน โข อายสฺมา อานนฺโท ภควนฺตํ เอตทโวจ “อิธาหํ ภนฺเต ปุพฺพณฺหสมยํ นิวาเสตฺวา ปตฺตจีวรมาทาย เวสาลิํ ปิณฺฑาย ปาวิสิํ, อทฺทสํ ขฺวาหํ ภนฺเต สมฺพหุเล ลิจฺฉวิกุมารเก สนฺถาคาเร อุปาสนํ กโรนฺเต ทูรโตว สุขุเมน ตาฬจฺฉิคฺคเฬน อสนํ อติปาเตนฺเต โปงฺขานุโปงฺขํ อวิราธิตํ, ทิสฺวาน เม เอตทโหสิ ‘สิกฺขิตา วติเม ลิจฺฉวิกุมารกา, สุสิกฺขิตา วติเม ลิจฺฉวิกุมารกา, ยตฺร หิ นาม ทูรโตว สุขุเมน ตาฬจฺฉิคฺคเฬน อสนํ อติปาเตสฺสนฺติ โปงฺขานุโปงฺขํ อวิราธิตนฺ’ติ”.}

\csromanmatikaanchor{mtk.298}
\csromantocmarkref{subparagraph}{6. อนฺทการสุตฺต}{mtk.298}
\bookmark[dest={mtk.298},level=5]{6. อนฺทการสุตฺต}
\prose{ตํ กิํ มญฺญสิ อานนฺท, กตมํ นุ โข ทุกฺกรตรํ วา ทุรภิสมฺภวตรํ วา, โย ทูรโตว สุขุเมน ตาฬจฺฉิคฺคเฬน อสนํ อติปาเตยฺย โปงฺขานุโปงฺขํ อวิราธิตํ, โย วา สตฺตธา ภินฺนสฺส วาลสฺส โกฏิยา โกฏิํ ปฏิวิชฺเฌยฺยาติ.\csromanspacer{} เอตเทว ภนฺเต ทุกฺกรตรญฺเจว ทุรภิสมฺภวตรญฺจ, \csromanfolio{396}โย วา\footnote{โย (สี)} สตฺตธา ภินฺนสฺส วาลสฺส โกฏิยา โกฏิํ ปฏิวิชฺเฌยฺยาติ.\csromanspacer{} อถ โข\footnote{อถ โข เต (สฺยา, กํ)} อานนฺท ทุปฺปฏิวิชฺฌตรํ ปฏิวิชฺฌนฺติ, เย “อิทํ ทุกฺขนฺ”ติ ยถาภูตํ ปฏิวิชฺฌนฺติ ฯเปฯ “อยํ ทุกฺขนิโรธคามินี ปฏิปทา”ติ ยถาภูตํ ปฏิวิชฺฌนฺติ.}

\prose{ตสฺมาติหานนฺท “อิทํ ทุกฺขนฺ”ติ โยโค กรณีโย ฯเปฯ “อยํ ทุกฺขนิโรธคามินี ปฏิปทา”ติ โยโค กรณีโยติ.\csromanspacer{}\csromanspacer{}ปญฺจมํ.}

\csromanheader{2}{6. อนฺธการสุตฺต}
\proseitem{1116}{อตฺถิ ภิกฺขเว โลกนฺตริกา อฆา อสํวุตา อนฺธการา อนฺธการติมิสา, ยตฺถมิเมสํ จนฺทิมสูริยานํ เอวํมหิทฺธิกานํ เอวํมหานุภาวานํ อาภาย นานุโภนฺตีติ.}

\prose{เอวํ วุตฺเต อญฺญตโร ภิกฺขุ ภควนฺตํ เอตทโวจ “มหา วต โส ภนฺเต อนฺธกาโร, สุมหา วต โส ภนฺเต อนฺธกาโร, อตฺถิ นุ โข ภนฺเต เอตมฺหา อนฺธการา อญฺโญ อนฺธกาโร มหนฺตตโร จ ภยานกตโร จา”ติ.\csromanspacer{} อตฺถิ โข ภิกฺขุ เอตมฺหา อนฺธการา อญฺโญ อนฺธกาโร มหนฺตตโร จ ภยานกตโร จาติ.}

\prose{กตโม ปน ภนฺเต เอตมฺหา อนฺธการา อญฺโญ อนฺธกาโร มหนฺตตโร จ ภยานกตโร จาติ.\csromanspacer{} เย หิ เกจิ ภิกฺขุ สมณา วา พฺราหฺมณา วา “อิทํ ทุกฺขนฺ”ติ ยถาภูตํ นปฺปชานนฺติ ฯเปฯ “อยํ ทุกฺขนิโรธคามินี ปฏิปทา”ติ ยถาภูตํ นปฺปชานนฺติ.\csromanspacer{} เต ชาติสํวตฺตนิเกสุ สงฺขาเรสุ อภิรมนฺติ ฯเปฯ อภิรตา ฯเปฯ อภิสงฺขโรนฺติ ฯเปฯ อภิสงฺขริตฺวา ชาตนฺธการมฺปิ ปปตนฺติ, ชรนฺธการมฺปิ ปปตนฺติ, มรณนฺธการมฺปิ ปปตนฺติ, โสกปริเทวทุกฺขโทมนสฺสุปายาสนฺธการมฺปิ ปปตนฺติ.\csromanspacer{} เต น ปริมุจฺจนฺติ ชาติยา ชราย มรเณน โสเกหิ ปริเทเวหิ ทุกฺเขหิ โทมนสฺเสหิ อุปายาเสหิ, น ปริมุจฺจนฺติ ทุกฺขสฺมาติ วทามิ.}

\prose{เย จ โข เกจิ ภิกฺขุ สมณา วา พฺราหฺมณา วา “อิทํ ทุกฺขนฺ”ติ ยถาภูตํ ปชานนฺติ ฯเปฯ “อยํ ทุกฺขนิโรธคามินี ปฏิปทา”ติ ยถาภูตํ ปชานนฺติ.\csromanspacer{} เต ชาติสํวตฺตนิเกสุ สงฺขาเรสุ นาภิรมนฺติ ฯเปฯ อนภิรตา ฯเปฯ นาภิสงฺขโรนฺติ ฯเปฯ อนภิสงฺขริตฺวา ชาตนฺธการมฺปิ}

\csromanheader{2}{12. สจฺจสํยุตฺต}
\csromanmatikaanchor{mtk.299}
\csromantocmarkref{subparagraph}{7-8. ฉิคฺคฬยุคสุตฺตานิ}{mtk.299}
\bookmark[dest={mtk.299},level=5]{7-8. ฉิคฺคฬยุคสุตฺตานิ}
\prose{\csromanfolio{397}นปฺปปตนฺติ, ชรนฺธการมฺปิ นปฺปปตนฺติ, มรณนฺธการมฺปิ นปฺปปตนฺติ, โสกปริเทวทุกฺขโทมนสฺสุปายาสนฺธการมฺปิ นปฺปปตนฺติ.\csromanspacer{} เต ปริมุจฺจนฺติ ชาติยา ชราย มรเณน โสเกหิ ปริเทเวหิ ทุกฺเขหิ โทมนสฺเสหิ อุปายาเสหิ, ปริมุจฺจนฺติ ทุกฺขสฺมาติ วทามิ.}

\prose{ตสฺมาติห ภิกฺขเว “อิทํ ทุกฺขนฺ”ติ โยโค กรณีโย ฯเปฯ “อยํ ทุกฺขนิโรธคามินี ปฏิปทา”ติ โยโค กรณีโยติ.\csromanspacer{}\csromanspacer{}ฉฏฺฐํ.}

\csromanheader{2}{7. ปฐมฉิคฺคฬยุคสุตฺต}
\proseitem{1117}{เสยฺยถาปิ ภิกฺขเว ปุริโส มหาสมุทฺเท เอกจฺฉิคฺคฬํ ยุคํ ปกฺขิเปยฺย, ตตฺราปิสฺส กาโณ กจฺฉโป, โส วสฺสสตสฺส วสฺสสตสฺส อจฺจเยน สกิํ สกิํ อุมฺมุชฺเชยฺย.\csromanspacer{} ตํ กิํ มญฺญถ ภิกฺขเว, อปิ นุ โข กาโณ กจฺฉโป วสฺสสตสฺส วสฺสสตสฺส อจฺจเยน สกิํ สกิํ อุมฺมุชฺชนฺโต อมุสฺมิํ เอกจฺฉิคฺคเฬ ยุเค คีวํ ปเวเสยฺยาติ.\csromanspacer{} ยทิ นูน ภนฺเต กทาจิ กรหจิ ทีฆสฺส อทฺธุโน อจฺจเยนาติ.}

\prose{ขิปฺปตรํ โข โส ภิกฺขเว กาโณ กจฺฉโป วสฺสสตสฺส วสฺสสตสฺส อจฺจเยน สกิํ สกิํ อุมฺมุชฺชนฺโต อมุสฺมิํ เอกจฺฉิคฺคเฬ ยุเค คีวํ ปเวเสยฺย, น เตฺววาหํ ภิกฺขเว สกิํ วินิปาตคเตน พาเลน\footnote{วินีตคเตน พหุเลน (ก)} มนุสฺสตฺตํ วทามิ.}

\prose{ตํ กิสฺส เหตุ, น เหตฺถ ภิกฺขเว อตฺถิ ธมฺมจริยา สมจริยา กุสลกิริยา ปุญฺญกิริยา, อญฺญมญฺญขาทิกา เอตฺถ ภิกฺขเว วตฺตติ ทุพฺพลขาทิกา.\csromanspacer{} ตํ กิสฺส เหตุ, อทิฏฺฐตฺตา ภิกฺขเว จตุนฺนํ อริยสจฺจานํ.\csromanspacer{} กตเมสํ จตุนฺนํ, ทุกฺขสฺส อริยสจฺจสฺส ฯเปฯ ทุกฺขนิโรธคามินิยา ปฏิปทาย อริยสจฺจสฺส.}

\prose{ตสฺมาติห ภิกฺขเว “อิทํ ทุกฺขนฺ”ติ โยโค กรณีโย ฯเปฯ “อยํ ทุกฺขนิโรธคามินี ปฏิปทา”ติ โยโค กรณีโยติ.\csromanspacer{}\csromanspacer{}สตฺตมํ.}

\csromanheader{2}{8. ทุติยฉิคฺคฬยุคสุตฺต}
\csromanmatikaanchor{mtk.300}
\csromantocmarkref{subparagraph}{9-10. สิเนรุปพฺพตราชสุตฺตานิ}{mtk.300}
\bookmark[dest={mtk.300},level=5]{9-10. สิเนรุปพฺพตราชสุตฺตานิ}
\proseitem{1118}{\csromanfolio{398}เสยฺยถาปิ ภิกฺขเว อยํ มหาปถวี เอโกทกา อสฺส, ตตฺร ปุริโส เอกจฺฉิคฺคฬํ ยุคํ ปกฺขิเปยฺย, ตเมนํ ปุรตฺถิโม วาโต ปจฺฉิเมน สํหเรยฺย.\csromanspacer{} ปจฺฉิโม วาโต ปุรตฺถิเมน สํหเรยฺย.\csromanspacer{} อุตฺตโร วาโต ทกฺขิเณน สํหเรยฺย.\csromanspacer{} ทกฺขิโณ วาโต อุตฺตเรน สํหเรยฺย.\csromanspacer{} ตตฺรสฺส กาโณ กจฺฉโป, โส วสฺสสตสฺส วสฺสสตสฺส อจฺจเยน สกิํ สกิํ อุมฺมุชฺเชยฺย.\csromanspacer{} ตํ กิํ มญฺญถ ภิกฺขเว, อปิ นุ โข กาโณ กจฺฉโป วสฺสสตสฺส วสฺสสตสฺส อจฺจเยน สกิํ สกิํ อุมฺมุชฺชนฺโต อมุสฺมิํ เอกจฺฉิคฺคเฬ ยุเค คีวํ ปเวเสยฺยาติ.\csromanspacer{} อธิจฺจมิทํ ภนฺเต, ยํ โส กาโณ กจฺฉโป วสฺสสตสฺส วสฺสสตสฺส อจฺจเยน สกิํ สกิํ อุมฺมุชฺชนฺโต อมุสฺมิํ เอกจฺฉิคฺคเฬ ยุเค คีวํ ปเวเสยฺยาติ.}

\prose{เอวํ อธิจฺจมิทํ ภิกฺขเว ยํ มนุสฺสตฺตํ ลภติ, เอวํ อธิจฺจมิทํ ภิกฺขเว ยํ ตถาคโต โลเก อุปฺปชฺชติ อรหํ สมฺมาสมฺพุทฺโธ, เอวํ อธิจฺจมิทํ ภิกฺขเว ยํ ตถาคตปฺปเวทิโต ธมฺมวินโย โลเก ทิพฺพติ, ตสฺสิทํ\footnote{ตยิทํ (?)} ภิกฺขเว มนุสฺสตฺตํ ลทฺธํ ตถาคโต โลเก อุปฺปนฺโน อรหํ สมฺมาสมฺพุทฺโธ, ตถาคตปฺปเวทิโต จ ธมฺมวินโย โลเก ทิพฺพติ.}

\prose{ตสฺมาติห ภิกฺขเว “อิทํ ทุกฺขนฺ”ติ โยโค กรณีโย ฯเปฯ “อยํ ทุกฺขนิโรธคามินี ปฏิปทา”ติ โยโค กรณีโยติ.\csromanspacer{}\csromanspacer{}อฏฺฐมํ.}

\csromanheader{2}{9. ปฐมสิเนรุปพฺพตราชสุตฺต}
\proseitem{1119}{เสยฺยถาปิ ภิกฺขเว ปุริโส สิเนรุสฺส ปพฺพตราชสฺส สตฺต มุคฺคมตฺติโย ปาสาณสกฺขรา อุปนิกฺขิเปยฺย.\csromanspacer{} ตํ กิํ มญฺญถ ภิกฺขเว, กตมํ นุ โข พหุตรํ, ยา วา\footnote{ยา จ} สตฺต มุคฺคมตฺติโย ปาสาณสกฺขรา อุปนิกฺขิตฺตา, โย วา\footnote{โย จ (สฺยา, กํ, อิ, ก) สํ 1. 356 ปิฏฺเฐปิ.} สิเนรุปพฺพตราชาติ.\csromanspacer{} เอตเทว ภนฺเต พหุตรํ ยทิทํ สิเนรุปพฺพตราชา, อปฺปมตฺติกา สตฺต มุคฺคมตฺติโย ปาสาณสกฺขรา อุปนิกฺขิตฺตา, สงฺขมฺปิ น อุเปนฺติ อุปนิธมฺปิ น อุเปนฺติ กลภาคมฺปิ น อุเปนฺติ, สิเนรุปพฺพตราชานํ อุปนิธาย สตฺต มุคฺคมตฺติโย ปาสาณสกฺขรา อุปนิกฺขิตฺตาติ.\csromanspacer{} เอวเมว โข ภิกฺขเว อริยสาวกสฺส ทิฏฺฐิสมฺปนฺนสฺส}

\csromanheader{2}{12. สจฺจสํยุตฺต}
\prose{\csromanfolio{399}ปุคฺคลสฺส อภิสเมตาวิโน เอตเทว พหุตรํ ทุกฺขํ, ยทิทํ ปริกฺขีณํ ปริยาทินฺนํ อปฺปมตฺตกํ อวสิฏฺฐํ, สงฺขมฺปิ น อุเปติ อุปนิธมฺปิ น อุเปติ กลภาคมฺปิ น อุเปติ, ปุริมํ ทุกฺขกฺขนฺธํ ปริกฺขีณํ ปริยาทินฺนํ อุปนิธาย, ยทิทํ สตฺตกฺขตฺตุปรมตา.\csromanspacer{} โย “อิทํ ทุกฺขนฺ”ติ ยถาภูตํ ปชานาติ. “อยํ ทุกฺขนิโรธคามินี ปฏิปทา”ติ ยถาภูตํ ปชานาติ.}

\prose{ตสฺมาติห ภิกฺขเว “อิทํ ทุกฺขนฺ”ติ โยโค กรณีโย ฯเปฯ “อยํ ทุกฺขนิโรธคามินี ปฏิปทา”ติ โยโค กรณีโยติ.\csromanspacer{}\csromanspacer{}นวมํ.}

\csromanheader{2}{10. ทุติยสิเนรุปพฺพตราชสุตฺต}
\proseitem{1120}{เสยฺยถาปิ ภิกฺขเว สิเนรุปพฺพตราชายํ ปริกฺขยํ ปริยาทานํ คจฺเฉยฺย ฐเปตฺวา สตฺต มุคฺคมตฺติโย ปาสาณสกฺขรา.\csromanspacer{} ตํ กิํ มญฺญถ ภิกฺขเว, กตมํ นุ โข พหุตรํ, ยํ วา สิเนรุสฺส ปพฺพตราชสฺส ปริกฺขีณํ ปริยาทินฺนํ, ยา วา สตฺต มุคฺคมตฺติโย ปาสาณสกฺขรา อวสิฏฺฐาติ.\csromanspacer{} เอตเทว ภนฺเต พหุตรํ สิเนรุสฺส ปพฺพตราชสฺส, ยทิทํ ปริกฺขีณํ ปริยาทินฺนํ, อปฺปมตฺติกา สตฺต มุคฺคมตฺติโย ปาสาณสกฺขรา อวสิฏฺฐา, สงฺขมฺปิ น อุเปนฺติ อุปนิธมฺปิ น อุเปนฺติ กลภาคมฺปิ น อุเปนฺติ, สิเนรุสฺส ปพฺพตราชสฺส ปริกฺขีณํ ปริยาทินฺนํ, อุปนิธาย สตฺต มุคฺคมตฺติโย ปาสาณสกฺขรา อวสิฏฺฐาติ.\csromanspacer{} เอวเมว โข ภิกฺขเว อริยสาวกสฺส ทิฏฺฐิสมฺปนฺนสฺส ปุคฺคลสฺส อภิสเมตาวิโน เอตเทว พหุตรํ ทุกฺขํ, ยทิทํ ปริกฺขีณํ ปริยาทินฺนํ อปฺปมตฺตกํ อวสิฏฺฐํ, สงฺขมฺปิ น อุเปติ อุปนิธมฺปิ น อุเปติ กลภาคมฺปิ น อุเปติ, ปุริมํ ทุกฺขกฺขนฺธํ ปริกฺขีณํ ปริยาทินฺนํ อุปนิธาย, ยทิทํ สตฺตกฺขตฺตุปรมตา.\csromanspacer{} โย “อิทํ ทุกฺขนฺ”ติ ยถาภูตํ ปชานาติ ฯเปฯ “อยํ ทุกฺขนิโรธคามินี ปฏิปทา”ติ ยถาภูตํ ปชานาติ.}

\prose{ตสฺมาติห ภิกฺขเว “อิทํ ทุกฺขนฺ”ติ โยโค กรณีโย ฯเปฯ “อยํ ทุกฺขนิโรธคามินี ปฏิปทา”ติ โยโค กรณีโยติ.\csromanspacer{}\csromanspacer{}ทสมํ.\csromanspacer{} ปปาตวคฺโค ปญฺจโม.}

\titlehead{\textbf{ตสฺสุทฺทานํ}}
\csromangathameasure{จินฺตา ปปาโต ปริฬาโห, กูฏํ วาลนฺธกาโร จ. \\ ฉิคฺคเฬน จ เทฺว วุตฺตา, สิเนรุ อปเร ทุเวติ.}
\csromangathagroup{\csromangathastanza{จินฺตา ปปาโต ปริฬาโห, กูฏํ วาลนฺธกาโร จ. \\ ฉิคฺคเฬน จ เทฺว วุตฺตา, สิเนรุ อปเร ทุเวติ.}}

\csromanmatikaanchor{mtk.301}
\csromantocmarknopage{subparagraph}{6. อภิสมยวคฺค}{mtk.301}
\bookmark[dest={mtk.301},level=5]{6. อภิสมยวคฺค}
\csromanheader{6}{6. อภิสมยวคฺค}
\csromanmatikaanchor{mtk.302}
\csromantocmarkref{subparagraph}{1. นขสิขสุตฺต}{mtk.302}
\bookmark[dest={mtk.302},level=5]{1. นขสิขสุตฺต}
\csromanheader{6}{1. นขสิขสุตฺต}
\proseitem{1121}{\csromanfolio{400}อถ โข ภควา ปริตฺตํ นขสิขายํ ปํสุํ อาโรเปตฺวา ภิกฺขู อามนฺเตสิ “ตํ กิํ มญฺญถ ภิกฺขเว, กตมํ นุ โข พหุตรํ โย วายํ มยา ปริตฺโต นขสิขายํ ปํสุ อาโรปิโต, อยํ วา มหาปถวี”ติ.\csromanspacer{} เอตเทว ภนฺเต พหุตรํ ยทิทํ มหาปถวี, อปฺปมตฺตกายํ ภควตา ปริตฺโต นขสิขายํ ปํสุ อาโรปิโต, สงฺขมฺปิ น อุเปติ อุปนิธมฺปิ น อุเปติ กลภาคมฺปิ น อุเปติ, มหาปถวิํ อุปนิธาย ภควตา ปริตฺโต นขสิขายํ ปํสุ อาโรปิโตติ.\csromanspacer{} เอวเมว โข ภิกฺขเว อริยสาวกสฺส ทิฏฺฐิสมฺปนฺนสฺส ปุคฺคลสฺส อภิสเมตาวิโน เอตเทว พหุตรํ ทุกฺขํ, ยทิทํ ปริกฺขีณํ ปริยาทินฺนํ, อปฺปมตฺตกํ อวสิฏฺฐํ, สงฺขมฺปิ น อุเปติ อุปนิธมฺปิ น อุเปติ กลภาคมฺปิ น อุเปติ, ปุริมํ ทุกฺขกฺขนฺธํ ปริกฺขีณํ ปริยาทินฺนํ อุปนิธาย, ยทิทํ สตฺตกฺขตฺตุปรมตา.\csromanspacer{} โย “อิทํ ทุกฺขนฺ”ติ ยถาภูตํ ปชานาติ ฯเปฯ “อยํ ทุกฺขนิโรธคามินี ปฏิปทา”ติ ยถาภูตํ ปชานาติ.}

\prose{ตสฺมาติห ภิกฺขเว “อิทํ ทุกฺขนฺ”ติ โยโค กรณีโย ฯเปฯ “อยํ ทุกฺขนิโรธคามินี ปฏิปทา”ติ โยโค กรณีโย.\csromanspacer{}\csromanspacer{}ปฐมํ.}

\csromanmatikaanchor{mtk.303}
\csromantocmarkref{subparagraph}{2. โปกฺขรณีสุตฺต}{mtk.303}
\bookmark[dest={mtk.303},level=5]{2. โปกฺขรณีสุตฺต}
\csromanheader{6}{2. โปกฺขรณีสุตฺต}
\proseitem{1122}{เสยฺยถาปิ ภิกฺขเว โปกฺขรณี ปญฺญาสโยชนานิ อายาเมน ปญฺญาสโยชนานิ วิตฺถาเรน ปญฺญาสโยชนานิ อุพฺเพเธน ปุณฺณา อุทกสฺส สมติตฺติกา กากเปยฺยา, ตโต ปุริโส กุสคฺเคน อุทกํ อุทฺธเรยฺย.\csromanspacer{} ตํ กิํ มญฺญถ ภิกฺขเว, กตมํ นุ โข พหุตรํ, ยํ วา กุสคฺเคน อุพฺภตํ, ยํ วา โปกฺขรณิยา อุทกนฺติ.\csromanspacer{} เอตเทว ภนฺเต พหุตรํ ยทิทํ โปกฺขรณิยา อุทกํ, อปฺปมตฺตกํ กุสคฺเคน อุทกํ อุพฺภตํ, สงฺขมฺปิ น อุเปติ อุปนิธมฺปิ น อุเปติ กลภาคมฺปิ น อุเปติ, โปกฺขรณิยา อุทกํ อุปนิธาย กุสคฺเคน อุทกํ อุพฺภตนฺติ.\csromanspacer{} เอวเมว โข ภิกฺขเว อริยสาวกสฺส ฯเปฯ โยโค กรณีโยติ.\csromanspacer{}\csromanspacer{}ทุติยํ.}

\csromanheader{2}{12. สจฺจสํยุตฺต}
\csromanheader{2}{3. ปฐมสํเภชฺชสุตฺต}
\csromanmatikaanchor{mtk.304}
\csromanmatikaanchor{mtk.305}
\csromantocmarkref{subparagraph}{3-4. สํเภชฺชสุตฺตานิ}{mtk.304}
\bookmark[dest={mtk.304},level=5]{3-4. สํเภชฺชสุตฺตานิ}
\csromantocmarkref{subparagraph}{5-6. มหาปถวีสุตฺตานิ}{mtk.305}
\bookmark[dest={mtk.305},level=5]{5-6. มหาปถวีสุตฺตานิ}
\proseitem{1123}{\csromanfolio{401}เสยฺยถาปิ ภิกฺขเว ยตฺถิมา มหานทิโย สํสนฺทนฺติ สเมนฺติ.\csromanspacer{} เสยฺยถิทํ, คงฺคา ยมุนา อจิรวตี สรภู มหี, ตโต ปุริโส เทฺว วา ตีณิ วา อุทกผุสิตานิ อุทฺธเรยฺย.\csromanspacer{} ตํ กิํ มญฺญถ ภิกฺขเว, กตมํ นุ โข พหุตรํ, ยานิ เทฺว วา ตีณิ วา อุทกผุสิตานิ อุพฺภตานิ, ยํ วา สํเภชฺช-อุทกนฺติ.\csromanspacer{} เอตเทว ภนฺเต พหุตรํ ยทิทํ สํเภชฺช-อุทกํ, อปฺปมตฺตกานิ เทฺว วา ตีณิ วา อุทกผุสิตานิ อุพฺภตานิ, สงฺขมฺปิ น อุเปนฺติ อุปนิธมฺปิ น อุเปนฺติ กลภาคมฺปิ น อุเปนฺติ, สํเภชฺช-อุทกํ อุปนิธาย เทฺว วา ตีณิ วา อุทกผุสิตานิ อุพฺภตานีติ.\csromanspacer{} เอวเมว โข ภิกฺขเว อริยสาวกสฺส ฯเปฯ โยโค กรณีโยติ.\csromanspacer{}\csromanspacer{}ตติยํ.}

\csromanheader{2}{4. ทุติยสํเภชฺชสุตฺต}
\proseitem{1124}{เสยฺยถาปิ ภิกฺขเว ยตฺถิมา มหานทิโย สํสนฺทนฺติ สเมนฺติ.\csromanspacer{} เสยฺยถิทํ, คงฺคา ยมุนา อจิรวตี สรภู มหี, ตํ อุทกํ ปริกฺขยํ ปริยาทานํ คจฺเฉยฺย ฐเปตฺวา เทฺว วา ตีณิ วา อุทกผุสิตานิ.\csromanspacer{} ตํ กิํ มญฺญถ ภิกฺขเว, กตมํ นุ โข พหุตรํ, ยํ วา สํเภชฺช-อุทกํ ปริกฺขีณํ ปริยาทินฺนํ, ยานิ เทฺว วา ตีณิ วา อุทกผุสิตานิ อวสิฏฺฐานีติ.\csromanspacer{} เอตเทว ภนฺเต พหุตรํ สํเภชฺช-อุทกํ, ยทิทํ ปริกฺขีณํ ปริยาทินฺนํ, อปฺปมตฺตกานิ เทฺว วา ตีณิ วา อุทกผุสิตานิ อวสิฏฺฐานิ, สงฺขมฺปิ น อุเปนฺติ อุปนิธมฺปิ น อุเปนฺติ กลภาคมฺปิ น อุเปนฺติ, สํเภชฺช-อุทกํ ปริกฺขีณํ ปริยาทินฺนํ อุปนิธาย เทฺว วา ตีณิ วา อุทกผุสิตานิ อวสิฏฺฐานีติ.\csromanspacer{} เอวเมว โข ภิกฺขเว อริยสาวกสฺส ฯเปฯ โยโค กรณีโยติ.\csromanspacer{}\csromanspacer{}จตุตฺถํ.}

\csromanheader{2}{5. ปฐมมหาปถวีสุตฺต}
\proseitem{1125}{เสยฺยถาปิ ภิกฺขเว ปุริโส มหาปถวิยา สตฺต โกลฏฺฐิมตฺติโย คุฬิกา อุปนิกฺขิเปยฺย.\csromanspacer{} ตํ กิํ มญฺญถ ภิกฺขเว, กตมํ นุ โข พหุตรํ, ยา วา สตฺต โกลฏฺฐิมตฺติโย คุฬิกา อุปนิกฺขิตฺตา, อยํ วา มหาปถวีติ.\csromanspacer{} เอตเทว ภนฺเต พหุตรํ ยทิทํ มหาปถวี, อปฺปมตฺติกา สตฺต โกลฏฺฐิมตฺติโย คุฬิกา อุปนิกฺขิตฺตา, สงฺขมฺปิ น อุเปนฺติ อุปนิธมฺปิ น อุเปนฺติ กลภาคมฺปิ น อุเปนฺติ, มหาปถวิํ อุปนิธาย สตฺต โกลฏฺฐิมตฺติโย คุฬิกา อุปนิกฺขิตฺตาติ.\csromanspacer{} เอวเมว โข ภิกฺขเว อริยสาวกสฺส ฯเปฯ โยโค กรณีโยติ.\csromanspacer{}\csromanspacer{}ปญฺจมํ.}

\csromanheader{2}{6. ทุติยมหาปถวีสุตฺต}
\csromanmatikaanchor{mtk.306}
\csromantocmarkref{subparagraph}{7-8. มหาสมุทฺทสุตฺตานิ}{mtk.306}
\bookmark[dest={mtk.306},level=5]{7-8. มหาสมุทฺทสุตฺตานิ}
\proseitem{1126}{\csromanfolio{402}เสยฺยถาปิ ภิกฺขเว มหาปถวี ปริกฺขยํ ปริยาทานํ คจฺเฉยฺย ฐเปตฺวา สตฺต โกลฏฺฐิมตฺติโย คุฬิกา.\csromanspacer{} ตํ กิํ มญฺญถ ภิกฺขเว, กตมํ นุ โข พหุตรํ, ยํ วา มหาปถวิยา ปริกฺขีณํ ปริยาทินฺนํ, ยา วา สตฺต โกลฏฺฐิมตฺติโย คุฬิกา อวสิฏฺฐาติ.\csromanspacer{} เอตเทว ภนฺเต พหุตรํ มหาปถวิยา ยทิทํ ปริกฺขีณํ ปริยาทินฺนํ, อปฺปมตฺติกา สตฺต โกลฏฺฐิมตฺติโย คุฬิกา อวสิฏฺฐา, สงฺขมฺปิ น อุเปนฺติ อุปนิธมฺปิ น อุเปนฺติ กลภาคมฺปิ น อุเปนฺติ, มหาปถวี ปริกฺขีณํ ปริยาทินฺนํ อุปนิธาย สตฺต โกลฏฺฐิมตฺติโย คุฬิกา อวสิฏฺฐาติ.\csromanspacer{} เอวเมว โข ภิกฺขเว อริยสาวกสฺส ฯเปฯ โยโค กรณีโยติ.\csromanspacer{}\csromanspacer{}ฉฏฺฐํ.}

\csromanheader{2}{7. ปฐมมหาสมุทฺทสุตฺต}
\proseitem{1127}{เสยฺยถาปิ ภิกฺขเว ปุริโส มหาสมุทฺทโต เทฺว วา ตีณิ วา อุทกผุสิตานิ อุทฺธริตานิ.\csromanspacer{} ตํ กิํ มญฺญถ ภิกฺขเว, กตมํ นุ โข พหุตรํ, ยานิ เทฺว วา ตีณิ วา อุทกผุสิตานิ อุพฺภตานิ, ยํ วา มหาสมุทฺเท อุทกนฺติ.\csromanspacer{} เอตเทว ภนฺเต พหุตรํ ยทิทํ มหาสมุทฺเท อุทกํ, อปฺปมตฺตกานิ เทฺว วา ตีณิ วา อุทกผุสิตานิ อุพฺภตานิ, สงฺขมฺปิ น อุเปนฺติ อุปนิธมฺปิ น อุเปนฺติ กลภาคมฺปิ น อุเปนฺติ, มหาสมุทฺเท อุทกํ อุปนิธาย เทฺว วา ตีณิ วา อุทกผุสิตานิ อุพฺภตานีติ.\csromanspacer{} เอวเมว โข ภิกฺขเว อริยสาวกสฺส ฯเปฯ โยโค กรณีโยติ.\csromanspacer{}\csromanspacer{}สตฺตมํ.}

\csromanheader{2}{8. ทุติยมหาสมุทฺทสุตฺต}
\proseitem{1128}{เสยฺยถาปิ ภิกฺขเว มหาสมุทฺเท อุทกํ ปริกฺขยํ\footnote{มหาสมุทฺโท ปริกฺขยํ (สี, สฺยา, กํ) สํ 1. 355 ปิฏฺเฐปิ.} ปริยาทานํ คจฺเฉยฺย ฐเปตฺวา เทฺว วา ตีณิ วา อุทกผุสิตานิ.\csromanspacer{} ตํ กิํ มญฺญถ ภิกฺขเว, กตมํ นุ โข พหุตรํ, ยํ วา มหาสมุทฺเท อุทกํ ปริกฺขีณํ ปริยาทินฺนํ, ยานิ เทฺว วา ตีณิ วา อุทกผุสิตานิ อวสิฏฺฐานีติ.\csromanspacer{} เอตเทว ภนฺเต พหุตรํ มหาสมุทฺเท อุทกํ ยทิทํ ปริกฺขีณํ ปริยาทินฺนํ, อปฺปมตฺตกานิ เทฺว วา ตีณิ วา อุทกผุสิตานิ อวสิฏฺฐานิ, สงฺขมฺปิ น อุเปนฺติ อุปนิธมฺปิ น อุเปนฺติ กลภาคมฺปิ น อุเปนฺติ, มหาสมุทฺเท อุทกํ}

\csromanheader{2}{12. สจฺจสํยุตฺต}
\csromanmatikaanchor{mtk.307}
\csromantocmarkref{subparagraph}{9-10. ปพฺพตูปมสุตฺตานิ}{mtk.307}
\bookmark[dest={mtk.307},level=5]{9-10. ปพฺพตูปมสุตฺตานิ}
\prose{\csromanfolio{403}ปริกฺขีณํ ปริยาทินฺนํ อุปนิธาย เทฺว วา ตีณิ วา อุทกผุสิตานิ อวสิฏฺฐานีติ.\csromanspacer{} เอวเมว โข ภิกฺขเว อริยสาวกสฺส ฯเปฯ โยโค กรณีโยติ.\csromanspacer{}\csromanspacer{}อฏฺฐมํ.}

\csromanheader{2}{9. ปฐมปพฺพตูปมสุตฺต}
\proseitem{1129}{เสยฺยถาปิ ภิกฺขเว ปุริโส หิมวโต ปพฺพตราชสฺส สตฺต สาสปมตฺติโย ปาสาณสกฺขรา อุปนิกฺขิเปยฺย.\csromanspacer{} ตํ กิํ มญฺญถ ภิกฺขเว, กตมํ นุ โข พหุตรํ, ยา วา สตฺต สาสปมตฺติโย ปาสาณสกฺขรา อุปนิกฺขิตฺตา, อยํ วา หิมวา ปพฺพตราชาติ.\csromanspacer{} เอตเทว ภนฺเต พหุตรํ ยทิทํ หิมวา ปพฺพตราชา, อปฺปมตฺติกา สตฺต สาสปมตฺติโย ปาสาณสกฺขรา อุปนิกฺขิตฺตา, สงฺขมฺปิ น อุเปนฺติ อุปนิธมฺปิ น อุเปนฺติ กลภาคมฺปิ น อุเปนฺติ, หิมวนฺตํ ปพฺพตราชานํ อุปนิธาย สตฺต สาสปมตฺติโย ปาสาณสกฺขรา อุปนิกฺขิตฺตาติ.\csromanspacer{} เอวเมว โข ภิกฺขเว อริยสาวกสฺส ฯเปฯ โยโค กรณีโยติ.\csromanspacer{}\csromanspacer{}นวมํ.}

\csromanheader{2}{10. ทุติยปพฺพตูปมสุตฺต}
\proseitem{1130}{เสยฺยถาปิ ภิกฺขเว หิมวา ปพฺพตราชา ปริกฺขยํ ปริยาทานํ คจฺเฉยฺย ฐเปตฺวา สตฺต สาสปมตฺติโย ปาสาณสกฺขรา.\csromanspacer{} ตํ กิํ มญฺญถ ภิกฺขเว, กตมํ นุ โข พหุตรํ, ยํ วา หิมวโต ปพฺพตราชสฺส ปริกฺขีณํ ปริยาทินฺนํ, ยา วา สตฺต สาสปมตฺติโย ปาสาณสกฺขรา อวสิฏฺฐาติ.\csromanspacer{} เอตเทว ภนฺเต พหุตรํ หิมวโต ปพฺพตราชสฺส ยทิทํ ปริกฺขีณํ ปริยาทินฺนํ, อปฺปมตฺติกา สตฺต สาสปมตฺติโย ปาสาณสกฺขรา อวสิฏฺฐา, สงฺขมฺปิ น อุเปนฺติ อุปนิธมฺปิ น อุเปนฺติ กลภาคมฺปิ น อุเปนฺติ, หิมวโต ปพฺพตราชสฺส ปริกฺขีณํ ปริยาทินฺนํ, อุปนิธาย สตฺต สาสปมตฺติโย ปาสาณสกฺขรา อวสิฏฺฐาติ.\csromanspacer{} เอวเมว โข ภิกฺขเว อริยสาวกสฺส ทิฏฺฐิสมฺปนฺนสฺส ปุคฺคลสฺส อภิสเมตาวิโน เอตเทว พหุตรํ ทุกฺขํ, ยทิทํ ปริกฺขีณํ ปริยาทินฺนํ, อปฺปมตฺตกํ อวสิฏฺฐํ, สงฺขมฺปิ น อุเปติ อุปนิธมฺปิ น อุเปติ กลภาคมฺปิ น อุเปติ, ปุริมํ ทุกฺขกฺขนฺธํ ปริกฺขีณํ ปริยาทินฺนํ อุปนิธาย, ยทิทํ สตฺตกฺขตฺตุปรมตา. “โย อิทํ ทุกฺขนฺ”ติ ยถาภูตํ ปชานาติ ฯเปฯ “อยํ ทุกฺขนิโรธคามินี ปฏิปทา”ติ ยถาภูตํ ปชานาติ.}

\csromanmatikaanchor{mtk.308}
\csromanmatikaanchor{mtk.309}
\csromantocmarknopage{subparagraph}{7. ปฐม-อามกธญฺญเปยฺยาลวคฺค}{mtk.308}
\bookmark[dest={mtk.308},level=5]{7. ปฐม-อามกธญฺญเปยฺยาลวคฺค}
\csromantocmarkref{subparagraph}{1-10. อญฺญตฺรสุตฺตาทีนิ}{mtk.309}
\bookmark[dest={mtk.309},level=5]{1-10. อญฺญตฺรสุตฺตาทีนิ}
\prose{\csromanfolio{404}ตสฺมาติห ภิกฺขเว “อิทํ ทุกฺขนฺ”ติ โยโค กรณีโย ฯเปฯ “อยํ ทุกฺขนิโรธคามินี ปฏิปทา”ติ โยโค กรณีโยติ.\csromanspacer{}\csromanspacer{}ทสมํ.\csromanspacer{} อภิสมยวคฺโค ฉฏฺโฐ.}

\titlehead{\textbf{ตสฺสุทฺทานํ}}
\csromangathameasure{นขสิขา โปกฺขรณี, สํเภชฺช อปเร ทุเว. \\ ปถวี เทฺว สมุทฺทา เทฺว, เทฺวมา จ ปพฺพตูปมาติ.}
\csromangathagroup{\csromangathastanza{นขสิขา โปกฺขรณี, สํเภชฺช อปเร ทุเว. \\ ปถวี เทฺว สมุทฺทา เทฺว, เทฺวมา จ ปพฺพตูปมาติ.}}

\csromanheader{6}{7. ปฐม-อามกธญฺญเปยฺยาลวคฺค}
\csromanheader{2}{1. อญฺญตฺรสุตฺต}
\proseitem{1131}{อถ โข ภควา ปริตฺตํ นขสิขายํ ปํสุํ อาโรเปตฺวา ภิกฺขู อามนฺเตสิ “ตํ กิํ มญฺญถ ภิกฺขเว, กตมํ นุ โข พหุตรํ, โย วายํ มยา ปริตฺโต นขสิขายํ ปํสุ อาโรปิโต, อยํ วา มหาปถวี”ติ.\csromanspacer{} เอตเทว ภนฺเต พหุตรํ ยทิทํ มหาปถวี, อปฺปมตฺตกายํ ภควตา ปริตฺโต นขสิขายํ ปํสุ อาโรปิโต, สงฺขมฺปิ น อุเปติ อุปนิธมฺปิ น อุเปติ กลภาคมฺปิ น อุเปติ, มหาปถวิํ อุปนิธาย ภควตา ปริตฺโต นขสิขายํ ปํสุ อาโรปิโตติ.\csromanspacer{} เอวเมว โข ภิกฺขเว อปฺปมตฺตกา เต สตฺตา, เย มนุสฺเสสุ ปจฺจาชายนฺติ.\csromanspacer{} อถ โข เอเตว พหุตรา สตฺตา เย อญฺญตฺร มนุสฺเสหิ\footnote{มนุสฺเสสุ (อิ, ก)} ปจฺจาชายนฺติ.\csromanspacer{} ตํ กิสฺส เหตุ, อทิฏฺฐตฺตา ภิกฺขเว จตุนฺนํ อริยสจฺจานํ.\csromanspacer{} กตเมสํ จตุนฺนํ, ทุกฺขสฺส อริยสจฺจสฺส ฯเปฯ ทุกฺขนิโรธคามินิยา ปฏิปทาย อริยสจฺจสฺส.}

\prose{ตสฺมาติห ภิกฺขเว “อิทํ ทุกฺขนฺ”ติ โยโค กรณีโย ฯเปฯ “อยํ ทุกฺขนิโรธคามินี ปฏิปทา”ติ โยโค กรณีโยติ.\csromanspacer{}\csromanspacer{}ปฐมํ.}

\csromanheader{2}{2. ปจฺจนฺตสุตฺต}
\proseitem{1132}{อถ โข ภควา ปริตฺตํ นขสิขายํ ปํสุํ อาโรเปตฺวา ภิกฺขู อามนฺเตสิ “ตํ กิํ มญฺญถ ภิกฺขเว, กตมํ นุ โข พหุตรํ, โย}

\csromanheader{2}{12. สจฺจสํยุตฺต}
\prose{\csromanfolio{405}วายํ มยา ปริตฺโต นขสิขายํ ปํสุ อาโรปิโต, อยํ วา มหาปถวี”ติ.\csromanspacer{} เอตเทว ภนฺเต พหุตรํ ยทิทํ มหาปถวี, อปฺปมตฺตกายํ ภควตา ปริตฺโต นขสิขายํ ปํสุ อาโรปิโต, สงฺขมฺปิ น อุเปติ อุปนิธมฺปิ น อุเปติ กลภาคมฺปิ น อุเปติ, มหาปถวิํ อุปนิธาย ภควตา ปริตฺโต นขสิขายํ ปํสุ อาโรปิโตติ.}

\prose{เอวเมว โข ภิกฺขเว อปฺปมตฺตกา เต สตฺตา, เย มชฺฌิเมสุ ชนปเทสุ ปจฺจาชายนฺติ, อถ โข เอเตว พหุตรา สตฺตา, เย ปจฺจนฺติเมสุ ชนปเทสุ ปจฺจาชายนฺติ อวิญฺญาตาเรสุ มิลกฺเขสุ\footnote{มิลกฺขูสุ (สฺยา, กํ, ก)} ฯเปฯ.\csromanspacer{}\csromanspacer{}ทุติยํ.}

\csromanheader{2}{3. ปญฺญาสุตฺต}
\proseitem{1133}{เอวเมว โข ภิกฺขเว อปฺปกา เต สตฺตา, เย ปน อริเยน ปญฺญาจกฺขุนา สมนฺนาคตา, อถ โข เอเตว พหุตรา สตฺตา, เย อวิชฺชาคตา สมฺมุฬฺหา ฯเปฯ.\csromanspacer{}\csromanspacer{}ตติยํ.}

\csromanheader{2}{4. สุราเมรยสุตฺต}
\proseitem{1134}{เอวเมว โข ภิกฺขเว อปฺปกา เต สตฺตา, เย สุราเมรยมชฺชปฺปมาทฏฺฐานา ปฏิวิรตา, อถ โข เอเตว พหุตรา สตฺตา, เย สุราเมรยมชฺชปฺปมาทฏฺฐานา ปฏิวิรตา ฯเปฯ.\csromanspacer{}\csromanspacer{}จตุตฺถํ.}

\csromanheader{2}{5. โอทกสุตฺต}
\proseitem{1135}{เอวเมว โข ภิกฺขเว อปฺปกา เต สตฺตา, เย ถลชา, อถ โข เอเตว พหุตรา สตฺตา, เย อุทกชา.\csromanspacer{} ตํ กิสฺส เหตุ ฯเปฯ.\csromanspacer{}\csromanspacer{}ปญฺจมํ.}

\csromanheader{2}{6. มตฺเตยฺยสุตฺต}
\proseitem{1136}{เอวเมว โข ภิกฺขเว อปฺปกา เต สตฺตา, เย มตฺเตยฺยา, อถ โข เอเตว พหุตรา สตฺตา, เย อมตฺเตยฺยา ฯเปฯ.\csromanspacer{}\csromanspacer{}ฉฏฺฐํ.}

\csromanheader{2}{7. เปตฺเตยฺยสุตฺต}
\csromanmatikaanchor{mtk.310}
\csromanmatikaanchor{mtk.311}
\csromantocmarknopage{subparagraph}{8. ทุติย-อามกธญฺญเปยฺยาลวคฺค}{mtk.310}
\bookmark[dest={mtk.310},level=5]{8. ทุติย-อามกธญฺญเปยฺยาลวคฺค}
\csromantocmarkref{subparagraph}{1-10. ปาณาติปาตสุตฺตาทีนิ}{mtk.311}
\bookmark[dest={mtk.311},level=5]{1-10. ปาณาติปาตสุตฺตาทีนิ}
\csromangathameasure{เอวเมว โข ภิกฺขเว อปฺปกา เต สตฺตา, เย เปตฺเตยฺยา, อถ โข เอเตว พหุตรา สตฺตา, เย อเปตฺเตยฺยา ฯเปฯ.สตฺตมํ.}
\csromangathagroup{\csromangathastanza{\csromanfolio{406}เอวเมว โข ภิกฺขเว อปฺปกา เต สตฺตา, เย เปตฺเตยฺยา, อถ โข เอเตว พหุตรา สตฺตา, เย อเปตฺเตยฺยา ฯเปฯ.\csromanspacer{}\csromanspacer{}สตฺตมํ.}}

\csromanheader{2}{8. สามญฺญสุตฺต}
\csromangathameasure{เอวเมว โข ภิกฺขเว อปฺปกา เต สตฺตา, เย สามญฺญา, อถ โข เอเตว พหุตรา สตฺตา, เย อสามญฺญา ฯเปฯ.อฏฺฐมํ.}
\csromangathagroup{\csromangathastanza{เอวเมว โข ภิกฺขเว อปฺปกา เต สตฺตา, เย สามญฺญา, อถ โข เอเตว พหุตรา สตฺตา, เย อสามญฺญา ฯเปฯ.\csromanspacer{}\csromanspacer{}อฏฺฐมํ.}}

\csromanheader{2}{9. พฺรหฺมญฺญสุตฺต}
\csromangathameasure{เอวเมว โข ภิกฺขเว อปฺปกา เต สตฺตา, เย พฺรหฺมญฺญา อถ โข เอเตว พหุตรา สตฺตา, เย อพฺรหฺมญฺญา ฯเปฯ.นวมํ.}
\csromangathagroup{\csromangathastanza{เอวเมว โข ภิกฺขเว อปฺปกา เต สตฺตา, เย พฺรหฺมญฺญา อถ โข เอเตว พหุตรา สตฺตา, เย อพฺรหฺมญฺญา ฯเปฯ.\csromanspacer{}\csromanspacer{}นวมํ.}}

\csromanheader{2}{10. ปจายิกสุตฺต}
\proseitem{1140}{เอวเมว โข ภิกฺขเว อปฺปกา เต สตฺตา, เย กุเล เชฏฺฐาปจายิโน.\csromanspacer{} อถ โข เอเตว พหุตรา สตฺตา, เย กุเล อเชฏฺฐาปจายิโนติ\footnote{อกุเล เชฏฺฐาปจายิโนติ (สฺยา, กํ)}.\csromanspacer{}\csromanspacer{}ทสมํ.\csromanspacer{} ปฐม-อามกธญฺญเปยฺยาลวคฺโค สตฺตโม.}

\titlehead{\textbf{ตสฺสุทฺทานํ}}
\csromangathameasure{อญฺญตฺร ปจฺจนฺตํ ปญฺญา, สุราเมรย-โอทกา. \\ มตฺเตยฺย เปตฺเตยฺยา จาปิ, สามญฺญํ พฺรหฺมปจายิกนฺติ.}
\csromangathagroup{\csromangathastanza{อญฺญตฺร ปจฺจนฺตํ ปญฺญา, สุราเมรย-โอทกา. \\ มตฺเตยฺย เปตฺเตยฺยา จาปิ, สามญฺญํ พฺรหฺมปจายิกนฺติ.}}

\csromanheader{6}{8. ทุติย-อามกธญฺญเปยฺยาลวคฺค}
\csromanheader{2}{1. ปาณาติปาตสุตฺต}
\csromangathameasure{เอวเมว โข ภิกฺขเว อปฺปกา เต สตฺตา, เย ปาณาติปาตา ปฏิวิรตา. อถ โข เอเตว พหุตรา สตฺตา, เย ปาณาติปาตา อปฺปฏิวิรตา. ตํ กิสฺส เหตุ ฯเปฯ.ปฐมํ.}
\csromangathagroup{\csromangathastanza{เอวเมว โข ภิกฺขเว อปฺปกา เต สตฺตา, เย ปาณาติปาตา ปฏิวิรตา.\csromanspacer{} อถ โข เอเตว พหุตรา สตฺตา, เย ปาณาติปาตา อปฺปฏิวิรตา.\csromanspacer{} ตํ กิสฺส เหตุ ฯเปฯ.\csromanspacer{}\csromanspacer{}ปฐมํ.}}

\csromanheader{2}{12. สจฺจสํยุตฺต}
\csromanheader{2}{2. อทินฺนาทานสุตฺต}
\proseitem{1142}{\csromanfolio{407}เอวเมว โข ภิกฺขเว อปฺปกา เต สตฺตา, เย อทินฺนาทานา ปฏิวิรตา.\csromanspacer{} อถ โข เอเตว พหุตรา สตฺตา, เย อทินฺนาทานา อปฺปฏิวิรตา ฯเปฯ.\csromanspacer{}\csromanspacer{}ทุติยํ.}

\csromanheader{2}{3. กาเมสุมิจฺฉาจารสุตฺต}
\proseitem{1143}{เอวเมว โข ภิกฺขเว อปฺปกา เต สตฺตา, เย กาเมสุมิจฺฉาจารา ปฏิวิรตา.\csromanspacer{} อถ โข เอเตว พหุตรา สตฺตา, เย กาเมสุ มิจฺฉาจารา อปฺปฏิวิรตา ฯเปฯ.\csromanspacer{}\csromanspacer{}ตติยํ.}

\csromanheader{2}{4. มุสาวาทสุตฺต}
\proseitem{1144}{เอวเมว โข ภิกฺขเว อปฺปกา เต สตฺตา, เย มุสาวาทา ปฏิวิรตา.\csromanspacer{} อถ โข เอเตว พหุตรา สตฺตา, เย มุสาวาทา อปฺปฏิวิรตา ฯเปฯ.\csromanspacer{}\csromanspacer{}จตุตฺถํ.}

\csromanheader{2}{5. เปสุญฺญสุตฺต}
\proseitem{1145}{เอวเมว โข ภิกฺขเว อปฺปกา เต สตฺตา, เย ปิสุณาย วาจาย ปฏิวิรตา.\csromanspacer{} อถ โข เอเตว พหุตรา สตฺตา, เย ปิสุณาย วาจาย อปฺปฏิวิรตา ฯเปฯ.\csromanspacer{}\csromanspacer{}ปญฺจมํ.}

\csromanheader{2}{6. ผรุสวาจาสุตฺต}
\proseitem{1146}{เอวเมว โข ภิกฺขเว อปฺปกา เต สตฺตา, เย ผรุสาย วาจาย ปฏิวิรตา.\csromanspacer{} อถ โข เอเตว พหุตรา สตฺตา, เย ผรุสาย วาจาย อปฺปฏิวิรตา ฯเปฯ.\csromanspacer{}\csromanspacer{}ฉฏฺฐํ.}

\csromanheader{2}{7. สมฺผปฺปลาปสุตฺต}
\proseitem{1147}{เอวเมว โข ภิกฺขเว อปฺปกา เต สตฺตา, เย สมฺผปฺปลาปา ปฏิวิรตา.\csromanspacer{} อถ โข เอเตว พหุตรา สตฺตา, เย สมฺผปฺปลาปา อปฺปฏิวิรตา ฯเปฯ.\csromanspacer{}\csromanspacer{}สตฺตมํ.}

\csromanheader{2}{8. พีชคามสุตฺต}
\csromanmatikaanchor{mtk.312}
\csromanmatikaanchor{mtk.313}
\csromantocmarknopage{subparagraph}{9. ตติย-อามกธญฺญเปยฺยาลวคฺค}{mtk.312}
\bookmark[dest={mtk.312},level=5]{9. ตติย-อามกธญฺญเปยฺยาลวคฺค}
\csromantocmarkref{subparagraph}{1-10. นจฺจคีตสุตฺตาทีนิ}{mtk.313}
\bookmark[dest={mtk.313},level=5]{1-10. นจฺจคีตสุตฺตาทีนิ}
\proseitem{1148}{\csromanfolio{408}เอวเมว โข ภิกฺขเว อปฺปกา เต สตฺตา, เย พีชคามภูตคามสมารมฺภา\footnote{พีชคามภูตคามสมารพฺภา (ก)} ปฏิวิรตา.\csromanspacer{} อถ โข เอเตว พหุตรา สตฺตา, เย พีชคามภูตคามสมารมฺภา อปฺปฏิวิรตา ฯเปฯ.\csromanspacer{}\csromanspacer{}อฏฺฐมํ.}

\csromanheader{2}{9. วิกาลโภชนสุตฺต}
\proseitem{1149}{เอวเมว โข ภิกฺขเว อปฺปกา เต สตฺตา, เย วิกาลโภชนา ปฏิวิรตา.\csromanspacer{} อถ โข เอเตว พหุตรา สตฺตา, เย วิกาลโภชนา อปฺปฏิวิรตา ฯเปฯ.\csromanspacer{}\csromanspacer{}นวมํ.}

\csromanheader{2}{10. คนฺธวิเลปนสุตฺต}
\proseitem{1150}{เอวเมว โข ภิกฺขเว อปฺปกา เต สตฺตา, เย มาลาคนฺธวิเลปนธารณมณฺฑนวิภูสนฏฺฐานา ปฏิวิรตา.\csromanspacer{} อถ โข เอเตว พหุตรา สตฺตา, เย มาลาคนฺธวิเลปนธารมณฺฑนวิภูสนฏฺฐานา อปฺปฏิวิรตา.\csromanspacer{} ทสมํ.\csromanspacer{} ทุติย-อามกธญฺญเปยฺยาลวคฺโค อฏฺฐโม.}

\titlehead{\textbf{ตสฺสุทฺทานํ}}
\csromangathameasure{ปาณํ อทินฺนํ กาเมสุ, มุสาวาทญฺจ เปสุญฺญํ. \\ ผรุสํ สมฺผปฺปลาปํ, พีชญฺจ วิกาลํ คนฺธนฺติ.}
\csromangathagroup{\csromangathastanza{ปาณํ อทินฺนํ กาเมสุ, มุสาวาทญฺจ เปสุญฺญํ. \\ ผรุสํ สมฺผปฺปลาปํ, พีชญฺจ วิกาลํ คนฺธนฺติ.}}

\csromanheader{6}{9. ตติย-อามกธญฺญเปยฺยาลวคฺค}
\csromanheader{2}{1. นจฺจคีตสุตฺต}
\proseitem{1151}{เอวเมว โข ภิกฺขเว อปฺปกา เต สตฺตา, เย นจฺจคีตวาทิตวิสูกทสฺสนา ปฏิวิรตา.\csromanspacer{} อถ โข เอเตว พหุตรา สตฺตา, เย นจฺจ คีต วาทิต วิสูก ทสฺสนา อปฺปฏิวิรตา.\csromanspacer{} ตํ กิสฺส เหตุ ฯเปฯ.\csromanspacer{}\csromanspacer{}ปฐมํ.}

\csromanheader{2}{12. สจฺจสํยุตฺต}
\csromanheader{2}{2. อุจฺจาสยนสุตฺต}
\proseitem{1152}{\csromanfolio{409}เอวเมว โข ภิกฺขเว อปฺปกา เต สตฺตา, เย อุจฺจาสยนมหาสยนา ปฏิวิรตา.\csromanspacer{} อถ โข เอเตว พหุตรา สตฺตา, เย อุจฺจาสยนมหาสยนา อปฺปฏิวิรตา ฯเปฯ.\csromanspacer{}\csromanspacer{}ทุติยํ.}

\csromanheader{2}{3. ชาตรูปรชตสุตฺต}
\proseitem{1153}{เอวเมว โข ภิกฺขเว อปฺปกา เต สตฺตา, เย ชาตรูปรชตปฏิคฺคหณา ปฏิวิรตา.\csromanspacer{} อถ โข เอเตว พหุตรา สตฺตา, เย ชาตรูปรชตปฏิคฺคหณา อปฺปฏิวิรตา ฯเปฯ.\csromanspacer{}\csromanspacer{}ตติยํ.}

\csromanheader{2}{4. อามกธญฺญสุตฺต}
\proseitem{1154}{เอวเมว โข ภิกฺขเว อปฺปกา เต สตฺตา, เย อามกธญฺญปฏิคฺคหณา ปฏิวิรตา.\csromanspacer{} อถ โข เอเตว พหุตรา สตฺตา, เย อามกธญฺญปฏิคฺคหณา อปฺปฏิวิรตา ฯเปฯ.\csromanspacer{}\csromanspacer{}จตุตฺถํ.}

\csromanheader{2}{5. อามกมํสสุตฺต}
\proseitem{1155}{เอวเมว โข ภิกฺขเว อปฺปกา เต สตฺตา, เย อามกมํสปฏิคฺคหณา ปฏิวิรตา.\csromanspacer{} อถ โข เอเตว พหุตรา สตฺตา, เย อามกมํสปฏิคฺคหณา อปฺปฏิวิรตา ฯเปฯ.\csromanspacer{}\csromanspacer{}ปญฺจมํ.}

\csromanheader{2}{6. กุมาริกสุตฺต}
\proseitem{1156}{เอวเมว โข ภิกฺขเว อปฺปกา เต สตฺตา, เย อิตฺถิกุมาริกปฏิคฺคหณา\footnote{อิตฺถิกุมาริกาปฏิคฺคหณา (ก)} ปฏิวิรตา.\csromanspacer{} อถ โข เอเตว พหุตรา สตฺตา, เย อิตฺถิกุมาริกปฏิคฺคหณา อปฺปฏิวิรตา ฯเปฯ.\csromanspacer{}\csromanspacer{}ฉฏฺฐํ.}

\csromanheader{2}{7. ทาสิทาสสุตฺต}
\proseitem{1157}{เอวเมว โข ภิกฺขเว อปฺปกา เต สตฺตา, เย ทาสิทาสปฏิคฺคหณา\footnote{ทาสีทาสปฏิคฺคหณา (สฺยา, กํ, อิ)} ปฏิวิรตา.\csromanspacer{} อถ โข เอเตว พหุตรา สตฺตา, เย ทาสิทาสปฏิคฺคหณา อปฺปฏิวิรตา ฯเปฯ.\csromanspacer{}\csromanspacer{}สตฺตมํ.}

\csromanheader{2}{8. อเชฬกสุตฺต}
\csromanmatikaanchor{mtk.314}
\csromanmatikaanchor{mtk.315}
\csromantocmarknopage{subparagraph}{10. จตฺตุตฺถ-อามกธญฺญเปยฺยาลวคฺค}{mtk.314}
\bookmark[dest={mtk.314},level=5]{10. จตฺตุตฺถ-อามกธญฺญเปยฺยาลวคฺค}
\csromantocmarkref{subparagraph}{1-11. เขตฺตวตฺถุสุตฺตาทีนิ}{mtk.315}
\bookmark[dest={mtk.315},level=5]{1-11. เขตฺตวตฺถุสุตฺตาทีนิ}
\proseitem{1158}{\csromanfolio{410}เอวเมว โข ภิกฺขเว อปฺปกา เต สตฺตา, เย อเชฬกปฏิคฺคหณา ปฏิวิรตา.\csromanspacer{} อถ โข เอเตว พหุตรา สตฺตา, เย อเชฬกปฏิคฺคหณา อปฺปฏิวิรตา ฯเปฯ.\csromanspacer{}\csromanspacer{}อฏฺฐมํ.}

\csromanheader{2}{9. กุกฺกุฏสูกรสุตฺต}
\proseitem{1159}{เอวเมว โข ภิกฺขเว อปฺปกา เต สตฺตา, เย กุกฺกุฏสูกรปฏิคฺคหณา ปฏิวิรตา.\csromanspacer{} อถ โข เอเตว พหุตรา สตฺตา, เย กุกฺกุฏสูกรปฏิคฺคหณา อปฺปฏิวิรตา ฯเปฯ.\csromanspacer{}\csromanspacer{}นวมํ.}

\csromanheader{2}{10. หตฺถิควสฺสสุตฺต}
\proseitem{1160}{เอวเมว โข ภิกฺขเว อปฺปกา เต สตฺตา, เย หตฺถิควสฺสวฬวปฏิคฺคหณา\footnote{หตฺถิควสฺสวฬวาปฏิคฺคหณา (สฺยา, กํ, อิ, ก)} ปฏิวิรตา.\csromanspacer{} อถ โข เอเตว พหุตรา สตฺตา, เย หตฺถิควสฺสวฬาวปฏิคฺคหณา อปฺปฏิวิรตา.\csromanspacer{}\csromanspacer{}ทสมํ.\csromanspacer{} ตติย-อามกธญฺญเปยฺยาลวคฺโค นวโม.}

\titlehead{\textbf{ตสฺสุทฺทานํ}}
\csromangathameasure{นจฺจํ สยนํ รชตํ, ธญฺญํ มํสํ กุมาริกา. \\ ทาสี อเชฬกญฺเจว, กุกฺกุฏสูกรหตฺถีติ.}
\csromangathagroup{\csromangathastanza{นจฺจํ สยนํ รชตํ, ธญฺญํ มํสํ กุมาริกา. \\ ทาสี อเชฬกญฺเจว, กุกฺกุฏสูกรหตฺถีติ.}}

\chapterhead{10. จตุตฺถ-อามกธญฺญเปยฺยาลวคฺค}
\csromanheader{2}{1. เขตฺตวตฺถุสุตฺต}
\proseitem{1161}{เอวเมว โข ภิกฺขเว อปฺปกา เต สตฺตา, เย เขตฺตวตฺถุปฏิคฺคหณา ปฏิวิรตา.\csromanspacer{} อถ โข เอเตว พหุตรา สตฺตา, เย เขตฺตวตฺถุปฏิคฺคหณา อปฺปฏิวิรตา ฯเปฯ.\csromanspacer{}\csromanspacer{}ปฐมํ.}

\csromanheader{2}{12. สจฺจสํยุตฺต}
\csromanheader{2}{2. กยวิกฺกยสุตฺต}
\proseitem{1162}{\csromanfolio{411}เอวเมว โข ภิกฺขเว อปฺปกา เต สตฺตา, เย กยวิกฺกยา ปฏิวิรตา.\csromanspacer{} อถ โข เอเตว พหุตรา สตฺตา, เย กยวิกฺกยา อปฺปฏิวิรตา ฯเปฯ.\csromanspacer{}\csromanspacer{}ทุติยํ.}

\csromanheader{2}{3. ทูเตยฺยสุตฺต}
\proseitem{1163}{เอวเมว โข ภิกฺขเว อปฺปกา เต สตฺตา, เย ทูเตยฺยปหินคมนานุโยคา ปฏิวิรตา.\csromanspacer{} อถ โข เอเตว พหุตรา สตฺตา, เย ทูเตยฺยปหินคมนานุโยคา อปฺปฏิวิรตา ฯเปฯ.\csromanspacer{}\csromanspacer{}ตติยํ.}

\csromanheader{2}{4. ตุลากูฏสุตฺต}
\proseitem{1164}{เอวเมว โข ภิกฺขเว อปฺปกา เต สตฺตา, เย ตุลากูฏกํสกูฏมานกูฏา ปฏิวิรตา.\csromanspacer{} อถ โข เอเตว พหุตรา สตฺตา, เย ตุลากูฏกํสกูฏมานกูฏา อปฺปฏิวิรตา ฯเปฯ.\csromanspacer{}\csromanspacer{}จตุตฺถํ.}

\csromanheader{2}{5. อุกฺโกฏนสุตฺต}
\proseitem{1165}{เอวเมว โข ภิกฺขเว อปฺปกา เต สตฺตา, เย อุกฺโกฏนวญฺจนนิกติสาจิโยคา\footnote{อุกฺโกฏนวญฺจนนิกติสาวิโยคา (สฺยา, กํ, อิ, ก)} ปฏิวิรตา, อถ โข เอเตว พหุตรา สตฺตา, เย อุกฺโกฏนวญฺจนนิกติสาจิโยคา อปฺปฏิวิรตา ฯเปฯ.\csromanspacer{}\csromanspacer{}ปญฺจมํ.}

\csromanheader{2}{6-11. เฉทนาทิสุตฺต}
\proseitem{1166-1171}{เอวเมว โข ภิกฺขเว อปฺปกา เต สตฺตา, เย เฉทนวธพนฺธนวิปราโมส-อาโลปสหสาการา\footnote{สาหสาการา (ก)} ปฏิวิรตา.\csromanspacer{} อถ โข เอเตว พหุตรา สตฺตา, เย เฉทนวธพนฺธนวิปราโมส-อาโลปสหสาการา อปฺปฏิวิรตา.\csromanspacer{} ตํ กิสฺส เหตุ, อทิฏฺฐตฺตา ภิกฺขเว จตุนฺนํ อริยสจฺจานํ.\csromanspacer{} กตเมสํ จตุนฺนํ, ทุกฺขสฺส อริยสจฺจสฺส ฯเปฯ ทุกฺขนิโรธคามินิยา ปฏิปทาย อริยสจฺจสฺส.}

\prose{ตสฺมาติห ภิกฺขเว “อิทํ ทุกฺขนฺ”ติ โยโค กรณีโย ฯเปฯ “อยํ ทุกฺขนิโรธคามินี ปฏิปทา”ติ โยโค กรณีโยติ.\csromanspacer{}\csromanspacer{}เอกาทสมํ.\csromanspacer{} จตุตฺถ-อามกธญฺญเปยฺยาลวคฺโค ทสโม.}

\titlehead{\textbf{ตสฺสุทฺทานํ}}
\csromanmatikaanchor{mtk.316}
\csromanmatikaanchor{mtk.317}
\csromantocmarknopage{subparagraph}{11. ปญฺจคติเปยฺยาลวคฺค}{mtk.316}
\bookmark[dest={mtk.316},level=5]{11. ปญฺจคติเปยฺยาลวคฺค}
\csromantocmarkref{subparagraph}{1-30. มนุสฺสจุตินิรยสุตฺตาทีนิ}{mtk.317}
\bookmark[dest={mtk.317},level=5]{1-30. มนุสฺสจุตินิรยสุตฺตาทีนิ}
\csromangathameasure{เขตฺตํ กายํ ทูเตยฺยญฺจ, ตุลากูฏํ อุกฺโกฏนํ. \\ เฉทนํ วธพนฺธนํ, วิปราโลปํ สาหสนฺติ.}
\csromangathagroup{\csromangathastanza{\csromanfolio{412}เขตฺตํ กายํ ทูเตยฺยญฺจ, ตุลากูฏํ อุกฺโกฏนํ. \\ เฉทนํ วธพนฺธนํ, วิปราโลปํ สาหสนฺติ.}}

\csromanheader{6}{11. ปญฺจคติเปยฺยาลวคฺค}
\csromanheader{2}{1. มนุสฺสจุตินิรยสุตฺต}
\proseitem{1172}{อถ โข ภควา ปริตฺตํ นขสิขายํ ปํสุํ อาโรเปตฺวา ภิกฺขู อามนฺเตสิ “ตํ กิํ มญฺญถ ภิกฺขเว, กตมํ นุ โข พหุตรํ, โย วายํ มยา ปริตฺโต นขสิขายํ ปํสุ อาโรปิโต, อยํ วา มหาปถวี”ติ.\csromanspacer{} เอตเทว ภนฺเต พหุตรํ ยทิทํ มหาปถวี, อปฺปมตฺตกายํ ภควตา ปริตฺโต นขสิขายํ ปํสุ อาโรปิโต, สงฺขมฺปิ น อุเปติ อุปนิธมฺปิ น อุเปติ กลภาคมฺปิ น อุเปติ, มหาปถวิํ อุปนิธาย ภควตา ปริตฺโต นขสิขายํ ปํสุ อาโรปิโตติ.\csromanspacer{} เอวเมว โข ภิกฺขเว อปฺปกา เต สตฺตา, เย มนุสฺสา จุตา มนุสฺเสสุ ปจฺจาชายนฺติ.\csromanspacer{} อถ โข เอเตว พหุตรา สตฺตา, เย มนุสฺสา จุตา นิรเย ปจฺจาชายนฺติ ฯเปฯ.\csromanspacer{}\csromanspacer{}ปฐมํ.}

\csromanheader{2}{2. มนุสฺสจุติติรจฺฉนสุตฺต}
\proseitem{1173}{เอวเมว โข ภิกฺขเว อปฺปกา เต สตฺตา, เย มนุสฺสา จุตา มนุสฺเสสุ ปจฺจาชายนฺติ.\csromanspacer{} อถ โข เอเตว พหุตรา สตฺตา, เย มนุสฺสา จุตา ติรจฺฉนโยนิยา ปจฺจาชายนฺติ ฯเปฯ.\csromanspacer{}\csromanspacer{}ทุติยํ.}

\csromanheader{2}{3. มนุสฺสจุติเปตฺติวิสยสุตฺต}
\proseitem{1174}{เอวเมว โข ภิกฺขเว อปฺปกา เต สตฺตา, เย มนุสฺสา จุตา มนุสฺเสสุ ปจฺจาชายนฺติ.\csromanspacer{} อถ โข เอเตว พหุตรา สตฺตา, เย มนุสฺสา จุตา เปตฺติวิสเย ปจฺจาชายนฺติ ฯเปฯ.\csromanspacer{}\csromanspacer{}ตติยํ.}

\prose{\textbf{4-5-6.}\csromanspacer{}\textbf{ มนุสฺสจุติเทวนิรยาทิสุตฺต}}

\proseitem{1175-1177}{เอวเมว โข ภิกฺขเว อปฺปกา เต สตฺตา, เย มนุสฺสา จุตา เทเวสุ ปจฺจาชายนฺติ.\csromanspacer{} อถ โข เอเตว พหุตรา สตฺตา,}

\csromanheader{2}{12. สจฺจสํยุตฺต}
\prose{\csromanfolio{413}เย มนุสฺสา สุตา นิรเย ปจฺจชายนฺติ ฯเปฯ ติรจฺฉานโยนิยา ปจฺจชายนฺติ ฯเปฯ เปตฺติวิสเย ปจฺจาชายนฺติ ฯเปฯ.\csromanspacer{} ฉฏฺฐํ.}

\csromanheader{2}{7-9. เทวจุตินิรยาทิสุตฺต}
\proseitem{1178-1180}{เอวเมว โข ภิกฺขเว อปฺปกา เต สตฺตา, เย เทวา จุตา เทเวสุ ปจฺจชายนฺติ.\csromanspacer{} อถ โข เอเตว พหุตรา สตฺตา, เย เทวา จุตา นิรเย ปจฺจาชายนฺติ ฯเปฯ ติรจฺฉานโยนิยา ปจฺจาชายนฺติ ฯเปฯ เปตฺติวิสเย ปจฺจาชายนฺติ ฯเปฯ.\csromanspacer{}\csromanspacer{}นวมํ.}

\csromanheader{2}{10-12. เทวมนุสฺสนิรยสุตฺต}
\proseitem{1181-1183}{เอวเมว โข ภิกฺขเว อปฺปกา เต สตฺตา, เย เทวา จุตา มนุสฺเสสุ ปจฺจาชายนฺติ.\csromanspacer{} อถ โข เอเตว พหุตรา สตฺตา, เย เทวา จุตา นิรเย ปจฺจาชายนฺติ ฯเปฯ ติรจฺฉานโยนิยา ปจฺจาชายนฺติ ฯเปฯ เปตฺติวิสเย ปจฺจาชายนฺติ.\csromanspacer{}\csromanspacer{}ทฺวาทสมํ.}

\csromanheader{2}{13-15. นิรยมนุสฺสนิรยาทิสุตฺต}
\proseitem{1184-1186}{เอวเมว โข ภิกฺขเว อปฺปกา เต สตฺตา, เย นิรยา จุตา มนุสฺเสสุ ปจฺจาชายนฺติ.\csromanspacer{} อถ โข เอเตว พหุตรา สตฺตา, เย นิรยา จุตา นิรเย ปจฺจาชายนฺติ ฯเปฯ ติรจฺฉานโยนิยา ปจฺจาชายนฺติ ฯเปฯ เปตฺติวิสเย ปจฺจาชายนฺติ ฯเปฯ.\csromanspacer{}\csromanspacer{}ปนฺนรสมํ.}

\csromanheader{2}{16-18. นิรยเทวนิรยาทิสุตฺต}
\proseitem{1187-1179}{เอวเมว โข ภิกฺขเว อปฺปกา เต สตฺตา, เย นิรยา จุตา เทเวสุ ปจฺจาชายนฺติ.\csromanspacer{} อถ โข เอเตว พหุตรา สตฺตา, เย นิรยา จุตา นิรเย ปจฺจาชายนฺติ ฯเปฯ ติรจฺฉานโยนิยา ปจฺจาชายนฺติ ฯเปฯ เปตฺติวิสเย ปจฺจาชายนฺติ ฯเปฯ.\csromanspacer{}\csromanspacer{}อฏฺฐารสมํ.}

\csromanheader{2}{19-21. ติรจฺฉานมนุสฺสนิรยาทิสุตฺต}
\proseitem{1190-1192}{เอวเมว โข ภิกฺขเว อปฺปกา เต สตฺตา, เย ติรจฺฉานโยนิยา จุตา มนุสฺเสสุ ปจฺจาชายนฺติ.\csromanspacer{} อถ โข เอเตว พหุตรา สตฺตา, เย ติรจฺฉานโยนิยา จุตา นิรเย ปจฺจาชายนฺติ ฯเปฯ \csromanfolio{414}ติรจฺฉานโยนิยา ปจฺจาชายนฺติ ฯเปฯ เปตฺติวิสเย ปจฺจาชายนฺติ ฯเปฯ.\csromanspacer{}\csromanspacer{}เอกวีสติมํ.}

\csromanheader{2}{22-24. ติรจฺฉานเทวนิรยาทิสุตฺต}
\proseitem{1193-1195}{เอวเมว โข ภิกฺขเว อปฺปกา เต สตฺตา, เย ติรจฺฉานโยนิยา จุตา เทเวสุ ปจฺจาชายนฺติ.\csromanspacer{} อถ โข เอเตว พหุตรา สตฺตา, เย ติรจฺฉานโยนิยา จุตา นิรเย ปจฺจาชายนฺติ ฯเปฯ ติรจฺฉานโยนิยา ปจฺจาชายนฺติ ฯเปฯ เปตฺติวิสเย ปจฺจาชายนฺติ ฯเปฯ.\csromanspacer{} จตุวีสติมํ.}

\csromanheader{2}{25-27. เปตฺติมนุสฺสนิรยาทิสุตฺต}
\proseitem{1196-1198}{เอวเมว โข ภิกฺขเว อปฺปกา เต สตฺตา, เย เปตฺติวิสยา จุตา มนุสฺเสสุ ปจฺจาชายนฺติ.\csromanspacer{} อถ โข เอเตว พหุตรา สตฺตา, เย เปตฺติวิสยา จุตา นิรเย ปจฺจาชายนฺติ ฯเปฯ ติรจฺฉานโยนิยา ปจฺจาชายนฺติ ฯเปฯ เปตฺติวิสเย ปจฺจาชายนฺติ ฯเปฯ.\csromanspacer{}\csromanspacer{}สตฺตวีสติมํ.}

\csromanheader{2}{28-29. เปตฺติเทวนิรยาทิสุตฺต}
\proseitem{1199-1200}{เอวเมว โข ภิกฺขเว อปฺปกา เต สตฺตา, เย เปตฺติวิสยา จุตา เทเวสุ ปจฺจาชายนฺติ.\csromanspacer{} อถ โข เอเตว พหุตรา สตฺตา, เย เปตฺติวิสยา จุตา นิรเย ปจฺจาชายนฺติ ฯเปฯ.\csromanspacer{} เอวเมว โข ภิกฺขเว อปฺปกา เต สตฺตา, เย เปตฺติวิสยา จุตา เทเวสุ ปจฺจาชายนฺติ.\csromanspacer{} อถ โข เอเตว พหุตรา สตฺตา, เย เปตฺติวิสยา จุตา ติรจฺฉานโยนิยา ปจฺจาชายนฺติ ฯเปฯ.\csromanspacer{}\csromanspacer{}เอกูนติํสติมํ.}

\csromanheader{2}{30. เปตฺติเทวเปตฺติวิสยสุตฺต}
\proseitem{1201}{เอวเมว โข ภิกฺขเว อปฺปกา เต สตฺตา, เย เปตฺติวิสยา จุตา เทเวสุ ปจฺจาชายนฺติ.\csromanspacer{} อถ โข เอเตว พหุตรา สตฺตา, เย เปตฺติวิสยา จุตา เปตฺติวิสเย ปจฺจาชายนฺติ.\csromanspacer{} ตํ กิสฺส เหตุ, อทิฏฺฐตฺตา ภิกฺขเว จตุนฺนํ อริยสจฺจานํ.\csromanspacer{} กตเมสํ จตุนฺนํ, ทุกฺขสฺส อริยสจฺจสฺส ทุกฺขสมุทยสฺส อริยสจฺจสฺส ทุกฺขนิโรธสฺส อริยสจฺจสฺส ทุกฺขนิโรธคามินิยา ปฏิปทาย อริยสจฺจสฺส.}

\csromanheader{2}{12. สจฺจสํยุตฺต}
\prose{\csromanfolio{415}ตสฺมาติห ภิกฺขเว “อิทํ ทุกฺขนฺ”ติ โยโค กรณีโย, “อยํ ทุกฺขสมุทโย”ติ โยโค กรณีโย, “อยํ ทุกฺขนิโรโธ”ติ โยโค กรณีโย, “อยํ ทุกฺขนิโรธคามินี ปฏิปทา”ติ โยโค กรณีโยติ.\csromanspacer{} อิทมโวจ ภควา.\csromanspacer{} อตฺตมนา เต ภิกฺขู ภควโต ภาสิตํ อภินนฺทุนฺติ.\csromanspacer{}\csromanspacer{}ติํสติมํ.\csromanspacer{} ปญฺจคติเปยฺยาลวคฺโค เอกาทสโม.}

\titlehead{\textbf{ตสฺสุทฺทานํ}}
\csromangathameasure{มนุสฺสโต จุตา ฉาปิ, เทวา จุตา นิรยโต. \\ ติรจฺฉาเปตฺติวิสยา, ติํสมตฺโต คติวคฺโคติ. \\ สจฺจสํยุตฺตํ ทฺวาทสมํ. มหาวคฺโค ปญฺจโม.}
\csromangathagroup{\csromangathastanza{มนุสฺสโต จุตา ฉาปิ, เทวา จุตา นิรยโต. \\ ติรจฺฉาเปตฺติวิสยา, ติํสมตฺโต คติวคฺโคติ. \\ สจฺจสํยุตฺตํ ทฺวาทสมํ. มหาวคฺโค ปญฺจโม.}}

\titlehead{\textbf{ตสฺสุทฺทานํ}}
\csromangathameasure{มคฺคโพชฺฌงฺคํ สติยา, อินฺทฺริยํ สมฺมปฺปธานํ. \\ พลิทฺธิปาทานุรุทฺธา, ฌานานาปาน สํยุตฺตํ. \\ โสตาปตฺติ สจฺจญฺจาติ, มหาวคฺโคติ วุจฺจตีติ.}
\csromangathagroup{\csromangathastanza{มคฺคโพชฺฌงฺคํ สติยา, อินฺทฺริยํ สมฺมปฺปธานํ. \\ พลิทฺธิปาทานุรุทฺธา, ฌานานาปาน สํยุตฺตํ. \\ โสตาปตฺติ สจฺจญฺจาติ, มหาวคฺโคติ วุจฺจตีติ.}}

\nitthitam{\textbf{มหาวคฺคสํยุตฺตปาฬิ นิฏฺฐิตา.}}
